\documentclass[12pt,a4paper]{article}
\usepackage{fontspec}
\setmainfont{FreeSerif} % или DejaVuSerif, LiberationSerif
\usepackage[utf8]{inputenc}
\usepackage[russian]{babel}
\usepackage{amsmath}
\usepackage{amssymb}
\usepackage{geometry}
\usepackage{float}
\usepackage{graphicx}
\usepackage{booktabs}
\usepackage{longtable}
\usepackage{hyperref}
\usepackage{tikz}
\geometry{margin=2cm}

\title{Лабораторная работа 1\\Градиентный спуск для тестовых функций}
\author{}
\date{}

\begin{document}
\maketitle

\section{Постановка задачи}
В работе исследуются две реализации градиентного спуска:
\begin{itemize}
\item \texttt{gradDest} --- метод с начальным шагом $t$ и внутренним циклом дробления шага до достижения монотонного уменьшения значения функции;
\item \texttt{fastGradDest} --- метод с одномерным поиском шага методом золотого сечения.
\end{itemize}

Оба метода останавливаются по одному из трёх критериев: малой норме градиента, малому изменению точки или малому изменению значения функции. Для функции Adjiman используется стандартная benchmark-область $[-1, 2] \times [-1, 1]$, поэтому вычисления выполняются с проекцией на допустимый прямоугольник. Это важно, поскольку вне этой области функция не имеет конечного глобального минимума.

\section{Тестовые функции и точные минимумы}
\begin{longtable}{|c|c|c|c|c|}
\caption{Выбранные тестовые функции и их эталонные минимумы.}\label{tab:analytical}\\
\hline
Функция & Формула & Область & Точка минимума & Минимум \\
\hline
\endfirsthead
\hline
Функция & Формула & Область & Точка минимума & Минимум \\
\hline
\endhead
Adjiman & $f(x, y) = \cos(x)\sin(y) - \frac{x}{1 + y^2}$ & $[-1, 2] \times [-1, 1]$ & (2.000000, 0.105783) & -2.0218067834 \\
Ackley 2 & $f(x, y) = -200 \exp \left(-0.02 \sqrt{x^2 + y^2}\right)$ & $[-32, 32] \times [-32, 32]$ & (0.000000, 0.000000) & -200.0000000000 \\
Schaffer N.2 & $f(x, y) = 0.5 + \frac{\sin^2(x^2 - y^2) - 0.5}{(1 + 0.001(x^2 + y^2))^2}$ & $[-100, 100] \times [-100, 100]$ & (0.000000, 0.000000) & 0.0000000000 \\
Rosenbrock & $f(x, y) = (1 - x)^2 + b\,(y - x^2)^2$ & $[-5, 5] \times [-5, 5]$ & (1.000000, 1.000000) & 0.0000000000 \\
\hline
\end{longtable}


Для функции Adjiman минимум стандартной benchmark-версии находится на границе по $x$ в точке
\[
x^\star = 2, \qquad
y^\star \approx 0.1057834694517169,
\]
где $y^\star$ является корнем уравнения
\[
\cos(2)\cos(y) + \frac{4y}{(1+y^2)^2} = 0.
\]
Численно это даёт значение
\[
f(x^\star, y^\star) \approx -2.0218067833597866.
\]

Для функции Ackley 2 глобальный минимум равен $-200$ в точке $(0,0)$, а для Schaffer N.2 глобальный минимум равен $0$ в точке $(0,0)$.

\section{Описание алгоритмов}
Исправленная реализация \texttt{gradDest} использует внутренний цикл подбора шага: если пробный переход не уменьшает значение функции, шаг делится пополам до тех пор, пока не будет найден убывающий шаг или пока шаг не станет слишком малым. Это устраняет логическую ошибку прежней версии, где шаг уменьшался, но текущая точка не обновлялась, что могло приводить к зависанию.

В реализации \texttt{fastGradDest} устранено лишнее повторное вычисление градиента в начале итерации, убраны дублирующиеся проверки критерия по изменению функции, а параметр \texttt{gradNormEps} больше не используется как стартовая граница для одномерного поиска. Для поиска правой границы отрезка теперь используется отдельный начальный шаг \texttt{initialStep}.

\section{Эксперименты}
\subsection{Влияние начальной точки}
\begin{longtable}{|c|c|c|c|c|c|c|}
\caption{Влияние начальной точки на сходимость методов.}\label{tab:starts}\\
\hline
Функция & Метод & Старт & Найденная точка & $f(x)$ & Итерации & Вызовы $f$ \\
\hline
\endfirsthead
\hline
Функция & Метод & Старт & Найденная точка & $f(x)$ & Итерации & Вызовы $f$ \\
\hline
\endhead
Adjiman & gradDest & (-1.000000, -1.000000) & (0.636164, -1.000000) & -0.994945 & 8 & 48 \\
Adjiman & fastGradDest & (-1.000000, -1.000000) & (0.636268, -1.000000) & -0.994945 & 2 & 74 \\
Adjiman & gradDest & (2.000000, 2.000000) & (2.000000, 0.106212) & -2.021806 & 47 & 328 \\
Adjiman & fastGradDest & (2.000000, 2.000000) & (2.000000, 0.105783) & -2.021807 & 2 & 71 \\
Adjiman & gradDest & (5.000000, 5.000000) & (2.000000, 0.106212) & -2.021806 & 47 & 328 \\
Adjiman & fastGradDest & (5.000000, 5.000000) & (2.000000, 0.105783) & -2.021807 & 2 & 71 \\
Adjiman & gradDest & (0.500000, 0.500000) & (2.000000, 0.105354) & -2.021806 & 51 & 355 \\
Adjiman & fastGradDest & (0.500000, 0.500000) & (2.000000, 0.105783) & -2.021807 & 4 & 142 \\
Adjiman & gradDest & (-0.562992, -0.090146) & (0.636115, -1.000000) & -0.994945 & 7 & 42 \\
Adjiman & fastGradDest & (-0.562992, -0.090146) & (0.636268, -1.000000) & -0.994945 & 2 & 74 \\
Ackley 2 & gradDest & (-1.000000, -1.000000) & (0.000000, 0.000000) & -200.000000 & 15 & 259 \\
Ackley 2 & fastGradDest & (-1.000000, -1.000000) & (-0.000000, -0.000000) & -199.999999 & 32 & 1088 \\
Ackley 2 & gradDest & (2.000000, 2.000000) & (-0.000000, -0.000000) & -200.000000 & 13 & 248 \\
Ackley 2 & fastGradDest & (2.000000, 2.000000) & (-0.000000, -0.000000) & -199.999999 & 31 & 1057 \\
Ackley 2 & gradDest & (5.000000, 5.000000) & (-0.000000, -0.000000) & -200.000000 & 12 & 224 \\
Ackley 2 & fastGradDest & (5.000000, 5.000000) & (0.000000, 0.000000) & -199.999999 & 32 & 1093 \\
Ackley 2 & gradDest & (0.500000, 0.500000) & (0.000000, 0.000000) & -200.000000 & 16 & 310 \\
Ackley 2 & fastGradDest & (0.500000, 0.500000) & (0.000000, 0.000000) & -199.999999 & 32 & 1088 \\
Ackley 2 & gradDest & (5.193736, -19.535712) & (0.000000, 0.000000) & -200.000000 & 25 & 386 \\
Ackley 2 & fastGradDest & (5.193736, -19.535712) & (0.000000, -0.000000) & -199.999999 & 36 & 1234 \\
Schaffer N.2 & gradDest & (-1.000000, -1.000000) & (-0.111557, -0.111557) & 0.000025 & 1097 & 6582 \\
Schaffer N.2 & fastGradDest & (-1.000000, -1.000000) & (0.000000, 0.000000) & 0.000000 & 1 & 63 \\
Schaffer N.2 & gradDest & (2.000000, 2.000000) & (0.111503, 0.111503) & 0.000025 & 1448 & 8688 \\
Schaffer N.2 & fastGradDest & (2.000000, 2.000000) & (-0.000000, -0.000000) & 0.000000 & 1 & 63 \\
Schaffer N.2 & gradDest & (5.000000, 5.000000) & (0.111541, 0.111541) & 0.000025 & 1938 & 11628 \\
Schaffer N.2 & fastGradDest & (5.000000, 5.000000) & (-0.000000, -0.000000) & 0.000000 & 1 & 63 \\
Schaffer N.2 & gradDest & (0.500000, 0.500000) & (0.111477, 0.111477) & 0.000025 & 750 & 4500 \\
Schaffer N.2 & fastGradDest & (0.500000, 0.500000) & (-0.000000, -0.000000) & 0.000000 & 1 & 63 \\
Schaffer N.2 & gradDest & (-70.365759, 23.193057) & (-70.404952, 23.165415) & 0.488142 & 7 & 86 \\
Schaffer N.2 & fastGradDest & (-70.365759, 23.193057) & (-72.654701, 22.438465) & 0.489130 & 2 & 68 \\
Rosenbrock & gradDest & (-1.000000, -1.000000) & (1.609039, 2.592703) & 0.372295 & 50 & 777 \\
Rosenbrock & fastGradDest & (-1.000000, -1.000000) & (-5.000000, 5.000000) & 40036.000000 & 10000 & 354997 \\
Rosenbrock & gradDest & (2.000000, 2.000000) & (1.168319, 1.366354) & 0.028523 & 4264 & 66711 \\
Rosenbrock & fastGradDest & (2.000000, 2.000000) & (5.000000, 5.000000) & 40016.000000 & 10000 & 355000 \\
Rosenbrock & gradDest & (5.000000, 5.000000) & (1.578995, 2.497126) & 0.336757 & 7098 & 117474 \\
Rosenbrock & fastGradDest & (5.000000, 5.000000) & (5.000000, 5.000000) & 40016.000000 & 10000 & 355000 \\
Rosenbrock & gradDest & (0.500000, 0.500000) & (0.936054, 0.875519) & 0.004135 & 1069 & 15528 \\
Rosenbrock & fastGradDest & (0.500000, 0.500000) & (-5.000000, 5.000000) & 40036.000000 & 10000 & 355000 \\
Rosenbrock & gradDest & (4.790501, -0.437017) & (1.578995, 2.497126) & 0.336757 & 7100 & 117486 \\
Rosenbrock & fastGradDest & (4.790501, -0.437017) & (5.000000, 5.000000) & 40016.000000 & 10000 & 355003 \\
\hline
\end{longtable}


\subsection{Влияние гиперпараметров}
\begin{longtable}{|c|c|c|c|c|c|}
\caption{Влияние начального шага $t$ на метод с дроблением шага.}\label{tab:grad_tuning}\\
\hline
Функция & $t$ & $f(x)$ & Итерации & Вызовы $f$ & Причина остановки \\
\hline
\endfirsthead
\hline
Функция & $t$ & $f(x)$ & Итерации & Вызовы $f$ & Причина остановки \\
\hline
\endhead
Adjiman & 0.1 & -0.994944 & 97 & 582 & func\_shift \\
Adjiman & 1 & -0.994945 & 8 & 48 & func\_shift \\
Adjiman & 10 & -2.021807 & 8 & 80 & func\_shift \\
Adjiman & 100 & -2.021807 & 13 & 170 & func\_shift \\
Ackley 2 & 0.1 & -200.000000 & 15 & 197 & point\_shift \\
Ackley 2 & 1 & -200.000000 & 15 & 259 & point\_shift \\
Ackley 2 & 10 & -199.999999 & 38 & 912 & func\_shift \\
Ackley 2 & 100 & -200.000000 & 12 & 295 & point\_shift \\
Schaffer N.2 & 0.1 & 0.000250 & 5209 & 31254 & func\_shift \\
Schaffer N.2 & 1 & 0.000025 & 1097 & 6582 & func\_shift \\
Schaffer N.2 & 10 & 0.000002 & 167 & 1002 & func\_shift \\
Schaffer N.2 & 100 & 0.000000 & 21 & 126 & func\_shift \\
Rosenbrock & 0.1 & 0.007455 & 2484 & 29718 & func\_shift \\
Rosenbrock & 1 & 0.372295 & 50 & 777 & func\_shift \\
Rosenbrock & 10 & 0.003195 & 1009 & 17725 & func\_shift \\
Rosenbrock & 100 & 0.001120 & 47 & 1017 & func\_shift \\
\hline
\end{longtable}

\begin{longtable}{|c|c|c|c|c|c|}
\caption{Влияние точности одномерного поиска на метод с линейным поиском.}\label{tab:fast_tuning}\\
\hline
Функция & $\varepsilon_{ls}$ & $f(x)$ & Итерации & Вызовы $f$ & Причина остановки \\
\hline
\endfirsthead
\hline
Функция & $\varepsilon_{ls}$ & $f(x)$ & Итерации & Вызовы $f$ & Причина остановки \\
\hline
\endhead
Adjiman & 1e-03 & -0.994945 & 2 & 54 & func\_shift \\
Adjiman & 1e-05 & -0.994945 & 2 & 74 & func\_shift \\
Adjiman & 1e-07 & -0.994945 & 2 & 90 & point\_shift \\
Ackley 2 & 1e-03 & -199.996820 & 10000 & 250000 & max\_iter \\
Ackley 2 & 1e-05 & -199.999999 & 32 & 1088 & func\_shift \\
Ackley 2 & 1e-07 & -200.000000 & 2 & 88 & point\_shift \\
Schaffer N.2 & 1e-03 & 0.000000 & 1 & 53 & grad\_norm \\
Schaffer N.2 & 1e-05 & 0.000000 & 1 & 63 & grad\_norm \\
Schaffer N.2 & 1e-07 & 0.000000 & 1 & 72 & grad\_norm \\
Rosenbrock & 1e-03 & 40036.000000 & 10000 & 259998 & max\_iter \\
Rosenbrock & 1e-05 & 40036.000000 & 10000 & 354997 & max\_iter \\
Rosenbrock & 1e-07 & 40036.000000 & 10000 & 449998 & max\_iter \\
\hline
\end{longtable}


По таблицам обычно видно две устойчивые тенденции:
\begin{itemize}
\item слишком большой начальный шаг в \texttt{gradDest} не ломает метод благодаря дроблению шага, но увеличивает число вычислений функции;
\item уменьшение точности одномерного поиска в \texttt{fastGradDest} снижает стоимость одной итерации, но может ухудшать итоговую точность найденного шага.
\end{itemize}

\subsection{Влияние овражности (функция Розенброка)}
Функция Розенброка $f(x,y) = (1-x)^2 + b(y-x^2)^2$ известна своей овражной структурой.
Параметр $b$ определяет крутизну оврага. В таблице~\ref{tab:rosenbrock} приведены результаты для различных $b$.
\begin{longtable}{|c|c|c|c|c|c|c|}
\caption{Влияние параметра овражности $b$ на сходимость.}\label{tab:rosenbrock}\\
\hline
Параметр $b$ & Метод & Найденная точка & $f(x)$ & Итерации & Вызовы $f$ & Причина \\
\hline
\endfirsthead
\hline
Параметр $b$ & Метод & Найденная точка & $f(x)$ & Итерации & Вызовы $f$ & Причина \\
\hline
\endhead
$1$ & gradDest & (0.998938, 0.997836) & 0.000001 & 76 & 641 & func\_shift \\
$1$ & fastGradDest & (-5.000000, 5.000000) & 436.000000 & 10000 & 354997 & max\_iter \\
$10$ & gradDest & (0.990671, 0.981571) & 0.000087 & 290 & 3321 & func\_shift \\
$10$ & fastGradDest & (-5.000000, 5.000000) & 4036.000000 & 10000 & 354997 & max\_iter \\
$100$ & gradDest & (1.166439, 1.362032) & 0.027913 & 3401 & 53138 & func\_shift \\
$100$ & fastGradDest & (-5.000000, 5.000000) & 40036.000000 & 10000 & 354997 & max\_iter \\
$1000$ & gradDest & (1.488698, 2.216151) & 0.238831 & 408 & 7735 & func\_shift \\
$1000$ & fastGradDest & (-5.000000, 5.000000) & 400036.000000 & 10000 & 354997 & max\_iter \\
\hline
\end{longtable}


На рисунке~\ref{fig:rosenbrock} показаны траектории методов для $b=100$.
\begin{figure}[H]
\centering
\begin{tikzpicture}[scale=0.9]
\draw[fill=orange!3, draw=black] (0,0) rectangle (12.0000,8.0000);
\draw[gray!35, line width=0.15pt] (1.2203,7.6277) -- (1.2027,7.7288);
\draw[gray!35, line width=0.15pt] (1.2027,7.7288) -- (1.1787,7.8644);
\draw[gray!35, line width=0.15pt] (1.1787,7.8644) -- (1.1544,8.0000);
\draw[gray!35, line width=0.15pt] (1.4237,6.5700) -- (1.4102,6.6441);
\draw[gray!35, line width=0.15pt] (1.4102,6.6441) -- (1.3852,6.7797);
\draw[gray!35, line width=0.15pt] (1.3852,6.7797) -- (1.3599,6.9153);
\draw[gray!35, line width=0.15pt] (1.3599,6.9153) -- (1.3342,7.0508);
\draw[gray!35, line width=0.15pt] (1.3342,7.0508) -- (1.3081,7.1864);
\draw[gray!35, line width=0.15pt] (1.3081,7.1864) -- (1.2816,7.3220);
\draw[gray!35, line width=0.15pt] (1.2816,7.3220) -- (1.2547,7.4576);
\draw[gray!35, line width=0.15pt] (1.2547,7.4576) -- (1.2273,7.5932);
\draw[gray!35, line width=0.15pt] (1.2203,7.6277) -- (1.2273,7.5932);
\draw[gray!35, line width=0.15pt] (1.6271,5.5580) -- (1.6269,5.5593);
\draw[gray!35, line width=0.15pt] (1.6269,5.5593) -- (1.6009,5.6949);
\draw[gray!35, line width=0.15pt] (1.6009,5.6949) -- (1.5745,5.8305);
\draw[gray!35, line width=0.15pt] (1.5745,5.8305) -- (1.5478,5.9661);
\draw[gray!35, line width=0.15pt] (1.5478,5.9661) -- (1.5207,6.1017);
\draw[gray!35, line width=0.15pt] (1.5207,6.1017) -- (1.4931,6.2373);
\draw[gray!35, line width=0.15pt] (1.4931,6.2373) -- (1.4652,6.3729);
\draw[gray!35, line width=0.15pt] (1.4652,6.3729) -- (1.4367,6.5085);
\draw[gray!35, line width=0.15pt] (1.4237,6.5700) -- (1.4367,6.5085);
\draw[gray!35, line width=0.15pt] (1.8305,4.5923) -- (1.8269,4.6102);
\draw[gray!35, line width=0.15pt] (1.8269,4.6102) -- (1.7996,4.7458);
\draw[gray!35, line width=0.15pt] (1.7996,4.7458) -- (1.7718,4.8814);
\draw[gray!35, line width=0.15pt] (1.7718,4.8814) -- (1.7437,5.0169);
\draw[gray!35, line width=0.15pt] (1.7437,5.0169) -- (1.7152,5.1525);
\draw[gray!35, line width=0.15pt] (1.7152,5.1525) -- (1.6862,5.2881);
\draw[gray!35, line width=0.15pt] (1.6862,5.2881) -- (1.6567,5.4237);
\draw[gray!35, line width=0.15pt] (1.6271,5.5580) -- (1.6567,5.4237);
\draw[gray!35, line width=0.15pt] (2.0339,3.6725) -- (2.0076,3.7966);
\draw[gray!35, line width=0.15pt] (2.0076,3.7966) -- (1.9785,3.9322);
\draw[gray!35, line width=0.15pt] (1.9785,3.9322) -- (1.9490,4.0678);
\draw[gray!35, line width=0.15pt] (1.9490,4.0678) -- (1.9190,4.2034);
\draw[gray!35, line width=0.15pt] (1.9190,4.2034) -- (1.8886,4.3390);
\draw[gray!35, line width=0.15pt] (1.8886,4.3390) -- (1.8577,4.4746);
\draw[gray!35, line width=0.15pt] (1.8305,4.5923) -- (1.8577,4.4746);
\draw[gray!35, line width=0.15pt] (2.2373,2.7988) -- (2.2264,2.8475);
\draw[gray!35, line width=0.15pt] (2.2264,2.8475) -- (2.1959,2.9831);
\draw[gray!35, line width=0.15pt] (2.1959,2.9831) -- (2.1650,3.1186);
\draw[gray!35, line width=0.15pt] (2.1650,3.1186) -- (2.1337,3.2542);
\draw[gray!35, line width=0.15pt] (2.1337,3.2542) -- (2.1019,3.3898);
\draw[gray!35, line width=0.15pt] (2.1019,3.3898) -- (2.0695,3.5254);
\draw[gray!35, line width=0.15pt] (2.0695,3.5254) -- (2.0367,3.6610);
\draw[gray!35, line width=0.15pt] (2.0339,3.6725) -- (2.0367,3.6610);
\draw[gray!35, line width=0.15pt] (2.4407,1.9710) -- (2.4258,2.0339);
\draw[gray!35, line width=0.15pt] (2.4258,2.0339) -- (2.3935,2.1695);
\draw[gray!35, line width=0.15pt] (2.3935,2.1695) -- (2.3608,2.3051);
\draw[gray!35, line width=0.15pt] (2.3608,2.3051) -- (2.3275,2.4407);
\draw[gray!35, line width=0.15pt] (2.3275,2.4407) -- (2.2938,2.5763);
\draw[gray!35, line width=0.15pt] (2.2938,2.5763) -- (2.2595,2.7119);
\draw[gray!35, line width=0.15pt] (2.2373,2.7988) -- (2.2595,2.7119);
\draw[gray!35, line width=0.15pt] (2.6441,1.1891) -- (2.6362,1.2203);
\draw[gray!35, line width=0.15pt] (2.6362,1.2203) -- (2.6021,1.3559);
\draw[gray!35, line width=0.15pt] (2.6021,1.3559) -- (2.5674,1.4915);
\draw[gray!35, line width=0.15pt] (2.5674,1.4915) -- (2.5322,1.6271);
\draw[gray!35, line width=0.15pt] (2.5322,1.6271) -- (2.4965,1.7627);
\draw[gray!35, line width=0.15pt] (2.4965,1.7627) -- (2.4603,1.8983);
\draw[gray!35, line width=0.15pt] (2.4407,1.9710) -- (2.4603,1.8983);
\draw[gray!35, line width=0.15pt] (2.8475,0.4533) -- (2.8236,0.5424);
\draw[gray!35, line width=0.15pt] (2.8236,0.5424) -- (2.7870,0.6780);
\draw[gray!35, line width=0.15pt] (2.7870,0.6780) -- (2.7499,0.8136);
\draw[gray!35, line width=0.15pt] (2.7499,0.8136) -- (2.7122,0.9492);
\draw[gray!35, line width=0.15pt] (2.7122,0.9492) -- (2.6739,1.0847);
\draw[gray!35, line width=0.15pt] (2.6441,1.1891) -- (2.6739,1.0847);
\draw[gray!35, line width=0.15pt] (2.9827,0.0000) -- (2.9430,0.1356);
\draw[gray!35, line width=0.15pt] (2.9430,0.1356) -- (2.9026,0.2712);
\draw[gray!35, line width=0.15pt] (2.9026,0.2712) -- (2.8616,0.4068);
\draw[gray!35, line width=0.15pt] (2.8475,0.4533) -- (2.8616,0.4068);
\draw[gray!35, line width=0.15pt] (9.0183,0.0000) -- (9.0581,0.1356);
\draw[gray!35, line width=0.15pt] (9.0581,0.1356) -- (9.0985,0.2712);
\draw[gray!35, line width=0.15pt] (9.0985,0.2712) -- (9.1395,0.4068);
\draw[gray!35, line width=0.15pt] (9.1395,0.4068) -- (9.1525,0.4496);
\draw[gray!35, line width=0.15pt] (9.1525,0.4496) -- (9.1774,0.5424);
\draw[gray!35, line width=0.15pt] (9.1774,0.5424) -- (9.2140,0.6780);
\draw[gray!35, line width=0.15pt] (9.2140,0.6780) -- (9.2512,0.8136);
\draw[gray!35, line width=0.15pt] (9.2512,0.8136) -- (9.2889,0.9492);
\draw[gray!35, line width=0.15pt] (9.2889,0.9492) -- (9.3272,1.0847);
\draw[gray!35, line width=0.15pt] (9.3272,1.0847) -- (9.3559,1.1852);
\draw[gray!35, line width=0.15pt] (9.3559,1.1852) -- (9.3648,1.2203);
\draw[gray!35, line width=0.15pt] (9.3648,1.2203) -- (9.3990,1.3559);
\draw[gray!35, line width=0.15pt] (9.3990,1.3559) -- (9.4336,1.4915);
\draw[gray!35, line width=0.15pt] (9.4336,1.4915) -- (9.4688,1.6271);
\draw[gray!35, line width=0.15pt] (9.4688,1.6271) -- (9.5045,1.7627);
\draw[gray!35, line width=0.15pt] (9.5045,1.7627) -- (9.5408,1.8983);
\draw[gray!35, line width=0.15pt] (9.5408,1.8983) -- (9.5593,1.9668);
\draw[gray!35, line width=0.15pt] (9.5593,1.9668) -- (9.5752,2.0339);
\draw[gray!35, line width=0.15pt] (9.5752,2.0339) -- (9.6075,2.1695);
\draw[gray!35, line width=0.15pt] (9.6075,2.1695) -- (9.6403,2.3051);
\draw[gray!35, line width=0.15pt] (9.6403,2.3051) -- (9.6736,2.4407);
\draw[gray!35, line width=0.15pt] (9.6736,2.4407) -- (9.7073,2.5763);
\draw[gray!35, line width=0.15pt] (9.7073,2.5763) -- (9.7416,2.7119);
\draw[gray!35, line width=0.15pt] (9.7416,2.7119) -- (9.7627,2.7944);
\draw[gray!35, line width=0.15pt] (9.7627,2.7944) -- (9.7746,2.8475);
\draw[gray!35, line width=0.15pt] (9.7746,2.8475) -- (9.8051,2.9831);
\draw[gray!35, line width=0.15pt] (9.8051,2.9831) -- (9.8360,3.1186);
\draw[gray!35, line width=0.15pt] (9.8360,3.1186) -- (9.8674,3.2542);
\draw[gray!35, line width=0.15pt] (9.8674,3.2542) -- (9.8992,3.3898);
\draw[gray!35, line width=0.15pt] (9.8992,3.3898) -- (9.9316,3.5254);
\draw[gray!35, line width=0.15pt] (9.9316,3.5254) -- (9.9645,3.6610);
\draw[gray!35, line width=0.15pt] (9.9645,3.6610) -- (9.9661,3.6678);
\draw[gray!35, line width=0.15pt] (9.9661,3.6678) -- (9.9934,3.7966);
\draw[gray!35, line width=0.15pt] (9.9934,3.7966) -- (10.0225,3.9322);
\draw[gray!35, line width=0.15pt] (10.0225,3.9322) -- (10.0521,4.0678);
\draw[gray!35, line width=0.15pt] (10.0521,4.0678) -- (10.0820,4.2034);
\draw[gray!35, line width=0.15pt] (10.0820,4.2034) -- (10.1125,4.3390);
\draw[gray!35, line width=0.15pt] (10.1125,4.3390) -- (10.1434,4.4746);
\draw[gray!35, line width=0.15pt] (10.1434,4.4746) -- (10.1695,4.5875);
\draw[gray!35, line width=0.15pt] (10.1695,4.5875) -- (10.1741,4.6102);
\draw[gray!35, line width=0.15pt] (10.1741,4.6102) -- (10.2014,4.7458);
\draw[gray!35, line width=0.15pt] (10.2014,4.7458) -- (10.2292,4.8814);
\draw[gray!35, line width=0.15pt] (10.2292,4.8814) -- (10.2573,5.0169);
\draw[gray!35, line width=0.15pt] (10.2573,5.0169) -- (10.2859,5.1525);
\draw[gray!35, line width=0.15pt] (10.2859,5.1525) -- (10.3149,5.2881);
\draw[gray!35, line width=0.15pt] (10.3149,5.2881) -- (10.3444,5.4237);
\draw[gray!35, line width=0.15pt] (10.3444,5.4237) -- (10.3729,5.5529);
\draw[gray!35, line width=0.15pt] (10.3729,5.5529) -- (10.3741,5.5593);
\draw[gray!35, line width=0.15pt] (10.3741,5.5593) -- (10.4001,5.6949);
\draw[gray!35, line width=0.15pt] (10.4001,5.6949) -- (10.4265,5.8305);
\draw[gray!35, line width=0.15pt] (10.4265,5.8305) -- (10.4533,5.9661);
\draw[gray!35, line width=0.15pt] (10.4533,5.9661) -- (10.4804,6.1017);
\draw[gray!35, line width=0.15pt] (10.4804,6.1017) -- (10.5080,6.2373);
\draw[gray!35, line width=0.15pt] (10.5080,6.2373) -- (10.5360,6.3729);
\draw[gray!35, line width=0.15pt] (10.5360,6.3729) -- (10.5644,6.5085);
\draw[gray!35, line width=0.15pt] (10.5644,6.5085) -- (10.5763,6.5646);
\draw[gray!35, line width=0.15pt] (10.5763,6.5646) -- (10.5908,6.6441);
\draw[gray!35, line width=0.15pt] (10.5908,6.6441) -- (10.6158,6.7797);
\draw[gray!35, line width=0.15pt] (10.6158,6.7797) -- (10.6412,6.9153);
\draw[gray!35, line width=0.15pt] (10.6412,6.9153) -- (10.6669,7.0508);
\draw[gray!35, line width=0.15pt] (10.6669,7.0508) -- (10.6930,7.1864);
\draw[gray!35, line width=0.15pt] (10.6930,7.1864) -- (10.7195,7.3220);
\draw[gray!35, line width=0.15pt] (10.7195,7.3220) -- (10.7464,7.4576);
\draw[gray!35, line width=0.15pt] (10.7464,7.4576) -- (10.7738,7.5932);
\draw[gray!35, line width=0.15pt] (10.7738,7.5932) -- (10.7797,7.6221);
\draw[gray!35, line width=0.15pt] (10.7797,7.6221) -- (10.7983,7.7288);
\draw[gray!35, line width=0.15pt] (10.7983,7.7288) -- (10.8223,7.8644);
\draw[gray!35, line width=0.15pt] (10.8223,7.8644) -- (10.8466,8.0000);
\draw[gray!45, line width=0.15pt] (0.6102,7.3155) -- (0.6091,7.3220);
\draw[gray!45, line width=0.15pt] (0.6091,7.3220) -- (0.5878,7.4576);
\draw[gray!45, line width=0.15pt] (0.5878,7.4576) -- (0.5663,7.5932);
\draw[gray!45, line width=0.15pt] (0.5663,7.5932) -- (0.5446,7.7288);
\draw[gray!45, line width=0.15pt] (0.5446,7.7288) -- (0.5226,7.8644);
\draw[gray!45, line width=0.15pt] (0.5226,7.8644) -- (0.5005,8.0000);
\draw[gray!45, line width=0.15pt] (0.8136,6.1200) -- (0.7944,6.2373);
\draw[gray!45, line width=0.15pt] (0.7944,6.2373) -- (0.7721,6.3729);
\draw[gray!45, line width=0.15pt] (0.7721,6.3729) -- (0.7495,6.5085);
\draw[gray!45, line width=0.15pt] (0.7495,6.5085) -- (0.7267,6.6441);
\draw[gray!45, line width=0.15pt] (0.7267,6.6441) -- (0.7037,6.7797);
\draw[gray!45, line width=0.15pt] (0.7037,6.7797) -- (0.6804,6.9153);
\draw[gray!45, line width=0.15pt] (0.6804,6.9153) -- (0.6569,7.0508);
\draw[gray!45, line width=0.15pt] (0.6569,7.0508) -- (0.6331,7.1864);
\draw[gray!45, line width=0.15pt] (0.6102,7.3155) -- (0.6331,7.1864);
\draw[gray!45, line width=0.15pt] (1.0169,4.9706) -- (1.0091,5.0169);
\draw[gray!45, line width=0.15pt] (1.0091,5.0169) -- (0.9859,5.1525);
\draw[gray!45, line width=0.15pt] (0.9859,5.1525) -- (0.9625,5.2881);
\draw[gray!45, line width=0.15pt] (0.9625,5.2881) -- (0.9389,5.4237);
\draw[gray!45, line width=0.15pt] (0.9389,5.4237) -- (0.9150,5.5593);
\draw[gray!45, line width=0.15pt] (0.9150,5.5593) -- (0.8909,5.6949);
\draw[gray!45, line width=0.15pt] (0.8909,5.6949) -- (0.8665,5.8305);
\draw[gray!45, line width=0.15pt] (0.8665,5.8305) -- (0.8419,5.9661);
\draw[gray!45, line width=0.15pt] (0.8419,5.9661) -- (0.8169,6.1017);
\draw[gray!45, line width=0.15pt] (0.8136,6.1200) -- (0.8169,6.1017);
\draw[gray!45, line width=0.15pt] (1.2203,3.8670) -- (1.2088,3.9322);
\draw[gray!45, line width=0.15pt] (1.2088,3.9322) -- (1.1846,4.0678);
\draw[gray!45, line width=0.15pt] (1.1846,4.0678) -- (1.1601,4.2034);
\draw[gray!45, line width=0.15pt] (1.1601,4.2034) -- (1.1354,4.3390);
\draw[gray!45, line width=0.15pt] (1.1354,4.3390) -- (1.1105,4.4746);
\draw[gray!45, line width=0.15pt] (1.1105,4.4746) -- (1.0853,4.6102);
\draw[gray!45, line width=0.15pt] (1.0853,4.6102) -- (1.0598,4.7458);
\draw[gray!45, line width=0.15pt] (1.0598,4.7458) -- (1.0340,4.8814);
\draw[gray!45, line width=0.15pt] (1.0169,4.9706) -- (1.0340,4.8814);
\draw[gray!45, line width=0.15pt] (1.4237,2.8094) -- (1.4167,2.8475);
\draw[gray!45, line width=0.15pt] (1.4167,2.8475) -- (1.3914,2.9831);
\draw[gray!45, line width=0.15pt] (1.3914,2.9831) -- (1.3659,3.1186);
\draw[gray!45, line width=0.15pt] (1.3659,3.1186) -- (1.3402,3.2542);
\draw[gray!45, line width=0.15pt] (1.3402,3.2542) -- (1.3142,3.3898);
\draw[gray!45, line width=0.15pt] (1.3142,3.3898) -- (1.2879,3.5254);
\draw[gray!45, line width=0.15pt] (1.2879,3.5254) -- (1.2613,3.6610);
\draw[gray!45, line width=0.15pt] (1.2613,3.6610) -- (1.2344,3.7966);
\draw[gray!45, line width=0.15pt] (1.2203,3.8670) -- (1.2344,3.7966);
\draw[gray!45, line width=0.15pt] (1.6271,1.7978) -- (1.6076,1.8983);
\draw[gray!45, line width=0.15pt] (1.6076,1.8983) -- (1.5810,2.0339);
\draw[gray!45, line width=0.15pt] (1.5810,2.0339) -- (1.5542,2.1695);
\draw[gray!45, line width=0.15pt] (1.5542,2.1695) -- (1.5271,2.3051);
\draw[gray!45, line width=0.15pt] (1.5271,2.3051) -- (1.4997,2.4407);
\draw[gray!45, line width=0.15pt] (1.4997,2.4407) -- (1.4720,2.5763);
\draw[gray!45, line width=0.15pt] (1.4720,2.5763) -- (1.4440,2.7119);
\draw[gray!45, line width=0.15pt] (1.4237,2.8094) -- (1.4440,2.7119);
\draw[gray!45, line width=0.15pt] (1.8305,0.8321) -- (1.8066,0.9492);
\draw[gray!45, line width=0.15pt] (1.8066,0.9492) -- (1.7787,1.0847);
\draw[gray!45, line width=0.15pt] (1.7787,1.0847) -- (1.7505,1.2203);
\draw[gray!45, line width=0.15pt] (1.7505,1.2203) -- (1.7221,1.3559);
\draw[gray!45, line width=0.15pt] (1.7221,1.3559) -- (1.6933,1.4915);
\draw[gray!45, line width=0.15pt] (1.6933,1.4915) -- (1.6642,1.6271);
\draw[gray!45, line width=0.15pt] (1.6642,1.6271) -- (1.6348,1.7627);
\draw[gray!45, line width=0.15pt] (1.6271,1.7978) -- (1.6348,1.7627);
\draw[gray!45, line width=0.15pt] (2.0151,0.0000) -- (1.9858,0.1356);
\draw[gray!45, line width=0.15pt] (1.9858,0.1356) -- (1.9563,0.2712);
\draw[gray!45, line width=0.15pt] (1.9563,0.2712) -- (1.9264,0.4068);
\draw[gray!45, line width=0.15pt] (1.9264,0.4068) -- (1.8962,0.5424);
\draw[gray!45, line width=0.15pt] (1.8962,0.5424) -- (1.8656,0.6780);
\draw[gray!45, line width=0.15pt] (1.8656,0.6780) -- (1.8347,0.8136);
\draw[gray!45, line width=0.15pt] (1.8305,0.8321) -- (1.8347,0.8136);
\draw[gray!45, line width=0.15pt] (9.9856,0.0000) -- (10.0149,0.1356);
\draw[gray!45, line width=0.15pt] (10.0149,0.1356) -- (10.0445,0.2712);
\draw[gray!45, line width=0.15pt] (10.0445,0.2712) -- (10.0744,0.4068);
\draw[gray!45, line width=0.15pt] (10.0744,0.4068) -- (10.1046,0.5424);
\draw[gray!45, line width=0.15pt] (10.1046,0.5424) -- (10.1352,0.6780);
\draw[gray!45, line width=0.15pt] (10.1352,0.6780) -- (10.1661,0.8136);
\draw[gray!45, line width=0.15pt] (10.1661,0.8136) -- (10.1695,0.8286);
\draw[gray!45, line width=0.15pt] (10.1695,0.8286) -- (10.1941,0.9492);
\draw[gray!45, line width=0.15pt] (10.1941,0.9492) -- (10.2220,1.0847);
\draw[gray!45, line width=0.15pt] (10.2220,1.0847) -- (10.2502,1.2203);
\draw[gray!45, line width=0.15pt] (10.2502,1.2203) -- (10.2787,1.3559);
\draw[gray!45, line width=0.15pt] (10.2787,1.3559) -- (10.3075,1.4915);
\draw[gray!45, line width=0.15pt] (10.3075,1.4915) -- (10.3366,1.6271);
\draw[gray!45, line width=0.15pt] (10.3366,1.6271) -- (10.3660,1.7627);
\draw[gray!45, line width=0.15pt] (10.3660,1.7627) -- (10.3729,1.7941);
\draw[gray!45, line width=0.15pt] (10.3729,1.7941) -- (10.3931,1.8983);
\draw[gray!45, line width=0.15pt] (10.3931,1.8983) -- (10.4197,2.0339);
\draw[gray!45, line width=0.15pt] (10.4197,2.0339) -- (10.4465,2.1695);
\draw[gray!45, line width=0.15pt] (10.4465,2.1695) -- (10.4736,2.3051);
\draw[gray!45, line width=0.15pt] (10.4736,2.3051) -- (10.5011,2.4407);
\draw[gray!45, line width=0.15pt] (10.5011,2.4407) -- (10.5288,2.5763);
\draw[gray!45, line width=0.15pt] (10.5288,2.5763) -- (10.5568,2.7119);
\draw[gray!45, line width=0.15pt] (10.5568,2.7119) -- (10.5763,2.8056);
\draw[gray!45, line width=0.15pt] (10.5763,2.8056) -- (10.5840,2.8475);
\draw[gray!45, line width=0.15pt] (10.5840,2.8475) -- (10.6093,2.9831);
\draw[gray!45, line width=0.15pt] (10.6093,2.9831) -- (10.6348,3.1186);
\draw[gray!45, line width=0.15pt] (10.6348,3.1186) -- (10.6606,3.2542);
\draw[gray!45, line width=0.15pt] (10.6606,3.2542) -- (10.6866,3.3898);
\draw[gray!45, line width=0.15pt] (10.6866,3.3898) -- (10.7129,3.5254);
\draw[gray!45, line width=0.15pt] (10.7129,3.5254) -- (10.7395,3.6610);
\draw[gray!45, line width=0.15pt] (10.7395,3.6610) -- (10.7664,3.7966);
\draw[gray!45, line width=0.15pt] (10.7664,3.7966) -- (10.7797,3.8630);
\draw[gray!45, line width=0.15pt] (10.7797,3.8630) -- (10.7919,3.9322);
\draw[gray!45, line width=0.15pt] (10.7919,3.9322) -- (10.8162,4.0678);
\draw[gray!45, line width=0.15pt] (10.8162,4.0678) -- (10.8406,4.2034);
\draw[gray!45, line width=0.15pt] (10.8406,4.2034) -- (10.8653,4.3390);
\draw[gray!45, line width=0.15pt] (10.8653,4.3390) -- (10.8903,4.4746);
\draw[gray!45, line width=0.15pt] (10.8903,4.4746) -- (10.9155,4.6102);
\draw[gray!45, line width=0.15pt] (10.9155,4.6102) -- (10.9410,4.7458);
\draw[gray!45, line width=0.15pt] (10.9410,4.7458) -- (10.9668,4.8814);
\draw[gray!45, line width=0.15pt] (10.9668,4.8814) -- (10.9831,4.9664);
\draw[gray!45, line width=0.15pt] (10.9831,4.9664) -- (10.9917,5.0169);
\draw[gray!45, line width=0.15pt] (10.9917,5.0169) -- (11.0148,5.1525);
\draw[gray!45, line width=0.15pt] (11.0148,5.1525) -- (11.0382,5.2881);
\draw[gray!45, line width=0.15pt] (11.0382,5.2881) -- (11.0618,5.4237);
\draw[gray!45, line width=0.15pt] (11.0618,5.4237) -- (11.0857,5.5593);
\draw[gray!45, line width=0.15pt] (11.0857,5.5593) -- (11.1099,5.6949);
\draw[gray!45, line width=0.15pt] (11.1099,5.6949) -- (11.1343,5.8305);
\draw[gray!45, line width=0.15pt] (11.1343,5.8305) -- (11.1589,5.9661);
\draw[gray!45, line width=0.15pt] (11.1589,5.9661) -- (11.1839,6.1017);
\draw[gray!45, line width=0.15pt] (11.1839,6.1017) -- (11.1864,6.1157);
\draw[gray!45, line width=0.15pt] (11.1864,6.1157) -- (11.2063,6.2373);
\draw[gray!45, line width=0.15pt] (11.2063,6.2373) -- (11.2287,6.3729);
\draw[gray!45, line width=0.15pt] (11.2287,6.3729) -- (11.2512,6.5085);
\draw[gray!45, line width=0.15pt] (11.2512,6.5085) -- (11.2740,6.6441);
\draw[gray!45, line width=0.15pt] (11.2740,6.6441) -- (11.2971,6.7797);
\draw[gray!45, line width=0.15pt] (11.2971,6.7797) -- (11.3204,6.9153);
\draw[gray!45, line width=0.15pt] (11.3204,6.9153) -- (11.3439,7.0508);
\draw[gray!45, line width=0.15pt] (11.3439,7.0508) -- (11.3677,7.1864);
\draw[gray!45, line width=0.15pt] (11.3677,7.1864) -- (11.3898,7.3110);
\draw[gray!45, line width=0.15pt] (11.3898,7.3110) -- (11.3916,7.3220);
\draw[gray!45, line width=0.15pt] (11.3916,7.3220) -- (11.4129,7.4576);
\draw[gray!45, line width=0.15pt] (11.4129,7.4576) -- (11.4344,7.5932);
\draw[gray!45, line width=0.15pt] (11.4344,7.5932) -- (11.4561,7.7288);
\draw[gray!45, line width=0.15pt] (11.4561,7.7288) -- (11.4781,7.8644);
\draw[gray!45, line width=0.15pt] (11.4781,7.8644) -- (11.5003,8.0000);
\draw[gray!55, line width=0.15pt] (0.2034,6.9593) -- (0.1899,7.0508);
\draw[gray!55, line width=0.15pt] (0.1899,7.0508) -- (0.1699,7.1864);
\draw[gray!55, line width=0.15pt] (0.1699,7.1864) -- (0.1497,7.3220);
\draw[gray!55, line width=0.15pt] (0.1497,7.3220) -- (0.1293,7.4576);
\draw[gray!55, line width=0.15pt] (0.1293,7.4576) -- (0.1088,7.5932);
\draw[gray!55, line width=0.15pt] (0.1088,7.5932) -- (0.0881,7.7288);
\draw[gray!55, line width=0.15pt] (0.0881,7.7288) -- (0.0672,7.8644);
\draw[gray!55, line width=0.15pt] (0.0672,7.8644) -- (0.0461,8.0000);
\draw[gray!55, line width=0.15pt] (0.4068,5.6718) -- (0.4033,5.6949);
\draw[gray!55, line width=0.15pt] (0.4033,5.6949) -- (0.3826,5.8305);
\draw[gray!55, line width=0.15pt] (0.3826,5.8305) -- (0.3617,5.9661);
\draw[gray!55, line width=0.15pt] (0.3617,5.9661) -- (0.3407,6.1017);
\draw[gray!55, line width=0.15pt] (0.3407,6.1017) -- (0.3195,6.2373);
\draw[gray!55, line width=0.15pt] (0.3195,6.2373) -- (0.2981,6.3729);
\draw[gray!55, line width=0.15pt] (0.2981,6.3729) -- (0.2765,6.5085);
\draw[gray!55, line width=0.15pt] (0.2765,6.5085) -- (0.2547,6.6441);
\draw[gray!55, line width=0.15pt] (0.2547,6.6441) -- (0.2328,6.7797);
\draw[gray!55, line width=0.15pt] (0.2328,6.7797) -- (0.2106,6.9153);
\draw[gray!55, line width=0.15pt] (0.2034,6.9593) -- (0.2106,6.9153);
\draw[gray!55, line width=0.15pt] (0.6102,4.4305) -- (0.6032,4.4746);
\draw[gray!55, line width=0.15pt] (0.6032,4.4746) -- (0.5817,4.6102);
\draw[gray!55, line width=0.15pt] (0.5817,4.6102) -- (0.5600,4.7458);
\draw[gray!55, line width=0.15pt] (0.5600,4.7458) -- (0.5381,4.8814);
\draw[gray!55, line width=0.15pt] (0.5381,4.8814) -- (0.5161,5.0169);
\draw[gray!55, line width=0.15pt] (0.5161,5.0169) -- (0.4938,5.1525);
\draw[gray!55, line width=0.15pt] (0.4938,5.1525) -- (0.4714,5.2881);
\draw[gray!55, line width=0.15pt] (0.4714,5.2881) -- (0.4487,5.4237);
\draw[gray!55, line width=0.15pt] (0.4487,5.4237) -- (0.4259,5.5593);
\draw[gray!55, line width=0.15pt] (0.4068,5.6718) -- (0.4259,5.5593);
\draw[gray!55, line width=0.15pt] (0.8136,3.2350) -- (0.8104,3.2542);
\draw[gray!55, line width=0.15pt] (0.8104,3.2542) -- (0.7881,3.3898);
\draw[gray!55, line width=0.15pt] (0.7881,3.3898) -- (0.7655,3.5254);
\draw[gray!55, line width=0.15pt] (0.7655,3.5254) -- (0.7428,3.6610);
\draw[gray!55, line width=0.15pt] (0.7428,3.6610) -- (0.7199,3.7966);
\draw[gray!55, line width=0.15pt] (0.7199,3.7966) -- (0.6968,3.9322);
\draw[gray!55, line width=0.15pt] (0.6968,3.9322) -- (0.6735,4.0678);
\draw[gray!55, line width=0.15pt] (0.6735,4.0678) -- (0.6500,4.2034);
\draw[gray!55, line width=0.15pt] (0.6500,4.2034) -- (0.6263,4.3390);
\draw[gray!55, line width=0.15pt] (0.6102,4.4305) -- (0.6263,4.3390);
\draw[gray!55, line width=0.15pt] (1.0169,2.0856) -- (1.0026,2.1695);
\draw[gray!55, line width=0.15pt] (1.0026,2.1695) -- (0.9792,2.3051);
\draw[gray!55, line width=0.15pt] (0.9792,2.3051) -- (0.9557,2.4407);
\draw[gray!55, line width=0.15pt] (0.9557,2.4407) -- (0.9319,2.5763);
\draw[gray!55, line width=0.15pt] (0.9319,2.5763) -- (0.9080,2.7119);
\draw[gray!55, line width=0.15pt] (0.9080,2.7119) -- (0.8838,2.8475);
\draw[gray!55, line width=0.15pt] (0.8838,2.8475) -- (0.8594,2.9831);
\draw[gray!55, line width=0.15pt] (0.8594,2.9831) -- (0.8348,3.1186);
\draw[gray!55, line width=0.15pt] (0.8136,3.2350) -- (0.8348,3.1186);
\draw[gray!55, line width=0.15pt] (1.2203,0.9821) -- (1.2020,1.0847);
\draw[gray!55, line width=0.15pt] (1.2020,1.0847) -- (1.1776,1.2203);
\draw[gray!55, line width=0.15pt] (1.1776,1.2203) -- (1.1530,1.3559);
\draw[gray!55, line width=0.15pt] (1.1530,1.3559) -- (1.1282,1.4915);
\draw[gray!55, line width=0.15pt] (1.1282,1.4915) -- (1.1032,1.6271);
\draw[gray!55, line width=0.15pt] (1.1032,1.6271) -- (1.0779,1.7627);
\draw[gray!55, line width=0.15pt] (1.0779,1.7627) -- (1.0525,1.8983);
\draw[gray!55, line width=0.15pt] (1.0525,1.8983) -- (1.0268,2.0339);
\draw[gray!55, line width=0.15pt] (1.0169,2.0856) -- (1.0268,2.0339);
\draw[gray!55, line width=0.15pt] (1.4096,0.0000) -- (1.3842,0.1356);
\draw[gray!55, line width=0.15pt] (1.3842,0.1356) -- (1.3585,0.2712);
\draw[gray!55, line width=0.15pt] (1.3585,0.2712) -- (1.3327,0.4068);
\draw[gray!55, line width=0.15pt] (1.3327,0.4068) -- (1.3066,0.5424);
\draw[gray!55, line width=0.15pt] (1.3066,0.5424) -- (1.2802,0.6780);
\draw[gray!55, line width=0.15pt] (1.2802,0.6780) -- (1.2537,0.8136);
\draw[gray!55, line width=0.15pt] (1.2537,0.8136) -- (1.2269,0.9492);
\draw[gray!55, line width=0.15pt] (1.2203,0.9821) -- (1.2269,0.9492);
\draw[gray!55, line width=0.15pt] (10.5909,0.0000) -- (10.6164,0.1356);
\draw[gray!55, line width=0.15pt] (10.6164,0.1356) -- (10.6421,0.2712);
\draw[gray!55, line width=0.15pt] (10.6421,0.2712) -- (10.6679,0.4068);
\draw[gray!55, line width=0.15pt] (10.6679,0.4068) -- (10.6941,0.5424);
\draw[gray!55, line width=0.15pt] (10.6941,0.5424) -- (10.7204,0.6780);
\draw[gray!55, line width=0.15pt] (10.7204,0.6780) -- (10.7470,0.8136);
\draw[gray!55, line width=0.15pt] (10.7470,0.8136) -- (10.7738,0.9492);
\draw[gray!55, line width=0.15pt] (10.7738,0.9492) -- (10.7797,0.9788);
\draw[gray!55, line width=0.15pt] (10.7797,0.9788) -- (10.7986,1.0847);
\draw[gray!55, line width=0.15pt] (10.7986,1.0847) -- (10.8230,1.2203);
\draw[gray!55, line width=0.15pt] (10.8230,1.2203) -- (10.8476,1.3559);
\draw[gray!55, line width=0.15pt] (10.8476,1.3559) -- (10.8724,1.4915);
\draw[gray!55, line width=0.15pt] (10.8724,1.4915) -- (10.8974,1.6271);
\draw[gray!55, line width=0.15pt] (10.8974,1.6271) -- (10.9227,1.7627);
\draw[gray!55, line width=0.15pt] (10.9227,1.7627) -- (10.9482,1.8983);
\draw[gray!55, line width=0.15pt] (10.9482,1.8983) -- (10.9739,2.0339);
\draw[gray!55, line width=0.15pt] (10.9739,2.0339) -- (10.9831,2.0822);
\draw[gray!55, line width=0.15pt] (10.9831,2.0822) -- (10.9980,2.1695);
\draw[gray!55, line width=0.15pt] (10.9980,2.1695) -- (11.0214,2.3051);
\draw[gray!55, line width=0.15pt] (11.0214,2.3051) -- (11.0449,2.4407);
\draw[gray!55, line width=0.15pt] (11.0449,2.4407) -- (11.0687,2.5763);
\draw[gray!55, line width=0.15pt] (11.0687,2.5763) -- (11.0927,2.7119);
\draw[gray!55, line width=0.15pt] (11.0927,2.7119) -- (11.1168,2.8475);
\draw[gray!55, line width=0.15pt] (11.1168,2.8475) -- (11.1412,2.9831);
\draw[gray!55, line width=0.15pt] (11.1412,2.9831) -- (11.1658,3.1186);
\draw[gray!55, line width=0.15pt] (11.1658,3.1186) -- (11.1864,3.2315);
\draw[gray!55, line width=0.15pt] (11.1864,3.2315) -- (11.1902,3.2542);
\draw[gray!55, line width=0.15pt] (11.1902,3.2542) -- (11.2125,3.3898);
\draw[gray!55, line width=0.15pt] (11.2125,3.3898) -- (11.2351,3.5254);
\draw[gray!55, line width=0.15pt] (11.2351,3.5254) -- (11.2578,3.6610);
\draw[gray!55, line width=0.15pt] (11.2578,3.6610) -- (11.2807,3.7966);
\draw[gray!55, line width=0.15pt] (11.2807,3.7966) -- (11.3038,3.9322);
\draw[gray!55, line width=0.15pt] (11.3038,3.9322) -- (11.3271,4.0678);
\draw[gray!55, line width=0.15pt] (11.3271,4.0678) -- (11.3506,4.2034);
\draw[gray!55, line width=0.15pt] (11.3506,4.2034) -- (11.3744,4.3390);
\draw[gray!55, line width=0.15pt] (11.3744,4.3390) -- (11.3898,4.4268);
\draw[gray!55, line width=0.15pt] (11.3898,4.4268) -- (11.3974,4.4746);
\draw[gray!55, line width=0.15pt] (11.3974,4.4746) -- (11.4189,4.6102);
\draw[gray!55, line width=0.15pt] (11.4189,4.6102) -- (11.4406,4.7458);
\draw[gray!55, line width=0.15pt] (11.4406,4.7458) -- (11.4625,4.8814);
\draw[gray!55, line width=0.15pt] (11.4625,4.8814) -- (11.4846,5.0169);
\draw[gray!55, line width=0.15pt] (11.4846,5.0169) -- (11.5068,5.1525);
\draw[gray!55, line width=0.15pt] (11.5068,5.1525) -- (11.5293,5.2881);
\draw[gray!55, line width=0.15pt] (11.5293,5.2881) -- (11.5519,5.4237);
\draw[gray!55, line width=0.15pt] (11.5519,5.4237) -- (11.5748,5.5593);
\draw[gray!55, line width=0.15pt] (11.5748,5.5593) -- (11.5932,5.6680);
\draw[gray!55, line width=0.15pt] (11.5932,5.6680) -- (11.5973,5.6949);
\draw[gray!55, line width=0.15pt] (11.5973,5.6949) -- (11.6180,5.8305);
\draw[gray!55, line width=0.15pt] (11.6180,5.8305) -- (11.6389,5.9661);
\draw[gray!55, line width=0.15pt] (11.6389,5.9661) -- (11.6599,6.1017);
\draw[gray!55, line width=0.15pt] (11.6599,6.1017) -- (11.6811,6.2373);
\draw[gray!55, line width=0.15pt] (11.6811,6.2373) -- (11.7025,6.3729);
\draw[gray!55, line width=0.15pt] (11.7025,6.3729) -- (11.7241,6.5085);
\draw[gray!55, line width=0.15pt] (11.7241,6.5085) -- (11.7459,6.6441);
\draw[gray!55, line width=0.15pt] (11.7459,6.6441) -- (11.7679,6.7797);
\draw[gray!55, line width=0.15pt] (11.7679,6.7797) -- (11.7900,6.9153);
\draw[gray!55, line width=0.15pt] (11.7900,6.9153) -- (11.7966,6.9553);
\draw[gray!55, line width=0.15pt] (11.7966,6.9553) -- (11.8106,7.0508);
\draw[gray!55, line width=0.15pt] (11.8106,7.0508) -- (11.8307,7.1864);
\draw[gray!55, line width=0.15pt] (11.8307,7.1864) -- (11.8509,7.3220);
\draw[gray!55, line width=0.15pt] (11.8509,7.3220) -- (11.8713,7.4576);
\draw[gray!55, line width=0.15pt] (11.8713,7.4576) -- (11.8918,7.5932);
\draw[gray!55, line width=0.15pt] (11.8918,7.5932) -- (11.9125,7.7288);
\draw[gray!55, line width=0.15pt] (11.9125,7.7288) -- (11.9334,7.8644);
\draw[gray!55, line width=0.15pt] (11.9334,7.8644) -- (11.9545,8.0000);
\draw[gray!65, line width=0.15pt] (0.2034,4.5272) -- (0.1911,4.6102);
\draw[gray!65, line width=0.15pt] (0.1911,4.6102) -- (0.1710,4.7458);
\draw[gray!65, line width=0.15pt] (0.1710,4.7458) -- (0.1508,4.8814);
\draw[gray!65, line width=0.15pt] (0.1508,4.8814) -- (0.1304,5.0169);
\draw[gray!65, line width=0.15pt] (0.1304,5.0169) -- (0.1098,5.1525);
\draw[gray!65, line width=0.15pt] (0.1098,5.1525) -- (0.0891,5.2881);
\draw[gray!65, line width=0.15pt] (0.0891,5.2881) -- (0.0683,5.4237);
\draw[gray!65, line width=0.15pt] (0.0683,5.4237) -- (0.0473,5.5593);
\draw[gray!65, line width=0.15pt] (0.0473,5.5593) -- (0.0261,5.6949);
\draw[gray!65, line width=0.15pt] (0.0261,5.6949) -- (0.0047,5.8305);
\draw[gray!65, line width=0.15pt] (0.0000,5.8605) -- (0.0047,5.8305);
\draw[gray!65, line width=0.15pt] (0.4068,3.2398) -- (0.4046,3.2542);
\draw[gray!65, line width=0.15pt] (0.4046,3.2542) -- (0.3838,3.3898);
\draw[gray!65, line width=0.15pt] (0.3838,3.3898) -- (0.3629,3.5254);
\draw[gray!65, line width=0.15pt] (0.3629,3.5254) -- (0.3418,3.6610);
\draw[gray!65, line width=0.15pt] (0.3418,3.6610) -- (0.3206,3.7966);
\draw[gray!65, line width=0.15pt] (0.3206,3.7966) -- (0.2992,3.9322);
\draw[gray!65, line width=0.15pt] (0.2992,3.9322) -- (0.2776,4.0678);
\draw[gray!65, line width=0.15pt] (0.2776,4.0678) -- (0.2559,4.2034);
\draw[gray!65, line width=0.15pt] (0.2559,4.2034) -- (0.2340,4.3390);
\draw[gray!65, line width=0.15pt] (0.2340,4.3390) -- (0.2120,4.4746);
\draw[gray!65, line width=0.15pt] (0.2034,4.5272) -- (0.2120,4.4746);
\draw[gray!65, line width=0.15pt] (0.6102,1.9985) -- (0.6045,2.0339);
\draw[gray!65, line width=0.15pt] (0.6045,2.0339) -- (0.5829,2.1695);
\draw[gray!65, line width=0.15pt] (0.5829,2.1695) -- (0.5612,2.3051);
\draw[gray!65, line width=0.15pt] (0.5612,2.3051) -- (0.5393,2.4407);
\draw[gray!65, line width=0.15pt] (0.5393,2.4407) -- (0.5172,2.5763);
\draw[gray!65, line width=0.15pt] (0.5172,2.5763) -- (0.4950,2.7119);
\draw[gray!65, line width=0.15pt] (0.4950,2.7119) -- (0.4726,2.8475);
\draw[gray!65, line width=0.15pt] (0.4726,2.8475) -- (0.4500,2.9831);
\draw[gray!65, line width=0.15pt] (0.4500,2.9831) -- (0.4273,3.1186);
\draw[gray!65, line width=0.15pt] (0.4068,3.2398) -- (0.4273,3.1186);
\draw[gray!65, line width=0.15pt] (0.8136,0.8031) -- (0.8118,0.8136);
\draw[gray!65, line width=0.15pt] (0.8118,0.8136) -- (0.7894,0.9492);
\draw[gray!65, line width=0.15pt] (0.7894,0.9492) -- (0.7668,1.0847);
\draw[gray!65, line width=0.15pt] (0.7668,1.0847) -- (0.7441,1.2203);
\draw[gray!65, line width=0.15pt] (0.7441,1.2203) -- (0.7211,1.3559);
\draw[gray!65, line width=0.15pt] (0.7211,1.3559) -- (0.6981,1.4915);
\draw[gray!65, line width=0.15pt] (0.6981,1.4915) -- (0.6748,1.6271);
\draw[gray!65, line width=0.15pt] (0.6748,1.6271) -- (0.6514,1.7627);
\draw[gray!65, line width=0.15pt] (0.6514,1.7627) -- (0.6277,1.8983);
\draw[gray!65, line width=0.15pt] (0.6102,1.9985) -- (0.6277,1.8983);
\draw[gray!65, line width=0.15pt] (0.9570,0.0000) -- (0.9332,0.1356);
\draw[gray!65, line width=0.15pt] (0.9332,0.1356) -- (0.9093,0.2712);
\draw[gray!65, line width=0.15pt] (0.9093,0.2712) -- (0.8851,0.4068);
\draw[gray!65, line width=0.15pt] (0.8851,0.4068) -- (0.8608,0.5424);
\draw[gray!65, line width=0.15pt] (0.8608,0.5424) -- (0.8363,0.6780);
\draw[gray!65, line width=0.15pt] (0.8136,0.8031) -- (0.8363,0.6780);
\draw[gray!65, line width=0.15pt] (11.0435,0.0000) -- (11.0673,0.1356);
\draw[gray!65, line width=0.15pt] (11.0673,0.1356) -- (11.0913,0.2712);
\draw[gray!65, line width=0.15pt] (11.0913,0.2712) -- (11.1154,0.4068);
\draw[gray!65, line width=0.15pt] (11.1154,0.4068) -- (11.1397,0.5424);
\draw[gray!65, line width=0.15pt] (11.1397,0.5424) -- (11.1642,0.6780);
\draw[gray!65, line width=0.15pt] (11.1642,0.6780) -- (11.1864,0.8000);
\draw[gray!65, line width=0.15pt] (11.1864,0.8000) -- (11.1887,0.8136);
\draw[gray!65, line width=0.15pt] (11.1887,0.8136) -- (11.2111,0.9492);
\draw[gray!65, line width=0.15pt] (11.2111,0.9492) -- (11.2337,1.0847);
\draw[gray!65, line width=0.15pt] (11.2337,1.0847) -- (11.2565,1.2203);
\draw[gray!65, line width=0.15pt] (11.2565,1.2203) -- (11.2794,1.3559);
\draw[gray!65, line width=0.15pt] (11.2794,1.3559) -- (11.3025,1.4915);
\draw[gray!65, line width=0.15pt] (11.3025,1.4915) -- (11.3257,1.6271);
\draw[gray!65, line width=0.15pt] (11.3257,1.6271) -- (11.3492,1.7627);
\draw[gray!65, line width=0.15pt] (11.3492,1.7627) -- (11.3728,1.8983);
\draw[gray!65, line width=0.15pt] (11.3728,1.8983) -- (11.3898,1.9954);
\draw[gray!65, line width=0.15pt] (11.3898,1.9954) -- (11.3960,2.0339);
\draw[gray!65, line width=0.15pt] (11.3960,2.0339) -- (11.4176,2.1695);
\draw[gray!65, line width=0.15pt] (11.4176,2.1695) -- (11.4393,2.3051);
\draw[gray!65, line width=0.15pt] (11.4393,2.3051) -- (11.4612,2.4407);
\draw[gray!65, line width=0.15pt] (11.4612,2.4407) -- (11.4833,2.5763);
\draw[gray!65, line width=0.15pt] (11.4833,2.5763) -- (11.5055,2.7119);
\draw[gray!65, line width=0.15pt] (11.5055,2.7119) -- (11.5280,2.8475);
\draw[gray!65, line width=0.15pt] (11.5280,2.8475) -- (11.5505,2.9831);
\draw[gray!65, line width=0.15pt] (11.5505,2.9831) -- (11.5733,3.1186);
\draw[gray!65, line width=0.15pt] (11.5733,3.1186) -- (11.5932,3.2366);
\draw[gray!65, line width=0.15pt] (11.5932,3.2366) -- (11.5959,3.2542);
\draw[gray!65, line width=0.15pt] (11.5959,3.2542) -- (11.6167,3.3898);
\draw[gray!65, line width=0.15pt] (11.6167,3.3898) -- (11.6376,3.5254);
\draw[gray!65, line width=0.15pt] (11.6376,3.5254) -- (11.6587,3.6610);
\draw[gray!65, line width=0.15pt] (11.6587,3.6610) -- (11.6800,3.7966);
\draw[gray!65, line width=0.15pt] (11.6800,3.7966) -- (11.7014,3.9322);
\draw[gray!65, line width=0.15pt] (11.7014,3.9322) -- (11.7229,4.0678);
\draw[gray!65, line width=0.15pt] (11.7229,4.0678) -- (11.7446,4.2034);
\draw[gray!65, line width=0.15pt] (11.7446,4.2034) -- (11.7665,4.3390);
\draw[gray!65, line width=0.15pt] (11.7665,4.3390) -- (11.7886,4.4746);
\draw[gray!65, line width=0.15pt] (11.7886,4.4746) -- (11.7966,4.5238);
\draw[gray!65, line width=0.15pt] (11.7966,4.5238) -- (11.8094,4.6102);
\draw[gray!65, line width=0.15pt] (11.8094,4.6102) -- (11.8295,4.7458);
\draw[gray!65, line width=0.15pt] (11.8295,4.7458) -- (11.8497,4.8814);
\draw[gray!65, line width=0.15pt] (11.8497,4.8814) -- (11.8701,5.0169);
\draw[gray!65, line width=0.15pt] (11.8701,5.0169) -- (11.8907,5.1525);
\draw[gray!65, line width=0.15pt] (11.8907,5.1525) -- (11.9114,5.2881);
\draw[gray!65, line width=0.15pt] (11.9114,5.2881) -- (11.9323,5.4237);
\draw[gray!65, line width=0.15pt] (11.9323,5.4237) -- (11.9533,5.5593);
\draw[gray!65, line width=0.15pt] (11.9533,5.5593) -- (11.9745,5.6949);
\draw[gray!65, line width=0.15pt] (11.9745,5.6949) -- (11.9958,5.8305);
\draw[gray!65, line width=0.15pt] (11.9958,5.8305) -- (12.0000,5.8570);
\draw[gray!75, line width=0.15pt] (0.2034,2.3848) -- (0.1951,2.4407);
\draw[gray!75, line width=0.15pt] (0.1951,2.4407) -- (0.1750,2.5763);
\draw[gray!75, line width=0.15pt] (0.1750,2.5763) -- (0.1547,2.7119);
\draw[gray!75, line width=0.15pt] (0.1547,2.7119) -- (0.1343,2.8475);
\draw[gray!75, line width=0.15pt] (0.1343,2.8475) -- (0.1138,2.9831);
\draw[gray!75, line width=0.15pt] (0.1138,2.9831) -- (0.0931,3.1186);
\draw[gray!75, line width=0.15pt] (0.0931,3.1186) -- (0.0723,3.2542);
\draw[gray!75, line width=0.15pt] (0.0723,3.2542) -- (0.0513,3.3898);
\draw[gray!75, line width=0.15pt] (0.0513,3.3898) -- (0.0302,3.5254);
\draw[gray!75, line width=0.15pt] (0.0302,3.5254) -- (0.0090,3.6610);
\draw[gray!75, line width=0.15pt] (0.0000,3.7180) -- (0.0090,3.6610);
\draw[gray!75, line width=0.15pt] (0.4068,1.0974) -- (0.3879,1.2203);
\draw[gray!75, line width=0.15pt] (0.3879,1.2203) -- (0.3669,1.3559);
\draw[gray!75, line width=0.15pt] (0.3669,1.3559) -- (0.3459,1.4915);
\draw[gray!75, line width=0.15pt] (0.3459,1.4915) -- (0.3246,1.6271);
\draw[gray!75, line width=0.15pt] (0.3246,1.6271) -- (0.3033,1.7627);
\draw[gray!75, line width=0.15pt] (0.3033,1.7627) -- (0.2818,1.8983);
\draw[gray!75, line width=0.15pt] (0.2818,1.8983) -- (0.2601,2.0339);
\draw[gray!75, line width=0.15pt] (0.2601,2.0339) -- (0.2383,2.1695);
\draw[gray!75, line width=0.15pt] (0.2383,2.1695) -- (0.2163,2.3051);
\draw[gray!75, line width=0.15pt] (0.2034,2.3848) -- (0.2163,2.3051);
\draw[gray!75, line width=0.15pt] (0.5872,0.0000) -- (0.5654,0.1356);
\draw[gray!75, line width=0.15pt] (0.5654,0.1356) -- (0.5435,0.2712);
\draw[gray!75, line width=0.15pt] (0.5435,0.2712) -- (0.5215,0.4068);
\draw[gray!75, line width=0.15pt] (0.5215,0.4068) -- (0.4993,0.5424);
\draw[gray!75, line width=0.15pt] (0.4993,0.5424) -- (0.4769,0.6780);
\draw[gray!75, line width=0.15pt] (0.4769,0.6780) -- (0.4544,0.8136);
\draw[gray!75, line width=0.15pt] (0.4544,0.8136) -- (0.4317,0.9492);
\draw[gray!75, line width=0.15pt] (0.4317,0.9492) -- (0.4089,1.0847);
\draw[gray!75, line width=0.15pt] (0.4068,1.0974) -- (0.4089,1.0847);
\draw[gray!75, line width=0.15pt] (11.4133,0.0000) -- (11.4350,0.1356);
\draw[gray!75, line width=0.15pt] (11.4350,0.1356) -- (11.4569,0.2712);
\draw[gray!75, line width=0.15pt] (11.4569,0.2712) -- (11.4790,0.4068);
\draw[gray!75, line width=0.15pt] (11.4790,0.4068) -- (11.5012,0.5424);
\draw[gray!75, line width=0.15pt] (11.5012,0.5424) -- (11.5236,0.6780);
\draw[gray!75, line width=0.15pt] (11.5236,0.6780) -- (11.5461,0.8136);
\draw[gray!75, line width=0.15pt] (11.5461,0.8136) -- (11.5688,0.9492);
\draw[gray!75, line width=0.15pt] (11.5688,0.9492) -- (11.5916,1.0847);
\draw[gray!75, line width=0.15pt] (11.5916,1.0847) -- (11.5932,1.0944);
\draw[gray!75, line width=0.15pt] (11.5932,1.0944) -- (11.6126,1.2203);
\draw[gray!75, line width=0.15pt] (11.6126,1.2203) -- (11.6335,1.3559);
\draw[gray!75, line width=0.15pt] (11.6335,1.3559) -- (11.6546,1.4915);
\draw[gray!75, line width=0.15pt] (11.6546,1.4915) -- (11.6758,1.6271);
\draw[gray!75, line width=0.15pt] (11.6758,1.6271) -- (11.6972,1.7627);
\draw[gray!75, line width=0.15pt] (11.6972,1.7627) -- (11.7187,1.8983);
\draw[gray!75, line width=0.15pt] (11.7187,1.8983) -- (11.7404,2.0339);
\draw[gray!75, line width=0.15pt] (11.7404,2.0339) -- (11.7622,2.1695);
\draw[gray!75, line width=0.15pt] (11.7622,2.1695) -- (11.7841,2.3051);
\draw[gray!75, line width=0.15pt] (11.7841,2.3051) -- (11.7966,2.3817);
\draw[gray!75, line width=0.15pt] (11.7966,2.3817) -- (11.8053,2.4407);
\draw[gray!75, line width=0.15pt] (11.8053,2.4407) -- (11.8255,2.5763);
\draw[gray!75, line width=0.15pt] (11.8255,2.5763) -- (11.8458,2.7119);
\draw[gray!75, line width=0.15pt] (11.8458,2.7119) -- (11.8662,2.8475);
\draw[gray!75, line width=0.15pt] (11.8662,2.8475) -- (11.8867,2.9831);
\draw[gray!75, line width=0.15pt] (11.8867,2.9831) -- (11.9074,3.1186);
\draw[gray!75, line width=0.15pt] (11.9074,3.1186) -- (11.9282,3.2542);
\draw[gray!75, line width=0.15pt] (11.9282,3.2542) -- (11.9492,3.3898);
\draw[gray!75, line width=0.15pt] (11.9492,3.3898) -- (11.9703,3.5254);
\draw[gray!75, line width=0.15pt] (11.9703,3.5254) -- (11.9915,3.6610);
\draw[gray!75, line width=0.15pt] (11.9915,3.6610) -- (12.0000,3.7149);
\draw[gray!85, line width=0.15pt] (0.2034,0.4479) -- (0.1894,0.5424);
\draw[gray!85, line width=0.15pt] (0.1894,0.5424) -- (0.1691,0.6780);
\draw[gray!85, line width=0.15pt] (0.1691,0.6780) -- (0.1488,0.8136);
\draw[gray!85, line width=0.15pt] (0.1488,0.8136) -- (0.1283,0.9492);
\draw[gray!85, line width=0.15pt] (0.1283,0.9492) -- (0.1077,1.0847);
\draw[gray!85, line width=0.15pt] (0.1077,1.0847) -- (0.0870,1.2203);
\draw[gray!85, line width=0.15pt] (0.0870,1.2203) -- (0.0662,1.3559);
\draw[gray!85, line width=0.15pt] (0.0662,1.3559) -- (0.0452,1.4915);
\draw[gray!85, line width=0.15pt] (0.0452,1.4915) -- (0.0241,1.6271);
\draw[gray!85, line width=0.15pt] (0.0241,1.6271) -- (0.0029,1.7627);
\draw[gray!85, line width=0.15pt] (0.0000,1.7811) -- (0.0029,1.7627);
\draw[gray!85, line width=0.15pt] (0.2755,0.0000) -- (0.2538,0.1356);
\draw[gray!85, line width=0.15pt] (0.2538,0.1356) -- (0.2320,0.2712);
\draw[gray!85, line width=0.15pt] (0.2320,0.2712) -- (0.2101,0.4068);
\draw[gray!85, line width=0.15pt] (0.2034,0.4479) -- (0.2101,0.4068);
\draw[gray!85, line width=0.15pt] (11.7249,0.0000) -- (11.7466,0.1356);
\draw[gray!85, line width=0.15pt] (11.7466,0.1356) -- (11.7684,0.2712);
\draw[gray!85, line width=0.15pt] (11.7684,0.2712) -- (11.7904,0.4068);
\draw[gray!85, line width=0.15pt] (11.7904,0.4068) -- (11.7966,0.4451);
\draw[gray!85, line width=0.15pt] (11.7966,0.4451) -- (11.8111,0.5424);
\draw[gray!85, line width=0.15pt] (11.8111,0.5424) -- (11.8313,0.6780);
\draw[gray!85, line width=0.15pt] (11.8313,0.6780) -- (11.8516,0.8136);
\draw[gray!85, line width=0.15pt] (11.8516,0.8136) -- (11.8721,0.9492);
\draw[gray!85, line width=0.15pt] (11.8721,0.9492) -- (11.8927,1.0847);
\draw[gray!85, line width=0.15pt] (11.8927,1.0847) -- (11.9134,1.2203);
\draw[gray!85, line width=0.15pt] (11.9134,1.2203) -- (11.9342,1.3559);
\draw[gray!85, line width=0.15pt] (11.9342,1.3559) -- (11.9552,1.4915);
\draw[gray!85, line width=0.15pt] (11.9552,1.4915) -- (11.9763,1.6271);
\draw[gray!85, line width=0.15pt] (11.9763,1.6271) -- (11.9976,1.7627);
\draw[gray!85, line width=0.15pt] (11.9976,1.7627) -- (12.0000,1.7782);
\draw[blue!80!black, line width=1.0pt] (3.6000,2.4000) -- (6.4160,2.8687) -- (4.4337,4.7866) -- (8.0038,5.6872) -- (3.7372,6.5364) -- (3.9108,6.5665) -- (3.8417,6.5527) -- (3.8635,6.5561) -- (3.8562,6.5542) -- (3.8590,6.5541) -- (3.8577,6.5526) -- (3.8594,6.5523) -- (3.8584,6.5497) -- (3.8608,6.5495) -- (3.8598,6.5481) -- (3.8628,6.5475) -- (3.8609,6.5465) -- (3.8628,6.5456) -- (3.8609,6.5441) -- (3.8631,6.5439) -- (3.8623,6.5425) -- (3.8650,6.5418) -- (3.8633,6.5409) -- (3.8668,6.5391) -- (3.8643,6.5380) -- (3.8666,6.5372) -- (3.8653,6.5363) -- (3.8683,6.5345) -- (3.8664,6.5335) -- (3.8683,6.5326) -- (3.8664,6.5310) -- (3.8686,6.5309) -- (3.8669,6.5281) -- (3.8700,6.5281) -- (3.8686,6.5266) -- (3.8704,6.5263) -- (3.8693,6.5237) -- (3.8718,6.5235) -- (3.8708,6.5221) -- (3.8737,6.5214) -- (3.8719,6.5205) -- (3.8755,6.5187) -- (3.8729,6.5176) -- (3.8752,6.5168) -- (3.8728,6.5151) -- (3.8754,6.5150) -- (3.8743,6.5136) -- (3.8774,6.5129) -- (3.8755,6.5120) -- (3.8773,6.5111) -- (3.8756,6.5096) -- (3.8777,6.5093) -- (3.8763,6.5066) -- (3.8790,6.5065) -- (3.8779,6.5051) -- (3.8810,6.5045) -- (3.8791,6.5035) -- (3.8809,6.5026) -- (3.8792,6.5011) -- (3.8813,6.5008) -- (3.8800,6.4981) -- (3.8826,6.4980) -- (3.8816,6.4966) -- (3.8845,6.4959) -- (3.8828,6.4950) -- (3.8862,6.4932) -- (3.8838,6.4921) -- (3.8859,6.4913) -- (3.8837,6.4897) -- (3.8862,6.4895) -- (3.8853,6.4881) -- (3.8880,6.4874) -- (3.8864,6.4865) -- (3.8896,6.4847) -- (3.8875,6.4837) -- (3.8894,6.4828) -- (3.8876,6.4812) -- (3.8897,6.4810) -- (3.8884,6.4783) -- (3.8910,6.4782) -- (3.8901,6.4768) -- (3.8929,6.4761) -- (3.8912,6.4752) -- (3.8945,6.4733) -- (3.8923,6.4723) -- (3.8943,6.4714) -- (3.8924,6.4699) -- (3.8945,6.4697) -- (3.8933,6.4670) -- (3.8959,6.4668) -- (3.8949,6.4654) -- (3.8977,6.4647) -- (3.8961,6.4638) -- (3.8992,6.4620) -- (3.8972,6.4610) -- (3.9008,6.4592) -- (3.8983,6.4581) -- (3.9004,6.4573) -- (3.8984,6.4557) -- (3.9006,6.4555) -- (3.8992,6.4528) -- (3.9020,6.4527) -- (3.9010,6.4512) -- (3.9038,6.4506) -- (3.9022,6.4496) -- (3.9053,6.4478) -- (3.9033,6.4468) -- (3.9069,6.4450) -- (3.9045,6.4439) -- (3.9065,6.4431) -- (3.9046,6.4415) -- (3.9067,6.4413) -- (3.9055,6.4386) -- (3.9080,6.4385) -- (3.9064,6.4357) -- (3.9093,6.4356) -- (3.9083,6.4342) -- (3.9111,6.4335) -- (3.9095,6.4326) -- (3.9126,6.4307) -- (3.9107,6.4297) -- (3.9141,6.4279) -- (3.9118,6.4269) -- (3.9137,6.4260) -- (3.9121,6.4245) -- (7.8436,5.9179) -- (7.8563,5.9150) -- (7.8556,5.9149) -- (7.8561,5.9145) -- (7.8553,5.9144) -- (7.8559,5.9140) -- (7.8550,5.9139) -- (7.8557,5.9135) -- (7.8547,5.9134) -- (7.8555,5.9129) -- (7.8544,5.9129) -- (7.8553,5.9124) -- (7.8541,5.9124) -- (7.8552,5.9119) -- (7.8538,5.9119) -- (7.8550,5.9114) -- (7.8535,5.9115) -- (7.8549,5.9109) -- (7.8531,5.9110) -- (7.8547,5.9104) -- (7.8528,5.9105) -- (7.8546,5.9098) -- (7.8524,5.9100) -- (7.8545,5.9093) -- (7.8533,5.9094) -- (7.8525,5.9052) -- (7.8513,5.9053) -- (7.8495,5.8969) -- (7.8473,5.8972) -- (7.8472,5.8951) -- (7.8454,5.8952) -- (7.8470,5.8946) -- (7.8451,5.8947) -- (7.8469,5.8940) -- (7.8448,5.8942) -- (7.8467,5.8935) -- (7.8444,5.8938) -- (7.8466,5.8930) -- (7.8441,5.8933) -- (7.8465,5.8925) -- (7.8451,5.8927) -- (7.8431,5.8843) -- (7.8412,5.8846) -- (7.8404,5.8804) -- (7.8380,5.8807) -- (7.8403,5.8799) -- (7.8377,5.8802) -- (7.8401,5.8794) -- (7.8374,5.8797) -- (7.8399,5.8789) -- (7.8370,5.8793) -- (7.8398,5.8784) -- (7.8382,5.8786) -- (7.8359,5.8704) -- (7.8328,5.8708) -- (7.8357,5.8699) -- (7.8325,5.8704) -- (7.8356,5.8694) -- (7.8322,5.8699) -- (7.8354,5.8689) -- (7.8318,5.8694) -- (7.8352,5.8684) -- (7.8334,5.8687) -- (7.8284,5.8525) -- (7.8230,5.8534) -- (7.8281,5.8520) -- (7.8227,5.8529) -- (7.8279,5.8515) -- (7.8225,5.8525) -- (7.8277,5.8510) -- (7.8222,5.8520) -- (7.8275,5.8505) -- (7.8219,5.8515) -- (7.8273,5.8501) -- (7.8217,5.8510) -- (7.8270,5.8496) -- (7.8214,5.8505) -- (7.8268,5.8491) -- (7.8212,5.8500) -- (7.8266,5.8486) -- (7.8209,5.8496) -- (7.8263,5.8481) -- (7.8207,5.8491) -- (7.8261,5.8476) -- (7.8205,5.8486) -- (7.8259,5.8471) -- (7.8202,5.8481) -- (7.8256,5.8466) -- (7.8200,5.8476) -- (7.8254,5.8462) -- (7.8197,5.8471) -- (7.8251,5.8457) -- (7.8195,5.8466) -- (7.8249,5.8452) -- (7.8193,5.8462) -- (7.8246,5.8447) -- (7.8191,5.8457) -- (7.8244,5.8442) -- (7.8188,5.8452) -- (7.8241,5.8438) -- (7.8186,5.8447) -- (7.8239,5.8433) -- (7.8184,5.8442) -- (7.8236,5.8428) -- (7.8182,5.8437) -- (7.8233,5.8423) -- (7.8180,5.8432) -- (7.8231,5.8418) -- (7.8178,5.8427) -- (7.8228,5.8414) -- (7.8176,5.8422) -- (7.8225,5.8409) -- (7.8173,5.8418) -- (7.8223,5.8404) -- (7.8171,5.8413) -- (7.8220,5.8399) -- (7.8169,5.8408) -- (7.8217,5.8395) -- (7.8167,5.8403) -- (7.8214,5.8390) -- (7.8165,5.8398) -- (7.8212,5.8385) -- (7.8163,5.8393) -- (7.8209,5.8380) -- (7.8161,5.8388) -- (7.8206,5.8376) -- (7.8160,5.8383) -- (7.8203,5.8371) -- (7.8158,5.8378) -- (7.8200,5.8366) -- (7.8156,5.8373) -- (7.8197,5.8362) -- (7.8154,5.8368) -- (7.8195,5.8357) -- (7.8152,5.8364) -- (7.8192,5.8352) -- (7.8150,5.8359) -- (7.8189,5.8347) -- (7.8148,5.8354) -- (7.8186,5.8343) -- (7.8146,5.8349) -- (7.8183,5.8338) -- (7.8144,5.8344) -- (7.8180,5.8333) -- (7.8143,5.8339) -- (7.8177,5.8329) -- (7.8141,5.8334) -- (7.8174,5.8324) -- (7.8139,5.8329) -- (7.8172,5.8319) -- (7.8137,5.8324) -- (7.8169,5.8315) -- (7.8135,5.8319) -- (7.8166,5.8310) -- (7.8133,5.8315) -- (7.8163,5.8305) -- (7.8131,5.8310) -- (7.8160,5.8301) -- (7.8130,5.8305) -- (7.8157,5.8296) -- (7.8128,5.8300) -- (7.8154,5.8291) -- (7.8126,5.8295) -- (7.8151,5.8287) -- (7.8124,5.8290) -- (7.8149,5.8282) -- (7.8122,5.8285) -- (7.8146,5.8277) -- (7.8120,5.8280) -- (7.8143,5.8273) -- (7.8118,5.8275) -- (7.8140,5.8268) -- (7.8116,5.8271) -- (7.8137,5.8263) -- (7.8114,5.8266) -- (7.8134,5.8259) -- (7.8113,5.8261) -- (7.8132,5.8254) -- (7.8111,5.8256) -- (7.8129,5.8249) -- (7.8109,5.8251) -- (7.8126,5.8245) -- (7.8107,5.8246) -- (7.8123,5.8240) -- (7.8105,5.8241) -- (7.8120,5.8235) -- (7.8103,5.8237) -- (7.8118,5.8231) -- (7.8101,5.8232) -- (7.8115,5.8226) -- (7.8099,5.8227) -- (7.8112,5.8221) -- (7.8097,5.8222) -- (7.8110,5.8217) -- (7.8095,5.8217) -- (7.8107,5.8212) -- (7.8093,5.8212) -- (7.8104,5.8207) -- (7.8091,5.8208) -- (7.8101,5.8203) -- (7.8089,5.8203) -- (7.8099,5.8198) -- (7.8087,5.8198) -- (7.8096,5.8193) -- (7.8085,5.8193) -- (7.8093,5.8189) -- (7.8083,5.8188) -- (7.8091,5.8184) -- (7.8081,5.8184) -- (7.8088,5.8179) -- (7.8078,5.8179) -- (7.8086,5.8175) -- (7.8076,5.8174) -- (7.8083,5.8170) -- (7.8074,5.8169) -- (7.8080,5.8165) -- (7.8072,5.8164) -- (7.8078,5.8161) -- (7.8070,5.8160) -- (7.8075,5.8156) -- (7.8068,5.8155) -- (7.8073,5.8151) -- (7.8066,5.8150) -- (7.8070,5.8146) -- (7.8063,5.8145) -- (7.8068,5.8142) -- (7.8061,5.8141) -- (7.8065,5.8137) -- (7.8059,5.8136) -- (7.8063,5.8132) -- (7.8057,5.8131) -- (7.8060,5.8128) -- (7.8055,5.8126) -- (7.8058,5.8123) -- (7.8052,5.8122) -- (7.8055,5.8118) -- (7.8050,5.8117) -- (7.8055,5.8110) -- (7.8043,5.8110) -- (7.8052,5.8106) -- (7.8041,5.8106) -- (7.8050,5.8101) -- (7.8039,5.8101) -- (7.8047,5.8097) -- (7.8037,5.8096) -- (7.8044,5.8092) -- (7.8035,5.8091) -- (7.8042,5.8087) -- (7.8033,5.8087) -- (7.8039,5.8083) -- (7.8031,5.8082) -- (7.8036,5.8078) -- (7.8029,5.8077) -- (7.8034,5.8073) -- (7.8027,5.8072) -- (7.8031,5.8069) -- (7.8025,5.8067) -- (7.8029,5.8064) -- (7.8023,5.8063) -- (7.8026,5.8059) -- (7.8021,5.8058) -- (7.8024,5.8055) -- (7.8018,5.8053) -- (7.8021,5.8050) -- (7.8016,5.8049) -- (7.8021,5.8042) -- (7.8009,5.8042) -- (7.8018,5.8038) -- (7.8008,5.8037) -- (7.8015,5.8033) -- (7.8006,5.8033) -- (7.8013,5.8028) -- (7.8004,5.8028) -- (7.8010,5.8024) -- (7.8002,5.8023) -- (7.8007,5.8019) -- (7.8000,5.8018) -- (7.8005,5.8015) -- (7.7997,5.8014) -- (7.8002,5.8010) -- (7.7995,5.8009) -- (7.8000,5.8005) -- (7.7993,5.8004) -- (7.7997,5.8001) -- (7.7991,5.7999) -- (7.7994,5.7996) -- (7.7989,5.7995) -- (7.7992,5.7991) -- (7.7987,5.7990) -- (7.7992,5.7984) -- (7.7980,5.7984) -- (7.7989,5.7979) -- (7.7978,5.7979) -- (7.7986,5.7974) -- (7.7976,5.7974) -- (7.7983,5.7970) -- (7.7974,5.7969) -- (7.7981,5.7965) -- (7.7972,5.7965) -- (7.7978,5.7961) -- (7.7970,5.7960) -- (7.7975,5.7956) -- (7.7968,5.7955) -- (7.7973,5.7951) -- (7.7966,5.7950) -- (7.7970,5.7947) -- (7.7964,5.7946) -- (7.7967,5.7942) -- (7.7962,5.7941) -- (7.7965,5.7938) -- (7.7960,5.7936) -- (7.7965,5.7930) -- (7.7953,5.7930) -- (7.7962,5.7925) -- (7.7952,5.7925) -- (7.7959,5.7921) -- (7.7950,5.7920) -- (7.7956,5.7916) -- (7.7948,5.7916) -- (7.7953,5.7912) -- (7.7946,5.7911) -- (7.7951,5.7907) -- (7.7944,5.7906) -- (7.7948,5.7903) -- (7.7942,5.7901) -- (7.7946,5.7898) -- (7.7940,5.7897) -- (7.7943,5.7893) -- (7.7938,5.7892) -- (7.7940,5.7889) -- (7.7936,5.7887) -- (7.7940,5.7881) -- (7.7929,5.7881) -- (7.7937,5.7876) -- (7.7927,5.7876) -- (7.7934,5.7872) -- (7.7925,5.7871) -- (7.7932,5.7867) -- (7.7924,5.7867) -- (7.7929,5.7863) -- (7.7922,5.7862) -- (7.7926,5.7858) -- (7.7920,5.7857) -- (7.7923,5.7854) -- (7.7918,5.7852) -- (7.7921,5.7849) -- (7.7916,5.7848) -- (7.7918,5.7845) -- (7.7913,5.7843) -- (7.7918,5.7837) -- (7.7907,5.7837) -- (7.7915,5.7832) -- (7.7905,5.7832) -- (7.7912,5.7828) -- (7.7903,5.7827) -- (7.7909,5.7823) -- (7.7902,5.7823) -- (7.7907,5.7819) -- (7.7900,5.7818) -- (7.7904,5.7814) -- (7.7898,5.7813) -- (7.7901,5.7810) -- (7.7896,5.7808) -- (7.7899,5.7805) -- (7.7894,5.7804) -- (7.7896,5.7801) -- (7.7887,5.7797) -- (7.7895,5.7793) -- (7.7885,5.7793) -- (7.7892,5.7788) -- (7.7883,5.7788) -- (7.7890,5.7784) -- (7.7882,5.7783) -- (7.7887,5.7779) -- (7.7880,5.7778) -- (7.7884,5.7775) -- (7.7878,5.7774) -- (7.7881,5.7770) -- (7.7876,5.7769) -- (7.7879,5.7766) -- (7.7874,5.7764) -- (7.7876,5.7761) -- (7.7867,5.7758) -- (7.7876,5.7754) -- (7.7866,5.7753) -- (7.7873,5.7749) -- (7.7864,5.7749) -- (7.7870,5.7745) -- (7.7862,5.7744) -- (7.7867,5.7740) -- (7.7860,5.7739) -- (7.7864,5.7736) -- (7.7858,5.7735) -- (7.7862,5.7731) -- (7.7856,5.7730) -- (7.7859,5.7727) -- (7.7854,5.7725) -- (7.7859,5.7719) -- (7.7848,5.7719) -- (7.7856,5.7715) -- (7.7846,5.7714) -- (7.7853,5.7710) -- (7.7844,5.7709) -- (7.7850,5.7706) -- (7.7843,5.7705) -- (7.7847,5.7701) -- (7.7841,5.7700) -- (7.7844,5.7697) -- (7.7839,5.7695) -- (7.7842,5.7692) -- (7.7837,5.7691) -- (7.7841,5.7685) -- (7.7830,5.7684) -- (7.7838,5.7680) -- (7.7829,5.7680) -- (7.7835,5.7676) -- (7.7827,5.7675) -- (7.7832,5.7671) -- (7.7825,5.7670) -- (7.7830,5.7667) -- (7.7823,5.7666) -- (7.7827,5.7662) -- (7.7821,5.7661) -- (7.7824,5.7658) -- (7.7819,5.7656) -- (7.7824,5.7650) -- (7.7813,5.7650) -- (7.7821,5.7646) -- (7.7812,5.7645) -- (7.7818,5.7641) -- (7.7810,5.7641) -- (7.7815,5.7637) -- (7.7808,5.7636) -- (7.7812,5.7632) -- (7.7806,5.7631) -- (7.7809,5.7628) -- (7.7804,5.7627) -- (7.7807,5.7623) -- (7.7798,5.7620) -- (7.7806,5.7616) -- (7.7796,5.7616) -- (7.7803,5.7611) -- (7.7794,5.7611) -- (7.7800,5.7607) -- (7.7793,5.7606) -- (7.7797,5.7603) -- (7.7791,5.7601) -- (7.7795,5.7598) -- (7.7789,5.7597) -- (7.7792,5.7594) -- (7.7787,5.7592) -- (7.7791,5.7586) -- (7.7781,5.7586) -- (7.7788,5.7582) -- (7.7779,5.7581) -- (7.7785,5.7577) -- (7.7778,5.7577) -- (7.7782,5.7573) -- (7.7776,5.7572) -- (7.7780,5.7568) -- (7.7774,5.7567) -- (7.7777,5.7564) -- (7.7772,5.7563) -- (7.7777,5.7556) -- (7.7766,5.7556) -- (7.7773,5.7552) -- (7.7764,5.7552) -- (7.7770,5.7548) -- (7.7763,5.7547) -- (7.7767,5.7543) -- (7.7761,5.7542) -- (7.7765,5.7539) -- (7.7759,5.7538) -- (7.7762,5.7534) -- (7.7757,5.7533) -- (7.7761,5.7527) -- (7.7751,5.7527) -- (7.7758,5.7523) -- (7.7749,5.7522) -- (7.7755,5.7518) -- (7.7748,5.7517) -- (7.7752,5.7514) -- (7.7746,5.7513) -- (7.7750,5.7509) -- (7.7744,5.7508) -- (7.7747,5.7505) -- (7.7738,5.7502) -- (7.7746,5.7497) -- (7.7736,5.7497) -- (7.7743,5.7493) -- (7.7735,5.7492) -- (7.7740,5.7489) -- (7.7733,5.7488) -- (7.7737,5.7484) -- (7.7731,5.7483) -- (7.7734,5.7480) -- (7.7729,5.7478) -- (7.7734,5.7472) -- (7.7723,5.7472) -- (7.7731,5.7468) -- (7.7722,5.7468) -- (7.7728,5.7464) -- (7.7720,5.7463) -- (7.7725,5.7459) -- (7.7718,5.7458) -- (7.7722,5.7455) -- (7.7717,5.7454) -- (7.7719,5.7450) -- (7.7710,5.7448) -- (7.7719,5.7443) -- (7.7709,5.7443) -- (7.7715,5.7439) -- (7.7707,5.7438) -- (7.7712,5.7434) -- (7.7706,5.7433) -- (7.7710,5.7430) -- (7.7704,5.7429) -- (7.7707,5.7426) -- (7.7702,5.7424) -- (7.7706,5.7418) -- (7.7696,5.7418) -- (7.7703,5.7414) -- (7.7694,5.7413) -- (7.7700,5.7409) -- (7.7693,5.7409) -- (7.7697,5.7405) -- (7.7691,5.7404) -- (7.7694,5.7401) -- (7.7689,5.7399) -- (7.7694,5.7393) -- (7.7683,5.7393) -- (7.7691,5.7389) -- (7.7682,5.7388) -- (7.7687,5.7385) -- (7.7680,5.7384) -- (7.7684,5.7380) -- (7.7678,5.7379) -- (7.7682,5.7376) -- (7.7677,5.7375) -- (7.7681,5.7368) -- (7.7670,5.7368) -- (7.7678,5.7364) -- (7.7669,5.7364) -- (7.7675,5.7360) -- (7.7668,5.7359) -- (7.7672,5.7356) -- (7.7666,5.7354) -- (7.7669,5.7351) -- (7.7664,5.7350) -- (7.7669,5.7344) -- (7.7658,5.7344) -- (7.7665,5.7340) -- (7.7657,5.7339) -- (7.7662,5.7335) -- (7.7655,5.7334) -- (7.7659,5.7331) -- (7.7653,5.7330) -- (7.7656,5.7327) -- (7.7652,5.7325) -- (7.7656,5.7319) -- (7.7646,5.7319) -- (7.7653,5.7315) -- (7.7644,5.7314) -- (7.7649,5.7311) -- (7.7643,5.7310) -- (7.7646,5.7306) -- (7.7641,5.7305) -- (7.7644,5.7302) -- (7.7634,5.7299) -- (7.7643,5.7295) -- (7.7633,5.7294) -- (7.7640,5.7290) -- (7.7632,5.7290) -- (7.7637,5.7286) -- (7.7630,5.7285) -- (7.7634,5.7282) -- (7.7629,5.7280) -- (7.7631,5.7277) -- (7.7622,5.7274) -- (7.7630,5.7270) -- (7.7621,5.7270) -- (7.7627,5.7266) -- (7.7620,5.7265) -- (7.7624,5.7261) -- (7.7618,5.7260) -- (7.7621,5.7257) -- (7.7616,5.7256) -- (7.7620,5.7250) -- (7.7610,5.7250) -- (7.7617,5.7246) -- (7.7609,5.7245) -- (7.7614,5.7241) -- (7.7607,5.7240) -- (7.7611,5.7237) -- (7.7606,5.7236) -- (7.7608,5.7233) -- (7.7600,5.7230) -- (7.7607,5.7225) -- (7.7598,5.7225) -- (7.7604,5.7221) -- (7.7597,5.7220) -- (7.7601,5.7217) -- (7.7595,5.7216) -- (7.7598,5.7213) -- (7.7589,5.7210) -- (7.7597,5.7205) -- (7.7588,5.7205) -- (7.7594,5.7201) -- (7.7586,5.7200) -- (7.7591,5.7197) -- (7.7585,5.7196) -- (7.7588,5.7192) -- (7.7583,5.7191) -- (7.7587,5.7185) -- (7.7577,5.7185) -- (7.7584,5.7181) -- (7.7576,5.7180) -- (7.7581,5.7177) -- (7.7574,5.7176) -- (7.7578,5.7172) -- (7.7573,5.7171) -- (7.7577,5.7165) -- (7.7567,5.7165) -- (7.7574,5.7161) -- (7.7566,5.7160) -- (7.7571,5.7157) -- (7.7564,5.7156) -- (7.7568,5.7152) -- (7.7562,5.7151) -- (7.7565,5.7148) -- (7.7556,5.7145) -- (7.7564,5.7141) -- (7.7555,5.7140) -- (7.7561,5.7137) -- (7.7554,5.7136) -- (7.7558,5.7133) -- (7.7552,5.7131) -- (7.7555,5.7128) -- (7.7546,5.7125) -- (7.7554,5.7121) -- (7.7545,5.7121) -- (7.7551,5.7117) -- (7.7544,5.7116) -- (7.7547,5.7113) -- (7.7542,5.7111) -- (7.7545,5.7108) -- (7.7536,5.7105) -- (7.7544,5.7101) -- (7.7535,5.7101) -- (7.7540,5.7097) -- (7.7533,5.7096) -- (7.7537,5.7093) -- (7.7532,5.7091) -- (7.7534,5.7088) -- (7.7526,5.7085) -- (7.7533,5.7081) -- (7.7525,5.7081) -- (7.7530,5.7077) -- (7.7523,5.7076) -- (7.7527,5.7073) -- (7.7522,5.7072) -- (7.7524,5.7069) -- (7.7516,5.7066) -- (7.7523,5.7061) -- (7.7515,5.7061) -- (7.7520,5.7057) -- (7.7513,5.7056) -- (7.7517,5.7053) -- (7.7512,5.7052) -- (7.7516,5.7046) -- (7.7506,5.7046) -- (7.7513,5.7042) -- (7.7505,5.7041) -- (7.7509,5.7037) -- (7.7503,5.7036) -- (7.7506,5.7033) -- (7.7502,5.7032) -- (7.7506,5.7026) -- (7.7496,5.7026) -- (7.7502,5.7022) -- (7.7495,5.7021) -- (7.7499,5.7018) -- (7.7493,5.7017) -- (7.7496,5.7014) -- (7.7487,5.7011) -- (7.7495,5.7006) -- (7.7486,5.7006) -- (7.7492,5.7002) -- (7.7485,5.7001) -- (7.7489,5.6998) -- (7.7483,5.6997) -- (7.7486,5.6994) -- (7.7477,5.6991) -- (7.7485,5.6987) -- (7.7476,5.6986) -- (7.7481,5.6983) -- (7.7475,5.6982) -- (7.7478,5.6978) -- (7.7473,5.6977) -- (7.7478,5.6971) -- (7.7467,5.6971) -- (7.7474,5.6967) -- (7.7466,5.6966) -- (7.7471,5.6963) -- (7.7465,5.6962) -- (7.7468,5.6959) -- (7.7459,5.6956) -- (7.7467,5.6952) -- (7.7458,5.6951) -- (7.7464,5.6948) -- (7.7457,5.6947) -- (7.7460,5.6943) -- (7.7455,5.6942) -- (7.7458,5.6939) -- (7.7449,5.6936) -- (7.7456,5.6932) -- (7.7448,5.6931) -- (7.7453,5.6928) -- (7.7447,5.6927) -- (7.7450,5.6924) -- (7.7445,5.6922) -- (7.7449,5.6917) -- (7.7440,5.6916) -- (7.7446,5.6913) -- (7.7439,5.6912) -- (7.7443,5.6908) -- (7.7437,5.6907) -- (7.7440,5.6904) -- (7.7431,5.6901) -- (7.7438,5.6897) -- (7.7430,5.6897) -- (7.7435,5.6893) -- (7.7429,5.6892) -- (7.7432,5.6889) -- (7.7427,5.6888) -- (7.7431,5.6882) -- (7.7422,5.6881) -- (7.7428,5.6878) -- (7.7421,5.6877) -- (7.7424,5.6874) -- (7.7419,5.6872) -- (7.7424,5.6866) -- (7.7413,5.6866) -- (7.7420,5.6862) -- (7.7412,5.6862) -- (7.7417,5.6858) -- (7.7411,5.6857) -- (7.7414,5.6854) -- (7.7405,5.6851) -- (7.7413,5.6847) -- (7.7404,5.6847) -- (7.7409,5.6843) -- (7.7403,5.6842) -- (7.7406,5.6839) -- (7.7401,5.6838) -- (7.7405,5.6832) -- (7.7396,5.6832) -- (7.7402,5.6828) -- (7.7395,5.6827) -- (7.7399,5.6824) -- (7.7393,5.6822) -- (7.7398,5.6817) -- (7.7388,5.6817) -- (7.7394,5.6813) -- (7.7387,5.6812) -- (7.7391,5.6808) -- (7.7385,5.6807) -- (7.7388,5.6804) -- (7.7379,5.6801) -- (7.7387,5.6797) -- (7.7379,5.6797) -- (7.7384,5.6793) -- (7.7377,5.6792) -- (7.7380,5.6789) -- (7.7371,5.6786) -- (7.7379,5.6782) -- (7.7370,5.6782) -- (7.7376,5.6778) -- (7.7369,5.6777) -- (7.7373,5.6774) -- (7.7368,5.6773) -- (7.7372,5.6767) -- (7.7362,5.6767) -- (7.7368,5.6763) -- (7.7361,5.6762) -- (7.7365,5.6759) -- (7.7360,5.6758) -- (7.7364,5.6752) -- (7.7354,5.6752) -- (7.7361,5.6748) -- (7.7353,5.6747) -- (7.7357,5.6744) -- (7.7352,5.6743) -- (7.7354,5.6740) -- (7.7346,5.6737) -- (7.7353,5.6733) -- (7.7345,5.6732) -- (7.7350,5.6729) -- (7.7344,5.6728) -- (7.7347,5.6725) -- (7.7338,5.6722) -- (7.7345,5.6718) -- (7.7338,5.6717) -- (7.7342,5.6714) -- (7.7336,5.6712) -- (7.7339,5.6709) -- (7.7330,5.6707) -- (7.7338,5.6703) -- (7.7330,5.6702) -- (7.7334,5.6699) -- (7.7328,5.6697) -- (7.7331,5.6694) -- (7.7322,5.6692) -- (7.7330,5.6687) -- (7.7322,5.6687) -- (7.7327,5.6683) -- (7.7321,5.6682) -- (7.7323,5.6679) -- (7.7315,5.6677) -- (7.7322,5.6672) -- (7.7314,5.6672) -- (7.7319,5.6668) -- (7.7313,5.6667) -- (7.7316,5.6664) -- (7.7307,5.6662) -- (7.7314,5.6657) -- (7.7306,5.6657) -- (7.7311,5.6653) -- (7.7305,5.6652) -- (7.7308,5.6649) -- (7.7299,5.6647) -- (7.7307,5.6643) -- (7.7298,5.6642) -- (7.7303,5.6639) -- (7.7297,5.6637) -- (7.7300,5.6634) -- (7.7291,5.6632) -- (7.7299,5.6628) -- (7.7291,5.6627) -- (7.7295,5.6624) -- (7.7290,5.6623) -- (7.7292,5.6620) -- (7.7284,5.6617) -- (7.7291,5.6613) -- (7.7283,5.6612) -- (7.7287,5.6609) -- (7.7282,5.6608) -- (7.7284,5.6605) -- (7.7276,5.6602) -- (7.7283,5.6598) -- (7.7275,5.6597) -- (7.7280,5.6594) -- (7.7274,5.6593) -- (7.7279,5.6587) -- (7.7269,5.6587) -- (7.7275,5.6583) -- (7.7268,5.6582) -- (7.7272,5.6579) -- (7.7267,5.6578) -- (7.7271,5.6572) -- (7.7261,5.6572) -- (7.7267,5.6568) -- (7.7260,5.6567) -- (7.7264,5.6564) -- (7.7259,5.6563) -- (7.7263,5.6557) -- (7.7254,5.6557) -- (7.7259,5.6553) -- (7.7253,5.6552) -- (7.7256,5.6549) -- (7.7251,5.6548) -- (7.7255,5.6542) -- (7.7246,5.6542) -- (7.7251,5.6538) -- (7.7245,5.6537) -- (7.7248,5.6534) -- (7.7239,5.6532) -- (7.7247,5.6528) -- (7.7239,5.6527) -- (7.7243,5.6524) -- (7.7238,5.6523) -- (7.7240,5.6520) -- (7.7232,5.6517) -- (7.7239,5.6513) -- (7.7231,5.6512) -- (7.7235,5.6509) -- (7.7230,5.6508) -- (7.7234,5.6502) -- (7.7225,5.6502) -- (7.7231,5.6498) -- (7.7224,5.6497) -- (7.7227,5.6494) -- (7.7222,5.6493) -- (7.7226,5.6487) -- (7.7217,5.6487) -- (7.7222,5.6483) -- (7.7216,5.6482) -- (7.7219,5.6479) -- (7.7211,5.6477) -- (7.7218,5.6473) -- (7.7210,5.6472) -- (7.7214,5.6469) -- (7.7209,5.6468) -- (7.7211,5.6465) -- (7.7203,5.6462) -- (7.7210,5.6458) -- (7.7203,5.6457) -- (7.7206,5.6454) -- (7.7201,5.6453) -- (7.7205,5.6447) -- (7.7196,5.6447) -- (7.7202,5.6443) -- (7.7195,5.6442) -- (7.7198,5.6439) -- (7.7190,5.6437) -- (7.7197,5.6433) -- (7.7189,5.6432) -- (7.7194,5.6429) -- (7.7188,5.6428) -- (7.7190,5.6425) -- (7.7183,5.6422) -- (7.7189,5.6418) -- (7.7182,5.6417) -- (7.7185,5.6414) -- (7.7181,5.6413) -- (7.7184,5.6407) -- (7.7176,5.6407) -- (7.7181,5.6403) -- (7.7175,5.6402) -- (7.7177,5.6399) -- (7.7169,5.6397) -- (7.7176,5.6393) -- (7.7168,5.6392) -- (7.7172,5.6389) -- (7.7167,5.6388) -- (7.7172,5.6382) -- (7.7162,5.6382) -- (7.7168,5.6378) -- (7.7161,5.6377) -- (7.7164,5.6374) -- (7.7155,5.6372) -- (7.7163,5.6367) -- (7.7155,5.6367) -- (7.7159,5.6364) -- (7.7154,5.6363) -- (7.7159,5.6357) -- (7.7149,5.6357) -- (7.7155,5.6353) -- (7.7148,5.6352) -- (7.7151,5.6349) -- (7.7147,5.6348) -- (7.7150,5.6342) -- (7.7142,5.6342) -- (7.7146,5.6338) -- (7.7141,5.6337) -- (7.7143,5.6335) -- (7.7135,5.6332) -- (7.7142,5.6328) -- (7.7135,5.6327) -- (7.7138,5.6324) -- (7.7133,5.6323) -- (7.7137,5.6317) -- (7.7129,5.6317) -- (7.7133,5.6313) -- (7.7128,5.6312) -- (7.7130,5.6310) -- (7.7122,5.6307) -- (7.7128,5.6303) -- (7.7122,5.6302) -- (7.7125,5.6299) -- (7.7120,5.6298) -- (7.7124,5.6292) -- (7.7115,5.6292) -- (7.7120,5.6288) -- (7.7114,5.6287) -- (7.7117,5.6285) -- (7.7109,5.6282) -- (7.7115,5.6278) -- (7.7109,5.6277) -- (7.7112,5.6274) -- (7.7103,5.6272) -- (7.7111,5.6267) -- (7.7102,5.6267) -- (7.7107,5.6264) -- (7.7101,5.6262) -- (7.7106,5.6257) -- (7.7096,5.6257) -- (7.7102,5.6253) -- (7.7096,5.6252) -- (7.7099,5.6249) -- (7.7090,5.6247) -- (7.7097,5.6243) -- (7.7090,5.6242) -- (7.7094,5.6239) -- (7.7089,5.6238) -- (7.7093,5.6232) -- (7.7084,5.6232) -- (7.7089,5.6228) -- (7.7083,5.6227) -- (7.7085,5.6224) -- (7.7077,5.6222) -- (7.7084,5.6218) -- (7.7077,5.6217) -- (7.7080,5.6214) -- (7.7076,5.6213) -- (7.7079,5.6207) -- (7.7071,5.6207) -- (7.7075,5.6203) -- (7.7070,5.6202) -- (7.7074,5.6197) -- (7.7065,5.6197) -- (7.7070,5.6193) -- (7.7064,5.6192) -- (7.7067,5.6189) -- (7.7059,5.6186) -- (7.7065,5.6183) -- (7.7058,5.6182) -- (7.7062,5.6179) -- (7.7057,5.6178) -- (7.7061,5.6172) -- (7.7052,5.6172) -- (7.7057,5.6168) -- (7.7051,5.6167) -- (7.7056,5.6162) -- (7.7046,5.6162) -- (7.7052,5.6158) -- (7.7046,5.6157) -- (7.7049,5.6154) -- (7.7040,5.6151) -- (7.7047,5.6148) -- (7.7040,5.6147) -- (7.7043,5.6144) -- (7.7039,5.6143) -- (7.7042,5.6137) -- (7.7034,5.6137) -- (7.7038,5.6133) -- (7.7033,5.6132) -- (7.7037,5.6127) -- (7.7028,5.6127) -- (7.7033,5.6123) -- (7.7027,5.6122) -- (7.7030,5.6119) -- (7.7022,5.6116) -- (7.7028,5.6113) -- (7.7021,5.6112) -- (7.7025,5.6109) -- (7.7016,5.6106) -- (7.7023,5.6102) -- (7.7016,5.6102) -- (7.7020,5.6099) -- (7.7015,5.6097) -- (7.7019,5.6092) -- (7.7010,5.6092) -- (7.7015,5.6088) -- (7.7009,5.6087) -- (7.7014,5.6082) -- (7.7004,5.6082) -- (7.7010,5.6078) -- (7.7003,5.6077) -- (7.7006,5.6074) -- (7.6998,5.6071) -- (7.7005,5.6068) -- (7.6998,5.6067) -- (7.7001,5.6064) -- (7.6992,5.6061) -- (7.7000,5.6057) -- (7.6992,5.6057) -- (7.6996,5.6054) -- (7.6991,5.6052) -- (7.6995,5.6047) -- (7.6986,5.6047) -- (7.6991,5.6043) -- (7.6985,5.6042) -- (7.6990,5.6037) -- (7.6980,5.6037) -- (7.6986,5.6033) -- (7.6980,5.6032) -- (7.6982,5.6029) -- (7.6975,5.6026) -- (7.6980,5.6023) -- (7.6974,5.6022) -- (7.6977,5.6019) -- (7.6969,5.6016) -- (7.6975,5.6013) -- (7.6968,5.6012) -- (7.6972,5.6009) -- (7.6963,5.6006) -- (7.6970,5.6002) -- (7.6963,5.6002) -- (7.6967,5.5999) -- (7.6962,5.5997) -- (7.6965,5.5992) -- (7.6957,5.5992) -- (7.6961,5.5988) -- (7.6956,5.5987) -- (7.6960,5.5982) -- (7.6952,5.5982) -- (7.6956,5.5978) -- (7.6951,5.5977) -- (7.6955,5.5972) -- (7.6946,5.5971) -- (7.6951,5.5968) -- (7.6945,5.5967) -- (7.6948,5.5964) -- (7.6940,5.5961) -- (7.6946,5.5958) -- (7.6940,5.5957) -- (7.6943,5.5954) -- (7.6935,5.5951) -- (7.6941,5.5948) -- (7.6934,5.5947) -- (7.6937,5.5944) -- (7.6929,5.5941) -- (7.6936,5.5937) -- (7.6929,5.5937) -- (7.6932,5.5934) -- (7.6923,5.5931) -- (7.6930,5.5927) -- (7.6923,5.5927) -- (7.6927,5.5924) -- (7.6922,5.5922) -- (7.6925,5.5917) -- (7.6918,5.5917) -- (7.6922,5.5913) -- (7.6917,5.5912) -- (7.6920,5.5907) -- (7.6912,5.5907) -- (7.6916,5.5903) -- (7.6911,5.5902) -- (7.6915,5.5897) -- (7.6907,5.5897) -- (7.6911,5.5893) -- (7.6906,5.5892) -- (7.6910,5.5887) -- (7.6901,5.5887) -- (7.6906,5.5883) -- (7.6900,5.5882) -- (7.6905,5.5877) -- (7.6896,5.5877) -- (7.6900,5.5873) -- (7.6895,5.5872) -- (7.6899,5.5867) -- (7.6890,5.5866) -- (7.6895,5.5863) -- (7.6890,5.5862) -- (7.6894,5.5856) -- (7.6885,5.5856) -- (7.6890,5.5853) -- (7.6884,5.5852) -- (7.6889,5.5846) -- (7.6879,5.5846) -- (7.6885,5.5843) -- (7.6879,5.5842) -- (7.6881,5.5839) -- (7.6874,5.5836) -- (7.6879,5.5833) -- (7.6873,5.5832) -- (7.6876,5.5829) -- (7.6869,5.5826) -- (7.6874,5.5823) -- (7.6868,5.5822) -- (7.6871,5.5819) -- (7.6863,5.5816) -- (7.6869,5.5813) -- (7.6863,5.5812) -- (7.6865,5.5809) -- (7.6858,5.5806) -- (7.6863,5.5803) -- (7.6857,5.5802) -- (7.6860,5.5799) -- (7.6852,5.5796) -- (7.6858,5.5793) -- (7.6852,5.5792) -- (7.6855,5.5789) -- (7.6847,5.5786) -- (7.6853,5.5783) -- (7.6847,5.5782) -- (7.6849,5.5779) -- (7.6842,5.5776) -- (7.6847,5.5773) -- (7.6841,5.5772) -- (7.6844,5.5769) -- (7.6837,5.5766) -- (7.6842,5.5763) -- (7.6836,5.5762) -- (7.6841,5.5756) -- (7.6831,5.5756) -- (7.6837,5.5753) -- (7.6831,5.5752) -- (7.6836,5.5746) -- (7.6826,5.5746) -- (7.6831,5.5743) -- (7.6826,5.5742) -- (7.6830,5.5737) -- (7.6821,5.5737) -- (7.6826,5.5733) -- (7.6820,5.5732) -- (7.6825,5.5727) -- (7.6816,5.5727) -- (7.6820,5.5723) -- (7.6815,5.5722) -- (7.6819,5.5717) -- (7.6811,5.5717) -- (7.6815,5.5713) -- (7.6810,5.5712) -- (7.6814,5.5707) -- (7.6805,5.5707) -- (7.6810,5.5703) -- (7.6805,5.5702) -- (7.6808,5.5697) -- (7.6800,5.5697) -- (7.6804,5.5694) -- (7.6799,5.5692) -- (7.6803,5.5687) -- (7.6795,5.5687) -- (7.6799,5.5684) -- (7.6794,5.5682) -- (7.6797,5.5677) -- (7.6790,5.5677) -- (7.6793,5.5674) -- (7.6785,5.5671) -- (7.6792,5.5667) -- (7.6785,5.5667) -- (7.6788,5.5664) -- (7.6780,5.5661) -- (7.6786,5.5658) -- (7.6780,5.5657) -- (7.6783,5.5654) -- (7.6775,5.5652) -- (7.6781,5.5648) -- (7.6775,5.5647) -- (7.6777,5.5644) -- (7.6770,5.5642) -- (7.6775,5.5638) -- (7.6769,5.5637) -- (7.6774,5.5632) -- (7.6765,5.5632) -- (7.6770,5.5628) -- (7.6764,5.5627) -- (7.6768,5.5622) -- (7.6760,5.5622) -- (7.6764,5.5618) -- (7.6759,5.5617) -- (7.6763,5.5612) -- (7.6755,5.5612) -- (7.6759,5.5609) -- (7.6754,5.5608) -- (7.6757,5.5602) -- (7.6750,5.5602) -- (7.6753,5.5599) -- (7.6745,5.5597) -- (7.6752,5.5593) -- (7.6745,5.5592) -- (7.6748,5.5589) -- (7.6740,5.5587) -- (7.6746,5.5583) -- (7.6740,5.5582) -- (7.6742,5.5579) -- (7.6735,5.5577) -- (7.6740,5.5573) -- (7.6735,5.5572) -- (7.6739,5.5567) -- (7.6730,5.5567) -- (7.6735,5.5563) -- (7.6730,5.5562) -- (7.6733,5.5557) -- (7.6725,5.5557) -- (7.6729,5.5554) -- (7.6725,5.5553) -- (7.6728,5.5547) -- (7.6720,5.5547) -- (7.6724,5.5544) -- (7.6715,5.5542) -- (7.6722,5.5538) -- (7.6715,5.5537) -- (7.6718,5.5534) -- (7.6710,5.5532) -- (7.6716,5.5528) -- (7.6710,5.5527) -- (7.6713,5.5525) -- (7.6706,5.5522) -- (7.6711,5.5518) -- (7.6705,5.5518) -- (7.6709,5.5512) -- (7.6701,5.5512) -- (7.6705,5.5509) -- (7.6700,5.5508) -- (7.6704,5.5503) -- (7.6696,5.5502) -- (7.6700,5.5499) -- (7.6691,5.5497) -- (7.6698,5.5493) -- (7.6691,5.5492) -- (7.6694,5.5489) -- (7.6686,5.5487) -- (7.6692,5.5483) -- (7.6686,5.5483) -- (7.6691,5.5477) -- (7.6682,5.5477) -- (7.6686,5.5474) -- (7.6681,5.5473) -- (7.6685,5.5467) -- (7.6677,5.5467) -- (7.6681,5.5464) -- (7.6676,5.5463) -- (7.6679,5.5458) -- (7.6672,5.5457) -- (7.6675,5.5454) -- (7.6667,5.5452) -- (7.6673,5.5448) -- (7.6667,5.5448) -- (7.6670,5.5445) -- (7.6663,5.5442) -- (7.6668,5.5439) -- (7.6662,5.5438) -- (7.6666,5.5433) -- (7.6658,5.5432) -- (7.6662,5.5429) -- (7.6658,5.5428) -- (7.6660,5.5423) -- (7.6653,5.5423) -- (7.6657,5.5420) -- (7.6649,5.5417) -- (7.6655,5.5414) -- (7.6649,5.5413) -- (7.6653,5.5407) -- (7.6644,5.5407) -- (7.6649,5.5404) -- (7.6644,5.5403) -- (7.6647,5.5398) -- (7.6639,5.5398) -- (7.6643,5.5394) -- (7.6634,5.5392) -- (7.6641,5.5388) -- (7.6635,5.5388) -- (7.6638,5.5385) -- (7.6630,5.5382) -- (7.6636,5.5379) -- (7.6630,5.5378) -- (7.6634,5.5373) -- (7.6626,5.5373) -- (7.6630,5.5369) -- (7.6625,5.5368) -- (7.6628,5.5363) -- (7.6621,5.5363) -- (7.6624,5.5360) -- (7.6616,5.5357) -- (7.6622,5.5354) -- (7.6616,5.5353) -- (7.6621,5.5348) -- (7.6612,5.5348) -- (7.6617,5.5344) -- (7.6612,5.5343) -- (7.6615,5.5338) -- (7.6607,5.5338) -- (7.6611,5.5335) -- (7.6602,5.5333) -- (7.6609,5.5329) -- (7.6603,5.5328) -- (7.6605,5.5325) -- (7.6598,5.5323) -- (7.6603,5.5319) -- (7.6598,5.5318) -- (7.6602,5.5313) -- (7.6594,5.5313) -- (7.6597,5.5310) -- (7.6589,5.5308) -- (7.6596,5.5304) -- (7.6589,5.5303) -- (7.6592,5.5300) -- (7.6585,5.5298) -- (7.6590,5.5294) -- (7.6584,5.5294) -- (7.6588,5.5288) -- (7.6580,5.5288) -- (7.6584,5.5285) -- (7.6575,5.5283) -- (7.6582,5.5279) -- (7.6576,5.5278) -- (7.6578,5.5276) -- (7.6571,5.5273) -- (7.6576,5.5270) -- (7.6571,5.5269) -- (7.6575,5.5264) -- (7.6567,5.5263) -- (7.6570,5.5260) -- (7.6562,5.5258) -- (7.6569,5.5254) -- (7.6562,5.5254) -- (7.6565,5.5251) -- (7.6558,5.5248) -- (7.6563,5.5245) -- (7.6558,5.5244) -- (7.6561,5.5239) -- (7.6554,5.5239) -- (7.6557,5.5236) -- (7.6549,5.5233) -- (7.6555,5.5230) -- (7.6549,5.5229) -- (7.6554,5.5224) -- (7.6545,5.5223) -- (7.6549,5.5220) -- (7.6544,5.5219) -- (7.6547,5.5214) -- (7.6540,5.5214) -- (7.6543,5.5211) -- (7.6536,5.5208) -- (7.6541,5.5205) -- (7.6536,5.5204) -- (7.6540,5.5199) -- (7.6532,5.5199) -- (7.6535,5.5196) -- (7.6527,5.5193) -- (7.6533,5.5190) -- (7.6527,5.5189) -- (7.6530,5.5186) -- (7.6523,5.5184) -- (7.6528,5.5180) -- (7.6523,5.5180) -- (7.6526,5.5174) -- (7.6519,5.5174) -- (7.6522,5.5171) -- (7.6514,5.5169) -- (7.6520,5.5165) -- (7.6514,5.5164) -- (7.6518,5.5159) -- (7.6510,5.5159) -- (7.6514,5.5156) -- (7.6505,5.5154) -- (7.6512,5.5150) -- (7.6506,5.5149) -- (7.6511,5.5144) -- (7.6501,5.5144) -- (7.6506,5.5141) -- (7.6501,5.5140) -- (7.6504,5.5135) -- (7.6497,5.5134) -- (7.6500,5.5132) -- (7.6493,5.5129) -- (7.6498,5.5126) -- (7.6493,5.5125) -- (7.6496,5.5120) -- (7.6489,5.5119) -- (7.6492,5.5117) -- (7.6484,5.5114) -- (7.6490,5.5111) -- (7.6484,5.5110) -- (7.6489,5.5105) -- (7.6480,5.5105) -- (7.6484,5.5101) -- (7.6475,5.5099) -- (7.6482,5.5096) -- (7.6476,5.5095) -- (7.6478,5.5092) -- (7.6472,5.5090) -- (7.6476,5.5086) -- (7.6471,5.5085) -- (7.6474,5.5080) -- (7.6468,5.5080) -- (7.6470,5.5077) -- (7.6463,5.5075) -- (7.6468,5.5071) -- (7.6463,5.5070) -- (7.6466,5.5065) -- (7.6459,5.5065) -- (7.6462,5.5062) -- (7.6455,5.5060) -- (7.6460,5.5056) -- (7.6455,5.5056) -- (7.6459,5.5050) -- (7.6451,5.5050) -- (7.6454,5.5047) -- (7.6446,5.5045) -- (7.6452,5.5041) -- (7.6447,5.5041) -- (7.6451,5.5035) -- (7.6443,5.5035) -- (7.6446,5.5032) -- (7.6438,5.5030) -- (7.6444,5.5026) -- (7.6438,5.5026) -- (7.6443,5.5020) -- (7.6434,5.5020) -- (7.6438,5.5017) -- (7.6430,5.5015) -- (7.6436,5.5012) -- (7.6430,5.5011) -- (7.6435,5.5006) -- (7.6426,5.5006) -- (7.6430,5.5002) -- (7.6421,5.5000) -- (7.6428,5.4997) -- (7.6422,5.4996) -- (7.6427,5.4991) -- (7.6418,5.4991) -- (7.6422,5.4988) -- (7.6413,5.4986) -- (7.6420,5.4982) -- (7.6414,5.4981) -- (7.6416,5.4978) -- (7.6410,5.4976) -- (7.6414,5.4973) -- (7.6409,5.4972) -- (7.6412,5.4967) -- (7.6406,5.4966) -- (7.6408,5.4964) -- (7.6402,5.4961) -- (7.6406,5.4958) -- (7.6401,5.4957) -- (7.6404,5.4952) -- (7.6398,5.4952) -- (7.6400,5.4949) -- (7.6393,5.4946) -- (7.6398,5.4943) -- (7.6389,5.4941) -- (7.6396,5.4937) -- (7.6390,5.4937) -- (7.6394,5.4931) -- (7.6385,5.4931) -- (7.6389,5.4928) -- (7.6381,5.4926) -- (7.6387,5.4923) -- (7.6381,5.4922) -- (7.6386,5.4917) -- (7.6377,5.4917) -- (7.6381,5.4914) -- (7.6373,5.4912) -- (7.6379,5.4908) -- (7.6373,5.4907) -- (7.6378,5.4902) -- (7.6369,5.4902) -- (7.6373,5.4899) -- (7.6365,5.4897) -- (7.6371,5.4893) -- (7.6365,5.4892) -- (7.6369,5.4887) -- (7.6362,5.4887) -- (7.6365,5.4884) -- (7.6357,5.4882) -- (7.6363,5.4879) -- (7.6357,5.4878) -- (7.6361,5.4873) -- (7.6354,5.4872) -- (7.6357,5.4870) -- (7.6349,5.4867) -- (7.6355,5.4864) -- (7.6350,5.4863) -- (7.6353,5.4858) -- (7.6346,5.4858) -- (7.6349,5.4855) -- (7.6342,5.4853) -- (7.6346,5.4849) -- (7.6342,5.4848) -- (7.6344,5.4844) -- (7.6338,5.4843) -- (7.6340,5.4840) -- (7.6334,5.4838) -- (7.6338,5.4835) -- (7.6329,5.4833) -- (7.6336,5.4829) -- (7.6330,5.4828) -- (7.6334,5.4823) -- (7.6326,5.4823) -- (7.6330,5.4820) -- (7.6322,5.4818) -- (7.6328,5.4814) -- (7.6322,5.4814) -- (7.6326,5.4809) -- (7.6318,5.4808) -- (7.6322,5.4806) -- (7.6314,5.4803) -- (7.6319,5.4800) -- (7.6314,5.4799) -- (7.6317,5.4794) -- (7.6311,5.4794) -- (7.6313,5.4791) -- (7.6307,5.4789) -- (7.6311,5.4785) -- (7.6302,5.4784) -- (7.6309,5.4780) -- (7.6303,5.4779) -- (7.6307,5.4774) -- (7.6299,5.4774) -- (7.6303,5.4771) -- (7.6295,5.4769) -- (7.6300,5.4765) -- (7.6295,5.4765) -- (7.6299,5.4760) -- (7.6291,5.4759) -- (7.6294,5.4757) -- (7.6287,5.4754) -- (7.6292,5.4751) -- (7.6287,5.4750) -- (7.6290,5.4745) -- (7.6284,5.4745) -- (7.6289,5.4739) -- (7.6280,5.4740) -- (7.6284,5.4737) -- (7.6276,5.4734) -- (7.6282,5.4731) -- (7.6276,5.4730) -- (7.6280,5.4725) -- (7.6272,5.4725) -- (7.6275,5.4722) -- (7.6268,5.4720) -- (7.6273,5.4717) -- (7.6268,5.4716) -- (7.6271,5.4711) -- (7.6265,5.4710) -- (7.6269,5.4705) -- (7.6261,5.4705) -- (7.6265,5.4702) -- (7.6257,5.4700) -- (7.6263,5.4697) -- (7.6257,5.4696) -- (7.6261,5.4691) -- (7.6254,5.4691) -- (7.6256,5.4688) -- (7.6250,5.4685) -- (7.6254,5.4682) -- (7.6250,5.4681) -- (7.6252,5.4677) -- (7.6246,5.4676) -- (7.6250,5.4671) -- (7.6242,5.4671) -- (7.6246,5.4668) -- (7.6238,5.4666) -- (7.6243,5.4663) -- (7.6239,5.4662) -- (7.6241,5.4657) -- (7.6235,5.4656) -- (7.6237,5.4654) -- (7.6231,5.4651) -- (7.6235,5.4648) -- (7.6227,5.4646) -- (7.6233,5.4643) -- (7.6227,5.4642) -- (7.6231,5.4637) -- (7.6224,5.4637) -- (7.6227,5.4634) -- (7.6220,5.4632) -- (7.6224,5.4629) -- (7.6216,5.4627) -- (7.6222,5.4623) -- (7.6216,5.4622) -- (7.6220,5.4617) -- (7.6213,5.4617) -- (7.6216,5.4614) -- (7.6209,5.4612) -- (7.6213,5.4609) -- (7.6209,5.4608) -- (7.6211,5.4603) -- (7.6205,5.4603) -- (7.6210,5.4598) -- (7.6202,5.4598) -- (7.6205,5.4595) -- (7.6198,5.4592) -- (7.6203,5.4589) -- (7.6198,5.4588) -- (7.6201,5.4584) -- (7.6194,5.4583) -- (7.6199,5.4578) -- (7.6191,5.4578) -- (7.6194,5.4575) -- (7.6187,5.4573) -- (7.6192,5.4570) -- (7.6187,5.4569) -- (7.6190,5.4564) -- (7.6184,5.4564) -- (7.6188,5.4558) -- (7.6180,5.4558) -- (7.6183,5.4556) -- (7.6176,5.4553) -- (7.6181,5.4550) -- (7.6176,5.4549) -- (7.6179,5.4545) -- (7.6173,5.4544) -- (7.6177,5.4539) -- (7.6169,5.4539) -- (7.6172,5.4536) -- (7.6165,5.4534) -- (7.6170,5.4531) -- (7.6165,5.4530) -- (7.6168,5.4525) -- (7.6162,5.4525) -- (7.6166,5.4519) -- (7.6158,5.4519) -- (7.6162,5.4517) -- (7.6154,5.4514) -- (7.6159,5.4511) -- (7.6155,5.4510) -- (7.6157,5.4506) -- (7.6151,5.4505) -- (7.6155,5.4500) -- (7.6148,5.4500) -- (7.6151,5.4497) -- (7.6144,5.4495) -- (7.6148,5.4492) -- (7.6144,5.4491) -- (7.6146,5.4486) -- (7.6140,5.4486) -- (7.6144,5.4481) -- (7.6137,5.4481) -- (7.6140,5.4478) -- (7.6133,5.4475) -- (7.6137,5.4472) -- (7.6129,5.4470) -- (7.6135,5.4467) -- (7.6130,5.4466) -- (7.6133,5.4461) -- (7.6126,5.4461) -- (7.6129,5.4458) -- (7.6123,5.4456) -- (7.6126,5.4453) -- (7.6119,5.4451) -- (7.6124,5.4448) -- (7.6119,5.4447) -- (7.6122,5.4442) -- (7.6116,5.4442) -- (7.6120,5.4437) -- (7.6112,5.4437) -- (7.6115,5.4434) -- (7.6108,5.4432) -- (7.6113,5.4428) -- (7.6109,5.4428) -- (7.6111,5.4423) -- (7.6105,5.4422) -- (7.6109,5.4418) -- (7.6102,5.4417) -- (7.6104,5.4415) -- (7.6098,5.4412) -- (7.6102,5.4409) -- (7.6094,5.4407) -- (7.6100,5.4404) -- (7.6095,5.4403) -- (7.6098,5.4398) -- (7.6091,5.4398) -- (7.6096,5.4393) -- (7.6088,5.4393) -- (7.6091,5.4390) -- (7.6084,5.4388) -- (7.6089,5.4385) -- (7.6080,5.4383) -- (7.6086,5.4379) -- (7.6081,5.4379) -- (7.6084,5.4374) -- (7.6078,5.4374) -- (7.6080,5.4371) -- (7.6074,5.4369) -- (7.6077,5.4366) -- (7.6070,5.4364) -- (7.6075,5.4360) -- (7.6071,5.4360) -- (7.6073,5.4355) -- (7.6067,5.4354) -- (7.6071,5.4350) -- (7.6064,5.4349) -- (7.6066,5.4347) -- (7.6060,5.4344) -- (7.6064,5.4341) -- (7.6056,5.4339) -- (7.6062,5.4336) -- (7.6057,5.4335) -- (7.6059,5.4331) -- (7.6054,5.4330) -- (7.6057,5.4325) -- (7.6050,5.4325) -- (7.6053,5.4322) -- (7.6046,5.4320) -- (7.6050,5.4317) -- (7.6043,5.4315) -- (7.6048,5.4312) -- (7.6043,5.4311) -- (7.6046,5.4306) -- (7.6040,5.4306) -- (7.6044,5.4301) -- (7.6037,5.4301) -- (7.6039,5.4298) -- (7.6033,5.4296) -- (7.6037,5.4293) -- (7.6029,5.4291) -- (7.6034,5.4288) -- (7.6030,5.4287) -- (7.6032,5.4282) -- (7.6026,5.4282) -- (7.6030,5.4277) -- (7.6023,5.4277) -- (7.6026,5.4274) -- (7.6020,5.4272) -- (7.6023,5.4269) -- (7.6016,5.4267) -- (7.6021,5.4264) -- (7.6016,5.4263) -- (7.6019,5.4258) -- (7.6013,5.4258) -- (7.6017,5.4253) -- (7.6010,5.4253) -- (7.6012,5.4250) -- (7.6006,5.4248) -- (7.6010,5.4245) -- (7.6003,5.4243) -- (7.6007,5.4240) -- (7.5999,5.4238) -- (7.6005,5.4234) -- (7.6000,5.4234) -- (7.6003,5.4229) -- (7.5996,5.4229) -- (7.6001,5.4224) -- (7.5993,5.4224) -- (7.5996,5.4221) -- (7.5989,5.4219) -- (7.5993,5.4216) -- (7.5986,5.4214) -- (7.5991,5.4210) -- (7.5986,5.4210) -- (7.5989,5.4205) -- (7.5983,5.4205) -- (7.5987,5.4200) -- (7.5980,5.4200) -- (7.5982,5.4197) -- (7.5976,5.4195) -- (7.5980,5.4192) -- (7.5973,5.4190) -- (7.5977,5.4187) -- (7.5969,5.4185) -- (7.5975,5.4181) -- (7.5970,5.4181) -- (7.5973,5.4176) -- (7.5967,5.4176) -- (7.5971,5.4171) -- (7.5963,5.4171) -- (7.5966,5.4168) -- (7.5960,5.4166) -- (7.5964,5.4163) -- (7.5956,5.4161) -- (7.5961,5.4158) -- (7.5953,5.4156) -- (7.5959,5.4152) -- (7.5954,5.4152) -- (7.5957,5.4147) -- (7.5950,5.4147) -- (7.5954,5.4142) -- (7.5947,5.4142) -- (7.5950,5.4139) -- (7.5944,5.4137) -- (7.5947,5.4134) -- (7.5940,5.4132) -- (7.5945,5.4129) -- (7.5936,5.4127) -- (7.5942,5.4124) -- (7.5938,5.4123) -- (7.5940,5.4118) -- (7.5934,5.4118) -- (7.5938,5.4113) -- (7.5931,5.4113) -- (7.5934,5.4110) -- (7.5928,5.4108) -- (7.5931,5.4105) -- (7.5924,5.4103) -- (7.5929,5.4100) -- (7.5921,5.4098) -- (7.5926,5.4095) -- (7.5921,5.4094) -- (7.5924,5.4090) -- (7.5918,5.4089) -- (7.5922,5.4085) -- (7.5915,5.4084) -- (7.5920,5.4079) -- (7.5912,5.4079) -- (7.5915,5.4077) -- (7.5908,5.4074) -- (7.5912,5.4072) -- (7.5905,5.4070) -- (7.5910,5.4066) -- (7.5901,5.4065) -- (7.5907,5.4061) -- (7.5902,5.4061) -- (7.5905,5.4056) -- (7.5899,5.4056) -- (7.5903,5.4051) -- (7.5896,5.4051) -- (7.5898,5.4048) -- (7.5893,5.4046) -- (7.5896,5.4043) -- (7.5889,5.4041) -- (7.5893,5.4038) -- (7.5886,5.4036) -- (7.5891,5.4033) -- (7.5886,5.4032) -- (7.5888,5.4028) -- (7.5883,5.4027) -- (7.5886,5.4023) -- (7.5880,5.4022) -- (7.5884,5.4017) -- (7.5877,5.4017) -- (7.5879,5.4015) -- (7.5874,5.4012) -- (7.5877,5.4010) -- (7.5870,5.4007) -- (7.5874,5.4005) -- (7.5867,5.4003) -- (7.5872,5.3999) -- (7.5863,5.3998) -- (7.5869,5.3994) -- (7.5864,5.3994) -- (7.5867,5.3989) -- (7.5861,5.3989) -- (7.5865,5.3984) -- (7.5858,5.3984) -- (7.5863,5.3979) -- (7.5855,5.3979) -- (7.5858,5.3976) -- (7.5852,5.3974) -- (7.5855,5.3971) -- (7.5848,5.3969) -- (7.5853,5.3966) -- (7.5845,5.3964) -- (7.5850,5.3961) -- (7.5846,5.3960) -- (7.5848,5.3956) -- (7.5843,5.3955) -- (7.5845,5.3951) -- (7.5840,5.3951) -- (7.5843,5.3946) -- (7.5836,5.3946) -- (7.5841,5.3941) -- (7.5833,5.3941) -- (7.5836,5.3938) -- (7.5830,5.3936) -- (7.5834,5.3933) -- (7.5827,5.3931) -- (7.5831,5.3928) -- (7.5823,5.3926) -- (7.5829,5.3923) -- (7.5824,5.3922) -- (7.5826,5.3918) -- (7.5821,5.3917) -- (7.5824,5.3913) -- (7.5818,5.3913) -- (7.5821,5.3908) -- (7.5815,5.3908) -- (7.5819,5.3903) -- (7.5812,5.3903) -- (7.5814,5.3900) -- (7.5809,5.3898) -- (7.5812,5.3895) -- (7.5805,5.3893) -- (7.5809,5.3890) -- (7.5802,5.3888) -- (7.5807,5.3885) -- (7.5799,5.3884) -- (7.5804,5.3880) -- (7.5800,5.3879) -- (7.5802,5.3875) -- (7.5797,5.3875) -- (7.5800,5.3870) -- (7.5794,5.3870) -- (7.5797,5.3865) -- (7.5791,5.3865) -- (7.5795,5.3860) -- (7.5788,5.3860) -- (7.5790,5.3858) -- (7.5784,5.3855) -- (7.5788,5.3853) -- (7.5781,5.3850) -- (7.5785,5.3848) -- (7.5778,5.3846) -- (7.5782,5.3843) -- (7.5775,5.3841) -- (7.5780,5.3838) -- (7.5776,5.3837) -- (7.5777,5.3833) -- (7.5773,5.3832) -- (7.5775,5.3828) -- (7.5770,5.3827) -- (7.5773,5.3823) -- (7.5767,5.3822) -- (7.5770,5.3818) -- (7.5764,5.3817) -- (7.5768,5.3813) -- (7.5761,5.3813) -- (7.5763,5.3810) -- (7.5757,5.3808) -- (7.5761,5.3805) -- (7.5754,5.3803) -- (7.5758,5.3800) -- (7.5751,5.3798) -- (7.5755,5.3795) -- (7.5748,5.3794) -- (7.5753,5.3790) -- (7.5745,5.3789) -- (7.5750,5.3785) -- (7.5746,5.3785) -- (7.5748,5.3780) -- (7.5743,5.3780) -- (7.5745,5.3776) -- (7.5740,5.3775) -- (7.5743,5.3771) -- (7.5737,5.3770) -- (7.5741,5.3766) -- (7.5734,5.3765) -- (7.5738,5.3761) -- (7.5731,5.3761) -- (7.5734,5.3758) -- (7.5728,5.3756) -- (7.5731,5.3753) -- (7.5725,5.3751) -- (7.5728,5.3748) -- (7.5722,5.3746) -- (7.5726,5.3743) -- (7.5718,5.3742) -- (7.5723,5.3739) -- (7.5715,5.3737) -- (7.5721,5.3734) -- (7.5716,5.3733) -- (7.5718,5.3729) -- (7.5713,5.3728) -- (7.5716,5.3724) -- (7.5710,5.3723) -- (7.5713,5.3719) -- (7.5708,5.3719) -- (7.5711,5.3714) -- (7.5705,5.3714) -- (7.5708,5.3709) -- (7.5702,5.3709) -- (7.5706,5.3704) -- (7.5699,5.3704) -- (7.5701,5.3702) -- (7.5696,5.3699) -- (7.5698,5.3697) -- (7.5693,5.3695) -- (7.5696,5.3692) -- (7.5690,5.3690) -- (7.5693,5.3687) -- (7.5687,5.3685) -- (7.5691,5.3682) -- (7.5684,5.3680) -- (7.5688,5.3677) -- (7.5680,5.3676) -- (7.5685,5.3673) -- (7.5677,5.3671) -- (7.5683,5.3668) -- (7.5679,5.3667) -- (7.5680,5.3663) -- (7.5676,5.3662) -- (7.5678,5.3658) -- (7.5673,5.3658) -- (7.5675,5.3653) -- (7.5670,5.3653) -- (7.5673,5.3648) -- (7.5667,5.3648) -- (7.5670,5.3643) -- (7.5664,5.3643) -- (7.5668,5.3639) -- (7.5661,5.3639) -- (7.5666,5.3634) -- (7.5658,5.3634) -- (7.5661,5.3631) -- (7.5655,5.3629) -- (7.5658,5.3627) -- (7.5652,5.3624) -- (7.5655,5.3622) -- (7.5650,5.3620) -- (7.5653,5.3617) -- (7.5647,5.3615) -- (7.5650,5.3612) -- (7.5644,5.3610) -- (7.5648,5.3607) -- (7.5641,5.3605) -- (7.5645,5.3603) -- (7.5638,5.3601) -- (7.5642,5.3598) -- (7.5635,5.3596) -- (7.5640,5.3593) -- (7.5632,5.3591) -- (7.5637,5.3588) -- (7.5633,5.3587) -- (7.5635,5.3583) -- (7.5630,5.3583) -- (7.5632,5.3579) -- (7.5627,5.3578) -- (7.5629,5.3574) -- (7.5624,5.3573) -- (7.5627,5.3569) -- (7.5622,5.3569) -- (7.5624,5.3564) -- (7.5619,5.3564) -- (7.5622,5.3559) -- (7.5616,5.3559) -- (7.5619,5.3555) -- (7.5613,5.3554) -- (7.5617,5.3550) -- (7.5610,5.3550) -- (7.5614,5.3545) -- (7.5607,5.3545) -- (7.5612,5.3540) -- (7.5605,5.3540) -- (7.5609,5.3536) -- (7.5602,5.3536) -- (7.5604,5.3533) -- (7.5599,5.3531) -- (7.5602,5.3528) -- (7.5596,5.3526) -- (7.5599,5.3524) -- (7.5593,5.3522) -- (7.5596,5.3519) -- (7.5590,5.3517) -- (7.5594,5.3514) -- (7.5588,5.3512) -- (7.5591,5.3510) -- (7.5585,5.3508) -- (7.5588,5.3505) -- (7.5582,5.3503) -- (7.5586,5.3500) -- (7.5579,5.3498) -- (7.5583,5.3495) -- (7.5576,5.3494) -- (7.5580,5.3491) -- (7.5573,5.3489) -- (7.5578,5.3486) -- (7.5571,5.3484) -- (7.5575,5.3481) -- (7.5568,5.3480) -- (7.5573,5.3477) -- (7.5565,5.3475) -- (7.5570,5.3472) -- (7.5562,5.3470) -- (7.5567,5.3467) -- (7.5559,5.3466) -- (7.5565,5.3463) -- (7.5561,5.3462) -- (7.5562,5.3458) -- (7.5558,5.3457) -- (7.5559,5.3453) -- (7.5555,5.3453) -- (7.5557,5.3448) -- (7.5552,5.3448) -- (7.5554,5.3444) -- (7.5550,5.3443) -- (7.5552,5.3439) -- (7.5547,5.3439) -- (7.5549,5.3434) -- (7.5544,5.3434) -- (7.5546,5.3430) -- (7.5541,5.3429) -- (7.5544,5.3425) -- (7.5539,5.3425) -- (7.5541,5.3420) -- (7.5536,5.3420) -- (7.5538,5.3416) -- (7.5533,5.3415) -- (7.5536,5.3411) -- (7.5531,5.3411) -- (7.5533,5.3406) -- (7.5528,5.3406) -- (7.5530,5.3402) -- (7.5525,5.3401) -- (7.5528,5.3397) -- (7.5522,5.3397) -- (7.5525,5.3393) -- (7.5520,5.3392) -- (7.5523,5.3388) -- (7.5517,5.3388) -- (7.5520,5.3383) -- (7.5514,5.3383) -- (7.5517,5.3379) -- (7.5512,5.3378) -- (7.5515,5.3374) -- (7.5509,5.3374) -- (7.5512,5.3369) -- (7.5506,5.3369) -- (7.5509,5.3365) -- (7.5504,5.3364) -- (7.5507,5.3360) -- (7.5501,5.3360) -- (7.5504,5.3356) -- (7.5498,5.3355) -- (7.5501,5.3351) -- (7.5496,5.3351) -- (7.5499,5.3346) -- (7.5493,5.3346) -- (7.5496,5.3342) -- (7.5490,5.3341) -- (7.5493,5.3337) -- (7.5488,5.3337) -- (7.5491,5.3333) -- (7.5485,5.3332) -- (7.5488,5.3328) -- (7.5482,5.3328) -- (7.5485,5.3323) -- (7.5480,5.3323) -- (7.5483,5.3319) -- (7.5477,5.3319) -- (7.5480,5.3314) -- (7.5474,5.3314) -- (7.5477,5.3310) -- (7.5472,5.3309) -- (7.5474,5.3305) -- (7.5469,5.3305) -- (7.5472,5.3301) -- (7.5466,5.3300) -- (7.5469,5.3296) -- (7.5464,5.3296) -- (7.5466,5.3291) -- (7.5461,5.3291) -- (7.5464,5.3287) -- (7.5459,5.3286) -- (7.5461,5.3282) -- (7.5456,5.3282) -- (7.5458,5.3278) -- (7.5453,5.3277) -- (7.5456,5.3273) -- (7.5451,5.3273) -- (7.5453,5.3269) -- (7.5448,5.3268) -- (7.5450,5.3264) -- (7.5445,5.3264) -- (7.5447,5.3260) -- (7.5443,5.3259) -- (7.5445,5.3255) -- (7.5440,5.3255) -- (7.5442,5.3251) -- (7.5438,5.3250) -- (7.5439,5.3246) -- (7.5435,5.3245) -- (7.5437,5.3242) -- (7.5432,5.3241) -- (7.5434,5.3237) -- (7.5426,5.3236) -- (7.5431,5.3233) -- (7.5423,5.3231) -- (7.5428,5.3228) -- (7.5421,5.3227) -- (7.5426,5.3224) -- (7.5418,5.3222) -- (7.5423,5.3219) -- (7.5416,5.3217) -- (7.5420,5.3215) -- (7.5413,5.3213) -- (7.5418,5.3210) -- (7.5411,5.3208) -- (7.5415,5.3206) -- (7.5408,5.3204) -- (7.5412,5.3201) -- (7.5406,5.3199) -- (7.5410,5.3197) -- (7.5403,5.3195) -- (7.5407,5.3192) -- (7.5401,5.3190) -- (7.5404,5.3188) -- (7.5398,5.3186) -- (7.5401,5.3183) -- (7.5396,5.3181) -- (7.5399,5.3179) -- (7.5393,5.3177) -- (7.5396,5.3174) -- (7.5391,5.3172) -- (7.5393,5.3170) -- (7.5388,5.3168) -- (7.5393,5.3163) -- (7.5386,5.3163) -- (7.5390,5.3159) -- (7.5383,5.3159) -- (7.5387,5.3154) -- (7.5381,5.3154) -- (7.5384,5.3150) -- (7.5378,5.3150) -- (7.5382,5.3145) -- (7.5376,5.3145) -- (7.5379,5.3141) -- (7.5373,5.3141) -- (7.5376,5.3136) -- (7.5371,5.3136) -- (7.5373,5.3132) -- (7.5368,5.3132) -- (7.5370,5.3128) -- (7.5366,5.3127) -- (7.5368,5.3123) -- (7.5363,5.3123) -- (7.5365,5.3119) -- (7.5361,5.3118) -- (7.5362,5.3114) -- (7.5354,5.3113) -- (7.5359,5.3110) -- (7.5352,5.3109) -- (7.5357,5.3106) -- (7.5350,5.3104) -- (7.5354,5.3101) -- (7.5347,5.3099) -- (7.5351,5.3097) -- (7.5345,5.3095) -- (7.5349,5.3092) -- (7.5343,5.3090) -- (7.5346,5.3088) -- (7.5340,5.3086) -- (7.5343,5.3084) -- (7.5338,5.3082) -- (7.5340,5.3079) -- (7.5335,5.3077) -- (7.5338,5.3075) -- (7.5328,5.3070) -- (7.5335,5.3067) -- (7.5330,5.3066) -- (7.5332,5.3063) -- (7.5328,5.3062) -- (7.5331,5.3055) -- (7.5322,5.3055) -- (7.5325,5.3053) -- (7.5320,5.3051) -- (7.5322,5.3048) -- (7.5313,5.3044) -- (7.5319,5.3041) -- (7.5315,5.3040) -- (7.5316,5.3036) -- (7.5308,5.3035) -- (7.5314,5.3032) -- (7.5306,5.3030) -- (7.5311,5.3028) -- (7.5304,5.3026) -- (7.5308,5.3023) -- (7.5302,5.3021) -- (7.5305,5.3019) -- (7.5299,5.3017) -- (7.5303,5.3014) -- (7.5297,5.3013) -- (7.5300,5.3010) -- (7.5295,5.3008) -- (7.5297,5.3006) -- (7.5288,5.3002) -- (7.5294,5.2998) -- (7.5290,5.2998) -- (7.5292,5.2994) -- (7.5283,5.2993) -- (7.5289,5.2989) -- (7.5281,5.2988) -- (7.5286,5.2985) -- (7.5279,5.2984) -- (7.5283,5.2981) -- (7.5277,5.2979) -- (7.5280,5.2977) -- (7.5275,5.2975) -- (7.5278,5.2972) -- (7.5272,5.2970) -- (7.5275,5.2968) -- (7.5270,5.2966) -- (7.5275,5.2961) -- (7.5268,5.2961) -- (7.5272,5.2957) -- (7.5265,5.2957) -- (7.5269,5.2953) -- (7.5263,5.2953) -- (7.5266,5.2949) -- (7.5261,5.2948) -- (7.5263,5.2944) -- (7.5258,5.2944) -- (7.5260,5.2940) -- (7.5256,5.2940) -- (7.5258,5.2932) -- (7.5250,5.2933) -- (7.5253,5.2930) -- (7.5243,5.2926) -- (7.5250,5.2923) -- (7.5246,5.2922) -- (7.5248,5.2915) -- (7.5240,5.2915) -- (7.5243,5.2913) -- (7.5233,5.2909) -- (7.5240,5.2905) -- (7.5235,5.2905) -- (7.5237,5.2901) -- (7.5229,5.2900) -- (7.5234,5.2897) -- (7.5227,5.2895) -- (7.5231,5.2893) -- (7.5225,5.2891) -- (7.5229,5.2888) -- (7.5223,5.2887) -- (7.5226,5.2884) -- (7.5221,5.2882) -- (7.5223,5.2880) -- (7.5214,5.2876) -- (7.5220,5.2873) -- (7.5216,5.2872) -- (7.5218,5.2865) -- (7.5211,5.2865) -- (7.5215,5.2861) -- (7.5208,5.2861) -- (7.5212,5.2856) -- (7.5206,5.2856) -- (7.5209,5.2852) -- (7.5204,5.2852) -- (7.5206,5.2848) -- (7.5202,5.2848) -- (7.5203,5.2844) -- (7.5199,5.2843) -- (7.5202,5.2836) -- (7.5194,5.2837) -- (7.5196,5.2834) -- (7.5187,5.2830) -- (7.5193,5.2827) -- (7.5185,5.2826) -- (7.5190,5.2823) -- (7.5183,5.2821) -- (7.5187,5.2819) -- (7.5181,5.2817) -- (7.5185,5.2814) -- (7.5179,5.2813) -- (7.5182,5.2810) -- (7.5177,5.2808) -- (7.5181,5.2804) -- (7.5175,5.2804) -- (7.5178,5.2800) -- (7.5173,5.2800) -- (7.5175,5.2796) -- (7.5170,5.2795) -- (7.5172,5.2791) -- (7.5168,5.2791) -- (7.5171,5.2784) -- (7.5163,5.2784) -- (7.5165,5.2782) -- (7.5156,5.2778) -- (7.5162,5.2775) -- (7.5154,5.2773) -- (7.5159,5.2770) -- (7.5152,5.2769) -- (7.5156,5.2766) -- (7.5150,5.2765) -- (7.5154,5.2762) -- (7.5148,5.2760) -- (7.5151,5.2758) -- (7.5142,5.2754) -- (7.5148,5.2751) -- (7.5140,5.2750) -- (7.5145,5.2747) -- (7.5138,5.2745) -- (7.5142,5.2742) -- (7.5136,5.2741) -- (7.5139,5.2738) -- (7.5134,5.2736) -- (7.5137,5.2734) -- (7.5128,5.2730) -- (7.5134,5.2727) -- (7.5126,5.2726) -- (7.5131,5.2723) -- (7.5124,5.2721) -- (7.5128,5.2719) -- (7.5122,5.2717) -- (7.5125,5.2715) -- (7.5120,5.2713) -- (7.5123,5.2710) -- (7.5114,5.2706) -- (7.5119,5.2703) -- (7.5112,5.2702) -- (7.5117,5.2699) -- (7.5110,5.2698) -- (7.5114,5.2695) -- (7.5108,5.2693) -- (7.5111,5.2691) -- (7.5106,5.2689) -- (7.5111,5.2685) -- (7.5104,5.2685) -- (7.5107,5.2681) -- (7.5102,5.2680) -- (7.5104,5.2677) -- (7.5100,5.2676) -- (7.5101,5.2673) -- (7.5094,5.2671) -- (7.5098,5.2669) -- (7.5092,5.2667) -- (7.5096,5.2665) -- (7.5090,5.2663) -- (7.5093,5.2660) -- (7.5088,5.2658) -- (7.5092,5.2654) -- (7.5086,5.2654) -- (7.5089,5.2650) -- (7.5084,5.2650) -- (7.5086,5.2646) -- (7.5082,5.2646) -- (7.5084,5.2639) -- (7.5077,5.2639) -- (7.5081,5.2635) -- (7.5075,5.2635) -- (7.5078,5.2631) -- (7.5072,5.2631) -- (7.5075,5.2627) -- (7.5070,5.2627) -- (7.5072,5.2623) -- (7.5064,5.2622) -- (7.5069,5.2619) -- (7.5063,5.2617) -- (7.5066,5.2615) -- (7.5061,5.2613) -- (7.5063,5.2611) -- (7.5054,5.2607) -- (7.5060,5.2604) -- (7.5053,5.2603) -- (7.5057,5.2600) -- (7.5051,5.2598) -- (7.5055,5.2596) -- (7.5049,5.2594) -- (7.5052,5.2592) -- (7.5043,5.2588) -- (7.5049,5.2585) -- (7.5042,5.2583) -- (7.5046,5.2581) -- (7.5040,5.2579) -- (7.5043,5.2577) -- (7.5038,5.2575) -- (7.5040,5.2573) -- (7.5032,5.2569) -- (7.5037,5.2566) -- (7.5030,5.2564) -- (7.5034,5.2562) -- (7.5029,5.2560) -- (7.5032,5.2558) -- (7.5027,5.2556) -- (7.5031,5.2551) -- (7.5025,5.2552) -- (7.5028,5.2548) -- (7.5023,5.2547) -- (7.5025,5.2544) -- (7.5021,5.2543) -- (7.5023,5.2536) -- (7.5016,5.2537) -- (7.5020,5.2532) -- (7.5014,5.2533) -- (7.5016,5.2529) -- (7.5011,5.2528) -- (7.5013,5.2525) -- (7.5005,5.2524) -- (7.5010,5.2521) -- (7.5004,5.2519) -- (7.5007,5.2517) -- (7.5002,5.2515) -- (7.5005,5.2513) -- (7.4996,5.2509) -- (7.5002,5.2506) -- (7.4995,5.2505) -- (7.4999,5.2502) -- (7.4993,5.2500) -- (7.4996,5.2498) -- (7.4991,5.2496) -- (7.4995,5.2492) -- (7.4989,5.2492) -- (7.4992,5.2488) -- (7.4987,5.2488) -- (7.4989,5.2484) -- (7.4981,5.2483) -- (7.4986,5.2480) -- (7.4980,5.2479) -- (7.4983,5.2477) -- (7.4978,5.2475) -- (7.4981,5.2473) -- (7.4972,5.2469) -- (7.4977,5.2466) -- (7.4971,5.2464) -- (7.4975,5.2462) -- (7.4969,5.2460) -- (7.4972,5.2458) -- (7.4963,5.2454) -- (7.4969,5.2451) -- (7.4961,5.2450) -- (7.4966,5.2447) -- (7.4960,5.2446) -- (7.4963,5.2443) -- (7.4958,5.2441) -- (7.4962,5.2437) -- (7.4956,5.2437) -- (7.4959,5.2433) -- (7.4954,5.2433) -- (7.4956,5.2430) -- (7.4949,5.2429) -- (7.4953,5.2426) -- (7.4947,5.2424) -- (7.4950,5.2422) -- (7.4945,5.2420) -- (7.4950,5.2416) -- (7.4944,5.2416) -- (7.4946,5.2412) -- (7.4942,5.2412) -- (7.4943,5.2408) -- (7.4936,5.2407) -- (7.4940,5.2405) -- (7.4934,5.2403) -- (7.4938,5.2401) -- (7.4933,5.2399) -- (7.4937,5.2395) -- (7.4931,5.2395) -- (7.4934,5.2391) -- (7.4929,5.2391) -- (7.4930,5.2387) -- (7.4923,5.2386) -- (7.4928,5.2383) -- (7.4922,5.2382) -- (7.4925,5.2380) -- (7.4920,5.2378) -- (7.4924,5.2374) -- (7.4918,5.2374) -- (7.4921,5.2370) -- (7.4916,5.2370) -- (7.4918,5.2366) -- (7.4911,5.2365) -- (7.4915,5.2362) -- (7.4909,5.2361) -- (7.4912,5.2359) -- (7.4903,5.2355) -- (7.4909,5.2352) -- (7.4902,5.2351) -- (7.4906,5.2348) -- (7.4901,5.2346) -- (7.4903,5.2344) -- (7.4895,5.2340) -- (7.4900,5.2337) -- (7.4893,5.2336) -- (7.4897,5.2334) -- (7.4892,5.2332) -- (7.4895,5.2330) -- (7.4886,5.2326) -- (7.4891,5.2323) -- (7.4885,5.2322) -- (7.4889,5.2319) -- (7.4883,5.2318) -- (7.4886,5.2315) -- (7.4878,5.2312) -- (7.4883,5.2309) -- (7.4876,5.2307) -- (7.4880,5.2305) -- (7.4875,5.2303) -- (7.4879,5.2299) -- (7.4873,5.2299) -- (7.4876,5.2295) -- (7.4871,5.2295) -- (7.4873,5.2292) -- (7.4866,5.2291) -- (7.4870,5.2288) -- (7.4864,5.2286) -- (7.4867,5.2284) -- (7.4858,5.2280) -- (7.4864,5.2277) -- (7.4857,5.2276) -- (7.4861,5.2274) -- (7.4856,5.2272) -- (7.4858,5.2270) -- (7.4850,5.2266) -- (7.4855,5.2263) -- (7.4849,5.2262) -- (7.4852,5.2260) -- (7.4848,5.2258) -- (7.4852,5.2254) -- (7.4846,5.2254) -- (7.4848,5.2250) -- (7.4844,5.2250) -- (7.4846,5.2243) -- (7.4839,5.2244) -- (7.4843,5.2240) -- (7.4838,5.2240) -- (7.4839,5.2236) -- (7.4832,5.2235) -- (7.4836,5.2232) -- (7.4831,5.2231) -- (7.4834,5.2229) -- (7.4829,5.2227) -- (7.4833,5.2223) -- (7.4827,5.2223) -- (7.4829,5.2219) -- (7.4825,5.2219) -- (7.4827,5.2212) -- (7.4821,5.2213) -- (7.4824,5.2209) -- (7.4819,5.2209) -- (7.4821,5.2205) -- (7.4813,5.2204) -- (7.4818,5.2202) -- (7.4812,5.2200) -- (7.4815,5.2198) -- (7.4807,5.2194) -- (7.4812,5.2191) -- (7.4806,5.2190) -- (7.4809,5.2188) -- (7.4804,5.2186) -- (7.4808,5.2182) -- (7.4802,5.2182) -- (7.4805,5.2179) -- (7.4801,5.2178) -- (7.4803,5.2172) -- (7.4796,5.2172) -- (7.4799,5.2168) -- (7.4794,5.2168) -- (7.4796,5.2165) -- (7.4789,5.2164) -- (7.4793,5.2161) -- (7.4788,5.2160) -- (7.4790,5.2157) -- (7.4782,5.2154) -- (7.4787,5.2151) -- (7.4781,5.2150) -- (7.4784,5.2147) -- (7.4780,5.2145) -- (7.4783,5.2142) -- (7.4778,5.2142) -- (7.4780,5.2138) -- (7.4772,5.2137) -- (7.4777,5.2134) -- (7.4771,5.2133) -- (7.4774,5.2131) -- (7.4766,5.2127) -- (7.4771,5.2124) -- (7.4765,5.2123) -- (7.4768,5.2121) -- (7.4763,5.2119) -- (7.4767,5.2115) -- (7.4762,5.2115) -- (7.4764,5.2112) -- (7.4760,5.2111) -- (7.4762,5.2105) -- (7.4755,5.2105) -- (7.4758,5.2101) -- (7.4754,5.2101) -- (7.4756,5.2095) -- (7.4749,5.2095) -- (7.4752,5.2091) -- (7.4747,5.2091) -- (7.4749,5.2088) -- (7.4742,5.2087) -- (7.4746,5.2084) -- (7.4741,5.2083) -- (7.4743,5.2081) -- (7.4736,5.2077) -- (7.4740,5.2074) -- (7.4735,5.2073) -- (7.4737,5.2071) -- (7.4729,5.2067) -- (7.4734,5.2064) -- (7.4728,5.2063) -- (7.4731,5.2061) -- (7.4727,5.2059) -- (7.4731,5.2055) -- (7.4725,5.2055) -- (7.4727,5.2052) -- (7.4720,5.2051) -- (7.4724,5.2048) -- (7.4719,5.2047) -- (7.4721,5.2045) -- (7.4713,5.2041) -- (7.4718,5.2038) -- (7.4713,5.2037) -- (7.4715,5.2035) -- (7.4707,5.2031) -- (7.4712,5.2028) -- (7.4706,5.2027) -- (7.4709,5.2025) -- (7.4705,5.2023) -- (7.4709,5.2019) -- (7.4704,5.2019) -- (7.4705,5.2016) -- (7.4698,5.2015) -- (7.4702,5.2012) -- (7.4697,5.2011) -- (7.4699,5.2009) -- (7.4692,5.2005) -- (7.4696,5.2003) -- (7.4691,5.2001) -- (7.4694,5.1999) -- (7.4686,5.1995) -- (7.4690,5.1993) -- (7.4685,5.1991) -- (7.4688,5.1989) -- (7.4680,5.1986) -- (7.4684,5.1983) -- (7.4679,5.1982) -- (7.4682,5.1979) -- (7.4673,5.1976) -- (7.4678,5.1973) -- (7.4673,5.1972) -- (7.4676,5.1970) -- (7.4667,5.1966) -- (7.4673,5.1963) -- (7.4667,5.1962) -- (7.4670,5.1960) -- (7.4661,5.1956) -- (7.4667,5.1954) -- (7.4661,5.1952) -- (7.4664,5.1950) -- (7.4655,5.1947) -- (7.4661,5.1944) -- (7.4655,5.1943) -- (7.4658,5.1940) -- (7.4654,5.1939) -- (7.4657,5.1935) -- (7.4652,5.1935) -- (7.4653,5.1932) -- (7.4647,5.1931) -- (7.4650,5.1928) -- (7.4646,5.1927) -- (7.4650,5.1923) -- (7.4644,5.1923) -- (7.4646,5.1920) -- (7.4639,5.1919) -- (7.4643,5.1916) -- (7.4638,5.1915) -- (7.4641,5.1913) -- (7.4633,5.1909) -- (7.4637,5.1906) -- (7.4632,5.1905) -- (7.4635,5.1903) -- (7.4627,5.1899) -- (7.4631,5.1897) -- (7.4626,5.1895) -- (7.4631,5.1891) -- (7.4625,5.1892) -- (7.4627,5.1888) -- (7.4623,5.1888) -- (7.4625,5.1882) -- (7.4619,5.1882) -- (7.4621,5.1879) -- (7.4617,5.1878) -- (7.4619,5.1872) -- (7.4613,5.1872) -- (7.4615,5.1869) -- (7.4608,5.1868) -- (7.4612,5.1866) -- (7.4607,5.1864) -- (7.4609,5.1862) -- (7.4602,5.1859) -- (7.4606,5.1856) -- (7.4601,5.1855) -- (7.4606,5.1851) -- (7.4600,5.1851) -- (7.4602,5.1848) -- (7.4595,5.1847) -- (7.4599,5.1844) -- (7.4594,5.1843) -- (7.4596,5.1841) -- (7.4589,5.1837) -- (7.4593,5.1835) -- (7.4588,5.1833) -- (7.4592,5.1829) -- (7.4587,5.1829) -- (7.4589,5.1826) -- (7.4582,5.1825) -- (7.4586,5.1823) -- (7.4581,5.1821) -- (7.4585,5.1817) -- (7.4580,5.1818) -- (7.4582,5.1814) -- (7.4574,5.1814) -- (7.4579,5.1811) -- (7.4574,5.1810) -- (7.4578,5.1806) -- (7.4572,5.1806) -- (7.4574,5.1802) -- (7.4571,5.1802) -- (7.4572,5.1796) -- (7.4567,5.1796) -- (7.4568,5.1793) -- (7.4561,5.1792) -- (7.4565,5.1790) -- (7.4561,5.1788) -- (7.4565,5.1784) -- (7.4559,5.1785) -- (7.4561,5.1781) -- (7.4554,5.1780) -- (7.4558,5.1778) -- (7.4553,5.1777) -- (7.4557,5.1773) -- (7.4552,5.1773) -- (7.4554,5.1770) -- (7.4547,5.1769) -- (7.4551,5.1766) -- (7.4546,5.1765) -- (7.4550,5.1761) -- (7.4545,5.1761) -- (7.4546,5.1758) -- (7.4540,5.1757) -- (7.4544,5.1755) -- (7.4539,5.1753) -- (7.4543,5.1749) -- (7.4538,5.1749) -- (7.4539,5.1746) -- (7.4533,5.1745) -- (7.4536,5.1743) -- (7.4532,5.1741) -- (7.4535,5.1738) -- (7.4531,5.1738) -- (7.4532,5.1735) -- (7.4526,5.1734) -- (7.4529,5.1731) -- (7.4525,5.1730) -- (7.4528,5.1726) -- (7.4523,5.1726) -- (7.4526,5.1720) -- (7.4519,5.1720) -- (7.4522,5.1717) -- (7.4518,5.1717) -- (7.4519,5.1711) -- (7.4514,5.1711) -- (7.4516,5.1708) -- (7.4509,5.1707) -- (7.4513,5.1705) -- (7.4508,5.1703) -- (7.4512,5.1699) -- (7.4507,5.1700) -- (7.4508,5.1696) -- (7.4502,5.1696) -- (7.4505,5.1693) -- (7.4501,5.1692) -- (7.4505,5.1688) -- (7.4500,5.1688) -- (7.4501,5.1685) -- (7.4495,5.1684) -- (7.4498,5.1682) -- (7.4490,5.1679) -- (7.4495,5.1676) -- (7.4490,5.1675) -- (7.4492,5.1673) -- (7.4485,5.1669) -- (7.4489,5.1667) -- (7.4485,5.1665) -- (7.4488,5.1662) -- (7.4483,5.1662) -- (7.4485,5.1659) -- (7.4479,5.1658) -- (7.4482,5.1655) -- (7.4478,5.1654) -- (7.4481,5.1650) -- (7.4476,5.1650) -- (7.4479,5.1644) -- (7.4472,5.1645) -- (7.4475,5.1641) -- (7.4471,5.1641) -- (7.4472,5.1635) -- (7.4467,5.1635) -- (7.4469,5.1632) -- (7.4462,5.1631) -- (7.4466,5.1629) -- (7.4462,5.1628) -- (7.4465,5.1624) -- (7.4460,5.1624) -- (7.4462,5.1618) -- (7.4456,5.1618) -- (7.4458,5.1615) -- (7.4451,5.1615) -- (7.4455,5.1612) -- (7.4451,5.1611) -- (7.4454,5.1607) -- (7.4449,5.1607) -- (7.4451,5.1604) -- (7.4445,5.1603) -- (7.4448,5.1601) -- (7.4440,5.1598) -- (7.4445,5.1595) -- (7.4440,5.1594) -- (7.4442,5.1592) -- (7.4435,5.1588) -- (7.4439,5.1586) -- (7.4435,5.1585) -- (7.4438,5.1581) -- (7.4433,5.1581) -- (7.4436,5.1575) -- (7.4430,5.1576) -- (7.4432,5.1572) -- (7.4428,5.1572) -- (7.4429,5.1566) -- (7.4424,5.1566) -- (7.4426,5.1563) -- (7.4420,5.1562) -- (7.4423,5.1560) -- (7.4415,5.1557) -- (7.4420,5.1555) -- (7.4415,5.1553) -- (7.4419,5.1550) -- (7.4414,5.1550) -- (7.4415,5.1547) -- (7.4409,5.1546) -- (7.4412,5.1544) -- (7.4404,5.1541) -- (7.4409,5.1538) -- (7.4404,5.1537) -- (7.4406,5.1535) -- (7.4400,5.1531) -- (7.4403,5.1529) -- (7.4399,5.1528) -- (7.4402,5.1524) -- (7.4398,5.1524) -- (7.4400,5.1518) -- (7.4394,5.1519) -- (7.4396,5.1516) -- (7.4389,5.1515) -- (7.4393,5.1512) -- (7.4389,5.1511) -- (7.4392,5.1508) -- (7.4388,5.1507) -- (7.4389,5.1502) -- (7.4384,5.1502) -- (7.4386,5.1499) -- (7.4379,5.1498) -- (7.4383,5.1496) -- (7.4378,5.1494) -- (7.4382,5.1491) -- (7.4377,5.1491) -- (7.4379,5.1485) -- (7.4373,5.1486) -- (7.4375,5.1482) -- (7.4369,5.1482) -- (7.4372,5.1479) -- (7.4368,5.1478) -- (7.4371,5.1475) -- (7.4367,5.1474) -- (7.4369,5.1469) -- (7.4363,5.1469) -- (7.4365,5.1466) -- (7.4359,5.1465) -- (7.4362,5.1463) -- (7.4358,5.1461) -- (7.4361,5.1458) -- (7.4357,5.1458) -- (7.4358,5.1452) -- (7.4353,5.1453) -- (7.4355,5.1450) -- (7.4348,5.1449) -- (7.4352,5.1447) -- (7.4344,5.1444) -- (7.4349,5.1441) -- (7.4344,5.1440) -- (7.4348,5.1436) -- (7.4343,5.1436) -- (7.4344,5.1433) -- (7.4338,5.1432) -- (7.4341,5.1430) -- (7.4334,5.1427) -- (7.4338,5.1425) -- (7.4334,5.1423) -- (7.4337,5.1420) -- (7.4333,5.1420) -- (7.4335,5.1414) -- (7.4329,5.1415) -- (7.4331,5.1412) -- (7.4324,5.1411) -- (7.4328,5.1409) -- (7.4324,5.1407) -- (7.4327,5.1404) -- (7.4323,5.1404) -- (7.4324,5.1398) -- (7.4319,5.1398) -- (7.4320,5.1395) -- (7.4315,5.1395) -- (7.4318,5.1392) -- (7.4310,5.1389) -- (7.4314,5.1387) -- (7.4310,5.1386) -- (7.4314,5.1382) -- (7.4309,5.1382) -- (7.4311,5.1377) -- (7.4305,5.1377) -- (7.4307,5.1374) -- (7.4301,5.1373) -- (7.4304,5.1371) -- (7.4297,5.1368) -- (7.4301,5.1366) -- (7.4297,5.1364) -- (7.4300,5.1361) -- (7.4296,5.1361) -- (7.4298,5.1355) -- (7.4292,5.1356) -- (7.4294,5.1353) -- (7.4287,5.1352) -- (7.4291,5.1350) -- (7.4287,5.1348) -- (7.4290,5.1345) -- (7.4286,5.1345) -- (7.4287,5.1339) -- (7.4282,5.1340) -- (7.4283,5.1337) -- (7.4278,5.1336) -- (7.4281,5.1334) -- (7.4274,5.1331) -- (7.4277,5.1328) -- (7.4274,5.1327) -- (7.4276,5.1324) -- (7.4272,5.1324) -- (7.4274,5.1318) -- (7.4269,5.1318) -- (7.4270,5.1316) -- (7.4265,5.1315) -- (7.4267,5.1313) -- (7.4261,5.1309) -- (7.4264,5.1307) -- (7.4257,5.1304) -- (7.4261,5.1302) -- (7.4257,5.1301) -- (7.4260,5.1297) -- (7.4256,5.1297) -- (7.4258,5.1292) -- (7.4252,5.1292) -- (7.4254,5.1289) -- (7.4248,5.1288) -- (7.4251,5.1286) -- (7.4244,5.1283) -- (7.4248,5.1281) -- (7.4244,5.1280) -- (7.4247,5.1276) -- (7.4242,5.1276) -- (7.4244,5.1271) -- (7.4239,5.1271) -- (7.4240,5.1268) -- (7.4235,5.1267) -- (7.4237,5.1265) -- (7.4231,5.1262) -- (7.4234,5.1260) -- (7.4231,5.1259) -- (7.4233,5.1256) -- (7.4229,5.1255) -- (7.4231,5.1250) -- (7.4226,5.1250) -- (7.4228,5.1245) -- (7.4222,5.1245) -- (7.4224,5.1242) -- (7.4218,5.1241) -- (7.4221,5.1239) -- (7.4214,5.1236) -- (7.4218,5.1234) -- (7.4214,5.1233) -- (7.4217,5.1230) -- (7.4213,5.1229) -- (7.4214,5.1224) -- (7.4210,5.1224) -- (7.4212,5.1219) -- (7.4206,5.1219) -- (7.4208,5.1216) -- (7.4202,5.1216) -- (7.4205,5.1214) -- (7.4198,5.1211) -- (7.4202,5.1208) -- (7.4198,5.1207) -- (7.4201,5.1204) -- (7.4197,5.1204) -- (7.4198,5.1199) -- (7.4193,5.1199) -- (7.4195,5.1193) -- (7.4190,5.1194) -- (7.4191,5.1191) -- (7.4186,5.1190) -- (7.4188,5.1188) -- (7.4182,5.1185) -- (7.4185,5.1183) -- (7.4178,5.1180) -- (7.4182,5.1178) -- (7.4178,5.1176) -- (7.4181,5.1173) -- (7.4177,5.1173) -- (7.4179,5.1168) -- (7.4174,5.1168) -- (7.4176,5.1163) -- (7.4171,5.1163) -- (7.4172,5.1160) -- (7.4166,5.1159) -- (7.4169,5.1157) -- (7.4163,5.1154) -- (7.4166,5.1152) -- (7.4159,5.1149) -- (7.4163,5.1147) -- (7.4159,5.1146) -- (7.4162,5.1143) -- (7.4158,5.1143) -- (7.4159,5.1137) -- (7.4155,5.1138) -- (7.4157,5.1132) -- (7.4151,5.1133) -- (7.4153,5.1130) -- (7.4147,5.1129) -- (7.4150,5.1127) -- (7.4144,5.1124) -- (7.4147,5.1122) -- (7.4140,5.1119) -- (7.4144,5.1117) -- (7.4140,5.1116) -- (7.4143,5.1113) -- (7.4139,5.1112) -- (7.4140,5.1107) -- (7.4136,5.1107) -- (7.4137,5.1102) -- (7.4132,5.1102) -- (7.4134,5.1097) -- (7.4129,5.1097) -- (7.4130,5.1095) -- (7.4125,5.1094) -- (7.4128,5.1092) -- (7.4121,5.1089) -- (7.4124,5.1087) -- (7.4118,5.1084) -- (7.4121,5.1082) -- (7.4114,5.1079) -- (7.4118,5.1077) -- (7.4114,5.1076) -- (7.4117,5.1073) -- (7.4113,5.1072) -- (7.4115,5.1068) -- (7.4110,5.1068) -- (7.4112,5.1062) -- (7.4107,5.1063) -- (7.4109,5.1057) -- (7.4104,5.1058) -- (7.4105,5.1055) -- (7.4100,5.1054) -- (7.4102,5.1052) -- (7.4096,5.1049) -- (7.4099,5.1047) -- (7.4093,5.1045) -- (7.4096,5.1042) -- (7.4089,5.1040) -- (7.4093,5.1037) -- (7.4089,5.1036) -- (7.4092,5.1033) -- (7.4088,5.1033) -- (7.4089,5.1028) -- (7.4085,5.1028) -- (7.4086,5.1023) -- (7.4082,5.1023) -- (7.4084,5.1018) -- (7.4079,5.1018) -- (7.4081,5.1013) -- (7.4076,5.1014) -- (7.4077,5.1011) -- (7.4072,5.1010) -- (7.4074,5.1008) -- (7.4068,5.1005) -- (7.4071,5.1003) -- (7.4065,5.1000) -- (7.4068,5.0998) -- (7.4062,5.0996) -- (7.4065,5.0994) -- (7.4058,5.0991) -- (7.4062,5.0989) -- (7.4058,5.0987) -- (7.4061,5.0984) -- (7.4057,5.0984) -- (7.4058,5.0980) -- (7.4054,5.0979) -- (7.4055,5.0975) -- (7.4051,5.0975) -- (7.4052,5.0970) -- (7.4048,5.0970) -- (7.4049,5.0965) -- (7.4045,5.0965) -- (7.4047,5.0960) -- (7.4042,5.0960) -- (7.4043,5.0958) -- (7.4038,5.0957) -- (7.4042,5.0953) -- (7.4037,5.0954) -- (7.4039,5.0948) -- (7.4034,5.0949) -- (7.4035,5.0946) -- (7.4030,5.0945) -- (7.4033,5.0944) -- (7.4027,5.0941) -- (7.4030,5.0939) -- (7.4024,5.0936) -- (7.4027,5.0934) -- (7.4021,5.0931) -- (7.4024,5.0929) -- (7.4017,5.0926) -- (7.4021,5.0924) -- (7.4014,5.0922) -- (7.4018,5.0920) -- (7.4011,5.0917) -- (7.4015,5.0915) -- (7.4008,5.0912) -- (7.4012,5.0910) -- (7.4008,5.0909) -- (7.4010,5.0906) -- (7.4004,5.0906) -- (7.4007,5.0903) -- (7.4000,5.0901) -- (7.4004,5.0899) -- (7.3997,5.0896);
\fill[blue!80!black] (3.6000,2.4000) circle (1.2pt);
\fill[blue!80!black] (6.4160,2.8687) circle (1.2pt);
\fill[blue!80!black] (4.4337,4.7866) circle (1.2pt);
\fill[blue!80!black] (8.0038,5.6872) circle (1.2pt);
\fill[blue!80!black] (3.7372,6.5364) circle (1.2pt);
\fill[blue!80!black] (3.9108,6.5665) circle (1.2pt);
\fill[blue!80!black] (3.8417,6.5527) circle (1.2pt);
\fill[blue!80!black] (3.8635,6.5561) circle (1.2pt);
\fill[blue!80!black] (3.8562,6.5542) circle (1.2pt);
\fill[blue!80!black] (3.8590,6.5541) circle (1.2pt);
\fill[blue!80!black] (3.8577,6.5526) circle (1.2pt);
\fill[blue!80!black] (3.8594,6.5523) circle (1.2pt);
\fill[blue!80!black] (3.8584,6.5497) circle (1.2pt);
\fill[blue!80!black] (3.8608,6.5495) circle (1.2pt);
\fill[blue!80!black] (3.8598,6.5481) circle (1.2pt);
\fill[blue!80!black] (3.8628,6.5475) circle (1.2pt);
\fill[blue!80!black] (3.8609,6.5465) circle (1.2pt);
\fill[blue!80!black] (3.8628,6.5456) circle (1.2pt);
\fill[blue!80!black] (3.8609,6.5441) circle (1.2pt);
\fill[blue!80!black] (3.8631,6.5439) circle (1.2pt);
\fill[blue!80!black] (3.8623,6.5425) circle (1.2pt);
\fill[blue!80!black] (3.8650,6.5418) circle (1.2pt);
\fill[blue!80!black] (3.8633,6.5409) circle (1.2pt);
\fill[blue!80!black] (3.8668,6.5391) circle (1.2pt);
\fill[blue!80!black] (3.8643,6.5380) circle (1.2pt);
\fill[blue!80!black] (3.8666,6.5372) circle (1.2pt);
\fill[blue!80!black] (3.8653,6.5363) circle (1.2pt);
\fill[blue!80!black] (3.8683,6.5345) circle (1.2pt);
\fill[blue!80!black] (3.8664,6.5335) circle (1.2pt);
\fill[blue!80!black] (3.8683,6.5326) circle (1.2pt);
\fill[blue!80!black] (3.8664,6.5310) circle (1.2pt);
\fill[blue!80!black] (3.8686,6.5309) circle (1.2pt);
\fill[blue!80!black] (3.8669,6.5281) circle (1.2pt);
\fill[blue!80!black] (3.8700,6.5281) circle (1.2pt);
\fill[blue!80!black] (3.8686,6.5266) circle (1.2pt);
\fill[blue!80!black] (3.8704,6.5263) circle (1.2pt);
\fill[blue!80!black] (3.8693,6.5237) circle (1.2pt);
\fill[blue!80!black] (3.8718,6.5235) circle (1.2pt);
\fill[blue!80!black] (3.8708,6.5221) circle (1.2pt);
\fill[blue!80!black] (3.8737,6.5214) circle (1.2pt);
\fill[blue!80!black] (3.8719,6.5205) circle (1.2pt);
\fill[blue!80!black] (3.8755,6.5187) circle (1.2pt);
\fill[blue!80!black] (3.8729,6.5176) circle (1.2pt);
\fill[blue!80!black] (3.8752,6.5168) circle (1.2pt);
\fill[blue!80!black] (3.8728,6.5151) circle (1.2pt);
\fill[blue!80!black] (3.8754,6.5150) circle (1.2pt);
\fill[blue!80!black] (3.8743,6.5136) circle (1.2pt);
\fill[blue!80!black] (3.8774,6.5129) circle (1.2pt);
\fill[blue!80!black] (3.8755,6.5120) circle (1.2pt);
\fill[blue!80!black] (3.8773,6.5111) circle (1.2pt);
\fill[blue!80!black] (3.8756,6.5096) circle (1.2pt);
\fill[blue!80!black] (3.8777,6.5093) circle (1.2pt);
\fill[blue!80!black] (3.8763,6.5066) circle (1.2pt);
\fill[blue!80!black] (3.8790,6.5065) circle (1.2pt);
\fill[blue!80!black] (3.8779,6.5051) circle (1.2pt);
\fill[blue!80!black] (3.8810,6.5045) circle (1.2pt);
\fill[blue!80!black] (3.8791,6.5035) circle (1.2pt);
\fill[blue!80!black] (3.8809,6.5026) circle (1.2pt);
\fill[blue!80!black] (3.8792,6.5011) circle (1.2pt);
\fill[blue!80!black] (3.8813,6.5008) circle (1.2pt);
\fill[blue!80!black] (3.8800,6.4981) circle (1.2pt);
\fill[blue!80!black] (3.8826,6.4980) circle (1.2pt);
\fill[blue!80!black] (3.8816,6.4966) circle (1.2pt);
\fill[blue!80!black] (3.8845,6.4959) circle (1.2pt);
\fill[blue!80!black] (3.8828,6.4950) circle (1.2pt);
\fill[blue!80!black] (3.8862,6.4932) circle (1.2pt);
\fill[blue!80!black] (3.8838,6.4921) circle (1.2pt);
\fill[blue!80!black] (3.8859,6.4913) circle (1.2pt);
\fill[blue!80!black] (3.8837,6.4897) circle (1.2pt);
\fill[blue!80!black] (3.8862,6.4895) circle (1.2pt);
\fill[blue!80!black] (3.8853,6.4881) circle (1.2pt);
\fill[blue!80!black] (3.8880,6.4874) circle (1.2pt);
\fill[blue!80!black] (3.8864,6.4865) circle (1.2pt);
\fill[blue!80!black] (3.8896,6.4847) circle (1.2pt);
\fill[blue!80!black] (3.8875,6.4837) circle (1.2pt);
\fill[blue!80!black] (3.8894,6.4828) circle (1.2pt);
\fill[blue!80!black] (3.8876,6.4812) circle (1.2pt);
\fill[blue!80!black] (3.8897,6.4810) circle (1.2pt);
\fill[blue!80!black] (3.8884,6.4783) circle (1.2pt);
\fill[blue!80!black] (3.8910,6.4782) circle (1.2pt);
\fill[blue!80!black] (3.8901,6.4768) circle (1.2pt);
\fill[blue!80!black] (3.8929,6.4761) circle (1.2pt);
\fill[blue!80!black] (3.8912,6.4752) circle (1.2pt);
\fill[blue!80!black] (3.8945,6.4733) circle (1.2pt);
\fill[blue!80!black] (3.8923,6.4723) circle (1.2pt);
\fill[blue!80!black] (3.8943,6.4714) circle (1.2pt);
\fill[blue!80!black] (3.8924,6.4699) circle (1.2pt);
\fill[blue!80!black] (3.8945,6.4697) circle (1.2pt);
\fill[blue!80!black] (3.8933,6.4670) circle (1.2pt);
\fill[blue!80!black] (3.8959,6.4668) circle (1.2pt);
\fill[blue!80!black] (3.8949,6.4654) circle (1.2pt);
\fill[blue!80!black] (3.8977,6.4647) circle (1.2pt);
\fill[blue!80!black] (3.8961,6.4638) circle (1.2pt);
\fill[blue!80!black] (3.8992,6.4620) circle (1.2pt);
\fill[blue!80!black] (3.8972,6.4610) circle (1.2pt);
\fill[blue!80!black] (3.9008,6.4592) circle (1.2pt);
\fill[blue!80!black] (3.8983,6.4581) circle (1.2pt);
\fill[blue!80!black] (3.9004,6.4573) circle (1.2pt);
\fill[blue!80!black] (3.8984,6.4557) circle (1.2pt);
\fill[blue!80!black] (3.9006,6.4555) circle (1.2pt);
\fill[blue!80!black] (3.8992,6.4528) circle (1.2pt);
\fill[blue!80!black] (3.9020,6.4527) circle (1.2pt);
\fill[blue!80!black] (3.9010,6.4512) circle (1.2pt);
\fill[blue!80!black] (3.9038,6.4506) circle (1.2pt);
\fill[blue!80!black] (3.9022,6.4496) circle (1.2pt);
\fill[blue!80!black] (3.9053,6.4478) circle (1.2pt);
\fill[blue!80!black] (3.9033,6.4468) circle (1.2pt);
\fill[blue!80!black] (3.9069,6.4450) circle (1.2pt);
\fill[blue!80!black] (3.9045,6.4439) circle (1.2pt);
\fill[blue!80!black] (3.9065,6.4431) circle (1.2pt);
\fill[blue!80!black] (3.9046,6.4415) circle (1.2pt);
\fill[blue!80!black] (3.9067,6.4413) circle (1.2pt);
\fill[blue!80!black] (3.9055,6.4386) circle (1.2pt);
\fill[blue!80!black] (3.9080,6.4385) circle (1.2pt);
\fill[blue!80!black] (3.9064,6.4357) circle (1.2pt);
\fill[blue!80!black] (3.9093,6.4356) circle (1.2pt);
\fill[blue!80!black] (3.9083,6.4342) circle (1.2pt);
\fill[blue!80!black] (3.9111,6.4335) circle (1.2pt);
\fill[blue!80!black] (3.9095,6.4326) circle (1.2pt);
\fill[blue!80!black] (3.9126,6.4307) circle (1.2pt);
\fill[blue!80!black] (3.9107,6.4297) circle (1.2pt);
\fill[blue!80!black] (3.9141,6.4279) circle (1.2pt);
\fill[blue!80!black] (3.9118,6.4269) circle (1.2pt);
\fill[blue!80!black] (3.9137,6.4260) circle (1.2pt);
\fill[blue!80!black] (3.9121,6.4245) circle (1.2pt);
\fill[blue!80!black] (7.8436,5.9179) circle (1.2pt);
\fill[blue!80!black] (7.8563,5.9150) circle (1.2pt);
\fill[blue!80!black] (7.8556,5.9149) circle (1.2pt);
\fill[blue!80!black] (7.8561,5.9145) circle (1.2pt);
\fill[blue!80!black] (7.8553,5.9144) circle (1.2pt);
\fill[blue!80!black] (7.8559,5.9140) circle (1.2pt);
\fill[blue!80!black] (7.8550,5.9139) circle (1.2pt);
\fill[blue!80!black] (7.8557,5.9135) circle (1.2pt);
\fill[blue!80!black] (7.8547,5.9134) circle (1.2pt);
\fill[blue!80!black] (7.8555,5.9129) circle (1.2pt);
\fill[blue!80!black] (7.8544,5.9129) circle (1.2pt);
\fill[blue!80!black] (7.8553,5.9124) circle (1.2pt);
\fill[blue!80!black] (7.8541,5.9124) circle (1.2pt);
\fill[blue!80!black] (7.8552,5.9119) circle (1.2pt);
\fill[blue!80!black] (7.8538,5.9119) circle (1.2pt);
\fill[blue!80!black] (7.8550,5.9114) circle (1.2pt);
\fill[blue!80!black] (7.8535,5.9115) circle (1.2pt);
\fill[blue!80!black] (7.8549,5.9109) circle (1.2pt);
\fill[blue!80!black] (7.8531,5.9110) circle (1.2pt);
\fill[blue!80!black] (7.8547,5.9104) circle (1.2pt);
\fill[blue!80!black] (7.8528,5.9105) circle (1.2pt);
\fill[blue!80!black] (7.8546,5.9098) circle (1.2pt);
\fill[blue!80!black] (7.8524,5.9100) circle (1.2pt);
\fill[blue!80!black] (7.8545,5.9093) circle (1.2pt);
\fill[blue!80!black] (7.8533,5.9094) circle (1.2pt);
\fill[blue!80!black] (7.8525,5.9052) circle (1.2pt);
\fill[blue!80!black] (7.8513,5.9053) circle (1.2pt);
\fill[blue!80!black] (7.8495,5.8969) circle (1.2pt);
\fill[blue!80!black] (7.8473,5.8972) circle (1.2pt);
\fill[blue!80!black] (7.8472,5.8951) circle (1.2pt);
\fill[blue!80!black] (7.8454,5.8952) circle (1.2pt);
\fill[blue!80!black] (7.8470,5.8946) circle (1.2pt);
\fill[blue!80!black] (7.8451,5.8947) circle (1.2pt);
\fill[blue!80!black] (7.8469,5.8940) circle (1.2pt);
\fill[blue!80!black] (7.8448,5.8942) circle (1.2pt);
\fill[blue!80!black] (7.8467,5.8935) circle (1.2pt);
\fill[blue!80!black] (7.8444,5.8938) circle (1.2pt);
\fill[blue!80!black] (7.8466,5.8930) circle (1.2pt);
\fill[blue!80!black] (7.8441,5.8933) circle (1.2pt);
\fill[blue!80!black] (7.8465,5.8925) circle (1.2pt);
\fill[blue!80!black] (7.8451,5.8927) circle (1.2pt);
\fill[blue!80!black] (7.8431,5.8843) circle (1.2pt);
\fill[blue!80!black] (7.8412,5.8846) circle (1.2pt);
\fill[blue!80!black] (7.8404,5.8804) circle (1.2pt);
\fill[blue!80!black] (7.8380,5.8807) circle (1.2pt);
\fill[blue!80!black] (7.8403,5.8799) circle (1.2pt);
\fill[blue!80!black] (7.8377,5.8802) circle (1.2pt);
\fill[blue!80!black] (7.8401,5.8794) circle (1.2pt);
\fill[blue!80!black] (7.8374,5.8797) circle (1.2pt);
\fill[blue!80!black] (7.8399,5.8789) circle (1.2pt);
\fill[blue!80!black] (7.8370,5.8793) circle (1.2pt);
\fill[blue!80!black] (7.8398,5.8784) circle (1.2pt);
\fill[blue!80!black] (7.8382,5.8786) circle (1.2pt);
\fill[blue!80!black] (7.8359,5.8704) circle (1.2pt);
\fill[blue!80!black] (7.8328,5.8708) circle (1.2pt);
\fill[blue!80!black] (7.8357,5.8699) circle (1.2pt);
\fill[blue!80!black] (7.8325,5.8704) circle (1.2pt);
\fill[blue!80!black] (7.8356,5.8694) circle (1.2pt);
\fill[blue!80!black] (7.8322,5.8699) circle (1.2pt);
\fill[blue!80!black] (7.8354,5.8689) circle (1.2pt);
\fill[blue!80!black] (7.8318,5.8694) circle (1.2pt);
\fill[blue!80!black] (7.8352,5.8684) circle (1.2pt);
\fill[blue!80!black] (7.8334,5.8687) circle (1.2pt);
\fill[blue!80!black] (7.8284,5.8525) circle (1.2pt);
\fill[blue!80!black] (7.8230,5.8534) circle (1.2pt);
\fill[blue!80!black] (7.8281,5.8520) circle (1.2pt);
\fill[blue!80!black] (7.8227,5.8529) circle (1.2pt);
\fill[blue!80!black] (7.8279,5.8515) circle (1.2pt);
\fill[blue!80!black] (7.8225,5.8525) circle (1.2pt);
\fill[blue!80!black] (7.8277,5.8510) circle (1.2pt);
\fill[blue!80!black] (7.8222,5.8520) circle (1.2pt);
\fill[blue!80!black] (7.8275,5.8505) circle (1.2pt);
\fill[blue!80!black] (7.8219,5.8515) circle (1.2pt);
\fill[blue!80!black] (7.8273,5.8501) circle (1.2pt);
\fill[blue!80!black] (7.8217,5.8510) circle (1.2pt);
\fill[blue!80!black] (7.8270,5.8496) circle (1.2pt);
\fill[blue!80!black] (7.8214,5.8505) circle (1.2pt);
\fill[blue!80!black] (7.8268,5.8491) circle (1.2pt);
\fill[blue!80!black] (7.8212,5.8500) circle (1.2pt);
\fill[blue!80!black] (7.8266,5.8486) circle (1.2pt);
\fill[blue!80!black] (7.8209,5.8496) circle (1.2pt);
\fill[blue!80!black] (7.8263,5.8481) circle (1.2pt);
\fill[blue!80!black] (7.8207,5.8491) circle (1.2pt);
\fill[blue!80!black] (7.8261,5.8476) circle (1.2pt);
\fill[blue!80!black] (7.8205,5.8486) circle (1.2pt);
\fill[blue!80!black] (7.8259,5.8471) circle (1.2pt);
\fill[blue!80!black] (7.8202,5.8481) circle (1.2pt);
\fill[blue!80!black] (7.8256,5.8466) circle (1.2pt);
\fill[blue!80!black] (7.8200,5.8476) circle (1.2pt);
\fill[blue!80!black] (7.8254,5.8462) circle (1.2pt);
\fill[blue!80!black] (7.8197,5.8471) circle (1.2pt);
\fill[blue!80!black] (7.8251,5.8457) circle (1.2pt);
\fill[blue!80!black] (7.8195,5.8466) circle (1.2pt);
\fill[blue!80!black] (7.8249,5.8452) circle (1.2pt);
\fill[blue!80!black] (7.8193,5.8462) circle (1.2pt);
\fill[blue!80!black] (7.8246,5.8447) circle (1.2pt);
\fill[blue!80!black] (7.8191,5.8457) circle (1.2pt);
\fill[blue!80!black] (7.8244,5.8442) circle (1.2pt);
\fill[blue!80!black] (7.8188,5.8452) circle (1.2pt);
\fill[blue!80!black] (7.8241,5.8438) circle (1.2pt);
\fill[blue!80!black] (7.8186,5.8447) circle (1.2pt);
\fill[blue!80!black] (7.8239,5.8433) circle (1.2pt);
\fill[blue!80!black] (7.8184,5.8442) circle (1.2pt);
\fill[blue!80!black] (7.8236,5.8428) circle (1.2pt);
\fill[blue!80!black] (7.8182,5.8437) circle (1.2pt);
\fill[blue!80!black] (7.8233,5.8423) circle (1.2pt);
\fill[blue!80!black] (7.8180,5.8432) circle (1.2pt);
\fill[blue!80!black] (7.8231,5.8418) circle (1.2pt);
\fill[blue!80!black] (7.8178,5.8427) circle (1.2pt);
\fill[blue!80!black] (7.8228,5.8414) circle (1.2pt);
\fill[blue!80!black] (7.8176,5.8422) circle (1.2pt);
\fill[blue!80!black] (7.8225,5.8409) circle (1.2pt);
\fill[blue!80!black] (7.8173,5.8418) circle (1.2pt);
\fill[blue!80!black] (7.8223,5.8404) circle (1.2pt);
\fill[blue!80!black] (7.8171,5.8413) circle (1.2pt);
\fill[blue!80!black] (7.8220,5.8399) circle (1.2pt);
\fill[blue!80!black] (7.8169,5.8408) circle (1.2pt);
\fill[blue!80!black] (7.8217,5.8395) circle (1.2pt);
\fill[blue!80!black] (7.8167,5.8403) circle (1.2pt);
\fill[blue!80!black] (7.8214,5.8390) circle (1.2pt);
\fill[blue!80!black] (7.8165,5.8398) circle (1.2pt);
\fill[blue!80!black] (7.8212,5.8385) circle (1.2pt);
\fill[blue!80!black] (7.8163,5.8393) circle (1.2pt);
\fill[blue!80!black] (7.8209,5.8380) circle (1.2pt);
\fill[blue!80!black] (7.8161,5.8388) circle (1.2pt);
\fill[blue!80!black] (7.8206,5.8376) circle (1.2pt);
\fill[blue!80!black] (7.8160,5.8383) circle (1.2pt);
\fill[blue!80!black] (7.8203,5.8371) circle (1.2pt);
\fill[blue!80!black] (7.8158,5.8378) circle (1.2pt);
\fill[blue!80!black] (7.8200,5.8366) circle (1.2pt);
\fill[blue!80!black] (7.8156,5.8373) circle (1.2pt);
\fill[blue!80!black] (7.8197,5.8362) circle (1.2pt);
\fill[blue!80!black] (7.8154,5.8368) circle (1.2pt);
\fill[blue!80!black] (7.8195,5.8357) circle (1.2pt);
\fill[blue!80!black] (7.8152,5.8364) circle (1.2pt);
\fill[blue!80!black] (7.8192,5.8352) circle (1.2pt);
\fill[blue!80!black] (7.8150,5.8359) circle (1.2pt);
\fill[blue!80!black] (7.8189,5.8347) circle (1.2pt);
\fill[blue!80!black] (7.8148,5.8354) circle (1.2pt);
\fill[blue!80!black] (7.8186,5.8343) circle (1.2pt);
\fill[blue!80!black] (7.8146,5.8349) circle (1.2pt);
\fill[blue!80!black] (7.8183,5.8338) circle (1.2pt);
\fill[blue!80!black] (7.8144,5.8344) circle (1.2pt);
\fill[blue!80!black] (7.8180,5.8333) circle (1.2pt);
\fill[blue!80!black] (7.8143,5.8339) circle (1.2pt);
\fill[blue!80!black] (7.8177,5.8329) circle (1.2pt);
\fill[blue!80!black] (7.8141,5.8334) circle (1.2pt);
\fill[blue!80!black] (7.8174,5.8324) circle (1.2pt);
\fill[blue!80!black] (7.8139,5.8329) circle (1.2pt);
\fill[blue!80!black] (7.8172,5.8319) circle (1.2pt);
\fill[blue!80!black] (7.8137,5.8324) circle (1.2pt);
\fill[blue!80!black] (7.8169,5.8315) circle (1.2pt);
\fill[blue!80!black] (7.8135,5.8319) circle (1.2pt);
\fill[blue!80!black] (7.8166,5.8310) circle (1.2pt);
\fill[blue!80!black] (7.8133,5.8315) circle (1.2pt);
\fill[blue!80!black] (7.8163,5.8305) circle (1.2pt);
\fill[blue!80!black] (7.8131,5.8310) circle (1.2pt);
\fill[blue!80!black] (7.8160,5.8301) circle (1.2pt);
\fill[blue!80!black] (7.8130,5.8305) circle (1.2pt);
\fill[blue!80!black] (7.8157,5.8296) circle (1.2pt);
\fill[blue!80!black] (7.8128,5.8300) circle (1.2pt);
\fill[blue!80!black] (7.8154,5.8291) circle (1.2pt);
\fill[blue!80!black] (7.8126,5.8295) circle (1.2pt);
\fill[blue!80!black] (7.8151,5.8287) circle (1.2pt);
\fill[blue!80!black] (7.8124,5.8290) circle (1.2pt);
\fill[blue!80!black] (7.8149,5.8282) circle (1.2pt);
\fill[blue!80!black] (7.8122,5.8285) circle (1.2pt);
\fill[blue!80!black] (7.8146,5.8277) circle (1.2pt);
\fill[blue!80!black] (7.8120,5.8280) circle (1.2pt);
\fill[blue!80!black] (7.8143,5.8273) circle (1.2pt);
\fill[blue!80!black] (7.8118,5.8275) circle (1.2pt);
\fill[blue!80!black] (7.8140,5.8268) circle (1.2pt);
\fill[blue!80!black] (7.8116,5.8271) circle (1.2pt);
\fill[blue!80!black] (7.8137,5.8263) circle (1.2pt);
\fill[blue!80!black] (7.8114,5.8266) circle (1.2pt);
\fill[blue!80!black] (7.8134,5.8259) circle (1.2pt);
\fill[blue!80!black] (7.8113,5.8261) circle (1.2pt);
\fill[blue!80!black] (7.8132,5.8254) circle (1.2pt);
\fill[blue!80!black] (7.8111,5.8256) circle (1.2pt);
\fill[blue!80!black] (7.8129,5.8249) circle (1.2pt);
\fill[blue!80!black] (7.8109,5.8251) circle (1.2pt);
\fill[blue!80!black] (7.8126,5.8245) circle (1.2pt);
\fill[blue!80!black] (7.8107,5.8246) circle (1.2pt);
\fill[blue!80!black] (7.8123,5.8240) circle (1.2pt);
\fill[blue!80!black] (7.8105,5.8241) circle (1.2pt);
\fill[blue!80!black] (7.8120,5.8235) circle (1.2pt);
\fill[blue!80!black] (7.8103,5.8237) circle (1.2pt);
\fill[blue!80!black] (7.8118,5.8231) circle (1.2pt);
\fill[blue!80!black] (7.8101,5.8232) circle (1.2pt);
\fill[blue!80!black] (7.8115,5.8226) circle (1.2pt);
\fill[blue!80!black] (7.8099,5.8227) circle (1.2pt);
\fill[blue!80!black] (7.8112,5.8221) circle (1.2pt);
\fill[blue!80!black] (7.8097,5.8222) circle (1.2pt);
\fill[blue!80!black] (7.8110,5.8217) circle (1.2pt);
\fill[blue!80!black] (7.8095,5.8217) circle (1.2pt);
\fill[blue!80!black] (7.8107,5.8212) circle (1.2pt);
\fill[blue!80!black] (7.8093,5.8212) circle (1.2pt);
\fill[blue!80!black] (7.8104,5.8207) circle (1.2pt);
\fill[blue!80!black] (7.8091,5.8208) circle (1.2pt);
\fill[blue!80!black] (7.8101,5.8203) circle (1.2pt);
\fill[blue!80!black] (7.8089,5.8203) circle (1.2pt);
\fill[blue!80!black] (7.8099,5.8198) circle (1.2pt);
\fill[blue!80!black] (7.8087,5.8198) circle (1.2pt);
\fill[blue!80!black] (7.8096,5.8193) circle (1.2pt);
\fill[blue!80!black] (7.8085,5.8193) circle (1.2pt);
\fill[blue!80!black] (7.8093,5.8189) circle (1.2pt);
\fill[blue!80!black] (7.8083,5.8188) circle (1.2pt);
\fill[blue!80!black] (7.8091,5.8184) circle (1.2pt);
\fill[blue!80!black] (7.8081,5.8184) circle (1.2pt);
\fill[blue!80!black] (7.8088,5.8179) circle (1.2pt);
\fill[blue!80!black] (7.8078,5.8179) circle (1.2pt);
\fill[blue!80!black] (7.8086,5.8175) circle (1.2pt);
\fill[blue!80!black] (7.8076,5.8174) circle (1.2pt);
\fill[blue!80!black] (7.8083,5.8170) circle (1.2pt);
\fill[blue!80!black] (7.8074,5.8169) circle (1.2pt);
\fill[blue!80!black] (7.8080,5.8165) circle (1.2pt);
\fill[blue!80!black] (7.8072,5.8164) circle (1.2pt);
\fill[blue!80!black] (7.8078,5.8161) circle (1.2pt);
\fill[blue!80!black] (7.8070,5.8160) circle (1.2pt);
\fill[blue!80!black] (7.8075,5.8156) circle (1.2pt);
\fill[blue!80!black] (7.8068,5.8155) circle (1.2pt);
\fill[blue!80!black] (7.8073,5.8151) circle (1.2pt);
\fill[blue!80!black] (7.8066,5.8150) circle (1.2pt);
\fill[blue!80!black] (7.8070,5.8146) circle (1.2pt);
\fill[blue!80!black] (7.8063,5.8145) circle (1.2pt);
\fill[blue!80!black] (7.8068,5.8142) circle (1.2pt);
\fill[blue!80!black] (7.8061,5.8141) circle (1.2pt);
\fill[blue!80!black] (7.8065,5.8137) circle (1.2pt);
\fill[blue!80!black] (7.8059,5.8136) circle (1.2pt);
\fill[blue!80!black] (7.8063,5.8132) circle (1.2pt);
\fill[blue!80!black] (7.8057,5.8131) circle (1.2pt);
\fill[blue!80!black] (7.8060,5.8128) circle (1.2pt);
\fill[blue!80!black] (7.8055,5.8126) circle (1.2pt);
\fill[blue!80!black] (7.8058,5.8123) circle (1.2pt);
\fill[blue!80!black] (7.8052,5.8122) circle (1.2pt);
\fill[blue!80!black] (7.8055,5.8118) circle (1.2pt);
\fill[blue!80!black] (7.8050,5.8117) circle (1.2pt);
\fill[blue!80!black] (7.8055,5.8110) circle (1.2pt);
\fill[blue!80!black] (7.8043,5.8110) circle (1.2pt);
\fill[blue!80!black] (7.8052,5.8106) circle (1.2pt);
\fill[blue!80!black] (7.8041,5.8106) circle (1.2pt);
\fill[blue!80!black] (7.8050,5.8101) circle (1.2pt);
\fill[blue!80!black] (7.8039,5.8101) circle (1.2pt);
\fill[blue!80!black] (7.8047,5.8097) circle (1.2pt);
\fill[blue!80!black] (7.8037,5.8096) circle (1.2pt);
\fill[blue!80!black] (7.8044,5.8092) circle (1.2pt);
\fill[blue!80!black] (7.8035,5.8091) circle (1.2pt);
\fill[blue!80!black] (7.8042,5.8087) circle (1.2pt);
\fill[blue!80!black] (7.8033,5.8087) circle (1.2pt);
\fill[blue!80!black] (7.8039,5.8083) circle (1.2pt);
\fill[blue!80!black] (7.8031,5.8082) circle (1.2pt);
\fill[blue!80!black] (7.8036,5.8078) circle (1.2pt);
\fill[blue!80!black] (7.8029,5.8077) circle (1.2pt);
\fill[blue!80!black] (7.8034,5.8073) circle (1.2pt);
\fill[blue!80!black] (7.8027,5.8072) circle (1.2pt);
\fill[blue!80!black] (7.8031,5.8069) circle (1.2pt);
\fill[blue!80!black] (7.8025,5.8067) circle (1.2pt);
\fill[blue!80!black] (7.8029,5.8064) circle (1.2pt);
\fill[blue!80!black] (7.8023,5.8063) circle (1.2pt);
\fill[blue!80!black] (7.8026,5.8059) circle (1.2pt);
\fill[blue!80!black] (7.8021,5.8058) circle (1.2pt);
\fill[blue!80!black] (7.8024,5.8055) circle (1.2pt);
\fill[blue!80!black] (7.8018,5.8053) circle (1.2pt);
\fill[blue!80!black] (7.8021,5.8050) circle (1.2pt);
\fill[blue!80!black] (7.8016,5.8049) circle (1.2pt);
\fill[blue!80!black] (7.8021,5.8042) circle (1.2pt);
\fill[blue!80!black] (7.8009,5.8042) circle (1.2pt);
\fill[blue!80!black] (7.8018,5.8038) circle (1.2pt);
\fill[blue!80!black] (7.8008,5.8037) circle (1.2pt);
\fill[blue!80!black] (7.8015,5.8033) circle (1.2pt);
\fill[blue!80!black] (7.8006,5.8033) circle (1.2pt);
\fill[blue!80!black] (7.8013,5.8028) circle (1.2pt);
\fill[blue!80!black] (7.8004,5.8028) circle (1.2pt);
\fill[blue!80!black] (7.8010,5.8024) circle (1.2pt);
\fill[blue!80!black] (7.8002,5.8023) circle (1.2pt);
\fill[blue!80!black] (7.8007,5.8019) circle (1.2pt);
\fill[blue!80!black] (7.8000,5.8018) circle (1.2pt);
\fill[blue!80!black] (7.8005,5.8015) circle (1.2pt);
\fill[blue!80!black] (7.7997,5.8014) circle (1.2pt);
\fill[blue!80!black] (7.8002,5.8010) circle (1.2pt);
\fill[blue!80!black] (7.7995,5.8009) circle (1.2pt);
\fill[blue!80!black] (7.8000,5.8005) circle (1.2pt);
\fill[blue!80!black] (7.7993,5.8004) circle (1.2pt);
\fill[blue!80!black] (7.7997,5.8001) circle (1.2pt);
\fill[blue!80!black] (7.7991,5.7999) circle (1.2pt);
\fill[blue!80!black] (7.7994,5.7996) circle (1.2pt);
\fill[blue!80!black] (7.7989,5.7995) circle (1.2pt);
\fill[blue!80!black] (7.7992,5.7991) circle (1.2pt);
\fill[blue!80!black] (7.7987,5.7990) circle (1.2pt);
\fill[blue!80!black] (7.7992,5.7984) circle (1.2pt);
\fill[blue!80!black] (7.7980,5.7984) circle (1.2pt);
\fill[blue!80!black] (7.7989,5.7979) circle (1.2pt);
\fill[blue!80!black] (7.7978,5.7979) circle (1.2pt);
\fill[blue!80!black] (7.7986,5.7974) circle (1.2pt);
\fill[blue!80!black] (7.7976,5.7974) circle (1.2pt);
\fill[blue!80!black] (7.7983,5.7970) circle (1.2pt);
\fill[blue!80!black] (7.7974,5.7969) circle (1.2pt);
\fill[blue!80!black] (7.7981,5.7965) circle (1.2pt);
\fill[blue!80!black] (7.7972,5.7965) circle (1.2pt);
\fill[blue!80!black] (7.7978,5.7961) circle (1.2pt);
\fill[blue!80!black] (7.7970,5.7960) circle (1.2pt);
\fill[blue!80!black] (7.7975,5.7956) circle (1.2pt);
\fill[blue!80!black] (7.7968,5.7955) circle (1.2pt);
\fill[blue!80!black] (7.7973,5.7951) circle (1.2pt);
\fill[blue!80!black] (7.7966,5.7950) circle (1.2pt);
\fill[blue!80!black] (7.7970,5.7947) circle (1.2pt);
\fill[blue!80!black] (7.7964,5.7946) circle (1.2pt);
\fill[blue!80!black] (7.7967,5.7942) circle (1.2pt);
\fill[blue!80!black] (7.7962,5.7941) circle (1.2pt);
\fill[blue!80!black] (7.7965,5.7938) circle (1.2pt);
\fill[blue!80!black] (7.7960,5.7936) circle (1.2pt);
\fill[blue!80!black] (7.7965,5.7930) circle (1.2pt);
\fill[blue!80!black] (7.7953,5.7930) circle (1.2pt);
\fill[blue!80!black] (7.7962,5.7925) circle (1.2pt);
\fill[blue!80!black] (7.7952,5.7925) circle (1.2pt);
\fill[blue!80!black] (7.7959,5.7921) circle (1.2pt);
\fill[blue!80!black] (7.7950,5.7920) circle (1.2pt);
\fill[blue!80!black] (7.7956,5.7916) circle (1.2pt);
\fill[blue!80!black] (7.7948,5.7916) circle (1.2pt);
\fill[blue!80!black] (7.7953,5.7912) circle (1.2pt);
\fill[blue!80!black] (7.7946,5.7911) circle (1.2pt);
\fill[blue!80!black] (7.7951,5.7907) circle (1.2pt);
\fill[blue!80!black] (7.7944,5.7906) circle (1.2pt);
\fill[blue!80!black] (7.7948,5.7903) circle (1.2pt);
\fill[blue!80!black] (7.7942,5.7901) circle (1.2pt);
\fill[blue!80!black] (7.7946,5.7898) circle (1.2pt);
\fill[blue!80!black] (7.7940,5.7897) circle (1.2pt);
\fill[blue!80!black] (7.7943,5.7893) circle (1.2pt);
\fill[blue!80!black] (7.7938,5.7892) circle (1.2pt);
\fill[blue!80!black] (7.7940,5.7889) circle (1.2pt);
\fill[blue!80!black] (7.7936,5.7887) circle (1.2pt);
\fill[blue!80!black] (7.7940,5.7881) circle (1.2pt);
\fill[blue!80!black] (7.7929,5.7881) circle (1.2pt);
\fill[blue!80!black] (7.7937,5.7876) circle (1.2pt);
\fill[blue!80!black] (7.7927,5.7876) circle (1.2pt);
\fill[blue!80!black] (7.7934,5.7872) circle (1.2pt);
\fill[blue!80!black] (7.7925,5.7871) circle (1.2pt);
\fill[blue!80!black] (7.7932,5.7867) circle (1.2pt);
\fill[blue!80!black] (7.7924,5.7867) circle (1.2pt);
\fill[blue!80!black] (7.7929,5.7863) circle (1.2pt);
\fill[blue!80!black] (7.7922,5.7862) circle (1.2pt);
\fill[blue!80!black] (7.7926,5.7858) circle (1.2pt);
\fill[blue!80!black] (7.7920,5.7857) circle (1.2pt);
\fill[blue!80!black] (7.7923,5.7854) circle (1.2pt);
\fill[blue!80!black] (7.7918,5.7852) circle (1.2pt);
\fill[blue!80!black] (7.7921,5.7849) circle (1.2pt);
\fill[blue!80!black] (7.7916,5.7848) circle (1.2pt);
\fill[blue!80!black] (7.7918,5.7845) circle (1.2pt);
\fill[blue!80!black] (7.7913,5.7843) circle (1.2pt);
\fill[blue!80!black] (7.7918,5.7837) circle (1.2pt);
\fill[blue!80!black] (7.7907,5.7837) circle (1.2pt);
\fill[blue!80!black] (7.7915,5.7832) circle (1.2pt);
\fill[blue!80!black] (7.7905,5.7832) circle (1.2pt);
\fill[blue!80!black] (7.7912,5.7828) circle (1.2pt);
\fill[blue!80!black] (7.7903,5.7827) circle (1.2pt);
\fill[blue!80!black] (7.7909,5.7823) circle (1.2pt);
\fill[blue!80!black] (7.7902,5.7823) circle (1.2pt);
\fill[blue!80!black] (7.7907,5.7819) circle (1.2pt);
\fill[blue!80!black] (7.7900,5.7818) circle (1.2pt);
\fill[blue!80!black] (7.7904,5.7814) circle (1.2pt);
\fill[blue!80!black] (7.7898,5.7813) circle (1.2pt);
\fill[blue!80!black] (7.7901,5.7810) circle (1.2pt);
\fill[blue!80!black] (7.7896,5.7808) circle (1.2pt);
\fill[blue!80!black] (7.7899,5.7805) circle (1.2pt);
\fill[blue!80!black] (7.7894,5.7804) circle (1.2pt);
\fill[blue!80!black] (7.7896,5.7801) circle (1.2pt);
\fill[blue!80!black] (7.7887,5.7797) circle (1.2pt);
\fill[blue!80!black] (7.7895,5.7793) circle (1.2pt);
\fill[blue!80!black] (7.7885,5.7793) circle (1.2pt);
\fill[blue!80!black] (7.7892,5.7788) circle (1.2pt);
\fill[blue!80!black] (7.7883,5.7788) circle (1.2pt);
\fill[blue!80!black] (7.7890,5.7784) circle (1.2pt);
\fill[blue!80!black] (7.7882,5.7783) circle (1.2pt);
\fill[blue!80!black] (7.7887,5.7779) circle (1.2pt);
\fill[blue!80!black] (7.7880,5.7778) circle (1.2pt);
\fill[blue!80!black] (7.7884,5.7775) circle (1.2pt);
\fill[blue!80!black] (7.7878,5.7774) circle (1.2pt);
\fill[blue!80!black] (7.7881,5.7770) circle (1.2pt);
\fill[blue!80!black] (7.7876,5.7769) circle (1.2pt);
\fill[blue!80!black] (7.7879,5.7766) circle (1.2pt);
\fill[blue!80!black] (7.7874,5.7764) circle (1.2pt);
\fill[blue!80!black] (7.7876,5.7761) circle (1.2pt);
\fill[blue!80!black] (7.7867,5.7758) circle (1.2pt);
\fill[blue!80!black] (7.7876,5.7754) circle (1.2pt);
\fill[blue!80!black] (7.7866,5.7753) circle (1.2pt);
\fill[blue!80!black] (7.7873,5.7749) circle (1.2pt);
\fill[blue!80!black] (7.7864,5.7749) circle (1.2pt);
\fill[blue!80!black] (7.7870,5.7745) circle (1.2pt);
\fill[blue!80!black] (7.7862,5.7744) circle (1.2pt);
\fill[blue!80!black] (7.7867,5.7740) circle (1.2pt);
\fill[blue!80!black] (7.7860,5.7739) circle (1.2pt);
\fill[blue!80!black] (7.7864,5.7736) circle (1.2pt);
\fill[blue!80!black] (7.7858,5.7735) circle (1.2pt);
\fill[blue!80!black] (7.7862,5.7731) circle (1.2pt);
\fill[blue!80!black] (7.7856,5.7730) circle (1.2pt);
\fill[blue!80!black] (7.7859,5.7727) circle (1.2pt);
\fill[blue!80!black] (7.7854,5.7725) circle (1.2pt);
\fill[blue!80!black] (7.7859,5.7719) circle (1.2pt);
\fill[blue!80!black] (7.7848,5.7719) circle (1.2pt);
\fill[blue!80!black] (7.7856,5.7715) circle (1.2pt);
\fill[blue!80!black] (7.7846,5.7714) circle (1.2pt);
\fill[blue!80!black] (7.7853,5.7710) circle (1.2pt);
\fill[blue!80!black] (7.7844,5.7709) circle (1.2pt);
\fill[blue!80!black] (7.7850,5.7706) circle (1.2pt);
\fill[blue!80!black] (7.7843,5.7705) circle (1.2pt);
\fill[blue!80!black] (7.7847,5.7701) circle (1.2pt);
\fill[blue!80!black] (7.7841,5.7700) circle (1.2pt);
\fill[blue!80!black] (7.7844,5.7697) circle (1.2pt);
\fill[blue!80!black] (7.7839,5.7695) circle (1.2pt);
\fill[blue!80!black] (7.7842,5.7692) circle (1.2pt);
\fill[blue!80!black] (7.7837,5.7691) circle (1.2pt);
\fill[blue!80!black] (7.7841,5.7685) circle (1.2pt);
\fill[blue!80!black] (7.7830,5.7684) circle (1.2pt);
\fill[blue!80!black] (7.7838,5.7680) circle (1.2pt);
\fill[blue!80!black] (7.7829,5.7680) circle (1.2pt);
\fill[blue!80!black] (7.7835,5.7676) circle (1.2pt);
\fill[blue!80!black] (7.7827,5.7675) circle (1.2pt);
\fill[blue!80!black] (7.7832,5.7671) circle (1.2pt);
\fill[blue!80!black] (7.7825,5.7670) circle (1.2pt);
\fill[blue!80!black] (7.7830,5.7667) circle (1.2pt);
\fill[blue!80!black] (7.7823,5.7666) circle (1.2pt);
\fill[blue!80!black] (7.7827,5.7662) circle (1.2pt);
\fill[blue!80!black] (7.7821,5.7661) circle (1.2pt);
\fill[blue!80!black] (7.7824,5.7658) circle (1.2pt);
\fill[blue!80!black] (7.7819,5.7656) circle (1.2pt);
\fill[blue!80!black] (7.7824,5.7650) circle (1.2pt);
\fill[blue!80!black] (7.7813,5.7650) circle (1.2pt);
\fill[blue!80!black] (7.7821,5.7646) circle (1.2pt);
\fill[blue!80!black] (7.7812,5.7645) circle (1.2pt);
\fill[blue!80!black] (7.7818,5.7641) circle (1.2pt);
\fill[blue!80!black] (7.7810,5.7641) circle (1.2pt);
\fill[blue!80!black] (7.7815,5.7637) circle (1.2pt);
\fill[blue!80!black] (7.7808,5.7636) circle (1.2pt);
\fill[blue!80!black] (7.7812,5.7632) circle (1.2pt);
\fill[blue!80!black] (7.7806,5.7631) circle (1.2pt);
\fill[blue!80!black] (7.7809,5.7628) circle (1.2pt);
\fill[blue!80!black] (7.7804,5.7627) circle (1.2pt);
\fill[blue!80!black] (7.7807,5.7623) circle (1.2pt);
\fill[blue!80!black] (7.7798,5.7620) circle (1.2pt);
\fill[blue!80!black] (7.7806,5.7616) circle (1.2pt);
\fill[blue!80!black] (7.7796,5.7616) circle (1.2pt);
\fill[blue!80!black] (7.7803,5.7611) circle (1.2pt);
\fill[blue!80!black] (7.7794,5.7611) circle (1.2pt);
\fill[blue!80!black] (7.7800,5.7607) circle (1.2pt);
\fill[blue!80!black] (7.7793,5.7606) circle (1.2pt);
\fill[blue!80!black] (7.7797,5.7603) circle (1.2pt);
\fill[blue!80!black] (7.7791,5.7601) circle (1.2pt);
\fill[blue!80!black] (7.7795,5.7598) circle (1.2pt);
\fill[blue!80!black] (7.7789,5.7597) circle (1.2pt);
\fill[blue!80!black] (7.7792,5.7594) circle (1.2pt);
\fill[blue!80!black] (7.7787,5.7592) circle (1.2pt);
\fill[blue!80!black] (7.7791,5.7586) circle (1.2pt);
\fill[blue!80!black] (7.7781,5.7586) circle (1.2pt);
\fill[blue!80!black] (7.7788,5.7582) circle (1.2pt);
\fill[blue!80!black] (7.7779,5.7581) circle (1.2pt);
\fill[blue!80!black] (7.7785,5.7577) circle (1.2pt);
\fill[blue!80!black] (7.7778,5.7577) circle (1.2pt);
\fill[blue!80!black] (7.7782,5.7573) circle (1.2pt);
\fill[blue!80!black] (7.7776,5.7572) circle (1.2pt);
\fill[blue!80!black] (7.7780,5.7568) circle (1.2pt);
\fill[blue!80!black] (7.7774,5.7567) circle (1.2pt);
\fill[blue!80!black] (7.7777,5.7564) circle (1.2pt);
\fill[blue!80!black] (7.7772,5.7563) circle (1.2pt);
\fill[blue!80!black] (7.7777,5.7556) circle (1.2pt);
\fill[blue!80!black] (7.7766,5.7556) circle (1.2pt);
\fill[blue!80!black] (7.7773,5.7552) circle (1.2pt);
\fill[blue!80!black] (7.7764,5.7552) circle (1.2pt);
\fill[blue!80!black] (7.7770,5.7548) circle (1.2pt);
\fill[blue!80!black] (7.7763,5.7547) circle (1.2pt);
\fill[blue!80!black] (7.7767,5.7543) circle (1.2pt);
\fill[blue!80!black] (7.7761,5.7542) circle (1.2pt);
\fill[blue!80!black] (7.7765,5.7539) circle (1.2pt);
\fill[blue!80!black] (7.7759,5.7538) circle (1.2pt);
\fill[blue!80!black] (7.7762,5.7534) circle (1.2pt);
\fill[blue!80!black] (7.7757,5.7533) circle (1.2pt);
\fill[blue!80!black] (7.7761,5.7527) circle (1.2pt);
\fill[blue!80!black] (7.7751,5.7527) circle (1.2pt);
\fill[blue!80!black] (7.7758,5.7523) circle (1.2pt);
\fill[blue!80!black] (7.7749,5.7522) circle (1.2pt);
\fill[blue!80!black] (7.7755,5.7518) circle (1.2pt);
\fill[blue!80!black] (7.7748,5.7517) circle (1.2pt);
\fill[blue!80!black] (7.7752,5.7514) circle (1.2pt);
\fill[blue!80!black] (7.7746,5.7513) circle (1.2pt);
\fill[blue!80!black] (7.7750,5.7509) circle (1.2pt);
\fill[blue!80!black] (7.7744,5.7508) circle (1.2pt);
\fill[blue!80!black] (7.7747,5.7505) circle (1.2pt);
\fill[blue!80!black] (7.7738,5.7502) circle (1.2pt);
\fill[blue!80!black] (7.7746,5.7497) circle (1.2pt);
\fill[blue!80!black] (7.7736,5.7497) circle (1.2pt);
\fill[blue!80!black] (7.7743,5.7493) circle (1.2pt);
\fill[blue!80!black] (7.7735,5.7492) circle (1.2pt);
\fill[blue!80!black] (7.7740,5.7489) circle (1.2pt);
\fill[blue!80!black] (7.7733,5.7488) circle (1.2pt);
\fill[blue!80!black] (7.7737,5.7484) circle (1.2pt);
\fill[blue!80!black] (7.7731,5.7483) circle (1.2pt);
\fill[blue!80!black] (7.7734,5.7480) circle (1.2pt);
\fill[blue!80!black] (7.7729,5.7478) circle (1.2pt);
\fill[blue!80!black] (7.7734,5.7472) circle (1.2pt);
\fill[blue!80!black] (7.7723,5.7472) circle (1.2pt);
\fill[blue!80!black] (7.7731,5.7468) circle (1.2pt);
\fill[blue!80!black] (7.7722,5.7468) circle (1.2pt);
\fill[blue!80!black] (7.7728,5.7464) circle (1.2pt);
\fill[blue!80!black] (7.7720,5.7463) circle (1.2pt);
\fill[blue!80!black] (7.7725,5.7459) circle (1.2pt);
\fill[blue!80!black] (7.7718,5.7458) circle (1.2pt);
\fill[blue!80!black] (7.7722,5.7455) circle (1.2pt);
\fill[blue!80!black] (7.7717,5.7454) circle (1.2pt);
\fill[blue!80!black] (7.7719,5.7450) circle (1.2pt);
\fill[blue!80!black] (7.7710,5.7448) circle (1.2pt);
\fill[blue!80!black] (7.7719,5.7443) circle (1.2pt);
\fill[blue!80!black] (7.7709,5.7443) circle (1.2pt);
\fill[blue!80!black] (7.7715,5.7439) circle (1.2pt);
\fill[blue!80!black] (7.7707,5.7438) circle (1.2pt);
\fill[blue!80!black] (7.7712,5.7434) circle (1.2pt);
\fill[blue!80!black] (7.7706,5.7433) circle (1.2pt);
\fill[blue!80!black] (7.7710,5.7430) circle (1.2pt);
\fill[blue!80!black] (7.7704,5.7429) circle (1.2pt);
\fill[blue!80!black] (7.7707,5.7426) circle (1.2pt);
\fill[blue!80!black] (7.7702,5.7424) circle (1.2pt);
\fill[blue!80!black] (7.7706,5.7418) circle (1.2pt);
\fill[blue!80!black] (7.7696,5.7418) circle (1.2pt);
\fill[blue!80!black] (7.7703,5.7414) circle (1.2pt);
\fill[blue!80!black] (7.7694,5.7413) circle (1.2pt);
\fill[blue!80!black] (7.7700,5.7409) circle (1.2pt);
\fill[blue!80!black] (7.7693,5.7409) circle (1.2pt);
\fill[blue!80!black] (7.7697,5.7405) circle (1.2pt);
\fill[blue!80!black] (7.7691,5.7404) circle (1.2pt);
\fill[blue!80!black] (7.7694,5.7401) circle (1.2pt);
\fill[blue!80!black] (7.7689,5.7399) circle (1.2pt);
\fill[blue!80!black] (7.7694,5.7393) circle (1.2pt);
\fill[blue!80!black] (7.7683,5.7393) circle (1.2pt);
\fill[blue!80!black] (7.7691,5.7389) circle (1.2pt);
\fill[blue!80!black] (7.7682,5.7388) circle (1.2pt);
\fill[blue!80!black] (7.7687,5.7385) circle (1.2pt);
\fill[blue!80!black] (7.7680,5.7384) circle (1.2pt);
\fill[blue!80!black] (7.7684,5.7380) circle (1.2pt);
\fill[blue!80!black] (7.7678,5.7379) circle (1.2pt);
\fill[blue!80!black] (7.7682,5.7376) circle (1.2pt);
\fill[blue!80!black] (7.7677,5.7375) circle (1.2pt);
\fill[blue!80!black] (7.7681,5.7368) circle (1.2pt);
\fill[blue!80!black] (7.7670,5.7368) circle (1.2pt);
\fill[blue!80!black] (7.7678,5.7364) circle (1.2pt);
\fill[blue!80!black] (7.7669,5.7364) circle (1.2pt);
\fill[blue!80!black] (7.7675,5.7360) circle (1.2pt);
\fill[blue!80!black] (7.7668,5.7359) circle (1.2pt);
\fill[blue!80!black] (7.7672,5.7356) circle (1.2pt);
\fill[blue!80!black] (7.7666,5.7354) circle (1.2pt);
\fill[blue!80!black] (7.7669,5.7351) circle (1.2pt);
\fill[blue!80!black] (7.7664,5.7350) circle (1.2pt);
\fill[blue!80!black] (7.7669,5.7344) circle (1.2pt);
\fill[blue!80!black] (7.7658,5.7344) circle (1.2pt);
\fill[blue!80!black] (7.7665,5.7340) circle (1.2pt);
\fill[blue!80!black] (7.7657,5.7339) circle (1.2pt);
\fill[blue!80!black] (7.7662,5.7335) circle (1.2pt);
\fill[blue!80!black] (7.7655,5.7334) circle (1.2pt);
\fill[blue!80!black] (7.7659,5.7331) circle (1.2pt);
\fill[blue!80!black] (7.7653,5.7330) circle (1.2pt);
\fill[blue!80!black] (7.7656,5.7327) circle (1.2pt);
\fill[blue!80!black] (7.7652,5.7325) circle (1.2pt);
\fill[blue!80!black] (7.7656,5.7319) circle (1.2pt);
\fill[blue!80!black] (7.7646,5.7319) circle (1.2pt);
\fill[blue!80!black] (7.7653,5.7315) circle (1.2pt);
\fill[blue!80!black] (7.7644,5.7314) circle (1.2pt);
\fill[blue!80!black] (7.7649,5.7311) circle (1.2pt);
\fill[blue!80!black] (7.7643,5.7310) circle (1.2pt);
\fill[blue!80!black] (7.7646,5.7306) circle (1.2pt);
\fill[blue!80!black] (7.7641,5.7305) circle (1.2pt);
\fill[blue!80!black] (7.7644,5.7302) circle (1.2pt);
\fill[blue!80!black] (7.7634,5.7299) circle (1.2pt);
\fill[blue!80!black] (7.7643,5.7295) circle (1.2pt);
\fill[blue!80!black] (7.7633,5.7294) circle (1.2pt);
\fill[blue!80!black] (7.7640,5.7290) circle (1.2pt);
\fill[blue!80!black] (7.7632,5.7290) circle (1.2pt);
\fill[blue!80!black] (7.7637,5.7286) circle (1.2pt);
\fill[blue!80!black] (7.7630,5.7285) circle (1.2pt);
\fill[blue!80!black] (7.7634,5.7282) circle (1.2pt);
\fill[blue!80!black] (7.7629,5.7280) circle (1.2pt);
\fill[blue!80!black] (7.7631,5.7277) circle (1.2pt);
\fill[blue!80!black] (7.7622,5.7274) circle (1.2pt);
\fill[blue!80!black] (7.7630,5.7270) circle (1.2pt);
\fill[blue!80!black] (7.7621,5.7270) circle (1.2pt);
\fill[blue!80!black] (7.7627,5.7266) circle (1.2pt);
\fill[blue!80!black] (7.7620,5.7265) circle (1.2pt);
\fill[blue!80!black] (7.7624,5.7261) circle (1.2pt);
\fill[blue!80!black] (7.7618,5.7260) circle (1.2pt);
\fill[blue!80!black] (7.7621,5.7257) circle (1.2pt);
\fill[blue!80!black] (7.7616,5.7256) circle (1.2pt);
\fill[blue!80!black] (7.7620,5.7250) circle (1.2pt);
\fill[blue!80!black] (7.7610,5.7250) circle (1.2pt);
\fill[blue!80!black] (7.7617,5.7246) circle (1.2pt);
\fill[blue!80!black] (7.7609,5.7245) circle (1.2pt);
\fill[blue!80!black] (7.7614,5.7241) circle (1.2pt);
\fill[blue!80!black] (7.7607,5.7240) circle (1.2pt);
\fill[blue!80!black] (7.7611,5.7237) circle (1.2pt);
\fill[blue!80!black] (7.7606,5.7236) circle (1.2pt);
\fill[blue!80!black] (7.7608,5.7233) circle (1.2pt);
\fill[blue!80!black] (7.7600,5.7230) circle (1.2pt);
\fill[blue!80!black] (7.7607,5.7225) circle (1.2pt);
\fill[blue!80!black] (7.7598,5.7225) circle (1.2pt);
\fill[blue!80!black] (7.7604,5.7221) circle (1.2pt);
\fill[blue!80!black] (7.7597,5.7220) circle (1.2pt);
\fill[blue!80!black] (7.7601,5.7217) circle (1.2pt);
\fill[blue!80!black] (7.7595,5.7216) circle (1.2pt);
\fill[blue!80!black] (7.7598,5.7213) circle (1.2pt);
\fill[blue!80!black] (7.7589,5.7210) circle (1.2pt);
\fill[blue!80!black] (7.7597,5.7205) circle (1.2pt);
\fill[blue!80!black] (7.7588,5.7205) circle (1.2pt);
\fill[blue!80!black] (7.7594,5.7201) circle (1.2pt);
\fill[blue!80!black] (7.7586,5.7200) circle (1.2pt);
\fill[blue!80!black] (7.7591,5.7197) circle (1.2pt);
\fill[blue!80!black] (7.7585,5.7196) circle (1.2pt);
\fill[blue!80!black] (7.7588,5.7192) circle (1.2pt);
\fill[blue!80!black] (7.7583,5.7191) circle (1.2pt);
\fill[blue!80!black] (7.7587,5.7185) circle (1.2pt);
\fill[blue!80!black] (7.7577,5.7185) circle (1.2pt);
\fill[blue!80!black] (7.7584,5.7181) circle (1.2pt);
\fill[blue!80!black] (7.7576,5.7180) circle (1.2pt);
\fill[blue!80!black] (7.7581,5.7177) circle (1.2pt);
\fill[blue!80!black] (7.7574,5.7176) circle (1.2pt);
\fill[blue!80!black] (7.7578,5.7172) circle (1.2pt);
\fill[blue!80!black] (7.7573,5.7171) circle (1.2pt);
\fill[blue!80!black] (7.7577,5.7165) circle (1.2pt);
\fill[blue!80!black] (7.7567,5.7165) circle (1.2pt);
\fill[blue!80!black] (7.7574,5.7161) circle (1.2pt);
\fill[blue!80!black] (7.7566,5.7160) circle (1.2pt);
\fill[blue!80!black] (7.7571,5.7157) circle (1.2pt);
\fill[blue!80!black] (7.7564,5.7156) circle (1.2pt);
\fill[blue!80!black] (7.7568,5.7152) circle (1.2pt);
\fill[blue!80!black] (7.7562,5.7151) circle (1.2pt);
\fill[blue!80!black] (7.7565,5.7148) circle (1.2pt);
\fill[blue!80!black] (7.7556,5.7145) circle (1.2pt);
\fill[blue!80!black] (7.7564,5.7141) circle (1.2pt);
\fill[blue!80!black] (7.7555,5.7140) circle (1.2pt);
\fill[blue!80!black] (7.7561,5.7137) circle (1.2pt);
\fill[blue!80!black] (7.7554,5.7136) circle (1.2pt);
\fill[blue!80!black] (7.7558,5.7133) circle (1.2pt);
\fill[blue!80!black] (7.7552,5.7131) circle (1.2pt);
\fill[blue!80!black] (7.7555,5.7128) circle (1.2pt);
\fill[blue!80!black] (7.7546,5.7125) circle (1.2pt);
\fill[blue!80!black] (7.7554,5.7121) circle (1.2pt);
\fill[blue!80!black] (7.7545,5.7121) circle (1.2pt);
\fill[blue!80!black] (7.7551,5.7117) circle (1.2pt);
\fill[blue!80!black] (7.7544,5.7116) circle (1.2pt);
\fill[blue!80!black] (7.7547,5.7113) circle (1.2pt);
\fill[blue!80!black] (7.7542,5.7111) circle (1.2pt);
\fill[blue!80!black] (7.7545,5.7108) circle (1.2pt);
\fill[blue!80!black] (7.7536,5.7105) circle (1.2pt);
\fill[blue!80!black] (7.7544,5.7101) circle (1.2pt);
\fill[blue!80!black] (7.7535,5.7101) circle (1.2pt);
\fill[blue!80!black] (7.7540,5.7097) circle (1.2pt);
\fill[blue!80!black] (7.7533,5.7096) circle (1.2pt);
\fill[blue!80!black] (7.7537,5.7093) circle (1.2pt);
\fill[blue!80!black] (7.7532,5.7091) circle (1.2pt);
\fill[blue!80!black] (7.7534,5.7088) circle (1.2pt);
\fill[blue!80!black] (7.7526,5.7085) circle (1.2pt);
\fill[blue!80!black] (7.7533,5.7081) circle (1.2pt);
\fill[blue!80!black] (7.7525,5.7081) circle (1.2pt);
\fill[blue!80!black] (7.7530,5.7077) circle (1.2pt);
\fill[blue!80!black] (7.7523,5.7076) circle (1.2pt);
\fill[blue!80!black] (7.7527,5.7073) circle (1.2pt);
\fill[blue!80!black] (7.7522,5.7072) circle (1.2pt);
\fill[blue!80!black] (7.7524,5.7069) circle (1.2pt);
\fill[blue!80!black] (7.7516,5.7066) circle (1.2pt);
\fill[blue!80!black] (7.7523,5.7061) circle (1.2pt);
\fill[blue!80!black] (7.7515,5.7061) circle (1.2pt);
\fill[blue!80!black] (7.7520,5.7057) circle (1.2pt);
\fill[blue!80!black] (7.7513,5.7056) circle (1.2pt);
\fill[blue!80!black] (7.7517,5.7053) circle (1.2pt);
\fill[blue!80!black] (7.7512,5.7052) circle (1.2pt);
\fill[blue!80!black] (7.7516,5.7046) circle (1.2pt);
\fill[blue!80!black] (7.7506,5.7046) circle (1.2pt);
\fill[blue!80!black] (7.7513,5.7042) circle (1.2pt);
\fill[blue!80!black] (7.7505,5.7041) circle (1.2pt);
\fill[blue!80!black] (7.7509,5.7037) circle (1.2pt);
\fill[blue!80!black] (7.7503,5.7036) circle (1.2pt);
\fill[blue!80!black] (7.7506,5.7033) circle (1.2pt);
\fill[blue!80!black] (7.7502,5.7032) circle (1.2pt);
\fill[blue!80!black] (7.7506,5.7026) circle (1.2pt);
\fill[blue!80!black] (7.7496,5.7026) circle (1.2pt);
\fill[blue!80!black] (7.7502,5.7022) circle (1.2pt);
\fill[blue!80!black] (7.7495,5.7021) circle (1.2pt);
\fill[blue!80!black] (7.7499,5.7018) circle (1.2pt);
\fill[blue!80!black] (7.7493,5.7017) circle (1.2pt);
\fill[blue!80!black] (7.7496,5.7014) circle (1.2pt);
\fill[blue!80!black] (7.7487,5.7011) circle (1.2pt);
\fill[blue!80!black] (7.7495,5.7006) circle (1.2pt);
\fill[blue!80!black] (7.7486,5.7006) circle (1.2pt);
\fill[blue!80!black] (7.7492,5.7002) circle (1.2pt);
\fill[blue!80!black] (7.7485,5.7001) circle (1.2pt);
\fill[blue!80!black] (7.7489,5.6998) circle (1.2pt);
\fill[blue!80!black] (7.7483,5.6997) circle (1.2pt);
\fill[blue!80!black] (7.7486,5.6994) circle (1.2pt);
\fill[blue!80!black] (7.7477,5.6991) circle (1.2pt);
\fill[blue!80!black] (7.7485,5.6987) circle (1.2pt);
\fill[blue!80!black] (7.7476,5.6986) circle (1.2pt);
\fill[blue!80!black] (7.7481,5.6983) circle (1.2pt);
\fill[blue!80!black] (7.7475,5.6982) circle (1.2pt);
\fill[blue!80!black] (7.7478,5.6978) circle (1.2pt);
\fill[blue!80!black] (7.7473,5.6977) circle (1.2pt);
\fill[blue!80!black] (7.7478,5.6971) circle (1.2pt);
\fill[blue!80!black] (7.7467,5.6971) circle (1.2pt);
\fill[blue!80!black] (7.7474,5.6967) circle (1.2pt);
\fill[blue!80!black] (7.7466,5.6966) circle (1.2pt);
\fill[blue!80!black] (7.7471,5.6963) circle (1.2pt);
\fill[blue!80!black] (7.7465,5.6962) circle (1.2pt);
\fill[blue!80!black] (7.7468,5.6959) circle (1.2pt);
\fill[blue!80!black] (7.7459,5.6956) circle (1.2pt);
\fill[blue!80!black] (7.7467,5.6952) circle (1.2pt);
\fill[blue!80!black] (7.7458,5.6951) circle (1.2pt);
\fill[blue!80!black] (7.7464,5.6948) circle (1.2pt);
\fill[blue!80!black] (7.7457,5.6947) circle (1.2pt);
\fill[blue!80!black] (7.7460,5.6943) circle (1.2pt);
\fill[blue!80!black] (7.7455,5.6942) circle (1.2pt);
\fill[blue!80!black] (7.7458,5.6939) circle (1.2pt);
\fill[blue!80!black] (7.7449,5.6936) circle (1.2pt);
\fill[blue!80!black] (7.7456,5.6932) circle (1.2pt);
\fill[blue!80!black] (7.7448,5.6931) circle (1.2pt);
\fill[blue!80!black] (7.7453,5.6928) circle (1.2pt);
\fill[blue!80!black] (7.7447,5.6927) circle (1.2pt);
\fill[blue!80!black] (7.7450,5.6924) circle (1.2pt);
\fill[blue!80!black] (7.7445,5.6922) circle (1.2pt);
\fill[blue!80!black] (7.7449,5.6917) circle (1.2pt);
\fill[blue!80!black] (7.7440,5.6916) circle (1.2pt);
\fill[blue!80!black] (7.7446,5.6913) circle (1.2pt);
\fill[blue!80!black] (7.7439,5.6912) circle (1.2pt);
\fill[blue!80!black] (7.7443,5.6908) circle (1.2pt);
\fill[blue!80!black] (7.7437,5.6907) circle (1.2pt);
\fill[blue!80!black] (7.7440,5.6904) circle (1.2pt);
\fill[blue!80!black] (7.7431,5.6901) circle (1.2pt);
\fill[blue!80!black] (7.7438,5.6897) circle (1.2pt);
\fill[blue!80!black] (7.7430,5.6897) circle (1.2pt);
\fill[blue!80!black] (7.7435,5.6893) circle (1.2pt);
\fill[blue!80!black] (7.7429,5.6892) circle (1.2pt);
\fill[blue!80!black] (7.7432,5.6889) circle (1.2pt);
\fill[blue!80!black] (7.7427,5.6888) circle (1.2pt);
\fill[blue!80!black] (7.7431,5.6882) circle (1.2pt);
\fill[blue!80!black] (7.7422,5.6881) circle (1.2pt);
\fill[blue!80!black] (7.7428,5.6878) circle (1.2pt);
\fill[blue!80!black] (7.7421,5.6877) circle (1.2pt);
\fill[blue!80!black] (7.7424,5.6874) circle (1.2pt);
\fill[blue!80!black] (7.7419,5.6872) circle (1.2pt);
\fill[blue!80!black] (7.7424,5.6866) circle (1.2pt);
\fill[blue!80!black] (7.7413,5.6866) circle (1.2pt);
\fill[blue!80!black] (7.7420,5.6862) circle (1.2pt);
\fill[blue!80!black] (7.7412,5.6862) circle (1.2pt);
\fill[blue!80!black] (7.7417,5.6858) circle (1.2pt);
\fill[blue!80!black] (7.7411,5.6857) circle (1.2pt);
\fill[blue!80!black] (7.7414,5.6854) circle (1.2pt);
\fill[blue!80!black] (7.7405,5.6851) circle (1.2pt);
\fill[blue!80!black] (7.7413,5.6847) circle (1.2pt);
\fill[blue!80!black] (7.7404,5.6847) circle (1.2pt);
\fill[blue!80!black] (7.7409,5.6843) circle (1.2pt);
\fill[blue!80!black] (7.7403,5.6842) circle (1.2pt);
\fill[blue!80!black] (7.7406,5.6839) circle (1.2pt);
\fill[blue!80!black] (7.7401,5.6838) circle (1.2pt);
\fill[blue!80!black] (7.7405,5.6832) circle (1.2pt);
\fill[blue!80!black] (7.7396,5.6832) circle (1.2pt);
\fill[blue!80!black] (7.7402,5.6828) circle (1.2pt);
\fill[blue!80!black] (7.7395,5.6827) circle (1.2pt);
\fill[blue!80!black] (7.7399,5.6824) circle (1.2pt);
\fill[blue!80!black] (7.7393,5.6822) circle (1.2pt);
\fill[blue!80!black] (7.7398,5.6817) circle (1.2pt);
\fill[blue!80!black] (7.7388,5.6817) circle (1.2pt);
\fill[blue!80!black] (7.7394,5.6813) circle (1.2pt);
\fill[blue!80!black] (7.7387,5.6812) circle (1.2pt);
\fill[blue!80!black] (7.7391,5.6808) circle (1.2pt);
\fill[blue!80!black] (7.7385,5.6807) circle (1.2pt);
\fill[blue!80!black] (7.7388,5.6804) circle (1.2pt);
\fill[blue!80!black] (7.7379,5.6801) circle (1.2pt);
\fill[blue!80!black] (7.7387,5.6797) circle (1.2pt);
\fill[blue!80!black] (7.7379,5.6797) circle (1.2pt);
\fill[blue!80!black] (7.7384,5.6793) circle (1.2pt);
\fill[blue!80!black] (7.7377,5.6792) circle (1.2pt);
\fill[blue!80!black] (7.7380,5.6789) circle (1.2pt);
\fill[blue!80!black] (7.7371,5.6786) circle (1.2pt);
\fill[blue!80!black] (7.7379,5.6782) circle (1.2pt);
\fill[blue!80!black] (7.7370,5.6782) circle (1.2pt);
\fill[blue!80!black] (7.7376,5.6778) circle (1.2pt);
\fill[blue!80!black] (7.7369,5.6777) circle (1.2pt);
\fill[blue!80!black] (7.7373,5.6774) circle (1.2pt);
\fill[blue!80!black] (7.7368,5.6773) circle (1.2pt);
\fill[blue!80!black] (7.7372,5.6767) circle (1.2pt);
\fill[blue!80!black] (7.7362,5.6767) circle (1.2pt);
\fill[blue!80!black] (7.7368,5.6763) circle (1.2pt);
\fill[blue!80!black] (7.7361,5.6762) circle (1.2pt);
\fill[blue!80!black] (7.7365,5.6759) circle (1.2pt);
\fill[blue!80!black] (7.7360,5.6758) circle (1.2pt);
\fill[blue!80!black] (7.7364,5.6752) circle (1.2pt);
\fill[blue!80!black] (7.7354,5.6752) circle (1.2pt);
\fill[blue!80!black] (7.7361,5.6748) circle (1.2pt);
\fill[blue!80!black] (7.7353,5.6747) circle (1.2pt);
\fill[blue!80!black] (7.7357,5.6744) circle (1.2pt);
\fill[blue!80!black] (7.7352,5.6743) circle (1.2pt);
\fill[blue!80!black] (7.7354,5.6740) circle (1.2pt);
\fill[blue!80!black] (7.7346,5.6737) circle (1.2pt);
\fill[blue!80!black] (7.7353,5.6733) circle (1.2pt);
\fill[blue!80!black] (7.7345,5.6732) circle (1.2pt);
\fill[blue!80!black] (7.7350,5.6729) circle (1.2pt);
\fill[blue!80!black] (7.7344,5.6728) circle (1.2pt);
\fill[blue!80!black] (7.7347,5.6725) circle (1.2pt);
\fill[blue!80!black] (7.7338,5.6722) circle (1.2pt);
\fill[blue!80!black] (7.7345,5.6718) circle (1.2pt);
\fill[blue!80!black] (7.7338,5.6717) circle (1.2pt);
\fill[blue!80!black] (7.7342,5.6714) circle (1.2pt);
\fill[blue!80!black] (7.7336,5.6712) circle (1.2pt);
\fill[blue!80!black] (7.7339,5.6709) circle (1.2pt);
\fill[blue!80!black] (7.7330,5.6707) circle (1.2pt);
\fill[blue!80!black] (7.7338,5.6703) circle (1.2pt);
\fill[blue!80!black] (7.7330,5.6702) circle (1.2pt);
\fill[blue!80!black] (7.7334,5.6699) circle (1.2pt);
\fill[blue!80!black] (7.7328,5.6697) circle (1.2pt);
\fill[blue!80!black] (7.7331,5.6694) circle (1.2pt);
\fill[blue!80!black] (7.7322,5.6692) circle (1.2pt);
\fill[blue!80!black] (7.7330,5.6687) circle (1.2pt);
\fill[blue!80!black] (7.7322,5.6687) circle (1.2pt);
\fill[blue!80!black] (7.7327,5.6683) circle (1.2pt);
\fill[blue!80!black] (7.7321,5.6682) circle (1.2pt);
\fill[blue!80!black] (7.7323,5.6679) circle (1.2pt);
\fill[blue!80!black] (7.7315,5.6677) circle (1.2pt);
\fill[blue!80!black] (7.7322,5.6672) circle (1.2pt);
\fill[blue!80!black] (7.7314,5.6672) circle (1.2pt);
\fill[blue!80!black] (7.7319,5.6668) circle (1.2pt);
\fill[blue!80!black] (7.7313,5.6667) circle (1.2pt);
\fill[blue!80!black] (7.7316,5.6664) circle (1.2pt);
\fill[blue!80!black] (7.7307,5.6662) circle (1.2pt);
\fill[blue!80!black] (7.7314,5.6657) circle (1.2pt);
\fill[blue!80!black] (7.7306,5.6657) circle (1.2pt);
\fill[blue!80!black] (7.7311,5.6653) circle (1.2pt);
\fill[blue!80!black] (7.7305,5.6652) circle (1.2pt);
\fill[blue!80!black] (7.7308,5.6649) circle (1.2pt);
\fill[blue!80!black] (7.7299,5.6647) circle (1.2pt);
\fill[blue!80!black] (7.7307,5.6643) circle (1.2pt);
\fill[blue!80!black] (7.7298,5.6642) circle (1.2pt);
\fill[blue!80!black] (7.7303,5.6639) circle (1.2pt);
\fill[blue!80!black] (7.7297,5.6637) circle (1.2pt);
\fill[blue!80!black] (7.7300,5.6634) circle (1.2pt);
\fill[blue!80!black] (7.7291,5.6632) circle (1.2pt);
\fill[blue!80!black] (7.7299,5.6628) circle (1.2pt);
\fill[blue!80!black] (7.7291,5.6627) circle (1.2pt);
\fill[blue!80!black] (7.7295,5.6624) circle (1.2pt);
\fill[blue!80!black] (7.7290,5.6623) circle (1.2pt);
\fill[blue!80!black] (7.7292,5.6620) circle (1.2pt);
\fill[blue!80!black] (7.7284,5.6617) circle (1.2pt);
\fill[blue!80!black] (7.7291,5.6613) circle (1.2pt);
\fill[blue!80!black] (7.7283,5.6612) circle (1.2pt);
\fill[blue!80!black] (7.7287,5.6609) circle (1.2pt);
\fill[blue!80!black] (7.7282,5.6608) circle (1.2pt);
\fill[blue!80!black] (7.7284,5.6605) circle (1.2pt);
\fill[blue!80!black] (7.7276,5.6602) circle (1.2pt);
\fill[blue!80!black] (7.7283,5.6598) circle (1.2pt);
\fill[blue!80!black] (7.7275,5.6597) circle (1.2pt);
\fill[blue!80!black] (7.7280,5.6594) circle (1.2pt);
\fill[blue!80!black] (7.7274,5.6593) circle (1.2pt);
\fill[blue!80!black] (7.7279,5.6587) circle (1.2pt);
\fill[blue!80!black] (7.7269,5.6587) circle (1.2pt);
\fill[blue!80!black] (7.7275,5.6583) circle (1.2pt);
\fill[blue!80!black] (7.7268,5.6582) circle (1.2pt);
\fill[blue!80!black] (7.7272,5.6579) circle (1.2pt);
\fill[blue!80!black] (7.7267,5.6578) circle (1.2pt);
\fill[blue!80!black] (7.7271,5.6572) circle (1.2pt);
\fill[blue!80!black] (7.7261,5.6572) circle (1.2pt);
\fill[blue!80!black] (7.7267,5.6568) circle (1.2pt);
\fill[blue!80!black] (7.7260,5.6567) circle (1.2pt);
\fill[blue!80!black] (7.7264,5.6564) circle (1.2pt);
\fill[blue!80!black] (7.7259,5.6563) circle (1.2pt);
\fill[blue!80!black] (7.7263,5.6557) circle (1.2pt);
\fill[blue!80!black] (7.7254,5.6557) circle (1.2pt);
\fill[blue!80!black] (7.7259,5.6553) circle (1.2pt);
\fill[blue!80!black] (7.7253,5.6552) circle (1.2pt);
\fill[blue!80!black] (7.7256,5.6549) circle (1.2pt);
\fill[blue!80!black] (7.7251,5.6548) circle (1.2pt);
\fill[blue!80!black] (7.7255,5.6542) circle (1.2pt);
\fill[blue!80!black] (7.7246,5.6542) circle (1.2pt);
\fill[blue!80!black] (7.7251,5.6538) circle (1.2pt);
\fill[blue!80!black] (7.7245,5.6537) circle (1.2pt);
\fill[blue!80!black] (7.7248,5.6534) circle (1.2pt);
\fill[blue!80!black] (7.7239,5.6532) circle (1.2pt);
\fill[blue!80!black] (7.7247,5.6528) circle (1.2pt);
\fill[blue!80!black] (7.7239,5.6527) circle (1.2pt);
\fill[blue!80!black] (7.7243,5.6524) circle (1.2pt);
\fill[blue!80!black] (7.7238,5.6523) circle (1.2pt);
\fill[blue!80!black] (7.7240,5.6520) circle (1.2pt);
\fill[blue!80!black] (7.7232,5.6517) circle (1.2pt);
\fill[blue!80!black] (7.7239,5.6513) circle (1.2pt);
\fill[blue!80!black] (7.7231,5.6512) circle (1.2pt);
\fill[blue!80!black] (7.7235,5.6509) circle (1.2pt);
\fill[blue!80!black] (7.7230,5.6508) circle (1.2pt);
\fill[blue!80!black] (7.7234,5.6502) circle (1.2pt);
\fill[blue!80!black] (7.7225,5.6502) circle (1.2pt);
\fill[blue!80!black] (7.7231,5.6498) circle (1.2pt);
\fill[blue!80!black] (7.7224,5.6497) circle (1.2pt);
\fill[blue!80!black] (7.7227,5.6494) circle (1.2pt);
\fill[blue!80!black] (7.7222,5.6493) circle (1.2pt);
\fill[blue!80!black] (7.7226,5.6487) circle (1.2pt);
\fill[blue!80!black] (7.7217,5.6487) circle (1.2pt);
\fill[blue!80!black] (7.7222,5.6483) circle (1.2pt);
\fill[blue!80!black] (7.7216,5.6482) circle (1.2pt);
\fill[blue!80!black] (7.7219,5.6479) circle (1.2pt);
\fill[blue!80!black] (7.7211,5.6477) circle (1.2pt);
\fill[blue!80!black] (7.7218,5.6473) circle (1.2pt);
\fill[blue!80!black] (7.7210,5.6472) circle (1.2pt);
\fill[blue!80!black] (7.7214,5.6469) circle (1.2pt);
\fill[blue!80!black] (7.7209,5.6468) circle (1.2pt);
\fill[blue!80!black] (7.7211,5.6465) circle (1.2pt);
\fill[blue!80!black] (7.7203,5.6462) circle (1.2pt);
\fill[blue!80!black] (7.7210,5.6458) circle (1.2pt);
\fill[blue!80!black] (7.7203,5.6457) circle (1.2pt);
\fill[blue!80!black] (7.7206,5.6454) circle (1.2pt);
\fill[blue!80!black] (7.7201,5.6453) circle (1.2pt);
\fill[blue!80!black] (7.7205,5.6447) circle (1.2pt);
\fill[blue!80!black] (7.7196,5.6447) circle (1.2pt);
\fill[blue!80!black] (7.7202,5.6443) circle (1.2pt);
\fill[blue!80!black] (7.7195,5.6442) circle (1.2pt);
\fill[blue!80!black] (7.7198,5.6439) circle (1.2pt);
\fill[blue!80!black] (7.7190,5.6437) circle (1.2pt);
\fill[blue!80!black] (7.7197,5.6433) circle (1.2pt);
\fill[blue!80!black] (7.7189,5.6432) circle (1.2pt);
\fill[blue!80!black] (7.7194,5.6429) circle (1.2pt);
\fill[blue!80!black] (7.7188,5.6428) circle (1.2pt);
\fill[blue!80!black] (7.7190,5.6425) circle (1.2pt);
\fill[blue!80!black] (7.7183,5.6422) circle (1.2pt);
\fill[blue!80!black] (7.7189,5.6418) circle (1.2pt);
\fill[blue!80!black] (7.7182,5.6417) circle (1.2pt);
\fill[blue!80!black] (7.7185,5.6414) circle (1.2pt);
\fill[blue!80!black] (7.7181,5.6413) circle (1.2pt);
\fill[blue!80!black] (7.7184,5.6407) circle (1.2pt);
\fill[blue!80!black] (7.7176,5.6407) circle (1.2pt);
\fill[blue!80!black] (7.7181,5.6403) circle (1.2pt);
\fill[blue!80!black] (7.7175,5.6402) circle (1.2pt);
\fill[blue!80!black] (7.7177,5.6399) circle (1.2pt);
\fill[blue!80!black] (7.7169,5.6397) circle (1.2pt);
\fill[blue!80!black] (7.7176,5.6393) circle (1.2pt);
\fill[blue!80!black] (7.7168,5.6392) circle (1.2pt);
\fill[blue!80!black] (7.7172,5.6389) circle (1.2pt);
\fill[blue!80!black] (7.7167,5.6388) circle (1.2pt);
\fill[blue!80!black] (7.7172,5.6382) circle (1.2pt);
\fill[blue!80!black] (7.7162,5.6382) circle (1.2pt);
\fill[blue!80!black] (7.7168,5.6378) circle (1.2pt);
\fill[blue!80!black] (7.7161,5.6377) circle (1.2pt);
\fill[blue!80!black] (7.7164,5.6374) circle (1.2pt);
\fill[blue!80!black] (7.7155,5.6372) circle (1.2pt);
\fill[blue!80!black] (7.7163,5.6367) circle (1.2pt);
\fill[blue!80!black] (7.7155,5.6367) circle (1.2pt);
\fill[blue!80!black] (7.7159,5.6364) circle (1.2pt);
\fill[blue!80!black] (7.7154,5.6363) circle (1.2pt);
\fill[blue!80!black] (7.7159,5.6357) circle (1.2pt);
\fill[blue!80!black] (7.7149,5.6357) circle (1.2pt);
\fill[blue!80!black] (7.7155,5.6353) circle (1.2pt);
\fill[blue!80!black] (7.7148,5.6352) circle (1.2pt);
\fill[blue!80!black] (7.7151,5.6349) circle (1.2pt);
\fill[blue!80!black] (7.7147,5.6348) circle (1.2pt);
\fill[blue!80!black] (7.7150,5.6342) circle (1.2pt);
\fill[blue!80!black] (7.7142,5.6342) circle (1.2pt);
\fill[blue!80!black] (7.7146,5.6338) circle (1.2pt);
\fill[blue!80!black] (7.7141,5.6337) circle (1.2pt);
\fill[blue!80!black] (7.7143,5.6335) circle (1.2pt);
\fill[blue!80!black] (7.7135,5.6332) circle (1.2pt);
\fill[blue!80!black] (7.7142,5.6328) circle (1.2pt);
\fill[blue!80!black] (7.7135,5.6327) circle (1.2pt);
\fill[blue!80!black] (7.7138,5.6324) circle (1.2pt);
\fill[blue!80!black] (7.7133,5.6323) circle (1.2pt);
\fill[blue!80!black] (7.7137,5.6317) circle (1.2pt);
\fill[blue!80!black] (7.7129,5.6317) circle (1.2pt);
\fill[blue!80!black] (7.7133,5.6313) circle (1.2pt);
\fill[blue!80!black] (7.7128,5.6312) circle (1.2pt);
\fill[blue!80!black] (7.7130,5.6310) circle (1.2pt);
\fill[blue!80!black] (7.7122,5.6307) circle (1.2pt);
\fill[blue!80!black] (7.7128,5.6303) circle (1.2pt);
\fill[blue!80!black] (7.7122,5.6302) circle (1.2pt);
\fill[blue!80!black] (7.7125,5.6299) circle (1.2pt);
\fill[blue!80!black] (7.7120,5.6298) circle (1.2pt);
\fill[blue!80!black] (7.7124,5.6292) circle (1.2pt);
\fill[blue!80!black] (7.7115,5.6292) circle (1.2pt);
\fill[blue!80!black] (7.7120,5.6288) circle (1.2pt);
\fill[blue!80!black] (7.7114,5.6287) circle (1.2pt);
\fill[blue!80!black] (7.7117,5.6285) circle (1.2pt);
\fill[blue!80!black] (7.7109,5.6282) circle (1.2pt);
\fill[blue!80!black] (7.7115,5.6278) circle (1.2pt);
\fill[blue!80!black] (7.7109,5.6277) circle (1.2pt);
\fill[blue!80!black] (7.7112,5.6274) circle (1.2pt);
\fill[blue!80!black] (7.7103,5.6272) circle (1.2pt);
\fill[blue!80!black] (7.7111,5.6267) circle (1.2pt);
\fill[blue!80!black] (7.7102,5.6267) circle (1.2pt);
\fill[blue!80!black] (7.7107,5.6264) circle (1.2pt);
\fill[blue!80!black] (7.7101,5.6262) circle (1.2pt);
\fill[blue!80!black] (7.7106,5.6257) circle (1.2pt);
\fill[blue!80!black] (7.7096,5.6257) circle (1.2pt);
\fill[blue!80!black] (7.7102,5.6253) circle (1.2pt);
\fill[blue!80!black] (7.7096,5.6252) circle (1.2pt);
\fill[blue!80!black] (7.7099,5.6249) circle (1.2pt);
\fill[blue!80!black] (7.7090,5.6247) circle (1.2pt);
\fill[blue!80!black] (7.7097,5.6243) circle (1.2pt);
\fill[blue!80!black] (7.7090,5.6242) circle (1.2pt);
\fill[blue!80!black] (7.7094,5.6239) circle (1.2pt);
\fill[blue!80!black] (7.7089,5.6238) circle (1.2pt);
\fill[blue!80!black] (7.7093,5.6232) circle (1.2pt);
\fill[blue!80!black] (7.7084,5.6232) circle (1.2pt);
\fill[blue!80!black] (7.7089,5.6228) circle (1.2pt);
\fill[blue!80!black] (7.7083,5.6227) circle (1.2pt);
\fill[blue!80!black] (7.7085,5.6224) circle (1.2pt);
\fill[blue!80!black] (7.7077,5.6222) circle (1.2pt);
\fill[blue!80!black] (7.7084,5.6218) circle (1.2pt);
\fill[blue!80!black] (7.7077,5.6217) circle (1.2pt);
\fill[blue!80!black] (7.7080,5.6214) circle (1.2pt);
\fill[blue!80!black] (7.7076,5.6213) circle (1.2pt);
\fill[blue!80!black] (7.7079,5.6207) circle (1.2pt);
\fill[blue!80!black] (7.7071,5.6207) circle (1.2pt);
\fill[blue!80!black] (7.7075,5.6203) circle (1.2pt);
\fill[blue!80!black] (7.7070,5.6202) circle (1.2pt);
\fill[blue!80!black] (7.7074,5.6197) circle (1.2pt);
\fill[blue!80!black] (7.7065,5.6197) circle (1.2pt);
\fill[blue!80!black] (7.7070,5.6193) circle (1.2pt);
\fill[blue!80!black] (7.7064,5.6192) circle (1.2pt);
\fill[blue!80!black] (7.7067,5.6189) circle (1.2pt);
\fill[blue!80!black] (7.7059,5.6186) circle (1.2pt);
\fill[blue!80!black] (7.7065,5.6183) circle (1.2pt);
\fill[blue!80!black] (7.7058,5.6182) circle (1.2pt);
\fill[blue!80!black] (7.7062,5.6179) circle (1.2pt);
\fill[blue!80!black] (7.7057,5.6178) circle (1.2pt);
\fill[blue!80!black] (7.7061,5.6172) circle (1.2pt);
\fill[blue!80!black] (7.7052,5.6172) circle (1.2pt);
\fill[blue!80!black] (7.7057,5.6168) circle (1.2pt);
\fill[blue!80!black] (7.7051,5.6167) circle (1.2pt);
\fill[blue!80!black] (7.7056,5.6162) circle (1.2pt);
\fill[blue!80!black] (7.7046,5.6162) circle (1.2pt);
\fill[blue!80!black] (7.7052,5.6158) circle (1.2pt);
\fill[blue!80!black] (7.7046,5.6157) circle (1.2pt);
\fill[blue!80!black] (7.7049,5.6154) circle (1.2pt);
\fill[blue!80!black] (7.7040,5.6151) circle (1.2pt);
\fill[blue!80!black] (7.7047,5.6148) circle (1.2pt);
\fill[blue!80!black] (7.7040,5.6147) circle (1.2pt);
\fill[blue!80!black] (7.7043,5.6144) circle (1.2pt);
\fill[blue!80!black] (7.7039,5.6143) circle (1.2pt);
\fill[blue!80!black] (7.7042,5.6137) circle (1.2pt);
\fill[blue!80!black] (7.7034,5.6137) circle (1.2pt);
\fill[blue!80!black] (7.7038,5.6133) circle (1.2pt);
\fill[blue!80!black] (7.7033,5.6132) circle (1.2pt);
\fill[blue!80!black] (7.7037,5.6127) circle (1.2pt);
\fill[blue!80!black] (7.7028,5.6127) circle (1.2pt);
\fill[blue!80!black] (7.7033,5.6123) circle (1.2pt);
\fill[blue!80!black] (7.7027,5.6122) circle (1.2pt);
\fill[blue!80!black] (7.7030,5.6119) circle (1.2pt);
\fill[blue!80!black] (7.7022,5.6116) circle (1.2pt);
\fill[blue!80!black] (7.7028,5.6113) circle (1.2pt);
\fill[blue!80!black] (7.7021,5.6112) circle (1.2pt);
\fill[blue!80!black] (7.7025,5.6109) circle (1.2pt);
\fill[blue!80!black] (7.7016,5.6106) circle (1.2pt);
\fill[blue!80!black] (7.7023,5.6102) circle (1.2pt);
\fill[blue!80!black] (7.7016,5.6102) circle (1.2pt);
\fill[blue!80!black] (7.7020,5.6099) circle (1.2pt);
\fill[blue!80!black] (7.7015,5.6097) circle (1.2pt);
\fill[blue!80!black] (7.7019,5.6092) circle (1.2pt);
\fill[blue!80!black] (7.7010,5.6092) circle (1.2pt);
\fill[blue!80!black] (7.7015,5.6088) circle (1.2pt);
\fill[blue!80!black] (7.7009,5.6087) circle (1.2pt);
\fill[blue!80!black] (7.7014,5.6082) circle (1.2pt);
\fill[blue!80!black] (7.7004,5.6082) circle (1.2pt);
\fill[blue!80!black] (7.7010,5.6078) circle (1.2pt);
\fill[blue!80!black] (7.7003,5.6077) circle (1.2pt);
\fill[blue!80!black] (7.7006,5.6074) circle (1.2pt);
\fill[blue!80!black] (7.6998,5.6071) circle (1.2pt);
\fill[blue!80!black] (7.7005,5.6068) circle (1.2pt);
\fill[blue!80!black] (7.6998,5.6067) circle (1.2pt);
\fill[blue!80!black] (7.7001,5.6064) circle (1.2pt);
\fill[blue!80!black] (7.6992,5.6061) circle (1.2pt);
\fill[blue!80!black] (7.7000,5.6057) circle (1.2pt);
\fill[blue!80!black] (7.6992,5.6057) circle (1.2pt);
\fill[blue!80!black] (7.6996,5.6054) circle (1.2pt);
\fill[blue!80!black] (7.6991,5.6052) circle (1.2pt);
\fill[blue!80!black] (7.6995,5.6047) circle (1.2pt);
\fill[blue!80!black] (7.6986,5.6047) circle (1.2pt);
\fill[blue!80!black] (7.6991,5.6043) circle (1.2pt);
\fill[blue!80!black] (7.6985,5.6042) circle (1.2pt);
\fill[blue!80!black] (7.6990,5.6037) circle (1.2pt);
\fill[blue!80!black] (7.6980,5.6037) circle (1.2pt);
\fill[blue!80!black] (7.6986,5.6033) circle (1.2pt);
\fill[blue!80!black] (7.6980,5.6032) circle (1.2pt);
\fill[blue!80!black] (7.6982,5.6029) circle (1.2pt);
\fill[blue!80!black] (7.6975,5.6026) circle (1.2pt);
\fill[blue!80!black] (7.6980,5.6023) circle (1.2pt);
\fill[blue!80!black] (7.6974,5.6022) circle (1.2pt);
\fill[blue!80!black] (7.6977,5.6019) circle (1.2pt);
\fill[blue!80!black] (7.6969,5.6016) circle (1.2pt);
\fill[blue!80!black] (7.6975,5.6013) circle (1.2pt);
\fill[blue!80!black] (7.6968,5.6012) circle (1.2pt);
\fill[blue!80!black] (7.6972,5.6009) circle (1.2pt);
\fill[blue!80!black] (7.6963,5.6006) circle (1.2pt);
\fill[blue!80!black] (7.6970,5.6002) circle (1.2pt);
\fill[blue!80!black] (7.6963,5.6002) circle (1.2pt);
\fill[blue!80!black] (7.6967,5.5999) circle (1.2pt);
\fill[blue!80!black] (7.6962,5.5997) circle (1.2pt);
\fill[blue!80!black] (7.6965,5.5992) circle (1.2pt);
\fill[blue!80!black] (7.6957,5.5992) circle (1.2pt);
\fill[blue!80!black] (7.6961,5.5988) circle (1.2pt);
\fill[blue!80!black] (7.6956,5.5987) circle (1.2pt);
\fill[blue!80!black] (7.6960,5.5982) circle (1.2pt);
\fill[blue!80!black] (7.6952,5.5982) circle (1.2pt);
\fill[blue!80!black] (7.6956,5.5978) circle (1.2pt);
\fill[blue!80!black] (7.6951,5.5977) circle (1.2pt);
\fill[blue!80!black] (7.6955,5.5972) circle (1.2pt);
\fill[blue!80!black] (7.6946,5.5971) circle (1.2pt);
\fill[blue!80!black] (7.6951,5.5968) circle (1.2pt);
\fill[blue!80!black] (7.6945,5.5967) circle (1.2pt);
\fill[blue!80!black] (7.6948,5.5964) circle (1.2pt);
\fill[blue!80!black] (7.6940,5.5961) circle (1.2pt);
\fill[blue!80!black] (7.6946,5.5958) circle (1.2pt);
\fill[blue!80!black] (7.6940,5.5957) circle (1.2pt);
\fill[blue!80!black] (7.6943,5.5954) circle (1.2pt);
\fill[blue!80!black] (7.6935,5.5951) circle (1.2pt);
\fill[blue!80!black] (7.6941,5.5948) circle (1.2pt);
\fill[blue!80!black] (7.6934,5.5947) circle (1.2pt);
\fill[blue!80!black] (7.6937,5.5944) circle (1.2pt);
\fill[blue!80!black] (7.6929,5.5941) circle (1.2pt);
\fill[blue!80!black] (7.6936,5.5937) circle (1.2pt);
\fill[blue!80!black] (7.6929,5.5937) circle (1.2pt);
\fill[blue!80!black] (7.6932,5.5934) circle (1.2pt);
\fill[blue!80!black] (7.6923,5.5931) circle (1.2pt);
\fill[blue!80!black] (7.6930,5.5927) circle (1.2pt);
\fill[blue!80!black] (7.6923,5.5927) circle (1.2pt);
\fill[blue!80!black] (7.6927,5.5924) circle (1.2pt);
\fill[blue!80!black] (7.6922,5.5922) circle (1.2pt);
\fill[blue!80!black] (7.6925,5.5917) circle (1.2pt);
\fill[blue!80!black] (7.6918,5.5917) circle (1.2pt);
\fill[blue!80!black] (7.6922,5.5913) circle (1.2pt);
\fill[blue!80!black] (7.6917,5.5912) circle (1.2pt);
\fill[blue!80!black] (7.6920,5.5907) circle (1.2pt);
\fill[blue!80!black] (7.6912,5.5907) circle (1.2pt);
\fill[blue!80!black] (7.6916,5.5903) circle (1.2pt);
\fill[blue!80!black] (7.6911,5.5902) circle (1.2pt);
\fill[blue!80!black] (7.6915,5.5897) circle (1.2pt);
\fill[blue!80!black] (7.6907,5.5897) circle (1.2pt);
\fill[blue!80!black] (7.6911,5.5893) circle (1.2pt);
\fill[blue!80!black] (7.6906,5.5892) circle (1.2pt);
\fill[blue!80!black] (7.6910,5.5887) circle (1.2pt);
\fill[blue!80!black] (7.6901,5.5887) circle (1.2pt);
\fill[blue!80!black] (7.6906,5.5883) circle (1.2pt);
\fill[blue!80!black] (7.6900,5.5882) circle (1.2pt);
\fill[blue!80!black] (7.6905,5.5877) circle (1.2pt);
\fill[blue!80!black] (7.6896,5.5877) circle (1.2pt);
\fill[blue!80!black] (7.6900,5.5873) circle (1.2pt);
\fill[blue!80!black] (7.6895,5.5872) circle (1.2pt);
\fill[blue!80!black] (7.6899,5.5867) circle (1.2pt);
\fill[blue!80!black] (7.6890,5.5866) circle (1.2pt);
\fill[blue!80!black] (7.6895,5.5863) circle (1.2pt);
\fill[blue!80!black] (7.6890,5.5862) circle (1.2pt);
\fill[blue!80!black] (7.6894,5.5856) circle (1.2pt);
\fill[blue!80!black] (7.6885,5.5856) circle (1.2pt);
\fill[blue!80!black] (7.6890,5.5853) circle (1.2pt);
\fill[blue!80!black] (7.6884,5.5852) circle (1.2pt);
\fill[blue!80!black] (7.6889,5.5846) circle (1.2pt);
\fill[blue!80!black] (7.6879,5.5846) circle (1.2pt);
\fill[blue!80!black] (7.6885,5.5843) circle (1.2pt);
\fill[blue!80!black] (7.6879,5.5842) circle (1.2pt);
\fill[blue!80!black] (7.6881,5.5839) circle (1.2pt);
\fill[blue!80!black] (7.6874,5.5836) circle (1.2pt);
\fill[blue!80!black] (7.6879,5.5833) circle (1.2pt);
\fill[blue!80!black] (7.6873,5.5832) circle (1.2pt);
\fill[blue!80!black] (7.6876,5.5829) circle (1.2pt);
\fill[blue!80!black] (7.6869,5.5826) circle (1.2pt);
\fill[blue!80!black] (7.6874,5.5823) circle (1.2pt);
\fill[blue!80!black] (7.6868,5.5822) circle (1.2pt);
\fill[blue!80!black] (7.6871,5.5819) circle (1.2pt);
\fill[blue!80!black] (7.6863,5.5816) circle (1.2pt);
\fill[blue!80!black] (7.6869,5.5813) circle (1.2pt);
\fill[blue!80!black] (7.6863,5.5812) circle (1.2pt);
\fill[blue!80!black] (7.6865,5.5809) circle (1.2pt);
\fill[blue!80!black] (7.6858,5.5806) circle (1.2pt);
\fill[blue!80!black] (7.6863,5.5803) circle (1.2pt);
\fill[blue!80!black] (7.6857,5.5802) circle (1.2pt);
\fill[blue!80!black] (7.6860,5.5799) circle (1.2pt);
\fill[blue!80!black] (7.6852,5.5796) circle (1.2pt);
\fill[blue!80!black] (7.6858,5.5793) circle (1.2pt);
\fill[blue!80!black] (7.6852,5.5792) circle (1.2pt);
\fill[blue!80!black] (7.6855,5.5789) circle (1.2pt);
\fill[blue!80!black] (7.6847,5.5786) circle (1.2pt);
\fill[blue!80!black] (7.6853,5.5783) circle (1.2pt);
\fill[blue!80!black] (7.6847,5.5782) circle (1.2pt);
\fill[blue!80!black] (7.6849,5.5779) circle (1.2pt);
\fill[blue!80!black] (7.6842,5.5776) circle (1.2pt);
\fill[blue!80!black] (7.6847,5.5773) circle (1.2pt);
\fill[blue!80!black] (7.6841,5.5772) circle (1.2pt);
\fill[blue!80!black] (7.6844,5.5769) circle (1.2pt);
\fill[blue!80!black] (7.6837,5.5766) circle (1.2pt);
\fill[blue!80!black] (7.6842,5.5763) circle (1.2pt);
\fill[blue!80!black] (7.6836,5.5762) circle (1.2pt);
\fill[blue!80!black] (7.6841,5.5756) circle (1.2pt);
\fill[blue!80!black] (7.6831,5.5756) circle (1.2pt);
\fill[blue!80!black] (7.6837,5.5753) circle (1.2pt);
\fill[blue!80!black] (7.6831,5.5752) circle (1.2pt);
\fill[blue!80!black] (7.6836,5.5746) circle (1.2pt);
\fill[blue!80!black] (7.6826,5.5746) circle (1.2pt);
\fill[blue!80!black] (7.6831,5.5743) circle (1.2pt);
\fill[blue!80!black] (7.6826,5.5742) circle (1.2pt);
\fill[blue!80!black] (7.6830,5.5737) circle (1.2pt);
\fill[blue!80!black] (7.6821,5.5737) circle (1.2pt);
\fill[blue!80!black] (7.6826,5.5733) circle (1.2pt);
\fill[blue!80!black] (7.6820,5.5732) circle (1.2pt);
\fill[blue!80!black] (7.6825,5.5727) circle (1.2pt);
\fill[blue!80!black] (7.6816,5.5727) circle (1.2pt);
\fill[blue!80!black] (7.6820,5.5723) circle (1.2pt);
\fill[blue!80!black] (7.6815,5.5722) circle (1.2pt);
\fill[blue!80!black] (7.6819,5.5717) circle (1.2pt);
\fill[blue!80!black] (7.6811,5.5717) circle (1.2pt);
\fill[blue!80!black] (7.6815,5.5713) circle (1.2pt);
\fill[blue!80!black] (7.6810,5.5712) circle (1.2pt);
\fill[blue!80!black] (7.6814,5.5707) circle (1.2pt);
\fill[blue!80!black] (7.6805,5.5707) circle (1.2pt);
\fill[blue!80!black] (7.6810,5.5703) circle (1.2pt);
\fill[blue!80!black] (7.6805,5.5702) circle (1.2pt);
\fill[blue!80!black] (7.6808,5.5697) circle (1.2pt);
\fill[blue!80!black] (7.6800,5.5697) circle (1.2pt);
\fill[blue!80!black] (7.6804,5.5694) circle (1.2pt);
\fill[blue!80!black] (7.6799,5.5692) circle (1.2pt);
\fill[blue!80!black] (7.6803,5.5687) circle (1.2pt);
\fill[blue!80!black] (7.6795,5.5687) circle (1.2pt);
\fill[blue!80!black] (7.6799,5.5684) circle (1.2pt);
\fill[blue!80!black] (7.6794,5.5682) circle (1.2pt);
\fill[blue!80!black] (7.6797,5.5677) circle (1.2pt);
\fill[blue!80!black] (7.6790,5.5677) circle (1.2pt);
\fill[blue!80!black] (7.6793,5.5674) circle (1.2pt);
\fill[blue!80!black] (7.6785,5.5671) circle (1.2pt);
\fill[blue!80!black] (7.6792,5.5667) circle (1.2pt);
\fill[blue!80!black] (7.6785,5.5667) circle (1.2pt);
\fill[blue!80!black] (7.6788,5.5664) circle (1.2pt);
\fill[blue!80!black] (7.6780,5.5661) circle (1.2pt);
\fill[blue!80!black] (7.6786,5.5658) circle (1.2pt);
\fill[blue!80!black] (7.6780,5.5657) circle (1.2pt);
\fill[blue!80!black] (7.6783,5.5654) circle (1.2pt);
\fill[blue!80!black] (7.6775,5.5652) circle (1.2pt);
\fill[blue!80!black] (7.6781,5.5648) circle (1.2pt);
\fill[blue!80!black] (7.6775,5.5647) circle (1.2pt);
\fill[blue!80!black] (7.6777,5.5644) circle (1.2pt);
\fill[blue!80!black] (7.6770,5.5642) circle (1.2pt);
\fill[blue!80!black] (7.6775,5.5638) circle (1.2pt);
\fill[blue!80!black] (7.6769,5.5637) circle (1.2pt);
\fill[blue!80!black] (7.6774,5.5632) circle (1.2pt);
\fill[blue!80!black] (7.6765,5.5632) circle (1.2pt);
\fill[blue!80!black] (7.6770,5.5628) circle (1.2pt);
\fill[blue!80!black] (7.6764,5.5627) circle (1.2pt);
\fill[blue!80!black] (7.6768,5.5622) circle (1.2pt);
\fill[blue!80!black] (7.6760,5.5622) circle (1.2pt);
\fill[blue!80!black] (7.6764,5.5618) circle (1.2pt);
\fill[blue!80!black] (7.6759,5.5617) circle (1.2pt);
\fill[blue!80!black] (7.6763,5.5612) circle (1.2pt);
\fill[blue!80!black] (7.6755,5.5612) circle (1.2pt);
\fill[blue!80!black] (7.6759,5.5609) circle (1.2pt);
\fill[blue!80!black] (7.6754,5.5608) circle (1.2pt);
\fill[blue!80!black] (7.6757,5.5602) circle (1.2pt);
\fill[blue!80!black] (7.6750,5.5602) circle (1.2pt);
\fill[blue!80!black] (7.6753,5.5599) circle (1.2pt);
\fill[blue!80!black] (7.6745,5.5597) circle (1.2pt);
\fill[blue!80!black] (7.6752,5.5593) circle (1.2pt);
\fill[blue!80!black] (7.6745,5.5592) circle (1.2pt);
\fill[blue!80!black] (7.6748,5.5589) circle (1.2pt);
\fill[blue!80!black] (7.6740,5.5587) circle (1.2pt);
\fill[blue!80!black] (7.6746,5.5583) circle (1.2pt);
\fill[blue!80!black] (7.6740,5.5582) circle (1.2pt);
\fill[blue!80!black] (7.6742,5.5579) circle (1.2pt);
\fill[blue!80!black] (7.6735,5.5577) circle (1.2pt);
\fill[blue!80!black] (7.6740,5.5573) circle (1.2pt);
\fill[blue!80!black] (7.6735,5.5572) circle (1.2pt);
\fill[blue!80!black] (7.6739,5.5567) circle (1.2pt);
\fill[blue!80!black] (7.6730,5.5567) circle (1.2pt);
\fill[blue!80!black] (7.6735,5.5563) circle (1.2pt);
\fill[blue!80!black] (7.6730,5.5562) circle (1.2pt);
\fill[blue!80!black] (7.6733,5.5557) circle (1.2pt);
\fill[blue!80!black] (7.6725,5.5557) circle (1.2pt);
\fill[blue!80!black] (7.6729,5.5554) circle (1.2pt);
\fill[blue!80!black] (7.6725,5.5553) circle (1.2pt);
\fill[blue!80!black] (7.6728,5.5547) circle (1.2pt);
\fill[blue!80!black] (7.6720,5.5547) circle (1.2pt);
\fill[blue!80!black] (7.6724,5.5544) circle (1.2pt);
\fill[blue!80!black] (7.6715,5.5542) circle (1.2pt);
\fill[blue!80!black] (7.6722,5.5538) circle (1.2pt);
\fill[blue!80!black] (7.6715,5.5537) circle (1.2pt);
\fill[blue!80!black] (7.6718,5.5534) circle (1.2pt);
\fill[blue!80!black] (7.6710,5.5532) circle (1.2pt);
\fill[blue!80!black] (7.6716,5.5528) circle (1.2pt);
\fill[blue!80!black] (7.6710,5.5527) circle (1.2pt);
\fill[blue!80!black] (7.6713,5.5525) circle (1.2pt);
\fill[blue!80!black] (7.6706,5.5522) circle (1.2pt);
\fill[blue!80!black] (7.6711,5.5518) circle (1.2pt);
\fill[blue!80!black] (7.6705,5.5518) circle (1.2pt);
\fill[blue!80!black] (7.6709,5.5512) circle (1.2pt);
\fill[blue!80!black] (7.6701,5.5512) circle (1.2pt);
\fill[blue!80!black] (7.6705,5.5509) circle (1.2pt);
\fill[blue!80!black] (7.6700,5.5508) circle (1.2pt);
\fill[blue!80!black] (7.6704,5.5503) circle (1.2pt);
\fill[blue!80!black] (7.6696,5.5502) circle (1.2pt);
\fill[blue!80!black] (7.6700,5.5499) circle (1.2pt);
\fill[blue!80!black] (7.6691,5.5497) circle (1.2pt);
\fill[blue!80!black] (7.6698,5.5493) circle (1.2pt);
\fill[blue!80!black] (7.6691,5.5492) circle (1.2pt);
\fill[blue!80!black] (7.6694,5.5489) circle (1.2pt);
\fill[blue!80!black] (7.6686,5.5487) circle (1.2pt);
\fill[blue!80!black] (7.6692,5.5483) circle (1.2pt);
\fill[blue!80!black] (7.6686,5.5483) circle (1.2pt);
\fill[blue!80!black] (7.6691,5.5477) circle (1.2pt);
\fill[blue!80!black] (7.6682,5.5477) circle (1.2pt);
\fill[blue!80!black] (7.6686,5.5474) circle (1.2pt);
\fill[blue!80!black] (7.6681,5.5473) circle (1.2pt);
\fill[blue!80!black] (7.6685,5.5467) circle (1.2pt);
\fill[blue!80!black] (7.6677,5.5467) circle (1.2pt);
\fill[blue!80!black] (7.6681,5.5464) circle (1.2pt);
\fill[blue!80!black] (7.6676,5.5463) circle (1.2pt);
\fill[blue!80!black] (7.6679,5.5458) circle (1.2pt);
\fill[blue!80!black] (7.6672,5.5457) circle (1.2pt);
\fill[blue!80!black] (7.6675,5.5454) circle (1.2pt);
\fill[blue!80!black] (7.6667,5.5452) circle (1.2pt);
\fill[blue!80!black] (7.6673,5.5448) circle (1.2pt);
\fill[blue!80!black] (7.6667,5.5448) circle (1.2pt);
\fill[blue!80!black] (7.6670,5.5445) circle (1.2pt);
\fill[blue!80!black] (7.6663,5.5442) circle (1.2pt);
\fill[blue!80!black] (7.6668,5.5439) circle (1.2pt);
\fill[blue!80!black] (7.6662,5.5438) circle (1.2pt);
\fill[blue!80!black] (7.6666,5.5433) circle (1.2pt);
\fill[blue!80!black] (7.6658,5.5432) circle (1.2pt);
\fill[blue!80!black] (7.6662,5.5429) circle (1.2pt);
\fill[blue!80!black] (7.6658,5.5428) circle (1.2pt);
\fill[blue!80!black] (7.6660,5.5423) circle (1.2pt);
\fill[blue!80!black] (7.6653,5.5423) circle (1.2pt);
\fill[blue!80!black] (7.6657,5.5420) circle (1.2pt);
\fill[blue!80!black] (7.6649,5.5417) circle (1.2pt);
\fill[blue!80!black] (7.6655,5.5414) circle (1.2pt);
\fill[blue!80!black] (7.6649,5.5413) circle (1.2pt);
\fill[blue!80!black] (7.6653,5.5407) circle (1.2pt);
\fill[blue!80!black] (7.6644,5.5407) circle (1.2pt);
\fill[blue!80!black] (7.6649,5.5404) circle (1.2pt);
\fill[blue!80!black] (7.6644,5.5403) circle (1.2pt);
\fill[blue!80!black] (7.6647,5.5398) circle (1.2pt);
\fill[blue!80!black] (7.6639,5.5398) circle (1.2pt);
\fill[blue!80!black] (7.6643,5.5394) circle (1.2pt);
\fill[blue!80!black] (7.6634,5.5392) circle (1.2pt);
\fill[blue!80!black] (7.6641,5.5388) circle (1.2pt);
\fill[blue!80!black] (7.6635,5.5388) circle (1.2pt);
\fill[blue!80!black] (7.6638,5.5385) circle (1.2pt);
\fill[blue!80!black] (7.6630,5.5382) circle (1.2pt);
\fill[blue!80!black] (7.6636,5.5379) circle (1.2pt);
\fill[blue!80!black] (7.6630,5.5378) circle (1.2pt);
\fill[blue!80!black] (7.6634,5.5373) circle (1.2pt);
\fill[blue!80!black] (7.6626,5.5373) circle (1.2pt);
\fill[blue!80!black] (7.6630,5.5369) circle (1.2pt);
\fill[blue!80!black] (7.6625,5.5368) circle (1.2pt);
\fill[blue!80!black] (7.6628,5.5363) circle (1.2pt);
\fill[blue!80!black] (7.6621,5.5363) circle (1.2pt);
\fill[blue!80!black] (7.6624,5.5360) circle (1.2pt);
\fill[blue!80!black] (7.6616,5.5357) circle (1.2pt);
\fill[blue!80!black] (7.6622,5.5354) circle (1.2pt);
\fill[blue!80!black] (7.6616,5.5353) circle (1.2pt);
\fill[blue!80!black] (7.6621,5.5348) circle (1.2pt);
\fill[blue!80!black] (7.6612,5.5348) circle (1.2pt);
\fill[blue!80!black] (7.6617,5.5344) circle (1.2pt);
\fill[blue!80!black] (7.6612,5.5343) circle (1.2pt);
\fill[blue!80!black] (7.6615,5.5338) circle (1.2pt);
\fill[blue!80!black] (7.6607,5.5338) circle (1.2pt);
\fill[blue!80!black] (7.6611,5.5335) circle (1.2pt);
\fill[blue!80!black] (7.6602,5.5333) circle (1.2pt);
\fill[blue!80!black] (7.6609,5.5329) circle (1.2pt);
\fill[blue!80!black] (7.6603,5.5328) circle (1.2pt);
\fill[blue!80!black] (7.6605,5.5325) circle (1.2pt);
\fill[blue!80!black] (7.6598,5.5323) circle (1.2pt);
\fill[blue!80!black] (7.6603,5.5319) circle (1.2pt);
\fill[blue!80!black] (7.6598,5.5318) circle (1.2pt);
\fill[blue!80!black] (7.6602,5.5313) circle (1.2pt);
\fill[blue!80!black] (7.6594,5.5313) circle (1.2pt);
\fill[blue!80!black] (7.6597,5.5310) circle (1.2pt);
\fill[blue!80!black] (7.6589,5.5308) circle (1.2pt);
\fill[blue!80!black] (7.6596,5.5304) circle (1.2pt);
\fill[blue!80!black] (7.6589,5.5303) circle (1.2pt);
\fill[blue!80!black] (7.6592,5.5300) circle (1.2pt);
\fill[blue!80!black] (7.6585,5.5298) circle (1.2pt);
\fill[blue!80!black] (7.6590,5.5294) circle (1.2pt);
\fill[blue!80!black] (7.6584,5.5294) circle (1.2pt);
\fill[blue!80!black] (7.6588,5.5288) circle (1.2pt);
\fill[blue!80!black] (7.6580,5.5288) circle (1.2pt);
\fill[blue!80!black] (7.6584,5.5285) circle (1.2pt);
\fill[blue!80!black] (7.6575,5.5283) circle (1.2pt);
\fill[blue!80!black] (7.6582,5.5279) circle (1.2pt);
\fill[blue!80!black] (7.6576,5.5278) circle (1.2pt);
\fill[blue!80!black] (7.6578,5.5276) circle (1.2pt);
\fill[blue!80!black] (7.6571,5.5273) circle (1.2pt);
\fill[blue!80!black] (7.6576,5.5270) circle (1.2pt);
\fill[blue!80!black] (7.6571,5.5269) circle (1.2pt);
\fill[blue!80!black] (7.6575,5.5264) circle (1.2pt);
\fill[blue!80!black] (7.6567,5.5263) circle (1.2pt);
\fill[blue!80!black] (7.6570,5.5260) circle (1.2pt);
\fill[blue!80!black] (7.6562,5.5258) circle (1.2pt);
\fill[blue!80!black] (7.6569,5.5254) circle (1.2pt);
\fill[blue!80!black] (7.6562,5.5254) circle (1.2pt);
\fill[blue!80!black] (7.6565,5.5251) circle (1.2pt);
\fill[blue!80!black] (7.6558,5.5248) circle (1.2pt);
\fill[blue!80!black] (7.6563,5.5245) circle (1.2pt);
\fill[blue!80!black] (7.6558,5.5244) circle (1.2pt);
\fill[blue!80!black] (7.6561,5.5239) circle (1.2pt);
\fill[blue!80!black] (7.6554,5.5239) circle (1.2pt);
\fill[blue!80!black] (7.6557,5.5236) circle (1.2pt);
\fill[blue!80!black] (7.6549,5.5233) circle (1.2pt);
\fill[blue!80!black] (7.6555,5.5230) circle (1.2pt);
\fill[blue!80!black] (7.6549,5.5229) circle (1.2pt);
\fill[blue!80!black] (7.6554,5.5224) circle (1.2pt);
\fill[blue!80!black] (7.6545,5.5223) circle (1.2pt);
\fill[blue!80!black] (7.6549,5.5220) circle (1.2pt);
\fill[blue!80!black] (7.6544,5.5219) circle (1.2pt);
\fill[blue!80!black] (7.6547,5.5214) circle (1.2pt);
\fill[blue!80!black] (7.6540,5.5214) circle (1.2pt);
\fill[blue!80!black] (7.6543,5.5211) circle (1.2pt);
\fill[blue!80!black] (7.6536,5.5208) circle (1.2pt);
\fill[blue!80!black] (7.6541,5.5205) circle (1.2pt);
\fill[blue!80!black] (7.6536,5.5204) circle (1.2pt);
\fill[blue!80!black] (7.6540,5.5199) circle (1.2pt);
\fill[blue!80!black] (7.6532,5.5199) circle (1.2pt);
\fill[blue!80!black] (7.6535,5.5196) circle (1.2pt);
\fill[blue!80!black] (7.6527,5.5193) circle (1.2pt);
\fill[blue!80!black] (7.6533,5.5190) circle (1.2pt);
\fill[blue!80!black] (7.6527,5.5189) circle (1.2pt);
\fill[blue!80!black] (7.6530,5.5186) circle (1.2pt);
\fill[blue!80!black] (7.6523,5.5184) circle (1.2pt);
\fill[blue!80!black] (7.6528,5.5180) circle (1.2pt);
\fill[blue!80!black] (7.6523,5.5180) circle (1.2pt);
\fill[blue!80!black] (7.6526,5.5174) circle (1.2pt);
\fill[blue!80!black] (7.6519,5.5174) circle (1.2pt);
\fill[blue!80!black] (7.6522,5.5171) circle (1.2pt);
\fill[blue!80!black] (7.6514,5.5169) circle (1.2pt);
\fill[blue!80!black] (7.6520,5.5165) circle (1.2pt);
\fill[blue!80!black] (7.6514,5.5164) circle (1.2pt);
\fill[blue!80!black] (7.6518,5.5159) circle (1.2pt);
\fill[blue!80!black] (7.6510,5.5159) circle (1.2pt);
\fill[blue!80!black] (7.6514,5.5156) circle (1.2pt);
\fill[blue!80!black] (7.6505,5.5154) circle (1.2pt);
\fill[blue!80!black] (7.6512,5.5150) circle (1.2pt);
\fill[blue!80!black] (7.6506,5.5149) circle (1.2pt);
\fill[blue!80!black] (7.6511,5.5144) circle (1.2pt);
\fill[blue!80!black] (7.6501,5.5144) circle (1.2pt);
\fill[blue!80!black] (7.6506,5.5141) circle (1.2pt);
\fill[blue!80!black] (7.6501,5.5140) circle (1.2pt);
\fill[blue!80!black] (7.6504,5.5135) circle (1.2pt);
\fill[blue!80!black] (7.6497,5.5134) circle (1.2pt);
\fill[blue!80!black] (7.6500,5.5132) circle (1.2pt);
\fill[blue!80!black] (7.6493,5.5129) circle (1.2pt);
\fill[blue!80!black] (7.6498,5.5126) circle (1.2pt);
\fill[blue!80!black] (7.6493,5.5125) circle (1.2pt);
\fill[blue!80!black] (7.6496,5.5120) circle (1.2pt);
\fill[blue!80!black] (7.6489,5.5119) circle (1.2pt);
\fill[blue!80!black] (7.6492,5.5117) circle (1.2pt);
\fill[blue!80!black] (7.6484,5.5114) circle (1.2pt);
\fill[blue!80!black] (7.6490,5.5111) circle (1.2pt);
\fill[blue!80!black] (7.6484,5.5110) circle (1.2pt);
\fill[blue!80!black] (7.6489,5.5105) circle (1.2pt);
\fill[blue!80!black] (7.6480,5.5105) circle (1.2pt);
\fill[blue!80!black] (7.6484,5.5101) circle (1.2pt);
\fill[blue!80!black] (7.6475,5.5099) circle (1.2pt);
\fill[blue!80!black] (7.6482,5.5096) circle (1.2pt);
\fill[blue!80!black] (7.6476,5.5095) circle (1.2pt);
\fill[blue!80!black] (7.6478,5.5092) circle (1.2pt);
\fill[blue!80!black] (7.6472,5.5090) circle (1.2pt);
\fill[blue!80!black] (7.6476,5.5086) circle (1.2pt);
\fill[blue!80!black] (7.6471,5.5085) circle (1.2pt);
\fill[blue!80!black] (7.6474,5.5080) circle (1.2pt);
\fill[blue!80!black] (7.6468,5.5080) circle (1.2pt);
\fill[blue!80!black] (7.6470,5.5077) circle (1.2pt);
\fill[blue!80!black] (7.6463,5.5075) circle (1.2pt);
\fill[blue!80!black] (7.6468,5.5071) circle (1.2pt);
\fill[blue!80!black] (7.6463,5.5070) circle (1.2pt);
\fill[blue!80!black] (7.6466,5.5065) circle (1.2pt);
\fill[blue!80!black] (7.6459,5.5065) circle (1.2pt);
\fill[blue!80!black] (7.6462,5.5062) circle (1.2pt);
\fill[blue!80!black] (7.6455,5.5060) circle (1.2pt);
\fill[blue!80!black] (7.6460,5.5056) circle (1.2pt);
\fill[blue!80!black] (7.6455,5.5056) circle (1.2pt);
\fill[blue!80!black] (7.6459,5.5050) circle (1.2pt);
\fill[blue!80!black] (7.6451,5.5050) circle (1.2pt);
\fill[blue!80!black] (7.6454,5.5047) circle (1.2pt);
\fill[blue!80!black] (7.6446,5.5045) circle (1.2pt);
\fill[blue!80!black] (7.6452,5.5041) circle (1.2pt);
\fill[blue!80!black] (7.6447,5.5041) circle (1.2pt);
\fill[blue!80!black] (7.6451,5.5035) circle (1.2pt);
\fill[blue!80!black] (7.6443,5.5035) circle (1.2pt);
\fill[blue!80!black] (7.6446,5.5032) circle (1.2pt);
\fill[blue!80!black] (7.6438,5.5030) circle (1.2pt);
\fill[blue!80!black] (7.6444,5.5026) circle (1.2pt);
\fill[blue!80!black] (7.6438,5.5026) circle (1.2pt);
\fill[blue!80!black] (7.6443,5.5020) circle (1.2pt);
\fill[blue!80!black] (7.6434,5.5020) circle (1.2pt);
\fill[blue!80!black] (7.6438,5.5017) circle (1.2pt);
\fill[blue!80!black] (7.6430,5.5015) circle (1.2pt);
\fill[blue!80!black] (7.6436,5.5012) circle (1.2pt);
\fill[blue!80!black] (7.6430,5.5011) circle (1.2pt);
\fill[blue!80!black] (7.6435,5.5006) circle (1.2pt);
\fill[blue!80!black] (7.6426,5.5006) circle (1.2pt);
\fill[blue!80!black] (7.6430,5.5002) circle (1.2pt);
\fill[blue!80!black] (7.6421,5.5000) circle (1.2pt);
\fill[blue!80!black] (7.6428,5.4997) circle (1.2pt);
\fill[blue!80!black] (7.6422,5.4996) circle (1.2pt);
\fill[blue!80!black] (7.6427,5.4991) circle (1.2pt);
\fill[blue!80!black] (7.6418,5.4991) circle (1.2pt);
\fill[blue!80!black] (7.6422,5.4988) circle (1.2pt);
\fill[blue!80!black] (7.6413,5.4986) circle (1.2pt);
\fill[blue!80!black] (7.6420,5.4982) circle (1.2pt);
\fill[blue!80!black] (7.6414,5.4981) circle (1.2pt);
\fill[blue!80!black] (7.6416,5.4978) circle (1.2pt);
\fill[blue!80!black] (7.6410,5.4976) circle (1.2pt);
\fill[blue!80!black] (7.6414,5.4973) circle (1.2pt);
\fill[blue!80!black] (7.6409,5.4972) circle (1.2pt);
\fill[blue!80!black] (7.6412,5.4967) circle (1.2pt);
\fill[blue!80!black] (7.6406,5.4966) circle (1.2pt);
\fill[blue!80!black] (7.6408,5.4964) circle (1.2pt);
\fill[blue!80!black] (7.6402,5.4961) circle (1.2pt);
\fill[blue!80!black] (7.6406,5.4958) circle (1.2pt);
\fill[blue!80!black] (7.6401,5.4957) circle (1.2pt);
\fill[blue!80!black] (7.6404,5.4952) circle (1.2pt);
\fill[blue!80!black] (7.6398,5.4952) circle (1.2pt);
\fill[blue!80!black] (7.6400,5.4949) circle (1.2pt);
\fill[blue!80!black] (7.6393,5.4946) circle (1.2pt);
\fill[blue!80!black] (7.6398,5.4943) circle (1.2pt);
\fill[blue!80!black] (7.6389,5.4941) circle (1.2pt);
\fill[blue!80!black] (7.6396,5.4937) circle (1.2pt);
\fill[blue!80!black] (7.6390,5.4937) circle (1.2pt);
\fill[blue!80!black] (7.6394,5.4931) circle (1.2pt);
\fill[blue!80!black] (7.6385,5.4931) circle (1.2pt);
\fill[blue!80!black] (7.6389,5.4928) circle (1.2pt);
\fill[blue!80!black] (7.6381,5.4926) circle (1.2pt);
\fill[blue!80!black] (7.6387,5.4923) circle (1.2pt);
\fill[blue!80!black] (7.6381,5.4922) circle (1.2pt);
\fill[blue!80!black] (7.6386,5.4917) circle (1.2pt);
\fill[blue!80!black] (7.6377,5.4917) circle (1.2pt);
\fill[blue!80!black] (7.6381,5.4914) circle (1.2pt);
\fill[blue!80!black] (7.6373,5.4912) circle (1.2pt);
\fill[blue!80!black] (7.6379,5.4908) circle (1.2pt);
\fill[blue!80!black] (7.6373,5.4907) circle (1.2pt);
\fill[blue!80!black] (7.6378,5.4902) circle (1.2pt);
\fill[blue!80!black] (7.6369,5.4902) circle (1.2pt);
\fill[blue!80!black] (7.6373,5.4899) circle (1.2pt);
\fill[blue!80!black] (7.6365,5.4897) circle (1.2pt);
\fill[blue!80!black] (7.6371,5.4893) circle (1.2pt);
\fill[blue!80!black] (7.6365,5.4892) circle (1.2pt);
\fill[blue!80!black] (7.6369,5.4887) circle (1.2pt);
\fill[blue!80!black] (7.6362,5.4887) circle (1.2pt);
\fill[blue!80!black] (7.6365,5.4884) circle (1.2pt);
\fill[blue!80!black] (7.6357,5.4882) circle (1.2pt);
\fill[blue!80!black] (7.6363,5.4879) circle (1.2pt);
\fill[blue!80!black] (7.6357,5.4878) circle (1.2pt);
\fill[blue!80!black] (7.6361,5.4873) circle (1.2pt);
\fill[blue!80!black] (7.6354,5.4872) circle (1.2pt);
\fill[blue!80!black] (7.6357,5.4870) circle (1.2pt);
\fill[blue!80!black] (7.6349,5.4867) circle (1.2pt);
\fill[blue!80!black] (7.6355,5.4864) circle (1.2pt);
\fill[blue!80!black] (7.6350,5.4863) circle (1.2pt);
\fill[blue!80!black] (7.6353,5.4858) circle (1.2pt);
\fill[blue!80!black] (7.6346,5.4858) circle (1.2pt);
\fill[blue!80!black] (7.6349,5.4855) circle (1.2pt);
\fill[blue!80!black] (7.6342,5.4853) circle (1.2pt);
\fill[blue!80!black] (7.6346,5.4849) circle (1.2pt);
\fill[blue!80!black] (7.6342,5.4848) circle (1.2pt);
\fill[blue!80!black] (7.6344,5.4844) circle (1.2pt);
\fill[blue!80!black] (7.6338,5.4843) circle (1.2pt);
\fill[blue!80!black] (7.6340,5.4840) circle (1.2pt);
\fill[blue!80!black] (7.6334,5.4838) circle (1.2pt);
\fill[blue!80!black] (7.6338,5.4835) circle (1.2pt);
\fill[blue!80!black] (7.6329,5.4833) circle (1.2pt);
\fill[blue!80!black] (7.6336,5.4829) circle (1.2pt);
\fill[blue!80!black] (7.6330,5.4828) circle (1.2pt);
\fill[blue!80!black] (7.6334,5.4823) circle (1.2pt);
\fill[blue!80!black] (7.6326,5.4823) circle (1.2pt);
\fill[blue!80!black] (7.6330,5.4820) circle (1.2pt);
\fill[blue!80!black] (7.6322,5.4818) circle (1.2pt);
\fill[blue!80!black] (7.6328,5.4814) circle (1.2pt);
\fill[blue!80!black] (7.6322,5.4814) circle (1.2pt);
\fill[blue!80!black] (7.6326,5.4809) circle (1.2pt);
\fill[blue!80!black] (7.6318,5.4808) circle (1.2pt);
\fill[blue!80!black] (7.6322,5.4806) circle (1.2pt);
\fill[blue!80!black] (7.6314,5.4803) circle (1.2pt);
\fill[blue!80!black] (7.6319,5.4800) circle (1.2pt);
\fill[blue!80!black] (7.6314,5.4799) circle (1.2pt);
\fill[blue!80!black] (7.6317,5.4794) circle (1.2pt);
\fill[blue!80!black] (7.6311,5.4794) circle (1.2pt);
\fill[blue!80!black] (7.6313,5.4791) circle (1.2pt);
\fill[blue!80!black] (7.6307,5.4789) circle (1.2pt);
\fill[blue!80!black] (7.6311,5.4785) circle (1.2pt);
\fill[blue!80!black] (7.6302,5.4784) circle (1.2pt);
\fill[blue!80!black] (7.6309,5.4780) circle (1.2pt);
\fill[blue!80!black] (7.6303,5.4779) circle (1.2pt);
\fill[blue!80!black] (7.6307,5.4774) circle (1.2pt);
\fill[blue!80!black] (7.6299,5.4774) circle (1.2pt);
\fill[blue!80!black] (7.6303,5.4771) circle (1.2pt);
\fill[blue!80!black] (7.6295,5.4769) circle (1.2pt);
\fill[blue!80!black] (7.6300,5.4765) circle (1.2pt);
\fill[blue!80!black] (7.6295,5.4765) circle (1.2pt);
\fill[blue!80!black] (7.6299,5.4760) circle (1.2pt);
\fill[blue!80!black] (7.6291,5.4759) circle (1.2pt);
\fill[blue!80!black] (7.6294,5.4757) circle (1.2pt);
\fill[blue!80!black] (7.6287,5.4754) circle (1.2pt);
\fill[blue!80!black] (7.6292,5.4751) circle (1.2pt);
\fill[blue!80!black] (7.6287,5.4750) circle (1.2pt);
\fill[blue!80!black] (7.6290,5.4745) circle (1.2pt);
\fill[blue!80!black] (7.6284,5.4745) circle (1.2pt);
\fill[blue!80!black] (7.6289,5.4739) circle (1.2pt);
\fill[blue!80!black] (7.6280,5.4740) circle (1.2pt);
\fill[blue!80!black] (7.6284,5.4737) circle (1.2pt);
\fill[blue!80!black] (7.6276,5.4734) circle (1.2pt);
\fill[blue!80!black] (7.6282,5.4731) circle (1.2pt);
\fill[blue!80!black] (7.6276,5.4730) circle (1.2pt);
\fill[blue!80!black] (7.6280,5.4725) circle (1.2pt);
\fill[blue!80!black] (7.6272,5.4725) circle (1.2pt);
\fill[blue!80!black] (7.6275,5.4722) circle (1.2pt);
\fill[blue!80!black] (7.6268,5.4720) circle (1.2pt);
\fill[blue!80!black] (7.6273,5.4717) circle (1.2pt);
\fill[blue!80!black] (7.6268,5.4716) circle (1.2pt);
\fill[blue!80!black] (7.6271,5.4711) circle (1.2pt);
\fill[blue!80!black] (7.6265,5.4710) circle (1.2pt);
\fill[blue!80!black] (7.6269,5.4705) circle (1.2pt);
\fill[blue!80!black] (7.6261,5.4705) circle (1.2pt);
\fill[blue!80!black] (7.6265,5.4702) circle (1.2pt);
\fill[blue!80!black] (7.6257,5.4700) circle (1.2pt);
\fill[blue!80!black] (7.6263,5.4697) circle (1.2pt);
\fill[blue!80!black] (7.6257,5.4696) circle (1.2pt);
\fill[blue!80!black] (7.6261,5.4691) circle (1.2pt);
\fill[blue!80!black] (7.6254,5.4691) circle (1.2pt);
\fill[blue!80!black] (7.6256,5.4688) circle (1.2pt);
\fill[blue!80!black] (7.6250,5.4685) circle (1.2pt);
\fill[blue!80!black] (7.6254,5.4682) circle (1.2pt);
\fill[blue!80!black] (7.6250,5.4681) circle (1.2pt);
\fill[blue!80!black] (7.6252,5.4677) circle (1.2pt);
\fill[blue!80!black] (7.6246,5.4676) circle (1.2pt);
\fill[blue!80!black] (7.6250,5.4671) circle (1.2pt);
\fill[blue!80!black] (7.6242,5.4671) circle (1.2pt);
\fill[blue!80!black] (7.6246,5.4668) circle (1.2pt);
\fill[blue!80!black] (7.6238,5.4666) circle (1.2pt);
\fill[blue!80!black] (7.6243,5.4663) circle (1.2pt);
\fill[blue!80!black] (7.6239,5.4662) circle (1.2pt);
\fill[blue!80!black] (7.6241,5.4657) circle (1.2pt);
\fill[blue!80!black] (7.6235,5.4656) circle (1.2pt);
\fill[blue!80!black] (7.6237,5.4654) circle (1.2pt);
\fill[blue!80!black] (7.6231,5.4651) circle (1.2pt);
\fill[blue!80!black] (7.6235,5.4648) circle (1.2pt);
\fill[blue!80!black] (7.6227,5.4646) circle (1.2pt);
\fill[blue!80!black] (7.6233,5.4643) circle (1.2pt);
\fill[blue!80!black] (7.6227,5.4642) circle (1.2pt);
\fill[blue!80!black] (7.6231,5.4637) circle (1.2pt);
\fill[blue!80!black] (7.6224,5.4637) circle (1.2pt);
\fill[blue!80!black] (7.6227,5.4634) circle (1.2pt);
\fill[blue!80!black] (7.6220,5.4632) circle (1.2pt);
\fill[blue!80!black] (7.6224,5.4629) circle (1.2pt);
\fill[blue!80!black] (7.6216,5.4627) circle (1.2pt);
\fill[blue!80!black] (7.6222,5.4623) circle (1.2pt);
\fill[blue!80!black] (7.6216,5.4622) circle (1.2pt);
\fill[blue!80!black] (7.6220,5.4617) circle (1.2pt);
\fill[blue!80!black] (7.6213,5.4617) circle (1.2pt);
\fill[blue!80!black] (7.6216,5.4614) circle (1.2pt);
\fill[blue!80!black] (7.6209,5.4612) circle (1.2pt);
\fill[blue!80!black] (7.6213,5.4609) circle (1.2pt);
\fill[blue!80!black] (7.6209,5.4608) circle (1.2pt);
\fill[blue!80!black] (7.6211,5.4603) circle (1.2pt);
\fill[blue!80!black] (7.6205,5.4603) circle (1.2pt);
\fill[blue!80!black] (7.6210,5.4598) circle (1.2pt);
\fill[blue!80!black] (7.6202,5.4598) circle (1.2pt);
\fill[blue!80!black] (7.6205,5.4595) circle (1.2pt);
\fill[blue!80!black] (7.6198,5.4592) circle (1.2pt);
\fill[blue!80!black] (7.6203,5.4589) circle (1.2pt);
\fill[blue!80!black] (7.6198,5.4588) circle (1.2pt);
\fill[blue!80!black] (7.6201,5.4584) circle (1.2pt);
\fill[blue!80!black] (7.6194,5.4583) circle (1.2pt);
\fill[blue!80!black] (7.6199,5.4578) circle (1.2pt);
\fill[blue!80!black] (7.6191,5.4578) circle (1.2pt);
\fill[blue!80!black] (7.6194,5.4575) circle (1.2pt);
\fill[blue!80!black] (7.6187,5.4573) circle (1.2pt);
\fill[blue!80!black] (7.6192,5.4570) circle (1.2pt);
\fill[blue!80!black] (7.6187,5.4569) circle (1.2pt);
\fill[blue!80!black] (7.6190,5.4564) circle (1.2pt);
\fill[blue!80!black] (7.6184,5.4564) circle (1.2pt);
\fill[blue!80!black] (7.6188,5.4558) circle (1.2pt);
\fill[blue!80!black] (7.6180,5.4558) circle (1.2pt);
\fill[blue!80!black] (7.6183,5.4556) circle (1.2pt);
\fill[blue!80!black] (7.6176,5.4553) circle (1.2pt);
\fill[blue!80!black] (7.6181,5.4550) circle (1.2pt);
\fill[blue!80!black] (7.6176,5.4549) circle (1.2pt);
\fill[blue!80!black] (7.6179,5.4545) circle (1.2pt);
\fill[blue!80!black] (7.6173,5.4544) circle (1.2pt);
\fill[blue!80!black] (7.6177,5.4539) circle (1.2pt);
\fill[blue!80!black] (7.6169,5.4539) circle (1.2pt);
\fill[blue!80!black] (7.6172,5.4536) circle (1.2pt);
\fill[blue!80!black] (7.6165,5.4534) circle (1.2pt);
\fill[blue!80!black] (7.6170,5.4531) circle (1.2pt);
\fill[blue!80!black] (7.6165,5.4530) circle (1.2pt);
\fill[blue!80!black] (7.6168,5.4525) circle (1.2pt);
\fill[blue!80!black] (7.6162,5.4525) circle (1.2pt);
\fill[blue!80!black] (7.6166,5.4519) circle (1.2pt);
\fill[blue!80!black] (7.6158,5.4519) circle (1.2pt);
\fill[blue!80!black] (7.6162,5.4517) circle (1.2pt);
\fill[blue!80!black] (7.6154,5.4514) circle (1.2pt);
\fill[blue!80!black] (7.6159,5.4511) circle (1.2pt);
\fill[blue!80!black] (7.6155,5.4510) circle (1.2pt);
\fill[blue!80!black] (7.6157,5.4506) circle (1.2pt);
\fill[blue!80!black] (7.6151,5.4505) circle (1.2pt);
\fill[blue!80!black] (7.6155,5.4500) circle (1.2pt);
\fill[blue!80!black] (7.6148,5.4500) circle (1.2pt);
\fill[blue!80!black] (7.6151,5.4497) circle (1.2pt);
\fill[blue!80!black] (7.6144,5.4495) circle (1.2pt);
\fill[blue!80!black] (7.6148,5.4492) circle (1.2pt);
\fill[blue!80!black] (7.6144,5.4491) circle (1.2pt);
\fill[blue!80!black] (7.6146,5.4486) circle (1.2pt);
\fill[blue!80!black] (7.6140,5.4486) circle (1.2pt);
\fill[blue!80!black] (7.6144,5.4481) circle (1.2pt);
\fill[blue!80!black] (7.6137,5.4481) circle (1.2pt);
\fill[blue!80!black] (7.6140,5.4478) circle (1.2pt);
\fill[blue!80!black] (7.6133,5.4475) circle (1.2pt);
\fill[blue!80!black] (7.6137,5.4472) circle (1.2pt);
\fill[blue!80!black] (7.6129,5.4470) circle (1.2pt);
\fill[blue!80!black] (7.6135,5.4467) circle (1.2pt);
\fill[blue!80!black] (7.6130,5.4466) circle (1.2pt);
\fill[blue!80!black] (7.6133,5.4461) circle (1.2pt);
\fill[blue!80!black] (7.6126,5.4461) circle (1.2pt);
\fill[blue!80!black] (7.6129,5.4458) circle (1.2pt);
\fill[blue!80!black] (7.6123,5.4456) circle (1.2pt);
\fill[blue!80!black] (7.6126,5.4453) circle (1.2pt);
\fill[blue!80!black] (7.6119,5.4451) circle (1.2pt);
\fill[blue!80!black] (7.6124,5.4448) circle (1.2pt);
\fill[blue!80!black] (7.6119,5.4447) circle (1.2pt);
\fill[blue!80!black] (7.6122,5.4442) circle (1.2pt);
\fill[blue!80!black] (7.6116,5.4442) circle (1.2pt);
\fill[blue!80!black] (7.6120,5.4437) circle (1.2pt);
\fill[blue!80!black] (7.6112,5.4437) circle (1.2pt);
\fill[blue!80!black] (7.6115,5.4434) circle (1.2pt);
\fill[blue!80!black] (7.6108,5.4432) circle (1.2pt);
\fill[blue!80!black] (7.6113,5.4428) circle (1.2pt);
\fill[blue!80!black] (7.6109,5.4428) circle (1.2pt);
\fill[blue!80!black] (7.6111,5.4423) circle (1.2pt);
\fill[blue!80!black] (7.6105,5.4422) circle (1.2pt);
\fill[blue!80!black] (7.6109,5.4418) circle (1.2pt);
\fill[blue!80!black] (7.6102,5.4417) circle (1.2pt);
\fill[blue!80!black] (7.6104,5.4415) circle (1.2pt);
\fill[blue!80!black] (7.6098,5.4412) circle (1.2pt);
\fill[blue!80!black] (7.6102,5.4409) circle (1.2pt);
\fill[blue!80!black] (7.6094,5.4407) circle (1.2pt);
\fill[blue!80!black] (7.6100,5.4404) circle (1.2pt);
\fill[blue!80!black] (7.6095,5.4403) circle (1.2pt);
\fill[blue!80!black] (7.6098,5.4398) circle (1.2pt);
\fill[blue!80!black] (7.6091,5.4398) circle (1.2pt);
\fill[blue!80!black] (7.6096,5.4393) circle (1.2pt);
\fill[blue!80!black] (7.6088,5.4393) circle (1.2pt);
\fill[blue!80!black] (7.6091,5.4390) circle (1.2pt);
\fill[blue!80!black] (7.6084,5.4388) circle (1.2pt);
\fill[blue!80!black] (7.6089,5.4385) circle (1.2pt);
\fill[blue!80!black] (7.6080,5.4383) circle (1.2pt);
\fill[blue!80!black] (7.6086,5.4379) circle (1.2pt);
\fill[blue!80!black] (7.6081,5.4379) circle (1.2pt);
\fill[blue!80!black] (7.6084,5.4374) circle (1.2pt);
\fill[blue!80!black] (7.6078,5.4374) circle (1.2pt);
\fill[blue!80!black] (7.6080,5.4371) circle (1.2pt);
\fill[blue!80!black] (7.6074,5.4369) circle (1.2pt);
\fill[blue!80!black] (7.6077,5.4366) circle (1.2pt);
\fill[blue!80!black] (7.6070,5.4364) circle (1.2pt);
\fill[blue!80!black] (7.6075,5.4360) circle (1.2pt);
\fill[blue!80!black] (7.6071,5.4360) circle (1.2pt);
\fill[blue!80!black] (7.6073,5.4355) circle (1.2pt);
\fill[blue!80!black] (7.6067,5.4354) circle (1.2pt);
\fill[blue!80!black] (7.6071,5.4350) circle (1.2pt);
\fill[blue!80!black] (7.6064,5.4349) circle (1.2pt);
\fill[blue!80!black] (7.6066,5.4347) circle (1.2pt);
\fill[blue!80!black] (7.6060,5.4344) circle (1.2pt);
\fill[blue!80!black] (7.6064,5.4341) circle (1.2pt);
\fill[blue!80!black] (7.6056,5.4339) circle (1.2pt);
\fill[blue!80!black] (7.6062,5.4336) circle (1.2pt);
\fill[blue!80!black] (7.6057,5.4335) circle (1.2pt);
\fill[blue!80!black] (7.6059,5.4331) circle (1.2pt);
\fill[blue!80!black] (7.6054,5.4330) circle (1.2pt);
\fill[blue!80!black] (7.6057,5.4325) circle (1.2pt);
\fill[blue!80!black] (7.6050,5.4325) circle (1.2pt);
\fill[blue!80!black] (7.6053,5.4322) circle (1.2pt);
\fill[blue!80!black] (7.6046,5.4320) circle (1.2pt);
\fill[blue!80!black] (7.6050,5.4317) circle (1.2pt);
\fill[blue!80!black] (7.6043,5.4315) circle (1.2pt);
\fill[blue!80!black] (7.6048,5.4312) circle (1.2pt);
\fill[blue!80!black] (7.6043,5.4311) circle (1.2pt);
\fill[blue!80!black] (7.6046,5.4306) circle (1.2pt);
\fill[blue!80!black] (7.6040,5.4306) circle (1.2pt);
\fill[blue!80!black] (7.6044,5.4301) circle (1.2pt);
\fill[blue!80!black] (7.6037,5.4301) circle (1.2pt);
\fill[blue!80!black] (7.6039,5.4298) circle (1.2pt);
\fill[blue!80!black] (7.6033,5.4296) circle (1.2pt);
\fill[blue!80!black] (7.6037,5.4293) circle (1.2pt);
\fill[blue!80!black] (7.6029,5.4291) circle (1.2pt);
\fill[blue!80!black] (7.6034,5.4288) circle (1.2pt);
\fill[blue!80!black] (7.6030,5.4287) circle (1.2pt);
\fill[blue!80!black] (7.6032,5.4282) circle (1.2pt);
\fill[blue!80!black] (7.6026,5.4282) circle (1.2pt);
\fill[blue!80!black] (7.6030,5.4277) circle (1.2pt);
\fill[blue!80!black] (7.6023,5.4277) circle (1.2pt);
\fill[blue!80!black] (7.6026,5.4274) circle (1.2pt);
\fill[blue!80!black] (7.6020,5.4272) circle (1.2pt);
\fill[blue!80!black] (7.6023,5.4269) circle (1.2pt);
\fill[blue!80!black] (7.6016,5.4267) circle (1.2pt);
\fill[blue!80!black] (7.6021,5.4264) circle (1.2pt);
\fill[blue!80!black] (7.6016,5.4263) circle (1.2pt);
\fill[blue!80!black] (7.6019,5.4258) circle (1.2pt);
\fill[blue!80!black] (7.6013,5.4258) circle (1.2pt);
\fill[blue!80!black] (7.6017,5.4253) circle (1.2pt);
\fill[blue!80!black] (7.6010,5.4253) circle (1.2pt);
\fill[blue!80!black] (7.6012,5.4250) circle (1.2pt);
\fill[blue!80!black] (7.6006,5.4248) circle (1.2pt);
\fill[blue!80!black] (7.6010,5.4245) circle (1.2pt);
\fill[blue!80!black] (7.6003,5.4243) circle (1.2pt);
\fill[blue!80!black] (7.6007,5.4240) circle (1.2pt);
\fill[blue!80!black] (7.5999,5.4238) circle (1.2pt);
\fill[blue!80!black] (7.6005,5.4234) circle (1.2pt);
\fill[blue!80!black] (7.6000,5.4234) circle (1.2pt);
\fill[blue!80!black] (7.6003,5.4229) circle (1.2pt);
\fill[blue!80!black] (7.5996,5.4229) circle (1.2pt);
\fill[blue!80!black] (7.6001,5.4224) circle (1.2pt);
\fill[blue!80!black] (7.5993,5.4224) circle (1.2pt);
\fill[blue!80!black] (7.5996,5.4221) circle (1.2pt);
\fill[blue!80!black] (7.5989,5.4219) circle (1.2pt);
\fill[blue!80!black] (7.5993,5.4216) circle (1.2pt);
\fill[blue!80!black] (7.5986,5.4214) circle (1.2pt);
\fill[blue!80!black] (7.5991,5.4210) circle (1.2pt);
\fill[blue!80!black] (7.5986,5.4210) circle (1.2pt);
\fill[blue!80!black] (7.5989,5.4205) circle (1.2pt);
\fill[blue!80!black] (7.5983,5.4205) circle (1.2pt);
\fill[blue!80!black] (7.5987,5.4200) circle (1.2pt);
\fill[blue!80!black] (7.5980,5.4200) circle (1.2pt);
\fill[blue!80!black] (7.5982,5.4197) circle (1.2pt);
\fill[blue!80!black] (7.5976,5.4195) circle (1.2pt);
\fill[blue!80!black] (7.5980,5.4192) circle (1.2pt);
\fill[blue!80!black] (7.5973,5.4190) circle (1.2pt);
\fill[blue!80!black] (7.5977,5.4187) circle (1.2pt);
\fill[blue!80!black] (7.5969,5.4185) circle (1.2pt);
\fill[blue!80!black] (7.5975,5.4181) circle (1.2pt);
\fill[blue!80!black] (7.5970,5.4181) circle (1.2pt);
\fill[blue!80!black] (7.5973,5.4176) circle (1.2pt);
\fill[blue!80!black] (7.5967,5.4176) circle (1.2pt);
\fill[blue!80!black] (7.5971,5.4171) circle (1.2pt);
\fill[blue!80!black] (7.5963,5.4171) circle (1.2pt);
\fill[blue!80!black] (7.5966,5.4168) circle (1.2pt);
\fill[blue!80!black] (7.5960,5.4166) circle (1.2pt);
\fill[blue!80!black] (7.5964,5.4163) circle (1.2pt);
\fill[blue!80!black] (7.5956,5.4161) circle (1.2pt);
\fill[blue!80!black] (7.5961,5.4158) circle (1.2pt);
\fill[blue!80!black] (7.5953,5.4156) circle (1.2pt);
\fill[blue!80!black] (7.5959,5.4152) circle (1.2pt);
\fill[blue!80!black] (7.5954,5.4152) circle (1.2pt);
\fill[blue!80!black] (7.5957,5.4147) circle (1.2pt);
\fill[blue!80!black] (7.5950,5.4147) circle (1.2pt);
\fill[blue!80!black] (7.5954,5.4142) circle (1.2pt);
\fill[blue!80!black] (7.5947,5.4142) circle (1.2pt);
\fill[blue!80!black] (7.5950,5.4139) circle (1.2pt);
\fill[blue!80!black] (7.5944,5.4137) circle (1.2pt);
\fill[blue!80!black] (7.5947,5.4134) circle (1.2pt);
\fill[blue!80!black] (7.5940,5.4132) circle (1.2pt);
\fill[blue!80!black] (7.5945,5.4129) circle (1.2pt);
\fill[blue!80!black] (7.5936,5.4127) circle (1.2pt);
\fill[blue!80!black] (7.5942,5.4124) circle (1.2pt);
\fill[blue!80!black] (7.5938,5.4123) circle (1.2pt);
\fill[blue!80!black] (7.5940,5.4118) circle (1.2pt);
\fill[blue!80!black] (7.5934,5.4118) circle (1.2pt);
\fill[blue!80!black] (7.5938,5.4113) circle (1.2pt);
\fill[blue!80!black] (7.5931,5.4113) circle (1.2pt);
\fill[blue!80!black] (7.5934,5.4110) circle (1.2pt);
\fill[blue!80!black] (7.5928,5.4108) circle (1.2pt);
\fill[blue!80!black] (7.5931,5.4105) circle (1.2pt);
\fill[blue!80!black] (7.5924,5.4103) circle (1.2pt);
\fill[blue!80!black] (7.5929,5.4100) circle (1.2pt);
\fill[blue!80!black] (7.5921,5.4098) circle (1.2pt);
\fill[blue!80!black] (7.5926,5.4095) circle (1.2pt);
\fill[blue!80!black] (7.5921,5.4094) circle (1.2pt);
\fill[blue!80!black] (7.5924,5.4090) circle (1.2pt);
\fill[blue!80!black] (7.5918,5.4089) circle (1.2pt);
\fill[blue!80!black] (7.5922,5.4085) circle (1.2pt);
\fill[blue!80!black] (7.5915,5.4084) circle (1.2pt);
\fill[blue!80!black] (7.5920,5.4079) circle (1.2pt);
\fill[blue!80!black] (7.5912,5.4079) circle (1.2pt);
\fill[blue!80!black] (7.5915,5.4077) circle (1.2pt);
\fill[blue!80!black] (7.5908,5.4074) circle (1.2pt);
\fill[blue!80!black] (7.5912,5.4072) circle (1.2pt);
\fill[blue!80!black] (7.5905,5.4070) circle (1.2pt);
\fill[blue!80!black] (7.5910,5.4066) circle (1.2pt);
\fill[blue!80!black] (7.5901,5.4065) circle (1.2pt);
\fill[blue!80!black] (7.5907,5.4061) circle (1.2pt);
\fill[blue!80!black] (7.5902,5.4061) circle (1.2pt);
\fill[blue!80!black] (7.5905,5.4056) circle (1.2pt);
\fill[blue!80!black] (7.5899,5.4056) circle (1.2pt);
\fill[blue!80!black] (7.5903,5.4051) circle (1.2pt);
\fill[blue!80!black] (7.5896,5.4051) circle (1.2pt);
\fill[blue!80!black] (7.5898,5.4048) circle (1.2pt);
\fill[blue!80!black] (7.5893,5.4046) circle (1.2pt);
\fill[blue!80!black] (7.5896,5.4043) circle (1.2pt);
\fill[blue!80!black] (7.5889,5.4041) circle (1.2pt);
\fill[blue!80!black] (7.5893,5.4038) circle (1.2pt);
\fill[blue!80!black] (7.5886,5.4036) circle (1.2pt);
\fill[blue!80!black] (7.5891,5.4033) circle (1.2pt);
\fill[blue!80!black] (7.5886,5.4032) circle (1.2pt);
\fill[blue!80!black] (7.5888,5.4028) circle (1.2pt);
\fill[blue!80!black] (7.5883,5.4027) circle (1.2pt);
\fill[blue!80!black] (7.5886,5.4023) circle (1.2pt);
\fill[blue!80!black] (7.5880,5.4022) circle (1.2pt);
\fill[blue!80!black] (7.5884,5.4017) circle (1.2pt);
\fill[blue!80!black] (7.5877,5.4017) circle (1.2pt);
\fill[blue!80!black] (7.5879,5.4015) circle (1.2pt);
\fill[blue!80!black] (7.5874,5.4012) circle (1.2pt);
\fill[blue!80!black] (7.5877,5.4010) circle (1.2pt);
\fill[blue!80!black] (7.5870,5.4007) circle (1.2pt);
\fill[blue!80!black] (7.5874,5.4005) circle (1.2pt);
\fill[blue!80!black] (7.5867,5.4003) circle (1.2pt);
\fill[blue!80!black] (7.5872,5.3999) circle (1.2pt);
\fill[blue!80!black] (7.5863,5.3998) circle (1.2pt);
\fill[blue!80!black] (7.5869,5.3994) circle (1.2pt);
\fill[blue!80!black] (7.5864,5.3994) circle (1.2pt);
\fill[blue!80!black] (7.5867,5.3989) circle (1.2pt);
\fill[blue!80!black] (7.5861,5.3989) circle (1.2pt);
\fill[blue!80!black] (7.5865,5.3984) circle (1.2pt);
\fill[blue!80!black] (7.5858,5.3984) circle (1.2pt);
\fill[blue!80!black] (7.5863,5.3979) circle (1.2pt);
\fill[blue!80!black] (7.5855,5.3979) circle (1.2pt);
\fill[blue!80!black] (7.5858,5.3976) circle (1.2pt);
\fill[blue!80!black] (7.5852,5.3974) circle (1.2pt);
\fill[blue!80!black] (7.5855,5.3971) circle (1.2pt);
\fill[blue!80!black] (7.5848,5.3969) circle (1.2pt);
\fill[blue!80!black] (7.5853,5.3966) circle (1.2pt);
\fill[blue!80!black] (7.5845,5.3964) circle (1.2pt);
\fill[blue!80!black] (7.5850,5.3961) circle (1.2pt);
\fill[blue!80!black] (7.5846,5.3960) circle (1.2pt);
\fill[blue!80!black] (7.5848,5.3956) circle (1.2pt);
\fill[blue!80!black] (7.5843,5.3955) circle (1.2pt);
\fill[blue!80!black] (7.5845,5.3951) circle (1.2pt);
\fill[blue!80!black] (7.5840,5.3951) circle (1.2pt);
\fill[blue!80!black] (7.5843,5.3946) circle (1.2pt);
\fill[blue!80!black] (7.5836,5.3946) circle (1.2pt);
\fill[blue!80!black] (7.5841,5.3941) circle (1.2pt);
\fill[blue!80!black] (7.5833,5.3941) circle (1.2pt);
\fill[blue!80!black] (7.5836,5.3938) circle (1.2pt);
\fill[blue!80!black] (7.5830,5.3936) circle (1.2pt);
\fill[blue!80!black] (7.5834,5.3933) circle (1.2pt);
\fill[blue!80!black] (7.5827,5.3931) circle (1.2pt);
\fill[blue!80!black] (7.5831,5.3928) circle (1.2pt);
\fill[blue!80!black] (7.5823,5.3926) circle (1.2pt);
\fill[blue!80!black] (7.5829,5.3923) circle (1.2pt);
\fill[blue!80!black] (7.5824,5.3922) circle (1.2pt);
\fill[blue!80!black] (7.5826,5.3918) circle (1.2pt);
\fill[blue!80!black] (7.5821,5.3917) circle (1.2pt);
\fill[blue!80!black] (7.5824,5.3913) circle (1.2pt);
\fill[blue!80!black] (7.5818,5.3913) circle (1.2pt);
\fill[blue!80!black] (7.5821,5.3908) circle (1.2pt);
\fill[blue!80!black] (7.5815,5.3908) circle (1.2pt);
\fill[blue!80!black] (7.5819,5.3903) circle (1.2pt);
\fill[blue!80!black] (7.5812,5.3903) circle (1.2pt);
\fill[blue!80!black] (7.5814,5.3900) circle (1.2pt);
\fill[blue!80!black] (7.5809,5.3898) circle (1.2pt);
\fill[blue!80!black] (7.5812,5.3895) circle (1.2pt);
\fill[blue!80!black] (7.5805,5.3893) circle (1.2pt);
\fill[blue!80!black] (7.5809,5.3890) circle (1.2pt);
\fill[blue!80!black] (7.5802,5.3888) circle (1.2pt);
\fill[blue!80!black] (7.5807,5.3885) circle (1.2pt);
\fill[blue!80!black] (7.5799,5.3884) circle (1.2pt);
\fill[blue!80!black] (7.5804,5.3880) circle (1.2pt);
\fill[blue!80!black] (7.5800,5.3879) circle (1.2pt);
\fill[blue!80!black] (7.5802,5.3875) circle (1.2pt);
\fill[blue!80!black] (7.5797,5.3875) circle (1.2pt);
\fill[blue!80!black] (7.5800,5.3870) circle (1.2pt);
\fill[blue!80!black] (7.5794,5.3870) circle (1.2pt);
\fill[blue!80!black] (7.5797,5.3865) circle (1.2pt);
\fill[blue!80!black] (7.5791,5.3865) circle (1.2pt);
\fill[blue!80!black] (7.5795,5.3860) circle (1.2pt);
\fill[blue!80!black] (7.5788,5.3860) circle (1.2pt);
\fill[blue!80!black] (7.5790,5.3858) circle (1.2pt);
\fill[blue!80!black] (7.5784,5.3855) circle (1.2pt);
\fill[blue!80!black] (7.5788,5.3853) circle (1.2pt);
\fill[blue!80!black] (7.5781,5.3850) circle (1.2pt);
\fill[blue!80!black] (7.5785,5.3848) circle (1.2pt);
\fill[blue!80!black] (7.5778,5.3846) circle (1.2pt);
\fill[blue!80!black] (7.5782,5.3843) circle (1.2pt);
\fill[blue!80!black] (7.5775,5.3841) circle (1.2pt);
\fill[blue!80!black] (7.5780,5.3838) circle (1.2pt);
\fill[blue!80!black] (7.5776,5.3837) circle (1.2pt);
\fill[blue!80!black] (7.5777,5.3833) circle (1.2pt);
\fill[blue!80!black] (7.5773,5.3832) circle (1.2pt);
\fill[blue!80!black] (7.5775,5.3828) circle (1.2pt);
\fill[blue!80!black] (7.5770,5.3827) circle (1.2pt);
\fill[blue!80!black] (7.5773,5.3823) circle (1.2pt);
\fill[blue!80!black] (7.5767,5.3822) circle (1.2pt);
\fill[blue!80!black] (7.5770,5.3818) circle (1.2pt);
\fill[blue!80!black] (7.5764,5.3817) circle (1.2pt);
\fill[blue!80!black] (7.5768,5.3813) circle (1.2pt);
\fill[blue!80!black] (7.5761,5.3813) circle (1.2pt);
\fill[blue!80!black] (7.5763,5.3810) circle (1.2pt);
\fill[blue!80!black] (7.5757,5.3808) circle (1.2pt);
\fill[blue!80!black] (7.5761,5.3805) circle (1.2pt);
\fill[blue!80!black] (7.5754,5.3803) circle (1.2pt);
\fill[blue!80!black] (7.5758,5.3800) circle (1.2pt);
\fill[blue!80!black] (7.5751,5.3798) circle (1.2pt);
\fill[blue!80!black] (7.5755,5.3795) circle (1.2pt);
\fill[blue!80!black] (7.5748,5.3794) circle (1.2pt);
\fill[blue!80!black] (7.5753,5.3790) circle (1.2pt);
\fill[blue!80!black] (7.5745,5.3789) circle (1.2pt);
\fill[blue!80!black] (7.5750,5.3785) circle (1.2pt);
\fill[blue!80!black] (7.5746,5.3785) circle (1.2pt);
\fill[blue!80!black] (7.5748,5.3780) circle (1.2pt);
\fill[blue!80!black] (7.5743,5.3780) circle (1.2pt);
\fill[blue!80!black] (7.5745,5.3776) circle (1.2pt);
\fill[blue!80!black] (7.5740,5.3775) circle (1.2pt);
\fill[blue!80!black] (7.5743,5.3771) circle (1.2pt);
\fill[blue!80!black] (7.5737,5.3770) circle (1.2pt);
\fill[blue!80!black] (7.5741,5.3766) circle (1.2pt);
\fill[blue!80!black] (7.5734,5.3765) circle (1.2pt);
\fill[blue!80!black] (7.5738,5.3761) circle (1.2pt);
\fill[blue!80!black] (7.5731,5.3761) circle (1.2pt);
\fill[blue!80!black] (7.5734,5.3758) circle (1.2pt);
\fill[blue!80!black] (7.5728,5.3756) circle (1.2pt);
\fill[blue!80!black] (7.5731,5.3753) circle (1.2pt);
\fill[blue!80!black] (7.5725,5.3751) circle (1.2pt);
\fill[blue!80!black] (7.5728,5.3748) circle (1.2pt);
\fill[blue!80!black] (7.5722,5.3746) circle (1.2pt);
\fill[blue!80!black] (7.5726,5.3743) circle (1.2pt);
\fill[blue!80!black] (7.5718,5.3742) circle (1.2pt);
\fill[blue!80!black] (7.5723,5.3739) circle (1.2pt);
\fill[blue!80!black] (7.5715,5.3737) circle (1.2pt);
\fill[blue!80!black] (7.5721,5.3734) circle (1.2pt);
\fill[blue!80!black] (7.5716,5.3733) circle (1.2pt);
\fill[blue!80!black] (7.5718,5.3729) circle (1.2pt);
\fill[blue!80!black] (7.5713,5.3728) circle (1.2pt);
\fill[blue!80!black] (7.5716,5.3724) circle (1.2pt);
\fill[blue!80!black] (7.5710,5.3723) circle (1.2pt);
\fill[blue!80!black] (7.5713,5.3719) circle (1.2pt);
\fill[blue!80!black] (7.5708,5.3719) circle (1.2pt);
\fill[blue!80!black] (7.5711,5.3714) circle (1.2pt);
\fill[blue!80!black] (7.5705,5.3714) circle (1.2pt);
\fill[blue!80!black] (7.5708,5.3709) circle (1.2pt);
\fill[blue!80!black] (7.5702,5.3709) circle (1.2pt);
\fill[blue!80!black] (7.5706,5.3704) circle (1.2pt);
\fill[blue!80!black] (7.5699,5.3704) circle (1.2pt);
\fill[blue!80!black] (7.5701,5.3702) circle (1.2pt);
\fill[blue!80!black] (7.5696,5.3699) circle (1.2pt);
\fill[blue!80!black] (7.5698,5.3697) circle (1.2pt);
\fill[blue!80!black] (7.5693,5.3695) circle (1.2pt);
\fill[blue!80!black] (7.5696,5.3692) circle (1.2pt);
\fill[blue!80!black] (7.5690,5.3690) circle (1.2pt);
\fill[blue!80!black] (7.5693,5.3687) circle (1.2pt);
\fill[blue!80!black] (7.5687,5.3685) circle (1.2pt);
\fill[blue!80!black] (7.5691,5.3682) circle (1.2pt);
\fill[blue!80!black] (7.5684,5.3680) circle (1.2pt);
\fill[blue!80!black] (7.5688,5.3677) circle (1.2pt);
\fill[blue!80!black] (7.5680,5.3676) circle (1.2pt);
\fill[blue!80!black] (7.5685,5.3673) circle (1.2pt);
\fill[blue!80!black] (7.5677,5.3671) circle (1.2pt);
\fill[blue!80!black] (7.5683,5.3668) circle (1.2pt);
\fill[blue!80!black] (7.5679,5.3667) circle (1.2pt);
\fill[blue!80!black] (7.5680,5.3663) circle (1.2pt);
\fill[blue!80!black] (7.5676,5.3662) circle (1.2pt);
\fill[blue!80!black] (7.5678,5.3658) circle (1.2pt);
\fill[blue!80!black] (7.5673,5.3658) circle (1.2pt);
\fill[blue!80!black] (7.5675,5.3653) circle (1.2pt);
\fill[blue!80!black] (7.5670,5.3653) circle (1.2pt);
\fill[blue!80!black] (7.5673,5.3648) circle (1.2pt);
\fill[blue!80!black] (7.5667,5.3648) circle (1.2pt);
\fill[blue!80!black] (7.5670,5.3643) circle (1.2pt);
\fill[blue!80!black] (7.5664,5.3643) circle (1.2pt);
\fill[blue!80!black] (7.5668,5.3639) circle (1.2pt);
\fill[blue!80!black] (7.5661,5.3639) circle (1.2pt);
\fill[blue!80!black] (7.5666,5.3634) circle (1.2pt);
\fill[blue!80!black] (7.5658,5.3634) circle (1.2pt);
\fill[blue!80!black] (7.5661,5.3631) circle (1.2pt);
\fill[blue!80!black] (7.5655,5.3629) circle (1.2pt);
\fill[blue!80!black] (7.5658,5.3627) circle (1.2pt);
\fill[blue!80!black] (7.5652,5.3624) circle (1.2pt);
\fill[blue!80!black] (7.5655,5.3622) circle (1.2pt);
\fill[blue!80!black] (7.5650,5.3620) circle (1.2pt);
\fill[blue!80!black] (7.5653,5.3617) circle (1.2pt);
\fill[blue!80!black] (7.5647,5.3615) circle (1.2pt);
\fill[blue!80!black] (7.5650,5.3612) circle (1.2pt);
\fill[blue!80!black] (7.5644,5.3610) circle (1.2pt);
\fill[blue!80!black] (7.5648,5.3607) circle (1.2pt);
\fill[blue!80!black] (7.5641,5.3605) circle (1.2pt);
\fill[blue!80!black] (7.5645,5.3603) circle (1.2pt);
\fill[blue!80!black] (7.5638,5.3601) circle (1.2pt);
\fill[blue!80!black] (7.5642,5.3598) circle (1.2pt);
\fill[blue!80!black] (7.5635,5.3596) circle (1.2pt);
\fill[blue!80!black] (7.5640,5.3593) circle (1.2pt);
\fill[blue!80!black] (7.5632,5.3591) circle (1.2pt);
\fill[blue!80!black] (7.5637,5.3588) circle (1.2pt);
\fill[blue!80!black] (7.5633,5.3587) circle (1.2pt);
\fill[blue!80!black] (7.5635,5.3583) circle (1.2pt);
\fill[blue!80!black] (7.5630,5.3583) circle (1.2pt);
\fill[blue!80!black] (7.5632,5.3579) circle (1.2pt);
\fill[blue!80!black] (7.5627,5.3578) circle (1.2pt);
\fill[blue!80!black] (7.5629,5.3574) circle (1.2pt);
\fill[blue!80!black] (7.5624,5.3573) circle (1.2pt);
\fill[blue!80!black] (7.5627,5.3569) circle (1.2pt);
\fill[blue!80!black] (7.5622,5.3569) circle (1.2pt);
\fill[blue!80!black] (7.5624,5.3564) circle (1.2pt);
\fill[blue!80!black] (7.5619,5.3564) circle (1.2pt);
\fill[blue!80!black] (7.5622,5.3559) circle (1.2pt);
\fill[blue!80!black] (7.5616,5.3559) circle (1.2pt);
\fill[blue!80!black] (7.5619,5.3555) circle (1.2pt);
\fill[blue!80!black] (7.5613,5.3554) circle (1.2pt);
\fill[blue!80!black] (7.5617,5.3550) circle (1.2pt);
\fill[blue!80!black] (7.5610,5.3550) circle (1.2pt);
\fill[blue!80!black] (7.5614,5.3545) circle (1.2pt);
\fill[blue!80!black] (7.5607,5.3545) circle (1.2pt);
\fill[blue!80!black] (7.5612,5.3540) circle (1.2pt);
\fill[blue!80!black] (7.5605,5.3540) circle (1.2pt);
\fill[blue!80!black] (7.5609,5.3536) circle (1.2pt);
\fill[blue!80!black] (7.5602,5.3536) circle (1.2pt);
\fill[blue!80!black] (7.5604,5.3533) circle (1.2pt);
\fill[blue!80!black] (7.5599,5.3531) circle (1.2pt);
\fill[blue!80!black] (7.5602,5.3528) circle (1.2pt);
\fill[blue!80!black] (7.5596,5.3526) circle (1.2pt);
\fill[blue!80!black] (7.5599,5.3524) circle (1.2pt);
\fill[blue!80!black] (7.5593,5.3522) circle (1.2pt);
\fill[blue!80!black] (7.5596,5.3519) circle (1.2pt);
\fill[blue!80!black] (7.5590,5.3517) circle (1.2pt);
\fill[blue!80!black] (7.5594,5.3514) circle (1.2pt);
\fill[blue!80!black] (7.5588,5.3512) circle (1.2pt);
\fill[blue!80!black] (7.5591,5.3510) circle (1.2pt);
\fill[blue!80!black] (7.5585,5.3508) circle (1.2pt);
\fill[blue!80!black] (7.5588,5.3505) circle (1.2pt);
\fill[blue!80!black] (7.5582,5.3503) circle (1.2pt);
\fill[blue!80!black] (7.5586,5.3500) circle (1.2pt);
\fill[blue!80!black] (7.5579,5.3498) circle (1.2pt);
\fill[blue!80!black] (7.5583,5.3495) circle (1.2pt);
\fill[blue!80!black] (7.5576,5.3494) circle (1.2pt);
\fill[blue!80!black] (7.5580,5.3491) circle (1.2pt);
\fill[blue!80!black] (7.5573,5.3489) circle (1.2pt);
\fill[blue!80!black] (7.5578,5.3486) circle (1.2pt);
\fill[blue!80!black] (7.5571,5.3484) circle (1.2pt);
\fill[blue!80!black] (7.5575,5.3481) circle (1.2pt);
\fill[blue!80!black] (7.5568,5.3480) circle (1.2pt);
\fill[blue!80!black] (7.5573,5.3477) circle (1.2pt);
\fill[blue!80!black] (7.5565,5.3475) circle (1.2pt);
\fill[blue!80!black] (7.5570,5.3472) circle (1.2pt);
\fill[blue!80!black] (7.5562,5.3470) circle (1.2pt);
\fill[blue!80!black] (7.5567,5.3467) circle (1.2pt);
\fill[blue!80!black] (7.5559,5.3466) circle (1.2pt);
\fill[blue!80!black] (7.5565,5.3463) circle (1.2pt);
\fill[blue!80!black] (7.5561,5.3462) circle (1.2pt);
\fill[blue!80!black] (7.5562,5.3458) circle (1.2pt);
\fill[blue!80!black] (7.5558,5.3457) circle (1.2pt);
\fill[blue!80!black] (7.5559,5.3453) circle (1.2pt);
\fill[blue!80!black] (7.5555,5.3453) circle (1.2pt);
\fill[blue!80!black] (7.5557,5.3448) circle (1.2pt);
\fill[blue!80!black] (7.5552,5.3448) circle (1.2pt);
\fill[blue!80!black] (7.5554,5.3444) circle (1.2pt);
\fill[blue!80!black] (7.5550,5.3443) circle (1.2pt);
\fill[blue!80!black] (7.5552,5.3439) circle (1.2pt);
\fill[blue!80!black] (7.5547,5.3439) circle (1.2pt);
\fill[blue!80!black] (7.5549,5.3434) circle (1.2pt);
\fill[blue!80!black] (7.5544,5.3434) circle (1.2pt);
\fill[blue!80!black] (7.5546,5.3430) circle (1.2pt);
\fill[blue!80!black] (7.5541,5.3429) circle (1.2pt);
\fill[blue!80!black] (7.5544,5.3425) circle (1.2pt);
\fill[blue!80!black] (7.5539,5.3425) circle (1.2pt);
\fill[blue!80!black] (7.5541,5.3420) circle (1.2pt);
\fill[blue!80!black] (7.5536,5.3420) circle (1.2pt);
\fill[blue!80!black] (7.5538,5.3416) circle (1.2pt);
\fill[blue!80!black] (7.5533,5.3415) circle (1.2pt);
\fill[blue!80!black] (7.5536,5.3411) circle (1.2pt);
\fill[blue!80!black] (7.5531,5.3411) circle (1.2pt);
\fill[blue!80!black] (7.5533,5.3406) circle (1.2pt);
\fill[blue!80!black] (7.5528,5.3406) circle (1.2pt);
\fill[blue!80!black] (7.5530,5.3402) circle (1.2pt);
\fill[blue!80!black] (7.5525,5.3401) circle (1.2pt);
\fill[blue!80!black] (7.5528,5.3397) circle (1.2pt);
\fill[blue!80!black] (7.5522,5.3397) circle (1.2pt);
\fill[blue!80!black] (7.5525,5.3393) circle (1.2pt);
\fill[blue!80!black] (7.5520,5.3392) circle (1.2pt);
\fill[blue!80!black] (7.5523,5.3388) circle (1.2pt);
\fill[blue!80!black] (7.5517,5.3388) circle (1.2pt);
\fill[blue!80!black] (7.5520,5.3383) circle (1.2pt);
\fill[blue!80!black] (7.5514,5.3383) circle (1.2pt);
\fill[blue!80!black] (7.5517,5.3379) circle (1.2pt);
\fill[blue!80!black] (7.5512,5.3378) circle (1.2pt);
\fill[blue!80!black] (7.5515,5.3374) circle (1.2pt);
\fill[blue!80!black] (7.5509,5.3374) circle (1.2pt);
\fill[blue!80!black] (7.5512,5.3369) circle (1.2pt);
\fill[blue!80!black] (7.5506,5.3369) circle (1.2pt);
\fill[blue!80!black] (7.5509,5.3365) circle (1.2pt);
\fill[blue!80!black] (7.5504,5.3364) circle (1.2pt);
\fill[blue!80!black] (7.5507,5.3360) circle (1.2pt);
\fill[blue!80!black] (7.5501,5.3360) circle (1.2pt);
\fill[blue!80!black] (7.5504,5.3356) circle (1.2pt);
\fill[blue!80!black] (7.5498,5.3355) circle (1.2pt);
\fill[blue!80!black] (7.5501,5.3351) circle (1.2pt);
\fill[blue!80!black] (7.5496,5.3351) circle (1.2pt);
\fill[blue!80!black] (7.5499,5.3346) circle (1.2pt);
\fill[blue!80!black] (7.5493,5.3346) circle (1.2pt);
\fill[blue!80!black] (7.5496,5.3342) circle (1.2pt);
\fill[blue!80!black] (7.5490,5.3341) circle (1.2pt);
\fill[blue!80!black] (7.5493,5.3337) circle (1.2pt);
\fill[blue!80!black] (7.5488,5.3337) circle (1.2pt);
\fill[blue!80!black] (7.5491,5.3333) circle (1.2pt);
\fill[blue!80!black] (7.5485,5.3332) circle (1.2pt);
\fill[blue!80!black] (7.5488,5.3328) circle (1.2pt);
\fill[blue!80!black] (7.5482,5.3328) circle (1.2pt);
\fill[blue!80!black] (7.5485,5.3323) circle (1.2pt);
\fill[blue!80!black] (7.5480,5.3323) circle (1.2pt);
\fill[blue!80!black] (7.5483,5.3319) circle (1.2pt);
\fill[blue!80!black] (7.5477,5.3319) circle (1.2pt);
\fill[blue!80!black] (7.5480,5.3314) circle (1.2pt);
\fill[blue!80!black] (7.5474,5.3314) circle (1.2pt);
\fill[blue!80!black] (7.5477,5.3310) circle (1.2pt);
\fill[blue!80!black] (7.5472,5.3309) circle (1.2pt);
\fill[blue!80!black] (7.5474,5.3305) circle (1.2pt);
\fill[blue!80!black] (7.5469,5.3305) circle (1.2pt);
\fill[blue!80!black] (7.5472,5.3301) circle (1.2pt);
\fill[blue!80!black] (7.5466,5.3300) circle (1.2pt);
\fill[blue!80!black] (7.5469,5.3296) circle (1.2pt);
\fill[blue!80!black] (7.5464,5.3296) circle (1.2pt);
\fill[blue!80!black] (7.5466,5.3291) circle (1.2pt);
\fill[blue!80!black] (7.5461,5.3291) circle (1.2pt);
\fill[blue!80!black] (7.5464,5.3287) circle (1.2pt);
\fill[blue!80!black] (7.5459,5.3286) circle (1.2pt);
\fill[blue!80!black] (7.5461,5.3282) circle (1.2pt);
\fill[blue!80!black] (7.5456,5.3282) circle (1.2pt);
\fill[blue!80!black] (7.5458,5.3278) circle (1.2pt);
\fill[blue!80!black] (7.5453,5.3277) circle (1.2pt);
\fill[blue!80!black] (7.5456,5.3273) circle (1.2pt);
\fill[blue!80!black] (7.5451,5.3273) circle (1.2pt);
\fill[blue!80!black] (7.5453,5.3269) circle (1.2pt);
\fill[blue!80!black] (7.5448,5.3268) circle (1.2pt);
\fill[blue!80!black] (7.5450,5.3264) circle (1.2pt);
\fill[blue!80!black] (7.5445,5.3264) circle (1.2pt);
\fill[blue!80!black] (7.5447,5.3260) circle (1.2pt);
\fill[blue!80!black] (7.5443,5.3259) circle (1.2pt);
\fill[blue!80!black] (7.5445,5.3255) circle (1.2pt);
\fill[blue!80!black] (7.5440,5.3255) circle (1.2pt);
\fill[blue!80!black] (7.5442,5.3251) circle (1.2pt);
\fill[blue!80!black] (7.5438,5.3250) circle (1.2pt);
\fill[blue!80!black] (7.5439,5.3246) circle (1.2pt);
\fill[blue!80!black] (7.5435,5.3245) circle (1.2pt);
\fill[blue!80!black] (7.5437,5.3242) circle (1.2pt);
\fill[blue!80!black] (7.5432,5.3241) circle (1.2pt);
\fill[blue!80!black] (7.5434,5.3237) circle (1.2pt);
\fill[blue!80!black] (7.5426,5.3236) circle (1.2pt);
\fill[blue!80!black] (7.5431,5.3233) circle (1.2pt);
\fill[blue!80!black] (7.5423,5.3231) circle (1.2pt);
\fill[blue!80!black] (7.5428,5.3228) circle (1.2pt);
\fill[blue!80!black] (7.5421,5.3227) circle (1.2pt);
\fill[blue!80!black] (7.5426,5.3224) circle (1.2pt);
\fill[blue!80!black] (7.5418,5.3222) circle (1.2pt);
\fill[blue!80!black] (7.5423,5.3219) circle (1.2pt);
\fill[blue!80!black] (7.5416,5.3217) circle (1.2pt);
\fill[blue!80!black] (7.5420,5.3215) circle (1.2pt);
\fill[blue!80!black] (7.5413,5.3213) circle (1.2pt);
\fill[blue!80!black] (7.5418,5.3210) circle (1.2pt);
\fill[blue!80!black] (7.5411,5.3208) circle (1.2pt);
\fill[blue!80!black] (7.5415,5.3206) circle (1.2pt);
\fill[blue!80!black] (7.5408,5.3204) circle (1.2pt);
\fill[blue!80!black] (7.5412,5.3201) circle (1.2pt);
\fill[blue!80!black] (7.5406,5.3199) circle (1.2pt);
\fill[blue!80!black] (7.5410,5.3197) circle (1.2pt);
\fill[blue!80!black] (7.5403,5.3195) circle (1.2pt);
\fill[blue!80!black] (7.5407,5.3192) circle (1.2pt);
\fill[blue!80!black] (7.5401,5.3190) circle (1.2pt);
\fill[blue!80!black] (7.5404,5.3188) circle (1.2pt);
\fill[blue!80!black] (7.5398,5.3186) circle (1.2pt);
\fill[blue!80!black] (7.5401,5.3183) circle (1.2pt);
\fill[blue!80!black] (7.5396,5.3181) circle (1.2pt);
\fill[blue!80!black] (7.5399,5.3179) circle (1.2pt);
\fill[blue!80!black] (7.5393,5.3177) circle (1.2pt);
\fill[blue!80!black] (7.5396,5.3174) circle (1.2pt);
\fill[blue!80!black] (7.5391,5.3172) circle (1.2pt);
\fill[blue!80!black] (7.5393,5.3170) circle (1.2pt);
\fill[blue!80!black] (7.5388,5.3168) circle (1.2pt);
\fill[blue!80!black] (7.5393,5.3163) circle (1.2pt);
\fill[blue!80!black] (7.5386,5.3163) circle (1.2pt);
\fill[blue!80!black] (7.5390,5.3159) circle (1.2pt);
\fill[blue!80!black] (7.5383,5.3159) circle (1.2pt);
\fill[blue!80!black] (7.5387,5.3154) circle (1.2pt);
\fill[blue!80!black] (7.5381,5.3154) circle (1.2pt);
\fill[blue!80!black] (7.5384,5.3150) circle (1.2pt);
\fill[blue!80!black] (7.5378,5.3150) circle (1.2pt);
\fill[blue!80!black] (7.5382,5.3145) circle (1.2pt);
\fill[blue!80!black] (7.5376,5.3145) circle (1.2pt);
\fill[blue!80!black] (7.5379,5.3141) circle (1.2pt);
\fill[blue!80!black] (7.5373,5.3141) circle (1.2pt);
\fill[blue!80!black] (7.5376,5.3136) circle (1.2pt);
\fill[blue!80!black] (7.5371,5.3136) circle (1.2pt);
\fill[blue!80!black] (7.5373,5.3132) circle (1.2pt);
\fill[blue!80!black] (7.5368,5.3132) circle (1.2pt);
\fill[blue!80!black] (7.5370,5.3128) circle (1.2pt);
\fill[blue!80!black] (7.5366,5.3127) circle (1.2pt);
\fill[blue!80!black] (7.5368,5.3123) circle (1.2pt);
\fill[blue!80!black] (7.5363,5.3123) circle (1.2pt);
\fill[blue!80!black] (7.5365,5.3119) circle (1.2pt);
\fill[blue!80!black] (7.5361,5.3118) circle (1.2pt);
\fill[blue!80!black] (7.5362,5.3114) circle (1.2pt);
\fill[blue!80!black] (7.5354,5.3113) circle (1.2pt);
\fill[blue!80!black] (7.5359,5.3110) circle (1.2pt);
\fill[blue!80!black] (7.5352,5.3109) circle (1.2pt);
\fill[blue!80!black] (7.5357,5.3106) circle (1.2pt);
\fill[blue!80!black] (7.5350,5.3104) circle (1.2pt);
\fill[blue!80!black] (7.5354,5.3101) circle (1.2pt);
\fill[blue!80!black] (7.5347,5.3099) circle (1.2pt);
\fill[blue!80!black] (7.5351,5.3097) circle (1.2pt);
\fill[blue!80!black] (7.5345,5.3095) circle (1.2pt);
\fill[blue!80!black] (7.5349,5.3092) circle (1.2pt);
\fill[blue!80!black] (7.5343,5.3090) circle (1.2pt);
\fill[blue!80!black] (7.5346,5.3088) circle (1.2pt);
\fill[blue!80!black] (7.5340,5.3086) circle (1.2pt);
\fill[blue!80!black] (7.5343,5.3084) circle (1.2pt);
\fill[blue!80!black] (7.5338,5.3082) circle (1.2pt);
\fill[blue!80!black] (7.5340,5.3079) circle (1.2pt);
\fill[blue!80!black] (7.5335,5.3077) circle (1.2pt);
\fill[blue!80!black] (7.5338,5.3075) circle (1.2pt);
\fill[blue!80!black] (7.5328,5.3070) circle (1.2pt);
\fill[blue!80!black] (7.5335,5.3067) circle (1.2pt);
\fill[blue!80!black] (7.5330,5.3066) circle (1.2pt);
\fill[blue!80!black] (7.5332,5.3063) circle (1.2pt);
\fill[blue!80!black] (7.5328,5.3062) circle (1.2pt);
\fill[blue!80!black] (7.5331,5.3055) circle (1.2pt);
\fill[blue!80!black] (7.5322,5.3055) circle (1.2pt);
\fill[blue!80!black] (7.5325,5.3053) circle (1.2pt);
\fill[blue!80!black] (7.5320,5.3051) circle (1.2pt);
\fill[blue!80!black] (7.5322,5.3048) circle (1.2pt);
\fill[blue!80!black] (7.5313,5.3044) circle (1.2pt);
\fill[blue!80!black] (7.5319,5.3041) circle (1.2pt);
\fill[blue!80!black] (7.5315,5.3040) circle (1.2pt);
\fill[blue!80!black] (7.5316,5.3036) circle (1.2pt);
\fill[blue!80!black] (7.5308,5.3035) circle (1.2pt);
\fill[blue!80!black] (7.5314,5.3032) circle (1.2pt);
\fill[blue!80!black] (7.5306,5.3030) circle (1.2pt);
\fill[blue!80!black] (7.5311,5.3028) circle (1.2pt);
\fill[blue!80!black] (7.5304,5.3026) circle (1.2pt);
\fill[blue!80!black] (7.5308,5.3023) circle (1.2pt);
\fill[blue!80!black] (7.5302,5.3021) circle (1.2pt);
\fill[blue!80!black] (7.5305,5.3019) circle (1.2pt);
\fill[blue!80!black] (7.5299,5.3017) circle (1.2pt);
\fill[blue!80!black] (7.5303,5.3014) circle (1.2pt);
\fill[blue!80!black] (7.5297,5.3013) circle (1.2pt);
\fill[blue!80!black] (7.5300,5.3010) circle (1.2pt);
\fill[blue!80!black] (7.5295,5.3008) circle (1.2pt);
\fill[blue!80!black] (7.5297,5.3006) circle (1.2pt);
\fill[blue!80!black] (7.5288,5.3002) circle (1.2pt);
\fill[blue!80!black] (7.5294,5.2998) circle (1.2pt);
\fill[blue!80!black] (7.5290,5.2998) circle (1.2pt);
\fill[blue!80!black] (7.5292,5.2994) circle (1.2pt);
\fill[blue!80!black] (7.5283,5.2993) circle (1.2pt);
\fill[blue!80!black] (7.5289,5.2989) circle (1.2pt);
\fill[blue!80!black] (7.5281,5.2988) circle (1.2pt);
\fill[blue!80!black] (7.5286,5.2985) circle (1.2pt);
\fill[blue!80!black] (7.5279,5.2984) circle (1.2pt);
\fill[blue!80!black] (7.5283,5.2981) circle (1.2pt);
\fill[blue!80!black] (7.5277,5.2979) circle (1.2pt);
\fill[blue!80!black] (7.5280,5.2977) circle (1.2pt);
\fill[blue!80!black] (7.5275,5.2975) circle (1.2pt);
\fill[blue!80!black] (7.5278,5.2972) circle (1.2pt);
\fill[blue!80!black] (7.5272,5.2970) circle (1.2pt);
\fill[blue!80!black] (7.5275,5.2968) circle (1.2pt);
\fill[blue!80!black] (7.5270,5.2966) circle (1.2pt);
\fill[blue!80!black] (7.5275,5.2961) circle (1.2pt);
\fill[blue!80!black] (7.5268,5.2961) circle (1.2pt);
\fill[blue!80!black] (7.5272,5.2957) circle (1.2pt);
\fill[blue!80!black] (7.5265,5.2957) circle (1.2pt);
\fill[blue!80!black] (7.5269,5.2953) circle (1.2pt);
\fill[blue!80!black] (7.5263,5.2953) circle (1.2pt);
\fill[blue!80!black] (7.5266,5.2949) circle (1.2pt);
\fill[blue!80!black] (7.5261,5.2948) circle (1.2pt);
\fill[blue!80!black] (7.5263,5.2944) circle (1.2pt);
\fill[blue!80!black] (7.5258,5.2944) circle (1.2pt);
\fill[blue!80!black] (7.5260,5.2940) circle (1.2pt);
\fill[blue!80!black] (7.5256,5.2940) circle (1.2pt);
\fill[blue!80!black] (7.5258,5.2932) circle (1.2pt);
\fill[blue!80!black] (7.5250,5.2933) circle (1.2pt);
\fill[blue!80!black] (7.5253,5.2930) circle (1.2pt);
\fill[blue!80!black] (7.5243,5.2926) circle (1.2pt);
\fill[blue!80!black] (7.5250,5.2923) circle (1.2pt);
\fill[blue!80!black] (7.5246,5.2922) circle (1.2pt);
\fill[blue!80!black] (7.5248,5.2915) circle (1.2pt);
\fill[blue!80!black] (7.5240,5.2915) circle (1.2pt);
\fill[blue!80!black] (7.5243,5.2913) circle (1.2pt);
\fill[blue!80!black] (7.5233,5.2909) circle (1.2pt);
\fill[blue!80!black] (7.5240,5.2905) circle (1.2pt);
\fill[blue!80!black] (7.5235,5.2905) circle (1.2pt);
\fill[blue!80!black] (7.5237,5.2901) circle (1.2pt);
\fill[blue!80!black] (7.5229,5.2900) circle (1.2pt);
\fill[blue!80!black] (7.5234,5.2897) circle (1.2pt);
\fill[blue!80!black] (7.5227,5.2895) circle (1.2pt);
\fill[blue!80!black] (7.5231,5.2893) circle (1.2pt);
\fill[blue!80!black] (7.5225,5.2891) circle (1.2pt);
\fill[blue!80!black] (7.5229,5.2888) circle (1.2pt);
\fill[blue!80!black] (7.5223,5.2887) circle (1.2pt);
\fill[blue!80!black] (7.5226,5.2884) circle (1.2pt);
\fill[blue!80!black] (7.5221,5.2882) circle (1.2pt);
\fill[blue!80!black] (7.5223,5.2880) circle (1.2pt);
\fill[blue!80!black] (7.5214,5.2876) circle (1.2pt);
\fill[blue!80!black] (7.5220,5.2873) circle (1.2pt);
\fill[blue!80!black] (7.5216,5.2872) circle (1.2pt);
\fill[blue!80!black] (7.5218,5.2865) circle (1.2pt);
\fill[blue!80!black] (7.5211,5.2865) circle (1.2pt);
\fill[blue!80!black] (7.5215,5.2861) circle (1.2pt);
\fill[blue!80!black] (7.5208,5.2861) circle (1.2pt);
\fill[blue!80!black] (7.5212,5.2856) circle (1.2pt);
\fill[blue!80!black] (7.5206,5.2856) circle (1.2pt);
\fill[blue!80!black] (7.5209,5.2852) circle (1.2pt);
\fill[blue!80!black] (7.5204,5.2852) circle (1.2pt);
\fill[blue!80!black] (7.5206,5.2848) circle (1.2pt);
\fill[blue!80!black] (7.5202,5.2848) circle (1.2pt);
\fill[blue!80!black] (7.5203,5.2844) circle (1.2pt);
\fill[blue!80!black] (7.5199,5.2843) circle (1.2pt);
\fill[blue!80!black] (7.5202,5.2836) circle (1.2pt);
\fill[blue!80!black] (7.5194,5.2837) circle (1.2pt);
\fill[blue!80!black] (7.5196,5.2834) circle (1.2pt);
\fill[blue!80!black] (7.5187,5.2830) circle (1.2pt);
\fill[blue!80!black] (7.5193,5.2827) circle (1.2pt);
\fill[blue!80!black] (7.5185,5.2826) circle (1.2pt);
\fill[blue!80!black] (7.5190,5.2823) circle (1.2pt);
\fill[blue!80!black] (7.5183,5.2821) circle (1.2pt);
\fill[blue!80!black] (7.5187,5.2819) circle (1.2pt);
\fill[blue!80!black] (7.5181,5.2817) circle (1.2pt);
\fill[blue!80!black] (7.5185,5.2814) circle (1.2pt);
\fill[blue!80!black] (7.5179,5.2813) circle (1.2pt);
\fill[blue!80!black] (7.5182,5.2810) circle (1.2pt);
\fill[blue!80!black] (7.5177,5.2808) circle (1.2pt);
\fill[blue!80!black] (7.5181,5.2804) circle (1.2pt);
\fill[blue!80!black] (7.5175,5.2804) circle (1.2pt);
\fill[blue!80!black] (7.5178,5.2800) circle (1.2pt);
\fill[blue!80!black] (7.5173,5.2800) circle (1.2pt);
\fill[blue!80!black] (7.5175,5.2796) circle (1.2pt);
\fill[blue!80!black] (7.5170,5.2795) circle (1.2pt);
\fill[blue!80!black] (7.5172,5.2791) circle (1.2pt);
\fill[blue!80!black] (7.5168,5.2791) circle (1.2pt);
\fill[blue!80!black] (7.5171,5.2784) circle (1.2pt);
\fill[blue!80!black] (7.5163,5.2784) circle (1.2pt);
\fill[blue!80!black] (7.5165,5.2782) circle (1.2pt);
\fill[blue!80!black] (7.5156,5.2778) circle (1.2pt);
\fill[blue!80!black] (7.5162,5.2775) circle (1.2pt);
\fill[blue!80!black] (7.5154,5.2773) circle (1.2pt);
\fill[blue!80!black] (7.5159,5.2770) circle (1.2pt);
\fill[blue!80!black] (7.5152,5.2769) circle (1.2pt);
\fill[blue!80!black] (7.5156,5.2766) circle (1.2pt);
\fill[blue!80!black] (7.5150,5.2765) circle (1.2pt);
\fill[blue!80!black] (7.5154,5.2762) circle (1.2pt);
\fill[blue!80!black] (7.5148,5.2760) circle (1.2pt);
\fill[blue!80!black] (7.5151,5.2758) circle (1.2pt);
\fill[blue!80!black] (7.5142,5.2754) circle (1.2pt);
\fill[blue!80!black] (7.5148,5.2751) circle (1.2pt);
\fill[blue!80!black] (7.5140,5.2750) circle (1.2pt);
\fill[blue!80!black] (7.5145,5.2747) circle (1.2pt);
\fill[blue!80!black] (7.5138,5.2745) circle (1.2pt);
\fill[blue!80!black] (7.5142,5.2742) circle (1.2pt);
\fill[blue!80!black] (7.5136,5.2741) circle (1.2pt);
\fill[blue!80!black] (7.5139,5.2738) circle (1.2pt);
\fill[blue!80!black] (7.5134,5.2736) circle (1.2pt);
\fill[blue!80!black] (7.5137,5.2734) circle (1.2pt);
\fill[blue!80!black] (7.5128,5.2730) circle (1.2pt);
\fill[blue!80!black] (7.5134,5.2727) circle (1.2pt);
\fill[blue!80!black] (7.5126,5.2726) circle (1.2pt);
\fill[blue!80!black] (7.5131,5.2723) circle (1.2pt);
\fill[blue!80!black] (7.5124,5.2721) circle (1.2pt);
\fill[blue!80!black] (7.5128,5.2719) circle (1.2pt);
\fill[blue!80!black] (7.5122,5.2717) circle (1.2pt);
\fill[blue!80!black] (7.5125,5.2715) circle (1.2pt);
\fill[blue!80!black] (7.5120,5.2713) circle (1.2pt);
\fill[blue!80!black] (7.5123,5.2710) circle (1.2pt);
\fill[blue!80!black] (7.5114,5.2706) circle (1.2pt);
\fill[blue!80!black] (7.5119,5.2703) circle (1.2pt);
\fill[blue!80!black] (7.5112,5.2702) circle (1.2pt);
\fill[blue!80!black] (7.5117,5.2699) circle (1.2pt);
\fill[blue!80!black] (7.5110,5.2698) circle (1.2pt);
\fill[blue!80!black] (7.5114,5.2695) circle (1.2pt);
\fill[blue!80!black] (7.5108,5.2693) circle (1.2pt);
\fill[blue!80!black] (7.5111,5.2691) circle (1.2pt);
\fill[blue!80!black] (7.5106,5.2689) circle (1.2pt);
\fill[blue!80!black] (7.5111,5.2685) circle (1.2pt);
\fill[blue!80!black] (7.5104,5.2685) circle (1.2pt);
\fill[blue!80!black] (7.5107,5.2681) circle (1.2pt);
\fill[blue!80!black] (7.5102,5.2680) circle (1.2pt);
\fill[blue!80!black] (7.5104,5.2677) circle (1.2pt);
\fill[blue!80!black] (7.5100,5.2676) circle (1.2pt);
\fill[blue!80!black] (7.5101,5.2673) circle (1.2pt);
\fill[blue!80!black] (7.5094,5.2671) circle (1.2pt);
\fill[blue!80!black] (7.5098,5.2669) circle (1.2pt);
\fill[blue!80!black] (7.5092,5.2667) circle (1.2pt);
\fill[blue!80!black] (7.5096,5.2665) circle (1.2pt);
\fill[blue!80!black] (7.5090,5.2663) circle (1.2pt);
\fill[blue!80!black] (7.5093,5.2660) circle (1.2pt);
\fill[blue!80!black] (7.5088,5.2658) circle (1.2pt);
\fill[blue!80!black] (7.5092,5.2654) circle (1.2pt);
\fill[blue!80!black] (7.5086,5.2654) circle (1.2pt);
\fill[blue!80!black] (7.5089,5.2650) circle (1.2pt);
\fill[blue!80!black] (7.5084,5.2650) circle (1.2pt);
\fill[blue!80!black] (7.5086,5.2646) circle (1.2pt);
\fill[blue!80!black] (7.5082,5.2646) circle (1.2pt);
\fill[blue!80!black] (7.5084,5.2639) circle (1.2pt);
\fill[blue!80!black] (7.5077,5.2639) circle (1.2pt);
\fill[blue!80!black] (7.5081,5.2635) circle (1.2pt);
\fill[blue!80!black] (7.5075,5.2635) circle (1.2pt);
\fill[blue!80!black] (7.5078,5.2631) circle (1.2pt);
\fill[blue!80!black] (7.5072,5.2631) circle (1.2pt);
\fill[blue!80!black] (7.5075,5.2627) circle (1.2pt);
\fill[blue!80!black] (7.5070,5.2627) circle (1.2pt);
\fill[blue!80!black] (7.5072,5.2623) circle (1.2pt);
\fill[blue!80!black] (7.5064,5.2622) circle (1.2pt);
\fill[blue!80!black] (7.5069,5.2619) circle (1.2pt);
\fill[blue!80!black] (7.5063,5.2617) circle (1.2pt);
\fill[blue!80!black] (7.5066,5.2615) circle (1.2pt);
\fill[blue!80!black] (7.5061,5.2613) circle (1.2pt);
\fill[blue!80!black] (7.5063,5.2611) circle (1.2pt);
\fill[blue!80!black] (7.5054,5.2607) circle (1.2pt);
\fill[blue!80!black] (7.5060,5.2604) circle (1.2pt);
\fill[blue!80!black] (7.5053,5.2603) circle (1.2pt);
\fill[blue!80!black] (7.5057,5.2600) circle (1.2pt);
\fill[blue!80!black] (7.5051,5.2598) circle (1.2pt);
\fill[blue!80!black] (7.5055,5.2596) circle (1.2pt);
\fill[blue!80!black] (7.5049,5.2594) circle (1.2pt);
\fill[blue!80!black] (7.5052,5.2592) circle (1.2pt);
\fill[blue!80!black] (7.5043,5.2588) circle (1.2pt);
\fill[blue!80!black] (7.5049,5.2585) circle (1.2pt);
\fill[blue!80!black] (7.5042,5.2583) circle (1.2pt);
\fill[blue!80!black] (7.5046,5.2581) circle (1.2pt);
\fill[blue!80!black] (7.5040,5.2579) circle (1.2pt);
\fill[blue!80!black] (7.5043,5.2577) circle (1.2pt);
\fill[blue!80!black] (7.5038,5.2575) circle (1.2pt);
\fill[blue!80!black] (7.5040,5.2573) circle (1.2pt);
\fill[blue!80!black] (7.5032,5.2569) circle (1.2pt);
\fill[blue!80!black] (7.5037,5.2566) circle (1.2pt);
\fill[blue!80!black] (7.5030,5.2564) circle (1.2pt);
\fill[blue!80!black] (7.5034,5.2562) circle (1.2pt);
\fill[blue!80!black] (7.5029,5.2560) circle (1.2pt);
\fill[blue!80!black] (7.5032,5.2558) circle (1.2pt);
\fill[blue!80!black] (7.5027,5.2556) circle (1.2pt);
\fill[blue!80!black] (7.5031,5.2551) circle (1.2pt);
\fill[blue!80!black] (7.5025,5.2552) circle (1.2pt);
\fill[blue!80!black] (7.5028,5.2548) circle (1.2pt);
\fill[blue!80!black] (7.5023,5.2547) circle (1.2pt);
\fill[blue!80!black] (7.5025,5.2544) circle (1.2pt);
\fill[blue!80!black] (7.5021,5.2543) circle (1.2pt);
\fill[blue!80!black] (7.5023,5.2536) circle (1.2pt);
\fill[blue!80!black] (7.5016,5.2537) circle (1.2pt);
\fill[blue!80!black] (7.5020,5.2532) circle (1.2pt);
\fill[blue!80!black] (7.5014,5.2533) circle (1.2pt);
\fill[blue!80!black] (7.5016,5.2529) circle (1.2pt);
\fill[blue!80!black] (7.5011,5.2528) circle (1.2pt);
\fill[blue!80!black] (7.5013,5.2525) circle (1.2pt);
\fill[blue!80!black] (7.5005,5.2524) circle (1.2pt);
\fill[blue!80!black] (7.5010,5.2521) circle (1.2pt);
\fill[blue!80!black] (7.5004,5.2519) circle (1.2pt);
\fill[blue!80!black] (7.5007,5.2517) circle (1.2pt);
\fill[blue!80!black] (7.5002,5.2515) circle (1.2pt);
\fill[blue!80!black] (7.5005,5.2513) circle (1.2pt);
\fill[blue!80!black] (7.4996,5.2509) circle (1.2pt);
\fill[blue!80!black] (7.5002,5.2506) circle (1.2pt);
\fill[blue!80!black] (7.4995,5.2505) circle (1.2pt);
\fill[blue!80!black] (7.4999,5.2502) circle (1.2pt);
\fill[blue!80!black] (7.4993,5.2500) circle (1.2pt);
\fill[blue!80!black] (7.4996,5.2498) circle (1.2pt);
\fill[blue!80!black] (7.4991,5.2496) circle (1.2pt);
\fill[blue!80!black] (7.4995,5.2492) circle (1.2pt);
\fill[blue!80!black] (7.4989,5.2492) circle (1.2pt);
\fill[blue!80!black] (7.4992,5.2488) circle (1.2pt);
\fill[blue!80!black] (7.4987,5.2488) circle (1.2pt);
\fill[blue!80!black] (7.4989,5.2484) circle (1.2pt);
\fill[blue!80!black] (7.4981,5.2483) circle (1.2pt);
\fill[blue!80!black] (7.4986,5.2480) circle (1.2pt);
\fill[blue!80!black] (7.4980,5.2479) circle (1.2pt);
\fill[blue!80!black] (7.4983,5.2477) circle (1.2pt);
\fill[blue!80!black] (7.4978,5.2475) circle (1.2pt);
\fill[blue!80!black] (7.4981,5.2473) circle (1.2pt);
\fill[blue!80!black] (7.4972,5.2469) circle (1.2pt);
\fill[blue!80!black] (7.4977,5.2466) circle (1.2pt);
\fill[blue!80!black] (7.4971,5.2464) circle (1.2pt);
\fill[blue!80!black] (7.4975,5.2462) circle (1.2pt);
\fill[blue!80!black] (7.4969,5.2460) circle (1.2pt);
\fill[blue!80!black] (7.4972,5.2458) circle (1.2pt);
\fill[blue!80!black] (7.4963,5.2454) circle (1.2pt);
\fill[blue!80!black] (7.4969,5.2451) circle (1.2pt);
\fill[blue!80!black] (7.4961,5.2450) circle (1.2pt);
\fill[blue!80!black] (7.4966,5.2447) circle (1.2pt);
\fill[blue!80!black] (7.4960,5.2446) circle (1.2pt);
\fill[blue!80!black] (7.4963,5.2443) circle (1.2pt);
\fill[blue!80!black] (7.4958,5.2441) circle (1.2pt);
\fill[blue!80!black] (7.4962,5.2437) circle (1.2pt);
\fill[blue!80!black] (7.4956,5.2437) circle (1.2pt);
\fill[blue!80!black] (7.4959,5.2433) circle (1.2pt);
\fill[blue!80!black] (7.4954,5.2433) circle (1.2pt);
\fill[blue!80!black] (7.4956,5.2430) circle (1.2pt);
\fill[blue!80!black] (7.4949,5.2429) circle (1.2pt);
\fill[blue!80!black] (7.4953,5.2426) circle (1.2pt);
\fill[blue!80!black] (7.4947,5.2424) circle (1.2pt);
\fill[blue!80!black] (7.4950,5.2422) circle (1.2pt);
\fill[blue!80!black] (7.4945,5.2420) circle (1.2pt);
\fill[blue!80!black] (7.4950,5.2416) circle (1.2pt);
\fill[blue!80!black] (7.4944,5.2416) circle (1.2pt);
\fill[blue!80!black] (7.4946,5.2412) circle (1.2pt);
\fill[blue!80!black] (7.4942,5.2412) circle (1.2pt);
\fill[blue!80!black] (7.4943,5.2408) circle (1.2pt);
\fill[blue!80!black] (7.4936,5.2407) circle (1.2pt);
\fill[blue!80!black] (7.4940,5.2405) circle (1.2pt);
\fill[blue!80!black] (7.4934,5.2403) circle (1.2pt);
\fill[blue!80!black] (7.4938,5.2401) circle (1.2pt);
\fill[blue!80!black] (7.4933,5.2399) circle (1.2pt);
\fill[blue!80!black] (7.4937,5.2395) circle (1.2pt);
\fill[blue!80!black] (7.4931,5.2395) circle (1.2pt);
\fill[blue!80!black] (7.4934,5.2391) circle (1.2pt);
\fill[blue!80!black] (7.4929,5.2391) circle (1.2pt);
\fill[blue!80!black] (7.4930,5.2387) circle (1.2pt);
\fill[blue!80!black] (7.4923,5.2386) circle (1.2pt);
\fill[blue!80!black] (7.4928,5.2383) circle (1.2pt);
\fill[blue!80!black] (7.4922,5.2382) circle (1.2pt);
\fill[blue!80!black] (7.4925,5.2380) circle (1.2pt);
\fill[blue!80!black] (7.4920,5.2378) circle (1.2pt);
\fill[blue!80!black] (7.4924,5.2374) circle (1.2pt);
\fill[blue!80!black] (7.4918,5.2374) circle (1.2pt);
\fill[blue!80!black] (7.4921,5.2370) circle (1.2pt);
\fill[blue!80!black] (7.4916,5.2370) circle (1.2pt);
\fill[blue!80!black] (7.4918,5.2366) circle (1.2pt);
\fill[blue!80!black] (7.4911,5.2365) circle (1.2pt);
\fill[blue!80!black] (7.4915,5.2362) circle (1.2pt);
\fill[blue!80!black] (7.4909,5.2361) circle (1.2pt);
\fill[blue!80!black] (7.4912,5.2359) circle (1.2pt);
\fill[blue!80!black] (7.4903,5.2355) circle (1.2pt);
\fill[blue!80!black] (7.4909,5.2352) circle (1.2pt);
\fill[blue!80!black] (7.4902,5.2351) circle (1.2pt);
\fill[blue!80!black] (7.4906,5.2348) circle (1.2pt);
\fill[blue!80!black] (7.4901,5.2346) circle (1.2pt);
\fill[blue!80!black] (7.4903,5.2344) circle (1.2pt);
\fill[blue!80!black] (7.4895,5.2340) circle (1.2pt);
\fill[blue!80!black] (7.4900,5.2337) circle (1.2pt);
\fill[blue!80!black] (7.4893,5.2336) circle (1.2pt);
\fill[blue!80!black] (7.4897,5.2334) circle (1.2pt);
\fill[blue!80!black] (7.4892,5.2332) circle (1.2pt);
\fill[blue!80!black] (7.4895,5.2330) circle (1.2pt);
\fill[blue!80!black] (7.4886,5.2326) circle (1.2pt);
\fill[blue!80!black] (7.4891,5.2323) circle (1.2pt);
\fill[blue!80!black] (7.4885,5.2322) circle (1.2pt);
\fill[blue!80!black] (7.4889,5.2319) circle (1.2pt);
\fill[blue!80!black] (7.4883,5.2318) circle (1.2pt);
\fill[blue!80!black] (7.4886,5.2315) circle (1.2pt);
\fill[blue!80!black] (7.4878,5.2312) circle (1.2pt);
\fill[blue!80!black] (7.4883,5.2309) circle (1.2pt);
\fill[blue!80!black] (7.4876,5.2307) circle (1.2pt);
\fill[blue!80!black] (7.4880,5.2305) circle (1.2pt);
\fill[blue!80!black] (7.4875,5.2303) circle (1.2pt);
\fill[blue!80!black] (7.4879,5.2299) circle (1.2pt);
\fill[blue!80!black] (7.4873,5.2299) circle (1.2pt);
\fill[blue!80!black] (7.4876,5.2295) circle (1.2pt);
\fill[blue!80!black] (7.4871,5.2295) circle (1.2pt);
\fill[blue!80!black] (7.4873,5.2292) circle (1.2pt);
\fill[blue!80!black] (7.4866,5.2291) circle (1.2pt);
\fill[blue!80!black] (7.4870,5.2288) circle (1.2pt);
\fill[blue!80!black] (7.4864,5.2286) circle (1.2pt);
\fill[blue!80!black] (7.4867,5.2284) circle (1.2pt);
\fill[blue!80!black] (7.4858,5.2280) circle (1.2pt);
\fill[blue!80!black] (7.4864,5.2277) circle (1.2pt);
\fill[blue!80!black] (7.4857,5.2276) circle (1.2pt);
\fill[blue!80!black] (7.4861,5.2274) circle (1.2pt);
\fill[blue!80!black] (7.4856,5.2272) circle (1.2pt);
\fill[blue!80!black] (7.4858,5.2270) circle (1.2pt);
\fill[blue!80!black] (7.4850,5.2266) circle (1.2pt);
\fill[blue!80!black] (7.4855,5.2263) circle (1.2pt);
\fill[blue!80!black] (7.4849,5.2262) circle (1.2pt);
\fill[blue!80!black] (7.4852,5.2260) circle (1.2pt);
\fill[blue!80!black] (7.4848,5.2258) circle (1.2pt);
\fill[blue!80!black] (7.4852,5.2254) circle (1.2pt);
\fill[blue!80!black] (7.4846,5.2254) circle (1.2pt);
\fill[blue!80!black] (7.4848,5.2250) circle (1.2pt);
\fill[blue!80!black] (7.4844,5.2250) circle (1.2pt);
\fill[blue!80!black] (7.4846,5.2243) circle (1.2pt);
\fill[blue!80!black] (7.4839,5.2244) circle (1.2pt);
\fill[blue!80!black] (7.4843,5.2240) circle (1.2pt);
\fill[blue!80!black] (7.4838,5.2240) circle (1.2pt);
\fill[blue!80!black] (7.4839,5.2236) circle (1.2pt);
\fill[blue!80!black] (7.4832,5.2235) circle (1.2pt);
\fill[blue!80!black] (7.4836,5.2232) circle (1.2pt);
\fill[blue!80!black] (7.4831,5.2231) circle (1.2pt);
\fill[blue!80!black] (7.4834,5.2229) circle (1.2pt);
\fill[blue!80!black] (7.4829,5.2227) circle (1.2pt);
\fill[blue!80!black] (7.4833,5.2223) circle (1.2pt);
\fill[blue!80!black] (7.4827,5.2223) circle (1.2pt);
\fill[blue!80!black] (7.4829,5.2219) circle (1.2pt);
\fill[blue!80!black] (7.4825,5.2219) circle (1.2pt);
\fill[blue!80!black] (7.4827,5.2212) circle (1.2pt);
\fill[blue!80!black] (7.4821,5.2213) circle (1.2pt);
\fill[blue!80!black] (7.4824,5.2209) circle (1.2pt);
\fill[blue!80!black] (7.4819,5.2209) circle (1.2pt);
\fill[blue!80!black] (7.4821,5.2205) circle (1.2pt);
\fill[blue!80!black] (7.4813,5.2204) circle (1.2pt);
\fill[blue!80!black] (7.4818,5.2202) circle (1.2pt);
\fill[blue!80!black] (7.4812,5.2200) circle (1.2pt);
\fill[blue!80!black] (7.4815,5.2198) circle (1.2pt);
\fill[blue!80!black] (7.4807,5.2194) circle (1.2pt);
\fill[blue!80!black] (7.4812,5.2191) circle (1.2pt);
\fill[blue!80!black] (7.4806,5.2190) circle (1.2pt);
\fill[blue!80!black] (7.4809,5.2188) circle (1.2pt);
\fill[blue!80!black] (7.4804,5.2186) circle (1.2pt);
\fill[blue!80!black] (7.4808,5.2182) circle (1.2pt);
\fill[blue!80!black] (7.4802,5.2182) circle (1.2pt);
\fill[blue!80!black] (7.4805,5.2179) circle (1.2pt);
\fill[blue!80!black] (7.4801,5.2178) circle (1.2pt);
\fill[blue!80!black] (7.4803,5.2172) circle (1.2pt);
\fill[blue!80!black] (7.4796,5.2172) circle (1.2pt);
\fill[blue!80!black] (7.4799,5.2168) circle (1.2pt);
\fill[blue!80!black] (7.4794,5.2168) circle (1.2pt);
\fill[blue!80!black] (7.4796,5.2165) circle (1.2pt);
\fill[blue!80!black] (7.4789,5.2164) circle (1.2pt);
\fill[blue!80!black] (7.4793,5.2161) circle (1.2pt);
\fill[blue!80!black] (7.4788,5.2160) circle (1.2pt);
\fill[blue!80!black] (7.4790,5.2157) circle (1.2pt);
\fill[blue!80!black] (7.4782,5.2154) circle (1.2pt);
\fill[blue!80!black] (7.4787,5.2151) circle (1.2pt);
\fill[blue!80!black] (7.4781,5.2150) circle (1.2pt);
\fill[blue!80!black] (7.4784,5.2147) circle (1.2pt);
\fill[blue!80!black] (7.4780,5.2145) circle (1.2pt);
\fill[blue!80!black] (7.4783,5.2142) circle (1.2pt);
\fill[blue!80!black] (7.4778,5.2142) circle (1.2pt);
\fill[blue!80!black] (7.4780,5.2138) circle (1.2pt);
\fill[blue!80!black] (7.4772,5.2137) circle (1.2pt);
\fill[blue!80!black] (7.4777,5.2134) circle (1.2pt);
\fill[blue!80!black] (7.4771,5.2133) circle (1.2pt);
\fill[blue!80!black] (7.4774,5.2131) circle (1.2pt);
\fill[blue!80!black] (7.4766,5.2127) circle (1.2pt);
\fill[blue!80!black] (7.4771,5.2124) circle (1.2pt);
\fill[blue!80!black] (7.4765,5.2123) circle (1.2pt);
\fill[blue!80!black] (7.4768,5.2121) circle (1.2pt);
\fill[blue!80!black] (7.4763,5.2119) circle (1.2pt);
\fill[blue!80!black] (7.4767,5.2115) circle (1.2pt);
\fill[blue!80!black] (7.4762,5.2115) circle (1.2pt);
\fill[blue!80!black] (7.4764,5.2112) circle (1.2pt);
\fill[blue!80!black] (7.4760,5.2111) circle (1.2pt);
\fill[blue!80!black] (7.4762,5.2105) circle (1.2pt);
\fill[blue!80!black] (7.4755,5.2105) circle (1.2pt);
\fill[blue!80!black] (7.4758,5.2101) circle (1.2pt);
\fill[blue!80!black] (7.4754,5.2101) circle (1.2pt);
\fill[blue!80!black] (7.4756,5.2095) circle (1.2pt);
\fill[blue!80!black] (7.4749,5.2095) circle (1.2pt);
\fill[blue!80!black] (7.4752,5.2091) circle (1.2pt);
\fill[blue!80!black] (7.4747,5.2091) circle (1.2pt);
\fill[blue!80!black] (7.4749,5.2088) circle (1.2pt);
\fill[blue!80!black] (7.4742,5.2087) circle (1.2pt);
\fill[blue!80!black] (7.4746,5.2084) circle (1.2pt);
\fill[blue!80!black] (7.4741,5.2083) circle (1.2pt);
\fill[blue!80!black] (7.4743,5.2081) circle (1.2pt);
\fill[blue!80!black] (7.4736,5.2077) circle (1.2pt);
\fill[blue!80!black] (7.4740,5.2074) circle (1.2pt);
\fill[blue!80!black] (7.4735,5.2073) circle (1.2pt);
\fill[blue!80!black] (7.4737,5.2071) circle (1.2pt);
\fill[blue!80!black] (7.4729,5.2067) circle (1.2pt);
\fill[blue!80!black] (7.4734,5.2064) circle (1.2pt);
\fill[blue!80!black] (7.4728,5.2063) circle (1.2pt);
\fill[blue!80!black] (7.4731,5.2061) circle (1.2pt);
\fill[blue!80!black] (7.4727,5.2059) circle (1.2pt);
\fill[blue!80!black] (7.4731,5.2055) circle (1.2pt);
\fill[blue!80!black] (7.4725,5.2055) circle (1.2pt);
\fill[blue!80!black] (7.4727,5.2052) circle (1.2pt);
\fill[blue!80!black] (7.4720,5.2051) circle (1.2pt);
\fill[blue!80!black] (7.4724,5.2048) circle (1.2pt);
\fill[blue!80!black] (7.4719,5.2047) circle (1.2pt);
\fill[blue!80!black] (7.4721,5.2045) circle (1.2pt);
\fill[blue!80!black] (7.4713,5.2041) circle (1.2pt);
\fill[blue!80!black] (7.4718,5.2038) circle (1.2pt);
\fill[blue!80!black] (7.4713,5.2037) circle (1.2pt);
\fill[blue!80!black] (7.4715,5.2035) circle (1.2pt);
\fill[blue!80!black] (7.4707,5.2031) circle (1.2pt);
\fill[blue!80!black] (7.4712,5.2028) circle (1.2pt);
\fill[blue!80!black] (7.4706,5.2027) circle (1.2pt);
\fill[blue!80!black] (7.4709,5.2025) circle (1.2pt);
\fill[blue!80!black] (7.4705,5.2023) circle (1.2pt);
\fill[blue!80!black] (7.4709,5.2019) circle (1.2pt);
\fill[blue!80!black] (7.4704,5.2019) circle (1.2pt);
\fill[blue!80!black] (7.4705,5.2016) circle (1.2pt);
\fill[blue!80!black] (7.4698,5.2015) circle (1.2pt);
\fill[blue!80!black] (7.4702,5.2012) circle (1.2pt);
\fill[blue!80!black] (7.4697,5.2011) circle (1.2pt);
\fill[blue!80!black] (7.4699,5.2009) circle (1.2pt);
\fill[blue!80!black] (7.4692,5.2005) circle (1.2pt);
\fill[blue!80!black] (7.4696,5.2003) circle (1.2pt);
\fill[blue!80!black] (7.4691,5.2001) circle (1.2pt);
\fill[blue!80!black] (7.4694,5.1999) circle (1.2pt);
\fill[blue!80!black] (7.4686,5.1995) circle (1.2pt);
\fill[blue!80!black] (7.4690,5.1993) circle (1.2pt);
\fill[blue!80!black] (7.4685,5.1991) circle (1.2pt);
\fill[blue!80!black] (7.4688,5.1989) circle (1.2pt);
\fill[blue!80!black] (7.4680,5.1986) circle (1.2pt);
\fill[blue!80!black] (7.4684,5.1983) circle (1.2pt);
\fill[blue!80!black] (7.4679,5.1982) circle (1.2pt);
\fill[blue!80!black] (7.4682,5.1979) circle (1.2pt);
\fill[blue!80!black] (7.4673,5.1976) circle (1.2pt);
\fill[blue!80!black] (7.4678,5.1973) circle (1.2pt);
\fill[blue!80!black] (7.4673,5.1972) circle (1.2pt);
\fill[blue!80!black] (7.4676,5.1970) circle (1.2pt);
\fill[blue!80!black] (7.4667,5.1966) circle (1.2pt);
\fill[blue!80!black] (7.4673,5.1963) circle (1.2pt);
\fill[blue!80!black] (7.4667,5.1962) circle (1.2pt);
\fill[blue!80!black] (7.4670,5.1960) circle (1.2pt);
\fill[blue!80!black] (7.4661,5.1956) circle (1.2pt);
\fill[blue!80!black] (7.4667,5.1954) circle (1.2pt);
\fill[blue!80!black] (7.4661,5.1952) circle (1.2pt);
\fill[blue!80!black] (7.4664,5.1950) circle (1.2pt);
\fill[blue!80!black] (7.4655,5.1947) circle (1.2pt);
\fill[blue!80!black] (7.4661,5.1944) circle (1.2pt);
\fill[blue!80!black] (7.4655,5.1943) circle (1.2pt);
\fill[blue!80!black] (7.4658,5.1940) circle (1.2pt);
\fill[blue!80!black] (7.4654,5.1939) circle (1.2pt);
\fill[blue!80!black] (7.4657,5.1935) circle (1.2pt);
\fill[blue!80!black] (7.4652,5.1935) circle (1.2pt);
\fill[blue!80!black] (7.4653,5.1932) circle (1.2pt);
\fill[blue!80!black] (7.4647,5.1931) circle (1.2pt);
\fill[blue!80!black] (7.4650,5.1928) circle (1.2pt);
\fill[blue!80!black] (7.4646,5.1927) circle (1.2pt);
\fill[blue!80!black] (7.4650,5.1923) circle (1.2pt);
\fill[blue!80!black] (7.4644,5.1923) circle (1.2pt);
\fill[blue!80!black] (7.4646,5.1920) circle (1.2pt);
\fill[blue!80!black] (7.4639,5.1919) circle (1.2pt);
\fill[blue!80!black] (7.4643,5.1916) circle (1.2pt);
\fill[blue!80!black] (7.4638,5.1915) circle (1.2pt);
\fill[blue!80!black] (7.4641,5.1913) circle (1.2pt);
\fill[blue!80!black] (7.4633,5.1909) circle (1.2pt);
\fill[blue!80!black] (7.4637,5.1906) circle (1.2pt);
\fill[blue!80!black] (7.4632,5.1905) circle (1.2pt);
\fill[blue!80!black] (7.4635,5.1903) circle (1.2pt);
\fill[blue!80!black] (7.4627,5.1899) circle (1.2pt);
\fill[blue!80!black] (7.4631,5.1897) circle (1.2pt);
\fill[blue!80!black] (7.4626,5.1895) circle (1.2pt);
\fill[blue!80!black] (7.4631,5.1891) circle (1.2pt);
\fill[blue!80!black] (7.4625,5.1892) circle (1.2pt);
\fill[blue!80!black] (7.4627,5.1888) circle (1.2pt);
\fill[blue!80!black] (7.4623,5.1888) circle (1.2pt);
\fill[blue!80!black] (7.4625,5.1882) circle (1.2pt);
\fill[blue!80!black] (7.4619,5.1882) circle (1.2pt);
\fill[blue!80!black] (7.4621,5.1879) circle (1.2pt);
\fill[blue!80!black] (7.4617,5.1878) circle (1.2pt);
\fill[blue!80!black] (7.4619,5.1872) circle (1.2pt);
\fill[blue!80!black] (7.4613,5.1872) circle (1.2pt);
\fill[blue!80!black] (7.4615,5.1869) circle (1.2pt);
\fill[blue!80!black] (7.4608,5.1868) circle (1.2pt);
\fill[blue!80!black] (7.4612,5.1866) circle (1.2pt);
\fill[blue!80!black] (7.4607,5.1864) circle (1.2pt);
\fill[blue!80!black] (7.4609,5.1862) circle (1.2pt);
\fill[blue!80!black] (7.4602,5.1859) circle (1.2pt);
\fill[blue!80!black] (7.4606,5.1856) circle (1.2pt);
\fill[blue!80!black] (7.4601,5.1855) circle (1.2pt);
\fill[blue!80!black] (7.4606,5.1851) circle (1.2pt);
\fill[blue!80!black] (7.4600,5.1851) circle (1.2pt);
\fill[blue!80!black] (7.4602,5.1848) circle (1.2pt);
\fill[blue!80!black] (7.4595,5.1847) circle (1.2pt);
\fill[blue!80!black] (7.4599,5.1844) circle (1.2pt);
\fill[blue!80!black] (7.4594,5.1843) circle (1.2pt);
\fill[blue!80!black] (7.4596,5.1841) circle (1.2pt);
\fill[blue!80!black] (7.4589,5.1837) circle (1.2pt);
\fill[blue!80!black] (7.4593,5.1835) circle (1.2pt);
\fill[blue!80!black] (7.4588,5.1833) circle (1.2pt);
\fill[blue!80!black] (7.4592,5.1829) circle (1.2pt);
\fill[blue!80!black] (7.4587,5.1829) circle (1.2pt);
\fill[blue!80!black] (7.4589,5.1826) circle (1.2pt);
\fill[blue!80!black] (7.4582,5.1825) circle (1.2pt);
\fill[blue!80!black] (7.4586,5.1823) circle (1.2pt);
\fill[blue!80!black] (7.4581,5.1821) circle (1.2pt);
\fill[blue!80!black] (7.4585,5.1817) circle (1.2pt);
\fill[blue!80!black] (7.4580,5.1818) circle (1.2pt);
\fill[blue!80!black] (7.4582,5.1814) circle (1.2pt);
\fill[blue!80!black] (7.4574,5.1814) circle (1.2pt);
\fill[blue!80!black] (7.4579,5.1811) circle (1.2pt);
\fill[blue!80!black] (7.4574,5.1810) circle (1.2pt);
\fill[blue!80!black] (7.4578,5.1806) circle (1.2pt);
\fill[blue!80!black] (7.4572,5.1806) circle (1.2pt);
\fill[blue!80!black] (7.4574,5.1802) circle (1.2pt);
\fill[blue!80!black] (7.4571,5.1802) circle (1.2pt);
\fill[blue!80!black] (7.4572,5.1796) circle (1.2pt);
\fill[blue!80!black] (7.4567,5.1796) circle (1.2pt);
\fill[blue!80!black] (7.4568,5.1793) circle (1.2pt);
\fill[blue!80!black] (7.4561,5.1792) circle (1.2pt);
\fill[blue!80!black] (7.4565,5.1790) circle (1.2pt);
\fill[blue!80!black] (7.4561,5.1788) circle (1.2pt);
\fill[blue!80!black] (7.4565,5.1784) circle (1.2pt);
\fill[blue!80!black] (7.4559,5.1785) circle (1.2pt);
\fill[blue!80!black] (7.4561,5.1781) circle (1.2pt);
\fill[blue!80!black] (7.4554,5.1780) circle (1.2pt);
\fill[blue!80!black] (7.4558,5.1778) circle (1.2pt);
\fill[blue!80!black] (7.4553,5.1777) circle (1.2pt);
\fill[blue!80!black] (7.4557,5.1773) circle (1.2pt);
\fill[blue!80!black] (7.4552,5.1773) circle (1.2pt);
\fill[blue!80!black] (7.4554,5.1770) circle (1.2pt);
\fill[blue!80!black] (7.4547,5.1769) circle (1.2pt);
\fill[blue!80!black] (7.4551,5.1766) circle (1.2pt);
\fill[blue!80!black] (7.4546,5.1765) circle (1.2pt);
\fill[blue!80!black] (7.4550,5.1761) circle (1.2pt);
\fill[blue!80!black] (7.4545,5.1761) circle (1.2pt);
\fill[blue!80!black] (7.4546,5.1758) circle (1.2pt);
\fill[blue!80!black] (7.4540,5.1757) circle (1.2pt);
\fill[blue!80!black] (7.4544,5.1755) circle (1.2pt);
\fill[blue!80!black] (7.4539,5.1753) circle (1.2pt);
\fill[blue!80!black] (7.4543,5.1749) circle (1.2pt);
\fill[blue!80!black] (7.4538,5.1749) circle (1.2pt);
\fill[blue!80!black] (7.4539,5.1746) circle (1.2pt);
\fill[blue!80!black] (7.4533,5.1745) circle (1.2pt);
\fill[blue!80!black] (7.4536,5.1743) circle (1.2pt);
\fill[blue!80!black] (7.4532,5.1741) circle (1.2pt);
\fill[blue!80!black] (7.4535,5.1738) circle (1.2pt);
\fill[blue!80!black] (7.4531,5.1738) circle (1.2pt);
\fill[blue!80!black] (7.4532,5.1735) circle (1.2pt);
\fill[blue!80!black] (7.4526,5.1734) circle (1.2pt);
\fill[blue!80!black] (7.4529,5.1731) circle (1.2pt);
\fill[blue!80!black] (7.4525,5.1730) circle (1.2pt);
\fill[blue!80!black] (7.4528,5.1726) circle (1.2pt);
\fill[blue!80!black] (7.4523,5.1726) circle (1.2pt);
\fill[blue!80!black] (7.4526,5.1720) circle (1.2pt);
\fill[blue!80!black] (7.4519,5.1720) circle (1.2pt);
\fill[blue!80!black] (7.4522,5.1717) circle (1.2pt);
\fill[blue!80!black] (7.4518,5.1717) circle (1.2pt);
\fill[blue!80!black] (7.4519,5.1711) circle (1.2pt);
\fill[blue!80!black] (7.4514,5.1711) circle (1.2pt);
\fill[blue!80!black] (7.4516,5.1708) circle (1.2pt);
\fill[blue!80!black] (7.4509,5.1707) circle (1.2pt);
\fill[blue!80!black] (7.4513,5.1705) circle (1.2pt);
\fill[blue!80!black] (7.4508,5.1703) circle (1.2pt);
\fill[blue!80!black] (7.4512,5.1699) circle (1.2pt);
\fill[blue!80!black] (7.4507,5.1700) circle (1.2pt);
\fill[blue!80!black] (7.4508,5.1696) circle (1.2pt);
\fill[blue!80!black] (7.4502,5.1696) circle (1.2pt);
\fill[blue!80!black] (7.4505,5.1693) circle (1.2pt);
\fill[blue!80!black] (7.4501,5.1692) circle (1.2pt);
\fill[blue!80!black] (7.4505,5.1688) circle (1.2pt);
\fill[blue!80!black] (7.4500,5.1688) circle (1.2pt);
\fill[blue!80!black] (7.4501,5.1685) circle (1.2pt);
\fill[blue!80!black] (7.4495,5.1684) circle (1.2pt);
\fill[blue!80!black] (7.4498,5.1682) circle (1.2pt);
\fill[blue!80!black] (7.4490,5.1679) circle (1.2pt);
\fill[blue!80!black] (7.4495,5.1676) circle (1.2pt);
\fill[blue!80!black] (7.4490,5.1675) circle (1.2pt);
\fill[blue!80!black] (7.4492,5.1673) circle (1.2pt);
\fill[blue!80!black] (7.4485,5.1669) circle (1.2pt);
\fill[blue!80!black] (7.4489,5.1667) circle (1.2pt);
\fill[blue!80!black] (7.4485,5.1665) circle (1.2pt);
\fill[blue!80!black] (7.4488,5.1662) circle (1.2pt);
\fill[blue!80!black] (7.4483,5.1662) circle (1.2pt);
\fill[blue!80!black] (7.4485,5.1659) circle (1.2pt);
\fill[blue!80!black] (7.4479,5.1658) circle (1.2pt);
\fill[blue!80!black] (7.4482,5.1655) circle (1.2pt);
\fill[blue!80!black] (7.4478,5.1654) circle (1.2pt);
\fill[blue!80!black] (7.4481,5.1650) circle (1.2pt);
\fill[blue!80!black] (7.4476,5.1650) circle (1.2pt);
\fill[blue!80!black] (7.4479,5.1644) circle (1.2pt);
\fill[blue!80!black] (7.4472,5.1645) circle (1.2pt);
\fill[blue!80!black] (7.4475,5.1641) circle (1.2pt);
\fill[blue!80!black] (7.4471,5.1641) circle (1.2pt);
\fill[blue!80!black] (7.4472,5.1635) circle (1.2pt);
\fill[blue!80!black] (7.4467,5.1635) circle (1.2pt);
\fill[blue!80!black] (7.4469,5.1632) circle (1.2pt);
\fill[blue!80!black] (7.4462,5.1631) circle (1.2pt);
\fill[blue!80!black] (7.4466,5.1629) circle (1.2pt);
\fill[blue!80!black] (7.4462,5.1628) circle (1.2pt);
\fill[blue!80!black] (7.4465,5.1624) circle (1.2pt);
\fill[blue!80!black] (7.4460,5.1624) circle (1.2pt);
\fill[blue!80!black] (7.4462,5.1618) circle (1.2pt);
\fill[blue!80!black] (7.4456,5.1618) circle (1.2pt);
\fill[blue!80!black] (7.4458,5.1615) circle (1.2pt);
\fill[blue!80!black] (7.4451,5.1615) circle (1.2pt);
\fill[blue!80!black] (7.4455,5.1612) circle (1.2pt);
\fill[blue!80!black] (7.4451,5.1611) circle (1.2pt);
\fill[blue!80!black] (7.4454,5.1607) circle (1.2pt);
\fill[blue!80!black] (7.4449,5.1607) circle (1.2pt);
\fill[blue!80!black] (7.4451,5.1604) circle (1.2pt);
\fill[blue!80!black] (7.4445,5.1603) circle (1.2pt);
\fill[blue!80!black] (7.4448,5.1601) circle (1.2pt);
\fill[blue!80!black] (7.4440,5.1598) circle (1.2pt);
\fill[blue!80!black] (7.4445,5.1595) circle (1.2pt);
\fill[blue!80!black] (7.4440,5.1594) circle (1.2pt);
\fill[blue!80!black] (7.4442,5.1592) circle (1.2pt);
\fill[blue!80!black] (7.4435,5.1588) circle (1.2pt);
\fill[blue!80!black] (7.4439,5.1586) circle (1.2pt);
\fill[blue!80!black] (7.4435,5.1585) circle (1.2pt);
\fill[blue!80!black] (7.4438,5.1581) circle (1.2pt);
\fill[blue!80!black] (7.4433,5.1581) circle (1.2pt);
\fill[blue!80!black] (7.4436,5.1575) circle (1.2pt);
\fill[blue!80!black] (7.4430,5.1576) circle (1.2pt);
\fill[blue!80!black] (7.4432,5.1572) circle (1.2pt);
\fill[blue!80!black] (7.4428,5.1572) circle (1.2pt);
\fill[blue!80!black] (7.4429,5.1566) circle (1.2pt);
\fill[blue!80!black] (7.4424,5.1566) circle (1.2pt);
\fill[blue!80!black] (7.4426,5.1563) circle (1.2pt);
\fill[blue!80!black] (7.4420,5.1562) circle (1.2pt);
\fill[blue!80!black] (7.4423,5.1560) circle (1.2pt);
\fill[blue!80!black] (7.4415,5.1557) circle (1.2pt);
\fill[blue!80!black] (7.4420,5.1555) circle (1.2pt);
\fill[blue!80!black] (7.4415,5.1553) circle (1.2pt);
\fill[blue!80!black] (7.4419,5.1550) circle (1.2pt);
\fill[blue!80!black] (7.4414,5.1550) circle (1.2pt);
\fill[blue!80!black] (7.4415,5.1547) circle (1.2pt);
\fill[blue!80!black] (7.4409,5.1546) circle (1.2pt);
\fill[blue!80!black] (7.4412,5.1544) circle (1.2pt);
\fill[blue!80!black] (7.4404,5.1541) circle (1.2pt);
\fill[blue!80!black] (7.4409,5.1538) circle (1.2pt);
\fill[blue!80!black] (7.4404,5.1537) circle (1.2pt);
\fill[blue!80!black] (7.4406,5.1535) circle (1.2pt);
\fill[blue!80!black] (7.4400,5.1531) circle (1.2pt);
\fill[blue!80!black] (7.4403,5.1529) circle (1.2pt);
\fill[blue!80!black] (7.4399,5.1528) circle (1.2pt);
\fill[blue!80!black] (7.4402,5.1524) circle (1.2pt);
\fill[blue!80!black] (7.4398,5.1524) circle (1.2pt);
\fill[blue!80!black] (7.4400,5.1518) circle (1.2pt);
\fill[blue!80!black] (7.4394,5.1519) circle (1.2pt);
\fill[blue!80!black] (7.4396,5.1516) circle (1.2pt);
\fill[blue!80!black] (7.4389,5.1515) circle (1.2pt);
\fill[blue!80!black] (7.4393,5.1512) circle (1.2pt);
\fill[blue!80!black] (7.4389,5.1511) circle (1.2pt);
\fill[blue!80!black] (7.4392,5.1508) circle (1.2pt);
\fill[blue!80!black] (7.4388,5.1507) circle (1.2pt);
\fill[blue!80!black] (7.4389,5.1502) circle (1.2pt);
\fill[blue!80!black] (7.4384,5.1502) circle (1.2pt);
\fill[blue!80!black] (7.4386,5.1499) circle (1.2pt);
\fill[blue!80!black] (7.4379,5.1498) circle (1.2pt);
\fill[blue!80!black] (7.4383,5.1496) circle (1.2pt);
\fill[blue!80!black] (7.4378,5.1494) circle (1.2pt);
\fill[blue!80!black] (7.4382,5.1491) circle (1.2pt);
\fill[blue!80!black] (7.4377,5.1491) circle (1.2pt);
\fill[blue!80!black] (7.4379,5.1485) circle (1.2pt);
\fill[blue!80!black] (7.4373,5.1486) circle (1.2pt);
\fill[blue!80!black] (7.4375,5.1482) circle (1.2pt);
\fill[blue!80!black] (7.4369,5.1482) circle (1.2pt);
\fill[blue!80!black] (7.4372,5.1479) circle (1.2pt);
\fill[blue!80!black] (7.4368,5.1478) circle (1.2pt);
\fill[blue!80!black] (7.4371,5.1475) circle (1.2pt);
\fill[blue!80!black] (7.4367,5.1474) circle (1.2pt);
\fill[blue!80!black] (7.4369,5.1469) circle (1.2pt);
\fill[blue!80!black] (7.4363,5.1469) circle (1.2pt);
\fill[blue!80!black] (7.4365,5.1466) circle (1.2pt);
\fill[blue!80!black] (7.4359,5.1465) circle (1.2pt);
\fill[blue!80!black] (7.4362,5.1463) circle (1.2pt);
\fill[blue!80!black] (7.4358,5.1461) circle (1.2pt);
\fill[blue!80!black] (7.4361,5.1458) circle (1.2pt);
\fill[blue!80!black] (7.4357,5.1458) circle (1.2pt);
\fill[blue!80!black] (7.4358,5.1452) circle (1.2pt);
\fill[blue!80!black] (7.4353,5.1453) circle (1.2pt);
\fill[blue!80!black] (7.4355,5.1450) circle (1.2pt);
\fill[blue!80!black] (7.4348,5.1449) circle (1.2pt);
\fill[blue!80!black] (7.4352,5.1447) circle (1.2pt);
\fill[blue!80!black] (7.4344,5.1444) circle (1.2pt);
\fill[blue!80!black] (7.4349,5.1441) circle (1.2pt);
\fill[blue!80!black] (7.4344,5.1440) circle (1.2pt);
\fill[blue!80!black] (7.4348,5.1436) circle (1.2pt);
\fill[blue!80!black] (7.4343,5.1436) circle (1.2pt);
\fill[blue!80!black] (7.4344,5.1433) circle (1.2pt);
\fill[blue!80!black] (7.4338,5.1432) circle (1.2pt);
\fill[blue!80!black] (7.4341,5.1430) circle (1.2pt);
\fill[blue!80!black] (7.4334,5.1427) circle (1.2pt);
\fill[blue!80!black] (7.4338,5.1425) circle (1.2pt);
\fill[blue!80!black] (7.4334,5.1423) circle (1.2pt);
\fill[blue!80!black] (7.4337,5.1420) circle (1.2pt);
\fill[blue!80!black] (7.4333,5.1420) circle (1.2pt);
\fill[blue!80!black] (7.4335,5.1414) circle (1.2pt);
\fill[blue!80!black] (7.4329,5.1415) circle (1.2pt);
\fill[blue!80!black] (7.4331,5.1412) circle (1.2pt);
\fill[blue!80!black] (7.4324,5.1411) circle (1.2pt);
\fill[blue!80!black] (7.4328,5.1409) circle (1.2pt);
\fill[blue!80!black] (7.4324,5.1407) circle (1.2pt);
\fill[blue!80!black] (7.4327,5.1404) circle (1.2pt);
\fill[blue!80!black] (7.4323,5.1404) circle (1.2pt);
\fill[blue!80!black] (7.4324,5.1398) circle (1.2pt);
\fill[blue!80!black] (7.4319,5.1398) circle (1.2pt);
\fill[blue!80!black] (7.4320,5.1395) circle (1.2pt);
\fill[blue!80!black] (7.4315,5.1395) circle (1.2pt);
\fill[blue!80!black] (7.4318,5.1392) circle (1.2pt);
\fill[blue!80!black] (7.4310,5.1389) circle (1.2pt);
\fill[blue!80!black] (7.4314,5.1387) circle (1.2pt);
\fill[blue!80!black] (7.4310,5.1386) circle (1.2pt);
\fill[blue!80!black] (7.4314,5.1382) circle (1.2pt);
\fill[blue!80!black] (7.4309,5.1382) circle (1.2pt);
\fill[blue!80!black] (7.4311,5.1377) circle (1.2pt);
\fill[blue!80!black] (7.4305,5.1377) circle (1.2pt);
\fill[blue!80!black] (7.4307,5.1374) circle (1.2pt);
\fill[blue!80!black] (7.4301,5.1373) circle (1.2pt);
\fill[blue!80!black] (7.4304,5.1371) circle (1.2pt);
\fill[blue!80!black] (7.4297,5.1368) circle (1.2pt);
\fill[blue!80!black] (7.4301,5.1366) circle (1.2pt);
\fill[blue!80!black] (7.4297,5.1364) circle (1.2pt);
\fill[blue!80!black] (7.4300,5.1361) circle (1.2pt);
\fill[blue!80!black] (7.4296,5.1361) circle (1.2pt);
\fill[blue!80!black] (7.4298,5.1355) circle (1.2pt);
\fill[blue!80!black] (7.4292,5.1356) circle (1.2pt);
\fill[blue!80!black] (7.4294,5.1353) circle (1.2pt);
\fill[blue!80!black] (7.4287,5.1352) circle (1.2pt);
\fill[blue!80!black] (7.4291,5.1350) circle (1.2pt);
\fill[blue!80!black] (7.4287,5.1348) circle (1.2pt);
\fill[blue!80!black] (7.4290,5.1345) circle (1.2pt);
\fill[blue!80!black] (7.4286,5.1345) circle (1.2pt);
\fill[blue!80!black] (7.4287,5.1339) circle (1.2pt);
\fill[blue!80!black] (7.4282,5.1340) circle (1.2pt);
\fill[blue!80!black] (7.4283,5.1337) circle (1.2pt);
\fill[blue!80!black] (7.4278,5.1336) circle (1.2pt);
\fill[blue!80!black] (7.4281,5.1334) circle (1.2pt);
\fill[blue!80!black] (7.4274,5.1331) circle (1.2pt);
\fill[blue!80!black] (7.4277,5.1328) circle (1.2pt);
\fill[blue!80!black] (7.4274,5.1327) circle (1.2pt);
\fill[blue!80!black] (7.4276,5.1324) circle (1.2pt);
\fill[blue!80!black] (7.4272,5.1324) circle (1.2pt);
\fill[blue!80!black] (7.4274,5.1318) circle (1.2pt);
\fill[blue!80!black] (7.4269,5.1318) circle (1.2pt);
\fill[blue!80!black] (7.4270,5.1316) circle (1.2pt);
\fill[blue!80!black] (7.4265,5.1315) circle (1.2pt);
\fill[blue!80!black] (7.4267,5.1313) circle (1.2pt);
\fill[blue!80!black] (7.4261,5.1309) circle (1.2pt);
\fill[blue!80!black] (7.4264,5.1307) circle (1.2pt);
\fill[blue!80!black] (7.4257,5.1304) circle (1.2pt);
\fill[blue!80!black] (7.4261,5.1302) circle (1.2pt);
\fill[blue!80!black] (7.4257,5.1301) circle (1.2pt);
\fill[blue!80!black] (7.4260,5.1297) circle (1.2pt);
\fill[blue!80!black] (7.4256,5.1297) circle (1.2pt);
\fill[blue!80!black] (7.4258,5.1292) circle (1.2pt);
\fill[blue!80!black] (7.4252,5.1292) circle (1.2pt);
\fill[blue!80!black] (7.4254,5.1289) circle (1.2pt);
\fill[blue!80!black] (7.4248,5.1288) circle (1.2pt);
\fill[blue!80!black] (7.4251,5.1286) circle (1.2pt);
\fill[blue!80!black] (7.4244,5.1283) circle (1.2pt);
\fill[blue!80!black] (7.4248,5.1281) circle (1.2pt);
\fill[blue!80!black] (7.4244,5.1280) circle (1.2pt);
\fill[blue!80!black] (7.4247,5.1276) circle (1.2pt);
\fill[blue!80!black] (7.4242,5.1276) circle (1.2pt);
\fill[blue!80!black] (7.4244,5.1271) circle (1.2pt);
\fill[blue!80!black] (7.4239,5.1271) circle (1.2pt);
\fill[blue!80!black] (7.4240,5.1268) circle (1.2pt);
\fill[blue!80!black] (7.4235,5.1267) circle (1.2pt);
\fill[blue!80!black] (7.4237,5.1265) circle (1.2pt);
\fill[blue!80!black] (7.4231,5.1262) circle (1.2pt);
\fill[blue!80!black] (7.4234,5.1260) circle (1.2pt);
\fill[blue!80!black] (7.4231,5.1259) circle (1.2pt);
\fill[blue!80!black] (7.4233,5.1256) circle (1.2pt);
\fill[blue!80!black] (7.4229,5.1255) circle (1.2pt);
\fill[blue!80!black] (7.4231,5.1250) circle (1.2pt);
\fill[blue!80!black] (7.4226,5.1250) circle (1.2pt);
\fill[blue!80!black] (7.4228,5.1245) circle (1.2pt);
\fill[blue!80!black] (7.4222,5.1245) circle (1.2pt);
\fill[blue!80!black] (7.4224,5.1242) circle (1.2pt);
\fill[blue!80!black] (7.4218,5.1241) circle (1.2pt);
\fill[blue!80!black] (7.4221,5.1239) circle (1.2pt);
\fill[blue!80!black] (7.4214,5.1236) circle (1.2pt);
\fill[blue!80!black] (7.4218,5.1234) circle (1.2pt);
\fill[blue!80!black] (7.4214,5.1233) circle (1.2pt);
\fill[blue!80!black] (7.4217,5.1230) circle (1.2pt);
\fill[blue!80!black] (7.4213,5.1229) circle (1.2pt);
\fill[blue!80!black] (7.4214,5.1224) circle (1.2pt);
\fill[blue!80!black] (7.4210,5.1224) circle (1.2pt);
\fill[blue!80!black] (7.4212,5.1219) circle (1.2pt);
\fill[blue!80!black] (7.4206,5.1219) circle (1.2pt);
\fill[blue!80!black] (7.4208,5.1216) circle (1.2pt);
\fill[blue!80!black] (7.4202,5.1216) circle (1.2pt);
\fill[blue!80!black] (7.4205,5.1214) circle (1.2pt);
\fill[blue!80!black] (7.4198,5.1211) circle (1.2pt);
\fill[blue!80!black] (7.4202,5.1208) circle (1.2pt);
\fill[blue!80!black] (7.4198,5.1207) circle (1.2pt);
\fill[blue!80!black] (7.4201,5.1204) circle (1.2pt);
\fill[blue!80!black] (7.4197,5.1204) circle (1.2pt);
\fill[blue!80!black] (7.4198,5.1199) circle (1.2pt);
\fill[blue!80!black] (7.4193,5.1199) circle (1.2pt);
\fill[blue!80!black] (7.4195,5.1193) circle (1.2pt);
\fill[blue!80!black] (7.4190,5.1194) circle (1.2pt);
\fill[blue!80!black] (7.4191,5.1191) circle (1.2pt);
\fill[blue!80!black] (7.4186,5.1190) circle (1.2pt);
\fill[blue!80!black] (7.4188,5.1188) circle (1.2pt);
\fill[blue!80!black] (7.4182,5.1185) circle (1.2pt);
\fill[blue!80!black] (7.4185,5.1183) circle (1.2pt);
\fill[blue!80!black] (7.4178,5.1180) circle (1.2pt);
\fill[blue!80!black] (7.4182,5.1178) circle (1.2pt);
\fill[blue!80!black] (7.4178,5.1176) circle (1.2pt);
\fill[blue!80!black] (7.4181,5.1173) circle (1.2pt);
\fill[blue!80!black] (7.4177,5.1173) circle (1.2pt);
\fill[blue!80!black] (7.4179,5.1168) circle (1.2pt);
\fill[blue!80!black] (7.4174,5.1168) circle (1.2pt);
\fill[blue!80!black] (7.4176,5.1163) circle (1.2pt);
\fill[blue!80!black] (7.4171,5.1163) circle (1.2pt);
\fill[blue!80!black] (7.4172,5.1160) circle (1.2pt);
\fill[blue!80!black] (7.4166,5.1159) circle (1.2pt);
\fill[blue!80!black] (7.4169,5.1157) circle (1.2pt);
\fill[blue!80!black] (7.4163,5.1154) circle (1.2pt);
\fill[blue!80!black] (7.4166,5.1152) circle (1.2pt);
\fill[blue!80!black] (7.4159,5.1149) circle (1.2pt);
\fill[blue!80!black] (7.4163,5.1147) circle (1.2pt);
\fill[blue!80!black] (7.4159,5.1146) circle (1.2pt);
\fill[blue!80!black] (7.4162,5.1143) circle (1.2pt);
\fill[blue!80!black] (7.4158,5.1143) circle (1.2pt);
\fill[blue!80!black] (7.4159,5.1137) circle (1.2pt);
\fill[blue!80!black] (7.4155,5.1138) circle (1.2pt);
\fill[blue!80!black] (7.4157,5.1132) circle (1.2pt);
\fill[blue!80!black] (7.4151,5.1133) circle (1.2pt);
\fill[blue!80!black] (7.4153,5.1130) circle (1.2pt);
\fill[blue!80!black] (7.4147,5.1129) circle (1.2pt);
\fill[blue!80!black] (7.4150,5.1127) circle (1.2pt);
\fill[blue!80!black] (7.4144,5.1124) circle (1.2pt);
\fill[blue!80!black] (7.4147,5.1122) circle (1.2pt);
\fill[blue!80!black] (7.4140,5.1119) circle (1.2pt);
\fill[blue!80!black] (7.4144,5.1117) circle (1.2pt);
\fill[blue!80!black] (7.4140,5.1116) circle (1.2pt);
\fill[blue!80!black] (7.4143,5.1113) circle (1.2pt);
\fill[blue!80!black] (7.4139,5.1112) circle (1.2pt);
\fill[blue!80!black] (7.4140,5.1107) circle (1.2pt);
\fill[blue!80!black] (7.4136,5.1107) circle (1.2pt);
\fill[blue!80!black] (7.4137,5.1102) circle (1.2pt);
\fill[blue!80!black] (7.4132,5.1102) circle (1.2pt);
\fill[blue!80!black] (7.4134,5.1097) circle (1.2pt);
\fill[blue!80!black] (7.4129,5.1097) circle (1.2pt);
\fill[blue!80!black] (7.4130,5.1095) circle (1.2pt);
\fill[blue!80!black] (7.4125,5.1094) circle (1.2pt);
\fill[blue!80!black] (7.4128,5.1092) circle (1.2pt);
\fill[blue!80!black] (7.4121,5.1089) circle (1.2pt);
\fill[blue!80!black] (7.4124,5.1087) circle (1.2pt);
\fill[blue!80!black] (7.4118,5.1084) circle (1.2pt);
\fill[blue!80!black] (7.4121,5.1082) circle (1.2pt);
\fill[blue!80!black] (7.4114,5.1079) circle (1.2pt);
\fill[blue!80!black] (7.4118,5.1077) circle (1.2pt);
\fill[blue!80!black] (7.4114,5.1076) circle (1.2pt);
\fill[blue!80!black] (7.4117,5.1073) circle (1.2pt);
\fill[blue!80!black] (7.4113,5.1072) circle (1.2pt);
\fill[blue!80!black] (7.4115,5.1068) circle (1.2pt);
\fill[blue!80!black] (7.4110,5.1068) circle (1.2pt);
\fill[blue!80!black] (7.4112,5.1062) circle (1.2pt);
\fill[blue!80!black] (7.4107,5.1063) circle (1.2pt);
\fill[blue!80!black] (7.4109,5.1057) circle (1.2pt);
\fill[blue!80!black] (7.4104,5.1058) circle (1.2pt);
\fill[blue!80!black] (7.4105,5.1055) circle (1.2pt);
\fill[blue!80!black] (7.4100,5.1054) circle (1.2pt);
\fill[blue!80!black] (7.4102,5.1052) circle (1.2pt);
\fill[blue!80!black] (7.4096,5.1049) circle (1.2pt);
\fill[blue!80!black] (7.4099,5.1047) circle (1.2pt);
\fill[blue!80!black] (7.4093,5.1045) circle (1.2pt);
\fill[blue!80!black] (7.4096,5.1042) circle (1.2pt);
\fill[blue!80!black] (7.4089,5.1040) circle (1.2pt);
\fill[blue!80!black] (7.4093,5.1037) circle (1.2pt);
\fill[blue!80!black] (7.4089,5.1036) circle (1.2pt);
\fill[blue!80!black] (7.4092,5.1033) circle (1.2pt);
\fill[blue!80!black] (7.4088,5.1033) circle (1.2pt);
\fill[blue!80!black] (7.4089,5.1028) circle (1.2pt);
\fill[blue!80!black] (7.4085,5.1028) circle (1.2pt);
\fill[blue!80!black] (7.4086,5.1023) circle (1.2pt);
\fill[blue!80!black] (7.4082,5.1023) circle (1.2pt);
\fill[blue!80!black] (7.4084,5.1018) circle (1.2pt);
\fill[blue!80!black] (7.4079,5.1018) circle (1.2pt);
\fill[blue!80!black] (7.4081,5.1013) circle (1.2pt);
\fill[blue!80!black] (7.4076,5.1014) circle (1.2pt);
\fill[blue!80!black] (7.4077,5.1011) circle (1.2pt);
\fill[blue!80!black] (7.4072,5.1010) circle (1.2pt);
\fill[blue!80!black] (7.4074,5.1008) circle (1.2pt);
\fill[blue!80!black] (7.4068,5.1005) circle (1.2pt);
\fill[blue!80!black] (7.4071,5.1003) circle (1.2pt);
\fill[blue!80!black] (7.4065,5.1000) circle (1.2pt);
\fill[blue!80!black] (7.4068,5.0998) circle (1.2pt);
\fill[blue!80!black] (7.4062,5.0996) circle (1.2pt);
\fill[blue!80!black] (7.4065,5.0994) circle (1.2pt);
\fill[blue!80!black] (7.4058,5.0991) circle (1.2pt);
\fill[blue!80!black] (7.4062,5.0989) circle (1.2pt);
\fill[blue!80!black] (7.4058,5.0987) circle (1.2pt);
\fill[blue!80!black] (7.4061,5.0984) circle (1.2pt);
\fill[blue!80!black] (7.4057,5.0984) circle (1.2pt);
\fill[blue!80!black] (7.4058,5.0980) circle (1.2pt);
\fill[blue!80!black] (7.4054,5.0979) circle (1.2pt);
\fill[blue!80!black] (7.4055,5.0975) circle (1.2pt);
\fill[blue!80!black] (7.4051,5.0975) circle (1.2pt);
\fill[blue!80!black] (7.4052,5.0970) circle (1.2pt);
\fill[blue!80!black] (7.4048,5.0970) circle (1.2pt);
\fill[blue!80!black] (7.4049,5.0965) circle (1.2pt);
\fill[blue!80!black] (7.4045,5.0965) circle (1.2pt);
\fill[blue!80!black] (7.4047,5.0960) circle (1.2pt);
\fill[blue!80!black] (7.4042,5.0960) circle (1.2pt);
\fill[blue!80!black] (7.4043,5.0958) circle (1.2pt);
\fill[blue!80!black] (7.4038,5.0957) circle (1.2pt);
\fill[blue!80!black] (7.4042,5.0953) circle (1.2pt);
\fill[blue!80!black] (7.4037,5.0954) circle (1.2pt);
\fill[blue!80!black] (7.4039,5.0948) circle (1.2pt);
\fill[blue!80!black] (7.4034,5.0949) circle (1.2pt);
\fill[blue!80!black] (7.4035,5.0946) circle (1.2pt);
\fill[blue!80!black] (7.4030,5.0945) circle (1.2pt);
\fill[blue!80!black] (7.4033,5.0944) circle (1.2pt);
\fill[blue!80!black] (7.4027,5.0941) circle (1.2pt);
\fill[blue!80!black] (7.4030,5.0939) circle (1.2pt);
\fill[blue!80!black] (7.4024,5.0936) circle (1.2pt);
\fill[blue!80!black] (7.4027,5.0934) circle (1.2pt);
\fill[blue!80!black] (7.4021,5.0931) circle (1.2pt);
\fill[blue!80!black] (7.4024,5.0929) circle (1.2pt);
\fill[blue!80!black] (7.4017,5.0926) circle (1.2pt);
\fill[blue!80!black] (7.4021,5.0924) circle (1.2pt);
\fill[blue!80!black] (7.4014,5.0922) circle (1.2pt);
\fill[blue!80!black] (7.4018,5.0920) circle (1.2pt);
\fill[blue!80!black] (7.4011,5.0917) circle (1.2pt);
\fill[blue!80!black] (7.4015,5.0915) circle (1.2pt);
\fill[blue!80!black] (7.4008,5.0912) circle (1.2pt);
\fill[blue!80!black] (7.4012,5.0910) circle (1.2pt);
\fill[blue!80!black] (7.4008,5.0909) circle (1.2pt);
\fill[blue!80!black] (7.4010,5.0906) circle (1.2pt);
\fill[blue!80!black] (7.4004,5.0906) circle (1.2pt);
\fill[blue!80!black] (7.4007,5.0903) circle (1.2pt);
\fill[blue!80!black] (7.4000,5.0901) circle (1.2pt);
\fill[blue!80!black] (7.4004,5.0899) circle (1.2pt);
\fill[blue!80!black] (7.3997,5.0896) circle (1.2pt);
\draw[red!80!black, line width=1.0pt] (3.6000,2.4000) -- (12.0000,8.0000) -- (0.0000,8.0000) -- (12.0000,8.0000) -- (0.0000,8.0000) -- (12.0000,8.0000) -- (0.0000,8.0000) -- (12.0000,8.0000) -- (0.0000,8.0000) -- (12.0000,8.0000) -- (0.0000,8.0000) -- (12.0000,8.0000) -- (0.0000,8.0000) -- (12.0000,8.0000) -- (0.0000,8.0000) -- (12.0000,8.0000) -- (0.0000,8.0000) -- (12.0000,8.0000) -- (0.0000,8.0000) -- (12.0000,8.0000) -- (0.0000,8.0000) -- (12.0000,8.0000) -- (0.0000,8.0000) -- (12.0000,8.0000) -- (0.0000,8.0000) -- (12.0000,8.0000) -- (0.0000,8.0000) -- (12.0000,8.0000) -- (0.0000,8.0000) -- (12.0000,8.0000) -- (0.0000,8.0000) -- (12.0000,8.0000) -- (0.0000,8.0000) -- (12.0000,8.0000) -- (0.0000,8.0000) -- (12.0000,8.0000) -- (0.0000,8.0000) -- (12.0000,8.0000) -- (0.0000,8.0000) -- (12.0000,8.0000) -- (0.0000,8.0000) -- (12.0000,8.0000) -- (0.0000,8.0000) -- (12.0000,8.0000) -- (0.0000,8.0000) -- (12.0000,8.0000) -- (0.0000,8.0000) -- (12.0000,8.0000) -- (0.0000,8.0000) -- (12.0000,8.0000) -- (0.0000,8.0000) -- (12.0000,8.0000) -- (0.0000,8.0000) -- (12.0000,8.0000) -- (0.0000,8.0000) -- (12.0000,8.0000) -- (0.0000,8.0000) -- (12.0000,8.0000) -- (0.0000,8.0000) -- (12.0000,8.0000) -- (0.0000,8.0000) -- (12.0000,8.0000) -- (0.0000,8.0000) -- (12.0000,8.0000) -- (0.0000,8.0000) -- (12.0000,8.0000) -- (0.0000,8.0000) -- (12.0000,8.0000) -- (0.0000,8.0000) -- (12.0000,8.0000) -- (0.0000,8.0000) -- (12.0000,8.0000) -- (0.0000,8.0000) -- (12.0000,8.0000) -- (0.0000,8.0000) -- (12.0000,8.0000) -- (0.0000,8.0000) -- (12.0000,8.0000) -- (0.0000,8.0000) -- (12.0000,8.0000) -- (0.0000,8.0000) -- (12.0000,8.0000) -- (0.0000,8.0000) -- (12.0000,8.0000) -- (0.0000,8.0000) -- (12.0000,8.0000) -- (0.0000,8.0000) -- (12.0000,8.0000) -- (0.0000,8.0000) -- (12.0000,8.0000) -- (0.0000,8.0000) -- (12.0000,8.0000) -- (0.0000,8.0000) -- (12.0000,8.0000) -- (0.0000,8.0000) -- (12.0000,8.0000) -- (0.0000,8.0000) -- (12.0000,8.0000) -- (0.0000,8.0000) -- (12.0000,8.0000) -- (0.0000,8.0000) -- (12.0000,8.0000) -- (0.0000,8.0000) -- (12.0000,8.0000) -- (0.0000,8.0000) -- (12.0000,8.0000) -- (0.0000,8.0000) -- (12.0000,8.0000) -- (0.0000,8.0000) -- (12.0000,8.0000) -- (0.0000,8.0000) -- (12.0000,8.0000) -- (0.0000,8.0000) -- (12.0000,8.0000) -- (0.0000,8.0000) -- (12.0000,8.0000) -- (0.0000,8.0000) -- (12.0000,8.0000) -- (0.0000,8.0000) -- (12.0000,8.0000) -- (0.0000,8.0000) -- (12.0000,8.0000) -- (0.0000,8.0000) -- (12.0000,8.0000) -- (0.0000,8.0000) -- (12.0000,8.0000) -- (0.0000,8.0000) -- (12.0000,8.0000) -- (0.0000,8.0000) -- (12.0000,8.0000) -- (0.0000,8.0000) -- (12.0000,8.0000) -- (0.0000,8.0000) -- (12.0000,8.0000) -- (0.0000,8.0000) -- (12.0000,8.0000) -- (0.0000,8.0000) -- (12.0000,8.0000) -- (0.0000,8.0000) -- (12.0000,8.0000) -- (0.0000,8.0000) -- (12.0000,8.0000) -- (0.0000,8.0000) -- (12.0000,8.0000) -- (0.0000,8.0000) -- (12.0000,8.0000) -- (0.0000,8.0000) -- (12.0000,8.0000) -- (0.0000,8.0000) -- (12.0000,8.0000) -- (0.0000,8.0000) -- (12.0000,8.0000) -- (0.0000,8.0000) -- (12.0000,8.0000) -- (0.0000,8.0000) -- (12.0000,8.0000) -- (0.0000,8.0000) -- (12.0000,8.0000) -- (0.0000,8.0000) -- (12.0000,8.0000) -- (0.0000,8.0000) -- (12.0000,8.0000) -- (0.0000,8.0000) -- (12.0000,8.0000) -- (0.0000,8.0000) -- (12.0000,8.0000) -- (0.0000,8.0000) -- (12.0000,8.0000) -- (0.0000,8.0000) -- (12.0000,8.0000) -- (0.0000,8.0000) -- (12.0000,8.0000) -- (0.0000,8.0000) -- (12.0000,8.0000) -- (0.0000,8.0000) -- (12.0000,8.0000) -- (0.0000,8.0000) -- (12.0000,8.0000) -- (0.0000,8.0000) -- (12.0000,8.0000) -- (0.0000,8.0000) -- (12.0000,8.0000) -- (0.0000,8.0000) -- (12.0000,8.0000) -- (0.0000,8.0000) -- (12.0000,8.0000) -- (0.0000,8.0000) -- (12.0000,8.0000) -- (0.0000,8.0000) -- (12.0000,8.0000) -- (0.0000,8.0000) -- (12.0000,8.0000) -- (0.0000,8.0000) -- (12.0000,8.0000) -- (0.0000,8.0000) -- (12.0000,8.0000) -- (0.0000,8.0000) -- (12.0000,8.0000) -- (0.0000,8.0000) -- (12.0000,8.0000) -- (0.0000,8.0000) -- (12.0000,8.0000) -- (0.0000,8.0000) -- (12.0000,8.0000) -- (0.0000,8.0000) -- (12.0000,8.0000) -- (0.0000,8.0000) -- (12.0000,8.0000) -- (0.0000,8.0000) -- (12.0000,8.0000) -- (0.0000,8.0000) -- (12.0000,8.0000) -- (0.0000,8.0000) -- (12.0000,8.0000) -- (0.0000,8.0000) -- (12.0000,8.0000) -- (0.0000,8.0000) -- (12.0000,8.0000) -- (0.0000,8.0000) -- (12.0000,8.0000) -- (0.0000,8.0000) -- (12.0000,8.0000) -- (0.0000,8.0000) -- (12.0000,8.0000) -- (0.0000,8.0000) -- (12.0000,8.0000) -- (0.0000,8.0000) -- (12.0000,8.0000) -- (0.0000,8.0000) -- (12.0000,8.0000) -- (0.0000,8.0000) -- (12.0000,8.0000) -- (0.0000,8.0000) -- (12.0000,8.0000) -- (0.0000,8.0000) -- (12.0000,8.0000) -- (0.0000,8.0000) -- (12.0000,8.0000) -- (0.0000,8.0000) -- (12.0000,8.0000) -- (0.0000,8.0000) -- (12.0000,8.0000) -- (0.0000,8.0000) -- (12.0000,8.0000) -- (0.0000,8.0000) -- (12.0000,8.0000) -- (0.0000,8.0000) -- (12.0000,8.0000) -- (0.0000,8.0000) -- (12.0000,8.0000) -- (0.0000,8.0000) -- (12.0000,8.0000) -- (0.0000,8.0000) -- (12.0000,8.0000) -- (0.0000,8.0000) -- (12.0000,8.0000) -- (0.0000,8.0000) -- (12.0000,8.0000) -- (0.0000,8.0000) -- (12.0000,8.0000) -- (0.0000,8.0000) -- (12.0000,8.0000) -- (0.0000,8.0000) -- (12.0000,8.0000) -- (0.0000,8.0000) -- (12.0000,8.0000) -- (0.0000,8.0000) -- (12.0000,8.0000) -- (0.0000,8.0000) -- (12.0000,8.0000) -- (0.0000,8.0000) -- (12.0000,8.0000) -- (0.0000,8.0000) -- (12.0000,8.0000) -- (0.0000,8.0000) -- (12.0000,8.0000) -- (0.0000,8.0000) -- (12.0000,8.0000) -- (0.0000,8.0000) -- (12.0000,8.0000) -- (0.0000,8.0000) -- (12.0000,8.0000) -- (0.0000,8.0000) -- (12.0000,8.0000) -- (0.0000,8.0000) -- (12.0000,8.0000) -- (0.0000,8.0000) -- (12.0000,8.0000) -- (0.0000,8.0000) -- (12.0000,8.0000) -- (0.0000,8.0000) -- (12.0000,8.0000) -- (0.0000,8.0000) -- (12.0000,8.0000) -- (0.0000,8.0000) -- (12.0000,8.0000) -- (0.0000,8.0000) -- (12.0000,8.0000) -- (0.0000,8.0000) -- (12.0000,8.0000) -- (0.0000,8.0000) -- (12.0000,8.0000) -- (0.0000,8.0000) -- (12.0000,8.0000) -- (0.0000,8.0000) -- (12.0000,8.0000) -- (0.0000,8.0000) -- (12.0000,8.0000) -- (0.0000,8.0000) -- (12.0000,8.0000) -- (0.0000,8.0000) -- (12.0000,8.0000) -- (0.0000,8.0000) -- (12.0000,8.0000) -- (0.0000,8.0000) -- (12.0000,8.0000) -- (0.0000,8.0000) -- (12.0000,8.0000) -- (0.0000,8.0000) -- (12.0000,8.0000) -- (0.0000,8.0000) -- (12.0000,8.0000) -- (0.0000,8.0000) -- (12.0000,8.0000) -- (0.0000,8.0000) -- (12.0000,8.0000) -- (0.0000,8.0000) -- (12.0000,8.0000) -- (0.0000,8.0000) -- (12.0000,8.0000) -- (0.0000,8.0000) -- (12.0000,8.0000) -- (0.0000,8.0000) -- (12.0000,8.0000) -- (0.0000,8.0000) -- (12.0000,8.0000) -- (0.0000,8.0000) -- (12.0000,8.0000) -- (0.0000,8.0000) -- (12.0000,8.0000) -- (0.0000,8.0000) -- (12.0000,8.0000) -- (0.0000,8.0000) -- (12.0000,8.0000) -- (0.0000,8.0000) -- (12.0000,8.0000) -- (0.0000,8.0000) -- (12.0000,8.0000) -- (0.0000,8.0000) -- (12.0000,8.0000) -- (0.0000,8.0000) -- (12.0000,8.0000) -- (0.0000,8.0000) -- (12.0000,8.0000) -- (0.0000,8.0000) -- (12.0000,8.0000) -- (0.0000,8.0000) -- (12.0000,8.0000) -- (0.0000,8.0000) -- (12.0000,8.0000) -- (0.0000,8.0000) -- (12.0000,8.0000) -- (0.0000,8.0000) -- (12.0000,8.0000) -- (0.0000,8.0000) -- (12.0000,8.0000) -- (0.0000,8.0000) -- (12.0000,8.0000) -- (0.0000,8.0000) -- (12.0000,8.0000) -- (0.0000,8.0000) -- (12.0000,8.0000) -- (0.0000,8.0000) -- (12.0000,8.0000) -- (0.0000,8.0000) -- (12.0000,8.0000) -- (0.0000,8.0000) -- (12.0000,8.0000) -- (0.0000,8.0000) -- (12.0000,8.0000) -- (0.0000,8.0000) -- (12.0000,8.0000) -- (0.0000,8.0000) -- (12.0000,8.0000) -- (0.0000,8.0000) -- (12.0000,8.0000) -- (0.0000,8.0000) -- (12.0000,8.0000) -- (0.0000,8.0000) -- (12.0000,8.0000) -- (0.0000,8.0000) -- (12.0000,8.0000) -- (0.0000,8.0000) -- (12.0000,8.0000) -- (0.0000,8.0000) -- (12.0000,8.0000) -- (0.0000,8.0000) -- (12.0000,8.0000) -- (0.0000,8.0000) -- (12.0000,8.0000) -- (0.0000,8.0000) -- (12.0000,8.0000) -- (0.0000,8.0000) -- (12.0000,8.0000) -- (0.0000,8.0000) -- (12.0000,8.0000) -- (0.0000,8.0000) -- (12.0000,8.0000) -- (0.0000,8.0000) -- (12.0000,8.0000) -- (0.0000,8.0000) -- (12.0000,8.0000) -- (0.0000,8.0000) -- (12.0000,8.0000) -- (0.0000,8.0000) -- (12.0000,8.0000) -- (0.0000,8.0000) -- (12.0000,8.0000) -- (0.0000,8.0000) -- (12.0000,8.0000) -- (0.0000,8.0000) -- (12.0000,8.0000) -- (0.0000,8.0000) -- (12.0000,8.0000) -- (0.0000,8.0000) -- (12.0000,8.0000) -- (0.0000,8.0000) -- (12.0000,8.0000) -- (0.0000,8.0000) -- (12.0000,8.0000) -- (0.0000,8.0000) -- (12.0000,8.0000) -- (0.0000,8.0000) -- (12.0000,8.0000) -- (0.0000,8.0000) -- (12.0000,8.0000) -- (0.0000,8.0000) -- (12.0000,8.0000) -- (0.0000,8.0000) -- (12.0000,8.0000) -- (0.0000,8.0000) -- (12.0000,8.0000) -- (0.0000,8.0000) -- (12.0000,8.0000) -- (0.0000,8.0000) -- (12.0000,8.0000) -- (0.0000,8.0000) -- (12.0000,8.0000) -- (0.0000,8.0000) -- (12.0000,8.0000) -- (0.0000,8.0000) -- (12.0000,8.0000) -- (0.0000,8.0000) -- (12.0000,8.0000) -- (0.0000,8.0000) -- (12.0000,8.0000) -- (0.0000,8.0000) -- (12.0000,8.0000) -- (0.0000,8.0000) -- (12.0000,8.0000) -- (0.0000,8.0000) -- (12.0000,8.0000) -- (0.0000,8.0000) -- (12.0000,8.0000) -- (0.0000,8.0000) -- (12.0000,8.0000) -- (0.0000,8.0000) -- (12.0000,8.0000) -- (0.0000,8.0000) -- (12.0000,8.0000) -- (0.0000,8.0000) -- (12.0000,8.0000) -- (0.0000,8.0000) -- (12.0000,8.0000) -- (0.0000,8.0000) -- (12.0000,8.0000) -- (0.0000,8.0000) -- (12.0000,8.0000) -- (0.0000,8.0000) -- (12.0000,8.0000) -- (0.0000,8.0000) -- (12.0000,8.0000) -- (0.0000,8.0000) -- (12.0000,8.0000) -- (0.0000,8.0000) -- (12.0000,8.0000) -- (0.0000,8.0000) -- (12.0000,8.0000) -- (0.0000,8.0000) -- (12.0000,8.0000) -- (0.0000,8.0000) -- (12.0000,8.0000) -- (0.0000,8.0000) -- (12.0000,8.0000) -- (0.0000,8.0000) -- (12.0000,8.0000) -- (0.0000,8.0000) -- (12.0000,8.0000) -- (0.0000,8.0000) -- (12.0000,8.0000) -- (0.0000,8.0000) -- (12.0000,8.0000) -- (0.0000,8.0000) -- (12.0000,8.0000) -- (0.0000,8.0000) -- (12.0000,8.0000) -- (0.0000,8.0000) -- (12.0000,8.0000) -- (0.0000,8.0000) -- (12.0000,8.0000) -- (0.0000,8.0000) -- (12.0000,8.0000) -- (0.0000,8.0000) -- (12.0000,8.0000) -- (0.0000,8.0000) -- (12.0000,8.0000) -- (0.0000,8.0000) -- (12.0000,8.0000) -- (0.0000,8.0000) -- (12.0000,8.0000) -- (0.0000,8.0000) -- (12.0000,8.0000) -- (0.0000,8.0000) -- (12.0000,8.0000) -- (0.0000,8.0000) -- (12.0000,8.0000) -- (0.0000,8.0000) -- (12.0000,8.0000) -- (0.0000,8.0000) -- (12.0000,8.0000) -- (0.0000,8.0000) -- (12.0000,8.0000) -- (0.0000,8.0000) -- (12.0000,8.0000) -- (0.0000,8.0000) -- (12.0000,8.0000) -- (0.0000,8.0000) -- (12.0000,8.0000) -- (0.0000,8.0000) -- (12.0000,8.0000) -- (0.0000,8.0000) -- (12.0000,8.0000) -- (0.0000,8.0000) -- (12.0000,8.0000) -- (0.0000,8.0000) -- (12.0000,8.0000) -- (0.0000,8.0000) -- (12.0000,8.0000) -- (0.0000,8.0000) -- (12.0000,8.0000) -- (0.0000,8.0000) -- (12.0000,8.0000) -- (0.0000,8.0000) -- (12.0000,8.0000) -- (0.0000,8.0000) -- (12.0000,8.0000) -- (0.0000,8.0000) -- (12.0000,8.0000) -- (0.0000,8.0000) -- (12.0000,8.0000) -- (0.0000,8.0000) -- (12.0000,8.0000) -- (0.0000,8.0000) -- (12.0000,8.0000) -- (0.0000,8.0000) -- (12.0000,8.0000) -- (0.0000,8.0000) -- (12.0000,8.0000) -- (0.0000,8.0000) -- (12.0000,8.0000) -- (0.0000,8.0000) -- (12.0000,8.0000) -- (0.0000,8.0000) -- (12.0000,8.0000) -- (0.0000,8.0000) -- (12.0000,8.0000) -- (0.0000,8.0000) -- (12.0000,8.0000) -- (0.0000,8.0000) -- (12.0000,8.0000) -- (0.0000,8.0000) -- (12.0000,8.0000) -- (0.0000,8.0000) -- (12.0000,8.0000) -- (0.0000,8.0000) -- (12.0000,8.0000) -- (0.0000,8.0000) -- (12.0000,8.0000) -- (0.0000,8.0000) -- (12.0000,8.0000) -- (0.0000,8.0000) -- (12.0000,8.0000) -- (0.0000,8.0000) -- (12.0000,8.0000) -- (0.0000,8.0000) -- (12.0000,8.0000) -- (0.0000,8.0000) -- (12.0000,8.0000) -- (0.0000,8.0000) -- (12.0000,8.0000) -- (0.0000,8.0000) -- (12.0000,8.0000) -- (0.0000,8.0000) -- (12.0000,8.0000) -- (0.0000,8.0000) -- (12.0000,8.0000) -- (0.0000,8.0000) -- (12.0000,8.0000) -- (0.0000,8.0000) -- (12.0000,8.0000) -- (0.0000,8.0000) -- (12.0000,8.0000) -- (0.0000,8.0000) -- (12.0000,8.0000) -- (0.0000,8.0000) -- (12.0000,8.0000) -- (0.0000,8.0000) -- (12.0000,8.0000) -- (0.0000,8.0000) -- (12.0000,8.0000) -- (0.0000,8.0000) -- (12.0000,8.0000) -- (0.0000,8.0000) -- (12.0000,8.0000) -- (0.0000,8.0000) -- (12.0000,8.0000) -- (0.0000,8.0000) -- (12.0000,8.0000) -- (0.0000,8.0000) -- (12.0000,8.0000) -- (0.0000,8.0000) -- (12.0000,8.0000) -- (0.0000,8.0000) -- (12.0000,8.0000) -- (0.0000,8.0000) -- (12.0000,8.0000) -- (0.0000,8.0000) -- (12.0000,8.0000) -- (0.0000,8.0000) -- (12.0000,8.0000) -- (0.0000,8.0000) -- (12.0000,8.0000) -- (0.0000,8.0000) -- (12.0000,8.0000) -- (0.0000,8.0000) -- (12.0000,8.0000) -- (0.0000,8.0000) -- (12.0000,8.0000) -- (0.0000,8.0000) -- (12.0000,8.0000) -- (0.0000,8.0000) -- (12.0000,8.0000) -- (0.0000,8.0000) -- (12.0000,8.0000) -- (0.0000,8.0000) -- (12.0000,8.0000) -- (0.0000,8.0000) -- (12.0000,8.0000) -- (0.0000,8.0000) -- (12.0000,8.0000) -- (0.0000,8.0000) -- (12.0000,8.0000) -- (0.0000,8.0000) -- (12.0000,8.0000) -- (0.0000,8.0000) -- (12.0000,8.0000) -- (0.0000,8.0000) -- (12.0000,8.0000) -- (0.0000,8.0000) -- (12.0000,8.0000) -- (0.0000,8.0000) -- (12.0000,8.0000) -- (0.0000,8.0000) -- (12.0000,8.0000) -- (0.0000,8.0000) -- (12.0000,8.0000) -- (0.0000,8.0000) -- (12.0000,8.0000) -- (0.0000,8.0000) -- (12.0000,8.0000) -- (0.0000,8.0000) -- (12.0000,8.0000) -- (0.0000,8.0000) -- (12.0000,8.0000) -- (0.0000,8.0000) -- (12.0000,8.0000) -- (0.0000,8.0000) -- (12.0000,8.0000) -- (0.0000,8.0000) -- (12.0000,8.0000) -- (0.0000,8.0000) -- (12.0000,8.0000) -- (0.0000,8.0000) -- (12.0000,8.0000) -- (0.0000,8.0000) -- (12.0000,8.0000) -- (0.0000,8.0000) -- (12.0000,8.0000) -- (0.0000,8.0000) -- (12.0000,8.0000) -- (0.0000,8.0000) -- (12.0000,8.0000) -- (0.0000,8.0000) -- (12.0000,8.0000) -- (0.0000,8.0000) -- (12.0000,8.0000) -- (0.0000,8.0000) -- (12.0000,8.0000) -- (0.0000,8.0000) -- (12.0000,8.0000) -- (0.0000,8.0000) -- (12.0000,8.0000) -- (0.0000,8.0000) -- (12.0000,8.0000) -- (0.0000,8.0000) -- (12.0000,8.0000) -- (0.0000,8.0000) -- (12.0000,8.0000) -- (0.0000,8.0000) -- (12.0000,8.0000) -- (0.0000,8.0000) -- (12.0000,8.0000) -- (0.0000,8.0000) -- (12.0000,8.0000) -- (0.0000,8.0000) -- (12.0000,8.0000) -- (0.0000,8.0000) -- (12.0000,8.0000) -- (0.0000,8.0000) -- (12.0000,8.0000) -- (0.0000,8.0000) -- (12.0000,8.0000) -- (0.0000,8.0000) -- (12.0000,8.0000) -- (0.0000,8.0000) -- (12.0000,8.0000) -- (0.0000,8.0000) -- (12.0000,8.0000) -- (0.0000,8.0000) -- (12.0000,8.0000) -- (0.0000,8.0000) -- (12.0000,8.0000) -- (0.0000,8.0000) -- (12.0000,8.0000) -- (0.0000,8.0000) -- (12.0000,8.0000) -- (0.0000,8.0000) -- (12.0000,8.0000) -- (0.0000,8.0000) -- (12.0000,8.0000) -- (0.0000,8.0000) -- (12.0000,8.0000) -- (0.0000,8.0000) -- (12.0000,8.0000) -- (0.0000,8.0000) -- (12.0000,8.0000) -- (0.0000,8.0000) -- (12.0000,8.0000) -- (0.0000,8.0000) -- (12.0000,8.0000) -- (0.0000,8.0000) -- (12.0000,8.0000) -- (0.0000,8.0000) -- (12.0000,8.0000) -- (0.0000,8.0000) -- (12.0000,8.0000) -- (0.0000,8.0000) -- (12.0000,8.0000) -- (0.0000,8.0000) -- (12.0000,8.0000) -- (0.0000,8.0000) -- (12.0000,8.0000) -- (0.0000,8.0000) -- (12.0000,8.0000) -- (0.0000,8.0000) -- (12.0000,8.0000) -- (0.0000,8.0000) -- (12.0000,8.0000) -- (0.0000,8.0000) -- (12.0000,8.0000) -- (0.0000,8.0000) -- (12.0000,8.0000) -- (0.0000,8.0000) -- (12.0000,8.0000) -- (0.0000,8.0000) -- (12.0000,8.0000) -- (0.0000,8.0000) -- (12.0000,8.0000) -- (0.0000,8.0000) -- (12.0000,8.0000) -- (0.0000,8.0000) -- (12.0000,8.0000) -- (0.0000,8.0000) -- (12.0000,8.0000) -- (0.0000,8.0000) -- (12.0000,8.0000) -- (0.0000,8.0000) -- (12.0000,8.0000) -- (0.0000,8.0000) -- (12.0000,8.0000) -- (0.0000,8.0000) -- (12.0000,8.0000) -- (0.0000,8.0000) -- (12.0000,8.0000) -- (0.0000,8.0000) -- (12.0000,8.0000) -- (0.0000,8.0000) -- (12.0000,8.0000) -- (0.0000,8.0000) -- (12.0000,8.0000) -- (0.0000,8.0000) -- (12.0000,8.0000) -- (0.0000,8.0000) -- (12.0000,8.0000) -- (0.0000,8.0000) -- (12.0000,8.0000) -- (0.0000,8.0000) -- (12.0000,8.0000) -- (0.0000,8.0000) -- (12.0000,8.0000) -- (0.0000,8.0000) -- (12.0000,8.0000) -- (0.0000,8.0000) -- (12.0000,8.0000) -- (0.0000,8.0000) -- (12.0000,8.0000) -- (0.0000,8.0000) -- (12.0000,8.0000) -- (0.0000,8.0000) -- (12.0000,8.0000) -- (0.0000,8.0000) -- (12.0000,8.0000) -- (0.0000,8.0000) -- (12.0000,8.0000) -- (0.0000,8.0000) -- (12.0000,8.0000) -- (0.0000,8.0000) -- (12.0000,8.0000) -- (0.0000,8.0000) -- (12.0000,8.0000) -- (0.0000,8.0000) -- (12.0000,8.0000) -- (0.0000,8.0000) -- (12.0000,8.0000) -- (0.0000,8.0000) -- (12.0000,8.0000) -- (0.0000,8.0000) -- (12.0000,8.0000) -- (0.0000,8.0000) -- (12.0000,8.0000) -- (0.0000,8.0000) -- (12.0000,8.0000) -- (0.0000,8.0000) -- (12.0000,8.0000) -- (0.0000,8.0000) -- (12.0000,8.0000) -- (0.0000,8.0000) -- (12.0000,8.0000) -- (0.0000,8.0000) -- (12.0000,8.0000) -- (0.0000,8.0000) -- (12.0000,8.0000) -- (0.0000,8.0000) -- (12.0000,8.0000) -- (0.0000,8.0000) -- (12.0000,8.0000) -- (0.0000,8.0000) -- (12.0000,8.0000) -- (0.0000,8.0000) -- (12.0000,8.0000) -- (0.0000,8.0000) -- (12.0000,8.0000) -- (0.0000,8.0000) -- (12.0000,8.0000) -- (0.0000,8.0000) -- (12.0000,8.0000) -- (0.0000,8.0000) -- (12.0000,8.0000) -- (0.0000,8.0000) -- (12.0000,8.0000) -- (0.0000,8.0000) -- (12.0000,8.0000) -- (0.0000,8.0000) -- (12.0000,8.0000) -- (0.0000,8.0000) -- (12.0000,8.0000) -- (0.0000,8.0000) -- (12.0000,8.0000) -- (0.0000,8.0000) -- (12.0000,8.0000) -- (0.0000,8.0000) -- (12.0000,8.0000) -- (0.0000,8.0000) -- (12.0000,8.0000) -- (0.0000,8.0000) -- (12.0000,8.0000) -- (0.0000,8.0000) -- (12.0000,8.0000) -- (0.0000,8.0000) -- (12.0000,8.0000) -- (0.0000,8.0000) -- (12.0000,8.0000) -- (0.0000,8.0000) -- (12.0000,8.0000) -- (0.0000,8.0000) -- (12.0000,8.0000) -- (0.0000,8.0000) -- (12.0000,8.0000) -- (0.0000,8.0000) -- (12.0000,8.0000) -- (0.0000,8.0000) -- (12.0000,8.0000) -- (0.0000,8.0000) -- (12.0000,8.0000) -- (0.0000,8.0000) -- (12.0000,8.0000) -- (0.0000,8.0000) -- (12.0000,8.0000) -- (0.0000,8.0000) -- (12.0000,8.0000) -- (0.0000,8.0000) -- (12.0000,8.0000) -- (0.0000,8.0000) -- (12.0000,8.0000) -- (0.0000,8.0000) -- (12.0000,8.0000) -- (0.0000,8.0000) -- (12.0000,8.0000) -- (0.0000,8.0000) -- (12.0000,8.0000) -- (0.0000,8.0000) -- (12.0000,8.0000) -- (0.0000,8.0000) -- (12.0000,8.0000) -- (0.0000,8.0000) -- (12.0000,8.0000) -- (0.0000,8.0000) -- (12.0000,8.0000) -- (0.0000,8.0000) -- (12.0000,8.0000) -- (0.0000,8.0000) -- (12.0000,8.0000) -- (0.0000,8.0000) -- (12.0000,8.0000) -- (0.0000,8.0000) -- (12.0000,8.0000) -- (0.0000,8.0000) -- (12.0000,8.0000) -- (0.0000,8.0000) -- (12.0000,8.0000) -- (0.0000,8.0000) -- (12.0000,8.0000) -- (0.0000,8.0000) -- (12.0000,8.0000) -- (0.0000,8.0000) -- (12.0000,8.0000) -- (0.0000,8.0000) -- (12.0000,8.0000) -- (0.0000,8.0000) -- (12.0000,8.0000) -- (0.0000,8.0000) -- (12.0000,8.0000) -- (0.0000,8.0000) -- (12.0000,8.0000) -- (0.0000,8.0000) -- (12.0000,8.0000) -- (0.0000,8.0000) -- (12.0000,8.0000) -- (0.0000,8.0000) -- (12.0000,8.0000) -- (0.0000,8.0000) -- (12.0000,8.0000) -- (0.0000,8.0000) -- (12.0000,8.0000) -- (0.0000,8.0000) -- (12.0000,8.0000) -- (0.0000,8.0000) -- (12.0000,8.0000) -- (0.0000,8.0000) -- (12.0000,8.0000) -- (0.0000,8.0000) -- (12.0000,8.0000) -- (0.0000,8.0000) -- (12.0000,8.0000) -- (0.0000,8.0000) -- (12.0000,8.0000) -- (0.0000,8.0000) -- (12.0000,8.0000) -- (0.0000,8.0000) -- (12.0000,8.0000) -- (0.0000,8.0000) -- (12.0000,8.0000) -- (0.0000,8.0000) -- (12.0000,8.0000) -- (0.0000,8.0000) -- (12.0000,8.0000) -- (0.0000,8.0000) -- (12.0000,8.0000) -- (0.0000,8.0000) -- (12.0000,8.0000) -- (0.0000,8.0000) -- (12.0000,8.0000) -- (0.0000,8.0000) -- (12.0000,8.0000) -- (0.0000,8.0000) -- (12.0000,8.0000) -- (0.0000,8.0000) -- (12.0000,8.0000) -- (0.0000,8.0000) -- (12.0000,8.0000) -- (0.0000,8.0000) -- (12.0000,8.0000) -- (0.0000,8.0000) -- (12.0000,8.0000) -- (0.0000,8.0000) -- (12.0000,8.0000) -- (0.0000,8.0000) -- (12.0000,8.0000) -- (0.0000,8.0000) -- (12.0000,8.0000) -- (0.0000,8.0000) -- (12.0000,8.0000) -- (0.0000,8.0000) -- (12.0000,8.0000) -- (0.0000,8.0000) -- (12.0000,8.0000) -- (0.0000,8.0000) -- (12.0000,8.0000) -- (0.0000,8.0000) -- (12.0000,8.0000) -- (0.0000,8.0000) -- (12.0000,8.0000) -- (0.0000,8.0000) -- (12.0000,8.0000) -- (0.0000,8.0000) -- (12.0000,8.0000) -- (0.0000,8.0000) -- (12.0000,8.0000) -- (0.0000,8.0000) -- (12.0000,8.0000) -- (0.0000,8.0000) -- (12.0000,8.0000) -- (0.0000,8.0000) -- (12.0000,8.0000) -- (0.0000,8.0000) -- (12.0000,8.0000) -- (0.0000,8.0000) -- (12.0000,8.0000) -- (0.0000,8.0000) -- (12.0000,8.0000) -- (0.0000,8.0000) -- (12.0000,8.0000) -- (0.0000,8.0000) -- (12.0000,8.0000) -- (0.0000,8.0000) -- (12.0000,8.0000) -- (0.0000,8.0000) -- (12.0000,8.0000) -- (0.0000,8.0000) -- (12.0000,8.0000) -- (0.0000,8.0000) -- (12.0000,8.0000) -- (0.0000,8.0000) -- (12.0000,8.0000) -- (0.0000,8.0000) -- (12.0000,8.0000) -- (0.0000,8.0000) -- (12.0000,8.0000) -- (0.0000,8.0000) -- (12.0000,8.0000) -- (0.0000,8.0000) -- (12.0000,8.0000) -- (0.0000,8.0000) -- (12.0000,8.0000) -- (0.0000,8.0000) -- (12.0000,8.0000) -- (0.0000,8.0000) -- (12.0000,8.0000) -- (0.0000,8.0000) -- (12.0000,8.0000) -- (0.0000,8.0000) -- (12.0000,8.0000) -- (0.0000,8.0000) -- (12.0000,8.0000) -- (0.0000,8.0000) -- (12.0000,8.0000) -- (0.0000,8.0000) -- (12.0000,8.0000) -- (0.0000,8.0000) -- (12.0000,8.0000) -- (0.0000,8.0000) -- (12.0000,8.0000) -- (0.0000,8.0000) -- (12.0000,8.0000) -- (0.0000,8.0000) -- (12.0000,8.0000) -- (0.0000,8.0000) -- (12.0000,8.0000) -- (0.0000,8.0000) -- (12.0000,8.0000) -- (0.0000,8.0000) -- (12.0000,8.0000) -- (0.0000,8.0000) -- (12.0000,8.0000) -- (0.0000,8.0000) -- (12.0000,8.0000) -- (0.0000,8.0000) -- (12.0000,8.0000) -- (0.0000,8.0000) -- (12.0000,8.0000) -- (0.0000,8.0000) -- (12.0000,8.0000) -- (0.0000,8.0000) -- (12.0000,8.0000) -- (0.0000,8.0000) -- (12.0000,8.0000) -- (0.0000,8.0000) -- (12.0000,8.0000) -- (0.0000,8.0000) -- (12.0000,8.0000) -- (0.0000,8.0000) -- (12.0000,8.0000) -- (0.0000,8.0000) -- (12.0000,8.0000) -- (0.0000,8.0000) -- (12.0000,8.0000) -- (0.0000,8.0000) -- (12.0000,8.0000) -- (0.0000,8.0000) -- (12.0000,8.0000) -- (0.0000,8.0000) -- (12.0000,8.0000) -- (0.0000,8.0000) -- (12.0000,8.0000) -- (0.0000,8.0000) -- (12.0000,8.0000) -- (0.0000,8.0000) -- (12.0000,8.0000) -- (0.0000,8.0000) -- (12.0000,8.0000) -- (0.0000,8.0000) -- (12.0000,8.0000) -- (0.0000,8.0000) -- (12.0000,8.0000) -- (0.0000,8.0000) -- (12.0000,8.0000) -- (0.0000,8.0000) -- (12.0000,8.0000) -- (0.0000,8.0000) -- (12.0000,8.0000) -- (0.0000,8.0000) -- (12.0000,8.0000) -- (0.0000,8.0000) -- (12.0000,8.0000) -- (0.0000,8.0000) -- (12.0000,8.0000) -- (0.0000,8.0000) -- (12.0000,8.0000) -- (0.0000,8.0000) -- (12.0000,8.0000) -- (0.0000,8.0000) -- (12.0000,8.0000) -- (0.0000,8.0000) -- (12.0000,8.0000) -- (0.0000,8.0000) -- (12.0000,8.0000) -- (0.0000,8.0000) -- (12.0000,8.0000) -- (0.0000,8.0000) -- (12.0000,8.0000) -- (0.0000,8.0000) -- (12.0000,8.0000) -- (0.0000,8.0000) -- (12.0000,8.0000) -- (0.0000,8.0000) -- (12.0000,8.0000) -- (0.0000,8.0000) -- (12.0000,8.0000) -- (0.0000,8.0000) -- (12.0000,8.0000) -- (0.0000,8.0000) -- (12.0000,8.0000) -- (0.0000,8.0000) -- (12.0000,8.0000) -- (0.0000,8.0000) -- (12.0000,8.0000) -- (0.0000,8.0000) -- (12.0000,8.0000) -- (0.0000,8.0000) -- (12.0000,8.0000) -- (0.0000,8.0000) -- (12.0000,8.0000) -- (0.0000,8.0000) -- (12.0000,8.0000) -- (0.0000,8.0000) -- (12.0000,8.0000) -- (0.0000,8.0000) -- (12.0000,8.0000) -- (0.0000,8.0000) -- (12.0000,8.0000) -- (0.0000,8.0000) -- (12.0000,8.0000) -- (0.0000,8.0000) -- (12.0000,8.0000) -- (0.0000,8.0000) -- (12.0000,8.0000) -- (0.0000,8.0000) -- (12.0000,8.0000) -- (0.0000,8.0000) -- (12.0000,8.0000) -- (0.0000,8.0000) -- (12.0000,8.0000) -- (0.0000,8.0000) -- (12.0000,8.0000) -- (0.0000,8.0000) -- (12.0000,8.0000) -- (0.0000,8.0000) -- (12.0000,8.0000) -- (0.0000,8.0000) -- (12.0000,8.0000) -- (0.0000,8.0000) -- (12.0000,8.0000) -- (0.0000,8.0000) -- (12.0000,8.0000) -- (0.0000,8.0000) -- (12.0000,8.0000) -- (0.0000,8.0000) -- (12.0000,8.0000) -- (0.0000,8.0000) -- (12.0000,8.0000) -- (0.0000,8.0000) -- (12.0000,8.0000) -- (0.0000,8.0000) -- (12.0000,8.0000) -- (0.0000,8.0000) -- (12.0000,8.0000) -- (0.0000,8.0000) -- (12.0000,8.0000) -- (0.0000,8.0000) -- (12.0000,8.0000) -- (0.0000,8.0000) -- (12.0000,8.0000) -- (0.0000,8.0000) -- (12.0000,8.0000) -- (0.0000,8.0000) -- (12.0000,8.0000) -- (0.0000,8.0000) -- (12.0000,8.0000) -- (0.0000,8.0000) -- (12.0000,8.0000) -- (0.0000,8.0000) -- (12.0000,8.0000) -- (0.0000,8.0000) -- (12.0000,8.0000) -- (0.0000,8.0000) -- (12.0000,8.0000) -- (0.0000,8.0000) -- (12.0000,8.0000) -- (0.0000,8.0000) -- (12.0000,8.0000) -- (0.0000,8.0000) -- (12.0000,8.0000) -- (0.0000,8.0000) -- (12.0000,8.0000) -- (0.0000,8.0000) -- (12.0000,8.0000) -- (0.0000,8.0000) -- (12.0000,8.0000) -- (0.0000,8.0000) -- (12.0000,8.0000) -- (0.0000,8.0000) -- (12.0000,8.0000) -- (0.0000,8.0000) -- (12.0000,8.0000) -- (0.0000,8.0000) -- (12.0000,8.0000) -- (0.0000,8.0000) -- (12.0000,8.0000) -- (0.0000,8.0000) -- (12.0000,8.0000) -- (0.0000,8.0000) -- (12.0000,8.0000) -- (0.0000,8.0000) -- (12.0000,8.0000) -- (0.0000,8.0000) -- (12.0000,8.0000) -- (0.0000,8.0000) -- (12.0000,8.0000) -- (0.0000,8.0000) -- (12.0000,8.0000) -- (0.0000,8.0000) -- (12.0000,8.0000) -- (0.0000,8.0000) -- (12.0000,8.0000) -- (0.0000,8.0000) -- (12.0000,8.0000) -- (0.0000,8.0000) -- (12.0000,8.0000) -- (0.0000,8.0000) -- (12.0000,8.0000) -- (0.0000,8.0000) -- (12.0000,8.0000) -- (0.0000,8.0000) -- (12.0000,8.0000) -- (0.0000,8.0000) -- (12.0000,8.0000) -- (0.0000,8.0000) -- (12.0000,8.0000) -- (0.0000,8.0000) -- (12.0000,8.0000) -- (0.0000,8.0000) -- (12.0000,8.0000) -- (0.0000,8.0000) -- (12.0000,8.0000) -- (0.0000,8.0000) -- (12.0000,8.0000) -- (0.0000,8.0000) -- (12.0000,8.0000) -- (0.0000,8.0000) -- (12.0000,8.0000) -- (0.0000,8.0000) -- (12.0000,8.0000) -- (0.0000,8.0000) -- (12.0000,8.0000) -- (0.0000,8.0000) -- (12.0000,8.0000) -- (0.0000,8.0000) -- (12.0000,8.0000) -- (0.0000,8.0000) -- (12.0000,8.0000) -- (0.0000,8.0000) -- (12.0000,8.0000) -- (0.0000,8.0000) -- (12.0000,8.0000) -- (0.0000,8.0000) -- (12.0000,8.0000) -- (0.0000,8.0000) -- (12.0000,8.0000) -- (0.0000,8.0000) -- (12.0000,8.0000) -- (0.0000,8.0000) -- (12.0000,8.0000) -- (0.0000,8.0000) -- (12.0000,8.0000) -- (0.0000,8.0000) -- (12.0000,8.0000) -- (0.0000,8.0000) -- (12.0000,8.0000) -- (0.0000,8.0000) -- (12.0000,8.0000) -- (0.0000,8.0000) -- (12.0000,8.0000) -- (0.0000,8.0000) -- (12.0000,8.0000) -- (0.0000,8.0000) -- (12.0000,8.0000) -- (0.0000,8.0000) -- (12.0000,8.0000) -- (0.0000,8.0000) -- (12.0000,8.0000) -- (0.0000,8.0000) -- (12.0000,8.0000) -- (0.0000,8.0000) -- (12.0000,8.0000) -- (0.0000,8.0000) -- (12.0000,8.0000) -- (0.0000,8.0000) -- (12.0000,8.0000) -- (0.0000,8.0000) -- (12.0000,8.0000) -- (0.0000,8.0000) -- (12.0000,8.0000) -- (0.0000,8.0000) -- (12.0000,8.0000) -- (0.0000,8.0000) -- (12.0000,8.0000) -- (0.0000,8.0000) -- (12.0000,8.0000) -- (0.0000,8.0000) -- (12.0000,8.0000) -- (0.0000,8.0000) -- (12.0000,8.0000) -- (0.0000,8.0000) -- (12.0000,8.0000) -- (0.0000,8.0000) -- (12.0000,8.0000) -- (0.0000,8.0000) -- (12.0000,8.0000) -- (0.0000,8.0000) -- (12.0000,8.0000) -- (0.0000,8.0000) -- (12.0000,8.0000) -- (0.0000,8.0000) -- (12.0000,8.0000) -- (0.0000,8.0000) -- (12.0000,8.0000) -- (0.0000,8.0000) -- (12.0000,8.0000) -- (0.0000,8.0000) -- (12.0000,8.0000) -- (0.0000,8.0000) -- (12.0000,8.0000) -- (0.0000,8.0000) -- (12.0000,8.0000) -- (0.0000,8.0000) -- (12.0000,8.0000) -- (0.0000,8.0000) -- (12.0000,8.0000) -- (0.0000,8.0000) -- (12.0000,8.0000) -- (0.0000,8.0000) -- (12.0000,8.0000) -- (0.0000,8.0000) -- (12.0000,8.0000) -- (0.0000,8.0000) -- (12.0000,8.0000) -- (0.0000,8.0000) -- (12.0000,8.0000) -- (0.0000,8.0000) -- (12.0000,8.0000) -- (0.0000,8.0000) -- (12.0000,8.0000) -- (0.0000,8.0000) -- (12.0000,8.0000) -- (0.0000,8.0000) -- (12.0000,8.0000) -- (0.0000,8.0000) -- (12.0000,8.0000) -- (0.0000,8.0000) -- (12.0000,8.0000) -- (0.0000,8.0000) -- (12.0000,8.0000) -- (0.0000,8.0000) -- (12.0000,8.0000) -- (0.0000,8.0000) -- (12.0000,8.0000) -- (0.0000,8.0000) -- (12.0000,8.0000) -- (0.0000,8.0000) -- (12.0000,8.0000) -- (0.0000,8.0000) -- (12.0000,8.0000) -- (0.0000,8.0000) -- (12.0000,8.0000) -- (0.0000,8.0000) -- (12.0000,8.0000) -- (0.0000,8.0000) -- (12.0000,8.0000) -- (0.0000,8.0000) -- (12.0000,8.0000) -- (0.0000,8.0000) -- (12.0000,8.0000) -- (0.0000,8.0000) -- (12.0000,8.0000) -- (0.0000,8.0000) -- (12.0000,8.0000) -- (0.0000,8.0000) -- (12.0000,8.0000) -- (0.0000,8.0000) -- (12.0000,8.0000) -- (0.0000,8.0000) -- (12.0000,8.0000) -- (0.0000,8.0000) -- (12.0000,8.0000) -- (0.0000,8.0000) -- (12.0000,8.0000) -- (0.0000,8.0000) -- (12.0000,8.0000) -- (0.0000,8.0000) -- (12.0000,8.0000) -- (0.0000,8.0000) -- (12.0000,8.0000) -- (0.0000,8.0000) -- (12.0000,8.0000) -- (0.0000,8.0000) -- (12.0000,8.0000) -- (0.0000,8.0000) -- (12.0000,8.0000) -- (0.0000,8.0000) -- (12.0000,8.0000) -- (0.0000,8.0000) -- (12.0000,8.0000) -- (0.0000,8.0000) -- (12.0000,8.0000) -- (0.0000,8.0000) -- (12.0000,8.0000) -- (0.0000,8.0000) -- (12.0000,8.0000) -- (0.0000,8.0000) -- (12.0000,8.0000) -- (0.0000,8.0000) -- (12.0000,8.0000) -- (0.0000,8.0000) -- (12.0000,8.0000) -- (0.0000,8.0000) -- (12.0000,8.0000) -- (0.0000,8.0000) -- (12.0000,8.0000) -- (0.0000,8.0000) -- (12.0000,8.0000) -- (0.0000,8.0000) -- (12.0000,8.0000) -- (0.0000,8.0000) -- (12.0000,8.0000) -- (0.0000,8.0000) -- (12.0000,8.0000) -- (0.0000,8.0000) -- (12.0000,8.0000) -- (0.0000,8.0000) -- (12.0000,8.0000) -- (0.0000,8.0000) -- (12.0000,8.0000) -- (0.0000,8.0000) -- (12.0000,8.0000) -- (0.0000,8.0000) -- (12.0000,8.0000) -- (0.0000,8.0000) -- (12.0000,8.0000) -- (0.0000,8.0000) -- (12.0000,8.0000) -- (0.0000,8.0000) -- (12.0000,8.0000) -- (0.0000,8.0000) -- (12.0000,8.0000) -- (0.0000,8.0000) -- (12.0000,8.0000) -- (0.0000,8.0000) -- (12.0000,8.0000) -- (0.0000,8.0000) -- (12.0000,8.0000) -- (0.0000,8.0000) -- (12.0000,8.0000) -- (0.0000,8.0000) -- (12.0000,8.0000) -- (0.0000,8.0000) -- (12.0000,8.0000) -- (0.0000,8.0000) -- (12.0000,8.0000) -- (0.0000,8.0000) -- (12.0000,8.0000) -- (0.0000,8.0000) -- (12.0000,8.0000) -- (0.0000,8.0000) -- (12.0000,8.0000) -- (0.0000,8.0000) -- (12.0000,8.0000) -- (0.0000,8.0000) -- (12.0000,8.0000) -- (0.0000,8.0000) -- (12.0000,8.0000) -- (0.0000,8.0000) -- (12.0000,8.0000) -- (0.0000,8.0000) -- (12.0000,8.0000) -- (0.0000,8.0000) -- (12.0000,8.0000) -- (0.0000,8.0000) -- (12.0000,8.0000) -- (0.0000,8.0000) -- (12.0000,8.0000) -- (0.0000,8.0000) -- (12.0000,8.0000) -- (0.0000,8.0000) -- (12.0000,8.0000) -- (0.0000,8.0000) -- (12.0000,8.0000) -- (0.0000,8.0000) -- (12.0000,8.0000) -- (0.0000,8.0000) -- (12.0000,8.0000) -- (0.0000,8.0000) -- (12.0000,8.0000) -- (0.0000,8.0000) -- (12.0000,8.0000) -- (0.0000,8.0000) -- (12.0000,8.0000) -- (0.0000,8.0000) -- (12.0000,8.0000) -- (0.0000,8.0000) -- (12.0000,8.0000) -- (0.0000,8.0000) -- (12.0000,8.0000) -- (0.0000,8.0000) -- (12.0000,8.0000) -- (0.0000,8.0000) -- (12.0000,8.0000) -- (0.0000,8.0000) -- (12.0000,8.0000) -- (0.0000,8.0000) -- (12.0000,8.0000) -- (0.0000,8.0000) -- (12.0000,8.0000) -- (0.0000,8.0000) -- (12.0000,8.0000) -- (0.0000,8.0000) -- (12.0000,8.0000) -- (0.0000,8.0000) -- (12.0000,8.0000) -- (0.0000,8.0000) -- (12.0000,8.0000) -- (0.0000,8.0000) -- (12.0000,8.0000) -- (0.0000,8.0000) -- (12.0000,8.0000) -- (0.0000,8.0000) -- (12.0000,8.0000) -- (0.0000,8.0000) -- (12.0000,8.0000) -- (0.0000,8.0000) -- (12.0000,8.0000) -- (0.0000,8.0000) -- (12.0000,8.0000) -- (0.0000,8.0000) -- (12.0000,8.0000) -- (0.0000,8.0000) -- (12.0000,8.0000) -- (0.0000,8.0000) -- (12.0000,8.0000) -- (0.0000,8.0000) -- (12.0000,8.0000) -- (0.0000,8.0000) -- (12.0000,8.0000) -- (0.0000,8.0000) -- (12.0000,8.0000) -- (0.0000,8.0000) -- (12.0000,8.0000) -- (0.0000,8.0000) -- (12.0000,8.0000) -- (0.0000,8.0000) -- (12.0000,8.0000) -- (0.0000,8.0000) -- (12.0000,8.0000) -- (0.0000,8.0000) -- (12.0000,8.0000) -- (0.0000,8.0000) -- (12.0000,8.0000) -- (0.0000,8.0000) -- (12.0000,8.0000) -- (0.0000,8.0000) -- (12.0000,8.0000) -- (0.0000,8.0000) -- (12.0000,8.0000) -- (0.0000,8.0000) -- (12.0000,8.0000) -- (0.0000,8.0000) -- (12.0000,8.0000) -- (0.0000,8.0000) -- (12.0000,8.0000) -- (0.0000,8.0000) -- (12.0000,8.0000) -- (0.0000,8.0000) -- (12.0000,8.0000) -- (0.0000,8.0000) -- (12.0000,8.0000) -- (0.0000,8.0000) -- (12.0000,8.0000) -- (0.0000,8.0000) -- (12.0000,8.0000) -- (0.0000,8.0000) -- (12.0000,8.0000) -- (0.0000,8.0000) -- (12.0000,8.0000) -- (0.0000,8.0000) -- (12.0000,8.0000) -- (0.0000,8.0000) -- (12.0000,8.0000) -- (0.0000,8.0000) -- (12.0000,8.0000) -- (0.0000,8.0000) -- (12.0000,8.0000) -- (0.0000,8.0000) -- (12.0000,8.0000) -- (0.0000,8.0000) -- (12.0000,8.0000) -- (0.0000,8.0000) -- (12.0000,8.0000) -- (0.0000,8.0000) -- (12.0000,8.0000) -- (0.0000,8.0000) -- (12.0000,8.0000) -- (0.0000,8.0000) -- (12.0000,8.0000) -- (0.0000,8.0000) -- (12.0000,8.0000) -- (0.0000,8.0000) -- (12.0000,8.0000) -- (0.0000,8.0000) -- (12.0000,8.0000) -- (0.0000,8.0000) -- (12.0000,8.0000) -- (0.0000,8.0000) -- (12.0000,8.0000) -- (0.0000,8.0000) -- (12.0000,8.0000) -- (0.0000,8.0000) -- (12.0000,8.0000) -- (0.0000,8.0000) -- (12.0000,8.0000) -- (0.0000,8.0000) -- (12.0000,8.0000) -- (0.0000,8.0000) -- (12.0000,8.0000) -- (0.0000,8.0000) -- (12.0000,8.0000) -- (0.0000,8.0000) -- (12.0000,8.0000) -- (0.0000,8.0000) -- (12.0000,8.0000) -- (0.0000,8.0000) -- (12.0000,8.0000) -- (0.0000,8.0000) -- (12.0000,8.0000) -- (0.0000,8.0000) -- (12.0000,8.0000) -- (0.0000,8.0000) -- (12.0000,8.0000) -- (0.0000,8.0000) -- (12.0000,8.0000) -- (0.0000,8.0000) -- (12.0000,8.0000) -- (0.0000,8.0000) -- (12.0000,8.0000) -- (0.0000,8.0000) -- (12.0000,8.0000) -- (0.0000,8.0000) -- (12.0000,8.0000) -- (0.0000,8.0000) -- (12.0000,8.0000) -- (0.0000,8.0000) -- (12.0000,8.0000) -- (0.0000,8.0000) -- (12.0000,8.0000) -- (0.0000,8.0000) -- (12.0000,8.0000) -- (0.0000,8.0000) -- (12.0000,8.0000) -- (0.0000,8.0000) -- (12.0000,8.0000) -- (0.0000,8.0000) -- (12.0000,8.0000) -- (0.0000,8.0000) -- (12.0000,8.0000) -- (0.0000,8.0000) -- (12.0000,8.0000) -- (0.0000,8.0000) -- (12.0000,8.0000) -- (0.0000,8.0000) -- (12.0000,8.0000) -- (0.0000,8.0000) -- (12.0000,8.0000) -- (0.0000,8.0000) -- (12.0000,8.0000) -- (0.0000,8.0000) -- (12.0000,8.0000) -- (0.0000,8.0000) -- (12.0000,8.0000) -- (0.0000,8.0000) -- (12.0000,8.0000) -- (0.0000,8.0000) -- (12.0000,8.0000) -- (0.0000,8.0000) -- (12.0000,8.0000) -- (0.0000,8.0000) -- (12.0000,8.0000) -- (0.0000,8.0000) -- (12.0000,8.0000) -- (0.0000,8.0000) -- (12.0000,8.0000) -- (0.0000,8.0000) -- (12.0000,8.0000) -- (0.0000,8.0000) -- (12.0000,8.0000) -- (0.0000,8.0000) -- (12.0000,8.0000) -- (0.0000,8.0000) -- (12.0000,8.0000) -- (0.0000,8.0000) -- (12.0000,8.0000) -- (0.0000,8.0000) -- (12.0000,8.0000) -- (0.0000,8.0000) -- (12.0000,8.0000) -- (0.0000,8.0000) -- (12.0000,8.0000) -- (0.0000,8.0000) -- (12.0000,8.0000) -- (0.0000,8.0000) -- (12.0000,8.0000) -- (0.0000,8.0000) -- (12.0000,8.0000) -- (0.0000,8.0000) -- (12.0000,8.0000) -- (0.0000,8.0000) -- (12.0000,8.0000) -- (0.0000,8.0000) -- (12.0000,8.0000) -- (0.0000,8.0000) -- (12.0000,8.0000) -- (0.0000,8.0000) -- (12.0000,8.0000) -- (0.0000,8.0000) -- (12.0000,8.0000) -- (0.0000,8.0000) -- (12.0000,8.0000) -- (0.0000,8.0000) -- (12.0000,8.0000) -- (0.0000,8.0000) -- (12.0000,8.0000) -- (0.0000,8.0000) -- (12.0000,8.0000) -- (0.0000,8.0000) -- (12.0000,8.0000) -- (0.0000,8.0000) -- (12.0000,8.0000) -- (0.0000,8.0000) -- (12.0000,8.0000) -- (0.0000,8.0000) -- (12.0000,8.0000) -- (0.0000,8.0000) -- (12.0000,8.0000) -- (0.0000,8.0000) -- (12.0000,8.0000) -- (0.0000,8.0000) -- (12.0000,8.0000) -- (0.0000,8.0000) -- (12.0000,8.0000) -- (0.0000,8.0000) -- (12.0000,8.0000) -- (0.0000,8.0000) -- (12.0000,8.0000) -- (0.0000,8.0000) -- (12.0000,8.0000) -- (0.0000,8.0000) -- (12.0000,8.0000) -- (0.0000,8.0000) -- (12.0000,8.0000) -- (0.0000,8.0000) -- (12.0000,8.0000) -- (0.0000,8.0000) -- (12.0000,8.0000) -- (0.0000,8.0000) -- (12.0000,8.0000) -- (0.0000,8.0000) -- (12.0000,8.0000) -- (0.0000,8.0000) -- (12.0000,8.0000) -- (0.0000,8.0000) -- (12.0000,8.0000) -- (0.0000,8.0000) -- (12.0000,8.0000) -- (0.0000,8.0000) -- (12.0000,8.0000) -- (0.0000,8.0000) -- (12.0000,8.0000) -- (0.0000,8.0000) -- (12.0000,8.0000) -- (0.0000,8.0000) -- (12.0000,8.0000) -- (0.0000,8.0000) -- (12.0000,8.0000) -- (0.0000,8.0000) -- (12.0000,8.0000) -- (0.0000,8.0000) -- (12.0000,8.0000) -- (0.0000,8.0000) -- (12.0000,8.0000) -- (0.0000,8.0000) -- (12.0000,8.0000) -- (0.0000,8.0000) -- (12.0000,8.0000) -- (0.0000,8.0000) -- (12.0000,8.0000) -- (0.0000,8.0000) -- (12.0000,8.0000) -- (0.0000,8.0000) -- (12.0000,8.0000) -- (0.0000,8.0000) -- (12.0000,8.0000) -- (0.0000,8.0000) -- (12.0000,8.0000) -- (0.0000,8.0000) -- (12.0000,8.0000) -- (0.0000,8.0000) -- (12.0000,8.0000) -- (0.0000,8.0000) -- (12.0000,8.0000) -- (0.0000,8.0000) -- (12.0000,8.0000) -- (0.0000,8.0000) -- (12.0000,8.0000) -- (0.0000,8.0000) -- (12.0000,8.0000) -- (0.0000,8.0000) -- (12.0000,8.0000) -- (0.0000,8.0000) -- (12.0000,8.0000) -- (0.0000,8.0000) -- (12.0000,8.0000) -- (0.0000,8.0000) -- (12.0000,8.0000) -- (0.0000,8.0000) -- (12.0000,8.0000) -- (0.0000,8.0000) -- (12.0000,8.0000) -- (0.0000,8.0000) -- (12.0000,8.0000) -- (0.0000,8.0000) -- (12.0000,8.0000) -- (0.0000,8.0000) -- (12.0000,8.0000) -- (0.0000,8.0000) -- (12.0000,8.0000) -- (0.0000,8.0000) -- (12.0000,8.0000) -- (0.0000,8.0000) -- (12.0000,8.0000) -- (0.0000,8.0000) -- (12.0000,8.0000) -- (0.0000,8.0000) -- (12.0000,8.0000) -- (0.0000,8.0000) -- (12.0000,8.0000) -- (0.0000,8.0000) -- (12.0000,8.0000) -- (0.0000,8.0000) -- (12.0000,8.0000) -- (0.0000,8.0000) -- (12.0000,8.0000) -- (0.0000,8.0000) -- (12.0000,8.0000) -- (0.0000,8.0000) -- (12.0000,8.0000) -- (0.0000,8.0000) -- (12.0000,8.0000) -- (0.0000,8.0000) -- (12.0000,8.0000) -- (0.0000,8.0000) -- (12.0000,8.0000) -- (0.0000,8.0000) -- (12.0000,8.0000) -- (0.0000,8.0000) -- (12.0000,8.0000) -- (0.0000,8.0000) -- (12.0000,8.0000) -- (0.0000,8.0000) -- (12.0000,8.0000) -- (0.0000,8.0000) -- (12.0000,8.0000) -- (0.0000,8.0000) -- (12.0000,8.0000) -- (0.0000,8.0000) -- (12.0000,8.0000) -- (0.0000,8.0000) -- (12.0000,8.0000) -- (0.0000,8.0000) -- (12.0000,8.0000) -- (0.0000,8.0000) -- (12.0000,8.0000) -- (0.0000,8.0000) -- (12.0000,8.0000) -- (0.0000,8.0000) -- (12.0000,8.0000) -- (0.0000,8.0000) -- (12.0000,8.0000) -- (0.0000,8.0000) -- (12.0000,8.0000) -- (0.0000,8.0000) -- (12.0000,8.0000) -- (0.0000,8.0000) -- (12.0000,8.0000) -- (0.0000,8.0000) -- (12.0000,8.0000) -- (0.0000,8.0000) -- (12.0000,8.0000) -- (0.0000,8.0000) -- (12.0000,8.0000) -- (0.0000,8.0000) -- (12.0000,8.0000) -- (0.0000,8.0000) -- (12.0000,8.0000) -- (0.0000,8.0000) -- (12.0000,8.0000) -- (0.0000,8.0000) -- (12.0000,8.0000) -- (0.0000,8.0000) -- (12.0000,8.0000) -- (0.0000,8.0000) -- (12.0000,8.0000) -- (0.0000,8.0000) -- (12.0000,8.0000) -- (0.0000,8.0000) -- (12.0000,8.0000) -- (0.0000,8.0000) -- (12.0000,8.0000) -- (0.0000,8.0000) -- (12.0000,8.0000) -- (0.0000,8.0000) -- (12.0000,8.0000) -- (0.0000,8.0000) -- (12.0000,8.0000) -- (0.0000,8.0000) -- (12.0000,8.0000) -- (0.0000,8.0000) -- (12.0000,8.0000) -- (0.0000,8.0000) -- (12.0000,8.0000) -- (0.0000,8.0000) -- (12.0000,8.0000) -- (0.0000,8.0000) -- (12.0000,8.0000) -- (0.0000,8.0000) -- (12.0000,8.0000) -- (0.0000,8.0000) -- (12.0000,8.0000) -- (0.0000,8.0000) -- (12.0000,8.0000) -- (0.0000,8.0000) -- (12.0000,8.0000) -- (0.0000,8.0000) -- (12.0000,8.0000) -- (0.0000,8.0000) -- (12.0000,8.0000) -- (0.0000,8.0000) -- (12.0000,8.0000) -- (0.0000,8.0000) -- (12.0000,8.0000) -- (0.0000,8.0000) -- (12.0000,8.0000) -- (0.0000,8.0000) -- (12.0000,8.0000) -- (0.0000,8.0000) -- (12.0000,8.0000) -- (0.0000,8.0000) -- (12.0000,8.0000) -- (0.0000,8.0000) -- (12.0000,8.0000) -- (0.0000,8.0000) -- (12.0000,8.0000) -- (0.0000,8.0000) -- (12.0000,8.0000) -- (0.0000,8.0000) -- (12.0000,8.0000) -- (0.0000,8.0000) -- (12.0000,8.0000) -- (0.0000,8.0000) -- (12.0000,8.0000) -- (0.0000,8.0000) -- (12.0000,8.0000) -- (0.0000,8.0000) -- (12.0000,8.0000) -- (0.0000,8.0000) -- (12.0000,8.0000) -- (0.0000,8.0000) -- (12.0000,8.0000) -- (0.0000,8.0000) -- (12.0000,8.0000) -- (0.0000,8.0000) -- (12.0000,8.0000) -- (0.0000,8.0000) -- (12.0000,8.0000) -- (0.0000,8.0000) -- (12.0000,8.0000) -- (0.0000,8.0000) -- (12.0000,8.0000) -- (0.0000,8.0000) -- (12.0000,8.0000) -- (0.0000,8.0000) -- (12.0000,8.0000) -- (0.0000,8.0000) -- (12.0000,8.0000) -- (0.0000,8.0000) -- (12.0000,8.0000) -- (0.0000,8.0000) -- (12.0000,8.0000) -- (0.0000,8.0000) -- (12.0000,8.0000) -- (0.0000,8.0000) -- (12.0000,8.0000) -- (0.0000,8.0000) -- (12.0000,8.0000) -- (0.0000,8.0000) -- (12.0000,8.0000) -- (0.0000,8.0000) -- (12.0000,8.0000) -- (0.0000,8.0000) -- (12.0000,8.0000) -- (0.0000,8.0000) -- (12.0000,8.0000) -- (0.0000,8.0000) -- (12.0000,8.0000) -- (0.0000,8.0000) -- (12.0000,8.0000) -- (0.0000,8.0000) -- (12.0000,8.0000) -- (0.0000,8.0000) -- (12.0000,8.0000) -- (0.0000,8.0000) -- (12.0000,8.0000) -- (0.0000,8.0000) -- (12.0000,8.0000) -- (0.0000,8.0000) -- (12.0000,8.0000) -- (0.0000,8.0000) -- (12.0000,8.0000) -- (0.0000,8.0000) -- (12.0000,8.0000) -- (0.0000,8.0000) -- (12.0000,8.0000) -- (0.0000,8.0000) -- (12.0000,8.0000) -- (0.0000,8.0000) -- (12.0000,8.0000) -- (0.0000,8.0000) -- (12.0000,8.0000) -- (0.0000,8.0000) -- (12.0000,8.0000) -- (0.0000,8.0000) -- (12.0000,8.0000) -- (0.0000,8.0000) -- (12.0000,8.0000) -- (0.0000,8.0000) -- (12.0000,8.0000) -- (0.0000,8.0000) -- (12.0000,8.0000) -- (0.0000,8.0000) -- (12.0000,8.0000) -- (0.0000,8.0000) -- (12.0000,8.0000) -- (0.0000,8.0000) -- (12.0000,8.0000) -- (0.0000,8.0000) -- (12.0000,8.0000) -- (0.0000,8.0000) -- (12.0000,8.0000) -- (0.0000,8.0000) -- (12.0000,8.0000) -- (0.0000,8.0000) -- (12.0000,8.0000) -- (0.0000,8.0000) -- (12.0000,8.0000) -- (0.0000,8.0000) -- (12.0000,8.0000) -- (0.0000,8.0000) -- (12.0000,8.0000) -- (0.0000,8.0000) -- (12.0000,8.0000) -- (0.0000,8.0000) -- (12.0000,8.0000) -- (0.0000,8.0000) -- (12.0000,8.0000) -- (0.0000,8.0000) -- (12.0000,8.0000) -- (0.0000,8.0000) -- (12.0000,8.0000) -- (0.0000,8.0000) -- (12.0000,8.0000) -- (0.0000,8.0000) -- (12.0000,8.0000) -- (0.0000,8.0000) -- (12.0000,8.0000) -- (0.0000,8.0000) -- (12.0000,8.0000) -- (0.0000,8.0000) -- (12.0000,8.0000) -- (0.0000,8.0000) -- (12.0000,8.0000) -- (0.0000,8.0000) -- (12.0000,8.0000) -- (0.0000,8.0000) -- (12.0000,8.0000) -- (0.0000,8.0000) -- (12.0000,8.0000) -- (0.0000,8.0000) -- (12.0000,8.0000) -- (0.0000,8.0000) -- (12.0000,8.0000) -- (0.0000,8.0000) -- (12.0000,8.0000) -- (0.0000,8.0000) -- (12.0000,8.0000) -- (0.0000,8.0000) -- (12.0000,8.0000) -- (0.0000,8.0000) -- (12.0000,8.0000) -- (0.0000,8.0000) -- (12.0000,8.0000) -- (0.0000,8.0000) -- (12.0000,8.0000) -- (0.0000,8.0000) -- (12.0000,8.0000) -- (0.0000,8.0000) -- (12.0000,8.0000) -- (0.0000,8.0000) -- (12.0000,8.0000) -- (0.0000,8.0000) -- (12.0000,8.0000) -- (0.0000,8.0000) -- (12.0000,8.0000) -- (0.0000,8.0000) -- (12.0000,8.0000) -- (0.0000,8.0000) -- (12.0000,8.0000) -- (0.0000,8.0000) -- (12.0000,8.0000) -- (0.0000,8.0000) -- (12.0000,8.0000) -- (0.0000,8.0000) -- (12.0000,8.0000) -- (0.0000,8.0000) -- (12.0000,8.0000) -- (0.0000,8.0000) -- (12.0000,8.0000) -- (0.0000,8.0000) -- (12.0000,8.0000) -- (0.0000,8.0000) -- (12.0000,8.0000) -- (0.0000,8.0000) -- (12.0000,8.0000) -- (0.0000,8.0000) -- (12.0000,8.0000) -- (0.0000,8.0000) -- (12.0000,8.0000) -- (0.0000,8.0000) -- (12.0000,8.0000) -- (0.0000,8.0000) -- (12.0000,8.0000) -- (0.0000,8.0000) -- (12.0000,8.0000) -- (0.0000,8.0000) -- (12.0000,8.0000) -- (0.0000,8.0000) -- (12.0000,8.0000) -- (0.0000,8.0000) -- (12.0000,8.0000) -- (0.0000,8.0000) -- (12.0000,8.0000) -- (0.0000,8.0000) -- (12.0000,8.0000) -- (0.0000,8.0000) -- (12.0000,8.0000) -- (0.0000,8.0000) -- (12.0000,8.0000) -- (0.0000,8.0000) -- (12.0000,8.0000) -- (0.0000,8.0000) -- (12.0000,8.0000) -- (0.0000,8.0000) -- (12.0000,8.0000) -- (0.0000,8.0000) -- (12.0000,8.0000) -- (0.0000,8.0000) -- (12.0000,8.0000) -- (0.0000,8.0000) -- (12.0000,8.0000) -- (0.0000,8.0000) -- (12.0000,8.0000) -- (0.0000,8.0000) -- (12.0000,8.0000) -- (0.0000,8.0000) -- (12.0000,8.0000) -- (0.0000,8.0000) -- (12.0000,8.0000) -- (0.0000,8.0000) -- (12.0000,8.0000) -- (0.0000,8.0000) -- (12.0000,8.0000) -- (0.0000,8.0000) -- (12.0000,8.0000) -- (0.0000,8.0000) -- (12.0000,8.0000) -- (0.0000,8.0000) -- (12.0000,8.0000) -- (0.0000,8.0000) -- (12.0000,8.0000) -- (0.0000,8.0000) -- (12.0000,8.0000) -- (0.0000,8.0000) -- (12.0000,8.0000) -- (0.0000,8.0000) -- (12.0000,8.0000) -- (0.0000,8.0000) -- (12.0000,8.0000) -- (0.0000,8.0000) -- (12.0000,8.0000) -- (0.0000,8.0000) -- (12.0000,8.0000) -- (0.0000,8.0000) -- (12.0000,8.0000) -- (0.0000,8.0000) -- (12.0000,8.0000) -- (0.0000,8.0000) -- (12.0000,8.0000) -- (0.0000,8.0000) -- (12.0000,8.0000) -- (0.0000,8.0000) -- (12.0000,8.0000) -- (0.0000,8.0000) -- (12.0000,8.0000) -- (0.0000,8.0000) -- (12.0000,8.0000) -- (0.0000,8.0000) -- (12.0000,8.0000) -- (0.0000,8.0000) -- (12.0000,8.0000) -- (0.0000,8.0000) -- (12.0000,8.0000) -- (0.0000,8.0000) -- (12.0000,8.0000) -- (0.0000,8.0000) -- (12.0000,8.0000) -- (0.0000,8.0000) -- (12.0000,8.0000) -- (0.0000,8.0000) -- (12.0000,8.0000) -- (0.0000,8.0000) -- (12.0000,8.0000) -- (0.0000,8.0000) -- (12.0000,8.0000) -- (0.0000,8.0000) -- (12.0000,8.0000) -- (0.0000,8.0000) -- (12.0000,8.0000) -- (0.0000,8.0000) -- (12.0000,8.0000) -- (0.0000,8.0000) -- (12.0000,8.0000) -- (0.0000,8.0000) -- (12.0000,8.0000) -- (0.0000,8.0000) -- (12.0000,8.0000) -- (0.0000,8.0000) -- (12.0000,8.0000) -- (0.0000,8.0000) -- (12.0000,8.0000) -- (0.0000,8.0000) -- (12.0000,8.0000) -- (0.0000,8.0000) -- (12.0000,8.0000) -- (0.0000,8.0000) -- (12.0000,8.0000) -- (0.0000,8.0000) -- (12.0000,8.0000) -- (0.0000,8.0000) -- (12.0000,8.0000) -- (0.0000,8.0000) -- (12.0000,8.0000) -- (0.0000,8.0000) -- (12.0000,8.0000) -- (0.0000,8.0000) -- (12.0000,8.0000) -- (0.0000,8.0000) -- (12.0000,8.0000) -- (0.0000,8.0000) -- (12.0000,8.0000) -- (0.0000,8.0000) -- (12.0000,8.0000) -- (0.0000,8.0000) -- (12.0000,8.0000) -- (0.0000,8.0000) -- (12.0000,8.0000) -- (0.0000,8.0000) -- (12.0000,8.0000) -- (0.0000,8.0000) -- (12.0000,8.0000) -- (0.0000,8.0000) -- (12.0000,8.0000) -- (0.0000,8.0000) -- (12.0000,8.0000) -- (0.0000,8.0000) -- (12.0000,8.0000) -- (0.0000,8.0000) -- (12.0000,8.0000) -- (0.0000,8.0000) -- (12.0000,8.0000) -- (0.0000,8.0000) -- (12.0000,8.0000) -- (0.0000,8.0000) -- (12.0000,8.0000) -- (0.0000,8.0000) -- (12.0000,8.0000) -- (0.0000,8.0000) -- (12.0000,8.0000) -- (0.0000,8.0000) -- (12.0000,8.0000) -- (0.0000,8.0000) -- (12.0000,8.0000) -- (0.0000,8.0000) -- (12.0000,8.0000) -- (0.0000,8.0000) -- (12.0000,8.0000) -- (0.0000,8.0000) -- (12.0000,8.0000) -- (0.0000,8.0000) -- (12.0000,8.0000) -- (0.0000,8.0000) -- (12.0000,8.0000) -- (0.0000,8.0000) -- (12.0000,8.0000) -- (0.0000,8.0000) -- (12.0000,8.0000) -- (0.0000,8.0000) -- (12.0000,8.0000) -- (0.0000,8.0000) -- (12.0000,8.0000) -- (0.0000,8.0000) -- (12.0000,8.0000) -- (0.0000,8.0000) -- (12.0000,8.0000) -- (0.0000,8.0000) -- (12.0000,8.0000) -- (0.0000,8.0000) -- (12.0000,8.0000) -- (0.0000,8.0000) -- (12.0000,8.0000) -- (0.0000,8.0000) -- (12.0000,8.0000) -- (0.0000,8.0000) -- (12.0000,8.0000) -- (0.0000,8.0000) -- (12.0000,8.0000) -- (0.0000,8.0000) -- (12.0000,8.0000) -- (0.0000,8.0000) -- (12.0000,8.0000) -- (0.0000,8.0000) -- (12.0000,8.0000) -- (0.0000,8.0000) -- (12.0000,8.0000) -- (0.0000,8.0000) -- (12.0000,8.0000) -- (0.0000,8.0000) -- (12.0000,8.0000) -- (0.0000,8.0000) -- (12.0000,8.0000) -- (0.0000,8.0000) -- (12.0000,8.0000) -- (0.0000,8.0000) -- (12.0000,8.0000) -- (0.0000,8.0000) -- (12.0000,8.0000) -- (0.0000,8.0000) -- (12.0000,8.0000) -- (0.0000,8.0000) -- (12.0000,8.0000) -- (0.0000,8.0000) -- (12.0000,8.0000) -- (0.0000,8.0000) -- (12.0000,8.0000) -- (0.0000,8.0000) -- (12.0000,8.0000) -- (0.0000,8.0000) -- (12.0000,8.0000) -- (0.0000,8.0000) -- (12.0000,8.0000) -- (0.0000,8.0000) -- (12.0000,8.0000) -- (0.0000,8.0000) -- (12.0000,8.0000) -- (0.0000,8.0000) -- (12.0000,8.0000) -- (0.0000,8.0000) -- (12.0000,8.0000) -- (0.0000,8.0000) -- (12.0000,8.0000) -- (0.0000,8.0000) -- (12.0000,8.0000) -- (0.0000,8.0000) -- (12.0000,8.0000) -- (0.0000,8.0000) -- (12.0000,8.0000) -- (0.0000,8.0000) -- (12.0000,8.0000) -- (0.0000,8.0000) -- (12.0000,8.0000) -- (0.0000,8.0000) -- (12.0000,8.0000) -- (0.0000,8.0000) -- (12.0000,8.0000) -- (0.0000,8.0000) -- (12.0000,8.0000) -- (0.0000,8.0000) -- (12.0000,8.0000) -- (0.0000,8.0000) -- (12.0000,8.0000) -- (0.0000,8.0000) -- (12.0000,8.0000) -- (0.0000,8.0000) -- (12.0000,8.0000) -- (0.0000,8.0000) -- (12.0000,8.0000) -- (0.0000,8.0000) -- (12.0000,8.0000) -- (0.0000,8.0000) -- (12.0000,8.0000) -- (0.0000,8.0000) -- (12.0000,8.0000) -- (0.0000,8.0000) -- (12.0000,8.0000) -- (0.0000,8.0000) -- (12.0000,8.0000) -- (0.0000,8.0000) -- (12.0000,8.0000) -- (0.0000,8.0000) -- (12.0000,8.0000) -- (0.0000,8.0000) -- (12.0000,8.0000) -- (0.0000,8.0000) -- (12.0000,8.0000) -- (0.0000,8.0000) -- (12.0000,8.0000) -- (0.0000,8.0000) -- (12.0000,8.0000) -- (0.0000,8.0000) -- (12.0000,8.0000) -- (0.0000,8.0000) -- (12.0000,8.0000) -- (0.0000,8.0000) -- (12.0000,8.0000) -- (0.0000,8.0000) -- (12.0000,8.0000) -- (0.0000,8.0000) -- (12.0000,8.0000) -- (0.0000,8.0000) -- (12.0000,8.0000) -- (0.0000,8.0000) -- (12.0000,8.0000) -- (0.0000,8.0000) -- (12.0000,8.0000) -- (0.0000,8.0000) -- (12.0000,8.0000) -- (0.0000,8.0000) -- (12.0000,8.0000) -- (0.0000,8.0000) -- (12.0000,8.0000) -- (0.0000,8.0000) -- (12.0000,8.0000) -- (0.0000,8.0000) -- (12.0000,8.0000) -- (0.0000,8.0000) -- (12.0000,8.0000) -- (0.0000,8.0000) -- (12.0000,8.0000) -- (0.0000,8.0000) -- (12.0000,8.0000) -- (0.0000,8.0000) -- (12.0000,8.0000) -- (0.0000,8.0000) -- (12.0000,8.0000) -- (0.0000,8.0000) -- (12.0000,8.0000) -- (0.0000,8.0000) -- (12.0000,8.0000) -- (0.0000,8.0000) -- (12.0000,8.0000) -- (0.0000,8.0000) -- (12.0000,8.0000) -- (0.0000,8.0000) -- (12.0000,8.0000) -- (0.0000,8.0000) -- (12.0000,8.0000) -- (0.0000,8.0000) -- (12.0000,8.0000) -- (0.0000,8.0000) -- (12.0000,8.0000) -- (0.0000,8.0000) -- (12.0000,8.0000) -- (0.0000,8.0000) -- (12.0000,8.0000) -- (0.0000,8.0000) -- (12.0000,8.0000) -- (0.0000,8.0000) -- (12.0000,8.0000) -- (0.0000,8.0000) -- (12.0000,8.0000) -- (0.0000,8.0000) -- (12.0000,8.0000) -- (0.0000,8.0000) -- (12.0000,8.0000) -- (0.0000,8.0000) -- (12.0000,8.0000) -- (0.0000,8.0000) -- (12.0000,8.0000) -- (0.0000,8.0000) -- (12.0000,8.0000) -- (0.0000,8.0000) -- (12.0000,8.0000) -- (0.0000,8.0000) -- (12.0000,8.0000) -- (0.0000,8.0000) -- (12.0000,8.0000) -- (0.0000,8.0000) -- (12.0000,8.0000) -- (0.0000,8.0000) -- (12.0000,8.0000) -- (0.0000,8.0000) -- (12.0000,8.0000) -- (0.0000,8.0000) -- (12.0000,8.0000) -- (0.0000,8.0000) -- (12.0000,8.0000) -- (0.0000,8.0000) -- (12.0000,8.0000) -- (0.0000,8.0000) -- (12.0000,8.0000) -- (0.0000,8.0000) -- (12.0000,8.0000) -- (0.0000,8.0000) -- (12.0000,8.0000) -- (0.0000,8.0000) -- (12.0000,8.0000) -- (0.0000,8.0000) -- (12.0000,8.0000) -- (0.0000,8.0000) -- (12.0000,8.0000) -- (0.0000,8.0000) -- (12.0000,8.0000) -- (0.0000,8.0000) -- (12.0000,8.0000) -- (0.0000,8.0000) -- (12.0000,8.0000) -- (0.0000,8.0000) -- (12.0000,8.0000) -- (0.0000,8.0000) -- (12.0000,8.0000) -- (0.0000,8.0000) -- (12.0000,8.0000) -- (0.0000,8.0000) -- (12.0000,8.0000) -- (0.0000,8.0000) -- (12.0000,8.0000) -- (0.0000,8.0000) -- (12.0000,8.0000) -- (0.0000,8.0000) -- (12.0000,8.0000) -- (0.0000,8.0000) -- (12.0000,8.0000) -- (0.0000,8.0000) -- (12.0000,8.0000) -- (0.0000,8.0000) -- (12.0000,8.0000) -- (0.0000,8.0000) -- (12.0000,8.0000) -- (0.0000,8.0000) -- (12.0000,8.0000) -- (0.0000,8.0000) -- (12.0000,8.0000) -- (0.0000,8.0000) -- (12.0000,8.0000) -- (0.0000,8.0000) -- (12.0000,8.0000) -- (0.0000,8.0000) -- (12.0000,8.0000) -- (0.0000,8.0000) -- (12.0000,8.0000) -- (0.0000,8.0000) -- (12.0000,8.0000) -- (0.0000,8.0000) -- (12.0000,8.0000) -- (0.0000,8.0000) -- (12.0000,8.0000) -- (0.0000,8.0000) -- (12.0000,8.0000) -- (0.0000,8.0000) -- (12.0000,8.0000) -- (0.0000,8.0000) -- (12.0000,8.0000) -- (0.0000,8.0000) -- (12.0000,8.0000) -- (0.0000,8.0000) -- (12.0000,8.0000) -- (0.0000,8.0000) -- (12.0000,8.0000) -- (0.0000,8.0000) -- (12.0000,8.0000) -- (0.0000,8.0000) -- (12.0000,8.0000) -- (0.0000,8.0000) -- (12.0000,8.0000) -- (0.0000,8.0000) -- (12.0000,8.0000) -- (0.0000,8.0000) -- (12.0000,8.0000) -- (0.0000,8.0000) -- (12.0000,8.0000) -- (0.0000,8.0000) -- (12.0000,8.0000) -- (0.0000,8.0000) -- (12.0000,8.0000) -- (0.0000,8.0000) -- (12.0000,8.0000) -- (0.0000,8.0000) -- (12.0000,8.0000) -- (0.0000,8.0000) -- (12.0000,8.0000) -- (0.0000,8.0000) -- (12.0000,8.0000) -- (0.0000,8.0000) -- (12.0000,8.0000) -- (0.0000,8.0000) -- (12.0000,8.0000) -- (0.0000,8.0000) -- (12.0000,8.0000) -- (0.0000,8.0000) -- (12.0000,8.0000) -- (0.0000,8.0000) -- (12.0000,8.0000) -- (0.0000,8.0000) -- (12.0000,8.0000) -- (0.0000,8.0000) -- (12.0000,8.0000) -- (0.0000,8.0000) -- (12.0000,8.0000) -- (0.0000,8.0000) -- (12.0000,8.0000) -- (0.0000,8.0000) -- (12.0000,8.0000) -- (0.0000,8.0000) -- (12.0000,8.0000) -- (0.0000,8.0000) -- (12.0000,8.0000) -- (0.0000,8.0000) -- (12.0000,8.0000) -- (0.0000,8.0000) -- (12.0000,8.0000) -- (0.0000,8.0000) -- (12.0000,8.0000) -- (0.0000,8.0000) -- (12.0000,8.0000) -- (0.0000,8.0000) -- (12.0000,8.0000) -- (0.0000,8.0000) -- (12.0000,8.0000) -- (0.0000,8.0000) -- (12.0000,8.0000) -- (0.0000,8.0000) -- (12.0000,8.0000) -- (0.0000,8.0000) -- (12.0000,8.0000) -- (0.0000,8.0000) -- (12.0000,8.0000) -- (0.0000,8.0000) -- (12.0000,8.0000) -- (0.0000,8.0000) -- (12.0000,8.0000) -- (0.0000,8.0000) -- (12.0000,8.0000) -- (0.0000,8.0000) -- (12.0000,8.0000) -- (0.0000,8.0000) -- (12.0000,8.0000) -- (0.0000,8.0000) -- (12.0000,8.0000) -- (0.0000,8.0000) -- (12.0000,8.0000) -- (0.0000,8.0000) -- (12.0000,8.0000) -- (0.0000,8.0000) -- (12.0000,8.0000) -- (0.0000,8.0000) -- (12.0000,8.0000) -- (0.0000,8.0000) -- (12.0000,8.0000) -- (0.0000,8.0000) -- (12.0000,8.0000) -- (0.0000,8.0000) -- (12.0000,8.0000) -- (0.0000,8.0000) -- (12.0000,8.0000) -- (0.0000,8.0000) -- (12.0000,8.0000) -- (0.0000,8.0000) -- (12.0000,8.0000) -- (0.0000,8.0000) -- (12.0000,8.0000) -- (0.0000,8.0000) -- (12.0000,8.0000) -- (0.0000,8.0000) -- (12.0000,8.0000) -- (0.0000,8.0000) -- (12.0000,8.0000) -- (0.0000,8.0000) -- (12.0000,8.0000) -- (0.0000,8.0000) -- (12.0000,8.0000) -- (0.0000,8.0000) -- (12.0000,8.0000) -- (0.0000,8.0000) -- (12.0000,8.0000) -- (0.0000,8.0000) -- (12.0000,8.0000) -- (0.0000,8.0000) -- (12.0000,8.0000) -- (0.0000,8.0000) -- (12.0000,8.0000) -- (0.0000,8.0000) -- (12.0000,8.0000) -- (0.0000,8.0000) -- (12.0000,8.0000) -- (0.0000,8.0000) -- (12.0000,8.0000) -- (0.0000,8.0000) -- (12.0000,8.0000) -- (0.0000,8.0000) -- (12.0000,8.0000) -- (0.0000,8.0000) -- (12.0000,8.0000) -- (0.0000,8.0000) -- (12.0000,8.0000) -- (0.0000,8.0000) -- (12.0000,8.0000) -- (0.0000,8.0000) -- (12.0000,8.0000) -- (0.0000,8.0000) -- (12.0000,8.0000) -- (0.0000,8.0000) -- (12.0000,8.0000) -- (0.0000,8.0000) -- (12.0000,8.0000) -- (0.0000,8.0000) -- (12.0000,8.0000) -- (0.0000,8.0000) -- (12.0000,8.0000) -- (0.0000,8.0000) -- (12.0000,8.0000) -- (0.0000,8.0000) -- (12.0000,8.0000) -- (0.0000,8.0000) -- (12.0000,8.0000) -- (0.0000,8.0000) -- (12.0000,8.0000) -- (0.0000,8.0000) -- (12.0000,8.0000) -- (0.0000,8.0000) -- (12.0000,8.0000) -- (0.0000,8.0000) -- (12.0000,8.0000) -- (0.0000,8.0000) -- (12.0000,8.0000) -- (0.0000,8.0000) -- (12.0000,8.0000) -- (0.0000,8.0000) -- (12.0000,8.0000) -- (0.0000,8.0000) -- (12.0000,8.0000) -- (0.0000,8.0000) -- (12.0000,8.0000) -- (0.0000,8.0000) -- (12.0000,8.0000) -- (0.0000,8.0000) -- (12.0000,8.0000) -- (0.0000,8.0000) -- (12.0000,8.0000) -- (0.0000,8.0000) -- (12.0000,8.0000) -- (0.0000,8.0000) -- (12.0000,8.0000) -- (0.0000,8.0000) -- (12.0000,8.0000) -- (0.0000,8.0000) -- (12.0000,8.0000) -- (0.0000,8.0000) -- (12.0000,8.0000) -- (0.0000,8.0000) -- (12.0000,8.0000) -- (0.0000,8.0000) -- (12.0000,8.0000) -- (0.0000,8.0000) -- (12.0000,8.0000) -- (0.0000,8.0000) -- (12.0000,8.0000) -- (0.0000,8.0000) -- (12.0000,8.0000) -- (0.0000,8.0000) -- (12.0000,8.0000) -- (0.0000,8.0000) -- (12.0000,8.0000) -- (0.0000,8.0000) -- (12.0000,8.0000) -- (0.0000,8.0000) -- (12.0000,8.0000) -- (0.0000,8.0000) -- (12.0000,8.0000) -- (0.0000,8.0000) -- (12.0000,8.0000) -- (0.0000,8.0000) -- (12.0000,8.0000) -- (0.0000,8.0000) -- (12.0000,8.0000) -- (0.0000,8.0000) -- (12.0000,8.0000) -- (0.0000,8.0000) -- (12.0000,8.0000) -- (0.0000,8.0000) -- (12.0000,8.0000) -- (0.0000,8.0000) -- (12.0000,8.0000) -- (0.0000,8.0000) -- (12.0000,8.0000) -- (0.0000,8.0000) -- (12.0000,8.0000) -- (0.0000,8.0000) -- (12.0000,8.0000) -- (0.0000,8.0000) -- (12.0000,8.0000) -- (0.0000,8.0000) -- (12.0000,8.0000) -- (0.0000,8.0000) -- (12.0000,8.0000) -- (0.0000,8.0000) -- (12.0000,8.0000) -- (0.0000,8.0000) -- (12.0000,8.0000) -- (0.0000,8.0000) -- (12.0000,8.0000) -- (0.0000,8.0000) -- (12.0000,8.0000) -- (0.0000,8.0000) -- (12.0000,8.0000) -- (0.0000,8.0000) -- (12.0000,8.0000) -- (0.0000,8.0000) -- (12.0000,8.0000) -- (0.0000,8.0000) -- (12.0000,8.0000) -- (0.0000,8.0000) -- (12.0000,8.0000) -- (0.0000,8.0000) -- (12.0000,8.0000) -- (0.0000,8.0000) -- (12.0000,8.0000) -- (0.0000,8.0000) -- (12.0000,8.0000) -- (0.0000,8.0000) -- (12.0000,8.0000) -- (0.0000,8.0000) -- (12.0000,8.0000) -- (0.0000,8.0000) -- (12.0000,8.0000) -- (0.0000,8.0000) -- (12.0000,8.0000) -- (0.0000,8.0000) -- (12.0000,8.0000) -- (0.0000,8.0000) -- (12.0000,8.0000) -- (0.0000,8.0000) -- (12.0000,8.0000) -- (0.0000,8.0000) -- (12.0000,8.0000) -- (0.0000,8.0000) -- (12.0000,8.0000) -- (0.0000,8.0000) -- (12.0000,8.0000) -- (0.0000,8.0000) -- (12.0000,8.0000) -- (0.0000,8.0000) -- (12.0000,8.0000) -- (0.0000,8.0000) -- (12.0000,8.0000) -- (0.0000,8.0000) -- (12.0000,8.0000) -- (0.0000,8.0000) -- (12.0000,8.0000) -- (0.0000,8.0000) -- (12.0000,8.0000) -- (0.0000,8.0000) -- (12.0000,8.0000) -- (0.0000,8.0000) -- (12.0000,8.0000) -- (0.0000,8.0000) -- (12.0000,8.0000) -- (0.0000,8.0000) -- (12.0000,8.0000) -- (0.0000,8.0000) -- (12.0000,8.0000) -- (0.0000,8.0000) -- (12.0000,8.0000) -- (0.0000,8.0000) -- (12.0000,8.0000) -- (0.0000,8.0000) -- (12.0000,8.0000) -- (0.0000,8.0000) -- (12.0000,8.0000) -- (0.0000,8.0000) -- (12.0000,8.0000) -- (0.0000,8.0000) -- (12.0000,8.0000) -- (0.0000,8.0000) -- (12.0000,8.0000) -- (0.0000,8.0000) -- (12.0000,8.0000) -- (0.0000,8.0000) -- (12.0000,8.0000) -- (0.0000,8.0000) -- (12.0000,8.0000) -- (0.0000,8.0000) -- (12.0000,8.0000) -- (0.0000,8.0000) -- (12.0000,8.0000) -- (0.0000,8.0000) -- (12.0000,8.0000) -- (0.0000,8.0000) -- (12.0000,8.0000) -- (0.0000,8.0000) -- (12.0000,8.0000) -- (0.0000,8.0000) -- (12.0000,8.0000) -- (0.0000,8.0000) -- (12.0000,8.0000) -- (0.0000,8.0000) -- (12.0000,8.0000) -- (0.0000,8.0000) -- (12.0000,8.0000) -- (0.0000,8.0000) -- (12.0000,8.0000) -- (0.0000,8.0000) -- (12.0000,8.0000) -- (0.0000,8.0000) -- (12.0000,8.0000) -- (0.0000,8.0000) -- (12.0000,8.0000) -- (0.0000,8.0000) -- (12.0000,8.0000) -- (0.0000,8.0000) -- (12.0000,8.0000) -- (0.0000,8.0000) -- (12.0000,8.0000) -- (0.0000,8.0000) -- (12.0000,8.0000) -- (0.0000,8.0000) -- (12.0000,8.0000) -- (0.0000,8.0000) -- (12.0000,8.0000) -- (0.0000,8.0000) -- (12.0000,8.0000) -- (0.0000,8.0000) -- (12.0000,8.0000) -- (0.0000,8.0000) -- (12.0000,8.0000) -- (0.0000,8.0000) -- (12.0000,8.0000) -- (0.0000,8.0000) -- (12.0000,8.0000) -- (0.0000,8.0000) -- (12.0000,8.0000) -- (0.0000,8.0000) -- (12.0000,8.0000) -- (0.0000,8.0000) -- (12.0000,8.0000) -- (0.0000,8.0000) -- (12.0000,8.0000) -- (0.0000,8.0000) -- (12.0000,8.0000) -- (0.0000,8.0000) -- (12.0000,8.0000) -- (0.0000,8.0000) -- (12.0000,8.0000) -- (0.0000,8.0000) -- (12.0000,8.0000) -- (0.0000,8.0000) -- (12.0000,8.0000) -- (0.0000,8.0000) -- (12.0000,8.0000) -- (0.0000,8.0000) -- (12.0000,8.0000) -- (0.0000,8.0000) -- (12.0000,8.0000) -- (0.0000,8.0000) -- (12.0000,8.0000) -- (0.0000,8.0000) -- (12.0000,8.0000) -- (0.0000,8.0000) -- (12.0000,8.0000) -- (0.0000,8.0000) -- (12.0000,8.0000) -- (0.0000,8.0000) -- (12.0000,8.0000) -- (0.0000,8.0000) -- (12.0000,8.0000) -- (0.0000,8.0000) -- (12.0000,8.0000) -- (0.0000,8.0000) -- (12.0000,8.0000) -- (0.0000,8.0000) -- (12.0000,8.0000) -- (0.0000,8.0000) -- (12.0000,8.0000) -- (0.0000,8.0000) -- (12.0000,8.0000) -- (0.0000,8.0000) -- (12.0000,8.0000) -- (0.0000,8.0000) -- (12.0000,8.0000) -- (0.0000,8.0000) -- (12.0000,8.0000) -- (0.0000,8.0000) -- (12.0000,8.0000) -- (0.0000,8.0000) -- (12.0000,8.0000) -- (0.0000,8.0000) -- (12.0000,8.0000) -- (0.0000,8.0000) -- (12.0000,8.0000) -- (0.0000,8.0000) -- (12.0000,8.0000) -- (0.0000,8.0000) -- (12.0000,8.0000) -- (0.0000,8.0000) -- (12.0000,8.0000) -- (0.0000,8.0000) -- (12.0000,8.0000) -- (0.0000,8.0000) -- (12.0000,8.0000) -- (0.0000,8.0000) -- (12.0000,8.0000) -- (0.0000,8.0000) -- (12.0000,8.0000) -- (0.0000,8.0000) -- (12.0000,8.0000) -- (0.0000,8.0000) -- (12.0000,8.0000) -- (0.0000,8.0000) -- (12.0000,8.0000) -- (0.0000,8.0000) -- (12.0000,8.0000) -- (0.0000,8.0000) -- (12.0000,8.0000) -- (0.0000,8.0000) -- (12.0000,8.0000) -- (0.0000,8.0000) -- (12.0000,8.0000) -- (0.0000,8.0000) -- (12.0000,8.0000) -- (0.0000,8.0000) -- (12.0000,8.0000) -- (0.0000,8.0000) -- (12.0000,8.0000) -- (0.0000,8.0000) -- (12.0000,8.0000) -- (0.0000,8.0000) -- (12.0000,8.0000) -- (0.0000,8.0000) -- (12.0000,8.0000) -- (0.0000,8.0000) -- (12.0000,8.0000) -- (0.0000,8.0000) -- (12.0000,8.0000) -- (0.0000,8.0000) -- (12.0000,8.0000) -- (0.0000,8.0000) -- (12.0000,8.0000) -- (0.0000,8.0000) -- (12.0000,8.0000) -- (0.0000,8.0000) -- (12.0000,8.0000) -- (0.0000,8.0000) -- (12.0000,8.0000) -- (0.0000,8.0000) -- (12.0000,8.0000) -- (0.0000,8.0000) -- (12.0000,8.0000) -- (0.0000,8.0000) -- (12.0000,8.0000) -- (0.0000,8.0000) -- (12.0000,8.0000) -- (0.0000,8.0000) -- (12.0000,8.0000) -- (0.0000,8.0000) -- (12.0000,8.0000) -- (0.0000,8.0000) -- (12.0000,8.0000) -- (0.0000,8.0000) -- (12.0000,8.0000) -- (0.0000,8.0000) -- (12.0000,8.0000) -- (0.0000,8.0000) -- (12.0000,8.0000) -- (0.0000,8.0000) -- (12.0000,8.0000) -- (0.0000,8.0000) -- (12.0000,8.0000) -- (0.0000,8.0000) -- (12.0000,8.0000) -- (0.0000,8.0000) -- (12.0000,8.0000) -- (0.0000,8.0000) -- (12.0000,8.0000) -- (0.0000,8.0000) -- (12.0000,8.0000) -- (0.0000,8.0000) -- (12.0000,8.0000) -- (0.0000,8.0000) -- (12.0000,8.0000) -- (0.0000,8.0000) -- (12.0000,8.0000) -- (0.0000,8.0000) -- (12.0000,8.0000) -- (0.0000,8.0000) -- (12.0000,8.0000) -- (0.0000,8.0000) -- (12.0000,8.0000) -- (0.0000,8.0000) -- (12.0000,8.0000) -- (0.0000,8.0000) -- (12.0000,8.0000) -- (0.0000,8.0000) -- (12.0000,8.0000) -- (0.0000,8.0000) -- (12.0000,8.0000) -- (0.0000,8.0000) -- (12.0000,8.0000) -- (0.0000,8.0000) -- (12.0000,8.0000) -- (0.0000,8.0000) -- (12.0000,8.0000) -- (0.0000,8.0000) -- (12.0000,8.0000) -- (0.0000,8.0000) -- (12.0000,8.0000) -- (0.0000,8.0000) -- (12.0000,8.0000) -- (0.0000,8.0000) -- (12.0000,8.0000) -- (0.0000,8.0000) -- (12.0000,8.0000) -- (0.0000,8.0000) -- (12.0000,8.0000) -- (0.0000,8.0000) -- (12.0000,8.0000) -- (0.0000,8.0000) -- (12.0000,8.0000) -- (0.0000,8.0000) -- (12.0000,8.0000) -- (0.0000,8.0000) -- (12.0000,8.0000) -- (0.0000,8.0000) -- (12.0000,8.0000) -- (0.0000,8.0000) -- (12.0000,8.0000) -- (0.0000,8.0000) -- (12.0000,8.0000) -- (0.0000,8.0000) -- (12.0000,8.0000) -- (0.0000,8.0000) -- (12.0000,8.0000) -- (0.0000,8.0000) -- (12.0000,8.0000) -- (0.0000,8.0000) -- (12.0000,8.0000) -- (0.0000,8.0000) -- (12.0000,8.0000) -- (0.0000,8.0000) -- (12.0000,8.0000) -- (0.0000,8.0000) -- (12.0000,8.0000) -- (0.0000,8.0000) -- (12.0000,8.0000) -- (0.0000,8.0000) -- (12.0000,8.0000) -- (0.0000,8.0000) -- (12.0000,8.0000) -- (0.0000,8.0000) -- (12.0000,8.0000) -- (0.0000,8.0000) -- (12.0000,8.0000) -- (0.0000,8.0000) -- (12.0000,8.0000) -- (0.0000,8.0000) -- (12.0000,8.0000) -- (0.0000,8.0000) -- (12.0000,8.0000) -- (0.0000,8.0000) -- (12.0000,8.0000) -- (0.0000,8.0000) -- (12.0000,8.0000) -- (0.0000,8.0000) -- (12.0000,8.0000) -- (0.0000,8.0000) -- (12.0000,8.0000) -- (0.0000,8.0000) -- (12.0000,8.0000) -- (0.0000,8.0000) -- (12.0000,8.0000) -- (0.0000,8.0000) -- (12.0000,8.0000) -- (0.0000,8.0000) -- (12.0000,8.0000) -- (0.0000,8.0000) -- (12.0000,8.0000) -- (0.0000,8.0000) -- (12.0000,8.0000) -- (0.0000,8.0000) -- (12.0000,8.0000) -- (0.0000,8.0000) -- (12.0000,8.0000) -- (0.0000,8.0000) -- (12.0000,8.0000) -- (0.0000,8.0000) -- (12.0000,8.0000) -- (0.0000,8.0000) -- (12.0000,8.0000) -- (0.0000,8.0000) -- (12.0000,8.0000) -- (0.0000,8.0000) -- (12.0000,8.0000) -- (0.0000,8.0000) -- (12.0000,8.0000) -- (0.0000,8.0000) -- (12.0000,8.0000) -- (0.0000,8.0000) -- (12.0000,8.0000) -- (0.0000,8.0000) -- (12.0000,8.0000) -- (0.0000,8.0000) -- (12.0000,8.0000) -- (0.0000,8.0000) -- (12.0000,8.0000) -- (0.0000,8.0000) -- (12.0000,8.0000) -- (0.0000,8.0000) -- (12.0000,8.0000) -- (0.0000,8.0000) -- (12.0000,8.0000) -- (0.0000,8.0000) -- (12.0000,8.0000) -- (0.0000,8.0000) -- (12.0000,8.0000) -- (0.0000,8.0000) -- (12.0000,8.0000) -- (0.0000,8.0000) -- (12.0000,8.0000) -- (0.0000,8.0000) -- (12.0000,8.0000) -- (0.0000,8.0000) -- (12.0000,8.0000) -- (0.0000,8.0000) -- (12.0000,8.0000) -- (0.0000,8.0000) -- (12.0000,8.0000) -- (0.0000,8.0000) -- (12.0000,8.0000) -- (0.0000,8.0000) -- (12.0000,8.0000) -- (0.0000,8.0000) -- (12.0000,8.0000) -- (0.0000,8.0000) -- (12.0000,8.0000) -- (0.0000,8.0000) -- (12.0000,8.0000) -- (0.0000,8.0000) -- (12.0000,8.0000) -- (0.0000,8.0000) -- (12.0000,8.0000) -- (0.0000,8.0000) -- (12.0000,8.0000) -- (0.0000,8.0000) -- (12.0000,8.0000) -- (0.0000,8.0000) -- (12.0000,8.0000) -- (0.0000,8.0000) -- (12.0000,8.0000) -- (0.0000,8.0000) -- (12.0000,8.0000) -- (0.0000,8.0000) -- (12.0000,8.0000) -- (0.0000,8.0000) -- (12.0000,8.0000) -- (0.0000,8.0000) -- (12.0000,8.0000) -- (0.0000,8.0000) -- (12.0000,8.0000) -- (0.0000,8.0000) -- (12.0000,8.0000) -- (0.0000,8.0000) -- (12.0000,8.0000) -- (0.0000,8.0000) -- (12.0000,8.0000) -- (0.0000,8.0000) -- (12.0000,8.0000) -- (0.0000,8.0000) -- (12.0000,8.0000) -- (0.0000,8.0000) -- (12.0000,8.0000) -- (0.0000,8.0000) -- (12.0000,8.0000) -- (0.0000,8.0000) -- (12.0000,8.0000) -- (0.0000,8.0000) -- (12.0000,8.0000) -- (0.0000,8.0000) -- (12.0000,8.0000) -- (0.0000,8.0000) -- (12.0000,8.0000) -- (0.0000,8.0000) -- (12.0000,8.0000) -- (0.0000,8.0000) -- (12.0000,8.0000) -- (0.0000,8.0000) -- (12.0000,8.0000) -- (0.0000,8.0000) -- (12.0000,8.0000) -- (0.0000,8.0000) -- (12.0000,8.0000) -- (0.0000,8.0000) -- (12.0000,8.0000) -- (0.0000,8.0000) -- (12.0000,8.0000) -- (0.0000,8.0000) -- (12.0000,8.0000) -- (0.0000,8.0000) -- (12.0000,8.0000) -- (0.0000,8.0000) -- (12.0000,8.0000) -- (0.0000,8.0000) -- (12.0000,8.0000) -- (0.0000,8.0000) -- (12.0000,8.0000) -- (0.0000,8.0000) -- (12.0000,8.0000) -- (0.0000,8.0000) -- (12.0000,8.0000) -- (0.0000,8.0000) -- (12.0000,8.0000) -- (0.0000,8.0000) -- (12.0000,8.0000) -- (0.0000,8.0000) -- (12.0000,8.0000) -- (0.0000,8.0000) -- (12.0000,8.0000) -- (0.0000,8.0000) -- (12.0000,8.0000) -- (0.0000,8.0000) -- (12.0000,8.0000) -- (0.0000,8.0000) -- (12.0000,8.0000) -- (0.0000,8.0000) -- (12.0000,8.0000) -- (0.0000,8.0000) -- (12.0000,8.0000) -- (0.0000,8.0000) -- (12.0000,8.0000) -- (0.0000,8.0000) -- (12.0000,8.0000) -- (0.0000,8.0000) -- (12.0000,8.0000) -- (0.0000,8.0000) -- (12.0000,8.0000) -- (0.0000,8.0000) -- (12.0000,8.0000) -- (0.0000,8.0000) -- (12.0000,8.0000) -- (0.0000,8.0000) -- (12.0000,8.0000) -- (0.0000,8.0000) -- (12.0000,8.0000) -- (0.0000,8.0000) -- (12.0000,8.0000) -- (0.0000,8.0000) -- (12.0000,8.0000) -- (0.0000,8.0000) -- (12.0000,8.0000) -- (0.0000,8.0000) -- (12.0000,8.0000) -- (0.0000,8.0000) -- (12.0000,8.0000) -- (0.0000,8.0000) -- (12.0000,8.0000) -- (0.0000,8.0000) -- (12.0000,8.0000) -- (0.0000,8.0000) -- (12.0000,8.0000) -- (0.0000,8.0000) -- (12.0000,8.0000) -- (0.0000,8.0000) -- (12.0000,8.0000) -- (0.0000,8.0000) -- (12.0000,8.0000) -- (0.0000,8.0000) -- (12.0000,8.0000) -- (0.0000,8.0000) -- (12.0000,8.0000) -- (0.0000,8.0000) -- (12.0000,8.0000) -- (0.0000,8.0000) -- (12.0000,8.0000) -- (0.0000,8.0000) -- (12.0000,8.0000) -- (0.0000,8.0000) -- (12.0000,8.0000) -- (0.0000,8.0000) -- (12.0000,8.0000) -- (0.0000,8.0000) -- (12.0000,8.0000) -- (0.0000,8.0000) -- (12.0000,8.0000) -- (0.0000,8.0000) -- (12.0000,8.0000) -- (0.0000,8.0000) -- (12.0000,8.0000) -- (0.0000,8.0000) -- (12.0000,8.0000) -- (0.0000,8.0000) -- (12.0000,8.0000) -- (0.0000,8.0000) -- (12.0000,8.0000) -- (0.0000,8.0000) -- (12.0000,8.0000) -- (0.0000,8.0000) -- (12.0000,8.0000) -- (0.0000,8.0000) -- (12.0000,8.0000) -- (0.0000,8.0000) -- (12.0000,8.0000) -- (0.0000,8.0000) -- (12.0000,8.0000) -- (0.0000,8.0000) -- (12.0000,8.0000) -- (0.0000,8.0000) -- (12.0000,8.0000) -- (0.0000,8.0000) -- (12.0000,8.0000) -- (0.0000,8.0000) -- (12.0000,8.0000) -- (0.0000,8.0000) -- (12.0000,8.0000) -- (0.0000,8.0000) -- (12.0000,8.0000) -- (0.0000,8.0000) -- (12.0000,8.0000) -- (0.0000,8.0000) -- (12.0000,8.0000) -- (0.0000,8.0000) -- (12.0000,8.0000) -- (0.0000,8.0000) -- (12.0000,8.0000) -- (0.0000,8.0000) -- (12.0000,8.0000) -- (0.0000,8.0000) -- (12.0000,8.0000) -- (0.0000,8.0000) -- (12.0000,8.0000) -- (0.0000,8.0000) -- (12.0000,8.0000) -- (0.0000,8.0000) -- (12.0000,8.0000) -- (0.0000,8.0000) -- (12.0000,8.0000) -- (0.0000,8.0000) -- (12.0000,8.0000) -- (0.0000,8.0000) -- (12.0000,8.0000) -- (0.0000,8.0000) -- (12.0000,8.0000) -- (0.0000,8.0000) -- (12.0000,8.0000) -- (0.0000,8.0000) -- (12.0000,8.0000) -- (0.0000,8.0000) -- (12.0000,8.0000) -- (0.0000,8.0000) -- (12.0000,8.0000) -- (0.0000,8.0000) -- (12.0000,8.0000) -- (0.0000,8.0000) -- (12.0000,8.0000) -- (0.0000,8.0000) -- (12.0000,8.0000) -- (0.0000,8.0000) -- (12.0000,8.0000) -- (0.0000,8.0000) -- (12.0000,8.0000) -- (0.0000,8.0000) -- (12.0000,8.0000) -- (0.0000,8.0000) -- (12.0000,8.0000) -- (0.0000,8.0000) -- (12.0000,8.0000) -- (0.0000,8.0000) -- (12.0000,8.0000) -- (0.0000,8.0000) -- (12.0000,8.0000) -- (0.0000,8.0000) -- (12.0000,8.0000) -- (0.0000,8.0000) -- (12.0000,8.0000) -- (0.0000,8.0000) -- (12.0000,8.0000) -- (0.0000,8.0000) -- (12.0000,8.0000) -- (0.0000,8.0000) -- (12.0000,8.0000) -- (0.0000,8.0000) -- (12.0000,8.0000) -- (0.0000,8.0000) -- (12.0000,8.0000) -- (0.0000,8.0000) -- (12.0000,8.0000) -- (0.0000,8.0000) -- (12.0000,8.0000) -- (0.0000,8.0000) -- (12.0000,8.0000) -- (0.0000,8.0000) -- (12.0000,8.0000) -- (0.0000,8.0000) -- (12.0000,8.0000) -- (0.0000,8.0000) -- (12.0000,8.0000) -- (0.0000,8.0000) -- (12.0000,8.0000) -- (0.0000,8.0000) -- (12.0000,8.0000) -- (0.0000,8.0000) -- (12.0000,8.0000) -- (0.0000,8.0000) -- (12.0000,8.0000) -- (0.0000,8.0000) -- (12.0000,8.0000) -- (0.0000,8.0000) -- (12.0000,8.0000) -- (0.0000,8.0000) -- (12.0000,8.0000) -- (0.0000,8.0000) -- (12.0000,8.0000) -- (0.0000,8.0000) -- (12.0000,8.0000) -- (0.0000,8.0000) -- (12.0000,8.0000) -- (0.0000,8.0000) -- (12.0000,8.0000) -- (0.0000,8.0000) -- (12.0000,8.0000) -- (0.0000,8.0000) -- (12.0000,8.0000) -- (0.0000,8.0000) -- (12.0000,8.0000) -- (0.0000,8.0000) -- (12.0000,8.0000) -- (0.0000,8.0000) -- (12.0000,8.0000) -- (0.0000,8.0000) -- (12.0000,8.0000) -- (0.0000,8.0000) -- (12.0000,8.0000) -- (0.0000,8.0000) -- (12.0000,8.0000) -- (0.0000,8.0000) -- (12.0000,8.0000) -- (0.0000,8.0000) -- (12.0000,8.0000) -- (0.0000,8.0000) -- (12.0000,8.0000) -- (0.0000,8.0000) -- (12.0000,8.0000) -- (0.0000,8.0000) -- (12.0000,8.0000) -- (0.0000,8.0000) -- (12.0000,8.0000) -- (0.0000,8.0000) -- (12.0000,8.0000) -- (0.0000,8.0000) -- (12.0000,8.0000) -- (0.0000,8.0000) -- (12.0000,8.0000) -- (0.0000,8.0000) -- (12.0000,8.0000) -- (0.0000,8.0000) -- (12.0000,8.0000) -- (0.0000,8.0000) -- (12.0000,8.0000) -- (0.0000,8.0000) -- (12.0000,8.0000) -- (0.0000,8.0000) -- (12.0000,8.0000) -- (0.0000,8.0000) -- (12.0000,8.0000) -- (0.0000,8.0000) -- (12.0000,8.0000) -- (0.0000,8.0000) -- (12.0000,8.0000) -- (0.0000,8.0000) -- (12.0000,8.0000) -- (0.0000,8.0000) -- (12.0000,8.0000) -- (0.0000,8.0000) -- (12.0000,8.0000) -- (0.0000,8.0000) -- (12.0000,8.0000) -- (0.0000,8.0000) -- (12.0000,8.0000) -- (0.0000,8.0000) -- (12.0000,8.0000) -- (0.0000,8.0000) -- (12.0000,8.0000) -- (0.0000,8.0000) -- (12.0000,8.0000) -- (0.0000,8.0000) -- (12.0000,8.0000) -- (0.0000,8.0000) -- (12.0000,8.0000) -- (0.0000,8.0000) -- (12.0000,8.0000) -- (0.0000,8.0000) -- (12.0000,8.0000) -- (0.0000,8.0000) -- (12.0000,8.0000) -- (0.0000,8.0000) -- (12.0000,8.0000) -- (0.0000,8.0000) -- (12.0000,8.0000) -- (0.0000,8.0000) -- (12.0000,8.0000) -- (0.0000,8.0000) -- (12.0000,8.0000) -- (0.0000,8.0000) -- (12.0000,8.0000) -- (0.0000,8.0000) -- (12.0000,8.0000) -- (0.0000,8.0000) -- (12.0000,8.0000) -- (0.0000,8.0000) -- (12.0000,8.0000) -- (0.0000,8.0000) -- (12.0000,8.0000) -- (0.0000,8.0000) -- (12.0000,8.0000) -- (0.0000,8.0000) -- (12.0000,8.0000) -- (0.0000,8.0000) -- (12.0000,8.0000) -- (0.0000,8.0000) -- (12.0000,8.0000) -- (0.0000,8.0000) -- (12.0000,8.0000) -- (0.0000,8.0000) -- (12.0000,8.0000) -- (0.0000,8.0000) -- (12.0000,8.0000) -- (0.0000,8.0000) -- (12.0000,8.0000) -- (0.0000,8.0000) -- (12.0000,8.0000) -- (0.0000,8.0000) -- (12.0000,8.0000) -- (0.0000,8.0000) -- (12.0000,8.0000) -- (0.0000,8.0000) -- (12.0000,8.0000) -- (0.0000,8.0000) -- (12.0000,8.0000) -- (0.0000,8.0000) -- (12.0000,8.0000) -- (0.0000,8.0000) -- (12.0000,8.0000) -- (0.0000,8.0000) -- (12.0000,8.0000) -- (0.0000,8.0000) -- (12.0000,8.0000) -- (0.0000,8.0000) -- (12.0000,8.0000) -- (0.0000,8.0000) -- (12.0000,8.0000) -- (0.0000,8.0000) -- (12.0000,8.0000) -- (0.0000,8.0000) -- (12.0000,8.0000) -- (0.0000,8.0000) -- (12.0000,8.0000) -- (0.0000,8.0000) -- (12.0000,8.0000) -- (0.0000,8.0000) -- (12.0000,8.0000) -- (0.0000,8.0000) -- (12.0000,8.0000) -- (0.0000,8.0000) -- (12.0000,8.0000) -- (0.0000,8.0000) -- (12.0000,8.0000) -- (0.0000,8.0000) -- (12.0000,8.0000) -- (0.0000,8.0000) -- (12.0000,8.0000) -- (0.0000,8.0000) -- (12.0000,8.0000) -- (0.0000,8.0000) -- (12.0000,8.0000) -- (0.0000,8.0000) -- (12.0000,8.0000) -- (0.0000,8.0000) -- (12.0000,8.0000) -- (0.0000,8.0000) -- (12.0000,8.0000) -- (0.0000,8.0000) -- (12.0000,8.0000) -- (0.0000,8.0000) -- (12.0000,8.0000) -- (0.0000,8.0000) -- (12.0000,8.0000) -- (0.0000,8.0000) -- (12.0000,8.0000) -- (0.0000,8.0000) -- (12.0000,8.0000) -- (0.0000,8.0000) -- (12.0000,8.0000) -- (0.0000,8.0000) -- (12.0000,8.0000) -- (0.0000,8.0000) -- (12.0000,8.0000) -- (0.0000,8.0000) -- (12.0000,8.0000) -- (0.0000,8.0000) -- (12.0000,8.0000) -- (0.0000,8.0000) -- (12.0000,8.0000) -- (0.0000,8.0000) -- (12.0000,8.0000) -- (0.0000,8.0000) -- (12.0000,8.0000) -- (0.0000,8.0000) -- (12.0000,8.0000) -- (0.0000,8.0000) -- (12.0000,8.0000) -- (0.0000,8.0000) -- (12.0000,8.0000) -- (0.0000,8.0000) -- (12.0000,8.0000) -- (0.0000,8.0000) -- (12.0000,8.0000) -- (0.0000,8.0000) -- (12.0000,8.0000) -- (0.0000,8.0000) -- (12.0000,8.0000) -- (0.0000,8.0000) -- (12.0000,8.0000) -- (0.0000,8.0000) -- (12.0000,8.0000) -- (0.0000,8.0000) -- (12.0000,8.0000) -- (0.0000,8.0000) -- (12.0000,8.0000) -- (0.0000,8.0000) -- (12.0000,8.0000) -- (0.0000,8.0000) -- (12.0000,8.0000) -- (0.0000,8.0000) -- (12.0000,8.0000) -- (0.0000,8.0000) -- (12.0000,8.0000) -- (0.0000,8.0000) -- (12.0000,8.0000) -- (0.0000,8.0000) -- (12.0000,8.0000) -- (0.0000,8.0000) -- (12.0000,8.0000) -- (0.0000,8.0000) -- (12.0000,8.0000) -- (0.0000,8.0000) -- (12.0000,8.0000) -- (0.0000,8.0000) -- (12.0000,8.0000) -- (0.0000,8.0000) -- (12.0000,8.0000) -- (0.0000,8.0000) -- (12.0000,8.0000) -- (0.0000,8.0000) -- (12.0000,8.0000) -- (0.0000,8.0000) -- (12.0000,8.0000) -- (0.0000,8.0000) -- (12.0000,8.0000) -- (0.0000,8.0000) -- (12.0000,8.0000) -- (0.0000,8.0000) -- (12.0000,8.0000) -- (0.0000,8.0000) -- (12.0000,8.0000) -- (0.0000,8.0000) -- (12.0000,8.0000) -- (0.0000,8.0000) -- (12.0000,8.0000) -- (0.0000,8.0000) -- (12.0000,8.0000) -- (0.0000,8.0000) -- (12.0000,8.0000) -- (0.0000,8.0000) -- (12.0000,8.0000) -- (0.0000,8.0000) -- (12.0000,8.0000) -- (0.0000,8.0000) -- (12.0000,8.0000) -- (0.0000,8.0000) -- (12.0000,8.0000) -- (0.0000,8.0000) -- (12.0000,8.0000) -- (0.0000,8.0000) -- (12.0000,8.0000) -- (0.0000,8.0000) -- (12.0000,8.0000) -- (0.0000,8.0000) -- (12.0000,8.0000) -- (0.0000,8.0000) -- (12.0000,8.0000) -- (0.0000,8.0000) -- (12.0000,8.0000) -- (0.0000,8.0000) -- (12.0000,8.0000) -- (0.0000,8.0000) -- (12.0000,8.0000) -- (0.0000,8.0000) -- (12.0000,8.0000) -- (0.0000,8.0000) -- (12.0000,8.0000) -- (0.0000,8.0000) -- (12.0000,8.0000) -- (0.0000,8.0000) -- (12.0000,8.0000) -- (0.0000,8.0000) -- (12.0000,8.0000) -- (0.0000,8.0000) -- (12.0000,8.0000) -- (0.0000,8.0000) -- (12.0000,8.0000) -- (0.0000,8.0000) -- (12.0000,8.0000) -- (0.0000,8.0000) -- (12.0000,8.0000) -- (0.0000,8.0000) -- (12.0000,8.0000) -- (0.0000,8.0000) -- (12.0000,8.0000) -- (0.0000,8.0000) -- (12.0000,8.0000) -- (0.0000,8.0000) -- (12.0000,8.0000) -- (0.0000,8.0000) -- (12.0000,8.0000) -- (0.0000,8.0000) -- (12.0000,8.0000) -- (0.0000,8.0000) -- (12.0000,8.0000) -- (0.0000,8.0000) -- (12.0000,8.0000) -- (0.0000,8.0000) -- (12.0000,8.0000) -- (0.0000,8.0000) -- (12.0000,8.0000) -- (0.0000,8.0000) -- (12.0000,8.0000) -- (0.0000,8.0000) -- (12.0000,8.0000) -- (0.0000,8.0000) -- (12.0000,8.0000) -- (0.0000,8.0000) -- (12.0000,8.0000) -- (0.0000,8.0000) -- (12.0000,8.0000) -- (0.0000,8.0000) -- (12.0000,8.0000) -- (0.0000,8.0000) -- (12.0000,8.0000) -- (0.0000,8.0000) -- (12.0000,8.0000) -- (0.0000,8.0000) -- (12.0000,8.0000) -- (0.0000,8.0000) -- (12.0000,8.0000) -- (0.0000,8.0000) -- (12.0000,8.0000) -- (0.0000,8.0000) -- (12.0000,8.0000) -- (0.0000,8.0000) -- (12.0000,8.0000) -- (0.0000,8.0000) -- (12.0000,8.0000) -- (0.0000,8.0000) -- (12.0000,8.0000) -- (0.0000,8.0000) -- (12.0000,8.0000) -- (0.0000,8.0000) -- (12.0000,8.0000) -- (0.0000,8.0000) -- (12.0000,8.0000) -- (0.0000,8.0000) -- (12.0000,8.0000) -- (0.0000,8.0000) -- (12.0000,8.0000) -- (0.0000,8.0000) -- (12.0000,8.0000) -- (0.0000,8.0000) -- (12.0000,8.0000) -- (0.0000,8.0000) -- (12.0000,8.0000) -- (0.0000,8.0000) -- (12.0000,8.0000) -- (0.0000,8.0000) -- (12.0000,8.0000) -- (0.0000,8.0000) -- (12.0000,8.0000) -- (0.0000,8.0000) -- (12.0000,8.0000) -- (0.0000,8.0000) -- (12.0000,8.0000) -- (0.0000,8.0000) -- (12.0000,8.0000) -- (0.0000,8.0000) -- (12.0000,8.0000) -- (0.0000,8.0000) -- (12.0000,8.0000) -- (0.0000,8.0000) -- (12.0000,8.0000) -- (0.0000,8.0000) -- (12.0000,8.0000) -- (0.0000,8.0000) -- (12.0000,8.0000) -- (0.0000,8.0000) -- (12.0000,8.0000) -- (0.0000,8.0000) -- (12.0000,8.0000) -- (0.0000,8.0000) -- (12.0000,8.0000) -- (0.0000,8.0000) -- (12.0000,8.0000) -- (0.0000,8.0000) -- (12.0000,8.0000) -- (0.0000,8.0000) -- (12.0000,8.0000) -- (0.0000,8.0000) -- (12.0000,8.0000) -- (0.0000,8.0000) -- (12.0000,8.0000) -- (0.0000,8.0000) -- (12.0000,8.0000) -- (0.0000,8.0000) -- (12.0000,8.0000) -- (0.0000,8.0000) -- (12.0000,8.0000) -- (0.0000,8.0000) -- (12.0000,8.0000) -- (0.0000,8.0000) -- (12.0000,8.0000) -- (0.0000,8.0000) -- (12.0000,8.0000) -- (0.0000,8.0000) -- (12.0000,8.0000) -- (0.0000,8.0000) -- (12.0000,8.0000) -- (0.0000,8.0000) -- (12.0000,8.0000) -- (0.0000,8.0000) -- (12.0000,8.0000) -- (0.0000,8.0000) -- (12.0000,8.0000) -- (0.0000,8.0000) -- (12.0000,8.0000) -- (0.0000,8.0000) -- (12.0000,8.0000) -- (0.0000,8.0000) -- (12.0000,8.0000) -- (0.0000,8.0000) -- (12.0000,8.0000) -- (0.0000,8.0000) -- (12.0000,8.0000) -- (0.0000,8.0000) -- (12.0000,8.0000) -- (0.0000,8.0000) -- (12.0000,8.0000) -- (0.0000,8.0000) -- (12.0000,8.0000) -- (0.0000,8.0000) -- (12.0000,8.0000) -- (0.0000,8.0000) -- (12.0000,8.0000) -- (0.0000,8.0000) -- (12.0000,8.0000) -- (0.0000,8.0000) -- (12.0000,8.0000) -- (0.0000,8.0000) -- (12.0000,8.0000) -- (0.0000,8.0000) -- (12.0000,8.0000) -- (0.0000,8.0000) -- (12.0000,8.0000) -- (0.0000,8.0000) -- (12.0000,8.0000) -- (0.0000,8.0000) -- (12.0000,8.0000) -- (0.0000,8.0000) -- (12.0000,8.0000) -- (0.0000,8.0000) -- (12.0000,8.0000) -- (0.0000,8.0000) -- (12.0000,8.0000) -- (0.0000,8.0000) -- (12.0000,8.0000) -- (0.0000,8.0000) -- (12.0000,8.0000) -- (0.0000,8.0000) -- (12.0000,8.0000) -- (0.0000,8.0000) -- (12.0000,8.0000) -- (0.0000,8.0000) -- (12.0000,8.0000) -- (0.0000,8.0000) -- (12.0000,8.0000) -- (0.0000,8.0000) -- (12.0000,8.0000) -- (0.0000,8.0000) -- (12.0000,8.0000) -- (0.0000,8.0000) -- (12.0000,8.0000) -- (0.0000,8.0000) -- (12.0000,8.0000) -- (0.0000,8.0000) -- (12.0000,8.0000) -- (0.0000,8.0000) -- (12.0000,8.0000) -- (0.0000,8.0000) -- (12.0000,8.0000) -- (0.0000,8.0000) -- (12.0000,8.0000) -- (0.0000,8.0000) -- (12.0000,8.0000) -- (0.0000,8.0000) -- (12.0000,8.0000) -- (0.0000,8.0000) -- (12.0000,8.0000) -- (0.0000,8.0000) -- (12.0000,8.0000) -- (0.0000,8.0000) -- (12.0000,8.0000) -- (0.0000,8.0000) -- (12.0000,8.0000) -- (0.0000,8.0000) -- (12.0000,8.0000) -- (0.0000,8.0000) -- (12.0000,8.0000) -- (0.0000,8.0000) -- (12.0000,8.0000) -- (0.0000,8.0000) -- (12.0000,8.0000) -- (0.0000,8.0000) -- (12.0000,8.0000) -- (0.0000,8.0000) -- (12.0000,8.0000) -- (0.0000,8.0000) -- (12.0000,8.0000) -- (0.0000,8.0000) -- (12.0000,8.0000) -- (0.0000,8.0000) -- (12.0000,8.0000) -- (0.0000,8.0000) -- (12.0000,8.0000) -- (0.0000,8.0000) -- (12.0000,8.0000) -- (0.0000,8.0000) -- (12.0000,8.0000) -- (0.0000,8.0000) -- (12.0000,8.0000) -- (0.0000,8.0000) -- (12.0000,8.0000) -- (0.0000,8.0000) -- (12.0000,8.0000) -- (0.0000,8.0000) -- (12.0000,8.0000) -- (0.0000,8.0000) -- (12.0000,8.0000) -- (0.0000,8.0000) -- (12.0000,8.0000) -- (0.0000,8.0000) -- (12.0000,8.0000) -- (0.0000,8.0000) -- (12.0000,8.0000) -- (0.0000,8.0000) -- (12.0000,8.0000) -- (0.0000,8.0000) -- (12.0000,8.0000) -- (0.0000,8.0000) -- (12.0000,8.0000) -- (0.0000,8.0000) -- (12.0000,8.0000) -- (0.0000,8.0000) -- (12.0000,8.0000) -- (0.0000,8.0000) -- (12.0000,8.0000) -- (0.0000,8.0000) -- (12.0000,8.0000) -- (0.0000,8.0000) -- (12.0000,8.0000) -- (0.0000,8.0000) -- (12.0000,8.0000) -- (0.0000,8.0000) -- (12.0000,8.0000) -- (0.0000,8.0000) -- (12.0000,8.0000) -- (0.0000,8.0000) -- (12.0000,8.0000) -- (0.0000,8.0000) -- (12.0000,8.0000) -- (0.0000,8.0000) -- (12.0000,8.0000) -- (0.0000,8.0000) -- (12.0000,8.0000) -- (0.0000,8.0000) -- (12.0000,8.0000) -- (0.0000,8.0000) -- (12.0000,8.0000) -- (0.0000,8.0000) -- (12.0000,8.0000) -- (0.0000,8.0000) -- (12.0000,8.0000) -- (0.0000,8.0000) -- (12.0000,8.0000) -- (0.0000,8.0000) -- (12.0000,8.0000) -- (0.0000,8.0000) -- (12.0000,8.0000) -- (0.0000,8.0000) -- (12.0000,8.0000) -- (0.0000,8.0000) -- (12.0000,8.0000) -- (0.0000,8.0000) -- (12.0000,8.0000) -- (0.0000,8.0000) -- (12.0000,8.0000) -- (0.0000,8.0000) -- (12.0000,8.0000) -- (0.0000,8.0000) -- (12.0000,8.0000) -- (0.0000,8.0000) -- (12.0000,8.0000) -- (0.0000,8.0000) -- (12.0000,8.0000) -- (0.0000,8.0000) -- (12.0000,8.0000) -- (0.0000,8.0000) -- (12.0000,8.0000) -- (0.0000,8.0000) -- (12.0000,8.0000) -- (0.0000,8.0000) -- (12.0000,8.0000) -- (0.0000,8.0000) -- (12.0000,8.0000) -- (0.0000,8.0000) -- (12.0000,8.0000) -- (0.0000,8.0000) -- (12.0000,8.0000) -- (0.0000,8.0000) -- (12.0000,8.0000) -- (0.0000,8.0000) -- (12.0000,8.0000) -- (0.0000,8.0000) -- (12.0000,8.0000) -- (0.0000,8.0000) -- (12.0000,8.0000) -- (0.0000,8.0000) -- (12.0000,8.0000) -- (0.0000,8.0000) -- (12.0000,8.0000) -- (0.0000,8.0000) -- (12.0000,8.0000) -- (0.0000,8.0000) -- (12.0000,8.0000) -- (0.0000,8.0000) -- (12.0000,8.0000) -- (0.0000,8.0000) -- (12.0000,8.0000) -- (0.0000,8.0000) -- (12.0000,8.0000) -- (0.0000,8.0000) -- (12.0000,8.0000) -- (0.0000,8.0000) -- (12.0000,8.0000) -- (0.0000,8.0000) -- (12.0000,8.0000) -- (0.0000,8.0000) -- (12.0000,8.0000) -- (0.0000,8.0000) -- (12.0000,8.0000) -- (0.0000,8.0000) -- (12.0000,8.0000) -- (0.0000,8.0000) -- (12.0000,8.0000) -- (0.0000,8.0000) -- (12.0000,8.0000) -- (0.0000,8.0000) -- (12.0000,8.0000) -- (0.0000,8.0000) -- (12.0000,8.0000) -- (0.0000,8.0000) -- (12.0000,8.0000) -- (0.0000,8.0000) -- (12.0000,8.0000) -- (0.0000,8.0000) -- (12.0000,8.0000) -- (0.0000,8.0000) -- (12.0000,8.0000) -- (0.0000,8.0000) -- (12.0000,8.0000) -- (0.0000,8.0000) -- (12.0000,8.0000) -- (0.0000,8.0000) -- (12.0000,8.0000) -- (0.0000,8.0000) -- (12.0000,8.0000) -- (0.0000,8.0000) -- (12.0000,8.0000) -- (0.0000,8.0000) -- (12.0000,8.0000) -- (0.0000,8.0000) -- (12.0000,8.0000) -- (0.0000,8.0000) -- (12.0000,8.0000) -- (0.0000,8.0000) -- (12.0000,8.0000) -- (0.0000,8.0000) -- (12.0000,8.0000) -- (0.0000,8.0000) -- (12.0000,8.0000) -- (0.0000,8.0000) -- (12.0000,8.0000) -- (0.0000,8.0000) -- (12.0000,8.0000) -- (0.0000,8.0000) -- (12.0000,8.0000) -- (0.0000,8.0000) -- (12.0000,8.0000) -- (0.0000,8.0000) -- (12.0000,8.0000) -- (0.0000,8.0000) -- (12.0000,8.0000) -- (0.0000,8.0000) -- (12.0000,8.0000) -- (0.0000,8.0000) -- (12.0000,8.0000) -- (0.0000,8.0000) -- (12.0000,8.0000) -- (0.0000,8.0000) -- (12.0000,8.0000) -- (0.0000,8.0000) -- (12.0000,8.0000) -- (0.0000,8.0000) -- (12.0000,8.0000) -- (0.0000,8.0000) -- (12.0000,8.0000) -- (0.0000,8.0000) -- (12.0000,8.0000) -- (0.0000,8.0000) -- (12.0000,8.0000) -- (0.0000,8.0000) -- (12.0000,8.0000) -- (0.0000,8.0000) -- (12.0000,8.0000) -- (0.0000,8.0000) -- (12.0000,8.0000) -- (0.0000,8.0000) -- (12.0000,8.0000) -- (0.0000,8.0000) -- (12.0000,8.0000) -- (0.0000,8.0000) -- (12.0000,8.0000) -- (0.0000,8.0000) -- (12.0000,8.0000) -- (0.0000,8.0000) -- (12.0000,8.0000) -- (0.0000,8.0000) -- (12.0000,8.0000) -- (0.0000,8.0000) -- (12.0000,8.0000) -- (0.0000,8.0000) -- (12.0000,8.0000) -- (0.0000,8.0000) -- (12.0000,8.0000) -- (0.0000,8.0000) -- (12.0000,8.0000) -- (0.0000,8.0000) -- (12.0000,8.0000) -- (0.0000,8.0000) -- (12.0000,8.0000) -- (0.0000,8.0000) -- (12.0000,8.0000) -- (0.0000,8.0000) -- (12.0000,8.0000) -- (0.0000,8.0000) -- (12.0000,8.0000) -- (0.0000,8.0000) -- (12.0000,8.0000) -- (0.0000,8.0000) -- (12.0000,8.0000) -- (0.0000,8.0000) -- (12.0000,8.0000) -- (0.0000,8.0000) -- (12.0000,8.0000) -- (0.0000,8.0000) -- (12.0000,8.0000) -- (0.0000,8.0000) -- (12.0000,8.0000) -- (0.0000,8.0000) -- (12.0000,8.0000) -- (0.0000,8.0000) -- (12.0000,8.0000) -- (0.0000,8.0000) -- (12.0000,8.0000) -- (0.0000,8.0000) -- (12.0000,8.0000) -- (0.0000,8.0000) -- (12.0000,8.0000) -- (0.0000,8.0000) -- (12.0000,8.0000) -- (0.0000,8.0000) -- (12.0000,8.0000) -- (0.0000,8.0000) -- (12.0000,8.0000) -- (0.0000,8.0000) -- (12.0000,8.0000) -- (0.0000,8.0000) -- (12.0000,8.0000) -- (0.0000,8.0000) -- (12.0000,8.0000) -- (0.0000,8.0000) -- (12.0000,8.0000) -- (0.0000,8.0000) -- (12.0000,8.0000) -- (0.0000,8.0000) -- (12.0000,8.0000) -- (0.0000,8.0000) -- (12.0000,8.0000) -- (0.0000,8.0000) -- (12.0000,8.0000) -- (0.0000,8.0000) -- (12.0000,8.0000) -- (0.0000,8.0000) -- (12.0000,8.0000) -- (0.0000,8.0000) -- (12.0000,8.0000) -- (0.0000,8.0000) -- (12.0000,8.0000) -- (0.0000,8.0000) -- (12.0000,8.0000) -- (0.0000,8.0000) -- (12.0000,8.0000) -- (0.0000,8.0000) -- (12.0000,8.0000) -- (0.0000,8.0000) -- (12.0000,8.0000) -- (0.0000,8.0000) -- (12.0000,8.0000) -- (0.0000,8.0000) -- (12.0000,8.0000) -- (0.0000,8.0000) -- (12.0000,8.0000) -- (0.0000,8.0000) -- (12.0000,8.0000) -- (0.0000,8.0000) -- (12.0000,8.0000) -- (0.0000,8.0000) -- (12.0000,8.0000) -- (0.0000,8.0000) -- (12.0000,8.0000) -- (0.0000,8.0000) -- (12.0000,8.0000) -- (0.0000,8.0000) -- (12.0000,8.0000) -- (0.0000,8.0000) -- (12.0000,8.0000) -- (0.0000,8.0000) -- (12.0000,8.0000) -- (0.0000,8.0000) -- (12.0000,8.0000) -- (0.0000,8.0000) -- (12.0000,8.0000) -- (0.0000,8.0000) -- (12.0000,8.0000) -- (0.0000,8.0000) -- (12.0000,8.0000) -- (0.0000,8.0000) -- (12.0000,8.0000) -- (0.0000,8.0000) -- (12.0000,8.0000) -- (0.0000,8.0000) -- (12.0000,8.0000) -- (0.0000,8.0000) -- (12.0000,8.0000) -- (0.0000,8.0000) -- (12.0000,8.0000) -- (0.0000,8.0000) -- (12.0000,8.0000) -- (0.0000,8.0000) -- (12.0000,8.0000) -- (0.0000,8.0000) -- (12.0000,8.0000) -- (0.0000,8.0000) -- (12.0000,8.0000) -- (0.0000,8.0000) -- (12.0000,8.0000) -- (0.0000,8.0000) -- (12.0000,8.0000) -- (0.0000,8.0000) -- (12.0000,8.0000) -- (0.0000,8.0000) -- (12.0000,8.0000) -- (0.0000,8.0000) -- (12.0000,8.0000) -- (0.0000,8.0000) -- (12.0000,8.0000) -- (0.0000,8.0000) -- (12.0000,8.0000) -- (0.0000,8.0000) -- (12.0000,8.0000) -- (0.0000,8.0000) -- (12.0000,8.0000) -- (0.0000,8.0000) -- (12.0000,8.0000) -- (0.0000,8.0000) -- (12.0000,8.0000) -- (0.0000,8.0000) -- (12.0000,8.0000) -- (0.0000,8.0000) -- (12.0000,8.0000) -- (0.0000,8.0000) -- (12.0000,8.0000) -- (0.0000,8.0000) -- (12.0000,8.0000) -- (0.0000,8.0000) -- (12.0000,8.0000) -- (0.0000,8.0000) -- (12.0000,8.0000) -- (0.0000,8.0000) -- (12.0000,8.0000) -- (0.0000,8.0000) -- (12.0000,8.0000) -- (0.0000,8.0000) -- (12.0000,8.0000) -- (0.0000,8.0000) -- (12.0000,8.0000) -- (0.0000,8.0000) -- (12.0000,8.0000) -- (0.0000,8.0000) -- (12.0000,8.0000) -- (0.0000,8.0000) -- (12.0000,8.0000) -- (0.0000,8.0000) -- (12.0000,8.0000) -- (0.0000,8.0000) -- (12.0000,8.0000) -- (0.0000,8.0000) -- (12.0000,8.0000) -- (0.0000,8.0000) -- (12.0000,8.0000) -- (0.0000,8.0000) -- (12.0000,8.0000) -- (0.0000,8.0000) -- (12.0000,8.0000) -- (0.0000,8.0000) -- (12.0000,8.0000) -- (0.0000,8.0000) -- (12.0000,8.0000) -- (0.0000,8.0000) -- (12.0000,8.0000) -- (0.0000,8.0000) -- (12.0000,8.0000) -- (0.0000,8.0000) -- (12.0000,8.0000) -- (0.0000,8.0000) -- (12.0000,8.0000) -- (0.0000,8.0000) -- (12.0000,8.0000) -- (0.0000,8.0000) -- (12.0000,8.0000) -- (0.0000,8.0000) -- (12.0000,8.0000) -- (0.0000,8.0000) -- (12.0000,8.0000) -- (0.0000,8.0000) -- (12.0000,8.0000) -- (0.0000,8.0000) -- (12.0000,8.0000) -- (0.0000,8.0000) -- (12.0000,8.0000) -- (0.0000,8.0000) -- (12.0000,8.0000) -- (0.0000,8.0000) -- (12.0000,8.0000) -- (0.0000,8.0000) -- (12.0000,8.0000) -- (0.0000,8.0000) -- (12.0000,8.0000) -- (0.0000,8.0000) -- (12.0000,8.0000) -- (0.0000,8.0000) -- (12.0000,8.0000) -- (0.0000,8.0000) -- (12.0000,8.0000) -- (0.0000,8.0000) -- (12.0000,8.0000) -- (0.0000,8.0000) -- (12.0000,8.0000) -- (0.0000,8.0000) -- (12.0000,8.0000) -- (0.0000,8.0000) -- (12.0000,8.0000) -- (0.0000,8.0000) -- (12.0000,8.0000) -- (0.0000,8.0000) -- (12.0000,8.0000) -- (0.0000,8.0000) -- (12.0000,8.0000) -- (0.0000,8.0000) -- (12.0000,8.0000) -- (0.0000,8.0000) -- (12.0000,8.0000) -- (0.0000,8.0000) -- (12.0000,8.0000) -- (0.0000,8.0000) -- (12.0000,8.0000) -- (0.0000,8.0000) -- (12.0000,8.0000) -- (0.0000,8.0000) -- (12.0000,8.0000) -- (0.0000,8.0000) -- (12.0000,8.0000) -- (0.0000,8.0000) -- (12.0000,8.0000) -- (0.0000,8.0000) -- (12.0000,8.0000) -- (0.0000,8.0000) -- (12.0000,8.0000) -- (0.0000,8.0000) -- (12.0000,8.0000) -- (0.0000,8.0000) -- (12.0000,8.0000) -- (0.0000,8.0000) -- (12.0000,8.0000) -- (0.0000,8.0000) -- (12.0000,8.0000) -- (0.0000,8.0000) -- (12.0000,8.0000) -- (0.0000,8.0000) -- (12.0000,8.0000) -- (0.0000,8.0000) -- (12.0000,8.0000) -- (0.0000,8.0000) -- (12.0000,8.0000) -- (0.0000,8.0000) -- (12.0000,8.0000) -- (0.0000,8.0000) -- (12.0000,8.0000) -- (0.0000,8.0000) -- (12.0000,8.0000) -- (0.0000,8.0000) -- (12.0000,8.0000) -- (0.0000,8.0000) -- (12.0000,8.0000) -- (0.0000,8.0000) -- (12.0000,8.0000) -- (0.0000,8.0000) -- (12.0000,8.0000) -- (0.0000,8.0000) -- (12.0000,8.0000) -- (0.0000,8.0000) -- (12.0000,8.0000) -- (0.0000,8.0000) -- (12.0000,8.0000) -- (0.0000,8.0000) -- (12.0000,8.0000) -- (0.0000,8.0000) -- (12.0000,8.0000) -- (0.0000,8.0000) -- (12.0000,8.0000) -- (0.0000,8.0000) -- (12.0000,8.0000) -- (0.0000,8.0000) -- (12.0000,8.0000) -- (0.0000,8.0000) -- (12.0000,8.0000) -- (0.0000,8.0000) -- (12.0000,8.0000) -- (0.0000,8.0000) -- (12.0000,8.0000) -- (0.0000,8.0000) -- (12.0000,8.0000) -- (0.0000,8.0000) -- (12.0000,8.0000) -- (0.0000,8.0000) -- (12.0000,8.0000) -- (0.0000,8.0000) -- (12.0000,8.0000) -- (0.0000,8.0000) -- (12.0000,8.0000) -- (0.0000,8.0000) -- (12.0000,8.0000) -- (0.0000,8.0000) -- (12.0000,8.0000) -- (0.0000,8.0000) -- (12.0000,8.0000) -- (0.0000,8.0000) -- (12.0000,8.0000) -- (0.0000,8.0000) -- (12.0000,8.0000) -- (0.0000,8.0000) -- (12.0000,8.0000) -- (0.0000,8.0000) -- (12.0000,8.0000) -- (0.0000,8.0000) -- (12.0000,8.0000) -- (0.0000,8.0000) -- (12.0000,8.0000) -- (0.0000,8.0000) -- (12.0000,8.0000) -- (0.0000,8.0000) -- (12.0000,8.0000) -- (0.0000,8.0000) -- (12.0000,8.0000) -- (0.0000,8.0000) -- (12.0000,8.0000) -- (0.0000,8.0000) -- (12.0000,8.0000) -- (0.0000,8.0000) -- (12.0000,8.0000) -- (0.0000,8.0000) -- (12.0000,8.0000) -- (0.0000,8.0000) -- (12.0000,8.0000) -- (0.0000,8.0000) -- (12.0000,8.0000) -- (0.0000,8.0000) -- (12.0000,8.0000) -- (0.0000,8.0000) -- (12.0000,8.0000) -- (0.0000,8.0000) -- (12.0000,8.0000) -- (0.0000,8.0000) -- (12.0000,8.0000) -- (0.0000,8.0000) -- (12.0000,8.0000) -- (0.0000,8.0000) -- (12.0000,8.0000) -- (0.0000,8.0000) -- (12.0000,8.0000) -- (0.0000,8.0000) -- (12.0000,8.0000) -- (0.0000,8.0000) -- (12.0000,8.0000) -- (0.0000,8.0000) -- (12.0000,8.0000) -- (0.0000,8.0000) -- (12.0000,8.0000) -- (0.0000,8.0000) -- (12.0000,8.0000) -- (0.0000,8.0000) -- (12.0000,8.0000) -- (0.0000,8.0000) -- (12.0000,8.0000) -- (0.0000,8.0000) -- (12.0000,8.0000) -- (0.0000,8.0000) -- (12.0000,8.0000) -- (0.0000,8.0000) -- (12.0000,8.0000) -- (0.0000,8.0000) -- (12.0000,8.0000) -- (0.0000,8.0000) -- (12.0000,8.0000) -- (0.0000,8.0000) -- (12.0000,8.0000) -- (0.0000,8.0000) -- (12.0000,8.0000) -- (0.0000,8.0000) -- (12.0000,8.0000) -- (0.0000,8.0000) -- (12.0000,8.0000) -- (0.0000,8.0000) -- (12.0000,8.0000) -- (0.0000,8.0000) -- (12.0000,8.0000) -- (0.0000,8.0000) -- (12.0000,8.0000) -- (0.0000,8.0000) -- (12.0000,8.0000) -- (0.0000,8.0000) -- (12.0000,8.0000) -- (0.0000,8.0000) -- (12.0000,8.0000) -- (0.0000,8.0000) -- (12.0000,8.0000) -- (0.0000,8.0000) -- (12.0000,8.0000) -- (0.0000,8.0000) -- (12.0000,8.0000) -- (0.0000,8.0000) -- (12.0000,8.0000) -- (0.0000,8.0000) -- (12.0000,8.0000) -- (0.0000,8.0000) -- (12.0000,8.0000) -- (0.0000,8.0000) -- (12.0000,8.0000) -- (0.0000,8.0000) -- (12.0000,8.0000) -- (0.0000,8.0000) -- (12.0000,8.0000) -- (0.0000,8.0000) -- (12.0000,8.0000) -- (0.0000,8.0000) -- (12.0000,8.0000) -- (0.0000,8.0000) -- (12.0000,8.0000) -- (0.0000,8.0000) -- (12.0000,8.0000) -- (0.0000,8.0000) -- (12.0000,8.0000) -- (0.0000,8.0000) -- (12.0000,8.0000) -- (0.0000,8.0000) -- (12.0000,8.0000) -- (0.0000,8.0000) -- (12.0000,8.0000) -- (0.0000,8.0000) -- (12.0000,8.0000) -- (0.0000,8.0000) -- (12.0000,8.0000) -- (0.0000,8.0000) -- (12.0000,8.0000) -- (0.0000,8.0000) -- (12.0000,8.0000) -- (0.0000,8.0000) -- (12.0000,8.0000) -- (0.0000,8.0000) -- (12.0000,8.0000) -- (0.0000,8.0000) -- (12.0000,8.0000) -- (0.0000,8.0000) -- (12.0000,8.0000) -- (0.0000,8.0000) -- (12.0000,8.0000) -- (0.0000,8.0000) -- (12.0000,8.0000) -- (0.0000,8.0000) -- (12.0000,8.0000) -- (0.0000,8.0000) -- (12.0000,8.0000) -- (0.0000,8.0000) -- (12.0000,8.0000) -- (0.0000,8.0000) -- (12.0000,8.0000) -- (0.0000,8.0000) -- (12.0000,8.0000) -- (0.0000,8.0000) -- (12.0000,8.0000) -- (0.0000,8.0000) -- (12.0000,8.0000) -- (0.0000,8.0000) -- (12.0000,8.0000) -- (0.0000,8.0000) -- (12.0000,8.0000) -- (0.0000,8.0000) -- (12.0000,8.0000) -- (0.0000,8.0000) -- (12.0000,8.0000) -- (0.0000,8.0000) -- (12.0000,8.0000) -- (0.0000,8.0000) -- (12.0000,8.0000) -- (0.0000,8.0000) -- (12.0000,8.0000) -- (0.0000,8.0000) -- (12.0000,8.0000) -- (0.0000,8.0000) -- (12.0000,8.0000) -- (0.0000,8.0000) -- (12.0000,8.0000) -- (0.0000,8.0000) -- (12.0000,8.0000) -- (0.0000,8.0000) -- (12.0000,8.0000) -- (0.0000,8.0000) -- (12.0000,8.0000) -- (0.0000,8.0000) -- (12.0000,8.0000) -- (0.0000,8.0000) -- (12.0000,8.0000) -- (0.0000,8.0000) -- (12.0000,8.0000) -- (0.0000,8.0000) -- (12.0000,8.0000) -- (0.0000,8.0000) -- (12.0000,8.0000) -- (0.0000,8.0000) -- (12.0000,8.0000) -- (0.0000,8.0000) -- (12.0000,8.0000) -- (0.0000,8.0000) -- (12.0000,8.0000) -- (0.0000,8.0000) -- (12.0000,8.0000) -- (0.0000,8.0000) -- (12.0000,8.0000) -- (0.0000,8.0000) -- (12.0000,8.0000) -- (0.0000,8.0000) -- (12.0000,8.0000) -- (0.0000,8.0000) -- (12.0000,8.0000) -- (0.0000,8.0000) -- (12.0000,8.0000) -- (0.0000,8.0000) -- (12.0000,8.0000) -- (0.0000,8.0000) -- (12.0000,8.0000) -- (0.0000,8.0000) -- (12.0000,8.0000) -- (0.0000,8.0000) -- (12.0000,8.0000) -- (0.0000,8.0000) -- (12.0000,8.0000) -- (0.0000,8.0000) -- (12.0000,8.0000) -- (0.0000,8.0000) -- (12.0000,8.0000) -- (0.0000,8.0000) -- (12.0000,8.0000) -- (0.0000,8.0000) -- (12.0000,8.0000) -- (0.0000,8.0000) -- (12.0000,8.0000) -- (0.0000,8.0000) -- (12.0000,8.0000) -- (0.0000,8.0000) -- (12.0000,8.0000) -- (0.0000,8.0000) -- (12.0000,8.0000) -- (0.0000,8.0000) -- (12.0000,8.0000) -- (0.0000,8.0000) -- (12.0000,8.0000) -- (0.0000,8.0000) -- (12.0000,8.0000) -- (0.0000,8.0000) -- (12.0000,8.0000) -- (0.0000,8.0000) -- (12.0000,8.0000) -- (0.0000,8.0000) -- (12.0000,8.0000) -- (0.0000,8.0000) -- (12.0000,8.0000) -- (0.0000,8.0000) -- (12.0000,8.0000) -- (0.0000,8.0000) -- (12.0000,8.0000) -- (0.0000,8.0000) -- (12.0000,8.0000) -- (0.0000,8.0000) -- (12.0000,8.0000) -- (0.0000,8.0000) -- (12.0000,8.0000) -- (0.0000,8.0000) -- (12.0000,8.0000) -- (0.0000,8.0000) -- (12.0000,8.0000) -- (0.0000,8.0000) -- (12.0000,8.0000) -- (0.0000,8.0000) -- (12.0000,8.0000) -- (0.0000,8.0000) -- (12.0000,8.0000) -- (0.0000,8.0000) -- (12.0000,8.0000) -- (0.0000,8.0000) -- (12.0000,8.0000) -- (0.0000,8.0000) -- (12.0000,8.0000) -- (0.0000,8.0000) -- (12.0000,8.0000) -- (0.0000,8.0000) -- (12.0000,8.0000) -- (0.0000,8.0000) -- (12.0000,8.0000) -- (0.0000,8.0000) -- (12.0000,8.0000) -- (0.0000,8.0000) -- (12.0000,8.0000) -- (0.0000,8.0000) -- (12.0000,8.0000) -- (0.0000,8.0000) -- (12.0000,8.0000) -- (0.0000,8.0000) -- (12.0000,8.0000) -- (0.0000,8.0000) -- (12.0000,8.0000) -- (0.0000,8.0000) -- (12.0000,8.0000) -- (0.0000,8.0000) -- (12.0000,8.0000) -- (0.0000,8.0000) -- (12.0000,8.0000) -- (0.0000,8.0000) -- (12.0000,8.0000) -- (0.0000,8.0000) -- (12.0000,8.0000) -- (0.0000,8.0000) -- (12.0000,8.0000) -- (0.0000,8.0000) -- (12.0000,8.0000) -- (0.0000,8.0000) -- (12.0000,8.0000) -- (0.0000,8.0000) -- (12.0000,8.0000) -- (0.0000,8.0000) -- (12.0000,8.0000) -- (0.0000,8.0000) -- (12.0000,8.0000) -- (0.0000,8.0000) -- (12.0000,8.0000) -- (0.0000,8.0000) -- (12.0000,8.0000) -- (0.0000,8.0000) -- (12.0000,8.0000) -- (0.0000,8.0000) -- (12.0000,8.0000) -- (0.0000,8.0000) -- (12.0000,8.0000) -- (0.0000,8.0000) -- (12.0000,8.0000) -- (0.0000,8.0000) -- (12.0000,8.0000) -- (0.0000,8.0000) -- (12.0000,8.0000) -- (0.0000,8.0000) -- (12.0000,8.0000) -- (0.0000,8.0000) -- (12.0000,8.0000) -- (0.0000,8.0000) -- (12.0000,8.0000) -- (0.0000,8.0000) -- (12.0000,8.0000) -- (0.0000,8.0000) -- (12.0000,8.0000) -- (0.0000,8.0000) -- (12.0000,8.0000) -- (0.0000,8.0000) -- (12.0000,8.0000) -- (0.0000,8.0000) -- (12.0000,8.0000) -- (0.0000,8.0000) -- (12.0000,8.0000) -- (0.0000,8.0000) -- (12.0000,8.0000) -- (0.0000,8.0000) -- (12.0000,8.0000) -- (0.0000,8.0000) -- (12.0000,8.0000) -- (0.0000,8.0000) -- (12.0000,8.0000) -- (0.0000,8.0000) -- (12.0000,8.0000) -- (0.0000,8.0000) -- (12.0000,8.0000) -- (0.0000,8.0000) -- (12.0000,8.0000) -- (0.0000,8.0000) -- (12.0000,8.0000) -- (0.0000,8.0000) -- (12.0000,8.0000) -- (0.0000,8.0000) -- (12.0000,8.0000) -- (0.0000,8.0000) -- (12.0000,8.0000) -- (0.0000,8.0000) -- (12.0000,8.0000) -- (0.0000,8.0000) -- (12.0000,8.0000) -- (0.0000,8.0000) -- (12.0000,8.0000) -- (0.0000,8.0000) -- (12.0000,8.0000) -- (0.0000,8.0000) -- (12.0000,8.0000) -- (0.0000,8.0000) -- (12.0000,8.0000) -- (0.0000,8.0000) -- (12.0000,8.0000) -- (0.0000,8.0000) -- (12.0000,8.0000) -- (0.0000,8.0000) -- (12.0000,8.0000) -- (0.0000,8.0000) -- (12.0000,8.0000) -- (0.0000,8.0000) -- (12.0000,8.0000) -- (0.0000,8.0000) -- (12.0000,8.0000) -- (0.0000,8.0000) -- (12.0000,8.0000) -- (0.0000,8.0000) -- (12.0000,8.0000) -- (0.0000,8.0000) -- (12.0000,8.0000) -- (0.0000,8.0000) -- (12.0000,8.0000) -- (0.0000,8.0000) -- (12.0000,8.0000) -- (0.0000,8.0000) -- (12.0000,8.0000) -- (0.0000,8.0000) -- (12.0000,8.0000) -- (0.0000,8.0000) -- (12.0000,8.0000) -- (0.0000,8.0000) -- (12.0000,8.0000) -- (0.0000,8.0000) -- (12.0000,8.0000) -- (0.0000,8.0000) -- (12.0000,8.0000) -- (0.0000,8.0000) -- (12.0000,8.0000) -- (0.0000,8.0000) -- (12.0000,8.0000) -- (0.0000,8.0000) -- (12.0000,8.0000) -- (0.0000,8.0000) -- (12.0000,8.0000) -- (0.0000,8.0000) -- (12.0000,8.0000) -- (0.0000,8.0000) -- (12.0000,8.0000) -- (0.0000,8.0000) -- (12.0000,8.0000) -- (0.0000,8.0000) -- (12.0000,8.0000) -- (0.0000,8.0000) -- (12.0000,8.0000) -- (0.0000,8.0000) -- (12.0000,8.0000) -- (0.0000,8.0000) -- (12.0000,8.0000) -- (0.0000,8.0000) -- (12.0000,8.0000) -- (0.0000,8.0000) -- (12.0000,8.0000) -- (0.0000,8.0000) -- (12.0000,8.0000) -- (0.0000,8.0000) -- (12.0000,8.0000) -- (0.0000,8.0000) -- (12.0000,8.0000) -- (0.0000,8.0000) -- (12.0000,8.0000) -- (0.0000,8.0000) -- (12.0000,8.0000) -- (0.0000,8.0000) -- (12.0000,8.0000) -- (0.0000,8.0000) -- (12.0000,8.0000) -- (0.0000,8.0000) -- (12.0000,8.0000) -- (0.0000,8.0000) -- (12.0000,8.0000) -- (0.0000,8.0000) -- (12.0000,8.0000) -- (0.0000,8.0000) -- (12.0000,8.0000) -- (0.0000,8.0000) -- (12.0000,8.0000) -- (0.0000,8.0000) -- (12.0000,8.0000) -- (0.0000,8.0000) -- (12.0000,8.0000) -- (0.0000,8.0000) -- (12.0000,8.0000) -- (0.0000,8.0000) -- (12.0000,8.0000) -- (0.0000,8.0000) -- (12.0000,8.0000) -- (0.0000,8.0000) -- (12.0000,8.0000) -- (0.0000,8.0000) -- (12.0000,8.0000) -- (0.0000,8.0000) -- (12.0000,8.0000) -- (0.0000,8.0000) -- (12.0000,8.0000) -- (0.0000,8.0000) -- (12.0000,8.0000) -- (0.0000,8.0000) -- (12.0000,8.0000) -- (0.0000,8.0000) -- (12.0000,8.0000) -- (0.0000,8.0000) -- (12.0000,8.0000) -- (0.0000,8.0000) -- (12.0000,8.0000) -- (0.0000,8.0000) -- (12.0000,8.0000) -- (0.0000,8.0000) -- (12.0000,8.0000) -- (0.0000,8.0000) -- (12.0000,8.0000) -- (0.0000,8.0000) -- (12.0000,8.0000) -- (0.0000,8.0000) -- (12.0000,8.0000) -- (0.0000,8.0000) -- (12.0000,8.0000) -- (0.0000,8.0000) -- (12.0000,8.0000) -- (0.0000,8.0000) -- (12.0000,8.0000) -- (0.0000,8.0000) -- (12.0000,8.0000) -- (0.0000,8.0000) -- (12.0000,8.0000) -- (0.0000,8.0000) -- (12.0000,8.0000) -- (0.0000,8.0000) -- (12.0000,8.0000) -- (0.0000,8.0000) -- (12.0000,8.0000) -- (0.0000,8.0000) -- (12.0000,8.0000) -- (0.0000,8.0000) -- (12.0000,8.0000) -- (0.0000,8.0000) -- (12.0000,8.0000) -- (0.0000,8.0000) -- (12.0000,8.0000) -- (0.0000,8.0000) -- (12.0000,8.0000) -- (0.0000,8.0000) -- (12.0000,8.0000) -- (0.0000,8.0000) -- (12.0000,8.0000) -- (0.0000,8.0000) -- (12.0000,8.0000) -- (0.0000,8.0000) -- (12.0000,8.0000) -- (0.0000,8.0000) -- (12.0000,8.0000) -- (0.0000,8.0000) -- (12.0000,8.0000) -- (0.0000,8.0000) -- (12.0000,8.0000) -- (0.0000,8.0000) -- (12.0000,8.0000) -- (0.0000,8.0000) -- (12.0000,8.0000) -- (0.0000,8.0000) -- (12.0000,8.0000) -- (0.0000,8.0000) -- (12.0000,8.0000) -- (0.0000,8.0000) -- (12.0000,8.0000) -- (0.0000,8.0000) -- (12.0000,8.0000) -- (0.0000,8.0000) -- (12.0000,8.0000) -- (0.0000,8.0000) -- (12.0000,8.0000) -- (0.0000,8.0000) -- (12.0000,8.0000) -- (0.0000,8.0000) -- (12.0000,8.0000) -- (0.0000,8.0000) -- (12.0000,8.0000) -- (0.0000,8.0000) -- (12.0000,8.0000) -- (0.0000,8.0000) -- (12.0000,8.0000) -- (0.0000,8.0000) -- (12.0000,8.0000) -- (0.0000,8.0000) -- (12.0000,8.0000) -- (0.0000,8.0000) -- (12.0000,8.0000) -- (0.0000,8.0000) -- (12.0000,8.0000) -- (0.0000,8.0000) -- (12.0000,8.0000) -- (0.0000,8.0000) -- (12.0000,8.0000) -- (0.0000,8.0000) -- (12.0000,8.0000) -- (0.0000,8.0000) -- (12.0000,8.0000) -- (0.0000,8.0000) -- (12.0000,8.0000) -- (0.0000,8.0000) -- (12.0000,8.0000) -- (0.0000,8.0000) -- (12.0000,8.0000) -- (0.0000,8.0000) -- (12.0000,8.0000) -- (0.0000,8.0000) -- (12.0000,8.0000) -- (0.0000,8.0000) -- (12.0000,8.0000) -- (0.0000,8.0000) -- (12.0000,8.0000) -- (0.0000,8.0000) -- (12.0000,8.0000) -- (0.0000,8.0000) -- (12.0000,8.0000) -- (0.0000,8.0000) -- (12.0000,8.0000) -- (0.0000,8.0000) -- (12.0000,8.0000) -- (0.0000,8.0000) -- (12.0000,8.0000) -- (0.0000,8.0000) -- (12.0000,8.0000) -- (0.0000,8.0000) -- (12.0000,8.0000) -- (0.0000,8.0000) -- (12.0000,8.0000) -- (0.0000,8.0000) -- (12.0000,8.0000) -- (0.0000,8.0000) -- (12.0000,8.0000) -- (0.0000,8.0000) -- (12.0000,8.0000) -- (0.0000,8.0000) -- (12.0000,8.0000) -- (0.0000,8.0000) -- (12.0000,8.0000) -- (0.0000,8.0000) -- (12.0000,8.0000) -- (0.0000,8.0000) -- (12.0000,8.0000) -- (0.0000,8.0000) -- (12.0000,8.0000) -- (0.0000,8.0000) -- (12.0000,8.0000) -- (0.0000,8.0000) -- (12.0000,8.0000) -- (0.0000,8.0000) -- (12.0000,8.0000) -- (0.0000,8.0000) -- (12.0000,8.0000) -- (0.0000,8.0000) -- (12.0000,8.0000) -- (0.0000,8.0000) -- (12.0000,8.0000) -- (0.0000,8.0000) -- (12.0000,8.0000) -- (0.0000,8.0000) -- (12.0000,8.0000) -- (0.0000,8.0000) -- (12.0000,8.0000) -- (0.0000,8.0000) -- (12.0000,8.0000) -- (0.0000,8.0000) -- (12.0000,8.0000) -- (0.0000,8.0000) -- (12.0000,8.0000) -- (0.0000,8.0000) -- (12.0000,8.0000) -- (0.0000,8.0000) -- (12.0000,8.0000) -- (0.0000,8.0000) -- (12.0000,8.0000) -- (0.0000,8.0000) -- (12.0000,8.0000) -- (0.0000,8.0000) -- (12.0000,8.0000) -- (0.0000,8.0000) -- (12.0000,8.0000) -- (0.0000,8.0000) -- (12.0000,8.0000) -- (0.0000,8.0000) -- (12.0000,8.0000) -- (0.0000,8.0000) -- (12.0000,8.0000) -- (0.0000,8.0000) -- (12.0000,8.0000) -- (0.0000,8.0000) -- (12.0000,8.0000) -- (0.0000,8.0000) -- (12.0000,8.0000) -- (0.0000,8.0000) -- (12.0000,8.0000) -- (0.0000,8.0000) -- (12.0000,8.0000) -- (0.0000,8.0000) -- (12.0000,8.0000) -- (0.0000,8.0000) -- (12.0000,8.0000) -- (0.0000,8.0000) -- (12.0000,8.0000) -- (0.0000,8.0000) -- (12.0000,8.0000) -- (0.0000,8.0000) -- (12.0000,8.0000) -- (0.0000,8.0000) -- (12.0000,8.0000) -- (0.0000,8.0000) -- (12.0000,8.0000) -- (0.0000,8.0000) -- (12.0000,8.0000) -- (0.0000,8.0000) -- (12.0000,8.0000) -- (0.0000,8.0000) -- (12.0000,8.0000) -- (0.0000,8.0000) -- (12.0000,8.0000) -- (0.0000,8.0000) -- (12.0000,8.0000) -- (0.0000,8.0000) -- (12.0000,8.0000) -- (0.0000,8.0000) -- (12.0000,8.0000) -- (0.0000,8.0000) -- (12.0000,8.0000) -- (0.0000,8.0000) -- (12.0000,8.0000) -- (0.0000,8.0000) -- (12.0000,8.0000) -- (0.0000,8.0000) -- (12.0000,8.0000) -- (0.0000,8.0000) -- (12.0000,8.0000) -- (0.0000,8.0000) -- (12.0000,8.0000) -- (0.0000,8.0000) -- (12.0000,8.0000) -- (0.0000,8.0000) -- (12.0000,8.0000) -- (0.0000,8.0000) -- (12.0000,8.0000) -- (0.0000,8.0000) -- (12.0000,8.0000) -- (0.0000,8.0000) -- (12.0000,8.0000) -- (0.0000,8.0000) -- (12.0000,8.0000) -- (0.0000,8.0000) -- (12.0000,8.0000) -- (0.0000,8.0000) -- (12.0000,8.0000) -- (0.0000,8.0000) -- (12.0000,8.0000) -- (0.0000,8.0000) -- (12.0000,8.0000) -- (0.0000,8.0000) -- (12.0000,8.0000) -- (0.0000,8.0000) -- (12.0000,8.0000) -- (0.0000,8.0000) -- (12.0000,8.0000) -- (0.0000,8.0000) -- (12.0000,8.0000) -- (0.0000,8.0000) -- (12.0000,8.0000) -- (0.0000,8.0000) -- (12.0000,8.0000) -- (0.0000,8.0000) -- (12.0000,8.0000) -- (0.0000,8.0000) -- (12.0000,8.0000) -- (0.0000,8.0000) -- (12.0000,8.0000) -- (0.0000,8.0000) -- (12.0000,8.0000) -- (0.0000,8.0000) -- (12.0000,8.0000) -- (0.0000,8.0000) -- (12.0000,8.0000) -- (0.0000,8.0000) -- (12.0000,8.0000) -- (0.0000,8.0000) -- (12.0000,8.0000) -- (0.0000,8.0000) -- (12.0000,8.0000) -- (0.0000,8.0000) -- (12.0000,8.0000) -- (0.0000,8.0000) -- (12.0000,8.0000) -- (0.0000,8.0000) -- (12.0000,8.0000) -- (0.0000,8.0000) -- (12.0000,8.0000) -- (0.0000,8.0000) -- (12.0000,8.0000) -- (0.0000,8.0000) -- (12.0000,8.0000) -- (0.0000,8.0000) -- (12.0000,8.0000) -- (0.0000,8.0000) -- (12.0000,8.0000) -- (0.0000,8.0000) -- (12.0000,8.0000) -- (0.0000,8.0000) -- (12.0000,8.0000) -- (0.0000,8.0000) -- (12.0000,8.0000) -- (0.0000,8.0000) -- (12.0000,8.0000) -- (0.0000,8.0000) -- (12.0000,8.0000) -- (0.0000,8.0000) -- (12.0000,8.0000) -- (0.0000,8.0000) -- (12.0000,8.0000) -- (0.0000,8.0000) -- (12.0000,8.0000) -- (0.0000,8.0000) -- (12.0000,8.0000) -- (0.0000,8.0000) -- (12.0000,8.0000) -- (0.0000,8.0000) -- (12.0000,8.0000) -- (0.0000,8.0000) -- (12.0000,8.0000) -- (0.0000,8.0000) -- (12.0000,8.0000) -- (0.0000,8.0000) -- (12.0000,8.0000) -- (0.0000,8.0000) -- (12.0000,8.0000) -- (0.0000,8.0000) -- (12.0000,8.0000) -- (0.0000,8.0000) -- (12.0000,8.0000) -- (0.0000,8.0000) -- (12.0000,8.0000) -- (0.0000,8.0000) -- (12.0000,8.0000) -- (0.0000,8.0000) -- (12.0000,8.0000) -- (0.0000,8.0000) -- (12.0000,8.0000) -- (0.0000,8.0000) -- (12.0000,8.0000) -- (0.0000,8.0000) -- (12.0000,8.0000) -- (0.0000,8.0000) -- (12.0000,8.0000) -- (0.0000,8.0000) -- (12.0000,8.0000) -- (0.0000,8.0000) -- (12.0000,8.0000) -- (0.0000,8.0000) -- (12.0000,8.0000) -- (0.0000,8.0000) -- (12.0000,8.0000) -- (0.0000,8.0000) -- (12.0000,8.0000) -- (0.0000,8.0000) -- (12.0000,8.0000) -- (0.0000,8.0000) -- (12.0000,8.0000) -- (0.0000,8.0000) -- (12.0000,8.0000) -- (0.0000,8.0000) -- (12.0000,8.0000) -- (0.0000,8.0000) -- (12.0000,8.0000) -- (0.0000,8.0000) -- (12.0000,8.0000) -- (0.0000,8.0000) -- (12.0000,8.0000) -- (0.0000,8.0000) -- (12.0000,8.0000) -- (0.0000,8.0000) -- (12.0000,8.0000) -- (0.0000,8.0000) -- (12.0000,8.0000) -- (0.0000,8.0000) -- (12.0000,8.0000) -- (0.0000,8.0000) -- (12.0000,8.0000) -- (0.0000,8.0000) -- (12.0000,8.0000) -- (0.0000,8.0000) -- (12.0000,8.0000) -- (0.0000,8.0000) -- (12.0000,8.0000) -- (0.0000,8.0000) -- (12.0000,8.0000) -- (0.0000,8.0000) -- (12.0000,8.0000) -- (0.0000,8.0000) -- (12.0000,8.0000) -- (0.0000,8.0000) -- (12.0000,8.0000) -- (0.0000,8.0000) -- (12.0000,8.0000) -- (0.0000,8.0000) -- (12.0000,8.0000) -- (0.0000,8.0000) -- (12.0000,8.0000) -- (0.0000,8.0000) -- (12.0000,8.0000) -- (0.0000,8.0000) -- (12.0000,8.0000) -- (0.0000,8.0000) -- (12.0000,8.0000) -- (0.0000,8.0000) -- (12.0000,8.0000) -- (0.0000,8.0000) -- (12.0000,8.0000) -- (0.0000,8.0000) -- (12.0000,8.0000) -- (0.0000,8.0000) -- (12.0000,8.0000) -- (0.0000,8.0000) -- (12.0000,8.0000) -- (0.0000,8.0000) -- (12.0000,8.0000) -- (0.0000,8.0000) -- (12.0000,8.0000) -- (0.0000,8.0000) -- (12.0000,8.0000) -- (0.0000,8.0000) -- (12.0000,8.0000) -- (0.0000,8.0000) -- (12.0000,8.0000) -- (0.0000,8.0000) -- (12.0000,8.0000) -- (0.0000,8.0000) -- (12.0000,8.0000) -- (0.0000,8.0000) -- (12.0000,8.0000) -- (0.0000,8.0000) -- (12.0000,8.0000) -- (0.0000,8.0000) -- (12.0000,8.0000) -- (0.0000,8.0000) -- (12.0000,8.0000) -- (0.0000,8.0000) -- (12.0000,8.0000) -- (0.0000,8.0000) -- (12.0000,8.0000) -- (0.0000,8.0000) -- (12.0000,8.0000) -- (0.0000,8.0000) -- (12.0000,8.0000) -- (0.0000,8.0000) -- (12.0000,8.0000) -- (0.0000,8.0000) -- (12.0000,8.0000) -- (0.0000,8.0000) -- (12.0000,8.0000) -- (0.0000,8.0000) -- (12.0000,8.0000) -- (0.0000,8.0000) -- (12.0000,8.0000) -- (0.0000,8.0000) -- (12.0000,8.0000) -- (0.0000,8.0000) -- (12.0000,8.0000) -- (0.0000,8.0000) -- (12.0000,8.0000) -- (0.0000,8.0000) -- (12.0000,8.0000) -- (0.0000,8.0000) -- (12.0000,8.0000) -- (0.0000,8.0000) -- (12.0000,8.0000) -- (0.0000,8.0000) -- (12.0000,8.0000) -- (0.0000,8.0000) -- (12.0000,8.0000) -- (0.0000,8.0000) -- (12.0000,8.0000) -- (0.0000,8.0000) -- (12.0000,8.0000) -- (0.0000,8.0000) -- (12.0000,8.0000) -- (0.0000,8.0000) -- (12.0000,8.0000) -- (0.0000,8.0000) -- (12.0000,8.0000) -- (0.0000,8.0000) -- (12.0000,8.0000) -- (0.0000,8.0000) -- (12.0000,8.0000) -- (0.0000,8.0000) -- (12.0000,8.0000) -- (0.0000,8.0000) -- (12.0000,8.0000) -- (0.0000,8.0000) -- (12.0000,8.0000) -- (0.0000,8.0000) -- (12.0000,8.0000) -- (0.0000,8.0000) -- (12.0000,8.0000) -- (0.0000,8.0000) -- (12.0000,8.0000) -- (0.0000,8.0000) -- (12.0000,8.0000) -- (0.0000,8.0000) -- (12.0000,8.0000) -- (0.0000,8.0000) -- (12.0000,8.0000) -- (0.0000,8.0000) -- (12.0000,8.0000) -- (0.0000,8.0000) -- (12.0000,8.0000) -- (0.0000,8.0000) -- (12.0000,8.0000) -- (0.0000,8.0000) -- (12.0000,8.0000) -- (0.0000,8.0000) -- (12.0000,8.0000) -- (0.0000,8.0000) -- (12.0000,8.0000) -- (0.0000,8.0000) -- (12.0000,8.0000) -- (0.0000,8.0000) -- (12.0000,8.0000) -- (0.0000,8.0000) -- (12.0000,8.0000) -- (0.0000,8.0000) -- (12.0000,8.0000) -- (0.0000,8.0000) -- (12.0000,8.0000) -- (0.0000,8.0000) -- (12.0000,8.0000) -- (0.0000,8.0000) -- (12.0000,8.0000) -- (0.0000,8.0000) -- (12.0000,8.0000) -- (0.0000,8.0000) -- (12.0000,8.0000) -- (0.0000,8.0000) -- (12.0000,8.0000) -- (0.0000,8.0000) -- (12.0000,8.0000) -- (0.0000,8.0000) -- (12.0000,8.0000) -- (0.0000,8.0000) -- (12.0000,8.0000) -- (0.0000,8.0000) -- (12.0000,8.0000) -- (0.0000,8.0000) -- (12.0000,8.0000) -- (0.0000,8.0000) -- (12.0000,8.0000) -- (0.0000,8.0000) -- (12.0000,8.0000) -- (0.0000,8.0000) -- (12.0000,8.0000) -- (0.0000,8.0000) -- (12.0000,8.0000) -- (0.0000,8.0000) -- (12.0000,8.0000) -- (0.0000,8.0000) -- (12.0000,8.0000) -- (0.0000,8.0000) -- (12.0000,8.0000) -- (0.0000,8.0000) -- (12.0000,8.0000) -- (0.0000,8.0000) -- (12.0000,8.0000) -- (0.0000,8.0000) -- (12.0000,8.0000) -- (0.0000,8.0000) -- (12.0000,8.0000) -- (0.0000,8.0000) -- (12.0000,8.0000) -- (0.0000,8.0000) -- (12.0000,8.0000) -- (0.0000,8.0000) -- (12.0000,8.0000) -- (0.0000,8.0000) -- (12.0000,8.0000) -- (0.0000,8.0000) -- (12.0000,8.0000) -- (0.0000,8.0000) -- (12.0000,8.0000) -- (0.0000,8.0000) -- (12.0000,8.0000) -- (0.0000,8.0000) -- (12.0000,8.0000) -- (0.0000,8.0000) -- (12.0000,8.0000) -- (0.0000,8.0000) -- (12.0000,8.0000) -- (0.0000,8.0000) -- (12.0000,8.0000) -- (0.0000,8.0000) -- (12.0000,8.0000) -- (0.0000,8.0000) -- (12.0000,8.0000) -- (0.0000,8.0000) -- (12.0000,8.0000) -- (0.0000,8.0000) -- (12.0000,8.0000) -- (0.0000,8.0000) -- (12.0000,8.0000) -- (0.0000,8.0000) -- (12.0000,8.0000) -- (0.0000,8.0000) -- (12.0000,8.0000) -- (0.0000,8.0000) -- (12.0000,8.0000) -- (0.0000,8.0000) -- (12.0000,8.0000) -- (0.0000,8.0000) -- (12.0000,8.0000) -- (0.0000,8.0000) -- (12.0000,8.0000) -- (0.0000,8.0000) -- (12.0000,8.0000) -- (0.0000,8.0000) -- (12.0000,8.0000) -- (0.0000,8.0000) -- (12.0000,8.0000) -- (0.0000,8.0000) -- (12.0000,8.0000) -- (0.0000,8.0000) -- (12.0000,8.0000) -- (0.0000,8.0000) -- (12.0000,8.0000) -- (0.0000,8.0000) -- (12.0000,8.0000) -- (0.0000,8.0000) -- (12.0000,8.0000) -- (0.0000,8.0000) -- (12.0000,8.0000) -- (0.0000,8.0000) -- (12.0000,8.0000) -- (0.0000,8.0000) -- (12.0000,8.0000) -- (0.0000,8.0000) -- (12.0000,8.0000) -- (0.0000,8.0000) -- (12.0000,8.0000) -- (0.0000,8.0000) -- (12.0000,8.0000) -- (0.0000,8.0000) -- (12.0000,8.0000) -- (0.0000,8.0000) -- (12.0000,8.0000) -- (0.0000,8.0000) -- (12.0000,8.0000) -- (0.0000,8.0000) -- (12.0000,8.0000) -- (0.0000,8.0000) -- (12.0000,8.0000) -- (0.0000,8.0000) -- (12.0000,8.0000) -- (0.0000,8.0000) -- (12.0000,8.0000) -- (0.0000,8.0000) -- (12.0000,8.0000) -- (0.0000,8.0000) -- (12.0000,8.0000) -- (0.0000,8.0000) -- (12.0000,8.0000) -- (0.0000,8.0000) -- (12.0000,8.0000) -- (0.0000,8.0000) -- (12.0000,8.0000) -- (0.0000,8.0000) -- (12.0000,8.0000) -- (0.0000,8.0000) -- (12.0000,8.0000) -- (0.0000,8.0000) -- (12.0000,8.0000) -- (0.0000,8.0000) -- (12.0000,8.0000) -- (0.0000,8.0000) -- (12.0000,8.0000) -- (0.0000,8.0000) -- (12.0000,8.0000) -- (0.0000,8.0000) -- (12.0000,8.0000) -- (0.0000,8.0000) -- (12.0000,8.0000) -- (0.0000,8.0000) -- (12.0000,8.0000) -- (0.0000,8.0000) -- (12.0000,8.0000) -- (0.0000,8.0000) -- (12.0000,8.0000) -- (0.0000,8.0000) -- (12.0000,8.0000) -- (0.0000,8.0000) -- (12.0000,8.0000) -- (0.0000,8.0000) -- (12.0000,8.0000) -- (0.0000,8.0000) -- (12.0000,8.0000) -- (0.0000,8.0000) -- (12.0000,8.0000) -- (0.0000,8.0000) -- (12.0000,8.0000) -- (0.0000,8.0000) -- (12.0000,8.0000) -- (0.0000,8.0000) -- (12.0000,8.0000) -- (0.0000,8.0000) -- (12.0000,8.0000) -- (0.0000,8.0000) -- (12.0000,8.0000) -- (0.0000,8.0000) -- (12.0000,8.0000) -- (0.0000,8.0000) -- (12.0000,8.0000) -- (0.0000,8.0000) -- (12.0000,8.0000) -- (0.0000,8.0000) -- (12.0000,8.0000) -- (0.0000,8.0000) -- (12.0000,8.0000) -- (0.0000,8.0000) -- (12.0000,8.0000) -- (0.0000,8.0000) -- (12.0000,8.0000) -- (0.0000,8.0000) -- (12.0000,8.0000) -- (0.0000,8.0000) -- (12.0000,8.0000) -- (0.0000,8.0000) -- (12.0000,8.0000) -- (0.0000,8.0000) -- (12.0000,8.0000) -- (0.0000,8.0000) -- (12.0000,8.0000) -- (0.0000,8.0000) -- (12.0000,8.0000) -- (0.0000,8.0000) -- (12.0000,8.0000) -- (0.0000,8.0000) -- (12.0000,8.0000) -- (0.0000,8.0000) -- (12.0000,8.0000) -- (0.0000,8.0000) -- (12.0000,8.0000) -- (0.0000,8.0000) -- (12.0000,8.0000) -- (0.0000,8.0000) -- (12.0000,8.0000) -- (0.0000,8.0000) -- (12.0000,8.0000) -- (0.0000,8.0000) -- (12.0000,8.0000) -- (0.0000,8.0000) -- (12.0000,8.0000) -- (0.0000,8.0000) -- (12.0000,8.0000) -- (0.0000,8.0000) -- (12.0000,8.0000) -- (0.0000,8.0000) -- (12.0000,8.0000) -- (0.0000,8.0000) -- (12.0000,8.0000) -- (0.0000,8.0000) -- (12.0000,8.0000) -- (0.0000,8.0000) -- (12.0000,8.0000) -- (0.0000,8.0000) -- (12.0000,8.0000) -- (0.0000,8.0000) -- (12.0000,8.0000) -- (0.0000,8.0000) -- (12.0000,8.0000) -- (0.0000,8.0000) -- (12.0000,8.0000) -- (0.0000,8.0000) -- (12.0000,8.0000) -- (0.0000,8.0000) -- (12.0000,8.0000) -- (0.0000,8.0000) -- (12.0000,8.0000) -- (0.0000,8.0000) -- (12.0000,8.0000) -- (0.0000,8.0000) -- (12.0000,8.0000) -- (0.0000,8.0000) -- (12.0000,8.0000) -- (0.0000,8.0000) -- (12.0000,8.0000) -- (0.0000,8.0000) -- (12.0000,8.0000) -- (0.0000,8.0000) -- (12.0000,8.0000) -- (0.0000,8.0000) -- (12.0000,8.0000) -- (0.0000,8.0000) -- (12.0000,8.0000) -- (0.0000,8.0000) -- (12.0000,8.0000) -- (0.0000,8.0000) -- (12.0000,8.0000) -- (0.0000,8.0000) -- (12.0000,8.0000) -- (0.0000,8.0000) -- (12.0000,8.0000) -- (0.0000,8.0000) -- (12.0000,8.0000) -- (0.0000,8.0000) -- (12.0000,8.0000) -- (0.0000,8.0000) -- (12.0000,8.0000) -- (0.0000,8.0000) -- (12.0000,8.0000) -- (0.0000,8.0000) -- (12.0000,8.0000) -- (0.0000,8.0000) -- (12.0000,8.0000) -- (0.0000,8.0000) -- (12.0000,8.0000) -- (0.0000,8.0000) -- (12.0000,8.0000) -- (0.0000,8.0000) -- (12.0000,8.0000) -- (0.0000,8.0000) -- (12.0000,8.0000) -- (0.0000,8.0000) -- (12.0000,8.0000) -- (0.0000,8.0000) -- (12.0000,8.0000) -- (0.0000,8.0000) -- (12.0000,8.0000) -- (0.0000,8.0000) -- (12.0000,8.0000) -- (0.0000,8.0000) -- (12.0000,8.0000) -- (0.0000,8.0000) -- (12.0000,8.0000) -- (0.0000,8.0000) -- (12.0000,8.0000) -- (0.0000,8.0000) -- (12.0000,8.0000) -- (0.0000,8.0000) -- (12.0000,8.0000) -- (0.0000,8.0000) -- (12.0000,8.0000) -- (0.0000,8.0000) -- (12.0000,8.0000) -- (0.0000,8.0000) -- (12.0000,8.0000) -- (0.0000,8.0000) -- (12.0000,8.0000) -- (0.0000,8.0000) -- (12.0000,8.0000) -- (0.0000,8.0000) -- (12.0000,8.0000) -- (0.0000,8.0000) -- (12.0000,8.0000) -- (0.0000,8.0000) -- (12.0000,8.0000) -- (0.0000,8.0000) -- (12.0000,8.0000) -- (0.0000,8.0000) -- (12.0000,8.0000) -- (0.0000,8.0000) -- (12.0000,8.0000) -- (0.0000,8.0000) -- (12.0000,8.0000) -- (0.0000,8.0000) -- (12.0000,8.0000) -- (0.0000,8.0000) -- (12.0000,8.0000) -- (0.0000,8.0000) -- (12.0000,8.0000) -- (0.0000,8.0000) -- (12.0000,8.0000) -- (0.0000,8.0000) -- (12.0000,8.0000) -- (0.0000,8.0000) -- (12.0000,8.0000) -- (0.0000,8.0000) -- (12.0000,8.0000) -- (0.0000,8.0000) -- (12.0000,8.0000) -- (0.0000,8.0000) -- (12.0000,8.0000) -- (0.0000,8.0000) -- (12.0000,8.0000) -- (0.0000,8.0000) -- (12.0000,8.0000) -- (0.0000,8.0000) -- (12.0000,8.0000) -- (0.0000,8.0000) -- (12.0000,8.0000) -- (0.0000,8.0000) -- (12.0000,8.0000) -- (0.0000,8.0000) -- (12.0000,8.0000) -- (0.0000,8.0000) -- (12.0000,8.0000) -- (0.0000,8.0000) -- (12.0000,8.0000) -- (0.0000,8.0000) -- (12.0000,8.0000) -- (0.0000,8.0000) -- (12.0000,8.0000) -- (0.0000,8.0000) -- (12.0000,8.0000) -- (0.0000,8.0000) -- (12.0000,8.0000) -- (0.0000,8.0000) -- (12.0000,8.0000) -- (0.0000,8.0000) -- (12.0000,8.0000) -- (0.0000,8.0000) -- (12.0000,8.0000) -- (0.0000,8.0000) -- (12.0000,8.0000) -- (0.0000,8.0000) -- (12.0000,8.0000) -- (0.0000,8.0000) -- (12.0000,8.0000) -- (0.0000,8.0000) -- (12.0000,8.0000) -- (0.0000,8.0000) -- (12.0000,8.0000) -- (0.0000,8.0000) -- (12.0000,8.0000) -- (0.0000,8.0000) -- (12.0000,8.0000) -- (0.0000,8.0000) -- (12.0000,8.0000) -- (0.0000,8.0000) -- (12.0000,8.0000) -- (0.0000,8.0000) -- (12.0000,8.0000) -- (0.0000,8.0000) -- (12.0000,8.0000) -- (0.0000,8.0000) -- (12.0000,8.0000) -- (0.0000,8.0000) -- (12.0000,8.0000) -- (0.0000,8.0000) -- (12.0000,8.0000) -- (0.0000,8.0000) -- (12.0000,8.0000) -- (0.0000,8.0000) -- (12.0000,8.0000) -- (0.0000,8.0000) -- (12.0000,8.0000) -- (0.0000,8.0000) -- (12.0000,8.0000) -- (0.0000,8.0000) -- (12.0000,8.0000) -- (0.0000,8.0000) -- (12.0000,8.0000) -- (0.0000,8.0000) -- (12.0000,8.0000) -- (0.0000,8.0000) -- (12.0000,8.0000) -- (0.0000,8.0000) -- (12.0000,8.0000) -- (0.0000,8.0000) -- (12.0000,8.0000) -- (0.0000,8.0000) -- (12.0000,8.0000) -- (0.0000,8.0000) -- (12.0000,8.0000) -- (0.0000,8.0000) -- (12.0000,8.0000) -- (0.0000,8.0000) -- (12.0000,8.0000) -- (0.0000,8.0000) -- (12.0000,8.0000) -- (0.0000,8.0000) -- (12.0000,8.0000) -- (0.0000,8.0000) -- (12.0000,8.0000) -- (0.0000,8.0000) -- (12.0000,8.0000) -- (0.0000,8.0000) -- (12.0000,8.0000) -- (0.0000,8.0000) -- (12.0000,8.0000) -- (0.0000,8.0000) -- (12.0000,8.0000) -- (0.0000,8.0000) -- (12.0000,8.0000) -- (0.0000,8.0000) -- (12.0000,8.0000) -- (0.0000,8.0000) -- (12.0000,8.0000) -- (0.0000,8.0000) -- (12.0000,8.0000) -- (0.0000,8.0000) -- (12.0000,8.0000) -- (0.0000,8.0000) -- (12.0000,8.0000) -- (0.0000,8.0000) -- (12.0000,8.0000) -- (0.0000,8.0000) -- (12.0000,8.0000) -- (0.0000,8.0000) -- (12.0000,8.0000) -- (0.0000,8.0000) -- (12.0000,8.0000) -- (0.0000,8.0000) -- (12.0000,8.0000) -- (0.0000,8.0000) -- (12.0000,8.0000) -- (0.0000,8.0000) -- (12.0000,8.0000) -- (0.0000,8.0000) -- (12.0000,8.0000) -- (0.0000,8.0000) -- (12.0000,8.0000) -- (0.0000,8.0000) -- (12.0000,8.0000) -- (0.0000,8.0000) -- (12.0000,8.0000) -- (0.0000,8.0000) -- (12.0000,8.0000) -- (0.0000,8.0000) -- (12.0000,8.0000) -- (0.0000,8.0000) -- (12.0000,8.0000) -- (0.0000,8.0000) -- (12.0000,8.0000) -- (0.0000,8.0000) -- (12.0000,8.0000) -- (0.0000,8.0000) -- (12.0000,8.0000) -- (0.0000,8.0000) -- (12.0000,8.0000) -- (0.0000,8.0000) -- (12.0000,8.0000) -- (0.0000,8.0000) -- (12.0000,8.0000) -- (0.0000,8.0000) -- (12.0000,8.0000) -- (0.0000,8.0000) -- (12.0000,8.0000) -- (0.0000,8.0000) -- (12.0000,8.0000) -- (0.0000,8.0000) -- (12.0000,8.0000) -- (0.0000,8.0000) -- (12.0000,8.0000) -- (0.0000,8.0000) -- (12.0000,8.0000) -- (0.0000,8.0000) -- (12.0000,8.0000) -- (0.0000,8.0000) -- (12.0000,8.0000) -- (0.0000,8.0000) -- (12.0000,8.0000) -- (0.0000,8.0000) -- (12.0000,8.0000) -- (0.0000,8.0000) -- (12.0000,8.0000) -- (0.0000,8.0000) -- (12.0000,8.0000) -- (0.0000,8.0000) -- (12.0000,8.0000) -- (0.0000,8.0000) -- (12.0000,8.0000) -- (0.0000,8.0000) -- (12.0000,8.0000) -- (0.0000,8.0000) -- (12.0000,8.0000) -- (0.0000,8.0000) -- (12.0000,8.0000) -- (0.0000,8.0000) -- (12.0000,8.0000) -- (0.0000,8.0000) -- (12.0000,8.0000) -- (0.0000,8.0000) -- (12.0000,8.0000) -- (0.0000,8.0000) -- (12.0000,8.0000) -- (0.0000,8.0000) -- (12.0000,8.0000) -- (0.0000,8.0000) -- (12.0000,8.0000) -- (0.0000,8.0000) -- (12.0000,8.0000) -- (0.0000,8.0000) -- (12.0000,8.0000) -- (0.0000,8.0000) -- (12.0000,8.0000) -- (0.0000,8.0000) -- (12.0000,8.0000) -- (0.0000,8.0000) -- (12.0000,8.0000) -- (0.0000,8.0000) -- (12.0000,8.0000) -- (0.0000,8.0000) -- (12.0000,8.0000) -- (0.0000,8.0000) -- (12.0000,8.0000) -- (0.0000,8.0000) -- (12.0000,8.0000) -- (0.0000,8.0000) -- (12.0000,8.0000) -- (0.0000,8.0000) -- (12.0000,8.0000) -- (0.0000,8.0000) -- (12.0000,8.0000) -- (0.0000,8.0000) -- (12.0000,8.0000) -- (0.0000,8.0000) -- (12.0000,8.0000) -- (0.0000,8.0000) -- (12.0000,8.0000) -- (0.0000,8.0000) -- (12.0000,8.0000) -- (0.0000,8.0000) -- (12.0000,8.0000) -- (0.0000,8.0000) -- (12.0000,8.0000) -- (0.0000,8.0000) -- (12.0000,8.0000) -- (0.0000,8.0000) -- (12.0000,8.0000) -- (0.0000,8.0000) -- (12.0000,8.0000) -- (0.0000,8.0000) -- (12.0000,8.0000) -- (0.0000,8.0000) -- (12.0000,8.0000) -- (0.0000,8.0000) -- (12.0000,8.0000) -- (0.0000,8.0000) -- (12.0000,8.0000) -- (0.0000,8.0000) -- (12.0000,8.0000) -- (0.0000,8.0000) -- (12.0000,8.0000) -- (0.0000,8.0000) -- (12.0000,8.0000) -- (0.0000,8.0000) -- (12.0000,8.0000) -- (0.0000,8.0000) -- (12.0000,8.0000) -- (0.0000,8.0000) -- (12.0000,8.0000) -- (0.0000,8.0000) -- (12.0000,8.0000) -- (0.0000,8.0000) -- (12.0000,8.0000) -- (0.0000,8.0000) -- (12.0000,8.0000) -- (0.0000,8.0000) -- (12.0000,8.0000) -- (0.0000,8.0000) -- (12.0000,8.0000) -- (0.0000,8.0000) -- (12.0000,8.0000) -- (0.0000,8.0000) -- (12.0000,8.0000) -- (0.0000,8.0000) -- (12.0000,8.0000) -- (0.0000,8.0000) -- (12.0000,8.0000) -- (0.0000,8.0000) -- (12.0000,8.0000) -- (0.0000,8.0000) -- (12.0000,8.0000) -- (0.0000,8.0000) -- (12.0000,8.0000) -- (0.0000,8.0000) -- (12.0000,8.0000) -- (0.0000,8.0000) -- (12.0000,8.0000) -- (0.0000,8.0000) -- (12.0000,8.0000) -- (0.0000,8.0000) -- (12.0000,8.0000) -- (0.0000,8.0000) -- (12.0000,8.0000) -- (0.0000,8.0000) -- (12.0000,8.0000) -- (0.0000,8.0000) -- (12.0000,8.0000) -- (0.0000,8.0000) -- (12.0000,8.0000) -- (0.0000,8.0000) -- (12.0000,8.0000) -- (0.0000,8.0000) -- (12.0000,8.0000) -- (0.0000,8.0000) -- (12.0000,8.0000) -- (0.0000,8.0000) -- (12.0000,8.0000) -- (0.0000,8.0000) -- (12.0000,8.0000) -- (0.0000,8.0000) -- (12.0000,8.0000) -- (0.0000,8.0000) -- (12.0000,8.0000) -- (0.0000,8.0000) -- (12.0000,8.0000) -- (0.0000,8.0000) -- (12.0000,8.0000) -- (0.0000,8.0000) -- (12.0000,8.0000) -- (0.0000,8.0000) -- (12.0000,8.0000) -- (0.0000,8.0000) -- (12.0000,8.0000) -- (0.0000,8.0000) -- (12.0000,8.0000) -- (0.0000,8.0000) -- (12.0000,8.0000) -- (0.0000,8.0000) -- (12.0000,8.0000) -- (0.0000,8.0000) -- (12.0000,8.0000) -- (0.0000,8.0000) -- (12.0000,8.0000) -- (0.0000,8.0000) -- (12.0000,8.0000) -- (0.0000,8.0000) -- (12.0000,8.0000) -- (0.0000,8.0000) -- (12.0000,8.0000) -- (0.0000,8.0000) -- (12.0000,8.0000) -- (0.0000,8.0000) -- (12.0000,8.0000) -- (0.0000,8.0000) -- (12.0000,8.0000) -- (0.0000,8.0000) -- (12.0000,8.0000) -- (0.0000,8.0000) -- (12.0000,8.0000) -- (0.0000,8.0000) -- (12.0000,8.0000) -- (0.0000,8.0000) -- (12.0000,8.0000) -- (0.0000,8.0000) -- (12.0000,8.0000) -- (0.0000,8.0000) -- (12.0000,8.0000) -- (0.0000,8.0000) -- (12.0000,8.0000) -- (0.0000,8.0000) -- (12.0000,8.0000) -- (0.0000,8.0000) -- (12.0000,8.0000) -- (0.0000,8.0000) -- (12.0000,8.0000) -- (0.0000,8.0000) -- (12.0000,8.0000) -- (0.0000,8.0000) -- (12.0000,8.0000) -- (0.0000,8.0000) -- (12.0000,8.0000) -- (0.0000,8.0000) -- (12.0000,8.0000) -- (0.0000,8.0000) -- (12.0000,8.0000) -- (0.0000,8.0000) -- (12.0000,8.0000) -- (0.0000,8.0000) -- (12.0000,8.0000) -- (0.0000,8.0000) -- (12.0000,8.0000) -- (0.0000,8.0000) -- (12.0000,8.0000) -- (0.0000,8.0000) -- (12.0000,8.0000) -- (0.0000,8.0000) -- (12.0000,8.0000) -- (0.0000,8.0000) -- (12.0000,8.0000) -- (0.0000,8.0000) -- (12.0000,8.0000) -- (0.0000,8.0000) -- (12.0000,8.0000) -- (0.0000,8.0000) -- (12.0000,8.0000) -- (0.0000,8.0000) -- (12.0000,8.0000) -- (0.0000,8.0000) -- (12.0000,8.0000) -- (0.0000,8.0000) -- (12.0000,8.0000) -- (0.0000,8.0000) -- (12.0000,8.0000) -- (0.0000,8.0000) -- (12.0000,8.0000) -- (0.0000,8.0000) -- (12.0000,8.0000) -- (0.0000,8.0000) -- (12.0000,8.0000) -- (0.0000,8.0000) -- (12.0000,8.0000) -- (0.0000,8.0000) -- (12.0000,8.0000) -- (0.0000,8.0000) -- (12.0000,8.0000) -- (0.0000,8.0000) -- (12.0000,8.0000) -- (0.0000,8.0000) -- (12.0000,8.0000) -- (0.0000,8.0000) -- (12.0000,8.0000) -- (0.0000,8.0000) -- (12.0000,8.0000) -- (0.0000,8.0000) -- (12.0000,8.0000) -- (0.0000,8.0000) -- (12.0000,8.0000) -- (0.0000,8.0000) -- (12.0000,8.0000) -- (0.0000,8.0000) -- (12.0000,8.0000) -- (0.0000,8.0000) -- (12.0000,8.0000) -- (0.0000,8.0000) -- (12.0000,8.0000) -- (0.0000,8.0000) -- (12.0000,8.0000) -- (0.0000,8.0000) -- (12.0000,8.0000) -- (0.0000,8.0000) -- (12.0000,8.0000) -- (0.0000,8.0000) -- (12.0000,8.0000) -- (0.0000,8.0000) -- (12.0000,8.0000) -- (0.0000,8.0000) -- (12.0000,8.0000) -- (0.0000,8.0000) -- (12.0000,8.0000) -- (0.0000,8.0000) -- (12.0000,8.0000) -- (0.0000,8.0000) -- (12.0000,8.0000) -- (0.0000,8.0000) -- (12.0000,8.0000) -- (0.0000,8.0000) -- (12.0000,8.0000) -- (0.0000,8.0000) -- (12.0000,8.0000) -- (0.0000,8.0000) -- (12.0000,8.0000) -- (0.0000,8.0000) -- (12.0000,8.0000) -- (0.0000,8.0000) -- (12.0000,8.0000) -- (0.0000,8.0000) -- (12.0000,8.0000) -- (0.0000,8.0000) -- (12.0000,8.0000) -- (0.0000,8.0000) -- (12.0000,8.0000) -- (0.0000,8.0000) -- (12.0000,8.0000) -- (0.0000,8.0000) -- (12.0000,8.0000) -- (0.0000,8.0000) -- (12.0000,8.0000) -- (0.0000,8.0000) -- (12.0000,8.0000) -- (0.0000,8.0000) -- (12.0000,8.0000) -- (0.0000,8.0000) -- (12.0000,8.0000) -- (0.0000,8.0000) -- (12.0000,8.0000) -- (0.0000,8.0000) -- (12.0000,8.0000) -- (0.0000,8.0000) -- (12.0000,8.0000) -- (0.0000,8.0000) -- (12.0000,8.0000) -- (0.0000,8.0000) -- (12.0000,8.0000) -- (0.0000,8.0000) -- (12.0000,8.0000) -- (0.0000,8.0000) -- (12.0000,8.0000) -- (0.0000,8.0000) -- (12.0000,8.0000) -- (0.0000,8.0000) -- (12.0000,8.0000) -- (0.0000,8.0000) -- (12.0000,8.0000) -- (0.0000,8.0000) -- (12.0000,8.0000) -- (0.0000,8.0000) -- (12.0000,8.0000) -- (0.0000,8.0000) -- (12.0000,8.0000) -- (0.0000,8.0000) -- (12.0000,8.0000) -- (0.0000,8.0000) -- (12.0000,8.0000) -- (0.0000,8.0000) -- (12.0000,8.0000) -- (0.0000,8.0000) -- (12.0000,8.0000) -- (0.0000,8.0000) -- (12.0000,8.0000) -- (0.0000,8.0000) -- (12.0000,8.0000) -- (0.0000,8.0000) -- (12.0000,8.0000) -- (0.0000,8.0000) -- (12.0000,8.0000) -- (0.0000,8.0000) -- (12.0000,8.0000) -- (0.0000,8.0000) -- (12.0000,8.0000) -- (0.0000,8.0000) -- (12.0000,8.0000) -- (0.0000,8.0000) -- (12.0000,8.0000) -- (0.0000,8.0000) -- (12.0000,8.0000) -- (0.0000,8.0000) -- (12.0000,8.0000) -- (0.0000,8.0000) -- (12.0000,8.0000) -- (0.0000,8.0000) -- (12.0000,8.0000) -- (0.0000,8.0000) -- (12.0000,8.0000) -- (0.0000,8.0000) -- (12.0000,8.0000) -- (0.0000,8.0000) -- (12.0000,8.0000) -- (0.0000,8.0000) -- (12.0000,8.0000) -- (0.0000,8.0000) -- (12.0000,8.0000) -- (0.0000,8.0000) -- (12.0000,8.0000) -- (0.0000,8.0000) -- (12.0000,8.0000) -- (0.0000,8.0000) -- (12.0000,8.0000) -- (0.0000,8.0000) -- (12.0000,8.0000) -- (0.0000,8.0000) -- (12.0000,8.0000) -- (0.0000,8.0000) -- (12.0000,8.0000) -- (0.0000,8.0000) -- (12.0000,8.0000) -- (0.0000,8.0000) -- (12.0000,8.0000) -- (0.0000,8.0000) -- (12.0000,8.0000) -- (0.0000,8.0000) -- (12.0000,8.0000) -- (0.0000,8.0000) -- (12.0000,8.0000) -- (0.0000,8.0000) -- (12.0000,8.0000) -- (0.0000,8.0000) -- (12.0000,8.0000) -- (0.0000,8.0000) -- (12.0000,8.0000) -- (0.0000,8.0000) -- (12.0000,8.0000) -- (0.0000,8.0000) -- (12.0000,8.0000) -- (0.0000,8.0000) -- (12.0000,8.0000) -- (0.0000,8.0000) -- (12.0000,8.0000) -- (0.0000,8.0000) -- (12.0000,8.0000) -- (0.0000,8.0000) -- (12.0000,8.0000) -- (0.0000,8.0000) -- (12.0000,8.0000) -- (0.0000,8.0000) -- (12.0000,8.0000) -- (0.0000,8.0000) -- (12.0000,8.0000) -- (0.0000,8.0000) -- (12.0000,8.0000) -- (0.0000,8.0000) -- (12.0000,8.0000) -- (0.0000,8.0000) -- (12.0000,8.0000) -- (0.0000,8.0000) -- (12.0000,8.0000) -- (0.0000,8.0000) -- (12.0000,8.0000) -- (0.0000,8.0000) -- (12.0000,8.0000) -- (0.0000,8.0000) -- (12.0000,8.0000) -- (0.0000,8.0000) -- (12.0000,8.0000) -- (0.0000,8.0000) -- (12.0000,8.0000) -- (0.0000,8.0000) -- (12.0000,8.0000) -- (0.0000,8.0000) -- (12.0000,8.0000) -- (0.0000,8.0000) -- (12.0000,8.0000) -- (0.0000,8.0000) -- (12.0000,8.0000) -- (0.0000,8.0000) -- (12.0000,8.0000) -- (0.0000,8.0000) -- (12.0000,8.0000) -- (0.0000,8.0000) -- (12.0000,8.0000) -- (0.0000,8.0000) -- (12.0000,8.0000) -- (0.0000,8.0000) -- (12.0000,8.0000) -- (0.0000,8.0000) -- (12.0000,8.0000) -- (0.0000,8.0000) -- (12.0000,8.0000) -- (0.0000,8.0000) -- (12.0000,8.0000) -- (0.0000,8.0000) -- (12.0000,8.0000) -- (0.0000,8.0000) -- (12.0000,8.0000) -- (0.0000,8.0000) -- (12.0000,8.0000) -- (0.0000,8.0000) -- (12.0000,8.0000) -- (0.0000,8.0000) -- (12.0000,8.0000) -- (0.0000,8.0000) -- (12.0000,8.0000) -- (0.0000,8.0000) -- (12.0000,8.0000) -- (0.0000,8.0000) -- (12.0000,8.0000) -- (0.0000,8.0000) -- (12.0000,8.0000) -- (0.0000,8.0000) -- (12.0000,8.0000) -- (0.0000,8.0000) -- (12.0000,8.0000) -- (0.0000,8.0000) -- (12.0000,8.0000) -- (0.0000,8.0000) -- (12.0000,8.0000) -- (0.0000,8.0000) -- (12.0000,8.0000) -- (0.0000,8.0000) -- (12.0000,8.0000) -- (0.0000,8.0000) -- (12.0000,8.0000) -- (0.0000,8.0000) -- (12.0000,8.0000) -- (0.0000,8.0000) -- (12.0000,8.0000) -- (0.0000,8.0000) -- (12.0000,8.0000) -- (0.0000,8.0000) -- (12.0000,8.0000) -- (0.0000,8.0000) -- (12.0000,8.0000) -- (0.0000,8.0000) -- (12.0000,8.0000) -- (0.0000,8.0000) -- (12.0000,8.0000) -- (0.0000,8.0000) -- (12.0000,8.0000) -- (0.0000,8.0000) -- (12.0000,8.0000) -- (0.0000,8.0000) -- (12.0000,8.0000) -- (0.0000,8.0000) -- (12.0000,8.0000) -- (0.0000,8.0000) -- (12.0000,8.0000) -- (0.0000,8.0000) -- (12.0000,8.0000) -- (0.0000,8.0000) -- (12.0000,8.0000) -- (0.0000,8.0000) -- (12.0000,8.0000) -- (0.0000,8.0000) -- (12.0000,8.0000) -- (0.0000,8.0000) -- (12.0000,8.0000) -- (0.0000,8.0000) -- (12.0000,8.0000) -- (0.0000,8.0000) -- (12.0000,8.0000) -- (0.0000,8.0000) -- (12.0000,8.0000) -- (0.0000,8.0000) -- (12.0000,8.0000) -- (0.0000,8.0000) -- (12.0000,8.0000) -- (0.0000,8.0000) -- (12.0000,8.0000) -- (0.0000,8.0000) -- (12.0000,8.0000) -- (0.0000,8.0000) -- (12.0000,8.0000) -- (0.0000,8.0000) -- (12.0000,8.0000) -- (0.0000,8.0000) -- (12.0000,8.0000) -- (0.0000,8.0000) -- (12.0000,8.0000) -- (0.0000,8.0000) -- (12.0000,8.0000) -- (0.0000,8.0000) -- (12.0000,8.0000) -- (0.0000,8.0000) -- (12.0000,8.0000) -- (0.0000,8.0000) -- (12.0000,8.0000) -- (0.0000,8.0000) -- (12.0000,8.0000) -- (0.0000,8.0000) -- (12.0000,8.0000) -- (0.0000,8.0000) -- (12.0000,8.0000) -- (0.0000,8.0000) -- (12.0000,8.0000) -- (0.0000,8.0000) -- (12.0000,8.0000) -- (0.0000,8.0000) -- (12.0000,8.0000) -- (0.0000,8.0000) -- (12.0000,8.0000) -- (0.0000,8.0000) -- (12.0000,8.0000) -- (0.0000,8.0000) -- (12.0000,8.0000) -- (0.0000,8.0000) -- (12.0000,8.0000) -- (0.0000,8.0000) -- (12.0000,8.0000) -- (0.0000,8.0000) -- (12.0000,8.0000) -- (0.0000,8.0000) -- (12.0000,8.0000) -- (0.0000,8.0000) -- (12.0000,8.0000) -- (0.0000,8.0000) -- (12.0000,8.0000) -- (0.0000,8.0000) -- (12.0000,8.0000) -- (0.0000,8.0000) -- (12.0000,8.0000) -- (0.0000,8.0000) -- (12.0000,8.0000) -- (0.0000,8.0000) -- (12.0000,8.0000) -- (0.0000,8.0000) -- (12.0000,8.0000) -- (0.0000,8.0000) -- (12.0000,8.0000) -- (0.0000,8.0000) -- (12.0000,8.0000) -- (0.0000,8.0000) -- (12.0000,8.0000) -- (0.0000,8.0000) -- (12.0000,8.0000) -- (0.0000,8.0000) -- (12.0000,8.0000) -- (0.0000,8.0000) -- (12.0000,8.0000) -- (0.0000,8.0000) -- (12.0000,8.0000) -- (0.0000,8.0000) -- (12.0000,8.0000) -- (0.0000,8.0000) -- (12.0000,8.0000) -- (0.0000,8.0000) -- (12.0000,8.0000) -- (0.0000,8.0000) -- (12.0000,8.0000) -- (0.0000,8.0000) -- (12.0000,8.0000) -- (0.0000,8.0000) -- (12.0000,8.0000) -- (0.0000,8.0000) -- (12.0000,8.0000) -- (0.0000,8.0000) -- (12.0000,8.0000) -- (0.0000,8.0000) -- (12.0000,8.0000) -- (0.0000,8.0000) -- (12.0000,8.0000) -- (0.0000,8.0000) -- (12.0000,8.0000) -- (0.0000,8.0000) -- (12.0000,8.0000) -- (0.0000,8.0000) -- (12.0000,8.0000) -- (0.0000,8.0000) -- (12.0000,8.0000) -- (0.0000,8.0000) -- (12.0000,8.0000) -- (0.0000,8.0000) -- (12.0000,8.0000) -- (0.0000,8.0000) -- (12.0000,8.0000) -- (0.0000,8.0000) -- (12.0000,8.0000) -- (0.0000,8.0000) -- (12.0000,8.0000) -- (0.0000,8.0000) -- (12.0000,8.0000) -- (0.0000,8.0000) -- (12.0000,8.0000) -- (0.0000,8.0000) -- (12.0000,8.0000) -- (0.0000,8.0000) -- (12.0000,8.0000) -- (0.0000,8.0000) -- (12.0000,8.0000) -- (0.0000,8.0000) -- (12.0000,8.0000) -- (0.0000,8.0000) -- (12.0000,8.0000) -- (0.0000,8.0000) -- (12.0000,8.0000) -- (0.0000,8.0000) -- (12.0000,8.0000) -- (0.0000,8.0000) -- (12.0000,8.0000) -- (0.0000,8.0000) -- (12.0000,8.0000) -- (0.0000,8.0000) -- (12.0000,8.0000) -- (0.0000,8.0000) -- (12.0000,8.0000) -- (0.0000,8.0000) -- (12.0000,8.0000) -- (0.0000,8.0000) -- (12.0000,8.0000) -- (0.0000,8.0000) -- (12.0000,8.0000) -- (0.0000,8.0000) -- (12.0000,8.0000) -- (0.0000,8.0000) -- (12.0000,8.0000) -- (0.0000,8.0000) -- (12.0000,8.0000) -- (0.0000,8.0000) -- (12.0000,8.0000) -- (0.0000,8.0000) -- (12.0000,8.0000) -- (0.0000,8.0000) -- (12.0000,8.0000) -- (0.0000,8.0000) -- (12.0000,8.0000) -- (0.0000,8.0000) -- (12.0000,8.0000) -- (0.0000,8.0000) -- (12.0000,8.0000) -- (0.0000,8.0000) -- (12.0000,8.0000) -- (0.0000,8.0000) -- (12.0000,8.0000) -- (0.0000,8.0000) -- (12.0000,8.0000) -- (0.0000,8.0000) -- (12.0000,8.0000) -- (0.0000,8.0000) -- (12.0000,8.0000) -- (0.0000,8.0000) -- (12.0000,8.0000) -- (0.0000,8.0000) -- (12.0000,8.0000) -- (0.0000,8.0000) -- (12.0000,8.0000) -- (0.0000,8.0000) -- (12.0000,8.0000) -- (0.0000,8.0000) -- (12.0000,8.0000) -- (0.0000,8.0000) -- (12.0000,8.0000) -- (0.0000,8.0000) -- (12.0000,8.0000) -- (0.0000,8.0000) -- (12.0000,8.0000) -- (0.0000,8.0000) -- (12.0000,8.0000) -- (0.0000,8.0000) -- (12.0000,8.0000) -- (0.0000,8.0000) -- (12.0000,8.0000) -- (0.0000,8.0000) -- (12.0000,8.0000) -- (0.0000,8.0000) -- (12.0000,8.0000) -- (0.0000,8.0000) -- (12.0000,8.0000) -- (0.0000,8.0000) -- (12.0000,8.0000) -- (0.0000,8.0000) -- (12.0000,8.0000) -- (0.0000,8.0000) -- (12.0000,8.0000) -- (0.0000,8.0000) -- (12.0000,8.0000) -- (0.0000,8.0000) -- (12.0000,8.0000) -- (0.0000,8.0000) -- (12.0000,8.0000) -- (0.0000,8.0000) -- (12.0000,8.0000) -- (0.0000,8.0000) -- (12.0000,8.0000) -- (0.0000,8.0000) -- (12.0000,8.0000) -- (0.0000,8.0000) -- (12.0000,8.0000) -- (0.0000,8.0000) -- (12.0000,8.0000) -- (0.0000,8.0000) -- (12.0000,8.0000) -- (0.0000,8.0000) -- (12.0000,8.0000) -- (0.0000,8.0000) -- (12.0000,8.0000) -- (0.0000,8.0000) -- (12.0000,8.0000) -- (0.0000,8.0000) -- (12.0000,8.0000) -- (0.0000,8.0000) -- (12.0000,8.0000) -- (0.0000,8.0000) -- (12.0000,8.0000) -- (0.0000,8.0000) -- (12.0000,8.0000) -- (0.0000,8.0000) -- (12.0000,8.0000) -- (0.0000,8.0000) -- (12.0000,8.0000) -- (0.0000,8.0000) -- (12.0000,8.0000) -- (0.0000,8.0000) -- (12.0000,8.0000) -- (0.0000,8.0000) -- (12.0000,8.0000) -- (0.0000,8.0000) -- (12.0000,8.0000) -- (0.0000,8.0000) -- (12.0000,8.0000) -- (0.0000,8.0000) -- (12.0000,8.0000) -- (0.0000,8.0000) -- (12.0000,8.0000) -- (0.0000,8.0000) -- (12.0000,8.0000) -- (0.0000,8.0000) -- (12.0000,8.0000) -- (0.0000,8.0000) -- (12.0000,8.0000) -- (0.0000,8.0000) -- (12.0000,8.0000) -- (0.0000,8.0000) -- (12.0000,8.0000) -- (0.0000,8.0000) -- (12.0000,8.0000) -- (0.0000,8.0000) -- (12.0000,8.0000) -- (0.0000,8.0000) -- (12.0000,8.0000) -- (0.0000,8.0000) -- (12.0000,8.0000) -- (0.0000,8.0000) -- (12.0000,8.0000) -- (0.0000,8.0000) -- (12.0000,8.0000) -- (0.0000,8.0000) -- (12.0000,8.0000) -- (0.0000,8.0000) -- (12.0000,8.0000) -- (0.0000,8.0000) -- (12.0000,8.0000) -- (0.0000,8.0000) -- (12.0000,8.0000) -- (0.0000,8.0000) -- (12.0000,8.0000) -- (0.0000,8.0000) -- (12.0000,8.0000) -- (0.0000,8.0000) -- (12.0000,8.0000) -- (0.0000,8.0000) -- (12.0000,8.0000) -- (0.0000,8.0000) -- (12.0000,8.0000) -- (0.0000,8.0000) -- (12.0000,8.0000) -- (0.0000,8.0000) -- (12.0000,8.0000) -- (0.0000,8.0000) -- (12.0000,8.0000) -- (0.0000,8.0000) -- (12.0000,8.0000) -- (0.0000,8.0000) -- (12.0000,8.0000) -- (0.0000,8.0000) -- (12.0000,8.0000) -- (0.0000,8.0000) -- (12.0000,8.0000) -- (0.0000,8.0000) -- (12.0000,8.0000) -- (0.0000,8.0000) -- (12.0000,8.0000) -- (0.0000,8.0000) -- (12.0000,8.0000) -- (0.0000,8.0000) -- (12.0000,8.0000) -- (0.0000,8.0000) -- (12.0000,8.0000) -- (0.0000,8.0000) -- (12.0000,8.0000) -- (0.0000,8.0000) -- (12.0000,8.0000) -- (0.0000,8.0000) -- (12.0000,8.0000) -- (0.0000,8.0000) -- (12.0000,8.0000) -- (0.0000,8.0000) -- (12.0000,8.0000) -- (0.0000,8.0000) -- (12.0000,8.0000) -- (0.0000,8.0000) -- (12.0000,8.0000) -- (0.0000,8.0000) -- (12.0000,8.0000) -- (0.0000,8.0000) -- (12.0000,8.0000) -- (0.0000,8.0000) -- (12.0000,8.0000) -- (0.0000,8.0000) -- (12.0000,8.0000) -- (0.0000,8.0000) -- (12.0000,8.0000) -- (0.0000,8.0000) -- (12.0000,8.0000) -- (0.0000,8.0000) -- (12.0000,8.0000) -- (0.0000,8.0000) -- (12.0000,8.0000) -- (0.0000,8.0000) -- (12.0000,8.0000) -- (0.0000,8.0000) -- (12.0000,8.0000) -- (0.0000,8.0000) -- (12.0000,8.0000) -- (0.0000,8.0000) -- (12.0000,8.0000) -- (0.0000,8.0000) -- (12.0000,8.0000) -- (0.0000,8.0000) -- (12.0000,8.0000) -- (0.0000,8.0000) -- (12.0000,8.0000) -- (0.0000,8.0000) -- (12.0000,8.0000) -- (0.0000,8.0000) -- (12.0000,8.0000) -- (0.0000,8.0000) -- (12.0000,8.0000) -- (0.0000,8.0000) -- (12.0000,8.0000) -- (0.0000,8.0000) -- (12.0000,8.0000) -- (0.0000,8.0000) -- (12.0000,8.0000) -- (0.0000,8.0000) -- (12.0000,8.0000) -- (0.0000,8.0000) -- (12.0000,8.0000) -- (0.0000,8.0000) -- (12.0000,8.0000) -- (0.0000,8.0000) -- (12.0000,8.0000) -- (0.0000,8.0000) -- (12.0000,8.0000) -- (0.0000,8.0000) -- (12.0000,8.0000) -- (0.0000,8.0000) -- (12.0000,8.0000) -- (0.0000,8.0000) -- (12.0000,8.0000) -- (0.0000,8.0000) -- (12.0000,8.0000) -- (0.0000,8.0000) -- (12.0000,8.0000) -- (0.0000,8.0000) -- (12.0000,8.0000) -- (0.0000,8.0000) -- (12.0000,8.0000) -- (0.0000,8.0000) -- (12.0000,8.0000) -- (0.0000,8.0000) -- (12.0000,8.0000) -- (0.0000,8.0000) -- (12.0000,8.0000) -- (0.0000,8.0000) -- (12.0000,8.0000) -- (0.0000,8.0000) -- (12.0000,8.0000) -- (0.0000,8.0000) -- (12.0000,8.0000) -- (0.0000,8.0000) -- (12.0000,8.0000) -- (0.0000,8.0000) -- (12.0000,8.0000) -- (0.0000,8.0000) -- (12.0000,8.0000) -- (0.0000,8.0000) -- (12.0000,8.0000) -- (0.0000,8.0000) -- (12.0000,8.0000) -- (0.0000,8.0000) -- (12.0000,8.0000) -- (0.0000,8.0000) -- (12.0000,8.0000) -- (0.0000,8.0000) -- (12.0000,8.0000) -- (0.0000,8.0000) -- (12.0000,8.0000) -- (0.0000,8.0000) -- (12.0000,8.0000) -- (0.0000,8.0000) -- (12.0000,8.0000) -- (0.0000,8.0000) -- (12.0000,8.0000) -- (0.0000,8.0000) -- (12.0000,8.0000) -- (0.0000,8.0000) -- (12.0000,8.0000) -- (0.0000,8.0000) -- (12.0000,8.0000) -- (0.0000,8.0000) -- (12.0000,8.0000) -- (0.0000,8.0000) -- (12.0000,8.0000) -- (0.0000,8.0000) -- (12.0000,8.0000) -- (0.0000,8.0000) -- (12.0000,8.0000) -- (0.0000,8.0000) -- (12.0000,8.0000) -- (0.0000,8.0000) -- (12.0000,8.0000) -- (0.0000,8.0000) -- (12.0000,8.0000) -- (0.0000,8.0000) -- (12.0000,8.0000) -- (0.0000,8.0000) -- (12.0000,8.0000) -- (0.0000,8.0000) -- (12.0000,8.0000) -- (0.0000,8.0000) -- (12.0000,8.0000) -- (0.0000,8.0000) -- (12.0000,8.0000) -- (0.0000,8.0000) -- (12.0000,8.0000) -- (0.0000,8.0000) -- (12.0000,8.0000) -- (0.0000,8.0000) -- (12.0000,8.0000) -- (0.0000,8.0000) -- (12.0000,8.0000) -- (0.0000,8.0000) -- (12.0000,8.0000) -- (0.0000,8.0000) -- (12.0000,8.0000) -- (0.0000,8.0000) -- (12.0000,8.0000) -- (0.0000,8.0000) -- (12.0000,8.0000) -- (0.0000,8.0000) -- (12.0000,8.0000) -- (0.0000,8.0000) -- (12.0000,8.0000) -- (0.0000,8.0000) -- (12.0000,8.0000) -- (0.0000,8.0000) -- (12.0000,8.0000) -- (0.0000,8.0000) -- (12.0000,8.0000) -- (0.0000,8.0000) -- (12.0000,8.0000) -- (0.0000,8.0000) -- (12.0000,8.0000) -- (0.0000,8.0000) -- (12.0000,8.0000) -- (0.0000,8.0000) -- (12.0000,8.0000) -- (0.0000,8.0000) -- (12.0000,8.0000) -- (0.0000,8.0000) -- (12.0000,8.0000) -- (0.0000,8.0000) -- (12.0000,8.0000) -- (0.0000,8.0000) -- (12.0000,8.0000) -- (0.0000,8.0000) -- (12.0000,8.0000) -- (0.0000,8.0000) -- (12.0000,8.0000) -- (0.0000,8.0000) -- (12.0000,8.0000) -- (0.0000,8.0000) -- (12.0000,8.0000) -- (0.0000,8.0000) -- (12.0000,8.0000) -- (0.0000,8.0000) -- (12.0000,8.0000) -- (0.0000,8.0000) -- (12.0000,8.0000) -- (0.0000,8.0000) -- (12.0000,8.0000) -- (0.0000,8.0000) -- (12.0000,8.0000) -- (0.0000,8.0000) -- (12.0000,8.0000) -- (0.0000,8.0000) -- (12.0000,8.0000) -- (0.0000,8.0000) -- (12.0000,8.0000) -- (0.0000,8.0000) -- (12.0000,8.0000) -- (0.0000,8.0000) -- (12.0000,8.0000) -- (0.0000,8.0000) -- (12.0000,8.0000) -- (0.0000,8.0000) -- (12.0000,8.0000) -- (0.0000,8.0000) -- (12.0000,8.0000) -- (0.0000,8.0000) -- (12.0000,8.0000) -- (0.0000,8.0000) -- (12.0000,8.0000) -- (0.0000,8.0000) -- (12.0000,8.0000) -- (0.0000,8.0000) -- (12.0000,8.0000) -- (0.0000,8.0000) -- (12.0000,8.0000) -- (0.0000,8.0000) -- (12.0000,8.0000) -- (0.0000,8.0000) -- (12.0000,8.0000) -- (0.0000,8.0000) -- (12.0000,8.0000) -- (0.0000,8.0000) -- (12.0000,8.0000) -- (0.0000,8.0000) -- (12.0000,8.0000) -- (0.0000,8.0000) -- (12.0000,8.0000) -- (0.0000,8.0000) -- (12.0000,8.0000) -- (0.0000,8.0000) -- (12.0000,8.0000) -- (0.0000,8.0000) -- (12.0000,8.0000) -- (0.0000,8.0000) -- (12.0000,8.0000) -- (0.0000,8.0000) -- (12.0000,8.0000) -- (0.0000,8.0000) -- (12.0000,8.0000) -- (0.0000,8.0000) -- (12.0000,8.0000) -- (0.0000,8.0000) -- (12.0000,8.0000) -- (0.0000,8.0000) -- (12.0000,8.0000) -- (0.0000,8.0000) -- (12.0000,8.0000) -- (0.0000,8.0000) -- (12.0000,8.0000) -- (0.0000,8.0000) -- (12.0000,8.0000) -- (0.0000,8.0000) -- (12.0000,8.0000) -- (0.0000,8.0000) -- (12.0000,8.0000) -- (0.0000,8.0000) -- (12.0000,8.0000) -- (0.0000,8.0000) -- (12.0000,8.0000) -- (0.0000,8.0000) -- (12.0000,8.0000) -- (0.0000,8.0000) -- (12.0000,8.0000) -- (0.0000,8.0000) -- (12.0000,8.0000) -- (0.0000,8.0000) -- (12.0000,8.0000) -- (0.0000,8.0000) -- (12.0000,8.0000) -- (0.0000,8.0000) -- (12.0000,8.0000) -- (0.0000,8.0000) -- (12.0000,8.0000) -- (0.0000,8.0000) -- (12.0000,8.0000) -- (0.0000,8.0000) -- (12.0000,8.0000) -- (0.0000,8.0000) -- (12.0000,8.0000) -- (0.0000,8.0000) -- (12.0000,8.0000) -- (0.0000,8.0000) -- (12.0000,8.0000) -- (0.0000,8.0000) -- (12.0000,8.0000) -- (0.0000,8.0000) -- (12.0000,8.0000) -- (0.0000,8.0000) -- (12.0000,8.0000) -- (0.0000,8.0000) -- (12.0000,8.0000) -- (0.0000,8.0000) -- (12.0000,8.0000) -- (0.0000,8.0000) -- (12.0000,8.0000) -- (0.0000,8.0000) -- (12.0000,8.0000) -- (0.0000,8.0000) -- (12.0000,8.0000) -- (0.0000,8.0000) -- (12.0000,8.0000) -- (0.0000,8.0000) -- (12.0000,8.0000) -- (0.0000,8.0000) -- (12.0000,8.0000) -- (0.0000,8.0000) -- (12.0000,8.0000) -- (0.0000,8.0000) -- (12.0000,8.0000) -- (0.0000,8.0000) -- (12.0000,8.0000) -- (0.0000,8.0000) -- (12.0000,8.0000) -- (0.0000,8.0000) -- (12.0000,8.0000) -- (0.0000,8.0000) -- (12.0000,8.0000) -- (0.0000,8.0000) -- (12.0000,8.0000) -- (0.0000,8.0000) -- (12.0000,8.0000) -- (0.0000,8.0000) -- (12.0000,8.0000) -- (0.0000,8.0000) -- (12.0000,8.0000) -- (0.0000,8.0000) -- (12.0000,8.0000) -- (0.0000,8.0000) -- (12.0000,8.0000) -- (0.0000,8.0000) -- (12.0000,8.0000) -- (0.0000,8.0000) -- (12.0000,8.0000) -- (0.0000,8.0000) -- (12.0000,8.0000) -- (0.0000,8.0000) -- (12.0000,8.0000) -- (0.0000,8.0000) -- (12.0000,8.0000) -- (0.0000,8.0000) -- (12.0000,8.0000) -- (0.0000,8.0000) -- (12.0000,8.0000) -- (0.0000,8.0000) -- (12.0000,8.0000) -- (0.0000,8.0000) -- (12.0000,8.0000) -- (0.0000,8.0000) -- (12.0000,8.0000) -- (0.0000,8.0000) -- (12.0000,8.0000) -- (0.0000,8.0000) -- (12.0000,8.0000) -- (0.0000,8.0000) -- (12.0000,8.0000) -- (0.0000,8.0000) -- (12.0000,8.0000) -- (0.0000,8.0000) -- (12.0000,8.0000) -- (0.0000,8.0000) -- (12.0000,8.0000) -- (0.0000,8.0000) -- (12.0000,8.0000) -- (0.0000,8.0000) -- (12.0000,8.0000) -- (0.0000,8.0000) -- (12.0000,8.0000) -- (0.0000,8.0000) -- (12.0000,8.0000) -- (0.0000,8.0000) -- (12.0000,8.0000) -- (0.0000,8.0000) -- (12.0000,8.0000) -- (0.0000,8.0000) -- (12.0000,8.0000) -- (0.0000,8.0000) -- (12.0000,8.0000) -- (0.0000,8.0000) -- (12.0000,8.0000) -- (0.0000,8.0000) -- (12.0000,8.0000) -- (0.0000,8.0000) -- (12.0000,8.0000) -- (0.0000,8.0000) -- (12.0000,8.0000) -- (0.0000,8.0000) -- (12.0000,8.0000) -- (0.0000,8.0000) -- (12.0000,8.0000) -- (0.0000,8.0000) -- (12.0000,8.0000) -- (0.0000,8.0000) -- (12.0000,8.0000) -- (0.0000,8.0000) -- (12.0000,8.0000) -- (0.0000,8.0000) -- (12.0000,8.0000) -- (0.0000,8.0000) -- (12.0000,8.0000) -- (0.0000,8.0000) -- (12.0000,8.0000) -- (0.0000,8.0000) -- (12.0000,8.0000) -- (0.0000,8.0000) -- (12.0000,8.0000) -- (0.0000,8.0000) -- (12.0000,8.0000) -- (0.0000,8.0000) -- (12.0000,8.0000) -- (0.0000,8.0000) -- (12.0000,8.0000) -- (0.0000,8.0000) -- (12.0000,8.0000) -- (0.0000,8.0000) -- (12.0000,8.0000) -- (0.0000,8.0000) -- (12.0000,8.0000) -- (0.0000,8.0000) -- (12.0000,8.0000) -- (0.0000,8.0000) -- (12.0000,8.0000) -- (0.0000,8.0000) -- (12.0000,8.0000) -- (0.0000,8.0000) -- (12.0000,8.0000) -- (0.0000,8.0000) -- (12.0000,8.0000) -- (0.0000,8.0000) -- (12.0000,8.0000) -- (0.0000,8.0000) -- (12.0000,8.0000) -- (0.0000,8.0000) -- (12.0000,8.0000) -- (0.0000,8.0000) -- (12.0000,8.0000) -- (0.0000,8.0000) -- (12.0000,8.0000) -- (0.0000,8.0000) -- (12.0000,8.0000) -- (0.0000,8.0000) -- (12.0000,8.0000) -- (0.0000,8.0000) -- (12.0000,8.0000) -- (0.0000,8.0000) -- (12.0000,8.0000) -- (0.0000,8.0000) -- (12.0000,8.0000) -- (0.0000,8.0000) -- (12.0000,8.0000) -- (0.0000,8.0000) -- (12.0000,8.0000) -- (0.0000,8.0000) -- (12.0000,8.0000) -- (0.0000,8.0000) -- (12.0000,8.0000) -- (0.0000,8.0000) -- (12.0000,8.0000) -- (0.0000,8.0000) -- (12.0000,8.0000) -- (0.0000,8.0000) -- (12.0000,8.0000) -- (0.0000,8.0000) -- (12.0000,8.0000) -- (0.0000,8.0000) -- (12.0000,8.0000) -- (0.0000,8.0000) -- (12.0000,8.0000) -- (0.0000,8.0000) -- (12.0000,8.0000) -- (0.0000,8.0000) -- (12.0000,8.0000) -- (0.0000,8.0000) -- (12.0000,8.0000) -- (0.0000,8.0000) -- (12.0000,8.0000) -- (0.0000,8.0000) -- (12.0000,8.0000) -- (0.0000,8.0000) -- (12.0000,8.0000) -- (0.0000,8.0000) -- (12.0000,8.0000) -- (0.0000,8.0000) -- (12.0000,8.0000) -- (0.0000,8.0000) -- (12.0000,8.0000) -- (0.0000,8.0000) -- (12.0000,8.0000) -- (0.0000,8.0000) -- (12.0000,8.0000) -- (0.0000,8.0000) -- (12.0000,8.0000) -- (0.0000,8.0000) -- (12.0000,8.0000) -- (0.0000,8.0000) -- (12.0000,8.0000) -- (0.0000,8.0000) -- (12.0000,8.0000) -- (0.0000,8.0000) -- (12.0000,8.0000) -- (0.0000,8.0000) -- (12.0000,8.0000) -- (0.0000,8.0000) -- (12.0000,8.0000) -- (0.0000,8.0000) -- (12.0000,8.0000) -- (0.0000,8.0000) -- (12.0000,8.0000) -- (0.0000,8.0000) -- (12.0000,8.0000) -- (0.0000,8.0000) -- (12.0000,8.0000) -- (0.0000,8.0000) -- (12.0000,8.0000) -- (0.0000,8.0000) -- (12.0000,8.0000) -- (0.0000,8.0000) -- (12.0000,8.0000) -- (0.0000,8.0000) -- (12.0000,8.0000) -- (0.0000,8.0000) -- (12.0000,8.0000) -- (0.0000,8.0000) -- (12.0000,8.0000) -- (0.0000,8.0000) -- (12.0000,8.0000) -- (0.0000,8.0000) -- (12.0000,8.0000) -- (0.0000,8.0000) -- (12.0000,8.0000) -- (0.0000,8.0000) -- (12.0000,8.0000) -- (0.0000,8.0000) -- (12.0000,8.0000) -- (0.0000,8.0000) -- (12.0000,8.0000) -- (0.0000,8.0000) -- (12.0000,8.0000) -- (0.0000,8.0000) -- (12.0000,8.0000) -- (0.0000,8.0000) -- (12.0000,8.0000) -- (0.0000,8.0000) -- (12.0000,8.0000) -- (0.0000,8.0000) -- (12.0000,8.0000) -- (0.0000,8.0000) -- (12.0000,8.0000) -- (0.0000,8.0000) -- (12.0000,8.0000) -- (0.0000,8.0000) -- (12.0000,8.0000) -- (0.0000,8.0000) -- (12.0000,8.0000) -- (0.0000,8.0000) -- (12.0000,8.0000) -- (0.0000,8.0000) -- (12.0000,8.0000) -- (0.0000,8.0000) -- (12.0000,8.0000) -- (0.0000,8.0000) -- (12.0000,8.0000) -- (0.0000,8.0000) -- (12.0000,8.0000) -- (0.0000,8.0000) -- (12.0000,8.0000) -- (0.0000,8.0000) -- (12.0000,8.0000) -- (0.0000,8.0000) -- (12.0000,8.0000) -- (0.0000,8.0000) -- (12.0000,8.0000) -- (0.0000,8.0000) -- (12.0000,8.0000) -- (0.0000,8.0000) -- (12.0000,8.0000) -- (0.0000,8.0000) -- (12.0000,8.0000) -- (0.0000,8.0000) -- (12.0000,8.0000) -- (0.0000,8.0000) -- (12.0000,8.0000) -- (0.0000,8.0000) -- (12.0000,8.0000) -- (0.0000,8.0000) -- (12.0000,8.0000) -- (0.0000,8.0000) -- (12.0000,8.0000) -- (0.0000,8.0000) -- (12.0000,8.0000) -- (0.0000,8.0000) -- (12.0000,8.0000) -- (0.0000,8.0000) -- (12.0000,8.0000) -- (0.0000,8.0000) -- (12.0000,8.0000) -- (0.0000,8.0000) -- (12.0000,8.0000) -- (0.0000,8.0000) -- (12.0000,8.0000) -- (0.0000,8.0000) -- (12.0000,8.0000) -- (0.0000,8.0000) -- (12.0000,8.0000) -- (0.0000,8.0000) -- (12.0000,8.0000) -- (0.0000,8.0000) -- (12.0000,8.0000) -- (0.0000,8.0000) -- (12.0000,8.0000) -- (0.0000,8.0000) -- (12.0000,8.0000) -- (0.0000,8.0000) -- (12.0000,8.0000) -- (0.0000,8.0000) -- (12.0000,8.0000) -- (0.0000,8.0000) -- (12.0000,8.0000) -- (0.0000,8.0000) -- (12.0000,8.0000) -- (0.0000,8.0000) -- (12.0000,8.0000) -- (0.0000,8.0000) -- (12.0000,8.0000) -- (0.0000,8.0000) -- (12.0000,8.0000) -- (0.0000,8.0000) -- (12.0000,8.0000) -- (0.0000,8.0000) -- (12.0000,8.0000) -- (0.0000,8.0000) -- (12.0000,8.0000) -- (0.0000,8.0000) -- (12.0000,8.0000) -- (0.0000,8.0000) -- (12.0000,8.0000) -- (0.0000,8.0000) -- (12.0000,8.0000) -- (0.0000,8.0000) -- (12.0000,8.0000) -- (0.0000,8.0000) -- (12.0000,8.0000) -- (0.0000,8.0000) -- (12.0000,8.0000) -- (0.0000,8.0000) -- (12.0000,8.0000) -- (0.0000,8.0000) -- (12.0000,8.0000) -- (0.0000,8.0000) -- (12.0000,8.0000) -- (0.0000,8.0000) -- (12.0000,8.0000) -- (0.0000,8.0000) -- (12.0000,8.0000) -- (0.0000,8.0000) -- (12.0000,8.0000) -- (0.0000,8.0000) -- (12.0000,8.0000) -- (0.0000,8.0000) -- (12.0000,8.0000) -- (0.0000,8.0000) -- (12.0000,8.0000) -- (0.0000,8.0000) -- (12.0000,8.0000) -- (0.0000,8.0000) -- (12.0000,8.0000) -- (0.0000,8.0000) -- (12.0000,8.0000) -- (0.0000,8.0000) -- (12.0000,8.0000) -- (0.0000,8.0000) -- (12.0000,8.0000) -- (0.0000,8.0000) -- (12.0000,8.0000) -- (0.0000,8.0000) -- (12.0000,8.0000) -- (0.0000,8.0000) -- (12.0000,8.0000) -- (0.0000,8.0000) -- (12.0000,8.0000) -- (0.0000,8.0000) -- (12.0000,8.0000) -- (0.0000,8.0000) -- (12.0000,8.0000) -- (0.0000,8.0000) -- (12.0000,8.0000) -- (0.0000,8.0000) -- (12.0000,8.0000) -- (0.0000,8.0000) -- (12.0000,8.0000) -- (0.0000,8.0000) -- (12.0000,8.0000) -- (0.0000,8.0000) -- (12.0000,8.0000) -- (0.0000,8.0000) -- (12.0000,8.0000) -- (0.0000,8.0000) -- (12.0000,8.0000) -- (0.0000,8.0000) -- (12.0000,8.0000) -- (0.0000,8.0000) -- (12.0000,8.0000) -- (0.0000,8.0000) -- (12.0000,8.0000) -- (0.0000,8.0000) -- (12.0000,8.0000) -- (0.0000,8.0000) -- (12.0000,8.0000) -- (0.0000,8.0000) -- (12.0000,8.0000) -- (0.0000,8.0000) -- (12.0000,8.0000) -- (0.0000,8.0000) -- (12.0000,8.0000) -- (0.0000,8.0000) -- (12.0000,8.0000) -- (0.0000,8.0000) -- (12.0000,8.0000) -- (0.0000,8.0000) -- (12.0000,8.0000) -- (0.0000,8.0000) -- (12.0000,8.0000) -- (0.0000,8.0000) -- (12.0000,8.0000) -- (0.0000,8.0000) -- (12.0000,8.0000) -- (0.0000,8.0000) -- (12.0000,8.0000) -- (0.0000,8.0000) -- (12.0000,8.0000) -- (0.0000,8.0000) -- (12.0000,8.0000) -- (0.0000,8.0000) -- (12.0000,8.0000) -- (0.0000,8.0000) -- (12.0000,8.0000) -- (0.0000,8.0000) -- (12.0000,8.0000) -- (0.0000,8.0000) -- (12.0000,8.0000) -- (0.0000,8.0000) -- (12.0000,8.0000) -- (0.0000,8.0000) -- (12.0000,8.0000) -- (0.0000,8.0000) -- (12.0000,8.0000) -- (0.0000,8.0000) -- (12.0000,8.0000) -- (0.0000,8.0000) -- (12.0000,8.0000) -- (0.0000,8.0000) -- (12.0000,8.0000) -- (0.0000,8.0000) -- (12.0000,8.0000) -- (0.0000,8.0000) -- (12.0000,8.0000) -- (0.0000,8.0000) -- (12.0000,8.0000) -- (0.0000,8.0000) -- (12.0000,8.0000) -- (0.0000,8.0000) -- (12.0000,8.0000) -- (0.0000,8.0000) -- (12.0000,8.0000) -- (0.0000,8.0000) -- (12.0000,8.0000) -- (0.0000,8.0000) -- (12.0000,8.0000) -- (0.0000,8.0000) -- (12.0000,8.0000) -- (0.0000,8.0000) -- (12.0000,8.0000) -- (0.0000,8.0000) -- (12.0000,8.0000) -- (0.0000,8.0000) -- (12.0000,8.0000) -- (0.0000,8.0000) -- (12.0000,8.0000) -- (0.0000,8.0000) -- (12.0000,8.0000) -- (0.0000,8.0000) -- (12.0000,8.0000) -- (0.0000,8.0000) -- (12.0000,8.0000) -- (0.0000,8.0000) -- (12.0000,8.0000) -- (0.0000,8.0000) -- (12.0000,8.0000) -- (0.0000,8.0000) -- (12.0000,8.0000) -- (0.0000,8.0000) -- (12.0000,8.0000) -- (0.0000,8.0000) -- (12.0000,8.0000) -- (0.0000,8.0000) -- (12.0000,8.0000) -- (0.0000,8.0000) -- (12.0000,8.0000) -- (0.0000,8.0000) -- (12.0000,8.0000) -- (0.0000,8.0000) -- (12.0000,8.0000) -- (0.0000,8.0000) -- (12.0000,8.0000) -- (0.0000,8.0000) -- (12.0000,8.0000) -- (0.0000,8.0000) -- (12.0000,8.0000) -- (0.0000,8.0000) -- (12.0000,8.0000) -- (0.0000,8.0000) -- (12.0000,8.0000) -- (0.0000,8.0000) -- (12.0000,8.0000) -- (0.0000,8.0000) -- (12.0000,8.0000) -- (0.0000,8.0000) -- (12.0000,8.0000) -- (0.0000,8.0000) -- (12.0000,8.0000) -- (0.0000,8.0000) -- (12.0000,8.0000) -- (0.0000,8.0000) -- (12.0000,8.0000) -- (0.0000,8.0000) -- (12.0000,8.0000) -- (0.0000,8.0000) -- (12.0000,8.0000) -- (0.0000,8.0000) -- (12.0000,8.0000) -- (0.0000,8.0000) -- (12.0000,8.0000) -- (0.0000,8.0000) -- (12.0000,8.0000) -- (0.0000,8.0000) -- (12.0000,8.0000) -- (0.0000,8.0000) -- (12.0000,8.0000) -- (0.0000,8.0000) -- (12.0000,8.0000) -- (0.0000,8.0000) -- (12.0000,8.0000) -- (0.0000,8.0000) -- (12.0000,8.0000) -- (0.0000,8.0000) -- (12.0000,8.0000) -- (0.0000,8.0000) -- (12.0000,8.0000) -- (0.0000,8.0000) -- (12.0000,8.0000) -- (0.0000,8.0000) -- (12.0000,8.0000) -- (0.0000,8.0000) -- (12.0000,8.0000) -- (0.0000,8.0000) -- (12.0000,8.0000) -- (0.0000,8.0000) -- (12.0000,8.0000) -- (0.0000,8.0000) -- (12.0000,8.0000) -- (0.0000,8.0000) -- (12.0000,8.0000) -- (0.0000,8.0000) -- (12.0000,8.0000) -- (0.0000,8.0000) -- (12.0000,8.0000) -- (0.0000,8.0000) -- (12.0000,8.0000) -- (0.0000,8.0000) -- (12.0000,8.0000) -- (0.0000,8.0000) -- (12.0000,8.0000) -- (0.0000,8.0000) -- (12.0000,8.0000) -- (0.0000,8.0000) -- (12.0000,8.0000) -- (0.0000,8.0000) -- (12.0000,8.0000) -- (0.0000,8.0000) -- (12.0000,8.0000) -- (0.0000,8.0000) -- (12.0000,8.0000) -- (0.0000,8.0000) -- (12.0000,8.0000) -- (0.0000,8.0000) -- (12.0000,8.0000) -- (0.0000,8.0000) -- (12.0000,8.0000) -- (0.0000,8.0000) -- (12.0000,8.0000) -- (0.0000,8.0000) -- (12.0000,8.0000) -- (0.0000,8.0000) -- (12.0000,8.0000) -- (0.0000,8.0000) -- (12.0000,8.0000) -- (0.0000,8.0000) -- (12.0000,8.0000) -- (0.0000,8.0000) -- (12.0000,8.0000) -- (0.0000,8.0000) -- (12.0000,8.0000) -- (0.0000,8.0000) -- (12.0000,8.0000) -- (0.0000,8.0000) -- (12.0000,8.0000) -- (0.0000,8.0000) -- (12.0000,8.0000) -- (0.0000,8.0000) -- (12.0000,8.0000) -- (0.0000,8.0000) -- (12.0000,8.0000) -- (0.0000,8.0000) -- (12.0000,8.0000) -- (0.0000,8.0000) -- (12.0000,8.0000) -- (0.0000,8.0000) -- (12.0000,8.0000) -- (0.0000,8.0000) -- (12.0000,8.0000) -- (0.0000,8.0000) -- (12.0000,8.0000) -- (0.0000,8.0000) -- (12.0000,8.0000) -- (0.0000,8.0000) -- (12.0000,8.0000) -- (0.0000,8.0000) -- (12.0000,8.0000) -- (0.0000,8.0000) -- (12.0000,8.0000) -- (0.0000,8.0000) -- (12.0000,8.0000) -- (0.0000,8.0000) -- (12.0000,8.0000) -- (0.0000,8.0000) -- (12.0000,8.0000) -- (0.0000,8.0000) -- (12.0000,8.0000) -- (0.0000,8.0000) -- (12.0000,8.0000) -- (0.0000,8.0000) -- (12.0000,8.0000) -- (0.0000,8.0000) -- (12.0000,8.0000) -- (0.0000,8.0000) -- (12.0000,8.0000) -- (0.0000,8.0000) -- (12.0000,8.0000) -- (0.0000,8.0000) -- (12.0000,8.0000) -- (0.0000,8.0000) -- (12.0000,8.0000) -- (0.0000,8.0000) -- (12.0000,8.0000) -- (0.0000,8.0000) -- (12.0000,8.0000) -- (0.0000,8.0000) -- (12.0000,8.0000) -- (0.0000,8.0000) -- (12.0000,8.0000) -- (0.0000,8.0000) -- (12.0000,8.0000) -- (0.0000,8.0000) -- (12.0000,8.0000) -- (0.0000,8.0000) -- (12.0000,8.0000) -- (0.0000,8.0000) -- (12.0000,8.0000) -- (0.0000,8.0000) -- (12.0000,8.0000) -- (0.0000,8.0000) -- (12.0000,8.0000) -- (0.0000,8.0000) -- (12.0000,8.0000) -- (0.0000,8.0000) -- (12.0000,8.0000) -- (0.0000,8.0000) -- (12.0000,8.0000) -- (0.0000,8.0000) -- (12.0000,8.0000) -- (0.0000,8.0000) -- (12.0000,8.0000) -- (0.0000,8.0000) -- (12.0000,8.0000) -- (0.0000,8.0000) -- (12.0000,8.0000) -- (0.0000,8.0000) -- (12.0000,8.0000) -- (0.0000,8.0000) -- (12.0000,8.0000) -- (0.0000,8.0000) -- (12.0000,8.0000) -- (0.0000,8.0000) -- (12.0000,8.0000) -- (0.0000,8.0000) -- (12.0000,8.0000) -- (0.0000,8.0000) -- (12.0000,8.0000) -- (0.0000,8.0000) -- (12.0000,8.0000) -- (0.0000,8.0000) -- (12.0000,8.0000) -- (0.0000,8.0000) -- (12.0000,8.0000) -- (0.0000,8.0000) -- (12.0000,8.0000) -- (0.0000,8.0000) -- (12.0000,8.0000) -- (0.0000,8.0000) -- (12.0000,8.0000) -- (0.0000,8.0000) -- (12.0000,8.0000) -- (0.0000,8.0000) -- (12.0000,8.0000) -- (0.0000,8.0000) -- (12.0000,8.0000) -- (0.0000,8.0000) -- (12.0000,8.0000) -- (0.0000,8.0000) -- (12.0000,8.0000) -- (0.0000,8.0000) -- (12.0000,8.0000) -- (0.0000,8.0000) -- (12.0000,8.0000) -- (0.0000,8.0000) -- (12.0000,8.0000) -- (0.0000,8.0000) -- (12.0000,8.0000) -- (0.0000,8.0000) -- (12.0000,8.0000) -- (0.0000,8.0000) -- (12.0000,8.0000) -- (0.0000,8.0000) -- (12.0000,8.0000) -- (0.0000,8.0000) -- (12.0000,8.0000) -- (0.0000,8.0000) -- (12.0000,8.0000) -- (0.0000,8.0000) -- (12.0000,8.0000) -- (0.0000,8.0000) -- (12.0000,8.0000) -- (0.0000,8.0000) -- (12.0000,8.0000) -- (0.0000,8.0000) -- (12.0000,8.0000) -- (0.0000,8.0000) -- (12.0000,8.0000) -- (0.0000,8.0000) -- (12.0000,8.0000) -- (0.0000,8.0000) -- (12.0000,8.0000) -- (0.0000,8.0000) -- (12.0000,8.0000) -- (0.0000,8.0000) -- (12.0000,8.0000) -- (0.0000,8.0000) -- (12.0000,8.0000) -- (0.0000,8.0000) -- (12.0000,8.0000) -- (0.0000,8.0000) -- (12.0000,8.0000) -- (0.0000,8.0000) -- (12.0000,8.0000) -- (0.0000,8.0000) -- (12.0000,8.0000) -- (0.0000,8.0000) -- (12.0000,8.0000) -- (0.0000,8.0000) -- (12.0000,8.0000) -- (0.0000,8.0000) -- (12.0000,8.0000) -- (0.0000,8.0000) -- (12.0000,8.0000) -- (0.0000,8.0000) -- (12.0000,8.0000) -- (0.0000,8.0000) -- (12.0000,8.0000) -- (0.0000,8.0000) -- (12.0000,8.0000) -- (0.0000,8.0000) -- (12.0000,8.0000) -- (0.0000,8.0000) -- (12.0000,8.0000) -- (0.0000,8.0000) -- (12.0000,8.0000) -- (0.0000,8.0000) -- (12.0000,8.0000) -- (0.0000,8.0000) -- (12.0000,8.0000) -- (0.0000,8.0000) -- (12.0000,8.0000) -- (0.0000,8.0000) -- (12.0000,8.0000) -- (0.0000,8.0000) -- (12.0000,8.0000) -- (0.0000,8.0000) -- (12.0000,8.0000) -- (0.0000,8.0000) -- (12.0000,8.0000) -- (0.0000,8.0000) -- (12.0000,8.0000) -- (0.0000,8.0000) -- (12.0000,8.0000) -- (0.0000,8.0000) -- (12.0000,8.0000) -- (0.0000,8.0000) -- (12.0000,8.0000) -- (0.0000,8.0000) -- (12.0000,8.0000) -- (0.0000,8.0000) -- (12.0000,8.0000) -- (0.0000,8.0000) -- (12.0000,8.0000) -- (0.0000,8.0000) -- (12.0000,8.0000) -- (0.0000,8.0000) -- (12.0000,8.0000) -- (0.0000,8.0000) -- (12.0000,8.0000) -- (0.0000,8.0000) -- (12.0000,8.0000) -- (0.0000,8.0000) -- (12.0000,8.0000) -- (0.0000,8.0000) -- (12.0000,8.0000) -- (0.0000,8.0000) -- (12.0000,8.0000) -- (0.0000,8.0000) -- (12.0000,8.0000) -- (0.0000,8.0000) -- (12.0000,8.0000) -- (0.0000,8.0000) -- (12.0000,8.0000) -- (0.0000,8.0000) -- (12.0000,8.0000) -- (0.0000,8.0000) -- (12.0000,8.0000) -- (0.0000,8.0000) -- (12.0000,8.0000) -- (0.0000,8.0000) -- (12.0000,8.0000) -- (0.0000,8.0000) -- (12.0000,8.0000) -- (0.0000,8.0000) -- (12.0000,8.0000) -- (0.0000,8.0000) -- (12.0000,8.0000) -- (0.0000,8.0000) -- (12.0000,8.0000) -- (0.0000,8.0000) -- (12.0000,8.0000) -- (0.0000,8.0000) -- (12.0000,8.0000) -- (0.0000,8.0000) -- (12.0000,8.0000) -- (0.0000,8.0000) -- (12.0000,8.0000) -- (0.0000,8.0000) -- (12.0000,8.0000) -- (0.0000,8.0000) -- (12.0000,8.0000) -- (0.0000,8.0000) -- (12.0000,8.0000) -- (0.0000,8.0000) -- (12.0000,8.0000) -- (0.0000,8.0000) -- (12.0000,8.0000) -- (0.0000,8.0000) -- (12.0000,8.0000) -- (0.0000,8.0000) -- (12.0000,8.0000) -- (0.0000,8.0000) -- (12.0000,8.0000) -- (0.0000,8.0000) -- (12.0000,8.0000) -- (0.0000,8.0000) -- (12.0000,8.0000) -- (0.0000,8.0000) -- (12.0000,8.0000) -- (0.0000,8.0000) -- (12.0000,8.0000) -- (0.0000,8.0000) -- (12.0000,8.0000) -- (0.0000,8.0000) -- (12.0000,8.0000) -- (0.0000,8.0000) -- (12.0000,8.0000) -- (0.0000,8.0000) -- (12.0000,8.0000) -- (0.0000,8.0000) -- (12.0000,8.0000) -- (0.0000,8.0000) -- (12.0000,8.0000) -- (0.0000,8.0000) -- (12.0000,8.0000) -- (0.0000,8.0000) -- (12.0000,8.0000) -- (0.0000,8.0000) -- (12.0000,8.0000) -- (0.0000,8.0000) -- (12.0000,8.0000) -- (0.0000,8.0000) -- (12.0000,8.0000) -- (0.0000,8.0000) -- (12.0000,8.0000) -- (0.0000,8.0000) -- (12.0000,8.0000) -- (0.0000,8.0000) -- (12.0000,8.0000) -- (0.0000,8.0000) -- (12.0000,8.0000) -- (0.0000,8.0000) -- (12.0000,8.0000) -- (0.0000,8.0000) -- (12.0000,8.0000) -- (0.0000,8.0000) -- (12.0000,8.0000) -- (0.0000,8.0000) -- (12.0000,8.0000) -- (0.0000,8.0000) -- (12.0000,8.0000) -- (0.0000,8.0000) -- (12.0000,8.0000) -- (0.0000,8.0000) -- (12.0000,8.0000) -- (0.0000,8.0000) -- (12.0000,8.0000) -- (0.0000,8.0000) -- (12.0000,8.0000) -- (0.0000,8.0000) -- (12.0000,8.0000) -- (0.0000,8.0000) -- (12.0000,8.0000) -- (0.0000,8.0000) -- (12.0000,8.0000) -- (0.0000,8.0000) -- (12.0000,8.0000) -- (0.0000,8.0000) -- (12.0000,8.0000) -- (0.0000,8.0000) -- (12.0000,8.0000) -- (0.0000,8.0000) -- (12.0000,8.0000) -- (0.0000,8.0000) -- (12.0000,8.0000) -- (0.0000,8.0000) -- (12.0000,8.0000) -- (0.0000,8.0000) -- (12.0000,8.0000) -- (0.0000,8.0000) -- (12.0000,8.0000) -- (0.0000,8.0000) -- (12.0000,8.0000) -- (0.0000,8.0000) -- (12.0000,8.0000) -- (0.0000,8.0000) -- (12.0000,8.0000) -- (0.0000,8.0000) -- (12.0000,8.0000) -- (0.0000,8.0000) -- (12.0000,8.0000) -- (0.0000,8.0000) -- (12.0000,8.0000) -- (0.0000,8.0000) -- (12.0000,8.0000) -- (0.0000,8.0000) -- (12.0000,8.0000) -- (0.0000,8.0000) -- (12.0000,8.0000) -- (0.0000,8.0000) -- (12.0000,8.0000) -- (0.0000,8.0000) -- (12.0000,8.0000) -- (0.0000,8.0000) -- (12.0000,8.0000) -- (0.0000,8.0000) -- (12.0000,8.0000) -- (0.0000,8.0000) -- (12.0000,8.0000) -- (0.0000,8.0000) -- (12.0000,8.0000) -- (0.0000,8.0000) -- (12.0000,8.0000) -- (0.0000,8.0000) -- (12.0000,8.0000) -- (0.0000,8.0000) -- (12.0000,8.0000) -- (0.0000,8.0000) -- (12.0000,8.0000) -- (0.0000,8.0000) -- (12.0000,8.0000) -- (0.0000,8.0000) -- (12.0000,8.0000) -- (0.0000,8.0000) -- (12.0000,8.0000) -- (0.0000,8.0000) -- (12.0000,8.0000) -- (0.0000,8.0000) -- (12.0000,8.0000) -- (0.0000,8.0000) -- (12.0000,8.0000) -- (0.0000,8.0000) -- (12.0000,8.0000) -- (0.0000,8.0000) -- (12.0000,8.0000) -- (0.0000,8.0000) -- (12.0000,8.0000) -- (0.0000,8.0000) -- (12.0000,8.0000) -- (0.0000,8.0000) -- (12.0000,8.0000) -- (0.0000,8.0000) -- (12.0000,8.0000) -- (0.0000,8.0000) -- (12.0000,8.0000) -- (0.0000,8.0000) -- (12.0000,8.0000) -- (0.0000,8.0000) -- (12.0000,8.0000) -- (0.0000,8.0000) -- (12.0000,8.0000) -- (0.0000,8.0000) -- (12.0000,8.0000) -- (0.0000,8.0000) -- (12.0000,8.0000) -- (0.0000,8.0000) -- (12.0000,8.0000) -- (0.0000,8.0000) -- (12.0000,8.0000) -- (0.0000,8.0000) -- (12.0000,8.0000) -- (0.0000,8.0000) -- (12.0000,8.0000) -- (0.0000,8.0000) -- (12.0000,8.0000) -- (0.0000,8.0000) -- (12.0000,8.0000) -- (0.0000,8.0000) -- (12.0000,8.0000) -- (0.0000,8.0000) -- (12.0000,8.0000) -- (0.0000,8.0000) -- (12.0000,8.0000) -- (0.0000,8.0000) -- (12.0000,8.0000) -- (0.0000,8.0000) -- (12.0000,8.0000) -- (0.0000,8.0000) -- (12.0000,8.0000) -- (0.0000,8.0000) -- (12.0000,8.0000) -- (0.0000,8.0000) -- (12.0000,8.0000) -- (0.0000,8.0000) -- (12.0000,8.0000) -- (0.0000,8.0000) -- (12.0000,8.0000) -- (0.0000,8.0000) -- (12.0000,8.0000) -- (0.0000,8.0000) -- (12.0000,8.0000) -- (0.0000,8.0000) -- (12.0000,8.0000) -- (0.0000,8.0000) -- (12.0000,8.0000) -- (0.0000,8.0000) -- (12.0000,8.0000) -- (0.0000,8.0000) -- (12.0000,8.0000) -- (0.0000,8.0000) -- (12.0000,8.0000) -- (0.0000,8.0000) -- (12.0000,8.0000) -- (0.0000,8.0000) -- (12.0000,8.0000) -- (0.0000,8.0000) -- (12.0000,8.0000) -- (0.0000,8.0000) -- (12.0000,8.0000) -- (0.0000,8.0000) -- (12.0000,8.0000) -- (0.0000,8.0000) -- (12.0000,8.0000) -- (0.0000,8.0000) -- (12.0000,8.0000) -- (0.0000,8.0000) -- (12.0000,8.0000) -- (0.0000,8.0000) -- (12.0000,8.0000) -- (0.0000,8.0000) -- (12.0000,8.0000) -- (0.0000,8.0000) -- (12.0000,8.0000) -- (0.0000,8.0000) -- (12.0000,8.0000) -- (0.0000,8.0000) -- (12.0000,8.0000) -- (0.0000,8.0000) -- (12.0000,8.0000) -- (0.0000,8.0000) -- (12.0000,8.0000) -- (0.0000,8.0000) -- (12.0000,8.0000) -- (0.0000,8.0000) -- (12.0000,8.0000) -- (0.0000,8.0000) -- (12.0000,8.0000) -- (0.0000,8.0000) -- (12.0000,8.0000) -- (0.0000,8.0000) -- (12.0000,8.0000) -- (0.0000,8.0000) -- (12.0000,8.0000) -- (0.0000,8.0000) -- (12.0000,8.0000) -- (0.0000,8.0000) -- (12.0000,8.0000) -- (0.0000,8.0000) -- (12.0000,8.0000) -- (0.0000,8.0000) -- (12.0000,8.0000) -- (0.0000,8.0000) -- (12.0000,8.0000) -- (0.0000,8.0000) -- (12.0000,8.0000) -- (0.0000,8.0000) -- (12.0000,8.0000) -- (0.0000,8.0000) -- (12.0000,8.0000) -- (0.0000,8.0000) -- (12.0000,8.0000) -- (0.0000,8.0000) -- (12.0000,8.0000) -- (0.0000,8.0000) -- (12.0000,8.0000) -- (0.0000,8.0000) -- (12.0000,8.0000) -- (0.0000,8.0000) -- (12.0000,8.0000) -- (0.0000,8.0000) -- (12.0000,8.0000) -- (0.0000,8.0000) -- (12.0000,8.0000) -- (0.0000,8.0000) -- (12.0000,8.0000) -- (0.0000,8.0000) -- (12.0000,8.0000) -- (0.0000,8.0000) -- (12.0000,8.0000) -- (0.0000,8.0000) -- (12.0000,8.0000) -- (0.0000,8.0000) -- (12.0000,8.0000) -- (0.0000,8.0000) -- (12.0000,8.0000) -- (0.0000,8.0000) -- (12.0000,8.0000) -- (0.0000,8.0000) -- (12.0000,8.0000) -- (0.0000,8.0000) -- (12.0000,8.0000) -- (0.0000,8.0000) -- (12.0000,8.0000) -- (0.0000,8.0000) -- (12.0000,8.0000) -- (0.0000,8.0000) -- (12.0000,8.0000) -- (0.0000,8.0000) -- (12.0000,8.0000) -- (0.0000,8.0000) -- (12.0000,8.0000) -- (0.0000,8.0000) -- (12.0000,8.0000) -- (0.0000,8.0000) -- (12.0000,8.0000) -- (0.0000,8.0000) -- (12.0000,8.0000) -- (0.0000,8.0000) -- (12.0000,8.0000) -- (0.0000,8.0000) -- (12.0000,8.0000) -- (0.0000,8.0000) -- (12.0000,8.0000) -- (0.0000,8.0000) -- (12.0000,8.0000) -- (0.0000,8.0000) -- (12.0000,8.0000) -- (0.0000,8.0000) -- (12.0000,8.0000) -- (0.0000,8.0000) -- (12.0000,8.0000) -- (0.0000,8.0000) -- (12.0000,8.0000) -- (0.0000,8.0000) -- (12.0000,8.0000) -- (0.0000,8.0000) -- (12.0000,8.0000) -- (0.0000,8.0000) -- (12.0000,8.0000) -- (0.0000,8.0000) -- (12.0000,8.0000) -- (0.0000,8.0000) -- (12.0000,8.0000) -- (0.0000,8.0000) -- (12.0000,8.0000) -- (0.0000,8.0000) -- (12.0000,8.0000) -- (0.0000,8.0000) -- (12.0000,8.0000) -- (0.0000,8.0000) -- (12.0000,8.0000) -- (0.0000,8.0000) -- (12.0000,8.0000) -- (0.0000,8.0000) -- (12.0000,8.0000) -- (0.0000,8.0000) -- (12.0000,8.0000) -- (0.0000,8.0000) -- (12.0000,8.0000) -- (0.0000,8.0000) -- (12.0000,8.0000) -- (0.0000,8.0000) -- (12.0000,8.0000) -- (0.0000,8.0000) -- (12.0000,8.0000) -- (0.0000,8.0000) -- (12.0000,8.0000) -- (0.0000,8.0000) -- (12.0000,8.0000) -- (0.0000,8.0000) -- (12.0000,8.0000) -- (0.0000,8.0000) -- (12.0000,8.0000) -- (0.0000,8.0000) -- (12.0000,8.0000) -- (0.0000,8.0000) -- (12.0000,8.0000) -- (0.0000,8.0000) -- (12.0000,8.0000) -- (0.0000,8.0000) -- (12.0000,8.0000) -- (0.0000,8.0000) -- (12.0000,8.0000) -- (0.0000,8.0000) -- (12.0000,8.0000) -- (0.0000,8.0000) -- (12.0000,8.0000) -- (0.0000,8.0000) -- (12.0000,8.0000) -- (0.0000,8.0000) -- (12.0000,8.0000) -- (0.0000,8.0000) -- (12.0000,8.0000) -- (0.0000,8.0000) -- (12.0000,8.0000) -- (0.0000,8.0000) -- (12.0000,8.0000) -- (0.0000,8.0000) -- (12.0000,8.0000) -- (0.0000,8.0000) -- (12.0000,8.0000) -- (0.0000,8.0000) -- (12.0000,8.0000) -- (0.0000,8.0000) -- (12.0000,8.0000) -- (0.0000,8.0000) -- (12.0000,8.0000) -- (0.0000,8.0000) -- (12.0000,8.0000) -- (0.0000,8.0000) -- (12.0000,8.0000) -- (0.0000,8.0000) -- (12.0000,8.0000) -- (0.0000,8.0000) -- (12.0000,8.0000) -- (0.0000,8.0000) -- (12.0000,8.0000) -- (0.0000,8.0000) -- (12.0000,8.0000) -- (0.0000,8.0000) -- (12.0000,8.0000) -- (0.0000,8.0000) -- (12.0000,8.0000) -- (0.0000,8.0000) -- (12.0000,8.0000) -- (0.0000,8.0000) -- (12.0000,8.0000) -- (0.0000,8.0000) -- (12.0000,8.0000) -- (0.0000,8.0000) -- (12.0000,8.0000) -- (0.0000,8.0000) -- (12.0000,8.0000) -- (0.0000,8.0000) -- (12.0000,8.0000) -- (0.0000,8.0000) -- (12.0000,8.0000) -- (0.0000,8.0000) -- (12.0000,8.0000) -- (0.0000,8.0000) -- (12.0000,8.0000) -- (0.0000,8.0000) -- (12.0000,8.0000) -- (0.0000,8.0000) -- (12.0000,8.0000) -- (0.0000,8.0000) -- (12.0000,8.0000) -- (0.0000,8.0000) -- (12.0000,8.0000) -- (0.0000,8.0000) -- (12.0000,8.0000) -- (0.0000,8.0000) -- (12.0000,8.0000) -- (0.0000,8.0000) -- (12.0000,8.0000) -- (0.0000,8.0000) -- (12.0000,8.0000) -- (0.0000,8.0000) -- (12.0000,8.0000) -- (0.0000,8.0000) -- (12.0000,8.0000) -- (0.0000,8.0000) -- (12.0000,8.0000) -- (0.0000,8.0000) -- (12.0000,8.0000) -- (0.0000,8.0000) -- (12.0000,8.0000) -- (0.0000,8.0000) -- (12.0000,8.0000) -- (0.0000,8.0000) -- (12.0000,8.0000) -- (0.0000,8.0000) -- (12.0000,8.0000) -- (0.0000,8.0000) -- (12.0000,8.0000) -- (0.0000,8.0000) -- (12.0000,8.0000) -- (0.0000,8.0000) -- (12.0000,8.0000) -- (0.0000,8.0000) -- (12.0000,8.0000) -- (0.0000,8.0000) -- (12.0000,8.0000) -- (0.0000,8.0000) -- (12.0000,8.0000) -- (0.0000,8.0000) -- (12.0000,8.0000) -- (0.0000,8.0000) -- (12.0000,8.0000) -- (0.0000,8.0000) -- (12.0000,8.0000) -- (0.0000,8.0000) -- (12.0000,8.0000) -- (0.0000,8.0000) -- (12.0000,8.0000) -- (0.0000,8.0000) -- (12.0000,8.0000) -- (0.0000,8.0000) -- (12.0000,8.0000) -- (0.0000,8.0000) -- (12.0000,8.0000) -- (0.0000,8.0000) -- (12.0000,8.0000) -- (0.0000,8.0000) -- (12.0000,8.0000) -- (0.0000,8.0000) -- (12.0000,8.0000) -- (0.0000,8.0000) -- (12.0000,8.0000) -- (0.0000,8.0000) -- (12.0000,8.0000) -- (0.0000,8.0000) -- (12.0000,8.0000) -- (0.0000,8.0000) -- (12.0000,8.0000) -- (0.0000,8.0000) -- (12.0000,8.0000) -- (0.0000,8.0000) -- (12.0000,8.0000) -- (0.0000,8.0000) -- (12.0000,8.0000) -- (0.0000,8.0000) -- (12.0000,8.0000) -- (0.0000,8.0000) -- (12.0000,8.0000) -- (0.0000,8.0000) -- (12.0000,8.0000) -- (0.0000,8.0000) -- (12.0000,8.0000) -- (0.0000,8.0000) -- (12.0000,8.0000) -- (0.0000,8.0000) -- (12.0000,8.0000) -- (0.0000,8.0000) -- (12.0000,8.0000) -- (0.0000,8.0000) -- (12.0000,8.0000) -- (0.0000,8.0000) -- (12.0000,8.0000) -- (0.0000,8.0000) -- (12.0000,8.0000) -- (0.0000,8.0000) -- (12.0000,8.0000) -- (0.0000,8.0000) -- (12.0000,8.0000) -- (0.0000,8.0000) -- (12.0000,8.0000) -- (0.0000,8.0000) -- (12.0000,8.0000) -- (0.0000,8.0000) -- (12.0000,8.0000) -- (0.0000,8.0000) -- (12.0000,8.0000) -- (0.0000,8.0000) -- (12.0000,8.0000) -- (0.0000,8.0000) -- (12.0000,8.0000) -- (0.0000,8.0000) -- (12.0000,8.0000) -- (0.0000,8.0000) -- (12.0000,8.0000) -- (0.0000,8.0000) -- (12.0000,8.0000) -- (0.0000,8.0000) -- (12.0000,8.0000) -- (0.0000,8.0000) -- (12.0000,8.0000) -- (0.0000,8.0000) -- (12.0000,8.0000) -- (0.0000,8.0000) -- (12.0000,8.0000) -- (0.0000,8.0000) -- (12.0000,8.0000) -- (0.0000,8.0000) -- (12.0000,8.0000) -- (0.0000,8.0000) -- (12.0000,8.0000) -- (0.0000,8.0000) -- (12.0000,8.0000) -- (0.0000,8.0000) -- (12.0000,8.0000) -- (0.0000,8.0000) -- (12.0000,8.0000) -- (0.0000,8.0000) -- (12.0000,8.0000) -- (0.0000,8.0000) -- (12.0000,8.0000) -- (0.0000,8.0000) -- (12.0000,8.0000) -- (0.0000,8.0000) -- (12.0000,8.0000) -- (0.0000,8.0000) -- (12.0000,8.0000) -- (0.0000,8.0000) -- (12.0000,8.0000) -- (0.0000,8.0000) -- (12.0000,8.0000) -- (0.0000,8.0000) -- (12.0000,8.0000) -- (0.0000,8.0000) -- (12.0000,8.0000) -- (0.0000,8.0000) -- (12.0000,8.0000) -- (0.0000,8.0000) -- (12.0000,8.0000) -- (0.0000,8.0000) -- (12.0000,8.0000) -- (0.0000,8.0000) -- (12.0000,8.0000) -- (0.0000,8.0000) -- (12.0000,8.0000) -- (0.0000,8.0000) -- (12.0000,8.0000) -- (0.0000,8.0000) -- (12.0000,8.0000) -- (0.0000,8.0000) -- (12.0000,8.0000) -- (0.0000,8.0000) -- (12.0000,8.0000) -- (0.0000,8.0000) -- (12.0000,8.0000) -- (0.0000,8.0000) -- (12.0000,8.0000) -- (0.0000,8.0000) -- (12.0000,8.0000) -- (0.0000,8.0000) -- (12.0000,8.0000) -- (0.0000,8.0000) -- (12.0000,8.0000) -- (0.0000,8.0000) -- (12.0000,8.0000) -- (0.0000,8.0000) -- (12.0000,8.0000) -- (0.0000,8.0000) -- (12.0000,8.0000) -- (0.0000,8.0000) -- (12.0000,8.0000) -- (0.0000,8.0000) -- (12.0000,8.0000) -- (0.0000,8.0000) -- (12.0000,8.0000) -- (0.0000,8.0000) -- (12.0000,8.0000) -- (0.0000,8.0000) -- (12.0000,8.0000) -- (0.0000,8.0000) -- (12.0000,8.0000) -- (0.0000,8.0000) -- (12.0000,8.0000) -- (0.0000,8.0000) -- (12.0000,8.0000) -- (0.0000,8.0000) -- (12.0000,8.0000) -- (0.0000,8.0000) -- (12.0000,8.0000) -- (0.0000,8.0000) -- (12.0000,8.0000) -- (0.0000,8.0000) -- (12.0000,8.0000) -- (0.0000,8.0000) -- (12.0000,8.0000) -- (0.0000,8.0000) -- (12.0000,8.0000) -- (0.0000,8.0000) -- (12.0000,8.0000) -- (0.0000,8.0000) -- (12.0000,8.0000) -- (0.0000,8.0000) -- (12.0000,8.0000) -- (0.0000,8.0000) -- (12.0000,8.0000) -- (0.0000,8.0000) -- (12.0000,8.0000) -- (0.0000,8.0000) -- (12.0000,8.0000) -- (0.0000,8.0000) -- (12.0000,8.0000) -- (0.0000,8.0000) -- (12.0000,8.0000) -- (0.0000,8.0000) -- (12.0000,8.0000) -- (0.0000,8.0000) -- (12.0000,8.0000) -- (0.0000,8.0000) -- (12.0000,8.0000) -- (0.0000,8.0000) -- (12.0000,8.0000) -- (0.0000,8.0000) -- (12.0000,8.0000) -- (0.0000,8.0000) -- (12.0000,8.0000) -- (0.0000,8.0000) -- (12.0000,8.0000) -- (0.0000,8.0000) -- (12.0000,8.0000) -- (0.0000,8.0000) -- (12.0000,8.0000) -- (0.0000,8.0000) -- (12.0000,8.0000) -- (0.0000,8.0000) -- (12.0000,8.0000) -- (0.0000,8.0000) -- (12.0000,8.0000) -- (0.0000,8.0000) -- (12.0000,8.0000) -- (0.0000,8.0000) -- (12.0000,8.0000) -- (0.0000,8.0000) -- (12.0000,8.0000) -- (0.0000,8.0000) -- (12.0000,8.0000) -- (0.0000,8.0000) -- (12.0000,8.0000) -- (0.0000,8.0000) -- (12.0000,8.0000) -- (0.0000,8.0000) -- (12.0000,8.0000) -- (0.0000,8.0000) -- (12.0000,8.0000) -- (0.0000,8.0000) -- (12.0000,8.0000) -- (0.0000,8.0000) -- (12.0000,8.0000) -- (0.0000,8.0000) -- (12.0000,8.0000) -- (0.0000,8.0000) -- (12.0000,8.0000) -- (0.0000,8.0000) -- (12.0000,8.0000) -- (0.0000,8.0000) -- (12.0000,8.0000) -- (0.0000,8.0000) -- (12.0000,8.0000) -- (0.0000,8.0000) -- (12.0000,8.0000) -- (0.0000,8.0000) -- (12.0000,8.0000) -- (0.0000,8.0000) -- (12.0000,8.0000) -- (0.0000,8.0000) -- (12.0000,8.0000) -- (0.0000,8.0000) -- (12.0000,8.0000) -- (0.0000,8.0000) -- (12.0000,8.0000) -- (0.0000,8.0000) -- (12.0000,8.0000) -- (0.0000,8.0000) -- (12.0000,8.0000) -- (0.0000,8.0000) -- (12.0000,8.0000) -- (0.0000,8.0000) -- (12.0000,8.0000) -- (0.0000,8.0000) -- (12.0000,8.0000) -- (0.0000,8.0000) -- (12.0000,8.0000) -- (0.0000,8.0000) -- (12.0000,8.0000) -- (0.0000,8.0000) -- (12.0000,8.0000) -- (0.0000,8.0000) -- (12.0000,8.0000) -- (0.0000,8.0000) -- (12.0000,8.0000) -- (0.0000,8.0000) -- (12.0000,8.0000) -- (0.0000,8.0000) -- (12.0000,8.0000) -- (0.0000,8.0000) -- (12.0000,8.0000) -- (0.0000,8.0000) -- (12.0000,8.0000) -- (0.0000,8.0000) -- (12.0000,8.0000) -- (0.0000,8.0000) -- (12.0000,8.0000) -- (0.0000,8.0000) -- (12.0000,8.0000) -- (0.0000,8.0000) -- (12.0000,8.0000) -- (0.0000,8.0000) -- (12.0000,8.0000) -- (0.0000,8.0000) -- (12.0000,8.0000) -- (0.0000,8.0000) -- (12.0000,8.0000) -- (0.0000,8.0000) -- (12.0000,8.0000) -- (0.0000,8.0000) -- (12.0000,8.0000) -- (0.0000,8.0000) -- (12.0000,8.0000) -- (0.0000,8.0000) -- (12.0000,8.0000) -- (0.0000,8.0000) -- (12.0000,8.0000) -- (0.0000,8.0000) -- (12.0000,8.0000) -- (0.0000,8.0000) -- (12.0000,8.0000) -- (0.0000,8.0000) -- (12.0000,8.0000) -- (0.0000,8.0000) -- (12.0000,8.0000) -- (0.0000,8.0000) -- (12.0000,8.0000) -- (0.0000,8.0000) -- (12.0000,8.0000) -- (0.0000,8.0000) -- (12.0000,8.0000) -- (0.0000,8.0000) -- (12.0000,8.0000) -- (0.0000,8.0000) -- (12.0000,8.0000) -- (0.0000,8.0000) -- (12.0000,8.0000) -- (0.0000,8.0000) -- (12.0000,8.0000) -- (0.0000,8.0000) -- (12.0000,8.0000) -- (0.0000,8.0000) -- (12.0000,8.0000) -- (0.0000,8.0000) -- (12.0000,8.0000) -- (0.0000,8.0000) -- (12.0000,8.0000) -- (0.0000,8.0000) -- (12.0000,8.0000) -- (0.0000,8.0000) -- (12.0000,8.0000) -- (0.0000,8.0000) -- (12.0000,8.0000) -- (0.0000,8.0000) -- (12.0000,8.0000) -- (0.0000,8.0000) -- (12.0000,8.0000) -- (0.0000,8.0000) -- (12.0000,8.0000) -- (0.0000,8.0000) -- (12.0000,8.0000) -- (0.0000,8.0000) -- (12.0000,8.0000) -- (0.0000,8.0000) -- (12.0000,8.0000) -- (0.0000,8.0000) -- (12.0000,8.0000) -- (0.0000,8.0000) -- (12.0000,8.0000) -- (0.0000,8.0000) -- (12.0000,8.0000) -- (0.0000,8.0000) -- (12.0000,8.0000) -- (0.0000,8.0000) -- (12.0000,8.0000) -- (0.0000,8.0000) -- (12.0000,8.0000) -- (0.0000,8.0000) -- (12.0000,8.0000) -- (0.0000,8.0000) -- (12.0000,8.0000) -- (0.0000,8.0000) -- (12.0000,8.0000) -- (0.0000,8.0000) -- (12.0000,8.0000) -- (0.0000,8.0000) -- (12.0000,8.0000) -- (0.0000,8.0000) -- (12.0000,8.0000) -- (0.0000,8.0000) -- (12.0000,8.0000) -- (0.0000,8.0000) -- (12.0000,8.0000) -- (0.0000,8.0000) -- (12.0000,8.0000) -- (0.0000,8.0000) -- (12.0000,8.0000) -- (0.0000,8.0000) -- (12.0000,8.0000) -- (0.0000,8.0000) -- (12.0000,8.0000) -- (0.0000,8.0000) -- (12.0000,8.0000) -- (0.0000,8.0000) -- (12.0000,8.0000) -- (0.0000,8.0000) -- (12.0000,8.0000) -- (0.0000,8.0000) -- (12.0000,8.0000) -- (0.0000,8.0000) -- (12.0000,8.0000) -- (0.0000,8.0000) -- (12.0000,8.0000) -- (0.0000,8.0000) -- (12.0000,8.0000) -- (0.0000,8.0000) -- (12.0000,8.0000) -- (0.0000,8.0000) -- (12.0000,8.0000) -- (0.0000,8.0000) -- (12.0000,8.0000) -- (0.0000,8.0000) -- (12.0000,8.0000) -- (0.0000,8.0000) -- (12.0000,8.0000) -- (0.0000,8.0000) -- (12.0000,8.0000) -- (0.0000,8.0000) -- (12.0000,8.0000) -- (0.0000,8.0000) -- (12.0000,8.0000) -- (0.0000,8.0000) -- (12.0000,8.0000) -- (0.0000,8.0000) -- (12.0000,8.0000) -- (0.0000,8.0000) -- (12.0000,8.0000) -- (0.0000,8.0000) -- (12.0000,8.0000) -- (0.0000,8.0000) -- (12.0000,8.0000) -- (0.0000,8.0000) -- (12.0000,8.0000) -- (0.0000,8.0000) -- (12.0000,8.0000) -- (0.0000,8.0000) -- (12.0000,8.0000) -- (0.0000,8.0000) -- (12.0000,8.0000) -- (0.0000,8.0000) -- (12.0000,8.0000) -- (0.0000,8.0000) -- (12.0000,8.0000) -- (0.0000,8.0000) -- (12.0000,8.0000) -- (0.0000,8.0000) -- (12.0000,8.0000) -- (0.0000,8.0000) -- (12.0000,8.0000) -- (0.0000,8.0000) -- (12.0000,8.0000) -- (0.0000,8.0000) -- (12.0000,8.0000) -- (0.0000,8.0000) -- (12.0000,8.0000) -- (0.0000,8.0000) -- (12.0000,8.0000) -- (0.0000,8.0000) -- (12.0000,8.0000) -- (0.0000,8.0000) -- (12.0000,8.0000) -- (0.0000,8.0000) -- (12.0000,8.0000) -- (0.0000,8.0000) -- (12.0000,8.0000) -- (0.0000,8.0000) -- (12.0000,8.0000) -- (0.0000,8.0000) -- (12.0000,8.0000) -- (0.0000,8.0000) -- (12.0000,8.0000) -- (0.0000,8.0000) -- (12.0000,8.0000) -- (0.0000,8.0000) -- (12.0000,8.0000) -- (0.0000,8.0000) -- (12.0000,8.0000) -- (0.0000,8.0000) -- (12.0000,8.0000) -- (0.0000,8.0000) -- (12.0000,8.0000) -- (0.0000,8.0000) -- (12.0000,8.0000) -- (0.0000,8.0000) -- (12.0000,8.0000) -- (0.0000,8.0000) -- (12.0000,8.0000) -- (0.0000,8.0000) -- (12.0000,8.0000) -- (0.0000,8.0000) -- (12.0000,8.0000) -- (0.0000,8.0000) -- (12.0000,8.0000) -- (0.0000,8.0000) -- (12.0000,8.0000) -- (0.0000,8.0000) -- (12.0000,8.0000) -- (0.0000,8.0000) -- (12.0000,8.0000) -- (0.0000,8.0000) -- (12.0000,8.0000) -- (0.0000,8.0000) -- (12.0000,8.0000) -- (0.0000,8.0000) -- (12.0000,8.0000) -- (0.0000,8.0000) -- (12.0000,8.0000) -- (0.0000,8.0000) -- (12.0000,8.0000) -- (0.0000,8.0000) -- (12.0000,8.0000) -- (0.0000,8.0000) -- (12.0000,8.0000) -- (0.0000,8.0000) -- (12.0000,8.0000) -- (0.0000,8.0000) -- (12.0000,8.0000) -- (0.0000,8.0000) -- (12.0000,8.0000) -- (0.0000,8.0000) -- (12.0000,8.0000) -- (0.0000,8.0000) -- (12.0000,8.0000) -- (0.0000,8.0000) -- (12.0000,8.0000) -- (0.0000,8.0000) -- (12.0000,8.0000) -- (0.0000,8.0000) -- (12.0000,8.0000) -- (0.0000,8.0000) -- (12.0000,8.0000) -- (0.0000,8.0000) -- (12.0000,8.0000) -- (0.0000,8.0000) -- (12.0000,8.0000) -- (0.0000,8.0000) -- (12.0000,8.0000) -- (0.0000,8.0000) -- (12.0000,8.0000) -- (0.0000,8.0000) -- (12.0000,8.0000) -- (0.0000,8.0000) -- (12.0000,8.0000) -- (0.0000,8.0000) -- (12.0000,8.0000) -- (0.0000,8.0000) -- (12.0000,8.0000) -- (0.0000,8.0000) -- (12.0000,8.0000) -- (0.0000,8.0000) -- (12.0000,8.0000) -- (0.0000,8.0000) -- (12.0000,8.0000) -- (0.0000,8.0000) -- (12.0000,8.0000) -- (0.0000,8.0000) -- (12.0000,8.0000) -- (0.0000,8.0000) -- (12.0000,8.0000) -- (0.0000,8.0000) -- (12.0000,8.0000) -- (0.0000,8.0000) -- (12.0000,8.0000) -- (0.0000,8.0000) -- (12.0000,8.0000) -- (0.0000,8.0000) -- (12.0000,8.0000) -- (0.0000,8.0000) -- (12.0000,8.0000) -- (0.0000,8.0000) -- (12.0000,8.0000) -- (0.0000,8.0000) -- (12.0000,8.0000) -- (0.0000,8.0000) -- (12.0000,8.0000) -- (0.0000,8.0000) -- (12.0000,8.0000) -- (0.0000,8.0000) -- (12.0000,8.0000) -- (0.0000,8.0000) -- (12.0000,8.0000) -- (0.0000,8.0000) -- (12.0000,8.0000) -- (0.0000,8.0000) -- (12.0000,8.0000) -- (0.0000,8.0000) -- (12.0000,8.0000) -- (0.0000,8.0000) -- (12.0000,8.0000) -- (0.0000,8.0000) -- (12.0000,8.0000) -- (0.0000,8.0000) -- (12.0000,8.0000) -- (0.0000,8.0000) -- (12.0000,8.0000) -- (0.0000,8.0000) -- (12.0000,8.0000) -- (0.0000,8.0000) -- (12.0000,8.0000) -- (0.0000,8.0000) -- (12.0000,8.0000) -- (0.0000,8.0000) -- (12.0000,8.0000) -- (0.0000,8.0000) -- (12.0000,8.0000) -- (0.0000,8.0000) -- (12.0000,8.0000) -- (0.0000,8.0000) -- (12.0000,8.0000) -- (0.0000,8.0000) -- (12.0000,8.0000) -- (0.0000,8.0000) -- (12.0000,8.0000) -- (0.0000,8.0000) -- (12.0000,8.0000) -- (0.0000,8.0000) -- (12.0000,8.0000) -- (0.0000,8.0000) -- (12.0000,8.0000) -- (0.0000,8.0000) -- (12.0000,8.0000) -- (0.0000,8.0000) -- (12.0000,8.0000) -- (0.0000,8.0000) -- (12.0000,8.0000) -- (0.0000,8.0000) -- (12.0000,8.0000) -- (0.0000,8.0000) -- (12.0000,8.0000) -- (0.0000,8.0000) -- (12.0000,8.0000) -- (0.0000,8.0000) -- (12.0000,8.0000) -- (0.0000,8.0000) -- (12.0000,8.0000) -- (0.0000,8.0000) -- (12.0000,8.0000) -- (0.0000,8.0000) -- (12.0000,8.0000) -- (0.0000,8.0000) -- (12.0000,8.0000) -- (0.0000,8.0000) -- (12.0000,8.0000) -- (0.0000,8.0000) -- (12.0000,8.0000) -- (0.0000,8.0000) -- (12.0000,8.0000) -- (0.0000,8.0000) -- (12.0000,8.0000) -- (0.0000,8.0000) -- (12.0000,8.0000) -- (0.0000,8.0000) -- (12.0000,8.0000) -- (0.0000,8.0000) -- (12.0000,8.0000) -- (0.0000,8.0000) -- (12.0000,8.0000) -- (0.0000,8.0000) -- (12.0000,8.0000) -- (0.0000,8.0000) -- (12.0000,8.0000) -- (0.0000,8.0000) -- (12.0000,8.0000) -- (0.0000,8.0000) -- (12.0000,8.0000) -- (0.0000,8.0000) -- (12.0000,8.0000) -- (0.0000,8.0000) -- (12.0000,8.0000) -- (0.0000,8.0000) -- (12.0000,8.0000) -- (0.0000,8.0000) -- (12.0000,8.0000) -- (0.0000,8.0000) -- (12.0000,8.0000) -- (0.0000,8.0000) -- (12.0000,8.0000) -- (0.0000,8.0000) -- (12.0000,8.0000) -- (0.0000,8.0000) -- (12.0000,8.0000) -- (0.0000,8.0000) -- (12.0000,8.0000) -- (0.0000,8.0000) -- (12.0000,8.0000) -- (0.0000,8.0000) -- (12.0000,8.0000) -- (0.0000,8.0000) -- (12.0000,8.0000) -- (0.0000,8.0000) -- (12.0000,8.0000) -- (0.0000,8.0000) -- (12.0000,8.0000) -- (0.0000,8.0000) -- (12.0000,8.0000) -- (0.0000,8.0000) -- (12.0000,8.0000) -- (0.0000,8.0000) -- (12.0000,8.0000) -- (0.0000,8.0000) -- (12.0000,8.0000) -- (0.0000,8.0000) -- (12.0000,8.0000) -- (0.0000,8.0000) -- (12.0000,8.0000) -- (0.0000,8.0000) -- (12.0000,8.0000) -- (0.0000,8.0000) -- (12.0000,8.0000) -- (0.0000,8.0000) -- (12.0000,8.0000) -- (0.0000,8.0000) -- (12.0000,8.0000) -- (0.0000,8.0000) -- (12.0000,8.0000) -- (0.0000,8.0000) -- (12.0000,8.0000) -- (0.0000,8.0000) -- (12.0000,8.0000) -- (0.0000,8.0000) -- (12.0000,8.0000) -- (0.0000,8.0000) -- (12.0000,8.0000) -- (0.0000,8.0000) -- (12.0000,8.0000) -- (0.0000,8.0000) -- (12.0000,8.0000) -- (0.0000,8.0000) -- (12.0000,8.0000) -- (0.0000,8.0000) -- (12.0000,8.0000) -- (0.0000,8.0000) -- (12.0000,8.0000) -- (0.0000,8.0000) -- (12.0000,8.0000) -- (0.0000,8.0000) -- (12.0000,8.0000) -- (0.0000,8.0000) -- (12.0000,8.0000) -- (0.0000,8.0000) -- (12.0000,8.0000) -- (0.0000,8.0000) -- (12.0000,8.0000) -- (0.0000,8.0000) -- (12.0000,8.0000) -- (0.0000,8.0000) -- (12.0000,8.0000) -- (0.0000,8.0000) -- (12.0000,8.0000) -- (0.0000,8.0000) -- (12.0000,8.0000) -- (0.0000,8.0000) -- (12.0000,8.0000) -- (0.0000,8.0000) -- (12.0000,8.0000) -- (0.0000,8.0000) -- (12.0000,8.0000) -- (0.0000,8.0000) -- (12.0000,8.0000) -- (0.0000,8.0000) -- (12.0000,8.0000) -- (0.0000,8.0000) -- (12.0000,8.0000) -- (0.0000,8.0000) -- (12.0000,8.0000) -- (0.0000,8.0000) -- (12.0000,8.0000) -- (0.0000,8.0000) -- (12.0000,8.0000) -- (0.0000,8.0000) -- (12.0000,8.0000) -- (0.0000,8.0000) -- (12.0000,8.0000) -- (0.0000,8.0000) -- (12.0000,8.0000) -- (0.0000,8.0000) -- (12.0000,8.0000) -- (0.0000,8.0000) -- (12.0000,8.0000) -- (0.0000,8.0000) -- (12.0000,8.0000) -- (0.0000,8.0000) -- (12.0000,8.0000) -- (0.0000,8.0000) -- (12.0000,8.0000) -- (0.0000,8.0000) -- (12.0000,8.0000) -- (0.0000,8.0000) -- (12.0000,8.0000) -- (0.0000,8.0000) -- (12.0000,8.0000) -- (0.0000,8.0000) -- (12.0000,8.0000) -- (0.0000,8.0000) -- (12.0000,8.0000) -- (0.0000,8.0000) -- (12.0000,8.0000) -- (0.0000,8.0000) -- (12.0000,8.0000) -- (0.0000,8.0000) -- (12.0000,8.0000) -- (0.0000,8.0000) -- (12.0000,8.0000) -- (0.0000,8.0000) -- (12.0000,8.0000) -- (0.0000,8.0000) -- (12.0000,8.0000) -- (0.0000,8.0000) -- (12.0000,8.0000) -- (0.0000,8.0000) -- (12.0000,8.0000) -- (0.0000,8.0000) -- (12.0000,8.0000) -- (0.0000,8.0000) -- (12.0000,8.0000) -- (0.0000,8.0000) -- (12.0000,8.0000) -- (0.0000,8.0000) -- (12.0000,8.0000) -- (0.0000,8.0000) -- (12.0000,8.0000) -- (0.0000,8.0000) -- (12.0000,8.0000) -- (0.0000,8.0000) -- (12.0000,8.0000) -- (0.0000,8.0000) -- (12.0000,8.0000) -- (0.0000,8.0000) -- (12.0000,8.0000) -- (0.0000,8.0000) -- (12.0000,8.0000) -- (0.0000,8.0000) -- (12.0000,8.0000) -- (0.0000,8.0000) -- (12.0000,8.0000) -- (0.0000,8.0000) -- (12.0000,8.0000) -- (0.0000,8.0000) -- (12.0000,8.0000) -- (0.0000,8.0000) -- (12.0000,8.0000) -- (0.0000,8.0000) -- (12.0000,8.0000) -- (0.0000,8.0000) -- (12.0000,8.0000) -- (0.0000,8.0000) -- (12.0000,8.0000) -- (0.0000,8.0000) -- (12.0000,8.0000) -- (0.0000,8.0000) -- (12.0000,8.0000) -- (0.0000,8.0000) -- (12.0000,8.0000) -- (0.0000,8.0000) -- (12.0000,8.0000) -- (0.0000,8.0000) -- (12.0000,8.0000) -- (0.0000,8.0000) -- (12.0000,8.0000) -- (0.0000,8.0000) -- (12.0000,8.0000) -- (0.0000,8.0000) -- (12.0000,8.0000) -- (0.0000,8.0000) -- (12.0000,8.0000) -- (0.0000,8.0000) -- (12.0000,8.0000) -- (0.0000,8.0000) -- (12.0000,8.0000) -- (0.0000,8.0000) -- (12.0000,8.0000) -- (0.0000,8.0000) -- (12.0000,8.0000) -- (0.0000,8.0000) -- (12.0000,8.0000) -- (0.0000,8.0000) -- (12.0000,8.0000) -- (0.0000,8.0000) -- (12.0000,8.0000) -- (0.0000,8.0000) -- (12.0000,8.0000) -- (0.0000,8.0000) -- (12.0000,8.0000) -- (0.0000,8.0000) -- (12.0000,8.0000) -- (0.0000,8.0000) -- (12.0000,8.0000) -- (0.0000,8.0000) -- (12.0000,8.0000) -- (0.0000,8.0000) -- (12.0000,8.0000) -- (0.0000,8.0000) -- (12.0000,8.0000) -- (0.0000,8.0000) -- (12.0000,8.0000) -- (0.0000,8.0000) -- (12.0000,8.0000) -- (0.0000,8.0000) -- (12.0000,8.0000) -- (0.0000,8.0000) -- (12.0000,8.0000) -- (0.0000,8.0000) -- (12.0000,8.0000) -- (0.0000,8.0000) -- (12.0000,8.0000) -- (0.0000,8.0000) -- (12.0000,8.0000) -- (0.0000,8.0000) -- (12.0000,8.0000) -- (0.0000,8.0000) -- (12.0000,8.0000) -- (0.0000,8.0000) -- (12.0000,8.0000) -- (0.0000,8.0000) -- (12.0000,8.0000) -- (0.0000,8.0000) -- (12.0000,8.0000) -- (0.0000,8.0000) -- (12.0000,8.0000) -- (0.0000,8.0000) -- (12.0000,8.0000) -- (0.0000,8.0000) -- (12.0000,8.0000) -- (0.0000,8.0000) -- (12.0000,8.0000) -- (0.0000,8.0000) -- (12.0000,8.0000) -- (0.0000,8.0000) -- (12.0000,8.0000) -- (0.0000,8.0000) -- (12.0000,8.0000) -- (0.0000,8.0000) -- (12.0000,8.0000) -- (0.0000,8.0000) -- (12.0000,8.0000) -- (0.0000,8.0000) -- (12.0000,8.0000) -- (0.0000,8.0000) -- (12.0000,8.0000) -- (0.0000,8.0000) -- (12.0000,8.0000) -- (0.0000,8.0000) -- (12.0000,8.0000) -- (0.0000,8.0000) -- (12.0000,8.0000) -- (0.0000,8.0000) -- (12.0000,8.0000) -- (0.0000,8.0000) -- (12.0000,8.0000) -- (0.0000,8.0000) -- (12.0000,8.0000) -- (0.0000,8.0000) -- (12.0000,8.0000) -- (0.0000,8.0000) -- (12.0000,8.0000) -- (0.0000,8.0000) -- (12.0000,8.0000) -- (0.0000,8.0000) -- (12.0000,8.0000) -- (0.0000,8.0000) -- (12.0000,8.0000) -- (0.0000,8.0000) -- (12.0000,8.0000) -- (0.0000,8.0000) -- (12.0000,8.0000) -- (0.0000,8.0000) -- (12.0000,8.0000) -- (0.0000,8.0000) -- (12.0000,8.0000) -- (0.0000,8.0000) -- (12.0000,8.0000) -- (0.0000,8.0000) -- (12.0000,8.0000) -- (0.0000,8.0000) -- (12.0000,8.0000) -- (0.0000,8.0000) -- (12.0000,8.0000) -- (0.0000,8.0000) -- (12.0000,8.0000) -- (0.0000,8.0000) -- (12.0000,8.0000) -- (0.0000,8.0000) -- (12.0000,8.0000) -- (0.0000,8.0000) -- (12.0000,8.0000) -- (0.0000,8.0000) -- (12.0000,8.0000) -- (0.0000,8.0000) -- (12.0000,8.0000) -- (0.0000,8.0000) -- (12.0000,8.0000) -- (0.0000,8.0000) -- (12.0000,8.0000) -- (0.0000,8.0000) -- (12.0000,8.0000) -- (0.0000,8.0000) -- (12.0000,8.0000) -- (0.0000,8.0000) -- (12.0000,8.0000) -- (0.0000,8.0000) -- (12.0000,8.0000) -- (0.0000,8.0000) -- (12.0000,8.0000) -- (0.0000,8.0000) -- (12.0000,8.0000) -- (0.0000,8.0000) -- (12.0000,8.0000) -- (0.0000,8.0000) -- (12.0000,8.0000) -- (0.0000,8.0000) -- (12.0000,8.0000) -- (0.0000,8.0000) -- (12.0000,8.0000) -- (0.0000,8.0000) -- (12.0000,8.0000) -- (0.0000,8.0000) -- (12.0000,8.0000) -- (0.0000,8.0000) -- (12.0000,8.0000) -- (0.0000,8.0000) -- (12.0000,8.0000) -- (0.0000,8.0000) -- (12.0000,8.0000) -- (0.0000,8.0000) -- (12.0000,8.0000) -- (0.0000,8.0000) -- (12.0000,8.0000) -- (0.0000,8.0000) -- (12.0000,8.0000) -- (0.0000,8.0000) -- (12.0000,8.0000) -- (0.0000,8.0000) -- (12.0000,8.0000) -- (0.0000,8.0000) -- (12.0000,8.0000) -- (0.0000,8.0000) -- (12.0000,8.0000) -- (0.0000,8.0000) -- (12.0000,8.0000) -- (0.0000,8.0000) -- (12.0000,8.0000) -- (0.0000,8.0000) -- (12.0000,8.0000) -- (0.0000,8.0000) -- (12.0000,8.0000) -- (0.0000,8.0000) -- (12.0000,8.0000) -- (0.0000,8.0000) -- (12.0000,8.0000) -- (0.0000,8.0000) -- (12.0000,8.0000) -- (0.0000,8.0000) -- (12.0000,8.0000) -- (0.0000,8.0000) -- (12.0000,8.0000) -- (0.0000,8.0000) -- (12.0000,8.0000) -- (0.0000,8.0000) -- (12.0000,8.0000) -- (0.0000,8.0000) -- (12.0000,8.0000) -- (0.0000,8.0000) -- (12.0000,8.0000) -- (0.0000,8.0000) -- (12.0000,8.0000) -- (0.0000,8.0000) -- (12.0000,8.0000) -- (0.0000,8.0000) -- (12.0000,8.0000) -- (0.0000,8.0000) -- (12.0000,8.0000) -- (0.0000,8.0000) -- (12.0000,8.0000) -- (0.0000,8.0000) -- (12.0000,8.0000) -- (0.0000,8.0000) -- (12.0000,8.0000) -- (0.0000,8.0000) -- (12.0000,8.0000) -- (0.0000,8.0000) -- (12.0000,8.0000) -- (0.0000,8.0000) -- (12.0000,8.0000) -- (0.0000,8.0000) -- (12.0000,8.0000) -- (0.0000,8.0000) -- (12.0000,8.0000) -- (0.0000,8.0000) -- (12.0000,8.0000) -- (0.0000,8.0000) -- (12.0000,8.0000) -- (0.0000,8.0000) -- (12.0000,8.0000) -- (0.0000,8.0000) -- (12.0000,8.0000) -- (0.0000,8.0000) -- (12.0000,8.0000) -- (0.0000,8.0000) -- (12.0000,8.0000) -- (0.0000,8.0000) -- (12.0000,8.0000) -- (0.0000,8.0000) -- (12.0000,8.0000) -- (0.0000,8.0000) -- (12.0000,8.0000) -- (0.0000,8.0000) -- (12.0000,8.0000) -- (0.0000,8.0000) -- (12.0000,8.0000) -- (0.0000,8.0000) -- (12.0000,8.0000) -- (0.0000,8.0000) -- (12.0000,8.0000) -- (0.0000,8.0000) -- (12.0000,8.0000) -- (0.0000,8.0000) -- (12.0000,8.0000) -- (0.0000,8.0000) -- (12.0000,8.0000) -- (0.0000,8.0000) -- (12.0000,8.0000) -- (0.0000,8.0000) -- (12.0000,8.0000) -- (0.0000,8.0000) -- (12.0000,8.0000) -- (0.0000,8.0000) -- (12.0000,8.0000) -- (0.0000,8.0000) -- (12.0000,8.0000) -- (0.0000,8.0000) -- (12.0000,8.0000) -- (0.0000,8.0000) -- (12.0000,8.0000) -- (0.0000,8.0000) -- (12.0000,8.0000) -- (0.0000,8.0000) -- (12.0000,8.0000) -- (0.0000,8.0000) -- (12.0000,8.0000) -- (0.0000,8.0000) -- (12.0000,8.0000) -- (0.0000,8.0000) -- (12.0000,8.0000) -- (0.0000,8.0000) -- (12.0000,8.0000) -- (0.0000,8.0000) -- (12.0000,8.0000) -- (0.0000,8.0000) -- (12.0000,8.0000) -- (0.0000,8.0000) -- (12.0000,8.0000) -- (0.0000,8.0000) -- (12.0000,8.0000) -- (0.0000,8.0000) -- (12.0000,8.0000) -- (0.0000,8.0000) -- (12.0000,8.0000) -- (0.0000,8.0000) -- (12.0000,8.0000) -- (0.0000,8.0000) -- (12.0000,8.0000) -- (0.0000,8.0000) -- (12.0000,8.0000) -- (0.0000,8.0000) -- (12.0000,8.0000) -- (0.0000,8.0000) -- (12.0000,8.0000) -- (0.0000,8.0000) -- (12.0000,8.0000) -- (0.0000,8.0000) -- (12.0000,8.0000) -- (0.0000,8.0000) -- (12.0000,8.0000) -- (0.0000,8.0000) -- (12.0000,8.0000) -- (0.0000,8.0000) -- (12.0000,8.0000) -- (0.0000,8.0000) -- (12.0000,8.0000) -- (0.0000,8.0000) -- (12.0000,8.0000) -- (0.0000,8.0000) -- (12.0000,8.0000) -- (0.0000,8.0000) -- (12.0000,8.0000) -- (0.0000,8.0000) -- (12.0000,8.0000) -- (0.0000,8.0000) -- (12.0000,8.0000) -- (0.0000,8.0000) -- (12.0000,8.0000) -- (0.0000,8.0000) -- (12.0000,8.0000) -- (0.0000,8.0000) -- (12.0000,8.0000) -- (0.0000,8.0000) -- (12.0000,8.0000) -- (0.0000,8.0000) -- (12.0000,8.0000) -- (0.0000,8.0000) -- (12.0000,8.0000) -- (0.0000,8.0000) -- (12.0000,8.0000) -- (0.0000,8.0000) -- (12.0000,8.0000) -- (0.0000,8.0000) -- (12.0000,8.0000) -- (0.0000,8.0000) -- (12.0000,8.0000) -- (0.0000,8.0000) -- (12.0000,8.0000) -- (0.0000,8.0000) -- (12.0000,8.0000) -- (0.0000,8.0000) -- (12.0000,8.0000) -- (0.0000,8.0000) -- (12.0000,8.0000) -- (0.0000,8.0000) -- (12.0000,8.0000) -- (0.0000,8.0000) -- (12.0000,8.0000) -- (0.0000,8.0000) -- (12.0000,8.0000) -- (0.0000,8.0000) -- (12.0000,8.0000) -- (0.0000,8.0000) -- (12.0000,8.0000) -- (0.0000,8.0000) -- (12.0000,8.0000) -- (0.0000,8.0000) -- (12.0000,8.0000) -- (0.0000,8.0000);
\fill[red!80!black] (3.6000,2.4000) circle (1.2pt);
\fill[red!80!black] (12.0000,8.0000) circle (1.2pt);
\fill[red!80!black] (0.0000,8.0000) circle (1.2pt);
\fill[red!80!black] (12.0000,8.0000) circle (1.2pt);
\fill[red!80!black] (0.0000,8.0000) circle (1.2pt);
\fill[red!80!black] (12.0000,8.0000) circle (1.2pt);
\fill[red!80!black] (0.0000,8.0000) circle (1.2pt);
\fill[red!80!black] (12.0000,8.0000) circle (1.2pt);
\fill[red!80!black] (0.0000,8.0000) circle (1.2pt);
\fill[red!80!black] (12.0000,8.0000) circle (1.2pt);
\fill[red!80!black] (0.0000,8.0000) circle (1.2pt);
\fill[red!80!black] (12.0000,8.0000) circle (1.2pt);
\fill[red!80!black] (0.0000,8.0000) circle (1.2pt);
\fill[red!80!black] (12.0000,8.0000) circle (1.2pt);
\fill[red!80!black] (0.0000,8.0000) circle (1.2pt);
\fill[red!80!black] (12.0000,8.0000) circle (1.2pt);
\fill[red!80!black] (0.0000,8.0000) circle (1.2pt);
\fill[red!80!black] (12.0000,8.0000) circle (1.2pt);
\fill[red!80!black] (0.0000,8.0000) circle (1.2pt);
\fill[red!80!black] (12.0000,8.0000) circle (1.2pt);
\fill[red!80!black] (0.0000,8.0000) circle (1.2pt);
\fill[red!80!black] (12.0000,8.0000) circle (1.2pt);
\fill[red!80!black] (0.0000,8.0000) circle (1.2pt);
\fill[red!80!black] (12.0000,8.0000) circle (1.2pt);
\fill[red!80!black] (0.0000,8.0000) circle (1.2pt);
\fill[red!80!black] (12.0000,8.0000) circle (1.2pt);
\fill[red!80!black] (0.0000,8.0000) circle (1.2pt);
\fill[red!80!black] (12.0000,8.0000) circle (1.2pt);
\fill[red!80!black] (0.0000,8.0000) circle (1.2pt);
\fill[red!80!black] (12.0000,8.0000) circle (1.2pt);
\fill[red!80!black] (0.0000,8.0000) circle (1.2pt);
\fill[red!80!black] (12.0000,8.0000) circle (1.2pt);
\fill[red!80!black] (0.0000,8.0000) circle (1.2pt);
\fill[red!80!black] (12.0000,8.0000) circle (1.2pt);
\fill[red!80!black] (0.0000,8.0000) circle (1.2pt);
\fill[red!80!black] (12.0000,8.0000) circle (1.2pt);
\fill[red!80!black] (0.0000,8.0000) circle (1.2pt);
\fill[red!80!black] (12.0000,8.0000) circle (1.2pt);
\fill[red!80!black] (0.0000,8.0000) circle (1.2pt);
\fill[red!80!black] (12.0000,8.0000) circle (1.2pt);
\fill[red!80!black] (0.0000,8.0000) circle (1.2pt);
\fill[red!80!black] (12.0000,8.0000) circle (1.2pt);
\fill[red!80!black] (0.0000,8.0000) circle (1.2pt);
\fill[red!80!black] (12.0000,8.0000) circle (1.2pt);
\fill[red!80!black] (0.0000,8.0000) circle (1.2pt);
\fill[red!80!black] (12.0000,8.0000) circle (1.2pt);
\fill[red!80!black] (0.0000,8.0000) circle (1.2pt);
\fill[red!80!black] (12.0000,8.0000) circle (1.2pt);
\fill[red!80!black] (0.0000,8.0000) circle (1.2pt);
\fill[red!80!black] (12.0000,8.0000) circle (1.2pt);
\fill[red!80!black] (0.0000,8.0000) circle (1.2pt);
\fill[red!80!black] (12.0000,8.0000) circle (1.2pt);
\fill[red!80!black] (0.0000,8.0000) circle (1.2pt);
\fill[red!80!black] (12.0000,8.0000) circle (1.2pt);
\fill[red!80!black] (0.0000,8.0000) circle (1.2pt);
\fill[red!80!black] (12.0000,8.0000) circle (1.2pt);
\fill[red!80!black] (0.0000,8.0000) circle (1.2pt);
\fill[red!80!black] (12.0000,8.0000) circle (1.2pt);
\fill[red!80!black] (0.0000,8.0000) circle (1.2pt);
\fill[red!80!black] (12.0000,8.0000) circle (1.2pt);
\fill[red!80!black] (0.0000,8.0000) circle (1.2pt);
\fill[red!80!black] (12.0000,8.0000) circle (1.2pt);
\fill[red!80!black] (0.0000,8.0000) circle (1.2pt);
\fill[red!80!black] (12.0000,8.0000) circle (1.2pt);
\fill[red!80!black] (0.0000,8.0000) circle (1.2pt);
\fill[red!80!black] (12.0000,8.0000) circle (1.2pt);
\fill[red!80!black] (0.0000,8.0000) circle (1.2pt);
\fill[red!80!black] (12.0000,8.0000) circle (1.2pt);
\fill[red!80!black] (0.0000,8.0000) circle (1.2pt);
\fill[red!80!black] (12.0000,8.0000) circle (1.2pt);
\fill[red!80!black] (0.0000,8.0000) circle (1.2pt);
\fill[red!80!black] (12.0000,8.0000) circle (1.2pt);
\fill[red!80!black] (0.0000,8.0000) circle (1.2pt);
\fill[red!80!black] (12.0000,8.0000) circle (1.2pt);
\fill[red!80!black] (0.0000,8.0000) circle (1.2pt);
\fill[red!80!black] (12.0000,8.0000) circle (1.2pt);
\fill[red!80!black] (0.0000,8.0000) circle (1.2pt);
\fill[red!80!black] (12.0000,8.0000) circle (1.2pt);
\fill[red!80!black] (0.0000,8.0000) circle (1.2pt);
\fill[red!80!black] (12.0000,8.0000) circle (1.2pt);
\fill[red!80!black] (0.0000,8.0000) circle (1.2pt);
\fill[red!80!black] (12.0000,8.0000) circle (1.2pt);
\fill[red!80!black] (0.0000,8.0000) circle (1.2pt);
\fill[red!80!black] (12.0000,8.0000) circle (1.2pt);
\fill[red!80!black] (0.0000,8.0000) circle (1.2pt);
\fill[red!80!black] (12.0000,8.0000) circle (1.2pt);
\fill[red!80!black] (0.0000,8.0000) circle (1.2pt);
\fill[red!80!black] (12.0000,8.0000) circle (1.2pt);
\fill[red!80!black] (0.0000,8.0000) circle (1.2pt);
\fill[red!80!black] (12.0000,8.0000) circle (1.2pt);
\fill[red!80!black] (0.0000,8.0000) circle (1.2pt);
\fill[red!80!black] (12.0000,8.0000) circle (1.2pt);
\fill[red!80!black] (0.0000,8.0000) circle (1.2pt);
\fill[red!80!black] (12.0000,8.0000) circle (1.2pt);
\fill[red!80!black] (0.0000,8.0000) circle (1.2pt);
\fill[red!80!black] (12.0000,8.0000) circle (1.2pt);
\fill[red!80!black] (0.0000,8.0000) circle (1.2pt);
\fill[red!80!black] (12.0000,8.0000) circle (1.2pt);
\fill[red!80!black] (0.0000,8.0000) circle (1.2pt);
\fill[red!80!black] (12.0000,8.0000) circle (1.2pt);
\fill[red!80!black] (0.0000,8.0000) circle (1.2pt);
\fill[red!80!black] (12.0000,8.0000) circle (1.2pt);
\fill[red!80!black] (0.0000,8.0000) circle (1.2pt);
\fill[red!80!black] (12.0000,8.0000) circle (1.2pt);
\fill[red!80!black] (0.0000,8.0000) circle (1.2pt);
\fill[red!80!black] (12.0000,8.0000) circle (1.2pt);
\fill[red!80!black] (0.0000,8.0000) circle (1.2pt);
\fill[red!80!black] (12.0000,8.0000) circle (1.2pt);
\fill[red!80!black] (0.0000,8.0000) circle (1.2pt);
\fill[red!80!black] (12.0000,8.0000) circle (1.2pt);
\fill[red!80!black] (0.0000,8.0000) circle (1.2pt);
\fill[red!80!black] (12.0000,8.0000) circle (1.2pt);
\fill[red!80!black] (0.0000,8.0000) circle (1.2pt);
\fill[red!80!black] (12.0000,8.0000) circle (1.2pt);
\fill[red!80!black] (0.0000,8.0000) circle (1.2pt);
\fill[red!80!black] (12.0000,8.0000) circle (1.2pt);
\fill[red!80!black] (0.0000,8.0000) circle (1.2pt);
\fill[red!80!black] (12.0000,8.0000) circle (1.2pt);
\fill[red!80!black] (0.0000,8.0000) circle (1.2pt);
\fill[red!80!black] (12.0000,8.0000) circle (1.2pt);
\fill[red!80!black] (0.0000,8.0000) circle (1.2pt);
\fill[red!80!black] (12.0000,8.0000) circle (1.2pt);
\fill[red!80!black] (0.0000,8.0000) circle (1.2pt);
\fill[red!80!black] (12.0000,8.0000) circle (1.2pt);
\fill[red!80!black] (0.0000,8.0000) circle (1.2pt);
\fill[red!80!black] (12.0000,8.0000) circle (1.2pt);
\fill[red!80!black] (0.0000,8.0000) circle (1.2pt);
\fill[red!80!black] (12.0000,8.0000) circle (1.2pt);
\fill[red!80!black] (0.0000,8.0000) circle (1.2pt);
\fill[red!80!black] (12.0000,8.0000) circle (1.2pt);
\fill[red!80!black] (0.0000,8.0000) circle (1.2pt);
\fill[red!80!black] (12.0000,8.0000) circle (1.2pt);
\fill[red!80!black] (0.0000,8.0000) circle (1.2pt);
\fill[red!80!black] (12.0000,8.0000) circle (1.2pt);
\fill[red!80!black] (0.0000,8.0000) circle (1.2pt);
\fill[red!80!black] (12.0000,8.0000) circle (1.2pt);
\fill[red!80!black] (0.0000,8.0000) circle (1.2pt);
\fill[red!80!black] (12.0000,8.0000) circle (1.2pt);
\fill[red!80!black] (0.0000,8.0000) circle (1.2pt);
\fill[red!80!black] (12.0000,8.0000) circle (1.2pt);
\fill[red!80!black] (0.0000,8.0000) circle (1.2pt);
\fill[red!80!black] (12.0000,8.0000) circle (1.2pt);
\fill[red!80!black] (0.0000,8.0000) circle (1.2pt);
\fill[red!80!black] (12.0000,8.0000) circle (1.2pt);
\fill[red!80!black] (0.0000,8.0000) circle (1.2pt);
\fill[red!80!black] (12.0000,8.0000) circle (1.2pt);
\fill[red!80!black] (0.0000,8.0000) circle (1.2pt);
\fill[red!80!black] (12.0000,8.0000) circle (1.2pt);
\fill[red!80!black] (0.0000,8.0000) circle (1.2pt);
\fill[red!80!black] (12.0000,8.0000) circle (1.2pt);
\fill[red!80!black] (0.0000,8.0000) circle (1.2pt);
\fill[red!80!black] (12.0000,8.0000) circle (1.2pt);
\fill[red!80!black] (0.0000,8.0000) circle (1.2pt);
\fill[red!80!black] (12.0000,8.0000) circle (1.2pt);
\fill[red!80!black] (0.0000,8.0000) circle (1.2pt);
\fill[red!80!black] (12.0000,8.0000) circle (1.2pt);
\fill[red!80!black] (0.0000,8.0000) circle (1.2pt);
\fill[red!80!black] (12.0000,8.0000) circle (1.2pt);
\fill[red!80!black] (0.0000,8.0000) circle (1.2pt);
\fill[red!80!black] (12.0000,8.0000) circle (1.2pt);
\fill[red!80!black] (0.0000,8.0000) circle (1.2pt);
\fill[red!80!black] (12.0000,8.0000) circle (1.2pt);
\fill[red!80!black] (0.0000,8.0000) circle (1.2pt);
\fill[red!80!black] (12.0000,8.0000) circle (1.2pt);
\fill[red!80!black] (0.0000,8.0000) circle (1.2pt);
\fill[red!80!black] (12.0000,8.0000) circle (1.2pt);
\fill[red!80!black] (0.0000,8.0000) circle (1.2pt);
\fill[red!80!black] (12.0000,8.0000) circle (1.2pt);
\fill[red!80!black] (0.0000,8.0000) circle (1.2pt);
\fill[red!80!black] (12.0000,8.0000) circle (1.2pt);
\fill[red!80!black] (0.0000,8.0000) circle (1.2pt);
\fill[red!80!black] (12.0000,8.0000) circle (1.2pt);
\fill[red!80!black] (0.0000,8.0000) circle (1.2pt);
\fill[red!80!black] (12.0000,8.0000) circle (1.2pt);
\fill[red!80!black] (0.0000,8.0000) circle (1.2pt);
\fill[red!80!black] (12.0000,8.0000) circle (1.2pt);
\fill[red!80!black] (0.0000,8.0000) circle (1.2pt);
\fill[red!80!black] (12.0000,8.0000) circle (1.2pt);
\fill[red!80!black] (0.0000,8.0000) circle (1.2pt);
\fill[red!80!black] (12.0000,8.0000) circle (1.2pt);
\fill[red!80!black] (0.0000,8.0000) circle (1.2pt);
\fill[red!80!black] (12.0000,8.0000) circle (1.2pt);
\fill[red!80!black] (0.0000,8.0000) circle (1.2pt);
\fill[red!80!black] (12.0000,8.0000) circle (1.2pt);
\fill[red!80!black] (0.0000,8.0000) circle (1.2pt);
\fill[red!80!black] (12.0000,8.0000) circle (1.2pt);
\fill[red!80!black] (0.0000,8.0000) circle (1.2pt);
\fill[red!80!black] (12.0000,8.0000) circle (1.2pt);
\fill[red!80!black] (0.0000,8.0000) circle (1.2pt);
\fill[red!80!black] (12.0000,8.0000) circle (1.2pt);
\fill[red!80!black] (0.0000,8.0000) circle (1.2pt);
\fill[red!80!black] (12.0000,8.0000) circle (1.2pt);
\fill[red!80!black] (0.0000,8.0000) circle (1.2pt);
\fill[red!80!black] (12.0000,8.0000) circle (1.2pt);
\fill[red!80!black] (0.0000,8.0000) circle (1.2pt);
\fill[red!80!black] (12.0000,8.0000) circle (1.2pt);
\fill[red!80!black] (0.0000,8.0000) circle (1.2pt);
\fill[red!80!black] (12.0000,8.0000) circle (1.2pt);
\fill[red!80!black] (0.0000,8.0000) circle (1.2pt);
\fill[red!80!black] (12.0000,8.0000) circle (1.2pt);
\fill[red!80!black] (0.0000,8.0000) circle (1.2pt);
\fill[red!80!black] (12.0000,8.0000) circle (1.2pt);
\fill[red!80!black] (0.0000,8.0000) circle (1.2pt);
\fill[red!80!black] (12.0000,8.0000) circle (1.2pt);
\fill[red!80!black] (0.0000,8.0000) circle (1.2pt);
\fill[red!80!black] (12.0000,8.0000) circle (1.2pt);
\fill[red!80!black] (0.0000,8.0000) circle (1.2pt);
\fill[red!80!black] (12.0000,8.0000) circle (1.2pt);
\fill[red!80!black] (0.0000,8.0000) circle (1.2pt);
\fill[red!80!black] (12.0000,8.0000) circle (1.2pt);
\fill[red!80!black] (0.0000,8.0000) circle (1.2pt);
\fill[red!80!black] (12.0000,8.0000) circle (1.2pt);
\fill[red!80!black] (0.0000,8.0000) circle (1.2pt);
\fill[red!80!black] (12.0000,8.0000) circle (1.2pt);
\fill[red!80!black] (0.0000,8.0000) circle (1.2pt);
\fill[red!80!black] (12.0000,8.0000) circle (1.2pt);
\fill[red!80!black] (0.0000,8.0000) circle (1.2pt);
\fill[red!80!black] (12.0000,8.0000) circle (1.2pt);
\fill[red!80!black] (0.0000,8.0000) circle (1.2pt);
\fill[red!80!black] (12.0000,8.0000) circle (1.2pt);
\fill[red!80!black] (0.0000,8.0000) circle (1.2pt);
\fill[red!80!black] (12.0000,8.0000) circle (1.2pt);
\fill[red!80!black] (0.0000,8.0000) circle (1.2pt);
\fill[red!80!black] (12.0000,8.0000) circle (1.2pt);
\fill[red!80!black] (0.0000,8.0000) circle (1.2pt);
\fill[red!80!black] (12.0000,8.0000) circle (1.2pt);
\fill[red!80!black] (0.0000,8.0000) circle (1.2pt);
\fill[red!80!black] (12.0000,8.0000) circle (1.2pt);
\fill[red!80!black] (0.0000,8.0000) circle (1.2pt);
\fill[red!80!black] (12.0000,8.0000) circle (1.2pt);
\fill[red!80!black] (0.0000,8.0000) circle (1.2pt);
\fill[red!80!black] (12.0000,8.0000) circle (1.2pt);
\fill[red!80!black] (0.0000,8.0000) circle (1.2pt);
\fill[red!80!black] (12.0000,8.0000) circle (1.2pt);
\fill[red!80!black] (0.0000,8.0000) circle (1.2pt);
\fill[red!80!black] (12.0000,8.0000) circle (1.2pt);
\fill[red!80!black] (0.0000,8.0000) circle (1.2pt);
\fill[red!80!black] (12.0000,8.0000) circle (1.2pt);
\fill[red!80!black] (0.0000,8.0000) circle (1.2pt);
\fill[red!80!black] (12.0000,8.0000) circle (1.2pt);
\fill[red!80!black] (0.0000,8.0000) circle (1.2pt);
\fill[red!80!black] (12.0000,8.0000) circle (1.2pt);
\fill[red!80!black] (0.0000,8.0000) circle (1.2pt);
\fill[red!80!black] (12.0000,8.0000) circle (1.2pt);
\fill[red!80!black] (0.0000,8.0000) circle (1.2pt);
\fill[red!80!black] (12.0000,8.0000) circle (1.2pt);
\fill[red!80!black] (0.0000,8.0000) circle (1.2pt);
\fill[red!80!black] (12.0000,8.0000) circle (1.2pt);
\fill[red!80!black] (0.0000,8.0000) circle (1.2pt);
\fill[red!80!black] (12.0000,8.0000) circle (1.2pt);
\fill[red!80!black] (0.0000,8.0000) circle (1.2pt);
\fill[red!80!black] (12.0000,8.0000) circle (1.2pt);
\fill[red!80!black] (0.0000,8.0000) circle (1.2pt);
\fill[red!80!black] (12.0000,8.0000) circle (1.2pt);
\fill[red!80!black] (0.0000,8.0000) circle (1.2pt);
\fill[red!80!black] (12.0000,8.0000) circle (1.2pt);
\fill[red!80!black] (0.0000,8.0000) circle (1.2pt);
\fill[red!80!black] (12.0000,8.0000) circle (1.2pt);
\fill[red!80!black] (0.0000,8.0000) circle (1.2pt);
\fill[red!80!black] (12.0000,8.0000) circle (1.2pt);
\fill[red!80!black] (0.0000,8.0000) circle (1.2pt);
\fill[red!80!black] (12.0000,8.0000) circle (1.2pt);
\fill[red!80!black] (0.0000,8.0000) circle (1.2pt);
\fill[red!80!black] (12.0000,8.0000) circle (1.2pt);
\fill[red!80!black] (0.0000,8.0000) circle (1.2pt);
\fill[red!80!black] (12.0000,8.0000) circle (1.2pt);
\fill[red!80!black] (0.0000,8.0000) circle (1.2pt);
\fill[red!80!black] (12.0000,8.0000) circle (1.2pt);
\fill[red!80!black] (0.0000,8.0000) circle (1.2pt);
\fill[red!80!black] (12.0000,8.0000) circle (1.2pt);
\fill[red!80!black] (0.0000,8.0000) circle (1.2pt);
\fill[red!80!black] (12.0000,8.0000) circle (1.2pt);
\fill[red!80!black] (0.0000,8.0000) circle (1.2pt);
\fill[red!80!black] (12.0000,8.0000) circle (1.2pt);
\fill[red!80!black] (0.0000,8.0000) circle (1.2pt);
\fill[red!80!black] (12.0000,8.0000) circle (1.2pt);
\fill[red!80!black] (0.0000,8.0000) circle (1.2pt);
\fill[red!80!black] (12.0000,8.0000) circle (1.2pt);
\fill[red!80!black] (0.0000,8.0000) circle (1.2pt);
\fill[red!80!black] (12.0000,8.0000) circle (1.2pt);
\fill[red!80!black] (0.0000,8.0000) circle (1.2pt);
\fill[red!80!black] (12.0000,8.0000) circle (1.2pt);
\fill[red!80!black] (0.0000,8.0000) circle (1.2pt);
\fill[red!80!black] (12.0000,8.0000) circle (1.2pt);
\fill[red!80!black] (0.0000,8.0000) circle (1.2pt);
\fill[red!80!black] (12.0000,8.0000) circle (1.2pt);
\fill[red!80!black] (0.0000,8.0000) circle (1.2pt);
\fill[red!80!black] (12.0000,8.0000) circle (1.2pt);
\fill[red!80!black] (0.0000,8.0000) circle (1.2pt);
\fill[red!80!black] (12.0000,8.0000) circle (1.2pt);
\fill[red!80!black] (0.0000,8.0000) circle (1.2pt);
\fill[red!80!black] (12.0000,8.0000) circle (1.2pt);
\fill[red!80!black] (0.0000,8.0000) circle (1.2pt);
\fill[red!80!black] (12.0000,8.0000) circle (1.2pt);
\fill[red!80!black] (0.0000,8.0000) circle (1.2pt);
\fill[red!80!black] (12.0000,8.0000) circle (1.2pt);
\fill[red!80!black] (0.0000,8.0000) circle (1.2pt);
\fill[red!80!black] (12.0000,8.0000) circle (1.2pt);
\fill[red!80!black] (0.0000,8.0000) circle (1.2pt);
\fill[red!80!black] (12.0000,8.0000) circle (1.2pt);
\fill[red!80!black] (0.0000,8.0000) circle (1.2pt);
\fill[red!80!black] (12.0000,8.0000) circle (1.2pt);
\fill[red!80!black] (0.0000,8.0000) circle (1.2pt);
\fill[red!80!black] (12.0000,8.0000) circle (1.2pt);
\fill[red!80!black] (0.0000,8.0000) circle (1.2pt);
\fill[red!80!black] (12.0000,8.0000) circle (1.2pt);
\fill[red!80!black] (0.0000,8.0000) circle (1.2pt);
\fill[red!80!black] (12.0000,8.0000) circle (1.2pt);
\fill[red!80!black] (0.0000,8.0000) circle (1.2pt);
\fill[red!80!black] (12.0000,8.0000) circle (1.2pt);
\fill[red!80!black] (0.0000,8.0000) circle (1.2pt);
\fill[red!80!black] (12.0000,8.0000) circle (1.2pt);
\fill[red!80!black] (0.0000,8.0000) circle (1.2pt);
\fill[red!80!black] (12.0000,8.0000) circle (1.2pt);
\fill[red!80!black] (0.0000,8.0000) circle (1.2pt);
\fill[red!80!black] (12.0000,8.0000) circle (1.2pt);
\fill[red!80!black] (0.0000,8.0000) circle (1.2pt);
\fill[red!80!black] (12.0000,8.0000) circle (1.2pt);
\fill[red!80!black] (0.0000,8.0000) circle (1.2pt);
\fill[red!80!black] (12.0000,8.0000) circle (1.2pt);
\fill[red!80!black] (0.0000,8.0000) circle (1.2pt);
\fill[red!80!black] (12.0000,8.0000) circle (1.2pt);
\fill[red!80!black] (0.0000,8.0000) circle (1.2pt);
\fill[red!80!black] (12.0000,8.0000) circle (1.2pt);
\fill[red!80!black] (0.0000,8.0000) circle (1.2pt);
\fill[red!80!black] (12.0000,8.0000) circle (1.2pt);
\fill[red!80!black] (0.0000,8.0000) circle (1.2pt);
\fill[red!80!black] (12.0000,8.0000) circle (1.2pt);
\fill[red!80!black] (0.0000,8.0000) circle (1.2pt);
\fill[red!80!black] (12.0000,8.0000) circle (1.2pt);
\fill[red!80!black] (0.0000,8.0000) circle (1.2pt);
\fill[red!80!black] (12.0000,8.0000) circle (1.2pt);
\fill[red!80!black] (0.0000,8.0000) circle (1.2pt);
\fill[red!80!black] (12.0000,8.0000) circle (1.2pt);
\fill[red!80!black] (0.0000,8.0000) circle (1.2pt);
\fill[red!80!black] (12.0000,8.0000) circle (1.2pt);
\fill[red!80!black] (0.0000,8.0000) circle (1.2pt);
\fill[red!80!black] (12.0000,8.0000) circle (1.2pt);
\fill[red!80!black] (0.0000,8.0000) circle (1.2pt);
\fill[red!80!black] (12.0000,8.0000) circle (1.2pt);
\fill[red!80!black] (0.0000,8.0000) circle (1.2pt);
\fill[red!80!black] (12.0000,8.0000) circle (1.2pt);
\fill[red!80!black] (0.0000,8.0000) circle (1.2pt);
\fill[red!80!black] (12.0000,8.0000) circle (1.2pt);
\fill[red!80!black] (0.0000,8.0000) circle (1.2pt);
\fill[red!80!black] (12.0000,8.0000) circle (1.2pt);
\fill[red!80!black] (0.0000,8.0000) circle (1.2pt);
\fill[red!80!black] (12.0000,8.0000) circle (1.2pt);
\fill[red!80!black] (0.0000,8.0000) circle (1.2pt);
\fill[red!80!black] (12.0000,8.0000) circle (1.2pt);
\fill[red!80!black] (0.0000,8.0000) circle (1.2pt);
\fill[red!80!black] (12.0000,8.0000) circle (1.2pt);
\fill[red!80!black] (0.0000,8.0000) circle (1.2pt);
\fill[red!80!black] (12.0000,8.0000) circle (1.2pt);
\fill[red!80!black] (0.0000,8.0000) circle (1.2pt);
\fill[red!80!black] (12.0000,8.0000) circle (1.2pt);
\fill[red!80!black] (0.0000,8.0000) circle (1.2pt);
\fill[red!80!black] (12.0000,8.0000) circle (1.2pt);
\fill[red!80!black] (0.0000,8.0000) circle (1.2pt);
\fill[red!80!black] (12.0000,8.0000) circle (1.2pt);
\fill[red!80!black] (0.0000,8.0000) circle (1.2pt);
\fill[red!80!black] (12.0000,8.0000) circle (1.2pt);
\fill[red!80!black] (0.0000,8.0000) circle (1.2pt);
\fill[red!80!black] (12.0000,8.0000) circle (1.2pt);
\fill[red!80!black] (0.0000,8.0000) circle (1.2pt);
\fill[red!80!black] (12.0000,8.0000) circle (1.2pt);
\fill[red!80!black] (0.0000,8.0000) circle (1.2pt);
\fill[red!80!black] (12.0000,8.0000) circle (1.2pt);
\fill[red!80!black] (0.0000,8.0000) circle (1.2pt);
\fill[red!80!black] (12.0000,8.0000) circle (1.2pt);
\fill[red!80!black] (0.0000,8.0000) circle (1.2pt);
\fill[red!80!black] (12.0000,8.0000) circle (1.2pt);
\fill[red!80!black] (0.0000,8.0000) circle (1.2pt);
\fill[red!80!black] (12.0000,8.0000) circle (1.2pt);
\fill[red!80!black] (0.0000,8.0000) circle (1.2pt);
\fill[red!80!black] (12.0000,8.0000) circle (1.2pt);
\fill[red!80!black] (0.0000,8.0000) circle (1.2pt);
\fill[red!80!black] (12.0000,8.0000) circle (1.2pt);
\fill[red!80!black] (0.0000,8.0000) circle (1.2pt);
\fill[red!80!black] (12.0000,8.0000) circle (1.2pt);
\fill[red!80!black] (0.0000,8.0000) circle (1.2pt);
\fill[red!80!black] (12.0000,8.0000) circle (1.2pt);
\fill[red!80!black] (0.0000,8.0000) circle (1.2pt);
\fill[red!80!black] (12.0000,8.0000) circle (1.2pt);
\fill[red!80!black] (0.0000,8.0000) circle (1.2pt);
\fill[red!80!black] (12.0000,8.0000) circle (1.2pt);
\fill[red!80!black] (0.0000,8.0000) circle (1.2pt);
\fill[red!80!black] (12.0000,8.0000) circle (1.2pt);
\fill[red!80!black] (0.0000,8.0000) circle (1.2pt);
\fill[red!80!black] (12.0000,8.0000) circle (1.2pt);
\fill[red!80!black] (0.0000,8.0000) circle (1.2pt);
\fill[red!80!black] (12.0000,8.0000) circle (1.2pt);
\fill[red!80!black] (0.0000,8.0000) circle (1.2pt);
\fill[red!80!black] (12.0000,8.0000) circle (1.2pt);
\fill[red!80!black] (0.0000,8.0000) circle (1.2pt);
\fill[red!80!black] (12.0000,8.0000) circle (1.2pt);
\fill[red!80!black] (0.0000,8.0000) circle (1.2pt);
\fill[red!80!black] (12.0000,8.0000) circle (1.2pt);
\fill[red!80!black] (0.0000,8.0000) circle (1.2pt);
\fill[red!80!black] (12.0000,8.0000) circle (1.2pt);
\fill[red!80!black] (0.0000,8.0000) circle (1.2pt);
\fill[red!80!black] (12.0000,8.0000) circle (1.2pt);
\fill[red!80!black] (0.0000,8.0000) circle (1.2pt);
\fill[red!80!black] (12.0000,8.0000) circle (1.2pt);
\fill[red!80!black] (0.0000,8.0000) circle (1.2pt);
\fill[red!80!black] (12.0000,8.0000) circle (1.2pt);
\fill[red!80!black] (0.0000,8.0000) circle (1.2pt);
\fill[red!80!black] (12.0000,8.0000) circle (1.2pt);
\fill[red!80!black] (0.0000,8.0000) circle (1.2pt);
\fill[red!80!black] (12.0000,8.0000) circle (1.2pt);
\fill[red!80!black] (0.0000,8.0000) circle (1.2pt);
\fill[red!80!black] (12.0000,8.0000) circle (1.2pt);
\fill[red!80!black] (0.0000,8.0000) circle (1.2pt);
\fill[red!80!black] (12.0000,8.0000) circle (1.2pt);
\fill[red!80!black] (0.0000,8.0000) circle (1.2pt);
\fill[red!80!black] (12.0000,8.0000) circle (1.2pt);
\fill[red!80!black] (0.0000,8.0000) circle (1.2pt);
\fill[red!80!black] (12.0000,8.0000) circle (1.2pt);
\fill[red!80!black] (0.0000,8.0000) circle (1.2pt);
\fill[red!80!black] (12.0000,8.0000) circle (1.2pt);
\fill[red!80!black] (0.0000,8.0000) circle (1.2pt);
\fill[red!80!black] (12.0000,8.0000) circle (1.2pt);
\fill[red!80!black] (0.0000,8.0000) circle (1.2pt);
\fill[red!80!black] (12.0000,8.0000) circle (1.2pt);
\fill[red!80!black] (0.0000,8.0000) circle (1.2pt);
\fill[red!80!black] (12.0000,8.0000) circle (1.2pt);
\fill[red!80!black] (0.0000,8.0000) circle (1.2pt);
\fill[red!80!black] (12.0000,8.0000) circle (1.2pt);
\fill[red!80!black] (0.0000,8.0000) circle (1.2pt);
\fill[red!80!black] (12.0000,8.0000) circle (1.2pt);
\fill[red!80!black] (0.0000,8.0000) circle (1.2pt);
\fill[red!80!black] (12.0000,8.0000) circle (1.2pt);
\fill[red!80!black] (0.0000,8.0000) circle (1.2pt);
\fill[red!80!black] (12.0000,8.0000) circle (1.2pt);
\fill[red!80!black] (0.0000,8.0000) circle (1.2pt);
\fill[red!80!black] (12.0000,8.0000) circle (1.2pt);
\fill[red!80!black] (0.0000,8.0000) circle (1.2pt);
\fill[red!80!black] (12.0000,8.0000) circle (1.2pt);
\fill[red!80!black] (0.0000,8.0000) circle (1.2pt);
\fill[red!80!black] (12.0000,8.0000) circle (1.2pt);
\fill[red!80!black] (0.0000,8.0000) circle (1.2pt);
\fill[red!80!black] (12.0000,8.0000) circle (1.2pt);
\fill[red!80!black] (0.0000,8.0000) circle (1.2pt);
\fill[red!80!black] (12.0000,8.0000) circle (1.2pt);
\fill[red!80!black] (0.0000,8.0000) circle (1.2pt);
\fill[red!80!black] (12.0000,8.0000) circle (1.2pt);
\fill[red!80!black] (0.0000,8.0000) circle (1.2pt);
\fill[red!80!black] (12.0000,8.0000) circle (1.2pt);
\fill[red!80!black] (0.0000,8.0000) circle (1.2pt);
\fill[red!80!black] (12.0000,8.0000) circle (1.2pt);
\fill[red!80!black] (0.0000,8.0000) circle (1.2pt);
\fill[red!80!black] (12.0000,8.0000) circle (1.2pt);
\fill[red!80!black] (0.0000,8.0000) circle (1.2pt);
\fill[red!80!black] (12.0000,8.0000) circle (1.2pt);
\fill[red!80!black] (0.0000,8.0000) circle (1.2pt);
\fill[red!80!black] (12.0000,8.0000) circle (1.2pt);
\fill[red!80!black] (0.0000,8.0000) circle (1.2pt);
\fill[red!80!black] (12.0000,8.0000) circle (1.2pt);
\fill[red!80!black] (0.0000,8.0000) circle (1.2pt);
\fill[red!80!black] (12.0000,8.0000) circle (1.2pt);
\fill[red!80!black] (0.0000,8.0000) circle (1.2pt);
\fill[red!80!black] (12.0000,8.0000) circle (1.2pt);
\fill[red!80!black] (0.0000,8.0000) circle (1.2pt);
\fill[red!80!black] (12.0000,8.0000) circle (1.2pt);
\fill[red!80!black] (0.0000,8.0000) circle (1.2pt);
\fill[red!80!black] (12.0000,8.0000) circle (1.2pt);
\fill[red!80!black] (0.0000,8.0000) circle (1.2pt);
\fill[red!80!black] (12.0000,8.0000) circle (1.2pt);
\fill[red!80!black] (0.0000,8.0000) circle (1.2pt);
\fill[red!80!black] (12.0000,8.0000) circle (1.2pt);
\fill[red!80!black] (0.0000,8.0000) circle (1.2pt);
\fill[red!80!black] (12.0000,8.0000) circle (1.2pt);
\fill[red!80!black] (0.0000,8.0000) circle (1.2pt);
\fill[red!80!black] (12.0000,8.0000) circle (1.2pt);
\fill[red!80!black] (0.0000,8.0000) circle (1.2pt);
\fill[red!80!black] (12.0000,8.0000) circle (1.2pt);
\fill[red!80!black] (0.0000,8.0000) circle (1.2pt);
\fill[red!80!black] (12.0000,8.0000) circle (1.2pt);
\fill[red!80!black] (0.0000,8.0000) circle (1.2pt);
\fill[red!80!black] (12.0000,8.0000) circle (1.2pt);
\fill[red!80!black] (0.0000,8.0000) circle (1.2pt);
\fill[red!80!black] (12.0000,8.0000) circle (1.2pt);
\fill[red!80!black] (0.0000,8.0000) circle (1.2pt);
\fill[red!80!black] (12.0000,8.0000) circle (1.2pt);
\fill[red!80!black] (0.0000,8.0000) circle (1.2pt);
\fill[red!80!black] (12.0000,8.0000) circle (1.2pt);
\fill[red!80!black] (0.0000,8.0000) circle (1.2pt);
\fill[red!80!black] (12.0000,8.0000) circle (1.2pt);
\fill[red!80!black] (0.0000,8.0000) circle (1.2pt);
\fill[red!80!black] (12.0000,8.0000) circle (1.2pt);
\fill[red!80!black] (0.0000,8.0000) circle (1.2pt);
\fill[red!80!black] (12.0000,8.0000) circle (1.2pt);
\fill[red!80!black] (0.0000,8.0000) circle (1.2pt);
\fill[red!80!black] (12.0000,8.0000) circle (1.2pt);
\fill[red!80!black] (0.0000,8.0000) circle (1.2pt);
\fill[red!80!black] (12.0000,8.0000) circle (1.2pt);
\fill[red!80!black] (0.0000,8.0000) circle (1.2pt);
\fill[red!80!black] (12.0000,8.0000) circle (1.2pt);
\fill[red!80!black] (0.0000,8.0000) circle (1.2pt);
\fill[red!80!black] (12.0000,8.0000) circle (1.2pt);
\fill[red!80!black] (0.0000,8.0000) circle (1.2pt);
\fill[red!80!black] (12.0000,8.0000) circle (1.2pt);
\fill[red!80!black] (0.0000,8.0000) circle (1.2pt);
\fill[red!80!black] (12.0000,8.0000) circle (1.2pt);
\fill[red!80!black] (0.0000,8.0000) circle (1.2pt);
\fill[red!80!black] (12.0000,8.0000) circle (1.2pt);
\fill[red!80!black] (0.0000,8.0000) circle (1.2pt);
\fill[red!80!black] (12.0000,8.0000) circle (1.2pt);
\fill[red!80!black] (0.0000,8.0000) circle (1.2pt);
\fill[red!80!black] (12.0000,8.0000) circle (1.2pt);
\fill[red!80!black] (0.0000,8.0000) circle (1.2pt);
\fill[red!80!black] (12.0000,8.0000) circle (1.2pt);
\fill[red!80!black] (0.0000,8.0000) circle (1.2pt);
\fill[red!80!black] (12.0000,8.0000) circle (1.2pt);
\fill[red!80!black] (0.0000,8.0000) circle (1.2pt);
\fill[red!80!black] (12.0000,8.0000) circle (1.2pt);
\fill[red!80!black] (0.0000,8.0000) circle (1.2pt);
\fill[red!80!black] (12.0000,8.0000) circle (1.2pt);
\fill[red!80!black] (0.0000,8.0000) circle (1.2pt);
\fill[red!80!black] (12.0000,8.0000) circle (1.2pt);
\fill[red!80!black] (0.0000,8.0000) circle (1.2pt);
\fill[red!80!black] (12.0000,8.0000) circle (1.2pt);
\fill[red!80!black] (0.0000,8.0000) circle (1.2pt);
\fill[red!80!black] (12.0000,8.0000) circle (1.2pt);
\fill[red!80!black] (0.0000,8.0000) circle (1.2pt);
\fill[red!80!black] (12.0000,8.0000) circle (1.2pt);
\fill[red!80!black] (0.0000,8.0000) circle (1.2pt);
\fill[red!80!black] (12.0000,8.0000) circle (1.2pt);
\fill[red!80!black] (0.0000,8.0000) circle (1.2pt);
\fill[red!80!black] (12.0000,8.0000) circle (1.2pt);
\fill[red!80!black] (0.0000,8.0000) circle (1.2pt);
\fill[red!80!black] (12.0000,8.0000) circle (1.2pt);
\fill[red!80!black] (0.0000,8.0000) circle (1.2pt);
\fill[red!80!black] (12.0000,8.0000) circle (1.2pt);
\fill[red!80!black] (0.0000,8.0000) circle (1.2pt);
\fill[red!80!black] (12.0000,8.0000) circle (1.2pt);
\fill[red!80!black] (0.0000,8.0000) circle (1.2pt);
\fill[red!80!black] (12.0000,8.0000) circle (1.2pt);
\fill[red!80!black] (0.0000,8.0000) circle (1.2pt);
\fill[red!80!black] (12.0000,8.0000) circle (1.2pt);
\fill[red!80!black] (0.0000,8.0000) circle (1.2pt);
\fill[red!80!black] (12.0000,8.0000) circle (1.2pt);
\fill[red!80!black] (0.0000,8.0000) circle (1.2pt);
\fill[red!80!black] (12.0000,8.0000) circle (1.2pt);
\fill[red!80!black] (0.0000,8.0000) circle (1.2pt);
\fill[red!80!black] (12.0000,8.0000) circle (1.2pt);
\fill[red!80!black] (0.0000,8.0000) circle (1.2pt);
\fill[red!80!black] (12.0000,8.0000) circle (1.2pt);
\fill[red!80!black] (0.0000,8.0000) circle (1.2pt);
\fill[red!80!black] (12.0000,8.0000) circle (1.2pt);
\fill[red!80!black] (0.0000,8.0000) circle (1.2pt);
\fill[red!80!black] (12.0000,8.0000) circle (1.2pt);
\fill[red!80!black] (0.0000,8.0000) circle (1.2pt);
\fill[red!80!black] (12.0000,8.0000) circle (1.2pt);
\fill[red!80!black] (0.0000,8.0000) circle (1.2pt);
\fill[red!80!black] (12.0000,8.0000) circle (1.2pt);
\fill[red!80!black] (0.0000,8.0000) circle (1.2pt);
\fill[red!80!black] (12.0000,8.0000) circle (1.2pt);
\fill[red!80!black] (0.0000,8.0000) circle (1.2pt);
\fill[red!80!black] (12.0000,8.0000) circle (1.2pt);
\fill[red!80!black] (0.0000,8.0000) circle (1.2pt);
\fill[red!80!black] (12.0000,8.0000) circle (1.2pt);
\fill[red!80!black] (0.0000,8.0000) circle (1.2pt);
\fill[red!80!black] (12.0000,8.0000) circle (1.2pt);
\fill[red!80!black] (0.0000,8.0000) circle (1.2pt);
\fill[red!80!black] (12.0000,8.0000) circle (1.2pt);
\fill[red!80!black] (0.0000,8.0000) circle (1.2pt);
\fill[red!80!black] (12.0000,8.0000) circle (1.2pt);
\fill[red!80!black] (0.0000,8.0000) circle (1.2pt);
\fill[red!80!black] (12.0000,8.0000) circle (1.2pt);
\fill[red!80!black] (0.0000,8.0000) circle (1.2pt);
\fill[red!80!black] (12.0000,8.0000) circle (1.2pt);
\fill[red!80!black] (0.0000,8.0000) circle (1.2pt);
\fill[red!80!black] (12.0000,8.0000) circle (1.2pt);
\fill[red!80!black] (0.0000,8.0000) circle (1.2pt);
\fill[red!80!black] (12.0000,8.0000) circle (1.2pt);
\fill[red!80!black] (0.0000,8.0000) circle (1.2pt);
\fill[red!80!black] (12.0000,8.0000) circle (1.2pt);
\fill[red!80!black] (0.0000,8.0000) circle (1.2pt);
\fill[red!80!black] (12.0000,8.0000) circle (1.2pt);
\fill[red!80!black] (0.0000,8.0000) circle (1.2pt);
\fill[red!80!black] (12.0000,8.0000) circle (1.2pt);
\fill[red!80!black] (0.0000,8.0000) circle (1.2pt);
\fill[red!80!black] (12.0000,8.0000) circle (1.2pt);
\fill[red!80!black] (0.0000,8.0000) circle (1.2pt);
\fill[red!80!black] (12.0000,8.0000) circle (1.2pt);
\fill[red!80!black] (0.0000,8.0000) circle (1.2pt);
\fill[red!80!black] (12.0000,8.0000) circle (1.2pt);
\fill[red!80!black] (0.0000,8.0000) circle (1.2pt);
\fill[red!80!black] (12.0000,8.0000) circle (1.2pt);
\fill[red!80!black] (0.0000,8.0000) circle (1.2pt);
\fill[red!80!black] (12.0000,8.0000) circle (1.2pt);
\fill[red!80!black] (0.0000,8.0000) circle (1.2pt);
\fill[red!80!black] (12.0000,8.0000) circle (1.2pt);
\fill[red!80!black] (0.0000,8.0000) circle (1.2pt);
\fill[red!80!black] (12.0000,8.0000) circle (1.2pt);
\fill[red!80!black] (0.0000,8.0000) circle (1.2pt);
\fill[red!80!black] (12.0000,8.0000) circle (1.2pt);
\fill[red!80!black] (0.0000,8.0000) circle (1.2pt);
\fill[red!80!black] (12.0000,8.0000) circle (1.2pt);
\fill[red!80!black] (0.0000,8.0000) circle (1.2pt);
\fill[red!80!black] (12.0000,8.0000) circle (1.2pt);
\fill[red!80!black] (0.0000,8.0000) circle (1.2pt);
\fill[red!80!black] (12.0000,8.0000) circle (1.2pt);
\fill[red!80!black] (0.0000,8.0000) circle (1.2pt);
\fill[red!80!black] (12.0000,8.0000) circle (1.2pt);
\fill[red!80!black] (0.0000,8.0000) circle (1.2pt);
\fill[red!80!black] (12.0000,8.0000) circle (1.2pt);
\fill[red!80!black] (0.0000,8.0000) circle (1.2pt);
\fill[red!80!black] (12.0000,8.0000) circle (1.2pt);
\fill[red!80!black] (0.0000,8.0000) circle (1.2pt);
\fill[red!80!black] (12.0000,8.0000) circle (1.2pt);
\fill[red!80!black] (0.0000,8.0000) circle (1.2pt);
\fill[red!80!black] (12.0000,8.0000) circle (1.2pt);
\fill[red!80!black] (0.0000,8.0000) circle (1.2pt);
\fill[red!80!black] (12.0000,8.0000) circle (1.2pt);
\fill[red!80!black] (0.0000,8.0000) circle (1.2pt);
\fill[red!80!black] (12.0000,8.0000) circle (1.2pt);
\fill[red!80!black] (0.0000,8.0000) circle (1.2pt);
\fill[red!80!black] (12.0000,8.0000) circle (1.2pt);
\fill[red!80!black] (0.0000,8.0000) circle (1.2pt);
\fill[red!80!black] (12.0000,8.0000) circle (1.2pt);
\fill[red!80!black] (0.0000,8.0000) circle (1.2pt);
\fill[red!80!black] (12.0000,8.0000) circle (1.2pt);
\fill[red!80!black] (0.0000,8.0000) circle (1.2pt);
\fill[red!80!black] (12.0000,8.0000) circle (1.2pt);
\fill[red!80!black] (0.0000,8.0000) circle (1.2pt);
\fill[red!80!black] (12.0000,8.0000) circle (1.2pt);
\fill[red!80!black] (0.0000,8.0000) circle (1.2pt);
\fill[red!80!black] (12.0000,8.0000) circle (1.2pt);
\fill[red!80!black] (0.0000,8.0000) circle (1.2pt);
\fill[red!80!black] (12.0000,8.0000) circle (1.2pt);
\fill[red!80!black] (0.0000,8.0000) circle (1.2pt);
\fill[red!80!black] (12.0000,8.0000) circle (1.2pt);
\fill[red!80!black] (0.0000,8.0000) circle (1.2pt);
\fill[red!80!black] (12.0000,8.0000) circle (1.2pt);
\fill[red!80!black] (0.0000,8.0000) circle (1.2pt);
\fill[red!80!black] (12.0000,8.0000) circle (1.2pt);
\fill[red!80!black] (0.0000,8.0000) circle (1.2pt);
\fill[red!80!black] (12.0000,8.0000) circle (1.2pt);
\fill[red!80!black] (0.0000,8.0000) circle (1.2pt);
\fill[red!80!black] (12.0000,8.0000) circle (1.2pt);
\fill[red!80!black] (0.0000,8.0000) circle (1.2pt);
\fill[red!80!black] (12.0000,8.0000) circle (1.2pt);
\fill[red!80!black] (0.0000,8.0000) circle (1.2pt);
\fill[red!80!black] (12.0000,8.0000) circle (1.2pt);
\fill[red!80!black] (0.0000,8.0000) circle (1.2pt);
\fill[red!80!black] (12.0000,8.0000) circle (1.2pt);
\fill[red!80!black] (0.0000,8.0000) circle (1.2pt);
\fill[red!80!black] (12.0000,8.0000) circle (1.2pt);
\fill[red!80!black] (0.0000,8.0000) circle (1.2pt);
\fill[red!80!black] (12.0000,8.0000) circle (1.2pt);
\fill[red!80!black] (0.0000,8.0000) circle (1.2pt);
\fill[red!80!black] (12.0000,8.0000) circle (1.2pt);
\fill[red!80!black] (0.0000,8.0000) circle (1.2pt);
\fill[red!80!black] (12.0000,8.0000) circle (1.2pt);
\fill[red!80!black] (0.0000,8.0000) circle (1.2pt);
\fill[red!80!black] (12.0000,8.0000) circle (1.2pt);
\fill[red!80!black] (0.0000,8.0000) circle (1.2pt);
\fill[red!80!black] (12.0000,8.0000) circle (1.2pt);
\fill[red!80!black] (0.0000,8.0000) circle (1.2pt);
\fill[red!80!black] (12.0000,8.0000) circle (1.2pt);
\fill[red!80!black] (0.0000,8.0000) circle (1.2pt);
\fill[red!80!black] (12.0000,8.0000) circle (1.2pt);
\fill[red!80!black] (0.0000,8.0000) circle (1.2pt);
\fill[red!80!black] (12.0000,8.0000) circle (1.2pt);
\fill[red!80!black] (0.0000,8.0000) circle (1.2pt);
\fill[red!80!black] (12.0000,8.0000) circle (1.2pt);
\fill[red!80!black] (0.0000,8.0000) circle (1.2pt);
\fill[red!80!black] (12.0000,8.0000) circle (1.2pt);
\fill[red!80!black] (0.0000,8.0000) circle (1.2pt);
\fill[red!80!black] (12.0000,8.0000) circle (1.2pt);
\fill[red!80!black] (0.0000,8.0000) circle (1.2pt);
\fill[red!80!black] (12.0000,8.0000) circle (1.2pt);
\fill[red!80!black] (0.0000,8.0000) circle (1.2pt);
\fill[red!80!black] (12.0000,8.0000) circle (1.2pt);
\fill[red!80!black] (0.0000,8.0000) circle (1.2pt);
\fill[red!80!black] (12.0000,8.0000) circle (1.2pt);
\fill[red!80!black] (0.0000,8.0000) circle (1.2pt);
\fill[red!80!black] (12.0000,8.0000) circle (1.2pt);
\fill[red!80!black] (0.0000,8.0000) circle (1.2pt);
\fill[red!80!black] (12.0000,8.0000) circle (1.2pt);
\fill[red!80!black] (0.0000,8.0000) circle (1.2pt);
\fill[red!80!black] (12.0000,8.0000) circle (1.2pt);
\fill[red!80!black] (0.0000,8.0000) circle (1.2pt);
\fill[red!80!black] (12.0000,8.0000) circle (1.2pt);
\fill[red!80!black] (0.0000,8.0000) circle (1.2pt);
\fill[red!80!black] (12.0000,8.0000) circle (1.2pt);
\fill[red!80!black] (0.0000,8.0000) circle (1.2pt);
\fill[red!80!black] (12.0000,8.0000) circle (1.2pt);
\fill[red!80!black] (0.0000,8.0000) circle (1.2pt);
\fill[red!80!black] (12.0000,8.0000) circle (1.2pt);
\fill[red!80!black] (0.0000,8.0000) circle (1.2pt);
\fill[red!80!black] (12.0000,8.0000) circle (1.2pt);
\fill[red!80!black] (0.0000,8.0000) circle (1.2pt);
\fill[red!80!black] (12.0000,8.0000) circle (1.2pt);
\fill[red!80!black] (0.0000,8.0000) circle (1.2pt);
\fill[red!80!black] (12.0000,8.0000) circle (1.2pt);
\fill[red!80!black] (0.0000,8.0000) circle (1.2pt);
\fill[red!80!black] (12.0000,8.0000) circle (1.2pt);
\fill[red!80!black] (0.0000,8.0000) circle (1.2pt);
\fill[red!80!black] (12.0000,8.0000) circle (1.2pt);
\fill[red!80!black] (0.0000,8.0000) circle (1.2pt);
\fill[red!80!black] (12.0000,8.0000) circle (1.2pt);
\fill[red!80!black] (0.0000,8.0000) circle (1.2pt);
\fill[red!80!black] (12.0000,8.0000) circle (1.2pt);
\fill[red!80!black] (0.0000,8.0000) circle (1.2pt);
\fill[red!80!black] (12.0000,8.0000) circle (1.2pt);
\fill[red!80!black] (0.0000,8.0000) circle (1.2pt);
\fill[red!80!black] (12.0000,8.0000) circle (1.2pt);
\fill[red!80!black] (0.0000,8.0000) circle (1.2pt);
\fill[red!80!black] (12.0000,8.0000) circle (1.2pt);
\fill[red!80!black] (0.0000,8.0000) circle (1.2pt);
\fill[red!80!black] (12.0000,8.0000) circle (1.2pt);
\fill[red!80!black] (0.0000,8.0000) circle (1.2pt);
\fill[red!80!black] (12.0000,8.0000) circle (1.2pt);
\fill[red!80!black] (0.0000,8.0000) circle (1.2pt);
\fill[red!80!black] (12.0000,8.0000) circle (1.2pt);
\fill[red!80!black] (0.0000,8.0000) circle (1.2pt);
\fill[red!80!black] (12.0000,8.0000) circle (1.2pt);
\fill[red!80!black] (0.0000,8.0000) circle (1.2pt);
\fill[red!80!black] (12.0000,8.0000) circle (1.2pt);
\fill[red!80!black] (0.0000,8.0000) circle (1.2pt);
\fill[red!80!black] (12.0000,8.0000) circle (1.2pt);
\fill[red!80!black] (0.0000,8.0000) circle (1.2pt);
\fill[red!80!black] (12.0000,8.0000) circle (1.2pt);
\fill[red!80!black] (0.0000,8.0000) circle (1.2pt);
\fill[red!80!black] (12.0000,8.0000) circle (1.2pt);
\fill[red!80!black] (0.0000,8.0000) circle (1.2pt);
\fill[red!80!black] (12.0000,8.0000) circle (1.2pt);
\fill[red!80!black] (0.0000,8.0000) circle (1.2pt);
\fill[red!80!black] (12.0000,8.0000) circle (1.2pt);
\fill[red!80!black] (0.0000,8.0000) circle (1.2pt);
\fill[red!80!black] (12.0000,8.0000) circle (1.2pt);
\fill[red!80!black] (0.0000,8.0000) circle (1.2pt);
\fill[red!80!black] (12.0000,8.0000) circle (1.2pt);
\fill[red!80!black] (0.0000,8.0000) circle (1.2pt);
\fill[red!80!black] (12.0000,8.0000) circle (1.2pt);
\fill[red!80!black] (0.0000,8.0000) circle (1.2pt);
\fill[red!80!black] (12.0000,8.0000) circle (1.2pt);
\fill[red!80!black] (0.0000,8.0000) circle (1.2pt);
\fill[red!80!black] (12.0000,8.0000) circle (1.2pt);
\fill[red!80!black] (0.0000,8.0000) circle (1.2pt);
\fill[red!80!black] (12.0000,8.0000) circle (1.2pt);
\fill[red!80!black] (0.0000,8.0000) circle (1.2pt);
\fill[red!80!black] (12.0000,8.0000) circle (1.2pt);
\fill[red!80!black] (0.0000,8.0000) circle (1.2pt);
\fill[red!80!black] (12.0000,8.0000) circle (1.2pt);
\fill[red!80!black] (0.0000,8.0000) circle (1.2pt);
\fill[red!80!black] (12.0000,8.0000) circle (1.2pt);
\fill[red!80!black] (0.0000,8.0000) circle (1.2pt);
\fill[red!80!black] (12.0000,8.0000) circle (1.2pt);
\fill[red!80!black] (0.0000,8.0000) circle (1.2pt);
\fill[red!80!black] (12.0000,8.0000) circle (1.2pt);
\fill[red!80!black] (0.0000,8.0000) circle (1.2pt);
\fill[red!80!black] (12.0000,8.0000) circle (1.2pt);
\fill[red!80!black] (0.0000,8.0000) circle (1.2pt);
\fill[red!80!black] (12.0000,8.0000) circle (1.2pt);
\fill[red!80!black] (0.0000,8.0000) circle (1.2pt);
\fill[red!80!black] (12.0000,8.0000) circle (1.2pt);
\fill[red!80!black] (0.0000,8.0000) circle (1.2pt);
\fill[red!80!black] (12.0000,8.0000) circle (1.2pt);
\fill[red!80!black] (0.0000,8.0000) circle (1.2pt);
\fill[red!80!black] (12.0000,8.0000) circle (1.2pt);
\fill[red!80!black] (0.0000,8.0000) circle (1.2pt);
\fill[red!80!black] (12.0000,8.0000) circle (1.2pt);
\fill[red!80!black] (0.0000,8.0000) circle (1.2pt);
\fill[red!80!black] (12.0000,8.0000) circle (1.2pt);
\fill[red!80!black] (0.0000,8.0000) circle (1.2pt);
\fill[red!80!black] (12.0000,8.0000) circle (1.2pt);
\fill[red!80!black] (0.0000,8.0000) circle (1.2pt);
\fill[red!80!black] (12.0000,8.0000) circle (1.2pt);
\fill[red!80!black] (0.0000,8.0000) circle (1.2pt);
\fill[red!80!black] (12.0000,8.0000) circle (1.2pt);
\fill[red!80!black] (0.0000,8.0000) circle (1.2pt);
\fill[red!80!black] (12.0000,8.0000) circle (1.2pt);
\fill[red!80!black] (0.0000,8.0000) circle (1.2pt);
\fill[red!80!black] (12.0000,8.0000) circle (1.2pt);
\fill[red!80!black] (0.0000,8.0000) circle (1.2pt);
\fill[red!80!black] (12.0000,8.0000) circle (1.2pt);
\fill[red!80!black] (0.0000,8.0000) circle (1.2pt);
\fill[red!80!black] (12.0000,8.0000) circle (1.2pt);
\fill[red!80!black] (0.0000,8.0000) circle (1.2pt);
\fill[red!80!black] (12.0000,8.0000) circle (1.2pt);
\fill[red!80!black] (0.0000,8.0000) circle (1.2pt);
\fill[red!80!black] (12.0000,8.0000) circle (1.2pt);
\fill[red!80!black] (0.0000,8.0000) circle (1.2pt);
\fill[red!80!black] (12.0000,8.0000) circle (1.2pt);
\fill[red!80!black] (0.0000,8.0000) circle (1.2pt);
\fill[red!80!black] (12.0000,8.0000) circle (1.2pt);
\fill[red!80!black] (0.0000,8.0000) circle (1.2pt);
\fill[red!80!black] (12.0000,8.0000) circle (1.2pt);
\fill[red!80!black] (0.0000,8.0000) circle (1.2pt);
\fill[red!80!black] (12.0000,8.0000) circle (1.2pt);
\fill[red!80!black] (0.0000,8.0000) circle (1.2pt);
\fill[red!80!black] (12.0000,8.0000) circle (1.2pt);
\fill[red!80!black] (0.0000,8.0000) circle (1.2pt);
\fill[red!80!black] (12.0000,8.0000) circle (1.2pt);
\fill[red!80!black] (0.0000,8.0000) circle (1.2pt);
\fill[red!80!black] (12.0000,8.0000) circle (1.2pt);
\fill[red!80!black] (0.0000,8.0000) circle (1.2pt);
\fill[red!80!black] (12.0000,8.0000) circle (1.2pt);
\fill[red!80!black] (0.0000,8.0000) circle (1.2pt);
\fill[red!80!black] (12.0000,8.0000) circle (1.2pt);
\fill[red!80!black] (0.0000,8.0000) circle (1.2pt);
\fill[red!80!black] (12.0000,8.0000) circle (1.2pt);
\fill[red!80!black] (0.0000,8.0000) circle (1.2pt);
\fill[red!80!black] (12.0000,8.0000) circle (1.2pt);
\fill[red!80!black] (0.0000,8.0000) circle (1.2pt);
\fill[red!80!black] (12.0000,8.0000) circle (1.2pt);
\fill[red!80!black] (0.0000,8.0000) circle (1.2pt);
\fill[red!80!black] (12.0000,8.0000) circle (1.2pt);
\fill[red!80!black] (0.0000,8.0000) circle (1.2pt);
\fill[red!80!black] (12.0000,8.0000) circle (1.2pt);
\fill[red!80!black] (0.0000,8.0000) circle (1.2pt);
\fill[red!80!black] (12.0000,8.0000) circle (1.2pt);
\fill[red!80!black] (0.0000,8.0000) circle (1.2pt);
\fill[red!80!black] (12.0000,8.0000) circle (1.2pt);
\fill[red!80!black] (0.0000,8.0000) circle (1.2pt);
\fill[red!80!black] (12.0000,8.0000) circle (1.2pt);
\fill[red!80!black] (0.0000,8.0000) circle (1.2pt);
\fill[red!80!black] (12.0000,8.0000) circle (1.2pt);
\fill[red!80!black] (0.0000,8.0000) circle (1.2pt);
\fill[red!80!black] (12.0000,8.0000) circle (1.2pt);
\fill[red!80!black] (0.0000,8.0000) circle (1.2pt);
\fill[red!80!black] (12.0000,8.0000) circle (1.2pt);
\fill[red!80!black] (0.0000,8.0000) circle (1.2pt);
\fill[red!80!black] (12.0000,8.0000) circle (1.2pt);
\fill[red!80!black] (0.0000,8.0000) circle (1.2pt);
\fill[red!80!black] (12.0000,8.0000) circle (1.2pt);
\fill[red!80!black] (0.0000,8.0000) circle (1.2pt);
\fill[red!80!black] (12.0000,8.0000) circle (1.2pt);
\fill[red!80!black] (0.0000,8.0000) circle (1.2pt);
\fill[red!80!black] (12.0000,8.0000) circle (1.2pt);
\fill[red!80!black] (0.0000,8.0000) circle (1.2pt);
\fill[red!80!black] (12.0000,8.0000) circle (1.2pt);
\fill[red!80!black] (0.0000,8.0000) circle (1.2pt);
\fill[red!80!black] (12.0000,8.0000) circle (1.2pt);
\fill[red!80!black] (0.0000,8.0000) circle (1.2pt);
\fill[red!80!black] (12.0000,8.0000) circle (1.2pt);
\fill[red!80!black] (0.0000,8.0000) circle (1.2pt);
\fill[red!80!black] (12.0000,8.0000) circle (1.2pt);
\fill[red!80!black] (0.0000,8.0000) circle (1.2pt);
\fill[red!80!black] (12.0000,8.0000) circle (1.2pt);
\fill[red!80!black] (0.0000,8.0000) circle (1.2pt);
\fill[red!80!black] (12.0000,8.0000) circle (1.2pt);
\fill[red!80!black] (0.0000,8.0000) circle (1.2pt);
\fill[red!80!black] (12.0000,8.0000) circle (1.2pt);
\fill[red!80!black] (0.0000,8.0000) circle (1.2pt);
\fill[red!80!black] (12.0000,8.0000) circle (1.2pt);
\fill[red!80!black] (0.0000,8.0000) circle (1.2pt);
\fill[red!80!black] (12.0000,8.0000) circle (1.2pt);
\fill[red!80!black] (0.0000,8.0000) circle (1.2pt);
\fill[red!80!black] (12.0000,8.0000) circle (1.2pt);
\fill[red!80!black] (0.0000,8.0000) circle (1.2pt);
\fill[red!80!black] (12.0000,8.0000) circle (1.2pt);
\fill[red!80!black] (0.0000,8.0000) circle (1.2pt);
\fill[red!80!black] (12.0000,8.0000) circle (1.2pt);
\fill[red!80!black] (0.0000,8.0000) circle (1.2pt);
\fill[red!80!black] (12.0000,8.0000) circle (1.2pt);
\fill[red!80!black] (0.0000,8.0000) circle (1.2pt);
\fill[red!80!black] (12.0000,8.0000) circle (1.2pt);
\fill[red!80!black] (0.0000,8.0000) circle (1.2pt);
\fill[red!80!black] (12.0000,8.0000) circle (1.2pt);
\fill[red!80!black] (0.0000,8.0000) circle (1.2pt);
\fill[red!80!black] (12.0000,8.0000) circle (1.2pt);
\fill[red!80!black] (0.0000,8.0000) circle (1.2pt);
\fill[red!80!black] (12.0000,8.0000) circle (1.2pt);
\fill[red!80!black] (0.0000,8.0000) circle (1.2pt);
\fill[red!80!black] (12.0000,8.0000) circle (1.2pt);
\fill[red!80!black] (0.0000,8.0000) circle (1.2pt);
\fill[red!80!black] (12.0000,8.0000) circle (1.2pt);
\fill[red!80!black] (0.0000,8.0000) circle (1.2pt);
\fill[red!80!black] (12.0000,8.0000) circle (1.2pt);
\fill[red!80!black] (0.0000,8.0000) circle (1.2pt);
\fill[red!80!black] (12.0000,8.0000) circle (1.2pt);
\fill[red!80!black] (0.0000,8.0000) circle (1.2pt);
\fill[red!80!black] (12.0000,8.0000) circle (1.2pt);
\fill[red!80!black] (0.0000,8.0000) circle (1.2pt);
\fill[red!80!black] (12.0000,8.0000) circle (1.2pt);
\fill[red!80!black] (0.0000,8.0000) circle (1.2pt);
\fill[red!80!black] (12.0000,8.0000) circle (1.2pt);
\fill[red!80!black] (0.0000,8.0000) circle (1.2pt);
\fill[red!80!black] (12.0000,8.0000) circle (1.2pt);
\fill[red!80!black] (0.0000,8.0000) circle (1.2pt);
\fill[red!80!black] (12.0000,8.0000) circle (1.2pt);
\fill[red!80!black] (0.0000,8.0000) circle (1.2pt);
\fill[red!80!black] (12.0000,8.0000) circle (1.2pt);
\fill[red!80!black] (0.0000,8.0000) circle (1.2pt);
\fill[red!80!black] (12.0000,8.0000) circle (1.2pt);
\fill[red!80!black] (0.0000,8.0000) circle (1.2pt);
\fill[red!80!black] (12.0000,8.0000) circle (1.2pt);
\fill[red!80!black] (0.0000,8.0000) circle (1.2pt);
\fill[red!80!black] (12.0000,8.0000) circle (1.2pt);
\fill[red!80!black] (0.0000,8.0000) circle (1.2pt);
\fill[red!80!black] (12.0000,8.0000) circle (1.2pt);
\fill[red!80!black] (0.0000,8.0000) circle (1.2pt);
\fill[red!80!black] (12.0000,8.0000) circle (1.2pt);
\fill[red!80!black] (0.0000,8.0000) circle (1.2pt);
\fill[red!80!black] (12.0000,8.0000) circle (1.2pt);
\fill[red!80!black] (0.0000,8.0000) circle (1.2pt);
\fill[red!80!black] (12.0000,8.0000) circle (1.2pt);
\fill[red!80!black] (0.0000,8.0000) circle (1.2pt);
\fill[red!80!black] (12.0000,8.0000) circle (1.2pt);
\fill[red!80!black] (0.0000,8.0000) circle (1.2pt);
\fill[red!80!black] (12.0000,8.0000) circle (1.2pt);
\fill[red!80!black] (0.0000,8.0000) circle (1.2pt);
\fill[red!80!black] (12.0000,8.0000) circle (1.2pt);
\fill[red!80!black] (0.0000,8.0000) circle (1.2pt);
\fill[red!80!black] (12.0000,8.0000) circle (1.2pt);
\fill[red!80!black] (0.0000,8.0000) circle (1.2pt);
\fill[red!80!black] (12.0000,8.0000) circle (1.2pt);
\fill[red!80!black] (0.0000,8.0000) circle (1.2pt);
\fill[red!80!black] (12.0000,8.0000) circle (1.2pt);
\fill[red!80!black] (0.0000,8.0000) circle (1.2pt);
\fill[red!80!black] (12.0000,8.0000) circle (1.2pt);
\fill[red!80!black] (0.0000,8.0000) circle (1.2pt);
\fill[red!80!black] (12.0000,8.0000) circle (1.2pt);
\fill[red!80!black] (0.0000,8.0000) circle (1.2pt);
\fill[red!80!black] (12.0000,8.0000) circle (1.2pt);
\fill[red!80!black] (0.0000,8.0000) circle (1.2pt);
\fill[red!80!black] (12.0000,8.0000) circle (1.2pt);
\fill[red!80!black] (0.0000,8.0000) circle (1.2pt);
\fill[red!80!black] (12.0000,8.0000) circle (1.2pt);
\fill[red!80!black] (0.0000,8.0000) circle (1.2pt);
\fill[red!80!black] (12.0000,8.0000) circle (1.2pt);
\fill[red!80!black] (0.0000,8.0000) circle (1.2pt);
\fill[red!80!black] (12.0000,8.0000) circle (1.2pt);
\fill[red!80!black] (0.0000,8.0000) circle (1.2pt);
\fill[red!80!black] (12.0000,8.0000) circle (1.2pt);
\fill[red!80!black] (0.0000,8.0000) circle (1.2pt);
\fill[red!80!black] (12.0000,8.0000) circle (1.2pt);
\fill[red!80!black] (0.0000,8.0000) circle (1.2pt);
\fill[red!80!black] (12.0000,8.0000) circle (1.2pt);
\fill[red!80!black] (0.0000,8.0000) circle (1.2pt);
\fill[red!80!black] (12.0000,8.0000) circle (1.2pt);
\fill[red!80!black] (0.0000,8.0000) circle (1.2pt);
\fill[red!80!black] (12.0000,8.0000) circle (1.2pt);
\fill[red!80!black] (0.0000,8.0000) circle (1.2pt);
\fill[red!80!black] (12.0000,8.0000) circle (1.2pt);
\fill[red!80!black] (0.0000,8.0000) circle (1.2pt);
\fill[red!80!black] (12.0000,8.0000) circle (1.2pt);
\fill[red!80!black] (0.0000,8.0000) circle (1.2pt);
\fill[red!80!black] (12.0000,8.0000) circle (1.2pt);
\fill[red!80!black] (0.0000,8.0000) circle (1.2pt);
\fill[red!80!black] (12.0000,8.0000) circle (1.2pt);
\fill[red!80!black] (0.0000,8.0000) circle (1.2pt);
\fill[red!80!black] (12.0000,8.0000) circle (1.2pt);
\fill[red!80!black] (0.0000,8.0000) circle (1.2pt);
\fill[red!80!black] (12.0000,8.0000) circle (1.2pt);
\fill[red!80!black] (0.0000,8.0000) circle (1.2pt);
\fill[red!80!black] (12.0000,8.0000) circle (1.2pt);
\fill[red!80!black] (0.0000,8.0000) circle (1.2pt);
\fill[red!80!black] (12.0000,8.0000) circle (1.2pt);
\fill[red!80!black] (0.0000,8.0000) circle (1.2pt);
\fill[red!80!black] (12.0000,8.0000) circle (1.2pt);
\fill[red!80!black] (0.0000,8.0000) circle (1.2pt);
\fill[red!80!black] (12.0000,8.0000) circle (1.2pt);
\fill[red!80!black] (0.0000,8.0000) circle (1.2pt);
\fill[red!80!black] (12.0000,8.0000) circle (1.2pt);
\fill[red!80!black] (0.0000,8.0000) circle (1.2pt);
\fill[red!80!black] (12.0000,8.0000) circle (1.2pt);
\fill[red!80!black] (0.0000,8.0000) circle (1.2pt);
\fill[red!80!black] (12.0000,8.0000) circle (1.2pt);
\fill[red!80!black] (0.0000,8.0000) circle (1.2pt);
\fill[red!80!black] (12.0000,8.0000) circle (1.2pt);
\fill[red!80!black] (0.0000,8.0000) circle (1.2pt);
\fill[red!80!black] (12.0000,8.0000) circle (1.2pt);
\fill[red!80!black] (0.0000,8.0000) circle (1.2pt);
\fill[red!80!black] (12.0000,8.0000) circle (1.2pt);
\fill[red!80!black] (0.0000,8.0000) circle (1.2pt);
\fill[red!80!black] (12.0000,8.0000) circle (1.2pt);
\fill[red!80!black] (0.0000,8.0000) circle (1.2pt);
\fill[red!80!black] (12.0000,8.0000) circle (1.2pt);
\fill[red!80!black] (0.0000,8.0000) circle (1.2pt);
\fill[red!80!black] (12.0000,8.0000) circle (1.2pt);
\fill[red!80!black] (0.0000,8.0000) circle (1.2pt);
\fill[red!80!black] (12.0000,8.0000) circle (1.2pt);
\fill[red!80!black] (0.0000,8.0000) circle (1.2pt);
\fill[red!80!black] (12.0000,8.0000) circle (1.2pt);
\fill[red!80!black] (0.0000,8.0000) circle (1.2pt);
\fill[red!80!black] (12.0000,8.0000) circle (1.2pt);
\fill[red!80!black] (0.0000,8.0000) circle (1.2pt);
\fill[red!80!black] (12.0000,8.0000) circle (1.2pt);
\fill[red!80!black] (0.0000,8.0000) circle (1.2pt);
\fill[red!80!black] (12.0000,8.0000) circle (1.2pt);
\fill[red!80!black] (0.0000,8.0000) circle (1.2pt);
\fill[red!80!black] (12.0000,8.0000) circle (1.2pt);
\fill[red!80!black] (0.0000,8.0000) circle (1.2pt);
\fill[red!80!black] (12.0000,8.0000) circle (1.2pt);
\fill[red!80!black] (0.0000,8.0000) circle (1.2pt);
\fill[red!80!black] (12.0000,8.0000) circle (1.2pt);
\fill[red!80!black] (0.0000,8.0000) circle (1.2pt);
\fill[red!80!black] (12.0000,8.0000) circle (1.2pt);
\fill[red!80!black] (0.0000,8.0000) circle (1.2pt);
\fill[red!80!black] (12.0000,8.0000) circle (1.2pt);
\fill[red!80!black] (0.0000,8.0000) circle (1.2pt);
\fill[red!80!black] (12.0000,8.0000) circle (1.2pt);
\fill[red!80!black] (0.0000,8.0000) circle (1.2pt);
\fill[red!80!black] (12.0000,8.0000) circle (1.2pt);
\fill[red!80!black] (0.0000,8.0000) circle (1.2pt);
\fill[red!80!black] (12.0000,8.0000) circle (1.2pt);
\fill[red!80!black] (0.0000,8.0000) circle (1.2pt);
\fill[red!80!black] (12.0000,8.0000) circle (1.2pt);
\fill[red!80!black] (0.0000,8.0000) circle (1.2pt);
\fill[red!80!black] (12.0000,8.0000) circle (1.2pt);
\fill[red!80!black] (0.0000,8.0000) circle (1.2pt);
\fill[red!80!black] (12.0000,8.0000) circle (1.2pt);
\fill[red!80!black] (0.0000,8.0000) circle (1.2pt);
\fill[red!80!black] (12.0000,8.0000) circle (1.2pt);
\fill[red!80!black] (0.0000,8.0000) circle (1.2pt);
\fill[red!80!black] (12.0000,8.0000) circle (1.2pt);
\fill[red!80!black] (0.0000,8.0000) circle (1.2pt);
\fill[red!80!black] (12.0000,8.0000) circle (1.2pt);
\fill[red!80!black] (0.0000,8.0000) circle (1.2pt);
\fill[red!80!black] (12.0000,8.0000) circle (1.2pt);
\fill[red!80!black] (0.0000,8.0000) circle (1.2pt);
\fill[red!80!black] (12.0000,8.0000) circle (1.2pt);
\fill[red!80!black] (0.0000,8.0000) circle (1.2pt);
\fill[red!80!black] (12.0000,8.0000) circle (1.2pt);
\fill[red!80!black] (0.0000,8.0000) circle (1.2pt);
\fill[red!80!black] (12.0000,8.0000) circle (1.2pt);
\fill[red!80!black] (0.0000,8.0000) circle (1.2pt);
\fill[red!80!black] (12.0000,8.0000) circle (1.2pt);
\fill[red!80!black] (0.0000,8.0000) circle (1.2pt);
\fill[red!80!black] (12.0000,8.0000) circle (1.2pt);
\fill[red!80!black] (0.0000,8.0000) circle (1.2pt);
\fill[red!80!black] (12.0000,8.0000) circle (1.2pt);
\fill[red!80!black] (0.0000,8.0000) circle (1.2pt);
\fill[red!80!black] (12.0000,8.0000) circle (1.2pt);
\fill[red!80!black] (0.0000,8.0000) circle (1.2pt);
\fill[red!80!black] (12.0000,8.0000) circle (1.2pt);
\fill[red!80!black] (0.0000,8.0000) circle (1.2pt);
\fill[red!80!black] (12.0000,8.0000) circle (1.2pt);
\fill[red!80!black] (0.0000,8.0000) circle (1.2pt);
\fill[red!80!black] (12.0000,8.0000) circle (1.2pt);
\fill[red!80!black] (0.0000,8.0000) circle (1.2pt);
\fill[red!80!black] (12.0000,8.0000) circle (1.2pt);
\fill[red!80!black] (0.0000,8.0000) circle (1.2pt);
\fill[red!80!black] (12.0000,8.0000) circle (1.2pt);
\fill[red!80!black] (0.0000,8.0000) circle (1.2pt);
\fill[red!80!black] (12.0000,8.0000) circle (1.2pt);
\fill[red!80!black] (0.0000,8.0000) circle (1.2pt);
\fill[red!80!black] (12.0000,8.0000) circle (1.2pt);
\fill[red!80!black] (0.0000,8.0000) circle (1.2pt);
\fill[red!80!black] (12.0000,8.0000) circle (1.2pt);
\fill[red!80!black] (0.0000,8.0000) circle (1.2pt);
\fill[red!80!black] (12.0000,8.0000) circle (1.2pt);
\fill[red!80!black] (0.0000,8.0000) circle (1.2pt);
\fill[red!80!black] (12.0000,8.0000) circle (1.2pt);
\fill[red!80!black] (0.0000,8.0000) circle (1.2pt);
\fill[red!80!black] (12.0000,8.0000) circle (1.2pt);
\fill[red!80!black] (0.0000,8.0000) circle (1.2pt);
\fill[red!80!black] (12.0000,8.0000) circle (1.2pt);
\fill[red!80!black] (0.0000,8.0000) circle (1.2pt);
\fill[red!80!black] (12.0000,8.0000) circle (1.2pt);
\fill[red!80!black] (0.0000,8.0000) circle (1.2pt);
\fill[red!80!black] (12.0000,8.0000) circle (1.2pt);
\fill[red!80!black] (0.0000,8.0000) circle (1.2pt);
\fill[red!80!black] (12.0000,8.0000) circle (1.2pt);
\fill[red!80!black] (0.0000,8.0000) circle (1.2pt);
\fill[red!80!black] (12.0000,8.0000) circle (1.2pt);
\fill[red!80!black] (0.0000,8.0000) circle (1.2pt);
\fill[red!80!black] (12.0000,8.0000) circle (1.2pt);
\fill[red!80!black] (0.0000,8.0000) circle (1.2pt);
\fill[red!80!black] (12.0000,8.0000) circle (1.2pt);
\fill[red!80!black] (0.0000,8.0000) circle (1.2pt);
\fill[red!80!black] (12.0000,8.0000) circle (1.2pt);
\fill[red!80!black] (0.0000,8.0000) circle (1.2pt);
\fill[red!80!black] (12.0000,8.0000) circle (1.2pt);
\fill[red!80!black] (0.0000,8.0000) circle (1.2pt);
\fill[red!80!black] (12.0000,8.0000) circle (1.2pt);
\fill[red!80!black] (0.0000,8.0000) circle (1.2pt);
\fill[red!80!black] (12.0000,8.0000) circle (1.2pt);
\fill[red!80!black] (0.0000,8.0000) circle (1.2pt);
\fill[red!80!black] (12.0000,8.0000) circle (1.2pt);
\fill[red!80!black] (0.0000,8.0000) circle (1.2pt);
\fill[red!80!black] (12.0000,8.0000) circle (1.2pt);
\fill[red!80!black] (0.0000,8.0000) circle (1.2pt);
\fill[red!80!black] (12.0000,8.0000) circle (1.2pt);
\fill[red!80!black] (0.0000,8.0000) circle (1.2pt);
\fill[red!80!black] (12.0000,8.0000) circle (1.2pt);
\fill[red!80!black] (0.0000,8.0000) circle (1.2pt);
\fill[red!80!black] (12.0000,8.0000) circle (1.2pt);
\fill[red!80!black] (0.0000,8.0000) circle (1.2pt);
\fill[red!80!black] (12.0000,8.0000) circle (1.2pt);
\fill[red!80!black] (0.0000,8.0000) circle (1.2pt);
\fill[red!80!black] (12.0000,8.0000) circle (1.2pt);
\fill[red!80!black] (0.0000,8.0000) circle (1.2pt);
\fill[red!80!black] (12.0000,8.0000) circle (1.2pt);
\fill[red!80!black] (0.0000,8.0000) circle (1.2pt);
\fill[red!80!black] (12.0000,8.0000) circle (1.2pt);
\fill[red!80!black] (0.0000,8.0000) circle (1.2pt);
\fill[red!80!black] (12.0000,8.0000) circle (1.2pt);
\fill[red!80!black] (0.0000,8.0000) circle (1.2pt);
\fill[red!80!black] (12.0000,8.0000) circle (1.2pt);
\fill[red!80!black] (0.0000,8.0000) circle (1.2pt);
\fill[red!80!black] (12.0000,8.0000) circle (1.2pt);
\fill[red!80!black] (0.0000,8.0000) circle (1.2pt);
\fill[red!80!black] (12.0000,8.0000) circle (1.2pt);
\fill[red!80!black] (0.0000,8.0000) circle (1.2pt);
\fill[red!80!black] (12.0000,8.0000) circle (1.2pt);
\fill[red!80!black] (0.0000,8.0000) circle (1.2pt);
\fill[red!80!black] (12.0000,8.0000) circle (1.2pt);
\fill[red!80!black] (0.0000,8.0000) circle (1.2pt);
\fill[red!80!black] (12.0000,8.0000) circle (1.2pt);
\fill[red!80!black] (0.0000,8.0000) circle (1.2pt);
\fill[red!80!black] (12.0000,8.0000) circle (1.2pt);
\fill[red!80!black] (0.0000,8.0000) circle (1.2pt);
\fill[red!80!black] (12.0000,8.0000) circle (1.2pt);
\fill[red!80!black] (0.0000,8.0000) circle (1.2pt);
\fill[red!80!black] (12.0000,8.0000) circle (1.2pt);
\fill[red!80!black] (0.0000,8.0000) circle (1.2pt);
\fill[red!80!black] (12.0000,8.0000) circle (1.2pt);
\fill[red!80!black] (0.0000,8.0000) circle (1.2pt);
\fill[red!80!black] (12.0000,8.0000) circle (1.2pt);
\fill[red!80!black] (0.0000,8.0000) circle (1.2pt);
\fill[red!80!black] (12.0000,8.0000) circle (1.2pt);
\fill[red!80!black] (0.0000,8.0000) circle (1.2pt);
\fill[red!80!black] (12.0000,8.0000) circle (1.2pt);
\fill[red!80!black] (0.0000,8.0000) circle (1.2pt);
\fill[red!80!black] (12.0000,8.0000) circle (1.2pt);
\fill[red!80!black] (0.0000,8.0000) circle (1.2pt);
\fill[red!80!black] (12.0000,8.0000) circle (1.2pt);
\fill[red!80!black] (0.0000,8.0000) circle (1.2pt);
\fill[red!80!black] (12.0000,8.0000) circle (1.2pt);
\fill[red!80!black] (0.0000,8.0000) circle (1.2pt);
\fill[red!80!black] (12.0000,8.0000) circle (1.2pt);
\fill[red!80!black] (0.0000,8.0000) circle (1.2pt);
\fill[red!80!black] (12.0000,8.0000) circle (1.2pt);
\fill[red!80!black] (0.0000,8.0000) circle (1.2pt);
\fill[red!80!black] (12.0000,8.0000) circle (1.2pt);
\fill[red!80!black] (0.0000,8.0000) circle (1.2pt);
\fill[red!80!black] (12.0000,8.0000) circle (1.2pt);
\fill[red!80!black] (0.0000,8.0000) circle (1.2pt);
\fill[red!80!black] (12.0000,8.0000) circle (1.2pt);
\fill[red!80!black] (0.0000,8.0000) circle (1.2pt);
\fill[red!80!black] (12.0000,8.0000) circle (1.2pt);
\fill[red!80!black] (0.0000,8.0000) circle (1.2pt);
\fill[red!80!black] (12.0000,8.0000) circle (1.2pt);
\fill[red!80!black] (0.0000,8.0000) circle (1.2pt);
\fill[red!80!black] (12.0000,8.0000) circle (1.2pt);
\fill[red!80!black] (0.0000,8.0000) circle (1.2pt);
\fill[red!80!black] (12.0000,8.0000) circle (1.2pt);
\fill[red!80!black] (0.0000,8.0000) circle (1.2pt);
\fill[red!80!black] (12.0000,8.0000) circle (1.2pt);
\fill[red!80!black] (0.0000,8.0000) circle (1.2pt);
\fill[red!80!black] (12.0000,8.0000) circle (1.2pt);
\fill[red!80!black] (0.0000,8.0000) circle (1.2pt);
\fill[red!80!black] (12.0000,8.0000) circle (1.2pt);
\fill[red!80!black] (0.0000,8.0000) circle (1.2pt);
\fill[red!80!black] (12.0000,8.0000) circle (1.2pt);
\fill[red!80!black] (0.0000,8.0000) circle (1.2pt);
\fill[red!80!black] (12.0000,8.0000) circle (1.2pt);
\fill[red!80!black] (0.0000,8.0000) circle (1.2pt);
\fill[red!80!black] (12.0000,8.0000) circle (1.2pt);
\fill[red!80!black] (0.0000,8.0000) circle (1.2pt);
\fill[red!80!black] (12.0000,8.0000) circle (1.2pt);
\fill[red!80!black] (0.0000,8.0000) circle (1.2pt);
\fill[red!80!black] (12.0000,8.0000) circle (1.2pt);
\fill[red!80!black] (0.0000,8.0000) circle (1.2pt);
\fill[red!80!black] (12.0000,8.0000) circle (1.2pt);
\fill[red!80!black] (0.0000,8.0000) circle (1.2pt);
\fill[red!80!black] (12.0000,8.0000) circle (1.2pt);
\fill[red!80!black] (0.0000,8.0000) circle (1.2pt);
\fill[red!80!black] (12.0000,8.0000) circle (1.2pt);
\fill[red!80!black] (0.0000,8.0000) circle (1.2pt);
\fill[red!80!black] (12.0000,8.0000) circle (1.2pt);
\fill[red!80!black] (0.0000,8.0000) circle (1.2pt);
\fill[red!80!black] (12.0000,8.0000) circle (1.2pt);
\fill[red!80!black] (0.0000,8.0000) circle (1.2pt);
\fill[red!80!black] (12.0000,8.0000) circle (1.2pt);
\fill[red!80!black] (0.0000,8.0000) circle (1.2pt);
\fill[red!80!black] (12.0000,8.0000) circle (1.2pt);
\fill[red!80!black] (0.0000,8.0000) circle (1.2pt);
\fill[red!80!black] (12.0000,8.0000) circle (1.2pt);
\fill[red!80!black] (0.0000,8.0000) circle (1.2pt);
\fill[red!80!black] (12.0000,8.0000) circle (1.2pt);
\fill[red!80!black] (0.0000,8.0000) circle (1.2pt);
\fill[red!80!black] (12.0000,8.0000) circle (1.2pt);
\fill[red!80!black] (0.0000,8.0000) circle (1.2pt);
\fill[red!80!black] (12.0000,8.0000) circle (1.2pt);
\fill[red!80!black] (0.0000,8.0000) circle (1.2pt);
\fill[red!80!black] (12.0000,8.0000) circle (1.2pt);
\fill[red!80!black] (0.0000,8.0000) circle (1.2pt);
\fill[red!80!black] (12.0000,8.0000) circle (1.2pt);
\fill[red!80!black] (0.0000,8.0000) circle (1.2pt);
\fill[red!80!black] (12.0000,8.0000) circle (1.2pt);
\fill[red!80!black] (0.0000,8.0000) circle (1.2pt);
\fill[red!80!black] (12.0000,8.0000) circle (1.2pt);
\fill[red!80!black] (0.0000,8.0000) circle (1.2pt);
\fill[red!80!black] (12.0000,8.0000) circle (1.2pt);
\fill[red!80!black] (0.0000,8.0000) circle (1.2pt);
\fill[red!80!black] (12.0000,8.0000) circle (1.2pt);
\fill[red!80!black] (0.0000,8.0000) circle (1.2pt);
\fill[red!80!black] (12.0000,8.0000) circle (1.2pt);
\fill[red!80!black] (0.0000,8.0000) circle (1.2pt);
\fill[red!80!black] (12.0000,8.0000) circle (1.2pt);
\fill[red!80!black] (0.0000,8.0000) circle (1.2pt);
\fill[red!80!black] (12.0000,8.0000) circle (1.2pt);
\fill[red!80!black] (0.0000,8.0000) circle (1.2pt);
\fill[red!80!black] (12.0000,8.0000) circle (1.2pt);
\fill[red!80!black] (0.0000,8.0000) circle (1.2pt);
\fill[red!80!black] (12.0000,8.0000) circle (1.2pt);
\fill[red!80!black] (0.0000,8.0000) circle (1.2pt);
\fill[red!80!black] (12.0000,8.0000) circle (1.2pt);
\fill[red!80!black] (0.0000,8.0000) circle (1.2pt);
\fill[red!80!black] (12.0000,8.0000) circle (1.2pt);
\fill[red!80!black] (0.0000,8.0000) circle (1.2pt);
\fill[red!80!black] (12.0000,8.0000) circle (1.2pt);
\fill[red!80!black] (0.0000,8.0000) circle (1.2pt);
\fill[red!80!black] (12.0000,8.0000) circle (1.2pt);
\fill[red!80!black] (0.0000,8.0000) circle (1.2pt);
\fill[red!80!black] (12.0000,8.0000) circle (1.2pt);
\fill[red!80!black] (0.0000,8.0000) circle (1.2pt);
\fill[red!80!black] (12.0000,8.0000) circle (1.2pt);
\fill[red!80!black] (0.0000,8.0000) circle (1.2pt);
\fill[red!80!black] (12.0000,8.0000) circle (1.2pt);
\fill[red!80!black] (0.0000,8.0000) circle (1.2pt);
\fill[red!80!black] (12.0000,8.0000) circle (1.2pt);
\fill[red!80!black] (0.0000,8.0000) circle (1.2pt);
\fill[red!80!black] (12.0000,8.0000) circle (1.2pt);
\fill[red!80!black] (0.0000,8.0000) circle (1.2pt);
\fill[red!80!black] (12.0000,8.0000) circle (1.2pt);
\fill[red!80!black] (0.0000,8.0000) circle (1.2pt);
\fill[red!80!black] (12.0000,8.0000) circle (1.2pt);
\fill[red!80!black] (0.0000,8.0000) circle (1.2pt);
\fill[red!80!black] (12.0000,8.0000) circle (1.2pt);
\fill[red!80!black] (0.0000,8.0000) circle (1.2pt);
\fill[red!80!black] (12.0000,8.0000) circle (1.2pt);
\fill[red!80!black] (0.0000,8.0000) circle (1.2pt);
\fill[red!80!black] (12.0000,8.0000) circle (1.2pt);
\fill[red!80!black] (0.0000,8.0000) circle (1.2pt);
\fill[red!80!black] (12.0000,8.0000) circle (1.2pt);
\fill[red!80!black] (0.0000,8.0000) circle (1.2pt);
\fill[red!80!black] (12.0000,8.0000) circle (1.2pt);
\fill[red!80!black] (0.0000,8.0000) circle (1.2pt);
\fill[red!80!black] (12.0000,8.0000) circle (1.2pt);
\fill[red!80!black] (0.0000,8.0000) circle (1.2pt);
\fill[red!80!black] (12.0000,8.0000) circle (1.2pt);
\fill[red!80!black] (0.0000,8.0000) circle (1.2pt);
\fill[red!80!black] (12.0000,8.0000) circle (1.2pt);
\fill[red!80!black] (0.0000,8.0000) circle (1.2pt);
\fill[red!80!black] (12.0000,8.0000) circle (1.2pt);
\fill[red!80!black] (0.0000,8.0000) circle (1.2pt);
\fill[red!80!black] (12.0000,8.0000) circle (1.2pt);
\fill[red!80!black] (0.0000,8.0000) circle (1.2pt);
\fill[red!80!black] (12.0000,8.0000) circle (1.2pt);
\fill[red!80!black] (0.0000,8.0000) circle (1.2pt);
\fill[red!80!black] (12.0000,8.0000) circle (1.2pt);
\fill[red!80!black] (0.0000,8.0000) circle (1.2pt);
\fill[red!80!black] (12.0000,8.0000) circle (1.2pt);
\fill[red!80!black] (0.0000,8.0000) circle (1.2pt);
\fill[red!80!black] (12.0000,8.0000) circle (1.2pt);
\fill[red!80!black] (0.0000,8.0000) circle (1.2pt);
\fill[red!80!black] (12.0000,8.0000) circle (1.2pt);
\fill[red!80!black] (0.0000,8.0000) circle (1.2pt);
\fill[red!80!black] (12.0000,8.0000) circle (1.2pt);
\fill[red!80!black] (0.0000,8.0000) circle (1.2pt);
\fill[red!80!black] (12.0000,8.0000) circle (1.2pt);
\fill[red!80!black] (0.0000,8.0000) circle (1.2pt);
\fill[red!80!black] (12.0000,8.0000) circle (1.2pt);
\fill[red!80!black] (0.0000,8.0000) circle (1.2pt);
\fill[red!80!black] (12.0000,8.0000) circle (1.2pt);
\fill[red!80!black] (0.0000,8.0000) circle (1.2pt);
\fill[red!80!black] (12.0000,8.0000) circle (1.2pt);
\fill[red!80!black] (0.0000,8.0000) circle (1.2pt);
\fill[red!80!black] (12.0000,8.0000) circle (1.2pt);
\fill[red!80!black] (0.0000,8.0000) circle (1.2pt);
\fill[red!80!black] (12.0000,8.0000) circle (1.2pt);
\fill[red!80!black] (0.0000,8.0000) circle (1.2pt);
\fill[red!80!black] (12.0000,8.0000) circle (1.2pt);
\fill[red!80!black] (0.0000,8.0000) circle (1.2pt);
\fill[red!80!black] (12.0000,8.0000) circle (1.2pt);
\fill[red!80!black] (0.0000,8.0000) circle (1.2pt);
\fill[red!80!black] (12.0000,8.0000) circle (1.2pt);
\fill[red!80!black] (0.0000,8.0000) circle (1.2pt);
\fill[red!80!black] (12.0000,8.0000) circle (1.2pt);
\fill[red!80!black] (0.0000,8.0000) circle (1.2pt);
\fill[red!80!black] (12.0000,8.0000) circle (1.2pt);
\fill[red!80!black] (0.0000,8.0000) circle (1.2pt);
\fill[red!80!black] (12.0000,8.0000) circle (1.2pt);
\fill[red!80!black] (0.0000,8.0000) circle (1.2pt);
\fill[red!80!black] (12.0000,8.0000) circle (1.2pt);
\fill[red!80!black] (0.0000,8.0000) circle (1.2pt);
\fill[red!80!black] (12.0000,8.0000) circle (1.2pt);
\fill[red!80!black] (0.0000,8.0000) circle (1.2pt);
\fill[red!80!black] (12.0000,8.0000) circle (1.2pt);
\fill[red!80!black] (0.0000,8.0000) circle (1.2pt);
\fill[red!80!black] (12.0000,8.0000) circle (1.2pt);
\fill[red!80!black] (0.0000,8.0000) circle (1.2pt);
\fill[red!80!black] (12.0000,8.0000) circle (1.2pt);
\fill[red!80!black] (0.0000,8.0000) circle (1.2pt);
\fill[red!80!black] (12.0000,8.0000) circle (1.2pt);
\fill[red!80!black] (0.0000,8.0000) circle (1.2pt);
\fill[red!80!black] (12.0000,8.0000) circle (1.2pt);
\fill[red!80!black] (0.0000,8.0000) circle (1.2pt);
\fill[red!80!black] (12.0000,8.0000) circle (1.2pt);
\fill[red!80!black] (0.0000,8.0000) circle (1.2pt);
\fill[red!80!black] (12.0000,8.0000) circle (1.2pt);
\fill[red!80!black] (0.0000,8.0000) circle (1.2pt);
\fill[red!80!black] (12.0000,8.0000) circle (1.2pt);
\fill[red!80!black] (0.0000,8.0000) circle (1.2pt);
\fill[red!80!black] (12.0000,8.0000) circle (1.2pt);
\fill[red!80!black] (0.0000,8.0000) circle (1.2pt);
\fill[red!80!black] (12.0000,8.0000) circle (1.2pt);
\fill[red!80!black] (0.0000,8.0000) circle (1.2pt);
\fill[red!80!black] (12.0000,8.0000) circle (1.2pt);
\fill[red!80!black] (0.0000,8.0000) circle (1.2pt);
\fill[red!80!black] (12.0000,8.0000) circle (1.2pt);
\fill[red!80!black] (0.0000,8.0000) circle (1.2pt);
\fill[red!80!black] (12.0000,8.0000) circle (1.2pt);
\fill[red!80!black] (0.0000,8.0000) circle (1.2pt);
\fill[red!80!black] (12.0000,8.0000) circle (1.2pt);
\fill[red!80!black] (0.0000,8.0000) circle (1.2pt);
\fill[red!80!black] (12.0000,8.0000) circle (1.2pt);
\fill[red!80!black] (0.0000,8.0000) circle (1.2pt);
\fill[red!80!black] (12.0000,8.0000) circle (1.2pt);
\fill[red!80!black] (0.0000,8.0000) circle (1.2pt);
\fill[red!80!black] (12.0000,8.0000) circle (1.2pt);
\fill[red!80!black] (0.0000,8.0000) circle (1.2pt);
\fill[red!80!black] (12.0000,8.0000) circle (1.2pt);
\fill[red!80!black] (0.0000,8.0000) circle (1.2pt);
\fill[red!80!black] (12.0000,8.0000) circle (1.2pt);
\fill[red!80!black] (0.0000,8.0000) circle (1.2pt);
\fill[red!80!black] (12.0000,8.0000) circle (1.2pt);
\fill[red!80!black] (0.0000,8.0000) circle (1.2pt);
\fill[red!80!black] (12.0000,8.0000) circle (1.2pt);
\fill[red!80!black] (0.0000,8.0000) circle (1.2pt);
\fill[red!80!black] (12.0000,8.0000) circle (1.2pt);
\fill[red!80!black] (0.0000,8.0000) circle (1.2pt);
\fill[red!80!black] (12.0000,8.0000) circle (1.2pt);
\fill[red!80!black] (0.0000,8.0000) circle (1.2pt);
\fill[red!80!black] (12.0000,8.0000) circle (1.2pt);
\fill[red!80!black] (0.0000,8.0000) circle (1.2pt);
\fill[red!80!black] (12.0000,8.0000) circle (1.2pt);
\fill[red!80!black] (0.0000,8.0000) circle (1.2pt);
\fill[red!80!black] (12.0000,8.0000) circle (1.2pt);
\fill[red!80!black] (0.0000,8.0000) circle (1.2pt);
\fill[red!80!black] (12.0000,8.0000) circle (1.2pt);
\fill[red!80!black] (0.0000,8.0000) circle (1.2pt);
\fill[red!80!black] (12.0000,8.0000) circle (1.2pt);
\fill[red!80!black] (0.0000,8.0000) circle (1.2pt);
\fill[red!80!black] (12.0000,8.0000) circle (1.2pt);
\fill[red!80!black] (0.0000,8.0000) circle (1.2pt);
\fill[red!80!black] (12.0000,8.0000) circle (1.2pt);
\fill[red!80!black] (0.0000,8.0000) circle (1.2pt);
\fill[red!80!black] (12.0000,8.0000) circle (1.2pt);
\fill[red!80!black] (0.0000,8.0000) circle (1.2pt);
\fill[red!80!black] (12.0000,8.0000) circle (1.2pt);
\fill[red!80!black] (0.0000,8.0000) circle (1.2pt);
\fill[red!80!black] (12.0000,8.0000) circle (1.2pt);
\fill[red!80!black] (0.0000,8.0000) circle (1.2pt);
\fill[red!80!black] (12.0000,8.0000) circle (1.2pt);
\fill[red!80!black] (0.0000,8.0000) circle (1.2pt);
\fill[red!80!black] (12.0000,8.0000) circle (1.2pt);
\fill[red!80!black] (0.0000,8.0000) circle (1.2pt);
\fill[red!80!black] (12.0000,8.0000) circle (1.2pt);
\fill[red!80!black] (0.0000,8.0000) circle (1.2pt);
\fill[red!80!black] (12.0000,8.0000) circle (1.2pt);
\fill[red!80!black] (0.0000,8.0000) circle (1.2pt);
\fill[red!80!black] (12.0000,8.0000) circle (1.2pt);
\fill[red!80!black] (0.0000,8.0000) circle (1.2pt);
\fill[red!80!black] (12.0000,8.0000) circle (1.2pt);
\fill[red!80!black] (0.0000,8.0000) circle (1.2pt);
\fill[red!80!black] (12.0000,8.0000) circle (1.2pt);
\fill[red!80!black] (0.0000,8.0000) circle (1.2pt);
\fill[red!80!black] (12.0000,8.0000) circle (1.2pt);
\fill[red!80!black] (0.0000,8.0000) circle (1.2pt);
\fill[red!80!black] (12.0000,8.0000) circle (1.2pt);
\fill[red!80!black] (0.0000,8.0000) circle (1.2pt);
\fill[red!80!black] (12.0000,8.0000) circle (1.2pt);
\fill[red!80!black] (0.0000,8.0000) circle (1.2pt);
\fill[red!80!black] (12.0000,8.0000) circle (1.2pt);
\fill[red!80!black] (0.0000,8.0000) circle (1.2pt);
\fill[red!80!black] (12.0000,8.0000) circle (1.2pt);
\fill[red!80!black] (0.0000,8.0000) circle (1.2pt);
\fill[red!80!black] (12.0000,8.0000) circle (1.2pt);
\fill[red!80!black] (0.0000,8.0000) circle (1.2pt);
\fill[red!80!black] (12.0000,8.0000) circle (1.2pt);
\fill[red!80!black] (0.0000,8.0000) circle (1.2pt);
\fill[red!80!black] (12.0000,8.0000) circle (1.2pt);
\fill[red!80!black] (0.0000,8.0000) circle (1.2pt);
\fill[red!80!black] (12.0000,8.0000) circle (1.2pt);
\fill[red!80!black] (0.0000,8.0000) circle (1.2pt);
\fill[red!80!black] (12.0000,8.0000) circle (1.2pt);
\fill[red!80!black] (0.0000,8.0000) circle (1.2pt);
\fill[red!80!black] (12.0000,8.0000) circle (1.2pt);
\fill[red!80!black] (0.0000,8.0000) circle (1.2pt);
\fill[red!80!black] (12.0000,8.0000) circle (1.2pt);
\fill[red!80!black] (0.0000,8.0000) circle (1.2pt);
\fill[red!80!black] (12.0000,8.0000) circle (1.2pt);
\fill[red!80!black] (0.0000,8.0000) circle (1.2pt);
\fill[red!80!black] (12.0000,8.0000) circle (1.2pt);
\fill[red!80!black] (0.0000,8.0000) circle (1.2pt);
\fill[red!80!black] (12.0000,8.0000) circle (1.2pt);
\fill[red!80!black] (0.0000,8.0000) circle (1.2pt);
\fill[red!80!black] (12.0000,8.0000) circle (1.2pt);
\fill[red!80!black] (0.0000,8.0000) circle (1.2pt);
\fill[red!80!black] (12.0000,8.0000) circle (1.2pt);
\fill[red!80!black] (0.0000,8.0000) circle (1.2pt);
\fill[red!80!black] (12.0000,8.0000) circle (1.2pt);
\fill[red!80!black] (0.0000,8.0000) circle (1.2pt);
\fill[red!80!black] (12.0000,8.0000) circle (1.2pt);
\fill[red!80!black] (0.0000,8.0000) circle (1.2pt);
\fill[red!80!black] (12.0000,8.0000) circle (1.2pt);
\fill[red!80!black] (0.0000,8.0000) circle (1.2pt);
\fill[red!80!black] (12.0000,8.0000) circle (1.2pt);
\fill[red!80!black] (0.0000,8.0000) circle (1.2pt);
\fill[red!80!black] (12.0000,8.0000) circle (1.2pt);
\fill[red!80!black] (0.0000,8.0000) circle (1.2pt);
\fill[red!80!black] (12.0000,8.0000) circle (1.2pt);
\fill[red!80!black] (0.0000,8.0000) circle (1.2pt);
\fill[red!80!black] (12.0000,8.0000) circle (1.2pt);
\fill[red!80!black] (0.0000,8.0000) circle (1.2pt);
\fill[red!80!black] (12.0000,8.0000) circle (1.2pt);
\fill[red!80!black] (0.0000,8.0000) circle (1.2pt);
\fill[red!80!black] (12.0000,8.0000) circle (1.2pt);
\fill[red!80!black] (0.0000,8.0000) circle (1.2pt);
\fill[red!80!black] (12.0000,8.0000) circle (1.2pt);
\fill[red!80!black] (0.0000,8.0000) circle (1.2pt);
\fill[red!80!black] (12.0000,8.0000) circle (1.2pt);
\fill[red!80!black] (0.0000,8.0000) circle (1.2pt);
\fill[red!80!black] (12.0000,8.0000) circle (1.2pt);
\fill[red!80!black] (0.0000,8.0000) circle (1.2pt);
\fill[red!80!black] (12.0000,8.0000) circle (1.2pt);
\fill[red!80!black] (0.0000,8.0000) circle (1.2pt);
\fill[red!80!black] (12.0000,8.0000) circle (1.2pt);
\fill[red!80!black] (0.0000,8.0000) circle (1.2pt);
\fill[red!80!black] (12.0000,8.0000) circle (1.2pt);
\fill[red!80!black] (0.0000,8.0000) circle (1.2pt);
\fill[red!80!black] (12.0000,8.0000) circle (1.2pt);
\fill[red!80!black] (0.0000,8.0000) circle (1.2pt);
\fill[red!80!black] (12.0000,8.0000) circle (1.2pt);
\fill[red!80!black] (0.0000,8.0000) circle (1.2pt);
\fill[red!80!black] (12.0000,8.0000) circle (1.2pt);
\fill[red!80!black] (0.0000,8.0000) circle (1.2pt);
\fill[red!80!black] (12.0000,8.0000) circle (1.2pt);
\fill[red!80!black] (0.0000,8.0000) circle (1.2pt);
\fill[red!80!black] (12.0000,8.0000) circle (1.2pt);
\fill[red!80!black] (0.0000,8.0000) circle (1.2pt);
\fill[red!80!black] (12.0000,8.0000) circle (1.2pt);
\fill[red!80!black] (0.0000,8.0000) circle (1.2pt);
\fill[red!80!black] (12.0000,8.0000) circle (1.2pt);
\fill[red!80!black] (0.0000,8.0000) circle (1.2pt);
\fill[red!80!black] (12.0000,8.0000) circle (1.2pt);
\fill[red!80!black] (0.0000,8.0000) circle (1.2pt);
\fill[red!80!black] (12.0000,8.0000) circle (1.2pt);
\fill[red!80!black] (0.0000,8.0000) circle (1.2pt);
\fill[red!80!black] (12.0000,8.0000) circle (1.2pt);
\fill[red!80!black] (0.0000,8.0000) circle (1.2pt);
\fill[red!80!black] (12.0000,8.0000) circle (1.2pt);
\fill[red!80!black] (0.0000,8.0000) circle (1.2pt);
\fill[red!80!black] (12.0000,8.0000) circle (1.2pt);
\fill[red!80!black] (0.0000,8.0000) circle (1.2pt);
\fill[red!80!black] (12.0000,8.0000) circle (1.2pt);
\fill[red!80!black] (0.0000,8.0000) circle (1.2pt);
\fill[red!80!black] (12.0000,8.0000) circle (1.2pt);
\fill[red!80!black] (0.0000,8.0000) circle (1.2pt);
\fill[red!80!black] (12.0000,8.0000) circle (1.2pt);
\fill[red!80!black] (0.0000,8.0000) circle (1.2pt);
\fill[red!80!black] (12.0000,8.0000) circle (1.2pt);
\fill[red!80!black] (0.0000,8.0000) circle (1.2pt);
\fill[red!80!black] (12.0000,8.0000) circle (1.2pt);
\fill[red!80!black] (0.0000,8.0000) circle (1.2pt);
\fill[red!80!black] (12.0000,8.0000) circle (1.2pt);
\fill[red!80!black] (0.0000,8.0000) circle (1.2pt);
\fill[red!80!black] (12.0000,8.0000) circle (1.2pt);
\fill[red!80!black] (0.0000,8.0000) circle (1.2pt);
\fill[red!80!black] (12.0000,8.0000) circle (1.2pt);
\fill[red!80!black] (0.0000,8.0000) circle (1.2pt);
\fill[red!80!black] (12.0000,8.0000) circle (1.2pt);
\fill[red!80!black] (0.0000,8.0000) circle (1.2pt);
\fill[red!80!black] (12.0000,8.0000) circle (1.2pt);
\fill[red!80!black] (0.0000,8.0000) circle (1.2pt);
\fill[red!80!black] (12.0000,8.0000) circle (1.2pt);
\fill[red!80!black] (0.0000,8.0000) circle (1.2pt);
\fill[red!80!black] (12.0000,8.0000) circle (1.2pt);
\fill[red!80!black] (0.0000,8.0000) circle (1.2pt);
\fill[red!80!black] (12.0000,8.0000) circle (1.2pt);
\fill[red!80!black] (0.0000,8.0000) circle (1.2pt);
\fill[red!80!black] (12.0000,8.0000) circle (1.2pt);
\fill[red!80!black] (0.0000,8.0000) circle (1.2pt);
\fill[red!80!black] (12.0000,8.0000) circle (1.2pt);
\fill[red!80!black] (0.0000,8.0000) circle (1.2pt);
\fill[red!80!black] (12.0000,8.0000) circle (1.2pt);
\fill[red!80!black] (0.0000,8.0000) circle (1.2pt);
\fill[red!80!black] (12.0000,8.0000) circle (1.2pt);
\fill[red!80!black] (0.0000,8.0000) circle (1.2pt);
\fill[red!80!black] (12.0000,8.0000) circle (1.2pt);
\fill[red!80!black] (0.0000,8.0000) circle (1.2pt);
\fill[red!80!black] (12.0000,8.0000) circle (1.2pt);
\fill[red!80!black] (0.0000,8.0000) circle (1.2pt);
\fill[red!80!black] (12.0000,8.0000) circle (1.2pt);
\fill[red!80!black] (0.0000,8.0000) circle (1.2pt);
\fill[red!80!black] (12.0000,8.0000) circle (1.2pt);
\fill[red!80!black] (0.0000,8.0000) circle (1.2pt);
\fill[red!80!black] (12.0000,8.0000) circle (1.2pt);
\fill[red!80!black] (0.0000,8.0000) circle (1.2pt);
\fill[red!80!black] (12.0000,8.0000) circle (1.2pt);
\fill[red!80!black] (0.0000,8.0000) circle (1.2pt);
\fill[red!80!black] (12.0000,8.0000) circle (1.2pt);
\fill[red!80!black] (0.0000,8.0000) circle (1.2pt);
\fill[red!80!black] (12.0000,8.0000) circle (1.2pt);
\fill[red!80!black] (0.0000,8.0000) circle (1.2pt);
\fill[red!80!black] (12.0000,8.0000) circle (1.2pt);
\fill[red!80!black] (0.0000,8.0000) circle (1.2pt);
\fill[red!80!black] (12.0000,8.0000) circle (1.2pt);
\fill[red!80!black] (0.0000,8.0000) circle (1.2pt);
\fill[red!80!black] (12.0000,8.0000) circle (1.2pt);
\fill[red!80!black] (0.0000,8.0000) circle (1.2pt);
\fill[red!80!black] (12.0000,8.0000) circle (1.2pt);
\fill[red!80!black] (0.0000,8.0000) circle (1.2pt);
\fill[red!80!black] (12.0000,8.0000) circle (1.2pt);
\fill[red!80!black] (0.0000,8.0000) circle (1.2pt);
\fill[red!80!black] (12.0000,8.0000) circle (1.2pt);
\fill[red!80!black] (0.0000,8.0000) circle (1.2pt);
\fill[red!80!black] (12.0000,8.0000) circle (1.2pt);
\fill[red!80!black] (0.0000,8.0000) circle (1.2pt);
\fill[red!80!black] (12.0000,8.0000) circle (1.2pt);
\fill[red!80!black] (0.0000,8.0000) circle (1.2pt);
\fill[red!80!black] (12.0000,8.0000) circle (1.2pt);
\fill[red!80!black] (0.0000,8.0000) circle (1.2pt);
\fill[red!80!black] (12.0000,8.0000) circle (1.2pt);
\fill[red!80!black] (0.0000,8.0000) circle (1.2pt);
\fill[red!80!black] (12.0000,8.0000) circle (1.2pt);
\fill[red!80!black] (0.0000,8.0000) circle (1.2pt);
\fill[red!80!black] (12.0000,8.0000) circle (1.2pt);
\fill[red!80!black] (0.0000,8.0000) circle (1.2pt);
\fill[red!80!black] (12.0000,8.0000) circle (1.2pt);
\fill[red!80!black] (0.0000,8.0000) circle (1.2pt);
\fill[red!80!black] (12.0000,8.0000) circle (1.2pt);
\fill[red!80!black] (0.0000,8.0000) circle (1.2pt);
\fill[red!80!black] (12.0000,8.0000) circle (1.2pt);
\fill[red!80!black] (0.0000,8.0000) circle (1.2pt);
\fill[red!80!black] (12.0000,8.0000) circle (1.2pt);
\fill[red!80!black] (0.0000,8.0000) circle (1.2pt);
\fill[red!80!black] (12.0000,8.0000) circle (1.2pt);
\fill[red!80!black] (0.0000,8.0000) circle (1.2pt);
\fill[red!80!black] (12.0000,8.0000) circle (1.2pt);
\fill[red!80!black] (0.0000,8.0000) circle (1.2pt);
\fill[red!80!black] (12.0000,8.0000) circle (1.2pt);
\fill[red!80!black] (0.0000,8.0000) circle (1.2pt);
\fill[red!80!black] (12.0000,8.0000) circle (1.2pt);
\fill[red!80!black] (0.0000,8.0000) circle (1.2pt);
\fill[red!80!black] (12.0000,8.0000) circle (1.2pt);
\fill[red!80!black] (0.0000,8.0000) circle (1.2pt);
\fill[red!80!black] (12.0000,8.0000) circle (1.2pt);
\fill[red!80!black] (0.0000,8.0000) circle (1.2pt);
\fill[red!80!black] (12.0000,8.0000) circle (1.2pt);
\fill[red!80!black] (0.0000,8.0000) circle (1.2pt);
\fill[red!80!black] (12.0000,8.0000) circle (1.2pt);
\fill[red!80!black] (0.0000,8.0000) circle (1.2pt);
\fill[red!80!black] (12.0000,8.0000) circle (1.2pt);
\fill[red!80!black] (0.0000,8.0000) circle (1.2pt);
\fill[red!80!black] (12.0000,8.0000) circle (1.2pt);
\fill[red!80!black] (0.0000,8.0000) circle (1.2pt);
\fill[red!80!black] (12.0000,8.0000) circle (1.2pt);
\fill[red!80!black] (0.0000,8.0000) circle (1.2pt);
\fill[red!80!black] (12.0000,8.0000) circle (1.2pt);
\fill[red!80!black] (0.0000,8.0000) circle (1.2pt);
\fill[red!80!black] (12.0000,8.0000) circle (1.2pt);
\fill[red!80!black] (0.0000,8.0000) circle (1.2pt);
\fill[red!80!black] (12.0000,8.0000) circle (1.2pt);
\fill[red!80!black] (0.0000,8.0000) circle (1.2pt);
\fill[red!80!black] (12.0000,8.0000) circle (1.2pt);
\fill[red!80!black] (0.0000,8.0000) circle (1.2pt);
\fill[red!80!black] (12.0000,8.0000) circle (1.2pt);
\fill[red!80!black] (0.0000,8.0000) circle (1.2pt);
\fill[red!80!black] (12.0000,8.0000) circle (1.2pt);
\fill[red!80!black] (0.0000,8.0000) circle (1.2pt);
\fill[red!80!black] (12.0000,8.0000) circle (1.2pt);
\fill[red!80!black] (0.0000,8.0000) circle (1.2pt);
\fill[red!80!black] (12.0000,8.0000) circle (1.2pt);
\fill[red!80!black] (0.0000,8.0000) circle (1.2pt);
\fill[red!80!black] (12.0000,8.0000) circle (1.2pt);
\fill[red!80!black] (0.0000,8.0000) circle (1.2pt);
\fill[red!80!black] (12.0000,8.0000) circle (1.2pt);
\fill[red!80!black] (0.0000,8.0000) circle (1.2pt);
\fill[red!80!black] (12.0000,8.0000) circle (1.2pt);
\fill[red!80!black] (0.0000,8.0000) circle (1.2pt);
\fill[red!80!black] (12.0000,8.0000) circle (1.2pt);
\fill[red!80!black] (0.0000,8.0000) circle (1.2pt);
\fill[red!80!black] (12.0000,8.0000) circle (1.2pt);
\fill[red!80!black] (0.0000,8.0000) circle (1.2pt);
\fill[red!80!black] (12.0000,8.0000) circle (1.2pt);
\fill[red!80!black] (0.0000,8.0000) circle (1.2pt);
\fill[red!80!black] (12.0000,8.0000) circle (1.2pt);
\fill[red!80!black] (0.0000,8.0000) circle (1.2pt);
\fill[red!80!black] (12.0000,8.0000) circle (1.2pt);
\fill[red!80!black] (0.0000,8.0000) circle (1.2pt);
\fill[red!80!black] (12.0000,8.0000) circle (1.2pt);
\fill[red!80!black] (0.0000,8.0000) circle (1.2pt);
\fill[red!80!black] (12.0000,8.0000) circle (1.2pt);
\fill[red!80!black] (0.0000,8.0000) circle (1.2pt);
\fill[red!80!black] (12.0000,8.0000) circle (1.2pt);
\fill[red!80!black] (0.0000,8.0000) circle (1.2pt);
\fill[red!80!black] (12.0000,8.0000) circle (1.2pt);
\fill[red!80!black] (0.0000,8.0000) circle (1.2pt);
\fill[red!80!black] (12.0000,8.0000) circle (1.2pt);
\fill[red!80!black] (0.0000,8.0000) circle (1.2pt);
\fill[red!80!black] (12.0000,8.0000) circle (1.2pt);
\fill[red!80!black] (0.0000,8.0000) circle (1.2pt);
\fill[red!80!black] (12.0000,8.0000) circle (1.2pt);
\fill[red!80!black] (0.0000,8.0000) circle (1.2pt);
\fill[red!80!black] (12.0000,8.0000) circle (1.2pt);
\fill[red!80!black] (0.0000,8.0000) circle (1.2pt);
\fill[red!80!black] (12.0000,8.0000) circle (1.2pt);
\fill[red!80!black] (0.0000,8.0000) circle (1.2pt);
\fill[red!80!black] (12.0000,8.0000) circle (1.2pt);
\fill[red!80!black] (0.0000,8.0000) circle (1.2pt);
\fill[red!80!black] (12.0000,8.0000) circle (1.2pt);
\fill[red!80!black] (0.0000,8.0000) circle (1.2pt);
\fill[red!80!black] (12.0000,8.0000) circle (1.2pt);
\fill[red!80!black] (0.0000,8.0000) circle (1.2pt);
\fill[red!80!black] (12.0000,8.0000) circle (1.2pt);
\fill[red!80!black] (0.0000,8.0000) circle (1.2pt);
\fill[red!80!black] (12.0000,8.0000) circle (1.2pt);
\fill[red!80!black] (0.0000,8.0000) circle (1.2pt);
\fill[red!80!black] (12.0000,8.0000) circle (1.2pt);
\fill[red!80!black] (0.0000,8.0000) circle (1.2pt);
\fill[red!80!black] (12.0000,8.0000) circle (1.2pt);
\fill[red!80!black] (0.0000,8.0000) circle (1.2pt);
\fill[red!80!black] (12.0000,8.0000) circle (1.2pt);
\fill[red!80!black] (0.0000,8.0000) circle (1.2pt);
\fill[red!80!black] (12.0000,8.0000) circle (1.2pt);
\fill[red!80!black] (0.0000,8.0000) circle (1.2pt);
\fill[red!80!black] (12.0000,8.0000) circle (1.2pt);
\fill[red!80!black] (0.0000,8.0000) circle (1.2pt);
\fill[red!80!black] (12.0000,8.0000) circle (1.2pt);
\fill[red!80!black] (0.0000,8.0000) circle (1.2pt);
\fill[red!80!black] (12.0000,8.0000) circle (1.2pt);
\fill[red!80!black] (0.0000,8.0000) circle (1.2pt);
\fill[red!80!black] (12.0000,8.0000) circle (1.2pt);
\fill[red!80!black] (0.0000,8.0000) circle (1.2pt);
\fill[red!80!black] (12.0000,8.0000) circle (1.2pt);
\fill[red!80!black] (0.0000,8.0000) circle (1.2pt);
\fill[red!80!black] (12.0000,8.0000) circle (1.2pt);
\fill[red!80!black] (0.0000,8.0000) circle (1.2pt);
\fill[red!80!black] (12.0000,8.0000) circle (1.2pt);
\fill[red!80!black] (0.0000,8.0000) circle (1.2pt);
\fill[red!80!black] (12.0000,8.0000) circle (1.2pt);
\fill[red!80!black] (0.0000,8.0000) circle (1.2pt);
\fill[red!80!black] (12.0000,8.0000) circle (1.2pt);
\fill[red!80!black] (0.0000,8.0000) circle (1.2pt);
\fill[red!80!black] (12.0000,8.0000) circle (1.2pt);
\fill[red!80!black] (0.0000,8.0000) circle (1.2pt);
\fill[red!80!black] (12.0000,8.0000) circle (1.2pt);
\fill[red!80!black] (0.0000,8.0000) circle (1.2pt);
\fill[red!80!black] (12.0000,8.0000) circle (1.2pt);
\fill[red!80!black] (0.0000,8.0000) circle (1.2pt);
\fill[red!80!black] (12.0000,8.0000) circle (1.2pt);
\fill[red!80!black] (0.0000,8.0000) circle (1.2pt);
\fill[red!80!black] (12.0000,8.0000) circle (1.2pt);
\fill[red!80!black] (0.0000,8.0000) circle (1.2pt);
\fill[red!80!black] (12.0000,8.0000) circle (1.2pt);
\fill[red!80!black] (0.0000,8.0000) circle (1.2pt);
\fill[red!80!black] (12.0000,8.0000) circle (1.2pt);
\fill[red!80!black] (0.0000,8.0000) circle (1.2pt);
\fill[red!80!black] (12.0000,8.0000) circle (1.2pt);
\fill[red!80!black] (0.0000,8.0000) circle (1.2pt);
\fill[red!80!black] (12.0000,8.0000) circle (1.2pt);
\fill[red!80!black] (0.0000,8.0000) circle (1.2pt);
\fill[red!80!black] (12.0000,8.0000) circle (1.2pt);
\fill[red!80!black] (0.0000,8.0000) circle (1.2pt);
\fill[red!80!black] (12.0000,8.0000) circle (1.2pt);
\fill[red!80!black] (0.0000,8.0000) circle (1.2pt);
\fill[red!80!black] (12.0000,8.0000) circle (1.2pt);
\fill[red!80!black] (0.0000,8.0000) circle (1.2pt);
\fill[red!80!black] (12.0000,8.0000) circle (1.2pt);
\fill[red!80!black] (0.0000,8.0000) circle (1.2pt);
\fill[red!80!black] (12.0000,8.0000) circle (1.2pt);
\fill[red!80!black] (0.0000,8.0000) circle (1.2pt);
\fill[red!80!black] (12.0000,8.0000) circle (1.2pt);
\fill[red!80!black] (0.0000,8.0000) circle (1.2pt);
\fill[red!80!black] (12.0000,8.0000) circle (1.2pt);
\fill[red!80!black] (0.0000,8.0000) circle (1.2pt);
\fill[red!80!black] (12.0000,8.0000) circle (1.2pt);
\fill[red!80!black] (0.0000,8.0000) circle (1.2pt);
\fill[red!80!black] (12.0000,8.0000) circle (1.2pt);
\fill[red!80!black] (0.0000,8.0000) circle (1.2pt);
\fill[red!80!black] (12.0000,8.0000) circle (1.2pt);
\fill[red!80!black] (0.0000,8.0000) circle (1.2pt);
\fill[red!80!black] (12.0000,8.0000) circle (1.2pt);
\fill[red!80!black] (0.0000,8.0000) circle (1.2pt);
\fill[red!80!black] (12.0000,8.0000) circle (1.2pt);
\fill[red!80!black] (0.0000,8.0000) circle (1.2pt);
\fill[red!80!black] (12.0000,8.0000) circle (1.2pt);
\fill[red!80!black] (0.0000,8.0000) circle (1.2pt);
\fill[red!80!black] (12.0000,8.0000) circle (1.2pt);
\fill[red!80!black] (0.0000,8.0000) circle (1.2pt);
\fill[red!80!black] (12.0000,8.0000) circle (1.2pt);
\fill[red!80!black] (0.0000,8.0000) circle (1.2pt);
\fill[red!80!black] (12.0000,8.0000) circle (1.2pt);
\fill[red!80!black] (0.0000,8.0000) circle (1.2pt);
\fill[red!80!black] (12.0000,8.0000) circle (1.2pt);
\fill[red!80!black] (0.0000,8.0000) circle (1.2pt);
\fill[red!80!black] (12.0000,8.0000) circle (1.2pt);
\fill[red!80!black] (0.0000,8.0000) circle (1.2pt);
\fill[red!80!black] (12.0000,8.0000) circle (1.2pt);
\fill[red!80!black] (0.0000,8.0000) circle (1.2pt);
\fill[red!80!black] (12.0000,8.0000) circle (1.2pt);
\fill[red!80!black] (0.0000,8.0000) circle (1.2pt);
\fill[red!80!black] (12.0000,8.0000) circle (1.2pt);
\fill[red!80!black] (0.0000,8.0000) circle (1.2pt);
\fill[red!80!black] (12.0000,8.0000) circle (1.2pt);
\fill[red!80!black] (0.0000,8.0000) circle (1.2pt);
\fill[red!80!black] (12.0000,8.0000) circle (1.2pt);
\fill[red!80!black] (0.0000,8.0000) circle (1.2pt);
\fill[red!80!black] (12.0000,8.0000) circle (1.2pt);
\fill[red!80!black] (0.0000,8.0000) circle (1.2pt);
\fill[red!80!black] (12.0000,8.0000) circle (1.2pt);
\fill[red!80!black] (0.0000,8.0000) circle (1.2pt);
\fill[red!80!black] (12.0000,8.0000) circle (1.2pt);
\fill[red!80!black] (0.0000,8.0000) circle (1.2pt);
\fill[red!80!black] (12.0000,8.0000) circle (1.2pt);
\fill[red!80!black] (0.0000,8.0000) circle (1.2pt);
\fill[red!80!black] (12.0000,8.0000) circle (1.2pt);
\fill[red!80!black] (0.0000,8.0000) circle (1.2pt);
\fill[red!80!black] (12.0000,8.0000) circle (1.2pt);
\fill[red!80!black] (0.0000,8.0000) circle (1.2pt);
\fill[red!80!black] (12.0000,8.0000) circle (1.2pt);
\fill[red!80!black] (0.0000,8.0000) circle (1.2pt);
\fill[red!80!black] (12.0000,8.0000) circle (1.2pt);
\fill[red!80!black] (0.0000,8.0000) circle (1.2pt);
\fill[red!80!black] (12.0000,8.0000) circle (1.2pt);
\fill[red!80!black] (0.0000,8.0000) circle (1.2pt);
\fill[red!80!black] (12.0000,8.0000) circle (1.2pt);
\fill[red!80!black] (0.0000,8.0000) circle (1.2pt);
\fill[red!80!black] (12.0000,8.0000) circle (1.2pt);
\fill[red!80!black] (0.0000,8.0000) circle (1.2pt);
\fill[red!80!black] (12.0000,8.0000) circle (1.2pt);
\fill[red!80!black] (0.0000,8.0000) circle (1.2pt);
\fill[red!80!black] (12.0000,8.0000) circle (1.2pt);
\fill[red!80!black] (0.0000,8.0000) circle (1.2pt);
\fill[red!80!black] (12.0000,8.0000) circle (1.2pt);
\fill[red!80!black] (0.0000,8.0000) circle (1.2pt);
\fill[red!80!black] (12.0000,8.0000) circle (1.2pt);
\fill[red!80!black] (0.0000,8.0000) circle (1.2pt);
\fill[red!80!black] (12.0000,8.0000) circle (1.2pt);
\fill[red!80!black] (0.0000,8.0000) circle (1.2pt);
\fill[red!80!black] (12.0000,8.0000) circle (1.2pt);
\fill[red!80!black] (0.0000,8.0000) circle (1.2pt);
\fill[red!80!black] (12.0000,8.0000) circle (1.2pt);
\fill[red!80!black] (0.0000,8.0000) circle (1.2pt);
\fill[red!80!black] (12.0000,8.0000) circle (1.2pt);
\fill[red!80!black] (0.0000,8.0000) circle (1.2pt);
\fill[red!80!black] (12.0000,8.0000) circle (1.2pt);
\fill[red!80!black] (0.0000,8.0000) circle (1.2pt);
\fill[red!80!black] (12.0000,8.0000) circle (1.2pt);
\fill[red!80!black] (0.0000,8.0000) circle (1.2pt);
\fill[red!80!black] (12.0000,8.0000) circle (1.2pt);
\fill[red!80!black] (0.0000,8.0000) circle (1.2pt);
\fill[red!80!black] (12.0000,8.0000) circle (1.2pt);
\fill[red!80!black] (0.0000,8.0000) circle (1.2pt);
\fill[red!80!black] (12.0000,8.0000) circle (1.2pt);
\fill[red!80!black] (0.0000,8.0000) circle (1.2pt);
\fill[red!80!black] (12.0000,8.0000) circle (1.2pt);
\fill[red!80!black] (0.0000,8.0000) circle (1.2pt);
\fill[red!80!black] (12.0000,8.0000) circle (1.2pt);
\fill[red!80!black] (0.0000,8.0000) circle (1.2pt);
\fill[red!80!black] (12.0000,8.0000) circle (1.2pt);
\fill[red!80!black] (0.0000,8.0000) circle (1.2pt);
\fill[red!80!black] (12.0000,8.0000) circle (1.2pt);
\fill[red!80!black] (0.0000,8.0000) circle (1.2pt);
\fill[red!80!black] (12.0000,8.0000) circle (1.2pt);
\fill[red!80!black] (0.0000,8.0000) circle (1.2pt);
\fill[red!80!black] (12.0000,8.0000) circle (1.2pt);
\fill[red!80!black] (0.0000,8.0000) circle (1.2pt);
\fill[red!80!black] (12.0000,8.0000) circle (1.2pt);
\fill[red!80!black] (0.0000,8.0000) circle (1.2pt);
\fill[red!80!black] (12.0000,8.0000) circle (1.2pt);
\fill[red!80!black] (0.0000,8.0000) circle (1.2pt);
\fill[red!80!black] (12.0000,8.0000) circle (1.2pt);
\fill[red!80!black] (0.0000,8.0000) circle (1.2pt);
\fill[red!80!black] (12.0000,8.0000) circle (1.2pt);
\fill[red!80!black] (0.0000,8.0000) circle (1.2pt);
\fill[red!80!black] (12.0000,8.0000) circle (1.2pt);
\fill[red!80!black] (0.0000,8.0000) circle (1.2pt);
\fill[red!80!black] (12.0000,8.0000) circle (1.2pt);
\fill[red!80!black] (0.0000,8.0000) circle (1.2pt);
\fill[red!80!black] (12.0000,8.0000) circle (1.2pt);
\fill[red!80!black] (0.0000,8.0000) circle (1.2pt);
\fill[red!80!black] (12.0000,8.0000) circle (1.2pt);
\fill[red!80!black] (0.0000,8.0000) circle (1.2pt);
\fill[red!80!black] (12.0000,8.0000) circle (1.2pt);
\fill[red!80!black] (0.0000,8.0000) circle (1.2pt);
\fill[red!80!black] (12.0000,8.0000) circle (1.2pt);
\fill[red!80!black] (0.0000,8.0000) circle (1.2pt);
\fill[red!80!black] (12.0000,8.0000) circle (1.2pt);
\fill[red!80!black] (0.0000,8.0000) circle (1.2pt);
\fill[red!80!black] (12.0000,8.0000) circle (1.2pt);
\fill[red!80!black] (0.0000,8.0000) circle (1.2pt);
\fill[red!80!black] (12.0000,8.0000) circle (1.2pt);
\fill[red!80!black] (0.0000,8.0000) circle (1.2pt);
\fill[red!80!black] (12.0000,8.0000) circle (1.2pt);
\fill[red!80!black] (0.0000,8.0000) circle (1.2pt);
\fill[red!80!black] (12.0000,8.0000) circle (1.2pt);
\fill[red!80!black] (0.0000,8.0000) circle (1.2pt);
\fill[red!80!black] (12.0000,8.0000) circle (1.2pt);
\fill[red!80!black] (0.0000,8.0000) circle (1.2pt);
\fill[red!80!black] (12.0000,8.0000) circle (1.2pt);
\fill[red!80!black] (0.0000,8.0000) circle (1.2pt);
\fill[red!80!black] (12.0000,8.0000) circle (1.2pt);
\fill[red!80!black] (0.0000,8.0000) circle (1.2pt);
\fill[red!80!black] (12.0000,8.0000) circle (1.2pt);
\fill[red!80!black] (0.0000,8.0000) circle (1.2pt);
\fill[red!80!black] (12.0000,8.0000) circle (1.2pt);
\fill[red!80!black] (0.0000,8.0000) circle (1.2pt);
\fill[red!80!black] (12.0000,8.0000) circle (1.2pt);
\fill[red!80!black] (0.0000,8.0000) circle (1.2pt);
\fill[red!80!black] (12.0000,8.0000) circle (1.2pt);
\fill[red!80!black] (0.0000,8.0000) circle (1.2pt);
\fill[red!80!black] (12.0000,8.0000) circle (1.2pt);
\fill[red!80!black] (0.0000,8.0000) circle (1.2pt);
\fill[red!80!black] (12.0000,8.0000) circle (1.2pt);
\fill[red!80!black] (0.0000,8.0000) circle (1.2pt);
\fill[red!80!black] (12.0000,8.0000) circle (1.2pt);
\fill[red!80!black] (0.0000,8.0000) circle (1.2pt);
\fill[red!80!black] (12.0000,8.0000) circle (1.2pt);
\fill[red!80!black] (0.0000,8.0000) circle (1.2pt);
\fill[red!80!black] (12.0000,8.0000) circle (1.2pt);
\fill[red!80!black] (0.0000,8.0000) circle (1.2pt);
\fill[red!80!black] (12.0000,8.0000) circle (1.2pt);
\fill[red!80!black] (0.0000,8.0000) circle (1.2pt);
\fill[red!80!black] (12.0000,8.0000) circle (1.2pt);
\fill[red!80!black] (0.0000,8.0000) circle (1.2pt);
\fill[red!80!black] (12.0000,8.0000) circle (1.2pt);
\fill[red!80!black] (0.0000,8.0000) circle (1.2pt);
\fill[red!80!black] (12.0000,8.0000) circle (1.2pt);
\fill[red!80!black] (0.0000,8.0000) circle (1.2pt);
\fill[red!80!black] (12.0000,8.0000) circle (1.2pt);
\fill[red!80!black] (0.0000,8.0000) circle (1.2pt);
\fill[red!80!black] (12.0000,8.0000) circle (1.2pt);
\fill[red!80!black] (0.0000,8.0000) circle (1.2pt);
\fill[red!80!black] (12.0000,8.0000) circle (1.2pt);
\fill[red!80!black] (0.0000,8.0000) circle (1.2pt);
\fill[red!80!black] (12.0000,8.0000) circle (1.2pt);
\fill[red!80!black] (0.0000,8.0000) circle (1.2pt);
\fill[red!80!black] (12.0000,8.0000) circle (1.2pt);
\fill[red!80!black] (0.0000,8.0000) circle (1.2pt);
\fill[red!80!black] (12.0000,8.0000) circle (1.2pt);
\fill[red!80!black] (0.0000,8.0000) circle (1.2pt);
\fill[red!80!black] (12.0000,8.0000) circle (1.2pt);
\fill[red!80!black] (0.0000,8.0000) circle (1.2pt);
\fill[red!80!black] (12.0000,8.0000) circle (1.2pt);
\fill[red!80!black] (0.0000,8.0000) circle (1.2pt);
\fill[red!80!black] (12.0000,8.0000) circle (1.2pt);
\fill[red!80!black] (0.0000,8.0000) circle (1.2pt);
\fill[red!80!black] (12.0000,8.0000) circle (1.2pt);
\fill[red!80!black] (0.0000,8.0000) circle (1.2pt);
\fill[red!80!black] (12.0000,8.0000) circle (1.2pt);
\fill[red!80!black] (0.0000,8.0000) circle (1.2pt);
\fill[red!80!black] (12.0000,8.0000) circle (1.2pt);
\fill[red!80!black] (0.0000,8.0000) circle (1.2pt);
\fill[red!80!black] (12.0000,8.0000) circle (1.2pt);
\fill[red!80!black] (0.0000,8.0000) circle (1.2pt);
\fill[red!80!black] (12.0000,8.0000) circle (1.2pt);
\fill[red!80!black] (0.0000,8.0000) circle (1.2pt);
\fill[red!80!black] (12.0000,8.0000) circle (1.2pt);
\fill[red!80!black] (0.0000,8.0000) circle (1.2pt);
\fill[red!80!black] (12.0000,8.0000) circle (1.2pt);
\fill[red!80!black] (0.0000,8.0000) circle (1.2pt);
\fill[red!80!black] (12.0000,8.0000) circle (1.2pt);
\fill[red!80!black] (0.0000,8.0000) circle (1.2pt);
\fill[red!80!black] (12.0000,8.0000) circle (1.2pt);
\fill[red!80!black] (0.0000,8.0000) circle (1.2pt);
\fill[red!80!black] (12.0000,8.0000) circle (1.2pt);
\fill[red!80!black] (0.0000,8.0000) circle (1.2pt);
\fill[red!80!black] (12.0000,8.0000) circle (1.2pt);
\fill[red!80!black] (0.0000,8.0000) circle (1.2pt);
\fill[red!80!black] (12.0000,8.0000) circle (1.2pt);
\fill[red!80!black] (0.0000,8.0000) circle (1.2pt);
\fill[red!80!black] (12.0000,8.0000) circle (1.2pt);
\fill[red!80!black] (0.0000,8.0000) circle (1.2pt);
\fill[red!80!black] (12.0000,8.0000) circle (1.2pt);
\fill[red!80!black] (0.0000,8.0000) circle (1.2pt);
\fill[red!80!black] (12.0000,8.0000) circle (1.2pt);
\fill[red!80!black] (0.0000,8.0000) circle (1.2pt);
\fill[red!80!black] (12.0000,8.0000) circle (1.2pt);
\fill[red!80!black] (0.0000,8.0000) circle (1.2pt);
\fill[red!80!black] (12.0000,8.0000) circle (1.2pt);
\fill[red!80!black] (0.0000,8.0000) circle (1.2pt);
\fill[red!80!black] (12.0000,8.0000) circle (1.2pt);
\fill[red!80!black] (0.0000,8.0000) circle (1.2pt);
\fill[red!80!black] (12.0000,8.0000) circle (1.2pt);
\fill[red!80!black] (0.0000,8.0000) circle (1.2pt);
\fill[red!80!black] (12.0000,8.0000) circle (1.2pt);
\fill[red!80!black] (0.0000,8.0000) circle (1.2pt);
\fill[red!80!black] (12.0000,8.0000) circle (1.2pt);
\fill[red!80!black] (0.0000,8.0000) circle (1.2pt);
\fill[red!80!black] (12.0000,8.0000) circle (1.2pt);
\fill[red!80!black] (0.0000,8.0000) circle (1.2pt);
\fill[red!80!black] (12.0000,8.0000) circle (1.2pt);
\fill[red!80!black] (0.0000,8.0000) circle (1.2pt);
\fill[red!80!black] (12.0000,8.0000) circle (1.2pt);
\fill[red!80!black] (0.0000,8.0000) circle (1.2pt);
\fill[red!80!black] (12.0000,8.0000) circle (1.2pt);
\fill[red!80!black] (0.0000,8.0000) circle (1.2pt);
\fill[red!80!black] (12.0000,8.0000) circle (1.2pt);
\fill[red!80!black] (0.0000,8.0000) circle (1.2pt);
\fill[red!80!black] (12.0000,8.0000) circle (1.2pt);
\fill[red!80!black] (0.0000,8.0000) circle (1.2pt);
\fill[red!80!black] (12.0000,8.0000) circle (1.2pt);
\fill[red!80!black] (0.0000,8.0000) circle (1.2pt);
\fill[red!80!black] (12.0000,8.0000) circle (1.2pt);
\fill[red!80!black] (0.0000,8.0000) circle (1.2pt);
\fill[red!80!black] (12.0000,8.0000) circle (1.2pt);
\fill[red!80!black] (0.0000,8.0000) circle (1.2pt);
\fill[red!80!black] (12.0000,8.0000) circle (1.2pt);
\fill[red!80!black] (0.0000,8.0000) circle (1.2pt);
\fill[red!80!black] (12.0000,8.0000) circle (1.2pt);
\fill[red!80!black] (0.0000,8.0000) circle (1.2pt);
\fill[red!80!black] (12.0000,8.0000) circle (1.2pt);
\fill[red!80!black] (0.0000,8.0000) circle (1.2pt);
\fill[red!80!black] (12.0000,8.0000) circle (1.2pt);
\fill[red!80!black] (0.0000,8.0000) circle (1.2pt);
\fill[red!80!black] (12.0000,8.0000) circle (1.2pt);
\fill[red!80!black] (0.0000,8.0000) circle (1.2pt);
\fill[red!80!black] (12.0000,8.0000) circle (1.2pt);
\fill[red!80!black] (0.0000,8.0000) circle (1.2pt);
\fill[red!80!black] (12.0000,8.0000) circle (1.2pt);
\fill[red!80!black] (0.0000,8.0000) circle (1.2pt);
\fill[red!80!black] (12.0000,8.0000) circle (1.2pt);
\fill[red!80!black] (0.0000,8.0000) circle (1.2pt);
\fill[red!80!black] (12.0000,8.0000) circle (1.2pt);
\fill[red!80!black] (0.0000,8.0000) circle (1.2pt);
\fill[red!80!black] (12.0000,8.0000) circle (1.2pt);
\fill[red!80!black] (0.0000,8.0000) circle (1.2pt);
\fill[red!80!black] (12.0000,8.0000) circle (1.2pt);
\fill[red!80!black] (0.0000,8.0000) circle (1.2pt);
\fill[red!80!black] (12.0000,8.0000) circle (1.2pt);
\fill[red!80!black] (0.0000,8.0000) circle (1.2pt);
\fill[red!80!black] (12.0000,8.0000) circle (1.2pt);
\fill[red!80!black] (0.0000,8.0000) circle (1.2pt);
\fill[red!80!black] (12.0000,8.0000) circle (1.2pt);
\fill[red!80!black] (0.0000,8.0000) circle (1.2pt);
\fill[red!80!black] (12.0000,8.0000) circle (1.2pt);
\fill[red!80!black] (0.0000,8.0000) circle (1.2pt);
\fill[red!80!black] (12.0000,8.0000) circle (1.2pt);
\fill[red!80!black] (0.0000,8.0000) circle (1.2pt);
\fill[red!80!black] (12.0000,8.0000) circle (1.2pt);
\fill[red!80!black] (0.0000,8.0000) circle (1.2pt);
\fill[red!80!black] (12.0000,8.0000) circle (1.2pt);
\fill[red!80!black] (0.0000,8.0000) circle (1.2pt);
\fill[red!80!black] (12.0000,8.0000) circle (1.2pt);
\fill[red!80!black] (0.0000,8.0000) circle (1.2pt);
\fill[red!80!black] (12.0000,8.0000) circle (1.2pt);
\fill[red!80!black] (0.0000,8.0000) circle (1.2pt);
\fill[red!80!black] (12.0000,8.0000) circle (1.2pt);
\fill[red!80!black] (0.0000,8.0000) circle (1.2pt);
\fill[red!80!black] (12.0000,8.0000) circle (1.2pt);
\fill[red!80!black] (0.0000,8.0000) circle (1.2pt);
\fill[red!80!black] (12.0000,8.0000) circle (1.2pt);
\fill[red!80!black] (0.0000,8.0000) circle (1.2pt);
\fill[red!80!black] (12.0000,8.0000) circle (1.2pt);
\fill[red!80!black] (0.0000,8.0000) circle (1.2pt);
\fill[red!80!black] (12.0000,8.0000) circle (1.2pt);
\fill[red!80!black] (0.0000,8.0000) circle (1.2pt);
\fill[red!80!black] (12.0000,8.0000) circle (1.2pt);
\fill[red!80!black] (0.0000,8.0000) circle (1.2pt);
\fill[red!80!black] (12.0000,8.0000) circle (1.2pt);
\fill[red!80!black] (0.0000,8.0000) circle (1.2pt);
\fill[red!80!black] (12.0000,8.0000) circle (1.2pt);
\fill[red!80!black] (0.0000,8.0000) circle (1.2pt);
\fill[red!80!black] (12.0000,8.0000) circle (1.2pt);
\fill[red!80!black] (0.0000,8.0000) circle (1.2pt);
\fill[red!80!black] (12.0000,8.0000) circle (1.2pt);
\fill[red!80!black] (0.0000,8.0000) circle (1.2pt);
\fill[red!80!black] (12.0000,8.0000) circle (1.2pt);
\fill[red!80!black] (0.0000,8.0000) circle (1.2pt);
\fill[red!80!black] (12.0000,8.0000) circle (1.2pt);
\fill[red!80!black] (0.0000,8.0000) circle (1.2pt);
\fill[red!80!black] (12.0000,8.0000) circle (1.2pt);
\fill[red!80!black] (0.0000,8.0000) circle (1.2pt);
\fill[red!80!black] (12.0000,8.0000) circle (1.2pt);
\fill[red!80!black] (0.0000,8.0000) circle (1.2pt);
\fill[red!80!black] (12.0000,8.0000) circle (1.2pt);
\fill[red!80!black] (0.0000,8.0000) circle (1.2pt);
\fill[red!80!black] (12.0000,8.0000) circle (1.2pt);
\fill[red!80!black] (0.0000,8.0000) circle (1.2pt);
\fill[red!80!black] (12.0000,8.0000) circle (1.2pt);
\fill[red!80!black] (0.0000,8.0000) circle (1.2pt);
\fill[red!80!black] (12.0000,8.0000) circle (1.2pt);
\fill[red!80!black] (0.0000,8.0000) circle (1.2pt);
\fill[red!80!black] (12.0000,8.0000) circle (1.2pt);
\fill[red!80!black] (0.0000,8.0000) circle (1.2pt);
\fill[red!80!black] (12.0000,8.0000) circle (1.2pt);
\fill[red!80!black] (0.0000,8.0000) circle (1.2pt);
\fill[red!80!black] (12.0000,8.0000) circle (1.2pt);
\fill[red!80!black] (0.0000,8.0000) circle (1.2pt);
\fill[red!80!black] (12.0000,8.0000) circle (1.2pt);
\fill[red!80!black] (0.0000,8.0000) circle (1.2pt);
\fill[red!80!black] (12.0000,8.0000) circle (1.2pt);
\fill[red!80!black] (0.0000,8.0000) circle (1.2pt);
\fill[red!80!black] (12.0000,8.0000) circle (1.2pt);
\fill[red!80!black] (0.0000,8.0000) circle (1.2pt);
\fill[red!80!black] (12.0000,8.0000) circle (1.2pt);
\fill[red!80!black] (0.0000,8.0000) circle (1.2pt);
\fill[red!80!black] (12.0000,8.0000) circle (1.2pt);
\fill[red!80!black] (0.0000,8.0000) circle (1.2pt);
\fill[red!80!black] (12.0000,8.0000) circle (1.2pt);
\fill[red!80!black] (0.0000,8.0000) circle (1.2pt);
\fill[red!80!black] (12.0000,8.0000) circle (1.2pt);
\fill[red!80!black] (0.0000,8.0000) circle (1.2pt);
\fill[red!80!black] (12.0000,8.0000) circle (1.2pt);
\fill[red!80!black] (0.0000,8.0000) circle (1.2pt);
\fill[red!80!black] (12.0000,8.0000) circle (1.2pt);
\fill[red!80!black] (0.0000,8.0000) circle (1.2pt);
\fill[red!80!black] (12.0000,8.0000) circle (1.2pt);
\fill[red!80!black] (0.0000,8.0000) circle (1.2pt);
\fill[red!80!black] (12.0000,8.0000) circle (1.2pt);
\fill[red!80!black] (0.0000,8.0000) circle (1.2pt);
\fill[red!80!black] (12.0000,8.0000) circle (1.2pt);
\fill[red!80!black] (0.0000,8.0000) circle (1.2pt);
\fill[red!80!black] (12.0000,8.0000) circle (1.2pt);
\fill[red!80!black] (0.0000,8.0000) circle (1.2pt);
\fill[red!80!black] (12.0000,8.0000) circle (1.2pt);
\fill[red!80!black] (0.0000,8.0000) circle (1.2pt);
\fill[red!80!black] (12.0000,8.0000) circle (1.2pt);
\fill[red!80!black] (0.0000,8.0000) circle (1.2pt);
\fill[red!80!black] (12.0000,8.0000) circle (1.2pt);
\fill[red!80!black] (0.0000,8.0000) circle (1.2pt);
\fill[red!80!black] (12.0000,8.0000) circle (1.2pt);
\fill[red!80!black] (0.0000,8.0000) circle (1.2pt);
\fill[red!80!black] (12.0000,8.0000) circle (1.2pt);
\fill[red!80!black] (0.0000,8.0000) circle (1.2pt);
\fill[red!80!black] (12.0000,8.0000) circle (1.2pt);
\fill[red!80!black] (0.0000,8.0000) circle (1.2pt);
\fill[red!80!black] (12.0000,8.0000) circle (1.2pt);
\fill[red!80!black] (0.0000,8.0000) circle (1.2pt);
\fill[red!80!black] (12.0000,8.0000) circle (1.2pt);
\fill[red!80!black] (0.0000,8.0000) circle (1.2pt);
\fill[red!80!black] (12.0000,8.0000) circle (1.2pt);
\fill[red!80!black] (0.0000,8.0000) circle (1.2pt);
\fill[red!80!black] (12.0000,8.0000) circle (1.2pt);
\fill[red!80!black] (0.0000,8.0000) circle (1.2pt);
\fill[red!80!black] (12.0000,8.0000) circle (1.2pt);
\fill[red!80!black] (0.0000,8.0000) circle (1.2pt);
\fill[red!80!black] (12.0000,8.0000) circle (1.2pt);
\fill[red!80!black] (0.0000,8.0000) circle (1.2pt);
\fill[red!80!black] (12.0000,8.0000) circle (1.2pt);
\fill[red!80!black] (0.0000,8.0000) circle (1.2pt);
\fill[red!80!black] (12.0000,8.0000) circle (1.2pt);
\fill[red!80!black] (0.0000,8.0000) circle (1.2pt);
\fill[red!80!black] (12.0000,8.0000) circle (1.2pt);
\fill[red!80!black] (0.0000,8.0000) circle (1.2pt);
\fill[red!80!black] (12.0000,8.0000) circle (1.2pt);
\fill[red!80!black] (0.0000,8.0000) circle (1.2pt);
\fill[red!80!black] (12.0000,8.0000) circle (1.2pt);
\fill[red!80!black] (0.0000,8.0000) circle (1.2pt);
\fill[red!80!black] (12.0000,8.0000) circle (1.2pt);
\fill[red!80!black] (0.0000,8.0000) circle (1.2pt);
\fill[red!80!black] (12.0000,8.0000) circle (1.2pt);
\fill[red!80!black] (0.0000,8.0000) circle (1.2pt);
\fill[red!80!black] (12.0000,8.0000) circle (1.2pt);
\fill[red!80!black] (0.0000,8.0000) circle (1.2pt);
\fill[red!80!black] (12.0000,8.0000) circle (1.2pt);
\fill[red!80!black] (0.0000,8.0000) circle (1.2pt);
\fill[red!80!black] (12.0000,8.0000) circle (1.2pt);
\fill[red!80!black] (0.0000,8.0000) circle (1.2pt);
\fill[red!80!black] (12.0000,8.0000) circle (1.2pt);
\fill[red!80!black] (0.0000,8.0000) circle (1.2pt);
\fill[red!80!black] (12.0000,8.0000) circle (1.2pt);
\fill[red!80!black] (0.0000,8.0000) circle (1.2pt);
\fill[red!80!black] (12.0000,8.0000) circle (1.2pt);
\fill[red!80!black] (0.0000,8.0000) circle (1.2pt);
\fill[red!80!black] (12.0000,8.0000) circle (1.2pt);
\fill[red!80!black] (0.0000,8.0000) circle (1.2pt);
\fill[red!80!black] (12.0000,8.0000) circle (1.2pt);
\fill[red!80!black] (0.0000,8.0000) circle (1.2pt);
\fill[red!80!black] (12.0000,8.0000) circle (1.2pt);
\fill[red!80!black] (0.0000,8.0000) circle (1.2pt);
\fill[red!80!black] (12.0000,8.0000) circle (1.2pt);
\fill[red!80!black] (0.0000,8.0000) circle (1.2pt);
\fill[red!80!black] (12.0000,8.0000) circle (1.2pt);
\fill[red!80!black] (0.0000,8.0000) circle (1.2pt);
\fill[red!80!black] (12.0000,8.0000) circle (1.2pt);
\fill[red!80!black] (0.0000,8.0000) circle (1.2pt);
\fill[red!80!black] (12.0000,8.0000) circle (1.2pt);
\fill[red!80!black] (0.0000,8.0000) circle (1.2pt);
\fill[red!80!black] (12.0000,8.0000) circle (1.2pt);
\fill[red!80!black] (0.0000,8.0000) circle (1.2pt);
\fill[red!80!black] (12.0000,8.0000) circle (1.2pt);
\fill[red!80!black] (0.0000,8.0000) circle (1.2pt);
\fill[red!80!black] (12.0000,8.0000) circle (1.2pt);
\fill[red!80!black] (0.0000,8.0000) circle (1.2pt);
\fill[red!80!black] (12.0000,8.0000) circle (1.2pt);
\fill[red!80!black] (0.0000,8.0000) circle (1.2pt);
\fill[red!80!black] (12.0000,8.0000) circle (1.2pt);
\fill[red!80!black] (0.0000,8.0000) circle (1.2pt);
\fill[red!80!black] (12.0000,8.0000) circle (1.2pt);
\fill[red!80!black] (0.0000,8.0000) circle (1.2pt);
\fill[red!80!black] (12.0000,8.0000) circle (1.2pt);
\fill[red!80!black] (0.0000,8.0000) circle (1.2pt);
\fill[red!80!black] (12.0000,8.0000) circle (1.2pt);
\fill[red!80!black] (0.0000,8.0000) circle (1.2pt);
\fill[red!80!black] (12.0000,8.0000) circle (1.2pt);
\fill[red!80!black] (0.0000,8.0000) circle (1.2pt);
\fill[red!80!black] (12.0000,8.0000) circle (1.2pt);
\fill[red!80!black] (0.0000,8.0000) circle (1.2pt);
\fill[red!80!black] (12.0000,8.0000) circle (1.2pt);
\fill[red!80!black] (0.0000,8.0000) circle (1.2pt);
\fill[red!80!black] (12.0000,8.0000) circle (1.2pt);
\fill[red!80!black] (0.0000,8.0000) circle (1.2pt);
\fill[red!80!black] (12.0000,8.0000) circle (1.2pt);
\fill[red!80!black] (0.0000,8.0000) circle (1.2pt);
\fill[red!80!black] (12.0000,8.0000) circle (1.2pt);
\fill[red!80!black] (0.0000,8.0000) circle (1.2pt);
\fill[red!80!black] (12.0000,8.0000) circle (1.2pt);
\fill[red!80!black] (0.0000,8.0000) circle (1.2pt);
\fill[red!80!black] (12.0000,8.0000) circle (1.2pt);
\fill[red!80!black] (0.0000,8.0000) circle (1.2pt);
\fill[red!80!black] (12.0000,8.0000) circle (1.2pt);
\fill[red!80!black] (0.0000,8.0000) circle (1.2pt);
\fill[red!80!black] (12.0000,8.0000) circle (1.2pt);
\fill[red!80!black] (0.0000,8.0000) circle (1.2pt);
\fill[red!80!black] (12.0000,8.0000) circle (1.2pt);
\fill[red!80!black] (0.0000,8.0000) circle (1.2pt);
\fill[red!80!black] (12.0000,8.0000) circle (1.2pt);
\fill[red!80!black] (0.0000,8.0000) circle (1.2pt);
\fill[red!80!black] (12.0000,8.0000) circle (1.2pt);
\fill[red!80!black] (0.0000,8.0000) circle (1.2pt);
\fill[red!80!black] (12.0000,8.0000) circle (1.2pt);
\fill[red!80!black] (0.0000,8.0000) circle (1.2pt);
\fill[red!80!black] (12.0000,8.0000) circle (1.2pt);
\fill[red!80!black] (0.0000,8.0000) circle (1.2pt);
\fill[red!80!black] (12.0000,8.0000) circle (1.2pt);
\fill[red!80!black] (0.0000,8.0000) circle (1.2pt);
\fill[red!80!black] (12.0000,8.0000) circle (1.2pt);
\fill[red!80!black] (0.0000,8.0000) circle (1.2pt);
\fill[red!80!black] (12.0000,8.0000) circle (1.2pt);
\fill[red!80!black] (0.0000,8.0000) circle (1.2pt);
\fill[red!80!black] (12.0000,8.0000) circle (1.2pt);
\fill[red!80!black] (0.0000,8.0000) circle (1.2pt);
\fill[red!80!black] (12.0000,8.0000) circle (1.2pt);
\fill[red!80!black] (0.0000,8.0000) circle (1.2pt);
\fill[red!80!black] (12.0000,8.0000) circle (1.2pt);
\fill[red!80!black] (0.0000,8.0000) circle (1.2pt);
\fill[red!80!black] (12.0000,8.0000) circle (1.2pt);
\fill[red!80!black] (0.0000,8.0000) circle (1.2pt);
\fill[red!80!black] (12.0000,8.0000) circle (1.2pt);
\fill[red!80!black] (0.0000,8.0000) circle (1.2pt);
\fill[red!80!black] (12.0000,8.0000) circle (1.2pt);
\fill[red!80!black] (0.0000,8.0000) circle (1.2pt);
\fill[red!80!black] (12.0000,8.0000) circle (1.2pt);
\fill[red!80!black] (0.0000,8.0000) circle (1.2pt);
\fill[red!80!black] (12.0000,8.0000) circle (1.2pt);
\fill[red!80!black] (0.0000,8.0000) circle (1.2pt);
\fill[red!80!black] (12.0000,8.0000) circle (1.2pt);
\fill[red!80!black] (0.0000,8.0000) circle (1.2pt);
\fill[red!80!black] (12.0000,8.0000) circle (1.2pt);
\fill[red!80!black] (0.0000,8.0000) circle (1.2pt);
\fill[red!80!black] (12.0000,8.0000) circle (1.2pt);
\fill[red!80!black] (0.0000,8.0000) circle (1.2pt);
\fill[red!80!black] (12.0000,8.0000) circle (1.2pt);
\fill[red!80!black] (0.0000,8.0000) circle (1.2pt);
\fill[red!80!black] (12.0000,8.0000) circle (1.2pt);
\fill[red!80!black] (0.0000,8.0000) circle (1.2pt);
\fill[red!80!black] (12.0000,8.0000) circle (1.2pt);
\fill[red!80!black] (0.0000,8.0000) circle (1.2pt);
\fill[red!80!black] (12.0000,8.0000) circle (1.2pt);
\fill[red!80!black] (0.0000,8.0000) circle (1.2pt);
\fill[red!80!black] (12.0000,8.0000) circle (1.2pt);
\fill[red!80!black] (0.0000,8.0000) circle (1.2pt);
\fill[red!80!black] (12.0000,8.0000) circle (1.2pt);
\fill[red!80!black] (0.0000,8.0000) circle (1.2pt);
\fill[red!80!black] (12.0000,8.0000) circle (1.2pt);
\fill[red!80!black] (0.0000,8.0000) circle (1.2pt);
\fill[red!80!black] (12.0000,8.0000) circle (1.2pt);
\fill[red!80!black] (0.0000,8.0000) circle (1.2pt);
\fill[red!80!black] (12.0000,8.0000) circle (1.2pt);
\fill[red!80!black] (0.0000,8.0000) circle (1.2pt);
\fill[red!80!black] (12.0000,8.0000) circle (1.2pt);
\fill[red!80!black] (0.0000,8.0000) circle (1.2pt);
\fill[red!80!black] (12.0000,8.0000) circle (1.2pt);
\fill[red!80!black] (0.0000,8.0000) circle (1.2pt);
\fill[red!80!black] (12.0000,8.0000) circle (1.2pt);
\fill[red!80!black] (0.0000,8.0000) circle (1.2pt);
\fill[red!80!black] (12.0000,8.0000) circle (1.2pt);
\fill[red!80!black] (0.0000,8.0000) circle (1.2pt);
\fill[red!80!black] (12.0000,8.0000) circle (1.2pt);
\fill[red!80!black] (0.0000,8.0000) circle (1.2pt);
\fill[red!80!black] (12.0000,8.0000) circle (1.2pt);
\fill[red!80!black] (0.0000,8.0000) circle (1.2pt);
\fill[red!80!black] (12.0000,8.0000) circle (1.2pt);
\fill[red!80!black] (0.0000,8.0000) circle (1.2pt);
\fill[red!80!black] (12.0000,8.0000) circle (1.2pt);
\fill[red!80!black] (0.0000,8.0000) circle (1.2pt);
\fill[red!80!black] (12.0000,8.0000) circle (1.2pt);
\fill[red!80!black] (0.0000,8.0000) circle (1.2pt);
\fill[red!80!black] (12.0000,8.0000) circle (1.2pt);
\fill[red!80!black] (0.0000,8.0000) circle (1.2pt);
\fill[red!80!black] (12.0000,8.0000) circle (1.2pt);
\fill[red!80!black] (0.0000,8.0000) circle (1.2pt);
\fill[red!80!black] (12.0000,8.0000) circle (1.2pt);
\fill[red!80!black] (0.0000,8.0000) circle (1.2pt);
\fill[red!80!black] (12.0000,8.0000) circle (1.2pt);
\fill[red!80!black] (0.0000,8.0000) circle (1.2pt);
\fill[red!80!black] (12.0000,8.0000) circle (1.2pt);
\fill[red!80!black] (0.0000,8.0000) circle (1.2pt);
\fill[red!80!black] (12.0000,8.0000) circle (1.2pt);
\fill[red!80!black] (0.0000,8.0000) circle (1.2pt);
\fill[red!80!black] (12.0000,8.0000) circle (1.2pt);
\fill[red!80!black] (0.0000,8.0000) circle (1.2pt);
\fill[red!80!black] (12.0000,8.0000) circle (1.2pt);
\fill[red!80!black] (0.0000,8.0000) circle (1.2pt);
\fill[red!80!black] (12.0000,8.0000) circle (1.2pt);
\fill[red!80!black] (0.0000,8.0000) circle (1.2pt);
\fill[red!80!black] (12.0000,8.0000) circle (1.2pt);
\fill[red!80!black] (0.0000,8.0000) circle (1.2pt);
\fill[red!80!black] (12.0000,8.0000) circle (1.2pt);
\fill[red!80!black] (0.0000,8.0000) circle (1.2pt);
\fill[red!80!black] (12.0000,8.0000) circle (1.2pt);
\fill[red!80!black] (0.0000,8.0000) circle (1.2pt);
\fill[red!80!black] (12.0000,8.0000) circle (1.2pt);
\fill[red!80!black] (0.0000,8.0000) circle (1.2pt);
\fill[red!80!black] (12.0000,8.0000) circle (1.2pt);
\fill[red!80!black] (0.0000,8.0000) circle (1.2pt);
\fill[red!80!black] (12.0000,8.0000) circle (1.2pt);
\fill[red!80!black] (0.0000,8.0000) circle (1.2pt);
\fill[red!80!black] (12.0000,8.0000) circle (1.2pt);
\fill[red!80!black] (0.0000,8.0000) circle (1.2pt);
\fill[red!80!black] (12.0000,8.0000) circle (1.2pt);
\fill[red!80!black] (0.0000,8.0000) circle (1.2pt);
\fill[red!80!black] (12.0000,8.0000) circle (1.2pt);
\fill[red!80!black] (0.0000,8.0000) circle (1.2pt);
\fill[red!80!black] (12.0000,8.0000) circle (1.2pt);
\fill[red!80!black] (0.0000,8.0000) circle (1.2pt);
\fill[red!80!black] (12.0000,8.0000) circle (1.2pt);
\fill[red!80!black] (0.0000,8.0000) circle (1.2pt);
\fill[red!80!black] (12.0000,8.0000) circle (1.2pt);
\fill[red!80!black] (0.0000,8.0000) circle (1.2pt);
\fill[red!80!black] (12.0000,8.0000) circle (1.2pt);
\fill[red!80!black] (0.0000,8.0000) circle (1.2pt);
\fill[red!80!black] (12.0000,8.0000) circle (1.2pt);
\fill[red!80!black] (0.0000,8.0000) circle (1.2pt);
\fill[red!80!black] (12.0000,8.0000) circle (1.2pt);
\fill[red!80!black] (0.0000,8.0000) circle (1.2pt);
\fill[red!80!black] (12.0000,8.0000) circle (1.2pt);
\fill[red!80!black] (0.0000,8.0000) circle (1.2pt);
\fill[red!80!black] (12.0000,8.0000) circle (1.2pt);
\fill[red!80!black] (0.0000,8.0000) circle (1.2pt);
\fill[red!80!black] (12.0000,8.0000) circle (1.2pt);
\fill[red!80!black] (0.0000,8.0000) circle (1.2pt);
\fill[red!80!black] (12.0000,8.0000) circle (1.2pt);
\fill[red!80!black] (0.0000,8.0000) circle (1.2pt);
\fill[red!80!black] (12.0000,8.0000) circle (1.2pt);
\fill[red!80!black] (0.0000,8.0000) circle (1.2pt);
\fill[red!80!black] (12.0000,8.0000) circle (1.2pt);
\fill[red!80!black] (0.0000,8.0000) circle (1.2pt);
\fill[red!80!black] (12.0000,8.0000) circle (1.2pt);
\fill[red!80!black] (0.0000,8.0000) circle (1.2pt);
\fill[red!80!black] (12.0000,8.0000) circle (1.2pt);
\fill[red!80!black] (0.0000,8.0000) circle (1.2pt);
\fill[red!80!black] (12.0000,8.0000) circle (1.2pt);
\fill[red!80!black] (0.0000,8.0000) circle (1.2pt);
\fill[red!80!black] (12.0000,8.0000) circle (1.2pt);
\fill[red!80!black] (0.0000,8.0000) circle (1.2pt);
\fill[red!80!black] (12.0000,8.0000) circle (1.2pt);
\fill[red!80!black] (0.0000,8.0000) circle (1.2pt);
\fill[red!80!black] (12.0000,8.0000) circle (1.2pt);
\fill[red!80!black] (0.0000,8.0000) circle (1.2pt);
\fill[red!80!black] (12.0000,8.0000) circle (1.2pt);
\fill[red!80!black] (0.0000,8.0000) circle (1.2pt);
\fill[red!80!black] (12.0000,8.0000) circle (1.2pt);
\fill[red!80!black] (0.0000,8.0000) circle (1.2pt);
\fill[red!80!black] (12.0000,8.0000) circle (1.2pt);
\fill[red!80!black] (0.0000,8.0000) circle (1.2pt);
\fill[red!80!black] (12.0000,8.0000) circle (1.2pt);
\fill[red!80!black] (0.0000,8.0000) circle (1.2pt);
\fill[red!80!black] (12.0000,8.0000) circle (1.2pt);
\fill[red!80!black] (0.0000,8.0000) circle (1.2pt);
\fill[red!80!black] (12.0000,8.0000) circle (1.2pt);
\fill[red!80!black] (0.0000,8.0000) circle (1.2pt);
\fill[red!80!black] (12.0000,8.0000) circle (1.2pt);
\fill[red!80!black] (0.0000,8.0000) circle (1.2pt);
\fill[red!80!black] (12.0000,8.0000) circle (1.2pt);
\fill[red!80!black] (0.0000,8.0000) circle (1.2pt);
\fill[red!80!black] (12.0000,8.0000) circle (1.2pt);
\fill[red!80!black] (0.0000,8.0000) circle (1.2pt);
\fill[red!80!black] (12.0000,8.0000) circle (1.2pt);
\fill[red!80!black] (0.0000,8.0000) circle (1.2pt);
\fill[red!80!black] (12.0000,8.0000) circle (1.2pt);
\fill[red!80!black] (0.0000,8.0000) circle (1.2pt);
\fill[red!80!black] (12.0000,8.0000) circle (1.2pt);
\fill[red!80!black] (0.0000,8.0000) circle (1.2pt);
\fill[red!80!black] (12.0000,8.0000) circle (1.2pt);
\fill[red!80!black] (0.0000,8.0000) circle (1.2pt);
\fill[red!80!black] (12.0000,8.0000) circle (1.2pt);
\fill[red!80!black] (0.0000,8.0000) circle (1.2pt);
\fill[red!80!black] (12.0000,8.0000) circle (1.2pt);
\fill[red!80!black] (0.0000,8.0000) circle (1.2pt);
\fill[red!80!black] (12.0000,8.0000) circle (1.2pt);
\fill[red!80!black] (0.0000,8.0000) circle (1.2pt);
\fill[red!80!black] (12.0000,8.0000) circle (1.2pt);
\fill[red!80!black] (0.0000,8.0000) circle (1.2pt);
\fill[red!80!black] (12.0000,8.0000) circle (1.2pt);
\fill[red!80!black] (0.0000,8.0000) circle (1.2pt);
\fill[red!80!black] (12.0000,8.0000) circle (1.2pt);
\fill[red!80!black] (0.0000,8.0000) circle (1.2pt);
\fill[red!80!black] (12.0000,8.0000) circle (1.2pt);
\fill[red!80!black] (0.0000,8.0000) circle (1.2pt);
\fill[red!80!black] (12.0000,8.0000) circle (1.2pt);
\fill[red!80!black] (0.0000,8.0000) circle (1.2pt);
\fill[red!80!black] (12.0000,8.0000) circle (1.2pt);
\fill[red!80!black] (0.0000,8.0000) circle (1.2pt);
\fill[red!80!black] (12.0000,8.0000) circle (1.2pt);
\fill[red!80!black] (0.0000,8.0000) circle (1.2pt);
\fill[red!80!black] (12.0000,8.0000) circle (1.2pt);
\fill[red!80!black] (0.0000,8.0000) circle (1.2pt);
\fill[red!80!black] (12.0000,8.0000) circle (1.2pt);
\fill[red!80!black] (0.0000,8.0000) circle (1.2pt);
\fill[red!80!black] (12.0000,8.0000) circle (1.2pt);
\fill[red!80!black] (0.0000,8.0000) circle (1.2pt);
\fill[red!80!black] (12.0000,8.0000) circle (1.2pt);
\fill[red!80!black] (0.0000,8.0000) circle (1.2pt);
\fill[red!80!black] (12.0000,8.0000) circle (1.2pt);
\fill[red!80!black] (0.0000,8.0000) circle (1.2pt);
\fill[red!80!black] (12.0000,8.0000) circle (1.2pt);
\fill[red!80!black] (0.0000,8.0000) circle (1.2pt);
\fill[red!80!black] (12.0000,8.0000) circle (1.2pt);
\fill[red!80!black] (0.0000,8.0000) circle (1.2pt);
\fill[red!80!black] (12.0000,8.0000) circle (1.2pt);
\fill[red!80!black] (0.0000,8.0000) circle (1.2pt);
\fill[red!80!black] (12.0000,8.0000) circle (1.2pt);
\fill[red!80!black] (0.0000,8.0000) circle (1.2pt);
\fill[red!80!black] (12.0000,8.0000) circle (1.2pt);
\fill[red!80!black] (0.0000,8.0000) circle (1.2pt);
\fill[red!80!black] (12.0000,8.0000) circle (1.2pt);
\fill[red!80!black] (0.0000,8.0000) circle (1.2pt);
\fill[red!80!black] (12.0000,8.0000) circle (1.2pt);
\fill[red!80!black] (0.0000,8.0000) circle (1.2pt);
\fill[red!80!black] (12.0000,8.0000) circle (1.2pt);
\fill[red!80!black] (0.0000,8.0000) circle (1.2pt);
\fill[red!80!black] (12.0000,8.0000) circle (1.2pt);
\fill[red!80!black] (0.0000,8.0000) circle (1.2pt);
\fill[red!80!black] (12.0000,8.0000) circle (1.2pt);
\fill[red!80!black] (0.0000,8.0000) circle (1.2pt);
\fill[red!80!black] (12.0000,8.0000) circle (1.2pt);
\fill[red!80!black] (0.0000,8.0000) circle (1.2pt);
\fill[red!80!black] (12.0000,8.0000) circle (1.2pt);
\fill[red!80!black] (0.0000,8.0000) circle (1.2pt);
\fill[red!80!black] (12.0000,8.0000) circle (1.2pt);
\fill[red!80!black] (0.0000,8.0000) circle (1.2pt);
\fill[red!80!black] (12.0000,8.0000) circle (1.2pt);
\fill[red!80!black] (0.0000,8.0000) circle (1.2pt);
\fill[red!80!black] (12.0000,8.0000) circle (1.2pt);
\fill[red!80!black] (0.0000,8.0000) circle (1.2pt);
\fill[red!80!black] (12.0000,8.0000) circle (1.2pt);
\fill[red!80!black] (0.0000,8.0000) circle (1.2pt);
\fill[red!80!black] (12.0000,8.0000) circle (1.2pt);
\fill[red!80!black] (0.0000,8.0000) circle (1.2pt);
\fill[red!80!black] (12.0000,8.0000) circle (1.2pt);
\fill[red!80!black] (0.0000,8.0000) circle (1.2pt);
\fill[red!80!black] (12.0000,8.0000) circle (1.2pt);
\fill[red!80!black] (0.0000,8.0000) circle (1.2pt);
\fill[red!80!black] (12.0000,8.0000) circle (1.2pt);
\fill[red!80!black] (0.0000,8.0000) circle (1.2pt);
\fill[red!80!black] (12.0000,8.0000) circle (1.2pt);
\fill[red!80!black] (0.0000,8.0000) circle (1.2pt);
\fill[red!80!black] (12.0000,8.0000) circle (1.2pt);
\fill[red!80!black] (0.0000,8.0000) circle (1.2pt);
\fill[red!80!black] (12.0000,8.0000) circle (1.2pt);
\fill[red!80!black] (0.0000,8.0000) circle (1.2pt);
\fill[red!80!black] (12.0000,8.0000) circle (1.2pt);
\fill[red!80!black] (0.0000,8.0000) circle (1.2pt);
\fill[red!80!black] (12.0000,8.0000) circle (1.2pt);
\fill[red!80!black] (0.0000,8.0000) circle (1.2pt);
\fill[red!80!black] (12.0000,8.0000) circle (1.2pt);
\fill[red!80!black] (0.0000,8.0000) circle (1.2pt);
\fill[red!80!black] (12.0000,8.0000) circle (1.2pt);
\fill[red!80!black] (0.0000,8.0000) circle (1.2pt);
\fill[red!80!black] (12.0000,8.0000) circle (1.2pt);
\fill[red!80!black] (0.0000,8.0000) circle (1.2pt);
\fill[red!80!black] (12.0000,8.0000) circle (1.2pt);
\fill[red!80!black] (0.0000,8.0000) circle (1.2pt);
\fill[red!80!black] (12.0000,8.0000) circle (1.2pt);
\fill[red!80!black] (0.0000,8.0000) circle (1.2pt);
\fill[red!80!black] (12.0000,8.0000) circle (1.2pt);
\fill[red!80!black] (0.0000,8.0000) circle (1.2pt);
\fill[red!80!black] (12.0000,8.0000) circle (1.2pt);
\fill[red!80!black] (0.0000,8.0000) circle (1.2pt);
\fill[red!80!black] (12.0000,8.0000) circle (1.2pt);
\fill[red!80!black] (0.0000,8.0000) circle (1.2pt);
\fill[red!80!black] (12.0000,8.0000) circle (1.2pt);
\fill[red!80!black] (0.0000,8.0000) circle (1.2pt);
\fill[red!80!black] (12.0000,8.0000) circle (1.2pt);
\fill[red!80!black] (0.0000,8.0000) circle (1.2pt);
\fill[red!80!black] (12.0000,8.0000) circle (1.2pt);
\fill[red!80!black] (0.0000,8.0000) circle (1.2pt);
\fill[red!80!black] (12.0000,8.0000) circle (1.2pt);
\fill[red!80!black] (0.0000,8.0000) circle (1.2pt);
\fill[red!80!black] (12.0000,8.0000) circle (1.2pt);
\fill[red!80!black] (0.0000,8.0000) circle (1.2pt);
\fill[red!80!black] (12.0000,8.0000) circle (1.2pt);
\fill[red!80!black] (0.0000,8.0000) circle (1.2pt);
\fill[red!80!black] (12.0000,8.0000) circle (1.2pt);
\fill[red!80!black] (0.0000,8.0000) circle (1.2pt);
\fill[red!80!black] (12.0000,8.0000) circle (1.2pt);
\fill[red!80!black] (0.0000,8.0000) circle (1.2pt);
\fill[red!80!black] (12.0000,8.0000) circle (1.2pt);
\fill[red!80!black] (0.0000,8.0000) circle (1.2pt);
\fill[red!80!black] (12.0000,8.0000) circle (1.2pt);
\fill[red!80!black] (0.0000,8.0000) circle (1.2pt);
\fill[red!80!black] (12.0000,8.0000) circle (1.2pt);
\fill[red!80!black] (0.0000,8.0000) circle (1.2pt);
\fill[red!80!black] (12.0000,8.0000) circle (1.2pt);
\fill[red!80!black] (0.0000,8.0000) circle (1.2pt);
\fill[red!80!black] (12.0000,8.0000) circle (1.2pt);
\fill[red!80!black] (0.0000,8.0000) circle (1.2pt);
\fill[red!80!black] (12.0000,8.0000) circle (1.2pt);
\fill[red!80!black] (0.0000,8.0000) circle (1.2pt);
\fill[red!80!black] (12.0000,8.0000) circle (1.2pt);
\fill[red!80!black] (0.0000,8.0000) circle (1.2pt);
\fill[red!80!black] (12.0000,8.0000) circle (1.2pt);
\fill[red!80!black] (0.0000,8.0000) circle (1.2pt);
\fill[red!80!black] (12.0000,8.0000) circle (1.2pt);
\fill[red!80!black] (0.0000,8.0000) circle (1.2pt);
\fill[red!80!black] (12.0000,8.0000) circle (1.2pt);
\fill[red!80!black] (0.0000,8.0000) circle (1.2pt);
\fill[red!80!black] (12.0000,8.0000) circle (1.2pt);
\fill[red!80!black] (0.0000,8.0000) circle (1.2pt);
\fill[red!80!black] (12.0000,8.0000) circle (1.2pt);
\fill[red!80!black] (0.0000,8.0000) circle (1.2pt);
\fill[red!80!black] (12.0000,8.0000) circle (1.2pt);
\fill[red!80!black] (0.0000,8.0000) circle (1.2pt);
\fill[red!80!black] (12.0000,8.0000) circle (1.2pt);
\fill[red!80!black] (0.0000,8.0000) circle (1.2pt);
\fill[red!80!black] (12.0000,8.0000) circle (1.2pt);
\fill[red!80!black] (0.0000,8.0000) circle (1.2pt);
\fill[red!80!black] (12.0000,8.0000) circle (1.2pt);
\fill[red!80!black] (0.0000,8.0000) circle (1.2pt);
\fill[red!80!black] (12.0000,8.0000) circle (1.2pt);
\fill[red!80!black] (0.0000,8.0000) circle (1.2pt);
\fill[red!80!black] (12.0000,8.0000) circle (1.2pt);
\fill[red!80!black] (0.0000,8.0000) circle (1.2pt);
\fill[red!80!black] (12.0000,8.0000) circle (1.2pt);
\fill[red!80!black] (0.0000,8.0000) circle (1.2pt);
\fill[red!80!black] (12.0000,8.0000) circle (1.2pt);
\fill[red!80!black] (0.0000,8.0000) circle (1.2pt);
\fill[red!80!black] (12.0000,8.0000) circle (1.2pt);
\fill[red!80!black] (0.0000,8.0000) circle (1.2pt);
\fill[red!80!black] (12.0000,8.0000) circle (1.2pt);
\fill[red!80!black] (0.0000,8.0000) circle (1.2pt);
\fill[red!80!black] (12.0000,8.0000) circle (1.2pt);
\fill[red!80!black] (0.0000,8.0000) circle (1.2pt);
\fill[red!80!black] (12.0000,8.0000) circle (1.2pt);
\fill[red!80!black] (0.0000,8.0000) circle (1.2pt);
\fill[red!80!black] (12.0000,8.0000) circle (1.2pt);
\fill[red!80!black] (0.0000,8.0000) circle (1.2pt);
\fill[red!80!black] (12.0000,8.0000) circle (1.2pt);
\fill[red!80!black] (0.0000,8.0000) circle (1.2pt);
\fill[red!80!black] (12.0000,8.0000) circle (1.2pt);
\fill[red!80!black] (0.0000,8.0000) circle (1.2pt);
\fill[red!80!black] (12.0000,8.0000) circle (1.2pt);
\fill[red!80!black] (0.0000,8.0000) circle (1.2pt);
\fill[red!80!black] (12.0000,8.0000) circle (1.2pt);
\fill[red!80!black] (0.0000,8.0000) circle (1.2pt);
\fill[red!80!black] (12.0000,8.0000) circle (1.2pt);
\fill[red!80!black] (0.0000,8.0000) circle (1.2pt);
\fill[red!80!black] (12.0000,8.0000) circle (1.2pt);
\fill[red!80!black] (0.0000,8.0000) circle (1.2pt);
\fill[red!80!black] (12.0000,8.0000) circle (1.2pt);
\fill[red!80!black] (0.0000,8.0000) circle (1.2pt);
\fill[red!80!black] (12.0000,8.0000) circle (1.2pt);
\fill[red!80!black] (0.0000,8.0000) circle (1.2pt);
\fill[red!80!black] (12.0000,8.0000) circle (1.2pt);
\fill[red!80!black] (0.0000,8.0000) circle (1.2pt);
\fill[red!80!black] (12.0000,8.0000) circle (1.2pt);
\fill[red!80!black] (0.0000,8.0000) circle (1.2pt);
\fill[red!80!black] (12.0000,8.0000) circle (1.2pt);
\fill[red!80!black] (0.0000,8.0000) circle (1.2pt);
\fill[red!80!black] (12.0000,8.0000) circle (1.2pt);
\fill[red!80!black] (0.0000,8.0000) circle (1.2pt);
\fill[red!80!black] (12.0000,8.0000) circle (1.2pt);
\fill[red!80!black] (0.0000,8.0000) circle (1.2pt);
\fill[red!80!black] (12.0000,8.0000) circle (1.2pt);
\fill[red!80!black] (0.0000,8.0000) circle (1.2pt);
\fill[red!80!black] (12.0000,8.0000) circle (1.2pt);
\fill[red!80!black] (0.0000,8.0000) circle (1.2pt);
\fill[red!80!black] (12.0000,8.0000) circle (1.2pt);
\fill[red!80!black] (0.0000,8.0000) circle (1.2pt);
\fill[red!80!black] (12.0000,8.0000) circle (1.2pt);
\fill[red!80!black] (0.0000,8.0000) circle (1.2pt);
\fill[red!80!black] (12.0000,8.0000) circle (1.2pt);
\fill[red!80!black] (0.0000,8.0000) circle (1.2pt);
\fill[red!80!black] (12.0000,8.0000) circle (1.2pt);
\fill[red!80!black] (0.0000,8.0000) circle (1.2pt);
\fill[red!80!black] (12.0000,8.0000) circle (1.2pt);
\fill[red!80!black] (0.0000,8.0000) circle (1.2pt);
\fill[red!80!black] (12.0000,8.0000) circle (1.2pt);
\fill[red!80!black] (0.0000,8.0000) circle (1.2pt);
\fill[red!80!black] (12.0000,8.0000) circle (1.2pt);
\fill[red!80!black] (0.0000,8.0000) circle (1.2pt);
\fill[red!80!black] (12.0000,8.0000) circle (1.2pt);
\fill[red!80!black] (0.0000,8.0000) circle (1.2pt);
\fill[red!80!black] (12.0000,8.0000) circle (1.2pt);
\fill[red!80!black] (0.0000,8.0000) circle (1.2pt);
\fill[red!80!black] (12.0000,8.0000) circle (1.2pt);
\fill[red!80!black] (0.0000,8.0000) circle (1.2pt);
\fill[red!80!black] (12.0000,8.0000) circle (1.2pt);
\fill[red!80!black] (0.0000,8.0000) circle (1.2pt);
\fill[red!80!black] (12.0000,8.0000) circle (1.2pt);
\fill[red!80!black] (0.0000,8.0000) circle (1.2pt);
\fill[red!80!black] (12.0000,8.0000) circle (1.2pt);
\fill[red!80!black] (0.0000,8.0000) circle (1.2pt);
\fill[red!80!black] (12.0000,8.0000) circle (1.2pt);
\fill[red!80!black] (0.0000,8.0000) circle (1.2pt);
\fill[red!80!black] (12.0000,8.0000) circle (1.2pt);
\fill[red!80!black] (0.0000,8.0000) circle (1.2pt);
\fill[red!80!black] (12.0000,8.0000) circle (1.2pt);
\fill[red!80!black] (0.0000,8.0000) circle (1.2pt);
\fill[red!80!black] (12.0000,8.0000) circle (1.2pt);
\fill[red!80!black] (0.0000,8.0000) circle (1.2pt);
\fill[red!80!black] (12.0000,8.0000) circle (1.2pt);
\fill[red!80!black] (0.0000,8.0000) circle (1.2pt);
\fill[red!80!black] (12.0000,8.0000) circle (1.2pt);
\fill[red!80!black] (0.0000,8.0000) circle (1.2pt);
\fill[red!80!black] (12.0000,8.0000) circle (1.2pt);
\fill[red!80!black] (0.0000,8.0000) circle (1.2pt);
\fill[red!80!black] (12.0000,8.0000) circle (1.2pt);
\fill[red!80!black] (0.0000,8.0000) circle (1.2pt);
\fill[red!80!black] (12.0000,8.0000) circle (1.2pt);
\fill[red!80!black] (0.0000,8.0000) circle (1.2pt);
\fill[red!80!black] (12.0000,8.0000) circle (1.2pt);
\fill[red!80!black] (0.0000,8.0000) circle (1.2pt);
\fill[red!80!black] (12.0000,8.0000) circle (1.2pt);
\fill[red!80!black] (0.0000,8.0000) circle (1.2pt);
\fill[red!80!black] (12.0000,8.0000) circle (1.2pt);
\fill[red!80!black] (0.0000,8.0000) circle (1.2pt);
\fill[red!80!black] (12.0000,8.0000) circle (1.2pt);
\fill[red!80!black] (0.0000,8.0000) circle (1.2pt);
\fill[red!80!black] (12.0000,8.0000) circle (1.2pt);
\fill[red!80!black] (0.0000,8.0000) circle (1.2pt);
\fill[red!80!black] (12.0000,8.0000) circle (1.2pt);
\fill[red!80!black] (0.0000,8.0000) circle (1.2pt);
\fill[red!80!black] (12.0000,8.0000) circle (1.2pt);
\fill[red!80!black] (0.0000,8.0000) circle (1.2pt);
\fill[red!80!black] (12.0000,8.0000) circle (1.2pt);
\fill[red!80!black] (0.0000,8.0000) circle (1.2pt);
\fill[red!80!black] (12.0000,8.0000) circle (1.2pt);
\fill[red!80!black] (0.0000,8.0000) circle (1.2pt);
\fill[red!80!black] (12.0000,8.0000) circle (1.2pt);
\fill[red!80!black] (0.0000,8.0000) circle (1.2pt);
\fill[red!80!black] (12.0000,8.0000) circle (1.2pt);
\fill[red!80!black] (0.0000,8.0000) circle (1.2pt);
\fill[red!80!black] (12.0000,8.0000) circle (1.2pt);
\fill[red!80!black] (0.0000,8.0000) circle (1.2pt);
\fill[red!80!black] (12.0000,8.0000) circle (1.2pt);
\fill[red!80!black] (0.0000,8.0000) circle (1.2pt);
\fill[red!80!black] (12.0000,8.0000) circle (1.2pt);
\fill[red!80!black] (0.0000,8.0000) circle (1.2pt);
\fill[red!80!black] (12.0000,8.0000) circle (1.2pt);
\fill[red!80!black] (0.0000,8.0000) circle (1.2pt);
\fill[red!80!black] (12.0000,8.0000) circle (1.2pt);
\fill[red!80!black] (0.0000,8.0000) circle (1.2pt);
\fill[red!80!black] (12.0000,8.0000) circle (1.2pt);
\fill[red!80!black] (0.0000,8.0000) circle (1.2pt);
\fill[red!80!black] (12.0000,8.0000) circle (1.2pt);
\fill[red!80!black] (0.0000,8.0000) circle (1.2pt);
\fill[red!80!black] (12.0000,8.0000) circle (1.2pt);
\fill[red!80!black] (0.0000,8.0000) circle (1.2pt);
\fill[red!80!black] (12.0000,8.0000) circle (1.2pt);
\fill[red!80!black] (0.0000,8.0000) circle (1.2pt);
\fill[red!80!black] (12.0000,8.0000) circle (1.2pt);
\fill[red!80!black] (0.0000,8.0000) circle (1.2pt);
\fill[red!80!black] (12.0000,8.0000) circle (1.2pt);
\fill[red!80!black] (0.0000,8.0000) circle (1.2pt);
\fill[red!80!black] (12.0000,8.0000) circle (1.2pt);
\fill[red!80!black] (0.0000,8.0000) circle (1.2pt);
\fill[red!80!black] (12.0000,8.0000) circle (1.2pt);
\fill[red!80!black] (0.0000,8.0000) circle (1.2pt);
\fill[red!80!black] (12.0000,8.0000) circle (1.2pt);
\fill[red!80!black] (0.0000,8.0000) circle (1.2pt);
\fill[red!80!black] (12.0000,8.0000) circle (1.2pt);
\fill[red!80!black] (0.0000,8.0000) circle (1.2pt);
\fill[red!80!black] (12.0000,8.0000) circle (1.2pt);
\fill[red!80!black] (0.0000,8.0000) circle (1.2pt);
\fill[red!80!black] (12.0000,8.0000) circle (1.2pt);
\fill[red!80!black] (0.0000,8.0000) circle (1.2pt);
\fill[red!80!black] (12.0000,8.0000) circle (1.2pt);
\fill[red!80!black] (0.0000,8.0000) circle (1.2pt);
\fill[red!80!black] (12.0000,8.0000) circle (1.2pt);
\fill[red!80!black] (0.0000,8.0000) circle (1.2pt);
\fill[red!80!black] (12.0000,8.0000) circle (1.2pt);
\fill[red!80!black] (0.0000,8.0000) circle (1.2pt);
\fill[red!80!black] (12.0000,8.0000) circle (1.2pt);
\fill[red!80!black] (0.0000,8.0000) circle (1.2pt);
\fill[red!80!black] (12.0000,8.0000) circle (1.2pt);
\fill[red!80!black] (0.0000,8.0000) circle (1.2pt);
\fill[red!80!black] (12.0000,8.0000) circle (1.2pt);
\fill[red!80!black] (0.0000,8.0000) circle (1.2pt);
\fill[red!80!black] (12.0000,8.0000) circle (1.2pt);
\fill[red!80!black] (0.0000,8.0000) circle (1.2pt);
\fill[red!80!black] (12.0000,8.0000) circle (1.2pt);
\fill[red!80!black] (0.0000,8.0000) circle (1.2pt);
\fill[red!80!black] (12.0000,8.0000) circle (1.2pt);
\fill[red!80!black] (0.0000,8.0000) circle (1.2pt);
\fill[red!80!black] (12.0000,8.0000) circle (1.2pt);
\fill[red!80!black] (0.0000,8.0000) circle (1.2pt);
\fill[red!80!black] (12.0000,8.0000) circle (1.2pt);
\fill[red!80!black] (0.0000,8.0000) circle (1.2pt);
\fill[red!80!black] (12.0000,8.0000) circle (1.2pt);
\fill[red!80!black] (0.0000,8.0000) circle (1.2pt);
\fill[red!80!black] (12.0000,8.0000) circle (1.2pt);
\fill[red!80!black] (0.0000,8.0000) circle (1.2pt);
\fill[red!80!black] (12.0000,8.0000) circle (1.2pt);
\fill[red!80!black] (0.0000,8.0000) circle (1.2pt);
\fill[red!80!black] (12.0000,8.0000) circle (1.2pt);
\fill[red!80!black] (0.0000,8.0000) circle (1.2pt);
\fill[red!80!black] (12.0000,8.0000) circle (1.2pt);
\fill[red!80!black] (0.0000,8.0000) circle (1.2pt);
\fill[red!80!black] (12.0000,8.0000) circle (1.2pt);
\fill[red!80!black] (0.0000,8.0000) circle (1.2pt);
\fill[red!80!black] (12.0000,8.0000) circle (1.2pt);
\fill[red!80!black] (0.0000,8.0000) circle (1.2pt);
\fill[red!80!black] (12.0000,8.0000) circle (1.2pt);
\fill[red!80!black] (0.0000,8.0000) circle (1.2pt);
\fill[red!80!black] (12.0000,8.0000) circle (1.2pt);
\fill[red!80!black] (0.0000,8.0000) circle (1.2pt);
\fill[red!80!black] (12.0000,8.0000) circle (1.2pt);
\fill[red!80!black] (0.0000,8.0000) circle (1.2pt);
\fill[red!80!black] (12.0000,8.0000) circle (1.2pt);
\fill[red!80!black] (0.0000,8.0000) circle (1.2pt);
\fill[red!80!black] (12.0000,8.0000) circle (1.2pt);
\fill[red!80!black] (0.0000,8.0000) circle (1.2pt);
\fill[red!80!black] (12.0000,8.0000) circle (1.2pt);
\fill[red!80!black] (0.0000,8.0000) circle (1.2pt);
\fill[red!80!black] (12.0000,8.0000) circle (1.2pt);
\fill[red!80!black] (0.0000,8.0000) circle (1.2pt);
\fill[red!80!black] (12.0000,8.0000) circle (1.2pt);
\fill[red!80!black] (0.0000,8.0000) circle (1.2pt);
\fill[red!80!black] (12.0000,8.0000) circle (1.2pt);
\fill[red!80!black] (0.0000,8.0000) circle (1.2pt);
\fill[red!80!black] (12.0000,8.0000) circle (1.2pt);
\fill[red!80!black] (0.0000,8.0000) circle (1.2pt);
\fill[red!80!black] (12.0000,8.0000) circle (1.2pt);
\fill[red!80!black] (0.0000,8.0000) circle (1.2pt);
\fill[red!80!black] (12.0000,8.0000) circle (1.2pt);
\fill[red!80!black] (0.0000,8.0000) circle (1.2pt);
\fill[red!80!black] (12.0000,8.0000) circle (1.2pt);
\fill[red!80!black] (0.0000,8.0000) circle (1.2pt);
\fill[red!80!black] (12.0000,8.0000) circle (1.2pt);
\fill[red!80!black] (0.0000,8.0000) circle (1.2pt);
\fill[red!80!black] (12.0000,8.0000) circle (1.2pt);
\fill[red!80!black] (0.0000,8.0000) circle (1.2pt);
\fill[red!80!black] (12.0000,8.0000) circle (1.2pt);
\fill[red!80!black] (0.0000,8.0000) circle (1.2pt);
\fill[red!80!black] (12.0000,8.0000) circle (1.2pt);
\fill[red!80!black] (0.0000,8.0000) circle (1.2pt);
\fill[red!80!black] (12.0000,8.0000) circle (1.2pt);
\fill[red!80!black] (0.0000,8.0000) circle (1.2pt);
\fill[red!80!black] (12.0000,8.0000) circle (1.2pt);
\fill[red!80!black] (0.0000,8.0000) circle (1.2pt);
\fill[red!80!black] (12.0000,8.0000) circle (1.2pt);
\fill[red!80!black] (0.0000,8.0000) circle (1.2pt);
\fill[red!80!black] (12.0000,8.0000) circle (1.2pt);
\fill[red!80!black] (0.0000,8.0000) circle (1.2pt);
\fill[red!80!black] (12.0000,8.0000) circle (1.2pt);
\fill[red!80!black] (0.0000,8.0000) circle (1.2pt);
\fill[red!80!black] (12.0000,8.0000) circle (1.2pt);
\fill[red!80!black] (0.0000,8.0000) circle (1.2pt);
\fill[red!80!black] (12.0000,8.0000) circle (1.2pt);
\fill[red!80!black] (0.0000,8.0000) circle (1.2pt);
\fill[red!80!black] (12.0000,8.0000) circle (1.2pt);
\fill[red!80!black] (0.0000,8.0000) circle (1.2pt);
\fill[red!80!black] (12.0000,8.0000) circle (1.2pt);
\fill[red!80!black] (0.0000,8.0000) circle (1.2pt);
\fill[red!80!black] (12.0000,8.0000) circle (1.2pt);
\fill[red!80!black] (0.0000,8.0000) circle (1.2pt);
\fill[red!80!black] (12.0000,8.0000) circle (1.2pt);
\fill[red!80!black] (0.0000,8.0000) circle (1.2pt);
\fill[red!80!black] (12.0000,8.0000) circle (1.2pt);
\fill[red!80!black] (0.0000,8.0000) circle (1.2pt);
\fill[red!80!black] (12.0000,8.0000) circle (1.2pt);
\fill[red!80!black] (0.0000,8.0000) circle (1.2pt);
\fill[red!80!black] (12.0000,8.0000) circle (1.2pt);
\fill[red!80!black] (0.0000,8.0000) circle (1.2pt);
\fill[red!80!black] (12.0000,8.0000) circle (1.2pt);
\fill[red!80!black] (0.0000,8.0000) circle (1.2pt);
\fill[red!80!black] (12.0000,8.0000) circle (1.2pt);
\fill[red!80!black] (0.0000,8.0000) circle (1.2pt);
\fill[red!80!black] (12.0000,8.0000) circle (1.2pt);
\fill[red!80!black] (0.0000,8.0000) circle (1.2pt);
\fill[red!80!black] (12.0000,8.0000) circle (1.2pt);
\fill[red!80!black] (0.0000,8.0000) circle (1.2pt);
\fill[red!80!black] (12.0000,8.0000) circle (1.2pt);
\fill[red!80!black] (0.0000,8.0000) circle (1.2pt);
\fill[red!80!black] (12.0000,8.0000) circle (1.2pt);
\fill[red!80!black] (0.0000,8.0000) circle (1.2pt);
\fill[red!80!black] (12.0000,8.0000) circle (1.2pt);
\fill[red!80!black] (0.0000,8.0000) circle (1.2pt);
\fill[red!80!black] (12.0000,8.0000) circle (1.2pt);
\fill[red!80!black] (0.0000,8.0000) circle (1.2pt);
\fill[red!80!black] (12.0000,8.0000) circle (1.2pt);
\fill[red!80!black] (0.0000,8.0000) circle (1.2pt);
\fill[red!80!black] (12.0000,8.0000) circle (1.2pt);
\fill[red!80!black] (0.0000,8.0000) circle (1.2pt);
\fill[red!80!black] (12.0000,8.0000) circle (1.2pt);
\fill[red!80!black] (0.0000,8.0000) circle (1.2pt);
\fill[red!80!black] (12.0000,8.0000) circle (1.2pt);
\fill[red!80!black] (0.0000,8.0000) circle (1.2pt);
\fill[red!80!black] (12.0000,8.0000) circle (1.2pt);
\fill[red!80!black] (0.0000,8.0000) circle (1.2pt);
\fill[red!80!black] (12.0000,8.0000) circle (1.2pt);
\fill[red!80!black] (0.0000,8.0000) circle (1.2pt);
\fill[red!80!black] (12.0000,8.0000) circle (1.2pt);
\fill[red!80!black] (0.0000,8.0000) circle (1.2pt);
\fill[red!80!black] (12.0000,8.0000) circle (1.2pt);
\fill[red!80!black] (0.0000,8.0000) circle (1.2pt);
\fill[red!80!black] (12.0000,8.0000) circle (1.2pt);
\fill[red!80!black] (0.0000,8.0000) circle (1.2pt);
\fill[red!80!black] (12.0000,8.0000) circle (1.2pt);
\fill[red!80!black] (0.0000,8.0000) circle (1.2pt);
\fill[red!80!black] (12.0000,8.0000) circle (1.2pt);
\fill[red!80!black] (0.0000,8.0000) circle (1.2pt);
\fill[red!80!black] (12.0000,8.0000) circle (1.2pt);
\fill[red!80!black] (0.0000,8.0000) circle (1.2pt);
\fill[red!80!black] (12.0000,8.0000) circle (1.2pt);
\fill[red!80!black] (0.0000,8.0000) circle (1.2pt);
\fill[red!80!black] (12.0000,8.0000) circle (1.2pt);
\fill[red!80!black] (0.0000,8.0000) circle (1.2pt);
\fill[red!80!black] (12.0000,8.0000) circle (1.2pt);
\fill[red!80!black] (0.0000,8.0000) circle (1.2pt);
\fill[red!80!black] (12.0000,8.0000) circle (1.2pt);
\fill[red!80!black] (0.0000,8.0000) circle (1.2pt);
\fill[red!80!black] (12.0000,8.0000) circle (1.2pt);
\fill[red!80!black] (0.0000,8.0000) circle (1.2pt);
\fill[red!80!black] (12.0000,8.0000) circle (1.2pt);
\fill[red!80!black] (0.0000,8.0000) circle (1.2pt);
\fill[red!80!black] (12.0000,8.0000) circle (1.2pt);
\fill[red!80!black] (0.0000,8.0000) circle (1.2pt);
\fill[red!80!black] (12.0000,8.0000) circle (1.2pt);
\fill[red!80!black] (0.0000,8.0000) circle (1.2pt);
\fill[red!80!black] (12.0000,8.0000) circle (1.2pt);
\fill[red!80!black] (0.0000,8.0000) circle (1.2pt);
\fill[red!80!black] (12.0000,8.0000) circle (1.2pt);
\fill[red!80!black] (0.0000,8.0000) circle (1.2pt);
\fill[red!80!black] (12.0000,8.0000) circle (1.2pt);
\fill[red!80!black] (0.0000,8.0000) circle (1.2pt);
\fill[red!80!black] (12.0000,8.0000) circle (1.2pt);
\fill[red!80!black] (0.0000,8.0000) circle (1.2pt);
\fill[red!80!black] (12.0000,8.0000) circle (1.2pt);
\fill[red!80!black] (0.0000,8.0000) circle (1.2pt);
\fill[red!80!black] (12.0000,8.0000) circle (1.2pt);
\fill[red!80!black] (0.0000,8.0000) circle (1.2pt);
\fill[red!80!black] (12.0000,8.0000) circle (1.2pt);
\fill[red!80!black] (0.0000,8.0000) circle (1.2pt);
\fill[red!80!black] (12.0000,8.0000) circle (1.2pt);
\fill[red!80!black] (0.0000,8.0000) circle (1.2pt);
\fill[red!80!black] (12.0000,8.0000) circle (1.2pt);
\fill[red!80!black] (0.0000,8.0000) circle (1.2pt);
\fill[red!80!black] (12.0000,8.0000) circle (1.2pt);
\fill[red!80!black] (0.0000,8.0000) circle (1.2pt);
\fill[red!80!black] (12.0000,8.0000) circle (1.2pt);
\fill[red!80!black] (0.0000,8.0000) circle (1.2pt);
\fill[red!80!black] (12.0000,8.0000) circle (1.2pt);
\fill[red!80!black] (0.0000,8.0000) circle (1.2pt);
\fill[red!80!black] (12.0000,8.0000) circle (1.2pt);
\fill[red!80!black] (0.0000,8.0000) circle (1.2pt);
\fill[red!80!black] (12.0000,8.0000) circle (1.2pt);
\fill[red!80!black] (0.0000,8.0000) circle (1.2pt);
\fill[red!80!black] (12.0000,8.0000) circle (1.2pt);
\fill[red!80!black] (0.0000,8.0000) circle (1.2pt);
\fill[red!80!black] (12.0000,8.0000) circle (1.2pt);
\fill[red!80!black] (0.0000,8.0000) circle (1.2pt);
\fill[red!80!black] (12.0000,8.0000) circle (1.2pt);
\fill[red!80!black] (0.0000,8.0000) circle (1.2pt);
\fill[red!80!black] (12.0000,8.0000) circle (1.2pt);
\fill[red!80!black] (0.0000,8.0000) circle (1.2pt);
\fill[red!80!black] (12.0000,8.0000) circle (1.2pt);
\fill[red!80!black] (0.0000,8.0000) circle (1.2pt);
\fill[red!80!black] (12.0000,8.0000) circle (1.2pt);
\fill[red!80!black] (0.0000,8.0000) circle (1.2pt);
\fill[red!80!black] (12.0000,8.0000) circle (1.2pt);
\fill[red!80!black] (0.0000,8.0000) circle (1.2pt);
\fill[red!80!black] (12.0000,8.0000) circle (1.2pt);
\fill[red!80!black] (0.0000,8.0000) circle (1.2pt);
\fill[red!80!black] (12.0000,8.0000) circle (1.2pt);
\fill[red!80!black] (0.0000,8.0000) circle (1.2pt);
\fill[red!80!black] (12.0000,8.0000) circle (1.2pt);
\fill[red!80!black] (0.0000,8.0000) circle (1.2pt);
\fill[red!80!black] (12.0000,8.0000) circle (1.2pt);
\fill[red!80!black] (0.0000,8.0000) circle (1.2pt);
\fill[red!80!black] (12.0000,8.0000) circle (1.2pt);
\fill[red!80!black] (0.0000,8.0000) circle (1.2pt);
\fill[red!80!black] (12.0000,8.0000) circle (1.2pt);
\fill[red!80!black] (0.0000,8.0000) circle (1.2pt);
\fill[red!80!black] (12.0000,8.0000) circle (1.2pt);
\fill[red!80!black] (0.0000,8.0000) circle (1.2pt);
\fill[red!80!black] (12.0000,8.0000) circle (1.2pt);
\fill[red!80!black] (0.0000,8.0000) circle (1.2pt);
\fill[red!80!black] (12.0000,8.0000) circle (1.2pt);
\fill[red!80!black] (0.0000,8.0000) circle (1.2pt);
\fill[red!80!black] (12.0000,8.0000) circle (1.2pt);
\fill[red!80!black] (0.0000,8.0000) circle (1.2pt);
\fill[red!80!black] (12.0000,8.0000) circle (1.2pt);
\fill[red!80!black] (0.0000,8.0000) circle (1.2pt);
\fill[red!80!black] (12.0000,8.0000) circle (1.2pt);
\fill[red!80!black] (0.0000,8.0000) circle (1.2pt);
\fill[red!80!black] (12.0000,8.0000) circle (1.2pt);
\fill[red!80!black] (0.0000,8.0000) circle (1.2pt);
\fill[red!80!black] (12.0000,8.0000) circle (1.2pt);
\fill[red!80!black] (0.0000,8.0000) circle (1.2pt);
\fill[red!80!black] (12.0000,8.0000) circle (1.2pt);
\fill[red!80!black] (0.0000,8.0000) circle (1.2pt);
\fill[red!80!black] (12.0000,8.0000) circle (1.2pt);
\fill[red!80!black] (0.0000,8.0000) circle (1.2pt);
\fill[red!80!black] (12.0000,8.0000) circle (1.2pt);
\fill[red!80!black] (0.0000,8.0000) circle (1.2pt);
\fill[red!80!black] (12.0000,8.0000) circle (1.2pt);
\fill[red!80!black] (0.0000,8.0000) circle (1.2pt);
\fill[red!80!black] (12.0000,8.0000) circle (1.2pt);
\fill[red!80!black] (0.0000,8.0000) circle (1.2pt);
\fill[red!80!black] (12.0000,8.0000) circle (1.2pt);
\fill[red!80!black] (0.0000,8.0000) circle (1.2pt);
\fill[red!80!black] (12.0000,8.0000) circle (1.2pt);
\fill[red!80!black] (0.0000,8.0000) circle (1.2pt);
\fill[red!80!black] (12.0000,8.0000) circle (1.2pt);
\fill[red!80!black] (0.0000,8.0000) circle (1.2pt);
\fill[red!80!black] (12.0000,8.0000) circle (1.2pt);
\fill[red!80!black] (0.0000,8.0000) circle (1.2pt);
\fill[red!80!black] (12.0000,8.0000) circle (1.2pt);
\fill[red!80!black] (0.0000,8.0000) circle (1.2pt);
\fill[red!80!black] (12.0000,8.0000) circle (1.2pt);
\fill[red!80!black] (0.0000,8.0000) circle (1.2pt);
\fill[red!80!black] (12.0000,8.0000) circle (1.2pt);
\fill[red!80!black] (0.0000,8.0000) circle (1.2pt);
\fill[red!80!black] (12.0000,8.0000) circle (1.2pt);
\fill[red!80!black] (0.0000,8.0000) circle (1.2pt);
\fill[red!80!black] (12.0000,8.0000) circle (1.2pt);
\fill[red!80!black] (0.0000,8.0000) circle (1.2pt);
\fill[red!80!black] (12.0000,8.0000) circle (1.2pt);
\fill[red!80!black] (0.0000,8.0000) circle (1.2pt);
\fill[red!80!black] (12.0000,8.0000) circle (1.2pt);
\fill[red!80!black] (0.0000,8.0000) circle (1.2pt);
\fill[red!80!black] (12.0000,8.0000) circle (1.2pt);
\fill[red!80!black] (0.0000,8.0000) circle (1.2pt);
\fill[red!80!black] (12.0000,8.0000) circle (1.2pt);
\fill[red!80!black] (0.0000,8.0000) circle (1.2pt);
\fill[red!80!black] (12.0000,8.0000) circle (1.2pt);
\fill[red!80!black] (0.0000,8.0000) circle (1.2pt);
\fill[red!80!black] (12.0000,8.0000) circle (1.2pt);
\fill[red!80!black] (0.0000,8.0000) circle (1.2pt);
\fill[red!80!black] (12.0000,8.0000) circle (1.2pt);
\fill[red!80!black] (0.0000,8.0000) circle (1.2pt);
\fill[red!80!black] (12.0000,8.0000) circle (1.2pt);
\fill[red!80!black] (0.0000,8.0000) circle (1.2pt);
\fill[red!80!black] (12.0000,8.0000) circle (1.2pt);
\fill[red!80!black] (0.0000,8.0000) circle (1.2pt);
\fill[red!80!black] (12.0000,8.0000) circle (1.2pt);
\fill[red!80!black] (0.0000,8.0000) circle (1.2pt);
\fill[red!80!black] (12.0000,8.0000) circle (1.2pt);
\fill[red!80!black] (0.0000,8.0000) circle (1.2pt);
\fill[red!80!black] (12.0000,8.0000) circle (1.2pt);
\fill[red!80!black] (0.0000,8.0000) circle (1.2pt);
\fill[red!80!black] (12.0000,8.0000) circle (1.2pt);
\fill[red!80!black] (0.0000,8.0000) circle (1.2pt);
\fill[red!80!black] (12.0000,8.0000) circle (1.2pt);
\fill[red!80!black] (0.0000,8.0000) circle (1.2pt);
\fill[red!80!black] (12.0000,8.0000) circle (1.2pt);
\fill[red!80!black] (0.0000,8.0000) circle (1.2pt);
\fill[red!80!black] (12.0000,8.0000) circle (1.2pt);
\fill[red!80!black] (0.0000,8.0000) circle (1.2pt);
\fill[red!80!black] (12.0000,8.0000) circle (1.2pt);
\fill[red!80!black] (0.0000,8.0000) circle (1.2pt);
\fill[red!80!black] (12.0000,8.0000) circle (1.2pt);
\fill[red!80!black] (0.0000,8.0000) circle (1.2pt);
\fill[red!80!black] (12.0000,8.0000) circle (1.2pt);
\fill[red!80!black] (0.0000,8.0000) circle (1.2pt);
\fill[red!80!black] (12.0000,8.0000) circle (1.2pt);
\fill[red!80!black] (0.0000,8.0000) circle (1.2pt);
\fill[red!80!black] (12.0000,8.0000) circle (1.2pt);
\fill[red!80!black] (0.0000,8.0000) circle (1.2pt);
\fill[red!80!black] (12.0000,8.0000) circle (1.2pt);
\fill[red!80!black] (0.0000,8.0000) circle (1.2pt);
\fill[red!80!black] (12.0000,8.0000) circle (1.2pt);
\fill[red!80!black] (0.0000,8.0000) circle (1.2pt);
\fill[red!80!black] (12.0000,8.0000) circle (1.2pt);
\fill[red!80!black] (0.0000,8.0000) circle (1.2pt);
\fill[red!80!black] (12.0000,8.0000) circle (1.2pt);
\fill[red!80!black] (0.0000,8.0000) circle (1.2pt);
\fill[red!80!black] (12.0000,8.0000) circle (1.2pt);
\fill[red!80!black] (0.0000,8.0000) circle (1.2pt);
\fill[red!80!black] (12.0000,8.0000) circle (1.2pt);
\fill[red!80!black] (0.0000,8.0000) circle (1.2pt);
\fill[red!80!black] (12.0000,8.0000) circle (1.2pt);
\fill[red!80!black] (0.0000,8.0000) circle (1.2pt);
\fill[red!80!black] (12.0000,8.0000) circle (1.2pt);
\fill[red!80!black] (0.0000,8.0000) circle (1.2pt);
\fill[red!80!black] (12.0000,8.0000) circle (1.2pt);
\fill[red!80!black] (0.0000,8.0000) circle (1.2pt);
\fill[red!80!black] (12.0000,8.0000) circle (1.2pt);
\fill[red!80!black] (0.0000,8.0000) circle (1.2pt);
\fill[red!80!black] (12.0000,8.0000) circle (1.2pt);
\fill[red!80!black] (0.0000,8.0000) circle (1.2pt);
\fill[red!80!black] (12.0000,8.0000) circle (1.2pt);
\fill[red!80!black] (0.0000,8.0000) circle (1.2pt);
\fill[red!80!black] (12.0000,8.0000) circle (1.2pt);
\fill[red!80!black] (0.0000,8.0000) circle (1.2pt);
\fill[red!80!black] (12.0000,8.0000) circle (1.2pt);
\fill[red!80!black] (0.0000,8.0000) circle (1.2pt);
\fill[red!80!black] (12.0000,8.0000) circle (1.2pt);
\fill[red!80!black] (0.0000,8.0000) circle (1.2pt);
\fill[red!80!black] (12.0000,8.0000) circle (1.2pt);
\fill[red!80!black] (0.0000,8.0000) circle (1.2pt);
\fill[red!80!black] (12.0000,8.0000) circle (1.2pt);
\fill[red!80!black] (0.0000,8.0000) circle (1.2pt);
\fill[red!80!black] (12.0000,8.0000) circle (1.2pt);
\fill[red!80!black] (0.0000,8.0000) circle (1.2pt);
\fill[red!80!black] (12.0000,8.0000) circle (1.2pt);
\fill[red!80!black] (0.0000,8.0000) circle (1.2pt);
\fill[red!80!black] (12.0000,8.0000) circle (1.2pt);
\fill[red!80!black] (0.0000,8.0000) circle (1.2pt);
\fill[red!80!black] (12.0000,8.0000) circle (1.2pt);
\fill[red!80!black] (0.0000,8.0000) circle (1.2pt);
\fill[red!80!black] (12.0000,8.0000) circle (1.2pt);
\fill[red!80!black] (0.0000,8.0000) circle (1.2pt);
\fill[red!80!black] (12.0000,8.0000) circle (1.2pt);
\fill[red!80!black] (0.0000,8.0000) circle (1.2pt);
\fill[red!80!black] (12.0000,8.0000) circle (1.2pt);
\fill[red!80!black] (0.0000,8.0000) circle (1.2pt);
\fill[red!80!black] (12.0000,8.0000) circle (1.2pt);
\fill[red!80!black] (0.0000,8.0000) circle (1.2pt);
\fill[red!80!black] (12.0000,8.0000) circle (1.2pt);
\fill[red!80!black] (0.0000,8.0000) circle (1.2pt);
\fill[red!80!black] (12.0000,8.0000) circle (1.2pt);
\fill[red!80!black] (0.0000,8.0000) circle (1.2pt);
\fill[red!80!black] (12.0000,8.0000) circle (1.2pt);
\fill[red!80!black] (0.0000,8.0000) circle (1.2pt);
\fill[red!80!black] (12.0000,8.0000) circle (1.2pt);
\fill[red!80!black] (0.0000,8.0000) circle (1.2pt);
\fill[red!80!black] (12.0000,8.0000) circle (1.2pt);
\fill[red!80!black] (0.0000,8.0000) circle (1.2pt);
\fill[red!80!black] (12.0000,8.0000) circle (1.2pt);
\fill[red!80!black] (0.0000,8.0000) circle (1.2pt);
\fill[red!80!black] (12.0000,8.0000) circle (1.2pt);
\fill[red!80!black] (0.0000,8.0000) circle (1.2pt);
\fill[red!80!black] (12.0000,8.0000) circle (1.2pt);
\fill[red!80!black] (0.0000,8.0000) circle (1.2pt);
\fill[red!80!black] (12.0000,8.0000) circle (1.2pt);
\fill[red!80!black] (0.0000,8.0000) circle (1.2pt);
\fill[red!80!black] (12.0000,8.0000) circle (1.2pt);
\fill[red!80!black] (0.0000,8.0000) circle (1.2pt);
\fill[red!80!black] (12.0000,8.0000) circle (1.2pt);
\fill[red!80!black] (0.0000,8.0000) circle (1.2pt);
\fill[red!80!black] (12.0000,8.0000) circle (1.2pt);
\fill[red!80!black] (0.0000,8.0000) circle (1.2pt);
\fill[red!80!black] (12.0000,8.0000) circle (1.2pt);
\fill[red!80!black] (0.0000,8.0000) circle (1.2pt);
\fill[red!80!black] (12.0000,8.0000) circle (1.2pt);
\fill[red!80!black] (0.0000,8.0000) circle (1.2pt);
\fill[red!80!black] (12.0000,8.0000) circle (1.2pt);
\fill[red!80!black] (0.0000,8.0000) circle (1.2pt);
\fill[red!80!black] (12.0000,8.0000) circle (1.2pt);
\fill[red!80!black] (0.0000,8.0000) circle (1.2pt);
\fill[red!80!black] (12.0000,8.0000) circle (1.2pt);
\fill[red!80!black] (0.0000,8.0000) circle (1.2pt);
\fill[red!80!black] (12.0000,8.0000) circle (1.2pt);
\fill[red!80!black] (0.0000,8.0000) circle (1.2pt);
\fill[red!80!black] (12.0000,8.0000) circle (1.2pt);
\fill[red!80!black] (0.0000,8.0000) circle (1.2pt);
\fill[red!80!black] (12.0000,8.0000) circle (1.2pt);
\fill[red!80!black] (0.0000,8.0000) circle (1.2pt);
\fill[red!80!black] (12.0000,8.0000) circle (1.2pt);
\fill[red!80!black] (0.0000,8.0000) circle (1.2pt);
\fill[red!80!black] (12.0000,8.0000) circle (1.2pt);
\fill[red!80!black] (0.0000,8.0000) circle (1.2pt);
\fill[red!80!black] (12.0000,8.0000) circle (1.2pt);
\fill[red!80!black] (0.0000,8.0000) circle (1.2pt);
\fill[red!80!black] (12.0000,8.0000) circle (1.2pt);
\fill[red!80!black] (0.0000,8.0000) circle (1.2pt);
\fill[red!80!black] (12.0000,8.0000) circle (1.2pt);
\fill[red!80!black] (0.0000,8.0000) circle (1.2pt);
\fill[red!80!black] (12.0000,8.0000) circle (1.2pt);
\fill[red!80!black] (0.0000,8.0000) circle (1.2pt);
\fill[red!80!black] (12.0000,8.0000) circle (1.2pt);
\fill[red!80!black] (0.0000,8.0000) circle (1.2pt);
\fill[red!80!black] (12.0000,8.0000) circle (1.2pt);
\fill[red!80!black] (0.0000,8.0000) circle (1.2pt);
\fill[red!80!black] (12.0000,8.0000) circle (1.2pt);
\fill[red!80!black] (0.0000,8.0000) circle (1.2pt);
\fill[red!80!black] (12.0000,8.0000) circle (1.2pt);
\fill[red!80!black] (0.0000,8.0000) circle (1.2pt);
\fill[red!80!black] (12.0000,8.0000) circle (1.2pt);
\fill[red!80!black] (0.0000,8.0000) circle (1.2pt);
\fill[red!80!black] (12.0000,8.0000) circle (1.2pt);
\fill[red!80!black] (0.0000,8.0000) circle (1.2pt);
\fill[red!80!black] (12.0000,8.0000) circle (1.2pt);
\fill[red!80!black] (0.0000,8.0000) circle (1.2pt);
\fill[red!80!black] (12.0000,8.0000) circle (1.2pt);
\fill[red!80!black] (0.0000,8.0000) circle (1.2pt);
\fill[red!80!black] (12.0000,8.0000) circle (1.2pt);
\fill[red!80!black] (0.0000,8.0000) circle (1.2pt);
\fill[red!80!black] (12.0000,8.0000) circle (1.2pt);
\fill[red!80!black] (0.0000,8.0000) circle (1.2pt);
\fill[red!80!black] (12.0000,8.0000) circle (1.2pt);
\fill[red!80!black] (0.0000,8.0000) circle (1.2pt);
\fill[red!80!black] (12.0000,8.0000) circle (1.2pt);
\fill[red!80!black] (0.0000,8.0000) circle (1.2pt);
\fill[red!80!black] (12.0000,8.0000) circle (1.2pt);
\fill[red!80!black] (0.0000,8.0000) circle (1.2pt);
\fill[red!80!black] (12.0000,8.0000) circle (1.2pt);
\fill[red!80!black] (0.0000,8.0000) circle (1.2pt);
\fill[red!80!black] (12.0000,8.0000) circle (1.2pt);
\fill[red!80!black] (0.0000,8.0000) circle (1.2pt);
\fill[red!80!black] (12.0000,8.0000) circle (1.2pt);
\fill[red!80!black] (0.0000,8.0000) circle (1.2pt);
\fill[red!80!black] (12.0000,8.0000) circle (1.2pt);
\fill[red!80!black] (0.0000,8.0000) circle (1.2pt);
\fill[red!80!black] (12.0000,8.0000) circle (1.2pt);
\fill[red!80!black] (0.0000,8.0000) circle (1.2pt);
\fill[red!80!black] (12.0000,8.0000) circle (1.2pt);
\fill[red!80!black] (0.0000,8.0000) circle (1.2pt);
\fill[red!80!black] (12.0000,8.0000) circle (1.2pt);
\fill[red!80!black] (0.0000,8.0000) circle (1.2pt);
\fill[red!80!black] (12.0000,8.0000) circle (1.2pt);
\fill[red!80!black] (0.0000,8.0000) circle (1.2pt);
\fill[red!80!black] (12.0000,8.0000) circle (1.2pt);
\fill[red!80!black] (0.0000,8.0000) circle (1.2pt);
\fill[red!80!black] (12.0000,8.0000) circle (1.2pt);
\fill[red!80!black] (0.0000,8.0000) circle (1.2pt);
\fill[red!80!black] (12.0000,8.0000) circle (1.2pt);
\fill[red!80!black] (0.0000,8.0000) circle (1.2pt);
\fill[red!80!black] (12.0000,8.0000) circle (1.2pt);
\fill[red!80!black] (0.0000,8.0000) circle (1.2pt);
\fill[red!80!black] (12.0000,8.0000) circle (1.2pt);
\fill[red!80!black] (0.0000,8.0000) circle (1.2pt);
\fill[red!80!black] (12.0000,8.0000) circle (1.2pt);
\fill[red!80!black] (0.0000,8.0000) circle (1.2pt);
\fill[red!80!black] (12.0000,8.0000) circle (1.2pt);
\fill[red!80!black] (0.0000,8.0000) circle (1.2pt);
\fill[red!80!black] (12.0000,8.0000) circle (1.2pt);
\fill[red!80!black] (0.0000,8.0000) circle (1.2pt);
\fill[red!80!black] (12.0000,8.0000) circle (1.2pt);
\fill[red!80!black] (0.0000,8.0000) circle (1.2pt);
\fill[red!80!black] (12.0000,8.0000) circle (1.2pt);
\fill[red!80!black] (0.0000,8.0000) circle (1.2pt);
\fill[red!80!black] (12.0000,8.0000) circle (1.2pt);
\fill[red!80!black] (0.0000,8.0000) circle (1.2pt);
\fill[red!80!black] (12.0000,8.0000) circle (1.2pt);
\fill[red!80!black] (0.0000,8.0000) circle (1.2pt);
\fill[red!80!black] (12.0000,8.0000) circle (1.2pt);
\fill[red!80!black] (0.0000,8.0000) circle (1.2pt);
\fill[red!80!black] (12.0000,8.0000) circle (1.2pt);
\fill[red!80!black] (0.0000,8.0000) circle (1.2pt);
\fill[red!80!black] (12.0000,8.0000) circle (1.2pt);
\fill[red!80!black] (0.0000,8.0000) circle (1.2pt);
\fill[red!80!black] (12.0000,8.0000) circle (1.2pt);
\fill[red!80!black] (0.0000,8.0000) circle (1.2pt);
\fill[red!80!black] (12.0000,8.0000) circle (1.2pt);
\fill[red!80!black] (0.0000,8.0000) circle (1.2pt);
\fill[red!80!black] (12.0000,8.0000) circle (1.2pt);
\fill[red!80!black] (0.0000,8.0000) circle (1.2pt);
\fill[red!80!black] (12.0000,8.0000) circle (1.2pt);
\fill[red!80!black] (0.0000,8.0000) circle (1.2pt);
\fill[red!80!black] (12.0000,8.0000) circle (1.2pt);
\fill[red!80!black] (0.0000,8.0000) circle (1.2pt);
\fill[red!80!black] (12.0000,8.0000) circle (1.2pt);
\fill[red!80!black] (0.0000,8.0000) circle (1.2pt);
\fill[red!80!black] (12.0000,8.0000) circle (1.2pt);
\fill[red!80!black] (0.0000,8.0000) circle (1.2pt);
\fill[red!80!black] (12.0000,8.0000) circle (1.2pt);
\fill[red!80!black] (0.0000,8.0000) circle (1.2pt);
\fill[red!80!black] (12.0000,8.0000) circle (1.2pt);
\fill[red!80!black] (0.0000,8.0000) circle (1.2pt);
\fill[red!80!black] (12.0000,8.0000) circle (1.2pt);
\fill[red!80!black] (0.0000,8.0000) circle (1.2pt);
\fill[red!80!black] (12.0000,8.0000) circle (1.2pt);
\fill[red!80!black] (0.0000,8.0000) circle (1.2pt);
\fill[red!80!black] (12.0000,8.0000) circle (1.2pt);
\fill[red!80!black] (0.0000,8.0000) circle (1.2pt);
\fill[red!80!black] (12.0000,8.0000) circle (1.2pt);
\fill[red!80!black] (0.0000,8.0000) circle (1.2pt);
\fill[red!80!black] (12.0000,8.0000) circle (1.2pt);
\fill[red!80!black] (0.0000,8.0000) circle (1.2pt);
\fill[red!80!black] (12.0000,8.0000) circle (1.2pt);
\fill[red!80!black] (0.0000,8.0000) circle (1.2pt);
\fill[red!80!black] (12.0000,8.0000) circle (1.2pt);
\fill[red!80!black] (0.0000,8.0000) circle (1.2pt);
\fill[red!80!black] (12.0000,8.0000) circle (1.2pt);
\fill[red!80!black] (0.0000,8.0000) circle (1.2pt);
\fill[red!80!black] (12.0000,8.0000) circle (1.2pt);
\fill[red!80!black] (0.0000,8.0000) circle (1.2pt);
\fill[red!80!black] (12.0000,8.0000) circle (1.2pt);
\fill[red!80!black] (0.0000,8.0000) circle (1.2pt);
\fill[red!80!black] (12.0000,8.0000) circle (1.2pt);
\fill[red!80!black] (0.0000,8.0000) circle (1.2pt);
\fill[red!80!black] (12.0000,8.0000) circle (1.2pt);
\fill[red!80!black] (0.0000,8.0000) circle (1.2pt);
\fill[red!80!black] (12.0000,8.0000) circle (1.2pt);
\fill[red!80!black] (0.0000,8.0000) circle (1.2pt);
\fill[red!80!black] (12.0000,8.0000) circle (1.2pt);
\fill[red!80!black] (0.0000,8.0000) circle (1.2pt);
\fill[red!80!black] (12.0000,8.0000) circle (1.2pt);
\fill[red!80!black] (0.0000,8.0000) circle (1.2pt);
\fill[red!80!black] (12.0000,8.0000) circle (1.2pt);
\fill[red!80!black] (0.0000,8.0000) circle (1.2pt);
\fill[red!80!black] (12.0000,8.0000) circle (1.2pt);
\fill[red!80!black] (0.0000,8.0000) circle (1.2pt);
\fill[red!80!black] (12.0000,8.0000) circle (1.2pt);
\fill[red!80!black] (0.0000,8.0000) circle (1.2pt);
\fill[red!80!black] (12.0000,8.0000) circle (1.2pt);
\fill[red!80!black] (0.0000,8.0000) circle (1.2pt);
\fill[red!80!black] (12.0000,8.0000) circle (1.2pt);
\fill[red!80!black] (0.0000,8.0000) circle (1.2pt);
\fill[red!80!black] (12.0000,8.0000) circle (1.2pt);
\fill[red!80!black] (0.0000,8.0000) circle (1.2pt);
\fill[red!80!black] (12.0000,8.0000) circle (1.2pt);
\fill[red!80!black] (0.0000,8.0000) circle (1.2pt);
\fill[red!80!black] (12.0000,8.0000) circle (1.2pt);
\fill[red!80!black] (0.0000,8.0000) circle (1.2pt);
\fill[red!80!black] (12.0000,8.0000) circle (1.2pt);
\fill[red!80!black] (0.0000,8.0000) circle (1.2pt);
\fill[red!80!black] (12.0000,8.0000) circle (1.2pt);
\fill[red!80!black] (0.0000,8.0000) circle (1.2pt);
\fill[red!80!black] (12.0000,8.0000) circle (1.2pt);
\fill[red!80!black] (0.0000,8.0000) circle (1.2pt);
\fill[red!80!black] (12.0000,8.0000) circle (1.2pt);
\fill[red!80!black] (0.0000,8.0000) circle (1.2pt);
\fill[red!80!black] (12.0000,8.0000) circle (1.2pt);
\fill[red!80!black] (0.0000,8.0000) circle (1.2pt);
\fill[red!80!black] (12.0000,8.0000) circle (1.2pt);
\fill[red!80!black] (0.0000,8.0000) circle (1.2pt);
\fill[red!80!black] (12.0000,8.0000) circle (1.2pt);
\fill[red!80!black] (0.0000,8.0000) circle (1.2pt);
\fill[red!80!black] (12.0000,8.0000) circle (1.2pt);
\fill[red!80!black] (0.0000,8.0000) circle (1.2pt);
\fill[red!80!black] (12.0000,8.0000) circle (1.2pt);
\fill[red!80!black] (0.0000,8.0000) circle (1.2pt);
\fill[red!80!black] (12.0000,8.0000) circle (1.2pt);
\fill[red!80!black] (0.0000,8.0000) circle (1.2pt);
\fill[red!80!black] (12.0000,8.0000) circle (1.2pt);
\fill[red!80!black] (0.0000,8.0000) circle (1.2pt);
\fill[red!80!black] (12.0000,8.0000) circle (1.2pt);
\fill[red!80!black] (0.0000,8.0000) circle (1.2pt);
\fill[red!80!black] (12.0000,8.0000) circle (1.2pt);
\fill[red!80!black] (0.0000,8.0000) circle (1.2pt);
\fill[red!80!black] (12.0000,8.0000) circle (1.2pt);
\fill[red!80!black] (0.0000,8.0000) circle (1.2pt);
\fill[red!80!black] (12.0000,8.0000) circle (1.2pt);
\fill[red!80!black] (0.0000,8.0000) circle (1.2pt);
\fill[red!80!black] (12.0000,8.0000) circle (1.2pt);
\fill[red!80!black] (0.0000,8.0000) circle (1.2pt);
\fill[red!80!black] (12.0000,8.0000) circle (1.2pt);
\fill[red!80!black] (0.0000,8.0000) circle (1.2pt);
\fill[red!80!black] (12.0000,8.0000) circle (1.2pt);
\fill[red!80!black] (0.0000,8.0000) circle (1.2pt);
\fill[red!80!black] (12.0000,8.0000) circle (1.2pt);
\fill[red!80!black] (0.0000,8.0000) circle (1.2pt);
\fill[red!80!black] (12.0000,8.0000) circle (1.2pt);
\fill[red!80!black] (0.0000,8.0000) circle (1.2pt);
\fill[red!80!black] (12.0000,8.0000) circle (1.2pt);
\fill[red!80!black] (0.0000,8.0000) circle (1.2pt);
\fill[red!80!black] (12.0000,8.0000) circle (1.2pt);
\fill[red!80!black] (0.0000,8.0000) circle (1.2pt);
\fill[red!80!black] (12.0000,8.0000) circle (1.2pt);
\fill[red!80!black] (0.0000,8.0000) circle (1.2pt);
\fill[red!80!black] (12.0000,8.0000) circle (1.2pt);
\fill[red!80!black] (0.0000,8.0000) circle (1.2pt);
\fill[red!80!black] (12.0000,8.0000) circle (1.2pt);
\fill[red!80!black] (0.0000,8.0000) circle (1.2pt);
\fill[red!80!black] (12.0000,8.0000) circle (1.2pt);
\fill[red!80!black] (0.0000,8.0000) circle (1.2pt);
\fill[red!80!black] (12.0000,8.0000) circle (1.2pt);
\fill[red!80!black] (0.0000,8.0000) circle (1.2pt);
\fill[red!80!black] (12.0000,8.0000) circle (1.2pt);
\fill[red!80!black] (0.0000,8.0000) circle (1.2pt);
\fill[red!80!black] (12.0000,8.0000) circle (1.2pt);
\fill[red!80!black] (0.0000,8.0000) circle (1.2pt);
\fill[red!80!black] (12.0000,8.0000) circle (1.2pt);
\fill[red!80!black] (0.0000,8.0000) circle (1.2pt);
\fill[red!80!black] (12.0000,8.0000) circle (1.2pt);
\fill[red!80!black] (0.0000,8.0000) circle (1.2pt);
\fill[red!80!black] (12.0000,8.0000) circle (1.2pt);
\fill[red!80!black] (0.0000,8.0000) circle (1.2pt);
\fill[red!80!black] (12.0000,8.0000) circle (1.2pt);
\fill[red!80!black] (0.0000,8.0000) circle (1.2pt);
\fill[red!80!black] (12.0000,8.0000) circle (1.2pt);
\fill[red!80!black] (0.0000,8.0000) circle (1.2pt);
\fill[red!80!black] (12.0000,8.0000) circle (1.2pt);
\fill[red!80!black] (0.0000,8.0000) circle (1.2pt);
\fill[red!80!black] (12.0000,8.0000) circle (1.2pt);
\fill[red!80!black] (0.0000,8.0000) circle (1.2pt);
\fill[red!80!black] (12.0000,8.0000) circle (1.2pt);
\fill[red!80!black] (0.0000,8.0000) circle (1.2pt);
\fill[red!80!black] (12.0000,8.0000) circle (1.2pt);
\fill[red!80!black] (0.0000,8.0000) circle (1.2pt);
\fill[red!80!black] (12.0000,8.0000) circle (1.2pt);
\fill[red!80!black] (0.0000,8.0000) circle (1.2pt);
\fill[red!80!black] (12.0000,8.0000) circle (1.2pt);
\fill[red!80!black] (0.0000,8.0000) circle (1.2pt);
\fill[red!80!black] (12.0000,8.0000) circle (1.2pt);
\fill[red!80!black] (0.0000,8.0000) circle (1.2pt);
\fill[red!80!black] (12.0000,8.0000) circle (1.2pt);
\fill[red!80!black] (0.0000,8.0000) circle (1.2pt);
\fill[red!80!black] (12.0000,8.0000) circle (1.2pt);
\fill[red!80!black] (0.0000,8.0000) circle (1.2pt);
\fill[red!80!black] (12.0000,8.0000) circle (1.2pt);
\fill[red!80!black] (0.0000,8.0000) circle (1.2pt);
\fill[red!80!black] (12.0000,8.0000) circle (1.2pt);
\fill[red!80!black] (0.0000,8.0000) circle (1.2pt);
\fill[red!80!black] (12.0000,8.0000) circle (1.2pt);
\fill[red!80!black] (0.0000,8.0000) circle (1.2pt);
\fill[red!80!black] (12.0000,8.0000) circle (1.2pt);
\fill[red!80!black] (0.0000,8.0000) circle (1.2pt);
\fill[red!80!black] (12.0000,8.0000) circle (1.2pt);
\fill[red!80!black] (0.0000,8.0000) circle (1.2pt);
\fill[red!80!black] (12.0000,8.0000) circle (1.2pt);
\fill[red!80!black] (0.0000,8.0000) circle (1.2pt);
\fill[red!80!black] (12.0000,8.0000) circle (1.2pt);
\fill[red!80!black] (0.0000,8.0000) circle (1.2pt);
\fill[red!80!black] (12.0000,8.0000) circle (1.2pt);
\fill[red!80!black] (0.0000,8.0000) circle (1.2pt);
\fill[red!80!black] (12.0000,8.0000) circle (1.2pt);
\fill[red!80!black] (0.0000,8.0000) circle (1.2pt);
\fill[red!80!black] (12.0000,8.0000) circle (1.2pt);
\fill[red!80!black] (0.0000,8.0000) circle (1.2pt);
\fill[red!80!black] (12.0000,8.0000) circle (1.2pt);
\fill[red!80!black] (0.0000,8.0000) circle (1.2pt);
\fill[red!80!black] (12.0000,8.0000) circle (1.2pt);
\fill[red!80!black] (0.0000,8.0000) circle (1.2pt);
\fill[red!80!black] (12.0000,8.0000) circle (1.2pt);
\fill[red!80!black] (0.0000,8.0000) circle (1.2pt);
\fill[red!80!black] (12.0000,8.0000) circle (1.2pt);
\fill[red!80!black] (0.0000,8.0000) circle (1.2pt);
\fill[red!80!black] (12.0000,8.0000) circle (1.2pt);
\fill[red!80!black] (0.0000,8.0000) circle (1.2pt);
\fill[red!80!black] (12.0000,8.0000) circle (1.2pt);
\fill[red!80!black] (0.0000,8.0000) circle (1.2pt);
\fill[red!80!black] (12.0000,8.0000) circle (1.2pt);
\fill[red!80!black] (0.0000,8.0000) circle (1.2pt);
\fill[red!80!black] (12.0000,8.0000) circle (1.2pt);
\fill[red!80!black] (0.0000,8.0000) circle (1.2pt);
\fill[red!80!black] (12.0000,8.0000) circle (1.2pt);
\fill[red!80!black] (0.0000,8.0000) circle (1.2pt);
\fill[red!80!black] (12.0000,8.0000) circle (1.2pt);
\fill[red!80!black] (0.0000,8.0000) circle (1.2pt);
\fill[red!80!black] (12.0000,8.0000) circle (1.2pt);
\fill[red!80!black] (0.0000,8.0000) circle (1.2pt);
\fill[red!80!black] (12.0000,8.0000) circle (1.2pt);
\fill[red!80!black] (0.0000,8.0000) circle (1.2pt);
\fill[red!80!black] (12.0000,8.0000) circle (1.2pt);
\fill[red!80!black] (0.0000,8.0000) circle (1.2pt);
\fill[red!80!black] (12.0000,8.0000) circle (1.2pt);
\fill[red!80!black] (0.0000,8.0000) circle (1.2pt);
\fill[red!80!black] (12.0000,8.0000) circle (1.2pt);
\fill[red!80!black] (0.0000,8.0000) circle (1.2pt);
\fill[red!80!black] (12.0000,8.0000) circle (1.2pt);
\fill[red!80!black] (0.0000,8.0000) circle (1.2pt);
\fill[red!80!black] (12.0000,8.0000) circle (1.2pt);
\fill[red!80!black] (0.0000,8.0000) circle (1.2pt);
\fill[red!80!black] (12.0000,8.0000) circle (1.2pt);
\fill[red!80!black] (0.0000,8.0000) circle (1.2pt);
\fill[red!80!black] (12.0000,8.0000) circle (1.2pt);
\fill[red!80!black] (0.0000,8.0000) circle (1.2pt);
\fill[red!80!black] (12.0000,8.0000) circle (1.2pt);
\fill[red!80!black] (0.0000,8.0000) circle (1.2pt);
\fill[red!80!black] (12.0000,8.0000) circle (1.2pt);
\fill[red!80!black] (0.0000,8.0000) circle (1.2pt);
\fill[red!80!black] (12.0000,8.0000) circle (1.2pt);
\fill[red!80!black] (0.0000,8.0000) circle (1.2pt);
\fill[red!80!black] (12.0000,8.0000) circle (1.2pt);
\fill[red!80!black] (0.0000,8.0000) circle (1.2pt);
\fill[red!80!black] (12.0000,8.0000) circle (1.2pt);
\fill[red!80!black] (0.0000,8.0000) circle (1.2pt);
\fill[red!80!black] (12.0000,8.0000) circle (1.2pt);
\fill[red!80!black] (0.0000,8.0000) circle (1.2pt);
\fill[red!80!black] (12.0000,8.0000) circle (1.2pt);
\fill[red!80!black] (0.0000,8.0000) circle (1.2pt);
\fill[red!80!black] (12.0000,8.0000) circle (1.2pt);
\fill[red!80!black] (0.0000,8.0000) circle (1.2pt);
\fill[red!80!black] (12.0000,8.0000) circle (1.2pt);
\fill[red!80!black] (0.0000,8.0000) circle (1.2pt);
\fill[red!80!black] (12.0000,8.0000) circle (1.2pt);
\fill[red!80!black] (0.0000,8.0000) circle (1.2pt);
\fill[red!80!black] (12.0000,8.0000) circle (1.2pt);
\fill[red!80!black] (0.0000,8.0000) circle (1.2pt);
\fill[red!80!black] (12.0000,8.0000) circle (1.2pt);
\fill[red!80!black] (0.0000,8.0000) circle (1.2pt);
\fill[red!80!black] (12.0000,8.0000) circle (1.2pt);
\fill[red!80!black] (0.0000,8.0000) circle (1.2pt);
\fill[red!80!black] (12.0000,8.0000) circle (1.2pt);
\fill[red!80!black] (0.0000,8.0000) circle (1.2pt);
\fill[red!80!black] (12.0000,8.0000) circle (1.2pt);
\fill[red!80!black] (0.0000,8.0000) circle (1.2pt);
\fill[red!80!black] (12.0000,8.0000) circle (1.2pt);
\fill[red!80!black] (0.0000,8.0000) circle (1.2pt);
\fill[red!80!black] (12.0000,8.0000) circle (1.2pt);
\fill[red!80!black] (0.0000,8.0000) circle (1.2pt);
\fill[red!80!black] (12.0000,8.0000) circle (1.2pt);
\fill[red!80!black] (0.0000,8.0000) circle (1.2pt);
\fill[red!80!black] (12.0000,8.0000) circle (1.2pt);
\fill[red!80!black] (0.0000,8.0000) circle (1.2pt);
\fill[red!80!black] (12.0000,8.0000) circle (1.2pt);
\fill[red!80!black] (0.0000,8.0000) circle (1.2pt);
\fill[red!80!black] (12.0000,8.0000) circle (1.2pt);
\fill[red!80!black] (0.0000,8.0000) circle (1.2pt);
\fill[red!80!black] (12.0000,8.0000) circle (1.2pt);
\fill[red!80!black] (0.0000,8.0000) circle (1.2pt);
\fill[red!80!black] (12.0000,8.0000) circle (1.2pt);
\fill[red!80!black] (0.0000,8.0000) circle (1.2pt);
\fill[red!80!black] (12.0000,8.0000) circle (1.2pt);
\fill[red!80!black] (0.0000,8.0000) circle (1.2pt);
\fill[red!80!black] (12.0000,8.0000) circle (1.2pt);
\fill[red!80!black] (0.0000,8.0000) circle (1.2pt);
\fill[red!80!black] (12.0000,8.0000) circle (1.2pt);
\fill[red!80!black] (0.0000,8.0000) circle (1.2pt);
\fill[red!80!black] (12.0000,8.0000) circle (1.2pt);
\fill[red!80!black] (0.0000,8.0000) circle (1.2pt);
\fill[red!80!black] (12.0000,8.0000) circle (1.2pt);
\fill[red!80!black] (0.0000,8.0000) circle (1.2pt);
\fill[red!80!black] (12.0000,8.0000) circle (1.2pt);
\fill[red!80!black] (0.0000,8.0000) circle (1.2pt);
\fill[red!80!black] (12.0000,8.0000) circle (1.2pt);
\fill[red!80!black] (0.0000,8.0000) circle (1.2pt);
\fill[red!80!black] (12.0000,8.0000) circle (1.2pt);
\fill[red!80!black] (0.0000,8.0000) circle (1.2pt);
\fill[red!80!black] (12.0000,8.0000) circle (1.2pt);
\fill[red!80!black] (0.0000,8.0000) circle (1.2pt);
\fill[red!80!black] (12.0000,8.0000) circle (1.2pt);
\fill[red!80!black] (0.0000,8.0000) circle (1.2pt);
\fill[red!80!black] (12.0000,8.0000) circle (1.2pt);
\fill[red!80!black] (0.0000,8.0000) circle (1.2pt);
\fill[red!80!black] (12.0000,8.0000) circle (1.2pt);
\fill[red!80!black] (0.0000,8.0000) circle (1.2pt);
\fill[red!80!black] (12.0000,8.0000) circle (1.2pt);
\fill[red!80!black] (0.0000,8.0000) circle (1.2pt);
\fill[red!80!black] (12.0000,8.0000) circle (1.2pt);
\fill[red!80!black] (0.0000,8.0000) circle (1.2pt);
\fill[red!80!black] (12.0000,8.0000) circle (1.2pt);
\fill[red!80!black] (0.0000,8.0000) circle (1.2pt);
\fill[red!80!black] (12.0000,8.0000) circle (1.2pt);
\fill[red!80!black] (0.0000,8.0000) circle (1.2pt);
\fill[red!80!black] (12.0000,8.0000) circle (1.2pt);
\fill[red!80!black] (0.0000,8.0000) circle (1.2pt);
\fill[red!80!black] (12.0000,8.0000) circle (1.2pt);
\fill[red!80!black] (0.0000,8.0000) circle (1.2pt);
\fill[red!80!black] (12.0000,8.0000) circle (1.2pt);
\fill[red!80!black] (0.0000,8.0000) circle (1.2pt);
\fill[red!80!black] (12.0000,8.0000) circle (1.2pt);
\fill[red!80!black] (0.0000,8.0000) circle (1.2pt);
\fill[red!80!black] (12.0000,8.0000) circle (1.2pt);
\fill[red!80!black] (0.0000,8.0000) circle (1.2pt);
\fill[red!80!black] (12.0000,8.0000) circle (1.2pt);
\fill[red!80!black] (0.0000,8.0000) circle (1.2pt);
\fill[red!80!black] (12.0000,8.0000) circle (1.2pt);
\fill[red!80!black] (0.0000,8.0000) circle (1.2pt);
\fill[red!80!black] (12.0000,8.0000) circle (1.2pt);
\fill[red!80!black] (0.0000,8.0000) circle (1.2pt);
\fill[red!80!black] (12.0000,8.0000) circle (1.2pt);
\fill[red!80!black] (0.0000,8.0000) circle (1.2pt);
\fill[red!80!black] (12.0000,8.0000) circle (1.2pt);
\fill[red!80!black] (0.0000,8.0000) circle (1.2pt);
\fill[red!80!black] (12.0000,8.0000) circle (1.2pt);
\fill[red!80!black] (0.0000,8.0000) circle (1.2pt);
\fill[red!80!black] (12.0000,8.0000) circle (1.2pt);
\fill[red!80!black] (0.0000,8.0000) circle (1.2pt);
\fill[red!80!black] (12.0000,8.0000) circle (1.2pt);
\fill[red!80!black] (0.0000,8.0000) circle (1.2pt);
\fill[red!80!black] (12.0000,8.0000) circle (1.2pt);
\fill[red!80!black] (0.0000,8.0000) circle (1.2pt);
\fill[red!80!black] (12.0000,8.0000) circle (1.2pt);
\fill[red!80!black] (0.0000,8.0000) circle (1.2pt);
\fill[red!80!black] (12.0000,8.0000) circle (1.2pt);
\fill[red!80!black] (0.0000,8.0000) circle (1.2pt);
\fill[red!80!black] (12.0000,8.0000) circle (1.2pt);
\fill[red!80!black] (0.0000,8.0000) circle (1.2pt);
\fill[red!80!black] (12.0000,8.0000) circle (1.2pt);
\fill[red!80!black] (0.0000,8.0000) circle (1.2pt);
\fill[red!80!black] (12.0000,8.0000) circle (1.2pt);
\fill[red!80!black] (0.0000,8.0000) circle (1.2pt);
\fill[red!80!black] (12.0000,8.0000) circle (1.2pt);
\fill[red!80!black] (0.0000,8.0000) circle (1.2pt);
\fill[red!80!black] (12.0000,8.0000) circle (1.2pt);
\fill[red!80!black] (0.0000,8.0000) circle (1.2pt);
\fill[red!80!black] (12.0000,8.0000) circle (1.2pt);
\fill[red!80!black] (0.0000,8.0000) circle (1.2pt);
\fill[red!80!black] (12.0000,8.0000) circle (1.2pt);
\fill[red!80!black] (0.0000,8.0000) circle (1.2pt);
\fill[red!80!black] (12.0000,8.0000) circle (1.2pt);
\fill[red!80!black] (0.0000,8.0000) circle (1.2pt);
\fill[red!80!black] (12.0000,8.0000) circle (1.2pt);
\fill[red!80!black] (0.0000,8.0000) circle (1.2pt);
\fill[red!80!black] (12.0000,8.0000) circle (1.2pt);
\fill[red!80!black] (0.0000,8.0000) circle (1.2pt);
\fill[red!80!black] (12.0000,8.0000) circle (1.2pt);
\fill[red!80!black] (0.0000,8.0000) circle (1.2pt);
\fill[red!80!black] (12.0000,8.0000) circle (1.2pt);
\fill[red!80!black] (0.0000,8.0000) circle (1.2pt);
\fill[red!80!black] (12.0000,8.0000) circle (1.2pt);
\fill[red!80!black] (0.0000,8.0000) circle (1.2pt);
\fill[red!80!black] (12.0000,8.0000) circle (1.2pt);
\fill[red!80!black] (0.0000,8.0000) circle (1.2pt);
\fill[red!80!black] (12.0000,8.0000) circle (1.2pt);
\fill[red!80!black] (0.0000,8.0000) circle (1.2pt);
\fill[red!80!black] (12.0000,8.0000) circle (1.2pt);
\fill[red!80!black] (0.0000,8.0000) circle (1.2pt);
\fill[red!80!black] (12.0000,8.0000) circle (1.2pt);
\fill[red!80!black] (0.0000,8.0000) circle (1.2pt);
\fill[red!80!black] (12.0000,8.0000) circle (1.2pt);
\fill[red!80!black] (0.0000,8.0000) circle (1.2pt);
\fill[red!80!black] (12.0000,8.0000) circle (1.2pt);
\fill[red!80!black] (0.0000,8.0000) circle (1.2pt);
\fill[red!80!black] (12.0000,8.0000) circle (1.2pt);
\fill[red!80!black] (0.0000,8.0000) circle (1.2pt);
\fill[red!80!black] (12.0000,8.0000) circle (1.2pt);
\fill[red!80!black] (0.0000,8.0000) circle (1.2pt);
\fill[red!80!black] (12.0000,8.0000) circle (1.2pt);
\fill[red!80!black] (0.0000,8.0000) circle (1.2pt);
\fill[red!80!black] (12.0000,8.0000) circle (1.2pt);
\fill[red!80!black] (0.0000,8.0000) circle (1.2pt);
\fill[red!80!black] (12.0000,8.0000) circle (1.2pt);
\fill[red!80!black] (0.0000,8.0000) circle (1.2pt);
\fill[red!80!black] (12.0000,8.0000) circle (1.2pt);
\fill[red!80!black] (0.0000,8.0000) circle (1.2pt);
\fill[red!80!black] (12.0000,8.0000) circle (1.2pt);
\fill[red!80!black] (0.0000,8.0000) circle (1.2pt);
\fill[red!80!black] (12.0000,8.0000) circle (1.2pt);
\fill[red!80!black] (0.0000,8.0000) circle (1.2pt);
\fill[red!80!black] (12.0000,8.0000) circle (1.2pt);
\fill[red!80!black] (0.0000,8.0000) circle (1.2pt);
\fill[red!80!black] (12.0000,8.0000) circle (1.2pt);
\fill[red!80!black] (0.0000,8.0000) circle (1.2pt);
\fill[red!80!black] (12.0000,8.0000) circle (1.2pt);
\fill[red!80!black] (0.0000,8.0000) circle (1.2pt);
\fill[red!80!black] (12.0000,8.0000) circle (1.2pt);
\fill[red!80!black] (0.0000,8.0000) circle (1.2pt);
\fill[red!80!black] (12.0000,8.0000) circle (1.2pt);
\fill[red!80!black] (0.0000,8.0000) circle (1.2pt);
\fill[red!80!black] (12.0000,8.0000) circle (1.2pt);
\fill[red!80!black] (0.0000,8.0000) circle (1.2pt);
\fill[red!80!black] (12.0000,8.0000) circle (1.2pt);
\fill[red!80!black] (0.0000,8.0000) circle (1.2pt);
\fill[red!80!black] (12.0000,8.0000) circle (1.2pt);
\fill[red!80!black] (0.0000,8.0000) circle (1.2pt);
\fill[red!80!black] (12.0000,8.0000) circle (1.2pt);
\fill[red!80!black] (0.0000,8.0000) circle (1.2pt);
\fill[red!80!black] (12.0000,8.0000) circle (1.2pt);
\fill[red!80!black] (0.0000,8.0000) circle (1.2pt);
\fill[red!80!black] (12.0000,8.0000) circle (1.2pt);
\fill[red!80!black] (0.0000,8.0000) circle (1.2pt);
\fill[red!80!black] (12.0000,8.0000) circle (1.2pt);
\fill[red!80!black] (0.0000,8.0000) circle (1.2pt);
\fill[red!80!black] (12.0000,8.0000) circle (1.2pt);
\fill[red!80!black] (0.0000,8.0000) circle (1.2pt);
\fill[red!80!black] (12.0000,8.0000) circle (1.2pt);
\fill[red!80!black] (0.0000,8.0000) circle (1.2pt);
\fill[red!80!black] (12.0000,8.0000) circle (1.2pt);
\fill[red!80!black] (0.0000,8.0000) circle (1.2pt);
\fill[red!80!black] (12.0000,8.0000) circle (1.2pt);
\fill[red!80!black] (0.0000,8.0000) circle (1.2pt);
\fill[red!80!black] (12.0000,8.0000) circle (1.2pt);
\fill[red!80!black] (0.0000,8.0000) circle (1.2pt);
\fill[red!80!black] (12.0000,8.0000) circle (1.2pt);
\fill[red!80!black] (0.0000,8.0000) circle (1.2pt);
\fill[red!80!black] (12.0000,8.0000) circle (1.2pt);
\fill[red!80!black] (0.0000,8.0000) circle (1.2pt);
\fill[red!80!black] (12.0000,8.0000) circle (1.2pt);
\fill[red!80!black] (0.0000,8.0000) circle (1.2pt);
\fill[red!80!black] (12.0000,8.0000) circle (1.2pt);
\fill[red!80!black] (0.0000,8.0000) circle (1.2pt);
\fill[red!80!black] (12.0000,8.0000) circle (1.2pt);
\fill[red!80!black] (0.0000,8.0000) circle (1.2pt);
\fill[red!80!black] (12.0000,8.0000) circle (1.2pt);
\fill[red!80!black] (0.0000,8.0000) circle (1.2pt);
\fill[red!80!black] (12.0000,8.0000) circle (1.2pt);
\fill[red!80!black] (0.0000,8.0000) circle (1.2pt);
\fill[red!80!black] (12.0000,8.0000) circle (1.2pt);
\fill[red!80!black] (0.0000,8.0000) circle (1.2pt);
\fill[red!80!black] (12.0000,8.0000) circle (1.2pt);
\fill[red!80!black] (0.0000,8.0000) circle (1.2pt);
\fill[red!80!black] (12.0000,8.0000) circle (1.2pt);
\fill[red!80!black] (0.0000,8.0000) circle (1.2pt);
\fill[red!80!black] (12.0000,8.0000) circle (1.2pt);
\fill[red!80!black] (0.0000,8.0000) circle (1.2pt);
\fill[red!80!black] (12.0000,8.0000) circle (1.2pt);
\fill[red!80!black] (0.0000,8.0000) circle (1.2pt);
\fill[red!80!black] (12.0000,8.0000) circle (1.2pt);
\fill[red!80!black] (0.0000,8.0000) circle (1.2pt);
\fill[red!80!black] (12.0000,8.0000) circle (1.2pt);
\fill[red!80!black] (0.0000,8.0000) circle (1.2pt);
\fill[red!80!black] (12.0000,8.0000) circle (1.2pt);
\fill[red!80!black] (0.0000,8.0000) circle (1.2pt);
\fill[red!80!black] (12.0000,8.0000) circle (1.2pt);
\fill[red!80!black] (0.0000,8.0000) circle (1.2pt);
\fill[red!80!black] (12.0000,8.0000) circle (1.2pt);
\fill[red!80!black] (0.0000,8.0000) circle (1.2pt);
\fill[red!80!black] (12.0000,8.0000) circle (1.2pt);
\fill[red!80!black] (0.0000,8.0000) circle (1.2pt);
\fill[red!80!black] (12.0000,8.0000) circle (1.2pt);
\fill[red!80!black] (0.0000,8.0000) circle (1.2pt);
\fill[red!80!black] (12.0000,8.0000) circle (1.2pt);
\fill[red!80!black] (0.0000,8.0000) circle (1.2pt);
\fill[red!80!black] (12.0000,8.0000) circle (1.2pt);
\fill[red!80!black] (0.0000,8.0000) circle (1.2pt);
\fill[red!80!black] (12.0000,8.0000) circle (1.2pt);
\fill[red!80!black] (0.0000,8.0000) circle (1.2pt);
\fill[red!80!black] (12.0000,8.0000) circle (1.2pt);
\fill[red!80!black] (0.0000,8.0000) circle (1.2pt);
\fill[red!80!black] (12.0000,8.0000) circle (1.2pt);
\fill[red!80!black] (0.0000,8.0000) circle (1.2pt);
\fill[red!80!black] (12.0000,8.0000) circle (1.2pt);
\fill[red!80!black] (0.0000,8.0000) circle (1.2pt);
\fill[red!80!black] (12.0000,8.0000) circle (1.2pt);
\fill[red!80!black] (0.0000,8.0000) circle (1.2pt);
\fill[red!80!black] (12.0000,8.0000) circle (1.2pt);
\fill[red!80!black] (0.0000,8.0000) circle (1.2pt);
\fill[red!80!black] (12.0000,8.0000) circle (1.2pt);
\fill[red!80!black] (0.0000,8.0000) circle (1.2pt);
\fill[red!80!black] (12.0000,8.0000) circle (1.2pt);
\fill[red!80!black] (0.0000,8.0000) circle (1.2pt);
\fill[red!80!black] (12.0000,8.0000) circle (1.2pt);
\fill[red!80!black] (0.0000,8.0000) circle (1.2pt);
\fill[red!80!black] (12.0000,8.0000) circle (1.2pt);
\fill[red!80!black] (0.0000,8.0000) circle (1.2pt);
\fill[red!80!black] (12.0000,8.0000) circle (1.2pt);
\fill[red!80!black] (0.0000,8.0000) circle (1.2pt);
\fill[red!80!black] (12.0000,8.0000) circle (1.2pt);
\fill[red!80!black] (0.0000,8.0000) circle (1.2pt);
\fill[red!80!black] (12.0000,8.0000) circle (1.2pt);
\fill[red!80!black] (0.0000,8.0000) circle (1.2pt);
\fill[red!80!black] (12.0000,8.0000) circle (1.2pt);
\fill[red!80!black] (0.0000,8.0000) circle (1.2pt);
\fill[red!80!black] (12.0000,8.0000) circle (1.2pt);
\fill[red!80!black] (0.0000,8.0000) circle (1.2pt);
\fill[red!80!black] (12.0000,8.0000) circle (1.2pt);
\fill[red!80!black] (0.0000,8.0000) circle (1.2pt);
\fill[red!80!black] (12.0000,8.0000) circle (1.2pt);
\fill[red!80!black] (0.0000,8.0000) circle (1.2pt);
\fill[red!80!black] (12.0000,8.0000) circle (1.2pt);
\fill[red!80!black] (0.0000,8.0000) circle (1.2pt);
\fill[red!80!black] (12.0000,8.0000) circle (1.2pt);
\fill[red!80!black] (0.0000,8.0000) circle (1.2pt);
\fill[red!80!black] (12.0000,8.0000) circle (1.2pt);
\fill[red!80!black] (0.0000,8.0000) circle (1.2pt);
\fill[red!80!black] (12.0000,8.0000) circle (1.2pt);
\fill[red!80!black] (0.0000,8.0000) circle (1.2pt);
\fill[red!80!black] (12.0000,8.0000) circle (1.2pt);
\fill[red!80!black] (0.0000,8.0000) circle (1.2pt);
\fill[red!80!black] (12.0000,8.0000) circle (1.2pt);
\fill[red!80!black] (0.0000,8.0000) circle (1.2pt);
\fill[red!80!black] (12.0000,8.0000) circle (1.2pt);
\fill[red!80!black] (0.0000,8.0000) circle (1.2pt);
\fill[red!80!black] (12.0000,8.0000) circle (1.2pt);
\fill[red!80!black] (0.0000,8.0000) circle (1.2pt);
\fill[red!80!black] (12.0000,8.0000) circle (1.2pt);
\fill[red!80!black] (0.0000,8.0000) circle (1.2pt);
\fill[red!80!black] (12.0000,8.0000) circle (1.2pt);
\fill[red!80!black] (0.0000,8.0000) circle (1.2pt);
\fill[red!80!black] (12.0000,8.0000) circle (1.2pt);
\fill[red!80!black] (0.0000,8.0000) circle (1.2pt);
\fill[red!80!black] (12.0000,8.0000) circle (1.2pt);
\fill[red!80!black] (0.0000,8.0000) circle (1.2pt);
\fill[red!80!black] (12.0000,8.0000) circle (1.2pt);
\fill[red!80!black] (0.0000,8.0000) circle (1.2pt);
\fill[red!80!black] (12.0000,8.0000) circle (1.2pt);
\fill[red!80!black] (0.0000,8.0000) circle (1.2pt);
\fill[red!80!black] (12.0000,8.0000) circle (1.2pt);
\fill[red!80!black] (0.0000,8.0000) circle (1.2pt);
\fill[red!80!black] (12.0000,8.0000) circle (1.2pt);
\fill[red!80!black] (0.0000,8.0000) circle (1.2pt);
\fill[red!80!black] (12.0000,8.0000) circle (1.2pt);
\fill[red!80!black] (0.0000,8.0000) circle (1.2pt);
\fill[red!80!black] (12.0000,8.0000) circle (1.2pt);
\fill[red!80!black] (0.0000,8.0000) circle (1.2pt);
\fill[red!80!black] (12.0000,8.0000) circle (1.2pt);
\fill[red!80!black] (0.0000,8.0000) circle (1.2pt);
\fill[red!80!black] (12.0000,8.0000) circle (1.2pt);
\fill[red!80!black] (0.0000,8.0000) circle (1.2pt);
\fill[red!80!black] (12.0000,8.0000) circle (1.2pt);
\fill[red!80!black] (0.0000,8.0000) circle (1.2pt);
\fill[red!80!black] (12.0000,8.0000) circle (1.2pt);
\fill[red!80!black] (0.0000,8.0000) circle (1.2pt);
\fill[red!80!black] (12.0000,8.0000) circle (1.2pt);
\fill[red!80!black] (0.0000,8.0000) circle (1.2pt);
\fill[red!80!black] (12.0000,8.0000) circle (1.2pt);
\fill[red!80!black] (0.0000,8.0000) circle (1.2pt);
\fill[red!80!black] (12.0000,8.0000) circle (1.2pt);
\fill[red!80!black] (0.0000,8.0000) circle (1.2pt);
\fill[red!80!black] (12.0000,8.0000) circle (1.2pt);
\fill[red!80!black] (0.0000,8.0000) circle (1.2pt);
\fill[red!80!black] (12.0000,8.0000) circle (1.2pt);
\fill[red!80!black] (0.0000,8.0000) circle (1.2pt);
\fill[red!80!black] (12.0000,8.0000) circle (1.2pt);
\fill[red!80!black] (0.0000,8.0000) circle (1.2pt);
\fill[red!80!black] (12.0000,8.0000) circle (1.2pt);
\fill[red!80!black] (0.0000,8.0000) circle (1.2pt);
\fill[red!80!black] (12.0000,8.0000) circle (1.2pt);
\fill[red!80!black] (0.0000,8.0000) circle (1.2pt);
\fill[red!80!black] (12.0000,8.0000) circle (1.2pt);
\fill[red!80!black] (0.0000,8.0000) circle (1.2pt);
\fill[red!80!black] (12.0000,8.0000) circle (1.2pt);
\fill[red!80!black] (0.0000,8.0000) circle (1.2pt);
\fill[red!80!black] (12.0000,8.0000) circle (1.2pt);
\fill[red!80!black] (0.0000,8.0000) circle (1.2pt);
\fill[red!80!black] (12.0000,8.0000) circle (1.2pt);
\fill[red!80!black] (0.0000,8.0000) circle (1.2pt);
\fill[red!80!black] (12.0000,8.0000) circle (1.2pt);
\fill[red!80!black] (0.0000,8.0000) circle (1.2pt);
\fill[red!80!black] (12.0000,8.0000) circle (1.2pt);
\fill[red!80!black] (0.0000,8.0000) circle (1.2pt);
\fill[red!80!black] (12.0000,8.0000) circle (1.2pt);
\fill[red!80!black] (0.0000,8.0000) circle (1.2pt);
\fill[red!80!black] (12.0000,8.0000) circle (1.2pt);
\fill[red!80!black] (0.0000,8.0000) circle (1.2pt);
\fill[red!80!black] (12.0000,8.0000) circle (1.2pt);
\fill[red!80!black] (0.0000,8.0000) circle (1.2pt);
\fill[red!80!black] (12.0000,8.0000) circle (1.2pt);
\fill[red!80!black] (0.0000,8.0000) circle (1.2pt);
\fill[red!80!black] (12.0000,8.0000) circle (1.2pt);
\fill[red!80!black] (0.0000,8.0000) circle (1.2pt);
\fill[red!80!black] (12.0000,8.0000) circle (1.2pt);
\fill[red!80!black] (0.0000,8.0000) circle (1.2pt);
\fill[red!80!black] (12.0000,8.0000) circle (1.2pt);
\fill[red!80!black] (0.0000,8.0000) circle (1.2pt);
\fill[red!80!black] (12.0000,8.0000) circle (1.2pt);
\fill[red!80!black] (0.0000,8.0000) circle (1.2pt);
\fill[red!80!black] (12.0000,8.0000) circle (1.2pt);
\fill[red!80!black] (0.0000,8.0000) circle (1.2pt);
\fill[red!80!black] (12.0000,8.0000) circle (1.2pt);
\fill[red!80!black] (0.0000,8.0000) circle (1.2pt);
\fill[red!80!black] (12.0000,8.0000) circle (1.2pt);
\fill[red!80!black] (0.0000,8.0000) circle (1.2pt);
\fill[red!80!black] (12.0000,8.0000) circle (1.2pt);
\fill[red!80!black] (0.0000,8.0000) circle (1.2pt);
\fill[red!80!black] (12.0000,8.0000) circle (1.2pt);
\fill[red!80!black] (0.0000,8.0000) circle (1.2pt);
\fill[red!80!black] (12.0000,8.0000) circle (1.2pt);
\fill[red!80!black] (0.0000,8.0000) circle (1.2pt);
\fill[red!80!black] (12.0000,8.0000) circle (1.2pt);
\fill[red!80!black] (0.0000,8.0000) circle (1.2pt);
\fill[red!80!black] (12.0000,8.0000) circle (1.2pt);
\fill[red!80!black] (0.0000,8.0000) circle (1.2pt);
\fill[red!80!black] (12.0000,8.0000) circle (1.2pt);
\fill[red!80!black] (0.0000,8.0000) circle (1.2pt);
\fill[red!80!black] (12.0000,8.0000) circle (1.2pt);
\fill[red!80!black] (0.0000,8.0000) circle (1.2pt);
\fill[red!80!black] (12.0000,8.0000) circle (1.2pt);
\fill[red!80!black] (0.0000,8.0000) circle (1.2pt);
\fill[red!80!black] (12.0000,8.0000) circle (1.2pt);
\fill[red!80!black] (0.0000,8.0000) circle (1.2pt);
\fill[red!80!black] (12.0000,8.0000) circle (1.2pt);
\fill[red!80!black] (0.0000,8.0000) circle (1.2pt);
\fill[red!80!black] (12.0000,8.0000) circle (1.2pt);
\fill[red!80!black] (0.0000,8.0000) circle (1.2pt);
\fill[red!80!black] (12.0000,8.0000) circle (1.2pt);
\fill[red!80!black] (0.0000,8.0000) circle (1.2pt);
\fill[red!80!black] (12.0000,8.0000) circle (1.2pt);
\fill[red!80!black] (0.0000,8.0000) circle (1.2pt);
\fill[red!80!black] (12.0000,8.0000) circle (1.2pt);
\fill[red!80!black] (0.0000,8.0000) circle (1.2pt);
\fill[red!80!black] (12.0000,8.0000) circle (1.2pt);
\fill[red!80!black] (0.0000,8.0000) circle (1.2pt);
\fill[red!80!black] (12.0000,8.0000) circle (1.2pt);
\fill[red!80!black] (0.0000,8.0000) circle (1.2pt);
\fill[red!80!black] (12.0000,8.0000) circle (1.2pt);
\fill[red!80!black] (0.0000,8.0000) circle (1.2pt);
\fill[red!80!black] (12.0000,8.0000) circle (1.2pt);
\fill[red!80!black] (0.0000,8.0000) circle (1.2pt);
\fill[red!80!black] (12.0000,8.0000) circle (1.2pt);
\fill[red!80!black] (0.0000,8.0000) circle (1.2pt);
\fill[red!80!black] (12.0000,8.0000) circle (1.2pt);
\fill[red!80!black] (0.0000,8.0000) circle (1.2pt);
\fill[red!80!black] (12.0000,8.0000) circle (1.2pt);
\fill[red!80!black] (0.0000,8.0000) circle (1.2pt);
\fill[red!80!black] (12.0000,8.0000) circle (1.2pt);
\fill[red!80!black] (0.0000,8.0000) circle (1.2pt);
\fill[red!80!black] (12.0000,8.0000) circle (1.2pt);
\fill[red!80!black] (0.0000,8.0000) circle (1.2pt);
\fill[red!80!black] (12.0000,8.0000) circle (1.2pt);
\fill[red!80!black] (0.0000,8.0000) circle (1.2pt);
\fill[red!80!black] (12.0000,8.0000) circle (1.2pt);
\fill[red!80!black] (0.0000,8.0000) circle (1.2pt);
\fill[red!80!black] (12.0000,8.0000) circle (1.2pt);
\fill[red!80!black] (0.0000,8.0000) circle (1.2pt);
\fill[red!80!black] (12.0000,8.0000) circle (1.2pt);
\fill[red!80!black] (0.0000,8.0000) circle (1.2pt);
\fill[red!80!black] (12.0000,8.0000) circle (1.2pt);
\fill[red!80!black] (0.0000,8.0000) circle (1.2pt);
\fill[red!80!black] (12.0000,8.0000) circle (1.2pt);
\fill[red!80!black] (0.0000,8.0000) circle (1.2pt);
\fill[red!80!black] (12.0000,8.0000) circle (1.2pt);
\fill[red!80!black] (0.0000,8.0000) circle (1.2pt);
\fill[red!80!black] (12.0000,8.0000) circle (1.2pt);
\fill[red!80!black] (0.0000,8.0000) circle (1.2pt);
\fill[red!80!black] (12.0000,8.0000) circle (1.2pt);
\fill[red!80!black] (0.0000,8.0000) circle (1.2pt);
\fill[red!80!black] (12.0000,8.0000) circle (1.2pt);
\fill[red!80!black] (0.0000,8.0000) circle (1.2pt);
\fill[red!80!black] (12.0000,8.0000) circle (1.2pt);
\fill[red!80!black] (0.0000,8.0000) circle (1.2pt);
\fill[red!80!black] (12.0000,8.0000) circle (1.2pt);
\fill[red!80!black] (0.0000,8.0000) circle (1.2pt);
\fill[red!80!black] (12.0000,8.0000) circle (1.2pt);
\fill[red!80!black] (0.0000,8.0000) circle (1.2pt);
\fill[red!80!black] (12.0000,8.0000) circle (1.2pt);
\fill[red!80!black] (0.0000,8.0000) circle (1.2pt);
\fill[red!80!black] (12.0000,8.0000) circle (1.2pt);
\fill[red!80!black] (0.0000,8.0000) circle (1.2pt);
\fill[red!80!black] (12.0000,8.0000) circle (1.2pt);
\fill[red!80!black] (0.0000,8.0000) circle (1.2pt);
\fill[red!80!black] (12.0000,8.0000) circle (1.2pt);
\fill[red!80!black] (0.0000,8.0000) circle (1.2pt);
\fill[red!80!black] (12.0000,8.0000) circle (1.2pt);
\fill[red!80!black] (0.0000,8.0000) circle (1.2pt);
\fill[red!80!black] (12.0000,8.0000) circle (1.2pt);
\fill[red!80!black] (0.0000,8.0000) circle (1.2pt);
\fill[red!80!black] (12.0000,8.0000) circle (1.2pt);
\fill[red!80!black] (0.0000,8.0000) circle (1.2pt);
\fill[red!80!black] (12.0000,8.0000) circle (1.2pt);
\fill[red!80!black] (0.0000,8.0000) circle (1.2pt);
\fill[red!80!black] (12.0000,8.0000) circle (1.2pt);
\fill[red!80!black] (0.0000,8.0000) circle (1.2pt);
\fill[red!80!black] (12.0000,8.0000) circle (1.2pt);
\fill[red!80!black] (0.0000,8.0000) circle (1.2pt);
\fill[red!80!black] (12.0000,8.0000) circle (1.2pt);
\fill[red!80!black] (0.0000,8.0000) circle (1.2pt);
\fill[red!80!black] (12.0000,8.0000) circle (1.2pt);
\fill[red!80!black] (0.0000,8.0000) circle (1.2pt);
\fill[red!80!black] (12.0000,8.0000) circle (1.2pt);
\fill[red!80!black] (0.0000,8.0000) circle (1.2pt);
\fill[red!80!black] (12.0000,8.0000) circle (1.2pt);
\fill[red!80!black] (0.0000,8.0000) circle (1.2pt);
\fill[red!80!black] (12.0000,8.0000) circle (1.2pt);
\fill[red!80!black] (0.0000,8.0000) circle (1.2pt);
\fill[red!80!black] (12.0000,8.0000) circle (1.2pt);
\fill[red!80!black] (0.0000,8.0000) circle (1.2pt);
\fill[red!80!black] (12.0000,8.0000) circle (1.2pt);
\fill[red!80!black] (0.0000,8.0000) circle (1.2pt);
\fill[red!80!black] (12.0000,8.0000) circle (1.2pt);
\fill[red!80!black] (0.0000,8.0000) circle (1.2pt);
\fill[red!80!black] (12.0000,8.0000) circle (1.2pt);
\fill[red!80!black] (0.0000,8.0000) circle (1.2pt);
\fill[red!80!black] (12.0000,8.0000) circle (1.2pt);
\fill[red!80!black] (0.0000,8.0000) circle (1.2pt);
\fill[red!80!black] (12.0000,8.0000) circle (1.2pt);
\fill[red!80!black] (0.0000,8.0000) circle (1.2pt);
\fill[red!80!black] (12.0000,8.0000) circle (1.2pt);
\fill[red!80!black] (0.0000,8.0000) circle (1.2pt);
\fill[red!80!black] (12.0000,8.0000) circle (1.2pt);
\fill[red!80!black] (0.0000,8.0000) circle (1.2pt);
\fill[red!80!black] (12.0000,8.0000) circle (1.2pt);
\fill[red!80!black] (0.0000,8.0000) circle (1.2pt);
\fill[red!80!black] (12.0000,8.0000) circle (1.2pt);
\fill[red!80!black] (0.0000,8.0000) circle (1.2pt);
\fill[red!80!black] (12.0000,8.0000) circle (1.2pt);
\fill[red!80!black] (0.0000,8.0000) circle (1.2pt);
\fill[red!80!black] (12.0000,8.0000) circle (1.2pt);
\fill[red!80!black] (0.0000,8.0000) circle (1.2pt);
\fill[red!80!black] (12.0000,8.0000) circle (1.2pt);
\fill[red!80!black] (0.0000,8.0000) circle (1.2pt);
\fill[red!80!black] (12.0000,8.0000) circle (1.2pt);
\fill[red!80!black] (0.0000,8.0000) circle (1.2pt);
\fill[red!80!black] (12.0000,8.0000) circle (1.2pt);
\fill[red!80!black] (0.0000,8.0000) circle (1.2pt);
\fill[red!80!black] (12.0000,8.0000) circle (1.2pt);
\fill[red!80!black] (0.0000,8.0000) circle (1.2pt);
\fill[red!80!black] (12.0000,8.0000) circle (1.2pt);
\fill[red!80!black] (0.0000,8.0000) circle (1.2pt);
\fill[red!80!black] (12.0000,8.0000) circle (1.2pt);
\fill[red!80!black] (0.0000,8.0000) circle (1.2pt);
\fill[red!80!black] (12.0000,8.0000) circle (1.2pt);
\fill[red!80!black] (0.0000,8.0000) circle (1.2pt);
\fill[red!80!black] (12.0000,8.0000) circle (1.2pt);
\fill[red!80!black] (0.0000,8.0000) circle (1.2pt);
\fill[red!80!black] (12.0000,8.0000) circle (1.2pt);
\fill[red!80!black] (0.0000,8.0000) circle (1.2pt);
\fill[red!80!black] (12.0000,8.0000) circle (1.2pt);
\fill[red!80!black] (0.0000,8.0000) circle (1.2pt);
\fill[red!80!black] (12.0000,8.0000) circle (1.2pt);
\fill[red!80!black] (0.0000,8.0000) circle (1.2pt);
\fill[red!80!black] (12.0000,8.0000) circle (1.2pt);
\fill[red!80!black] (0.0000,8.0000) circle (1.2pt);
\fill[red!80!black] (12.0000,8.0000) circle (1.2pt);
\fill[red!80!black] (0.0000,8.0000) circle (1.2pt);
\fill[red!80!black] (12.0000,8.0000) circle (1.2pt);
\fill[red!80!black] (0.0000,8.0000) circle (1.2pt);
\fill[red!80!black] (12.0000,8.0000) circle (1.2pt);
\fill[red!80!black] (0.0000,8.0000) circle (1.2pt);
\fill[red!80!black] (12.0000,8.0000) circle (1.2pt);
\fill[red!80!black] (0.0000,8.0000) circle (1.2pt);
\fill[red!80!black] (12.0000,8.0000) circle (1.2pt);
\fill[red!80!black] (0.0000,8.0000) circle (1.2pt);
\fill[red!80!black] (12.0000,8.0000) circle (1.2pt);
\fill[red!80!black] (0.0000,8.0000) circle (1.2pt);
\fill[red!80!black] (12.0000,8.0000) circle (1.2pt);
\fill[red!80!black] (0.0000,8.0000) circle (1.2pt);
\fill[red!80!black] (12.0000,8.0000) circle (1.2pt);
\fill[red!80!black] (0.0000,8.0000) circle (1.2pt);
\fill[red!80!black] (12.0000,8.0000) circle (1.2pt);
\fill[red!80!black] (0.0000,8.0000) circle (1.2pt);
\fill[red!80!black] (12.0000,8.0000) circle (1.2pt);
\fill[red!80!black] (0.0000,8.0000) circle (1.2pt);
\fill[red!80!black] (12.0000,8.0000) circle (1.2pt);
\fill[red!80!black] (0.0000,8.0000) circle (1.2pt);
\fill[red!80!black] (12.0000,8.0000) circle (1.2pt);
\fill[red!80!black] (0.0000,8.0000) circle (1.2pt);
\fill[red!80!black] (12.0000,8.0000) circle (1.2pt);
\fill[red!80!black] (0.0000,8.0000) circle (1.2pt);
\fill[red!80!black] (12.0000,8.0000) circle (1.2pt);
\fill[red!80!black] (0.0000,8.0000) circle (1.2pt);
\fill[red!80!black] (12.0000,8.0000) circle (1.2pt);
\fill[red!80!black] (0.0000,8.0000) circle (1.2pt);
\fill[red!80!black] (12.0000,8.0000) circle (1.2pt);
\fill[red!80!black] (0.0000,8.0000) circle (1.2pt);
\fill[red!80!black] (12.0000,8.0000) circle (1.2pt);
\fill[red!80!black] (0.0000,8.0000) circle (1.2pt);
\fill[red!80!black] (12.0000,8.0000) circle (1.2pt);
\fill[red!80!black] (0.0000,8.0000) circle (1.2pt);
\fill[red!80!black] (12.0000,8.0000) circle (1.2pt);
\fill[red!80!black] (0.0000,8.0000) circle (1.2pt);
\fill[red!80!black] (12.0000,8.0000) circle (1.2pt);
\fill[red!80!black] (0.0000,8.0000) circle (1.2pt);
\fill[red!80!black] (12.0000,8.0000) circle (1.2pt);
\fill[red!80!black] (0.0000,8.0000) circle (1.2pt);
\fill[red!80!black] (12.0000,8.0000) circle (1.2pt);
\fill[red!80!black] (0.0000,8.0000) circle (1.2pt);
\fill[red!80!black] (12.0000,8.0000) circle (1.2pt);
\fill[red!80!black] (0.0000,8.0000) circle (1.2pt);
\fill[red!80!black] (12.0000,8.0000) circle (1.2pt);
\fill[red!80!black] (0.0000,8.0000) circle (1.2pt);
\fill[red!80!black] (12.0000,8.0000) circle (1.2pt);
\fill[red!80!black] (0.0000,8.0000) circle (1.2pt);
\fill[red!80!black] (12.0000,8.0000) circle (1.2pt);
\fill[red!80!black] (0.0000,8.0000) circle (1.2pt);
\fill[red!80!black] (12.0000,8.0000) circle (1.2pt);
\fill[red!80!black] (0.0000,8.0000) circle (1.2pt);
\fill[red!80!black] (12.0000,8.0000) circle (1.2pt);
\fill[red!80!black] (0.0000,8.0000) circle (1.2pt);
\fill[red!80!black] (12.0000,8.0000) circle (1.2pt);
\fill[red!80!black] (0.0000,8.0000) circle (1.2pt);
\fill[red!80!black] (12.0000,8.0000) circle (1.2pt);
\fill[red!80!black] (0.0000,8.0000) circle (1.2pt);
\fill[red!80!black] (12.0000,8.0000) circle (1.2pt);
\fill[red!80!black] (0.0000,8.0000) circle (1.2pt);
\fill[red!80!black] (12.0000,8.0000) circle (1.2pt);
\fill[red!80!black] (0.0000,8.0000) circle (1.2pt);
\fill[red!80!black] (12.0000,8.0000) circle (1.2pt);
\fill[red!80!black] (0.0000,8.0000) circle (1.2pt);
\fill[red!80!black] (12.0000,8.0000) circle (1.2pt);
\fill[red!80!black] (0.0000,8.0000) circle (1.2pt);
\fill[red!80!black] (12.0000,8.0000) circle (1.2pt);
\fill[red!80!black] (0.0000,8.0000) circle (1.2pt);
\fill[red!80!black] (12.0000,8.0000) circle (1.2pt);
\fill[red!80!black] (0.0000,8.0000) circle (1.2pt);
\fill[red!80!black] (12.0000,8.0000) circle (1.2pt);
\fill[red!80!black] (0.0000,8.0000) circle (1.2pt);
\fill[red!80!black] (12.0000,8.0000) circle (1.2pt);
\fill[red!80!black] (0.0000,8.0000) circle (1.2pt);
\fill[red!80!black] (12.0000,8.0000) circle (1.2pt);
\fill[red!80!black] (0.0000,8.0000) circle (1.2pt);
\fill[red!80!black] (12.0000,8.0000) circle (1.2pt);
\fill[red!80!black] (0.0000,8.0000) circle (1.2pt);
\fill[red!80!black] (12.0000,8.0000) circle (1.2pt);
\fill[red!80!black] (0.0000,8.0000) circle (1.2pt);
\fill[red!80!black] (12.0000,8.0000) circle (1.2pt);
\fill[red!80!black] (0.0000,8.0000) circle (1.2pt);
\fill[red!80!black] (12.0000,8.0000) circle (1.2pt);
\fill[red!80!black] (0.0000,8.0000) circle (1.2pt);
\fill[red!80!black] (12.0000,8.0000) circle (1.2pt);
\fill[red!80!black] (0.0000,8.0000) circle (1.2pt);
\fill[red!80!black] (12.0000,8.0000) circle (1.2pt);
\fill[red!80!black] (0.0000,8.0000) circle (1.2pt);
\fill[red!80!black] (12.0000,8.0000) circle (1.2pt);
\fill[red!80!black] (0.0000,8.0000) circle (1.2pt);
\fill[red!80!black] (12.0000,8.0000) circle (1.2pt);
\fill[red!80!black] (0.0000,8.0000) circle (1.2pt);
\fill[red!80!black] (12.0000,8.0000) circle (1.2pt);
\fill[red!80!black] (0.0000,8.0000) circle (1.2pt);
\fill[red!80!black] (12.0000,8.0000) circle (1.2pt);
\fill[red!80!black] (0.0000,8.0000) circle (1.2pt);
\fill[red!80!black] (12.0000,8.0000) circle (1.2pt);
\fill[red!80!black] (0.0000,8.0000) circle (1.2pt);
\fill[red!80!black] (12.0000,8.0000) circle (1.2pt);
\fill[red!80!black] (0.0000,8.0000) circle (1.2pt);
\fill[red!80!black] (12.0000,8.0000) circle (1.2pt);
\fill[red!80!black] (0.0000,8.0000) circle (1.2pt);
\fill[red!80!black] (12.0000,8.0000) circle (1.2pt);
\fill[red!80!black] (0.0000,8.0000) circle (1.2pt);
\fill[red!80!black] (12.0000,8.0000) circle (1.2pt);
\fill[red!80!black] (0.0000,8.0000) circle (1.2pt);
\fill[red!80!black] (12.0000,8.0000) circle (1.2pt);
\fill[red!80!black] (0.0000,8.0000) circle (1.2pt);
\fill[red!80!black] (12.0000,8.0000) circle (1.2pt);
\fill[red!80!black] (0.0000,8.0000) circle (1.2pt);
\fill[red!80!black] (12.0000,8.0000) circle (1.2pt);
\fill[red!80!black] (0.0000,8.0000) circle (1.2pt);
\fill[red!80!black] (12.0000,8.0000) circle (1.2pt);
\fill[red!80!black] (0.0000,8.0000) circle (1.2pt);
\fill[red!80!black] (12.0000,8.0000) circle (1.2pt);
\fill[red!80!black] (0.0000,8.0000) circle (1.2pt);
\fill[red!80!black] (12.0000,8.0000) circle (1.2pt);
\fill[red!80!black] (0.0000,8.0000) circle (1.2pt);
\fill[red!80!black] (12.0000,8.0000) circle (1.2pt);
\fill[red!80!black] (0.0000,8.0000) circle (1.2pt);
\fill[red!80!black] (12.0000,8.0000) circle (1.2pt);
\fill[red!80!black] (0.0000,8.0000) circle (1.2pt);
\fill[red!80!black] (12.0000,8.0000) circle (1.2pt);
\fill[red!80!black] (0.0000,8.0000) circle (1.2pt);
\fill[red!80!black] (12.0000,8.0000) circle (1.2pt);
\fill[red!80!black] (0.0000,8.0000) circle (1.2pt);
\fill[red!80!black] (12.0000,8.0000) circle (1.2pt);
\fill[red!80!black] (0.0000,8.0000) circle (1.2pt);
\fill[red!80!black] (12.0000,8.0000) circle (1.2pt);
\fill[red!80!black] (0.0000,8.0000) circle (1.2pt);
\fill[red!80!black] (12.0000,8.0000) circle (1.2pt);
\fill[red!80!black] (0.0000,8.0000) circle (1.2pt);
\fill[red!80!black] (12.0000,8.0000) circle (1.2pt);
\fill[red!80!black] (0.0000,8.0000) circle (1.2pt);
\fill[red!80!black] (12.0000,8.0000) circle (1.2pt);
\fill[red!80!black] (0.0000,8.0000) circle (1.2pt);
\fill[red!80!black] (12.0000,8.0000) circle (1.2pt);
\fill[red!80!black] (0.0000,8.0000) circle (1.2pt);
\fill[red!80!black] (12.0000,8.0000) circle (1.2pt);
\fill[red!80!black] (0.0000,8.0000) circle (1.2pt);
\fill[red!80!black] (12.0000,8.0000) circle (1.2pt);
\fill[red!80!black] (0.0000,8.0000) circle (1.2pt);
\fill[red!80!black] (12.0000,8.0000) circle (1.2pt);
\fill[red!80!black] (0.0000,8.0000) circle (1.2pt);
\fill[red!80!black] (12.0000,8.0000) circle (1.2pt);
\fill[red!80!black] (0.0000,8.0000) circle (1.2pt);
\fill[red!80!black] (12.0000,8.0000) circle (1.2pt);
\fill[red!80!black] (0.0000,8.0000) circle (1.2pt);
\fill[red!80!black] (12.0000,8.0000) circle (1.2pt);
\fill[red!80!black] (0.0000,8.0000) circle (1.2pt);
\fill[red!80!black] (12.0000,8.0000) circle (1.2pt);
\fill[red!80!black] (0.0000,8.0000) circle (1.2pt);
\fill[red!80!black] (12.0000,8.0000) circle (1.2pt);
\fill[red!80!black] (0.0000,8.0000) circle (1.2pt);
\fill[red!80!black] (12.0000,8.0000) circle (1.2pt);
\fill[red!80!black] (0.0000,8.0000) circle (1.2pt);
\fill[red!80!black] (12.0000,8.0000) circle (1.2pt);
\fill[red!80!black] (0.0000,8.0000) circle (1.2pt);
\fill[red!80!black] (12.0000,8.0000) circle (1.2pt);
\fill[red!80!black] (0.0000,8.0000) circle (1.2pt);
\fill[red!80!black] (12.0000,8.0000) circle (1.2pt);
\fill[red!80!black] (0.0000,8.0000) circle (1.2pt);
\fill[red!80!black] (12.0000,8.0000) circle (1.2pt);
\fill[red!80!black] (0.0000,8.0000) circle (1.2pt);
\fill[red!80!black] (12.0000,8.0000) circle (1.2pt);
\fill[red!80!black] (0.0000,8.0000) circle (1.2pt);
\fill[red!80!black] (12.0000,8.0000) circle (1.2pt);
\fill[red!80!black] (0.0000,8.0000) circle (1.2pt);
\fill[red!80!black] (12.0000,8.0000) circle (1.2pt);
\fill[red!80!black] (0.0000,8.0000) circle (1.2pt);
\fill[red!80!black] (12.0000,8.0000) circle (1.2pt);
\fill[red!80!black] (0.0000,8.0000) circle (1.2pt);
\fill[red!80!black] (12.0000,8.0000) circle (1.2pt);
\fill[red!80!black] (0.0000,8.0000) circle (1.2pt);
\fill[red!80!black] (12.0000,8.0000) circle (1.2pt);
\fill[red!80!black] (0.0000,8.0000) circle (1.2pt);
\fill[red!80!black] (12.0000,8.0000) circle (1.2pt);
\fill[red!80!black] (0.0000,8.0000) circle (1.2pt);
\fill[red!80!black] (12.0000,8.0000) circle (1.2pt);
\fill[red!80!black] (0.0000,8.0000) circle (1.2pt);
\fill[red!80!black] (12.0000,8.0000) circle (1.2pt);
\fill[red!80!black] (0.0000,8.0000) circle (1.2pt);
\fill[red!80!black] (12.0000,8.0000) circle (1.2pt);
\fill[red!80!black] (0.0000,8.0000) circle (1.2pt);
\fill[red!80!black] (12.0000,8.0000) circle (1.2pt);
\fill[red!80!black] (0.0000,8.0000) circle (1.2pt);
\fill[red!80!black] (12.0000,8.0000) circle (1.2pt);
\fill[red!80!black] (0.0000,8.0000) circle (1.2pt);
\fill[red!80!black] (12.0000,8.0000) circle (1.2pt);
\fill[red!80!black] (0.0000,8.0000) circle (1.2pt);
\fill[red!80!black] (12.0000,8.0000) circle (1.2pt);
\fill[red!80!black] (0.0000,8.0000) circle (1.2pt);
\fill[red!80!black] (12.0000,8.0000) circle (1.2pt);
\fill[red!80!black] (0.0000,8.0000) circle (1.2pt);
\fill[red!80!black] (12.0000,8.0000) circle (1.2pt);
\fill[red!80!black] (0.0000,8.0000) circle (1.2pt);
\fill[red!80!black] (12.0000,8.0000) circle (1.2pt);
\fill[red!80!black] (0.0000,8.0000) circle (1.2pt);
\fill[red!80!black] (12.0000,8.0000) circle (1.2pt);
\fill[red!80!black] (0.0000,8.0000) circle (1.2pt);
\fill[red!80!black] (12.0000,8.0000) circle (1.2pt);
\fill[red!80!black] (0.0000,8.0000) circle (1.2pt);
\fill[red!80!black] (12.0000,8.0000) circle (1.2pt);
\fill[red!80!black] (0.0000,8.0000) circle (1.2pt);
\fill[red!80!black] (12.0000,8.0000) circle (1.2pt);
\fill[red!80!black] (0.0000,8.0000) circle (1.2pt);
\fill[red!80!black] (12.0000,8.0000) circle (1.2pt);
\fill[red!80!black] (0.0000,8.0000) circle (1.2pt);
\fill[red!80!black] (12.0000,8.0000) circle (1.2pt);
\fill[red!80!black] (0.0000,8.0000) circle (1.2pt);
\fill[red!80!black] (12.0000,8.0000) circle (1.2pt);
\fill[red!80!black] (0.0000,8.0000) circle (1.2pt);
\fill[red!80!black] (12.0000,8.0000) circle (1.2pt);
\fill[red!80!black] (0.0000,8.0000) circle (1.2pt);
\fill[red!80!black] (12.0000,8.0000) circle (1.2pt);
\fill[red!80!black] (0.0000,8.0000) circle (1.2pt);
\fill[red!80!black] (12.0000,8.0000) circle (1.2pt);
\fill[red!80!black] (0.0000,8.0000) circle (1.2pt);
\fill[red!80!black] (12.0000,8.0000) circle (1.2pt);
\fill[red!80!black] (0.0000,8.0000) circle (1.2pt);
\fill[red!80!black] (12.0000,8.0000) circle (1.2pt);
\fill[red!80!black] (0.0000,8.0000) circle (1.2pt);
\fill[red!80!black] (12.0000,8.0000) circle (1.2pt);
\fill[red!80!black] (0.0000,8.0000) circle (1.2pt);
\fill[red!80!black] (12.0000,8.0000) circle (1.2pt);
\fill[red!80!black] (0.0000,8.0000) circle (1.2pt);
\fill[red!80!black] (12.0000,8.0000) circle (1.2pt);
\fill[red!80!black] (0.0000,8.0000) circle (1.2pt);
\fill[red!80!black] (12.0000,8.0000) circle (1.2pt);
\fill[red!80!black] (0.0000,8.0000) circle (1.2pt);
\fill[red!80!black] (12.0000,8.0000) circle (1.2pt);
\fill[red!80!black] (0.0000,8.0000) circle (1.2pt);
\fill[red!80!black] (12.0000,8.0000) circle (1.2pt);
\fill[red!80!black] (0.0000,8.0000) circle (1.2pt);
\fill[red!80!black] (12.0000,8.0000) circle (1.2pt);
\fill[red!80!black] (0.0000,8.0000) circle (1.2pt);
\fill[red!80!black] (12.0000,8.0000) circle (1.2pt);
\fill[red!80!black] (0.0000,8.0000) circle (1.2pt);
\fill[red!80!black] (12.0000,8.0000) circle (1.2pt);
\fill[red!80!black] (0.0000,8.0000) circle (1.2pt);
\fill[red!80!black] (12.0000,8.0000) circle (1.2pt);
\fill[red!80!black] (0.0000,8.0000) circle (1.2pt);
\fill[red!80!black] (12.0000,8.0000) circle (1.2pt);
\fill[red!80!black] (0.0000,8.0000) circle (1.2pt);
\fill[red!80!black] (12.0000,8.0000) circle (1.2pt);
\fill[red!80!black] (0.0000,8.0000) circle (1.2pt);
\fill[red!80!black] (12.0000,8.0000) circle (1.2pt);
\fill[red!80!black] (0.0000,8.0000) circle (1.2pt);
\fill[red!80!black] (12.0000,8.0000) circle (1.2pt);
\fill[red!80!black] (0.0000,8.0000) circle (1.2pt);
\fill[red!80!black] (12.0000,8.0000) circle (1.2pt);
\fill[red!80!black] (0.0000,8.0000) circle (1.2pt);
\fill[red!80!black] (12.0000,8.0000) circle (1.2pt);
\fill[red!80!black] (0.0000,8.0000) circle (1.2pt);
\fill[red!80!black] (12.0000,8.0000) circle (1.2pt);
\fill[red!80!black] (0.0000,8.0000) circle (1.2pt);
\fill[red!80!black] (12.0000,8.0000) circle (1.2pt);
\fill[red!80!black] (0.0000,8.0000) circle (1.2pt);
\fill[red!80!black] (12.0000,8.0000) circle (1.2pt);
\fill[red!80!black] (0.0000,8.0000) circle (1.2pt);
\fill[red!80!black] (12.0000,8.0000) circle (1.2pt);
\fill[red!80!black] (0.0000,8.0000) circle (1.2pt);
\fill[red!80!black] (12.0000,8.0000) circle (1.2pt);
\fill[red!80!black] (0.0000,8.0000) circle (1.2pt);
\fill[red!80!black] (12.0000,8.0000) circle (1.2pt);
\fill[red!80!black] (0.0000,8.0000) circle (1.2pt);
\fill[red!80!black] (12.0000,8.0000) circle (1.2pt);
\fill[red!80!black] (0.0000,8.0000) circle (1.2pt);
\fill[red!80!black] (12.0000,8.0000) circle (1.2pt);
\fill[red!80!black] (0.0000,8.0000) circle (1.2pt);
\fill[red!80!black] (12.0000,8.0000) circle (1.2pt);
\fill[red!80!black] (0.0000,8.0000) circle (1.2pt);
\fill[red!80!black] (12.0000,8.0000) circle (1.2pt);
\fill[red!80!black] (0.0000,8.0000) circle (1.2pt);
\fill[red!80!black] (12.0000,8.0000) circle (1.2pt);
\fill[red!80!black] (0.0000,8.0000) circle (1.2pt);
\fill[red!80!black] (12.0000,8.0000) circle (1.2pt);
\fill[red!80!black] (0.0000,8.0000) circle (1.2pt);
\fill[red!80!black] (12.0000,8.0000) circle (1.2pt);
\fill[red!80!black] (0.0000,8.0000) circle (1.2pt);
\fill[red!80!black] (12.0000,8.0000) circle (1.2pt);
\fill[red!80!black] (0.0000,8.0000) circle (1.2pt);
\fill[red!80!black] (12.0000,8.0000) circle (1.2pt);
\fill[red!80!black] (0.0000,8.0000) circle (1.2pt);
\fill[red!80!black] (12.0000,8.0000) circle (1.2pt);
\fill[red!80!black] (0.0000,8.0000) circle (1.2pt);
\fill[red!80!black] (12.0000,8.0000) circle (1.2pt);
\fill[red!80!black] (0.0000,8.0000) circle (1.2pt);
\fill[red!80!black] (12.0000,8.0000) circle (1.2pt);
\fill[red!80!black] (0.0000,8.0000) circle (1.2pt);
\fill[red!80!black] (12.0000,8.0000) circle (1.2pt);
\fill[red!80!black] (0.0000,8.0000) circle (1.2pt);
\fill[red!80!black] (12.0000,8.0000) circle (1.2pt);
\fill[red!80!black] (0.0000,8.0000) circle (1.2pt);
\fill[red!80!black] (12.0000,8.0000) circle (1.2pt);
\fill[red!80!black] (0.0000,8.0000) circle (1.2pt);
\fill[red!80!black] (12.0000,8.0000) circle (1.2pt);
\fill[red!80!black] (0.0000,8.0000) circle (1.2pt);
\fill[red!80!black] (12.0000,8.0000) circle (1.2pt);
\fill[red!80!black] (0.0000,8.0000) circle (1.2pt);
\fill[red!80!black] (12.0000,8.0000) circle (1.2pt);
\fill[red!80!black] (0.0000,8.0000) circle (1.2pt);
\fill[red!80!black] (12.0000,8.0000) circle (1.2pt);
\fill[red!80!black] (0.0000,8.0000) circle (1.2pt);
\fill[red!80!black] (12.0000,8.0000) circle (1.2pt);
\fill[red!80!black] (0.0000,8.0000) circle (1.2pt);
\fill[red!80!black] (12.0000,8.0000) circle (1.2pt);
\fill[red!80!black] (0.0000,8.0000) circle (1.2pt);
\fill[red!80!black] (12.0000,8.0000) circle (1.2pt);
\fill[red!80!black] (0.0000,8.0000) circle (1.2pt);
\fill[red!80!black] (12.0000,8.0000) circle (1.2pt);
\fill[red!80!black] (0.0000,8.0000) circle (1.2pt);
\fill[red!80!black] (12.0000,8.0000) circle (1.2pt);
\fill[red!80!black] (0.0000,8.0000) circle (1.2pt);
\fill[red!80!black] (12.0000,8.0000) circle (1.2pt);
\fill[red!80!black] (0.0000,8.0000) circle (1.2pt);
\fill[red!80!black] (12.0000,8.0000) circle (1.2pt);
\fill[red!80!black] (0.0000,8.0000) circle (1.2pt);
\fill[red!80!black] (12.0000,8.0000) circle (1.2pt);
\fill[red!80!black] (0.0000,8.0000) circle (1.2pt);
\fill[red!80!black] (12.0000,8.0000) circle (1.2pt);
\fill[red!80!black] (0.0000,8.0000) circle (1.2pt);
\fill[red!80!black] (12.0000,8.0000) circle (1.2pt);
\fill[red!80!black] (0.0000,8.0000) circle (1.2pt);
\fill[red!80!black] (12.0000,8.0000) circle (1.2pt);
\fill[red!80!black] (0.0000,8.0000) circle (1.2pt);
\fill[red!80!black] (12.0000,8.0000) circle (1.2pt);
\fill[red!80!black] (0.0000,8.0000) circle (1.2pt);
\fill[red!80!black] (12.0000,8.0000) circle (1.2pt);
\fill[red!80!black] (0.0000,8.0000) circle (1.2pt);
\fill[red!80!black] (12.0000,8.0000) circle (1.2pt);
\fill[red!80!black] (0.0000,8.0000) circle (1.2pt);
\fill[red!80!black] (12.0000,8.0000) circle (1.2pt);
\fill[red!80!black] (0.0000,8.0000) circle (1.2pt);
\fill[red!80!black] (12.0000,8.0000) circle (1.2pt);
\fill[red!80!black] (0.0000,8.0000) circle (1.2pt);
\fill[red!80!black] (12.0000,8.0000) circle (1.2pt);
\fill[red!80!black] (0.0000,8.0000) circle (1.2pt);
\fill[red!80!black] (12.0000,8.0000) circle (1.2pt);
\fill[red!80!black] (0.0000,8.0000) circle (1.2pt);
\fill[red!80!black] (12.0000,8.0000) circle (1.2pt);
\fill[red!80!black] (0.0000,8.0000) circle (1.2pt);
\fill[red!80!black] (12.0000,8.0000) circle (1.2pt);
\fill[red!80!black] (0.0000,8.0000) circle (1.2pt);
\fill[red!80!black] (12.0000,8.0000) circle (1.2pt);
\fill[red!80!black] (0.0000,8.0000) circle (1.2pt);
\fill[red!80!black] (12.0000,8.0000) circle (1.2pt);
\fill[red!80!black] (0.0000,8.0000) circle (1.2pt);
\fill[red!80!black] (12.0000,8.0000) circle (1.2pt);
\fill[red!80!black] (0.0000,8.0000) circle (1.2pt);
\fill[red!80!black] (12.0000,8.0000) circle (1.2pt);
\fill[red!80!black] (0.0000,8.0000) circle (1.2pt);
\fill[red!80!black] (12.0000,8.0000) circle (1.2pt);
\fill[red!80!black] (0.0000,8.0000) circle (1.2pt);
\fill[red!80!black] (12.0000,8.0000) circle (1.2pt);
\fill[red!80!black] (0.0000,8.0000) circle (1.2pt);
\fill[red!80!black] (12.0000,8.0000) circle (1.2pt);
\fill[red!80!black] (0.0000,8.0000) circle (1.2pt);
\fill[red!80!black] (12.0000,8.0000) circle (1.2pt);
\fill[red!80!black] (0.0000,8.0000) circle (1.2pt);
\fill[red!80!black] (12.0000,8.0000) circle (1.2pt);
\fill[red!80!black] (0.0000,8.0000) circle (1.2pt);
\fill[red!80!black] (12.0000,8.0000) circle (1.2pt);
\fill[red!80!black] (0.0000,8.0000) circle (1.2pt);
\fill[red!80!black] (12.0000,8.0000) circle (1.2pt);
\fill[red!80!black] (0.0000,8.0000) circle (1.2pt);
\fill[red!80!black] (12.0000,8.0000) circle (1.2pt);
\fill[red!80!black] (0.0000,8.0000) circle (1.2pt);
\fill[red!80!black] (12.0000,8.0000) circle (1.2pt);
\fill[red!80!black] (0.0000,8.0000) circle (1.2pt);
\fill[red!80!black] (12.0000,8.0000) circle (1.2pt);
\fill[red!80!black] (0.0000,8.0000) circle (1.2pt);
\fill[red!80!black] (12.0000,8.0000) circle (1.2pt);
\fill[red!80!black] (0.0000,8.0000) circle (1.2pt);
\fill[red!80!black] (12.0000,8.0000) circle (1.2pt);
\fill[red!80!black] (0.0000,8.0000) circle (1.2pt);
\fill[red!80!black] (12.0000,8.0000) circle (1.2pt);
\fill[red!80!black] (0.0000,8.0000) circle (1.2pt);
\fill[red!80!black] (12.0000,8.0000) circle (1.2pt);
\fill[red!80!black] (0.0000,8.0000) circle (1.2pt);
\fill[red!80!black] (12.0000,8.0000) circle (1.2pt);
\fill[red!80!black] (0.0000,8.0000) circle (1.2pt);
\fill[red!80!black] (12.0000,8.0000) circle (1.2pt);
\fill[red!80!black] (0.0000,8.0000) circle (1.2pt);
\fill[red!80!black] (12.0000,8.0000) circle (1.2pt);
\fill[red!80!black] (0.0000,8.0000) circle (1.2pt);
\fill[red!80!black] (12.0000,8.0000) circle (1.2pt);
\fill[red!80!black] (0.0000,8.0000) circle (1.2pt);
\fill[red!80!black] (12.0000,8.0000) circle (1.2pt);
\fill[red!80!black] (0.0000,8.0000) circle (1.2pt);
\fill[red!80!black] (12.0000,8.0000) circle (1.2pt);
\fill[red!80!black] (0.0000,8.0000) circle (1.2pt);
\fill[red!80!black] (12.0000,8.0000) circle (1.2pt);
\fill[red!80!black] (0.0000,8.0000) circle (1.2pt);
\fill[red!80!black] (12.0000,8.0000) circle (1.2pt);
\fill[red!80!black] (0.0000,8.0000) circle (1.2pt);
\fill[red!80!black] (12.0000,8.0000) circle (1.2pt);
\fill[red!80!black] (0.0000,8.0000) circle (1.2pt);
\fill[red!80!black] (12.0000,8.0000) circle (1.2pt);
\fill[red!80!black] (0.0000,8.0000) circle (1.2pt);
\fill[red!80!black] (12.0000,8.0000) circle (1.2pt);
\fill[red!80!black] (0.0000,8.0000) circle (1.2pt);
\fill[red!80!black] (12.0000,8.0000) circle (1.2pt);
\fill[red!80!black] (0.0000,8.0000) circle (1.2pt);
\fill[red!80!black] (12.0000,8.0000) circle (1.2pt);
\fill[red!80!black] (0.0000,8.0000) circle (1.2pt);
\fill[red!80!black] (12.0000,8.0000) circle (1.2pt);
\fill[red!80!black] (0.0000,8.0000) circle (1.2pt);
\fill[red!80!black] (12.0000,8.0000) circle (1.2pt);
\fill[red!80!black] (0.0000,8.0000) circle (1.2pt);
\fill[red!80!black] (12.0000,8.0000) circle (1.2pt);
\fill[red!80!black] (0.0000,8.0000) circle (1.2pt);
\fill[red!80!black] (12.0000,8.0000) circle (1.2pt);
\fill[red!80!black] (0.0000,8.0000) circle (1.2pt);
\fill[red!80!black] (12.0000,8.0000) circle (1.2pt);
\fill[red!80!black] (0.0000,8.0000) circle (1.2pt);
\fill[red!80!black] (12.0000,8.0000) circle (1.2pt);
\fill[red!80!black] (0.0000,8.0000) circle (1.2pt);
\fill[red!80!black] (12.0000,8.0000) circle (1.2pt);
\fill[red!80!black] (0.0000,8.0000) circle (1.2pt);
\fill[red!80!black] (12.0000,8.0000) circle (1.2pt);
\fill[red!80!black] (0.0000,8.0000) circle (1.2pt);
\fill[red!80!black] (12.0000,8.0000) circle (1.2pt);
\fill[red!80!black] (0.0000,8.0000) circle (1.2pt);
\fill[red!80!black] (12.0000,8.0000) circle (1.2pt);
\fill[red!80!black] (0.0000,8.0000) circle (1.2pt);
\fill[red!80!black] (12.0000,8.0000) circle (1.2pt);
\fill[red!80!black] (0.0000,8.0000) circle (1.2pt);
\fill[red!80!black] (12.0000,8.0000) circle (1.2pt);
\fill[red!80!black] (0.0000,8.0000) circle (1.2pt);
\fill[red!80!black] (12.0000,8.0000) circle (1.2pt);
\fill[red!80!black] (0.0000,8.0000) circle (1.2pt);
\fill[red!80!black] (12.0000,8.0000) circle (1.2pt);
\fill[red!80!black] (0.0000,8.0000) circle (1.2pt);
\fill[red!80!black] (12.0000,8.0000) circle (1.2pt);
\fill[red!80!black] (0.0000,8.0000) circle (1.2pt);
\fill[red!80!black] (12.0000,8.0000) circle (1.2pt);
\fill[red!80!black] (0.0000,8.0000) circle (1.2pt);
\fill[red!80!black] (12.0000,8.0000) circle (1.2pt);
\fill[red!80!black] (0.0000,8.0000) circle (1.2pt);
\fill[red!80!black] (12.0000,8.0000) circle (1.2pt);
\fill[red!80!black] (0.0000,8.0000) circle (1.2pt);
\fill[red!80!black] (12.0000,8.0000) circle (1.2pt);
\fill[red!80!black] (0.0000,8.0000) circle (1.2pt);
\fill[red!80!black] (12.0000,8.0000) circle (1.2pt);
\fill[red!80!black] (0.0000,8.0000) circle (1.2pt);
\fill[red!80!black] (12.0000,8.0000) circle (1.2pt);
\fill[red!80!black] (0.0000,8.0000) circle (1.2pt);
\fill[red!80!black] (12.0000,8.0000) circle (1.2pt);
\fill[red!80!black] (0.0000,8.0000) circle (1.2pt);
\fill[red!80!black] (12.0000,8.0000) circle (1.2pt);
\fill[red!80!black] (0.0000,8.0000) circle (1.2pt);
\fill[red!80!black] (12.0000,8.0000) circle (1.2pt);
\fill[red!80!black] (0.0000,8.0000) circle (1.2pt);
\fill[red!80!black] (12.0000,8.0000) circle (1.2pt);
\fill[red!80!black] (0.0000,8.0000) circle (1.2pt);
\fill[red!80!black] (12.0000,8.0000) circle (1.2pt);
\fill[red!80!black] (0.0000,8.0000) circle (1.2pt);
\fill[red!80!black] (12.0000,8.0000) circle (1.2pt);
\fill[red!80!black] (0.0000,8.0000) circle (1.2pt);
\fill[red!80!black] (12.0000,8.0000) circle (1.2pt);
\fill[red!80!black] (0.0000,8.0000) circle (1.2pt);
\fill[red!80!black] (12.0000,8.0000) circle (1.2pt);
\fill[red!80!black] (0.0000,8.0000) circle (1.2pt);
\fill[red!80!black] (12.0000,8.0000) circle (1.2pt);
\fill[red!80!black] (0.0000,8.0000) circle (1.2pt);
\fill[red!80!black] (12.0000,8.0000) circle (1.2pt);
\fill[red!80!black] (0.0000,8.0000) circle (1.2pt);
\fill[red!80!black] (12.0000,8.0000) circle (1.2pt);
\fill[red!80!black] (0.0000,8.0000) circle (1.2pt);
\fill[red!80!black] (12.0000,8.0000) circle (1.2pt);
\fill[red!80!black] (0.0000,8.0000) circle (1.2pt);
\fill[red!80!black] (12.0000,8.0000) circle (1.2pt);
\fill[red!80!black] (0.0000,8.0000) circle (1.2pt);
\fill[red!80!black] (12.0000,8.0000) circle (1.2pt);
\fill[red!80!black] (0.0000,8.0000) circle (1.2pt);
\fill[red!80!black] (12.0000,8.0000) circle (1.2pt);
\fill[red!80!black] (0.0000,8.0000) circle (1.2pt);
\fill[red!80!black] (12.0000,8.0000) circle (1.2pt);
\fill[red!80!black] (0.0000,8.0000) circle (1.2pt);
\fill[red!80!black] (12.0000,8.0000) circle (1.2pt);
\fill[red!80!black] (0.0000,8.0000) circle (1.2pt);
\fill[red!80!black] (12.0000,8.0000) circle (1.2pt);
\fill[red!80!black] (0.0000,8.0000) circle (1.2pt);
\fill[red!80!black] (12.0000,8.0000) circle (1.2pt);
\fill[red!80!black] (0.0000,8.0000) circle (1.2pt);
\fill[red!80!black] (12.0000,8.0000) circle (1.2pt);
\fill[red!80!black] (0.0000,8.0000) circle (1.2pt);
\fill[red!80!black] (12.0000,8.0000) circle (1.2pt);
\fill[red!80!black] (0.0000,8.0000) circle (1.2pt);
\fill[red!80!black] (12.0000,8.0000) circle (1.2pt);
\fill[red!80!black] (0.0000,8.0000) circle (1.2pt);
\fill[red!80!black] (12.0000,8.0000) circle (1.2pt);
\fill[red!80!black] (0.0000,8.0000) circle (1.2pt);
\fill[red!80!black] (12.0000,8.0000) circle (1.2pt);
\fill[red!80!black] (0.0000,8.0000) circle (1.2pt);
\fill[red!80!black] (12.0000,8.0000) circle (1.2pt);
\fill[red!80!black] (0.0000,8.0000) circle (1.2pt);
\fill[red!80!black] (12.0000,8.0000) circle (1.2pt);
\fill[red!80!black] (0.0000,8.0000) circle (1.2pt);
\fill[red!80!black] (12.0000,8.0000) circle (1.2pt);
\fill[red!80!black] (0.0000,8.0000) circle (1.2pt);
\fill[red!80!black] (12.0000,8.0000) circle (1.2pt);
\fill[red!80!black] (0.0000,8.0000) circle (1.2pt);
\fill[red!80!black] (12.0000,8.0000) circle (1.2pt);
\fill[red!80!black] (0.0000,8.0000) circle (1.2pt);
\fill[red!80!black] (12.0000,8.0000) circle (1.2pt);
\fill[red!80!black] (0.0000,8.0000) circle (1.2pt);
\fill[red!80!black] (12.0000,8.0000) circle (1.2pt);
\fill[red!80!black] (0.0000,8.0000) circle (1.2pt);
\fill[red!80!black] (12.0000,8.0000) circle (1.2pt);
\fill[red!80!black] (0.0000,8.0000) circle (1.2pt);
\fill[red!80!black] (12.0000,8.0000) circle (1.2pt);
\fill[red!80!black] (0.0000,8.0000) circle (1.2pt);
\fill[red!80!black] (12.0000,8.0000) circle (1.2pt);
\fill[red!80!black] (0.0000,8.0000) circle (1.2pt);
\fill[red!80!black] (12.0000,8.0000) circle (1.2pt);
\fill[red!80!black] (0.0000,8.0000) circle (1.2pt);
\fill[red!80!black] (12.0000,8.0000) circle (1.2pt);
\fill[red!80!black] (0.0000,8.0000) circle (1.2pt);
\fill[red!80!black] (12.0000,8.0000) circle (1.2pt);
\fill[red!80!black] (0.0000,8.0000) circle (1.2pt);
\fill[red!80!black] (12.0000,8.0000) circle (1.2pt);
\fill[red!80!black] (0.0000,8.0000) circle (1.2pt);
\fill[red!80!black] (12.0000,8.0000) circle (1.2pt);
\fill[red!80!black] (0.0000,8.0000) circle (1.2pt);
\fill[red!80!black] (12.0000,8.0000) circle (1.2pt);
\fill[red!80!black] (0.0000,8.0000) circle (1.2pt);
\fill[red!80!black] (12.0000,8.0000) circle (1.2pt);
\fill[red!80!black] (0.0000,8.0000) circle (1.2pt);
\fill[red!80!black] (12.0000,8.0000) circle (1.2pt);
\fill[red!80!black] (0.0000,8.0000) circle (1.2pt);
\fill[red!80!black] (12.0000,8.0000) circle (1.2pt);
\fill[red!80!black] (0.0000,8.0000) circle (1.2pt);
\fill[red!80!black] (12.0000,8.0000) circle (1.2pt);
\fill[red!80!black] (0.0000,8.0000) circle (1.2pt);
\fill[red!80!black] (12.0000,8.0000) circle (1.2pt);
\fill[red!80!black] (0.0000,8.0000) circle (1.2pt);
\fill[red!80!black] (12.0000,8.0000) circle (1.2pt);
\fill[red!80!black] (0.0000,8.0000) circle (1.2pt);
\fill[red!80!black] (12.0000,8.0000) circle (1.2pt);
\fill[red!80!black] (0.0000,8.0000) circle (1.2pt);
\fill[red!80!black] (12.0000,8.0000) circle (1.2pt);
\fill[red!80!black] (0.0000,8.0000) circle (1.2pt);
\fill[red!80!black] (12.0000,8.0000) circle (1.2pt);
\fill[red!80!black] (0.0000,8.0000) circle (1.2pt);
\fill[red!80!black] (12.0000,8.0000) circle (1.2pt);
\fill[red!80!black] (0.0000,8.0000) circle (1.2pt);
\fill[red!80!black] (12.0000,8.0000) circle (1.2pt);
\fill[red!80!black] (0.0000,8.0000) circle (1.2pt);
\fill[red!80!black] (12.0000,8.0000) circle (1.2pt);
\fill[red!80!black] (0.0000,8.0000) circle (1.2pt);
\fill[red!80!black] (12.0000,8.0000) circle (1.2pt);
\fill[red!80!black] (0.0000,8.0000) circle (1.2pt);
\fill[red!80!black] (12.0000,8.0000) circle (1.2pt);
\fill[red!80!black] (0.0000,8.0000) circle (1.2pt);
\fill[red!80!black] (12.0000,8.0000) circle (1.2pt);
\fill[red!80!black] (0.0000,8.0000) circle (1.2pt);
\fill[red!80!black] (12.0000,8.0000) circle (1.2pt);
\fill[red!80!black] (0.0000,8.0000) circle (1.2pt);
\fill[red!80!black] (12.0000,8.0000) circle (1.2pt);
\fill[red!80!black] (0.0000,8.0000) circle (1.2pt);
\fill[red!80!black] (12.0000,8.0000) circle (1.2pt);
\fill[red!80!black] (0.0000,8.0000) circle (1.2pt);
\fill[red!80!black] (12.0000,8.0000) circle (1.2pt);
\fill[red!80!black] (0.0000,8.0000) circle (1.2pt);
\fill[red!80!black] (12.0000,8.0000) circle (1.2pt);
\fill[red!80!black] (0.0000,8.0000) circle (1.2pt);
\fill[red!80!black] (12.0000,8.0000) circle (1.2pt);
\fill[red!80!black] (0.0000,8.0000) circle (1.2pt);
\fill[red!80!black] (12.0000,8.0000) circle (1.2pt);
\fill[red!80!black] (0.0000,8.0000) circle (1.2pt);
\fill[red!80!black] (12.0000,8.0000) circle (1.2pt);
\fill[red!80!black] (0.0000,8.0000) circle (1.2pt);
\fill[red!80!black] (12.0000,8.0000) circle (1.2pt);
\fill[red!80!black] (0.0000,8.0000) circle (1.2pt);
\fill[red!80!black] (12.0000,8.0000) circle (1.2pt);
\fill[red!80!black] (0.0000,8.0000) circle (1.2pt);
\fill[red!80!black] (12.0000,8.0000) circle (1.2pt);
\fill[red!80!black] (0.0000,8.0000) circle (1.2pt);
\fill[red!80!black] (12.0000,8.0000) circle (1.2pt);
\fill[red!80!black] (0.0000,8.0000) circle (1.2pt);
\fill[red!80!black] (12.0000,8.0000) circle (1.2pt);
\fill[red!80!black] (0.0000,8.0000) circle (1.2pt);
\fill[red!80!black] (12.0000,8.0000) circle (1.2pt);
\fill[red!80!black] (0.0000,8.0000) circle (1.2pt);
\fill[red!80!black] (12.0000,8.0000) circle (1.2pt);
\fill[red!80!black] (0.0000,8.0000) circle (1.2pt);
\fill[red!80!black] (12.0000,8.0000) circle (1.2pt);
\fill[red!80!black] (0.0000,8.0000) circle (1.2pt);
\fill[red!80!black] (12.0000,8.0000) circle (1.2pt);
\fill[red!80!black] (0.0000,8.0000) circle (1.2pt);
\fill[red!80!black] (12.0000,8.0000) circle (1.2pt);
\fill[red!80!black] (0.0000,8.0000) circle (1.2pt);
\fill[red!80!black] (12.0000,8.0000) circle (1.2pt);
\fill[red!80!black] (0.0000,8.0000) circle (1.2pt);
\fill[red!80!black] (12.0000,8.0000) circle (1.2pt);
\fill[red!80!black] (0.0000,8.0000) circle (1.2pt);
\fill[red!80!black] (12.0000,8.0000) circle (1.2pt);
\fill[red!80!black] (0.0000,8.0000) circle (1.2pt);
\fill[red!80!black] (12.0000,8.0000) circle (1.2pt);
\fill[red!80!black] (0.0000,8.0000) circle (1.2pt);
\fill[red!80!black] (12.0000,8.0000) circle (1.2pt);
\fill[red!80!black] (0.0000,8.0000) circle (1.2pt);
\fill[red!80!black] (12.0000,8.0000) circle (1.2pt);
\fill[red!80!black] (0.0000,8.0000) circle (1.2pt);
\fill[red!80!black] (12.0000,8.0000) circle (1.2pt);
\fill[red!80!black] (0.0000,8.0000) circle (1.2pt);
\fill[red!80!black] (12.0000,8.0000) circle (1.2pt);
\fill[red!80!black] (0.0000,8.0000) circle (1.2pt);
\fill[red!80!black] (12.0000,8.0000) circle (1.2pt);
\fill[red!80!black] (0.0000,8.0000) circle (1.2pt);
\fill[red!80!black] (12.0000,8.0000) circle (1.2pt);
\fill[red!80!black] (0.0000,8.0000) circle (1.2pt);
\fill[red!80!black] (12.0000,8.0000) circle (1.2pt);
\fill[red!80!black] (0.0000,8.0000) circle (1.2pt);
\fill[red!80!black] (12.0000,8.0000) circle (1.2pt);
\fill[red!80!black] (0.0000,8.0000) circle (1.2pt);
\fill[red!80!black] (12.0000,8.0000) circle (1.2pt);
\fill[red!80!black] (0.0000,8.0000) circle (1.2pt);
\fill[red!80!black] (12.0000,8.0000) circle (1.2pt);
\fill[red!80!black] (0.0000,8.0000) circle (1.2pt);
\fill[red!80!black] (12.0000,8.0000) circle (1.2pt);
\fill[red!80!black] (0.0000,8.0000) circle (1.2pt);
\fill[red!80!black] (12.0000,8.0000) circle (1.2pt);
\fill[red!80!black] (0.0000,8.0000) circle (1.2pt);
\fill[red!80!black] (12.0000,8.0000) circle (1.2pt);
\fill[red!80!black] (0.0000,8.0000) circle (1.2pt);
\fill[red!80!black] (12.0000,8.0000) circle (1.2pt);
\fill[red!80!black] (0.0000,8.0000) circle (1.2pt);
\fill[red!80!black] (12.0000,8.0000) circle (1.2pt);
\fill[red!80!black] (0.0000,8.0000) circle (1.2pt);
\fill[red!80!black] (12.0000,8.0000) circle (1.2pt);
\fill[red!80!black] (0.0000,8.0000) circle (1.2pt);
\fill[red!80!black] (12.0000,8.0000) circle (1.2pt);
\fill[red!80!black] (0.0000,8.0000) circle (1.2pt);
\fill[red!80!black] (12.0000,8.0000) circle (1.2pt);
\fill[red!80!black] (0.0000,8.0000) circle (1.2pt);
\fill[red!80!black] (12.0000,8.0000) circle (1.2pt);
\fill[red!80!black] (0.0000,8.0000) circle (1.2pt);
\fill[red!80!black] (12.0000,8.0000) circle (1.2pt);
\fill[red!80!black] (0.0000,8.0000) circle (1.2pt);
\fill[red!80!black] (12.0000,8.0000) circle (1.2pt);
\fill[red!80!black] (0.0000,8.0000) circle (1.2pt);
\fill[red!80!black] (12.0000,8.0000) circle (1.2pt);
\fill[red!80!black] (0.0000,8.0000) circle (1.2pt);
\fill[red!80!black] (12.0000,8.0000) circle (1.2pt);
\fill[red!80!black] (0.0000,8.0000) circle (1.2pt);
\fill[red!80!black] (12.0000,8.0000) circle (1.2pt);
\fill[red!80!black] (0.0000,8.0000) circle (1.2pt);
\fill[red!80!black] (12.0000,8.0000) circle (1.2pt);
\fill[red!80!black] (0.0000,8.0000) circle (1.2pt);
\fill[red!80!black] (12.0000,8.0000) circle (1.2pt);
\fill[red!80!black] (0.0000,8.0000) circle (1.2pt);
\fill[red!80!black] (12.0000,8.0000) circle (1.2pt);
\fill[red!80!black] (0.0000,8.0000) circle (1.2pt);
\fill[red!80!black] (12.0000,8.0000) circle (1.2pt);
\fill[red!80!black] (0.0000,8.0000) circle (1.2pt);
\fill[red!80!black] (12.0000,8.0000) circle (1.2pt);
\fill[red!80!black] (0.0000,8.0000) circle (1.2pt);
\fill[red!80!black] (12.0000,8.0000) circle (1.2pt);
\fill[red!80!black] (0.0000,8.0000) circle (1.2pt);
\fill[red!80!black] (12.0000,8.0000) circle (1.2pt);
\fill[red!80!black] (0.0000,8.0000) circle (1.2pt);
\fill[red!80!black] (12.0000,8.0000) circle (1.2pt);
\fill[red!80!black] (0.0000,8.0000) circle (1.2pt);
\fill[red!80!black] (12.0000,8.0000) circle (1.2pt);
\fill[red!80!black] (0.0000,8.0000) circle (1.2pt);
\fill[red!80!black] (12.0000,8.0000) circle (1.2pt);
\fill[red!80!black] (0.0000,8.0000) circle (1.2pt);
\fill[red!80!black] (12.0000,8.0000) circle (1.2pt);
\fill[red!80!black] (0.0000,8.0000) circle (1.2pt);
\fill[red!80!black] (12.0000,8.0000) circle (1.2pt);
\fill[red!80!black] (0.0000,8.0000) circle (1.2pt);
\fill[red!80!black] (12.0000,8.0000) circle (1.2pt);
\fill[red!80!black] (0.0000,8.0000) circle (1.2pt);
\fill[red!80!black] (12.0000,8.0000) circle (1.2pt);
\fill[red!80!black] (0.0000,8.0000) circle (1.2pt);
\fill[red!80!black] (12.0000,8.0000) circle (1.2pt);
\fill[red!80!black] (0.0000,8.0000) circle (1.2pt);
\fill[red!80!black] (12.0000,8.0000) circle (1.2pt);
\fill[red!80!black] (0.0000,8.0000) circle (1.2pt);
\fill[red!80!black] (12.0000,8.0000) circle (1.2pt);
\fill[red!80!black] (0.0000,8.0000) circle (1.2pt);
\fill[red!80!black] (12.0000,8.0000) circle (1.2pt);
\fill[red!80!black] (0.0000,8.0000) circle (1.2pt);
\fill[red!80!black] (12.0000,8.0000) circle (1.2pt);
\fill[red!80!black] (0.0000,8.0000) circle (1.2pt);
\fill[red!80!black] (12.0000,8.0000) circle (1.2pt);
\fill[red!80!black] (0.0000,8.0000) circle (1.2pt);
\fill[red!80!black] (12.0000,8.0000) circle (1.2pt);
\fill[red!80!black] (0.0000,8.0000) circle (1.2pt);
\fill[red!80!black] (12.0000,8.0000) circle (1.2pt);
\fill[red!80!black] (0.0000,8.0000) circle (1.2pt);
\fill[red!80!black] (12.0000,8.0000) circle (1.2pt);
\fill[red!80!black] (0.0000,8.0000) circle (1.2pt);
\fill[red!80!black] (12.0000,8.0000) circle (1.2pt);
\fill[red!80!black] (0.0000,8.0000) circle (1.2pt);
\fill[red!80!black] (12.0000,8.0000) circle (1.2pt);
\fill[red!80!black] (0.0000,8.0000) circle (1.2pt);
\fill[red!80!black] (12.0000,8.0000) circle (1.2pt);
\fill[red!80!black] (0.0000,8.0000) circle (1.2pt);
\fill[red!80!black] (12.0000,8.0000) circle (1.2pt);
\fill[red!80!black] (0.0000,8.0000) circle (1.2pt);
\fill[red!80!black] (12.0000,8.0000) circle (1.2pt);
\fill[red!80!black] (0.0000,8.0000) circle (1.2pt);
\fill[red!80!black] (12.0000,8.0000) circle (1.2pt);
\fill[red!80!black] (0.0000,8.0000) circle (1.2pt);
\fill[red!80!black] (12.0000,8.0000) circle (1.2pt);
\fill[red!80!black] (0.0000,8.0000) circle (1.2pt);
\fill[red!80!black] (12.0000,8.0000) circle (1.2pt);
\fill[red!80!black] (0.0000,8.0000) circle (1.2pt);
\fill[red!80!black] (12.0000,8.0000) circle (1.2pt);
\fill[red!80!black] (0.0000,8.0000) circle (1.2pt);
\fill[red!80!black] (12.0000,8.0000) circle (1.2pt);
\fill[red!80!black] (0.0000,8.0000) circle (1.2pt);
\fill[red!80!black] (12.0000,8.0000) circle (1.2pt);
\fill[red!80!black] (0.0000,8.0000) circle (1.2pt);
\fill[red!80!black] (12.0000,8.0000) circle (1.2pt);
\fill[red!80!black] (0.0000,8.0000) circle (1.2pt);
\fill[red!80!black] (12.0000,8.0000) circle (1.2pt);
\fill[red!80!black] (0.0000,8.0000) circle (1.2pt);
\fill[red!80!black] (12.0000,8.0000) circle (1.2pt);
\fill[red!80!black] (0.0000,8.0000) circle (1.2pt);
\fill[red!80!black] (12.0000,8.0000) circle (1.2pt);
\fill[red!80!black] (0.0000,8.0000) circle (1.2pt);
\fill[red!80!black] (12.0000,8.0000) circle (1.2pt);
\fill[red!80!black] (0.0000,8.0000) circle (1.2pt);
\fill[red!80!black] (12.0000,8.0000) circle (1.2pt);
\fill[red!80!black] (0.0000,8.0000) circle (1.2pt);
\fill[red!80!black] (12.0000,8.0000) circle (1.2pt);
\fill[red!80!black] (0.0000,8.0000) circle (1.2pt);
\fill[red!80!black] (12.0000,8.0000) circle (1.2pt);
\fill[red!80!black] (0.0000,8.0000) circle (1.2pt);
\fill[red!80!black] (12.0000,8.0000) circle (1.2pt);
\fill[red!80!black] (0.0000,8.0000) circle (1.2pt);
\fill[red!80!black] (12.0000,8.0000) circle (1.2pt);
\fill[red!80!black] (0.0000,8.0000) circle (1.2pt);
\fill[red!80!black] (12.0000,8.0000) circle (1.2pt);
\fill[red!80!black] (0.0000,8.0000) circle (1.2pt);
\fill[red!80!black] (12.0000,8.0000) circle (1.2pt);
\fill[red!80!black] (0.0000,8.0000) circle (1.2pt);
\fill[red!80!black] (12.0000,8.0000) circle (1.2pt);
\fill[red!80!black] (0.0000,8.0000) circle (1.2pt);
\fill[red!80!black] (12.0000,8.0000) circle (1.2pt);
\fill[red!80!black] (0.0000,8.0000) circle (1.2pt);
\fill[red!80!black] (12.0000,8.0000) circle (1.2pt);
\fill[red!80!black] (0.0000,8.0000) circle (1.2pt);
\fill[red!80!black] (12.0000,8.0000) circle (1.2pt);
\fill[red!80!black] (0.0000,8.0000) circle (1.2pt);
\fill[red!80!black] (12.0000,8.0000) circle (1.2pt);
\fill[red!80!black] (0.0000,8.0000) circle (1.2pt);
\fill[red!80!black] (12.0000,8.0000) circle (1.2pt);
\fill[red!80!black] (0.0000,8.0000) circle (1.2pt);
\fill[red!80!black] (12.0000,8.0000) circle (1.2pt);
\fill[red!80!black] (0.0000,8.0000) circle (1.2pt);
\fill[red!80!black] (12.0000,8.0000) circle (1.2pt);
\fill[red!80!black] (0.0000,8.0000) circle (1.2pt);
\fill[red!80!black] (12.0000,8.0000) circle (1.2pt);
\fill[red!80!black] (0.0000,8.0000) circle (1.2pt);
\fill[red!80!black] (12.0000,8.0000) circle (1.2pt);
\fill[red!80!black] (0.0000,8.0000) circle (1.2pt);
\fill[red!80!black] (12.0000,8.0000) circle (1.2pt);
\fill[red!80!black] (0.0000,8.0000) circle (1.2pt);
\fill[red!80!black] (12.0000,8.0000) circle (1.2pt);
\fill[red!80!black] (0.0000,8.0000) circle (1.2pt);
\fill[red!80!black] (12.0000,8.0000) circle (1.2pt);
\fill[red!80!black] (0.0000,8.0000) circle (1.2pt);
\fill[red!80!black] (12.0000,8.0000) circle (1.2pt);
\fill[red!80!black] (0.0000,8.0000) circle (1.2pt);
\fill[red!80!black] (12.0000,8.0000) circle (1.2pt);
\fill[red!80!black] (0.0000,8.0000) circle (1.2pt);
\fill[red!80!black] (12.0000,8.0000) circle (1.2pt);
\fill[red!80!black] (0.0000,8.0000) circle (1.2pt);
\fill[red!80!black] (12.0000,8.0000) circle (1.2pt);
\fill[red!80!black] (0.0000,8.0000) circle (1.2pt);
\fill[red!80!black] (12.0000,8.0000) circle (1.2pt);
\fill[red!80!black] (0.0000,8.0000) circle (1.2pt);
\fill[red!80!black] (12.0000,8.0000) circle (1.2pt);
\fill[red!80!black] (0.0000,8.0000) circle (1.2pt);
\fill[red!80!black] (12.0000,8.0000) circle (1.2pt);
\fill[red!80!black] (0.0000,8.0000) circle (1.2pt);
\fill[red!80!black] (12.0000,8.0000) circle (1.2pt);
\fill[red!80!black] (0.0000,8.0000) circle (1.2pt);
\fill[red!80!black] (12.0000,8.0000) circle (1.2pt);
\fill[red!80!black] (0.0000,8.0000) circle (1.2pt);
\fill[red!80!black] (12.0000,8.0000) circle (1.2pt);
\fill[red!80!black] (0.0000,8.0000) circle (1.2pt);
\fill[red!80!black] (12.0000,8.0000) circle (1.2pt);
\fill[red!80!black] (0.0000,8.0000) circle (1.2pt);
\fill[red!80!black] (12.0000,8.0000) circle (1.2pt);
\fill[red!80!black] (0.0000,8.0000) circle (1.2pt);
\fill[red!80!black] (12.0000,8.0000) circle (1.2pt);
\fill[red!80!black] (0.0000,8.0000) circle (1.2pt);
\fill[red!80!black] (12.0000,8.0000) circle (1.2pt);
\fill[red!80!black] (0.0000,8.0000) circle (1.2pt);
\fill[red!80!black] (12.0000,8.0000) circle (1.2pt);
\fill[red!80!black] (0.0000,8.0000) circle (1.2pt);
\fill[red!80!black] (12.0000,8.0000) circle (1.2pt);
\fill[red!80!black] (0.0000,8.0000) circle (1.2pt);
\fill[red!80!black] (12.0000,8.0000) circle (1.2pt);
\fill[red!80!black] (0.0000,8.0000) circle (1.2pt);
\fill[red!80!black] (12.0000,8.0000) circle (1.2pt);
\fill[red!80!black] (0.0000,8.0000) circle (1.2pt);
\fill[red!80!black] (12.0000,8.0000) circle (1.2pt);
\fill[red!80!black] (0.0000,8.0000) circle (1.2pt);
\fill[red!80!black] (12.0000,8.0000) circle (1.2pt);
\fill[red!80!black] (0.0000,8.0000) circle (1.2pt);
\fill[red!80!black] (12.0000,8.0000) circle (1.2pt);
\fill[red!80!black] (0.0000,8.0000) circle (1.2pt);
\fill[red!80!black] (12.0000,8.0000) circle (1.2pt);
\fill[red!80!black] (0.0000,8.0000) circle (1.2pt);
\fill[red!80!black] (12.0000,8.0000) circle (1.2pt);
\fill[red!80!black] (0.0000,8.0000) circle (1.2pt);
\fill[red!80!black] (12.0000,8.0000) circle (1.2pt);
\fill[red!80!black] (0.0000,8.0000) circle (1.2pt);
\fill[red!80!black] (12.0000,8.0000) circle (1.2pt);
\fill[red!80!black] (0.0000,8.0000) circle (1.2pt);
\fill[red!80!black] (12.0000,8.0000) circle (1.2pt);
\fill[red!80!black] (0.0000,8.0000) circle (1.2pt);
\fill[red!80!black] (12.0000,8.0000) circle (1.2pt);
\fill[red!80!black] (0.0000,8.0000) circle (1.2pt);
\fill[red!80!black] (12.0000,8.0000) circle (1.2pt);
\fill[red!80!black] (0.0000,8.0000) circle (1.2pt);
\fill[red!80!black] (12.0000,8.0000) circle (1.2pt);
\fill[red!80!black] (0.0000,8.0000) circle (1.2pt);
\fill[red!80!black] (12.0000,8.0000) circle (1.2pt);
\fill[red!80!black] (0.0000,8.0000) circle (1.2pt);
\fill[red!80!black] (12.0000,8.0000) circle (1.2pt);
\fill[red!80!black] (0.0000,8.0000) circle (1.2pt);
\fill[red!80!black] (12.0000,8.0000) circle (1.2pt);
\fill[red!80!black] (0.0000,8.0000) circle (1.2pt);
\fill[red!80!black] (12.0000,8.0000) circle (1.2pt);
\fill[red!80!black] (0.0000,8.0000) circle (1.2pt);
\fill[red!80!black] (12.0000,8.0000) circle (1.2pt);
\fill[red!80!black] (0.0000,8.0000) circle (1.2pt);
\fill[red!80!black] (12.0000,8.0000) circle (1.2pt);
\fill[red!80!black] (0.0000,8.0000) circle (1.2pt);
\fill[red!80!black] (12.0000,8.0000) circle (1.2pt);
\fill[red!80!black] (0.0000,8.0000) circle (1.2pt);
\fill[red!80!black] (12.0000,8.0000) circle (1.2pt);
\fill[red!80!black] (0.0000,8.0000) circle (1.2pt);
\fill[red!80!black] (12.0000,8.0000) circle (1.2pt);
\fill[red!80!black] (0.0000,8.0000) circle (1.2pt);
\fill[red!80!black] (12.0000,8.0000) circle (1.2pt);
\fill[red!80!black] (0.0000,8.0000) circle (1.2pt);
\fill[red!80!black] (12.0000,8.0000) circle (1.2pt);
\fill[red!80!black] (0.0000,8.0000) circle (1.2pt);
\fill[red!80!black] (12.0000,8.0000) circle (1.2pt);
\fill[red!80!black] (0.0000,8.0000) circle (1.2pt);
\fill[red!80!black] (12.0000,8.0000) circle (1.2pt);
\fill[red!80!black] (0.0000,8.0000) circle (1.2pt);
\fill[red!80!black] (12.0000,8.0000) circle (1.2pt);
\fill[red!80!black] (0.0000,8.0000) circle (1.2pt);
\fill[red!80!black] (12.0000,8.0000) circle (1.2pt);
\fill[red!80!black] (0.0000,8.0000) circle (1.2pt);
\fill[red!80!black] (12.0000,8.0000) circle (1.2pt);
\fill[red!80!black] (0.0000,8.0000) circle (1.2pt);
\fill[red!80!black] (12.0000,8.0000) circle (1.2pt);
\fill[red!80!black] (0.0000,8.0000) circle (1.2pt);
\fill[red!80!black] (12.0000,8.0000) circle (1.2pt);
\fill[red!80!black] (0.0000,8.0000) circle (1.2pt);
\fill[red!80!black] (12.0000,8.0000) circle (1.2pt);
\fill[red!80!black] (0.0000,8.0000) circle (1.2pt);
\fill[red!80!black] (12.0000,8.0000) circle (1.2pt);
\fill[red!80!black] (0.0000,8.0000) circle (1.2pt);
\fill[red!80!black] (12.0000,8.0000) circle (1.2pt);
\fill[red!80!black] (0.0000,8.0000) circle (1.2pt);
\fill[red!80!black] (12.0000,8.0000) circle (1.2pt);
\fill[red!80!black] (0.0000,8.0000) circle (1.2pt);
\fill[red!80!black] (12.0000,8.0000) circle (1.2pt);
\fill[red!80!black] (0.0000,8.0000) circle (1.2pt);
\fill[red!80!black] (12.0000,8.0000) circle (1.2pt);
\fill[red!80!black] (0.0000,8.0000) circle (1.2pt);
\fill[red!80!black] (12.0000,8.0000) circle (1.2pt);
\fill[red!80!black] (0.0000,8.0000) circle (1.2pt);
\fill[red!80!black] (12.0000,8.0000) circle (1.2pt);
\fill[red!80!black] (0.0000,8.0000) circle (1.2pt);
\fill[red!80!black] (12.0000,8.0000) circle (1.2pt);
\fill[red!80!black] (0.0000,8.0000) circle (1.2pt);
\fill[red!80!black] (12.0000,8.0000) circle (1.2pt);
\fill[red!80!black] (0.0000,8.0000) circle (1.2pt);
\fill[red!80!black] (12.0000,8.0000) circle (1.2pt);
\fill[red!80!black] (0.0000,8.0000) circle (1.2pt);
\fill[red!80!black] (12.0000,8.0000) circle (1.2pt);
\fill[red!80!black] (0.0000,8.0000) circle (1.2pt);
\fill[red!80!black] (12.0000,8.0000) circle (1.2pt);
\fill[red!80!black] (0.0000,8.0000) circle (1.2pt);
\fill[red!80!black] (12.0000,8.0000) circle (1.2pt);
\fill[red!80!black] (0.0000,8.0000) circle (1.2pt);
\fill[red!80!black] (12.0000,8.0000) circle (1.2pt);
\fill[red!80!black] (0.0000,8.0000) circle (1.2pt);
\fill[red!80!black] (12.0000,8.0000) circle (1.2pt);
\fill[red!80!black] (0.0000,8.0000) circle (1.2pt);
\fill[red!80!black] (12.0000,8.0000) circle (1.2pt);
\fill[red!80!black] (0.0000,8.0000) circle (1.2pt);
\fill[red!80!black] (12.0000,8.0000) circle (1.2pt);
\fill[red!80!black] (0.0000,8.0000) circle (1.2pt);
\fill[red!80!black] (12.0000,8.0000) circle (1.2pt);
\fill[red!80!black] (0.0000,8.0000) circle (1.2pt);
\fill[red!80!black] (12.0000,8.0000) circle (1.2pt);
\fill[red!80!black] (0.0000,8.0000) circle (1.2pt);
\fill[red!80!black] (12.0000,8.0000) circle (1.2pt);
\fill[red!80!black] (0.0000,8.0000) circle (1.2pt);
\fill[red!80!black] (12.0000,8.0000) circle (1.2pt);
\fill[red!80!black] (0.0000,8.0000) circle (1.2pt);
\fill[red!80!black] (12.0000,8.0000) circle (1.2pt);
\fill[red!80!black] (0.0000,8.0000) circle (1.2pt);
\fill[red!80!black] (12.0000,8.0000) circle (1.2pt);
\fill[red!80!black] (0.0000,8.0000) circle (1.2pt);
\fill[red!80!black] (12.0000,8.0000) circle (1.2pt);
\fill[red!80!black] (0.0000,8.0000) circle (1.2pt);
\fill[red!80!black] (12.0000,8.0000) circle (1.2pt);
\fill[red!80!black] (0.0000,8.0000) circle (1.2pt);
\fill[red!80!black] (12.0000,8.0000) circle (1.2pt);
\fill[red!80!black] (0.0000,8.0000) circle (1.2pt);
\fill[red!80!black] (12.0000,8.0000) circle (1.2pt);
\fill[red!80!black] (0.0000,8.0000) circle (1.2pt);
\fill[red!80!black] (12.0000,8.0000) circle (1.2pt);
\fill[red!80!black] (0.0000,8.0000) circle (1.2pt);
\fill[red!80!black] (12.0000,8.0000) circle (1.2pt);
\fill[red!80!black] (0.0000,8.0000) circle (1.2pt);
\fill[red!80!black] (12.0000,8.0000) circle (1.2pt);
\fill[red!80!black] (0.0000,8.0000) circle (1.2pt);
\fill[red!80!black] (12.0000,8.0000) circle (1.2pt);
\fill[red!80!black] (0.0000,8.0000) circle (1.2pt);
\fill[red!80!black] (12.0000,8.0000) circle (1.2pt);
\fill[red!80!black] (0.0000,8.0000) circle (1.2pt);
\fill[red!80!black] (12.0000,8.0000) circle (1.2pt);
\fill[red!80!black] (0.0000,8.0000) circle (1.2pt);
\fill[red!80!black] (12.0000,8.0000) circle (1.2pt);
\fill[red!80!black] (0.0000,8.0000) circle (1.2pt);
\fill[red!80!black] (12.0000,8.0000) circle (1.2pt);
\fill[red!80!black] (0.0000,8.0000) circle (1.2pt);
\fill[red!80!black] (12.0000,8.0000) circle (1.2pt);
\fill[red!80!black] (0.0000,8.0000) circle (1.2pt);
\fill[red!80!black] (12.0000,8.0000) circle (1.2pt);
\fill[red!80!black] (0.0000,8.0000) circle (1.2pt);
\fill[red!80!black] (12.0000,8.0000) circle (1.2pt);
\fill[red!80!black] (0.0000,8.0000) circle (1.2pt);
\fill[red!80!black] (12.0000,8.0000) circle (1.2pt);
\fill[red!80!black] (0.0000,8.0000) circle (1.2pt);
\fill[red!80!black] (12.0000,8.0000) circle (1.2pt);
\fill[red!80!black] (0.0000,8.0000) circle (1.2pt);
\fill[red!80!black] (12.0000,8.0000) circle (1.2pt);
\fill[red!80!black] (0.0000,8.0000) circle (1.2pt);
\fill[red!80!black] (12.0000,8.0000) circle (1.2pt);
\fill[red!80!black] (0.0000,8.0000) circle (1.2pt);
\fill[red!80!black] (12.0000,8.0000) circle (1.2pt);
\fill[red!80!black] (0.0000,8.0000) circle (1.2pt);
\fill[red!80!black] (12.0000,8.0000) circle (1.2pt);
\fill[red!80!black] (0.0000,8.0000) circle (1.2pt);
\fill[red!80!black] (12.0000,8.0000) circle (1.2pt);
\fill[red!80!black] (0.0000,8.0000) circle (1.2pt);
\fill[red!80!black] (12.0000,8.0000) circle (1.2pt);
\fill[red!80!black] (0.0000,8.0000) circle (1.2pt);
\fill[red!80!black] (12.0000,8.0000) circle (1.2pt);
\fill[red!80!black] (0.0000,8.0000) circle (1.2pt);
\fill[red!80!black] (12.0000,8.0000) circle (1.2pt);
\fill[red!80!black] (0.0000,8.0000) circle (1.2pt);
\fill[red!80!black] (12.0000,8.0000) circle (1.2pt);
\fill[red!80!black] (0.0000,8.0000) circle (1.2pt);
\fill[red!80!black] (12.0000,8.0000) circle (1.2pt);
\fill[red!80!black] (0.0000,8.0000) circle (1.2pt);
\fill[red!80!black] (12.0000,8.0000) circle (1.2pt);
\fill[red!80!black] (0.0000,8.0000) circle (1.2pt);
\fill[red!80!black] (12.0000,8.0000) circle (1.2pt);
\fill[red!80!black] (0.0000,8.0000) circle (1.2pt);
\fill[red!80!black] (12.0000,8.0000) circle (1.2pt);
\fill[red!80!black] (0.0000,8.0000) circle (1.2pt);
\fill[red!80!black] (12.0000,8.0000) circle (1.2pt);
\fill[red!80!black] (0.0000,8.0000) circle (1.2pt);
\fill[red!80!black] (12.0000,8.0000) circle (1.2pt);
\fill[red!80!black] (0.0000,8.0000) circle (1.2pt);
\fill[red!80!black] (12.0000,8.0000) circle (1.2pt);
\fill[red!80!black] (0.0000,8.0000) circle (1.2pt);
\fill[red!80!black] (12.0000,8.0000) circle (1.2pt);
\fill[red!80!black] (0.0000,8.0000) circle (1.2pt);
\fill[red!80!black] (12.0000,8.0000) circle (1.2pt);
\fill[red!80!black] (0.0000,8.0000) circle (1.2pt);
\fill[red!80!black] (12.0000,8.0000) circle (1.2pt);
\fill[red!80!black] (0.0000,8.0000) circle (1.2pt);
\fill[red!80!black] (12.0000,8.0000) circle (1.2pt);
\fill[red!80!black] (0.0000,8.0000) circle (1.2pt);
\fill[red!80!black] (12.0000,8.0000) circle (1.2pt);
\fill[red!80!black] (0.0000,8.0000) circle (1.2pt);
\fill[red!80!black] (12.0000,8.0000) circle (1.2pt);
\fill[red!80!black] (0.0000,8.0000) circle (1.2pt);
\fill[red!80!black] (12.0000,8.0000) circle (1.2pt);
\fill[red!80!black] (0.0000,8.0000) circle (1.2pt);
\fill[red!80!black] (12.0000,8.0000) circle (1.2pt);
\fill[red!80!black] (0.0000,8.0000) circle (1.2pt);
\fill[red!80!black] (12.0000,8.0000) circle (1.2pt);
\fill[red!80!black] (0.0000,8.0000) circle (1.2pt);
\fill[red!80!black] (12.0000,8.0000) circle (1.2pt);
\fill[red!80!black] (0.0000,8.0000) circle (1.2pt);
\fill[red!80!black] (12.0000,8.0000) circle (1.2pt);
\fill[red!80!black] (0.0000,8.0000) circle (1.2pt);
\fill[red!80!black] (12.0000,8.0000) circle (1.2pt);
\fill[red!80!black] (0.0000,8.0000) circle (1.2pt);
\fill[red!80!black] (12.0000,8.0000) circle (1.2pt);
\fill[red!80!black] (0.0000,8.0000) circle (1.2pt);
\fill[red!80!black] (12.0000,8.0000) circle (1.2pt);
\fill[red!80!black] (0.0000,8.0000) circle (1.2pt);
\fill[red!80!black] (12.0000,8.0000) circle (1.2pt);
\fill[red!80!black] (0.0000,8.0000) circle (1.2pt);
\fill[red!80!black] (12.0000,8.0000) circle (1.2pt);
\fill[red!80!black] (0.0000,8.0000) circle (1.2pt);
\fill[red!80!black] (12.0000,8.0000) circle (1.2pt);
\fill[red!80!black] (0.0000,8.0000) circle (1.2pt);
\fill[red!80!black] (12.0000,8.0000) circle (1.2pt);
\fill[red!80!black] (0.0000,8.0000) circle (1.2pt);
\fill[red!80!black] (12.0000,8.0000) circle (1.2pt);
\fill[red!80!black] (0.0000,8.0000) circle (1.2pt);
\fill[red!80!black] (12.0000,8.0000) circle (1.2pt);
\fill[red!80!black] (0.0000,8.0000) circle (1.2pt);
\fill[red!80!black] (12.0000,8.0000) circle (1.2pt);
\fill[red!80!black] (0.0000,8.0000) circle (1.2pt);
\fill[red!80!black] (12.0000,8.0000) circle (1.2pt);
\fill[red!80!black] (0.0000,8.0000) circle (1.2pt);
\fill[red!80!black] (12.0000,8.0000) circle (1.2pt);
\fill[red!80!black] (0.0000,8.0000) circle (1.2pt);
\fill[red!80!black] (12.0000,8.0000) circle (1.2pt);
\fill[red!80!black] (0.0000,8.0000) circle (1.2pt);
\fill[red!80!black] (12.0000,8.0000) circle (1.2pt);
\fill[red!80!black] (0.0000,8.0000) circle (1.2pt);
\fill[red!80!black] (12.0000,8.0000) circle (1.2pt);
\fill[red!80!black] (0.0000,8.0000) circle (1.2pt);
\fill[red!80!black] (12.0000,8.0000) circle (1.2pt);
\fill[red!80!black] (0.0000,8.0000) circle (1.2pt);
\fill[red!80!black] (12.0000,8.0000) circle (1.2pt);
\fill[red!80!black] (0.0000,8.0000) circle (1.2pt);
\fill[red!80!black] (12.0000,8.0000) circle (1.2pt);
\fill[red!80!black] (0.0000,8.0000) circle (1.2pt);
\fill[red!80!black] (12.0000,8.0000) circle (1.2pt);
\fill[red!80!black] (0.0000,8.0000) circle (1.2pt);
\fill[red!80!black] (12.0000,8.0000) circle (1.2pt);
\fill[red!80!black] (0.0000,8.0000) circle (1.2pt);
\fill[red!80!black] (12.0000,8.0000) circle (1.2pt);
\fill[red!80!black] (0.0000,8.0000) circle (1.2pt);
\fill[red!80!black] (12.0000,8.0000) circle (1.2pt);
\fill[red!80!black] (0.0000,8.0000) circle (1.2pt);
\fill[red!80!black] (12.0000,8.0000) circle (1.2pt);
\fill[red!80!black] (0.0000,8.0000) circle (1.2pt);
\fill[red!80!black] (12.0000,8.0000) circle (1.2pt);
\fill[red!80!black] (0.0000,8.0000) circle (1.2pt);
\fill[red!80!black] (12.0000,8.0000) circle (1.2pt);
\fill[red!80!black] (0.0000,8.0000) circle (1.2pt);
\fill[red!80!black] (12.0000,8.0000) circle (1.2pt);
\fill[red!80!black] (0.0000,8.0000) circle (1.2pt);
\fill[red!80!black] (12.0000,8.0000) circle (1.2pt);
\fill[red!80!black] (0.0000,8.0000) circle (1.2pt);
\fill[red!80!black] (12.0000,8.0000) circle (1.2pt);
\fill[red!80!black] (0.0000,8.0000) circle (1.2pt);
\fill[red!80!black] (12.0000,8.0000) circle (1.2pt);
\fill[red!80!black] (0.0000,8.0000) circle (1.2pt);
\fill[red!80!black] (12.0000,8.0000) circle (1.2pt);
\fill[red!80!black] (0.0000,8.0000) circle (1.2pt);
\fill[red!80!black] (12.0000,8.0000) circle (1.2pt);
\fill[red!80!black] (0.0000,8.0000) circle (1.2pt);
\fill[red!80!black] (12.0000,8.0000) circle (1.2pt);
\fill[red!80!black] (0.0000,8.0000) circle (1.2pt);
\fill[red!80!black] (12.0000,8.0000) circle (1.2pt);
\fill[red!80!black] (0.0000,8.0000) circle (1.2pt);
\fill[red!80!black] (12.0000,8.0000) circle (1.2pt);
\fill[red!80!black] (0.0000,8.0000) circle (1.2pt);
\fill[red!80!black] (12.0000,8.0000) circle (1.2pt);
\fill[red!80!black] (0.0000,8.0000) circle (1.2pt);
\fill[red!80!black] (12.0000,8.0000) circle (1.2pt);
\fill[red!80!black] (0.0000,8.0000) circle (1.2pt);
\fill[red!80!black] (12.0000,8.0000) circle (1.2pt);
\fill[red!80!black] (0.0000,8.0000) circle (1.2pt);
\fill[red!80!black] (12.0000,8.0000) circle (1.2pt);
\fill[red!80!black] (0.0000,8.0000) circle (1.2pt);
\fill[red!80!black] (12.0000,8.0000) circle (1.2pt);
\fill[red!80!black] (0.0000,8.0000) circle (1.2pt);
\fill[red!80!black] (12.0000,8.0000) circle (1.2pt);
\fill[red!80!black] (0.0000,8.0000) circle (1.2pt);
\fill[red!80!black] (12.0000,8.0000) circle (1.2pt);
\fill[red!80!black] (0.0000,8.0000) circle (1.2pt);
\fill[red!80!black] (12.0000,8.0000) circle (1.2pt);
\fill[red!80!black] (0.0000,8.0000) circle (1.2pt);
\fill[red!80!black] (12.0000,8.0000) circle (1.2pt);
\fill[red!80!black] (0.0000,8.0000) circle (1.2pt);
\fill[red!80!black] (12.0000,8.0000) circle (1.2pt);
\fill[red!80!black] (0.0000,8.0000) circle (1.2pt);
\fill[red!80!black] (12.0000,8.0000) circle (1.2pt);
\fill[red!80!black] (0.0000,8.0000) circle (1.2pt);
\fill[red!80!black] (12.0000,8.0000) circle (1.2pt);
\fill[red!80!black] (0.0000,8.0000) circle (1.2pt);
\fill[red!80!black] (12.0000,8.0000) circle (1.2pt);
\fill[red!80!black] (0.0000,8.0000) circle (1.2pt);
\fill[red!80!black] (12.0000,8.0000) circle (1.2pt);
\fill[red!80!black] (0.0000,8.0000) circle (1.2pt);
\fill[red!80!black] (12.0000,8.0000) circle (1.2pt);
\fill[red!80!black] (0.0000,8.0000) circle (1.2pt);
\fill[red!80!black] (12.0000,8.0000) circle (1.2pt);
\fill[red!80!black] (0.0000,8.0000) circle (1.2pt);
\fill[red!80!black] (12.0000,8.0000) circle (1.2pt);
\fill[red!80!black] (0.0000,8.0000) circle (1.2pt);
\fill[red!80!black] (12.0000,8.0000) circle (1.2pt);
\fill[red!80!black] (0.0000,8.0000) circle (1.2pt);
\fill[red!80!black] (12.0000,8.0000) circle (1.2pt);
\fill[red!80!black] (0.0000,8.0000) circle (1.2pt);
\fill[red!80!black] (12.0000,8.0000) circle (1.2pt);
\fill[red!80!black] (0.0000,8.0000) circle (1.2pt);
\fill[red!80!black] (12.0000,8.0000) circle (1.2pt);
\fill[red!80!black] (0.0000,8.0000) circle (1.2pt);
\fill[red!80!black] (12.0000,8.0000) circle (1.2pt);
\fill[red!80!black] (0.0000,8.0000) circle (1.2pt);
\fill[red!80!black] (12.0000,8.0000) circle (1.2pt);
\fill[red!80!black] (0.0000,8.0000) circle (1.2pt);
\fill[red!80!black] (12.0000,8.0000) circle (1.2pt);
\fill[red!80!black] (0.0000,8.0000) circle (1.2pt);
\fill[red!80!black] (12.0000,8.0000) circle (1.2pt);
\fill[red!80!black] (0.0000,8.0000) circle (1.2pt);
\fill[red!80!black] (12.0000,8.0000) circle (1.2pt);
\fill[red!80!black] (0.0000,8.0000) circle (1.2pt);
\fill[red!80!black] (12.0000,8.0000) circle (1.2pt);
\fill[red!80!black] (0.0000,8.0000) circle (1.2pt);
\fill[red!80!black] (12.0000,8.0000) circle (1.2pt);
\fill[red!80!black] (0.0000,8.0000) circle (1.2pt);
\fill[red!80!black] (12.0000,8.0000) circle (1.2pt);
\fill[red!80!black] (0.0000,8.0000) circle (1.2pt);
\fill[red!80!black] (12.0000,8.0000) circle (1.2pt);
\fill[red!80!black] (0.0000,8.0000) circle (1.2pt);
\fill[red!80!black] (12.0000,8.0000) circle (1.2pt);
\fill[red!80!black] (0.0000,8.0000) circle (1.2pt);
\fill[red!80!black] (12.0000,8.0000) circle (1.2pt);
\fill[red!80!black] (0.0000,8.0000) circle (1.2pt);
\fill[red!80!black] (12.0000,8.0000) circle (1.2pt);
\fill[red!80!black] (0.0000,8.0000) circle (1.2pt);
\fill[red!80!black] (12.0000,8.0000) circle (1.2pt);
\fill[red!80!black] (0.0000,8.0000) circle (1.2pt);
\fill[red!80!black] (12.0000,8.0000) circle (1.2pt);
\fill[red!80!black] (0.0000,8.0000) circle (1.2pt);
\fill[red!80!black] (12.0000,8.0000) circle (1.2pt);
\fill[red!80!black] (0.0000,8.0000) circle (1.2pt);
\fill[red!80!black] (12.0000,8.0000) circle (1.2pt);
\fill[red!80!black] (0.0000,8.0000) circle (1.2pt);
\fill[red!80!black] (12.0000,8.0000) circle (1.2pt);
\fill[red!80!black] (0.0000,8.0000) circle (1.2pt);
\fill[red!80!black] (12.0000,8.0000) circle (1.2pt);
\fill[red!80!black] (0.0000,8.0000) circle (1.2pt);
\fill[red!80!black] (12.0000,8.0000) circle (1.2pt);
\fill[red!80!black] (0.0000,8.0000) circle (1.2pt);
\fill[red!80!black] (12.0000,8.0000) circle (1.2pt);
\fill[red!80!black] (0.0000,8.0000) circle (1.2pt);
\fill[red!80!black] (12.0000,8.0000) circle (1.2pt);
\fill[red!80!black] (0.0000,8.0000) circle (1.2pt);
\fill[red!80!black] (12.0000,8.0000) circle (1.2pt);
\fill[red!80!black] (0.0000,8.0000) circle (1.2pt);
\fill[red!80!black] (12.0000,8.0000) circle (1.2pt);
\fill[red!80!black] (0.0000,8.0000) circle (1.2pt);
\fill[red!80!black] (12.0000,8.0000) circle (1.2pt);
\fill[red!80!black] (0.0000,8.0000) circle (1.2pt);
\fill[red!80!black] (12.0000,8.0000) circle (1.2pt);
\fill[red!80!black] (0.0000,8.0000) circle (1.2pt);
\fill[red!80!black] (12.0000,8.0000) circle (1.2pt);
\fill[red!80!black] (0.0000,8.0000) circle (1.2pt);
\fill[red!80!black] (12.0000,8.0000) circle (1.2pt);
\fill[red!80!black] (0.0000,8.0000) circle (1.2pt);
\fill[red!80!black] (12.0000,8.0000) circle (1.2pt);
\fill[red!80!black] (0.0000,8.0000) circle (1.2pt);
\fill[red!80!black] (12.0000,8.0000) circle (1.2pt);
\fill[red!80!black] (0.0000,8.0000) circle (1.2pt);
\fill[red!80!black] (12.0000,8.0000) circle (1.2pt);
\fill[red!80!black] (0.0000,8.0000) circle (1.2pt);
\fill[red!80!black] (12.0000,8.0000) circle (1.2pt);
\fill[red!80!black] (0.0000,8.0000) circle (1.2pt);
\fill[red!80!black] (12.0000,8.0000) circle (1.2pt);
\fill[red!80!black] (0.0000,8.0000) circle (1.2pt);
\fill[red!80!black] (12.0000,8.0000) circle (1.2pt);
\fill[red!80!black] (0.0000,8.0000) circle (1.2pt);
\fill[red!80!black] (12.0000,8.0000) circle (1.2pt);
\fill[red!80!black] (0.0000,8.0000) circle (1.2pt);
\fill[red!80!black] (12.0000,8.0000) circle (1.2pt);
\fill[red!80!black] (0.0000,8.0000) circle (1.2pt);
\fill[red!80!black] (12.0000,8.0000) circle (1.2pt);
\fill[red!80!black] (0.0000,8.0000) circle (1.2pt);
\fill[red!80!black] (12.0000,8.0000) circle (1.2pt);
\fill[red!80!black] (0.0000,8.0000) circle (1.2pt);
\fill[red!80!black] (12.0000,8.0000) circle (1.2pt);
\fill[red!80!black] (0.0000,8.0000) circle (1.2pt);
\fill[red!80!black] (12.0000,8.0000) circle (1.2pt);
\fill[red!80!black] (0.0000,8.0000) circle (1.2pt);
\fill[red!80!black] (12.0000,8.0000) circle (1.2pt);
\fill[red!80!black] (0.0000,8.0000) circle (1.2pt);
\fill[red!80!black] (12.0000,8.0000) circle (1.2pt);
\fill[red!80!black] (0.0000,8.0000) circle (1.2pt);
\fill[red!80!black] (12.0000,8.0000) circle (1.2pt);
\fill[red!80!black] (0.0000,8.0000) circle (1.2pt);
\fill[red!80!black] (12.0000,8.0000) circle (1.2pt);
\fill[red!80!black] (0.0000,8.0000) circle (1.2pt);
\fill[red!80!black] (12.0000,8.0000) circle (1.2pt);
\fill[red!80!black] (0.0000,8.0000) circle (1.2pt);
\fill[red!80!black] (12.0000,8.0000) circle (1.2pt);
\fill[red!80!black] (0.0000,8.0000) circle (1.2pt);
\fill[red!80!black] (12.0000,8.0000) circle (1.2pt);
\fill[red!80!black] (0.0000,8.0000) circle (1.2pt);
\fill[red!80!black] (12.0000,8.0000) circle (1.2pt);
\fill[red!80!black] (0.0000,8.0000) circle (1.2pt);
\fill[red!80!black] (12.0000,8.0000) circle (1.2pt);
\fill[red!80!black] (0.0000,8.0000) circle (1.2pt);
\fill[red!80!black] (12.0000,8.0000) circle (1.2pt);
\fill[red!80!black] (0.0000,8.0000) circle (1.2pt);
\fill[red!80!black] (12.0000,8.0000) circle (1.2pt);
\fill[red!80!black] (0.0000,8.0000) circle (1.2pt);
\fill[red!80!black] (12.0000,8.0000) circle (1.2pt);
\fill[red!80!black] (0.0000,8.0000) circle (1.2pt);
\fill[red!80!black] (12.0000,8.0000) circle (1.2pt);
\fill[red!80!black] (0.0000,8.0000) circle (1.2pt);
\fill[red!80!black] (12.0000,8.0000) circle (1.2pt);
\fill[red!80!black] (0.0000,8.0000) circle (1.2pt);
\fill[red!80!black] (12.0000,8.0000) circle (1.2pt);
\fill[red!80!black] (0.0000,8.0000) circle (1.2pt);
\fill[red!80!black] (12.0000,8.0000) circle (1.2pt);
\fill[red!80!black] (0.0000,8.0000) circle (1.2pt);
\fill[red!80!black] (12.0000,8.0000) circle (1.2pt);
\fill[red!80!black] (0.0000,8.0000) circle (1.2pt);
\fill[red!80!black] (12.0000,8.0000) circle (1.2pt);
\fill[red!80!black] (0.0000,8.0000) circle (1.2pt);
\fill[red!80!black] (12.0000,8.0000) circle (1.2pt);
\fill[red!80!black] (0.0000,8.0000) circle (1.2pt);
\fill[red!80!black] (12.0000,8.0000) circle (1.2pt);
\fill[red!80!black] (0.0000,8.0000) circle (1.2pt);
\fill[red!80!black] (12.0000,8.0000) circle (1.2pt);
\fill[red!80!black] (0.0000,8.0000) circle (1.2pt);
\fill[red!80!black] (12.0000,8.0000) circle (1.2pt);
\fill[red!80!black] (0.0000,8.0000) circle (1.2pt);
\fill[red!80!black] (12.0000,8.0000) circle (1.2pt);
\fill[red!80!black] (0.0000,8.0000) circle (1.2pt);
\fill[red!80!black] (12.0000,8.0000) circle (1.2pt);
\fill[red!80!black] (0.0000,8.0000) circle (1.2pt);
\fill[red!80!black] (12.0000,8.0000) circle (1.2pt);
\fill[red!80!black] (0.0000,8.0000) circle (1.2pt);
\fill[red!80!black] (12.0000,8.0000) circle (1.2pt);
\fill[red!80!black] (0.0000,8.0000) circle (1.2pt);
\fill[red!80!black] (12.0000,8.0000) circle (1.2pt);
\fill[red!80!black] (0.0000,8.0000) circle (1.2pt);
\fill[red!80!black] (12.0000,8.0000) circle (1.2pt);
\fill[red!80!black] (0.0000,8.0000) circle (1.2pt);
\fill[red!80!black] (12.0000,8.0000) circle (1.2pt);
\fill[red!80!black] (0.0000,8.0000) circle (1.2pt);
\fill[red!80!black] (12.0000,8.0000) circle (1.2pt);
\fill[red!80!black] (0.0000,8.0000) circle (1.2pt);
\fill[red!80!black] (12.0000,8.0000) circle (1.2pt);
\fill[red!80!black] (0.0000,8.0000) circle (1.2pt);
\fill[red!80!black] (12.0000,8.0000) circle (1.2pt);
\fill[red!80!black] (0.0000,8.0000) circle (1.2pt);
\fill[red!80!black] (12.0000,8.0000) circle (1.2pt);
\fill[red!80!black] (0.0000,8.0000) circle (1.2pt);
\fill[red!80!black] (12.0000,8.0000) circle (1.2pt);
\fill[red!80!black] (0.0000,8.0000) circle (1.2pt);
\fill[red!80!black] (12.0000,8.0000) circle (1.2pt);
\fill[red!80!black] (0.0000,8.0000) circle (1.2pt);
\fill[red!80!black] (12.0000,8.0000) circle (1.2pt);
\fill[red!80!black] (0.0000,8.0000) circle (1.2pt);
\fill[red!80!black] (12.0000,8.0000) circle (1.2pt);
\fill[red!80!black] (0.0000,8.0000) circle (1.2pt);
\fill[red!80!black] (12.0000,8.0000) circle (1.2pt);
\fill[red!80!black] (0.0000,8.0000) circle (1.2pt);
\fill[red!80!black] (12.0000,8.0000) circle (1.2pt);
\fill[red!80!black] (0.0000,8.0000) circle (1.2pt);
\fill[red!80!black] (12.0000,8.0000) circle (1.2pt);
\fill[red!80!black] (0.0000,8.0000) circle (1.2pt);
\fill[red!80!black] (12.0000,8.0000) circle (1.2pt);
\fill[red!80!black] (0.0000,8.0000) circle (1.2pt);
\fill[red!80!black] (12.0000,8.0000) circle (1.2pt);
\fill[red!80!black] (0.0000,8.0000) circle (1.2pt);
\fill[red!80!black] (12.0000,8.0000) circle (1.2pt);
\fill[red!80!black] (0.0000,8.0000) circle (1.2pt);
\fill[red!80!black] (12.0000,8.0000) circle (1.2pt);
\fill[red!80!black] (0.0000,8.0000) circle (1.2pt);
\fill[red!80!black] (12.0000,8.0000) circle (1.2pt);
\fill[red!80!black] (0.0000,8.0000) circle (1.2pt);
\fill[red!80!black] (12.0000,8.0000) circle (1.2pt);
\fill[red!80!black] (0.0000,8.0000) circle (1.2pt);
\fill[red!80!black] (12.0000,8.0000) circle (1.2pt);
\fill[red!80!black] (0.0000,8.0000) circle (1.2pt);
\fill[red!80!black] (12.0000,8.0000) circle (1.2pt);
\fill[red!80!black] (0.0000,8.0000) circle (1.2pt);
\fill[red!80!black] (12.0000,8.0000) circle (1.2pt);
\fill[red!80!black] (0.0000,8.0000) circle (1.2pt);
\fill[red!80!black] (12.0000,8.0000) circle (1.2pt);
\fill[red!80!black] (0.0000,8.0000) circle (1.2pt);
\fill[red!80!black] (12.0000,8.0000) circle (1.2pt);
\fill[red!80!black] (0.0000,8.0000) circle (1.2pt);
\fill[red!80!black] (12.0000,8.0000) circle (1.2pt);
\fill[red!80!black] (0.0000,8.0000) circle (1.2pt);
\fill[red!80!black] (12.0000,8.0000) circle (1.2pt);
\fill[red!80!black] (0.0000,8.0000) circle (1.2pt);
\fill[red!80!black] (12.0000,8.0000) circle (1.2pt);
\fill[red!80!black] (0.0000,8.0000) circle (1.2pt);
\fill[red!80!black] (12.0000,8.0000) circle (1.2pt);
\fill[red!80!black] (0.0000,8.0000) circle (1.2pt);
\fill[red!80!black] (12.0000,8.0000) circle (1.2pt);
\fill[red!80!black] (0.0000,8.0000) circle (1.2pt);
\fill[red!80!black] (12.0000,8.0000) circle (1.2pt);
\fill[red!80!black] (0.0000,8.0000) circle (1.2pt);
\fill[red!80!black] (12.0000,8.0000) circle (1.2pt);
\fill[red!80!black] (0.0000,8.0000) circle (1.2pt);
\fill[red!80!black] (12.0000,8.0000) circle (1.2pt);
\fill[red!80!black] (0.0000,8.0000) circle (1.2pt);
\fill[red!80!black] (12.0000,8.0000) circle (1.2pt);
\fill[red!80!black] (0.0000,8.0000) circle (1.2pt);
\fill[red!80!black] (12.0000,8.0000) circle (1.2pt);
\fill[red!80!black] (0.0000,8.0000) circle (1.2pt);
\fill[red!80!black] (12.0000,8.0000) circle (1.2pt);
\fill[red!80!black] (0.0000,8.0000) circle (1.2pt);
\fill[red!80!black] (12.0000,8.0000) circle (1.2pt);
\fill[red!80!black] (0.0000,8.0000) circle (1.2pt);
\fill[red!80!black] (12.0000,8.0000) circle (1.2pt);
\fill[red!80!black] (0.0000,8.0000) circle (1.2pt);
\fill[red!80!black] (12.0000,8.0000) circle (1.2pt);
\fill[red!80!black] (0.0000,8.0000) circle (1.2pt);
\fill[red!80!black] (12.0000,8.0000) circle (1.2pt);
\fill[red!80!black] (0.0000,8.0000) circle (1.2pt);
\fill[red!80!black] (12.0000,8.0000) circle (1.2pt);
\fill[red!80!black] (0.0000,8.0000) circle (1.2pt);
\fill[red!80!black] (12.0000,8.0000) circle (1.2pt);
\fill[red!80!black] (0.0000,8.0000) circle (1.2pt);
\fill[red!80!black] (12.0000,8.0000) circle (1.2pt);
\fill[red!80!black] (0.0000,8.0000) circle (1.2pt);
\fill[red!80!black] (12.0000,8.0000) circle (1.2pt);
\fill[red!80!black] (0.0000,8.0000) circle (1.2pt);
\fill[red!80!black] (12.0000,8.0000) circle (1.2pt);
\fill[red!80!black] (0.0000,8.0000) circle (1.2pt);
\fill[red!80!black] (12.0000,8.0000) circle (1.2pt);
\fill[red!80!black] (0.0000,8.0000) circle (1.2pt);
\fill[red!80!black] (12.0000,8.0000) circle (1.2pt);
\fill[red!80!black] (0.0000,8.0000) circle (1.2pt);
\fill[red!80!black] (12.0000,8.0000) circle (1.2pt);
\fill[red!80!black] (0.0000,8.0000) circle (1.2pt);
\fill[red!80!black] (12.0000,8.0000) circle (1.2pt);
\fill[red!80!black] (0.0000,8.0000) circle (1.2pt);
\fill[red!80!black] (12.0000,8.0000) circle (1.2pt);
\fill[red!80!black] (0.0000,8.0000) circle (1.2pt);
\fill[red!80!black] (12.0000,8.0000) circle (1.2pt);
\fill[red!80!black] (0.0000,8.0000) circle (1.2pt);
\fill[red!80!black] (12.0000,8.0000) circle (1.2pt);
\fill[red!80!black] (0.0000,8.0000) circle (1.2pt);
\fill[red!80!black] (12.0000,8.0000) circle (1.2pt);
\fill[red!80!black] (0.0000,8.0000) circle (1.2pt);
\fill[red!80!black] (12.0000,8.0000) circle (1.2pt);
\fill[red!80!black] (0.0000,8.0000) circle (1.2pt);
\fill[red!80!black] (12.0000,8.0000) circle (1.2pt);
\fill[red!80!black] (0.0000,8.0000) circle (1.2pt);
\fill[red!80!black] (12.0000,8.0000) circle (1.2pt);
\fill[red!80!black] (0.0000,8.0000) circle (1.2pt);
\fill[red!80!black] (12.0000,8.0000) circle (1.2pt);
\fill[red!80!black] (0.0000,8.0000) circle (1.2pt);
\fill[red!80!black] (12.0000,8.0000) circle (1.2pt);
\fill[red!80!black] (0.0000,8.0000) circle (1.2pt);
\fill[red!80!black] (12.0000,8.0000) circle (1.2pt);
\fill[red!80!black] (0.0000,8.0000) circle (1.2pt);
\fill[red!80!black] (12.0000,8.0000) circle (1.2pt);
\fill[red!80!black] (0.0000,8.0000) circle (1.2pt);
\fill[red!80!black] (12.0000,8.0000) circle (1.2pt);
\fill[red!80!black] (0.0000,8.0000) circle (1.2pt);
\fill[red!80!black] (12.0000,8.0000) circle (1.2pt);
\fill[red!80!black] (0.0000,8.0000) circle (1.2pt);
\fill[red!80!black] (12.0000,8.0000) circle (1.2pt);
\fill[red!80!black] (0.0000,8.0000) circle (1.2pt);
\fill[red!80!black] (12.0000,8.0000) circle (1.2pt);
\fill[red!80!black] (0.0000,8.0000) circle (1.2pt);
\fill[red!80!black] (12.0000,8.0000) circle (1.2pt);
\fill[red!80!black] (0.0000,8.0000) circle (1.2pt);
\fill[red!80!black] (12.0000,8.0000) circle (1.2pt);
\fill[red!80!black] (0.0000,8.0000) circle (1.2pt);
\fill[red!80!black] (12.0000,8.0000) circle (1.2pt);
\fill[red!80!black] (0.0000,8.0000) circle (1.2pt);
\fill[red!80!black] (12.0000,8.0000) circle (1.2pt);
\fill[red!80!black] (0.0000,8.0000) circle (1.2pt);
\fill[red!80!black] (12.0000,8.0000) circle (1.2pt);
\fill[red!80!black] (0.0000,8.0000) circle (1.2pt);
\fill[red!80!black] (12.0000,8.0000) circle (1.2pt);
\fill[red!80!black] (0.0000,8.0000) circle (1.2pt);
\fill[red!80!black] (12.0000,8.0000) circle (1.2pt);
\fill[red!80!black] (0.0000,8.0000) circle (1.2pt);
\fill[red!80!black] (12.0000,8.0000) circle (1.2pt);
\fill[red!80!black] (0.0000,8.0000) circle (1.2pt);
\fill[red!80!black] (12.0000,8.0000) circle (1.2pt);
\fill[red!80!black] (0.0000,8.0000) circle (1.2pt);
\fill[red!80!black] (12.0000,8.0000) circle (1.2pt);
\fill[red!80!black] (0.0000,8.0000) circle (1.2pt);
\fill[red!80!black] (12.0000,8.0000) circle (1.2pt);
\fill[red!80!black] (0.0000,8.0000) circle (1.2pt);
\fill[red!80!black] (12.0000,8.0000) circle (1.2pt);
\fill[red!80!black] (0.0000,8.0000) circle (1.2pt);
\fill[red!80!black] (12.0000,8.0000) circle (1.2pt);
\fill[red!80!black] (0.0000,8.0000) circle (1.2pt);
\fill[red!80!black] (12.0000,8.0000) circle (1.2pt);
\fill[red!80!black] (0.0000,8.0000) circle (1.2pt);
\fill[red!80!black] (12.0000,8.0000) circle (1.2pt);
\fill[red!80!black] (0.0000,8.0000) circle (1.2pt);
\fill[red!80!black] (12.0000,8.0000) circle (1.2pt);
\fill[red!80!black] (0.0000,8.0000) circle (1.2pt);
\fill[red!80!black] (12.0000,8.0000) circle (1.2pt);
\fill[red!80!black] (0.0000,8.0000) circle (1.2pt);
\fill[red!80!black] (12.0000,8.0000) circle (1.2pt);
\fill[red!80!black] (0.0000,8.0000) circle (1.2pt);
\fill[red!80!black] (12.0000,8.0000) circle (1.2pt);
\fill[red!80!black] (0.0000,8.0000) circle (1.2pt);
\fill[red!80!black] (12.0000,8.0000) circle (1.2pt);
\fill[red!80!black] (0.0000,8.0000) circle (1.2pt);
\fill[red!80!black] (12.0000,8.0000) circle (1.2pt);
\fill[red!80!black] (0.0000,8.0000) circle (1.2pt);
\fill[red!80!black] (12.0000,8.0000) circle (1.2pt);
\fill[red!80!black] (0.0000,8.0000) circle (1.2pt);
\fill[red!80!black] (12.0000,8.0000) circle (1.2pt);
\fill[red!80!black] (0.0000,8.0000) circle (1.2pt);
\fill[red!80!black] (12.0000,8.0000) circle (1.2pt);
\fill[red!80!black] (0.0000,8.0000) circle (1.2pt);
\fill[red!80!black] (12.0000,8.0000) circle (1.2pt);
\fill[red!80!black] (0.0000,8.0000) circle (1.2pt);
\fill[red!80!black] (12.0000,8.0000) circle (1.2pt);
\fill[red!80!black] (0.0000,8.0000) circle (1.2pt);
\fill[red!80!black] (12.0000,8.0000) circle (1.2pt);
\fill[red!80!black] (0.0000,8.0000) circle (1.2pt);
\fill[red!80!black] (12.0000,8.0000) circle (1.2pt);
\fill[red!80!black] (0.0000,8.0000) circle (1.2pt);
\fill[red!80!black] (12.0000,8.0000) circle (1.2pt);
\fill[red!80!black] (0.0000,8.0000) circle (1.2pt);
\fill[red!80!black] (12.0000,8.0000) circle (1.2pt);
\fill[red!80!black] (0.0000,8.0000) circle (1.2pt);
\fill[red!80!black] (12.0000,8.0000) circle (1.2pt);
\fill[red!80!black] (0.0000,8.0000) circle (1.2pt);
\fill[red!80!black] (12.0000,8.0000) circle (1.2pt);
\fill[red!80!black] (0.0000,8.0000) circle (1.2pt);
\fill[red!80!black] (12.0000,8.0000) circle (1.2pt);
\fill[red!80!black] (0.0000,8.0000) circle (1.2pt);
\fill[red!80!black] (12.0000,8.0000) circle (1.2pt);
\fill[red!80!black] (0.0000,8.0000) circle (1.2pt);
\fill[red!80!black] (12.0000,8.0000) circle (1.2pt);
\fill[red!80!black] (0.0000,8.0000) circle (1.2pt);
\fill[red!80!black] (12.0000,8.0000) circle (1.2pt);
\fill[red!80!black] (0.0000,8.0000) circle (1.2pt);
\fill[red!80!black] (12.0000,8.0000) circle (1.2pt);
\fill[red!80!black] (0.0000,8.0000) circle (1.2pt);
\fill[red!80!black] (12.0000,8.0000) circle (1.2pt);
\fill[red!80!black] (0.0000,8.0000) circle (1.2pt);
\fill[red!80!black] (12.0000,8.0000) circle (1.2pt);
\fill[red!80!black] (0.0000,8.0000) circle (1.2pt);
\fill[red!80!black] (12.0000,8.0000) circle (1.2pt);
\fill[red!80!black] (0.0000,8.0000) circle (1.2pt);
\fill[red!80!black] (12.0000,8.0000) circle (1.2pt);
\fill[red!80!black] (0.0000,8.0000) circle (1.2pt);
\fill[red!80!black] (12.0000,8.0000) circle (1.2pt);
\fill[red!80!black] (0.0000,8.0000) circle (1.2pt);
\fill[red!80!black] (12.0000,8.0000) circle (1.2pt);
\fill[red!80!black] (0.0000,8.0000) circle (1.2pt);
\fill[red!80!black] (12.0000,8.0000) circle (1.2pt);
\fill[red!80!black] (0.0000,8.0000) circle (1.2pt);
\fill[red!80!black] (12.0000,8.0000) circle (1.2pt);
\fill[red!80!black] (0.0000,8.0000) circle (1.2pt);
\fill[red!80!black] (12.0000,8.0000) circle (1.2pt);
\fill[red!80!black] (0.0000,8.0000) circle (1.2pt);
\fill[red!80!black] (12.0000,8.0000) circle (1.2pt);
\fill[red!80!black] (0.0000,8.0000) circle (1.2pt);
\fill[red!80!black] (12.0000,8.0000) circle (1.2pt);
\fill[red!80!black] (0.0000,8.0000) circle (1.2pt);
\fill[red!80!black] (12.0000,8.0000) circle (1.2pt);
\fill[red!80!black] (0.0000,8.0000) circle (1.2pt);
\fill[red!80!black] (12.0000,8.0000) circle (1.2pt);
\fill[red!80!black] (0.0000,8.0000) circle (1.2pt);
\fill[red!80!black] (12.0000,8.0000) circle (1.2pt);
\fill[red!80!black] (0.0000,8.0000) circle (1.2pt);
\fill[red!80!black] (12.0000,8.0000) circle (1.2pt);
\fill[red!80!black] (0.0000,8.0000) circle (1.2pt);
\fill[red!80!black] (12.0000,8.0000) circle (1.2pt);
\fill[red!80!black] (0.0000,8.0000) circle (1.2pt);
\fill[red!80!black] (12.0000,8.0000) circle (1.2pt);
\fill[red!80!black] (0.0000,8.0000) circle (1.2pt);
\fill[red!80!black] (12.0000,8.0000) circle (1.2pt);
\fill[red!80!black] (0.0000,8.0000) circle (1.2pt);
\fill[red!80!black] (12.0000,8.0000) circle (1.2pt);
\fill[red!80!black] (0.0000,8.0000) circle (1.2pt);
\fill[red!80!black] (12.0000,8.0000) circle (1.2pt);
\fill[red!80!black] (0.0000,8.0000) circle (1.2pt);
\fill[red!80!black] (12.0000,8.0000) circle (1.2pt);
\fill[red!80!black] (0.0000,8.0000) circle (1.2pt);
\fill[red!80!black] (12.0000,8.0000) circle (1.2pt);
\fill[red!80!black] (0.0000,8.0000) circle (1.2pt);
\fill[red!80!black] (12.0000,8.0000) circle (1.2pt);
\fill[red!80!black] (0.0000,8.0000) circle (1.2pt);
\fill[red!80!black] (12.0000,8.0000) circle (1.2pt);
\fill[red!80!black] (0.0000,8.0000) circle (1.2pt);
\fill[red!80!black] (12.0000,8.0000) circle (1.2pt);
\fill[red!80!black] (0.0000,8.0000) circle (1.2pt);
\fill[red!80!black] (12.0000,8.0000) circle (1.2pt);
\fill[red!80!black] (0.0000,8.0000) circle (1.2pt);
\fill[red!80!black] (12.0000,8.0000) circle (1.2pt);
\fill[red!80!black] (0.0000,8.0000) circle (1.2pt);
\fill[red!80!black] (12.0000,8.0000) circle (1.2pt);
\fill[red!80!black] (0.0000,8.0000) circle (1.2pt);
\fill[red!80!black] (12.0000,8.0000) circle (1.2pt);
\fill[red!80!black] (0.0000,8.0000) circle (1.2pt);
\fill[red!80!black] (12.0000,8.0000) circle (1.2pt);
\fill[red!80!black] (0.0000,8.0000) circle (1.2pt);
\fill[red!80!black] (12.0000,8.0000) circle (1.2pt);
\fill[red!80!black] (0.0000,8.0000) circle (1.2pt);
\fill[red!80!black] (12.0000,8.0000) circle (1.2pt);
\fill[red!80!black] (0.0000,8.0000) circle (1.2pt);
\fill[red!80!black] (12.0000,8.0000) circle (1.2pt);
\fill[red!80!black] (0.0000,8.0000) circle (1.2pt);
\fill[red!80!black] (12.0000,8.0000) circle (1.2pt);
\fill[red!80!black] (0.0000,8.0000) circle (1.2pt);
\fill[red!80!black] (12.0000,8.0000) circle (1.2pt);
\fill[red!80!black] (0.0000,8.0000) circle (1.2pt);
\fill[red!80!black] (12.0000,8.0000) circle (1.2pt);
\fill[red!80!black] (0.0000,8.0000) circle (1.2pt);
\fill[red!80!black] (12.0000,8.0000) circle (1.2pt);
\fill[red!80!black] (0.0000,8.0000) circle (1.2pt);
\fill[red!80!black] (12.0000,8.0000) circle (1.2pt);
\fill[red!80!black] (0.0000,8.0000) circle (1.2pt);
\fill[red!80!black] (12.0000,8.0000) circle (1.2pt);
\fill[red!80!black] (0.0000,8.0000) circle (1.2pt);
\fill[red!80!black] (12.0000,8.0000) circle (1.2pt);
\fill[red!80!black] (0.0000,8.0000) circle (1.2pt);
\fill[red!80!black] (12.0000,8.0000) circle (1.2pt);
\fill[red!80!black] (0.0000,8.0000) circle (1.2pt);
\fill[red!80!black] (12.0000,8.0000) circle (1.2pt);
\fill[red!80!black] (0.0000,8.0000) circle (1.2pt);
\fill[red!80!black] (12.0000,8.0000) circle (1.2pt);
\fill[red!80!black] (0.0000,8.0000) circle (1.2pt);
\fill[red!80!black] (12.0000,8.0000) circle (1.2pt);
\fill[red!80!black] (0.0000,8.0000) circle (1.2pt);
\fill[red!80!black] (12.0000,8.0000) circle (1.2pt);
\fill[red!80!black] (0.0000,8.0000) circle (1.2pt);
\fill[red!80!black] (12.0000,8.0000) circle (1.2pt);
\fill[red!80!black] (0.0000,8.0000) circle (1.2pt);
\fill[red!80!black] (12.0000,8.0000) circle (1.2pt);
\fill[red!80!black] (0.0000,8.0000) circle (1.2pt);
\fill[red!80!black] (12.0000,8.0000) circle (1.2pt);
\fill[red!80!black] (0.0000,8.0000) circle (1.2pt);
\fill[red!80!black] (12.0000,8.0000) circle (1.2pt);
\fill[red!80!black] (0.0000,8.0000) circle (1.2pt);
\fill[red!80!black] (12.0000,8.0000) circle (1.2pt);
\fill[red!80!black] (0.0000,8.0000) circle (1.2pt);
\fill[red!80!black] (12.0000,8.0000) circle (1.2pt);
\fill[red!80!black] (0.0000,8.0000) circle (1.2pt);
\fill[red!80!black] (12.0000,8.0000) circle (1.2pt);
\fill[red!80!black] (0.0000,8.0000) circle (1.2pt);
\fill[red!80!black] (12.0000,8.0000) circle (1.2pt);
\fill[red!80!black] (0.0000,8.0000) circle (1.2pt);
\fill[red!80!black] (12.0000,8.0000) circle (1.2pt);
\fill[red!80!black] (0.0000,8.0000) circle (1.2pt);
\fill[red!80!black] (12.0000,8.0000) circle (1.2pt);
\fill[red!80!black] (0.0000,8.0000) circle (1.2pt);
\fill[red!80!black] (12.0000,8.0000) circle (1.2pt);
\fill[red!80!black] (0.0000,8.0000) circle (1.2pt);
\fill[red!80!black] (12.0000,8.0000) circle (1.2pt);
\fill[red!80!black] (0.0000,8.0000) circle (1.2pt);
\fill[red!80!black] (12.0000,8.0000) circle (1.2pt);
\fill[red!80!black] (0.0000,8.0000) circle (1.2pt);
\fill[red!80!black] (12.0000,8.0000) circle (1.2pt);
\fill[red!80!black] (0.0000,8.0000) circle (1.2pt);
\fill[red!80!black] (12.0000,8.0000) circle (1.2pt);
\fill[red!80!black] (0.0000,8.0000) circle (1.2pt);
\fill[red!80!black] (12.0000,8.0000) circle (1.2pt);
\fill[red!80!black] (0.0000,8.0000) circle (1.2pt);
\fill[red!80!black] (12.0000,8.0000) circle (1.2pt);
\fill[red!80!black] (0.0000,8.0000) circle (1.2pt);
\fill[red!80!black] (12.0000,8.0000) circle (1.2pt);
\fill[red!80!black] (0.0000,8.0000) circle (1.2pt);
\fill[red!80!black] (12.0000,8.0000) circle (1.2pt);
\fill[red!80!black] (0.0000,8.0000) circle (1.2pt);
\fill[red!80!black] (12.0000,8.0000) circle (1.2pt);
\fill[red!80!black] (0.0000,8.0000) circle (1.2pt);
\fill[red!80!black] (12.0000,8.0000) circle (1.2pt);
\fill[red!80!black] (0.0000,8.0000) circle (1.2pt);
\fill[red!80!black] (12.0000,8.0000) circle (1.2pt);
\fill[red!80!black] (0.0000,8.0000) circle (1.2pt);
\fill[red!80!black] (12.0000,8.0000) circle (1.2pt);
\fill[red!80!black] (0.0000,8.0000) circle (1.2pt);
\fill[red!80!black] (12.0000,8.0000) circle (1.2pt);
\fill[red!80!black] (0.0000,8.0000) circle (1.2pt);
\fill[red!80!black] (12.0000,8.0000) circle (1.2pt);
\fill[red!80!black] (0.0000,8.0000) circle (1.2pt);
\fill[red!80!black] (12.0000,8.0000) circle (1.2pt);
\fill[red!80!black] (0.0000,8.0000) circle (1.2pt);
\fill[red!80!black] (12.0000,8.0000) circle (1.2pt);
\fill[red!80!black] (0.0000,8.0000) circle (1.2pt);
\fill[red!80!black] (12.0000,8.0000) circle (1.2pt);
\fill[red!80!black] (0.0000,8.0000) circle (1.2pt);
\fill[red!80!black] (12.0000,8.0000) circle (1.2pt);
\fill[red!80!black] (0.0000,8.0000) circle (1.2pt);
\fill[red!80!black] (12.0000,8.0000) circle (1.2pt);
\fill[red!80!black] (0.0000,8.0000) circle (1.2pt);
\fill[red!80!black] (12.0000,8.0000) circle (1.2pt);
\fill[red!80!black] (0.0000,8.0000) circle (1.2pt);
\fill[red!80!black] (12.0000,8.0000) circle (1.2pt);
\fill[red!80!black] (0.0000,8.0000) circle (1.2pt);
\fill[red!80!black] (12.0000,8.0000) circle (1.2pt);
\fill[red!80!black] (0.0000,8.0000) circle (1.2pt);
\fill[red!80!black] (12.0000,8.0000) circle (1.2pt);
\fill[red!80!black] (0.0000,8.0000) circle (1.2pt);
\fill[red!80!black] (12.0000,8.0000) circle (1.2pt);
\fill[red!80!black] (0.0000,8.0000) circle (1.2pt);
\fill[red!80!black] (12.0000,8.0000) circle (1.2pt);
\fill[red!80!black] (0.0000,8.0000) circle (1.2pt);
\fill[red!80!black] (12.0000,8.0000) circle (1.2pt);
\fill[red!80!black] (0.0000,8.0000) circle (1.2pt);
\fill[red!80!black] (12.0000,8.0000) circle (1.2pt);
\fill[red!80!black] (0.0000,8.0000) circle (1.2pt);
\fill[red!80!black] (12.0000,8.0000) circle (1.2pt);
\fill[red!80!black] (0.0000,8.0000) circle (1.2pt);
\fill[red!80!black] (12.0000,8.0000) circle (1.2pt);
\fill[red!80!black] (0.0000,8.0000) circle (1.2pt);
\fill[red!80!black] (12.0000,8.0000) circle (1.2pt);
\fill[red!80!black] (0.0000,8.0000) circle (1.2pt);
\fill[red!80!black] (12.0000,8.0000) circle (1.2pt);
\fill[red!80!black] (0.0000,8.0000) circle (1.2pt);
\fill[red!80!black] (12.0000,8.0000) circle (1.2pt);
\fill[red!80!black] (0.0000,8.0000) circle (1.2pt);
\fill[red!80!black] (12.0000,8.0000) circle (1.2pt);
\fill[red!80!black] (0.0000,8.0000) circle (1.2pt);
\fill[red!80!black] (12.0000,8.0000) circle (1.2pt);
\fill[red!80!black] (0.0000,8.0000) circle (1.2pt);
\fill[red!80!black] (12.0000,8.0000) circle (1.2pt);
\fill[red!80!black] (0.0000,8.0000) circle (1.2pt);
\fill[red!80!black] (12.0000,8.0000) circle (1.2pt);
\fill[red!80!black] (0.0000,8.0000) circle (1.2pt);
\fill[red!80!black] (12.0000,8.0000) circle (1.2pt);
\fill[red!80!black] (0.0000,8.0000) circle (1.2pt);
\fill[red!80!black] (12.0000,8.0000) circle (1.2pt);
\fill[red!80!black] (0.0000,8.0000) circle (1.2pt);
\fill[red!80!black] (12.0000,8.0000) circle (1.2pt);
\fill[red!80!black] (0.0000,8.0000) circle (1.2pt);
\fill[red!80!black] (12.0000,8.0000) circle (1.2pt);
\fill[red!80!black] (0.0000,8.0000) circle (1.2pt);
\fill[red!80!black] (12.0000,8.0000) circle (1.2pt);
\fill[red!80!black] (0.0000,8.0000) circle (1.2pt);
\fill[red!80!black] (12.0000,8.0000) circle (1.2pt);
\fill[red!80!black] (0.0000,8.0000) circle (1.2pt);
\fill[red!80!black] (12.0000,8.0000) circle (1.2pt);
\fill[red!80!black] (0.0000,8.0000) circle (1.2pt);
\fill[red!80!black] (12.0000,8.0000) circle (1.2pt);
\fill[red!80!black] (0.0000,8.0000) circle (1.2pt);
\fill[red!80!black] (12.0000,8.0000) circle (1.2pt);
\fill[red!80!black] (0.0000,8.0000) circle (1.2pt);
\fill[red!80!black] (12.0000,8.0000) circle (1.2pt);
\fill[red!80!black] (0.0000,8.0000) circle (1.2pt);
\fill[red!80!black] (12.0000,8.0000) circle (1.2pt);
\fill[red!80!black] (0.0000,8.0000) circle (1.2pt);
\fill[red!80!black] (12.0000,8.0000) circle (1.2pt);
\fill[red!80!black] (0.0000,8.0000) circle (1.2pt);
\fill[red!80!black] (12.0000,8.0000) circle (1.2pt);
\fill[red!80!black] (0.0000,8.0000) circle (1.2pt);
\fill[red!80!black] (12.0000,8.0000) circle (1.2pt);
\fill[red!80!black] (0.0000,8.0000) circle (1.2pt);
\fill[red!80!black] (12.0000,8.0000) circle (1.2pt);
\fill[red!80!black] (0.0000,8.0000) circle (1.2pt);
\fill[red!80!black] (12.0000,8.0000) circle (1.2pt);
\fill[red!80!black] (0.0000,8.0000) circle (1.2pt);
\fill[red!80!black] (12.0000,8.0000) circle (1.2pt);
\fill[red!80!black] (0.0000,8.0000) circle (1.2pt);
\fill[red!80!black] (12.0000,8.0000) circle (1.2pt);
\fill[red!80!black] (0.0000,8.0000) circle (1.2pt);
\fill[red!80!black] (12.0000,8.0000) circle (1.2pt);
\fill[red!80!black] (0.0000,8.0000) circle (1.2pt);
\fill[red!80!black] (12.0000,8.0000) circle (1.2pt);
\fill[red!80!black] (0.0000,8.0000) circle (1.2pt);
\fill[red!80!black] (12.0000,8.0000) circle (1.2pt);
\fill[red!80!black] (0.0000,8.0000) circle (1.2pt);
\fill[red!80!black] (12.0000,8.0000) circle (1.2pt);
\fill[red!80!black] (0.0000,8.0000) circle (1.2pt);
\fill[red!80!black] (12.0000,8.0000) circle (1.2pt);
\fill[red!80!black] (0.0000,8.0000) circle (1.2pt);
\fill[red!80!black] (12.0000,8.0000) circle (1.2pt);
\fill[red!80!black] (0.0000,8.0000) circle (1.2pt);
\fill[red!80!black] (12.0000,8.0000) circle (1.2pt);
\fill[red!80!black] (0.0000,8.0000) circle (1.2pt);
\fill[red!80!black] (12.0000,8.0000) circle (1.2pt);
\fill[red!80!black] (0.0000,8.0000) circle (1.2pt);
\fill[red!80!black] (12.0000,8.0000) circle (1.2pt);
\fill[red!80!black] (0.0000,8.0000) circle (1.2pt);
\fill[red!80!black] (12.0000,8.0000) circle (1.2pt);
\fill[red!80!black] (0.0000,8.0000) circle (1.2pt);
\fill[red!80!black] (12.0000,8.0000) circle (1.2pt);
\fill[red!80!black] (0.0000,8.0000) circle (1.2pt);
\fill[red!80!black] (12.0000,8.0000) circle (1.2pt);
\fill[red!80!black] (0.0000,8.0000) circle (1.2pt);
\fill[red!80!black] (12.0000,8.0000) circle (1.2pt);
\fill[red!80!black] (0.0000,8.0000) circle (1.2pt);
\fill[red!80!black] (12.0000,8.0000) circle (1.2pt);
\fill[red!80!black] (0.0000,8.0000) circle (1.2pt);
\fill[red!80!black] (12.0000,8.0000) circle (1.2pt);
\fill[red!80!black] (0.0000,8.0000) circle (1.2pt);
\fill[red!80!black] (12.0000,8.0000) circle (1.2pt);
\fill[red!80!black] (0.0000,8.0000) circle (1.2pt);
\fill[red!80!black] (12.0000,8.0000) circle (1.2pt);
\fill[red!80!black] (0.0000,8.0000) circle (1.2pt);
\fill[red!80!black] (12.0000,8.0000) circle (1.2pt);
\fill[red!80!black] (0.0000,8.0000) circle (1.2pt);
\fill[red!80!black] (12.0000,8.0000) circle (1.2pt);
\fill[red!80!black] (0.0000,8.0000) circle (1.2pt);
\fill[red!80!black] (12.0000,8.0000) circle (1.2pt);
\fill[red!80!black] (0.0000,8.0000) circle (1.2pt);
\fill[red!80!black] (12.0000,8.0000) circle (1.2pt);
\fill[red!80!black] (0.0000,8.0000) circle (1.2pt);
\fill[red!80!black] (12.0000,8.0000) circle (1.2pt);
\fill[red!80!black] (0.0000,8.0000) circle (1.2pt);
\fill[red!80!black] (12.0000,8.0000) circle (1.2pt);
\fill[red!80!black] (0.0000,8.0000) circle (1.2pt);
\fill[red!80!black] (12.0000,8.0000) circle (1.2pt);
\fill[red!80!black] (0.0000,8.0000) circle (1.2pt);
\fill[red!80!black] (12.0000,8.0000) circle (1.2pt);
\fill[red!80!black] (0.0000,8.0000) circle (1.2pt);
\fill[red!80!black] (12.0000,8.0000) circle (1.2pt);
\fill[red!80!black] (0.0000,8.0000) circle (1.2pt);
\fill[red!80!black] (12.0000,8.0000) circle (1.2pt);
\fill[red!80!black] (0.0000,8.0000) circle (1.2pt);
\fill[red!80!black] (12.0000,8.0000) circle (1.2pt);
\fill[red!80!black] (0.0000,8.0000) circle (1.2pt);
\fill[red!80!black] (12.0000,8.0000) circle (1.2pt);
\fill[red!80!black] (0.0000,8.0000) circle (1.2pt);
\fill[red!80!black] (12.0000,8.0000) circle (1.2pt);
\fill[red!80!black] (0.0000,8.0000) circle (1.2pt);
\fill[red!80!black] (12.0000,8.0000) circle (1.2pt);
\fill[red!80!black] (0.0000,8.0000) circle (1.2pt);
\fill[red!80!black] (12.0000,8.0000) circle (1.2pt);
\fill[red!80!black] (0.0000,8.0000) circle (1.2pt);
\fill[red!80!black] (12.0000,8.0000) circle (1.2pt);
\fill[red!80!black] (0.0000,8.0000) circle (1.2pt);
\fill[red!80!black] (12.0000,8.0000) circle (1.2pt);
\fill[red!80!black] (0.0000,8.0000) circle (1.2pt);
\fill[red!80!black] (12.0000,8.0000) circle (1.2pt);
\fill[red!80!black] (0.0000,8.0000) circle (1.2pt);
\fill[red!80!black] (12.0000,8.0000) circle (1.2pt);
\fill[red!80!black] (0.0000,8.0000) circle (1.2pt);
\fill[red!80!black] (12.0000,8.0000) circle (1.2pt);
\fill[red!80!black] (0.0000,8.0000) circle (1.2pt);
\fill[red!80!black] (12.0000,8.0000) circle (1.2pt);
\fill[red!80!black] (0.0000,8.0000) circle (1.2pt);
\fill[red!80!black] (12.0000,8.0000) circle (1.2pt);
\fill[red!80!black] (0.0000,8.0000) circle (1.2pt);
\fill[red!80!black] (12.0000,8.0000) circle (1.2pt);
\fill[red!80!black] (0.0000,8.0000) circle (1.2pt);
\fill[red!80!black] (12.0000,8.0000) circle (1.2pt);
\fill[red!80!black] (0.0000,8.0000) circle (1.2pt);
\fill[red!80!black] (12.0000,8.0000) circle (1.2pt);
\fill[red!80!black] (0.0000,8.0000) circle (1.2pt);
\fill[red!80!black] (12.0000,8.0000) circle (1.2pt);
\fill[red!80!black] (0.0000,8.0000) circle (1.2pt);
\fill[red!80!black] (12.0000,8.0000) circle (1.2pt);
\fill[red!80!black] (0.0000,8.0000) circle (1.2pt);
\fill[red!80!black] (12.0000,8.0000) circle (1.2pt);
\fill[red!80!black] (0.0000,8.0000) circle (1.2pt);
\fill[red!80!black] (12.0000,8.0000) circle (1.2pt);
\fill[red!80!black] (0.0000,8.0000) circle (1.2pt);
\fill[red!80!black] (12.0000,8.0000) circle (1.2pt);
\fill[red!80!black] (0.0000,8.0000) circle (1.2pt);
\fill[red!80!black] (12.0000,8.0000) circle (1.2pt);
\fill[red!80!black] (0.0000,8.0000) circle (1.2pt);
\fill[red!80!black] (12.0000,8.0000) circle (1.2pt);
\fill[red!80!black] (0.0000,8.0000) circle (1.2pt);
\fill[red!80!black] (12.0000,8.0000) circle (1.2pt);
\fill[red!80!black] (0.0000,8.0000) circle (1.2pt);
\fill[red!80!black] (12.0000,8.0000) circle (1.2pt);
\fill[red!80!black] (0.0000,8.0000) circle (1.2pt);
\fill[red!80!black] (12.0000,8.0000) circle (1.2pt);
\fill[red!80!black] (0.0000,8.0000) circle (1.2pt);
\fill[red!80!black] (12.0000,8.0000) circle (1.2pt);
\fill[red!80!black] (0.0000,8.0000) circle (1.2pt);
\fill[red!80!black] (12.0000,8.0000) circle (1.2pt);
\fill[red!80!black] (0.0000,8.0000) circle (1.2pt);
\fill[red!80!black] (12.0000,8.0000) circle (1.2pt);
\fill[red!80!black] (0.0000,8.0000) circle (1.2pt);
\fill[red!80!black] (12.0000,8.0000) circle (1.2pt);
\fill[red!80!black] (0.0000,8.0000) circle (1.2pt);
\fill[red!80!black] (12.0000,8.0000) circle (1.2pt);
\fill[red!80!black] (0.0000,8.0000) circle (1.2pt);
\fill[red!80!black] (12.0000,8.0000) circle (1.2pt);
\fill[red!80!black] (0.0000,8.0000) circle (1.2pt);
\fill[red!80!black] (12.0000,8.0000) circle (1.2pt);
\fill[red!80!black] (0.0000,8.0000) circle (1.2pt);
\fill[red!80!black] (12.0000,8.0000) circle (1.2pt);
\fill[red!80!black] (0.0000,8.0000) circle (1.2pt);
\fill[red!80!black] (12.0000,8.0000) circle (1.2pt);
\fill[red!80!black] (0.0000,8.0000) circle (1.2pt);
\fill[red!80!black] (12.0000,8.0000) circle (1.2pt);
\fill[red!80!black] (0.0000,8.0000) circle (1.2pt);
\fill[red!80!black] (12.0000,8.0000) circle (1.2pt);
\fill[red!80!black] (0.0000,8.0000) circle (1.2pt);
\fill[red!80!black] (12.0000,8.0000) circle (1.2pt);
\fill[red!80!black] (0.0000,8.0000) circle (1.2pt);
\fill[red!80!black] (12.0000,8.0000) circle (1.2pt);
\fill[red!80!black] (0.0000,8.0000) circle (1.2pt);
\fill[red!80!black] (12.0000,8.0000) circle (1.2pt);
\fill[red!80!black] (0.0000,8.0000) circle (1.2pt);
\fill[red!80!black] (12.0000,8.0000) circle (1.2pt);
\fill[red!80!black] (0.0000,8.0000) circle (1.2pt);
\fill[red!80!black] (12.0000,8.0000) circle (1.2pt);
\fill[red!80!black] (0.0000,8.0000) circle (1.2pt);
\fill[red!80!black] (12.0000,8.0000) circle (1.2pt);
\fill[red!80!black] (0.0000,8.0000) circle (1.2pt);
\fill[red!80!black] (12.0000,8.0000) circle (1.2pt);
\fill[red!80!black] (0.0000,8.0000) circle (1.2pt);
\fill[red!80!black] (12.0000,8.0000) circle (1.2pt);
\fill[red!80!black] (0.0000,8.0000) circle (1.2pt);
\fill[red!80!black] (12.0000,8.0000) circle (1.2pt);
\fill[red!80!black] (0.0000,8.0000) circle (1.2pt);
\fill[red!80!black] (12.0000,8.0000) circle (1.2pt);
\fill[red!80!black] (0.0000,8.0000) circle (1.2pt);
\fill[red!80!black] (12.0000,8.0000) circle (1.2pt);
\fill[red!80!black] (0.0000,8.0000) circle (1.2pt);
\fill[red!80!black] (12.0000,8.0000) circle (1.2pt);
\fill[red!80!black] (0.0000,8.0000) circle (1.2pt);
\fill[red!80!black] (12.0000,8.0000) circle (1.2pt);
\fill[red!80!black] (0.0000,8.0000) circle (1.2pt);
\fill[red!80!black] (12.0000,8.0000) circle (1.2pt);
\fill[red!80!black] (0.0000,8.0000) circle (1.2pt);
\fill[red!80!black] (12.0000,8.0000) circle (1.2pt);
\fill[red!80!black] (0.0000,8.0000) circle (1.2pt);
\fill[red!80!black] (12.0000,8.0000) circle (1.2pt);
\fill[red!80!black] (0.0000,8.0000) circle (1.2pt);
\fill[red!80!black] (12.0000,8.0000) circle (1.2pt);
\fill[red!80!black] (0.0000,8.0000) circle (1.2pt);
\fill[red!80!black] (12.0000,8.0000) circle (1.2pt);
\fill[red!80!black] (0.0000,8.0000) circle (1.2pt);
\fill[red!80!black] (12.0000,8.0000) circle (1.2pt);
\fill[red!80!black] (0.0000,8.0000) circle (1.2pt);
\fill[red!80!black] (12.0000,8.0000) circle (1.2pt);
\fill[red!80!black] (0.0000,8.0000) circle (1.2pt);
\fill[red!80!black] (12.0000,8.0000) circle (1.2pt);
\fill[red!80!black] (0.0000,8.0000) circle (1.2pt);
\fill[red!80!black] (12.0000,8.0000) circle (1.2pt);
\fill[red!80!black] (0.0000,8.0000) circle (1.2pt);
\fill[red!80!black] (12.0000,8.0000) circle (1.2pt);
\fill[red!80!black] (0.0000,8.0000) circle (1.2pt);
\fill[red!80!black] (12.0000,8.0000) circle (1.2pt);
\fill[red!80!black] (0.0000,8.0000) circle (1.2pt);
\fill[red!80!black] (12.0000,8.0000) circle (1.2pt);
\fill[red!80!black] (0.0000,8.0000) circle (1.2pt);
\fill[red!80!black] (12.0000,8.0000) circle (1.2pt);
\fill[red!80!black] (0.0000,8.0000) circle (1.2pt);
\fill[red!80!black] (12.0000,8.0000) circle (1.2pt);
\fill[red!80!black] (0.0000,8.0000) circle (1.2pt);
\fill[red!80!black] (12.0000,8.0000) circle (1.2pt);
\fill[red!80!black] (0.0000,8.0000) circle (1.2pt);
\fill[red!80!black] (12.0000,8.0000) circle (1.2pt);
\fill[red!80!black] (0.0000,8.0000) circle (1.2pt);
\fill[red!80!black] (12.0000,8.0000) circle (1.2pt);
\fill[red!80!black] (0.0000,8.0000) circle (1.2pt);
\fill[red!80!black] (12.0000,8.0000) circle (1.2pt);
\fill[red!80!black] (0.0000,8.0000) circle (1.2pt);
\fill[red!80!black] (12.0000,8.0000) circle (1.2pt);
\fill[red!80!black] (0.0000,8.0000) circle (1.2pt);
\fill[red!80!black] (12.0000,8.0000) circle (1.2pt);
\fill[red!80!black] (0.0000,8.0000) circle (1.2pt);
\fill[red!80!black] (12.0000,8.0000) circle (1.2pt);
\fill[red!80!black] (0.0000,8.0000) circle (1.2pt);
\fill[red!80!black] (12.0000,8.0000) circle (1.2pt);
\fill[red!80!black] (0.0000,8.0000) circle (1.2pt);
\fill[red!80!black] (12.0000,8.0000) circle (1.2pt);
\fill[red!80!black] (0.0000,8.0000) circle (1.2pt);
\fill[red!80!black] (12.0000,8.0000) circle (1.2pt);
\fill[red!80!black] (0.0000,8.0000) circle (1.2pt);
\fill[red!80!black] (12.0000,8.0000) circle (1.2pt);
\fill[red!80!black] (0.0000,8.0000) circle (1.2pt);
\fill[red!80!black] (12.0000,8.0000) circle (1.2pt);
\fill[red!80!black] (0.0000,8.0000) circle (1.2pt);
\fill[red!80!black] (12.0000,8.0000) circle (1.2pt);
\fill[red!80!black] (0.0000,8.0000) circle (1.2pt);
\fill[red!80!black] (12.0000,8.0000) circle (1.2pt);
\fill[red!80!black] (0.0000,8.0000) circle (1.2pt);
\fill[red!80!black] (12.0000,8.0000) circle (1.2pt);
\fill[red!80!black] (0.0000,8.0000) circle (1.2pt);
\fill[red!80!black] (12.0000,8.0000) circle (1.2pt);
\fill[red!80!black] (0.0000,8.0000) circle (1.2pt);
\fill[red!80!black] (12.0000,8.0000) circle (1.2pt);
\fill[red!80!black] (0.0000,8.0000) circle (1.2pt);
\fill[red!80!black] (12.0000,8.0000) circle (1.2pt);
\fill[red!80!black] (0.0000,8.0000) circle (1.2pt);
\fill[red!80!black] (12.0000,8.0000) circle (1.2pt);
\fill[red!80!black] (0.0000,8.0000) circle (1.2pt);
\fill[red!80!black] (12.0000,8.0000) circle (1.2pt);
\fill[red!80!black] (0.0000,8.0000) circle (1.2pt);
\fill[red!80!black] (12.0000,8.0000) circle (1.2pt);
\fill[red!80!black] (0.0000,8.0000) circle (1.2pt);
\fill[red!80!black] (12.0000,8.0000) circle (1.2pt);
\fill[red!80!black] (0.0000,8.0000) circle (1.2pt);
\fill[red!80!black] (12.0000,8.0000) circle (1.2pt);
\fill[red!80!black] (0.0000,8.0000) circle (1.2pt);
\fill[red!80!black] (12.0000,8.0000) circle (1.2pt);
\fill[red!80!black] (0.0000,8.0000) circle (1.2pt);
\fill[red!80!black] (12.0000,8.0000) circle (1.2pt);
\fill[red!80!black] (0.0000,8.0000) circle (1.2pt);
\fill[red!80!black] (12.0000,8.0000) circle (1.2pt);
\fill[red!80!black] (0.0000,8.0000) circle (1.2pt);
\fill[red!80!black] (12.0000,8.0000) circle (1.2pt);
\fill[red!80!black] (0.0000,8.0000) circle (1.2pt);
\fill[red!80!black] (12.0000,8.0000) circle (1.2pt);
\fill[red!80!black] (0.0000,8.0000) circle (1.2pt);
\fill[red!80!black] (12.0000,8.0000) circle (1.2pt);
\fill[red!80!black] (0.0000,8.0000) circle (1.2pt);
\fill[red!80!black] (12.0000,8.0000) circle (1.2pt);
\fill[red!80!black] (0.0000,8.0000) circle (1.2pt);
\fill[red!80!black] (12.0000,8.0000) circle (1.2pt);
\fill[red!80!black] (0.0000,8.0000) circle (1.2pt);
\fill[red!80!black] (12.0000,8.0000) circle (1.2pt);
\fill[red!80!black] (0.0000,8.0000) circle (1.2pt);
\fill[red!80!black] (12.0000,8.0000) circle (1.2pt);
\fill[red!80!black] (0.0000,8.0000) circle (1.2pt);
\fill[red!80!black] (12.0000,8.0000) circle (1.2pt);
\fill[red!80!black] (0.0000,8.0000) circle (1.2pt);
\fill[red!80!black] (12.0000,8.0000) circle (1.2pt);
\fill[red!80!black] (0.0000,8.0000) circle (1.2pt);
\fill[red!80!black] (12.0000,8.0000) circle (1.2pt);
\fill[red!80!black] (0.0000,8.0000) circle (1.2pt);
\fill[red!80!black] (12.0000,8.0000) circle (1.2pt);
\fill[red!80!black] (0.0000,8.0000) circle (1.2pt);
\fill[red!80!black] (12.0000,8.0000) circle (1.2pt);
\fill[red!80!black] (0.0000,8.0000) circle (1.2pt);
\fill[red!80!black] (12.0000,8.0000) circle (1.2pt);
\fill[red!80!black] (0.0000,8.0000) circle (1.2pt);
\fill[red!80!black] (12.0000,8.0000) circle (1.2pt);
\fill[red!80!black] (0.0000,8.0000) circle (1.2pt);
\fill[red!80!black] (12.0000,8.0000) circle (1.2pt);
\fill[red!80!black] (0.0000,8.0000) circle (1.2pt);
\fill[red!80!black] (12.0000,8.0000) circle (1.2pt);
\fill[red!80!black] (0.0000,8.0000) circle (1.2pt);
\fill[red!80!black] (12.0000,8.0000) circle (1.2pt);
\fill[red!80!black] (0.0000,8.0000) circle (1.2pt);
\fill[red!80!black] (12.0000,8.0000) circle (1.2pt);
\fill[red!80!black] (0.0000,8.0000) circle (1.2pt);
\fill[red!80!black] (12.0000,8.0000) circle (1.2pt);
\fill[red!80!black] (0.0000,8.0000) circle (1.2pt);
\fill[red!80!black] (12.0000,8.0000) circle (1.2pt);
\fill[red!80!black] (0.0000,8.0000) circle (1.2pt);
\fill[red!80!black] (12.0000,8.0000) circle (1.2pt);
\fill[red!80!black] (0.0000,8.0000) circle (1.2pt);
\fill[red!80!black] (12.0000,8.0000) circle (1.2pt);
\fill[red!80!black] (0.0000,8.0000) circle (1.2pt);
\fill[red!80!black] (12.0000,8.0000) circle (1.2pt);
\fill[red!80!black] (0.0000,8.0000) circle (1.2pt);
\fill[red!80!black] (12.0000,8.0000) circle (1.2pt);
\fill[red!80!black] (0.0000,8.0000) circle (1.2pt);
\fill[red!80!black] (12.0000,8.0000) circle (1.2pt);
\fill[red!80!black] (0.0000,8.0000) circle (1.2pt);
\fill[red!80!black] (12.0000,8.0000) circle (1.2pt);
\fill[red!80!black] (0.0000,8.0000) circle (1.2pt);
\fill[red!80!black] (12.0000,8.0000) circle (1.2pt);
\fill[red!80!black] (0.0000,8.0000) circle (1.2pt);
\fill[red!80!black] (12.0000,8.0000) circle (1.2pt);
\fill[red!80!black] (0.0000,8.0000) circle (1.2pt);
\fill[red!80!black] (12.0000,8.0000) circle (1.2pt);
\fill[red!80!black] (0.0000,8.0000) circle (1.2pt);
\fill[red!80!black] (12.0000,8.0000) circle (1.2pt);
\fill[red!80!black] (0.0000,8.0000) circle (1.2pt);
\fill[red!80!black] (12.0000,8.0000) circle (1.2pt);
\fill[red!80!black] (0.0000,8.0000) circle (1.2pt);
\fill[red!80!black] (12.0000,8.0000) circle (1.2pt);
\fill[red!80!black] (0.0000,8.0000) circle (1.2pt);
\fill[red!80!black] (12.0000,8.0000) circle (1.2pt);
\fill[red!80!black] (0.0000,8.0000) circle (1.2pt);
\fill[red!80!black] (12.0000,8.0000) circle (1.2pt);
\fill[red!80!black] (0.0000,8.0000) circle (1.2pt);
\fill[red!80!black] (12.0000,8.0000) circle (1.2pt);
\fill[red!80!black] (0.0000,8.0000) circle (1.2pt);
\fill[red!80!black] (12.0000,8.0000) circle (1.2pt);
\fill[red!80!black] (0.0000,8.0000) circle (1.2pt);
\fill[red!80!black] (12.0000,8.0000) circle (1.2pt);
\fill[red!80!black] (0.0000,8.0000) circle (1.2pt);
\fill[red!80!black] (12.0000,8.0000) circle (1.2pt);
\fill[red!80!black] (0.0000,8.0000) circle (1.2pt);
\fill[red!80!black] (12.0000,8.0000) circle (1.2pt);
\fill[red!80!black] (0.0000,8.0000) circle (1.2pt);
\fill[red!80!black] (12.0000,8.0000) circle (1.2pt);
\fill[red!80!black] (0.0000,8.0000) circle (1.2pt);
\fill[red!80!black] (12.0000,8.0000) circle (1.2pt);
\fill[red!80!black] (0.0000,8.0000) circle (1.2pt);
\fill[red!80!black] (12.0000,8.0000) circle (1.2pt);
\fill[red!80!black] (0.0000,8.0000) circle (1.2pt);
\fill[red!80!black] (12.0000,8.0000) circle (1.2pt);
\fill[red!80!black] (0.0000,8.0000) circle (1.2pt);
\fill[red!80!black] (12.0000,8.0000) circle (1.2pt);
\fill[red!80!black] (0.0000,8.0000) circle (1.2pt);
\fill[red!80!black] (12.0000,8.0000) circle (1.2pt);
\fill[red!80!black] (0.0000,8.0000) circle (1.2pt);
\fill[red!80!black] (12.0000,8.0000) circle (1.2pt);
\fill[red!80!black] (0.0000,8.0000) circle (1.2pt);
\fill[red!80!black] (12.0000,8.0000) circle (1.2pt);
\fill[red!80!black] (0.0000,8.0000) circle (1.2pt);
\fill[red!80!black] (12.0000,8.0000) circle (1.2pt);
\fill[red!80!black] (0.0000,8.0000) circle (1.2pt);
\fill[red!80!black] (12.0000,8.0000) circle (1.2pt);
\fill[red!80!black] (0.0000,8.0000) circle (1.2pt);
\fill[red!80!black] (12.0000,8.0000) circle (1.2pt);
\fill[red!80!black] (0.0000,8.0000) circle (1.2pt);
\fill[red!80!black] (12.0000,8.0000) circle (1.2pt);
\fill[red!80!black] (0.0000,8.0000) circle (1.2pt);
\fill[red!80!black] (12.0000,8.0000) circle (1.2pt);
\fill[red!80!black] (0.0000,8.0000) circle (1.2pt);
\fill[red!80!black] (12.0000,8.0000) circle (1.2pt);
\fill[red!80!black] (0.0000,8.0000) circle (1.2pt);
\fill[red!80!black] (12.0000,8.0000) circle (1.2pt);
\fill[red!80!black] (0.0000,8.0000) circle (1.2pt);
\fill[red!80!black] (12.0000,8.0000) circle (1.2pt);
\fill[red!80!black] (0.0000,8.0000) circle (1.2pt);
\fill[red!80!black] (12.0000,8.0000) circle (1.2pt);
\fill[red!80!black] (0.0000,8.0000) circle (1.2pt);
\fill[red!80!black] (12.0000,8.0000) circle (1.2pt);
\fill[red!80!black] (0.0000,8.0000) circle (1.2pt);
\fill[red!80!black] (12.0000,8.0000) circle (1.2pt);
\fill[red!80!black] (0.0000,8.0000) circle (1.2pt);
\fill[red!80!black] (12.0000,8.0000) circle (1.2pt);
\fill[red!80!black] (0.0000,8.0000) circle (1.2pt);
\fill[red!80!black] (12.0000,8.0000) circle (1.2pt);
\fill[red!80!black] (0.0000,8.0000) circle (1.2pt);
\fill[red!80!black] (12.0000,8.0000) circle (1.2pt);
\fill[red!80!black] (0.0000,8.0000) circle (1.2pt);
\fill[red!80!black] (12.0000,8.0000) circle (1.2pt);
\fill[red!80!black] (0.0000,8.0000) circle (1.2pt);
\fill[red!80!black] (12.0000,8.0000) circle (1.2pt);
\fill[red!80!black] (0.0000,8.0000) circle (1.2pt);
\fill[red!80!black] (12.0000,8.0000) circle (1.2pt);
\fill[red!80!black] (0.0000,8.0000) circle (1.2pt);
\fill[red!80!black] (12.0000,8.0000) circle (1.2pt);
\fill[red!80!black] (0.0000,8.0000) circle (1.2pt);
\fill[red!80!black] (12.0000,8.0000) circle (1.2pt);
\fill[red!80!black] (0.0000,8.0000) circle (1.2pt);
\fill[red!80!black] (12.0000,8.0000) circle (1.2pt);
\fill[red!80!black] (0.0000,8.0000) circle (1.2pt);
\fill[red!80!black] (12.0000,8.0000) circle (1.2pt);
\fill[red!80!black] (0.0000,8.0000) circle (1.2pt);
\fill[red!80!black] (12.0000,8.0000) circle (1.2pt);
\fill[red!80!black] (0.0000,8.0000) circle (1.2pt);
\fill[red!80!black] (12.0000,8.0000) circle (1.2pt);
\fill[red!80!black] (0.0000,8.0000) circle (1.2pt);
\fill[red!80!black] (12.0000,8.0000) circle (1.2pt);
\fill[red!80!black] (0.0000,8.0000) circle (1.2pt);
\fill[red!80!black] (12.0000,8.0000) circle (1.2pt);
\fill[red!80!black] (0.0000,8.0000) circle (1.2pt);
\fill[red!80!black] (12.0000,8.0000) circle (1.2pt);
\fill[red!80!black] (0.0000,8.0000) circle (1.2pt);
\fill[red!80!black] (12.0000,8.0000) circle (1.2pt);
\fill[red!80!black] (0.0000,8.0000) circle (1.2pt);
\fill[red!80!black] (12.0000,8.0000) circle (1.2pt);
\fill[red!80!black] (0.0000,8.0000) circle (1.2pt);
\fill[red!80!black] (12.0000,8.0000) circle (1.2pt);
\fill[red!80!black] (0.0000,8.0000) circle (1.2pt);
\fill[red!80!black] (12.0000,8.0000) circle (1.2pt);
\fill[red!80!black] (0.0000,8.0000) circle (1.2pt);
\fill[red!80!black] (12.0000,8.0000) circle (1.2pt);
\fill[red!80!black] (0.0000,8.0000) circle (1.2pt);
\fill[red!80!black] (12.0000,8.0000) circle (1.2pt);
\fill[red!80!black] (0.0000,8.0000) circle (1.2pt);
\fill[red!80!black] (12.0000,8.0000) circle (1.2pt);
\fill[red!80!black] (0.0000,8.0000) circle (1.2pt);
\fill[red!80!black] (12.0000,8.0000) circle (1.2pt);
\fill[red!80!black] (0.0000,8.0000) circle (1.2pt);
\fill[red!80!black] (12.0000,8.0000) circle (1.2pt);
\fill[red!80!black] (0.0000,8.0000) circle (1.2pt);
\fill[red!80!black] (12.0000,8.0000) circle (1.2pt);
\fill[red!80!black] (0.0000,8.0000) circle (1.2pt);
\fill[red!80!black] (12.0000,8.0000) circle (1.2pt);
\fill[red!80!black] (0.0000,8.0000) circle (1.2pt);
\fill[red!80!black] (12.0000,8.0000) circle (1.2pt);
\fill[red!80!black] (0.0000,8.0000) circle (1.2pt);
\fill[red!80!black] (12.0000,8.0000) circle (1.2pt);
\fill[red!80!black] (0.0000,8.0000) circle (1.2pt);
\fill[red!80!black] (12.0000,8.0000) circle (1.2pt);
\fill[red!80!black] (0.0000,8.0000) circle (1.2pt);
\fill[red!80!black] (12.0000,8.0000) circle (1.2pt);
\fill[red!80!black] (0.0000,8.0000) circle (1.2pt);
\fill[red!80!black] (12.0000,8.0000) circle (1.2pt);
\fill[red!80!black] (0.0000,8.0000) circle (1.2pt);
\fill[red!80!black] (12.0000,8.0000) circle (1.2pt);
\fill[red!80!black] (0.0000,8.0000) circle (1.2pt);
\fill[red!80!black] (12.0000,8.0000) circle (1.2pt);
\fill[red!80!black] (0.0000,8.0000) circle (1.2pt);
\fill[red!80!black] (12.0000,8.0000) circle (1.2pt);
\fill[red!80!black] (0.0000,8.0000) circle (1.2pt);
\fill[red!80!black] (12.0000,8.0000) circle (1.2pt);
\fill[red!80!black] (0.0000,8.0000) circle (1.2pt);
\fill[red!80!black] (12.0000,8.0000) circle (1.2pt);
\fill[red!80!black] (0.0000,8.0000) circle (1.2pt);
\fill[red!80!black] (12.0000,8.0000) circle (1.2pt);
\fill[red!80!black] (0.0000,8.0000) circle (1.2pt);
\fill[red!80!black] (12.0000,8.0000) circle (1.2pt);
\fill[red!80!black] (0.0000,8.0000) circle (1.2pt);
\fill[red!80!black] (12.0000,8.0000) circle (1.2pt);
\fill[red!80!black] (0.0000,8.0000) circle (1.2pt);
\fill[red!80!black] (12.0000,8.0000) circle (1.2pt);
\fill[red!80!black] (0.0000,8.0000) circle (1.2pt);
\fill[red!80!black] (12.0000,8.0000) circle (1.2pt);
\fill[red!80!black] (0.0000,8.0000) circle (1.2pt);
\fill[red!80!black] (12.0000,8.0000) circle (1.2pt);
\fill[red!80!black] (0.0000,8.0000) circle (1.2pt);
\fill[red!80!black] (12.0000,8.0000) circle (1.2pt);
\fill[red!80!black] (0.0000,8.0000) circle (1.2pt);
\fill[red!80!black] (12.0000,8.0000) circle (1.2pt);
\fill[red!80!black] (0.0000,8.0000) circle (1.2pt);
\fill[red!80!black] (12.0000,8.0000) circle (1.2pt);
\fill[red!80!black] (0.0000,8.0000) circle (1.2pt);
\fill[red!80!black] (12.0000,8.0000) circle (1.2pt);
\fill[red!80!black] (0.0000,8.0000) circle (1.2pt);
\fill[red!80!black] (12.0000,8.0000) circle (1.2pt);
\fill[red!80!black] (0.0000,8.0000) circle (1.2pt);
\fill[red!80!black] (12.0000,8.0000) circle (1.2pt);
\fill[red!80!black] (0.0000,8.0000) circle (1.2pt);
\fill[red!80!black] (12.0000,8.0000) circle (1.2pt);
\fill[red!80!black] (0.0000,8.0000) circle (1.2pt);
\fill[red!80!black] (12.0000,8.0000) circle (1.2pt);
\fill[red!80!black] (0.0000,8.0000) circle (1.2pt);
\fill[red!80!black] (12.0000,8.0000) circle (1.2pt);
\fill[red!80!black] (0.0000,8.0000) circle (1.2pt);
\fill[red!80!black] (12.0000,8.0000) circle (1.2pt);
\fill[red!80!black] (0.0000,8.0000) circle (1.2pt);
\fill[red!80!black] (12.0000,8.0000) circle (1.2pt);
\fill[red!80!black] (0.0000,8.0000) circle (1.2pt);
\fill[red!80!black] (12.0000,8.0000) circle (1.2pt);
\fill[red!80!black] (0.0000,8.0000) circle (1.2pt);
\fill[red!80!black] (12.0000,8.0000) circle (1.2pt);
\fill[red!80!black] (0.0000,8.0000) circle (1.2pt);
\fill[red!80!black] (12.0000,8.0000) circle (1.2pt);
\fill[red!80!black] (0.0000,8.0000) circle (1.2pt);
\fill[red!80!black] (12.0000,8.0000) circle (1.2pt);
\fill[red!80!black] (0.0000,8.0000) circle (1.2pt);
\fill[red!80!black] (12.0000,8.0000) circle (1.2pt);
\fill[red!80!black] (0.0000,8.0000) circle (1.2pt);
\fill[red!80!black] (12.0000,8.0000) circle (1.2pt);
\fill[red!80!black] (0.0000,8.0000) circle (1.2pt);
\fill[red!80!black] (12.0000,8.0000) circle (1.2pt);
\fill[red!80!black] (0.0000,8.0000) circle (1.2pt);
\fill[red!80!black] (12.0000,8.0000) circle (1.2pt);
\fill[red!80!black] (0.0000,8.0000) circle (1.2pt);
\fill[red!80!black] (12.0000,8.0000) circle (1.2pt);
\fill[red!80!black] (0.0000,8.0000) circle (1.2pt);
\fill[red!80!black] (12.0000,8.0000) circle (1.2pt);
\fill[red!80!black] (0.0000,8.0000) circle (1.2pt);
\fill[red!80!black] (12.0000,8.0000) circle (1.2pt);
\fill[red!80!black] (0.0000,8.0000) circle (1.2pt);
\fill[red!80!black] (12.0000,8.0000) circle (1.2pt);
\fill[red!80!black] (0.0000,8.0000) circle (1.2pt);
\fill[red!80!black] (12.0000,8.0000) circle (1.2pt);
\fill[red!80!black] (0.0000,8.0000) circle (1.2pt);
\fill[red!80!black] (12.0000,8.0000) circle (1.2pt);
\fill[red!80!black] (0.0000,8.0000) circle (1.2pt);
\fill[red!80!black] (12.0000,8.0000) circle (1.2pt);
\fill[red!80!black] (0.0000,8.0000) circle (1.2pt);
\fill[red!80!black] (12.0000,8.0000) circle (1.2pt);
\fill[red!80!black] (0.0000,8.0000) circle (1.2pt);
\fill[red!80!black] (12.0000,8.0000) circle (1.2pt);
\fill[red!80!black] (0.0000,8.0000) circle (1.2pt);
\fill[red!80!black] (12.0000,8.0000) circle (1.2pt);
\fill[red!80!black] (0.0000,8.0000) circle (1.2pt);
\fill[red!80!black] (12.0000,8.0000) circle (1.2pt);
\fill[red!80!black] (0.0000,8.0000) circle (1.2pt);
\fill[red!80!black] (12.0000,8.0000) circle (1.2pt);
\fill[red!80!black] (0.0000,8.0000) circle (1.2pt);
\fill[red!80!black] (12.0000,8.0000) circle (1.2pt);
\fill[red!80!black] (0.0000,8.0000) circle (1.2pt);
\fill[red!80!black] (12.0000,8.0000) circle (1.2pt);
\fill[red!80!black] (0.0000,8.0000) circle (1.2pt);
\fill[red!80!black] (12.0000,8.0000) circle (1.2pt);
\fill[red!80!black] (0.0000,8.0000) circle (1.2pt);
\fill[red!80!black] (12.0000,8.0000) circle (1.2pt);
\fill[red!80!black] (0.0000,8.0000) circle (1.2pt);
\fill[red!80!black] (12.0000,8.0000) circle (1.2pt);
\fill[red!80!black] (0.0000,8.0000) circle (1.2pt);
\fill[red!80!black] (12.0000,8.0000) circle (1.2pt);
\fill[red!80!black] (0.0000,8.0000) circle (1.2pt);
\fill[red!80!black] (12.0000,8.0000) circle (1.2pt);
\fill[red!80!black] (0.0000,8.0000) circle (1.2pt);
\fill[red!80!black] (12.0000,8.0000) circle (1.2pt);
\fill[red!80!black] (0.0000,8.0000) circle (1.2pt);
\fill[red!80!black] (12.0000,8.0000) circle (1.2pt);
\fill[red!80!black] (0.0000,8.0000) circle (1.2pt);
\fill[red!80!black] (12.0000,8.0000) circle (1.2pt);
\fill[red!80!black] (0.0000,8.0000) circle (1.2pt);
\fill[red!80!black] (12.0000,8.0000) circle (1.2pt);
\fill[red!80!black] (0.0000,8.0000) circle (1.2pt);
\fill[red!80!black] (12.0000,8.0000) circle (1.2pt);
\fill[red!80!black] (0.0000,8.0000) circle (1.2pt);
\fill[red!80!black] (12.0000,8.0000) circle (1.2pt);
\fill[red!80!black] (0.0000,8.0000) circle (1.2pt);
\fill[red!80!black] (12.0000,8.0000) circle (1.2pt);
\fill[red!80!black] (0.0000,8.0000) circle (1.2pt);
\fill[red!80!black] (12.0000,8.0000) circle (1.2pt);
\fill[red!80!black] (0.0000,8.0000) circle (1.2pt);
\fill[red!80!black] (12.0000,8.0000) circle (1.2pt);
\fill[red!80!black] (0.0000,8.0000) circle (1.2pt);
\fill[red!80!black] (12.0000,8.0000) circle (1.2pt);
\fill[red!80!black] (0.0000,8.0000) circle (1.2pt);
\fill[red!80!black] (12.0000,8.0000) circle (1.2pt);
\fill[red!80!black] (0.0000,8.0000) circle (1.2pt);
\fill[red!80!black] (12.0000,8.0000) circle (1.2pt);
\fill[red!80!black] (0.0000,8.0000) circle (1.2pt);
\fill[red!80!black] (12.0000,8.0000) circle (1.2pt);
\fill[red!80!black] (0.0000,8.0000) circle (1.2pt);
\fill[red!80!black] (12.0000,8.0000) circle (1.2pt);
\fill[red!80!black] (0.0000,8.0000) circle (1.2pt);
\fill[red!80!black] (12.0000,8.0000) circle (1.2pt);
\fill[red!80!black] (0.0000,8.0000) circle (1.2pt);
\fill[red!80!black] (12.0000,8.0000) circle (1.2pt);
\fill[red!80!black] (0.0000,8.0000) circle (1.2pt);
\fill[red!80!black] (12.0000,8.0000) circle (1.2pt);
\fill[red!80!black] (0.0000,8.0000) circle (1.2pt);
\fill[red!80!black] (12.0000,8.0000) circle (1.2pt);
\fill[red!80!black] (0.0000,8.0000) circle (1.2pt);
\fill[red!80!black] (12.0000,8.0000) circle (1.2pt);
\fill[red!80!black] (0.0000,8.0000) circle (1.2pt);
\fill[red!80!black] (12.0000,8.0000) circle (1.2pt);
\fill[red!80!black] (0.0000,8.0000) circle (1.2pt);
\fill[red!80!black] (12.0000,8.0000) circle (1.2pt);
\fill[red!80!black] (0.0000,8.0000) circle (1.2pt);
\fill[red!80!black] (12.0000,8.0000) circle (1.2pt);
\fill[red!80!black] (0.0000,8.0000) circle (1.2pt);
\fill[red!80!black] (12.0000,8.0000) circle (1.2pt);
\fill[red!80!black] (0.0000,8.0000) circle (1.2pt);
\fill[red!80!black] (12.0000,8.0000) circle (1.2pt);
\fill[red!80!black] (0.0000,8.0000) circle (1.2pt);
\fill[red!80!black] (12.0000,8.0000) circle (1.2pt);
\fill[red!80!black] (0.0000,8.0000) circle (1.2pt);
\fill[red!80!black] (12.0000,8.0000) circle (1.2pt);
\fill[red!80!black] (0.0000,8.0000) circle (1.2pt);
\fill[red!80!black] (12.0000,8.0000) circle (1.2pt);
\fill[red!80!black] (0.0000,8.0000) circle (1.2pt);
\fill[red!80!black] (12.0000,8.0000) circle (1.2pt);
\fill[red!80!black] (0.0000,8.0000) circle (1.2pt);
\fill[red!80!black] (12.0000,8.0000) circle (1.2pt);
\fill[red!80!black] (0.0000,8.0000) circle (1.2pt);
\fill[red!80!black] (12.0000,8.0000) circle (1.2pt);
\fill[red!80!black] (0.0000,8.0000) circle (1.2pt);
\fill[red!80!black] (12.0000,8.0000) circle (1.2pt);
\fill[red!80!black] (0.0000,8.0000) circle (1.2pt);
\fill[red!80!black] (12.0000,8.0000) circle (1.2pt);
\fill[red!80!black] (0.0000,8.0000) circle (1.2pt);
\fill[red!80!black] (12.0000,8.0000) circle (1.2pt);
\fill[red!80!black] (0.0000,8.0000) circle (1.2pt);
\fill[red!80!black] (12.0000,8.0000) circle (1.2pt);
\fill[red!80!black] (0.0000,8.0000) circle (1.2pt);
\fill[red!80!black] (12.0000,8.0000) circle (1.2pt);
\fill[red!80!black] (0.0000,8.0000) circle (1.2pt);
\fill[red!80!black] (12.0000,8.0000) circle (1.2pt);
\fill[red!80!black] (0.0000,8.0000) circle (1.2pt);
\fill[red!80!black] (12.0000,8.0000) circle (1.2pt);
\fill[red!80!black] (0.0000,8.0000) circle (1.2pt);
\fill[red!80!black] (12.0000,8.0000) circle (1.2pt);
\fill[red!80!black] (0.0000,8.0000) circle (1.2pt);
\fill[red!80!black] (12.0000,8.0000) circle (1.2pt);
\fill[red!80!black] (0.0000,8.0000) circle (1.2pt);
\fill[red!80!black] (12.0000,8.0000) circle (1.2pt);
\fill[red!80!black] (0.0000,8.0000) circle (1.2pt);
\fill[red!80!black] (12.0000,8.0000) circle (1.2pt);
\fill[red!80!black] (0.0000,8.0000) circle (1.2pt);
\fill[red!80!black] (12.0000,8.0000) circle (1.2pt);
\fill[red!80!black] (0.0000,8.0000) circle (1.2pt);
\fill[red!80!black] (12.0000,8.0000) circle (1.2pt);
\fill[red!80!black] (0.0000,8.0000) circle (1.2pt);
\fill[red!80!black] (12.0000,8.0000) circle (1.2pt);
\fill[red!80!black] (0.0000,8.0000) circle (1.2pt);
\fill[red!80!black] (12.0000,8.0000) circle (1.2pt);
\fill[red!80!black] (0.0000,8.0000) circle (1.2pt);
\fill[red!80!black] (12.0000,8.0000) circle (1.2pt);
\fill[red!80!black] (0.0000,8.0000) circle (1.2pt);
\fill[red!80!black] (12.0000,8.0000) circle (1.2pt);
\fill[red!80!black] (0.0000,8.0000) circle (1.2pt);
\fill[red!80!black] (12.0000,8.0000) circle (1.2pt);
\fill[red!80!black] (0.0000,8.0000) circle (1.2pt);
\fill[red!80!black] (12.0000,8.0000) circle (1.2pt);
\fill[red!80!black] (0.0000,8.0000) circle (1.2pt);
\fill[red!80!black] (12.0000,8.0000) circle (1.2pt);
\fill[red!80!black] (0.0000,8.0000) circle (1.2pt);
\fill[red!80!black] (12.0000,8.0000) circle (1.2pt);
\fill[red!80!black] (0.0000,8.0000) circle (1.2pt);
\fill[red!80!black] (12.0000,8.0000) circle (1.2pt);
\fill[red!80!black] (0.0000,8.0000) circle (1.2pt);
\fill[red!80!black] (12.0000,8.0000) circle (1.2pt);
\fill[red!80!black] (0.0000,8.0000) circle (1.2pt);
\fill[red!80!black] (12.0000,8.0000) circle (1.2pt);
\fill[red!80!black] (0.0000,8.0000) circle (1.2pt);
\fill[red!80!black] (12.0000,8.0000) circle (1.2pt);
\fill[red!80!black] (0.0000,8.0000) circle (1.2pt);
\fill[red!80!black] (12.0000,8.0000) circle (1.2pt);
\fill[red!80!black] (0.0000,8.0000) circle (1.2pt);
\fill[red!80!black] (12.0000,8.0000) circle (1.2pt);
\fill[red!80!black] (0.0000,8.0000) circle (1.2pt);
\fill[red!80!black] (12.0000,8.0000) circle (1.2pt);
\fill[red!80!black] (0.0000,8.0000) circle (1.2pt);
\fill[red!80!black] (12.0000,8.0000) circle (1.2pt);
\fill[red!80!black] (0.0000,8.0000) circle (1.2pt);
\fill[red!80!black] (12.0000,8.0000) circle (1.2pt);
\fill[red!80!black] (0.0000,8.0000) circle (1.2pt);
\fill[red!80!black] (12.0000,8.0000) circle (1.2pt);
\fill[red!80!black] (0.0000,8.0000) circle (1.2pt);
\fill[red!80!black] (12.0000,8.0000) circle (1.2pt);
\fill[red!80!black] (0.0000,8.0000) circle (1.2pt);
\fill[red!80!black] (12.0000,8.0000) circle (1.2pt);
\fill[red!80!black] (0.0000,8.0000) circle (1.2pt);
\fill[red!80!black] (12.0000,8.0000) circle (1.2pt);
\fill[red!80!black] (0.0000,8.0000) circle (1.2pt);
\fill[red!80!black] (12.0000,8.0000) circle (1.2pt);
\fill[red!80!black] (0.0000,8.0000) circle (1.2pt);
\fill[red!80!black] (12.0000,8.0000) circle (1.2pt);
\fill[red!80!black] (0.0000,8.0000) circle (1.2pt);
\fill[red!80!black] (12.0000,8.0000) circle (1.2pt);
\fill[red!80!black] (0.0000,8.0000) circle (1.2pt);
\fill[red!80!black] (12.0000,8.0000) circle (1.2pt);
\fill[red!80!black] (0.0000,8.0000) circle (1.2pt);
\fill[red!80!black] (12.0000,8.0000) circle (1.2pt);
\fill[red!80!black] (0.0000,8.0000) circle (1.2pt);
\fill[red!80!black] (12.0000,8.0000) circle (1.2pt);
\fill[red!80!black] (0.0000,8.0000) circle (1.2pt);
\fill[red!80!black] (12.0000,8.0000) circle (1.2pt);
\fill[red!80!black] (0.0000,8.0000) circle (1.2pt);
\fill[red!80!black] (12.0000,8.0000) circle (1.2pt);
\fill[red!80!black] (0.0000,8.0000) circle (1.2pt);
\fill[red!80!black] (12.0000,8.0000) circle (1.2pt);
\fill[red!80!black] (0.0000,8.0000) circle (1.2pt);
\fill[red!80!black] (12.0000,8.0000) circle (1.2pt);
\fill[red!80!black] (0.0000,8.0000) circle (1.2pt);
\fill[red!80!black] (12.0000,8.0000) circle (1.2pt);
\fill[red!80!black] (0.0000,8.0000) circle (1.2pt);
\fill[red!80!black] (12.0000,8.0000) circle (1.2pt);
\fill[red!80!black] (0.0000,8.0000) circle (1.2pt);
\fill[red!80!black] (12.0000,8.0000) circle (1.2pt);
\fill[red!80!black] (0.0000,8.0000) circle (1.2pt);
\fill[red!80!black] (12.0000,8.0000) circle (1.2pt);
\fill[red!80!black] (0.0000,8.0000) circle (1.2pt);
\fill[red!80!black] (12.0000,8.0000) circle (1.2pt);
\fill[red!80!black] (0.0000,8.0000) circle (1.2pt);
\fill[red!80!black] (12.0000,8.0000) circle (1.2pt);
\fill[red!80!black] (0.0000,8.0000) circle (1.2pt);
\fill[red!80!black] (12.0000,8.0000) circle (1.2pt);
\fill[red!80!black] (0.0000,8.0000) circle (1.2pt);
\fill[red!80!black] (12.0000,8.0000) circle (1.2pt);
\fill[red!80!black] (0.0000,8.0000) circle (1.2pt);
\fill[red!80!black] (12.0000,8.0000) circle (1.2pt);
\fill[red!80!black] (0.0000,8.0000) circle (1.2pt);
\fill[red!80!black] (12.0000,8.0000) circle (1.2pt);
\fill[red!80!black] (0.0000,8.0000) circle (1.2pt);
\fill[red!80!black] (12.0000,8.0000) circle (1.2pt);
\fill[red!80!black] (0.0000,8.0000) circle (1.2pt);
\fill[red!80!black] (12.0000,8.0000) circle (1.2pt);
\fill[red!80!black] (0.0000,8.0000) circle (1.2pt);
\fill[red!80!black] (12.0000,8.0000) circle (1.2pt);
\fill[red!80!black] (0.0000,8.0000) circle (1.2pt);
\fill[red!80!black] (12.0000,8.0000) circle (1.2pt);
\fill[red!80!black] (0.0000,8.0000) circle (1.2pt);
\fill[red!80!black] (12.0000,8.0000) circle (1.2pt);
\fill[red!80!black] (0.0000,8.0000) circle (1.2pt);
\fill[red!80!black] (12.0000,8.0000) circle (1.2pt);
\fill[red!80!black] (0.0000,8.0000) circle (1.2pt);
\fill[red!80!black] (12.0000,8.0000) circle (1.2pt);
\fill[red!80!black] (0.0000,8.0000) circle (1.2pt);
\fill[red!80!black] (12.0000,8.0000) circle (1.2pt);
\fill[red!80!black] (0.0000,8.0000) circle (1.2pt);
\fill[red!80!black] (12.0000,8.0000) circle (1.2pt);
\fill[red!80!black] (0.0000,8.0000) circle (1.2pt);
\fill[red!80!black] (12.0000,8.0000) circle (1.2pt);
\fill[red!80!black] (0.0000,8.0000) circle (1.2pt);
\fill[red!80!black] (12.0000,8.0000) circle (1.2pt);
\fill[red!80!black] (0.0000,8.0000) circle (1.2pt);
\fill[red!80!black] (12.0000,8.0000) circle (1.2pt);
\fill[red!80!black] (0.0000,8.0000) circle (1.2pt);
\fill[red!80!black] (12.0000,8.0000) circle (1.2pt);
\fill[red!80!black] (0.0000,8.0000) circle (1.2pt);
\fill[red!80!black] (12.0000,8.0000) circle (1.2pt);
\fill[red!80!black] (0.0000,8.0000) circle (1.2pt);
\fill[red!80!black] (12.0000,8.0000) circle (1.2pt);
\fill[red!80!black] (0.0000,8.0000) circle (1.2pt);
\fill[red!80!black] (12.0000,8.0000) circle (1.2pt);
\fill[red!80!black] (0.0000,8.0000) circle (1.2pt);
\fill[red!80!black] (12.0000,8.0000) circle (1.2pt);
\fill[red!80!black] (0.0000,8.0000) circle (1.2pt);
\fill[red!80!black] (12.0000,8.0000) circle (1.2pt);
\fill[red!80!black] (0.0000,8.0000) circle (1.2pt);
\fill[red!80!black] (12.0000,8.0000) circle (1.2pt);
\fill[red!80!black] (0.0000,8.0000) circle (1.2pt);
\fill[red!80!black] (12.0000,8.0000) circle (1.2pt);
\fill[red!80!black] (0.0000,8.0000) circle (1.2pt);
\fill[red!80!black] (12.0000,8.0000) circle (1.2pt);
\fill[red!80!black] (0.0000,8.0000) circle (1.2pt);
\fill[red!80!black] (12.0000,8.0000) circle (1.2pt);
\fill[red!80!black] (0.0000,8.0000) circle (1.2pt);
\fill[red!80!black] (12.0000,8.0000) circle (1.2pt);
\fill[red!80!black] (0.0000,8.0000) circle (1.2pt);
\fill[red!80!black] (12.0000,8.0000) circle (1.2pt);
\fill[red!80!black] (0.0000,8.0000) circle (1.2pt);
\fill[red!80!black] (12.0000,8.0000) circle (1.2pt);
\fill[red!80!black] (0.0000,8.0000) circle (1.2pt);
\fill[red!80!black] (12.0000,8.0000) circle (1.2pt);
\fill[red!80!black] (0.0000,8.0000) circle (1.2pt);
\fill[red!80!black] (12.0000,8.0000) circle (1.2pt);
\fill[red!80!black] (0.0000,8.0000) circle (1.2pt);
\fill[red!80!black] (12.0000,8.0000) circle (1.2pt);
\fill[red!80!black] (0.0000,8.0000) circle (1.2pt);
\fill[red!80!black] (12.0000,8.0000) circle (1.2pt);
\fill[red!80!black] (0.0000,8.0000) circle (1.2pt);
\fill[red!80!black] (12.0000,8.0000) circle (1.2pt);
\fill[red!80!black] (0.0000,8.0000) circle (1.2pt);
\fill[red!80!black] (12.0000,8.0000) circle (1.2pt);
\fill[red!80!black] (0.0000,8.0000) circle (1.2pt);
\fill[red!80!black] (12.0000,8.0000) circle (1.2pt);
\fill[red!80!black] (0.0000,8.0000) circle (1.2pt);
\fill[red!80!black] (12.0000,8.0000) circle (1.2pt);
\fill[red!80!black] (0.0000,8.0000) circle (1.2pt);
\fill[red!80!black] (12.0000,8.0000) circle (1.2pt);
\fill[red!80!black] (0.0000,8.0000) circle (1.2pt);
\fill[red!80!black] (12.0000,8.0000) circle (1.2pt);
\fill[red!80!black] (0.0000,8.0000) circle (1.2pt);
\fill[red!80!black] (12.0000,8.0000) circle (1.2pt);
\fill[red!80!black] (0.0000,8.0000) circle (1.2pt);
\fill[red!80!black] (12.0000,8.0000) circle (1.2pt);
\fill[red!80!black] (0.0000,8.0000) circle (1.2pt);
\fill[red!80!black] (12.0000,8.0000) circle (1.2pt);
\fill[red!80!black] (0.0000,8.0000) circle (1.2pt);
\fill[red!80!black] (12.0000,8.0000) circle (1.2pt);
\fill[red!80!black] (0.0000,8.0000) circle (1.2pt);
\fill[red!80!black] (12.0000,8.0000) circle (1.2pt);
\fill[red!80!black] (0.0000,8.0000) circle (1.2pt);
\fill[red!80!black] (12.0000,8.0000) circle (1.2pt);
\fill[red!80!black] (0.0000,8.0000) circle (1.2pt);
\fill[red!80!black] (12.0000,8.0000) circle (1.2pt);
\fill[red!80!black] (0.0000,8.0000) circle (1.2pt);
\fill[red!80!black] (12.0000,8.0000) circle (1.2pt);
\fill[red!80!black] (0.0000,8.0000) circle (1.2pt);
\fill[red!80!black] (12.0000,8.0000) circle (1.2pt);
\fill[red!80!black] (0.0000,8.0000) circle (1.2pt);
\fill[red!80!black] (12.0000,8.0000) circle (1.2pt);
\fill[red!80!black] (0.0000,8.0000) circle (1.2pt);
\fill[red!80!black] (12.0000,8.0000) circle (1.2pt);
\fill[red!80!black] (0.0000,8.0000) circle (1.2pt);
\fill[red!80!black] (12.0000,8.0000) circle (1.2pt);
\fill[red!80!black] (0.0000,8.0000) circle (1.2pt);
\fill[red!80!black] (12.0000,8.0000) circle (1.2pt);
\fill[red!80!black] (0.0000,8.0000) circle (1.2pt);
\fill[red!80!black] (12.0000,8.0000) circle (1.2pt);
\fill[red!80!black] (0.0000,8.0000) circle (1.2pt);
\fill[red!80!black] (12.0000,8.0000) circle (1.2pt);
\fill[red!80!black] (0.0000,8.0000) circle (1.2pt);
\fill[red!80!black] (12.0000,8.0000) circle (1.2pt);
\fill[red!80!black] (0.0000,8.0000) circle (1.2pt);
\fill[red!80!black] (12.0000,8.0000) circle (1.2pt);
\fill[red!80!black] (0.0000,8.0000) circle (1.2pt);
\fill[red!80!black] (12.0000,8.0000) circle (1.2pt);
\fill[red!80!black] (0.0000,8.0000) circle (1.2pt);
\fill[red!80!black] (12.0000,8.0000) circle (1.2pt);
\fill[red!80!black] (0.0000,8.0000) circle (1.2pt);
\fill[red!80!black] (12.0000,8.0000) circle (1.2pt);
\fill[red!80!black] (0.0000,8.0000) circle (1.2pt);
\fill[red!80!black] (12.0000,8.0000) circle (1.2pt);
\fill[red!80!black] (0.0000,8.0000) circle (1.2pt);
\fill[red!80!black] (12.0000,8.0000) circle (1.2pt);
\fill[red!80!black] (0.0000,8.0000) circle (1.2pt);
\fill[red!80!black] (12.0000,8.0000) circle (1.2pt);
\fill[red!80!black] (0.0000,8.0000) circle (1.2pt);
\fill[red!80!black] (12.0000,8.0000) circle (1.2pt);
\fill[red!80!black] (0.0000,8.0000) circle (1.2pt);
\fill[red!80!black] (12.0000,8.0000) circle (1.2pt);
\fill[red!80!black] (0.0000,8.0000) circle (1.2pt);
\fill[red!80!black] (12.0000,8.0000) circle (1.2pt);
\fill[red!80!black] (0.0000,8.0000) circle (1.2pt);
\fill[red!80!black] (12.0000,8.0000) circle (1.2pt);
\fill[red!80!black] (0.0000,8.0000) circle (1.2pt);
\fill[red!80!black] (12.0000,8.0000) circle (1.2pt);
\fill[red!80!black] (0.0000,8.0000) circle (1.2pt);
\fill[red!80!black] (12.0000,8.0000) circle (1.2pt);
\fill[red!80!black] (0.0000,8.0000) circle (1.2pt);
\fill[red!80!black] (12.0000,8.0000) circle (1.2pt);
\fill[red!80!black] (0.0000,8.0000) circle (1.2pt);
\fill[red!80!black] (12.0000,8.0000) circle (1.2pt);
\fill[red!80!black] (0.0000,8.0000) circle (1.2pt);
\fill[red!80!black] (12.0000,8.0000) circle (1.2pt);
\fill[red!80!black] (0.0000,8.0000) circle (1.2pt);
\fill[red!80!black] (12.0000,8.0000) circle (1.2pt);
\fill[red!80!black] (0.0000,8.0000) circle (1.2pt);
\fill[red!80!black] (12.0000,8.0000) circle (1.2pt);
\fill[red!80!black] (0.0000,8.0000) circle (1.2pt);
\fill[red!80!black] (12.0000,8.0000) circle (1.2pt);
\fill[red!80!black] (0.0000,8.0000) circle (1.2pt);
\fill[red!80!black] (12.0000,8.0000) circle (1.2pt);
\fill[red!80!black] (0.0000,8.0000) circle (1.2pt);
\fill[red!80!black] (12.0000,8.0000) circle (1.2pt);
\fill[red!80!black] (0.0000,8.0000) circle (1.2pt);
\fill[red!80!black] (12.0000,8.0000) circle (1.2pt);
\fill[red!80!black] (0.0000,8.0000) circle (1.2pt);
\fill[red!80!black] (12.0000,8.0000) circle (1.2pt);
\fill[red!80!black] (0.0000,8.0000) circle (1.2pt);
\fill[red!80!black] (12.0000,8.0000) circle (1.2pt);
\fill[red!80!black] (0.0000,8.0000) circle (1.2pt);
\fill[red!80!black] (12.0000,8.0000) circle (1.2pt);
\fill[red!80!black] (0.0000,8.0000) circle (1.2pt);
\fill[red!80!black] (12.0000,8.0000) circle (1.2pt);
\fill[red!80!black] (0.0000,8.0000) circle (1.2pt);
\fill[red!80!black] (12.0000,8.0000) circle (1.2pt);
\fill[red!80!black] (0.0000,8.0000) circle (1.2pt);
\fill[red!80!black] (12.0000,8.0000) circle (1.2pt);
\fill[red!80!black] (0.0000,8.0000) circle (1.2pt);
\fill[red!80!black] (12.0000,8.0000) circle (1.2pt);
\fill[red!80!black] (0.0000,8.0000) circle (1.2pt);
\fill[red!80!black] (12.0000,8.0000) circle (1.2pt);
\fill[red!80!black] (0.0000,8.0000) circle (1.2pt);
\fill[red!80!black] (12.0000,8.0000) circle (1.2pt);
\fill[red!80!black] (0.0000,8.0000) circle (1.2pt);
\fill[red!80!black] (12.0000,8.0000) circle (1.2pt);
\fill[red!80!black] (0.0000,8.0000) circle (1.2pt);
\fill[red!80!black] (12.0000,8.0000) circle (1.2pt);
\fill[red!80!black] (0.0000,8.0000) circle (1.2pt);
\fill[red!80!black] (12.0000,8.0000) circle (1.2pt);
\fill[red!80!black] (0.0000,8.0000) circle (1.2pt);
\fill[red!80!black] (12.0000,8.0000) circle (1.2pt);
\fill[red!80!black] (0.0000,8.0000) circle (1.2pt);
\fill[red!80!black] (12.0000,8.0000) circle (1.2pt);
\fill[red!80!black] (0.0000,8.0000) circle (1.2pt);
\fill[red!80!black] (12.0000,8.0000) circle (1.2pt);
\fill[red!80!black] (0.0000,8.0000) circle (1.2pt);
\fill[red!80!black] (12.0000,8.0000) circle (1.2pt);
\fill[red!80!black] (0.0000,8.0000) circle (1.2pt);
\fill[red!80!black] (12.0000,8.0000) circle (1.2pt);
\fill[red!80!black] (0.0000,8.0000) circle (1.2pt);
\fill[red!80!black] (12.0000,8.0000) circle (1.2pt);
\fill[red!80!black] (0.0000,8.0000) circle (1.2pt);
\fill[red!80!black] (12.0000,8.0000) circle (1.2pt);
\fill[red!80!black] (0.0000,8.0000) circle (1.2pt);
\fill[red!80!black] (12.0000,8.0000) circle (1.2pt);
\fill[red!80!black] (0.0000,8.0000) circle (1.2pt);
\fill[red!80!black] (12.0000,8.0000) circle (1.2pt);
\fill[red!80!black] (0.0000,8.0000) circle (1.2pt);
\fill[red!80!black] (12.0000,8.0000) circle (1.2pt);
\fill[red!80!black] (0.0000,8.0000) circle (1.2pt);
\fill[red!80!black] (12.0000,8.0000) circle (1.2pt);
\fill[red!80!black] (0.0000,8.0000) circle (1.2pt);
\fill[red!80!black] (12.0000,8.0000) circle (1.2pt);
\fill[red!80!black] (0.0000,8.0000) circle (1.2pt);
\fill[red!80!black] (12.0000,8.0000) circle (1.2pt);
\fill[red!80!black] (0.0000,8.0000) circle (1.2pt);
\fill[red!80!black] (12.0000,8.0000) circle (1.2pt);
\fill[red!80!black] (0.0000,8.0000) circle (1.2pt);
\fill[red!80!black] (12.0000,8.0000) circle (1.2pt);
\fill[red!80!black] (0.0000,8.0000) circle (1.2pt);
\fill[red!80!black] (12.0000,8.0000) circle (1.2pt);
\fill[red!80!black] (0.0000,8.0000) circle (1.2pt);
\fill[red!80!black] (12.0000,8.0000) circle (1.2pt);
\fill[red!80!black] (0.0000,8.0000) circle (1.2pt);
\fill[red!80!black] (12.0000,8.0000) circle (1.2pt);
\fill[red!80!black] (0.0000,8.0000) circle (1.2pt);
\fill[red!80!black] (12.0000,8.0000) circle (1.2pt);
\fill[red!80!black] (0.0000,8.0000) circle (1.2pt);
\fill[red!80!black] (12.0000,8.0000) circle (1.2pt);
\fill[red!80!black] (0.0000,8.0000) circle (1.2pt);
\fill[red!80!black] (12.0000,8.0000) circle (1.2pt);
\fill[red!80!black] (0.0000,8.0000) circle (1.2pt);
\fill[red!80!black] (12.0000,8.0000) circle (1.2pt);
\fill[red!80!black] (0.0000,8.0000) circle (1.2pt);
\fill[red!80!black] (12.0000,8.0000) circle (1.2pt);
\fill[red!80!black] (0.0000,8.0000) circle (1.2pt);
\fill[red!80!black] (12.0000,8.0000) circle (1.2pt);
\fill[red!80!black] (0.0000,8.0000) circle (1.2pt);
\fill[red!80!black] (12.0000,8.0000) circle (1.2pt);
\fill[red!80!black] (0.0000,8.0000) circle (1.2pt);
\fill[red!80!black] (12.0000,8.0000) circle (1.2pt);
\fill[red!80!black] (0.0000,8.0000) circle (1.2pt);
\fill[red!80!black] (12.0000,8.0000) circle (1.2pt);
\fill[red!80!black] (0.0000,8.0000) circle (1.2pt);
\fill[red!80!black] (12.0000,8.0000) circle (1.2pt);
\fill[red!80!black] (0.0000,8.0000) circle (1.2pt);
\fill[red!80!black] (12.0000,8.0000) circle (1.2pt);
\fill[red!80!black] (0.0000,8.0000) circle (1.2pt);
\fill[red!80!black] (12.0000,8.0000) circle (1.2pt);
\fill[red!80!black] (0.0000,8.0000) circle (1.2pt);
\fill[red!80!black] (12.0000,8.0000) circle (1.2pt);
\fill[red!80!black] (0.0000,8.0000) circle (1.2pt);
\fill[red!80!black] (12.0000,8.0000) circle (1.2pt);
\fill[red!80!black] (0.0000,8.0000) circle (1.2pt);
\fill[red!80!black] (12.0000,8.0000) circle (1.2pt);
\fill[red!80!black] (0.0000,8.0000) circle (1.2pt);
\fill[red!80!black] (12.0000,8.0000) circle (1.2pt);
\fill[red!80!black] (0.0000,8.0000) circle (1.2pt);
\fill[red!80!black] (12.0000,8.0000) circle (1.2pt);
\fill[red!80!black] (0.0000,8.0000) circle (1.2pt);
\fill[red!80!black] (12.0000,8.0000) circle (1.2pt);
\fill[red!80!black] (0.0000,8.0000) circle (1.2pt);
\fill[red!80!black] (12.0000,8.0000) circle (1.2pt);
\fill[red!80!black] (0.0000,8.0000) circle (1.2pt);
\fill[red!80!black] (12.0000,8.0000) circle (1.2pt);
\fill[red!80!black] (0.0000,8.0000) circle (1.2pt);
\fill[red!80!black] (12.0000,8.0000) circle (1.2pt);
\fill[red!80!black] (0.0000,8.0000) circle (1.2pt);
\fill[red!80!black] (12.0000,8.0000) circle (1.2pt);
\fill[red!80!black] (0.0000,8.0000) circle (1.2pt);
\fill[red!80!black] (12.0000,8.0000) circle (1.2pt);
\fill[red!80!black] (0.0000,8.0000) circle (1.2pt);
\fill[red!80!black] (12.0000,8.0000) circle (1.2pt);
\fill[red!80!black] (0.0000,8.0000) circle (1.2pt);
\fill[red!80!black] (12.0000,8.0000) circle (1.2pt);
\fill[red!80!black] (0.0000,8.0000) circle (1.2pt);
\fill[red!80!black] (12.0000,8.0000) circle (1.2pt);
\fill[red!80!black] (0.0000,8.0000) circle (1.2pt);
\fill[red!80!black] (12.0000,8.0000) circle (1.2pt);
\fill[red!80!black] (0.0000,8.0000) circle (1.2pt);
\fill[red!80!black] (12.0000,8.0000) circle (1.2pt);
\fill[red!80!black] (0.0000,8.0000) circle (1.2pt);
\fill[red!80!black] (12.0000,8.0000) circle (1.2pt);
\fill[red!80!black] (0.0000,8.0000) circle (1.2pt);
\fill[red!80!black] (12.0000,8.0000) circle (1.2pt);
\fill[red!80!black] (0.0000,8.0000) circle (1.2pt);
\fill[red!80!black] (12.0000,8.0000) circle (1.2pt);
\fill[red!80!black] (0.0000,8.0000) circle (1.2pt);
\fill[red!80!black] (12.0000,8.0000) circle (1.2pt);
\fill[red!80!black] (0.0000,8.0000) circle (1.2pt);
\fill[red!80!black] (12.0000,8.0000) circle (1.2pt);
\fill[red!80!black] (0.0000,8.0000) circle (1.2pt);
\fill[red!80!black] (12.0000,8.0000) circle (1.2pt);
\fill[red!80!black] (0.0000,8.0000) circle (1.2pt);
\fill[red!80!black] (12.0000,8.0000) circle (1.2pt);
\fill[red!80!black] (0.0000,8.0000) circle (1.2pt);
\fill[red!80!black] (12.0000,8.0000) circle (1.2pt);
\fill[red!80!black] (0.0000,8.0000) circle (1.2pt);
\fill[red!80!black] (12.0000,8.0000) circle (1.2pt);
\fill[red!80!black] (0.0000,8.0000) circle (1.2pt);
\fill[red!80!black] (12.0000,8.0000) circle (1.2pt);
\fill[red!80!black] (0.0000,8.0000) circle (1.2pt);
\fill[red!80!black] (12.0000,8.0000) circle (1.2pt);
\fill[red!80!black] (0.0000,8.0000) circle (1.2pt);
\fill[red!80!black] (12.0000,8.0000) circle (1.2pt);
\fill[red!80!black] (0.0000,8.0000) circle (1.2pt);
\fill[red!80!black] (12.0000,8.0000) circle (1.2pt);
\fill[red!80!black] (0.0000,8.0000) circle (1.2pt);
\fill[red!80!black] (12.0000,8.0000) circle (1.2pt);
\fill[red!80!black] (0.0000,8.0000) circle (1.2pt);
\fill[red!80!black] (12.0000,8.0000) circle (1.2pt);
\fill[red!80!black] (0.0000,8.0000) circle (1.2pt);
\fill[red!80!black] (12.0000,8.0000) circle (1.2pt);
\fill[red!80!black] (0.0000,8.0000) circle (1.2pt);
\fill[red!80!black] (12.0000,8.0000) circle (1.2pt);
\fill[red!80!black] (0.0000,8.0000) circle (1.2pt);
\fill[red!80!black] (12.0000,8.0000) circle (1.2pt);
\fill[red!80!black] (0.0000,8.0000) circle (1.2pt);
\fill[red!80!black] (12.0000,8.0000) circle (1.2pt);
\fill[red!80!black] (0.0000,8.0000) circle (1.2pt);
\fill[red!80!black] (12.0000,8.0000) circle (1.2pt);
\fill[red!80!black] (0.0000,8.0000) circle (1.2pt);
\fill[red!80!black] (12.0000,8.0000) circle (1.2pt);
\fill[red!80!black] (0.0000,8.0000) circle (1.2pt);
\fill[red!80!black] (12.0000,8.0000) circle (1.2pt);
\fill[red!80!black] (0.0000,8.0000) circle (1.2pt);
\fill[red!80!black] (12.0000,8.0000) circle (1.2pt);
\fill[red!80!black] (0.0000,8.0000) circle (1.2pt);
\fill[red!80!black] (12.0000,8.0000) circle (1.2pt);
\fill[red!80!black] (0.0000,8.0000) circle (1.2pt);
\fill[red!80!black] (12.0000,8.0000) circle (1.2pt);
\fill[red!80!black] (0.0000,8.0000) circle (1.2pt);
\fill[red!80!black] (12.0000,8.0000) circle (1.2pt);
\fill[red!80!black] (0.0000,8.0000) circle (1.2pt);
\fill[red!80!black] (12.0000,8.0000) circle (1.2pt);
\fill[red!80!black] (0.0000,8.0000) circle (1.2pt);
\fill[red!80!black] (12.0000,8.0000) circle (1.2pt);
\fill[red!80!black] (0.0000,8.0000) circle (1.2pt);
\fill[red!80!black] (12.0000,8.0000) circle (1.2pt);
\fill[red!80!black] (0.0000,8.0000) circle (1.2pt);
\fill[red!80!black] (12.0000,8.0000) circle (1.2pt);
\fill[red!80!black] (0.0000,8.0000) circle (1.2pt);
\fill[red!80!black] (12.0000,8.0000) circle (1.2pt);
\fill[red!80!black] (0.0000,8.0000) circle (1.2pt);
\fill[red!80!black] (12.0000,8.0000) circle (1.2pt);
\fill[red!80!black] (0.0000,8.0000) circle (1.2pt);
\fill[red!80!black] (12.0000,8.0000) circle (1.2pt);
\fill[red!80!black] (0.0000,8.0000) circle (1.2pt);
\fill[red!80!black] (12.0000,8.0000) circle (1.2pt);
\fill[red!80!black] (0.0000,8.0000) circle (1.2pt);
\fill[red!80!black] (12.0000,8.0000) circle (1.2pt);
\fill[red!80!black] (0.0000,8.0000) circle (1.2pt);
\fill[red!80!black] (12.0000,8.0000) circle (1.2pt);
\fill[red!80!black] (0.0000,8.0000) circle (1.2pt);
\fill[red!80!black] (12.0000,8.0000) circle (1.2pt);
\fill[red!80!black] (0.0000,8.0000) circle (1.2pt);
\fill[red!80!black] (12.0000,8.0000) circle (1.2pt);
\fill[red!80!black] (0.0000,8.0000) circle (1.2pt);
\fill[red!80!black] (12.0000,8.0000) circle (1.2pt);
\fill[red!80!black] (0.0000,8.0000) circle (1.2pt);
\fill[red!80!black] (12.0000,8.0000) circle (1.2pt);
\fill[red!80!black] (0.0000,8.0000) circle (1.2pt);
\fill[red!80!black] (12.0000,8.0000) circle (1.2pt);
\fill[red!80!black] (0.0000,8.0000) circle (1.2pt);
\fill[red!80!black] (12.0000,8.0000) circle (1.2pt);
\fill[red!80!black] (0.0000,8.0000) circle (1.2pt);
\fill[red!80!black] (12.0000,8.0000) circle (1.2pt);
\fill[red!80!black] (0.0000,8.0000) circle (1.2pt);
\fill[red!80!black] (12.0000,8.0000) circle (1.2pt);
\fill[red!80!black] (0.0000,8.0000) circle (1.2pt);
\fill[red!80!black] (12.0000,8.0000) circle (1.2pt);
\fill[red!80!black] (0.0000,8.0000) circle (1.2pt);
\fill[red!80!black] (12.0000,8.0000) circle (1.2pt);
\fill[red!80!black] (0.0000,8.0000) circle (1.2pt);
\fill[red!80!black] (12.0000,8.0000) circle (1.2pt);
\fill[red!80!black] (0.0000,8.0000) circle (1.2pt);
\fill[red!80!black] (12.0000,8.0000) circle (1.2pt);
\fill[red!80!black] (0.0000,8.0000) circle (1.2pt);
\fill[red!80!black] (12.0000,8.0000) circle (1.2pt);
\fill[red!80!black] (0.0000,8.0000) circle (1.2pt);
\fill[red!80!black] (12.0000,8.0000) circle (1.2pt);
\fill[red!80!black] (0.0000,8.0000) circle (1.2pt);
\fill[red!80!black] (12.0000,8.0000) circle (1.2pt);
\fill[red!80!black] (0.0000,8.0000) circle (1.2pt);
\fill[red!80!black] (12.0000,8.0000) circle (1.2pt);
\fill[red!80!black] (0.0000,8.0000) circle (1.2pt);
\fill[red!80!black] (12.0000,8.0000) circle (1.2pt);
\fill[red!80!black] (0.0000,8.0000) circle (1.2pt);
\fill[red!80!black] (12.0000,8.0000) circle (1.2pt);
\fill[red!80!black] (0.0000,8.0000) circle (1.2pt);
\fill[red!80!black] (12.0000,8.0000) circle (1.2pt);
\fill[red!80!black] (0.0000,8.0000) circle (1.2pt);
\fill[red!80!black] (12.0000,8.0000) circle (1.2pt);
\fill[red!80!black] (0.0000,8.0000) circle (1.2pt);
\fill[red!80!black] (12.0000,8.0000) circle (1.2pt);
\fill[red!80!black] (0.0000,8.0000) circle (1.2pt);
\fill[red!80!black] (12.0000,8.0000) circle (1.2pt);
\fill[red!80!black] (0.0000,8.0000) circle (1.2pt);
\fill[red!80!black] (12.0000,8.0000) circle (1.2pt);
\fill[red!80!black] (0.0000,8.0000) circle (1.2pt);
\fill[red!80!black] (12.0000,8.0000) circle (1.2pt);
\fill[red!80!black] (0.0000,8.0000) circle (1.2pt);
\fill[red!80!black] (12.0000,8.0000) circle (1.2pt);
\fill[red!80!black] (0.0000,8.0000) circle (1.2pt);
\fill[red!80!black] (12.0000,8.0000) circle (1.2pt);
\fill[red!80!black] (0.0000,8.0000) circle (1.2pt);
\fill[red!80!black] (12.0000,8.0000) circle (1.2pt);
\fill[red!80!black] (0.0000,8.0000) circle (1.2pt);
\fill[red!80!black] (12.0000,8.0000) circle (1.2pt);
\fill[red!80!black] (0.0000,8.0000) circle (1.2pt);
\fill[red!80!black] (12.0000,8.0000) circle (1.2pt);
\fill[red!80!black] (0.0000,8.0000) circle (1.2pt);
\fill[red!80!black] (12.0000,8.0000) circle (1.2pt);
\fill[red!80!black] (0.0000,8.0000) circle (1.2pt);
\fill[red!80!black] (12.0000,8.0000) circle (1.2pt);
\fill[red!80!black] (0.0000,8.0000) circle (1.2pt);
\fill[red!80!black] (12.0000,8.0000) circle (1.2pt);
\fill[red!80!black] (0.0000,8.0000) circle (1.2pt);
\fill[red!80!black] (12.0000,8.0000) circle (1.2pt);
\fill[red!80!black] (0.0000,8.0000) circle (1.2pt);
\fill[red!80!black] (12.0000,8.0000) circle (1.2pt);
\fill[red!80!black] (0.0000,8.0000) circle (1.2pt);
\fill[red!80!black] (12.0000,8.0000) circle (1.2pt);
\fill[red!80!black] (0.0000,8.0000) circle (1.2pt);
\fill[red!80!black] (12.0000,8.0000) circle (1.2pt);
\fill[red!80!black] (0.0000,8.0000) circle (1.2pt);
\fill[red!80!black] (12.0000,8.0000) circle (1.2pt);
\fill[red!80!black] (0.0000,8.0000) circle (1.2pt);
\fill[red!80!black] (12.0000,8.0000) circle (1.2pt);
\fill[red!80!black] (0.0000,8.0000) circle (1.2pt);
\fill[red!80!black] (12.0000,8.0000) circle (1.2pt);
\fill[red!80!black] (0.0000,8.0000) circle (1.2pt);
\fill[red!80!black] (12.0000,8.0000) circle (1.2pt);
\fill[red!80!black] (0.0000,8.0000) circle (1.2pt);
\fill[red!80!black] (12.0000,8.0000) circle (1.2pt);
\fill[red!80!black] (0.0000,8.0000) circle (1.2pt);
\fill[red!80!black] (12.0000,8.0000) circle (1.2pt);
\fill[red!80!black] (0.0000,8.0000) circle (1.2pt);
\fill[red!80!black] (12.0000,8.0000) circle (1.2pt);
\fill[red!80!black] (0.0000,8.0000) circle (1.2pt);
\fill[red!80!black] (12.0000,8.0000) circle (1.2pt);
\fill[red!80!black] (0.0000,8.0000) circle (1.2pt);
\fill[red!80!black] (12.0000,8.0000) circle (1.2pt);
\fill[red!80!black] (0.0000,8.0000) circle (1.2pt);
\fill[red!80!black] (12.0000,8.0000) circle (1.2pt);
\fill[red!80!black] (0.0000,8.0000) circle (1.2pt);
\fill[red!80!black] (12.0000,8.0000) circle (1.2pt);
\fill[red!80!black] (0.0000,8.0000) circle (1.2pt);
\fill[red!80!black] (12.0000,8.0000) circle (1.2pt);
\fill[red!80!black] (0.0000,8.0000) circle (1.2pt);
\fill[red!80!black] (12.0000,8.0000) circle (1.2pt);
\fill[red!80!black] (0.0000,8.0000) circle (1.2pt);
\fill[red!80!black] (12.0000,8.0000) circle (1.2pt);
\fill[red!80!black] (0.0000,8.0000) circle (1.2pt);
\fill[red!80!black] (12.0000,8.0000) circle (1.2pt);
\fill[red!80!black] (0.0000,8.0000) circle (1.2pt);
\fill[red!80!black] (12.0000,8.0000) circle (1.2pt);
\fill[red!80!black] (0.0000,8.0000) circle (1.2pt);
\fill[red!80!black] (12.0000,8.0000) circle (1.2pt);
\fill[red!80!black] (0.0000,8.0000) circle (1.2pt);
\fill[red!80!black] (12.0000,8.0000) circle (1.2pt);
\fill[red!80!black] (0.0000,8.0000) circle (1.2pt);
\fill[red!80!black] (12.0000,8.0000) circle (1.2pt);
\fill[red!80!black] (0.0000,8.0000) circle (1.2pt);
\fill[red!80!black] (12.0000,8.0000) circle (1.2pt);
\fill[red!80!black] (0.0000,8.0000) circle (1.2pt);
\fill[red!80!black] (12.0000,8.0000) circle (1.2pt);
\fill[red!80!black] (0.0000,8.0000) circle (1.2pt);
\fill[red!80!black] (12.0000,8.0000) circle (1.2pt);
\fill[red!80!black] (0.0000,8.0000) circle (1.2pt);
\fill[red!80!black] (12.0000,8.0000) circle (1.2pt);
\fill[red!80!black] (0.0000,8.0000) circle (1.2pt);
\fill[red!80!black] (12.0000,8.0000) circle (1.2pt);
\fill[red!80!black] (0.0000,8.0000) circle (1.2pt);
\fill[red!80!black] (12.0000,8.0000) circle (1.2pt);
\fill[red!80!black] (0.0000,8.0000) circle (1.2pt);
\fill[red!80!black] (12.0000,8.0000) circle (1.2pt);
\fill[red!80!black] (0.0000,8.0000) circle (1.2pt);
\fill[red!80!black] (12.0000,8.0000) circle (1.2pt);
\fill[red!80!black] (0.0000,8.0000) circle (1.2pt);
\fill[red!80!black] (12.0000,8.0000) circle (1.2pt);
\fill[red!80!black] (0.0000,8.0000) circle (1.2pt);
\fill[red!80!black] (12.0000,8.0000) circle (1.2pt);
\fill[red!80!black] (0.0000,8.0000) circle (1.2pt);
\fill[red!80!black] (12.0000,8.0000) circle (1.2pt);
\fill[red!80!black] (0.0000,8.0000) circle (1.2pt);
\fill[red!80!black] (12.0000,8.0000) circle (1.2pt);
\fill[red!80!black] (0.0000,8.0000) circle (1.2pt);
\fill[red!80!black] (12.0000,8.0000) circle (1.2pt);
\fill[red!80!black] (0.0000,8.0000) circle (1.2pt);
\fill[red!80!black] (12.0000,8.0000) circle (1.2pt);
\fill[red!80!black] (0.0000,8.0000) circle (1.2pt);
\fill[red!80!black] (12.0000,8.0000) circle (1.2pt);
\fill[red!80!black] (0.0000,8.0000) circle (1.2pt);
\fill[red!80!black] (12.0000,8.0000) circle (1.2pt);
\fill[red!80!black] (0.0000,8.0000) circle (1.2pt);
\fill[red!80!black] (12.0000,8.0000) circle (1.2pt);
\fill[red!80!black] (0.0000,8.0000) circle (1.2pt);
\fill[red!80!black] (12.0000,8.0000) circle (1.2pt);
\fill[red!80!black] (0.0000,8.0000) circle (1.2pt);
\fill[red!80!black] (12.0000,8.0000) circle (1.2pt);
\fill[red!80!black] (0.0000,8.0000) circle (1.2pt);
\fill[red!80!black] (12.0000,8.0000) circle (1.2pt);
\fill[red!80!black] (0.0000,8.0000) circle (1.2pt);
\fill[red!80!black] (12.0000,8.0000) circle (1.2pt);
\fill[red!80!black] (0.0000,8.0000) circle (1.2pt);
\fill[red!80!black] (12.0000,8.0000) circle (1.2pt);
\fill[red!80!black] (0.0000,8.0000) circle (1.2pt);
\fill[red!80!black] (12.0000,8.0000) circle (1.2pt);
\fill[red!80!black] (0.0000,8.0000) circle (1.2pt);
\fill[red!80!black] (12.0000,8.0000) circle (1.2pt);
\fill[red!80!black] (0.0000,8.0000) circle (1.2pt);
\fill[red!80!black] (12.0000,8.0000) circle (1.2pt);
\fill[red!80!black] (0.0000,8.0000) circle (1.2pt);
\fill[red!80!black] (12.0000,8.0000) circle (1.2pt);
\fill[red!80!black] (0.0000,8.0000) circle (1.2pt);
\fill[red!80!black] (12.0000,8.0000) circle (1.2pt);
\fill[red!80!black] (0.0000,8.0000) circle (1.2pt);
\fill[red!80!black] (12.0000,8.0000) circle (1.2pt);
\fill[red!80!black] (0.0000,8.0000) circle (1.2pt);
\fill[red!80!black] (12.0000,8.0000) circle (1.2pt);
\fill[red!80!black] (0.0000,8.0000) circle (1.2pt);
\fill[red!80!black] (12.0000,8.0000) circle (1.2pt);
\fill[red!80!black] (0.0000,8.0000) circle (1.2pt);
\fill[red!80!black] (12.0000,8.0000) circle (1.2pt);
\fill[red!80!black] (0.0000,8.0000) circle (1.2pt);
\fill[red!80!black] (12.0000,8.0000) circle (1.2pt);
\fill[red!80!black] (0.0000,8.0000) circle (1.2pt);
\fill[red!80!black] (12.0000,8.0000) circle (1.2pt);
\fill[red!80!black] (0.0000,8.0000) circle (1.2pt);
\fill[red!80!black] (12.0000,8.0000) circle (1.2pt);
\fill[red!80!black] (0.0000,8.0000) circle (1.2pt);
\fill[red!80!black] (12.0000,8.0000) circle (1.2pt);
\fill[red!80!black] (0.0000,8.0000) circle (1.2pt);
\fill[red!80!black] (12.0000,8.0000) circle (1.2pt);
\fill[red!80!black] (0.0000,8.0000) circle (1.2pt);
\fill[red!80!black] (12.0000,8.0000) circle (1.2pt);
\fill[red!80!black] (0.0000,8.0000) circle (1.2pt);
\fill[red!80!black] (12.0000,8.0000) circle (1.2pt);
\fill[red!80!black] (0.0000,8.0000) circle (1.2pt);
\fill[red!80!black] (12.0000,8.0000) circle (1.2pt);
\fill[red!80!black] (0.0000,8.0000) circle (1.2pt);
\fill[red!80!black] (12.0000,8.0000) circle (1.2pt);
\fill[red!80!black] (0.0000,8.0000) circle (1.2pt);
\fill[red!80!black] (12.0000,8.0000) circle (1.2pt);
\fill[red!80!black] (0.0000,8.0000) circle (1.2pt);
\fill[red!80!black] (12.0000,8.0000) circle (1.2pt);
\fill[red!80!black] (0.0000,8.0000) circle (1.2pt);
\fill[red!80!black] (12.0000,8.0000) circle (1.2pt);
\fill[red!80!black] (0.0000,8.0000) circle (1.2pt);
\fill[red!80!black] (12.0000,8.0000) circle (1.2pt);
\fill[red!80!black] (0.0000,8.0000) circle (1.2pt);
\fill[red!80!black] (12.0000,8.0000) circle (1.2pt);
\fill[red!80!black] (0.0000,8.0000) circle (1.2pt);
\fill[red!80!black] (12.0000,8.0000) circle (1.2pt);
\fill[red!80!black] (0.0000,8.0000) circle (1.2pt);
\fill[red!80!black] (12.0000,8.0000) circle (1.2pt);
\fill[red!80!black] (0.0000,8.0000) circle (1.2pt);
\fill[red!80!black] (12.0000,8.0000) circle (1.2pt);
\fill[red!80!black] (0.0000,8.0000) circle (1.2pt);
\fill[red!80!black] (12.0000,8.0000) circle (1.2pt);
\fill[red!80!black] (0.0000,8.0000) circle (1.2pt);
\fill[red!80!black] (12.0000,8.0000) circle (1.2pt);
\fill[red!80!black] (0.0000,8.0000) circle (1.2pt);
\fill[red!80!black] (12.0000,8.0000) circle (1.2pt);
\fill[red!80!black] (0.0000,8.0000) circle (1.2pt);
\fill[red!80!black] (12.0000,8.0000) circle (1.2pt);
\fill[red!80!black] (0.0000,8.0000) circle (1.2pt);
\fill[red!80!black] (12.0000,8.0000) circle (1.2pt);
\fill[red!80!black] (0.0000,8.0000) circle (1.2pt);
\fill[red!80!black] (12.0000,8.0000) circle (1.2pt);
\fill[red!80!black] (0.0000,8.0000) circle (1.2pt);
\fill[red!80!black] (12.0000,8.0000) circle (1.2pt);
\fill[red!80!black] (0.0000,8.0000) circle (1.2pt);
\fill[red!80!black] (12.0000,8.0000) circle (1.2pt);
\fill[red!80!black] (0.0000,8.0000) circle (1.2pt);
\fill[red!80!black] (12.0000,8.0000) circle (1.2pt);
\fill[red!80!black] (0.0000,8.0000) circle (1.2pt);
\fill[red!80!black] (12.0000,8.0000) circle (1.2pt);
\fill[red!80!black] (0.0000,8.0000) circle (1.2pt);
\fill[red!80!black] (12.0000,8.0000) circle (1.2pt);
\fill[red!80!black] (0.0000,8.0000) circle (1.2pt);
\fill[red!80!black] (12.0000,8.0000) circle (1.2pt);
\fill[red!80!black] (0.0000,8.0000) circle (1.2pt);
\fill[red!80!black] (12.0000,8.0000) circle (1.2pt);
\fill[red!80!black] (0.0000,8.0000) circle (1.2pt);
\fill[red!80!black] (12.0000,8.0000) circle (1.2pt);
\fill[red!80!black] (0.0000,8.0000) circle (1.2pt);
\fill[red!80!black] (12.0000,8.0000) circle (1.2pt);
\fill[red!80!black] (0.0000,8.0000) circle (1.2pt);
\fill[red!80!black] (12.0000,8.0000) circle (1.2pt);
\fill[red!80!black] (0.0000,8.0000) circle (1.2pt);
\fill[red!80!black] (12.0000,8.0000) circle (1.2pt);
\fill[red!80!black] (0.0000,8.0000) circle (1.2pt);
\fill[red!80!black] (12.0000,8.0000) circle (1.2pt);
\fill[red!80!black] (0.0000,8.0000) circle (1.2pt);
\fill[red!80!black] (12.0000,8.0000) circle (1.2pt);
\fill[red!80!black] (0.0000,8.0000) circle (1.2pt);
\fill[red!80!black] (12.0000,8.0000) circle (1.2pt);
\fill[red!80!black] (0.0000,8.0000) circle (1.2pt);
\fill[red!80!black] (12.0000,8.0000) circle (1.2pt);
\fill[red!80!black] (0.0000,8.0000) circle (1.2pt);
\fill[red!80!black] (12.0000,8.0000) circle (1.2pt);
\fill[red!80!black] (0.0000,8.0000) circle (1.2pt);
\fill[red!80!black] (12.0000,8.0000) circle (1.2pt);
\fill[red!80!black] (0.0000,8.0000) circle (1.2pt);
\fill[red!80!black] (12.0000,8.0000) circle (1.2pt);
\fill[red!80!black] (0.0000,8.0000) circle (1.2pt);
\fill[red!80!black] (12.0000,8.0000) circle (1.2pt);
\fill[red!80!black] (0.0000,8.0000) circle (1.2pt);
\fill[red!80!black] (12.0000,8.0000) circle (1.2pt);
\fill[red!80!black] (0.0000,8.0000) circle (1.2pt);
\fill[red!80!black] (12.0000,8.0000) circle (1.2pt);
\fill[red!80!black] (0.0000,8.0000) circle (1.2pt);
\fill[red!80!black] (12.0000,8.0000) circle (1.2pt);
\fill[red!80!black] (0.0000,8.0000) circle (1.2pt);
\fill[red!80!black] (12.0000,8.0000) circle (1.2pt);
\fill[red!80!black] (0.0000,8.0000) circle (1.2pt);
\fill[red!80!black] (12.0000,8.0000) circle (1.2pt);
\fill[red!80!black] (0.0000,8.0000) circle (1.2pt);
\fill[red!80!black] (12.0000,8.0000) circle (1.2pt);
\fill[red!80!black] (0.0000,8.0000) circle (1.2pt);
\fill[red!80!black] (12.0000,8.0000) circle (1.2pt);
\fill[red!80!black] (0.0000,8.0000) circle (1.2pt);
\fill[red!80!black] (12.0000,8.0000) circle (1.2pt);
\fill[red!80!black] (0.0000,8.0000) circle (1.2pt);
\fill[red!80!black] (12.0000,8.0000) circle (1.2pt);
\fill[red!80!black] (0.0000,8.0000) circle (1.2pt);
\fill[red!80!black] (12.0000,8.0000) circle (1.2pt);
\fill[red!80!black] (0.0000,8.0000) circle (1.2pt);
\fill[red!80!black] (12.0000,8.0000) circle (1.2pt);
\fill[red!80!black] (0.0000,8.0000) circle (1.2pt);
\fill[red!80!black] (12.0000,8.0000) circle (1.2pt);
\fill[red!80!black] (0.0000,8.0000) circle (1.2pt);
\fill[red!80!black] (12.0000,8.0000) circle (1.2pt);
\fill[red!80!black] (0.0000,8.0000) circle (1.2pt);
\fill[red!80!black] (12.0000,8.0000) circle (1.2pt);
\fill[red!80!black] (0.0000,8.0000) circle (1.2pt);
\fill[red!80!black] (12.0000,8.0000) circle (1.2pt);
\fill[red!80!black] (0.0000,8.0000) circle (1.2pt);
\fill[red!80!black] (12.0000,8.0000) circle (1.2pt);
\fill[red!80!black] (0.0000,8.0000) circle (1.2pt);
\fill[red!80!black] (12.0000,8.0000) circle (1.2pt);
\fill[red!80!black] (0.0000,8.0000) circle (1.2pt);
\fill[red!80!black] (12.0000,8.0000) circle (1.2pt);
\fill[red!80!black] (0.0000,8.0000) circle (1.2pt);
\fill[red!80!black] (12.0000,8.0000) circle (1.2pt);
\fill[red!80!black] (0.0000,8.0000) circle (1.2pt);
\fill[red!80!black] (12.0000,8.0000) circle (1.2pt);
\fill[red!80!black] (0.0000,8.0000) circle (1.2pt);
\fill[red!80!black] (12.0000,8.0000) circle (1.2pt);
\fill[red!80!black] (0.0000,8.0000) circle (1.2pt);
\fill[red!80!black] (12.0000,8.0000) circle (1.2pt);
\fill[red!80!black] (0.0000,8.0000) circle (1.2pt);
\fill[red!80!black] (12.0000,8.0000) circle (1.2pt);
\fill[red!80!black] (0.0000,8.0000) circle (1.2pt);
\fill[red!80!black] (12.0000,8.0000) circle (1.2pt);
\fill[red!80!black] (0.0000,8.0000) circle (1.2pt);
\fill[red!80!black] (12.0000,8.0000) circle (1.2pt);
\fill[red!80!black] (0.0000,8.0000) circle (1.2pt);
\fill[red!80!black] (12.0000,8.0000) circle (1.2pt);
\fill[red!80!black] (0.0000,8.0000) circle (1.2pt);
\fill[red!80!black] (12.0000,8.0000) circle (1.2pt);
\fill[red!80!black] (0.0000,8.0000) circle (1.2pt);
\fill[red!80!black] (12.0000,8.0000) circle (1.2pt);
\fill[red!80!black] (0.0000,8.0000) circle (1.2pt);
\fill[red!80!black] (12.0000,8.0000) circle (1.2pt);
\fill[red!80!black] (0.0000,8.0000) circle (1.2pt);
\fill[red!80!black] (12.0000,8.0000) circle (1.2pt);
\fill[red!80!black] (0.0000,8.0000) circle (1.2pt);
\fill[red!80!black] (12.0000,8.0000) circle (1.2pt);
\fill[red!80!black] (0.0000,8.0000) circle (1.2pt);
\fill[red!80!black] (12.0000,8.0000) circle (1.2pt);
\fill[red!80!black] (0.0000,8.0000) circle (1.2pt);
\fill[red!80!black] (12.0000,8.0000) circle (1.2pt);
\fill[red!80!black] (0.0000,8.0000) circle (1.2pt);
\fill[red!80!black] (12.0000,8.0000) circle (1.2pt);
\fill[red!80!black] (0.0000,8.0000) circle (1.2pt);
\fill[red!80!black] (12.0000,8.0000) circle (1.2pt);
\fill[red!80!black] (0.0000,8.0000) circle (1.2pt);
\fill[red!80!black] (12.0000,8.0000) circle (1.2pt);
\fill[red!80!black] (0.0000,8.0000) circle (1.2pt);
\fill[red!80!black] (12.0000,8.0000) circle (1.2pt);
\fill[red!80!black] (0.0000,8.0000) circle (1.2pt);
\fill[red!80!black] (12.0000,8.0000) circle (1.2pt);
\fill[red!80!black] (0.0000,8.0000) circle (1.2pt);
\fill[red!80!black] (12.0000,8.0000) circle (1.2pt);
\fill[red!80!black] (0.0000,8.0000) circle (1.2pt);
\fill[red!80!black] (12.0000,8.0000) circle (1.2pt);
\fill[red!80!black] (0.0000,8.0000) circle (1.2pt);
\fill[red!80!black] (12.0000,8.0000) circle (1.2pt);
\fill[red!80!black] (0.0000,8.0000) circle (1.2pt);
\fill[red!80!black] (12.0000,8.0000) circle (1.2pt);
\fill[red!80!black] (0.0000,8.0000) circle (1.2pt);
\fill[red!80!black] (12.0000,8.0000) circle (1.2pt);
\fill[red!80!black] (0.0000,8.0000) circle (1.2pt);
\fill[red!80!black] (12.0000,8.0000) circle (1.2pt);
\fill[red!80!black] (0.0000,8.0000) circle (1.2pt);
\fill[red!80!black] (12.0000,8.0000) circle (1.2pt);
\fill[red!80!black] (0.0000,8.0000) circle (1.2pt);
\fill[red!80!black] (12.0000,8.0000) circle (1.2pt);
\fill[red!80!black] (0.0000,8.0000) circle (1.2pt);
\fill[red!80!black] (12.0000,8.0000) circle (1.2pt);
\fill[red!80!black] (0.0000,8.0000) circle (1.2pt);
\fill[red!80!black] (12.0000,8.0000) circle (1.2pt);
\fill[red!80!black] (0.0000,8.0000) circle (1.2pt);
\fill[red!80!black] (12.0000,8.0000) circle (1.2pt);
\fill[red!80!black] (0.0000,8.0000) circle (1.2pt);
\fill[red!80!black] (12.0000,8.0000) circle (1.2pt);
\fill[red!80!black] (0.0000,8.0000) circle (1.2pt);
\fill[red!80!black] (12.0000,8.0000) circle (1.2pt);
\fill[red!80!black] (0.0000,8.0000) circle (1.2pt);
\fill[red!80!black] (12.0000,8.0000) circle (1.2pt);
\fill[red!80!black] (0.0000,8.0000) circle (1.2pt);
\fill[red!80!black] (12.0000,8.0000) circle (1.2pt);
\fill[red!80!black] (0.0000,8.0000) circle (1.2pt);
\fill[red!80!black] (12.0000,8.0000) circle (1.2pt);
\fill[red!80!black] (0.0000,8.0000) circle (1.2pt);
\fill[red!80!black] (12.0000,8.0000) circle (1.2pt);
\fill[red!80!black] (0.0000,8.0000) circle (1.2pt);
\fill[red!80!black] (12.0000,8.0000) circle (1.2pt);
\fill[red!80!black] (0.0000,8.0000) circle (1.2pt);
\fill[red!80!black] (12.0000,8.0000) circle (1.2pt);
\fill[red!80!black] (0.0000,8.0000) circle (1.2pt);
\fill[red!80!black] (12.0000,8.0000) circle (1.2pt);
\fill[red!80!black] (0.0000,8.0000) circle (1.2pt);
\fill[red!80!black] (12.0000,8.0000) circle (1.2pt);
\fill[red!80!black] (0.0000,8.0000) circle (1.2pt);
\fill[red!80!black] (12.0000,8.0000) circle (1.2pt);
\fill[red!80!black] (0.0000,8.0000) circle (1.2pt);
\fill[red!80!black] (12.0000,8.0000) circle (1.2pt);
\fill[red!80!black] (0.0000,8.0000) circle (1.2pt);
\fill[red!80!black] (12.0000,8.0000) circle (1.2pt);
\fill[red!80!black] (0.0000,8.0000) circle (1.2pt);
\fill[red!80!black] (12.0000,8.0000) circle (1.2pt);
\fill[red!80!black] (0.0000,8.0000) circle (1.2pt);
\fill[red!80!black] (12.0000,8.0000) circle (1.2pt);
\fill[red!80!black] (0.0000,8.0000) circle (1.2pt);
\fill[red!80!black] (12.0000,8.0000) circle (1.2pt);
\fill[red!80!black] (0.0000,8.0000) circle (1.2pt);
\fill[red!80!black] (12.0000,8.0000) circle (1.2pt);
\fill[red!80!black] (0.0000,8.0000) circle (1.2pt);
\fill[red!80!black] (12.0000,8.0000) circle (1.2pt);
\fill[red!80!black] (0.0000,8.0000) circle (1.2pt);
\fill[red!80!black] (12.0000,8.0000) circle (1.2pt);
\fill[red!80!black] (0.0000,8.0000) circle (1.2pt);
\fill[red!80!black] (12.0000,8.0000) circle (1.2pt);
\fill[red!80!black] (0.0000,8.0000) circle (1.2pt);
\fill[red!80!black] (12.0000,8.0000) circle (1.2pt);
\fill[red!80!black] (0.0000,8.0000) circle (1.2pt);
\fill[red!80!black] (12.0000,8.0000) circle (1.2pt);
\fill[red!80!black] (0.0000,8.0000) circle (1.2pt);
\fill[red!80!black] (12.0000,8.0000) circle (1.2pt);
\fill[red!80!black] (0.0000,8.0000) circle (1.2pt);
\fill[red!80!black] (12.0000,8.0000) circle (1.2pt);
\fill[red!80!black] (0.0000,8.0000) circle (1.2pt);
\fill[red!80!black] (12.0000,8.0000) circle (1.2pt);
\fill[red!80!black] (0.0000,8.0000) circle (1.2pt);
\fill[red!80!black] (12.0000,8.0000) circle (1.2pt);
\fill[red!80!black] (0.0000,8.0000) circle (1.2pt);
\fill[red!80!black] (12.0000,8.0000) circle (1.2pt);
\fill[red!80!black] (0.0000,8.0000) circle (1.2pt);
\fill[red!80!black] (12.0000,8.0000) circle (1.2pt);
\fill[red!80!black] (0.0000,8.0000) circle (1.2pt);
\fill[red!80!black] (12.0000,8.0000) circle (1.2pt);
\fill[red!80!black] (0.0000,8.0000) circle (1.2pt);
\fill[red!80!black] (12.0000,8.0000) circle (1.2pt);
\fill[red!80!black] (0.0000,8.0000) circle (1.2pt);
\fill[red!80!black] (12.0000,8.0000) circle (1.2pt);
\fill[red!80!black] (0.0000,8.0000) circle (1.2pt);
\fill[red!80!black] (12.0000,8.0000) circle (1.2pt);
\fill[red!80!black] (0.0000,8.0000) circle (1.2pt);
\fill[red!80!black] (12.0000,8.0000) circle (1.2pt);
\fill[red!80!black] (0.0000,8.0000) circle (1.2pt);
\fill[red!80!black] (12.0000,8.0000) circle (1.2pt);
\fill[red!80!black] (0.0000,8.0000) circle (1.2pt);
\fill[red!80!black] (12.0000,8.0000) circle (1.2pt);
\fill[red!80!black] (0.0000,8.0000) circle (1.2pt);
\fill[red!80!black] (12.0000,8.0000) circle (1.2pt);
\fill[red!80!black] (0.0000,8.0000) circle (1.2pt);
\fill[red!80!black] (12.0000,8.0000) circle (1.2pt);
\fill[red!80!black] (0.0000,8.0000) circle (1.2pt);
\fill[red!80!black] (12.0000,8.0000) circle (1.2pt);
\fill[red!80!black] (0.0000,8.0000) circle (1.2pt);
\fill[red!80!black] (12.0000,8.0000) circle (1.2pt);
\fill[red!80!black] (0.0000,8.0000) circle (1.2pt);
\fill[red!80!black] (12.0000,8.0000) circle (1.2pt);
\fill[red!80!black] (0.0000,8.0000) circle (1.2pt);
\fill[red!80!black] (12.0000,8.0000) circle (1.2pt);
\fill[red!80!black] (0.0000,8.0000) circle (1.2pt);
\fill[red!80!black] (12.0000,8.0000) circle (1.2pt);
\fill[red!80!black] (0.0000,8.0000) circle (1.2pt);
\fill[red!80!black] (12.0000,8.0000) circle (1.2pt);
\fill[red!80!black] (0.0000,8.0000) circle (1.2pt);
\fill[red!80!black] (12.0000,8.0000) circle (1.2pt);
\fill[red!80!black] (0.0000,8.0000) circle (1.2pt);
\fill[red!80!black] (12.0000,8.0000) circle (1.2pt);
\fill[red!80!black] (0.0000,8.0000) circle (1.2pt);
\fill[red!80!black] (12.0000,8.0000) circle (1.2pt);
\fill[red!80!black] (0.0000,8.0000) circle (1.2pt);
\fill[red!80!black] (12.0000,8.0000) circle (1.2pt);
\fill[red!80!black] (0.0000,8.0000) circle (1.2pt);
\fill[red!80!black] (12.0000,8.0000) circle (1.2pt);
\fill[red!80!black] (0.0000,8.0000) circle (1.2pt);
\fill[red!80!black] (12.0000,8.0000) circle (1.2pt);
\fill[red!80!black] (0.0000,8.0000) circle (1.2pt);
\fill[red!80!black] (12.0000,8.0000) circle (1.2pt);
\fill[red!80!black] (0.0000,8.0000) circle (1.2pt);
\fill[red!80!black] (12.0000,8.0000) circle (1.2pt);
\fill[red!80!black] (0.0000,8.0000) circle (1.2pt);
\fill[red!80!black] (12.0000,8.0000) circle (1.2pt);
\fill[red!80!black] (0.0000,8.0000) circle (1.2pt);
\fill[red!80!black] (12.0000,8.0000) circle (1.2pt);
\fill[red!80!black] (0.0000,8.0000) circle (1.2pt);
\fill[red!80!black] (12.0000,8.0000) circle (1.2pt);
\fill[red!80!black] (0.0000,8.0000) circle (1.2pt);
\fill[red!80!black] (12.0000,8.0000) circle (1.2pt);
\fill[red!80!black] (0.0000,8.0000) circle (1.2pt);
\fill[red!80!black] (12.0000,8.0000) circle (1.2pt);
\fill[red!80!black] (0.0000,8.0000) circle (1.2pt);
\fill[red!80!black] (12.0000,8.0000) circle (1.2pt);
\fill[red!80!black] (0.0000,8.0000) circle (1.2pt);
\fill[red!80!black] (12.0000,8.0000) circle (1.2pt);
\fill[red!80!black] (0.0000,8.0000) circle (1.2pt);
\fill[red!80!black] (12.0000,8.0000) circle (1.2pt);
\fill[red!80!black] (0.0000,8.0000) circle (1.2pt);
\fill[red!80!black] (12.0000,8.0000) circle (1.2pt);
\fill[red!80!black] (0.0000,8.0000) circle (1.2pt);
\fill[red!80!black] (12.0000,8.0000) circle (1.2pt);
\fill[red!80!black] (0.0000,8.0000) circle (1.2pt);
\fill[red!80!black] (12.0000,8.0000) circle (1.2pt);
\fill[red!80!black] (0.0000,8.0000) circle (1.2pt);
\fill[red!80!black] (12.0000,8.0000) circle (1.2pt);
\fill[red!80!black] (0.0000,8.0000) circle (1.2pt);
\fill[red!80!black] (12.0000,8.0000) circle (1.2pt);
\fill[red!80!black] (0.0000,8.0000) circle (1.2pt);
\fill[red!80!black] (12.0000,8.0000) circle (1.2pt);
\fill[red!80!black] (0.0000,8.0000) circle (1.2pt);
\fill[red!80!black] (12.0000,8.0000) circle (1.2pt);
\fill[red!80!black] (0.0000,8.0000) circle (1.2pt);
\fill[red!80!black] (12.0000,8.0000) circle (1.2pt);
\fill[red!80!black] (0.0000,8.0000) circle (1.2pt);
\fill[red!80!black] (12.0000,8.0000) circle (1.2pt);
\fill[red!80!black] (0.0000,8.0000) circle (1.2pt);
\fill[red!80!black] (12.0000,8.0000) circle (1.2pt);
\fill[red!80!black] (0.0000,8.0000) circle (1.2pt);
\fill[red!80!black] (12.0000,8.0000) circle (1.2pt);
\fill[red!80!black] (0.0000,8.0000) circle (1.2pt);
\fill[red!80!black] (12.0000,8.0000) circle (1.2pt);
\fill[red!80!black] (0.0000,8.0000) circle (1.2pt);
\fill[red!80!black] (12.0000,8.0000) circle (1.2pt);
\fill[red!80!black] (0.0000,8.0000) circle (1.2pt);
\fill[red!80!black] (12.0000,8.0000) circle (1.2pt);
\fill[red!80!black] (0.0000,8.0000) circle (1.2pt);
\fill[red!80!black] (12.0000,8.0000) circle (1.2pt);
\fill[red!80!black] (0.0000,8.0000) circle (1.2pt);
\fill[red!80!black] (12.0000,8.0000) circle (1.2pt);
\fill[red!80!black] (0.0000,8.0000) circle (1.2pt);
\fill[red!80!black] (12.0000,8.0000) circle (1.2pt);
\fill[red!80!black] (0.0000,8.0000) circle (1.2pt);
\fill[red!80!black] (12.0000,8.0000) circle (1.2pt);
\fill[red!80!black] (0.0000,8.0000) circle (1.2pt);
\fill[red!80!black] (12.0000,8.0000) circle (1.2pt);
\fill[red!80!black] (0.0000,8.0000) circle (1.2pt);
\fill[red!80!black] (12.0000,8.0000) circle (1.2pt);
\fill[red!80!black] (0.0000,8.0000) circle (1.2pt);
\fill[red!80!black] (12.0000,8.0000) circle (1.2pt);
\fill[red!80!black] (0.0000,8.0000) circle (1.2pt);
\fill[red!80!black] (12.0000,8.0000) circle (1.2pt);
\fill[red!80!black] (0.0000,8.0000) circle (1.2pt);
\fill[red!80!black] (12.0000,8.0000) circle (1.2pt);
\fill[red!80!black] (0.0000,8.0000) circle (1.2pt);
\fill[red!80!black] (12.0000,8.0000) circle (1.2pt);
\fill[red!80!black] (0.0000,8.0000) circle (1.2pt);
\fill[red!80!black] (12.0000,8.0000) circle (1.2pt);
\fill[red!80!black] (0.0000,8.0000) circle (1.2pt);
\fill[red!80!black] (12.0000,8.0000) circle (1.2pt);
\fill[red!80!black] (0.0000,8.0000) circle (1.2pt);
\fill[red!80!black] (12.0000,8.0000) circle (1.2pt);
\fill[red!80!black] (0.0000,8.0000) circle (1.2pt);
\fill[red!80!black] (12.0000,8.0000) circle (1.2pt);
\fill[red!80!black] (0.0000,8.0000) circle (1.2pt);
\fill[red!80!black] (12.0000,8.0000) circle (1.2pt);
\fill[red!80!black] (0.0000,8.0000) circle (1.2pt);
\fill[red!80!black] (12.0000,8.0000) circle (1.2pt);
\fill[red!80!black] (0.0000,8.0000) circle (1.2pt);
\fill[red!80!black] (12.0000,8.0000) circle (1.2pt);
\fill[red!80!black] (0.0000,8.0000) circle (1.2pt);
\fill[red!80!black] (12.0000,8.0000) circle (1.2pt);
\fill[red!80!black] (0.0000,8.0000) circle (1.2pt);
\fill[red!80!black] (12.0000,8.0000) circle (1.2pt);
\fill[red!80!black] (0.0000,8.0000) circle (1.2pt);
\fill[red!80!black] (12.0000,8.0000) circle (1.2pt);
\fill[red!80!black] (0.0000,8.0000) circle (1.2pt);
\fill[red!80!black] (12.0000,8.0000) circle (1.2pt);
\fill[red!80!black] (0.0000,8.0000) circle (1.2pt);
\fill[red!80!black] (12.0000,8.0000) circle (1.2pt);
\fill[red!80!black] (0.0000,8.0000) circle (1.2pt);
\fill[red!80!black] (12.0000,8.0000) circle (1.2pt);
\fill[red!80!black] (0.0000,8.0000) circle (1.2pt);
\fill[red!80!black] (12.0000,8.0000) circle (1.2pt);
\fill[red!80!black] (0.0000,8.0000) circle (1.2pt);
\fill[red!80!black] (12.0000,8.0000) circle (1.2pt);
\fill[red!80!black] (0.0000,8.0000) circle (1.2pt);
\fill[red!80!black] (12.0000,8.0000) circle (1.2pt);
\fill[red!80!black] (0.0000,8.0000) circle (1.2pt);
\fill[red!80!black] (12.0000,8.0000) circle (1.2pt);
\fill[red!80!black] (0.0000,8.0000) circle (1.2pt);
\fill[red!80!black] (12.0000,8.0000) circle (1.2pt);
\fill[red!80!black] (0.0000,8.0000) circle (1.2pt);
\fill[red!80!black] (12.0000,8.0000) circle (1.2pt);
\fill[red!80!black] (0.0000,8.0000) circle (1.2pt);
\fill[red!80!black] (12.0000,8.0000) circle (1.2pt);
\fill[red!80!black] (0.0000,8.0000) circle (1.2pt);
\fill[red!80!black] (12.0000,8.0000) circle (1.2pt);
\fill[red!80!black] (0.0000,8.0000) circle (1.2pt);
\fill[red!80!black] (12.0000,8.0000) circle (1.2pt);
\fill[red!80!black] (0.0000,8.0000) circle (1.2pt);
\fill[red!80!black] (12.0000,8.0000) circle (1.2pt);
\fill[red!80!black] (0.0000,8.0000) circle (1.2pt);
\fill[red!80!black] (12.0000,8.0000) circle (1.2pt);
\fill[red!80!black] (0.0000,8.0000) circle (1.2pt);
\fill[red!80!black] (12.0000,8.0000) circle (1.2pt);
\fill[red!80!black] (0.0000,8.0000) circle (1.2pt);
\fill[red!80!black] (12.0000,8.0000) circle (1.2pt);
\fill[red!80!black] (0.0000,8.0000) circle (1.2pt);
\fill[red!80!black] (12.0000,8.0000) circle (1.2pt);
\fill[red!80!black] (0.0000,8.0000) circle (1.2pt);
\fill[red!80!black] (12.0000,8.0000) circle (1.2pt);
\fill[red!80!black] (0.0000,8.0000) circle (1.2pt);
\fill[red!80!black] (12.0000,8.0000) circle (1.2pt);
\fill[red!80!black] (0.0000,8.0000) circle (1.2pt);
\fill[red!80!black] (12.0000,8.0000) circle (1.2pt);
\fill[red!80!black] (0.0000,8.0000) circle (1.2pt);
\fill[red!80!black] (12.0000,8.0000) circle (1.2pt);
\fill[red!80!black] (0.0000,8.0000) circle (1.2pt);
\fill[red!80!black] (12.0000,8.0000) circle (1.2pt);
\fill[red!80!black] (0.0000,8.0000) circle (1.2pt);
\fill[red!80!black] (12.0000,8.0000) circle (1.2pt);
\fill[red!80!black] (0.0000,8.0000) circle (1.2pt);
\fill[red!80!black] (12.0000,8.0000) circle (1.2pt);
\fill[red!80!black] (0.0000,8.0000) circle (1.2pt);
\fill[red!80!black] (12.0000,8.0000) circle (1.2pt);
\fill[red!80!black] (0.0000,8.0000) circle (1.2pt);
\fill[red!80!black] (12.0000,8.0000) circle (1.2pt);
\fill[red!80!black] (0.0000,8.0000) circle (1.2pt);
\fill[red!80!black] (12.0000,8.0000) circle (1.2pt);
\fill[red!80!black] (0.0000,8.0000) circle (1.2pt);
\fill[red!80!black] (12.0000,8.0000) circle (1.2pt);
\fill[red!80!black] (0.0000,8.0000) circle (1.2pt);
\fill[red!80!black] (12.0000,8.0000) circle (1.2pt);
\fill[red!80!black] (0.0000,8.0000) circle (1.2pt);
\fill[red!80!black] (12.0000,8.0000) circle (1.2pt);
\fill[red!80!black] (0.0000,8.0000) circle (1.2pt);
\fill[red!80!black] (12.0000,8.0000) circle (1.2pt);
\fill[red!80!black] (0.0000,8.0000) circle (1.2pt);
\fill[red!80!black] (12.0000,8.0000) circle (1.2pt);
\fill[red!80!black] (0.0000,8.0000) circle (1.2pt);
\fill[red!80!black] (12.0000,8.0000) circle (1.2pt);
\fill[red!80!black] (0.0000,8.0000) circle (1.2pt);
\fill[red!80!black] (12.0000,8.0000) circle (1.2pt);
\fill[red!80!black] (0.0000,8.0000) circle (1.2pt);
\fill[red!80!black] (12.0000,8.0000) circle (1.2pt);
\fill[red!80!black] (0.0000,8.0000) circle (1.2pt);
\fill[red!80!black] (12.0000,8.0000) circle (1.2pt);
\fill[red!80!black] (0.0000,8.0000) circle (1.2pt);
\fill[red!80!black] (12.0000,8.0000) circle (1.2pt);
\fill[red!80!black] (0.0000,8.0000) circle (1.2pt);
\fill[red!80!black] (12.0000,8.0000) circle (1.2pt);
\fill[red!80!black] (0.0000,8.0000) circle (1.2pt);
\fill[red!80!black] (12.0000,8.0000) circle (1.2pt);
\fill[red!80!black] (0.0000,8.0000) circle (1.2pt);
\fill[red!80!black] (12.0000,8.0000) circle (1.2pt);
\fill[red!80!black] (0.0000,8.0000) circle (1.2pt);
\fill[red!80!black] (12.0000,8.0000) circle (1.2pt);
\fill[red!80!black] (0.0000,8.0000) circle (1.2pt);
\fill[red!80!black] (12.0000,8.0000) circle (1.2pt);
\fill[red!80!black] (0.0000,8.0000) circle (1.2pt);
\fill[red!80!black] (12.0000,8.0000) circle (1.2pt);
\fill[red!80!black] (0.0000,8.0000) circle (1.2pt);
\fill[red!80!black] (12.0000,8.0000) circle (1.2pt);
\fill[red!80!black] (0.0000,8.0000) circle (1.2pt);
\fill[red!80!black] (12.0000,8.0000) circle (1.2pt);
\fill[red!80!black] (0.0000,8.0000) circle (1.2pt);
\fill[red!80!black] (12.0000,8.0000) circle (1.2pt);
\fill[red!80!black] (0.0000,8.0000) circle (1.2pt);
\fill[red!80!black] (12.0000,8.0000) circle (1.2pt);
\fill[red!80!black] (0.0000,8.0000) circle (1.2pt);
\fill[red!80!black] (12.0000,8.0000) circle (1.2pt);
\fill[red!80!black] (0.0000,8.0000) circle (1.2pt);
\fill[red!80!black] (12.0000,8.0000) circle (1.2pt);
\fill[red!80!black] (0.0000,8.0000) circle (1.2pt);
\fill[red!80!black] (12.0000,8.0000) circle (1.2pt);
\fill[red!80!black] (0.0000,8.0000) circle (1.2pt);
\fill[red!80!black] (12.0000,8.0000) circle (1.2pt);
\fill[red!80!black] (0.0000,8.0000) circle (1.2pt);
\fill[red!80!black] (12.0000,8.0000) circle (1.2pt);
\fill[red!80!black] (0.0000,8.0000) circle (1.2pt);
\fill[red!80!black] (12.0000,8.0000) circle (1.2pt);
\fill[red!80!black] (0.0000,8.0000) circle (1.2pt);
\fill[red!80!black] (12.0000,8.0000) circle (1.2pt);
\fill[red!80!black] (0.0000,8.0000) circle (1.2pt);
\fill[red!80!black] (12.0000,8.0000) circle (1.2pt);
\fill[red!80!black] (0.0000,8.0000) circle (1.2pt);
\fill[red!80!black] (12.0000,8.0000) circle (1.2pt);
\fill[red!80!black] (0.0000,8.0000) circle (1.2pt);
\fill[red!80!black] (12.0000,8.0000) circle (1.2pt);
\fill[red!80!black] (0.0000,8.0000) circle (1.2pt);
\fill[red!80!black] (12.0000,8.0000) circle (1.2pt);
\fill[red!80!black] (0.0000,8.0000) circle (1.2pt);
\fill[red!80!black] (12.0000,8.0000) circle (1.2pt);
\fill[red!80!black] (0.0000,8.0000) circle (1.2pt);
\fill[red!80!black] (12.0000,8.0000) circle (1.2pt);
\fill[red!80!black] (0.0000,8.0000) circle (1.2pt);
\fill[red!80!black] (12.0000,8.0000) circle (1.2pt);
\fill[red!80!black] (0.0000,8.0000) circle (1.2pt);
\fill[red!80!black] (12.0000,8.0000) circle (1.2pt);
\fill[red!80!black] (0.0000,8.0000) circle (1.2pt);
\fill[red!80!black] (12.0000,8.0000) circle (1.2pt);
\fill[red!80!black] (0.0000,8.0000) circle (1.2pt);
\fill[red!80!black] (12.0000,8.0000) circle (1.2pt);
\fill[red!80!black] (0.0000,8.0000) circle (1.2pt);
\fill[red!80!black] (12.0000,8.0000) circle (1.2pt);
\fill[red!80!black] (0.0000,8.0000) circle (1.2pt);
\fill[red!80!black] (12.0000,8.0000) circle (1.2pt);
\fill[red!80!black] (0.0000,8.0000) circle (1.2pt);
\fill[red!80!black] (12.0000,8.0000) circle (1.2pt);
\fill[red!80!black] (0.0000,8.0000) circle (1.2pt);
\fill[red!80!black] (12.0000,8.0000) circle (1.2pt);
\fill[red!80!black] (0.0000,8.0000) circle (1.2pt);
\fill[red!80!black] (12.0000,8.0000) circle (1.2pt);
\fill[red!80!black] (0.0000,8.0000) circle (1.2pt);
\fill[red!80!black] (12.0000,8.0000) circle (1.2pt);
\fill[red!80!black] (0.0000,8.0000) circle (1.2pt);
\fill[red!80!black] (12.0000,8.0000) circle (1.2pt);
\fill[red!80!black] (0.0000,8.0000) circle (1.2pt);
\fill[red!80!black] (12.0000,8.0000) circle (1.2pt);
\fill[red!80!black] (0.0000,8.0000) circle (1.2pt);
\fill[red!80!black] (12.0000,8.0000) circle (1.2pt);
\fill[red!80!black] (0.0000,8.0000) circle (1.2pt);
\fill[red!80!black] (12.0000,8.0000) circle (1.2pt);
\fill[red!80!black] (0.0000,8.0000) circle (1.2pt);
\fill[red!80!black] (12.0000,8.0000) circle (1.2pt);
\fill[red!80!black] (0.0000,8.0000) circle (1.2pt);
\fill[red!80!black] (12.0000,8.0000) circle (1.2pt);
\fill[red!80!black] (0.0000,8.0000) circle (1.2pt);
\fill[red!80!black] (12.0000,8.0000) circle (1.2pt);
\fill[red!80!black] (0.0000,8.0000) circle (1.2pt);
\fill[red!80!black] (12.0000,8.0000) circle (1.2pt);
\fill[red!80!black] (0.0000,8.0000) circle (1.2pt);
\fill[red!80!black] (12.0000,8.0000) circle (1.2pt);
\fill[red!80!black] (0.0000,8.0000) circle (1.2pt);
\fill[red!80!black] (12.0000,8.0000) circle (1.2pt);
\fill[red!80!black] (0.0000,8.0000) circle (1.2pt);
\fill[red!80!black] (12.0000,8.0000) circle (1.2pt);
\fill[red!80!black] (0.0000,8.0000) circle (1.2pt);
\fill[red!80!black] (12.0000,8.0000) circle (1.2pt);
\fill[red!80!black] (0.0000,8.0000) circle (1.2pt);
\fill[red!80!black] (12.0000,8.0000) circle (1.2pt);
\fill[red!80!black] (0.0000,8.0000) circle (1.2pt);
\fill[red!80!black] (12.0000,8.0000) circle (1.2pt);
\fill[red!80!black] (0.0000,8.0000) circle (1.2pt);
\fill[red!80!black] (12.0000,8.0000) circle (1.2pt);
\fill[red!80!black] (0.0000,8.0000) circle (1.2pt);
\fill[red!80!black] (12.0000,8.0000) circle (1.2pt);
\fill[red!80!black] (0.0000,8.0000) circle (1.2pt);
\fill[red!80!black] (12.0000,8.0000) circle (1.2pt);
\fill[red!80!black] (0.0000,8.0000) circle (1.2pt);
\fill[red!80!black] (12.0000,8.0000) circle (1.2pt);
\fill[red!80!black] (0.0000,8.0000) circle (1.2pt);
\fill[red!80!black] (12.0000,8.0000) circle (1.2pt);
\fill[red!80!black] (0.0000,8.0000) circle (1.2pt);
\fill[red!80!black] (12.0000,8.0000) circle (1.2pt);
\fill[red!80!black] (0.0000,8.0000) circle (1.2pt);
\fill[red!80!black] (12.0000,8.0000) circle (1.2pt);
\fill[red!80!black] (0.0000,8.0000) circle (1.2pt);
\fill[red!80!black] (12.0000,8.0000) circle (1.2pt);
\fill[red!80!black] (0.0000,8.0000) circle (1.2pt);
\fill[red!80!black] (12.0000,8.0000) circle (1.2pt);
\fill[red!80!black] (0.0000,8.0000) circle (1.2pt);
\fill[red!80!black] (12.0000,8.0000) circle (1.2pt);
\fill[red!80!black] (0.0000,8.0000) circle (1.2pt);
\fill[red!80!black] (12.0000,8.0000) circle (1.2pt);
\fill[red!80!black] (0.0000,8.0000) circle (1.2pt);
\fill[red!80!black] (12.0000,8.0000) circle (1.2pt);
\fill[red!80!black] (0.0000,8.0000) circle (1.2pt);
\fill[red!80!black] (12.0000,8.0000) circle (1.2pt);
\fill[red!80!black] (0.0000,8.0000) circle (1.2pt);
\fill[red!80!black] (12.0000,8.0000) circle (1.2pt);
\fill[red!80!black] (0.0000,8.0000) circle (1.2pt);
\fill[red!80!black] (12.0000,8.0000) circle (1.2pt);
\fill[red!80!black] (0.0000,8.0000) circle (1.2pt);
\fill[red!80!black] (12.0000,8.0000) circle (1.2pt);
\fill[red!80!black] (0.0000,8.0000) circle (1.2pt);
\fill[red!80!black] (12.0000,8.0000) circle (1.2pt);
\fill[red!80!black] (0.0000,8.0000) circle (1.2pt);
\fill[red!80!black] (12.0000,8.0000) circle (1.2pt);
\fill[red!80!black] (0.0000,8.0000) circle (1.2pt);
\fill[red!80!black] (12.0000,8.0000) circle (1.2pt);
\fill[red!80!black] (0.0000,8.0000) circle (1.2pt);
\fill[red!80!black] (12.0000,8.0000) circle (1.2pt);
\fill[red!80!black] (0.0000,8.0000) circle (1.2pt);
\fill[red!80!black] (12.0000,8.0000) circle (1.2pt);
\fill[red!80!black] (0.0000,8.0000) circle (1.2pt);
\fill[red!80!black] (12.0000,8.0000) circle (1.2pt);
\fill[red!80!black] (0.0000,8.0000) circle (1.2pt);
\fill[red!80!black] (12.0000,8.0000) circle (1.2pt);
\fill[red!80!black] (0.0000,8.0000) circle (1.2pt);
\fill[red!80!black] (12.0000,8.0000) circle (1.2pt);
\fill[red!80!black] (0.0000,8.0000) circle (1.2pt);
\fill[red!80!black] (12.0000,8.0000) circle (1.2pt);
\fill[red!80!black] (0.0000,8.0000) circle (1.2pt);
\fill[red!80!black] (12.0000,8.0000) circle (1.2pt);
\fill[red!80!black] (0.0000,8.0000) circle (1.2pt);
\fill[red!80!black] (12.0000,8.0000) circle (1.2pt);
\fill[red!80!black] (0.0000,8.0000) circle (1.2pt);
\fill[red!80!black] (12.0000,8.0000) circle (1.2pt);
\fill[red!80!black] (0.0000,8.0000) circle (1.2pt);
\fill[red!80!black] (12.0000,8.0000) circle (1.2pt);
\fill[red!80!black] (0.0000,8.0000) circle (1.2pt);
\fill[red!80!black] (12.0000,8.0000) circle (1.2pt);
\fill[red!80!black] (0.0000,8.0000) circle (1.2pt);
\fill[red!80!black] (12.0000,8.0000) circle (1.2pt);
\fill[red!80!black] (0.0000,8.0000) circle (1.2pt);
\fill[red!80!black] (12.0000,8.0000) circle (1.2pt);
\fill[red!80!black] (0.0000,8.0000) circle (1.2pt);
\fill[red!80!black] (12.0000,8.0000) circle (1.2pt);
\fill[red!80!black] (0.0000,8.0000) circle (1.2pt);
\fill[red!80!black] (12.0000,8.0000) circle (1.2pt);
\fill[red!80!black] (0.0000,8.0000) circle (1.2pt);
\fill[red!80!black] (12.0000,8.0000) circle (1.2pt);
\fill[red!80!black] (0.0000,8.0000) circle (1.2pt);
\fill[red!80!black] (12.0000,8.0000) circle (1.2pt);
\fill[red!80!black] (0.0000,8.0000) circle (1.2pt);
\fill[red!80!black] (12.0000,8.0000) circle (1.2pt);
\fill[red!80!black] (0.0000,8.0000) circle (1.2pt);
\fill[red!80!black] (12.0000,8.0000) circle (1.2pt);
\fill[red!80!black] (0.0000,8.0000) circle (1.2pt);
\fill[red!80!black] (12.0000,8.0000) circle (1.2pt);
\fill[red!80!black] (0.0000,8.0000) circle (1.2pt);
\fill[red!80!black] (12.0000,8.0000) circle (1.2pt);
\fill[red!80!black] (0.0000,8.0000) circle (1.2pt);
\fill[red!80!black] (12.0000,8.0000) circle (1.2pt);
\fill[red!80!black] (0.0000,8.0000) circle (1.2pt);
\fill[red!80!black] (12.0000,8.0000) circle (1.2pt);
\fill[red!80!black] (0.0000,8.0000) circle (1.2pt);
\fill[red!80!black] (12.0000,8.0000) circle (1.2pt);
\fill[red!80!black] (0.0000,8.0000) circle (1.2pt);
\fill[red!80!black] (12.0000,8.0000) circle (1.2pt);
\fill[red!80!black] (0.0000,8.0000) circle (1.2pt);
\fill[red!80!black] (12.0000,8.0000) circle (1.2pt);
\fill[red!80!black] (0.0000,8.0000) circle (1.2pt);
\fill[red!80!black] (12.0000,8.0000) circle (1.2pt);
\fill[red!80!black] (0.0000,8.0000) circle (1.2pt);
\fill[red!80!black] (12.0000,8.0000) circle (1.2pt);
\fill[red!80!black] (0.0000,8.0000) circle (1.2pt);
\fill[red!80!black] (12.0000,8.0000) circle (1.2pt);
\fill[red!80!black] (0.0000,8.0000) circle (1.2pt);
\fill[red!80!black] (12.0000,8.0000) circle (1.2pt);
\fill[red!80!black] (0.0000,8.0000) circle (1.2pt);
\fill[red!80!black] (12.0000,8.0000) circle (1.2pt);
\fill[red!80!black] (0.0000,8.0000) circle (1.2pt);
\fill[red!80!black] (12.0000,8.0000) circle (1.2pt);
\fill[red!80!black] (0.0000,8.0000) circle (1.2pt);
\fill[red!80!black] (12.0000,8.0000) circle (1.2pt);
\fill[red!80!black] (0.0000,8.0000) circle (1.2pt);
\fill[red!80!black] (12.0000,8.0000) circle (1.2pt);
\fill[red!80!black] (0.0000,8.0000) circle (1.2pt);
\fill[red!80!black] (12.0000,8.0000) circle (1.2pt);
\fill[red!80!black] (0.0000,8.0000) circle (1.2pt);
\fill[red!80!black] (12.0000,8.0000) circle (1.2pt);
\fill[red!80!black] (0.0000,8.0000) circle (1.2pt);
\fill[red!80!black] (12.0000,8.0000) circle (1.2pt);
\fill[red!80!black] (0.0000,8.0000) circle (1.2pt);
\fill[red!80!black] (12.0000,8.0000) circle (1.2pt);
\fill[red!80!black] (0.0000,8.0000) circle (1.2pt);
\fill[red!80!black] (12.0000,8.0000) circle (1.2pt);
\fill[red!80!black] (0.0000,8.0000) circle (1.2pt);
\fill[red!80!black] (12.0000,8.0000) circle (1.2pt);
\fill[red!80!black] (0.0000,8.0000) circle (1.2pt);
\fill[red!80!black] (12.0000,8.0000) circle (1.2pt);
\fill[red!80!black] (0.0000,8.0000) circle (1.2pt);
\fill[red!80!black] (12.0000,8.0000) circle (1.2pt);
\fill[red!80!black] (0.0000,8.0000) circle (1.2pt);
\fill[red!80!black] (12.0000,8.0000) circle (1.2pt);
\fill[red!80!black] (0.0000,8.0000) circle (1.2pt);
\fill[red!80!black] (12.0000,8.0000) circle (1.2pt);
\fill[red!80!black] (0.0000,8.0000) circle (1.2pt);
\fill[red!80!black] (12.0000,8.0000) circle (1.2pt);
\fill[red!80!black] (0.0000,8.0000) circle (1.2pt);
\fill[red!80!black] (12.0000,8.0000) circle (1.2pt);
\fill[red!80!black] (0.0000,8.0000) circle (1.2pt);
\fill[red!80!black] (12.0000,8.0000) circle (1.2pt);
\fill[red!80!black] (0.0000,8.0000) circle (1.2pt);
\fill[red!80!black] (12.0000,8.0000) circle (1.2pt);
\fill[red!80!black] (0.0000,8.0000) circle (1.2pt);
\fill[red!80!black] (12.0000,8.0000) circle (1.2pt);
\fill[red!80!black] (0.0000,8.0000) circle (1.2pt);
\fill[red!80!black] (12.0000,8.0000) circle (1.2pt);
\fill[red!80!black] (0.0000,8.0000) circle (1.2pt);
\fill[red!80!black] (12.0000,8.0000) circle (1.2pt);
\fill[red!80!black] (0.0000,8.0000) circle (1.2pt);
\fill[red!80!black] (12.0000,8.0000) circle (1.2pt);
\fill[red!80!black] (0.0000,8.0000) circle (1.2pt);
\fill[red!80!black] (12.0000,8.0000) circle (1.2pt);
\fill[red!80!black] (0.0000,8.0000) circle (1.2pt);
\fill[red!80!black] (12.0000,8.0000) circle (1.2pt);
\fill[red!80!black] (0.0000,8.0000) circle (1.2pt);
\fill[red!80!black] (12.0000,8.0000) circle (1.2pt);
\fill[red!80!black] (0.0000,8.0000) circle (1.2pt);
\fill[red!80!black] (12.0000,8.0000) circle (1.2pt);
\fill[red!80!black] (0.0000,8.0000) circle (1.2pt);
\fill[red!80!black] (12.0000,8.0000) circle (1.2pt);
\fill[red!80!black] (0.0000,8.0000) circle (1.2pt);
\fill[red!80!black] (12.0000,8.0000) circle (1.2pt);
\fill[red!80!black] (0.0000,8.0000) circle (1.2pt);
\fill[red!80!black] (12.0000,8.0000) circle (1.2pt);
\fill[red!80!black] (0.0000,8.0000) circle (1.2pt);
\fill[red!80!black] (12.0000,8.0000) circle (1.2pt);
\fill[red!80!black] (0.0000,8.0000) circle (1.2pt);
\fill[red!80!black] (12.0000,8.0000) circle (1.2pt);
\fill[red!80!black] (0.0000,8.0000) circle (1.2pt);
\fill[red!80!black] (12.0000,8.0000) circle (1.2pt);
\fill[red!80!black] (0.0000,8.0000) circle (1.2pt);
\fill[red!80!black] (12.0000,8.0000) circle (1.2pt);
\fill[red!80!black] (0.0000,8.0000) circle (1.2pt);
\fill[red!80!black] (12.0000,8.0000) circle (1.2pt);
\fill[red!80!black] (0.0000,8.0000) circle (1.2pt);
\fill[red!80!black] (12.0000,8.0000) circle (1.2pt);
\fill[red!80!black] (0.0000,8.0000) circle (1.2pt);
\fill[red!80!black] (12.0000,8.0000) circle (1.2pt);
\fill[red!80!black] (0.0000,8.0000) circle (1.2pt);
\fill[red!80!black] (12.0000,8.0000) circle (1.2pt);
\fill[red!80!black] (0.0000,8.0000) circle (1.2pt);
\fill[red!80!black] (12.0000,8.0000) circle (1.2pt);
\fill[red!80!black] (0.0000,8.0000) circle (1.2pt);
\fill[red!80!black] (12.0000,8.0000) circle (1.2pt);
\fill[red!80!black] (0.0000,8.0000) circle (1.2pt);
\fill[red!80!black] (12.0000,8.0000) circle (1.2pt);
\fill[red!80!black] (0.0000,8.0000) circle (1.2pt);
\fill[red!80!black] (12.0000,8.0000) circle (1.2pt);
\fill[red!80!black] (0.0000,8.0000) circle (1.2pt);
\fill[red!80!black] (12.0000,8.0000) circle (1.2pt);
\fill[red!80!black] (0.0000,8.0000) circle (1.2pt);
\fill[red!80!black] (12.0000,8.0000) circle (1.2pt);
\fill[red!80!black] (0.0000,8.0000) circle (1.2pt);
\fill[red!80!black] (12.0000,8.0000) circle (1.2pt);
\fill[red!80!black] (0.0000,8.0000) circle (1.2pt);
\fill[red!80!black] (12.0000,8.0000) circle (1.2pt);
\fill[red!80!black] (0.0000,8.0000) circle (1.2pt);
\fill[red!80!black] (12.0000,8.0000) circle (1.2pt);
\fill[red!80!black] (0.0000,8.0000) circle (1.2pt);
\fill[red!80!black] (12.0000,8.0000) circle (1.2pt);
\fill[red!80!black] (0.0000,8.0000) circle (1.2pt);
\fill[red!80!black] (12.0000,8.0000) circle (1.2pt);
\fill[red!80!black] (0.0000,8.0000) circle (1.2pt);
\fill[red!80!black] (12.0000,8.0000) circle (1.2pt);
\fill[red!80!black] (0.0000,8.0000) circle (1.2pt);
\fill[red!80!black] (12.0000,8.0000) circle (1.2pt);
\fill[red!80!black] (0.0000,8.0000) circle (1.2pt);
\fill[red!80!black] (12.0000,8.0000) circle (1.2pt);
\fill[red!80!black] (0.0000,8.0000) circle (1.2pt);
\fill[red!80!black] (12.0000,8.0000) circle (1.2pt);
\fill[red!80!black] (0.0000,8.0000) circle (1.2pt);
\fill[red!80!black] (12.0000,8.0000) circle (1.2pt);
\fill[red!80!black] (0.0000,8.0000) circle (1.2pt);
\fill[red!80!black] (12.0000,8.0000) circle (1.2pt);
\fill[red!80!black] (0.0000,8.0000) circle (1.2pt);
\fill[red!80!black] (12.0000,8.0000) circle (1.2pt);
\fill[red!80!black] (0.0000,8.0000) circle (1.2pt);
\fill[red!80!black] (12.0000,8.0000) circle (1.2pt);
\fill[red!80!black] (0.0000,8.0000) circle (1.2pt);
\fill[red!80!black] (12.0000,8.0000) circle (1.2pt);
\fill[red!80!black] (0.0000,8.0000) circle (1.2pt);
\fill[red!80!black] (12.0000,8.0000) circle (1.2pt);
\fill[red!80!black] (0.0000,8.0000) circle (1.2pt);
\fill[red!80!black] (12.0000,8.0000) circle (1.2pt);
\fill[red!80!black] (0.0000,8.0000) circle (1.2pt);
\fill[red!80!black] (12.0000,8.0000) circle (1.2pt);
\fill[red!80!black] (0.0000,8.0000) circle (1.2pt);
\fill[red!80!black] (12.0000,8.0000) circle (1.2pt);
\fill[red!80!black] (0.0000,8.0000) circle (1.2pt);
\fill[red!80!black] (12.0000,8.0000) circle (1.2pt);
\fill[red!80!black] (0.0000,8.0000) circle (1.2pt);
\fill[red!80!black] (12.0000,8.0000) circle (1.2pt);
\fill[red!80!black] (0.0000,8.0000) circle (1.2pt);
\fill[red!80!black] (12.0000,8.0000) circle (1.2pt);
\fill[red!80!black] (0.0000,8.0000) circle (1.2pt);
\fill[red!80!black] (12.0000,8.0000) circle (1.2pt);
\fill[red!80!black] (0.0000,8.0000) circle (1.2pt);
\fill[red!80!black] (12.0000,8.0000) circle (1.2pt);
\fill[red!80!black] (0.0000,8.0000) circle (1.2pt);
\fill[red!80!black] (12.0000,8.0000) circle (1.2pt);
\fill[red!80!black] (0.0000,8.0000) circle (1.2pt);
\fill[red!80!black] (12.0000,8.0000) circle (1.2pt);
\fill[red!80!black] (0.0000,8.0000) circle (1.2pt);
\fill[red!80!black] (12.0000,8.0000) circle (1.2pt);
\fill[red!80!black] (0.0000,8.0000) circle (1.2pt);
\fill[red!80!black] (12.0000,8.0000) circle (1.2pt);
\fill[red!80!black] (0.0000,8.0000) circle (1.2pt);
\fill[red!80!black] (12.0000,8.0000) circle (1.2pt);
\fill[red!80!black] (0.0000,8.0000) circle (1.2pt);
\fill[red!80!black] (12.0000,8.0000) circle (1.2pt);
\fill[red!80!black] (0.0000,8.0000) circle (1.2pt);
\fill[red!80!black] (12.0000,8.0000) circle (1.2pt);
\fill[red!80!black] (0.0000,8.0000) circle (1.2pt);
\fill[red!80!black] (12.0000,8.0000) circle (1.2pt);
\fill[red!80!black] (0.0000,8.0000) circle (1.2pt);
\fill[red!80!black] (12.0000,8.0000) circle (1.2pt);
\fill[red!80!black] (0.0000,8.0000) circle (1.2pt);
\fill[red!80!black] (12.0000,8.0000) circle (1.2pt);
\fill[red!80!black] (0.0000,8.0000) circle (1.2pt);
\fill[red!80!black] (12.0000,8.0000) circle (1.2pt);
\fill[red!80!black] (0.0000,8.0000) circle (1.2pt);
\fill[red!80!black] (12.0000,8.0000) circle (1.2pt);
\fill[red!80!black] (0.0000,8.0000) circle (1.2pt);
\fill[red!80!black] (12.0000,8.0000) circle (1.2pt);
\fill[red!80!black] (0.0000,8.0000) circle (1.2pt);
\fill[red!80!black] (12.0000,8.0000) circle (1.2pt);
\fill[red!80!black] (0.0000,8.0000) circle (1.2pt);
\fill[red!80!black] (12.0000,8.0000) circle (1.2pt);
\fill[red!80!black] (0.0000,8.0000) circle (1.2pt);
\fill[red!80!black] (12.0000,8.0000) circle (1.2pt);
\fill[red!80!black] (0.0000,8.0000) circle (1.2pt);
\fill[red!80!black] (12.0000,8.0000) circle (1.2pt);
\fill[red!80!black] (0.0000,8.0000) circle (1.2pt);
\fill[red!80!black] (12.0000,8.0000) circle (1.2pt);
\fill[red!80!black] (0.0000,8.0000) circle (1.2pt);
\fill[red!80!black] (12.0000,8.0000) circle (1.2pt);
\fill[red!80!black] (0.0000,8.0000) circle (1.2pt);
\fill[red!80!black] (12.0000,8.0000) circle (1.2pt);
\fill[red!80!black] (0.0000,8.0000) circle (1.2pt);
\fill[red!80!black] (12.0000,8.0000) circle (1.2pt);
\fill[red!80!black] (0.0000,8.0000) circle (1.2pt);
\fill[red!80!black] (12.0000,8.0000) circle (1.2pt);
\fill[red!80!black] (0.0000,8.0000) circle (1.2pt);
\fill[red!80!black] (12.0000,8.0000) circle (1.2pt);
\fill[red!80!black] (0.0000,8.0000) circle (1.2pt);
\fill[red!80!black] (12.0000,8.0000) circle (1.2pt);
\fill[red!80!black] (0.0000,8.0000) circle (1.2pt);
\fill[red!80!black] (12.0000,8.0000) circle (1.2pt);
\fill[red!80!black] (0.0000,8.0000) circle (1.2pt);
\fill[red!80!black] (12.0000,8.0000) circle (1.2pt);
\fill[red!80!black] (0.0000,8.0000) circle (1.2pt);
\fill[red!80!black] (12.0000,8.0000) circle (1.2pt);
\fill[red!80!black] (0.0000,8.0000) circle (1.2pt);
\fill[red!80!black] (12.0000,8.0000) circle (1.2pt);
\fill[red!80!black] (0.0000,8.0000) circle (1.2pt);
\fill[red!80!black] (12.0000,8.0000) circle (1.2pt);
\fill[red!80!black] (0.0000,8.0000) circle (1.2pt);
\fill[red!80!black] (12.0000,8.0000) circle (1.2pt);
\fill[red!80!black] (0.0000,8.0000) circle (1.2pt);
\fill[red!80!black] (12.0000,8.0000) circle (1.2pt);
\fill[red!80!black] (0.0000,8.0000) circle (1.2pt);
\fill[red!80!black] (12.0000,8.0000) circle (1.2pt);
\fill[red!80!black] (0.0000,8.0000) circle (1.2pt);
\fill[red!80!black] (12.0000,8.0000) circle (1.2pt);
\fill[red!80!black] (0.0000,8.0000) circle (1.2pt);
\fill[red!80!black] (12.0000,8.0000) circle (1.2pt);
\fill[red!80!black] (0.0000,8.0000) circle (1.2pt);
\fill[red!80!black] (12.0000,8.0000) circle (1.2pt);
\fill[red!80!black] (0.0000,8.0000) circle (1.2pt);
\fill[red!80!black] (12.0000,8.0000) circle (1.2pt);
\fill[red!80!black] (0.0000,8.0000) circle (1.2pt);
\fill[red!80!black] (12.0000,8.0000) circle (1.2pt);
\fill[red!80!black] (0.0000,8.0000) circle (1.2pt);
\fill[red!80!black] (12.0000,8.0000) circle (1.2pt);
\fill[red!80!black] (0.0000,8.0000) circle (1.2pt);
\fill[red!80!black] (12.0000,8.0000) circle (1.2pt);
\fill[red!80!black] (0.0000,8.0000) circle (1.2pt);
\fill[red!80!black] (12.0000,8.0000) circle (1.2pt);
\fill[red!80!black] (0.0000,8.0000) circle (1.2pt);
\fill[red!80!black] (12.0000,8.0000) circle (1.2pt);
\fill[red!80!black] (0.0000,8.0000) circle (1.2pt);
\fill[red!80!black] (12.0000,8.0000) circle (1.2pt);
\fill[red!80!black] (0.0000,8.0000) circle (1.2pt);
\fill[red!80!black] (12.0000,8.0000) circle (1.2pt);
\fill[red!80!black] (0.0000,8.0000) circle (1.2pt);
\fill[red!80!black] (12.0000,8.0000) circle (1.2pt);
\fill[red!80!black] (0.0000,8.0000) circle (1.2pt);
\fill[red!80!black] (12.0000,8.0000) circle (1.2pt);
\fill[red!80!black] (0.0000,8.0000) circle (1.2pt);
\fill[red!80!black] (12.0000,8.0000) circle (1.2pt);
\fill[red!80!black] (0.0000,8.0000) circle (1.2pt);
\fill[red!80!black] (12.0000,8.0000) circle (1.2pt);
\fill[red!80!black] (0.0000,8.0000) circle (1.2pt);
\fill[red!80!black] (12.0000,8.0000) circle (1.2pt);
\fill[red!80!black] (0.0000,8.0000) circle (1.2pt);
\fill[red!80!black] (12.0000,8.0000) circle (1.2pt);
\fill[red!80!black] (0.0000,8.0000) circle (1.2pt);
\fill[red!80!black] (12.0000,8.0000) circle (1.2pt);
\fill[red!80!black] (0.0000,8.0000) circle (1.2pt);
\fill[red!80!black] (12.0000,8.0000) circle (1.2pt);
\fill[red!80!black] (0.0000,8.0000) circle (1.2pt);
\fill[red!80!black] (12.0000,8.0000) circle (1.2pt);
\fill[red!80!black] (0.0000,8.0000) circle (1.2pt);
\fill[red!80!black] (12.0000,8.0000) circle (1.2pt);
\fill[red!80!black] (0.0000,8.0000) circle (1.2pt);
\fill[red!80!black] (12.0000,8.0000) circle (1.2pt);
\fill[red!80!black] (0.0000,8.0000) circle (1.2pt);
\fill[red!80!black] (12.0000,8.0000) circle (1.2pt);
\fill[red!80!black] (0.0000,8.0000) circle (1.2pt);
\fill[red!80!black] (12.0000,8.0000) circle (1.2pt);
\fill[red!80!black] (0.0000,8.0000) circle (1.2pt);
\fill[red!80!black] (12.0000,8.0000) circle (1.2pt);
\fill[red!80!black] (0.0000,8.0000) circle (1.2pt);
\fill[red!80!black] (12.0000,8.0000) circle (1.2pt);
\fill[red!80!black] (0.0000,8.0000) circle (1.2pt);
\fill[red!80!black] (12.0000,8.0000) circle (1.2pt);
\fill[red!80!black] (0.0000,8.0000) circle (1.2pt);
\fill[red!80!black] (12.0000,8.0000) circle (1.2pt);
\fill[red!80!black] (0.0000,8.0000) circle (1.2pt);
\fill[red!80!black] (12.0000,8.0000) circle (1.2pt);
\fill[red!80!black] (0.0000,8.0000) circle (1.2pt);
\fill[red!80!black] (12.0000,8.0000) circle (1.2pt);
\fill[red!80!black] (0.0000,8.0000) circle (1.2pt);
\fill[red!80!black] (12.0000,8.0000) circle (1.2pt);
\fill[red!80!black] (0.0000,8.0000) circle (1.2pt);
\fill[red!80!black] (12.0000,8.0000) circle (1.2pt);
\fill[red!80!black] (0.0000,8.0000) circle (1.2pt);
\fill[red!80!black] (12.0000,8.0000) circle (1.2pt);
\fill[red!80!black] (0.0000,8.0000) circle (1.2pt);
\fill[red!80!black] (12.0000,8.0000) circle (1.2pt);
\fill[red!80!black] (0.0000,8.0000) circle (1.2pt);
\fill[red!80!black] (12.0000,8.0000) circle (1.2pt);
\fill[red!80!black] (0.0000,8.0000) circle (1.2pt);
\fill[red!80!black] (12.0000,8.0000) circle (1.2pt);
\fill[red!80!black] (0.0000,8.0000) circle (1.2pt);
\fill[red!80!black] (12.0000,8.0000) circle (1.2pt);
\fill[red!80!black] (0.0000,8.0000) circle (1.2pt);
\fill[red!80!black] (12.0000,8.0000) circle (1.2pt);
\fill[red!80!black] (0.0000,8.0000) circle (1.2pt);
\fill[red!80!black] (12.0000,8.0000) circle (1.2pt);
\fill[red!80!black] (0.0000,8.0000) circle (1.2pt);
\fill[red!80!black] (12.0000,8.0000) circle (1.2pt);
\fill[red!80!black] (0.0000,8.0000) circle (1.2pt);
\fill[red!80!black] (12.0000,8.0000) circle (1.2pt);
\fill[red!80!black] (0.0000,8.0000) circle (1.2pt);
\fill[red!80!black] (12.0000,8.0000) circle (1.2pt);
\fill[red!80!black] (0.0000,8.0000) circle (1.2pt);
\fill[red!80!black] (12.0000,8.0000) circle (1.2pt);
\fill[red!80!black] (0.0000,8.0000) circle (1.2pt);
\fill[red!80!black] (12.0000,8.0000) circle (1.2pt);
\fill[red!80!black] (0.0000,8.0000) circle (1.2pt);
\fill[red!80!black] (12.0000,8.0000) circle (1.2pt);
\fill[red!80!black] (0.0000,8.0000) circle (1.2pt);
\fill[red!80!black] (12.0000,8.0000) circle (1.2pt);
\fill[red!80!black] (0.0000,8.0000) circle (1.2pt);
\fill[red!80!black] (12.0000,8.0000) circle (1.2pt);
\fill[red!80!black] (0.0000,8.0000) circle (1.2pt);
\fill[red!80!black] (12.0000,8.0000) circle (1.2pt);
\fill[red!80!black] (0.0000,8.0000) circle (1.2pt);
\fill[red!80!black] (12.0000,8.0000) circle (1.2pt);
\fill[red!80!black] (0.0000,8.0000) circle (1.2pt);
\fill[red!80!black] (12.0000,8.0000) circle (1.2pt);
\fill[red!80!black] (0.0000,8.0000) circle (1.2pt);
\fill[red!80!black] (12.0000,8.0000) circle (1.2pt);
\fill[red!80!black] (0.0000,8.0000) circle (1.2pt);
\fill[red!80!black] (12.0000,8.0000) circle (1.2pt);
\fill[red!80!black] (0.0000,8.0000) circle (1.2pt);
\fill[red!80!black] (12.0000,8.0000) circle (1.2pt);
\fill[red!80!black] (0.0000,8.0000) circle (1.2pt);
\fill[red!80!black] (12.0000,8.0000) circle (1.2pt);
\fill[red!80!black] (0.0000,8.0000) circle (1.2pt);
\fill[red!80!black] (12.0000,8.0000) circle (1.2pt);
\fill[red!80!black] (0.0000,8.0000) circle (1.2pt);
\fill[red!80!black] (12.0000,8.0000) circle (1.2pt);
\fill[red!80!black] (0.0000,8.0000) circle (1.2pt);
\fill[red!80!black] (12.0000,8.0000) circle (1.2pt);
\fill[red!80!black] (0.0000,8.0000) circle (1.2pt);
\fill[red!80!black] (12.0000,8.0000) circle (1.2pt);
\fill[red!80!black] (0.0000,8.0000) circle (1.2pt);
\fill[red!80!black] (12.0000,8.0000) circle (1.2pt);
\fill[red!80!black] (0.0000,8.0000) circle (1.2pt);
\fill[red!80!black] (12.0000,8.0000) circle (1.2pt);
\fill[red!80!black] (0.0000,8.0000) circle (1.2pt);
\fill[red!80!black] (12.0000,8.0000) circle (1.2pt);
\fill[red!80!black] (0.0000,8.0000) circle (1.2pt);
\fill[red!80!black] (12.0000,8.0000) circle (1.2pt);
\fill[red!80!black] (0.0000,8.0000) circle (1.2pt);
\fill[red!80!black] (12.0000,8.0000) circle (1.2pt);
\fill[red!80!black] (0.0000,8.0000) circle (1.2pt);
\fill[red!80!black] (12.0000,8.0000) circle (1.2pt);
\fill[red!80!black] (0.0000,8.0000) circle (1.2pt);
\fill[red!80!black] (12.0000,8.0000) circle (1.2pt);
\fill[red!80!black] (0.0000,8.0000) circle (1.2pt);
\fill[red!80!black] (12.0000,8.0000) circle (1.2pt);
\fill[red!80!black] (0.0000,8.0000) circle (1.2pt);
\fill[red!80!black] (12.0000,8.0000) circle (1.2pt);
\fill[red!80!black] (0.0000,8.0000) circle (1.2pt);
\fill[red!80!black] (12.0000,8.0000) circle (1.2pt);
\fill[red!80!black] (0.0000,8.0000) circle (1.2pt);
\fill[red!80!black] (12.0000,8.0000) circle (1.2pt);
\fill[red!80!black] (0.0000,8.0000) circle (1.2pt);
\fill[red!80!black] (12.0000,8.0000) circle (1.2pt);
\fill[red!80!black] (0.0000,8.0000) circle (1.2pt);
\fill[red!80!black] (12.0000,8.0000) circle (1.2pt);
\fill[red!80!black] (0.0000,8.0000) circle (1.2pt);
\fill[red!80!black] (12.0000,8.0000) circle (1.2pt);
\fill[red!80!black] (0.0000,8.0000) circle (1.2pt);
\fill[red!80!black] (12.0000,8.0000) circle (1.2pt);
\fill[red!80!black] (0.0000,8.0000) circle (1.2pt);
\fill[red!80!black] (12.0000,8.0000) circle (1.2pt);
\fill[red!80!black] (0.0000,8.0000) circle (1.2pt);
\fill[red!80!black] (12.0000,8.0000) circle (1.2pt);
\fill[red!80!black] (0.0000,8.0000) circle (1.2pt);
\fill[red!80!black] (12.0000,8.0000) circle (1.2pt);
\fill[red!80!black] (0.0000,8.0000) circle (1.2pt);
\fill[red!80!black] (12.0000,8.0000) circle (1.2pt);
\fill[red!80!black] (0.0000,8.0000) circle (1.2pt);
\fill[red!80!black] (12.0000,8.0000) circle (1.2pt);
\fill[red!80!black] (0.0000,8.0000) circle (1.2pt);
\fill[red!80!black] (12.0000,8.0000) circle (1.2pt);
\fill[red!80!black] (0.0000,8.0000) circle (1.2pt);
\fill[red!80!black] (12.0000,8.0000) circle (1.2pt);
\fill[red!80!black] (0.0000,8.0000) circle (1.2pt);
\fill[red!80!black] (12.0000,8.0000) circle (1.2pt);
\fill[red!80!black] (0.0000,8.0000) circle (1.2pt);
\fill[red!80!black] (12.0000,8.0000) circle (1.2pt);
\fill[red!80!black] (0.0000,8.0000) circle (1.2pt);
\fill[red!80!black] (12.0000,8.0000) circle (1.2pt);
\fill[red!80!black] (0.0000,8.0000) circle (1.2pt);
\fill[red!80!black] (12.0000,8.0000) circle (1.2pt);
\fill[red!80!black] (0.0000,8.0000) circle (1.2pt);
\fill[red!80!black] (12.0000,8.0000) circle (1.2pt);
\fill[red!80!black] (0.0000,8.0000) circle (1.2pt);
\fill[red!80!black] (12.0000,8.0000) circle (1.2pt);
\fill[red!80!black] (0.0000,8.0000) circle (1.2pt);
\fill[red!80!black] (12.0000,8.0000) circle (1.2pt);
\fill[red!80!black] (0.0000,8.0000) circle (1.2pt);
\fill[red!80!black] (12.0000,8.0000) circle (1.2pt);
\fill[red!80!black] (0.0000,8.0000) circle (1.2pt);
\fill[red!80!black] (12.0000,8.0000) circle (1.2pt);
\fill[red!80!black] (0.0000,8.0000) circle (1.2pt);
\fill[red!80!black] (12.0000,8.0000) circle (1.2pt);
\fill[red!80!black] (0.0000,8.0000) circle (1.2pt);
\fill[red!80!black] (12.0000,8.0000) circle (1.2pt);
\fill[red!80!black] (0.0000,8.0000) circle (1.2pt);
\fill[red!80!black] (12.0000,8.0000) circle (1.2pt);
\fill[red!80!black] (0.0000,8.0000) circle (1.2pt);
\fill[red!80!black] (12.0000,8.0000) circle (1.2pt);
\fill[red!80!black] (0.0000,8.0000) circle (1.2pt);
\fill[red!80!black] (12.0000,8.0000) circle (1.2pt);
\fill[red!80!black] (0.0000,8.0000) circle (1.2pt);
\fill[red!80!black] (12.0000,8.0000) circle (1.2pt);
\fill[red!80!black] (0.0000,8.0000) circle (1.2pt);
\fill[red!80!black] (12.0000,8.0000) circle (1.2pt);
\fill[red!80!black] (0.0000,8.0000) circle (1.2pt);
\fill[red!80!black] (12.0000,8.0000) circle (1.2pt);
\fill[red!80!black] (0.0000,8.0000) circle (1.2pt);
\fill[red!80!black] (12.0000,8.0000) circle (1.2pt);
\fill[red!80!black] (0.0000,8.0000) circle (1.2pt);
\fill[red!80!black] (12.0000,8.0000) circle (1.2pt);
\fill[red!80!black] (0.0000,8.0000) circle (1.2pt);
\fill[red!80!black] (12.0000,8.0000) circle (1.2pt);
\fill[red!80!black] (0.0000,8.0000) circle (1.2pt);
\fill[red!80!black] (12.0000,8.0000) circle (1.2pt);
\fill[red!80!black] (0.0000,8.0000) circle (1.2pt);
\fill[red!80!black] (12.0000,8.0000) circle (1.2pt);
\fill[red!80!black] (0.0000,8.0000) circle (1.2pt);
\fill[red!80!black] (12.0000,8.0000) circle (1.2pt);
\fill[red!80!black] (0.0000,8.0000) circle (1.2pt);
\fill[red!80!black] (12.0000,8.0000) circle (1.2pt);
\fill[red!80!black] (0.0000,8.0000) circle (1.2pt);
\fill[red!80!black] (12.0000,8.0000) circle (1.2pt);
\fill[red!80!black] (0.0000,8.0000) circle (1.2pt);
\fill[red!80!black] (12.0000,8.0000) circle (1.2pt);
\fill[red!80!black] (0.0000,8.0000) circle (1.2pt);
\fill[red!80!black] (12.0000,8.0000) circle (1.2pt);
\fill[red!80!black] (0.0000,8.0000) circle (1.2pt);
\fill[red!80!black] (12.0000,8.0000) circle (1.2pt);
\fill[red!80!black] (0.0000,8.0000) circle (1.2pt);
\fill[red!80!black] (12.0000,8.0000) circle (1.2pt);
\fill[red!80!black] (0.0000,8.0000) circle (1.2pt);
\fill[red!80!black] (12.0000,8.0000) circle (1.2pt);
\fill[red!80!black] (0.0000,8.0000) circle (1.2pt);
\fill[red!80!black] (12.0000,8.0000) circle (1.2pt);
\fill[red!80!black] (0.0000,8.0000) circle (1.2pt);
\fill[red!80!black] (12.0000,8.0000) circle (1.2pt);
\fill[red!80!black] (0.0000,8.0000) circle (1.2pt);
\fill[red!80!black] (12.0000,8.0000) circle (1.2pt);
\fill[red!80!black] (0.0000,8.0000) circle (1.2pt);
\fill[red!80!black] (12.0000,8.0000) circle (1.2pt);
\fill[red!80!black] (0.0000,8.0000) circle (1.2pt);
\fill[red!80!black] (12.0000,8.0000) circle (1.2pt);
\fill[red!80!black] (0.0000,8.0000) circle (1.2pt);
\fill[red!80!black] (12.0000,8.0000) circle (1.2pt);
\fill[red!80!black] (0.0000,8.0000) circle (1.2pt);
\fill[red!80!black] (12.0000,8.0000) circle (1.2pt);
\fill[red!80!black] (0.0000,8.0000) circle (1.2pt);
\fill[red!80!black] (12.0000,8.0000) circle (1.2pt);
\fill[red!80!black] (0.0000,8.0000) circle (1.2pt);
\fill[red!80!black] (12.0000,8.0000) circle (1.2pt);
\fill[red!80!black] (0.0000,8.0000) circle (1.2pt);
\fill[red!80!black] (12.0000,8.0000) circle (1.2pt);
\fill[red!80!black] (0.0000,8.0000) circle (1.2pt);
\fill[red!80!black] (12.0000,8.0000) circle (1.2pt);
\fill[red!80!black] (0.0000,8.0000) circle (1.2pt);
\fill[red!80!black] (12.0000,8.0000) circle (1.2pt);
\fill[red!80!black] (0.0000,8.0000) circle (1.2pt);
\fill[red!80!black] (12.0000,8.0000) circle (1.2pt);
\fill[red!80!black] (0.0000,8.0000) circle (1.2pt);
\fill[red!80!black] (12.0000,8.0000) circle (1.2pt);
\fill[red!80!black] (0.0000,8.0000) circle (1.2pt);
\fill[red!80!black] (12.0000,8.0000) circle (1.2pt);
\fill[red!80!black] (0.0000,8.0000) circle (1.2pt);
\fill[red!80!black] (12.0000,8.0000) circle (1.2pt);
\fill[red!80!black] (0.0000,8.0000) circle (1.2pt);
\fill[red!80!black] (12.0000,8.0000) circle (1.2pt);
\fill[red!80!black] (0.0000,8.0000) circle (1.2pt);
\fill[red!80!black] (12.0000,8.0000) circle (1.2pt);
\fill[red!80!black] (0.0000,8.0000) circle (1.2pt);
\fill[red!80!black] (12.0000,8.0000) circle (1.2pt);
\fill[red!80!black] (0.0000,8.0000) circle (1.2pt);
\fill[red!80!black] (12.0000,8.0000) circle (1.2pt);
\fill[red!80!black] (0.0000,8.0000) circle (1.2pt);
\fill[red!80!black] (12.0000,8.0000) circle (1.2pt);
\fill[red!80!black] (0.0000,8.0000) circle (1.2pt);
\fill[red!80!black] (12.0000,8.0000) circle (1.2pt);
\fill[red!80!black] (0.0000,8.0000) circle (1.2pt);
\fill[red!80!black] (12.0000,8.0000) circle (1.2pt);
\fill[red!80!black] (0.0000,8.0000) circle (1.2pt);
\fill[red!80!black] (12.0000,8.0000) circle (1.2pt);
\fill[red!80!black] (0.0000,8.0000) circle (1.2pt);
\fill[red!80!black] (12.0000,8.0000) circle (1.2pt);
\fill[red!80!black] (0.0000,8.0000) circle (1.2pt);
\fill[red!80!black] (12.0000,8.0000) circle (1.2pt);
\fill[red!80!black] (0.0000,8.0000) circle (1.2pt);
\fill[red!80!black] (12.0000,8.0000) circle (1.2pt);
\fill[red!80!black] (0.0000,8.0000) circle (1.2pt);
\fill[red!80!black] (12.0000,8.0000) circle (1.2pt);
\fill[red!80!black] (0.0000,8.0000) circle (1.2pt);
\fill[red!80!black] (12.0000,8.0000) circle (1.2pt);
\fill[red!80!black] (0.0000,8.0000) circle (1.2pt);
\fill[red!80!black] (12.0000,8.0000) circle (1.2pt);
\fill[red!80!black] (0.0000,8.0000) circle (1.2pt);
\fill[red!80!black] (12.0000,8.0000) circle (1.2pt);
\fill[red!80!black] (0.0000,8.0000) circle (1.2pt);
\fill[red!80!black] (12.0000,8.0000) circle (1.2pt);
\fill[red!80!black] (0.0000,8.0000) circle (1.2pt);
\fill[red!80!black] (12.0000,8.0000) circle (1.2pt);
\fill[red!80!black] (0.0000,8.0000) circle (1.2pt);
\fill[red!80!black] (12.0000,8.0000) circle (1.2pt);
\fill[red!80!black] (0.0000,8.0000) circle (1.2pt);
\fill[red!80!black] (12.0000,8.0000) circle (1.2pt);
\fill[red!80!black] (0.0000,8.0000) circle (1.2pt);
\fill[red!80!black] (12.0000,8.0000) circle (1.2pt);
\fill[red!80!black] (0.0000,8.0000) circle (1.2pt);
\fill[red!80!black] (12.0000,8.0000) circle (1.2pt);
\fill[red!80!black] (0.0000,8.0000) circle (1.2pt);
\fill[red!80!black] (12.0000,8.0000) circle (1.2pt);
\fill[red!80!black] (0.0000,8.0000) circle (1.2pt);
\fill[red!80!black] (12.0000,8.0000) circle (1.2pt);
\fill[red!80!black] (0.0000,8.0000) circle (1.2pt);
\fill[red!80!black] (12.0000,8.0000) circle (1.2pt);
\fill[red!80!black] (0.0000,8.0000) circle (1.2pt);
\fill[red!80!black] (12.0000,8.0000) circle (1.2pt);
\fill[red!80!black] (0.0000,8.0000) circle (1.2pt);
\fill[red!80!black] (12.0000,8.0000) circle (1.2pt);
\fill[red!80!black] (0.0000,8.0000) circle (1.2pt);
\fill[red!80!black] (12.0000,8.0000) circle (1.2pt);
\fill[red!80!black] (0.0000,8.0000) circle (1.2pt);
\fill[red!80!black] (12.0000,8.0000) circle (1.2pt);
\fill[red!80!black] (0.0000,8.0000) circle (1.2pt);
\fill[red!80!black] (12.0000,8.0000) circle (1.2pt);
\fill[red!80!black] (0.0000,8.0000) circle (1.2pt);
\fill[red!80!black] (12.0000,8.0000) circle (1.2pt);
\fill[red!80!black] (0.0000,8.0000) circle (1.2pt);
\fill[red!80!black] (12.0000,8.0000) circle (1.2pt);
\fill[red!80!black] (0.0000,8.0000) circle (1.2pt);
\fill[red!80!black] (12.0000,8.0000) circle (1.2pt);
\fill[red!80!black] (0.0000,8.0000) circle (1.2pt);
\fill[red!80!black] (12.0000,8.0000) circle (1.2pt);
\fill[red!80!black] (0.0000,8.0000) circle (1.2pt);
\fill[red!80!black] (12.0000,8.0000) circle (1.2pt);
\fill[red!80!black] (0.0000,8.0000) circle (1.2pt);
\fill[red!80!black] (12.0000,8.0000) circle (1.2pt);
\fill[red!80!black] (0.0000,8.0000) circle (1.2pt);
\fill[red!80!black] (12.0000,8.0000) circle (1.2pt);
\fill[red!80!black] (0.0000,8.0000) circle (1.2pt);
\fill[red!80!black] (12.0000,8.0000) circle (1.2pt);
\fill[red!80!black] (0.0000,8.0000) circle (1.2pt);
\fill[red!80!black] (12.0000,8.0000) circle (1.2pt);
\fill[red!80!black] (0.0000,8.0000) circle (1.2pt);
\fill[red!80!black] (12.0000,8.0000) circle (1.2pt);
\fill[red!80!black] (0.0000,8.0000) circle (1.2pt);
\fill[red!80!black] (12.0000,8.0000) circle (1.2pt);
\fill[red!80!black] (0.0000,8.0000) circle (1.2pt);
\fill[red!80!black] (12.0000,8.0000) circle (1.2pt);
\fill[red!80!black] (0.0000,8.0000) circle (1.2pt);
\fill[red!80!black] (12.0000,8.0000) circle (1.2pt);
\fill[red!80!black] (0.0000,8.0000) circle (1.2pt);
\fill[red!80!black] (12.0000,8.0000) circle (1.2pt);
\fill[red!80!black] (0.0000,8.0000) circle (1.2pt);
\fill[red!80!black] (12.0000,8.0000) circle (1.2pt);
\fill[red!80!black] (0.0000,8.0000) circle (1.2pt);
\fill[red!80!black] (12.0000,8.0000) circle (1.2pt);
\fill[red!80!black] (0.0000,8.0000) circle (1.2pt);
\fill[red!80!black] (12.0000,8.0000) circle (1.2pt);
\fill[red!80!black] (0.0000,8.0000) circle (1.2pt);
\fill[red!80!black] (12.0000,8.0000) circle (1.2pt);
\fill[red!80!black] (0.0000,8.0000) circle (1.2pt);
\fill[red!80!black] (12.0000,8.0000) circle (1.2pt);
\fill[red!80!black] (0.0000,8.0000) circle (1.2pt);
\fill[red!80!black] (12.0000,8.0000) circle (1.2pt);
\fill[red!80!black] (0.0000,8.0000) circle (1.2pt);
\fill[red!80!black] (12.0000,8.0000) circle (1.2pt);
\fill[red!80!black] (0.0000,8.0000) circle (1.2pt);
\fill[red!80!black] (12.0000,8.0000) circle (1.2pt);
\fill[red!80!black] (0.0000,8.0000) circle (1.2pt);
\fill[red!80!black] (12.0000,8.0000) circle (1.2pt);
\fill[red!80!black] (0.0000,8.0000) circle (1.2pt);
\fill[red!80!black] (12.0000,8.0000) circle (1.2pt);
\fill[red!80!black] (0.0000,8.0000) circle (1.2pt);
\fill[red!80!black] (12.0000,8.0000) circle (1.2pt);
\fill[red!80!black] (0.0000,8.0000) circle (1.2pt);
\fill[red!80!black] (12.0000,8.0000) circle (1.2pt);
\fill[red!80!black] (0.0000,8.0000) circle (1.2pt);
\fill[red!80!black] (12.0000,8.0000) circle (1.2pt);
\fill[red!80!black] (0.0000,8.0000) circle (1.2pt);
\fill[red!80!black] (12.0000,8.0000) circle (1.2pt);
\fill[red!80!black] (0.0000,8.0000) circle (1.2pt);
\fill[red!80!black] (12.0000,8.0000) circle (1.2pt);
\fill[red!80!black] (0.0000,8.0000) circle (1.2pt);
\fill[red!80!black] (12.0000,8.0000) circle (1.2pt);
\fill[red!80!black] (0.0000,8.0000) circle (1.2pt);
\fill[red!80!black] (12.0000,8.0000) circle (1.2pt);
\fill[red!80!black] (0.0000,8.0000) circle (1.2pt);
\fill[red!80!black] (12.0000,8.0000) circle (1.2pt);
\fill[red!80!black] (0.0000,8.0000) circle (1.2pt);
\fill[red!80!black] (12.0000,8.0000) circle (1.2pt);
\fill[red!80!black] (0.0000,8.0000) circle (1.2pt);
\fill[red!80!black] (12.0000,8.0000) circle (1.2pt);
\fill[red!80!black] (0.0000,8.0000) circle (1.2pt);
\fill[red!80!black] (12.0000,8.0000) circle (1.2pt);
\fill[red!80!black] (0.0000,8.0000) circle (1.2pt);
\fill[red!80!black] (12.0000,8.0000) circle (1.2pt);
\fill[red!80!black] (0.0000,8.0000) circle (1.2pt);
\fill[red!80!black] (12.0000,8.0000) circle (1.2pt);
\fill[red!80!black] (0.0000,8.0000) circle (1.2pt);
\fill[red!80!black] (12.0000,8.0000) circle (1.2pt);
\fill[red!80!black] (0.0000,8.0000) circle (1.2pt);
\fill[red!80!black] (12.0000,8.0000) circle (1.2pt);
\fill[red!80!black] (0.0000,8.0000) circle (1.2pt);
\fill[red!80!black] (12.0000,8.0000) circle (1.2pt);
\fill[red!80!black] (0.0000,8.0000) circle (1.2pt);
\fill[red!80!black] (12.0000,8.0000) circle (1.2pt);
\fill[red!80!black] (0.0000,8.0000) circle (1.2pt);
\fill[red!80!black] (12.0000,8.0000) circle (1.2pt);
\fill[red!80!black] (0.0000,8.0000) circle (1.2pt);
\fill[red!80!black] (12.0000,8.0000) circle (1.2pt);
\fill[red!80!black] (0.0000,8.0000) circle (1.2pt);
\fill[red!80!black] (12.0000,8.0000) circle (1.2pt);
\fill[red!80!black] (0.0000,8.0000) circle (1.2pt);
\fill[red!80!black] (12.0000,8.0000) circle (1.2pt);
\fill[red!80!black] (0.0000,8.0000) circle (1.2pt);
\fill[red!80!black] (12.0000,8.0000) circle (1.2pt);
\fill[red!80!black] (0.0000,8.0000) circle (1.2pt);
\fill[red!80!black] (12.0000,8.0000) circle (1.2pt);
\fill[red!80!black] (0.0000,8.0000) circle (1.2pt);
\fill[red!80!black] (12.0000,8.0000) circle (1.2pt);
\fill[red!80!black] (0.0000,8.0000) circle (1.2pt);
\fill[red!80!black] (12.0000,8.0000) circle (1.2pt);
\fill[red!80!black] (0.0000,8.0000) circle (1.2pt);
\fill[red!80!black] (12.0000,8.0000) circle (1.2pt);
\fill[red!80!black] (0.0000,8.0000) circle (1.2pt);
\fill[red!80!black] (12.0000,8.0000) circle (1.2pt);
\fill[red!80!black] (0.0000,8.0000) circle (1.2pt);
\fill[red!80!black] (12.0000,8.0000) circle (1.2pt);
\fill[red!80!black] (0.0000,8.0000) circle (1.2pt);
\fill[red!80!black] (12.0000,8.0000) circle (1.2pt);
\fill[red!80!black] (0.0000,8.0000) circle (1.2pt);
\fill[red!80!black] (12.0000,8.0000) circle (1.2pt);
\fill[red!80!black] (0.0000,8.0000) circle (1.2pt);
\fill[red!80!black] (12.0000,8.0000) circle (1.2pt);
\fill[red!80!black] (0.0000,8.0000) circle (1.2pt);
\fill[red!80!black] (12.0000,8.0000) circle (1.2pt);
\fill[red!80!black] (0.0000,8.0000) circle (1.2pt);
\fill[red!80!black] (12.0000,8.0000) circle (1.2pt);
\fill[red!80!black] (0.0000,8.0000) circle (1.2pt);
\fill[red!80!black] (12.0000,8.0000) circle (1.2pt);
\fill[red!80!black] (0.0000,8.0000) circle (1.2pt);
\fill[red!80!black] (12.0000,8.0000) circle (1.2pt);
\fill[red!80!black] (0.0000,8.0000) circle (1.2pt);
\fill[red!80!black] (12.0000,8.0000) circle (1.2pt);
\fill[red!80!black] (0.0000,8.0000) circle (1.2pt);
\fill[red!80!black] (12.0000,8.0000) circle (1.2pt);
\fill[red!80!black] (0.0000,8.0000) circle (1.2pt);
\fill[red!80!black] (12.0000,8.0000) circle (1.2pt);
\fill[red!80!black] (0.0000,8.0000) circle (1.2pt);
\fill[red!80!black] (12.0000,8.0000) circle (1.2pt);
\fill[red!80!black] (0.0000,8.0000) circle (1.2pt);
\fill[red!80!black] (12.0000,8.0000) circle (1.2pt);
\fill[red!80!black] (0.0000,8.0000) circle (1.2pt);
\fill[red!80!black] (12.0000,8.0000) circle (1.2pt);
\fill[red!80!black] (0.0000,8.0000) circle (1.2pt);
\fill[red!80!black] (12.0000,8.0000) circle (1.2pt);
\fill[red!80!black] (0.0000,8.0000) circle (1.2pt);
\fill[red!80!black] (12.0000,8.0000) circle (1.2pt);
\fill[red!80!black] (0.0000,8.0000) circle (1.2pt);
\fill[red!80!black] (12.0000,8.0000) circle (1.2pt);
\fill[red!80!black] (0.0000,8.0000) circle (1.2pt);
\fill[red!80!black] (12.0000,8.0000) circle (1.2pt);
\fill[red!80!black] (0.0000,8.0000) circle (1.2pt);
\fill[red!80!black] (12.0000,8.0000) circle (1.2pt);
\fill[red!80!black] (0.0000,8.0000) circle (1.2pt);
\fill[red!80!black] (12.0000,8.0000) circle (1.2pt);
\fill[red!80!black] (0.0000,8.0000) circle (1.2pt);
\fill[red!80!black] (12.0000,8.0000) circle (1.2pt);
\fill[red!80!black] (0.0000,8.0000) circle (1.2pt);
\fill[red!80!black] (12.0000,8.0000) circle (1.2pt);
\fill[red!80!black] (0.0000,8.0000) circle (1.2pt);
\fill[red!80!black] (12.0000,8.0000) circle (1.2pt);
\fill[red!80!black] (0.0000,8.0000) circle (1.2pt);
\fill[red!80!black] (12.0000,8.0000) circle (1.2pt);
\fill[red!80!black] (0.0000,8.0000) circle (1.2pt);
\fill[red!80!black] (12.0000,8.0000) circle (1.2pt);
\fill[red!80!black] (0.0000,8.0000) circle (1.2pt);
\fill[red!80!black] (12.0000,8.0000) circle (1.2pt);
\fill[red!80!black] (0.0000,8.0000) circle (1.2pt);
\fill[red!80!black] (12.0000,8.0000) circle (1.2pt);
\fill[red!80!black] (0.0000,8.0000) circle (1.2pt);
\fill[red!80!black] (12.0000,8.0000) circle (1.2pt);
\fill[red!80!black] (0.0000,8.0000) circle (1.2pt);
\fill[red!80!black] (12.0000,8.0000) circle (1.2pt);
\fill[red!80!black] (0.0000,8.0000) circle (1.2pt);
\fill[red!80!black] (12.0000,8.0000) circle (1.2pt);
\fill[red!80!black] (0.0000,8.0000) circle (1.2pt);
\fill[red!80!black] (12.0000,8.0000) circle (1.2pt);
\fill[red!80!black] (0.0000,8.0000) circle (1.2pt);
\fill[red!80!black] (12.0000,8.0000) circle (1.2pt);
\fill[red!80!black] (0.0000,8.0000) circle (1.2pt);
\fill[red!80!black] (12.0000,8.0000) circle (1.2pt);
\fill[red!80!black] (0.0000,8.0000) circle (1.2pt);
\fill[red!80!black] (12.0000,8.0000) circle (1.2pt);
\fill[red!80!black] (0.0000,8.0000) circle (1.2pt);
\fill[red!80!black] (12.0000,8.0000) circle (1.2pt);
\fill[red!80!black] (0.0000,8.0000) circle (1.2pt);
\fill[red!80!black] (12.0000,8.0000) circle (1.2pt);
\fill[red!80!black] (0.0000,8.0000) circle (1.2pt);
\fill[red!80!black] (12.0000,8.0000) circle (1.2pt);
\fill[red!80!black] (0.0000,8.0000) circle (1.2pt);
\fill[red!80!black] (12.0000,8.0000) circle (1.2pt);
\fill[red!80!black] (0.0000,8.0000) circle (1.2pt);
\fill[red!80!black] (12.0000,8.0000) circle (1.2pt);
\fill[red!80!black] (0.0000,8.0000) circle (1.2pt);
\fill[red!80!black] (12.0000,8.0000) circle (1.2pt);
\fill[red!80!black] (0.0000,8.0000) circle (1.2pt);
\fill[red!80!black] (12.0000,8.0000) circle (1.2pt);
\fill[red!80!black] (0.0000,8.0000) circle (1.2pt);
\fill[red!80!black] (12.0000,8.0000) circle (1.2pt);
\fill[red!80!black] (0.0000,8.0000) circle (1.2pt);
\fill[red!80!black] (12.0000,8.0000) circle (1.2pt);
\fill[red!80!black] (0.0000,8.0000) circle (1.2pt);
\fill[red!80!black] (12.0000,8.0000) circle (1.2pt);
\fill[red!80!black] (0.0000,8.0000) circle (1.2pt);
\fill[red!80!black] (12.0000,8.0000) circle (1.2pt);
\fill[red!80!black] (0.0000,8.0000) circle (1.2pt);
\fill[red!80!black] (12.0000,8.0000) circle (1.2pt);
\fill[red!80!black] (0.0000,8.0000) circle (1.2pt);
\fill[red!80!black] (12.0000,8.0000) circle (1.2pt);
\fill[red!80!black] (0.0000,8.0000) circle (1.2pt);
\fill[red!80!black] (12.0000,8.0000) circle (1.2pt);
\fill[red!80!black] (0.0000,8.0000) circle (1.2pt);
\fill[red!80!black] (12.0000,8.0000) circle (1.2pt);
\fill[red!80!black] (0.0000,8.0000) circle (1.2pt);
\fill[red!80!black] (12.0000,8.0000) circle (1.2pt);
\fill[red!80!black] (0.0000,8.0000) circle (1.2pt);
\fill[red!80!black] (12.0000,8.0000) circle (1.2pt);
\fill[red!80!black] (0.0000,8.0000) circle (1.2pt);
\fill[red!80!black] (12.0000,8.0000) circle (1.2pt);
\fill[red!80!black] (0.0000,8.0000) circle (1.2pt);
\fill[red!80!black] (12.0000,8.0000) circle (1.2pt);
\fill[red!80!black] (0.0000,8.0000) circle (1.2pt);
\fill[red!80!black] (12.0000,8.0000) circle (1.2pt);
\fill[red!80!black] (0.0000,8.0000) circle (1.2pt);
\fill[red!80!black] (12.0000,8.0000) circle (1.2pt);
\fill[red!80!black] (0.0000,8.0000) circle (1.2pt);
\fill[red!80!black] (12.0000,8.0000) circle (1.2pt);
\fill[red!80!black] (0.0000,8.0000) circle (1.2pt);
\fill[red!80!black] (12.0000,8.0000) circle (1.2pt);
\fill[red!80!black] (0.0000,8.0000) circle (1.2pt);
\fill[red!80!black] (12.0000,8.0000) circle (1.2pt);
\fill[red!80!black] (0.0000,8.0000) circle (1.2pt);
\fill[red!80!black] (12.0000,8.0000) circle (1.2pt);
\fill[red!80!black] (0.0000,8.0000) circle (1.2pt);
\fill[red!80!black] (12.0000,8.0000) circle (1.2pt);
\fill[red!80!black] (0.0000,8.0000) circle (1.2pt);
\fill[red!80!black] (12.0000,8.0000) circle (1.2pt);
\fill[red!80!black] (0.0000,8.0000) circle (1.2pt);
\fill[red!80!black] (12.0000,8.0000) circle (1.2pt);
\fill[red!80!black] (0.0000,8.0000) circle (1.2pt);
\fill[red!80!black] (12.0000,8.0000) circle (1.2pt);
\fill[red!80!black] (0.0000,8.0000) circle (1.2pt);
\fill[red!80!black] (12.0000,8.0000) circle (1.2pt);
\fill[red!80!black] (0.0000,8.0000) circle (1.2pt);
\fill[red!80!black] (12.0000,8.0000) circle (1.2pt);
\fill[red!80!black] (0.0000,8.0000) circle (1.2pt);
\fill[red!80!black] (12.0000,8.0000) circle (1.2pt);
\fill[red!80!black] (0.0000,8.0000) circle (1.2pt);
\fill[red!80!black] (12.0000,8.0000) circle (1.2pt);
\fill[red!80!black] (0.0000,8.0000) circle (1.2pt);
\fill[red!80!black] (12.0000,8.0000) circle (1.2pt);
\fill[red!80!black] (0.0000,8.0000) circle (1.2pt);
\fill[red!80!black] (12.0000,8.0000) circle (1.2pt);
\fill[red!80!black] (0.0000,8.0000) circle (1.2pt);
\fill[red!80!black] (12.0000,8.0000) circle (1.2pt);
\fill[red!80!black] (0.0000,8.0000) circle (1.2pt);
\fill[red!80!black] (12.0000,8.0000) circle (1.2pt);
\fill[red!80!black] (0.0000,8.0000) circle (1.2pt);
\fill[red!80!black] (12.0000,8.0000) circle (1.2pt);
\fill[red!80!black] (0.0000,8.0000) circle (1.2pt);
\fill[red!80!black] (12.0000,8.0000) circle (1.2pt);
\fill[red!80!black] (0.0000,8.0000) circle (1.2pt);
\fill[red!80!black] (12.0000,8.0000) circle (1.2pt);
\fill[red!80!black] (0.0000,8.0000) circle (1.2pt);
\fill[red!80!black] (12.0000,8.0000) circle (1.2pt);
\fill[red!80!black] (0.0000,8.0000) circle (1.2pt);
\fill[red!80!black] (12.0000,8.0000) circle (1.2pt);
\fill[red!80!black] (0.0000,8.0000) circle (1.2pt);
\fill[red!80!black] (12.0000,8.0000) circle (1.2pt);
\fill[red!80!black] (0.0000,8.0000) circle (1.2pt);
\fill[red!80!black] (12.0000,8.0000) circle (1.2pt);
\fill[red!80!black] (0.0000,8.0000) circle (1.2pt);
\fill[red!80!black] (12.0000,8.0000) circle (1.2pt);
\fill[red!80!black] (0.0000,8.0000) circle (1.2pt);
\fill[red!80!black] (12.0000,8.0000) circle (1.2pt);
\fill[red!80!black] (0.0000,8.0000) circle (1.2pt);
\fill[red!80!black] (12.0000,8.0000) circle (1.2pt);
\fill[red!80!black] (0.0000,8.0000) circle (1.2pt);
\fill[red!80!black] (12.0000,8.0000) circle (1.2pt);
\fill[red!80!black] (0.0000,8.0000) circle (1.2pt);
\fill[red!80!black] (12.0000,8.0000) circle (1.2pt);
\fill[red!80!black] (0.0000,8.0000) circle (1.2pt);
\fill[red!80!black] (12.0000,8.0000) circle (1.2pt);
\fill[red!80!black] (0.0000,8.0000) circle (1.2pt);
\fill[red!80!black] (12.0000,8.0000) circle (1.2pt);
\fill[red!80!black] (0.0000,8.0000) circle (1.2pt);
\fill[red!80!black] (12.0000,8.0000) circle (1.2pt);
\fill[red!80!black] (0.0000,8.0000) circle (1.2pt);
\fill[red!80!black] (12.0000,8.0000) circle (1.2pt);
\fill[red!80!black] (0.0000,8.0000) circle (1.2pt);
\fill[red!80!black] (12.0000,8.0000) circle (1.2pt);
\fill[red!80!black] (0.0000,8.0000) circle (1.2pt);
\fill[red!80!black] (12.0000,8.0000) circle (1.2pt);
\fill[red!80!black] (0.0000,8.0000) circle (1.2pt);
\fill[red!80!black] (12.0000,8.0000) circle (1.2pt);
\fill[red!80!black] (0.0000,8.0000) circle (1.2pt);
\fill[red!80!black] (12.0000,8.0000) circle (1.2pt);
\fill[red!80!black] (0.0000,8.0000) circle (1.2pt);
\fill[red!80!black] (12.0000,8.0000) circle (1.2pt);
\fill[red!80!black] (0.0000,8.0000) circle (1.2pt);
\fill[red!80!black] (12.0000,8.0000) circle (1.2pt);
\fill[red!80!black] (0.0000,8.0000) circle (1.2pt);
\fill[red!80!black] (12.0000,8.0000) circle (1.2pt);
\fill[red!80!black] (0.0000,8.0000) circle (1.2pt);
\fill[red!80!black] (12.0000,8.0000) circle (1.2pt);
\fill[red!80!black] (0.0000,8.0000) circle (1.2pt);
\fill[red!80!black] (12.0000,8.0000) circle (1.2pt);
\fill[red!80!black] (0.0000,8.0000) circle (1.2pt);
\fill[red!80!black] (12.0000,8.0000) circle (1.2pt);
\fill[red!80!black] (0.0000,8.0000) circle (1.2pt);
\fill[red!80!black] (12.0000,8.0000) circle (1.2pt);
\fill[red!80!black] (0.0000,8.0000) circle (1.2pt);
\fill[red!80!black] (12.0000,8.0000) circle (1.2pt);
\fill[red!80!black] (0.0000,8.0000) circle (1.2pt);
\fill[red!80!black] (12.0000,8.0000) circle (1.2pt);
\fill[red!80!black] (0.0000,8.0000) circle (1.2pt);
\fill[red!80!black] (12.0000,8.0000) circle (1.2pt);
\fill[red!80!black] (0.0000,8.0000) circle (1.2pt);
\fill[red!80!black] (12.0000,8.0000) circle (1.2pt);
\fill[red!80!black] (0.0000,8.0000) circle (1.2pt);
\fill[red!80!black] (12.0000,8.0000) circle (1.2pt);
\fill[red!80!black] (0.0000,8.0000) circle (1.2pt);
\fill[red!80!black] (12.0000,8.0000) circle (1.2pt);
\fill[red!80!black] (0.0000,8.0000) circle (1.2pt);
\fill[red!80!black] (12.0000,8.0000) circle (1.2pt);
\fill[red!80!black] (0.0000,8.0000) circle (1.2pt);
\fill[red!80!black] (12.0000,8.0000) circle (1.2pt);
\fill[red!80!black] (0.0000,8.0000) circle (1.2pt);
\fill[red!80!black] (12.0000,8.0000) circle (1.2pt);
\fill[red!80!black] (0.0000,8.0000) circle (1.2pt);
\fill[red!80!black] (12.0000,8.0000) circle (1.2pt);
\fill[red!80!black] (0.0000,8.0000) circle (1.2pt);
\fill[red!80!black] (12.0000,8.0000) circle (1.2pt);
\fill[red!80!black] (0.0000,8.0000) circle (1.2pt);
\fill[red!80!black] (12.0000,8.0000) circle (1.2pt);
\fill[red!80!black] (0.0000,8.0000) circle (1.2pt);
\fill[red!80!black] (12.0000,8.0000) circle (1.2pt);
\fill[red!80!black] (0.0000,8.0000) circle (1.2pt);
\fill[red!80!black] (12.0000,8.0000) circle (1.2pt);
\fill[red!80!black] (0.0000,8.0000) circle (1.2pt);
\fill[red!80!black] (12.0000,8.0000) circle (1.2pt);
\fill[red!80!black] (0.0000,8.0000) circle (1.2pt);
\fill[red!80!black] (12.0000,8.0000) circle (1.2pt);
\fill[red!80!black] (0.0000,8.0000) circle (1.2pt);
\fill[red!80!black] (12.0000,8.0000) circle (1.2pt);
\fill[red!80!black] (0.0000,8.0000) circle (1.2pt);
\fill[red!80!black] (12.0000,8.0000) circle (1.2pt);
\fill[red!80!black] (0.0000,8.0000) circle (1.2pt);
\fill[red!80!black] (12.0000,8.0000) circle (1.2pt);
\fill[red!80!black] (0.0000,8.0000) circle (1.2pt);
\fill[red!80!black] (12.0000,8.0000) circle (1.2pt);
\fill[red!80!black] (0.0000,8.0000) circle (1.2pt);
\fill[red!80!black] (12.0000,8.0000) circle (1.2pt);
\fill[red!80!black] (0.0000,8.0000) circle (1.2pt);
\fill[red!80!black] (12.0000,8.0000) circle (1.2pt);
\fill[red!80!black] (0.0000,8.0000) circle (1.2pt);
\fill[red!80!black] (12.0000,8.0000) circle (1.2pt);
\fill[red!80!black] (0.0000,8.0000) circle (1.2pt);
\fill[red!80!black] (12.0000,8.0000) circle (1.2pt);
\fill[red!80!black] (0.0000,8.0000) circle (1.2pt);
\fill[red!80!black] (12.0000,8.0000) circle (1.2pt);
\fill[red!80!black] (0.0000,8.0000) circle (1.2pt);
\fill[red!80!black] (12.0000,8.0000) circle (1.2pt);
\fill[red!80!black] (0.0000,8.0000) circle (1.2pt);
\fill[red!80!black] (12.0000,8.0000) circle (1.2pt);
\fill[red!80!black] (0.0000,8.0000) circle (1.2pt);
\fill[red!80!black] (12.0000,8.0000) circle (1.2pt);
\fill[red!80!black] (0.0000,8.0000) circle (1.2pt);
\fill[red!80!black] (12.0000,8.0000) circle (1.2pt);
\fill[red!80!black] (0.0000,8.0000) circle (1.2pt);
\fill[red!80!black] (12.0000,8.0000) circle (1.2pt);
\fill[red!80!black] (0.0000,8.0000) circle (1.2pt);
\fill[red!80!black] (12.0000,8.0000) circle (1.2pt);
\fill[red!80!black] (0.0000,8.0000) circle (1.2pt);
\fill[red!80!black] (12.0000,8.0000) circle (1.2pt);
\fill[red!80!black] (0.0000,8.0000) circle (1.2pt);
\fill[red!80!black] (12.0000,8.0000) circle (1.2pt);
\fill[red!80!black] (0.0000,8.0000) circle (1.2pt);
\fill[red!80!black] (12.0000,8.0000) circle (1.2pt);
\fill[red!80!black] (0.0000,8.0000) circle (1.2pt);
\fill[red!80!black] (12.0000,8.0000) circle (1.2pt);
\fill[red!80!black] (0.0000,8.0000) circle (1.2pt);
\fill[red!80!black] (12.0000,8.0000) circle (1.2pt);
\fill[red!80!black] (0.0000,8.0000) circle (1.2pt);
\fill[red!80!black] (12.0000,8.0000) circle (1.2pt);
\fill[red!80!black] (0.0000,8.0000) circle (1.2pt);
\fill[red!80!black] (12.0000,8.0000) circle (1.2pt);
\fill[red!80!black] (0.0000,8.0000) circle (1.2pt);
\fill[red!80!black] (12.0000,8.0000) circle (1.2pt);
\fill[red!80!black] (0.0000,8.0000) circle (1.2pt);
\fill[red!80!black] (12.0000,8.0000) circle (1.2pt);
\fill[red!80!black] (0.0000,8.0000) circle (1.2pt);
\fill[red!80!black] (12.0000,8.0000) circle (1.2pt);
\fill[red!80!black] (0.0000,8.0000) circle (1.2pt);
\fill[red!80!black] (12.0000,8.0000) circle (1.2pt);
\fill[red!80!black] (0.0000,8.0000) circle (1.2pt);
\fill[red!80!black] (12.0000,8.0000) circle (1.2pt);
\fill[red!80!black] (0.0000,8.0000) circle (1.2pt);
\fill[red!80!black] (12.0000,8.0000) circle (1.2pt);
\fill[red!80!black] (0.0000,8.0000) circle (1.2pt);
\fill[red!80!black] (12.0000,8.0000) circle (1.2pt);
\fill[red!80!black] (0.0000,8.0000) circle (1.2pt);
\fill[red!80!black] (12.0000,8.0000) circle (1.2pt);
\fill[red!80!black] (0.0000,8.0000) circle (1.2pt);
\fill[red!80!black] (12.0000,8.0000) circle (1.2pt);
\fill[red!80!black] (0.0000,8.0000) circle (1.2pt);
\fill[red!80!black] (12.0000,8.0000) circle (1.2pt);
\fill[red!80!black] (0.0000,8.0000) circle (1.2pt);
\fill[red!80!black] (12.0000,8.0000) circle (1.2pt);
\fill[red!80!black] (0.0000,8.0000) circle (1.2pt);
\fill[red!80!black] (12.0000,8.0000) circle (1.2pt);
\fill[red!80!black] (0.0000,8.0000) circle (1.2pt);
\fill[red!80!black] (12.0000,8.0000) circle (1.2pt);
\fill[red!80!black] (0.0000,8.0000) circle (1.2pt);
\fill[red!80!black] (12.0000,8.0000) circle (1.2pt);
\fill[red!80!black] (0.0000,8.0000) circle (1.2pt);
\fill[red!80!black] (12.0000,8.0000) circle (1.2pt);
\fill[red!80!black] (0.0000,8.0000) circle (1.2pt);
\fill[red!80!black] (12.0000,8.0000) circle (1.2pt);
\fill[red!80!black] (0.0000,8.0000) circle (1.2pt);
\fill[red!80!black] (12.0000,8.0000) circle (1.2pt);
\fill[red!80!black] (0.0000,8.0000) circle (1.2pt);
\fill[red!80!black] (12.0000,8.0000) circle (1.2pt);
\fill[red!80!black] (0.0000,8.0000) circle (1.2pt);
\fill[red!80!black] (12.0000,8.0000) circle (1.2pt);
\fill[red!80!black] (0.0000,8.0000) circle (1.2pt);
\fill[red!80!black] (12.0000,8.0000) circle (1.2pt);
\fill[red!80!black] (0.0000,8.0000) circle (1.2pt);
\fill[red!80!black] (12.0000,8.0000) circle (1.2pt);
\fill[red!80!black] (0.0000,8.0000) circle (1.2pt);
\fill[red!80!black] (12.0000,8.0000) circle (1.2pt);
\fill[red!80!black] (0.0000,8.0000) circle (1.2pt);
\fill[red!80!black] (12.0000,8.0000) circle (1.2pt);
\fill[red!80!black] (0.0000,8.0000) circle (1.2pt);
\fill[red!80!black] (12.0000,8.0000) circle (1.2pt);
\fill[red!80!black] (0.0000,8.0000) circle (1.2pt);
\fill[red!80!black] (12.0000,8.0000) circle (1.2pt);
\fill[red!80!black] (0.0000,8.0000) circle (1.2pt);
\fill[red!80!black] (12.0000,8.0000) circle (1.2pt);
\fill[red!80!black] (0.0000,8.0000) circle (1.2pt);
\fill[red!80!black] (12.0000,8.0000) circle (1.2pt);
\fill[red!80!black] (0.0000,8.0000) circle (1.2pt);
\fill[red!80!black] (12.0000,8.0000) circle (1.2pt);
\fill[red!80!black] (0.0000,8.0000) circle (1.2pt);
\fill[red!80!black] (12.0000,8.0000) circle (1.2pt);
\fill[red!80!black] (0.0000,8.0000) circle (1.2pt);
\fill[red!80!black] (12.0000,8.0000) circle (1.2pt);
\fill[red!80!black] (0.0000,8.0000) circle (1.2pt);
\fill[red!80!black] (12.0000,8.0000) circle (1.2pt);
\fill[red!80!black] (0.0000,8.0000) circle (1.2pt);
\fill[red!80!black] (12.0000,8.0000) circle (1.2pt);
\fill[red!80!black] (0.0000,8.0000) circle (1.2pt);
\fill[red!80!black] (12.0000,8.0000) circle (1.2pt);
\fill[red!80!black] (0.0000,8.0000) circle (1.2pt);
\fill[red!80!black] (12.0000,8.0000) circle (1.2pt);
\fill[red!80!black] (0.0000,8.0000) circle (1.2pt);
\fill[red!80!black] (12.0000,8.0000) circle (1.2pt);
\fill[red!80!black] (0.0000,8.0000) circle (1.2pt);
\fill[red!80!black] (12.0000,8.0000) circle (1.2pt);
\fill[red!80!black] (0.0000,8.0000) circle (1.2pt);
\fill[red!80!black] (12.0000,8.0000) circle (1.2pt);
\fill[red!80!black] (0.0000,8.0000) circle (1.2pt);
\fill[red!80!black] (12.0000,8.0000) circle (1.2pt);
\fill[red!80!black] (0.0000,8.0000) circle (1.2pt);
\fill[red!80!black] (12.0000,8.0000) circle (1.2pt);
\fill[red!80!black] (0.0000,8.0000) circle (1.2pt);
\fill[red!80!black] (12.0000,8.0000) circle (1.2pt);
\fill[red!80!black] (0.0000,8.0000) circle (1.2pt);
\fill[red!80!black] (12.0000,8.0000) circle (1.2pt);
\fill[red!80!black] (0.0000,8.0000) circle (1.2pt);
\fill[red!80!black] (12.0000,8.0000) circle (1.2pt);
\fill[red!80!black] (0.0000,8.0000) circle (1.2pt);
\fill[red!80!black] (12.0000,8.0000) circle (1.2pt);
\fill[red!80!black] (0.0000,8.0000) circle (1.2pt);
\fill[red!80!black] (12.0000,8.0000) circle (1.2pt);
\fill[red!80!black] (0.0000,8.0000) circle (1.2pt);
\fill[red!80!black] (12.0000,8.0000) circle (1.2pt);
\fill[red!80!black] (0.0000,8.0000) circle (1.2pt);
\fill[red!80!black] (12.0000,8.0000) circle (1.2pt);
\fill[red!80!black] (0.0000,8.0000) circle (1.2pt);
\fill[red!80!black] (12.0000,8.0000) circle (1.2pt);
\fill[red!80!black] (0.0000,8.0000) circle (1.2pt);
\fill[red!80!black] (12.0000,8.0000) circle (1.2pt);
\fill[red!80!black] (0.0000,8.0000) circle (1.2pt);
\fill[red!80!black] (12.0000,8.0000) circle (1.2pt);
\fill[red!80!black] (0.0000,8.0000) circle (1.2pt);
\fill[red!80!black] (12.0000,8.0000) circle (1.2pt);
\fill[red!80!black] (0.0000,8.0000) circle (1.2pt);
\fill[red!80!black] (12.0000,8.0000) circle (1.2pt);
\fill[red!80!black] (0.0000,8.0000) circle (1.2pt);
\fill[red!80!black] (12.0000,8.0000) circle (1.2pt);
\fill[red!80!black] (0.0000,8.0000) circle (1.2pt);
\fill[red!80!black] (12.0000,8.0000) circle (1.2pt);
\fill[red!80!black] (0.0000,8.0000) circle (1.2pt);
\fill[red!80!black] (12.0000,8.0000) circle (1.2pt);
\fill[red!80!black] (0.0000,8.0000) circle (1.2pt);
\fill[red!80!black] (12.0000,8.0000) circle (1.2pt);
\fill[red!80!black] (0.0000,8.0000) circle (1.2pt);
\fill[red!80!black] (12.0000,8.0000) circle (1.2pt);
\fill[red!80!black] (0.0000,8.0000) circle (1.2pt);
\fill[red!80!black] (12.0000,8.0000) circle (1.2pt);
\fill[red!80!black] (0.0000,8.0000) circle (1.2pt);
\fill[red!80!black] (12.0000,8.0000) circle (1.2pt);
\fill[red!80!black] (0.0000,8.0000) circle (1.2pt);
\fill[red!80!black] (12.0000,8.0000) circle (1.2pt);
\fill[red!80!black] (0.0000,8.0000) circle (1.2pt);
\fill[red!80!black] (12.0000,8.0000) circle (1.2pt);
\fill[red!80!black] (0.0000,8.0000) circle (1.2pt);
\fill[red!80!black] (12.0000,8.0000) circle (1.2pt);
\fill[red!80!black] (0.0000,8.0000) circle (1.2pt);
\fill[red!80!black] (12.0000,8.0000) circle (1.2pt);
\fill[red!80!black] (0.0000,8.0000) circle (1.2pt);
\fill[red!80!black] (12.0000,8.0000) circle (1.2pt);
\fill[red!80!black] (0.0000,8.0000) circle (1.2pt);
\fill[red!80!black] (12.0000,8.0000) circle (1.2pt);
\fill[red!80!black] (0.0000,8.0000) circle (1.2pt);
\fill[red!80!black] (12.0000,8.0000) circle (1.2pt);
\fill[red!80!black] (0.0000,8.0000) circle (1.2pt);
\fill[red!80!black] (12.0000,8.0000) circle (1.2pt);
\fill[red!80!black] (0.0000,8.0000) circle (1.2pt);
\fill[red!80!black] (12.0000,8.0000) circle (1.2pt);
\fill[red!80!black] (0.0000,8.0000) circle (1.2pt);
\fill[red!80!black] (12.0000,8.0000) circle (1.2pt);
\fill[red!80!black] (0.0000,8.0000) circle (1.2pt);
\fill[red!80!black] (12.0000,8.0000) circle (1.2pt);
\fill[red!80!black] (0.0000,8.0000) circle (1.2pt);
\fill[red!80!black] (12.0000,8.0000) circle (1.2pt);
\fill[red!80!black] (0.0000,8.0000) circle (1.2pt);
\fill[red!80!black] (12.0000,8.0000) circle (1.2pt);
\fill[red!80!black] (0.0000,8.0000) circle (1.2pt);
\fill[red!80!black] (12.0000,8.0000) circle (1.2pt);
\fill[red!80!black] (0.0000,8.0000) circle (1.2pt);
\fill[red!80!black] (12.0000,8.0000) circle (1.2pt);
\fill[red!80!black] (0.0000,8.0000) circle (1.2pt);
\fill[red!80!black] (12.0000,8.0000) circle (1.2pt);
\fill[red!80!black] (0.0000,8.0000) circle (1.2pt);
\fill[red!80!black] (12.0000,8.0000) circle (1.2pt);
\fill[red!80!black] (0.0000,8.0000) circle (1.2pt);
\fill[red!80!black] (12.0000,8.0000) circle (1.2pt);
\fill[red!80!black] (0.0000,8.0000) circle (1.2pt);
\fill[red!80!black] (12.0000,8.0000) circle (1.2pt);
\fill[red!80!black] (0.0000,8.0000) circle (1.2pt);
\fill[red!80!black] (12.0000,8.0000) circle (1.2pt);
\fill[red!80!black] (0.0000,8.0000) circle (1.2pt);
\fill[red!80!black] (12.0000,8.0000) circle (1.2pt);
\fill[red!80!black] (0.0000,8.0000) circle (1.2pt);
\fill[red!80!black] (12.0000,8.0000) circle (1.2pt);
\fill[red!80!black] (0.0000,8.0000) circle (1.2pt);
\fill[red!80!black] (12.0000,8.0000) circle (1.2pt);
\fill[red!80!black] (0.0000,8.0000) circle (1.2pt);
\fill[red!80!black] (12.0000,8.0000) circle (1.2pt);
\fill[red!80!black] (0.0000,8.0000) circle (1.2pt);
\fill[red!80!black] (12.0000,8.0000) circle (1.2pt);
\fill[red!80!black] (0.0000,8.0000) circle (1.2pt);
\fill[red!80!black] (12.0000,8.0000) circle (1.2pt);
\fill[red!80!black] (0.0000,8.0000) circle (1.2pt);
\fill[red!80!black] (12.0000,8.0000) circle (1.2pt);
\fill[red!80!black] (0.0000,8.0000) circle (1.2pt);
\fill[red!80!black] (12.0000,8.0000) circle (1.2pt);
\fill[red!80!black] (0.0000,8.0000) circle (1.2pt);
\fill[red!80!black] (12.0000,8.0000) circle (1.2pt);
\fill[red!80!black] (0.0000,8.0000) circle (1.2pt);
\fill[red!80!black] (12.0000,8.0000) circle (1.2pt);
\fill[red!80!black] (0.0000,8.0000) circle (1.2pt);
\fill[red!80!black] (12.0000,8.0000) circle (1.2pt);
\fill[red!80!black] (0.0000,8.0000) circle (1.2pt);
\fill[red!80!black] (12.0000,8.0000) circle (1.2pt);
\fill[red!80!black] (0.0000,8.0000) circle (1.2pt);
\fill[red!80!black] (12.0000,8.0000) circle (1.2pt);
\fill[red!80!black] (0.0000,8.0000) circle (1.2pt);
\fill[red!80!black] (12.0000,8.0000) circle (1.2pt);
\fill[red!80!black] (0.0000,8.0000) circle (1.2pt);
\fill[red!80!black] (12.0000,8.0000) circle (1.2pt);
\fill[red!80!black] (0.0000,8.0000) circle (1.2pt);
\fill[red!80!black] (12.0000,8.0000) circle (1.2pt);
\fill[red!80!black] (0.0000,8.0000) circle (1.2pt);
\fill[red!80!black] (12.0000,8.0000) circle (1.2pt);
\fill[red!80!black] (0.0000,8.0000) circle (1.2pt);
\fill[red!80!black] (12.0000,8.0000) circle (1.2pt);
\fill[red!80!black] (0.0000,8.0000) circle (1.2pt);
\fill[red!80!black] (12.0000,8.0000) circle (1.2pt);
\fill[red!80!black] (0.0000,8.0000) circle (1.2pt);
\fill[red!80!black] (12.0000,8.0000) circle (1.2pt);
\fill[red!80!black] (0.0000,8.0000) circle (1.2pt);
\fill[red!80!black] (12.0000,8.0000) circle (1.2pt);
\fill[red!80!black] (0.0000,8.0000) circle (1.2pt);
\fill[red!80!black] (12.0000,8.0000) circle (1.2pt);
\fill[red!80!black] (0.0000,8.0000) circle (1.2pt);
\fill[red!80!black] (12.0000,8.0000) circle (1.2pt);
\fill[red!80!black] (0.0000,8.0000) circle (1.2pt);
\fill[red!80!black] (12.0000,8.0000) circle (1.2pt);
\fill[red!80!black] (0.0000,8.0000) circle (1.2pt);
\fill[red!80!black] (12.0000,8.0000) circle (1.2pt);
\fill[red!80!black] (0.0000,8.0000) circle (1.2pt);
\fill[red!80!black] (12.0000,8.0000) circle (1.2pt);
\fill[red!80!black] (0.0000,8.0000) circle (1.2pt);
\fill[red!80!black] (12.0000,8.0000) circle (1.2pt);
\fill[red!80!black] (0.0000,8.0000) circle (1.2pt);
\fill[red!80!black] (12.0000,8.0000) circle (1.2pt);
\fill[red!80!black] (0.0000,8.0000) circle (1.2pt);
\fill[red!80!black] (12.0000,8.0000) circle (1.2pt);
\fill[red!80!black] (0.0000,8.0000) circle (1.2pt);
\fill[red!80!black] (12.0000,8.0000) circle (1.2pt);
\fill[red!80!black] (0.0000,8.0000) circle (1.2pt);
\fill[red!80!black] (12.0000,8.0000) circle (1.2pt);
\fill[red!80!black] (0.0000,8.0000) circle (1.2pt);
\fill[red!80!black] (12.0000,8.0000) circle (1.2pt);
\fill[red!80!black] (0.0000,8.0000) circle (1.2pt);
\fill[red!80!black] (12.0000,8.0000) circle (1.2pt);
\fill[red!80!black] (0.0000,8.0000) circle (1.2pt);
\fill[red!80!black] (12.0000,8.0000) circle (1.2pt);
\fill[red!80!black] (0.0000,8.0000) circle (1.2pt);
\fill[red!80!black] (12.0000,8.0000) circle (1.2pt);
\fill[red!80!black] (0.0000,8.0000) circle (1.2pt);
\fill[red!80!black] (12.0000,8.0000) circle (1.2pt);
\fill[red!80!black] (0.0000,8.0000) circle (1.2pt);
\fill[red!80!black] (12.0000,8.0000) circle (1.2pt);
\fill[red!80!black] (0.0000,8.0000) circle (1.2pt);
\fill[red!80!black] (12.0000,8.0000) circle (1.2pt);
\fill[red!80!black] (0.0000,8.0000) circle (1.2pt);
\fill[red!80!black] (12.0000,8.0000) circle (1.2pt);
\fill[red!80!black] (0.0000,8.0000) circle (1.2pt);
\fill[red!80!black] (12.0000,8.0000) circle (1.2pt);
\fill[red!80!black] (0.0000,8.0000) circle (1.2pt);
\fill[red!80!black] (12.0000,8.0000) circle (1.2pt);
\fill[red!80!black] (0.0000,8.0000) circle (1.2pt);
\fill[red!80!black] (12.0000,8.0000) circle (1.2pt);
\fill[red!80!black] (0.0000,8.0000) circle (1.2pt);
\fill[red!80!black] (12.0000,8.0000) circle (1.2pt);
\fill[red!80!black] (0.0000,8.0000) circle (1.2pt);
\fill[red!80!black] (12.0000,8.0000) circle (1.2pt);
\fill[red!80!black] (0.0000,8.0000) circle (1.2pt);
\fill[red!80!black] (12.0000,8.0000) circle (1.2pt);
\fill[red!80!black] (0.0000,8.0000) circle (1.2pt);
\fill[red!80!black] (12.0000,8.0000) circle (1.2pt);
\fill[red!80!black] (0.0000,8.0000) circle (1.2pt);
\fill[red!80!black] (12.0000,8.0000) circle (1.2pt);
\fill[red!80!black] (0.0000,8.0000) circle (1.2pt);
\fill[red!80!black] (12.0000,8.0000) circle (1.2pt);
\fill[red!80!black] (0.0000,8.0000) circle (1.2pt);
\fill[red!80!black] (12.0000,8.0000) circle (1.2pt);
\fill[red!80!black] (0.0000,8.0000) circle (1.2pt);
\fill[red!80!black] (12.0000,8.0000) circle (1.2pt);
\fill[red!80!black] (0.0000,8.0000) circle (1.2pt);
\fill[red!80!black] (12.0000,8.0000) circle (1.2pt);
\fill[red!80!black] (0.0000,8.0000) circle (1.2pt);
\fill[red!80!black] (12.0000,8.0000) circle (1.2pt);
\fill[red!80!black] (0.0000,8.0000) circle (1.2pt);
\fill[red!80!black] (12.0000,8.0000) circle (1.2pt);
\fill[red!80!black] (0.0000,8.0000) circle (1.2pt);
\fill[red!80!black] (12.0000,8.0000) circle (1.2pt);
\fill[red!80!black] (0.0000,8.0000) circle (1.2pt);
\fill[red!80!black] (12.0000,8.0000) circle (1.2pt);
\fill[red!80!black] (0.0000,8.0000) circle (1.2pt);
\fill[red!80!black] (12.0000,8.0000) circle (1.2pt);
\fill[red!80!black] (0.0000,8.0000) circle (1.2pt);
\fill[red!80!black] (12.0000,8.0000) circle (1.2pt);
\fill[red!80!black] (0.0000,8.0000) circle (1.2pt);
\fill[red!80!black] (12.0000,8.0000) circle (1.2pt);
\fill[red!80!black] (0.0000,8.0000) circle (1.2pt);
\fill[red!80!black] (12.0000,8.0000) circle (1.2pt);
\fill[red!80!black] (0.0000,8.0000) circle (1.2pt);
\fill[red!80!black] (12.0000,8.0000) circle (1.2pt);
\fill[red!80!black] (0.0000,8.0000) circle (1.2pt);
\fill[red!80!black] (12.0000,8.0000) circle (1.2pt);
\fill[red!80!black] (0.0000,8.0000) circle (1.2pt);
\fill[red!80!black] (12.0000,8.0000) circle (1.2pt);
\fill[red!80!black] (0.0000,8.0000) circle (1.2pt);
\fill[red!80!black] (12.0000,8.0000) circle (1.2pt);
\fill[red!80!black] (0.0000,8.0000) circle (1.2pt);
\fill[red!80!black] (12.0000,8.0000) circle (1.2pt);
\fill[red!80!black] (0.0000,8.0000) circle (1.2pt);
\fill[red!80!black] (12.0000,8.0000) circle (1.2pt);
\fill[red!80!black] (0.0000,8.0000) circle (1.2pt);
\fill[red!80!black] (12.0000,8.0000) circle (1.2pt);
\fill[red!80!black] (0.0000,8.0000) circle (1.2pt);
\fill[red!80!black] (12.0000,8.0000) circle (1.2pt);
\fill[red!80!black] (0.0000,8.0000) circle (1.2pt);
\fill[red!80!black] (12.0000,8.0000) circle (1.2pt);
\fill[red!80!black] (0.0000,8.0000) circle (1.2pt);
\fill[red!80!black] (12.0000,8.0000) circle (1.2pt);
\fill[red!80!black] (0.0000,8.0000) circle (1.2pt);
\fill[red!80!black] (12.0000,8.0000) circle (1.2pt);
\fill[red!80!black] (0.0000,8.0000) circle (1.2pt);
\fill[red!80!black] (12.0000,8.0000) circle (1.2pt);
\fill[red!80!black] (0.0000,8.0000) circle (1.2pt);
\fill[red!80!black] (12.0000,8.0000) circle (1.2pt);
\fill[red!80!black] (0.0000,8.0000) circle (1.2pt);
\fill[red!80!black] (12.0000,8.0000) circle (1.2pt);
\fill[red!80!black] (0.0000,8.0000) circle (1.2pt);
\fill[red!80!black] (12.0000,8.0000) circle (1.2pt);
\fill[red!80!black] (0.0000,8.0000) circle (1.2pt);
\fill[red!80!black] (12.0000,8.0000) circle (1.2pt);
\fill[red!80!black] (0.0000,8.0000) circle (1.2pt);
\fill[red!80!black] (12.0000,8.0000) circle (1.2pt);
\fill[red!80!black] (0.0000,8.0000) circle (1.2pt);
\fill[red!80!black] (12.0000,8.0000) circle (1.2pt);
\fill[red!80!black] (0.0000,8.0000) circle (1.2pt);
\fill[red!80!black] (12.0000,8.0000) circle (1.2pt);
\fill[red!80!black] (0.0000,8.0000) circle (1.2pt);
\fill[red!80!black] (12.0000,8.0000) circle (1.2pt);
\fill[red!80!black] (0.0000,8.0000) circle (1.2pt);
\fill[red!80!black] (12.0000,8.0000) circle (1.2pt);
\fill[red!80!black] (0.0000,8.0000) circle (1.2pt);
\fill[red!80!black] (12.0000,8.0000) circle (1.2pt);
\fill[red!80!black] (0.0000,8.0000) circle (1.2pt);
\fill[red!80!black] (12.0000,8.0000) circle (1.2pt);
\fill[red!80!black] (0.0000,8.0000) circle (1.2pt);
\fill[red!80!black] (12.0000,8.0000) circle (1.2pt);
\fill[red!80!black] (0.0000,8.0000) circle (1.2pt);
\fill[red!80!black] (12.0000,8.0000) circle (1.2pt);
\fill[red!80!black] (0.0000,8.0000) circle (1.2pt);
\fill[red!80!black] (12.0000,8.0000) circle (1.2pt);
\fill[red!80!black] (0.0000,8.0000) circle (1.2pt);
\fill[red!80!black] (12.0000,8.0000) circle (1.2pt);
\fill[red!80!black] (0.0000,8.0000) circle (1.2pt);
\fill[red!80!black] (12.0000,8.0000) circle (1.2pt);
\fill[red!80!black] (0.0000,8.0000) circle (1.2pt);
\fill[red!80!black] (12.0000,8.0000) circle (1.2pt);
\fill[red!80!black] (0.0000,8.0000) circle (1.2pt);
\fill[red!80!black] (12.0000,8.0000) circle (1.2pt);
\fill[red!80!black] (0.0000,8.0000) circle (1.2pt);
\fill[red!80!black] (12.0000,8.0000) circle (1.2pt);
\fill[red!80!black] (0.0000,8.0000) circle (1.2pt);
\fill[red!80!black] (12.0000,8.0000) circle (1.2pt);
\fill[red!80!black] (0.0000,8.0000) circle (1.2pt);
\fill[red!80!black] (12.0000,8.0000) circle (1.2pt);
\fill[red!80!black] (0.0000,8.0000) circle (1.2pt);
\fill[red!80!black] (12.0000,8.0000) circle (1.2pt);
\fill[red!80!black] (0.0000,8.0000) circle (1.2pt);
\fill[red!80!black] (12.0000,8.0000) circle (1.2pt);
\fill[red!80!black] (0.0000,8.0000) circle (1.2pt);
\fill[red!80!black] (12.0000,8.0000) circle (1.2pt);
\fill[red!80!black] (0.0000,8.0000) circle (1.2pt);
\fill[red!80!black] (12.0000,8.0000) circle (1.2pt);
\fill[red!80!black] (0.0000,8.0000) circle (1.2pt);
\fill[red!80!black] (12.0000,8.0000) circle (1.2pt);
\fill[red!80!black] (0.0000,8.0000) circle (1.2pt);
\fill[red!80!black] (12.0000,8.0000) circle (1.2pt);
\fill[red!80!black] (0.0000,8.0000) circle (1.2pt);
\fill[red!80!black] (12.0000,8.0000) circle (1.2pt);
\fill[red!80!black] (0.0000,8.0000) circle (1.2pt);
\fill[red!80!black] (12.0000,8.0000) circle (1.2pt);
\fill[red!80!black] (0.0000,8.0000) circle (1.2pt);
\fill[red!80!black] (12.0000,8.0000) circle (1.2pt);
\fill[red!80!black] (0.0000,8.0000) circle (1.2pt);
\fill[red!80!black] (12.0000,8.0000) circle (1.2pt);
\fill[red!80!black] (0.0000,8.0000) circle (1.2pt);
\fill[red!80!black] (12.0000,8.0000) circle (1.2pt);
\fill[red!80!black] (0.0000,8.0000) circle (1.2pt);
\fill[red!80!black] (12.0000,8.0000) circle (1.2pt);
\fill[red!80!black] (0.0000,8.0000) circle (1.2pt);
\fill[red!80!black] (12.0000,8.0000) circle (1.2pt);
\fill[red!80!black] (0.0000,8.0000) circle (1.2pt);
\fill[red!80!black] (12.0000,8.0000) circle (1.2pt);
\fill[red!80!black] (0.0000,8.0000) circle (1.2pt);
\fill[red!80!black] (12.0000,8.0000) circle (1.2pt);
\fill[red!80!black] (0.0000,8.0000) circle (1.2pt);
\fill[red!80!black] (12.0000,8.0000) circle (1.2pt);
\fill[red!80!black] (0.0000,8.0000) circle (1.2pt);
\fill[red!80!black] (12.0000,8.0000) circle (1.2pt);
\fill[red!80!black] (0.0000,8.0000) circle (1.2pt);
\fill[red!80!black] (12.0000,8.0000) circle (1.2pt);
\fill[red!80!black] (0.0000,8.0000) circle (1.2pt);
\fill[red!80!black] (12.0000,8.0000) circle (1.2pt);
\fill[red!80!black] (0.0000,8.0000) circle (1.2pt);
\fill[red!80!black] (12.0000,8.0000) circle (1.2pt);
\fill[red!80!black] (0.0000,8.0000) circle (1.2pt);
\fill[red!80!black] (12.0000,8.0000) circle (1.2pt);
\fill[red!80!black] (0.0000,8.0000) circle (1.2pt);
\fill[red!80!black] (12.0000,8.0000) circle (1.2pt);
\fill[red!80!black] (0.0000,8.0000) circle (1.2pt);
\fill[red!80!black] (12.0000,8.0000) circle (1.2pt);
\fill[red!80!black] (0.0000,8.0000) circle (1.2pt);
\fill[red!80!black] (12.0000,8.0000) circle (1.2pt);
\fill[red!80!black] (0.0000,8.0000) circle (1.2pt);
\fill[red!80!black] (12.0000,8.0000) circle (1.2pt);
\fill[red!80!black] (0.0000,8.0000) circle (1.2pt);
\fill[red!80!black] (12.0000,8.0000) circle (1.2pt);
\fill[red!80!black] (0.0000,8.0000) circle (1.2pt);
\fill[red!80!black] (12.0000,8.0000) circle (1.2pt);
\fill[red!80!black] (0.0000,8.0000) circle (1.2pt);
\fill[red!80!black] (12.0000,8.0000) circle (1.2pt);
\fill[red!80!black] (0.0000,8.0000) circle (1.2pt);
\fill[red!80!black] (12.0000,8.0000) circle (1.2pt);
\fill[red!80!black] (0.0000,8.0000) circle (1.2pt);
\fill[red!80!black] (12.0000,8.0000) circle (1.2pt);
\fill[red!80!black] (0.0000,8.0000) circle (1.2pt);
\fill[red!80!black] (12.0000,8.0000) circle (1.2pt);
\fill[red!80!black] (0.0000,8.0000) circle (1.2pt);
\fill[red!80!black] (12.0000,8.0000) circle (1.2pt);
\fill[red!80!black] (0.0000,8.0000) circle (1.2pt);
\fill[red!80!black] (12.0000,8.0000) circle (1.2pt);
\fill[red!80!black] (0.0000,8.0000) circle (1.2pt);
\fill[red!80!black] (12.0000,8.0000) circle (1.2pt);
\fill[red!80!black] (0.0000,8.0000) circle (1.2pt);
\fill[red!80!black] (12.0000,8.0000) circle (1.2pt);
\fill[red!80!black] (0.0000,8.0000) circle (1.2pt);
\fill[red!80!black] (12.0000,8.0000) circle (1.2pt);
\fill[red!80!black] (0.0000,8.0000) circle (1.2pt);
\fill[red!80!black] (12.0000,8.0000) circle (1.2pt);
\fill[red!80!black] (0.0000,8.0000) circle (1.2pt);
\fill[red!80!black] (12.0000,8.0000) circle (1.2pt);
\fill[red!80!black] (0.0000,8.0000) circle (1.2pt);
\fill[red!80!black] (12.0000,8.0000) circle (1.2pt);
\fill[red!80!black] (0.0000,8.0000) circle (1.2pt);
\fill[red!80!black] (12.0000,8.0000) circle (1.2pt);
\fill[red!80!black] (0.0000,8.0000) circle (1.2pt);
\fill[red!80!black] (12.0000,8.0000) circle (1.2pt);
\fill[red!80!black] (0.0000,8.0000) circle (1.2pt);
\fill[red!80!black] (12.0000,8.0000) circle (1.2pt);
\fill[red!80!black] (0.0000,8.0000) circle (1.2pt);
\fill[red!80!black] (12.0000,8.0000) circle (1.2pt);
\fill[red!80!black] (0.0000,8.0000) circle (1.2pt);
\fill[red!80!black] (12.0000,8.0000) circle (1.2pt);
\fill[red!80!black] (0.0000,8.0000) circle (1.2pt);
\fill[red!80!black] (12.0000,8.0000) circle (1.2pt);
\fill[red!80!black] (0.0000,8.0000) circle (1.2pt);
\fill[red!80!black] (12.0000,8.0000) circle (1.2pt);
\fill[red!80!black] (0.0000,8.0000) circle (1.2pt);
\fill[red!80!black] (12.0000,8.0000) circle (1.2pt);
\fill[red!80!black] (0.0000,8.0000) circle (1.2pt);
\fill[red!80!black] (12.0000,8.0000) circle (1.2pt);
\fill[red!80!black] (0.0000,8.0000) circle (1.2pt);
\fill[red!80!black] (12.0000,8.0000) circle (1.2pt);
\fill[red!80!black] (0.0000,8.0000) circle (1.2pt);
\fill[red!80!black] (12.0000,8.0000) circle (1.2pt);
\fill[red!80!black] (0.0000,8.0000) circle (1.2pt);
\fill[red!80!black] (12.0000,8.0000) circle (1.2pt);
\fill[red!80!black] (0.0000,8.0000) circle (1.2pt);
\fill[red!80!black] (12.0000,8.0000) circle (1.2pt);
\fill[red!80!black] (0.0000,8.0000) circle (1.2pt);
\fill[red!80!black] (12.0000,8.0000) circle (1.2pt);
\fill[red!80!black] (0.0000,8.0000) circle (1.2pt);
\fill[red!80!black] (12.0000,8.0000) circle (1.2pt);
\fill[red!80!black] (0.0000,8.0000) circle (1.2pt);
\fill[red!80!black] (12.0000,8.0000) circle (1.2pt);
\fill[red!80!black] (0.0000,8.0000) circle (1.2pt);
\fill[red!80!black] (12.0000,8.0000) circle (1.2pt);
\fill[red!80!black] (0.0000,8.0000) circle (1.2pt);
\fill[red!80!black] (12.0000,8.0000) circle (1.2pt);
\fill[red!80!black] (0.0000,8.0000) circle (1.2pt);
\fill[red!80!black] (12.0000,8.0000) circle (1.2pt);
\fill[red!80!black] (0.0000,8.0000) circle (1.2pt);
\fill[red!80!black] (12.0000,8.0000) circle (1.2pt);
\fill[red!80!black] (0.0000,8.0000) circle (1.2pt);
\fill[red!80!black] (12.0000,8.0000) circle (1.2pt);
\fill[red!80!black] (0.0000,8.0000) circle (1.2pt);
\fill[red!80!black] (12.0000,8.0000) circle (1.2pt);
\fill[red!80!black] (0.0000,8.0000) circle (1.2pt);
\fill[red!80!black] (12.0000,8.0000) circle (1.2pt);
\fill[red!80!black] (0.0000,8.0000) circle (1.2pt);
\fill[red!80!black] (12.0000,8.0000) circle (1.2pt);
\fill[red!80!black] (0.0000,8.0000) circle (1.2pt);
\fill[red!80!black] (12.0000,8.0000) circle (1.2pt);
\fill[red!80!black] (0.0000,8.0000) circle (1.2pt);
\fill[red!80!black] (12.0000,8.0000) circle (1.2pt);
\fill[red!80!black] (0.0000,8.0000) circle (1.2pt);
\fill[red!80!black] (12.0000,8.0000) circle (1.2pt);
\fill[red!80!black] (0.0000,8.0000) circle (1.2pt);
\fill[red!80!black] (12.0000,8.0000) circle (1.2pt);
\fill[red!80!black] (0.0000,8.0000) circle (1.2pt);
\fill[red!80!black] (12.0000,8.0000) circle (1.2pt);
\fill[red!80!black] (0.0000,8.0000) circle (1.2pt);
\fill[red!80!black] (12.0000,8.0000) circle (1.2pt);
\fill[red!80!black] (0.0000,8.0000) circle (1.2pt);
\fill[red!80!black] (12.0000,8.0000) circle (1.2pt);
\fill[red!80!black] (0.0000,8.0000) circle (1.2pt);
\fill[red!80!black] (12.0000,8.0000) circle (1.2pt);
\fill[red!80!black] (0.0000,8.0000) circle (1.2pt);
\fill[red!80!black] (12.0000,8.0000) circle (1.2pt);
\fill[red!80!black] (0.0000,8.0000) circle (1.2pt);
\fill[red!80!black] (12.0000,8.0000) circle (1.2pt);
\fill[red!80!black] (0.0000,8.0000) circle (1.2pt);
\fill[red!80!black] (12.0000,8.0000) circle (1.2pt);
\fill[red!80!black] (0.0000,8.0000) circle (1.2pt);
\fill[red!80!black] (12.0000,8.0000) circle (1.2pt);
\fill[red!80!black] (0.0000,8.0000) circle (1.2pt);
\fill[red!80!black] (12.0000,8.0000) circle (1.2pt);
\fill[red!80!black] (0.0000,8.0000) circle (1.2pt);
\fill[red!80!black] (12.0000,8.0000) circle (1.2pt);
\fill[red!80!black] (0.0000,8.0000) circle (1.2pt);
\fill[red!80!black] (12.0000,8.0000) circle (1.2pt);
\fill[red!80!black] (0.0000,8.0000) circle (1.2pt);
\fill[red!80!black] (12.0000,8.0000) circle (1.2pt);
\fill[red!80!black] (0.0000,8.0000) circle (1.2pt);
\fill[red!80!black] (12.0000,8.0000) circle (1.2pt);
\fill[red!80!black] (0.0000,8.0000) circle (1.2pt);
\fill[red!80!black] (12.0000,8.0000) circle (1.2pt);
\fill[red!80!black] (0.0000,8.0000) circle (1.2pt);
\fill[red!80!black] (12.0000,8.0000) circle (1.2pt);
\fill[red!80!black] (0.0000,8.0000) circle (1.2pt);
\fill[red!80!black] (12.0000,8.0000) circle (1.2pt);
\fill[red!80!black] (0.0000,8.0000) circle (1.2pt);
\fill[red!80!black] (12.0000,8.0000) circle (1.2pt);
\fill[red!80!black] (0.0000,8.0000) circle (1.2pt);
\fill[red!80!black] (12.0000,8.0000) circle (1.2pt);
\fill[red!80!black] (0.0000,8.0000) circle (1.2pt);
\fill[red!80!black] (12.0000,8.0000) circle (1.2pt);
\fill[red!80!black] (0.0000,8.0000) circle (1.2pt);
\fill[red!80!black] (12.0000,8.0000) circle (1.2pt);
\fill[red!80!black] (0.0000,8.0000) circle (1.2pt);
\fill[red!80!black] (12.0000,8.0000) circle (1.2pt);
\fill[red!80!black] (0.0000,8.0000) circle (1.2pt);
\fill[red!80!black] (12.0000,8.0000) circle (1.2pt);
\fill[red!80!black] (0.0000,8.0000) circle (1.2pt);
\fill[red!80!black] (12.0000,8.0000) circle (1.2pt);
\fill[red!80!black] (0.0000,8.0000) circle (1.2pt);
\fill[red!80!black] (12.0000,8.0000) circle (1.2pt);
\fill[red!80!black] (0.0000,8.0000) circle (1.2pt);
\fill[red!80!black] (12.0000,8.0000) circle (1.2pt);
\fill[red!80!black] (0.0000,8.0000) circle (1.2pt);
\fill[red!80!black] (12.0000,8.0000) circle (1.2pt);
\fill[red!80!black] (0.0000,8.0000) circle (1.2pt);
\fill[red!80!black] (12.0000,8.0000) circle (1.2pt);
\fill[red!80!black] (0.0000,8.0000) circle (1.2pt);
\fill[red!80!black] (12.0000,8.0000) circle (1.2pt);
\fill[red!80!black] (0.0000,8.0000) circle (1.2pt);
\fill[red!80!black] (12.0000,8.0000) circle (1.2pt);
\fill[red!80!black] (0.0000,8.0000) circle (1.2pt);
\fill[red!80!black] (12.0000,8.0000) circle (1.2pt);
\fill[red!80!black] (0.0000,8.0000) circle (1.2pt);
\fill[red!80!black] (12.0000,8.0000) circle (1.2pt);
\fill[red!80!black] (0.0000,8.0000) circle (1.2pt);
\fill[red!80!black] (12.0000,8.0000) circle (1.2pt);
\fill[red!80!black] (0.0000,8.0000) circle (1.2pt);
\fill[red!80!black] (12.0000,8.0000) circle (1.2pt);
\fill[red!80!black] (0.0000,8.0000) circle (1.2pt);
\fill[red!80!black] (12.0000,8.0000) circle (1.2pt);
\fill[red!80!black] (0.0000,8.0000) circle (1.2pt);
\fill[red!80!black] (12.0000,8.0000) circle (1.2pt);
\fill[red!80!black] (0.0000,8.0000) circle (1.2pt);
\fill[red!80!black] (12.0000,8.0000) circle (1.2pt);
\fill[red!80!black] (0.0000,8.0000) circle (1.2pt);
\fill[red!80!black] (12.0000,8.0000) circle (1.2pt);
\fill[red!80!black] (0.0000,8.0000) circle (1.2pt);
\fill[red!80!black] (12.0000,8.0000) circle (1.2pt);
\fill[red!80!black] (0.0000,8.0000) circle (1.2pt);
\fill[red!80!black] (12.0000,8.0000) circle (1.2pt);
\fill[red!80!black] (0.0000,8.0000) circle (1.2pt);
\fill[red!80!black] (12.0000,8.0000) circle (1.2pt);
\fill[red!80!black] (0.0000,8.0000) circle (1.2pt);
\fill[red!80!black] (12.0000,8.0000) circle (1.2pt);
\fill[red!80!black] (0.0000,8.0000) circle (1.2pt);
\fill[red!80!black] (12.0000,8.0000) circle (1.2pt);
\fill[red!80!black] (0.0000,8.0000) circle (1.2pt);
\fill[red!80!black] (12.0000,8.0000) circle (1.2pt);
\fill[red!80!black] (0.0000,8.0000) circle (1.2pt);
\fill[red!80!black] (12.0000,8.0000) circle (1.2pt);
\fill[red!80!black] (0.0000,8.0000) circle (1.2pt);
\fill[red!80!black] (12.0000,8.0000) circle (1.2pt);
\fill[red!80!black] (0.0000,8.0000) circle (1.2pt);
\fill[red!80!black] (12.0000,8.0000) circle (1.2pt);
\fill[red!80!black] (0.0000,8.0000) circle (1.2pt);
\fill[red!80!black] (12.0000,8.0000) circle (1.2pt);
\fill[red!80!black] (0.0000,8.0000) circle (1.2pt);
\fill[red!80!black] (12.0000,8.0000) circle (1.2pt);
\fill[red!80!black] (0.0000,8.0000) circle (1.2pt);
\fill[red!80!black] (12.0000,8.0000) circle (1.2pt);
\fill[red!80!black] (0.0000,8.0000) circle (1.2pt);
\fill[red!80!black] (12.0000,8.0000) circle (1.2pt);
\fill[red!80!black] (0.0000,8.0000) circle (1.2pt);
\fill[red!80!black] (12.0000,8.0000) circle (1.2pt);
\fill[red!80!black] (0.0000,8.0000) circle (1.2pt);
\fill[red!80!black] (12.0000,8.0000) circle (1.2pt);
\fill[red!80!black] (0.0000,8.0000) circle (1.2pt);
\fill[red!80!black] (12.0000,8.0000) circle (1.2pt);
\fill[red!80!black] (0.0000,8.0000) circle (1.2pt);
\fill[red!80!black] (12.0000,8.0000) circle (1.2pt);
\fill[red!80!black] (0.0000,8.0000) circle (1.2pt);
\fill[red!80!black] (12.0000,8.0000) circle (1.2pt);
\fill[red!80!black] (0.0000,8.0000) circle (1.2pt);
\fill[red!80!black] (12.0000,8.0000) circle (1.2pt);
\fill[red!80!black] (0.0000,8.0000) circle (1.2pt);
\fill[red!80!black] (12.0000,8.0000) circle (1.2pt);
\fill[red!80!black] (0.0000,8.0000) circle (1.2pt);
\fill[red!80!black] (12.0000,8.0000) circle (1.2pt);
\fill[red!80!black] (0.0000,8.0000) circle (1.2pt);
\fill[red!80!black] (12.0000,8.0000) circle (1.2pt);
\fill[red!80!black] (0.0000,8.0000) circle (1.2pt);
\fill[red!80!black] (12.0000,8.0000) circle (1.2pt);
\fill[red!80!black] (0.0000,8.0000) circle (1.2pt);
\fill[red!80!black] (12.0000,8.0000) circle (1.2pt);
\fill[red!80!black] (0.0000,8.0000) circle (1.2pt);
\fill[red!80!black] (12.0000,8.0000) circle (1.2pt);
\fill[red!80!black] (0.0000,8.0000) circle (1.2pt);
\fill[red!80!black] (12.0000,8.0000) circle (1.2pt);
\fill[red!80!black] (0.0000,8.0000) circle (1.2pt);
\fill[red!80!black] (12.0000,8.0000) circle (1.2pt);
\fill[red!80!black] (0.0000,8.0000) circle (1.2pt);
\fill[red!80!black] (12.0000,8.0000) circle (1.2pt);
\fill[red!80!black] (0.0000,8.0000) circle (1.2pt);
\fill[red!80!black] (12.0000,8.0000) circle (1.2pt);
\fill[red!80!black] (0.0000,8.0000) circle (1.2pt);
\fill[red!80!black] (12.0000,8.0000) circle (1.2pt);
\fill[red!80!black] (0.0000,8.0000) circle (1.2pt);
\fill[red!80!black] (12.0000,8.0000) circle (1.2pt);
\fill[red!80!black] (0.0000,8.0000) circle (1.2pt);
\fill[red!80!black] (12.0000,8.0000) circle (1.2pt);
\fill[red!80!black] (0.0000,8.0000) circle (1.2pt);
\fill[red!80!black] (12.0000,8.0000) circle (1.2pt);
\fill[red!80!black] (0.0000,8.0000) circle (1.2pt);
\fill[red!80!black] (12.0000,8.0000) circle (1.2pt);
\fill[red!80!black] (0.0000,8.0000) circle (1.2pt);
\fill[red!80!black] (12.0000,8.0000) circle (1.2pt);
\fill[red!80!black] (0.0000,8.0000) circle (1.2pt);
\fill[red!80!black] (12.0000,8.0000) circle (1.2pt);
\fill[red!80!black] (0.0000,8.0000) circle (1.2pt);
\fill[red!80!black] (12.0000,8.0000) circle (1.2pt);
\fill[red!80!black] (0.0000,8.0000) circle (1.2pt);
\fill[red!80!black] (12.0000,8.0000) circle (1.2pt);
\fill[red!80!black] (0.0000,8.0000) circle (1.2pt);
\fill[red!80!black] (12.0000,8.0000) circle (1.2pt);
\fill[red!80!black] (0.0000,8.0000) circle (1.2pt);
\fill[red!80!black] (12.0000,8.0000) circle (1.2pt);
\fill[red!80!black] (0.0000,8.0000) circle (1.2pt);
\fill[red!80!black] (12.0000,8.0000) circle (1.2pt);
\fill[red!80!black] (0.0000,8.0000) circle (1.2pt);
\fill[red!80!black] (12.0000,8.0000) circle (1.2pt);
\fill[red!80!black] (0.0000,8.0000) circle (1.2pt);
\fill[red!80!black] (12.0000,8.0000) circle (1.2pt);
\fill[red!80!black] (0.0000,8.0000) circle (1.2pt);
\fill[red!80!black] (12.0000,8.0000) circle (1.2pt);
\fill[red!80!black] (0.0000,8.0000) circle (1.2pt);
\fill[red!80!black] (12.0000,8.0000) circle (1.2pt);
\fill[red!80!black] (0.0000,8.0000) circle (1.2pt);
\fill[red!80!black] (12.0000,8.0000) circle (1.2pt);
\fill[red!80!black] (0.0000,8.0000) circle (1.2pt);
\fill[red!80!black] (12.0000,8.0000) circle (1.2pt);
\fill[red!80!black] (0.0000,8.0000) circle (1.2pt);
\fill[red!80!black] (12.0000,8.0000) circle (1.2pt);
\fill[red!80!black] (0.0000,8.0000) circle (1.2pt);
\fill[red!80!black] (12.0000,8.0000) circle (1.2pt);
\fill[red!80!black] (0.0000,8.0000) circle (1.2pt);
\fill[red!80!black] (12.0000,8.0000) circle (1.2pt);
\fill[red!80!black] (0.0000,8.0000) circle (1.2pt);
\fill[red!80!black] (12.0000,8.0000) circle (1.2pt);
\fill[red!80!black] (0.0000,8.0000) circle (1.2pt);
\fill[red!80!black] (12.0000,8.0000) circle (1.2pt);
\fill[red!80!black] (0.0000,8.0000) circle (1.2pt);
\fill[red!80!black] (12.0000,8.0000) circle (1.2pt);
\fill[red!80!black] (0.0000,8.0000) circle (1.2pt);
\fill[red!80!black] (12.0000,8.0000) circle (1.2pt);
\fill[red!80!black] (0.0000,8.0000) circle (1.2pt);
\fill[red!80!black] (12.0000,8.0000) circle (1.2pt);
\fill[red!80!black] (0.0000,8.0000) circle (1.2pt);
\fill[red!80!black] (12.0000,8.0000) circle (1.2pt);
\fill[red!80!black] (0.0000,8.0000) circle (1.2pt);
\fill[red!80!black] (12.0000,8.0000) circle (1.2pt);
\fill[red!80!black] (0.0000,8.0000) circle (1.2pt);
\fill[red!80!black] (12.0000,8.0000) circle (1.2pt);
\fill[red!80!black] (0.0000,8.0000) circle (1.2pt);
\fill[red!80!black] (12.0000,8.0000) circle (1.2pt);
\fill[red!80!black] (0.0000,8.0000) circle (1.2pt);
\fill[red!80!black] (12.0000,8.0000) circle (1.2pt);
\fill[red!80!black] (0.0000,8.0000) circle (1.2pt);
\fill[red!80!black] (12.0000,8.0000) circle (1.2pt);
\fill[red!80!black] (0.0000,8.0000) circle (1.2pt);
\fill[red!80!black] (12.0000,8.0000) circle (1.2pt);
\fill[red!80!black] (0.0000,8.0000) circle (1.2pt);
\fill[red!80!black] (12.0000,8.0000) circle (1.2pt);
\fill[red!80!black] (0.0000,8.0000) circle (1.2pt);
\fill[red!80!black] (12.0000,8.0000) circle (1.2pt);
\fill[red!80!black] (0.0000,8.0000) circle (1.2pt);
\fill[red!80!black] (12.0000,8.0000) circle (1.2pt);
\fill[red!80!black] (0.0000,8.0000) circle (1.2pt);
\fill[red!80!black] (12.0000,8.0000) circle (1.2pt);
\fill[red!80!black] (0.0000,8.0000) circle (1.2pt);
\fill[red!80!black] (12.0000,8.0000) circle (1.2pt);
\fill[red!80!black] (0.0000,8.0000) circle (1.2pt);
\fill[red!80!black] (12.0000,8.0000) circle (1.2pt);
\fill[red!80!black] (0.0000,8.0000) circle (1.2pt);
\fill[red!80!black] (12.0000,8.0000) circle (1.2pt);
\fill[red!80!black] (0.0000,8.0000) circle (1.2pt);
\fill[red!80!black] (12.0000,8.0000) circle (1.2pt);
\fill[red!80!black] (0.0000,8.0000) circle (1.2pt);
\fill[red!80!black] (12.0000,8.0000) circle (1.2pt);
\fill[red!80!black] (0.0000,8.0000) circle (1.2pt);
\fill[red!80!black] (12.0000,8.0000) circle (1.2pt);
\fill[red!80!black] (0.0000,8.0000) circle (1.2pt);
\fill[red!80!black] (12.0000,8.0000) circle (1.2pt);
\fill[red!80!black] (0.0000,8.0000) circle (1.2pt);
\fill[red!80!black] (12.0000,8.0000) circle (1.2pt);
\fill[red!80!black] (0.0000,8.0000) circle (1.2pt);
\fill[red!80!black] (12.0000,8.0000) circle (1.2pt);
\fill[red!80!black] (0.0000,8.0000) circle (1.2pt);
\fill[red!80!black] (12.0000,8.0000) circle (1.2pt);
\fill[red!80!black] (0.0000,8.0000) circle (1.2pt);
\fill[red!80!black] (12.0000,8.0000) circle (1.2pt);
\fill[red!80!black] (0.0000,8.0000) circle (1.2pt);
\fill[red!80!black] (12.0000,8.0000) circle (1.2pt);
\fill[red!80!black] (0.0000,8.0000) circle (1.2pt);
\fill[red!80!black] (12.0000,8.0000) circle (1.2pt);
\fill[red!80!black] (0.0000,8.0000) circle (1.2pt);
\fill[red!80!black] (12.0000,8.0000) circle (1.2pt);
\fill[red!80!black] (0.0000,8.0000) circle (1.2pt);
\fill[red!80!black] (12.0000,8.0000) circle (1.2pt);
\fill[red!80!black] (0.0000,8.0000) circle (1.2pt);
\fill[red!80!black] (12.0000,8.0000) circle (1.2pt);
\fill[red!80!black] (0.0000,8.0000) circle (1.2pt);
\fill[red!80!black] (12.0000,8.0000) circle (1.2pt);
\fill[red!80!black] (0.0000,8.0000) circle (1.2pt);
\fill[red!80!black] (12.0000,8.0000) circle (1.2pt);
\fill[red!80!black] (0.0000,8.0000) circle (1.2pt);
\fill[red!80!black] (12.0000,8.0000) circle (1.2pt);
\fill[red!80!black] (0.0000,8.0000) circle (1.2pt);
\fill[red!80!black] (12.0000,8.0000) circle (1.2pt);
\fill[red!80!black] (0.0000,8.0000) circle (1.2pt);
\fill[red!80!black] (12.0000,8.0000) circle (1.2pt);
\fill[red!80!black] (0.0000,8.0000) circle (1.2pt);
\fill[red!80!black] (12.0000,8.0000) circle (1.2pt);
\fill[red!80!black] (0.0000,8.0000) circle (1.2pt);
\fill[red!80!black] (12.0000,8.0000) circle (1.2pt);
\fill[red!80!black] (0.0000,8.0000) circle (1.2pt);
\fill[red!80!black] (12.0000,8.0000) circle (1.2pt);
\fill[red!80!black] (0.0000,8.0000) circle (1.2pt);
\fill[red!80!black] (12.0000,8.0000) circle (1.2pt);
\fill[red!80!black] (0.0000,8.0000) circle (1.2pt);
\fill[red!80!black] (12.0000,8.0000) circle (1.2pt);
\fill[red!80!black] (0.0000,8.0000) circle (1.2pt);
\fill[red!80!black] (12.0000,8.0000) circle (1.2pt);
\fill[red!80!black] (0.0000,8.0000) circle (1.2pt);
\fill[red!80!black] (12.0000,8.0000) circle (1.2pt);
\fill[red!80!black] (0.0000,8.0000) circle (1.2pt);
\fill[red!80!black] (12.0000,8.0000) circle (1.2pt);
\fill[red!80!black] (0.0000,8.0000) circle (1.2pt);
\fill[red!80!black] (12.0000,8.0000) circle (1.2pt);
\fill[red!80!black] (0.0000,8.0000) circle (1.2pt);
\fill[red!80!black] (12.0000,8.0000) circle (1.2pt);
\fill[red!80!black] (0.0000,8.0000) circle (1.2pt);
\fill[red!80!black] (12.0000,8.0000) circle (1.2pt);
\fill[red!80!black] (0.0000,8.0000) circle (1.2pt);
\fill[red!80!black] (12.0000,8.0000) circle (1.2pt);
\fill[red!80!black] (0.0000,8.0000) circle (1.2pt);
\fill[red!80!black] (12.0000,8.0000) circle (1.2pt);
\fill[red!80!black] (0.0000,8.0000) circle (1.2pt);
\fill[red!80!black] (12.0000,8.0000) circle (1.2pt);
\fill[red!80!black] (0.0000,8.0000) circle (1.2pt);
\fill[red!80!black] (12.0000,8.0000) circle (1.2pt);
\fill[red!80!black] (0.0000,8.0000) circle (1.2pt);
\fill[red!80!black] (12.0000,8.0000) circle (1.2pt);
\fill[red!80!black] (0.0000,8.0000) circle (1.2pt);
\fill[red!80!black] (12.0000,8.0000) circle (1.2pt);
\fill[red!80!black] (0.0000,8.0000) circle (1.2pt);
\fill[red!80!black] (12.0000,8.0000) circle (1.2pt);
\fill[red!80!black] (0.0000,8.0000) circle (1.2pt);
\fill[red!80!black] (12.0000,8.0000) circle (1.2pt);
\fill[red!80!black] (0.0000,8.0000) circle (1.2pt);
\fill[red!80!black] (12.0000,8.0000) circle (1.2pt);
\fill[red!80!black] (0.0000,8.0000) circle (1.2pt);
\fill[red!80!black] (12.0000,8.0000) circle (1.2pt);
\fill[red!80!black] (0.0000,8.0000) circle (1.2pt);
\fill[red!80!black] (12.0000,8.0000) circle (1.2pt);
\fill[red!80!black] (0.0000,8.0000) circle (1.2pt);
\fill[red!80!black] (12.0000,8.0000) circle (1.2pt);
\fill[red!80!black] (0.0000,8.0000) circle (1.2pt);
\fill[red!80!black] (12.0000,8.0000) circle (1.2pt);
\fill[red!80!black] (0.0000,8.0000) circle (1.2pt);
\fill[red!80!black] (12.0000,8.0000) circle (1.2pt);
\fill[red!80!black] (0.0000,8.0000) circle (1.2pt);
\fill[red!80!black] (12.0000,8.0000) circle (1.2pt);
\fill[red!80!black] (0.0000,8.0000) circle (1.2pt);
\fill[red!80!black] (12.0000,8.0000) circle (1.2pt);
\fill[red!80!black] (0.0000,8.0000) circle (1.2pt);
\fill[red!80!black] (12.0000,8.0000) circle (1.2pt);
\fill[red!80!black] (0.0000,8.0000) circle (1.2pt);
\fill[red!80!black] (12.0000,8.0000) circle (1.2pt);
\fill[red!80!black] (0.0000,8.0000) circle (1.2pt);
\fill[red!80!black] (12.0000,8.0000) circle (1.2pt);
\fill[red!80!black] (0.0000,8.0000) circle (1.2pt);
\fill[red!80!black] (12.0000,8.0000) circle (1.2pt);
\fill[red!80!black] (0.0000,8.0000) circle (1.2pt);
\fill[red!80!black] (12.0000,8.0000) circle (1.2pt);
\fill[red!80!black] (0.0000,8.0000) circle (1.2pt);
\fill[red!80!black] (12.0000,8.0000) circle (1.2pt);
\fill[red!80!black] (0.0000,8.0000) circle (1.2pt);
\fill[red!80!black] (12.0000,8.0000) circle (1.2pt);
\fill[red!80!black] (0.0000,8.0000) circle (1.2pt);
\fill[red!80!black] (12.0000,8.0000) circle (1.2pt);
\fill[red!80!black] (0.0000,8.0000) circle (1.2pt);
\fill[red!80!black] (12.0000,8.0000) circle (1.2pt);
\fill[red!80!black] (0.0000,8.0000) circle (1.2pt);
\fill[red!80!black] (12.0000,8.0000) circle (1.2pt);
\fill[red!80!black] (0.0000,8.0000) circle (1.2pt);
\fill[red!80!black] (12.0000,8.0000) circle (1.2pt);
\fill[red!80!black] (0.0000,8.0000) circle (1.2pt);
\fill[red!80!black] (12.0000,8.0000) circle (1.2pt);
\fill[red!80!black] (0.0000,8.0000) circle (1.2pt);
\fill[red!80!black] (12.0000,8.0000) circle (1.2pt);
\fill[red!80!black] (0.0000,8.0000) circle (1.2pt);
\fill[red!80!black] (12.0000,8.0000) circle (1.2pt);
\fill[red!80!black] (0.0000,8.0000) circle (1.2pt);
\fill[red!80!black] (12.0000,8.0000) circle (1.2pt);
\fill[red!80!black] (0.0000,8.0000) circle (1.2pt);
\fill[red!80!black] (12.0000,8.0000) circle (1.2pt);
\fill[red!80!black] (0.0000,8.0000) circle (1.2pt);
\fill[red!80!black] (12.0000,8.0000) circle (1.2pt);
\fill[red!80!black] (0.0000,8.0000) circle (1.2pt);
\fill[red!80!black] (12.0000,8.0000) circle (1.2pt);
\fill[red!80!black] (0.0000,8.0000) circle (1.2pt);
\fill[red!80!black] (12.0000,8.0000) circle (1.2pt);
\fill[red!80!black] (0.0000,8.0000) circle (1.2pt);
\fill[red!80!black] (12.0000,8.0000) circle (1.2pt);
\fill[red!80!black] (0.0000,8.0000) circle (1.2pt);
\fill[red!80!black] (12.0000,8.0000) circle (1.2pt);
\fill[red!80!black] (0.0000,8.0000) circle (1.2pt);
\fill[red!80!black] (12.0000,8.0000) circle (1.2pt);
\fill[red!80!black] (0.0000,8.0000) circle (1.2pt);
\fill[red!80!black] (12.0000,8.0000) circle (1.2pt);
\fill[red!80!black] (0.0000,8.0000) circle (1.2pt);
\fill[red!80!black] (12.0000,8.0000) circle (1.2pt);
\fill[red!80!black] (0.0000,8.0000) circle (1.2pt);
\fill[red!80!black] (12.0000,8.0000) circle (1.2pt);
\fill[red!80!black] (0.0000,8.0000) circle (1.2pt);
\fill[red!80!black] (12.0000,8.0000) circle (1.2pt);
\fill[red!80!black] (0.0000,8.0000) circle (1.2pt);
\fill[red!80!black] (12.0000,8.0000) circle (1.2pt);
\fill[red!80!black] (0.0000,8.0000) circle (1.2pt);
\fill[red!80!black] (12.0000,8.0000) circle (1.2pt);
\fill[red!80!black] (0.0000,8.0000) circle (1.2pt);
\fill[red!80!black] (12.0000,8.0000) circle (1.2pt);
\fill[red!80!black] (0.0000,8.0000) circle (1.2pt);
\fill[red!80!black] (12.0000,8.0000) circle (1.2pt);
\fill[red!80!black] (0.0000,8.0000) circle (1.2pt);
\fill[red!80!black] (12.0000,8.0000) circle (1.2pt);
\fill[red!80!black] (0.0000,8.0000) circle (1.2pt);
\fill[red!80!black] (12.0000,8.0000) circle (1.2pt);
\fill[red!80!black] (0.0000,8.0000) circle (1.2pt);
\fill[red!80!black] (12.0000,8.0000) circle (1.2pt);
\fill[red!80!black] (0.0000,8.0000) circle (1.2pt);
\fill[red!80!black] (12.0000,8.0000) circle (1.2pt);
\fill[red!80!black] (0.0000,8.0000) circle (1.2pt);
\fill[red!80!black] (12.0000,8.0000) circle (1.2pt);
\fill[red!80!black] (0.0000,8.0000) circle (1.2pt);
\fill[red!80!black] (12.0000,8.0000) circle (1.2pt);
\fill[red!80!black] (0.0000,8.0000) circle (1.2pt);
\fill[red!80!black] (12.0000,8.0000) circle (1.2pt);
\fill[red!80!black] (0.0000,8.0000) circle (1.2pt);
\fill[red!80!black] (12.0000,8.0000) circle (1.2pt);
\fill[red!80!black] (0.0000,8.0000) circle (1.2pt);
\fill[red!80!black] (12.0000,8.0000) circle (1.2pt);
\fill[red!80!black] (0.0000,8.0000) circle (1.2pt);
\fill[red!80!black] (12.0000,8.0000) circle (1.2pt);
\fill[red!80!black] (0.0000,8.0000) circle (1.2pt);
\fill[red!80!black] (12.0000,8.0000) circle (1.2pt);
\fill[red!80!black] (0.0000,8.0000) circle (1.2pt);
\fill[red!80!black] (12.0000,8.0000) circle (1.2pt);
\fill[red!80!black] (0.0000,8.0000) circle (1.2pt);
\fill[red!80!black] (12.0000,8.0000) circle (1.2pt);
\fill[red!80!black] (0.0000,8.0000) circle (1.2pt);
\fill[red!80!black] (12.0000,8.0000) circle (1.2pt);
\fill[red!80!black] (0.0000,8.0000) circle (1.2pt);
\fill[red!80!black] (12.0000,8.0000) circle (1.2pt);
\fill[red!80!black] (0.0000,8.0000) circle (1.2pt);
\fill[red!80!black] (12.0000,8.0000) circle (1.2pt);
\fill[red!80!black] (0.0000,8.0000) circle (1.2pt);
\fill[red!80!black] (12.0000,8.0000) circle (1.2pt);
\fill[red!80!black] (0.0000,8.0000) circle (1.2pt);
\fill[red!80!black] (12.0000,8.0000) circle (1.2pt);
\fill[red!80!black] (0.0000,8.0000) circle (1.2pt);
\fill[red!80!black] (12.0000,8.0000) circle (1.2pt);
\fill[red!80!black] (0.0000,8.0000) circle (1.2pt);
\fill[red!80!black] (12.0000,8.0000) circle (1.2pt);
\fill[red!80!black] (0.0000,8.0000) circle (1.2pt);
\fill[red!80!black] (12.0000,8.0000) circle (1.2pt);
\fill[red!80!black] (0.0000,8.0000) circle (1.2pt);
\fill[red!80!black] (12.0000,8.0000) circle (1.2pt);
\fill[red!80!black] (0.0000,8.0000) circle (1.2pt);
\fill[red!80!black] (12.0000,8.0000) circle (1.2pt);
\fill[red!80!black] (0.0000,8.0000) circle (1.2pt);
\fill[red!80!black] (12.0000,8.0000) circle (1.2pt);
\fill[red!80!black] (0.0000,8.0000) circle (1.2pt);
\fill[red!80!black] (12.0000,8.0000) circle (1.2pt);
\fill[red!80!black] (0.0000,8.0000) circle (1.2pt);
\fill[red!80!black] (12.0000,8.0000) circle (1.2pt);
\fill[red!80!black] (0.0000,8.0000) circle (1.2pt);
\fill[red!80!black] (12.0000,8.0000) circle (1.2pt);
\fill[red!80!black] (0.0000,8.0000) circle (1.2pt);
\fill[red!80!black] (12.0000,8.0000) circle (1.2pt);
\fill[red!80!black] (0.0000,8.0000) circle (1.2pt);
\fill[red!80!black] (12.0000,8.0000) circle (1.2pt);
\fill[red!80!black] (0.0000,8.0000) circle (1.2pt);
\fill[red!80!black] (12.0000,8.0000) circle (1.2pt);
\fill[red!80!black] (0.0000,8.0000) circle (1.2pt);
\fill[red!80!black] (12.0000,8.0000) circle (1.2pt);
\fill[red!80!black] (0.0000,8.0000) circle (1.2pt);
\fill[red!80!black] (12.0000,8.0000) circle (1.2pt);
\fill[red!80!black] (0.0000,8.0000) circle (1.2pt);
\fill[red!80!black] (12.0000,8.0000) circle (1.2pt);
\fill[red!80!black] (0.0000,8.0000) circle (1.2pt);
\fill[red!80!black] (12.0000,8.0000) circle (1.2pt);
\fill[red!80!black] (0.0000,8.0000) circle (1.2pt);
\fill[red!80!black] (12.0000,8.0000) circle (1.2pt);
\fill[red!80!black] (0.0000,8.0000) circle (1.2pt);
\fill[red!80!black] (12.0000,8.0000) circle (1.2pt);
\fill[red!80!black] (0.0000,8.0000) circle (1.2pt);
\fill[red!80!black] (12.0000,8.0000) circle (1.2pt);
\fill[red!80!black] (0.0000,8.0000) circle (1.2pt);
\fill[red!80!black] (12.0000,8.0000) circle (1.2pt);
\fill[red!80!black] (0.0000,8.0000) circle (1.2pt);
\fill[red!80!black] (12.0000,8.0000) circle (1.2pt);
\fill[red!80!black] (0.0000,8.0000) circle (1.2pt);
\fill[red!80!black] (12.0000,8.0000) circle (1.2pt);
\fill[red!80!black] (0.0000,8.0000) circle (1.2pt);
\fill[red!80!black] (12.0000,8.0000) circle (1.2pt);
\fill[red!80!black] (0.0000,8.0000) circle (1.2pt);
\fill[red!80!black] (12.0000,8.0000) circle (1.2pt);
\fill[red!80!black] (0.0000,8.0000) circle (1.2pt);
\fill[red!80!black] (12.0000,8.0000) circle (1.2pt);
\fill[red!80!black] (0.0000,8.0000) circle (1.2pt);
\fill[red!80!black] (12.0000,8.0000) circle (1.2pt);
\fill[red!80!black] (0.0000,8.0000) circle (1.2pt);
\fill[red!80!black] (12.0000,8.0000) circle (1.2pt);
\fill[red!80!black] (0.0000,8.0000) circle (1.2pt);
\fill[red!80!black] (12.0000,8.0000) circle (1.2pt);
\fill[red!80!black] (0.0000,8.0000) circle (1.2pt);
\fill[red!80!black] (12.0000,8.0000) circle (1.2pt);
\fill[red!80!black] (0.0000,8.0000) circle (1.2pt);
\fill[red!80!black] (12.0000,8.0000) circle (1.2pt);
\fill[red!80!black] (0.0000,8.0000) circle (1.2pt);
\fill[red!80!black] (12.0000,8.0000) circle (1.2pt);
\fill[red!80!black] (0.0000,8.0000) circle (1.2pt);
\fill[red!80!black] (12.0000,8.0000) circle (1.2pt);
\fill[red!80!black] (0.0000,8.0000) circle (1.2pt);
\fill[red!80!black] (12.0000,8.0000) circle (1.2pt);
\fill[red!80!black] (0.0000,8.0000) circle (1.2pt);
\fill[red!80!black] (12.0000,8.0000) circle (1.2pt);
\fill[red!80!black] (0.0000,8.0000) circle (1.2pt);
\fill[red!80!black] (12.0000,8.0000) circle (1.2pt);
\fill[red!80!black] (0.0000,8.0000) circle (1.2pt);
\fill[red!80!black] (12.0000,8.0000) circle (1.2pt);
\fill[red!80!black] (0.0000,8.0000) circle (1.2pt);
\fill[red!80!black] (12.0000,8.0000) circle (1.2pt);
\fill[red!80!black] (0.0000,8.0000) circle (1.2pt);
\fill[red!80!black] (12.0000,8.0000) circle (1.2pt);
\fill[red!80!black] (0.0000,8.0000) circle (1.2pt);
\fill[red!80!black] (12.0000,8.0000) circle (1.2pt);
\fill[red!80!black] (0.0000,8.0000) circle (1.2pt);
\fill[red!80!black] (12.0000,8.0000) circle (1.2pt);
\fill[red!80!black] (0.0000,8.0000) circle (1.2pt);
\fill[red!80!black] (12.0000,8.0000) circle (1.2pt);
\fill[red!80!black] (0.0000,8.0000) circle (1.2pt);
\fill[red!80!black] (12.0000,8.0000) circle (1.2pt);
\fill[red!80!black] (0.0000,8.0000) circle (1.2pt);
\fill[red!80!black] (12.0000,8.0000) circle (1.2pt);
\fill[red!80!black] (0.0000,8.0000) circle (1.2pt);
\fill[red!80!black] (12.0000,8.0000) circle (1.2pt);
\fill[red!80!black] (0.0000,8.0000) circle (1.2pt);
\fill[red!80!black] (12.0000,8.0000) circle (1.2pt);
\fill[red!80!black] (0.0000,8.0000) circle (1.2pt);
\fill[red!80!black] (12.0000,8.0000) circle (1.2pt);
\fill[red!80!black] (0.0000,8.0000) circle (1.2pt);
\fill[red!80!black] (12.0000,8.0000) circle (1.2pt);
\fill[red!80!black] (0.0000,8.0000) circle (1.2pt);
\fill[red!80!black] (12.0000,8.0000) circle (1.2pt);
\fill[red!80!black] (0.0000,8.0000) circle (1.2pt);
\fill[red!80!black] (12.0000,8.0000) circle (1.2pt);
\fill[red!80!black] (0.0000,8.0000) circle (1.2pt);
\fill[red!80!black] (12.0000,8.0000) circle (1.2pt);
\fill[red!80!black] (0.0000,8.0000) circle (1.2pt);
\fill[red!80!black] (12.0000,8.0000) circle (1.2pt);
\fill[red!80!black] (0.0000,8.0000) circle (1.2pt);
\fill[red!80!black] (12.0000,8.0000) circle (1.2pt);
\fill[red!80!black] (0.0000,8.0000) circle (1.2pt);
\fill[red!80!black] (12.0000,8.0000) circle (1.2pt);
\fill[red!80!black] (0.0000,8.0000) circle (1.2pt);
\fill[red!80!black] (12.0000,8.0000) circle (1.2pt);
\fill[red!80!black] (0.0000,8.0000) circle (1.2pt);
\fill[red!80!black] (12.0000,8.0000) circle (1.2pt);
\fill[red!80!black] (0.0000,8.0000) circle (1.2pt);
\fill[red!80!black] (12.0000,8.0000) circle (1.2pt);
\fill[red!80!black] (0.0000,8.0000) circle (1.2pt);
\fill[red!80!black] (12.0000,8.0000) circle (1.2pt);
\fill[red!80!black] (0.0000,8.0000) circle (1.2pt);
\fill[red!80!black] (12.0000,8.0000) circle (1.2pt);
\fill[red!80!black] (0.0000,8.0000) circle (1.2pt);
\fill[red!80!black] (12.0000,8.0000) circle (1.2pt);
\fill[red!80!black] (0.0000,8.0000) circle (1.2pt);
\fill[red!80!black] (12.0000,8.0000) circle (1.2pt);
\fill[red!80!black] (0.0000,8.0000) circle (1.2pt);
\fill[red!80!black] (12.0000,8.0000) circle (1.2pt);
\fill[red!80!black] (0.0000,8.0000) circle (1.2pt);
\fill[red!80!black] (12.0000,8.0000) circle (1.2pt);
\fill[red!80!black] (0.0000,8.0000) circle (1.2pt);
\fill[red!80!black] (12.0000,8.0000) circle (1.2pt);
\fill[red!80!black] (0.0000,8.0000) circle (1.2pt);
\fill[red!80!black] (12.0000,8.0000) circle (1.2pt);
\fill[red!80!black] (0.0000,8.0000) circle (1.2pt);
\fill[red!80!black] (12.0000,8.0000) circle (1.2pt);
\fill[red!80!black] (0.0000,8.0000) circle (1.2pt);
\fill[red!80!black] (12.0000,8.0000) circle (1.2pt);
\fill[red!80!black] (0.0000,8.0000) circle (1.2pt);
\fill[red!80!black] (12.0000,8.0000) circle (1.2pt);
\fill[red!80!black] (0.0000,8.0000) circle (1.2pt);
\fill[red!80!black] (12.0000,8.0000) circle (1.2pt);
\fill[red!80!black] (0.0000,8.0000) circle (1.2pt);
\fill[red!80!black] (12.0000,8.0000) circle (1.2pt);
\fill[red!80!black] (0.0000,8.0000) circle (1.2pt);
\fill[red!80!black] (12.0000,8.0000) circle (1.2pt);
\fill[red!80!black] (0.0000,8.0000) circle (1.2pt);
\fill[red!80!black] (12.0000,8.0000) circle (1.2pt);
\fill[red!80!black] (0.0000,8.0000) circle (1.2pt);
\fill[red!80!black] (12.0000,8.0000) circle (1.2pt);
\fill[red!80!black] (0.0000,8.0000) circle (1.2pt);
\fill[red!80!black] (12.0000,8.0000) circle (1.2pt);
\fill[red!80!black] (0.0000,8.0000) circle (1.2pt);
\fill[red!80!black] (12.0000,8.0000) circle (1.2pt);
\fill[red!80!black] (0.0000,8.0000) circle (1.2pt);
\fill[red!80!black] (12.0000,8.0000) circle (1.2pt);
\fill[red!80!black] (0.0000,8.0000) circle (1.2pt);
\fill[red!80!black] (12.0000,8.0000) circle (1.2pt);
\fill[red!80!black] (0.0000,8.0000) circle (1.2pt);
\fill[red!80!black] (12.0000,8.0000) circle (1.2pt);
\fill[red!80!black] (0.0000,8.0000) circle (1.2pt);
\fill[red!80!black] (12.0000,8.0000) circle (1.2pt);
\fill[red!80!black] (0.0000,8.0000) circle (1.2pt);
\fill[red!80!black] (12.0000,8.0000) circle (1.2pt);
\fill[red!80!black] (0.0000,8.0000) circle (1.2pt);
\fill[red!80!black] (12.0000,8.0000) circle (1.2pt);
\fill[red!80!black] (0.0000,8.0000) circle (1.2pt);
\fill[red!80!black] (12.0000,8.0000) circle (1.2pt);
\fill[red!80!black] (0.0000,8.0000) circle (1.2pt);
\fill[red!80!black] (12.0000,8.0000) circle (1.2pt);
\fill[red!80!black] (0.0000,8.0000) circle (1.2pt);
\fill[red!80!black] (12.0000,8.0000) circle (1.2pt);
\fill[red!80!black] (0.0000,8.0000) circle (1.2pt);
\fill[red!80!black] (12.0000,8.0000) circle (1.2pt);
\fill[red!80!black] (0.0000,8.0000) circle (1.2pt);
\fill[red!80!black] (12.0000,8.0000) circle (1.2pt);
\fill[red!80!black] (0.0000,8.0000) circle (1.2pt);
\fill[red!80!black] (12.0000,8.0000) circle (1.2pt);
\fill[red!80!black] (0.0000,8.0000) circle (1.2pt);
\fill[red!80!black] (12.0000,8.0000) circle (1.2pt);
\fill[red!80!black] (0.0000,8.0000) circle (1.2pt);
\fill[red!80!black] (12.0000,8.0000) circle (1.2pt);
\fill[red!80!black] (0.0000,8.0000) circle (1.2pt);
\fill[red!80!black] (12.0000,8.0000) circle (1.2pt);
\fill[red!80!black] (0.0000,8.0000) circle (1.2pt);
\fill[red!80!black] (12.0000,8.0000) circle (1.2pt);
\fill[red!80!black] (0.0000,8.0000) circle (1.2pt);
\fill[red!80!black] (12.0000,8.0000) circle (1.2pt);
\fill[red!80!black] (0.0000,8.0000) circle (1.2pt);
\fill[red!80!black] (12.0000,8.0000) circle (1.2pt);
\fill[red!80!black] (0.0000,8.0000) circle (1.2pt);
\fill[red!80!black] (12.0000,8.0000) circle (1.2pt);
\fill[red!80!black] (0.0000,8.0000) circle (1.2pt);
\fill[red!80!black] (12.0000,8.0000) circle (1.2pt);
\fill[red!80!black] (0.0000,8.0000) circle (1.2pt);
\fill[red!80!black] (12.0000,8.0000) circle (1.2pt);
\fill[red!80!black] (0.0000,8.0000) circle (1.2pt);
\fill[red!80!black] (12.0000,8.0000) circle (1.2pt);
\fill[red!80!black] (0.0000,8.0000) circle (1.2pt);
\fill[red!80!black] (12.0000,8.0000) circle (1.2pt);
\fill[red!80!black] (0.0000,8.0000) circle (1.2pt);
\fill[red!80!black] (12.0000,8.0000) circle (1.2pt);
\fill[red!80!black] (0.0000,8.0000) circle (1.2pt);
\fill[red!80!black] (12.0000,8.0000) circle (1.2pt);
\fill[red!80!black] (0.0000,8.0000) circle (1.2pt);
\fill[red!80!black] (12.0000,8.0000) circle (1.2pt);
\fill[red!80!black] (0.0000,8.0000) circle (1.2pt);
\fill[red!80!black] (12.0000,8.0000) circle (1.2pt);
\fill[red!80!black] (0.0000,8.0000) circle (1.2pt);
\fill[red!80!black] (12.0000,8.0000) circle (1.2pt);
\fill[red!80!black] (0.0000,8.0000) circle (1.2pt);
\fill[red!80!black] (12.0000,8.0000) circle (1.2pt);
\fill[red!80!black] (0.0000,8.0000) circle (1.2pt);
\fill[red!80!black] (12.0000,8.0000) circle (1.2pt);
\fill[red!80!black] (0.0000,8.0000) circle (1.2pt);
\fill[red!80!black] (12.0000,8.0000) circle (1.2pt);
\fill[red!80!black] (0.0000,8.0000) circle (1.2pt);
\fill[red!80!black] (12.0000,8.0000) circle (1.2pt);
\fill[red!80!black] (0.0000,8.0000) circle (1.2pt);
\fill[red!80!black] (12.0000,8.0000) circle (1.2pt);
\fill[red!80!black] (0.0000,8.0000) circle (1.2pt);
\fill[red!80!black] (12.0000,8.0000) circle (1.2pt);
\fill[red!80!black] (0.0000,8.0000) circle (1.2pt);
\fill[red!80!black] (12.0000,8.0000) circle (1.2pt);
\fill[red!80!black] (0.0000,8.0000) circle (1.2pt);
\fill[red!80!black] (12.0000,8.0000) circle (1.2pt);
\fill[red!80!black] (0.0000,8.0000) circle (1.2pt);
\fill[red!80!black] (12.0000,8.0000) circle (1.2pt);
\fill[red!80!black] (0.0000,8.0000) circle (1.2pt);
\fill[red!80!black] (12.0000,8.0000) circle (1.2pt);
\fill[red!80!black] (0.0000,8.0000) circle (1.2pt);
\fill[red!80!black] (12.0000,8.0000) circle (1.2pt);
\fill[red!80!black] (0.0000,8.0000) circle (1.2pt);
\fill[red!80!black] (12.0000,8.0000) circle (1.2pt);
\fill[red!80!black] (0.0000,8.0000) circle (1.2pt);
\fill[red!80!black] (12.0000,8.0000) circle (1.2pt);
\fill[red!80!black] (0.0000,8.0000) circle (1.2pt);
\fill[red!80!black] (12.0000,8.0000) circle (1.2pt);
\fill[red!80!black] (0.0000,8.0000) circle (1.2pt);
\fill[red!80!black] (12.0000,8.0000) circle (1.2pt);
\fill[red!80!black] (0.0000,8.0000) circle (1.2pt);
\fill[red!80!black] (12.0000,8.0000) circle (1.2pt);
\fill[red!80!black] (0.0000,8.0000) circle (1.2pt);
\fill[red!80!black] (12.0000,8.0000) circle (1.2pt);
\fill[red!80!black] (0.0000,8.0000) circle (1.2pt);
\fill[red!80!black] (12.0000,8.0000) circle (1.2pt);
\fill[red!80!black] (0.0000,8.0000) circle (1.2pt);
\fill[red!80!black] (12.0000,8.0000) circle (1.2pt);
\fill[red!80!black] (0.0000,8.0000) circle (1.2pt);
\fill[red!80!black] (12.0000,8.0000) circle (1.2pt);
\fill[red!80!black] (0.0000,8.0000) circle (1.2pt);
\fill[red!80!black] (12.0000,8.0000) circle (1.2pt);
\fill[red!80!black] (0.0000,8.0000) circle (1.2pt);
\fill[red!80!black] (12.0000,8.0000) circle (1.2pt);
\fill[red!80!black] (0.0000,8.0000) circle (1.2pt);
\fill[red!80!black] (12.0000,8.0000) circle (1.2pt);
\fill[red!80!black] (0.0000,8.0000) circle (1.2pt);
\fill[red!80!black] (12.0000,8.0000) circle (1.2pt);
\fill[red!80!black] (0.0000,8.0000) circle (1.2pt);
\fill[red!80!black] (12.0000,8.0000) circle (1.2pt);
\fill[red!80!black] (0.0000,8.0000) circle (1.2pt);
\fill[red!80!black] (12.0000,8.0000) circle (1.2pt);
\fill[red!80!black] (0.0000,8.0000) circle (1.2pt);
\fill[red!80!black] (12.0000,8.0000) circle (1.2pt);
\fill[red!80!black] (0.0000,8.0000) circle (1.2pt);
\fill[red!80!black] (12.0000,8.0000) circle (1.2pt);
\fill[red!80!black] (0.0000,8.0000) circle (1.2pt);
\fill[red!80!black] (12.0000,8.0000) circle (1.2pt);
\fill[red!80!black] (0.0000,8.0000) circle (1.2pt);
\fill[red!80!black] (12.0000,8.0000) circle (1.2pt);
\fill[red!80!black] (0.0000,8.0000) circle (1.2pt);
\fill[red!80!black] (12.0000,8.0000) circle (1.2pt);
\fill[red!80!black] (0.0000,8.0000) circle (1.2pt);
\fill[red!80!black] (12.0000,8.0000) circle (1.2pt);
\fill[red!80!black] (0.0000,8.0000) circle (1.2pt);
\fill[red!80!black] (12.0000,8.0000) circle (1.2pt);
\fill[red!80!black] (0.0000,8.0000) circle (1.2pt);
\fill[red!80!black] (12.0000,8.0000) circle (1.2pt);
\fill[red!80!black] (0.0000,8.0000) circle (1.2pt);
\fill[red!80!black] (12.0000,8.0000) circle (1.2pt);
\fill[red!80!black] (0.0000,8.0000) circle (1.2pt);
\fill[red!80!black] (12.0000,8.0000) circle (1.2pt);
\fill[red!80!black] (0.0000,8.0000) circle (1.2pt);
\fill[red!80!black] (12.0000,8.0000) circle (1.2pt);
\fill[red!80!black] (0.0000,8.0000) circle (1.2pt);
\fill[red!80!black] (12.0000,8.0000) circle (1.2pt);
\fill[red!80!black] (0.0000,8.0000) circle (1.2pt);
\fill[red!80!black] (12.0000,8.0000) circle (1.2pt);
\fill[red!80!black] (0.0000,8.0000) circle (1.2pt);
\fill[red!80!black] (12.0000,8.0000) circle (1.2pt);
\fill[red!80!black] (0.0000,8.0000) circle (1.2pt);
\fill[red!80!black] (12.0000,8.0000) circle (1.2pt);
\fill[red!80!black] (0.0000,8.0000) circle (1.2pt);
\fill[red!80!black] (12.0000,8.0000) circle (1.2pt);
\fill[red!80!black] (0.0000,8.0000) circle (1.2pt);
\fill[red!80!black] (12.0000,8.0000) circle (1.2pt);
\fill[red!80!black] (0.0000,8.0000) circle (1.2pt);
\fill[red!80!black] (12.0000,8.0000) circle (1.2pt);
\fill[red!80!black] (0.0000,8.0000) circle (1.2pt);
\fill[red!80!black] (12.0000,8.0000) circle (1.2pt);
\fill[red!80!black] (0.0000,8.0000) circle (1.2pt);
\fill[red!80!black] (12.0000,8.0000) circle (1.2pt);
\fill[red!80!black] (0.0000,8.0000) circle (1.2pt);
\fill[red!80!black] (12.0000,8.0000) circle (1.2pt);
\fill[red!80!black] (0.0000,8.0000) circle (1.2pt);
\fill[red!80!black] (12.0000,8.0000) circle (1.2pt);
\fill[red!80!black] (0.0000,8.0000) circle (1.2pt);
\fill[red!80!black] (12.0000,8.0000) circle (1.2pt);
\fill[red!80!black] (0.0000,8.0000) circle (1.2pt);
\fill[red!80!black] (12.0000,8.0000) circle (1.2pt);
\fill[red!80!black] (0.0000,8.0000) circle (1.2pt);
\fill[red!80!black] (12.0000,8.0000) circle (1.2pt);
\fill[red!80!black] (0.0000,8.0000) circle (1.2pt);
\fill[red!80!black] (12.0000,8.0000) circle (1.2pt);
\fill[red!80!black] (0.0000,8.0000) circle (1.2pt);
\fill[red!80!black] (12.0000,8.0000) circle (1.2pt);
\fill[red!80!black] (0.0000,8.0000) circle (1.2pt);
\fill[red!80!black] (12.0000,8.0000) circle (1.2pt);
\fill[red!80!black] (0.0000,8.0000) circle (1.2pt);
\fill[red!80!black] (12.0000,8.0000) circle (1.2pt);
\fill[red!80!black] (0.0000,8.0000) circle (1.2pt);
\fill[red!80!black] (12.0000,8.0000) circle (1.2pt);
\fill[red!80!black] (0.0000,8.0000) circle (1.2pt);
\fill[red!80!black] (12.0000,8.0000) circle (1.2pt);
\fill[red!80!black] (0.0000,8.0000) circle (1.2pt);
\fill[red!80!black] (12.0000,8.0000) circle (1.2pt);
\fill[red!80!black] (0.0000,8.0000) circle (1.2pt);
\fill[red!80!black] (12.0000,8.0000) circle (1.2pt);
\fill[red!80!black] (0.0000,8.0000) circle (1.2pt);
\fill[red!80!black] (12.0000,8.0000) circle (1.2pt);
\fill[red!80!black] (0.0000,8.0000) circle (1.2pt);
\fill[red!80!black] (12.0000,8.0000) circle (1.2pt);
\fill[red!80!black] (0.0000,8.0000) circle (1.2pt);
\fill[red!80!black] (12.0000,8.0000) circle (1.2pt);
\fill[red!80!black] (0.0000,8.0000) circle (1.2pt);
\fill[red!80!black] (12.0000,8.0000) circle (1.2pt);
\fill[red!80!black] (0.0000,8.0000) circle (1.2pt);
\fill[red!80!black] (12.0000,8.0000) circle (1.2pt);
\fill[red!80!black] (0.0000,8.0000) circle (1.2pt);
\fill[red!80!black] (12.0000,8.0000) circle (1.2pt);
\fill[red!80!black] (0.0000,8.0000) circle (1.2pt);
\fill[red!80!black] (12.0000,8.0000) circle (1.2pt);
\fill[red!80!black] (0.0000,8.0000) circle (1.2pt);
\fill[red!80!black] (12.0000,8.0000) circle (1.2pt);
\fill[red!80!black] (0.0000,8.0000) circle (1.2pt);
\fill[red!80!black] (12.0000,8.0000) circle (1.2pt);
\fill[red!80!black] (0.0000,8.0000) circle (1.2pt);
\fill[red!80!black] (12.0000,8.0000) circle (1.2pt);
\fill[red!80!black] (0.0000,8.0000) circle (1.2pt);
\fill[red!80!black] (12.0000,8.0000) circle (1.2pt);
\fill[red!80!black] (0.0000,8.0000) circle (1.2pt);
\fill[red!80!black] (12.0000,8.0000) circle (1.2pt);
\fill[red!80!black] (0.0000,8.0000) circle (1.2pt);
\fill[red!80!black] (12.0000,8.0000) circle (1.2pt);
\fill[red!80!black] (0.0000,8.0000) circle (1.2pt);
\fill[red!80!black] (12.0000,8.0000) circle (1.2pt);
\fill[red!80!black] (0.0000,8.0000) circle (1.2pt);
\fill[red!80!black] (12.0000,8.0000) circle (1.2pt);
\fill[red!80!black] (0.0000,8.0000) circle (1.2pt);
\fill[red!80!black] (12.0000,8.0000) circle (1.2pt);
\fill[red!80!black] (0.0000,8.0000) circle (1.2pt);
\fill[red!80!black] (12.0000,8.0000) circle (1.2pt);
\fill[red!80!black] (0.0000,8.0000) circle (1.2pt);
\fill[red!80!black] (12.0000,8.0000) circle (1.2pt);
\fill[red!80!black] (0.0000,8.0000) circle (1.2pt);
\fill[red!80!black] (12.0000,8.0000) circle (1.2pt);
\fill[red!80!black] (0.0000,8.0000) circle (1.2pt);
\fill[red!80!black] (12.0000,8.0000) circle (1.2pt);
\fill[red!80!black] (0.0000,8.0000) circle (1.2pt);
\fill[red!80!black] (12.0000,8.0000) circle (1.2pt);
\fill[red!80!black] (0.0000,8.0000) circle (1.2pt);
\fill[red!80!black] (12.0000,8.0000) circle (1.2pt);
\fill[red!80!black] (0.0000,8.0000) circle (1.2pt);
\fill[red!80!black] (12.0000,8.0000) circle (1.2pt);
\fill[red!80!black] (0.0000,8.0000) circle (1.2pt);
\fill[red!80!black] (12.0000,8.0000) circle (1.2pt);
\fill[red!80!black] (0.0000,8.0000) circle (1.2pt);
\fill[red!80!black] (12.0000,8.0000) circle (1.2pt);
\fill[red!80!black] (0.0000,8.0000) circle (1.2pt);
\fill[red!80!black] (12.0000,8.0000) circle (1.2pt);
\fill[red!80!black] (0.0000,8.0000) circle (1.2pt);
\fill[red!80!black] (12.0000,8.0000) circle (1.2pt);
\fill[red!80!black] (0.0000,8.0000) circle (1.2pt);
\fill[red!80!black] (12.0000,8.0000) circle (1.2pt);
\fill[red!80!black] (0.0000,8.0000) circle (1.2pt);
\fill[red!80!black] (12.0000,8.0000) circle (1.2pt);
\fill[red!80!black] (0.0000,8.0000) circle (1.2pt);
\fill[red!80!black] (12.0000,8.0000) circle (1.2pt);
\fill[red!80!black] (0.0000,8.0000) circle (1.2pt);
\fill[red!80!black] (12.0000,8.0000) circle (1.2pt);
\fill[red!80!black] (0.0000,8.0000) circle (1.2pt);
\fill[red!80!black] (12.0000,8.0000) circle (1.2pt);
\fill[red!80!black] (0.0000,8.0000) circle (1.2pt);
\fill[red!80!black] (12.0000,8.0000) circle (1.2pt);
\fill[red!80!black] (0.0000,8.0000) circle (1.2pt);
\fill[red!80!black] (12.0000,8.0000) circle (1.2pt);
\fill[red!80!black] (0.0000,8.0000) circle (1.2pt);
\fill[red!80!black] (12.0000,8.0000) circle (1.2pt);
\fill[red!80!black] (0.0000,8.0000) circle (1.2pt);
\fill[red!80!black] (12.0000,8.0000) circle (1.2pt);
\fill[red!80!black] (0.0000,8.0000) circle (1.2pt);
\fill[red!80!black] (12.0000,8.0000) circle (1.2pt);
\fill[red!80!black] (0.0000,8.0000) circle (1.2pt);
\fill[red!80!black] (12.0000,8.0000) circle (1.2pt);
\fill[red!80!black] (0.0000,8.0000) circle (1.2pt);
\fill[red!80!black] (12.0000,8.0000) circle (1.2pt);
\fill[red!80!black] (0.0000,8.0000) circle (1.2pt);
\fill[red!80!black] (12.0000,8.0000) circle (1.2pt);
\fill[red!80!black] (0.0000,8.0000) circle (1.2pt);
\fill[red!80!black] (12.0000,8.0000) circle (1.2pt);
\fill[red!80!black] (0.0000,8.0000) circle (1.2pt);
\fill[red!80!black] (12.0000,8.0000) circle (1.2pt);
\fill[red!80!black] (0.0000,8.0000) circle (1.2pt);
\fill[red!80!black] (12.0000,8.0000) circle (1.2pt);
\fill[red!80!black] (0.0000,8.0000) circle (1.2pt);
\fill[red!80!black] (12.0000,8.0000) circle (1.2pt);
\fill[red!80!black] (0.0000,8.0000) circle (1.2pt);
\fill[red!80!black] (12.0000,8.0000) circle (1.2pt);
\fill[red!80!black] (0.0000,8.0000) circle (1.2pt);
\fill[red!80!black] (12.0000,8.0000) circle (1.2pt);
\fill[red!80!black] (0.0000,8.0000) circle (1.2pt);
\fill[red!80!black] (12.0000,8.0000) circle (1.2pt);
\fill[red!80!black] (0.0000,8.0000) circle (1.2pt);
\fill[red!80!black] (12.0000,8.0000) circle (1.2pt);
\fill[red!80!black] (0.0000,8.0000) circle (1.2pt);
\fill[red!80!black] (12.0000,8.0000) circle (1.2pt);
\fill[red!80!black] (0.0000,8.0000) circle (1.2pt);
\fill[red!80!black] (12.0000,8.0000) circle (1.2pt);
\fill[red!80!black] (0.0000,8.0000) circle (1.2pt);
\fill[red!80!black] (12.0000,8.0000) circle (1.2pt);
\fill[red!80!black] (0.0000,8.0000) circle (1.2pt);
\fill[red!80!black] (12.0000,8.0000) circle (1.2pt);
\fill[red!80!black] (0.0000,8.0000) circle (1.2pt);
\fill[red!80!black] (12.0000,8.0000) circle (1.2pt);
\fill[red!80!black] (0.0000,8.0000) circle (1.2pt);
\fill[red!80!black] (12.0000,8.0000) circle (1.2pt);
\fill[red!80!black] (0.0000,8.0000) circle (1.2pt);
\fill[red!80!black] (12.0000,8.0000) circle (1.2pt);
\fill[red!80!black] (0.0000,8.0000) circle (1.2pt);
\fill[red!80!black] (12.0000,8.0000) circle (1.2pt);
\fill[red!80!black] (0.0000,8.0000) circle (1.2pt);
\fill[red!80!black] (12.0000,8.0000) circle (1.2pt);
\fill[red!80!black] (0.0000,8.0000) circle (1.2pt);
\fill[red!80!black] (12.0000,8.0000) circle (1.2pt);
\fill[red!80!black] (0.0000,8.0000) circle (1.2pt);
\fill[red!80!black] (12.0000,8.0000) circle (1.2pt);
\fill[red!80!black] (0.0000,8.0000) circle (1.2pt);
\fill[red!80!black] (12.0000,8.0000) circle (1.2pt);
\fill[red!80!black] (0.0000,8.0000) circle (1.2pt);
\fill[red!80!black] (12.0000,8.0000) circle (1.2pt);
\fill[red!80!black] (0.0000,8.0000) circle (1.2pt);
\fill[red!80!black] (12.0000,8.0000) circle (1.2pt);
\fill[red!80!black] (0.0000,8.0000) circle (1.2pt);
\fill[red!80!black] (12.0000,8.0000) circle (1.2pt);
\fill[red!80!black] (0.0000,8.0000) circle (1.2pt);
\fill[red!80!black] (12.0000,8.0000) circle (1.2pt);
\fill[red!80!black] (0.0000,8.0000) circle (1.2pt);
\fill[red!80!black] (12.0000,8.0000) circle (1.2pt);
\fill[red!80!black] (0.0000,8.0000) circle (1.2pt);
\fill[red!80!black] (12.0000,8.0000) circle (1.2pt);
\fill[red!80!black] (0.0000,8.0000) circle (1.2pt);
\fill[red!80!black] (12.0000,8.0000) circle (1.2pt);
\fill[red!80!black] (0.0000,8.0000) circle (1.2pt);
\fill[red!80!black] (12.0000,8.0000) circle (1.2pt);
\fill[red!80!black] (0.0000,8.0000) circle (1.2pt);
\fill[red!80!black] (12.0000,8.0000) circle (1.2pt);
\fill[red!80!black] (0.0000,8.0000) circle (1.2pt);
\fill[red!80!black] (12.0000,8.0000) circle (1.2pt);
\fill[red!80!black] (0.0000,8.0000) circle (1.2pt);
\fill[red!80!black] (12.0000,8.0000) circle (1.2pt);
\fill[red!80!black] (0.0000,8.0000) circle (1.2pt);
\fill[red!80!black] (12.0000,8.0000) circle (1.2pt);
\fill[red!80!black] (0.0000,8.0000) circle (1.2pt);
\fill[red!80!black] (12.0000,8.0000) circle (1.2pt);
\fill[red!80!black] (0.0000,8.0000) circle (1.2pt);
\fill[red!80!black] (12.0000,8.0000) circle (1.2pt);
\fill[red!80!black] (0.0000,8.0000) circle (1.2pt);
\fill[red!80!black] (12.0000,8.0000) circle (1.2pt);
\fill[red!80!black] (0.0000,8.0000) circle (1.2pt);
\fill[red!80!black] (12.0000,8.0000) circle (1.2pt);
\fill[red!80!black] (0.0000,8.0000) circle (1.2pt);
\fill[red!80!black] (12.0000,8.0000) circle (1.2pt);
\fill[red!80!black] (0.0000,8.0000) circle (1.2pt);
\fill[red!80!black] (12.0000,8.0000) circle (1.2pt);
\fill[red!80!black] (0.0000,8.0000) circle (1.2pt);
\fill[red!80!black] (12.0000,8.0000) circle (1.2pt);
\fill[red!80!black] (0.0000,8.0000) circle (1.2pt);
\fill[red!80!black] (12.0000,8.0000) circle (1.2pt);
\fill[red!80!black] (0.0000,8.0000) circle (1.2pt);
\fill[red!80!black] (12.0000,8.0000) circle (1.2pt);
\fill[red!80!black] (0.0000,8.0000) circle (1.2pt);
\fill[red!80!black] (12.0000,8.0000) circle (1.2pt);
\fill[red!80!black] (0.0000,8.0000) circle (1.2pt);
\fill[red!80!black] (12.0000,8.0000) circle (1.2pt);
\fill[red!80!black] (0.0000,8.0000) circle (1.2pt);
\fill[red!80!black] (12.0000,8.0000) circle (1.2pt);
\fill[red!80!black] (0.0000,8.0000) circle (1.2pt);
\fill[red!80!black] (12.0000,8.0000) circle (1.2pt);
\fill[red!80!black] (0.0000,8.0000) circle (1.2pt);
\fill[red!80!black] (12.0000,8.0000) circle (1.2pt);
\fill[red!80!black] (0.0000,8.0000) circle (1.2pt);
\fill[red!80!black] (12.0000,8.0000) circle (1.2pt);
\fill[red!80!black] (0.0000,8.0000) circle (1.2pt);
\fill[red!80!black] (12.0000,8.0000) circle (1.2pt);
\fill[red!80!black] (0.0000,8.0000) circle (1.2pt);
\fill[red!80!black] (12.0000,8.0000) circle (1.2pt);
\fill[red!80!black] (0.0000,8.0000) circle (1.2pt);
\fill[red!80!black] (12.0000,8.0000) circle (1.2pt);
\fill[red!80!black] (0.0000,8.0000) circle (1.2pt);
\fill[red!80!black] (12.0000,8.0000) circle (1.2pt);
\fill[red!80!black] (0.0000,8.0000) circle (1.2pt);
\fill[red!80!black] (12.0000,8.0000) circle (1.2pt);
\fill[red!80!black] (0.0000,8.0000) circle (1.2pt);
\fill[red!80!black] (12.0000,8.0000) circle (1.2pt);
\fill[red!80!black] (0.0000,8.0000) circle (1.2pt);
\fill[red!80!black] (12.0000,8.0000) circle (1.2pt);
\fill[red!80!black] (0.0000,8.0000) circle (1.2pt);
\fill[red!80!black] (12.0000,8.0000) circle (1.2pt);
\fill[red!80!black] (0.0000,8.0000) circle (1.2pt);
\fill[red!80!black] (12.0000,8.0000) circle (1.2pt);
\fill[red!80!black] (0.0000,8.0000) circle (1.2pt);
\fill[red!80!black] (12.0000,8.0000) circle (1.2pt);
\fill[red!80!black] (0.0000,8.0000) circle (1.2pt);
\fill[red!80!black] (12.0000,8.0000) circle (1.2pt);
\fill[red!80!black] (0.0000,8.0000) circle (1.2pt);
\fill[red!80!black] (12.0000,8.0000) circle (1.2pt);
\fill[red!80!black] (0.0000,8.0000) circle (1.2pt);
\fill[red!80!black] (12.0000,8.0000) circle (1.2pt);
\fill[red!80!black] (0.0000,8.0000) circle (1.2pt);
\fill[red!80!black] (12.0000,8.0000) circle (1.2pt);
\fill[red!80!black] (0.0000,8.0000) circle (1.2pt);
\fill[red!80!black] (12.0000,8.0000) circle (1.2pt);
\fill[red!80!black] (0.0000,8.0000) circle (1.2pt);
\fill[red!80!black] (12.0000,8.0000) circle (1.2pt);
\fill[red!80!black] (0.0000,8.0000) circle (1.2pt);
\fill[red!80!black] (12.0000,8.0000) circle (1.2pt);
\fill[red!80!black] (0.0000,8.0000) circle (1.2pt);
\fill[red!80!black] (12.0000,8.0000) circle (1.2pt);
\fill[red!80!black] (0.0000,8.0000) circle (1.2pt);
\fill[red!80!black] (12.0000,8.0000) circle (1.2pt);
\fill[red!80!black] (0.0000,8.0000) circle (1.2pt);
\fill[red!80!black] (12.0000,8.0000) circle (1.2pt);
\fill[red!80!black] (0.0000,8.0000) circle (1.2pt);
\fill[red!80!black] (12.0000,8.0000) circle (1.2pt);
\fill[red!80!black] (0.0000,8.0000) circle (1.2pt);
\fill[red!80!black] (12.0000,8.0000) circle (1.2pt);
\fill[red!80!black] (0.0000,8.0000) circle (1.2pt);
\fill[red!80!black] (12.0000,8.0000) circle (1.2pt);
\fill[red!80!black] (0.0000,8.0000) circle (1.2pt);
\fill[red!80!black] (12.0000,8.0000) circle (1.2pt);
\fill[red!80!black] (0.0000,8.0000) circle (1.2pt);
\fill[red!80!black] (12.0000,8.0000) circle (1.2pt);
\fill[red!80!black] (0.0000,8.0000) circle (1.2pt);
\fill[red!80!black] (12.0000,8.0000) circle (1.2pt);
\fill[red!80!black] (0.0000,8.0000) circle (1.2pt);
\fill[red!80!black] (12.0000,8.0000) circle (1.2pt);
\fill[red!80!black] (0.0000,8.0000) circle (1.2pt);
\fill[red!80!black] (12.0000,8.0000) circle (1.2pt);
\fill[red!80!black] (0.0000,8.0000) circle (1.2pt);
\fill[red!80!black] (12.0000,8.0000) circle (1.2pt);
\fill[red!80!black] (0.0000,8.0000) circle (1.2pt);
\fill[red!80!black] (12.0000,8.0000) circle (1.2pt);
\fill[red!80!black] (0.0000,8.0000) circle (1.2pt);
\fill[red!80!black] (12.0000,8.0000) circle (1.2pt);
\fill[red!80!black] (0.0000,8.0000) circle (1.2pt);
\fill[red!80!black] (12.0000,8.0000) circle (1.2pt);
\fill[red!80!black] (0.0000,8.0000) circle (1.2pt);
\fill[red!80!black] (12.0000,8.0000) circle (1.2pt);
\fill[red!80!black] (0.0000,8.0000) circle (1.2pt);
\fill[red!80!black] (12.0000,8.0000) circle (1.2pt);
\fill[red!80!black] (0.0000,8.0000) circle (1.2pt);
\fill[red!80!black] (12.0000,8.0000) circle (1.2pt);
\fill[red!80!black] (0.0000,8.0000) circle (1.2pt);
\fill[red!80!black] (12.0000,8.0000) circle (1.2pt);
\fill[red!80!black] (0.0000,8.0000) circle (1.2pt);
\fill[red!80!black] (12.0000,8.0000) circle (1.2pt);
\fill[red!80!black] (0.0000,8.0000) circle (1.2pt);
\fill[red!80!black] (12.0000,8.0000) circle (1.2pt);
\fill[red!80!black] (0.0000,8.0000) circle (1.2pt);
\fill[red!80!black] (12.0000,8.0000) circle (1.2pt);
\fill[red!80!black] (0.0000,8.0000) circle (1.2pt);
\fill[red!80!black] (12.0000,8.0000) circle (1.2pt);
\fill[red!80!black] (0.0000,8.0000) circle (1.2pt);
\fill[red!80!black] (12.0000,8.0000) circle (1.2pt);
\fill[red!80!black] (0.0000,8.0000) circle (1.2pt);
\fill[red!80!black] (12.0000,8.0000) circle (1.2pt);
\fill[red!80!black] (0.0000,8.0000) circle (1.2pt);
\fill[red!80!black] (12.0000,8.0000) circle (1.2pt);
\fill[red!80!black] (0.0000,8.0000) circle (1.2pt);
\fill[red!80!black] (12.0000,8.0000) circle (1.2pt);
\fill[red!80!black] (0.0000,8.0000) circle (1.2pt);
\fill[red!80!black] (12.0000,8.0000) circle (1.2pt);
\fill[red!80!black] (0.0000,8.0000) circle (1.2pt);
\fill[red!80!black] (12.0000,8.0000) circle (1.2pt);
\fill[red!80!black] (0.0000,8.0000) circle (1.2pt);
\fill[red!80!black] (12.0000,8.0000) circle (1.2pt);
\fill[red!80!black] (0.0000,8.0000) circle (1.2pt);
\fill[red!80!black] (12.0000,8.0000) circle (1.2pt);
\fill[red!80!black] (0.0000,8.0000) circle (1.2pt);
\fill[red!80!black] (12.0000,8.0000) circle (1.2pt);
\fill[red!80!black] (0.0000,8.0000) circle (1.2pt);
\fill[red!80!black] (12.0000,8.0000) circle (1.2pt);
\fill[red!80!black] (0.0000,8.0000) circle (1.2pt);
\fill[red!80!black] (12.0000,8.0000) circle (1.2pt);
\fill[red!80!black] (0.0000,8.0000) circle (1.2pt);
\fill[red!80!black] (12.0000,8.0000) circle (1.2pt);
\fill[red!80!black] (0.0000,8.0000) circle (1.2pt);
\fill[red!80!black] (12.0000,8.0000) circle (1.2pt);
\fill[red!80!black] (0.0000,8.0000) circle (1.2pt);
\fill[red!80!black] (12.0000,8.0000) circle (1.2pt);
\fill[red!80!black] (0.0000,8.0000) circle (1.2pt);
\fill[red!80!black] (12.0000,8.0000) circle (1.2pt);
\fill[red!80!black] (0.0000,8.0000) circle (1.2pt);
\fill[red!80!black] (12.0000,8.0000) circle (1.2pt);
\fill[red!80!black] (0.0000,8.0000) circle (1.2pt);
\fill[red!80!black] (12.0000,8.0000) circle (1.2pt);
\fill[red!80!black] (0.0000,8.0000) circle (1.2pt);
\fill[red!80!black] (12.0000,8.0000) circle (1.2pt);
\fill[red!80!black] (0.0000,8.0000) circle (1.2pt);
\fill[red!80!black] (12.0000,8.0000) circle (1.2pt);
\fill[red!80!black] (0.0000,8.0000) circle (1.2pt);
\fill[red!80!black] (12.0000,8.0000) circle (1.2pt);
\fill[red!80!black] (0.0000,8.0000) circle (1.2pt);
\fill[red!80!black] (12.0000,8.0000) circle (1.2pt);
\fill[red!80!black] (0.0000,8.0000) circle (1.2pt);
\fill[red!80!black] (12.0000,8.0000) circle (1.2pt);
\fill[red!80!black] (0.0000,8.0000) circle (1.2pt);
\fill[red!80!black] (12.0000,8.0000) circle (1.2pt);
\fill[red!80!black] (0.0000,8.0000) circle (1.2pt);
\fill[red!80!black] (12.0000,8.0000) circle (1.2pt);
\fill[red!80!black] (0.0000,8.0000) circle (1.2pt);
\fill[red!80!black] (12.0000,8.0000) circle (1.2pt);
\fill[red!80!black] (0.0000,8.0000) circle (1.2pt);
\fill[red!80!black] (12.0000,8.0000) circle (1.2pt);
\fill[red!80!black] (0.0000,8.0000) circle (1.2pt);
\fill[red!80!black] (12.0000,8.0000) circle (1.2pt);
\fill[red!80!black] (0.0000,8.0000) circle (1.2pt);
\fill[red!80!black] (12.0000,8.0000) circle (1.2pt);
\fill[red!80!black] (0.0000,8.0000) circle (1.2pt);
\fill[red!80!black] (12.0000,8.0000) circle (1.2pt);
\fill[red!80!black] (0.0000,8.0000) circle (1.2pt);
\fill[red!80!black] (12.0000,8.0000) circle (1.2pt);
\fill[red!80!black] (0.0000,8.0000) circle (1.2pt);
\fill[red!80!black] (12.0000,8.0000) circle (1.2pt);
\fill[red!80!black] (0.0000,8.0000) circle (1.2pt);
\fill[red!80!black] (12.0000,8.0000) circle (1.2pt);
\fill[red!80!black] (0.0000,8.0000) circle (1.2pt);
\fill[red!80!black] (12.0000,8.0000) circle (1.2pt);
\fill[red!80!black] (0.0000,8.0000) circle (1.2pt);
\fill[red!80!black] (12.0000,8.0000) circle (1.2pt);
\fill[red!80!black] (0.0000,8.0000) circle (1.2pt);
\fill[red!80!black] (12.0000,8.0000) circle (1.2pt);
\fill[red!80!black] (0.0000,8.0000) circle (1.2pt);
\fill[red!80!black] (12.0000,8.0000) circle (1.2pt);
\fill[red!80!black] (0.0000,8.0000) circle (1.2pt);
\fill[red!80!black] (12.0000,8.0000) circle (1.2pt);
\fill[red!80!black] (0.0000,8.0000) circle (1.2pt);
\fill[red!80!black] (12.0000,8.0000) circle (1.2pt);
\fill[red!80!black] (0.0000,8.0000) circle (1.2pt);
\fill[red!80!black] (12.0000,8.0000) circle (1.2pt);
\fill[red!80!black] (0.0000,8.0000) circle (1.2pt);
\fill[red!80!black] (12.0000,8.0000) circle (1.2pt);
\fill[red!80!black] (0.0000,8.0000) circle (1.2pt);
\fill[red!80!black] (12.0000,8.0000) circle (1.2pt);
\fill[red!80!black] (0.0000,8.0000) circle (1.2pt);
\fill[red!80!black] (12.0000,8.0000) circle (1.2pt);
\fill[red!80!black] (0.0000,8.0000) circle (1.2pt);
\fill[red!80!black] (12.0000,8.0000) circle (1.2pt);
\fill[red!80!black] (0.0000,8.0000) circle (1.2pt);
\fill[red!80!black] (12.0000,8.0000) circle (1.2pt);
\fill[red!80!black] (0.0000,8.0000) circle (1.2pt);
\fill[red!80!black] (12.0000,8.0000) circle (1.2pt);
\fill[red!80!black] (0.0000,8.0000) circle (1.2pt);
\fill[red!80!black] (12.0000,8.0000) circle (1.2pt);
\fill[red!80!black] (0.0000,8.0000) circle (1.2pt);
\fill[red!80!black] (12.0000,8.0000) circle (1.2pt);
\fill[red!80!black] (0.0000,8.0000) circle (1.2pt);
\fill[red!80!black] (12.0000,8.0000) circle (1.2pt);
\fill[red!80!black] (0.0000,8.0000) circle (1.2pt);
\fill[red!80!black] (12.0000,8.0000) circle (1.2pt);
\fill[red!80!black] (0.0000,8.0000) circle (1.2pt);
\fill[red!80!black] (12.0000,8.0000) circle (1.2pt);
\fill[red!80!black] (0.0000,8.0000) circle (1.2pt);
\fill[red!80!black] (12.0000,8.0000) circle (1.2pt);
\fill[red!80!black] (0.0000,8.0000) circle (1.2pt);
\fill[red!80!black] (12.0000,8.0000) circle (1.2pt);
\fill[red!80!black] (0.0000,8.0000) circle (1.2pt);
\fill[red!80!black] (12.0000,8.0000) circle (1.2pt);
\fill[red!80!black] (0.0000,8.0000) circle (1.2pt);
\fill[red!80!black] (12.0000,8.0000) circle (1.2pt);
\fill[red!80!black] (0.0000,8.0000) circle (1.2pt);
\fill[red!80!black] (12.0000,8.0000) circle (1.2pt);
\fill[red!80!black] (0.0000,8.0000) circle (1.2pt);
\fill[red!80!black] (12.0000,8.0000) circle (1.2pt);
\fill[red!80!black] (0.0000,8.0000) circle (1.2pt);
\fill[red!80!black] (12.0000,8.0000) circle (1.2pt);
\fill[red!80!black] (0.0000,8.0000) circle (1.2pt);
\fill[red!80!black] (12.0000,8.0000) circle (1.2pt);
\fill[red!80!black] (0.0000,8.0000) circle (1.2pt);
\fill[red!80!black] (12.0000,8.0000) circle (1.2pt);
\fill[red!80!black] (0.0000,8.0000) circle (1.2pt);
\fill[red!80!black] (12.0000,8.0000) circle (1.2pt);
\fill[red!80!black] (0.0000,8.0000) circle (1.2pt);
\fill[red!80!black] (12.0000,8.0000) circle (1.2pt);
\fill[red!80!black] (0.0000,8.0000) circle (1.2pt);
\fill[red!80!black] (12.0000,8.0000) circle (1.2pt);
\fill[red!80!black] (0.0000,8.0000) circle (1.2pt);
\fill[red!80!black] (12.0000,8.0000) circle (1.2pt);
\fill[red!80!black] (0.0000,8.0000) circle (1.2pt);
\fill[red!80!black] (12.0000,8.0000) circle (1.2pt);
\fill[red!80!black] (0.0000,8.0000) circle (1.2pt);
\fill[red!80!black] (12.0000,8.0000) circle (1.2pt);
\fill[red!80!black] (0.0000,8.0000) circle (1.2pt);
\fill[red!80!black] (12.0000,8.0000) circle (1.2pt);
\fill[red!80!black] (0.0000,8.0000) circle (1.2pt);
\fill[red!80!black] (12.0000,8.0000) circle (1.2pt);
\fill[red!80!black] (0.0000,8.0000) circle (1.2pt);
\fill[red!80!black] (12.0000,8.0000) circle (1.2pt);
\fill[red!80!black] (0.0000,8.0000) circle (1.2pt);
\fill[red!80!black] (12.0000,8.0000) circle (1.2pt);
\fill[red!80!black] (0.0000,8.0000) circle (1.2pt);
\fill[red!80!black] (12.0000,8.0000) circle (1.2pt);
\fill[red!80!black] (0.0000,8.0000) circle (1.2pt);
\fill[red!80!black] (12.0000,8.0000) circle (1.2pt);
\fill[red!80!black] (0.0000,8.0000) circle (1.2pt);
\fill[red!80!black] (12.0000,8.0000) circle (1.2pt);
\fill[red!80!black] (0.0000,8.0000) circle (1.2pt);
\fill[red!80!black] (12.0000,8.0000) circle (1.2pt);
\fill[red!80!black] (0.0000,8.0000) circle (1.2pt);
\fill[red!80!black] (12.0000,8.0000) circle (1.2pt);
\fill[red!80!black] (0.0000,8.0000) circle (1.2pt);
\fill[red!80!black] (12.0000,8.0000) circle (1.2pt);
\fill[red!80!black] (0.0000,8.0000) circle (1.2pt);
\fill[red!80!black] (12.0000,8.0000) circle (1.2pt);
\fill[red!80!black] (0.0000,8.0000) circle (1.2pt);
\fill[red!80!black] (12.0000,8.0000) circle (1.2pt);
\fill[red!80!black] (0.0000,8.0000) circle (1.2pt);
\fill[red!80!black] (12.0000,8.0000) circle (1.2pt);
\fill[red!80!black] (0.0000,8.0000) circle (1.2pt);
\fill[red!80!black] (12.0000,8.0000) circle (1.2pt);
\fill[red!80!black] (0.0000,8.0000) circle (1.2pt);
\fill[red!80!black] (12.0000,8.0000) circle (1.2pt);
\fill[red!80!black] (0.0000,8.0000) circle (1.2pt);
\fill[red!80!black] (12.0000,8.0000) circle (1.2pt);
\fill[red!80!black] (0.0000,8.0000) circle (1.2pt);
\fill[red!80!black] (12.0000,8.0000) circle (1.2pt);
\fill[red!80!black] (0.0000,8.0000) circle (1.2pt);
\fill[red!80!black] (12.0000,8.0000) circle (1.2pt);
\fill[red!80!black] (0.0000,8.0000) circle (1.2pt);
\fill[red!80!black] (12.0000,8.0000) circle (1.2pt);
\fill[red!80!black] (0.0000,8.0000) circle (1.2pt);
\fill[red!80!black] (12.0000,8.0000) circle (1.2pt);
\fill[red!80!black] (0.0000,8.0000) circle (1.2pt);
\fill[red!80!black] (12.0000,8.0000) circle (1.2pt);
\fill[red!80!black] (0.0000,8.0000) circle (1.2pt);
\fill[red!80!black] (12.0000,8.0000) circle (1.2pt);
\fill[red!80!black] (0.0000,8.0000) circle (1.2pt);
\fill[red!80!black] (12.0000,8.0000) circle (1.2pt);
\fill[red!80!black] (0.0000,8.0000) circle (1.2pt);
\fill[red!80!black] (12.0000,8.0000) circle (1.2pt);
\fill[red!80!black] (0.0000,8.0000) circle (1.2pt);
\fill[red!80!black] (12.0000,8.0000) circle (1.2pt);
\fill[red!80!black] (0.0000,8.0000) circle (1.2pt);
\fill[red!80!black] (12.0000,8.0000) circle (1.2pt);
\fill[red!80!black] (0.0000,8.0000) circle (1.2pt);
\fill[red!80!black] (12.0000,8.0000) circle (1.2pt);
\fill[red!80!black] (0.0000,8.0000) circle (1.2pt);
\fill[red!80!black] (12.0000,8.0000) circle (1.2pt);
\fill[red!80!black] (0.0000,8.0000) circle (1.2pt);
\fill[red!80!black] (12.0000,8.0000) circle (1.2pt);
\fill[red!80!black] (0.0000,8.0000) circle (1.2pt);
\fill[red!80!black] (12.0000,8.0000) circle (1.2pt);
\fill[red!80!black] (0.0000,8.0000) circle (1.2pt);
\fill[red!80!black] (12.0000,8.0000) circle (1.2pt);
\fill[red!80!black] (0.0000,8.0000) circle (1.2pt);
\fill[red!80!black] (12.0000,8.0000) circle (1.2pt);
\fill[red!80!black] (0.0000,8.0000) circle (1.2pt);
\fill[red!80!black] (12.0000,8.0000) circle (1.2pt);
\fill[red!80!black] (0.0000,8.0000) circle (1.2pt);
\fill[red!80!black] (12.0000,8.0000) circle (1.2pt);
\fill[red!80!black] (0.0000,8.0000) circle (1.2pt);
\fill[red!80!black] (12.0000,8.0000) circle (1.2pt);
\fill[red!80!black] (0.0000,8.0000) circle (1.2pt);
\fill[red!80!black] (12.0000,8.0000) circle (1.2pt);
\fill[red!80!black] (0.0000,8.0000) circle (1.2pt);
\fill[red!80!black] (12.0000,8.0000) circle (1.2pt);
\fill[red!80!black] (0.0000,8.0000) circle (1.2pt);
\fill[red!80!black] (12.0000,8.0000) circle (1.2pt);
\fill[red!80!black] (0.0000,8.0000) circle (1.2pt);
\fill[red!80!black] (12.0000,8.0000) circle (1.2pt);
\fill[red!80!black] (0.0000,8.0000) circle (1.2pt);
\fill[red!80!black] (12.0000,8.0000) circle (1.2pt);
\fill[red!80!black] (0.0000,8.0000) circle (1.2pt);
\fill[red!80!black] (12.0000,8.0000) circle (1.2pt);
\fill[red!80!black] (0.0000,8.0000) circle (1.2pt);
\fill[red!80!black] (12.0000,8.0000) circle (1.2pt);
\fill[red!80!black] (0.0000,8.0000) circle (1.2pt);
\fill[red!80!black] (12.0000,8.0000) circle (1.2pt);
\fill[red!80!black] (0.0000,8.0000) circle (1.2pt);
\fill[red!80!black] (12.0000,8.0000) circle (1.2pt);
\fill[red!80!black] (0.0000,8.0000) circle (1.2pt);
\fill[red!80!black] (12.0000,8.0000) circle (1.2pt);
\fill[red!80!black] (0.0000,8.0000) circle (1.2pt);
\fill[red!80!black] (12.0000,8.0000) circle (1.2pt);
\fill[red!80!black] (0.0000,8.0000) circle (1.2pt);
\fill[red!80!black] (12.0000,8.0000) circle (1.2pt);
\fill[red!80!black] (0.0000,8.0000) circle (1.2pt);
\fill[red!80!black] (12.0000,8.0000) circle (1.2pt);
\fill[red!80!black] (0.0000,8.0000) circle (1.2pt);
\fill[red!80!black] (12.0000,8.0000) circle (1.2pt);
\fill[red!80!black] (0.0000,8.0000) circle (1.2pt);
\fill[red!80!black] (12.0000,8.0000) circle (1.2pt);
\fill[red!80!black] (0.0000,8.0000) circle (1.2pt);
\fill[red!80!black] (12.0000,8.0000) circle (1.2pt);
\fill[red!80!black] (0.0000,8.0000) circle (1.2pt);
\fill[red!80!black] (12.0000,8.0000) circle (1.2pt);
\fill[red!80!black] (0.0000,8.0000) circle (1.2pt);
\fill[red!80!black] (12.0000,8.0000) circle (1.2pt);
\fill[red!80!black] (0.0000,8.0000) circle (1.2pt);
\fill[red!80!black] (12.0000,8.0000) circle (1.2pt);
\fill[red!80!black] (0.0000,8.0000) circle (1.2pt);
\fill[red!80!black] (12.0000,8.0000) circle (1.2pt);
\fill[red!80!black] (0.0000,8.0000) circle (1.2pt);
\fill[red!80!black] (12.0000,8.0000) circle (1.2pt);
\fill[red!80!black] (0.0000,8.0000) circle (1.2pt);
\fill[red!80!black] (12.0000,8.0000) circle (1.2pt);
\fill[red!80!black] (0.0000,8.0000) circle (1.2pt);
\fill[red!80!black] (12.0000,8.0000) circle (1.2pt);
\fill[red!80!black] (0.0000,8.0000) circle (1.2pt);
\fill[red!80!black] (12.0000,8.0000) circle (1.2pt);
\fill[red!80!black] (0.0000,8.0000) circle (1.2pt);
\fill[red!80!black] (12.0000,8.0000) circle (1.2pt);
\fill[red!80!black] (0.0000,8.0000) circle (1.2pt);
\fill[red!80!black] (12.0000,8.0000) circle (1.2pt);
\fill[red!80!black] (0.0000,8.0000) circle (1.2pt);
\fill[red!80!black] (12.0000,8.0000) circle (1.2pt);
\fill[red!80!black] (0.0000,8.0000) circle (1.2pt);
\fill[red!80!black] (12.0000,8.0000) circle (1.2pt);
\fill[red!80!black] (0.0000,8.0000) circle (1.2pt);
\fill[red!80!black] (12.0000,8.0000) circle (1.2pt);
\fill[red!80!black] (0.0000,8.0000) circle (1.2pt);
\fill[red!80!black] (12.0000,8.0000) circle (1.2pt);
\fill[red!80!black] (0.0000,8.0000) circle (1.2pt);
\fill[red!80!black] (12.0000,8.0000) circle (1.2pt);
\fill[red!80!black] (0.0000,8.0000) circle (1.2pt);
\fill[red!80!black] (12.0000,8.0000) circle (1.2pt);
\fill[red!80!black] (0.0000,8.0000) circle (1.2pt);
\fill[red!80!black] (12.0000,8.0000) circle (1.2pt);
\fill[red!80!black] (0.0000,8.0000) circle (1.2pt);
\fill[red!80!black] (12.0000,8.0000) circle (1.2pt);
\fill[red!80!black] (0.0000,8.0000) circle (1.2pt);
\fill[red!80!black] (12.0000,8.0000) circle (1.2pt);
\fill[red!80!black] (0.0000,8.0000) circle (1.2pt);
\fill[red!80!black] (12.0000,8.0000) circle (1.2pt);
\fill[red!80!black] (0.0000,8.0000) circle (1.2pt);
\fill[red!80!black] (12.0000,8.0000) circle (1.2pt);
\fill[red!80!black] (0.0000,8.0000) circle (1.2pt);
\fill[red!80!black] (12.0000,8.0000) circle (1.2pt);
\fill[red!80!black] (0.0000,8.0000) circle (1.2pt);
\fill[red!80!black] (12.0000,8.0000) circle (1.2pt);
\fill[red!80!black] (0.0000,8.0000) circle (1.2pt);
\fill[red!80!black] (12.0000,8.0000) circle (1.2pt);
\fill[red!80!black] (0.0000,8.0000) circle (1.2pt);
\fill[red!80!black] (12.0000,8.0000) circle (1.2pt);
\fill[red!80!black] (0.0000,8.0000) circle (1.2pt);
\fill[red!80!black] (12.0000,8.0000) circle (1.2pt);
\fill[red!80!black] (0.0000,8.0000) circle (1.2pt);
\fill[red!80!black] (12.0000,8.0000) circle (1.2pt);
\fill[red!80!black] (0.0000,8.0000) circle (1.2pt);
\fill[red!80!black] (12.0000,8.0000) circle (1.2pt);
\fill[red!80!black] (0.0000,8.0000) circle (1.2pt);
\fill[red!80!black] (12.0000,8.0000) circle (1.2pt);
\fill[red!80!black] (0.0000,8.0000) circle (1.2pt);
\fill[red!80!black] (12.0000,8.0000) circle (1.2pt);
\fill[red!80!black] (0.0000,8.0000) circle (1.2pt);
\fill[red!80!black] (12.0000,8.0000) circle (1.2pt);
\fill[red!80!black] (0.0000,8.0000) circle (1.2pt);
\fill[red!80!black] (12.0000,8.0000) circle (1.2pt);
\fill[red!80!black] (0.0000,8.0000) circle (1.2pt);
\fill[red!80!black] (12.0000,8.0000) circle (1.2pt);
\fill[red!80!black] (0.0000,8.0000) circle (1.2pt);
\fill[red!80!black] (12.0000,8.0000) circle (1.2pt);
\fill[red!80!black] (0.0000,8.0000) circle (1.2pt);
\fill[red!80!black] (12.0000,8.0000) circle (1.2pt);
\fill[red!80!black] (0.0000,8.0000) circle (1.2pt);
\fill[red!80!black] (12.0000,8.0000) circle (1.2pt);
\fill[red!80!black] (0.0000,8.0000) circle (1.2pt);
\fill[red!80!black] (12.0000,8.0000) circle (1.2pt);
\fill[red!80!black] (0.0000,8.0000) circle (1.2pt);
\fill[red!80!black] (12.0000,8.0000) circle (1.2pt);
\fill[red!80!black] (0.0000,8.0000) circle (1.2pt);
\fill[red!80!black] (12.0000,8.0000) circle (1.2pt);
\fill[red!80!black] (0.0000,8.0000) circle (1.2pt);
\fill[red!80!black] (12.0000,8.0000) circle (1.2pt);
\fill[red!80!black] (0.0000,8.0000) circle (1.2pt);
\fill[red!80!black] (12.0000,8.0000) circle (1.2pt);
\fill[red!80!black] (0.0000,8.0000) circle (1.2pt);
\fill[red!80!black] (12.0000,8.0000) circle (1.2pt);
\fill[red!80!black] (0.0000,8.0000) circle (1.2pt);
\fill[red!80!black] (12.0000,8.0000) circle (1.2pt);
\fill[red!80!black] (0.0000,8.0000) circle (1.2pt);
\fill[red!80!black] (12.0000,8.0000) circle (1.2pt);
\fill[red!80!black] (0.0000,8.0000) circle (1.2pt);
\fill[red!80!black] (12.0000,8.0000) circle (1.2pt);
\fill[red!80!black] (0.0000,8.0000) circle (1.2pt);
\fill[red!80!black] (12.0000,8.0000) circle (1.2pt);
\fill[red!80!black] (0.0000,8.0000) circle (1.2pt);
\fill[red!80!black] (12.0000,8.0000) circle (1.2pt);
\fill[red!80!black] (0.0000,8.0000) circle (1.2pt);
\fill[red!80!black] (12.0000,8.0000) circle (1.2pt);
\fill[red!80!black] (0.0000,8.0000) circle (1.2pt);
\fill[red!80!black] (12.0000,8.0000) circle (1.2pt);
\fill[red!80!black] (0.0000,8.0000) circle (1.2pt);
\fill[red!80!black] (12.0000,8.0000) circle (1.2pt);
\fill[red!80!black] (0.0000,8.0000) circle (1.2pt);
\fill[red!80!black] (12.0000,8.0000) circle (1.2pt);
\fill[red!80!black] (0.0000,8.0000) circle (1.2pt);
\fill[red!80!black] (12.0000,8.0000) circle (1.2pt);
\fill[red!80!black] (0.0000,8.0000) circle (1.2pt);
\fill[red!80!black] (12.0000,8.0000) circle (1.2pt);
\fill[red!80!black] (0.0000,8.0000) circle (1.2pt);
\fill[red!80!black] (12.0000,8.0000) circle (1.2pt);
\fill[red!80!black] (0.0000,8.0000) circle (1.2pt);
\fill[red!80!black] (12.0000,8.0000) circle (1.2pt);
\fill[red!80!black] (0.0000,8.0000) circle (1.2pt);
\fill[red!80!black] (12.0000,8.0000) circle (1.2pt);
\fill[red!80!black] (0.0000,8.0000) circle (1.2pt);
\fill[red!80!black] (12.0000,8.0000) circle (1.2pt);
\fill[red!80!black] (0.0000,8.0000) circle (1.2pt);
\fill[red!80!black] (12.0000,8.0000) circle (1.2pt);
\fill[red!80!black] (0.0000,8.0000) circle (1.2pt);
\fill[red!80!black] (12.0000,8.0000) circle (1.2pt);
\fill[red!80!black] (0.0000,8.0000) circle (1.2pt);
\fill[red!80!black] (12.0000,8.0000) circle (1.2pt);
\fill[red!80!black] (0.0000,8.0000) circle (1.2pt);
\fill[red!80!black] (12.0000,8.0000) circle (1.2pt);
\fill[red!80!black] (0.0000,8.0000) circle (1.2pt);
\fill[red!80!black] (12.0000,8.0000) circle (1.2pt);
\fill[red!80!black] (0.0000,8.0000) circle (1.2pt);
\fill[red!80!black] (12.0000,8.0000) circle (1.2pt);
\fill[red!80!black] (0.0000,8.0000) circle (1.2pt);
\fill[red!80!black] (12.0000,8.0000) circle (1.2pt);
\fill[red!80!black] (0.0000,8.0000) circle (1.2pt);
\fill[red!80!black] (12.0000,8.0000) circle (1.2pt);
\fill[red!80!black] (0.0000,8.0000) circle (1.2pt);
\fill[red!80!black] (12.0000,8.0000) circle (1.2pt);
\fill[red!80!black] (0.0000,8.0000) circle (1.2pt);
\fill[red!80!black] (12.0000,8.0000) circle (1.2pt);
\fill[red!80!black] (0.0000,8.0000) circle (1.2pt);
\fill[red!80!black] (12.0000,8.0000) circle (1.2pt);
\fill[red!80!black] (0.0000,8.0000) circle (1.2pt);
\fill[red!80!black] (12.0000,8.0000) circle (1.2pt);
\fill[red!80!black] (0.0000,8.0000) circle (1.2pt);
\fill[red!80!black] (12.0000,8.0000) circle (1.2pt);
\fill[red!80!black] (0.0000,8.0000) circle (1.2pt);
\fill[red!80!black] (12.0000,8.0000) circle (1.2pt);
\fill[red!80!black] (0.0000,8.0000) circle (1.2pt);
\fill[red!80!black] (12.0000,8.0000) circle (1.2pt);
\fill[red!80!black] (0.0000,8.0000) circle (1.2pt);
\fill[red!80!black] (12.0000,8.0000) circle (1.2pt);
\fill[red!80!black] (0.0000,8.0000) circle (1.2pt);
\fill[red!80!black] (12.0000,8.0000) circle (1.2pt);
\fill[red!80!black] (0.0000,8.0000) circle (1.2pt);
\fill[red!80!black] (12.0000,8.0000) circle (1.2pt);
\fill[red!80!black] (0.0000,8.0000) circle (1.2pt);
\fill[red!80!black] (12.0000,8.0000) circle (1.2pt);
\fill[red!80!black] (0.0000,8.0000) circle (1.2pt);
\fill[red!80!black] (12.0000,8.0000) circle (1.2pt);
\fill[red!80!black] (0.0000,8.0000) circle (1.2pt);
\fill[red!80!black] (12.0000,8.0000) circle (1.2pt);
\fill[red!80!black] (0.0000,8.0000) circle (1.2pt);
\fill[red!80!black] (12.0000,8.0000) circle (1.2pt);
\fill[red!80!black] (0.0000,8.0000) circle (1.2pt);
\fill[red!80!black] (12.0000,8.0000) circle (1.2pt);
\fill[red!80!black] (0.0000,8.0000) circle (1.2pt);
\fill[red!80!black] (12.0000,8.0000) circle (1.2pt);
\fill[red!80!black] (0.0000,8.0000) circle (1.2pt);
\fill[red!80!black] (12.0000,8.0000) circle (1.2pt);
\fill[red!80!black] (0.0000,8.0000) circle (1.2pt);
\fill[red!80!black] (12.0000,8.0000) circle (1.2pt);
\fill[red!80!black] (0.0000,8.0000) circle (1.2pt);
\fill[red!80!black] (12.0000,8.0000) circle (1.2pt);
\fill[red!80!black] (0.0000,8.0000) circle (1.2pt);
\fill[red!80!black] (12.0000,8.0000) circle (1.2pt);
\fill[red!80!black] (0.0000,8.0000) circle (1.2pt);
\fill[red!80!black] (12.0000,8.0000) circle (1.2pt);
\fill[red!80!black] (0.0000,8.0000) circle (1.2pt);
\fill[red!80!black] (12.0000,8.0000) circle (1.2pt);
\fill[red!80!black] (0.0000,8.0000) circle (1.2pt);
\fill[red!80!black] (12.0000,8.0000) circle (1.2pt);
\fill[red!80!black] (0.0000,8.0000) circle (1.2pt);
\fill[red!80!black] (12.0000,8.0000) circle (1.2pt);
\fill[red!80!black] (0.0000,8.0000) circle (1.2pt);
\fill[red!80!black] (12.0000,8.0000) circle (1.2pt);
\fill[red!80!black] (0.0000,8.0000) circle (1.2pt);
\fill[red!80!black] (12.0000,8.0000) circle (1.2pt);
\fill[red!80!black] (0.0000,8.0000) circle (1.2pt);
\fill[red!80!black] (12.0000,8.0000) circle (1.2pt);
\fill[red!80!black] (0.0000,8.0000) circle (1.2pt);
\fill[red!80!black] (12.0000,8.0000) circle (1.2pt);
\fill[red!80!black] (0.0000,8.0000) circle (1.2pt);
\fill[red!80!black] (12.0000,8.0000) circle (1.2pt);
\fill[red!80!black] (0.0000,8.0000) circle (1.2pt);
\draw[->, thick] (0,0) -- (12.6000,0) node[right] {$x$};
\draw[->, thick] (0,0) -- (0,8.6000) node[above] {$y$};
\draw (0.0000,0) -- (0.0000,-0.12) node[below] {\small -5.00};
\draw (0,0.0000) -- (-0.12,0.0000) node[left] {\small -5.00};
\draw (2.4000,0) -- (2.4000,-0.12) node[below] {\small -3.00};
\draw (0,1.6000) -- (-0.12,1.6000) node[left] {\small -3.00};
\draw (4.8000,0) -- (4.8000,-0.12) node[below] {\small -1.00};
\draw (0,3.2000) -- (-0.12,3.2000) node[left] {\small -1.00};
\draw (7.2000,0) -- (7.2000,-0.12) node[below] {\small 1.00};
\draw (0,4.8000) -- (-0.12,4.8000) node[left] {\small 1.00};
\draw (9.6000,0) -- (9.6000,-0.12) node[below] {\small 3.00};
\draw (0,6.4000) -- (-0.12,6.4000) node[left] {\small 3.00};
\draw (12.0000,0) -- (12.0000,-0.12) node[below] {\small 5.00};
\draw (0,8.0000) -- (-0.12,8.0000) node[left] {\small 5.00};
\draw[fill=white, draw=black] (7.6,6.7) rectangle (11.8,7.9);
\draw[blue!80!black, line width=1.0pt] (7.9,7.45) -- (8.9,7.45);
\node[anchor=west] at (9.1,7.45) {\small gradDest};
\draw[red!80!black, line width=1.0pt] (7.9,7.05) -- (8.9,7.05);
\node[anchor=west] at (9.1,7.05) {\small fastGradDest};
\end{tikzpicture}

\caption{Функция Розенброка ($b=100$): линии уровня и траектории.}
\label{fig:rosenbrock}
\end{figure}

\section{Линии уровня и траектории}
На рисунках ниже показаны линии уровня и траектории обоих методов для старта $(-1, -1)$.

\begin{figure}[H]
\centering
\begin{tikzpicture}[scale=0.9]
\draw[fill=orange!3, draw=black] (0,0) rectangle (12.0000,8.0000);
\draw[gray!35, line width=0.15pt] (10.3729,3.6479) -- (10.3680,3.6610);
\draw[gray!35, line width=0.15pt] (10.3680,3.6610) -- (10.3343,3.7966);
\draw[gray!35, line width=0.15pt] (10.3343,3.7966) -- (10.3178,3.9322);
\draw[gray!35, line width=0.15pt] (10.3178,3.9322) -- (10.3166,4.0678);
\draw[gray!35, line width=0.15pt] (10.3166,4.0678) -- (10.3294,4.2034);
\draw[gray!35, line width=0.15pt] (10.3294,4.2034) -- (10.3547,4.3390);
\draw[gray!35, line width=0.15pt] (10.3547,4.3390) -- (10.3729,4.4056);
\draw[gray!35, line width=0.15pt] (10.5763,3.2781) -- (10.4953,3.3898);
\draw[gray!35, line width=0.15pt] (10.4953,3.3898) -- (10.4208,3.5254);
\draw[gray!35, line width=0.15pt] (10.3729,3.6479) -- (10.4208,3.5254);
\draw[gray!35, line width=0.15pt] (10.3913,4.4746) -- (10.3729,4.4056);
\draw[gray!35, line width=0.15pt] (10.3913,4.4746) -- (10.4382,4.6102);
\draw[gray!35, line width=0.15pt] (10.4382,4.6102) -- (10.4944,4.7458);
\draw[gray!35, line width=0.15pt] (10.4944,4.7458) -- (10.5589,4.8814);
\draw[gray!35, line width=0.15pt] (10.5589,4.8814) -- (10.5763,4.9138);
\draw[gray!35, line width=0.15pt] (10.7797,3.0685) -- (10.7224,3.1186);
\draw[gray!35, line width=0.15pt] (10.7224,3.1186) -- (10.5946,3.2542);
\draw[gray!35, line width=0.15pt] (10.5763,3.2781) -- (10.5946,3.2542);
\draw[gray!35, line width=0.15pt] (10.6307,5.0169) -- (10.5763,4.9138);
\draw[gray!35, line width=0.15pt] (10.6307,5.0169) -- (10.7090,5.1525);
\draw[gray!35, line width=0.15pt] (10.7090,5.1525) -- (10.7797,5.2667);
\draw[gray!35, line width=0.15pt] (10.9831,2.9137) -- (10.8837,2.9831);
\draw[gray!35, line width=0.15pt] (10.7797,3.0685) -- (10.8837,2.9831);
\draw[gray!35, line width=0.15pt] (10.7928,5.2881) -- (10.7797,5.2667);
\draw[gray!35, line width=0.15pt] (10.7928,5.2881) -- (10.8815,5.4237);
\draw[gray!35, line width=0.15pt] (10.8815,5.4237) -- (10.9741,5.5593);
\draw[gray!35, line width=0.15pt] (10.9741,5.5593) -- (10.9831,5.5718);
\draw[gray!35, line width=0.15pt] (11.1864,2.7894) -- (11.0846,2.8475);
\draw[gray!35, line width=0.15pt] (10.9831,2.9137) -- (11.0846,2.8475);
\draw[gray!35, line width=0.15pt] (11.0702,5.6949) -- (10.9831,5.5718);
\draw[gray!35, line width=0.15pt] (11.0702,5.6949) -- (11.1687,5.8305);
\draw[gray!35, line width=0.15pt] (11.1687,5.8305) -- (11.1864,5.8544);
\draw[gray!35, line width=0.15pt] (11.3898,2.6848) -- (11.3326,2.7119);
\draw[gray!35, line width=0.15pt] (11.1864,2.7894) -- (11.3326,2.7119);
\draw[gray!35, line width=0.15pt] (11.2693,5.9661) -- (11.1864,5.8544);
\draw[gray!35, line width=0.15pt] (11.2693,5.9661) -- (11.3712,6.1017);
\draw[gray!35, line width=0.15pt] (11.3712,6.1017) -- (11.3898,6.1263);
\draw[gray!35, line width=0.15pt] (11.3898,2.6848) -- (11.5932,2.5945);
\draw[gray!35, line width=0.15pt] (11.4740,6.2373) -- (11.3898,6.1263);
\draw[gray!35, line width=0.15pt] (11.4740,6.2373) -- (11.5770,6.3729);
\draw[gray!35, line width=0.15pt] (11.5770,6.3729) -- (11.5932,6.3942);
\draw[gray!35, line width=0.15pt] (11.7966,2.5150) -- (11.6367,2.5763);
\draw[gray!35, line width=0.15pt] (11.5932,2.5945) -- (11.6367,2.5763);
\draw[gray!35, line width=0.15pt] (11.6803,6.5085) -- (11.5932,6.3942);
\draw[gray!35, line width=0.15pt] (11.6803,6.5085) -- (11.7828,6.6441);
\draw[gray!35, line width=0.15pt] (11.7828,6.6441) -- (11.7966,6.6625);
\draw[gray!35, line width=0.15pt] (11.7966,2.5150) -- (12.0000,2.4429);
\draw[gray!35, line width=0.15pt] (11.8850,6.7797) -- (11.7966,6.6625);
\draw[gray!35, line width=0.15pt] (11.8850,6.7797) -- (11.9858,6.9153);
\draw[gray!35, line width=0.15pt] (11.9858,6.9153) -- (12.0000,6.9347);
\draw[gray!45, line width=0.15pt] (8.3390,2.9712) -- (8.3380,2.9831);
\draw[gray!45, line width=0.15pt] (8.3380,2.9831) -- (8.3368,3.1186);
\draw[gray!45, line width=0.15pt] (8.3368,3.1186) -- (8.3390,3.1517);
\draw[gray!45, line width=0.15pt] (8.5424,2.3067) -- (8.4680,2.4407);
\draw[gray!45, line width=0.15pt] (8.4680,2.4407) -- (8.4130,2.5763);
\draw[gray!45, line width=0.15pt] (8.4130,2.5763) -- (8.3745,2.7119);
\draw[gray!45, line width=0.15pt] (8.3745,2.7119) -- (8.3501,2.8475);
\draw[gray!45, line width=0.15pt] (8.3390,2.9712) -- (8.3501,2.8475);
\draw[gray!45, line width=0.15pt] (8.3455,3.2542) -- (8.3390,3.1517);
\draw[gray!45, line width=0.15pt] (8.3455,3.2542) -- (8.3634,3.3898);
\draw[gray!45, line width=0.15pt] (8.3634,3.3898) -- (8.3900,3.5254);
\draw[gray!45, line width=0.15pt] (8.3900,3.5254) -- (8.4247,3.6610);
\draw[gray!45, line width=0.15pt] (8.4247,3.6610) -- (8.4672,3.7966);
\draw[gray!45, line width=0.15pt] (8.4672,3.7966) -- (8.5171,3.9322);
\draw[gray!45, line width=0.15pt] (8.5171,3.9322) -- (8.5424,3.9912);
\draw[gray!45, line width=0.15pt] (8.7458,2.0690) -- (8.6465,2.1695);
\draw[gray!45, line width=0.15pt] (8.6465,2.1695) -- (8.5434,2.3051);
\draw[gray!45, line width=0.15pt] (8.5424,2.3067) -- (8.5434,2.3051);
\draw[gray!45, line width=0.15pt] (8.5742,4.0678) -- (8.5424,3.9912);
\draw[gray!45, line width=0.15pt] (8.5742,4.0678) -- (8.6380,4.2034);
\draw[gray!45, line width=0.15pt] (8.6380,4.2034) -- (8.7082,4.3390);
\draw[gray!45, line width=0.15pt] (8.7082,4.3390) -- (8.7458,4.4049);
\draw[gray!45, line width=0.15pt] (8.9492,1.9171) -- (8.7859,2.0339);
\draw[gray!45, line width=0.15pt] (8.7458,2.0690) -- (8.7859,2.0339);
\draw[gray!45, line width=0.15pt] (8.7845,4.4746) -- (8.7458,4.4049);
\draw[gray!45, line width=0.15pt] (8.7845,4.4746) -- (8.8662,4.6102);
\draw[gray!45, line width=0.15pt] (8.8662,4.6102) -- (8.9492,4.7393);
\draw[gray!45, line width=0.15pt] (9.1525,1.8099) -- (8.9805,1.8983);
\draw[gray!45, line width=0.15pt] (8.9492,1.9171) -- (8.9805,1.8983);
\draw[gray!45, line width=0.15pt] (8.9532,4.7458) -- (8.9492,4.7393);
\draw[gray!45, line width=0.15pt] (8.9532,4.7458) -- (9.0447,4.8814);
\draw[gray!45, line width=0.15pt] (9.0447,4.8814) -- (9.1404,5.0169);
\draw[gray!45, line width=0.15pt] (9.1404,5.0169) -- (9.1525,5.0333);
\draw[gray!45, line width=0.15pt] (9.3559,1.7295) -- (9.2655,1.7627);
\draw[gray!45, line width=0.15pt] (9.1525,1.8099) -- (9.2655,1.7627);
\draw[gray!45, line width=0.15pt] (9.2396,5.1525) -- (9.1525,5.0333);
\draw[gray!45, line width=0.15pt] (9.2396,5.1525) -- (9.3419,5.2881);
\draw[gray!45, line width=0.15pt] (9.3419,5.2881) -- (9.3559,5.3061);
\draw[gray!45, line width=0.15pt] (9.3559,1.7295) -- (9.5593,1.6672);
\draw[gray!45, line width=0.15pt] (9.4465,5.4237) -- (9.3559,5.3061);
\draw[gray!45, line width=0.15pt] (9.4465,5.4237) -- (9.5532,5.5593);
\draw[gray!45, line width=0.15pt] (9.5532,5.5593) -- (9.5593,5.5669);
\draw[gray!45, line width=0.15pt] (9.7627,1.6175) -- (9.7201,1.6271);
\draw[gray!45, line width=0.15pt] (9.5593,1.6672) -- (9.7201,1.6271);
\draw[gray!45, line width=0.15pt] (9.6612,5.6949) -- (9.5593,5.5669);
\draw[gray!45, line width=0.15pt] (9.6612,5.6949) -- (9.7627,5.8213);
\draw[gray!45, line width=0.15pt] (9.7627,1.6175) -- (9.9661,1.5768);
\draw[gray!45, line width=0.15pt] (9.7700,5.8305) -- (9.7627,5.8213);
\draw[gray!45, line width=0.15pt] (9.7700,5.8305) -- (9.8791,5.9661);
\draw[gray!45, line width=0.15pt] (9.8791,5.9661) -- (9.9661,6.0741);
\draw[gray!45, line width=0.15pt] (9.9661,1.5768) -- (10.1695,1.5428);
\draw[gray!45, line width=0.15pt] (9.9881,6.1017) -- (9.9661,6.0741);
\draw[gray!45, line width=0.15pt] (9.9881,6.1017) -- (10.0965,6.2373);
\draw[gray!45, line width=0.15pt] (10.0965,6.2373) -- (10.1695,6.3292);
\draw[gray!45, line width=0.15pt] (10.1695,1.5428) -- (10.3729,1.5138);
\draw[gray!45, line width=0.15pt] (10.2040,6.3729) -- (10.1695,6.3292);
\draw[gray!45, line width=0.15pt] (10.2040,6.3729) -- (10.3101,6.5085);
\draw[gray!45, line width=0.15pt] (10.3101,6.5085) -- (10.3729,6.5897);
\draw[gray!45, line width=0.15pt] (10.5763,1.4885) -- (10.5507,1.4915);
\draw[gray!45, line width=0.15pt] (10.3729,1.5138) -- (10.5507,1.4915);
\draw[gray!45, line width=0.15pt] (10.4147,6.6441) -- (10.3729,6.5897);
\draw[gray!45, line width=0.15pt] (10.4147,6.6441) -- (10.5173,6.7797);
\draw[gray!45, line width=0.15pt] (10.5173,6.7797) -- (10.5763,6.8590);
\draw[gray!45, line width=0.15pt] (10.5763,1.4885) -- (10.7797,1.4658);
\draw[gray!45, line width=0.15pt] (10.6180,6.9153) -- (10.5763,6.8590);
\draw[gray!45, line width=0.15pt] (10.6180,6.9153) -- (10.7163,7.0508);
\draw[gray!45, line width=0.15pt] (10.7163,7.0508) -- (10.7797,7.1403);
\draw[gray!45, line width=0.15pt] (10.7797,1.4658) -- (10.9831,1.4451);
\draw[gray!45, line width=0.15pt] (10.8123,7.1864) -- (10.7797,7.1403);
\draw[gray!45, line width=0.15pt] (10.8123,7.1864) -- (10.9060,7.3220);
\draw[gray!45, line width=0.15pt] (10.9060,7.3220) -- (10.9831,7.4369);
\draw[gray!45, line width=0.15pt] (10.9831,1.4451) -- (11.1864,1.4259);
\draw[gray!45, line width=0.15pt] (10.9970,7.4576) -- (10.9831,7.4369);
\draw[gray!45, line width=0.15pt] (10.9970,7.4576) -- (11.0857,7.5932);
\draw[gray!45, line width=0.15pt] (11.0857,7.5932) -- (11.1718,7.7288);
\draw[gray!45, line width=0.15pt] (11.1718,7.7288) -- (11.1864,7.7526);
\draw[gray!45, line width=0.15pt] (11.1864,1.4259) -- (11.3898,1.4077);
\draw[gray!45, line width=0.15pt] (11.2556,7.8644) -- (11.1864,7.7526);
\draw[gray!45, line width=0.15pt] (11.2556,7.8644) -- (11.3370,8.0000);
\draw[gray!45, line width=0.15pt] (11.3898,1.4077) -- (11.5932,1.3901);
\draw[gray!45, line width=0.15pt] (11.5932,1.3901) -- (11.7966,1.3728);
\draw[gray!45, line width=0.15pt] (12.0000,1.3555) -- (11.9950,1.3559);
\draw[gray!45, line width=0.15pt] (11.7966,1.3728) -- (11.9950,1.3559);
\draw[gray!55, line width=0.15pt] (3.0191,0.0000) -- (3.0508,0.0346);
\draw[gray!55, line width=0.15pt] (3.1474,0.1356) -- (3.0508,0.0346);
\draw[gray!55, line width=0.15pt] (3.1474,0.1356) -- (3.2542,0.2456);
\draw[gray!55, line width=0.15pt] (3.2801,0.2712) -- (3.2542,0.2456);
\draw[gray!55, line width=0.15pt] (3.2801,0.2712) -- (3.4193,0.4068);
\draw[gray!55, line width=0.15pt] (3.4193,0.4068) -- (3.4576,0.4436);
\draw[gray!55, line width=0.15pt] (3.5639,0.5424) -- (3.4576,0.4436);
\draw[gray!55, line width=0.15pt] (3.5639,0.5424) -- (3.6610,0.6323);
\draw[gray!55, line width=0.15pt] (3.7119,0.6780) -- (3.6610,0.6323);
\draw[gray!55, line width=0.15pt] (3.7119,0.6780) -- (3.8626,0.8136);
\draw[gray!55, line width=0.15pt] (3.8626,0.8136) -- (3.8644,0.8151);
\draw[gray!55, line width=0.15pt] (4.0180,0.9492) -- (3.8644,0.8151);
\draw[gray!55, line width=0.15pt] (4.0180,0.9492) -- (4.0678,0.9931);
\draw[gray!55, line width=0.15pt] (4.1740,1.0847) -- (4.0678,0.9931);
\draw[gray!55, line width=0.15pt] (4.1740,1.0847) -- (4.2712,1.1699);
\draw[gray!55, line width=0.15pt] (4.3296,1.2203) -- (4.2712,1.1699);
\draw[gray!55, line width=0.15pt] (4.3296,1.2203) -- (4.4746,1.3480);
\draw[gray!55, line width=0.15pt] (4.4837,1.3559) -- (4.4746,1.3480);
\draw[gray!55, line width=0.15pt] (4.4837,1.3559) -- (4.6369,1.4915);
\draw[gray!55, line width=0.15pt] (4.6369,1.4915) -- (4.6780,1.5288);
\draw[gray!55, line width=0.15pt] (4.7866,1.6271) -- (4.6780,1.5288);
\draw[gray!55, line width=0.15pt] (4.7866,1.6271) -- (4.8814,1.7157);
\draw[gray!55, line width=0.15pt] (4.9315,1.7627) -- (4.8814,1.7157);
\draw[gray!55, line width=0.15pt] (4.9315,1.7627) -- (5.0714,1.8983);
\draw[gray!55, line width=0.15pt] (5.0714,1.8983) -- (5.0847,1.9116);
\draw[gray!55, line width=0.15pt] (5.2071,2.0339) -- (5.0847,1.9116);
\draw[gray!55, line width=0.15pt] (5.2071,2.0339) -- (5.2881,2.1184);
\draw[gray!55, line width=0.15pt] (5.3365,2.1695) -- (5.2881,2.1184);
\draw[gray!55, line width=0.15pt] (5.3365,2.1695) -- (5.4606,2.3051);
\draw[gray!55, line width=0.15pt] (5.4606,2.3051) -- (5.4915,2.3398);
\draw[gray!55, line width=0.15pt] (5.5800,2.4407) -- (5.4915,2.3398);
\draw[gray!55, line width=0.15pt] (5.5800,2.4407) -- (5.6940,2.5763);
\draw[gray!55, line width=0.15pt] (5.6940,2.5763) -- (5.6949,2.5774);
\draw[gray!55, line width=0.15pt] (5.8049,2.7119) -- (5.6949,2.5774);
\draw[gray!55, line width=0.15pt] (5.8049,2.7119) -- (5.8983,2.8303);
\draw[gray!55, line width=0.15pt] (5.9115,2.8475) -- (5.8983,2.8303);
\draw[gray!55, line width=0.15pt] (5.9115,2.8475) -- (6.0159,2.9831);
\draw[gray!55, line width=0.15pt] (6.0159,2.9831) -- (6.1017,3.0968);
\draw[gray!55, line width=0.15pt] (6.1178,3.1186) -- (6.1017,3.0968);
\draw[gray!55, line width=0.15pt] (6.1178,3.1186) -- (6.2186,3.2542);
\draw[gray!55, line width=0.15pt] (6.2186,3.2542) -- (6.3051,3.3713);
\draw[gray!55, line width=0.15pt] (6.3184,3.3898) -- (6.3051,3.3713);
\draw[gray!55, line width=0.15pt] (6.3184,3.3898) -- (6.4184,3.5254);
\draw[gray!55, line width=0.15pt] (6.4184,3.5254) -- (6.5085,3.6471);
\draw[gray!55, line width=0.15pt] (6.5185,3.6610) -- (6.5085,3.6471);
\draw[gray!55, line width=0.15pt] (6.5185,3.6610) -- (6.6198,3.7966);
\draw[gray!55, line width=0.15pt] (6.6198,3.7966) -- (6.7119,3.9181);
\draw[gray!55, line width=0.15pt] (6.7223,3.9322) -- (6.7119,3.9181);
\draw[gray!55, line width=0.15pt] (6.7223,3.9322) -- (6.8266,4.0678);
\draw[gray!55, line width=0.15pt] (6.8266,4.0678) -- (6.9153,4.1806);
\draw[gray!55, line width=0.15pt] (6.9328,4.2034) -- (6.9153,4.1806);
\draw[gray!55, line width=0.15pt] (6.9328,4.2034) -- (7.0412,4.3390);
\draw[gray!55, line width=0.15pt] (7.0412,4.3390) -- (7.1186,4.4334);
\draw[gray!55, line width=0.15pt] (7.1518,4.4746) -- (7.1186,4.4334);
\draw[gray!55, line width=0.15pt] (7.1518,4.4746) -- (7.2647,4.6102);
\draw[gray!55, line width=0.15pt] (7.2647,4.6102) -- (7.3220,4.6771);
\draw[gray!55, line width=0.15pt] (7.3798,4.7458) -- (7.3220,4.6771);
\draw[gray!55, line width=0.15pt] (7.3798,4.7458) -- (7.4969,4.8814);
\draw[gray!55, line width=0.15pt] (7.4969,4.8814) -- (7.5254,4.9135);
\draw[gray!55, line width=0.15pt] (7.6158,5.0169) -- (7.5254,4.9135);
\draw[gray!55, line width=0.15pt] (7.6158,5.0169) -- (7.7288,5.1439);
\draw[gray!55, line width=0.15pt] (7.7364,5.1525) -- (7.7288,5.1439);
\draw[gray!55, line width=0.15pt] (7.7364,5.1525) -- (7.8579,5.2881);
\draw[gray!55, line width=0.15pt] (7.8579,5.2881) -- (7.9322,5.3699);
\draw[gray!55, line width=0.15pt] (7.9804,5.4237) -- (7.9322,5.3699);
\draw[gray!55, line width=0.15pt] (7.9804,5.4237) -- (8.1034,5.5593);
\draw[gray!55, line width=0.15pt] (8.1034,5.5593) -- (8.1356,5.5945);
\draw[gray!55, line width=0.15pt] (8.2262,5.6949) -- (8.1356,5.5945);
\draw[gray!55, line width=0.15pt] (8.2262,5.6949) -- (8.3390,5.8194);
\draw[gray!55, line width=0.15pt] (8.3489,5.8305) -- (8.3390,5.8194);
\draw[gray!55, line width=0.15pt] (8.3489,5.8305) -- (8.4704,5.9661);
\draw[gray!55, line width=0.15pt] (8.4704,5.9661) -- (8.5424,6.0466);
\draw[gray!55, line width=0.15pt] (8.5909,6.1017) -- (8.5424,6.0466);
\draw[gray!55, line width=0.15pt] (8.5909,6.1017) -- (8.7098,6.2373);
\draw[gray!55, line width=0.15pt] (8.7098,6.2373) -- (8.7458,6.2786);
\draw[gray!55, line width=0.15pt] (8.8267,6.3729) -- (8.7458,6.2786);
\draw[gray!55, line width=0.15pt] (8.8267,6.3729) -- (8.9416,6.5085);
\draw[gray!55, line width=0.15pt] (8.9416,6.5085) -- (8.9492,6.5175);
\draw[gray!55, line width=0.15pt] (9.0537,6.6441) -- (8.9492,6.5175);
\draw[gray!55, line width=0.15pt] (9.0537,6.6441) -- (9.1525,6.7662);
\draw[gray!55, line width=0.15pt] (9.1633,6.7797) -- (9.1525,6.7662);
\draw[gray!55, line width=0.15pt] (9.1633,6.7797) -- (9.2698,6.9153);
\draw[gray!55, line width=0.15pt] (9.2698,6.9153) -- (9.3559,7.0277);
\draw[gray!55, line width=0.15pt] (9.3735,7.0508) -- (9.3559,7.0277);
\draw[gray!55, line width=0.15pt] (9.3735,7.0508) -- (9.4738,7.1864);
\draw[gray!55, line width=0.15pt] (9.4738,7.1864) -- (9.5593,7.3054);
\draw[gray!55, line width=0.15pt] (9.5712,7.3220) -- (9.5593,7.3054);
\draw[gray!55, line width=0.15pt] (9.5712,7.3220) -- (9.6651,7.4576);
\draw[gray!55, line width=0.15pt] (9.6651,7.4576) -- (9.7559,7.5932);
\draw[gray!55, line width=0.15pt] (9.7559,7.5932) -- (9.7627,7.6037);
\draw[gray!55, line width=0.15pt] (9.8434,7.7288) -- (9.7627,7.6037);
\draw[gray!55, line width=0.15pt] (9.8434,7.7288) -- (9.9278,7.8644);
\draw[gray!55, line width=0.15pt] (9.9278,7.8644) -- (9.9661,7.9281);
\draw[gray!55, line width=0.15pt] (10.0091,8.0000) -- (9.9661,7.9281);
\draw[gray!55, line width=0.15pt] (11.3942,0.0000) -- (11.5932,0.0528);
\draw[gray!55, line width=0.15pt] (11.5932,0.0528) -- (11.7966,0.0997);
\draw[gray!55, line width=0.15pt] (11.7966,0.0997) -- (11.9761,0.1356);
\draw[gray!55, line width=0.15pt] (11.9761,0.1356) -- (12.0000,0.1400);
\draw[gray!65, line width=0.15pt] (1.0384,0.0000) -- (1.1389,0.1356);
\draw[gray!65, line width=0.15pt] (1.1389,0.1356) -- (1.2203,0.2418);
\draw[gray!65, line width=0.15pt] (1.2435,0.2712) -- (1.2203,0.2418);
\draw[gray!65, line width=0.15pt] (1.2435,0.2712) -- (1.3536,0.4068);
\draw[gray!65, line width=0.15pt] (1.3536,0.4068) -- (1.4237,0.4907);
\draw[gray!65, line width=0.15pt] (1.4680,0.5424) -- (1.4237,0.4907);
\draw[gray!65, line width=0.15pt] (1.4680,0.5424) -- (1.5873,0.6780);
\draw[gray!65, line width=0.15pt] (1.5873,0.6780) -- (1.6271,0.7222);
\draw[gray!65, line width=0.15pt] (1.7114,0.8136) -- (1.6271,0.7222);
\draw[gray!65, line width=0.15pt] (1.7114,0.8136) -- (1.8305,0.9403);
\draw[gray!65, line width=0.15pt] (1.8390,0.9492) -- (1.8305,0.9403);
\draw[gray!65, line width=0.15pt] (1.8390,0.9492) -- (1.9719,1.0847);
\draw[gray!65, line width=0.15pt] (1.9719,1.0847) -- (2.0339,1.1472);
\draw[gray!65, line width=0.15pt] (2.1081,1.2203) -- (2.0339,1.1472);
\draw[gray!65, line width=0.15pt] (2.1081,1.2203) -- (2.2373,1.3465);
\draw[gray!65, line width=0.15pt] (2.2471,1.3559) -- (2.2373,1.3465);
\draw[gray!65, line width=0.15pt] (2.2471,1.3559) -- (2.3900,1.4915);
\draw[gray!65, line width=0.15pt] (2.3900,1.4915) -- (2.4407,1.5394);
\draw[gray!65, line width=0.15pt] (2.5351,1.6271) -- (2.4407,1.5394);
\draw[gray!65, line width=0.15pt] (2.5351,1.6271) -- (2.6441,1.7283);
\draw[gray!65, line width=0.15pt] (2.6816,1.7627) -- (2.6441,1.7283);
\draw[gray!65, line width=0.15pt] (2.6816,1.7627) -- (2.8296,1.8983);
\draw[gray!65, line width=0.15pt] (2.8296,1.8983) -- (2.8475,1.9147);
\draw[gray!65, line width=0.15pt] (2.9787,2.0339) -- (2.8475,1.9147);
\draw[gray!65, line width=0.15pt] (2.9787,2.0339) -- (3.0508,2.0999);
\draw[gray!65, line width=0.15pt] (3.1277,2.1695) -- (3.0508,2.0999);
\draw[gray!65, line width=0.15pt] (3.1277,2.1695) -- (3.2542,2.2855);
\draw[gray!65, line width=0.15pt] (3.2758,2.3051) -- (3.2542,2.2855);
\draw[gray!65, line width=0.15pt] (3.2758,2.3051) -- (3.4235,2.4407);
\draw[gray!65, line width=0.15pt] (3.4235,2.4407) -- (3.4576,2.4724);
\draw[gray!65, line width=0.15pt] (3.5699,2.5763) -- (3.4576,2.4724);
\draw[gray!65, line width=0.15pt] (3.5699,2.5763) -- (3.6610,2.6619);
\draw[gray!65, line width=0.15pt] (3.7143,2.7119) -- (3.6610,2.6619);
\draw[gray!65, line width=0.15pt] (3.7143,2.7119) -- (3.8567,2.8475);
\draw[gray!65, line width=0.15pt] (3.8567,2.8475) -- (3.8644,2.8548);
\draw[gray!65, line width=0.15pt] (3.9977,2.9831) -- (3.8644,2.8548);
\draw[gray!65, line width=0.15pt] (3.9977,2.9831) -- (4.0678,3.0514);
\draw[gray!65, line width=0.15pt] (4.1364,3.1186) -- (4.0678,3.0514);
\draw[gray!65, line width=0.15pt] (4.1364,3.1186) -- (4.2712,3.2524);
\draw[gray!65, line width=0.15pt] (4.2731,3.2542) -- (4.2712,3.2524);
\draw[gray!65, line width=0.15pt] (4.2731,3.2542) -- (4.4084,3.3898);
\draw[gray!65, line width=0.15pt] (4.4084,3.3898) -- (4.4746,3.4567);
\draw[gray!65, line width=0.15pt] (4.5421,3.5254) -- (4.4746,3.4567);
\draw[gray!65, line width=0.15pt] (4.5421,3.5254) -- (4.6747,3.6610);
\draw[gray!65, line width=0.15pt] (4.6747,3.6610) -- (4.6780,3.6644);
\draw[gray!65, line width=0.15pt] (4.8066,3.7966) -- (4.6780,3.6644);
\draw[gray!65, line width=0.15pt] (4.8066,3.7966) -- (4.8814,3.8735);
\draw[gray!65, line width=0.15pt] (4.9380,3.9322) -- (4.8814,3.8735);
\draw[gray!65, line width=0.15pt] (4.9380,3.9322) -- (5.0693,4.0678);
\draw[gray!65, line width=0.15pt] (5.0693,4.0678) -- (5.0847,4.0836);
\draw[gray!65, line width=0.15pt] (5.2009,4.2034) -- (5.0847,4.0836);
\draw[gray!65, line width=0.15pt] (5.2009,4.2034) -- (5.2881,4.2927);
\draw[gray!65, line width=0.15pt] (5.3330,4.3390) -- (5.2881,4.2927);
\draw[gray!65, line width=0.15pt] (5.3330,4.3390) -- (5.4659,4.4746);
\draw[gray!65, line width=0.15pt] (5.4659,4.4746) -- (5.4915,4.5003);
\draw[gray!65, line width=0.15pt] (5.5997,4.6102) -- (5.4915,4.5003);
\draw[gray!65, line width=0.15pt] (5.5997,4.6102) -- (5.6949,4.7056);
\draw[gray!65, line width=0.15pt] (5.7346,4.7458) -- (5.6949,4.7056);
\draw[gray!65, line width=0.15pt] (5.7346,4.7458) -- (5.8705,4.8814);
\draw[gray!65, line width=0.15pt] (5.8705,4.8814) -- (5.8983,4.9087);
\draw[gray!65, line width=0.15pt] (6.0072,5.0169) -- (5.8983,4.9087);
\draw[gray!65, line width=0.15pt] (6.0072,5.0169) -- (6.1017,5.1096);
\draw[gray!65, line width=0.15pt] (6.1449,5.1525) -- (6.1017,5.1096);
\draw[gray!65, line width=0.15pt] (6.1449,5.1525) -- (6.2833,5.2881);
\draw[gray!65, line width=0.15pt] (6.2833,5.2881) -- (6.3051,5.3093);
\draw[gray!65, line width=0.15pt] (6.4217,5.4237) -- (6.3051,5.3093);
\draw[gray!65, line width=0.15pt] (6.4217,5.4237) -- (6.5085,5.5082);
\draw[gray!65, line width=0.15pt] (6.5603,5.5593) -- (6.5085,5.5082);
\draw[gray!65, line width=0.15pt] (6.5603,5.5593) -- (6.6986,5.6949);
\draw[gray!65, line width=0.15pt] (6.6986,5.6949) -- (6.7119,5.7079);
\draw[gray!65, line width=0.15pt] (6.8357,5.8305) -- (6.7119,5.7079);
\draw[gray!65, line width=0.15pt] (6.8357,5.8305) -- (6.9153,5.9092);
\draw[gray!65, line width=0.15pt] (6.9719,5.9661) -- (6.9153,5.9092);
\draw[gray!65, line width=0.15pt] (6.9719,5.9661) -- (7.1066,6.1017);
\draw[gray!65, line width=0.15pt] (7.1066,6.1017) -- (7.1186,6.1138);
\draw[gray!65, line width=0.15pt] (7.2388,6.2373) -- (7.1186,6.1138);
\draw[gray!65, line width=0.15pt] (7.2388,6.2373) -- (7.3220,6.3235);
\draw[gray!65, line width=0.15pt] (7.3690,6.3729) -- (7.3220,6.3235);
\draw[gray!65, line width=0.15pt] (7.3690,6.3729) -- (7.4965,6.5085);
\draw[gray!65, line width=0.15pt] (7.4965,6.5085) -- (7.5254,6.5397);
\draw[gray!65, line width=0.15pt] (7.6206,6.6441) -- (7.5254,6.5397);
\draw[gray!65, line width=0.15pt] (7.6206,6.6441) -- (7.7288,6.7647);
\draw[gray!65, line width=0.15pt] (7.7420,6.7797) -- (7.7288,6.7647);
\draw[gray!65, line width=0.15pt] (7.7420,6.7797) -- (7.8592,6.9153);
\draw[gray!65, line width=0.15pt] (7.8592,6.9153) -- (7.9322,7.0015);
\draw[gray!65, line width=0.15pt] (7.9732,7.0508) -- (7.9322,7.0015);
\draw[gray!65, line width=0.15pt] (7.9732,7.0508) -- (8.0833,7.1864);
\draw[gray!65, line width=0.15pt] (8.0833,7.1864) -- (8.1356,7.2527);
\draw[gray!65, line width=0.15pt] (8.1895,7.3220) -- (8.1356,7.2527);
\draw[gray!65, line width=0.15pt] (8.1895,7.3220) -- (8.2918,7.4576);
\draw[gray!65, line width=0.15pt] (8.2918,7.4576) -- (8.3390,7.5223);
\draw[gray!65, line width=0.15pt] (8.3901,7.5932) -- (8.3390,7.5223);
\draw[gray!65, line width=0.15pt] (8.3901,7.5932) -- (8.4844,7.7288);
\draw[gray!65, line width=0.15pt] (8.4844,7.7288) -- (8.5424,7.8152);
\draw[gray!65, line width=0.15pt] (8.5750,7.8644) -- (8.5424,7.8152);
\draw[gray!65, line width=0.15pt] (8.5750,7.8644) -- (8.6615,8.0000);
\draw[gray!75, line width=0.15pt] (0.0039,0.6780) -- (0.0000,0.6730);
\draw[gray!75, line width=0.15pt] (0.0039,0.6780) -- (0.1121,0.8136);
\draw[gray!75, line width=0.15pt] (0.1121,0.8136) -- (0.2034,0.9249);
\draw[gray!75, line width=0.15pt] (0.2236,0.9492) -- (0.2034,0.9249);
\draw[gray!75, line width=0.15pt] (0.2236,0.9492) -- (0.3395,1.0847);
\draw[gray!75, line width=0.15pt] (0.3395,1.0847) -- (0.4068,1.1618);
\draw[gray!75, line width=0.15pt] (0.4588,1.2203) -- (0.4068,1.1618);
\draw[gray!75, line width=0.15pt] (0.4588,1.2203) -- (0.5816,1.3559);
\draw[gray!75, line width=0.15pt] (0.5816,1.3559) -- (0.6102,1.3869);
\draw[gray!75, line width=0.15pt] (0.7081,1.4915) -- (0.6102,1.3869);
\draw[gray!75, line width=0.15pt] (0.7081,1.4915) -- (0.8136,1.6026);
\draw[gray!75, line width=0.15pt] (0.8372,1.6271) -- (0.8136,1.6026);
\draw[gray!75, line width=0.15pt] (0.8372,1.6271) -- (0.9693,1.7627);
\draw[gray!75, line width=0.15pt] (0.9693,1.7627) -- (1.0169,1.8111);
\draw[gray!75, line width=0.15pt] (1.1039,1.8983) -- (1.0169,1.8111);
\draw[gray!75, line width=0.15pt] (1.1039,1.8983) -- (1.2203,2.0145);
\draw[gray!75, line width=0.15pt] (1.2401,2.0339) -- (1.2203,2.0145);
\draw[gray!75, line width=0.15pt] (1.2401,2.0339) -- (1.3783,2.1695);
\draw[gray!75, line width=0.15pt] (1.3783,2.1695) -- (1.4237,2.2141);
\draw[gray!75, line width=0.15pt] (1.5176,2.3051) -- (1.4237,2.2141);
\draw[gray!75, line width=0.15pt] (1.5176,2.3051) -- (1.6271,2.4116);
\draw[gray!75, line width=0.15pt] (1.6574,2.4407) -- (1.6271,2.4116);
\draw[gray!75, line width=0.15pt] (1.6574,2.4407) -- (1.7977,2.5763);
\draw[gray!75, line width=0.15pt] (1.7977,2.5763) -- (1.8305,2.6082);
\draw[gray!75, line width=0.15pt] (1.9381,2.7119) -- (1.8305,2.6082);
\draw[gray!75, line width=0.15pt] (1.9381,2.7119) -- (2.0339,2.8050);
\draw[gray!75, line width=0.15pt] (2.0780,2.8475) -- (2.0339,2.8050);
\draw[gray!75, line width=0.15pt] (2.0780,2.8475) -- (2.2171,2.9831);
\draw[gray!75, line width=0.15pt] (2.2171,2.9831) -- (2.2373,3.0029);
\draw[gray!75, line width=0.15pt] (2.3557,3.1186) -- (2.2373,3.0029);
\draw[gray!75, line width=0.15pt] (2.3557,3.1186) -- (2.4407,3.2028);
\draw[gray!75, line width=0.15pt] (2.4931,3.2542) -- (2.4407,3.2028);
\draw[gray!75, line width=0.15pt] (2.4931,3.2542) -- (2.6294,3.3898);
\draw[gray!75, line width=0.15pt] (2.6294,3.3898) -- (2.6441,3.4046);
\draw[gray!75, line width=0.15pt] (2.7649,3.5254) -- (2.6441,3.4046);
\draw[gray!75, line width=0.15pt] (2.7649,3.5254) -- (2.8475,3.6089);
\draw[gray!75, line width=0.15pt] (2.8994,3.6610) -- (2.8475,3.6089);
\draw[gray!75, line width=0.15pt] (2.8994,3.6610) -- (3.0332,3.7966);
\draw[gray!75, line width=0.15pt] (3.0332,3.7966) -- (3.0508,3.8147);
\draw[gray!75, line width=0.15pt] (3.1665,3.9322) -- (3.0508,3.8147);
\draw[gray!75, line width=0.15pt] (3.1665,3.9322) -- (3.2542,4.0218);
\draw[gray!75, line width=0.15pt] (3.2996,4.0678) -- (3.2542,4.0218);
\draw[gray!75, line width=0.15pt] (3.2996,4.0678) -- (3.4328,4.2034);
\draw[gray!75, line width=0.15pt] (3.4328,4.2034) -- (3.4576,4.2287);
\draw[gray!75, line width=0.15pt] (3.5664,4.3390) -- (3.4576,4.2287);
\draw[gray!75, line width=0.15pt] (3.5664,4.3390) -- (3.6610,4.4345);
\draw[gray!75, line width=0.15pt] (3.7009,4.4746) -- (3.6610,4.4345);
\draw[gray!75, line width=0.15pt] (3.7009,4.4746) -- (3.8365,4.6102);
\draw[gray!75, line width=0.15pt] (3.8365,4.6102) -- (3.8644,4.6379);
\draw[gray!75, line width=0.15pt] (3.9733,4.7458) -- (3.8644,4.6379);
\draw[gray!75, line width=0.15pt] (3.9733,4.7458) -- (4.0678,4.8383);
\draw[gray!75, line width=0.15pt] (4.1119,4.8814) -- (4.0678,4.8383);
\draw[gray!75, line width=0.15pt] (4.1119,4.8814) -- (4.2523,5.0169);
\draw[gray!75, line width=0.15pt] (4.2523,5.0169) -- (4.2712,5.0351);
\draw[gray!75, line width=0.15pt] (4.3940,5.1525) -- (4.2712,5.0351);
\draw[gray!75, line width=0.15pt] (4.3940,5.1525) -- (4.4746,5.2285);
\draw[gray!75, line width=0.15pt] (4.5378,5.2881) -- (4.4746,5.2285);
\draw[gray!75, line width=0.15pt] (4.5378,5.2881) -- (4.6780,5.4187);
\draw[gray!75, line width=0.15pt] (4.6834,5.4237) -- (4.6780,5.4187);
\draw[gray!75, line width=0.15pt] (4.6834,5.4237) -- (4.8296,5.5593);
\draw[gray!75, line width=0.15pt] (4.8296,5.5593) -- (4.8814,5.6067);
\draw[gray!75, line width=0.15pt] (4.9771,5.6949) -- (4.8814,5.6067);
\draw[gray!75, line width=0.15pt] (4.9771,5.6949) -- (5.0847,5.7932);
\draw[gray!75, line width=0.15pt] (5.1253,5.8305) -- (5.0847,5.7932);
\draw[gray!75, line width=0.15pt] (5.1253,5.8305) -- (5.2737,5.9661);
\draw[gray!75, line width=0.15pt] (5.2737,5.9661) -- (5.2881,5.9793);
\draw[gray!75, line width=0.15pt] (5.4209,6.1017) -- (5.2881,5.9793);
\draw[gray!75, line width=0.15pt] (5.4209,6.1017) -- (5.4915,6.1665);
\draw[gray!75, line width=0.15pt] (5.5675,6.2373) -- (5.4915,6.1665);
\draw[gray!75, line width=0.15pt] (5.5675,6.2373) -- (5.6949,6.3559);
\draw[gray!75, line width=0.15pt] (5.7129,6.3729) -- (5.6949,6.3559);
\draw[gray!75, line width=0.15pt] (5.7129,6.3729) -- (5.8555,6.5085);
\draw[gray!75, line width=0.15pt] (5.8555,6.5085) -- (5.8983,6.5494);
\draw[gray!75, line width=0.15pt] (5.9955,6.6441) -- (5.8983,6.5494);
\draw[gray!75, line width=0.15pt] (5.9955,6.6441) -- (6.1017,6.7484);
\draw[gray!75, line width=0.15pt] (6.1329,6.7797) -- (6.1017,6.7484);
\draw[gray!75, line width=0.15pt] (6.1329,6.7797) -- (6.2664,6.9153);
\draw[gray!75, line width=0.15pt] (6.2664,6.9153) -- (6.3051,6.9551);
\draw[gray!75, line width=0.15pt] (6.3959,7.0508) -- (6.3051,6.9551);
\draw[gray!75, line width=0.15pt] (6.3959,7.0508) -- (6.5085,7.1716);
\draw[gray!75, line width=0.15pt] (6.5220,7.1864) -- (6.5085,7.1716);
\draw[gray!75, line width=0.15pt] (6.5220,7.1864) -- (6.6428,7.3220);
\draw[gray!75, line width=0.15pt] (6.6428,7.3220) -- (6.7119,7.4015);
\draw[gray!75, line width=0.15pt] (6.7596,7.4576) -- (6.7119,7.4015);
\draw[gray!75, line width=0.15pt] (6.7596,7.4576) -- (6.8717,7.5932);
\draw[gray!75, line width=0.15pt] (6.8717,7.5932) -- (6.9153,7.6475);
\draw[gray!75, line width=0.15pt] (6.9789,7.7288) -- (6.9153,7.6475);
\draw[gray!75, line width=0.15pt] (6.9789,7.7288) -- (7.0817,7.8644);
\draw[gray!75, line width=0.15pt] (7.0817,7.8644) -- (7.1186,7.9148);
\draw[gray!75, line width=0.15pt] (7.1795,8.0000) -- (7.1186,7.9148);
\draw[gray!85, line width=0.15pt] (0.0959,2.4407) -- (0.0000,2.3371);
\draw[gray!85, line width=0.15pt] (0.0959,2.4407) -- (0.2034,2.5575);
\draw[gray!85, line width=0.15pt] (0.2209,2.5763) -- (0.2034,2.5575);
\draw[gray!85, line width=0.15pt] (0.2209,2.5763) -- (0.3458,2.7119);
\draw[gray!85, line width=0.15pt] (0.3458,2.7119) -- (0.4068,2.7789);
\draw[gray!85, line width=0.15pt] (0.4700,2.8475) -- (0.4068,2.7789);
\draw[gray!85, line width=0.15pt] (0.4700,2.8475) -- (0.5930,2.9831);
\draw[gray!85, line width=0.15pt] (0.5930,2.9831) -- (0.6102,3.0024);
\draw[gray!85, line width=0.15pt] (0.7149,3.1186) -- (0.6102,3.0024);
\draw[gray!85, line width=0.15pt] (0.7149,3.1186) -- (0.8136,3.2304);
\draw[gray!85, line width=0.15pt] (0.8350,3.2542) -- (0.8136,3.2304);
\draw[gray!85, line width=0.15pt] (0.8350,3.2542) -- (0.9534,3.3898);
\draw[gray!85, line width=0.15pt] (0.9534,3.3898) -- (1.0169,3.4644);
\draw[gray!85, line width=0.15pt] (1.0698,3.5254) -- (1.0169,3.4644);
\draw[gray!85, line width=0.15pt] (1.0698,3.5254) -- (1.1843,3.6610);
\draw[gray!85, line width=0.15pt] (1.1843,3.6610) -- (1.2203,3.7050);
\draw[gray!85, line width=0.15pt] (1.2969,3.7966) -- (1.2203,3.7050);
\draw[gray!85, line width=0.15pt] (1.2969,3.7966) -- (1.4076,3.9322);
\draw[gray!85, line width=0.15pt] (1.4076,3.9322) -- (1.4237,3.9525);
\draw[gray!85, line width=0.15pt] (1.5168,4.0678) -- (1.4237,3.9525);
\draw[gray!85, line width=0.15pt] (1.5168,4.0678) -- (1.6248,4.2034);
\draw[gray!85, line width=0.15pt] (1.6248,4.2034) -- (1.6271,4.2064);
\draw[gray!85, line width=0.15pt] (1.7317,4.3390) -- (1.6271,4.2064);
\draw[gray!85, line width=0.15pt] (1.7317,4.3390) -- (1.8305,4.4647);
\draw[gray!85, line width=0.15pt] (1.8384,4.4746) -- (1.8305,4.4647);
\draw[gray!85, line width=0.15pt] (1.8384,4.4746) -- (1.9449,4.6102);
\draw[gray!85, line width=0.15pt] (1.9449,4.6102) -- (2.0339,4.7231);
\draw[gray!85, line width=0.15pt] (2.0522,4.7458) -- (2.0339,4.7231);
\draw[gray!85, line width=0.15pt] (2.0522,4.7458) -- (2.1604,4.8814);
\draw[gray!85, line width=0.15pt] (2.1604,4.8814) -- (2.2373,4.9764);
\draw[gray!85, line width=0.15pt] (2.2707,5.0169) -- (2.2373,4.9764);
\draw[gray!85, line width=0.15pt] (2.2707,5.0169) -- (2.3833,5.1525);
\draw[gray!85, line width=0.15pt] (2.3833,5.1525) -- (2.4407,5.2205);
\draw[gray!85, line width=0.15pt] (2.4990,5.2881) -- (2.4407,5.2205);
\draw[gray!85, line width=0.15pt] (2.4990,5.2881) -- (2.6184,5.4237);
\draw[gray!85, line width=0.15pt] (2.6184,5.4237) -- (2.6441,5.4524);
\draw[gray!85, line width=0.15pt] (2.7414,5.5593) -- (2.6441,5.4524);
\draw[gray!85, line width=0.15pt] (2.7414,5.5593) -- (2.8475,5.6718);
\draw[gray!85, line width=0.15pt] (2.8696,5.6949) -- (2.8475,5.6718);
\draw[gray!85, line width=0.15pt] (2.8696,5.6949) -- (3.0017,5.8305);
\draw[gray!85, line width=0.15pt] (3.0017,5.8305) -- (3.0508,5.8795);
\draw[gray!85, line width=0.15pt] (3.1385,5.9661) -- (3.0508,5.8795);
\draw[gray!85, line width=0.15pt] (3.1385,5.9661) -- (3.2542,6.0768);
\draw[gray!85, line width=0.15pt] (3.2804,6.1017) -- (3.2542,6.0768);
\draw[gray!85, line width=0.15pt] (3.2804,6.1017) -- (3.4259,6.2373);
\draw[gray!85, line width=0.15pt] (3.4259,6.2373) -- (3.4576,6.2662);
\draw[gray!85, line width=0.15pt] (3.5744,6.3729) -- (3.4576,6.2662);
\draw[gray!85, line width=0.15pt] (3.5744,6.3729) -- (3.6610,6.4500);
\draw[gray!85, line width=0.15pt] (3.7263,6.5085) -- (3.6610,6.4500);
\draw[gray!85, line width=0.15pt] (3.7263,6.5085) -- (3.8644,6.6298);
\draw[gray!85, line width=0.15pt] (3.8804,6.6441) -- (3.8644,6.6298);
\draw[gray!85, line width=0.15pt] (3.8804,6.6441) -- (4.0343,6.7797);
\draw[gray!85, line width=0.15pt] (4.0343,6.7797) -- (4.0678,6.8088);
\draw[gray!85, line width=0.15pt] (4.1877,6.9153) -- (4.0678,6.8088);
\draw[gray!85, line width=0.15pt] (4.1877,6.9153) -- (4.2712,6.9887);
\draw[gray!85, line width=0.15pt] (4.3403,7.0508) -- (4.2712,6.9887);
\draw[gray!85, line width=0.15pt] (4.3403,7.0508) -- (4.4746,7.1711);
\draw[gray!85, line width=0.15pt] (4.4913,7.1864) -- (4.4746,7.1711);
\draw[gray!85, line width=0.15pt] (4.4913,7.1864) -- (4.6378,7.3220);
\draw[gray!85, line width=0.15pt] (4.6378,7.3220) -- (4.6780,7.3594);
\draw[gray!85, line width=0.15pt] (4.7803,7.4576) -- (4.6780,7.3594);
\draw[gray!85, line width=0.15pt] (4.7803,7.4576) -- (4.8814,7.5554);
\draw[gray!85, line width=0.15pt] (4.9191,7.5932) -- (4.8814,7.5554);
\draw[gray!85, line width=0.15pt] (4.9191,7.5932) -- (5.0524,7.7288);
\draw[gray!85, line width=0.15pt] (5.0524,7.7288) -- (5.0847,7.7624);
\draw[gray!85, line width=0.15pt] (5.1795,7.8644) -- (5.0847,7.7624);
\draw[gray!85, line width=0.15pt] (5.1795,7.8644) -- (5.2881,7.9837);
\draw[gray!85, line width=0.15pt] (5.3024,8.0000) -- (5.2881,7.9837);
\draw[blue!80!black, line width=1.0pt] (0.0000,0.0000) -- (4.8323,0.0000) -- (6.1370,0.0000) -- (6.4231,0.0000) -- (6.5066,0.0000) -- (6.5327,0.0000) -- (6.5411,0.0000) -- (6.5438,0.0000) -- (6.5447,0.0000);
\fill[blue!80!black] (0.0000,0.0000) circle (1.2pt);
\fill[blue!80!black] (4.8323,0.0000) circle (1.2pt);
\fill[blue!80!black] (6.1370,0.0000) circle (1.2pt);
\fill[blue!80!black] (6.4231,0.0000) circle (1.2pt);
\fill[blue!80!black] (6.5066,0.0000) circle (1.2pt);
\fill[blue!80!black] (6.5327,0.0000) circle (1.2pt);
\fill[blue!80!black] (6.5411,0.0000) circle (1.2pt);
\fill[blue!80!black] (6.5438,0.0000) circle (1.2pt);
\fill[blue!80!black] (6.5447,0.0000) circle (1.2pt);
\draw[red!80!black, line width=1.0pt] (0.0000,0.0000) -- (6.5451,0.0000) -- (6.5451,0.0000);
\fill[red!80!black] (0.0000,0.0000) circle (1.2pt);
\fill[red!80!black] (6.5451,0.0000) circle (1.2pt);
\fill[red!80!black] (6.5451,0.0000) circle (1.2pt);
\draw[->, thick] (0,0) -- (12.6000,0) node[right] {$x$};
\draw[->, thick] (0,0) -- (0,8.6000) node[above] {$y$};
\draw (0.0000,0) -- (0.0000,-0.12) node[below] {\small -1.00};
\draw (0,0.0000) -- (-0.12,0.0000) node[left] {\small -1.00};
\draw (2.4000,0) -- (2.4000,-0.12) node[below] {\small -0.40};
\draw (0,1.6000) -- (-0.12,1.6000) node[left] {\small -0.60};
\draw (4.8000,0) -- (4.8000,-0.12) node[below] {\small 0.20};
\draw (0,3.2000) -- (-0.12,3.2000) node[left] {\small -0.20};
\draw (7.2000,0) -- (7.2000,-0.12) node[below] {\small 0.80};
\draw (0,4.8000) -- (-0.12,4.8000) node[left] {\small 0.20};
\draw (9.6000,0) -- (9.6000,-0.12) node[below] {\small 1.40};
\draw (0,6.4000) -- (-0.12,6.4000) node[left] {\small 0.60};
\draw (12.0000,0) -- (12.0000,-0.12) node[below] {\small 2.00};
\draw (0,8.0000) -- (-0.12,8.0000) node[left] {\small 1.00};
\draw[fill=white, draw=black] (7.6,6.7) rectangle (11.8,7.9);
\draw[blue!80!black, line width=1.0pt] (7.9,7.45) -- (8.9,7.45);
\node[anchor=west] at (9.1,7.45) {\small gradDest};
\draw[red!80!black, line width=1.0pt] (7.9,7.05) -- (8.9,7.05);
\node[anchor=west] at (9.1,7.05) {\small fastGradDest};
\end{tikzpicture}

\caption{Функция Adjiman.}
\end{figure}

\begin{figure}[H]
\centering
\begin{tikzpicture}[scale=0.9]
\draw[fill=orange!3, draw=black] (0,0) rectangle (12.0000,8.0000);
\draw[gray!35, line width=0.15pt] (5.0847,3.7909) -- (5.0811,3.7966);
\draw[gray!35, line width=0.15pt] (5.0811,3.7966) -- (5.0368,3.9322);
\draw[gray!35, line width=0.15pt] (5.0368,3.9322) -- (5.0368,4.0678);
\draw[gray!35, line width=0.15pt] (5.0368,4.0678) -- (5.0811,4.2034);
\draw[gray!35, line width=0.15pt] (5.0811,4.2034) -- (5.0847,4.2091);
\draw[gray!35, line width=0.15pt] (5.2881,3.5648) -- (5.1774,3.6610);
\draw[gray!35, line width=0.15pt] (5.0847,3.7909) -- (5.1774,3.6610);
\draw[gray!35, line width=0.15pt] (5.1774,4.3390) -- (5.0847,4.2091);
\draw[gray!35, line width=0.15pt] (5.1774,4.3390) -- (5.2881,4.4352);
\draw[gray!35, line width=0.15pt] (5.4915,3.4516) -- (5.3473,3.5254);
\draw[gray!35, line width=0.15pt] (5.2881,3.5648) -- (5.3473,3.5254);
\draw[gray!35, line width=0.15pt] (5.3473,4.4746) -- (5.2881,4.4352);
\draw[gray!35, line width=0.15pt] (5.3473,4.4746) -- (5.4915,4.5484);
\draw[gray!35, line width=0.15pt] (5.6949,3.3874) -- (5.6864,3.3898);
\draw[gray!35, line width=0.15pt] (5.4915,3.4516) -- (5.6864,3.3898);
\draw[gray!35, line width=0.15pt] (5.6864,4.6102) -- (5.4915,4.5484);
\draw[gray!35, line width=0.15pt] (5.6864,4.6102) -- (5.6949,4.6126);
\draw[gray!35, line width=0.15pt] (5.6949,3.3874) -- (5.8983,3.3579);
\draw[gray!35, line width=0.15pt] (5.8983,4.6421) -- (5.6949,4.6126);
\draw[gray!35, line width=0.15pt] (5.8983,3.3579) -- (6.1017,3.3579);
\draw[gray!35, line width=0.15pt] (6.1017,4.6421) -- (5.8983,4.6421);
\draw[gray!35, line width=0.15pt] (6.1017,3.3579) -- (6.3051,3.3874);
\draw[gray!35, line width=0.15pt] (6.3051,4.6126) -- (6.1017,4.6421);
\draw[gray!35, line width=0.15pt] (6.3051,3.3874) -- (6.3136,3.3898);
\draw[gray!35, line width=0.15pt] (6.3136,3.3898) -- (6.5085,3.4516);
\draw[gray!35, line width=0.15pt] (6.5085,4.5484) -- (6.3136,4.6102);
\draw[gray!35, line width=0.15pt] (6.3051,4.6126) -- (6.3136,4.6102);
\draw[gray!35, line width=0.15pt] (6.5085,3.4516) -- (6.6527,3.5254);
\draw[gray!35, line width=0.15pt] (6.6527,3.5254) -- (6.7119,3.5648);
\draw[gray!35, line width=0.15pt] (6.7119,4.4352) -- (6.6527,4.4746);
\draw[gray!35, line width=0.15pt] (6.5085,4.5484) -- (6.6527,4.4746);
\draw[gray!35, line width=0.15pt] (6.7119,3.5648) -- (6.8226,3.6610);
\draw[gray!35, line width=0.15pt] (6.8226,3.6610) -- (6.9153,3.7909);
\draw[gray!35, line width=0.15pt] (6.9153,4.2091) -- (6.8226,4.3390);
\draw[gray!35, line width=0.15pt] (6.7119,4.4352) -- (6.8226,4.3390);
\draw[gray!35, line width=0.15pt] (6.9153,3.7909) -- (6.9189,3.7966);
\draw[gray!35, line width=0.15pt] (6.9189,3.7966) -- (6.9632,3.9322);
\draw[gray!35, line width=0.15pt] (6.9632,3.9322) -- (6.9632,4.0678);
\draw[gray!35, line width=0.15pt] (6.9632,4.0678) -- (6.9189,4.2034);
\draw[gray!35, line width=0.15pt] (6.9153,4.2091) -- (6.9189,4.2034);
\draw[gray!45, line width=0.15pt] (4.2712,3.5215) -- (4.2685,3.5254);
\draw[gray!45, line width=0.15pt] (4.2685,3.5254) -- (4.1980,3.6610);
\draw[gray!45, line width=0.15pt] (4.1980,3.6610) -- (4.1525,3.7966);
\draw[gray!45, line width=0.15pt] (4.1525,3.7966) -- (4.1302,3.9322);
\draw[gray!45, line width=0.15pt] (4.1302,3.9322) -- (4.1302,4.0678);
\draw[gray!45, line width=0.15pt] (4.1302,4.0678) -- (4.1525,4.2034);
\draw[gray!45, line width=0.15pt] (4.1525,4.2034) -- (4.1980,4.3390);
\draw[gray!45, line width=0.15pt] (4.1980,4.3390) -- (4.2685,4.4746);
\draw[gray!45, line width=0.15pt] (4.2685,4.4746) -- (4.2712,4.4785);
\draw[gray!45, line width=0.15pt] (4.4746,3.2776) -- (4.3672,3.3898);
\draw[gray!45, line width=0.15pt] (4.2712,3.5215) -- (4.3672,3.3898);
\draw[gray!45, line width=0.15pt] (4.3672,4.6102) -- (4.2712,4.4785);
\draw[gray!45, line width=0.15pt] (4.3672,4.6102) -- (4.4746,4.7224);
\draw[gray!45, line width=0.15pt] (4.6780,3.1164) -- (4.6746,3.1186);
\draw[gray!45, line width=0.15pt] (4.6746,3.1186) -- (4.4992,3.2542);
\draw[gray!45, line width=0.15pt] (4.4746,3.2776) -- (4.4992,3.2542);
\draw[gray!45, line width=0.15pt] (4.4992,4.7458) -- (4.4746,4.7224);
\draw[gray!45, line width=0.15pt] (4.4992,4.7458) -- (4.6746,4.8814);
\draw[gray!45, line width=0.15pt] (4.6746,4.8814) -- (4.6780,4.8836);
\draw[gray!45, line width=0.15pt] (4.6780,3.1164) -- (4.8814,2.9995);
\draw[gray!45, line width=0.15pt] (4.8814,5.0005) -- (4.6780,4.8836);
\draw[gray!45, line width=0.15pt] (5.0847,2.9114) -- (4.9164,2.9831);
\draw[gray!45, line width=0.15pt] (4.8814,2.9995) -- (4.9164,2.9831);
\draw[gray!45, line width=0.15pt] (4.9164,5.0169) -- (4.8814,5.0005);
\draw[gray!45, line width=0.15pt] (4.9164,5.0169) -- (5.0847,5.0886);
\draw[gray!45, line width=0.15pt] (5.2881,2.8457) -- (5.2823,2.8475);
\draw[gray!45, line width=0.15pt] (5.0847,2.9114) -- (5.2823,2.8475);
\draw[gray!45, line width=0.15pt] (5.2823,5.1525) -- (5.0847,5.0886);
\draw[gray!45, line width=0.15pt] (5.2823,5.1525) -- (5.2881,5.1543);
\draw[gray!45, line width=0.15pt] (5.2881,2.8457) -- (5.4915,2.7986);
\draw[gray!45, line width=0.15pt] (5.4915,5.2014) -- (5.2881,5.1543);
\draw[gray!45, line width=0.15pt] (5.4915,2.7986) -- (5.6949,2.7683);
\draw[gray!45, line width=0.15pt] (5.6949,5.2317) -- (5.4915,5.2014);
\draw[gray!45, line width=0.15pt] (5.6949,2.7683) -- (5.8983,2.7535);
\draw[gray!45, line width=0.15pt] (5.8983,5.2465) -- (5.6949,5.2317);
\draw[gray!45, line width=0.15pt] (5.8983,2.7535) -- (6.1017,2.7535);
\draw[gray!45, line width=0.15pt] (6.1017,5.2465) -- (5.8983,5.2465);
\draw[gray!45, line width=0.15pt] (6.1017,2.7535) -- (6.3051,2.7683);
\draw[gray!45, line width=0.15pt] (6.3051,5.2317) -- (6.1017,5.2465);
\draw[gray!45, line width=0.15pt] (6.3051,2.7683) -- (6.5085,2.7986);
\draw[gray!45, line width=0.15pt] (6.5085,5.2014) -- (6.3051,5.2317);
\draw[gray!45, line width=0.15pt] (6.5085,2.7986) -- (6.7119,2.8457);
\draw[gray!45, line width=0.15pt] (6.7119,5.1543) -- (6.5085,5.2014);
\draw[gray!45, line width=0.15pt] (6.7119,2.8457) -- (6.7177,2.8475);
\draw[gray!45, line width=0.15pt] (6.7177,2.8475) -- (6.9153,2.9114);
\draw[gray!45, line width=0.15pt] (6.9153,5.0886) -- (6.7177,5.1525);
\draw[gray!45, line width=0.15pt] (6.7119,5.1543) -- (6.7177,5.1525);
\draw[gray!45, line width=0.15pt] (6.9153,2.9114) -- (7.0836,2.9831);
\draw[gray!45, line width=0.15pt] (7.0836,2.9831) -- (7.1186,2.9995);
\draw[gray!45, line width=0.15pt] (7.1186,5.0005) -- (7.0836,5.0169);
\draw[gray!45, line width=0.15pt] (6.9153,5.0886) -- (7.0836,5.0169);
\draw[gray!45, line width=0.15pt] (7.1186,2.9995) -- (7.3220,3.1164);
\draw[gray!45, line width=0.15pt] (7.3220,4.8836) -- (7.1186,5.0005);
\draw[gray!45, line width=0.15pt] (7.3220,3.1164) -- (7.3254,3.1186);
\draw[gray!45, line width=0.15pt] (7.3254,3.1186) -- (7.5008,3.2542);
\draw[gray!45, line width=0.15pt] (7.5008,3.2542) -- (7.5254,3.2776);
\draw[gray!45, line width=0.15pt] (7.5254,4.7224) -- (7.5008,4.7458);
\draw[gray!45, line width=0.15pt] (7.5008,4.7458) -- (7.3254,4.8814);
\draw[gray!45, line width=0.15pt] (7.3220,4.8836) -- (7.3254,4.8814);
\draw[gray!45, line width=0.15pt] (7.5254,3.2776) -- (7.6328,3.3898);
\draw[gray!45, line width=0.15pt] (7.6328,3.3898) -- (7.7288,3.5215);
\draw[gray!45, line width=0.15pt] (7.7288,4.4785) -- (7.6328,4.6102);
\draw[gray!45, line width=0.15pt] (7.5254,4.7224) -- (7.6328,4.6102);
\draw[gray!45, line width=0.15pt] (7.7288,3.5215) -- (7.7315,3.5254);
\draw[gray!45, line width=0.15pt] (7.7315,3.5254) -- (7.8020,3.6610);
\draw[gray!45, line width=0.15pt] (7.8020,3.6610) -- (7.8475,3.7966);
\draw[gray!45, line width=0.15pt] (7.8475,3.7966) -- (7.8698,3.9322);
\draw[gray!45, line width=0.15pt] (7.8698,3.9322) -- (7.8698,4.0678);
\draw[gray!45, line width=0.15pt] (7.8698,4.0678) -- (7.8475,4.2034);
\draw[gray!45, line width=0.15pt] (7.8475,4.2034) -- (7.8020,4.3390);
\draw[gray!45, line width=0.15pt] (7.8020,4.3390) -- (7.7315,4.4746);
\draw[gray!45, line width=0.15pt] (7.7288,4.4785) -- (7.7315,4.4746);
\draw[gray!55, line width=0.15pt] (3.2542,3.4408) -- (3.2168,3.5254);
\draw[gray!55, line width=0.15pt] (3.2168,3.5254) -- (3.1723,3.6610);
\draw[gray!55, line width=0.15pt] (3.1723,3.6610) -- (3.1432,3.7966);
\draw[gray!55, line width=0.15pt] (3.1432,3.7966) -- (3.1287,3.9322);
\draw[gray!55, line width=0.15pt] (3.1287,3.9322) -- (3.1287,4.0678);
\draw[gray!55, line width=0.15pt] (3.1287,4.0678) -- (3.1432,4.2034);
\draw[gray!55, line width=0.15pt] (3.1432,4.2034) -- (3.1723,4.3390);
\draw[gray!55, line width=0.15pt] (3.1723,4.3390) -- (3.2168,4.4746);
\draw[gray!55, line width=0.15pt] (3.2168,4.4746) -- (3.2542,4.5592);
\draw[gray!55, line width=0.15pt] (3.4576,3.1090) -- (3.4497,3.1186);
\draw[gray!55, line width=0.15pt] (3.4497,3.1186) -- (3.3540,3.2542);
\draw[gray!55, line width=0.15pt] (3.3540,3.2542) -- (3.2770,3.3898);
\draw[gray!55, line width=0.15pt] (3.2542,3.4408) -- (3.2770,3.3898);
\draw[gray!55, line width=0.15pt] (3.2770,4.6102) -- (3.2542,4.5592);
\draw[gray!55, line width=0.15pt] (3.2770,4.6102) -- (3.3540,4.7458);
\draw[gray!55, line width=0.15pt] (3.3540,4.7458) -- (3.4497,4.8814);
\draw[gray!55, line width=0.15pt] (3.4497,4.8814) -- (3.4576,4.8910);
\draw[gray!55, line width=0.15pt] (3.6610,2.8893) -- (3.5661,2.9831);
\draw[gray!55, line width=0.15pt] (3.4576,3.1090) -- (3.5661,2.9831);
\draw[gray!55, line width=0.15pt] (3.5661,5.0169) -- (3.4576,4.8910);
\draw[gray!55, line width=0.15pt] (3.5661,5.0169) -- (3.6610,5.1107);
\draw[gray!55, line width=0.15pt] (3.8644,2.7194) -- (3.7062,2.8475);
\draw[gray!55, line width=0.15pt] (3.6610,2.8893) -- (3.7062,2.8475);
\draw[gray!55, line width=0.15pt] (3.7062,5.1525) -- (3.6610,5.1107);
\draw[gray!55, line width=0.15pt] (3.7062,5.1525) -- (3.8644,5.2806);
\draw[gray!55, line width=0.15pt] (4.0678,2.5830) -- (3.8745,2.7119);
\draw[gray!55, line width=0.15pt] (3.8644,2.7194) -- (3.8745,2.7119);
\draw[gray!55, line width=0.15pt] (3.8745,5.2881) -- (3.8644,5.2806);
\draw[gray!55, line width=0.15pt] (3.8745,5.2881) -- (4.0678,5.4170);
\draw[gray!55, line width=0.15pt] (4.2712,2.4708) -- (4.0791,2.5763);
\draw[gray!55, line width=0.15pt] (4.0678,2.5830) -- (4.0791,2.5763);
\draw[gray!55, line width=0.15pt] (4.0791,5.4237) -- (4.0678,5.4170);
\draw[gray!55, line width=0.15pt] (4.0791,5.4237) -- (4.2712,5.5292);
\draw[gray!55, line width=0.15pt] (4.4746,2.3774) -- (4.3340,2.4407);
\draw[gray!55, line width=0.15pt] (4.2712,2.4708) -- (4.3340,2.4407);
\draw[gray!55, line width=0.15pt] (4.3340,5.5593) -- (4.2712,5.5292);
\draw[gray!55, line width=0.15pt] (4.3340,5.5593) -- (4.4746,5.6226);
\draw[gray!55, line width=0.15pt] (4.6780,2.2998) -- (4.6636,2.3051);
\draw[gray!55, line width=0.15pt] (4.4746,2.3774) -- (4.6636,2.3051);
\draw[gray!55, line width=0.15pt] (4.6636,5.6949) -- (4.4746,5.6226);
\draw[gray!55, line width=0.15pt] (4.6636,5.6949) -- (4.6780,5.7002);
\draw[gray!55, line width=0.15pt] (4.6780,2.2998) -- (4.8814,2.2360);
\draw[gray!55, line width=0.15pt] (4.8814,5.7640) -- (4.6780,5.7002);
\draw[gray!55, line width=0.15pt] (4.8814,2.2360) -- (5.0847,2.1847);
\draw[gray!55, line width=0.15pt] (5.0847,5.8153) -- (4.8814,5.7640);
\draw[gray!55, line width=0.15pt] (5.2881,2.1445) -- (5.1611,2.1695);
\draw[gray!55, line width=0.15pt] (5.0847,2.1847) -- (5.1611,2.1695);
\draw[gray!55, line width=0.15pt] (5.1611,5.8305) -- (5.0847,5.8153);
\draw[gray!55, line width=0.15pt] (5.1611,5.8305) -- (5.2881,5.8555);
\draw[gray!55, line width=0.15pt] (5.2881,2.1445) -- (5.4915,2.1149);
\draw[gray!55, line width=0.15pt] (5.4915,5.8851) -- (5.2881,5.8555);
\draw[gray!55, line width=0.15pt] (5.4915,2.1149) -- (5.6949,2.0954);
\draw[gray!55, line width=0.15pt] (5.6949,5.9046) -- (5.4915,5.8851);
\draw[gray!55, line width=0.15pt] (5.6949,2.0954) -- (5.8983,2.0858);
\draw[gray!55, line width=0.15pt] (5.8983,5.9142) -- (5.6949,5.9046);
\draw[gray!55, line width=0.15pt] (5.8983,2.0858) -- (6.1017,2.0858);
\draw[gray!55, line width=0.15pt] (6.1017,5.9142) -- (5.8983,5.9142);
\draw[gray!55, line width=0.15pt] (6.1017,2.0858) -- (6.3051,2.0954);
\draw[gray!55, line width=0.15pt] (6.3051,5.9046) -- (6.1017,5.9142);
\draw[gray!55, line width=0.15pt] (6.3051,2.0954) -- (6.5085,2.1149);
\draw[gray!55, line width=0.15pt] (6.5085,5.8851) -- (6.3051,5.9046);
\draw[gray!55, line width=0.15pt] (6.5085,2.1149) -- (6.7119,2.1445);
\draw[gray!55, line width=0.15pt] (6.7119,5.8555) -- (6.5085,5.8851);
\draw[gray!55, line width=0.15pt] (6.7119,2.1445) -- (6.8389,2.1695);
\draw[gray!55, line width=0.15pt] (6.8389,2.1695) -- (6.9153,2.1847);
\draw[gray!55, line width=0.15pt] (6.9153,5.8153) -- (6.8389,5.8305);
\draw[gray!55, line width=0.15pt] (6.7119,5.8555) -- (6.8389,5.8305);
\draw[gray!55, line width=0.15pt] (6.9153,2.1847) -- (7.1186,2.2360);
\draw[gray!55, line width=0.15pt] (7.1186,5.7640) -- (6.9153,5.8153);
\draw[gray!55, line width=0.15pt] (7.1186,2.2360) -- (7.3220,2.2998);
\draw[gray!55, line width=0.15pt] (7.3220,5.7002) -- (7.1186,5.7640);
\draw[gray!55, line width=0.15pt] (7.3220,2.2998) -- (7.3364,2.3051);
\draw[gray!55, line width=0.15pt] (7.3364,2.3051) -- (7.5254,2.3774);
\draw[gray!55, line width=0.15pt] (7.5254,5.6226) -- (7.3364,5.6949);
\draw[gray!55, line width=0.15pt] (7.3220,5.7002) -- (7.3364,5.6949);
\draw[gray!55, line width=0.15pt] (7.5254,2.3774) -- (7.6660,2.4407);
\draw[gray!55, line width=0.15pt] (7.6660,2.4407) -- (7.7288,2.4708);
\draw[gray!55, line width=0.15pt] (7.7288,5.5292) -- (7.6660,5.5593);
\draw[gray!55, line width=0.15pt] (7.5254,5.6226) -- (7.6660,5.5593);
\draw[gray!55, line width=0.15pt] (7.7288,2.4708) -- (7.9209,2.5763);
\draw[gray!55, line width=0.15pt] (7.9209,2.5763) -- (7.9322,2.5830);
\draw[gray!55, line width=0.15pt] (7.9322,5.4170) -- (7.9209,5.4237);
\draw[gray!55, line width=0.15pt] (7.7288,5.5292) -- (7.9209,5.4237);
\draw[gray!55, line width=0.15pt] (7.9322,2.5830) -- (8.1255,2.7119);
\draw[gray!55, line width=0.15pt] (8.1255,2.7119) -- (8.1356,2.7194);
\draw[gray!55, line width=0.15pt] (8.1356,5.2806) -- (8.1255,5.2881);
\draw[gray!55, line width=0.15pt] (7.9322,5.4170) -- (8.1255,5.2881);
\draw[gray!55, line width=0.15pt] (8.1356,2.7194) -- (8.2938,2.8475);
\draw[gray!55, line width=0.15pt] (8.2938,2.8475) -- (8.3390,2.8893);
\draw[gray!55, line width=0.15pt] (8.3390,5.1107) -- (8.2938,5.1525);
\draw[gray!55, line width=0.15pt] (8.1356,5.2806) -- (8.2938,5.1525);
\draw[gray!55, line width=0.15pt] (8.3390,2.8893) -- (8.4339,2.9831);
\draw[gray!55, line width=0.15pt] (8.4339,2.9831) -- (8.5424,3.1090);
\draw[gray!55, line width=0.15pt] (8.5424,4.8910) -- (8.4339,5.0169);
\draw[gray!55, line width=0.15pt] (8.3390,5.1107) -- (8.4339,5.0169);
\draw[gray!55, line width=0.15pt] (8.5424,3.1090) -- (8.5503,3.1186);
\draw[gray!55, line width=0.15pt] (8.5503,3.1186) -- (8.6460,3.2542);
\draw[gray!55, line width=0.15pt] (8.6460,3.2542) -- (8.7230,3.3898);
\draw[gray!55, line width=0.15pt] (8.7230,3.3898) -- (8.7458,3.4408);
\draw[gray!55, line width=0.15pt] (8.7458,4.5592) -- (8.7230,4.6102);
\draw[gray!55, line width=0.15pt] (8.7230,4.6102) -- (8.6460,4.7458);
\draw[gray!55, line width=0.15pt] (8.6460,4.7458) -- (8.5503,4.8814);
\draw[gray!55, line width=0.15pt] (8.5424,4.8910) -- (8.5503,4.8814);
\draw[gray!55, line width=0.15pt] (8.7458,3.4408) -- (8.7832,3.5254);
\draw[gray!55, line width=0.15pt] (8.7832,3.5254) -- (8.8277,3.6610);
\draw[gray!55, line width=0.15pt] (8.8277,3.6610) -- (8.8568,3.7966);
\draw[gray!55, line width=0.15pt] (8.8568,3.7966) -- (8.8713,3.9322);
\draw[gray!55, line width=0.15pt] (8.8713,3.9322) -- (8.8713,4.0678);
\draw[gray!55, line width=0.15pt] (8.8713,4.0678) -- (8.8568,4.2034);
\draw[gray!55, line width=0.15pt] (8.8568,4.2034) -- (8.8277,4.3390);
\draw[gray!55, line width=0.15pt] (8.8277,4.3390) -- (8.7832,4.4746);
\draw[gray!55, line width=0.15pt] (8.7458,4.5592) -- (8.7832,4.4746);
\draw[gray!65, line width=0.15pt] (2.0339,3.6995) -- (2.0189,3.7966);
\draw[gray!65, line width=0.15pt] (2.0189,3.7966) -- (2.0084,3.9322);
\draw[gray!65, line width=0.15pt] (2.0084,3.9322) -- (2.0084,4.0678);
\draw[gray!65, line width=0.15pt] (2.0084,4.0678) -- (2.0189,4.2034);
\draw[gray!65, line width=0.15pt] (2.0189,4.2034) -- (2.0339,4.3005);
\draw[gray!65, line width=0.15pt] (2.2373,3.1102) -- (2.2325,3.1186);
\draw[gray!65, line width=0.15pt] (2.2325,3.1186) -- (2.1669,3.2542);
\draw[gray!65, line width=0.15pt] (2.1669,3.2542) -- (2.1132,3.3898);
\draw[gray!65, line width=0.15pt] (2.1132,3.3898) -- (2.0710,3.5254);
\draw[gray!65, line width=0.15pt] (2.0710,3.5254) -- (2.0398,3.6610);
\draw[gray!65, line width=0.15pt] (2.0339,3.6995) -- (2.0398,3.6610);
\draw[gray!65, line width=0.15pt] (2.0398,4.3390) -- (2.0339,4.3005);
\draw[gray!65, line width=0.15pt] (2.0398,4.3390) -- (2.0710,4.4746);
\draw[gray!65, line width=0.15pt] (2.0710,4.4746) -- (2.1132,4.6102);
\draw[gray!65, line width=0.15pt] (2.1132,4.6102) -- (2.1669,4.7458);
\draw[gray!65, line width=0.15pt] (2.1669,4.7458) -- (2.2325,4.8814);
\draw[gray!65, line width=0.15pt] (2.2325,4.8814) -- (2.2373,4.8898);
\draw[gray!65, line width=0.15pt] (2.4407,2.7952) -- (2.4009,2.8475);
\draw[gray!65, line width=0.15pt] (2.4009,2.8475) -- (2.3099,2.9831);
\draw[gray!65, line width=0.15pt] (2.2373,3.1102) -- (2.3099,2.9831);
\draw[gray!65, line width=0.15pt] (2.3099,5.0169) -- (2.2373,4.8898);
\draw[gray!65, line width=0.15pt] (2.3099,5.0169) -- (2.4009,5.1525);
\draw[gray!65, line width=0.15pt] (2.4009,5.1525) -- (2.4407,5.2048);
\draw[gray!65, line width=0.15pt] (2.6441,2.5583) -- (2.6265,2.5763);
\draw[gray!65, line width=0.15pt] (2.6265,2.5763) -- (2.5059,2.7119);
\draw[gray!65, line width=0.15pt] (2.4407,2.7952) -- (2.5059,2.7119);
\draw[gray!65, line width=0.15pt] (2.5059,5.2881) -- (2.4407,5.2048);
\draw[gray!65, line width=0.15pt] (2.5059,5.2881) -- (2.6265,5.4237);
\draw[gray!65, line width=0.15pt] (2.6265,5.4237) -- (2.6441,5.4417);
\draw[gray!65, line width=0.15pt] (2.8475,2.3673) -- (2.7643,2.4407);
\draw[gray!65, line width=0.15pt] (2.6441,2.5583) -- (2.7643,2.4407);
\draw[gray!65, line width=0.15pt] (2.7643,5.5593) -- (2.6441,5.4417);
\draw[gray!65, line width=0.15pt] (2.7643,5.5593) -- (2.8475,5.6327);
\draw[gray!65, line width=0.15pt] (3.0508,2.2061) -- (2.9216,2.3051);
\draw[gray!65, line width=0.15pt] (2.8475,2.3673) -- (2.9216,2.3051);
\draw[gray!65, line width=0.15pt] (2.9216,5.6949) -- (2.8475,5.6327);
\draw[gray!65, line width=0.15pt] (2.9216,5.6949) -- (3.0508,5.7939);
\draw[gray!65, line width=0.15pt] (3.2542,2.0678) -- (3.1017,2.1695);
\draw[gray!65, line width=0.15pt] (3.0508,2.2061) -- (3.1017,2.1695);
\draw[gray!65, line width=0.15pt] (3.1017,5.8305) -- (3.0508,5.7939);
\draw[gray!65, line width=0.15pt] (3.1017,5.8305) -- (3.2542,5.9322);
\draw[gray!65, line width=0.15pt] (3.4576,1.9478) -- (3.3091,2.0339);
\draw[gray!65, line width=0.15pt] (3.2542,2.0678) -- (3.3091,2.0339);
\draw[gray!65, line width=0.15pt] (3.3091,5.9661) -- (3.2542,5.9322);
\draw[gray!65, line width=0.15pt] (3.3091,5.9661) -- (3.4576,6.0522);
\draw[gray!65, line width=0.15pt] (3.6610,1.8428) -- (3.5509,1.8983);
\draw[gray!65, line width=0.15pt] (3.4576,1.9478) -- (3.5509,1.8983);
\draw[gray!65, line width=0.15pt] (3.5509,6.1017) -- (3.4576,6.0522);
\draw[gray!65, line width=0.15pt] (3.5509,6.1017) -- (3.6610,6.1572);
\draw[gray!65, line width=0.15pt] (3.8644,1.7510) -- (3.8375,1.7627);
\draw[gray!65, line width=0.15pt] (3.6610,1.8428) -- (3.8375,1.7627);
\draw[gray!65, line width=0.15pt] (3.8375,6.2373) -- (3.6610,6.1572);
\draw[gray!65, line width=0.15pt] (3.8375,6.2373) -- (3.8644,6.2490);
\draw[gray!65, line width=0.15pt] (3.8644,1.7510) -- (4.0678,1.6706);
\draw[gray!65, line width=0.15pt] (4.0678,6.3294) -- (3.8644,6.2490);
\draw[gray!65, line width=0.15pt] (4.2712,1.6006) -- (4.1929,1.6271);
\draw[gray!65, line width=0.15pt] (4.0678,1.6706) -- (4.1929,1.6271);
\draw[gray!65, line width=0.15pt] (4.1929,6.3729) -- (4.0678,6.3294);
\draw[gray!65, line width=0.15pt] (4.1929,6.3729) -- (4.2712,6.3994);
\draw[gray!65, line width=0.15pt] (4.2712,1.6006) -- (4.4746,1.5400);
\draw[gray!65, line width=0.15pt] (4.4746,6.4600) -- (4.2712,6.3994);
\draw[gray!65, line width=0.15pt] (4.6780,1.4883) -- (4.6653,1.4915);
\draw[gray!65, line width=0.15pt] (4.4746,1.5400) -- (4.6653,1.4915);
\draw[gray!65, line width=0.15pt] (4.6653,6.5085) -- (4.4746,6.4600);
\draw[gray!65, line width=0.15pt] (4.6653,6.5085) -- (4.6780,6.5117);
\draw[gray!65, line width=0.15pt] (4.6780,1.4883) -- (4.8814,1.4446);
\draw[gray!65, line width=0.15pt] (4.8814,6.5554) -- (4.6780,6.5117);
\draw[gray!65, line width=0.15pt] (4.8814,1.4446) -- (5.0847,1.4088);
\draw[gray!65, line width=0.15pt] (5.0847,6.5912) -- (4.8814,6.5554);
\draw[gray!65, line width=0.15pt] (5.0847,1.4088) -- (5.2881,1.3807);
\draw[gray!65, line width=0.15pt] (5.2881,6.6193) -- (5.0847,6.5912);
\draw[gray!65, line width=0.15pt] (5.2881,1.3807) -- (5.4915,1.3599);
\draw[gray!65, line width=0.15pt] (5.4915,6.6401) -- (5.2881,6.6193);
\draw[gray!65, line width=0.15pt] (5.6949,1.3459) -- (5.5492,1.3559);
\draw[gray!65, line width=0.15pt] (5.4915,1.3599) -- (5.5492,1.3559);
\draw[gray!65, line width=0.15pt] (5.5492,6.6441) -- (5.4915,6.6401);
\draw[gray!65, line width=0.15pt] (5.5492,6.6441) -- (5.6949,6.6541);
\draw[gray!65, line width=0.15pt] (5.6949,1.3459) -- (5.8983,1.3389);
\draw[gray!65, line width=0.15pt] (5.8983,6.6611) -- (5.6949,6.6541);
\draw[gray!65, line width=0.15pt] (5.8983,1.3389) -- (6.1017,1.3389);
\draw[gray!65, line width=0.15pt] (6.1017,6.6611) -- (5.8983,6.6611);
\draw[gray!65, line width=0.15pt] (6.1017,1.3389) -- (6.3051,1.3459);
\draw[gray!65, line width=0.15pt] (6.3051,6.6541) -- (6.1017,6.6611);
\draw[gray!65, line width=0.15pt] (6.3051,1.3459) -- (6.4508,1.3559);
\draw[gray!65, line width=0.15pt] (6.4508,1.3559) -- (6.5085,1.3599);
\draw[gray!65, line width=0.15pt] (6.5085,6.6401) -- (6.4508,6.6441);
\draw[gray!65, line width=0.15pt] (6.3051,6.6541) -- (6.4508,6.6441);
\draw[gray!65, line width=0.15pt] (6.5085,1.3599) -- (6.7119,1.3807);
\draw[gray!65, line width=0.15pt] (6.7119,6.6193) -- (6.5085,6.6401);
\draw[gray!65, line width=0.15pt] (6.7119,1.3807) -- (6.9153,1.4088);
\draw[gray!65, line width=0.15pt] (6.9153,6.5912) -- (6.7119,6.6193);
\draw[gray!65, line width=0.15pt] (6.9153,1.4088) -- (7.1186,1.4446);
\draw[gray!65, line width=0.15pt] (7.1186,6.5554) -- (6.9153,6.5912);
\draw[gray!65, line width=0.15pt] (7.1186,1.4446) -- (7.3220,1.4883);
\draw[gray!65, line width=0.15pt] (7.3220,6.5117) -- (7.1186,6.5554);
\draw[gray!65, line width=0.15pt] (7.3220,1.4883) -- (7.3347,1.4915);
\draw[gray!65, line width=0.15pt] (7.3347,1.4915) -- (7.5254,1.5400);
\draw[gray!65, line width=0.15pt] (7.5254,6.4600) -- (7.3347,6.5085);
\draw[gray!65, line width=0.15pt] (7.3220,6.5117) -- (7.3347,6.5085);
\draw[gray!65, line width=0.15pt] (7.5254,1.5400) -- (7.7288,1.6006);
\draw[gray!65, line width=0.15pt] (7.7288,6.3994) -- (7.5254,6.4600);
\draw[gray!65, line width=0.15pt] (7.7288,1.6006) -- (7.8071,1.6271);
\draw[gray!65, line width=0.15pt] (7.8071,1.6271) -- (7.9322,1.6706);
\draw[gray!65, line width=0.15pt] (7.9322,6.3294) -- (7.8071,6.3729);
\draw[gray!65, line width=0.15pt] (7.7288,6.3994) -- (7.8071,6.3729);
\draw[gray!65, line width=0.15pt] (7.9322,1.6706) -- (8.1356,1.7510);
\draw[gray!65, line width=0.15pt] (8.1356,6.2490) -- (7.9322,6.3294);
\draw[gray!65, line width=0.15pt] (8.1356,1.7510) -- (8.1625,1.7627);
\draw[gray!65, line width=0.15pt] (8.1625,1.7627) -- (8.3390,1.8428);
\draw[gray!65, line width=0.15pt] (8.3390,6.1572) -- (8.1625,6.2373);
\draw[gray!65, line width=0.15pt] (8.1356,6.2490) -- (8.1625,6.2373);
\draw[gray!65, line width=0.15pt] (8.3390,1.8428) -- (8.4491,1.8983);
\draw[gray!65, line width=0.15pt] (8.4491,1.8983) -- (8.5424,1.9478);
\draw[gray!65, line width=0.15pt] (8.5424,6.0522) -- (8.4491,6.1017);
\draw[gray!65, line width=0.15pt] (8.3390,6.1572) -- (8.4491,6.1017);
\draw[gray!65, line width=0.15pt] (8.5424,1.9478) -- (8.6909,2.0339);
\draw[gray!65, line width=0.15pt] (8.6909,2.0339) -- (8.7458,2.0678);
\draw[gray!65, line width=0.15pt] (8.7458,5.9322) -- (8.6909,5.9661);
\draw[gray!65, line width=0.15pt] (8.5424,6.0522) -- (8.6909,5.9661);
\draw[gray!65, line width=0.15pt] (8.7458,2.0678) -- (8.8983,2.1695);
\draw[gray!65, line width=0.15pt] (8.8983,2.1695) -- (8.9492,2.2061);
\draw[gray!65, line width=0.15pt] (8.9492,5.7939) -- (8.8983,5.8305);
\draw[gray!65, line width=0.15pt] (8.7458,5.9322) -- (8.8983,5.8305);
\draw[gray!65, line width=0.15pt] (8.9492,2.2061) -- (9.0784,2.3051);
\draw[gray!65, line width=0.15pt] (9.0784,2.3051) -- (9.1525,2.3673);
\draw[gray!65, line width=0.15pt] (9.1525,5.6327) -- (9.0784,5.6949);
\draw[gray!65, line width=0.15pt] (8.9492,5.7939) -- (9.0784,5.6949);
\draw[gray!65, line width=0.15pt] (9.1525,2.3673) -- (9.2357,2.4407);
\draw[gray!65, line width=0.15pt] (9.2357,2.4407) -- (9.3559,2.5583);
\draw[gray!65, line width=0.15pt] (9.3559,5.4417) -- (9.2357,5.5593);
\draw[gray!65, line width=0.15pt] (9.1525,5.6327) -- (9.2357,5.5593);
\draw[gray!65, line width=0.15pt] (9.3559,2.5583) -- (9.3735,2.5763);
\draw[gray!65, line width=0.15pt] (9.3735,2.5763) -- (9.4941,2.7119);
\draw[gray!65, line width=0.15pt] (9.4941,2.7119) -- (9.5593,2.7952);
\draw[gray!65, line width=0.15pt] (9.5593,5.2048) -- (9.4941,5.2881);
\draw[gray!65, line width=0.15pt] (9.4941,5.2881) -- (9.3735,5.4237);
\draw[gray!65, line width=0.15pt] (9.3559,5.4417) -- (9.3735,5.4237);
\draw[gray!65, line width=0.15pt] (9.5593,2.7952) -- (9.5991,2.8475);
\draw[gray!65, line width=0.15pt] (9.5991,2.8475) -- (9.6901,2.9831);
\draw[gray!65, line width=0.15pt] (9.6901,2.9831) -- (9.7627,3.1102);
\draw[gray!65, line width=0.15pt] (9.7627,4.8898) -- (9.6901,5.0169);
\draw[gray!65, line width=0.15pt] (9.6901,5.0169) -- (9.5991,5.1525);
\draw[gray!65, line width=0.15pt] (9.5593,5.2048) -- (9.5991,5.1525);
\draw[gray!65, line width=0.15pt] (9.7627,3.1102) -- (9.7675,3.1186);
\draw[gray!65, line width=0.15pt] (9.7675,3.1186) -- (9.8331,3.2542);
\draw[gray!65, line width=0.15pt] (9.8331,3.2542) -- (9.8868,3.3898);
\draw[gray!65, line width=0.15pt] (9.8868,3.3898) -- (9.9290,3.5254);
\draw[gray!65, line width=0.15pt] (9.9290,3.5254) -- (9.9602,3.6610);
\draw[gray!65, line width=0.15pt] (9.9602,3.6610) -- (9.9661,3.6995);
\draw[gray!65, line width=0.15pt] (9.9661,4.3005) -- (9.9602,4.3390);
\draw[gray!65, line width=0.15pt] (9.9602,4.3390) -- (9.9290,4.4746);
\draw[gray!65, line width=0.15pt] (9.9290,4.4746) -- (9.8868,4.6102);
\draw[gray!65, line width=0.15pt] (9.8868,4.6102) -- (9.8331,4.7458);
\draw[gray!65, line width=0.15pt] (9.8331,4.7458) -- (9.7675,4.8814);
\draw[gray!65, line width=0.15pt] (9.7627,4.8898) -- (9.7675,4.8814);
\draw[gray!65, line width=0.15pt] (9.9661,3.6995) -- (9.9811,3.7966);
\draw[gray!65, line width=0.15pt] (9.9811,3.7966) -- (9.9916,3.9322);
\draw[gray!65, line width=0.15pt] (9.9916,3.9322) -- (9.9916,4.0678);
\draw[gray!65, line width=0.15pt] (9.9916,4.0678) -- (9.9811,4.2034);
\draw[gray!65, line width=0.15pt] (9.9661,4.3005) -- (9.9811,4.2034);
\draw[gray!75, line width=0.15pt] (0.8136,3.3966) -- (0.7831,3.5254);
\draw[gray!75, line width=0.15pt] (0.7831,3.5254) -- (0.7592,3.6610);
\draw[gray!75, line width=0.15pt] (0.7592,3.6610) -- (0.7433,3.7966);
\draw[gray!75, line width=0.15pt] (0.7433,3.7966) -- (0.7355,3.9322);
\draw[gray!75, line width=0.15pt] (0.7355,3.9322) -- (0.7355,4.0678);
\draw[gray!75, line width=0.15pt] (0.7355,4.0678) -- (0.7433,4.2034);
\draw[gray!75, line width=0.15pt] (0.7433,4.2034) -- (0.7592,4.3390);
\draw[gray!75, line width=0.15pt] (0.7592,4.3390) -- (0.7831,4.4746);
\draw[gray!75, line width=0.15pt] (0.7831,4.4746) -- (0.8136,4.6034);
\draw[gray!75, line width=0.15pt] (1.0169,2.8676) -- (0.9605,2.9831);
\draw[gray!75, line width=0.15pt] (0.9605,2.9831) -- (0.9032,3.1186);
\draw[gray!75, line width=0.15pt] (0.9032,3.1186) -- (0.8549,3.2542);
\draw[gray!75, line width=0.15pt] (0.8549,3.2542) -- (0.8151,3.3898);
\draw[gray!75, line width=0.15pt] (0.8136,3.3966) -- (0.8151,3.3898);
\draw[gray!75, line width=0.15pt] (0.8151,4.6102) -- (0.8136,4.6034);
\draw[gray!75, line width=0.15pt] (0.8151,4.6102) -- (0.8549,4.7458);
\draw[gray!75, line width=0.15pt] (0.8549,4.7458) -- (0.9032,4.8814);
\draw[gray!75, line width=0.15pt] (0.9032,4.8814) -- (0.9605,5.0169);
\draw[gray!75, line width=0.15pt] (0.9605,5.0169) -- (1.0169,5.1324);
\draw[gray!75, line width=0.15pt] (1.2203,2.5290) -- (1.1874,2.5763);
\draw[gray!75, line width=0.15pt] (1.1874,2.5763) -- (1.1020,2.7119);
\draw[gray!75, line width=0.15pt] (1.1020,2.7119) -- (1.0268,2.8475);
\draw[gray!75, line width=0.15pt] (1.0169,2.8676) -- (1.0268,2.8475);
\draw[gray!75, line width=0.15pt] (1.0268,5.1525) -- (1.0169,5.1324);
\draw[gray!75, line width=0.15pt] (1.0268,5.1525) -- (1.1020,5.2881);
\draw[gray!75, line width=0.15pt] (1.1020,5.2881) -- (1.1874,5.4237);
\draw[gray!75, line width=0.15pt] (1.1874,5.4237) -- (1.2203,5.4710);
\draw[gray!75, line width=0.15pt] (1.4237,2.2650) -- (1.3894,2.3051);
\draw[gray!75, line width=0.15pt] (1.3894,2.3051) -- (1.2828,2.4407);
\draw[gray!75, line width=0.15pt] (1.2203,2.5290) -- (1.2828,2.4407);
\draw[gray!75, line width=0.15pt] (1.2828,5.5593) -- (1.2203,5.4710);
\draw[gray!75, line width=0.15pt] (1.2828,5.5593) -- (1.3894,5.6949);
\draw[gray!75, line width=0.15pt] (1.3894,5.6949) -- (1.4237,5.7350);
\draw[gray!75, line width=0.15pt] (1.6271,2.0453) -- (1.5075,2.1695);
\draw[gray!75, line width=0.15pt] (1.4237,2.2650) -- (1.5075,2.1695);
\draw[gray!75, line width=0.15pt] (1.5075,5.8305) -- (1.4237,5.7350);
\draw[gray!75, line width=0.15pt] (1.5075,5.8305) -- (1.6271,5.9547);
\draw[gray!75, line width=0.15pt] (1.8305,1.8571) -- (1.7831,1.8983);
\draw[gray!75, line width=0.15pt] (1.7831,1.8983) -- (1.6385,2.0339);
\draw[gray!75, line width=0.15pt] (1.6271,2.0453) -- (1.6385,2.0339);
\draw[gray!75, line width=0.15pt] (1.6385,5.9661) -- (1.6271,5.9547);
\draw[gray!75, line width=0.15pt] (1.6385,5.9661) -- (1.7831,6.1017);
\draw[gray!75, line width=0.15pt] (1.7831,6.1017) -- (1.8305,6.1429);
\draw[gray!75, line width=0.15pt] (2.0339,1.6918) -- (1.9432,1.7627);
\draw[gray!75, line width=0.15pt] (1.8305,1.8571) -- (1.9432,1.7627);
\draw[gray!75, line width=0.15pt] (1.9432,6.2373) -- (1.8305,6.1429);
\draw[gray!75, line width=0.15pt] (1.9432,6.2373) -- (2.0339,6.3082);
\draw[gray!75, line width=0.15pt] (2.2373,1.5451) -- (2.1206,1.6271);
\draw[gray!75, line width=0.15pt] (2.0339,1.6918) -- (2.1206,1.6271);
\draw[gray!75, line width=0.15pt] (2.1206,6.3729) -- (2.0339,6.3082);
\draw[gray!75, line width=0.15pt] (2.1206,6.3729) -- (2.2373,6.4549);
\draw[gray!75, line width=0.15pt] (2.4407,1.4137) -- (2.3176,1.4915);
\draw[gray!75, line width=0.15pt] (2.2373,1.5451) -- (2.3176,1.4915);
\draw[gray!75, line width=0.15pt] (2.3176,6.5085) -- (2.2373,6.4549);
\draw[gray!75, line width=0.15pt] (2.3176,6.5085) -- (2.4407,6.5863);
\draw[gray!75, line width=0.15pt] (2.6441,1.2955) -- (2.5378,1.3559);
\draw[gray!75, line width=0.15pt] (2.4407,1.4137) -- (2.5378,1.3559);
\draw[gray!75, line width=0.15pt] (2.5378,6.6441) -- (2.4407,6.5863);
\draw[gray!75, line width=0.15pt] (2.5378,6.6441) -- (2.6441,6.7045);
\draw[gray!75, line width=0.15pt] (2.8475,1.1888) -- (2.7856,1.2203);
\draw[gray!75, line width=0.15pt] (2.6441,1.2955) -- (2.7856,1.2203);
\draw[gray!75, line width=0.15pt] (2.7856,6.7797) -- (2.6441,6.7045);
\draw[gray!75, line width=0.15pt] (2.7856,6.7797) -- (2.8475,6.8112);
\draw[gray!75, line width=0.15pt] (2.8475,1.1888) -- (3.0508,1.0923);
\draw[gray!75, line width=0.15pt] (3.0508,6.9077) -- (2.8475,6.8112);
\draw[gray!75, line width=0.15pt] (3.2542,1.0050) -- (3.0680,1.0847);
\draw[gray!75, line width=0.15pt] (3.0508,1.0923) -- (3.0680,1.0847);
\draw[gray!75, line width=0.15pt] (3.0680,6.9153) -- (3.0508,6.9077);
\draw[gray!75, line width=0.15pt] (3.0680,6.9153) -- (3.2542,6.9950);
\draw[gray!75, line width=0.15pt] (3.4576,0.9263) -- (3.3975,0.9492);
\draw[gray!75, line width=0.15pt] (3.2542,1.0050) -- (3.3975,0.9492);
\draw[gray!75, line width=0.15pt] (3.3975,7.0508) -- (3.2542,6.9950);
\draw[gray!75, line width=0.15pt] (3.3975,7.0508) -- (3.4576,7.0737);
\draw[gray!75, line width=0.15pt] (3.4576,0.9263) -- (3.6610,0.8552);
\draw[gray!75, line width=0.15pt] (3.6610,7.1448) -- (3.4576,7.0737);
\draw[gray!75, line width=0.15pt] (3.8644,0.7916) -- (3.7936,0.8136);
\draw[gray!75, line width=0.15pt] (3.6610,0.8552) -- (3.7936,0.8136);
\draw[gray!75, line width=0.15pt] (3.7936,7.1864) -- (3.6610,7.1448);
\draw[gray!75, line width=0.15pt] (3.7936,7.1864) -- (3.8644,7.2084);
\draw[gray!75, line width=0.15pt] (3.8644,0.7916) -- (4.0678,0.7347);
\draw[gray!75, line width=0.15pt] (4.0678,7.2653) -- (3.8644,7.2084);
\draw[gray!75, line width=0.15pt] (4.0678,0.7347) -- (4.2712,0.6846);
\draw[gray!75, line width=0.15pt] (4.2712,7.3154) -- (4.0678,7.2653);
\draw[gray!75, line width=0.15pt] (4.4746,0.6403) -- (4.3014,0.6780);
\draw[gray!75, line width=0.15pt] (4.2712,0.6846) -- (4.3014,0.6780);
\draw[gray!75, line width=0.15pt] (4.3014,7.3220) -- (4.2712,7.3154);
\draw[gray!75, line width=0.15pt] (4.3014,7.3220) -- (4.4746,7.3597);
\draw[gray!75, line width=0.15pt] (4.4746,0.6403) -- (4.6780,0.6021);
\draw[gray!75, line width=0.15pt] (4.6780,7.3979) -- (4.4746,7.3597);
\draw[gray!75, line width=0.15pt] (4.6780,0.6021) -- (4.8814,0.5699);
\draw[gray!75, line width=0.15pt] (4.8814,7.4301) -- (4.6780,7.3979);
\draw[gray!75, line width=0.15pt] (4.8814,0.5699) -- (5.0847,0.5434);
\draw[gray!75, line width=0.15pt] (5.0847,7.4566) -- (4.8814,7.4301);
\draw[gray!75, line width=0.15pt] (5.2881,0.5220) -- (5.0949,0.5424);
\draw[gray!75, line width=0.15pt] (5.0847,0.5434) -- (5.0949,0.5424);
\draw[gray!75, line width=0.15pt] (5.0949,7.4576) -- (5.0847,7.4566);
\draw[gray!75, line width=0.15pt] (5.0949,7.4576) -- (5.2881,7.4780);
\draw[gray!75, line width=0.15pt] (5.2881,0.5220) -- (5.4915,0.5061);
\draw[gray!75, line width=0.15pt] (5.4915,7.4939) -- (5.2881,7.4780);
\draw[gray!75, line width=0.15pt] (5.4915,0.5061) -- (5.6949,0.4956);
\draw[gray!75, line width=0.15pt] (5.6949,7.5044) -- (5.4915,7.4939);
\draw[gray!75, line width=0.15pt] (5.6949,0.4956) -- (5.8983,0.4903);
\draw[gray!75, line width=0.15pt] (5.8983,7.5097) -- (5.6949,7.5044);
\draw[gray!75, line width=0.15pt] (5.8983,0.4903) -- (6.1017,0.4903);
\draw[gray!75, line width=0.15pt] (6.1017,7.5097) -- (5.8983,7.5097);
\draw[gray!75, line width=0.15pt] (6.1017,0.4903) -- (6.3051,0.4956);
\draw[gray!75, line width=0.15pt] (6.3051,7.5044) -- (6.1017,7.5097);
\draw[gray!75, line width=0.15pt] (6.3051,0.4956) -- (6.5085,0.5061);
\draw[gray!75, line width=0.15pt] (6.5085,7.4939) -- (6.3051,7.5044);
\draw[gray!75, line width=0.15pt] (6.5085,0.5061) -- (6.7119,0.5220);
\draw[gray!75, line width=0.15pt] (6.7119,7.4780) -- (6.5085,7.4939);
\draw[gray!75, line width=0.15pt] (6.7119,0.5220) -- (6.9051,0.5424);
\draw[gray!75, line width=0.15pt] (6.9051,0.5424) -- (6.9153,0.5434);
\draw[gray!75, line width=0.15pt] (6.9153,7.4566) -- (6.9051,7.4576);
\draw[gray!75, line width=0.15pt] (6.7119,7.4780) -- (6.9051,7.4576);
\draw[gray!75, line width=0.15pt] (6.9153,0.5434) -- (7.1186,0.5699);
\draw[gray!75, line width=0.15pt] (7.1186,7.4301) -- (6.9153,7.4566);
\draw[gray!75, line width=0.15pt] (7.1186,0.5699) -- (7.3220,0.6021);
\draw[gray!75, line width=0.15pt] (7.3220,7.3979) -- (7.1186,7.4301);
\draw[gray!75, line width=0.15pt] (7.3220,0.6021) -- (7.5254,0.6403);
\draw[gray!75, line width=0.15pt] (7.5254,7.3597) -- (7.3220,7.3979);
\draw[gray!75, line width=0.15pt] (7.5254,0.6403) -- (7.6986,0.6780);
\draw[gray!75, line width=0.15pt] (7.6986,0.6780) -- (7.7288,0.6846);
\draw[gray!75, line width=0.15pt] (7.7288,7.3154) -- (7.6986,7.3220);
\draw[gray!75, line width=0.15pt] (7.5254,7.3597) -- (7.6986,7.3220);
\draw[gray!75, line width=0.15pt] (7.7288,0.6846) -- (7.9322,0.7347);
\draw[gray!75, line width=0.15pt] (7.9322,7.2653) -- (7.7288,7.3154);
\draw[gray!75, line width=0.15pt] (7.9322,0.7347) -- (8.1356,0.7916);
\draw[gray!75, line width=0.15pt] (8.1356,7.2084) -- (7.9322,7.2653);
\draw[gray!75, line width=0.15pt] (8.1356,0.7916) -- (8.2064,0.8136);
\draw[gray!75, line width=0.15pt] (8.2064,0.8136) -- (8.3390,0.8552);
\draw[gray!75, line width=0.15pt] (8.3390,7.1448) -- (8.2064,7.1864);
\draw[gray!75, line width=0.15pt] (8.1356,7.2084) -- (8.2064,7.1864);
\draw[gray!75, line width=0.15pt] (8.3390,0.8552) -- (8.5424,0.9263);
\draw[gray!75, line width=0.15pt] (8.5424,7.0737) -- (8.3390,7.1448);
\draw[gray!75, line width=0.15pt] (8.5424,0.9263) -- (8.6025,0.9492);
\draw[gray!75, line width=0.15pt] (8.6025,0.9492) -- (8.7458,1.0050);
\draw[gray!75, line width=0.15pt] (8.7458,6.9950) -- (8.6025,7.0508);
\draw[gray!75, line width=0.15pt] (8.5424,7.0737) -- (8.6025,7.0508);
\draw[gray!75, line width=0.15pt] (8.7458,1.0050) -- (8.9320,1.0847);
\draw[gray!75, line width=0.15pt] (8.9320,1.0847) -- (8.9492,1.0923);
\draw[gray!75, line width=0.15pt] (8.9492,6.9077) -- (8.9320,6.9153);
\draw[gray!75, line width=0.15pt] (8.7458,6.9950) -- (8.9320,6.9153);
\draw[gray!75, line width=0.15pt] (8.9492,1.0923) -- (9.1525,1.1888);
\draw[gray!75, line width=0.15pt] (9.1525,6.8112) -- (8.9492,6.9077);
\draw[gray!75, line width=0.15pt] (9.1525,1.1888) -- (9.2144,1.2203);
\draw[gray!75, line width=0.15pt] (9.2144,1.2203) -- (9.3559,1.2955);
\draw[gray!75, line width=0.15pt] (9.3559,6.7045) -- (9.2144,6.7797);
\draw[gray!75, line width=0.15pt] (9.1525,6.8112) -- (9.2144,6.7797);
\draw[gray!75, line width=0.15pt] (9.3559,1.2955) -- (9.4622,1.3559);
\draw[gray!75, line width=0.15pt] (9.4622,1.3559) -- (9.5593,1.4137);
\draw[gray!75, line width=0.15pt] (9.5593,6.5863) -- (9.4622,6.6441);
\draw[gray!75, line width=0.15pt] (9.3559,6.7045) -- (9.4622,6.6441);
\draw[gray!75, line width=0.15pt] (9.5593,1.4137) -- (9.6824,1.4915);
\draw[gray!75, line width=0.15pt] (9.6824,1.4915) -- (9.7627,1.5451);
\draw[gray!75, line width=0.15pt] (9.7627,6.4549) -- (9.6824,6.5085);
\draw[gray!75, line width=0.15pt] (9.5593,6.5863) -- (9.6824,6.5085);
\draw[gray!75, line width=0.15pt] (9.7627,1.5451) -- (9.8794,1.6271);
\draw[gray!75, line width=0.15pt] (9.8794,1.6271) -- (9.9661,1.6918);
\draw[gray!75, line width=0.15pt] (9.9661,6.3082) -- (9.8794,6.3729);
\draw[gray!75, line width=0.15pt] (9.7627,6.4549) -- (9.8794,6.3729);
\draw[gray!75, line width=0.15pt] (9.9661,1.6918) -- (10.0568,1.7627);
\draw[gray!75, line width=0.15pt] (10.0568,1.7627) -- (10.1695,1.8571);
\draw[gray!75, line width=0.15pt] (10.1695,6.1429) -- (10.0568,6.2373);
\draw[gray!75, line width=0.15pt] (9.9661,6.3082) -- (10.0568,6.2373);
\draw[gray!75, line width=0.15pt] (10.1695,1.8571) -- (10.2169,1.8983);
\draw[gray!75, line width=0.15pt] (10.2169,1.8983) -- (10.3615,2.0339);
\draw[gray!75, line width=0.15pt] (10.3615,2.0339) -- (10.3729,2.0453);
\draw[gray!75, line width=0.15pt] (10.3729,5.9547) -- (10.3615,5.9661);
\draw[gray!75, line width=0.15pt] (10.3615,5.9661) -- (10.2169,6.1017);
\draw[gray!75, line width=0.15pt] (10.1695,6.1429) -- (10.2169,6.1017);
\draw[gray!75, line width=0.15pt] (10.3729,2.0453) -- (10.4925,2.1695);
\draw[gray!75, line width=0.15pt] (10.4925,2.1695) -- (10.5763,2.2650);
\draw[gray!75, line width=0.15pt] (10.5763,5.7350) -- (10.4925,5.8305);
\draw[gray!75, line width=0.15pt] (10.3729,5.9547) -- (10.4925,5.8305);
\draw[gray!75, line width=0.15pt] (10.5763,2.2650) -- (10.6106,2.3051);
\draw[gray!75, line width=0.15pt] (10.6106,2.3051) -- (10.7172,2.4407);
\draw[gray!75, line width=0.15pt] (10.7172,2.4407) -- (10.7797,2.5290);
\draw[gray!75, line width=0.15pt] (10.7797,5.4710) -- (10.7172,5.5593);
\draw[gray!75, line width=0.15pt] (10.7172,5.5593) -- (10.6106,5.6949);
\draw[gray!75, line width=0.15pt] (10.5763,5.7350) -- (10.6106,5.6949);
\draw[gray!75, line width=0.15pt] (10.7797,2.5290) -- (10.8126,2.5763);
\draw[gray!75, line width=0.15pt] (10.8126,2.5763) -- (10.8980,2.7119);
\draw[gray!75, line width=0.15pt] (10.8980,2.7119) -- (10.9732,2.8475);
\draw[gray!75, line width=0.15pt] (10.9732,2.8475) -- (10.9831,2.8676);
\draw[gray!75, line width=0.15pt] (10.9831,5.1324) -- (10.9732,5.1525);
\draw[gray!75, line width=0.15pt] (10.9732,5.1525) -- (10.8980,5.2881);
\draw[gray!75, line width=0.15pt] (10.8980,5.2881) -- (10.8126,5.4237);
\draw[gray!75, line width=0.15pt] (10.7797,5.4710) -- (10.8126,5.4237);
\draw[gray!75, line width=0.15pt] (10.9831,2.8676) -- (11.0395,2.9831);
\draw[gray!75, line width=0.15pt] (11.0395,2.9831) -- (11.0968,3.1186);
\draw[gray!75, line width=0.15pt] (11.0968,3.1186) -- (11.1451,3.2542);
\draw[gray!75, line width=0.15pt] (11.1451,3.2542) -- (11.1849,3.3898);
\draw[gray!75, line width=0.15pt] (11.1849,3.3898) -- (11.1864,3.3966);
\draw[gray!75, line width=0.15pt] (11.1864,4.6034) -- (11.1849,4.6102);
\draw[gray!75, line width=0.15pt] (11.1849,4.6102) -- (11.1451,4.7458);
\draw[gray!75, line width=0.15pt] (11.1451,4.7458) -- (11.0968,4.8814);
\draw[gray!75, line width=0.15pt] (11.0968,4.8814) -- (11.0395,5.0169);
\draw[gray!75, line width=0.15pt] (10.9831,5.1324) -- (11.0395,5.0169);
\draw[gray!75, line width=0.15pt] (11.1864,3.3966) -- (11.2169,3.5254);
\draw[gray!75, line width=0.15pt] (11.2169,3.5254) -- (11.2408,3.6610);
\draw[gray!75, line width=0.15pt] (11.2408,3.6610) -- (11.2567,3.7966);
\draw[gray!75, line width=0.15pt] (11.2567,3.7966) -- (11.2645,3.9322);
\draw[gray!75, line width=0.15pt] (11.2645,3.9322) -- (11.2645,4.0678);
\draw[gray!75, line width=0.15pt] (11.2645,4.0678) -- (11.2567,4.2034);
\draw[gray!75, line width=0.15pt] (11.2567,4.2034) -- (11.2408,4.3390);
\draw[gray!75, line width=0.15pt] (11.2408,4.3390) -- (11.2169,4.4746);
\draw[gray!75, line width=0.15pt] (11.1864,4.6034) -- (11.2169,4.4746);
\draw[gray!85, line width=0.15pt] (0.2034,1.7110) -- (0.1575,1.7627);
\draw[gray!85, line width=0.15pt] (0.1575,1.7627) -- (0.0453,1.8983);
\draw[gray!85, line width=0.15pt] (0.0000,1.9573) -- (0.0453,1.8983);
\draw[gray!85, line width=0.15pt] (0.0453,6.1017) -- (0.0000,6.0427);
\draw[gray!85, line width=0.15pt] (0.0453,6.1017) -- (0.1575,6.2373);
\draw[gray!85, line width=0.15pt] (0.1575,6.2373) -- (0.2034,6.2890);
\draw[gray!85, line width=0.15pt] (0.4068,1.4958) -- (0.2792,1.6271);
\draw[gray!85, line width=0.15pt] (0.2034,1.7110) -- (0.2792,1.6271);
\draw[gray!85, line width=0.15pt] (0.2792,6.3729) -- (0.2034,6.2890);
\draw[gray!85, line width=0.15pt] (0.2792,6.3729) -- (0.4068,6.5042);
\draw[gray!85, line width=0.15pt] (0.6102,1.3051) -- (0.5534,1.3559);
\draw[gray!85, line width=0.15pt] (0.5534,1.3559) -- (0.4110,1.4915);
\draw[gray!85, line width=0.15pt] (0.4068,1.4958) -- (0.4110,1.4915);
\draw[gray!85, line width=0.15pt] (0.4110,6.5085) -- (0.4068,6.5042);
\draw[gray!85, line width=0.15pt] (0.4110,6.5085) -- (0.5534,6.6441);
\draw[gray!85, line width=0.15pt] (0.5534,6.6441) -- (0.6102,6.6949);
\draw[gray!85, line width=0.15pt] (0.8136,1.1331) -- (0.7075,1.2203);
\draw[gray!85, line width=0.15pt] (0.6102,1.3051) -- (0.7075,1.2203);
\draw[gray!85, line width=0.15pt] (0.7075,6.7797) -- (0.6102,6.6949);
\draw[gray!85, line width=0.15pt] (0.7075,6.7797) -- (0.8136,6.8669);
\draw[gray!85, line width=0.15pt] (1.0169,0.9770) -- (0.8743,1.0847);
\draw[gray!85, line width=0.15pt] (0.8136,1.1331) -- (0.8743,1.0847);
\draw[gray!85, line width=0.15pt] (0.8743,6.9153) -- (0.8136,6.8669);
\draw[gray!85, line width=0.15pt] (0.8743,6.9153) -- (1.0169,7.0230);
\draw[gray!85, line width=0.15pt] (1.2203,0.8343) -- (1.0551,0.9492);
\draw[gray!85, line width=0.15pt] (1.0169,0.9770) -- (1.0551,0.9492);
\draw[gray!85, line width=0.15pt] (1.0551,7.0508) -- (1.0169,7.0230);
\draw[gray!85, line width=0.15pt] (1.0551,7.0508) -- (1.2203,7.1657);
\draw[gray!85, line width=0.15pt] (1.4237,0.7034) -- (1.2515,0.8136);
\draw[gray!85, line width=0.15pt] (1.2203,0.8343) -- (1.2515,0.8136);
\draw[gray!85, line width=0.15pt] (1.2515,7.1864) -- (1.2203,7.1657);
\draw[gray!85, line width=0.15pt] (1.2515,7.1864) -- (1.4237,7.2966);
\draw[gray!85, line width=0.15pt] (1.6271,0.5829) -- (1.4654,0.6780);
\draw[gray!85, line width=0.15pt] (1.4237,0.7034) -- (1.4654,0.6780);
\draw[gray!85, line width=0.15pt] (1.4654,7.3220) -- (1.4237,7.2966);
\draw[gray!85, line width=0.15pt] (1.4654,7.3220) -- (1.6271,7.4171);
\draw[gray!85, line width=0.15pt] (1.8305,0.4717) -- (1.6997,0.5424);
\draw[gray!85, line width=0.15pt] (1.6271,0.5829) -- (1.6997,0.5424);
\draw[gray!85, line width=0.15pt] (1.6997,7.4576) -- (1.6271,7.4171);
\draw[gray!85, line width=0.15pt] (1.6997,7.4576) -- (1.8305,7.5283);
\draw[gray!85, line width=0.15pt] (2.0339,0.3689) -- (1.9576,0.4068);
\draw[gray!85, line width=0.15pt] (1.8305,0.4717) -- (1.9576,0.4068);
\draw[gray!85, line width=0.15pt] (1.9576,7.5932) -- (1.8305,7.5283);
\draw[gray!85, line width=0.15pt] (1.9576,7.5932) -- (2.0339,7.6311);
\draw[gray!85, line width=0.15pt] (2.0339,0.3689) -- (2.2373,0.2740);
\draw[gray!85, line width=0.15pt] (2.2373,7.7260) -- (2.0339,7.6311);
\draw[gray!85, line width=0.15pt] (2.4407,0.1861) -- (2.2437,0.2712);
\draw[gray!85, line width=0.15pt] (2.2373,0.2740) -- (2.2437,0.2712);
\draw[gray!85, line width=0.15pt] (2.2437,7.7288) -- (2.2373,7.7260);
\draw[gray!85, line width=0.15pt] (2.2437,7.7288) -- (2.4407,7.8139);
\draw[gray!85, line width=0.15pt] (2.6441,0.1050) -- (2.5666,0.1356);
\draw[gray!85, line width=0.15pt] (2.4407,0.1861) -- (2.5666,0.1356);
\draw[gray!85, line width=0.15pt] (2.5666,7.8644) -- (2.4407,7.8139);
\draw[gray!85, line width=0.15pt] (2.5666,7.8644) -- (2.6441,7.8950);
\draw[gray!85, line width=0.15pt] (2.6441,0.1050) -- (2.8475,0.0302);
\draw[gray!85, line width=0.15pt] (2.8475,7.9698) -- (2.6441,7.8950);
\draw[gray!85, line width=0.15pt] (2.8475,0.0302) -- (2.9360,0.0000);
\draw[gray!85, line width=0.15pt] (2.9360,8.0000) -- (2.8475,7.9698);
\draw[gray!85, line width=0.15pt] (9.0640,0.0000) -- (9.1525,0.0302);
\draw[gray!85, line width=0.15pt] (9.1525,7.9698) -- (9.0640,8.0000);
\draw[gray!85, line width=0.15pt] (9.1525,0.0302) -- (9.3559,0.1050);
\draw[gray!85, line width=0.15pt] (9.3559,7.8950) -- (9.1525,7.9698);
\draw[gray!85, line width=0.15pt] (9.3559,0.1050) -- (9.4334,0.1356);
\draw[gray!85, line width=0.15pt] (9.4334,0.1356) -- (9.5593,0.1861);
\draw[gray!85, line width=0.15pt] (9.5593,7.8139) -- (9.4334,7.8644);
\draw[gray!85, line width=0.15pt] (9.3559,7.8950) -- (9.4334,7.8644);
\draw[gray!85, line width=0.15pt] (9.5593,0.1861) -- (9.7563,0.2712);
\draw[gray!85, line width=0.15pt] (9.7563,0.2712) -- (9.7627,0.2740);
\draw[gray!85, line width=0.15pt] (9.7627,7.7260) -- (9.7563,7.7288);
\draw[gray!85, line width=0.15pt] (9.5593,7.8139) -- (9.7563,7.7288);
\draw[gray!85, line width=0.15pt] (9.7627,0.2740) -- (9.9661,0.3689);
\draw[gray!85, line width=0.15pt] (9.9661,7.6311) -- (9.7627,7.7260);
\draw[gray!85, line width=0.15pt] (9.9661,0.3689) -- (10.0424,0.4068);
\draw[gray!85, line width=0.15pt] (10.0424,0.4068) -- (10.1695,0.4717);
\draw[gray!85, line width=0.15pt] (10.1695,7.5283) -- (10.0424,7.5932);
\draw[gray!85, line width=0.15pt] (9.9661,7.6311) -- (10.0424,7.5932);
\draw[gray!85, line width=0.15pt] (10.1695,0.4717) -- (10.3003,0.5424);
\draw[gray!85, line width=0.15pt] (10.3003,0.5424) -- (10.3729,0.5829);
\draw[gray!85, line width=0.15pt] (10.3729,7.4171) -- (10.3003,7.4576);
\draw[gray!85, line width=0.15pt] (10.1695,7.5283) -- (10.3003,7.4576);
\draw[gray!85, line width=0.15pt] (10.3729,0.5829) -- (10.5346,0.6780);
\draw[gray!85, line width=0.15pt] (10.5346,0.6780) -- (10.5763,0.7034);
\draw[gray!85, line width=0.15pt] (10.5763,7.2966) -- (10.5346,7.3220);
\draw[gray!85, line width=0.15pt] (10.3729,7.4171) -- (10.5346,7.3220);
\draw[gray!85, line width=0.15pt] (10.5763,0.7034) -- (10.7485,0.8136);
\draw[gray!85, line width=0.15pt] (10.7485,0.8136) -- (10.7797,0.8343);
\draw[gray!85, line width=0.15pt] (10.7797,7.1657) -- (10.7485,7.1864);
\draw[gray!85, line width=0.15pt] (10.5763,7.2966) -- (10.7485,7.1864);
\draw[gray!85, line width=0.15pt] (10.7797,0.8343) -- (10.9449,0.9492);
\draw[gray!85, line width=0.15pt] (10.9449,0.9492) -- (10.9831,0.9770);
\draw[gray!85, line width=0.15pt] (10.9831,7.0230) -- (10.9449,7.0508);
\draw[gray!85, line width=0.15pt] (10.7797,7.1657) -- (10.9449,7.0508);
\draw[gray!85, line width=0.15pt] (10.9831,0.9770) -- (11.1257,1.0847);
\draw[gray!85, line width=0.15pt] (11.1257,1.0847) -- (11.1864,1.1331);
\draw[gray!85, line width=0.15pt] (11.1864,6.8669) -- (11.1257,6.9153);
\draw[gray!85, line width=0.15pt] (10.9831,7.0230) -- (11.1257,6.9153);
\draw[gray!85, line width=0.15pt] (11.1864,1.1331) -- (11.2925,1.2203);
\draw[gray!85, line width=0.15pt] (11.2925,1.2203) -- (11.3898,1.3051);
\draw[gray!85, line width=0.15pt] (11.3898,6.6949) -- (11.2925,6.7797);
\draw[gray!85, line width=0.15pt] (11.1864,6.8669) -- (11.2925,6.7797);
\draw[gray!85, line width=0.15pt] (11.3898,1.3051) -- (11.4466,1.3559);
\draw[gray!85, line width=0.15pt] (11.4466,1.3559) -- (11.5890,1.4915);
\draw[gray!85, line width=0.15pt] (11.5890,1.4915) -- (11.5932,1.4958);
\draw[gray!85, line width=0.15pt] (11.5932,6.5042) -- (11.5890,6.5085);
\draw[gray!85, line width=0.15pt] (11.5890,6.5085) -- (11.4466,6.6441);
\draw[gray!85, line width=0.15pt] (11.3898,6.6949) -- (11.4466,6.6441);
\draw[gray!85, line width=0.15pt] (11.5932,1.4958) -- (11.7208,1.6271);
\draw[gray!85, line width=0.15pt] (11.7208,1.6271) -- (11.7966,1.7110);
\draw[gray!85, line width=0.15pt] (11.7966,6.2890) -- (11.7208,6.3729);
\draw[gray!85, line width=0.15pt] (11.5932,6.5042) -- (11.7208,6.3729);
\draw[gray!85, line width=0.15pt] (11.7966,1.7110) -- (11.8425,1.7627);
\draw[gray!85, line width=0.15pt] (11.8425,1.7627) -- (11.9547,1.8983);
\draw[gray!85, line width=0.15pt] (11.9547,1.8983) -- (12.0000,1.9573);
\draw[gray!85, line width=0.15pt] (12.0000,6.0427) -- (11.9547,6.1017);
\draw[gray!85, line width=0.15pt] (11.9547,6.1017) -- (11.8425,6.2373);
\draw[gray!85, line width=0.15pt] (11.7966,6.2890) -- (11.8425,6.2373);
\draw[blue!80!black, line width=1.0pt] (5.8125,3.8750) -- (6.0703,4.0468) -- (5.9391,3.9594) -- (6.0048,4.0032) -- (5.9965,3.9977) -- (6.0006,4.0004) -- (5.9996,3.9997) -- (6.0001,4.0001) -- (6.0000,4.0000) -- (6.0000,4.0000) -- (6.0000,4.0000) -- (6.0000,4.0000) -- (6.0000,4.0000) -- (6.0000,4.0000) -- (6.0000,4.0000) -- (6.0000,4.0000);
\fill[blue!80!black] (5.8125,3.8750) circle (1.2pt);
\fill[blue!80!black] (6.0703,4.0468) circle (1.2pt);
\fill[blue!80!black] (5.9391,3.9594) circle (1.2pt);
\fill[blue!80!black] (6.0048,4.0032) circle (1.2pt);
\fill[blue!80!black] (5.9965,3.9977) circle (1.2pt);
\fill[blue!80!black] (6.0006,4.0004) circle (1.2pt);
\fill[blue!80!black] (5.9996,3.9997) circle (1.2pt);
\fill[blue!80!black] (6.0001,4.0001) circle (1.2pt);
\fill[blue!80!black] (6.0000,4.0000) circle (1.2pt);
\fill[blue!80!black] (6.0000,4.0000) circle (1.2pt);
\fill[blue!80!black] (6.0000,4.0000) circle (1.2pt);
\fill[blue!80!black] (6.0000,4.0000) circle (1.2pt);
\fill[blue!80!black] (6.0000,4.0000) circle (1.2pt);
\fill[blue!80!black] (6.0000,4.0000) circle (1.2pt);
\fill[blue!80!black] (6.0000,4.0000) circle (1.2pt);
\fill[blue!80!black] (6.0000,4.0000) circle (1.2pt);
\draw[red!80!black, line width=1.0pt] (5.8125,3.8750) -- (6.0000,4.0000) -- (6.0000,4.0000) -- (6.0000,4.0000) -- (6.0000,4.0000) -- (6.0000,4.0000) -- (6.0000,4.0000) -- (6.0000,4.0000) -- (6.0000,4.0000) -- (6.0000,4.0000) -- (6.0000,4.0000) -- (6.0000,4.0000) -- (6.0000,4.0000) -- (6.0000,4.0000) -- (6.0000,4.0000) -- (6.0000,4.0000) -- (6.0000,4.0000) -- (6.0000,4.0000) -- (6.0000,4.0000) -- (6.0000,4.0000) -- (6.0000,4.0000) -- (6.0000,4.0000) -- (6.0000,4.0000) -- (6.0000,4.0000) -- (6.0000,4.0000) -- (6.0000,4.0000) -- (6.0000,4.0000) -- (6.0000,4.0000) -- (6.0000,4.0000) -- (6.0000,4.0000) -- (6.0000,4.0000) -- (6.0000,4.0000) -- (6.0000,4.0000);
\fill[red!80!black] (5.8125,3.8750) circle (1.2pt);
\fill[red!80!black] (6.0000,4.0000) circle (1.2pt);
\fill[red!80!black] (6.0000,4.0000) circle (1.2pt);
\fill[red!80!black] (6.0000,4.0000) circle (1.2pt);
\fill[red!80!black] (6.0000,4.0000) circle (1.2pt);
\fill[red!80!black] (6.0000,4.0000) circle (1.2pt);
\fill[red!80!black] (6.0000,4.0000) circle (1.2pt);
\fill[red!80!black] (6.0000,4.0000) circle (1.2pt);
\fill[red!80!black] (6.0000,4.0000) circle (1.2pt);
\fill[red!80!black] (6.0000,4.0000) circle (1.2pt);
\fill[red!80!black] (6.0000,4.0000) circle (1.2pt);
\fill[red!80!black] (6.0000,4.0000) circle (1.2pt);
\fill[red!80!black] (6.0000,4.0000) circle (1.2pt);
\fill[red!80!black] (6.0000,4.0000) circle (1.2pt);
\fill[red!80!black] (6.0000,4.0000) circle (1.2pt);
\fill[red!80!black] (6.0000,4.0000) circle (1.2pt);
\fill[red!80!black] (6.0000,4.0000) circle (1.2pt);
\fill[red!80!black] (6.0000,4.0000) circle (1.2pt);
\fill[red!80!black] (6.0000,4.0000) circle (1.2pt);
\fill[red!80!black] (6.0000,4.0000) circle (1.2pt);
\fill[red!80!black] (6.0000,4.0000) circle (1.2pt);
\fill[red!80!black] (6.0000,4.0000) circle (1.2pt);
\fill[red!80!black] (6.0000,4.0000) circle (1.2pt);
\fill[red!80!black] (6.0000,4.0000) circle (1.2pt);
\fill[red!80!black] (6.0000,4.0000) circle (1.2pt);
\fill[red!80!black] (6.0000,4.0000) circle (1.2pt);
\fill[red!80!black] (6.0000,4.0000) circle (1.2pt);
\fill[red!80!black] (6.0000,4.0000) circle (1.2pt);
\fill[red!80!black] (6.0000,4.0000) circle (1.2pt);
\fill[red!80!black] (6.0000,4.0000) circle (1.2pt);
\fill[red!80!black] (6.0000,4.0000) circle (1.2pt);
\fill[red!80!black] (6.0000,4.0000) circle (1.2pt);
\fill[red!80!black] (6.0000,4.0000) circle (1.2pt);
\draw[->, thick] (0,0) -- (12.6000,0) node[right] {$x$};
\draw[->, thick] (0,0) -- (0,8.6000) node[above] {$y$};
\draw (0.0000,0) -- (0.0000,-0.12) node[below] {\small -32.00};
\draw (0,0.0000) -- (-0.12,0.0000) node[left] {\small -32.00};
\draw (2.4000,0) -- (2.4000,-0.12) node[below] {\small -19.20};
\draw (0,1.6000) -- (-0.12,1.6000) node[left] {\small -19.20};
\draw (4.8000,0) -- (4.8000,-0.12) node[below] {\small -6.40};
\draw (0,3.2000) -- (-0.12,3.2000) node[left] {\small -6.40};
\draw (7.2000,0) -- (7.2000,-0.12) node[below] {\small 6.40};
\draw (0,4.8000) -- (-0.12,4.8000) node[left] {\small 6.40};
\draw (9.6000,0) -- (9.6000,-0.12) node[below] {\small 19.20};
\draw (0,6.4000) -- (-0.12,6.4000) node[left] {\small 19.20};
\draw (12.0000,0) -- (12.0000,-0.12) node[below] {\small 32.00};
\draw (0,8.0000) -- (-0.12,8.0000) node[left] {\small 32.00};
\draw[fill=white, draw=black] (7.6,6.7) rectangle (11.8,7.9);
\draw[blue!80!black, line width=1.0pt] (7.9,7.45) -- (8.9,7.45);
\node[anchor=west] at (9.1,7.45) {\small gradDest};
\draw[red!80!black, line width=1.0pt] (7.9,7.05) -- (8.9,7.05);
\node[anchor=west] at (9.1,7.05) {\small fastGradDest};
\end{tikzpicture}

\caption{Функция Ackley 2.}
\end{figure}

\begin{figure}[H]
\centering
\begin{tikzpicture}[scale=0.9]
\draw[fill=orange!3, draw=black] (0,0) rectangle (12.0000,8.0000);
\draw[gray!35, line width=0.15pt] (5.4915,3.6498) -- (5.4747,3.6610);
\draw[gray!35, line width=0.15pt] (5.4747,3.6610) -- (5.4915,3.6621);
\draw[gray!35, line width=0.15pt] (5.4915,3.9261) -- (5.4342,3.9322);
\draw[gray!35, line width=0.15pt] (5.4342,3.9322) -- (5.4342,4.0678);
\draw[gray!35, line width=0.15pt] (5.4342,4.0678) -- (5.4915,4.0739);
\draw[gray!35, line width=0.15pt] (5.4915,4.3379) -- (5.4747,4.3390);
\draw[gray!35, line width=0.15pt] (5.4747,4.3390) -- (5.4915,4.3502);
\draw[gray!35, line width=0.15pt] (5.4915,3.6498) -- (5.4931,3.6610);
\draw[gray!35, line width=0.15pt] (5.4931,3.6610) -- (5.4915,3.6621);
\draw[gray!35, line width=0.15pt] (5.6949,3.7828) -- (5.6742,3.7966);
\draw[gray!35, line width=0.15pt] (5.4915,3.9261) -- (5.5021,3.9322);
\draw[gray!35, line width=0.15pt] (5.6742,3.7966) -- (5.6949,3.8124);
\draw[gray!35, line width=0.15pt] (5.5021,3.9322) -- (5.5021,4.0678);
\draw[gray!35, line width=0.15pt] (5.5021,4.0678) -- (5.4915,4.0739);
\draw[gray!35, line width=0.15pt] (5.6949,4.1876) -- (5.6742,4.2034);
\draw[gray!35, line width=0.15pt] (5.4915,4.3379) -- (5.4931,4.3390);
\draw[gray!35, line width=0.15pt] (5.6742,4.2034) -- (5.6949,4.2172);
\draw[gray!35, line width=0.15pt] (5.4915,4.3502) -- (5.4931,4.3390);
\draw[gray!35, line width=0.15pt] (5.8983,3.6228) -- (5.8891,3.6610);
\draw[gray!35, line width=0.15pt] (5.6949,3.7828) -- (5.7187,3.7966);
\draw[gray!35, line width=0.15pt] (5.8891,3.6610) -- (5.8983,3.6681);
\draw[gray!35, line width=0.15pt] (5.7187,3.7966) -- (5.6949,3.8124);
\draw[gray!35, line width=0.15pt] (5.8983,3.9090) -- (5.8634,3.9322);
\draw[gray!35, line width=0.15pt] (5.8634,3.9322) -- (5.8634,4.0678);
\draw[gray!35, line width=0.15pt] (5.6949,4.1876) -- (5.7187,4.2034);
\draw[gray!35, line width=0.15pt] (5.8634,4.0678) -- (5.8983,4.0910);
\draw[gray!35, line width=0.15pt] (5.7187,4.2034) -- (5.6949,4.2172);
\draw[gray!35, line width=0.15pt] (5.8983,4.3319) -- (5.8891,4.3390);
\draw[gray!35, line width=0.15pt] (5.8891,4.3390) -- (5.8983,4.3772);
\draw[gray!35, line width=0.15pt] (5.8983,3.6228) -- (6.1017,3.6228);
\draw[gray!35, line width=0.15pt] (6.1017,3.6681) -- (5.8983,3.6681);
\draw[gray!35, line width=0.15pt] (5.8983,3.9090) -- (6.1017,3.9090);
\draw[gray!35, line width=0.15pt] (6.1017,4.0910) -- (5.8983,4.0910);
\draw[gray!35, line width=0.15pt] (5.8983,4.3319) -- (6.1017,4.3319);
\draw[gray!35, line width=0.15pt] (6.1017,4.3772) -- (5.8983,4.3772);
\draw[gray!35, line width=0.15pt] (6.1017,3.6228) -- (6.1109,3.6610);
\draw[gray!35, line width=0.15pt] (6.1109,3.6610) -- (6.1017,3.6681);
\draw[gray!35, line width=0.15pt] (6.3051,3.7828) -- (6.2813,3.7966);
\draw[gray!35, line width=0.15pt] (6.1017,3.9090) -- (6.1366,3.9322);
\draw[gray!35, line width=0.15pt] (6.2813,3.7966) -- (6.3051,3.8124);
\draw[gray!35, line width=0.15pt] (6.1366,3.9322) -- (6.1366,4.0678);
\draw[gray!35, line width=0.15pt] (6.1366,4.0678) -- (6.1017,4.0910);
\draw[gray!35, line width=0.15pt] (6.3051,4.1876) -- (6.2813,4.2034);
\draw[gray!35, line width=0.15pt] (6.1017,4.3319) -- (6.1109,4.3390);
\draw[gray!35, line width=0.15pt] (6.2813,4.2034) -- (6.3051,4.2172);
\draw[gray!35, line width=0.15pt] (6.1017,4.3772) -- (6.1109,4.3390);
\draw[gray!35, line width=0.15pt] (6.5085,3.6498) -- (6.5069,3.6610);
\draw[gray!35, line width=0.15pt] (6.3051,3.7828) -- (6.3258,3.7966);
\draw[gray!35, line width=0.15pt] (6.5069,3.6610) -- (6.5085,3.6621);
\draw[gray!35, line width=0.15pt] (6.3258,3.7966) -- (6.3051,3.8124);
\draw[gray!35, line width=0.15pt] (6.5085,3.9261) -- (6.4979,3.9322);
\draw[gray!35, line width=0.15pt] (6.4979,3.9322) -- (6.4979,4.0678);
\draw[gray!35, line width=0.15pt] (6.3051,4.1876) -- (6.3258,4.2034);
\draw[gray!35, line width=0.15pt] (6.4979,4.0678) -- (6.5085,4.0739);
\draw[gray!35, line width=0.15pt] (6.3258,4.2034) -- (6.3051,4.2172);
\draw[gray!35, line width=0.15pt] (6.5085,4.3379) -- (6.5069,4.3390);
\draw[gray!35, line width=0.15pt] (6.5069,4.3390) -- (6.5085,4.3502);
\draw[gray!35, line width=0.15pt] (6.5085,3.6498) -- (6.5253,3.6610);
\draw[gray!35, line width=0.15pt] (6.5085,3.6621) -- (6.5253,3.6610);
\draw[gray!35, line width=0.15pt] (6.5085,3.9261) -- (6.5658,3.9322);
\draw[gray!35, line width=0.15pt] (6.5658,3.9322) -- (6.5658,4.0678);
\draw[gray!35, line width=0.15pt] (6.5085,4.0739) -- (6.5658,4.0678);
\draw[gray!35, line width=0.15pt] (6.5085,4.3379) -- (6.5253,4.3390);
\draw[gray!35, line width=0.15pt] (6.5085,4.3502) -- (6.5253,4.3390);
\draw[gray!45, line width=0.15pt] (5.2881,3.5068) -- (5.2603,3.5254);
\draw[gray!45, line width=0.15pt] (5.2603,3.5254) -- (5.2375,3.6610);
\draw[gray!45, line width=0.15pt] (5.2375,3.6610) -- (5.2881,3.6737);
\draw[gray!45, line width=0.15pt] (5.2881,3.9240) -- (5.2192,3.9322);
\draw[gray!45, line width=0.15pt] (5.2192,3.9322) -- (5.2192,4.0678);
\draw[gray!45, line width=0.15pt] (5.2192,4.0678) -- (5.2881,4.0760);
\draw[gray!45, line width=0.15pt] (5.2881,4.3263) -- (5.2375,4.3390);
\draw[gray!45, line width=0.15pt] (5.2375,4.3390) -- (5.2603,4.4746);
\draw[gray!45, line width=0.15pt] (5.2603,4.4746) -- (5.2881,4.4932);
\draw[gray!45, line width=0.15pt] (5.2881,3.5068) -- (5.4915,3.4916);
\draw[gray!45, line width=0.15pt] (5.4915,3.6859) -- (5.2881,3.6737);
\draw[gray!45, line width=0.15pt] (5.2881,3.9240) -- (5.4915,3.9031);
\draw[gray!45, line width=0.15pt] (5.4915,4.0969) -- (5.2881,4.0760);
\draw[gray!45, line width=0.15pt] (5.2881,4.3263) -- (5.4915,4.3141);
\draw[gray!45, line width=0.15pt] (5.4915,4.5084) -- (5.2881,4.4932);
\draw[gray!45, line width=0.15pt] (5.4915,3.4916) -- (5.5106,3.5254);
\draw[gray!45, line width=0.15pt] (5.5106,3.5254) -- (5.5289,3.6610);
\draw[gray!45, line width=0.15pt] (5.5289,3.6610) -- (5.4915,3.6859);
\draw[gray!45, line width=0.15pt] (5.6949,3.7612) -- (5.6418,3.7966);
\draw[gray!45, line width=0.15pt] (5.4915,3.9031) -- (5.5420,3.9322);
\draw[gray!45, line width=0.15pt] (5.6418,3.7966) -- (5.6949,3.8372);
\draw[gray!45, line width=0.15pt] (5.5420,3.9322) -- (5.5420,4.0678);
\draw[gray!45, line width=0.15pt] (5.5420,4.0678) -- (5.4915,4.0969);
\draw[gray!45, line width=0.15pt] (5.6949,4.1628) -- (5.6418,4.2034);
\draw[gray!45, line width=0.15pt] (5.4915,4.3141) -- (5.5289,4.3390);
\draw[gray!45, line width=0.15pt] (5.6418,4.2034) -- (5.6949,4.2388);
\draw[gray!45, line width=0.15pt] (5.5289,4.3390) -- (5.5106,4.4746);
\draw[gray!45, line width=0.15pt] (5.4915,4.5084) -- (5.5106,4.4746);
\draw[gray!45, line width=0.15pt] (5.8983,3.4795) -- (5.8860,3.5254);
\draw[gray!45, line width=0.15pt] (5.8860,3.5254) -- (5.8547,3.6610);
\draw[gray!45, line width=0.15pt] (5.6949,3.7612) -- (5.7558,3.7966);
\draw[gray!45, line width=0.15pt] (5.8547,3.6610) -- (5.8983,3.6947);
\draw[gray!45, line width=0.15pt] (5.7558,3.7966) -- (5.6949,3.8372);
\draw[gray!45, line width=0.15pt] (5.8983,3.8857) -- (5.8286,3.9322);
\draw[gray!45, line width=0.15pt] (5.8286,3.9322) -- (5.8286,4.0678);
\draw[gray!45, line width=0.15pt] (5.6949,4.1628) -- (5.7558,4.2034);
\draw[gray!45, line width=0.15pt] (5.8286,4.0678) -- (5.8983,4.1143);
\draw[gray!45, line width=0.15pt] (5.7558,4.2034) -- (5.6949,4.2388);
\draw[gray!45, line width=0.15pt] (5.8983,4.3053) -- (5.8547,4.3390);
\draw[gray!45, line width=0.15pt] (5.8547,4.3390) -- (5.8860,4.4746);
\draw[gray!45, line width=0.15pt] (5.8860,4.4746) -- (5.8983,4.5205);
\draw[gray!45, line width=0.15pt] (5.8983,3.4795) -- (6.1017,3.4795);
\draw[gray!45, line width=0.15pt] (6.1017,3.6947) -- (5.8983,3.6947);
\draw[gray!45, line width=0.15pt] (5.8983,3.8857) -- (6.1017,3.8857);
\draw[gray!45, line width=0.15pt] (6.1017,4.1143) -- (5.8983,4.1143);
\draw[gray!45, line width=0.15pt] (5.8983,4.3053) -- (6.1017,4.3053);
\draw[gray!45, line width=0.15pt] (6.1017,4.5205) -- (5.8983,4.5205);
\draw[gray!45, line width=0.15pt] (6.1017,3.4795) -- (6.1140,3.5254);
\draw[gray!45, line width=0.15pt] (6.1140,3.5254) -- (6.1453,3.6610);
\draw[gray!45, line width=0.15pt] (6.1453,3.6610) -- (6.1017,3.6947);
\draw[gray!45, line width=0.15pt] (6.3051,3.7612) -- (6.2442,3.7966);
\draw[gray!45, line width=0.15pt] (6.1017,3.8857) -- (6.1714,3.9322);
\draw[gray!45, line width=0.15pt] (6.2442,3.7966) -- (6.3051,3.8372);
\draw[gray!45, line width=0.15pt] (6.1714,3.9322) -- (6.1714,4.0678);
\draw[gray!45, line width=0.15pt] (6.1714,4.0678) -- (6.1017,4.1143);
\draw[gray!45, line width=0.15pt] (6.3051,4.1628) -- (6.2442,4.2034);
\draw[gray!45, line width=0.15pt] (6.1017,4.3053) -- (6.1453,4.3390);
\draw[gray!45, line width=0.15pt] (6.2442,4.2034) -- (6.3051,4.2388);
\draw[gray!45, line width=0.15pt] (6.1453,4.3390) -- (6.1140,4.4746);
\draw[gray!45, line width=0.15pt] (6.1017,4.5205) -- (6.1140,4.4746);
\draw[gray!45, line width=0.15pt] (6.5085,3.4916) -- (6.4894,3.5254);
\draw[gray!45, line width=0.15pt] (6.4894,3.5254) -- (6.4711,3.6610);
\draw[gray!45, line width=0.15pt] (6.3051,3.7612) -- (6.3582,3.7966);
\draw[gray!45, line width=0.15pt] (6.4711,3.6610) -- (6.5085,3.6859);
\draw[gray!45, line width=0.15pt] (6.3582,3.7966) -- (6.3051,3.8372);
\draw[gray!45, line width=0.15pt] (6.5085,3.9031) -- (6.4580,3.9322);
\draw[gray!45, line width=0.15pt] (6.4580,3.9322) -- (6.4580,4.0678);
\draw[gray!45, line width=0.15pt] (6.3051,4.1628) -- (6.3582,4.2034);
\draw[gray!45, line width=0.15pt] (6.4580,4.0678) -- (6.5085,4.0969);
\draw[gray!45, line width=0.15pt] (6.3582,4.2034) -- (6.3051,4.2388);
\draw[gray!45, line width=0.15pt] (6.5085,4.3141) -- (6.4711,4.3390);
\draw[gray!45, line width=0.15pt] (6.4711,4.3390) -- (6.4894,4.4746);
\draw[gray!45, line width=0.15pt] (6.4894,4.4746) -- (6.5085,4.5084);
\draw[gray!45, line width=0.15pt] (6.5085,3.4916) -- (6.7119,3.5068);
\draw[gray!45, line width=0.15pt] (6.7119,3.6737) -- (6.5085,3.6859);
\draw[gray!45, line width=0.15pt] (6.5085,3.9031) -- (6.7119,3.9240);
\draw[gray!45, line width=0.15pt] (6.7119,4.0760) -- (6.5085,4.0969);
\draw[gray!45, line width=0.15pt] (6.5085,4.3141) -- (6.7119,4.3263);
\draw[gray!45, line width=0.15pt] (6.7119,4.4932) -- (6.5085,4.5084);
\draw[gray!45, line width=0.15pt] (6.7119,3.5068) -- (6.7397,3.5254);
\draw[gray!45, line width=0.15pt] (6.7397,3.5254) -- (6.7625,3.6610);
\draw[gray!45, line width=0.15pt] (6.7119,3.6737) -- (6.7625,3.6610);
\draw[gray!45, line width=0.15pt] (6.7119,3.9240) -- (6.7808,3.9322);
\draw[gray!45, line width=0.15pt] (6.7808,3.9322) -- (6.7808,4.0678);
\draw[gray!45, line width=0.15pt] (6.7119,4.0760) -- (6.7808,4.0678);
\draw[gray!45, line width=0.15pt] (6.7119,4.3263) -- (6.7625,4.3390);
\draw[gray!45, line width=0.15pt] (6.7625,4.3390) -- (6.7397,4.4746);
\draw[gray!45, line width=0.15pt] (6.7119,4.4932) -- (6.7397,4.4746);
\draw[gray!55, line width=0.15pt] (4.4746,3.5242) -- (4.4732,3.5254);
\draw[gray!55, line width=0.15pt] (4.4732,3.5254) -- (4.4746,3.5390);
\draw[gray!55, line width=0.15pt] (4.4746,4.4610) -- (4.4732,4.4746);
\draw[gray!55, line width=0.15pt] (4.4732,4.4746) -- (4.4746,4.4758);
\draw[gray!55, line width=0.15pt] (4.4746,3.5242) -- (4.4768,3.5254);
\draw[gray!55, line width=0.15pt] (4.4746,3.5390) -- (4.4768,3.5254);
\draw[gray!55, line width=0.15pt] (4.4746,4.4610) -- (4.4768,4.4746);
\draw[gray!55, line width=0.15pt] (4.4746,4.4758) -- (4.4768,4.4746);
\draw[gray!55, line width=0.15pt] (4.8814,3.1818) -- (4.7728,3.2542);
\draw[gray!55, line width=0.15pt] (4.7728,3.2542) -- (4.8814,3.2662);
\draw[gray!55, line width=0.15pt] (4.8814,4.7338) -- (4.7728,4.7458);
\draw[gray!55, line width=0.15pt] (4.7728,4.7458) -- (4.8814,4.8182);
\draw[gray!55, line width=0.15pt] (4.8814,3.1818) -- (4.8993,3.2542);
\draw[gray!55, line width=0.15pt] (4.8993,3.2542) -- (4.8814,3.2662);
\draw[gray!55, line width=0.15pt] (5.0847,3.3601) -- (5.0402,3.3898);
\draw[gray!55, line width=0.15pt] (5.0402,3.3898) -- (5.0847,3.4382);
\draw[gray!55, line width=0.15pt] (5.0847,3.7611) -- (5.0727,3.7966);
\draw[gray!55, line width=0.15pt] (5.0727,3.7966) -- (5.0446,3.9322);
\draw[gray!55, line width=0.15pt] (5.0446,3.9322) -- (5.0446,4.0678);
\draw[gray!55, line width=0.15pt] (5.0446,4.0678) -- (5.0727,4.2034);
\draw[gray!55, line width=0.15pt] (5.0727,4.2034) -- (5.0847,4.2389);
\draw[gray!55, line width=0.15pt] (5.0847,4.5618) -- (5.0402,4.6102);
\draw[gray!55, line width=0.15pt] (4.8814,4.7338) -- (4.8993,4.7458);
\draw[gray!55, line width=0.15pt] (5.0402,4.6102) -- (5.0847,4.6399);
\draw[gray!55, line width=0.15pt] (4.8814,4.8182) -- (4.8993,4.7458);
\draw[gray!55, line width=0.15pt] (5.2881,2.9821) -- (5.2863,2.9831);
\draw[gray!55, line width=0.15pt] (5.2863,2.9831) -- (5.2881,2.9845);
\draw[gray!55, line width=0.15pt] (5.0847,3.3601) -- (5.1573,3.3898);
\draw[gray!55, line width=0.15pt] (5.1573,3.3898) -- (5.0847,3.4382);
\draw[gray!55, line width=0.15pt] (5.2881,3.4578) -- (5.1866,3.5254);
\draw[gray!55, line width=0.15pt] (5.1866,3.5254) -- (5.1396,3.6610);
\draw[gray!55, line width=0.15pt] (5.0847,3.7611) -- (5.0944,3.7966);
\draw[gray!55, line width=0.15pt] (5.1396,3.6610) -- (5.2881,3.6983);
\draw[gray!55, line width=0.15pt] (5.0944,3.7966) -- (5.2881,3.8985);
\draw[gray!55, line width=0.15pt] (5.2881,4.1015) -- (5.0944,4.2034);
\draw[gray!55, line width=0.15pt] (5.0944,4.2034) -- (5.0847,4.2389);
\draw[gray!55, line width=0.15pt] (5.2881,4.3017) -- (5.1396,4.3390);
\draw[gray!55, line width=0.15pt] (5.1396,4.3390) -- (5.1866,4.4746);
\draw[gray!55, line width=0.15pt] (5.0847,4.5618) -- (5.1573,4.6102);
\draw[gray!55, line width=0.15pt] (5.1866,4.4746) -- (5.2881,4.5422);
\draw[gray!55, line width=0.15pt] (5.0847,4.6399) -- (5.1573,4.6102);
\draw[gray!55, line width=0.15pt] (5.2881,5.0155) -- (5.2863,5.0169);
\draw[gray!55, line width=0.15pt] (5.2863,5.0169) -- (5.2881,5.0179);
\draw[gray!55, line width=0.15pt] (5.2881,2.9821) -- (5.3085,2.9831);
\draw[gray!55, line width=0.15pt] (5.2881,2.9845) -- (5.3085,2.9831);
\draw[gray!55, line width=0.15pt] (5.2881,3.4578) -- (5.4915,3.4264);
\draw[gray!55, line width=0.15pt] (5.4915,3.7098) -- (5.2881,3.6983);
\draw[gray!55, line width=0.15pt] (5.2881,3.8985) -- (5.4915,3.8801);
\draw[gray!55, line width=0.15pt] (5.4915,4.1199) -- (5.2881,4.1015);
\draw[gray!55, line width=0.15pt] (5.2881,4.3017) -- (5.4915,4.2902);
\draw[gray!55, line width=0.15pt] (5.4915,4.5736) -- (5.2881,4.5422);
\draw[gray!55, line width=0.15pt] (5.2881,5.0155) -- (5.3085,5.0169);
\draw[gray!55, line width=0.15pt] (5.2881,5.0179) -- (5.3085,5.0169);
\draw[gray!55, line width=0.15pt] (5.6949,3.3818) -- (5.6416,3.3898);
\draw[gray!55, line width=0.15pt] (5.4915,3.4264) -- (5.5475,3.5254);
\draw[gray!55, line width=0.15pt] (5.6416,3.3898) -- (5.6949,3.3963);
\draw[gray!55, line width=0.15pt] (5.5475,3.5254) -- (5.5647,3.6610);
\draw[gray!55, line width=0.15pt] (5.5647,3.6610) -- (5.4915,3.7098);
\draw[gray!55, line width=0.15pt] (5.6949,3.7396) -- (5.6094,3.7966);
\draw[gray!55, line width=0.15pt] (5.4915,3.8801) -- (5.5819,3.9322);
\draw[gray!55, line width=0.15pt] (5.6094,3.7966) -- (5.6949,3.8620);
\draw[gray!55, line width=0.15pt] (5.5819,3.9322) -- (5.5819,4.0678);
\draw[gray!55, line width=0.15pt] (5.5819,4.0678) -- (5.4915,4.1199);
\draw[gray!55, line width=0.15pt] (5.6949,4.1380) -- (5.6094,4.2034);
\draw[gray!55, line width=0.15pt] (5.4915,4.2902) -- (5.5647,4.3390);
\draw[gray!55, line width=0.15pt] (5.6094,4.2034) -- (5.6949,4.2604);
\draw[gray!55, line width=0.15pt] (5.5647,4.3390) -- (5.5475,4.4746);
\draw[gray!55, line width=0.15pt] (5.5475,4.4746) -- (5.4915,4.5736);
\draw[gray!55, line width=0.15pt] (5.6949,4.6037) -- (5.6416,4.6102);
\draw[gray!55, line width=0.15pt] (5.6416,4.6102) -- (5.6949,4.6182);
\draw[gray!55, line width=0.15pt] (5.6949,3.3818) -- (5.8983,3.3631);
\draw[gray!55, line width=0.15pt] (5.8478,3.5254) -- (5.6949,3.3963);
\draw[gray!55, line width=0.15pt] (5.8478,3.5254) -- (5.8202,3.6610);
\draw[gray!55, line width=0.15pt] (5.6949,3.7396) -- (5.7930,3.7966);
\draw[gray!55, line width=0.15pt] (5.8202,3.6610) -- (5.8983,3.7213);
\draw[gray!55, line width=0.15pt] (5.7930,3.7966) -- (5.6949,3.8620);
\draw[gray!55, line width=0.15pt] (5.8983,3.8625) -- (5.7937,3.9322);
\draw[gray!55, line width=0.15pt] (5.7937,3.9322) -- (5.7937,4.0678);
\draw[gray!55, line width=0.15pt] (5.6949,4.1380) -- (5.7930,4.2034);
\draw[gray!55, line width=0.15pt] (5.7937,4.0678) -- (5.8983,4.1375);
\draw[gray!55, line width=0.15pt] (5.7930,4.2034) -- (5.6949,4.2604);
\draw[gray!55, line width=0.15pt] (5.8983,4.2787) -- (5.8202,4.3390);
\draw[gray!55, line width=0.15pt] (5.8202,4.3390) -- (5.8478,4.4746);
\draw[gray!55, line width=0.15pt] (5.6949,4.6037) -- (5.8478,4.4746);
\draw[gray!55, line width=0.15pt] (5.8983,4.6369) -- (5.6949,4.6182);
\draw[gray!55, line width=0.15pt] (5.8983,3.3631) -- (6.1017,3.3631);
\draw[gray!55, line width=0.15pt] (6.1017,3.7213) -- (5.8983,3.7213);
\draw[gray!55, line width=0.15pt] (5.8983,3.8625) -- (6.1017,3.8625);
\draw[gray!55, line width=0.15pt] (6.1017,4.1375) -- (5.8983,4.1375);
\draw[gray!55, line width=0.15pt] (5.8983,4.2787) -- (6.1017,4.2787);
\draw[gray!55, line width=0.15pt] (6.1017,4.6369) -- (5.8983,4.6369);
\draw[gray!55, line width=0.15pt] (6.1017,3.3631) -- (6.3051,3.3818);
\draw[gray!55, line width=0.15pt] (6.3051,3.3963) -- (6.1522,3.5254);
\draw[gray!55, line width=0.15pt] (6.1522,3.5254) -- (6.1798,3.6610);
\draw[gray!55, line width=0.15pt] (6.1798,3.6610) -- (6.1017,3.7213);
\draw[gray!55, line width=0.15pt] (6.3051,3.7396) -- (6.2070,3.7966);
\draw[gray!55, line width=0.15pt] (6.1017,3.8625) -- (6.2063,3.9322);
\draw[gray!55, line width=0.15pt] (6.2070,3.7966) -- (6.3051,3.8620);
\draw[gray!55, line width=0.15pt] (6.2063,3.9322) -- (6.2063,4.0678);
\draw[gray!55, line width=0.15pt] (6.2063,4.0678) -- (6.1017,4.1375);
\draw[gray!55, line width=0.15pt] (6.3051,4.1380) -- (6.2070,4.2034);
\draw[gray!55, line width=0.15pt] (6.1017,4.2787) -- (6.1798,4.3390);
\draw[gray!55, line width=0.15pt] (6.2070,4.2034) -- (6.3051,4.2604);
\draw[gray!55, line width=0.15pt] (6.1798,4.3390) -- (6.1522,4.4746);
\draw[gray!55, line width=0.15pt] (6.1522,4.4746) -- (6.3051,4.6037);
\draw[gray!55, line width=0.15pt] (6.3051,4.6182) -- (6.1017,4.6369);
\draw[gray!55, line width=0.15pt] (6.3051,3.3818) -- (6.3584,3.3898);
\draw[gray!55, line width=0.15pt] (6.3584,3.3898) -- (6.3051,3.3963);
\draw[gray!55, line width=0.15pt] (6.5085,3.4264) -- (6.4525,3.5254);
\draw[gray!55, line width=0.15pt] (6.4525,3.5254) -- (6.4353,3.6610);
\draw[gray!55, line width=0.15pt] (6.3051,3.7396) -- (6.3906,3.7966);
\draw[gray!55, line width=0.15pt] (6.4353,3.6610) -- (6.5085,3.7098);
\draw[gray!55, line width=0.15pt] (6.3906,3.7966) -- (6.3051,3.8620);
\draw[gray!55, line width=0.15pt] (6.5085,3.8801) -- (6.4181,3.9322);
\draw[gray!55, line width=0.15pt] (6.4181,3.9322) -- (6.4181,4.0678);
\draw[gray!55, line width=0.15pt] (6.3051,4.1380) -- (6.3906,4.2034);
\draw[gray!55, line width=0.15pt] (6.4181,4.0678) -- (6.5085,4.1199);
\draw[gray!55, line width=0.15pt] (6.3906,4.2034) -- (6.3051,4.2604);
\draw[gray!55, line width=0.15pt] (6.5085,4.2902) -- (6.4353,4.3390);
\draw[gray!55, line width=0.15pt] (6.4353,4.3390) -- (6.4525,4.4746);
\draw[gray!55, line width=0.15pt] (6.3051,4.6037) -- (6.3584,4.6102);
\draw[gray!55, line width=0.15pt] (6.4525,4.4746) -- (6.5085,4.5736);
\draw[gray!55, line width=0.15pt] (6.3051,4.6182) -- (6.3584,4.6102);
\draw[gray!55, line width=0.15pt] (6.7119,2.9821) -- (6.6915,2.9831);
\draw[gray!55, line width=0.15pt] (6.6915,2.9831) -- (6.7119,2.9845);
\draw[gray!55, line width=0.15pt] (6.5085,3.4264) -- (6.7119,3.4578);
\draw[gray!55, line width=0.15pt] (6.7119,3.6983) -- (6.5085,3.7098);
\draw[gray!55, line width=0.15pt] (6.5085,3.8801) -- (6.7119,3.8985);
\draw[gray!55, line width=0.15pt] (6.7119,4.1015) -- (6.5085,4.1199);
\draw[gray!55, line width=0.15pt] (6.5085,4.2902) -- (6.7119,4.3017);
\draw[gray!55, line width=0.15pt] (6.7119,4.5422) -- (6.5085,4.5736);
\draw[gray!55, line width=0.15pt] (6.7119,5.0155) -- (6.6915,5.0169);
\draw[gray!55, line width=0.15pt] (6.6915,5.0169) -- (6.7119,5.0179);
\draw[gray!55, line width=0.15pt] (6.7119,2.9821) -- (6.7137,2.9831);
\draw[gray!55, line width=0.15pt] (6.7119,2.9845) -- (6.7137,2.9831);
\draw[gray!55, line width=0.15pt] (6.9153,3.3601) -- (6.8427,3.3898);
\draw[gray!55, line width=0.15pt] (6.7119,3.4578) -- (6.8134,3.5254);
\draw[gray!55, line width=0.15pt] (6.8427,3.3898) -- (6.9153,3.4382);
\draw[gray!55, line width=0.15pt] (6.8134,3.5254) -- (6.8604,3.6610);
\draw[gray!55, line width=0.15pt] (6.8604,3.6610) -- (6.7119,3.6983);
\draw[gray!55, line width=0.15pt] (6.9153,3.7611) -- (6.9056,3.7966);
\draw[gray!55, line width=0.15pt] (6.7119,3.8985) -- (6.9056,3.7966);
\draw[gray!55, line width=0.15pt] (6.9056,4.2034) -- (6.7119,4.1015);
\draw[gray!55, line width=0.15pt] (6.7119,4.3017) -- (6.8604,4.3390);
\draw[gray!55, line width=0.15pt] (6.9056,4.2034) -- (6.9153,4.2389);
\draw[gray!55, line width=0.15pt] (6.8604,4.3390) -- (6.8134,4.4746);
\draw[gray!55, line width=0.15pt] (6.8134,4.4746) -- (6.7119,4.5422);
\draw[gray!55, line width=0.15pt] (6.9153,4.5618) -- (6.8427,4.6102);
\draw[gray!55, line width=0.15pt] (6.8427,4.6102) -- (6.9153,4.6399);
\draw[gray!55, line width=0.15pt] (6.7119,5.0155) -- (6.7137,5.0169);
\draw[gray!55, line width=0.15pt] (6.7119,5.0179) -- (6.7137,5.0169);
\draw[gray!55, line width=0.15pt] (7.1186,3.1818) -- (7.1007,3.2542);
\draw[gray!55, line width=0.15pt] (6.9153,3.3601) -- (6.9598,3.3898);
\draw[gray!55, line width=0.15pt] (7.1007,3.2542) -- (7.1186,3.2662);
\draw[gray!55, line width=0.15pt] (6.9153,3.4382) -- (6.9598,3.3898);
\draw[gray!55, line width=0.15pt] (6.9153,3.7611) -- (6.9273,3.7966);
\draw[gray!55, line width=0.15pt] (6.9273,3.7966) -- (6.9554,3.9322);
\draw[gray!55, line width=0.15pt] (6.9554,3.9322) -- (6.9554,4.0678);
\draw[gray!55, line width=0.15pt] (6.9554,4.0678) -- (6.9273,4.2034);
\draw[gray!55, line width=0.15pt] (6.9153,4.2389) -- (6.9273,4.2034);
\draw[gray!55, line width=0.15pt] (6.9153,4.5618) -- (6.9598,4.6102);
\draw[gray!55, line width=0.15pt] (6.9598,4.6102) -- (6.9153,4.6399);
\draw[gray!55, line width=0.15pt] (7.1186,4.7338) -- (7.1007,4.7458);
\draw[gray!55, line width=0.15pt] (7.1007,4.7458) -- (7.1186,4.8182);
\draw[gray!55, line width=0.15pt] (7.1186,3.1818) -- (7.2272,3.2542);
\draw[gray!55, line width=0.15pt] (7.1186,3.2662) -- (7.2272,3.2542);
\draw[gray!55, line width=0.15pt] (7.1186,4.7338) -- (7.2272,4.7458);
\draw[gray!55, line width=0.15pt] (7.1186,4.8182) -- (7.2272,4.7458);
\draw[gray!55, line width=0.15pt] (7.5254,3.5242) -- (7.5232,3.5254);
\draw[gray!55, line width=0.15pt] (7.5232,3.5254) -- (7.5254,3.5390);
\draw[gray!55, line width=0.15pt] (7.5254,4.4610) -- (7.5232,4.4746);
\draw[gray!55, line width=0.15pt] (7.5232,4.4746) -- (7.5254,4.4758);
\draw[gray!55, line width=0.15pt] (7.5254,3.5242) -- (7.5268,3.5254);
\draw[gray!55, line width=0.15pt] (7.5254,3.5390) -- (7.5268,3.5254);
\draw[gray!55, line width=0.15pt] (7.5254,4.4610) -- (7.5268,4.4746);
\draw[gray!55, line width=0.15pt] (7.5254,4.4758) -- (7.5268,4.4746);
\draw[gray!65, line width=0.15pt] (2.6441,3.5185) -- (2.6367,3.5254);
\draw[gray!65, line width=0.15pt] (2.6367,3.5254) -- (2.6400,3.6610);
\draw[gray!65, line width=0.15pt] (2.6400,3.6610) -- (2.6441,3.6628);
\draw[gray!65, line width=0.15pt] (2.6441,4.3372) -- (2.6400,4.3390);
\draw[gray!65, line width=0.15pt] (2.6400,4.3390) -- (2.6367,4.4746);
\draw[gray!65, line width=0.15pt] (2.6367,4.4746) -- (2.6441,4.4815);
\draw[gray!65, line width=0.15pt] (2.6441,3.5185) -- (2.8475,3.5082);
\draw[gray!65, line width=0.15pt] (2.8475,3.6640) -- (2.6441,3.6628);
\draw[gray!65, line width=0.15pt] (2.6441,4.3372) -- (2.8475,4.3360);
\draw[gray!65, line width=0.15pt] (2.8475,4.4918) -- (2.6441,4.4815);
\draw[gray!65, line width=0.15pt] (3.0508,3.3672) -- (3.0064,3.3898);
\draw[gray!65, line width=0.15pt] (2.8475,3.5082) -- (2.8931,3.5254);
\draw[gray!65, line width=0.15pt] (3.0064,3.3898) -- (3.0508,3.4383);
\draw[gray!65, line width=0.15pt] (2.8931,3.5254) -- (2.8878,3.6610);
\draw[gray!65, line width=0.15pt] (2.8475,3.6640) -- (2.8878,3.6610);
\draw[gray!65, line width=0.15pt] (3.0508,3.8880) -- (2.9476,3.9322);
\draw[gray!65, line width=0.15pt] (2.9476,3.9322) -- (2.9476,4.0678);
\draw[gray!65, line width=0.15pt] (2.9476,4.0678) -- (3.0508,4.1120);
\draw[gray!65, line width=0.15pt] (2.8475,4.3360) -- (2.8878,4.3390);
\draw[gray!65, line width=0.15pt] (2.8878,4.3390) -- (2.8931,4.4746);
\draw[gray!65, line width=0.15pt] (2.8931,4.4746) -- (2.8475,4.4918);
\draw[gray!65, line width=0.15pt] (3.0508,4.5617) -- (3.0064,4.6102);
\draw[gray!65, line width=0.15pt] (3.0064,4.6102) -- (3.0508,4.6328);
\draw[gray!65, line width=0.15pt] (3.2542,3.0057) -- (3.2278,3.1186);
\draw[gray!65, line width=0.15pt] (3.2278,3.1186) -- (3.2077,3.2542);
\draw[gray!65, line width=0.15pt] (3.0508,3.3672) -- (3.0775,3.3898);
\draw[gray!65, line width=0.15pt] (3.2077,3.2542) -- (3.2542,3.2789);
\draw[gray!65, line width=0.15pt] (3.0508,3.4383) -- (3.0775,3.3898);
\draw[gray!65, line width=0.15pt] (3.0508,3.8880) -- (3.0776,3.9322);
\draw[gray!65, line width=0.15pt] (3.0776,3.9322) -- (3.0776,4.0678);
\draw[gray!65, line width=0.15pt] (3.0508,4.1120) -- (3.0776,4.0678);
\draw[gray!65, line width=0.15pt] (3.0508,4.5617) -- (3.0775,4.6102);
\draw[gray!65, line width=0.15pt] (3.0775,4.6102) -- (3.0508,4.6328);
\draw[gray!65, line width=0.15pt] (3.2542,4.7211) -- (3.2077,4.7458);
\draw[gray!65, line width=0.15pt] (3.2077,4.7458) -- (3.2278,4.8814);
\draw[gray!65, line width=0.15pt] (3.2278,4.8814) -- (3.2542,4.9943);
\draw[gray!65, line width=0.15pt] (3.2542,3.0057) -- (3.4576,3.0526);
\draw[gray!65, line width=0.15pt] (3.4576,3.2674) -- (3.2542,3.2789);
\draw[gray!65, line width=0.15pt] (3.2542,4.7211) -- (3.4576,4.7326);
\draw[gray!65, line width=0.15pt] (3.4576,4.9474) -- (3.2542,4.9943);
\draw[gray!65, line width=0.15pt] (3.6610,2.4149) -- (3.6223,2.4407);
\draw[gray!65, line width=0.15pt] (3.6223,2.4407) -- (3.6610,2.4506);
\draw[gray!65, line width=0.15pt] (3.6610,2.9484) -- (3.5533,2.9831);
\draw[gray!65, line width=0.15pt] (3.4576,3.0526) -- (3.6140,3.1186);
\draw[gray!65, line width=0.15pt] (3.5533,2.9831) -- (3.6610,3.0912);
\draw[gray!65, line width=0.15pt] (3.6140,3.1186) -- (3.6610,3.1436);
\draw[gray!65, line width=0.15pt] (3.6610,3.2955) -- (3.4576,3.2674);
\draw[gray!65, line width=0.15pt] (3.6610,3.4748) -- (3.5661,3.5254);
\draw[gray!65, line width=0.15pt] (3.5661,3.5254) -- (3.6028,3.6610);
\draw[gray!65, line width=0.15pt] (3.6028,3.6610) -- (3.6610,3.6870);
\draw[gray!65, line width=0.15pt] (3.6610,3.9260) -- (3.6489,3.9322);
\draw[gray!65, line width=0.15pt] (3.6489,3.9322) -- (3.6489,4.0678);
\draw[gray!65, line width=0.15pt] (3.6489,4.0678) -- (3.6610,4.0740);
\draw[gray!65, line width=0.15pt] (3.6610,4.3130) -- (3.6028,4.3390);
\draw[gray!65, line width=0.15pt] (3.6028,4.3390) -- (3.5661,4.4746);
\draw[gray!65, line width=0.15pt] (3.5661,4.4746) -- (3.6610,4.5252);
\draw[gray!65, line width=0.15pt] (3.4576,4.7326) -- (3.6610,4.7045);
\draw[gray!65, line width=0.15pt] (3.6610,4.8564) -- (3.6140,4.8814);
\draw[gray!65, line width=0.15pt] (3.6140,4.8814) -- (3.4576,4.9474);
\draw[gray!65, line width=0.15pt] (3.6610,4.9088) -- (3.5533,5.0169);
\draw[gray!65, line width=0.15pt] (3.5533,5.0169) -- (3.6610,5.0516);
\draw[gray!65, line width=0.15pt] (3.6610,5.5494) -- (3.6223,5.5593);
\draw[gray!65, line width=0.15pt] (3.6223,5.5593) -- (3.6610,5.5851);
\draw[gray!65, line width=0.15pt] (3.6610,2.4149) -- (3.6759,2.4407);
\draw[gray!65, line width=0.15pt] (3.6759,2.4407) -- (3.6610,2.4506);
\draw[gray!65, line width=0.15pt] (3.8644,2.5505) -- (3.8257,2.5763);
\draw[gray!65, line width=0.15pt] (3.8257,2.5763) -- (3.8638,2.7119);
\draw[gray!65, line width=0.15pt] (3.8638,2.7119) -- (3.8086,2.8475);
\draw[gray!65, line width=0.15pt] (3.6610,2.9484) -- (3.7056,2.9831);
\draw[gray!65, line width=0.15pt] (3.8086,2.8475) -- (3.8644,2.8795);
\draw[gray!65, line width=0.15pt] (3.6610,3.0912) -- (3.7056,2.9831);
\draw[gray!65, line width=0.15pt] (3.6610,3.1436) -- (3.7115,3.2542);
\draw[gray!65, line width=0.15pt] (3.7115,3.2542) -- (3.6610,3.2955);
\draw[gray!65, line width=0.15pt] (3.8644,3.3852) -- (3.8553,3.3898);
\draw[gray!65, line width=0.15pt] (3.6610,3.4748) -- (3.7122,3.5254);
\draw[gray!65, line width=0.15pt] (3.8553,3.3898) -- (3.8644,3.3933);
\draw[gray!65, line width=0.15pt] (3.7122,3.5254) -- (3.7011,3.6610);
\draw[gray!65, line width=0.15pt] (3.7011,3.6610) -- (3.6610,3.6870);
\draw[gray!65, line width=0.15pt] (3.8644,3.7418) -- (3.7839,3.7966);
\draw[gray!65, line width=0.15pt] (3.6610,3.9260) -- (3.6728,3.9322);
\draw[gray!65, line width=0.15pt] (3.7839,3.7966) -- (3.8644,3.8585);
\draw[gray!65, line width=0.15pt] (3.6728,3.9322) -- (3.6728,4.0678);
\draw[gray!65, line width=0.15pt] (3.6728,4.0678) -- (3.6610,4.0740);
\draw[gray!65, line width=0.15pt] (3.8644,4.1415) -- (3.7839,4.2034);
\draw[gray!65, line width=0.15pt] (3.6610,4.3130) -- (3.7011,4.3390);
\draw[gray!65, line width=0.15pt] (3.7839,4.2034) -- (3.8644,4.2582);
\draw[gray!65, line width=0.15pt] (3.7011,4.3390) -- (3.7122,4.4746);
\draw[gray!65, line width=0.15pt] (3.7122,4.4746) -- (3.6610,4.5252);
\draw[gray!65, line width=0.15pt] (3.8644,4.6067) -- (3.8553,4.6102);
\draw[gray!65, line width=0.15pt] (3.6610,4.7045) -- (3.7115,4.7458);
\draw[gray!65, line width=0.15pt] (3.8553,4.6102) -- (3.8644,4.6148);
\draw[gray!65, line width=0.15pt] (3.6610,4.8564) -- (3.7115,4.7458);
\draw[gray!65, line width=0.15pt] (3.6610,4.9088) -- (3.7056,5.0169);
\draw[gray!65, line width=0.15pt] (3.7056,5.0169) -- (3.6610,5.0516);
\draw[gray!65, line width=0.15pt] (3.8644,5.1205) -- (3.8086,5.1525);
\draw[gray!65, line width=0.15pt] (3.8086,5.1525) -- (3.8638,5.2881);
\draw[gray!65, line width=0.15pt] (3.8638,5.2881) -- (3.8257,5.4237);
\draw[gray!65, line width=0.15pt] (3.6610,5.5494) -- (3.6759,5.5593);
\draw[gray!65, line width=0.15pt] (3.8257,5.4237) -- (3.8644,5.4495);
\draw[gray!65, line width=0.15pt] (3.6610,5.5851) -- (3.6759,5.5593);
\draw[gray!65, line width=0.15pt] (3.8644,2.5505) -- (4.0678,2.5759);
\draw[gray!65, line width=0.15pt] (4.0678,2.8583) -- (3.8644,2.8795);
\draw[gray!65, line width=0.15pt] (3.8644,3.3852) -- (4.0678,3.3382);
\draw[gray!65, line width=0.15pt] (4.0678,3.4632) -- (3.8644,3.3933);
\draw[gray!65, line width=0.15pt] (3.8644,3.7418) -- (4.0678,3.6992);
\draw[gray!65, line width=0.15pt] (4.0319,3.9322) -- (3.8644,3.8585);
\draw[gray!65, line width=0.15pt] (4.0319,3.9322) -- (4.0319,4.0678);
\draw[gray!65, line width=0.15pt] (3.8644,4.1415) -- (4.0319,4.0678);
\draw[gray!65, line width=0.15pt] (4.0678,4.3008) -- (3.8644,4.2582);
\draw[gray!65, line width=0.15pt] (3.8644,4.6067) -- (4.0678,4.5368);
\draw[gray!65, line width=0.15pt] (4.0678,4.6618) -- (3.8644,4.6148);
\draw[gray!65, line width=0.15pt] (3.8644,5.1205) -- (4.0678,5.1417);
\draw[gray!65, line width=0.15pt] (4.0678,5.4241) -- (3.8644,5.4495);
\draw[gray!65, line width=0.15pt] (4.0678,2.5759) -- (4.2712,2.5391);
\draw[gray!65, line width=0.15pt] (4.2712,2.8875) -- (4.0678,2.8583);
\draw[gray!65, line width=0.15pt] (4.0678,3.3382) -- (4.2712,3.3769);
\draw[gray!65, line width=0.15pt] (4.2712,3.3988) -- (4.0678,3.4632);
\draw[gray!65, line width=0.15pt] (4.0678,3.6992) -- (4.2712,3.7357);
\draw[gray!65, line width=0.15pt] (4.2712,3.8642) -- (4.0942,3.9322);
\draw[gray!65, line width=0.15pt] (4.0942,3.9322) -- (4.0942,4.0678);
\draw[gray!65, line width=0.15pt] (4.0942,4.0678) -- (4.2712,4.1358);
\draw[gray!65, line width=0.15pt] (4.2712,4.2643) -- (4.0678,4.3008);
\draw[gray!65, line width=0.15pt] (4.0678,4.5368) -- (4.2712,4.6012);
\draw[gray!65, line width=0.15pt] (4.2712,4.6231) -- (4.0678,4.6618);
\draw[gray!65, line width=0.15pt] (4.0678,5.1417) -- (4.2712,5.1125);
\draw[gray!65, line width=0.15pt] (4.2712,5.4609) -- (4.0678,5.4241);
\draw[gray!65, line width=0.15pt] (4.4746,2.3688) -- (4.4226,2.4407);
\draw[gray!65, line width=0.15pt] (4.2712,2.5391) -- (4.3192,2.5763);
\draw[gray!65, line width=0.15pt] (4.4226,2.4407) -- (4.4746,2.4704);
\draw[gray!65, line width=0.15pt] (4.3192,2.5763) -- (4.2874,2.7119);
\draw[gray!65, line width=0.15pt] (4.2874,2.7119) -- (4.3313,2.8475);
\draw[gray!65, line width=0.15pt] (4.3313,2.8475) -- (4.2712,2.8875);
\draw[gray!65, line width=0.15pt] (4.4746,2.9297) -- (4.3945,2.9831);
\draw[gray!65, line width=0.15pt] (4.3945,2.9831) -- (4.4591,3.1186);
\draw[gray!65, line width=0.15pt] (4.4591,3.1186) -- (4.3875,3.2542);
\draw[gray!65, line width=0.15pt] (4.2712,3.3769) -- (4.2931,3.3898);
\draw[gray!65, line width=0.15pt] (4.3875,3.2542) -- (4.4746,3.3167);
\draw[gray!65, line width=0.15pt] (4.2931,3.3898) -- (4.2712,3.3988);
\draw[gray!65, line width=0.15pt] (4.4746,3.4478) -- (4.3849,3.5254);
\draw[gray!65, line width=0.15pt] (4.3849,3.5254) -- (4.3900,3.6610);
\draw[gray!65, line width=0.15pt] (4.2712,3.7357) -- (4.3446,3.7966);
\draw[gray!65, line width=0.15pt] (4.3900,3.6610) -- (4.4746,3.7058);
\draw[gray!65, line width=0.15pt] (4.3446,3.7966) -- (4.2712,3.8642);
\draw[gray!65, line width=0.15pt] (4.4746,3.9002) -- (4.4030,3.9322);
\draw[gray!65, line width=0.15pt] (4.4030,3.9322) -- (4.4030,4.0678);
\draw[gray!65, line width=0.15pt] (4.2712,4.1358) -- (4.3446,4.2034);
\draw[gray!65, line width=0.15pt] (4.4030,4.0678) -- (4.4746,4.0998);
\draw[gray!65, line width=0.15pt] (4.3446,4.2034) -- (4.2712,4.2643);
\draw[gray!65, line width=0.15pt] (4.4746,4.2942) -- (4.3900,4.3390);
\draw[gray!65, line width=0.15pt] (4.3900,4.3390) -- (4.3849,4.4746);
\draw[gray!65, line width=0.15pt] (4.2712,4.6012) -- (4.2931,4.6102);
\draw[gray!65, line width=0.15pt] (4.3849,4.4746) -- (4.4746,4.5522);
\draw[gray!65, line width=0.15pt] (4.2931,4.6102) -- (4.2712,4.6231);
\draw[gray!65, line width=0.15pt] (4.4746,4.6833) -- (4.3875,4.7458);
\draw[gray!65, line width=0.15pt] (4.3875,4.7458) -- (4.4591,4.8814);
\draw[gray!65, line width=0.15pt] (4.4591,4.8814) -- (4.3945,5.0169);
\draw[gray!65, line width=0.15pt] (4.2712,5.1125) -- (4.3313,5.1525);
\draw[gray!65, line width=0.15pt] (4.3945,5.0169) -- (4.4746,5.0703);
\draw[gray!65, line width=0.15pt] (4.3313,5.1525) -- (4.2874,5.2881);
\draw[gray!65, line width=0.15pt] (4.2874,5.2881) -- (4.3192,5.4237);
\draw[gray!65, line width=0.15pt] (4.3192,5.4237) -- (4.2712,5.4609);
\draw[gray!65, line width=0.15pt] (4.4746,5.5296) -- (4.4226,5.5593);
\draw[gray!65, line width=0.15pt] (4.4226,5.5593) -- (4.4746,5.6312);
\draw[gray!65, line width=0.15pt] (4.6780,2.1518) -- (4.5085,2.1695);
\draw[gray!65, line width=0.15pt] (4.5085,2.1695) -- (4.5789,2.3051);
\draw[gray!65, line width=0.15pt] (4.4746,2.3688) -- (4.6369,2.4407);
\draw[gray!65, line width=0.15pt] (4.5789,2.3051) -- (4.6780,2.4094);
\draw[gray!65, line width=0.15pt] (4.4746,2.4704) -- (4.6369,2.4407);
\draw[gray!65, line width=0.15pt] (4.4746,2.9297) -- (4.6780,2.9727);
\draw[gray!65, line width=0.15pt] (4.6780,3.2975) -- (4.4746,3.3167);
\draw[gray!65, line width=0.15pt] (4.4746,3.4478) -- (4.6186,3.5254);
\draw[gray!65, line width=0.15pt] (4.6186,3.5254) -- (4.5683,3.6610);
\draw[gray!65, line width=0.15pt] (4.4746,3.7058) -- (4.5683,3.6610);
\draw[gray!65, line width=0.15pt] (4.4746,3.9002) -- (4.5256,3.9322);
\draw[gray!65, line width=0.15pt] (4.5256,3.9322) -- (4.5256,4.0678);
\draw[gray!65, line width=0.15pt] (4.4746,4.0998) -- (4.5256,4.0678);
\draw[gray!65, line width=0.15pt] (4.4746,4.2942) -- (4.5683,4.3390);
\draw[gray!65, line width=0.15pt] (4.5683,4.3390) -- (4.6186,4.4746);
\draw[gray!65, line width=0.15pt] (4.4746,4.5522) -- (4.6186,4.4746);
\draw[gray!65, line width=0.15pt] (4.4746,4.6833) -- (4.6780,4.7025);
\draw[gray!65, line width=0.15pt] (4.6780,5.0273) -- (4.4746,5.0703);
\draw[gray!65, line width=0.15pt] (4.4746,5.5296) -- (4.6369,5.5593);
\draw[gray!65, line width=0.15pt] (4.6369,5.5593) -- (4.4746,5.6312);
\draw[gray!65, line width=0.15pt] (4.6780,5.5906) -- (4.5789,5.6949);
\draw[gray!65, line width=0.15pt] (4.5789,5.6949) -- (4.5085,5.8305);
\draw[gray!65, line width=0.15pt] (4.5085,5.8305) -- (4.6780,5.8482);
\draw[gray!65, line width=0.15pt] (4.6780,2.1518) -- (4.8814,2.1385);
\draw[gray!65, line width=0.15pt] (4.7153,2.4407) -- (4.6780,2.4094);
\draw[gray!65, line width=0.15pt] (4.7153,2.4407) -- (4.8814,2.4743);
\draw[gray!65, line width=0.15pt] (4.6780,2.9727) -- (4.8814,2.9250);
\draw[gray!65, line width=0.15pt] (4.8814,3.3114) -- (4.6780,3.2975);
\draw[gray!65, line width=0.15pt] (4.8814,3.4829) -- (4.7482,3.5254);
\draw[gray!65, line width=0.15pt] (4.7482,3.5254) -- (4.8430,3.6610);
\draw[gray!65, line width=0.15pt] (4.8430,3.6610) -- (4.8814,3.6742);
\draw[gray!65, line width=0.15pt] (4.8814,4.3258) -- (4.8430,4.3390);
\draw[gray!65, line width=0.15pt] (4.8430,4.3390) -- (4.7482,4.4746);
\draw[gray!65, line width=0.15pt] (4.7482,4.4746) -- (4.8814,4.5171);
\draw[gray!65, line width=0.15pt] (4.6780,4.7025) -- (4.8814,4.6886);
\draw[gray!65, line width=0.15pt] (4.8814,5.0750) -- (4.6780,5.0273);
\draw[gray!65, line width=0.15pt] (4.8814,5.5257) -- (4.7153,5.5593);
\draw[gray!65, line width=0.15pt] (4.6780,5.5906) -- (4.7153,5.5593);
\draw[gray!65, line width=0.15pt] (4.8814,5.8615) -- (4.6780,5.8482);
\draw[gray!65, line width=0.15pt] (5.0847,2.0043) -- (5.0507,2.0339);
\draw[gray!65, line width=0.15pt] (4.8814,2.1385) -- (4.9184,2.1695);
\draw[gray!65, line width=0.15pt] (5.0507,2.0339) -- (5.0847,2.0517);
\draw[gray!65, line width=0.15pt] (4.9184,2.1695) -- (4.9011,2.3051);
\draw[gray!65, line width=0.15pt] (4.9011,2.3051) -- (4.9432,2.4407);
\draw[gray!65, line width=0.15pt] (4.9432,2.4407) -- (4.8814,2.4743);
\draw[gray!65, line width=0.15pt] (5.0847,2.5702) -- (5.0778,2.5763);
\draw[gray!65, line width=0.15pt] (5.0778,2.5763) -- (5.0073,2.7119);
\draw[gray!65, line width=0.15pt] (5.0073,2.7119) -- (5.0653,2.8475);
\draw[gray!65, line width=0.15pt] (4.8814,2.9250) -- (4.9750,2.9831);
\draw[gray!65, line width=0.15pt] (5.0653,2.8475) -- (5.0847,2.8621);
\draw[gray!65, line width=0.15pt] (4.9750,2.9831) -- (4.9463,3.1186);
\draw[gray!65, line width=0.15pt] (4.9463,3.1186) -- (4.9671,3.2542);
\draw[gray!65, line width=0.15pt] (4.9671,3.2542) -- (4.8814,3.3114);
\draw[gray!65, line width=0.15pt] (5.0847,3.3214) -- (4.9821,3.3898);
\draw[gray!65, line width=0.15pt] (4.8814,3.4829) -- (5.0184,3.5254);
\draw[gray!65, line width=0.15pt] (4.9821,3.3898) -- (5.0847,3.5012);
\draw[gray!65, line width=0.15pt] (5.0184,3.5254) -- (5.0847,3.5886);
\draw[gray!65, line width=0.15pt] (5.0124,3.7966) -- (4.8814,3.6742);
\draw[gray!65, line width=0.15pt] (5.0124,3.7966) -- (4.9363,3.9322);
\draw[gray!65, line width=0.15pt] (4.9363,3.9322) -- (4.9363,4.0678);
\draw[gray!65, line width=0.15pt] (4.9363,4.0678) -- (5.0124,4.2034);
\draw[gray!65, line width=0.15pt] (4.8814,4.3258) -- (5.0124,4.2034);
\draw[gray!65, line width=0.15pt] (5.0847,4.4114) -- (5.0184,4.4746);
\draw[gray!65, line width=0.15pt] (5.0184,4.4746) -- (4.8814,4.5171);
\draw[gray!65, line width=0.15pt] (5.0847,4.4988) -- (4.9821,4.6102);
\draw[gray!65, line width=0.15pt] (4.8814,4.6886) -- (4.9671,4.7458);
\draw[gray!65, line width=0.15pt] (4.9821,4.6102) -- (5.0847,4.6786);
\draw[gray!65, line width=0.15pt] (4.9671,4.7458) -- (4.9463,4.8814);
\draw[gray!65, line width=0.15pt] (4.9463,4.8814) -- (4.9750,5.0169);
\draw[gray!65, line width=0.15pt] (4.9750,5.0169) -- (4.8814,5.0750);
\draw[gray!65, line width=0.15pt] (5.0847,5.1379) -- (5.0653,5.1525);
\draw[gray!65, line width=0.15pt] (5.0653,5.1525) -- (5.0073,5.2881);
\draw[gray!65, line width=0.15pt] (5.0073,5.2881) -- (5.0778,5.4237);
\draw[gray!65, line width=0.15pt] (4.8814,5.5257) -- (4.9432,5.5593);
\draw[gray!65, line width=0.15pt] (5.0778,5.4237) -- (5.0847,5.4298);
\draw[gray!65, line width=0.15pt] (4.9432,5.5593) -- (4.9011,5.6949);
\draw[gray!65, line width=0.15pt] (4.9011,5.6949) -- (4.9184,5.8305);
\draw[gray!65, line width=0.15pt] (4.9184,5.8305) -- (4.8814,5.8615);
\draw[gray!65, line width=0.15pt] (5.0847,5.9483) -- (5.0507,5.9661);
\draw[gray!65, line width=0.15pt] (5.0507,5.9661) -- (5.0847,5.9957);
\draw[gray!65, line width=0.15pt] (5.2881,1.7578) -- (5.2777,1.7627);
\draw[gray!65, line width=0.15pt] (5.2777,1.7627) -- (5.2623,1.8983);
\draw[gray!65, line width=0.15pt] (5.0847,2.0043) -- (5.1574,2.0339);
\draw[gray!65, line width=0.15pt] (5.2623,1.8983) -- (5.2881,1.9287);
\draw[gray!65, line width=0.15pt] (5.0847,2.0517) -- (5.1574,2.0339);
\draw[gray!65, line width=0.15pt] (5.2881,2.3774) -- (5.2122,2.4407);
\draw[gray!65, line width=0.15pt] (5.0847,2.5702) -- (5.0900,2.5763);
\draw[gray!65, line width=0.15pt] (5.2122,2.4407) -- (5.2881,2.4748);
\draw[gray!65, line width=0.15pt] (5.0900,2.5763) -- (5.1948,2.7119);
\draw[gray!65, line width=0.15pt] (5.1948,2.7119) -- (5.0983,2.8475);
\draw[gray!65, line width=0.15pt] (5.0983,2.8475) -- (5.0847,2.8621);
\draw[gray!65, line width=0.15pt] (5.2881,2.9233) -- (5.1718,2.9831);
\draw[gray!65, line width=0.15pt] (5.1718,2.9831) -- (5.2881,3.0791);
\draw[gray!65, line width=0.15pt] (5.2881,3.1654) -- (5.2243,3.2542);
\draw[gray!65, line width=0.15pt] (5.0847,3.3214) -- (5.2518,3.3898);
\draw[gray!65, line width=0.15pt] (5.2243,3.2542) -- (5.2881,3.3456);
\draw[gray!65, line width=0.15pt] (5.2518,3.3898) -- (5.0847,3.5012);
\draw[gray!65, line width=0.15pt] (5.2881,3.4087) -- (5.1130,3.5254);
\draw[gray!65, line width=0.15pt] (5.0847,3.5886) -- (5.1130,3.5254);
\draw[gray!65, line width=0.15pt] (5.2881,3.7229) -- (5.1428,3.7966);
\draw[gray!65, line width=0.15pt] (5.1428,3.7966) -- (5.2881,3.8730);
\draw[gray!65, line width=0.15pt] (5.2881,4.1270) -- (5.1428,4.2034);
\draw[gray!65, line width=0.15pt] (5.1428,4.2034) -- (5.2881,4.2771);
\draw[gray!65, line width=0.15pt] (5.1130,4.4746) -- (5.0847,4.4114);
\draw[gray!65, line width=0.15pt] (5.0847,4.4988) -- (5.2518,4.6102);
\draw[gray!65, line width=0.15pt] (5.1130,4.4746) -- (5.2881,4.5913);
\draw[gray!65, line width=0.15pt] (5.2518,4.6102) -- (5.0847,4.6786);
\draw[gray!65, line width=0.15pt] (5.2881,4.6544) -- (5.2243,4.7458);
\draw[gray!65, line width=0.15pt] (5.2243,4.7458) -- (5.2881,4.8346);
\draw[gray!65, line width=0.15pt] (5.2881,4.9209) -- (5.1718,5.0169);
\draw[gray!65, line width=0.15pt] (5.0847,5.1379) -- (5.0983,5.1525);
\draw[gray!65, line width=0.15pt] (5.1718,5.0169) -- (5.2881,5.0767);
\draw[gray!65, line width=0.15pt] (5.0983,5.1525) -- (5.1948,5.2881);
\draw[gray!65, line width=0.15pt] (5.1948,5.2881) -- (5.0900,5.4237);
\draw[gray!65, line width=0.15pt] (5.0900,5.4237) -- (5.0847,5.4298);
\draw[gray!65, line width=0.15pt] (5.2881,5.5252) -- (5.2122,5.5593);
\draw[gray!65, line width=0.15pt] (5.2122,5.5593) -- (5.2881,5.6226);
\draw[gray!65, line width=0.15pt] (5.0847,5.9483) -- (5.1574,5.9661);
\draw[gray!65, line width=0.15pt] (5.1574,5.9661) -- (5.0847,5.9957);
\draw[gray!65, line width=0.15pt] (5.2881,6.0713) -- (5.2623,6.1017);
\draw[gray!65, line width=0.15pt] (5.2623,6.1017) -- (5.2777,6.2373);
\draw[gray!65, line width=0.15pt] (5.2777,6.2373) -- (5.2881,6.2422);
\draw[gray!65, line width=0.15pt] (5.2881,1.7578) -- (5.4915,1.7600);
\draw[gray!65, line width=0.15pt] (5.4915,1.9252) -- (5.2881,1.9287);
\draw[gray!65, line width=0.15pt] (5.2881,2.3774) -- (5.4915,2.4019);
\draw[gray!65, line width=0.15pt] (5.4915,2.4674) -- (5.2881,2.4748);
\draw[gray!65, line width=0.15pt] (5.2881,2.9233) -- (5.4915,2.9267);
\draw[gray!65, line width=0.15pt] (5.4915,3.0455) -- (5.2881,3.0791);
\draw[gray!65, line width=0.15pt] (5.2881,3.1654) -- (5.4915,3.2287);
\draw[gray!65, line width=0.15pt] (5.3830,3.3898) -- (5.2881,3.3456);
\draw[gray!65, line width=0.15pt] (5.2881,3.4087) -- (5.3830,3.3898);
\draw[gray!65, line width=0.15pt] (5.4915,3.7337) -- (5.2881,3.7229);
\draw[gray!65, line width=0.15pt] (5.2881,3.8730) -- (5.4915,3.8572);
\draw[gray!65, line width=0.15pt] (5.4915,4.1428) -- (5.2881,4.1270);
\draw[gray!65, line width=0.15pt] (5.2881,4.2771) -- (5.4915,4.2663);
\draw[gray!65, line width=0.15pt] (5.3830,4.6102) -- (5.2881,4.5913);
\draw[gray!65, line width=0.15pt] (5.2881,4.6544) -- (5.3830,4.6102);
\draw[gray!65, line width=0.15pt] (5.4915,4.7713) -- (5.2881,4.8346);
\draw[gray!65, line width=0.15pt] (5.2881,4.9209) -- (5.4915,4.9545);
\draw[gray!65, line width=0.15pt] (5.4915,5.0733) -- (5.2881,5.0767);
\draw[gray!65, line width=0.15pt] (5.2881,5.5252) -- (5.4915,5.5326);
\draw[gray!65, line width=0.15pt] (5.4915,5.5981) -- (5.2881,5.6226);
\draw[gray!65, line width=0.15pt] (5.2881,6.0713) -- (5.4915,6.0748);
\draw[gray!65, line width=0.15pt] (5.4915,6.2400) -- (5.2881,6.2422);
\draw[gray!65, line width=0.15pt] (5.4915,1.7600) -- (5.4942,1.7627);
\draw[gray!65, line width=0.15pt] (5.4942,1.7627) -- (5.4960,1.8983);
\draw[gray!65, line width=0.15pt] (5.4915,1.9252) -- (5.4960,1.8983);
\draw[gray!65, line width=0.15pt] (5.4915,2.4019) -- (5.5305,2.4407);
\draw[gray!65, line width=0.15pt] (5.5305,2.4407) -- (5.4915,2.4674);
\draw[gray!65, line width=0.15pt] (5.6949,2.5226) -- (5.6127,2.5763);
\draw[gray!65, line width=0.15pt] (5.6127,2.5763) -- (5.5488,2.7119);
\draw[gray!65, line width=0.15pt] (5.5488,2.7119) -- (5.6036,2.8475);
\draw[gray!65, line width=0.15pt] (5.4915,2.9267) -- (5.5587,2.9831);
\draw[gray!65, line width=0.15pt] (5.6036,2.8475) -- (5.6949,2.8964);
\draw[gray!65, line width=0.15pt] (5.4915,3.0455) -- (5.5587,2.9831);
\draw[gray!65, line width=0.15pt] (5.4915,3.2287) -- (5.5113,3.2542);
\draw[gray!65, line width=0.15pt] (5.5113,3.2542) -- (5.6949,3.3416);
\draw[gray!65, line width=0.15pt] (5.6949,3.4285) -- (5.5843,3.5254);
\draw[gray!65, line width=0.15pt] (5.5843,3.5254) -- (5.6005,3.6610);
\draw[gray!65, line width=0.15pt] (5.6005,3.6610) -- (5.4915,3.7337);
\draw[gray!65, line width=0.15pt] (5.6949,3.7180) -- (5.5769,3.7966);
\draw[gray!65, line width=0.15pt] (5.4915,3.8572) -- (5.6218,3.9322);
\draw[gray!65, line width=0.15pt] (5.5769,3.7966) -- (5.6949,3.8868);
\draw[gray!65, line width=0.15pt] (5.6218,3.9322) -- (5.6218,4.0678);
\draw[gray!65, line width=0.15pt] (5.6218,4.0678) -- (5.4915,4.1428);
\draw[gray!65, line width=0.15pt] (5.6949,4.1132) -- (5.5769,4.2034);
\draw[gray!65, line width=0.15pt] (5.4915,4.2663) -- (5.6005,4.3390);
\draw[gray!65, line width=0.15pt] (5.5769,4.2034) -- (5.6949,4.2820);
\draw[gray!65, line width=0.15pt] (5.6005,4.3390) -- (5.5843,4.4746);
\draw[gray!65, line width=0.15pt] (5.5843,4.4746) -- (5.6949,4.5715);
\draw[gray!65, line width=0.15pt] (5.6949,4.6584) -- (5.5113,4.7458);
\draw[gray!65, line width=0.15pt] (5.4915,4.7713) -- (5.5113,4.7458);
\draw[gray!65, line width=0.15pt] (5.4915,4.9545) -- (5.5587,5.0169);
\draw[gray!65, line width=0.15pt] (5.5587,5.0169) -- (5.4915,5.0733);
\draw[gray!65, line width=0.15pt] (5.6949,5.1036) -- (5.6036,5.1525);
\draw[gray!65, line width=0.15pt] (5.6036,5.1525) -- (5.5488,5.2881);
\draw[gray!65, line width=0.15pt] (5.5488,5.2881) -- (5.6127,5.4237);
\draw[gray!65, line width=0.15pt] (5.4915,5.5326) -- (5.5305,5.5593);
\draw[gray!65, line width=0.15pt] (5.6127,5.4237) -- (5.6949,5.4774);
\draw[gray!65, line width=0.15pt] (5.4915,5.5981) -- (5.5305,5.5593);
\draw[gray!65, line width=0.15pt] (5.4915,6.0748) -- (5.4960,6.1017);
\draw[gray!65, line width=0.15pt] (5.4960,6.1017) -- (5.4942,6.2373);
\draw[gray!65, line width=0.15pt] (5.4915,6.2400) -- (5.4942,6.2373);
\draw[gray!65, line width=0.15pt] (5.8983,1.9650) -- (5.8320,2.0339);
\draw[gray!65, line width=0.15pt] (5.8320,2.0339) -- (5.8983,2.0518);
\draw[gray!65, line width=0.15pt] (5.8983,2.4326) -- (5.8890,2.4407);
\draw[gray!65, line width=0.15pt] (5.6949,2.5226) -- (5.7877,2.5763);
\draw[gray!65, line width=0.15pt] (5.8890,2.4407) -- (5.8983,2.4485);
\draw[gray!65, line width=0.15pt] (5.7877,2.5763) -- (5.8983,2.6879);
\draw[gray!65, line width=0.15pt] (5.8983,2.7295) -- (5.7963,2.8475);
\draw[gray!65, line width=0.15pt] (5.7963,2.8475) -- (5.6949,2.8964);
\draw[gray!65, line width=0.15pt] (5.8983,2.9353) -- (5.8503,2.9831);
\draw[gray!65, line width=0.15pt] (5.8503,2.9831) -- (5.8983,3.0171);
\draw[gray!65, line width=0.15pt] (5.6949,3.3416) -- (5.8983,3.2909);
\draw[gray!65, line width=0.15pt] (5.8095,3.5254) -- (5.6949,3.4285);
\draw[gray!65, line width=0.15pt] (5.8095,3.5254) -- (5.7857,3.6610);
\draw[gray!65, line width=0.15pt] (5.6949,3.7180) -- (5.8302,3.7966);
\draw[gray!65, line width=0.15pt] (5.7857,3.6610) -- (5.8983,3.7479);
\draw[gray!65, line width=0.15pt] (5.8302,3.7966) -- (5.6949,3.8868);
\draw[gray!65, line width=0.15pt] (5.8983,3.8392) -- (5.7588,3.9322);
\draw[gray!65, line width=0.15pt] (5.7588,3.9322) -- (5.7588,4.0678);
\draw[gray!65, line width=0.15pt] (5.6949,4.1132) -- (5.8302,4.2034);
\draw[gray!65, line width=0.15pt] (5.7588,4.0678) -- (5.8983,4.1608);
\draw[gray!65, line width=0.15pt] (5.8302,4.2034) -- (5.6949,4.2820);
\draw[gray!65, line width=0.15pt] (5.8983,4.2521) -- (5.7857,4.3390);
\draw[gray!65, line width=0.15pt] (5.7857,4.3390) -- (5.8095,4.4746);
\draw[gray!65, line width=0.15pt] (5.6949,4.5715) -- (5.8095,4.4746);
\draw[gray!65, line width=0.15pt] (5.8983,4.7091) -- (5.6949,4.6584);
\draw[gray!65, line width=0.15pt] (5.8983,4.9829) -- (5.8503,5.0169);
\draw[gray!65, line width=0.15pt] (5.6949,5.1036) -- (5.7963,5.1525);
\draw[gray!65, line width=0.15pt] (5.8503,5.0169) -- (5.8983,5.0647);
\draw[gray!65, line width=0.15pt] (5.7963,5.1525) -- (5.8983,5.2705);
\draw[gray!65, line width=0.15pt] (5.8983,5.3121) -- (5.7877,5.4237);
\draw[gray!65, line width=0.15pt] (5.7877,5.4237) -- (5.6949,5.4774);
\draw[gray!65, line width=0.15pt] (5.8983,5.5515) -- (5.8890,5.5593);
\draw[gray!65, line width=0.15pt] (5.8890,5.5593) -- (5.8983,5.5674);
\draw[gray!65, line width=0.15pt] (5.8983,5.9482) -- (5.8320,5.9661);
\draw[gray!65, line width=0.15pt] (5.8320,5.9661) -- (5.8983,6.0350);
\draw[gray!65, line width=0.15pt] (5.8983,1.9650) -- (6.1017,1.9650);
\draw[gray!65, line width=0.15pt] (6.1017,2.0518) -- (5.8983,2.0518);
\draw[gray!65, line width=0.15pt] (5.8983,2.4326) -- (6.1017,2.4326);
\draw[gray!65, line width=0.15pt] (6.1017,2.4485) -- (5.8983,2.4485);
\draw[gray!65, line width=0.15pt] (5.8983,2.6879) -- (6.1017,2.6879);
\draw[gray!65, line width=0.15pt] (6.1017,2.7295) -- (5.8983,2.7295);
\draw[gray!65, line width=0.15pt] (5.8983,2.9353) -- (6.1017,2.9353);
\draw[gray!65, line width=0.15pt] (6.1017,3.0171) -- (5.8983,3.0171);
\draw[gray!65, line width=0.15pt] (5.8983,3.2909) -- (6.1017,3.2909);
\draw[gray!65, line width=0.15pt] (6.1017,3.7479) -- (5.8983,3.7479);
\draw[gray!65, line width=0.15pt] (5.8983,3.8392) -- (6.1017,3.8392);
\draw[gray!65, line width=0.15pt] (6.1017,4.1608) -- (5.8983,4.1608);
\draw[gray!65, line width=0.15pt] (5.8983,4.2521) -- (6.1017,4.2521);
\draw[gray!65, line width=0.15pt] (6.1017,4.7091) -- (5.8983,4.7091);
\draw[gray!65, line width=0.15pt] (5.8983,4.9829) -- (6.1017,4.9829);
\draw[gray!65, line width=0.15pt] (6.1017,5.0647) -- (5.8983,5.0647);
\draw[gray!65, line width=0.15pt] (5.8983,5.2705) -- (6.1017,5.2705);
\draw[gray!65, line width=0.15pt] (6.1017,5.3121) -- (5.8983,5.3121);
\draw[gray!65, line width=0.15pt] (5.8983,5.5515) -- (6.1017,5.5515);
\draw[gray!65, line width=0.15pt] (6.1017,5.5674) -- (5.8983,5.5674);
\draw[gray!65, line width=0.15pt] (5.8983,5.9482) -- (6.1017,5.9482);
\draw[gray!65, line width=0.15pt] (6.1017,6.0350) -- (5.8983,6.0350);
\draw[gray!65, line width=0.15pt] (6.1017,1.9650) -- (6.1680,2.0339);
\draw[gray!65, line width=0.15pt] (6.1017,2.0518) -- (6.1680,2.0339);
\draw[gray!65, line width=0.15pt] (6.1017,2.4326) -- (6.1110,2.4407);
\draw[gray!65, line width=0.15pt] (6.1110,2.4407) -- (6.1017,2.4485);
\draw[gray!65, line width=0.15pt] (6.3051,2.5226) -- (6.2123,2.5763);
\draw[gray!65, line width=0.15pt] (6.1017,2.6879) -- (6.2123,2.5763);
\draw[gray!65, line width=0.15pt] (6.2037,2.8475) -- (6.1017,2.7295);
\draw[gray!65, line width=0.15pt] (6.1017,2.9353) -- (6.1497,2.9831);
\draw[gray!65, line width=0.15pt] (6.2037,2.8475) -- (6.3051,2.8964);
\draw[gray!65, line width=0.15pt] (6.1017,3.0171) -- (6.1497,2.9831);
\draw[gray!65, line width=0.15pt] (6.1017,3.2909) -- (6.3051,3.3416);
\draw[gray!65, line width=0.15pt] (6.3051,3.4285) -- (6.1905,3.5254);
\draw[gray!65, line width=0.15pt] (6.1905,3.5254) -- (6.2143,3.6610);
\draw[gray!65, line width=0.15pt] (6.2143,3.6610) -- (6.1017,3.7479);
\draw[gray!65, line width=0.15pt] (6.3051,3.7180) -- (6.1698,3.7966);
\draw[gray!65, line width=0.15pt] (6.1017,3.8392) -- (6.2412,3.9322);
\draw[gray!65, line width=0.15pt] (6.1698,3.7966) -- (6.3051,3.8868);
\draw[gray!65, line width=0.15pt] (6.2412,3.9322) -- (6.2412,4.0678);
\draw[gray!65, line width=0.15pt] (6.2412,4.0678) -- (6.1017,4.1608);
\draw[gray!65, line width=0.15pt] (6.3051,4.1132) -- (6.1698,4.2034);
\draw[gray!65, line width=0.15pt] (6.1017,4.2521) -- (6.2143,4.3390);
\draw[gray!65, line width=0.15pt] (6.1698,4.2034) -- (6.3051,4.2820);
\draw[gray!65, line width=0.15pt] (6.2143,4.3390) -- (6.1905,4.4746);
\draw[gray!65, line width=0.15pt] (6.1905,4.4746) -- (6.3051,4.5715);
\draw[gray!65, line width=0.15pt] (6.3051,4.6584) -- (6.1017,4.7091);
\draw[gray!65, line width=0.15pt] (6.1017,4.9829) -- (6.1497,5.0169);
\draw[gray!65, line width=0.15pt] (6.1497,5.0169) -- (6.1017,5.0647);
\draw[gray!65, line width=0.15pt] (6.3051,5.1036) -- (6.2037,5.1525);
\draw[gray!65, line width=0.15pt] (6.1017,5.2705) -- (6.2037,5.1525);
\draw[gray!65, line width=0.15pt] (6.2123,5.4237) -- (6.1017,5.3121);
\draw[gray!65, line width=0.15pt] (6.1017,5.5515) -- (6.1110,5.5593);
\draw[gray!65, line width=0.15pt] (6.2123,5.4237) -- (6.3051,5.4774);
\draw[gray!65, line width=0.15pt] (6.1017,5.5674) -- (6.1110,5.5593);
\draw[gray!65, line width=0.15pt] (6.1017,5.9482) -- (6.1680,5.9661);
\draw[gray!65, line width=0.15pt] (6.1017,6.0350) -- (6.1680,5.9661);
\draw[gray!65, line width=0.15pt] (6.5085,1.7600) -- (6.5058,1.7627);
\draw[gray!65, line width=0.15pt] (6.5058,1.7627) -- (6.5040,1.8983);
\draw[gray!65, line width=0.15pt] (6.5040,1.8983) -- (6.5085,1.9252);
\draw[gray!65, line width=0.15pt] (6.5085,2.4019) -- (6.4695,2.4407);
\draw[gray!65, line width=0.15pt] (6.3051,2.5226) -- (6.3873,2.5763);
\draw[gray!65, line width=0.15pt] (6.4695,2.4407) -- (6.5085,2.4674);
\draw[gray!65, line width=0.15pt] (6.3873,2.5763) -- (6.4512,2.7119);
\draw[gray!65, line width=0.15pt] (6.4512,2.7119) -- (6.3964,2.8475);
\draw[gray!65, line width=0.15pt] (6.3964,2.8475) -- (6.3051,2.8964);
\draw[gray!65, line width=0.15pt] (6.5085,2.9267) -- (6.4413,2.9831);
\draw[gray!65, line width=0.15pt] (6.4413,2.9831) -- (6.5085,3.0455);
\draw[gray!65, line width=0.15pt] (6.5085,3.2287) -- (6.4887,3.2542);
\draw[gray!65, line width=0.15pt] (6.3051,3.3416) -- (6.4887,3.2542);
\draw[gray!65, line width=0.15pt] (6.4157,3.5254) -- (6.3051,3.4285);
\draw[gray!65, line width=0.15pt] (6.4157,3.5254) -- (6.3995,3.6610);
\draw[gray!65, line width=0.15pt] (6.3051,3.7180) -- (6.4231,3.7966);
\draw[gray!65, line width=0.15pt] (6.3995,3.6610) -- (6.5085,3.7337);
\draw[gray!65, line width=0.15pt] (6.4231,3.7966) -- (6.3051,3.8868);
\draw[gray!65, line width=0.15pt] (6.5085,3.8572) -- (6.3782,3.9322);
\draw[gray!65, line width=0.15pt] (6.3782,3.9322) -- (6.3782,4.0678);
\draw[gray!65, line width=0.15pt] (6.3051,4.1132) -- (6.4231,4.2034);
\draw[gray!65, line width=0.15pt] (6.3782,4.0678) -- (6.5085,4.1428);
\draw[gray!65, line width=0.15pt] (6.4231,4.2034) -- (6.3051,4.2820);
\draw[gray!65, line width=0.15pt] (6.5085,4.2663) -- (6.3995,4.3390);
\draw[gray!65, line width=0.15pt] (6.3995,4.3390) -- (6.4157,4.4746);
\draw[gray!65, line width=0.15pt] (6.3051,4.5715) -- (6.4157,4.4746);
\draw[gray!65, line width=0.15pt] (6.4887,4.7458) -- (6.3051,4.6584);
\draw[gray!65, line width=0.15pt] (6.4887,4.7458) -- (6.5085,4.7713);
\draw[gray!65, line width=0.15pt] (6.5085,4.9545) -- (6.4413,5.0169);
\draw[gray!65, line width=0.15pt] (6.3051,5.1036) -- (6.3964,5.1525);
\draw[gray!65, line width=0.15pt] (6.4413,5.0169) -- (6.5085,5.0733);
\draw[gray!65, line width=0.15pt] (6.3964,5.1525) -- (6.4512,5.2881);
\draw[gray!65, line width=0.15pt] (6.4512,5.2881) -- (6.3873,5.4237);
\draw[gray!65, line width=0.15pt] (6.3873,5.4237) -- (6.3051,5.4774);
\draw[gray!65, line width=0.15pt] (6.5085,5.5326) -- (6.4695,5.5593);
\draw[gray!65, line width=0.15pt] (6.4695,5.5593) -- (6.5085,5.5981);
\draw[gray!65, line width=0.15pt] (6.5085,6.0748) -- (6.5040,6.1017);
\draw[gray!65, line width=0.15pt] (6.5040,6.1017) -- (6.5058,6.2373);
\draw[gray!65, line width=0.15pt] (6.5058,6.2373) -- (6.5085,6.2400);
\draw[gray!65, line width=0.15pt] (6.5085,1.7600) -- (6.7119,1.7578);
\draw[gray!65, line width=0.15pt] (6.7119,1.9287) -- (6.5085,1.9252);
\draw[gray!65, line width=0.15pt] (6.5085,2.4019) -- (6.7119,2.3774);
\draw[gray!65, line width=0.15pt] (6.7119,2.4748) -- (6.5085,2.4674);
\draw[gray!65, line width=0.15pt] (6.5085,2.9267) -- (6.7119,2.9233);
\draw[gray!65, line width=0.15pt] (6.7119,3.0791) -- (6.5085,3.0455);
\draw[gray!65, line width=0.15pt] (6.5085,3.2287) -- (6.7119,3.1654);
\draw[gray!65, line width=0.15pt] (6.7119,3.3456) -- (6.6170,3.3898);
\draw[gray!65, line width=0.15pt] (6.6170,3.3898) -- (6.7119,3.4087);
\draw[gray!65, line width=0.15pt] (6.7119,3.7229) -- (6.5085,3.7337);
\draw[gray!65, line width=0.15pt] (6.5085,3.8572) -- (6.7119,3.8730);
\draw[gray!65, line width=0.15pt] (6.7119,4.1270) -- (6.5085,4.1428);
\draw[gray!65, line width=0.15pt] (6.5085,4.2663) -- (6.7119,4.2771);
\draw[gray!65, line width=0.15pt] (6.7119,4.5913) -- (6.6170,4.6102);
\draw[gray!65, line width=0.15pt] (6.6170,4.6102) -- (6.7119,4.6544);
\draw[gray!65, line width=0.15pt] (6.7119,4.8346) -- (6.5085,4.7713);
\draw[gray!65, line width=0.15pt] (6.5085,4.9545) -- (6.7119,4.9209);
\draw[gray!65, line width=0.15pt] (6.7119,5.0767) -- (6.5085,5.0733);
\draw[gray!65, line width=0.15pt] (6.5085,5.5326) -- (6.7119,5.5252);
\draw[gray!65, line width=0.15pt] (6.7119,5.6226) -- (6.5085,5.5981);
\draw[gray!65, line width=0.15pt] (6.5085,6.0748) -- (6.7119,6.0713);
\draw[gray!65, line width=0.15pt] (6.7119,6.2422) -- (6.5085,6.2400);
\draw[gray!65, line width=0.15pt] (6.7119,1.7578) -- (6.7223,1.7627);
\draw[gray!65, line width=0.15pt] (6.7223,1.7627) -- (6.7377,1.8983);
\draw[gray!65, line width=0.15pt] (6.7377,1.8983) -- (6.7119,1.9287);
\draw[gray!65, line width=0.15pt] (6.9153,2.0043) -- (6.8426,2.0339);
\draw[gray!65, line width=0.15pt] (6.8426,2.0339) -- (6.9153,2.0517);
\draw[gray!65, line width=0.15pt] (6.7119,2.3774) -- (6.7878,2.4407);
\draw[gray!65, line width=0.15pt] (6.7878,2.4407) -- (6.7119,2.4748);
\draw[gray!65, line width=0.15pt] (6.9153,2.5702) -- (6.9100,2.5763);
\draw[gray!65, line width=0.15pt] (6.9100,2.5763) -- (6.8052,2.7119);
\draw[gray!65, line width=0.15pt] (6.8052,2.7119) -- (6.9017,2.8475);
\draw[gray!65, line width=0.15pt] (6.7119,2.9233) -- (6.8282,2.9831);
\draw[gray!65, line width=0.15pt] (6.9017,2.8475) -- (6.9153,2.8621);
\draw[gray!65, line width=0.15pt] (6.7119,3.0791) -- (6.8282,2.9831);
\draw[gray!65, line width=0.15pt] (6.7119,3.1654) -- (6.7757,3.2542);
\draw[gray!65, line width=0.15pt] (6.7757,3.2542) -- (6.7119,3.3456);
\draw[gray!65, line width=0.15pt] (6.9153,3.3214) -- (6.7482,3.3898);
\draw[gray!65, line width=0.15pt] (6.7119,3.4087) -- (6.8870,3.5254);
\draw[gray!65, line width=0.15pt] (6.7482,3.3898) -- (6.9153,3.5012);
\draw[gray!65, line width=0.15pt] (6.8870,3.5254) -- (6.9153,3.5886);
\draw[gray!65, line width=0.15pt] (6.8572,3.7966) -- (6.7119,3.7229);
\draw[gray!65, line width=0.15pt] (6.7119,3.8730) -- (6.8572,3.7966);
\draw[gray!65, line width=0.15pt] (6.8572,4.2034) -- (6.7119,4.1270);
\draw[gray!65, line width=0.15pt] (6.7119,4.2771) -- (6.8572,4.2034);
\draw[gray!65, line width=0.15pt] (6.9153,4.4114) -- (6.8870,4.4746);
\draw[gray!65, line width=0.15pt] (6.8870,4.4746) -- (6.7119,4.5913);
\draw[gray!65, line width=0.15pt] (6.9153,4.4988) -- (6.7482,4.6102);
\draw[gray!65, line width=0.15pt] (6.7119,4.6544) -- (6.7757,4.7458);
\draw[gray!65, line width=0.15pt] (6.7482,4.6102) -- (6.9153,4.6786);
\draw[gray!65, line width=0.15pt] (6.7119,4.8346) -- (6.7757,4.7458);
\draw[gray!65, line width=0.15pt] (6.7119,4.9209) -- (6.8282,5.0169);
\draw[gray!65, line width=0.15pt] (6.8282,5.0169) -- (6.7119,5.0767);
\draw[gray!65, line width=0.15pt] (6.9153,5.1379) -- (6.9017,5.1525);
\draw[gray!65, line width=0.15pt] (6.9017,5.1525) -- (6.8052,5.2881);
\draw[gray!65, line width=0.15pt] (6.8052,5.2881) -- (6.9100,5.4237);
\draw[gray!65, line width=0.15pt] (6.7119,5.5252) -- (6.7878,5.5593);
\draw[gray!65, line width=0.15pt] (6.9100,5.4237) -- (6.9153,5.4298);
\draw[gray!65, line width=0.15pt] (6.7119,5.6226) -- (6.7878,5.5593);
\draw[gray!65, line width=0.15pt] (6.9153,5.9483) -- (6.8426,5.9661);
\draw[gray!65, line width=0.15pt] (6.7119,6.0713) -- (6.7377,6.1017);
\draw[gray!65, line width=0.15pt] (6.8426,5.9661) -- (6.9153,5.9957);
\draw[gray!65, line width=0.15pt] (6.7377,6.1017) -- (6.7223,6.2373);
\draw[gray!65, line width=0.15pt] (6.7119,6.2422) -- (6.7223,6.2373);
\draw[gray!65, line width=0.15pt] (6.9153,2.0043) -- (6.9493,2.0339);
\draw[gray!65, line width=0.15pt] (6.9493,2.0339) -- (6.9153,2.0517);
\draw[gray!65, line width=0.15pt] (7.1186,2.1385) -- (7.0816,2.1695);
\draw[gray!65, line width=0.15pt] (7.0816,2.1695) -- (7.0989,2.3051);
\draw[gray!65, line width=0.15pt] (7.0989,2.3051) -- (7.0568,2.4407);
\draw[gray!65, line width=0.15pt] (6.9153,2.5702) -- (6.9222,2.5763);
\draw[gray!65, line width=0.15pt] (7.0568,2.4407) -- (7.1186,2.4743);
\draw[gray!65, line width=0.15pt] (6.9222,2.5763) -- (6.9927,2.7119);
\draw[gray!65, line width=0.15pt] (6.9927,2.7119) -- (6.9347,2.8475);
\draw[gray!65, line width=0.15pt] (6.9347,2.8475) -- (6.9153,2.8621);
\draw[gray!65, line width=0.15pt] (7.1186,2.9250) -- (7.0250,2.9831);
\draw[gray!65, line width=0.15pt] (7.0250,2.9831) -- (7.0537,3.1186);
\draw[gray!65, line width=0.15pt] (7.0537,3.1186) -- (7.0329,3.2542);
\draw[gray!65, line width=0.15pt] (6.9153,3.3214) -- (7.0179,3.3898);
\draw[gray!65, line width=0.15pt] (7.0329,3.2542) -- (7.1186,3.3114);
\draw[gray!65, line width=0.15pt] (7.0179,3.3898) -- (6.9153,3.5012);
\draw[gray!65, line width=0.15pt] (7.1186,3.4829) -- (6.9816,3.5254);
\draw[gray!65, line width=0.15pt] (6.9153,3.5886) -- (6.9816,3.5254);
\draw[gray!65, line width=0.15pt] (7.1186,3.6742) -- (6.9876,3.7966);
\draw[gray!65, line width=0.15pt] (6.9876,3.7966) -- (7.0637,3.9322);
\draw[gray!65, line width=0.15pt] (7.0637,3.9322) -- (7.0637,4.0678);
\draw[gray!65, line width=0.15pt] (7.0637,4.0678) -- (6.9876,4.2034);
\draw[gray!65, line width=0.15pt] (6.9876,4.2034) -- (7.1186,4.3258);
\draw[gray!65, line width=0.15pt] (6.9816,4.4746) -- (6.9153,4.4114);
\draw[gray!65, line width=0.15pt] (6.9153,4.4988) -- (7.0179,4.6102);
\draw[gray!65, line width=0.15pt] (6.9816,4.4746) -- (7.1186,4.5171);
\draw[gray!65, line width=0.15pt] (7.0179,4.6102) -- (6.9153,4.6786);
\draw[gray!65, line width=0.15pt] (7.1186,4.6886) -- (7.0329,4.7458);
\draw[gray!65, line width=0.15pt] (7.0329,4.7458) -- (7.0537,4.8814);
\draw[gray!65, line width=0.15pt] (7.0537,4.8814) -- (7.0250,5.0169);
\draw[gray!65, line width=0.15pt] (6.9153,5.1379) -- (6.9347,5.1525);
\draw[gray!65, line width=0.15pt] (7.0250,5.0169) -- (7.1186,5.0750);
\draw[gray!65, line width=0.15pt] (6.9347,5.1525) -- (6.9927,5.2881);
\draw[gray!65, line width=0.15pt] (6.9927,5.2881) -- (6.9222,5.4237);
\draw[gray!65, line width=0.15pt] (6.9222,5.4237) -- (6.9153,5.4298);
\draw[gray!65, line width=0.15pt] (7.1186,5.5257) -- (7.0568,5.5593);
\draw[gray!65, line width=0.15pt] (7.0568,5.5593) -- (7.0989,5.6949);
\draw[gray!65, line width=0.15pt] (7.0989,5.6949) -- (7.0816,5.8305);
\draw[gray!65, line width=0.15pt] (6.9153,5.9483) -- (6.9493,5.9661);
\draw[gray!65, line width=0.15pt] (7.0816,5.8305) -- (7.1186,5.8615);
\draw[gray!65, line width=0.15pt] (6.9153,5.9957) -- (6.9493,5.9661);
\draw[gray!65, line width=0.15pt] (7.1186,2.1385) -- (7.3220,2.1518);
\draw[gray!65, line width=0.15pt] (7.3220,2.4094) -- (7.2847,2.4407);
\draw[gray!65, line width=0.15pt] (7.1186,2.4743) -- (7.2847,2.4407);
\draw[gray!65, line width=0.15pt] (7.1186,2.9250) -- (7.3220,2.9727);
\draw[gray!65, line width=0.15pt] (7.3220,3.2975) -- (7.1186,3.3114);
\draw[gray!65, line width=0.15pt] (7.1186,3.4829) -- (7.2518,3.5254);
\draw[gray!65, line width=0.15pt] (7.2518,3.5254) -- (7.1570,3.6610);
\draw[gray!65, line width=0.15pt] (7.1186,3.6742) -- (7.1570,3.6610);
\draw[gray!65, line width=0.15pt] (7.1186,4.3258) -- (7.1570,4.3390);
\draw[gray!65, line width=0.15pt] (7.1570,4.3390) -- (7.2518,4.4746);
\draw[gray!65, line width=0.15pt] (7.1186,4.5171) -- (7.2518,4.4746);
\draw[gray!65, line width=0.15pt] (7.1186,4.6886) -- (7.3220,4.7025);
\draw[gray!65, line width=0.15pt] (7.3220,5.0273) -- (7.1186,5.0750);
\draw[gray!65, line width=0.15pt] (7.1186,5.5257) -- (7.2847,5.5593);
\draw[gray!65, line width=0.15pt] (7.2847,5.5593) -- (7.3220,5.5906);
\draw[gray!65, line width=0.15pt] (7.3220,5.8482) -- (7.1186,5.8615);
\draw[gray!65, line width=0.15pt] (7.3220,2.1518) -- (7.4915,2.1695);
\draw[gray!65, line width=0.15pt] (7.4915,2.1695) -- (7.4211,2.3051);
\draw[gray!65, line width=0.15pt] (7.4211,2.3051) -- (7.3220,2.4094);
\draw[gray!65, line width=0.15pt] (7.5254,2.3688) -- (7.3631,2.4407);
\draw[gray!65, line width=0.15pt] (7.3631,2.4407) -- (7.5254,2.4704);
\draw[gray!65, line width=0.15pt] (7.3220,2.9727) -- (7.5254,2.9297);
\draw[gray!65, line width=0.15pt] (7.5254,3.3167) -- (7.3220,3.2975);
\draw[gray!65, line width=0.15pt] (7.5254,3.4478) -- (7.3814,3.5254);
\draw[gray!65, line width=0.15pt] (7.3814,3.5254) -- (7.4317,3.6610);
\draw[gray!65, line width=0.15pt] (7.4317,3.6610) -- (7.5254,3.7058);
\draw[gray!65, line width=0.15pt] (7.5254,3.9002) -- (7.4744,3.9322);
\draw[gray!65, line width=0.15pt] (7.4744,3.9322) -- (7.4744,4.0678);
\draw[gray!65, line width=0.15pt] (7.4744,4.0678) -- (7.5254,4.0998);
\draw[gray!65, line width=0.15pt] (7.5254,4.2942) -- (7.4317,4.3390);
\draw[gray!65, line width=0.15pt] (7.4317,4.3390) -- (7.3814,4.4746);
\draw[gray!65, line width=0.15pt] (7.3814,4.4746) -- (7.5254,4.5522);
\draw[gray!65, line width=0.15pt] (7.3220,4.7025) -- (7.5254,4.6833);
\draw[gray!65, line width=0.15pt] (7.5254,5.0703) -- (7.3220,5.0273);
\draw[gray!65, line width=0.15pt] (7.5254,5.5296) -- (7.3631,5.5593);
\draw[gray!65, line width=0.15pt] (7.3220,5.5906) -- (7.4211,5.6949);
\draw[gray!65, line width=0.15pt] (7.3631,5.5593) -- (7.5254,5.6312);
\draw[gray!65, line width=0.15pt] (7.4211,5.6949) -- (7.4915,5.8305);
\draw[gray!65, line width=0.15pt] (7.3220,5.8482) -- (7.4915,5.8305);
\draw[gray!65, line width=0.15pt] (7.5254,2.3688) -- (7.5774,2.4407);
\draw[gray!65, line width=0.15pt] (7.5774,2.4407) -- (7.5254,2.4704);
\draw[gray!65, line width=0.15pt] (7.7288,2.5391) -- (7.6808,2.5763);
\draw[gray!65, line width=0.15pt] (7.6808,2.5763) -- (7.7126,2.7119);
\draw[gray!65, line width=0.15pt] (7.7126,2.7119) -- (7.6687,2.8475);
\draw[gray!65, line width=0.15pt] (7.5254,2.9297) -- (7.6055,2.9831);
\draw[gray!65, line width=0.15pt] (7.6687,2.8475) -- (7.7288,2.8875);
\draw[gray!65, line width=0.15pt] (7.6055,2.9831) -- (7.5409,3.1186);
\draw[gray!65, line width=0.15pt] (7.5409,3.1186) -- (7.6125,3.2542);
\draw[gray!65, line width=0.15pt] (7.6125,3.2542) -- (7.5254,3.3167);
\draw[gray!65, line width=0.15pt] (7.7288,3.3769) -- (7.7069,3.3898);
\draw[gray!65, line width=0.15pt] (7.5254,3.4478) -- (7.6151,3.5254);
\draw[gray!65, line width=0.15pt] (7.7069,3.3898) -- (7.7288,3.3988);
\draw[gray!65, line width=0.15pt] (7.6151,3.5254) -- (7.6100,3.6610);
\draw[gray!65, line width=0.15pt] (7.6100,3.6610) -- (7.5254,3.7058);
\draw[gray!65, line width=0.15pt] (7.7288,3.7357) -- (7.6554,3.7966);
\draw[gray!65, line width=0.15pt] (7.5254,3.9002) -- (7.5970,3.9322);
\draw[gray!65, line width=0.15pt] (7.6554,3.7966) -- (7.7288,3.8642);
\draw[gray!65, line width=0.15pt] (7.5970,3.9322) -- (7.5970,4.0678);
\draw[gray!65, line width=0.15pt] (7.5970,4.0678) -- (7.5254,4.0998);
\draw[gray!65, line width=0.15pt] (7.7288,4.1358) -- (7.6554,4.2034);
\draw[gray!65, line width=0.15pt] (7.5254,4.2942) -- (7.6100,4.3390);
\draw[gray!65, line width=0.15pt] (7.6554,4.2034) -- (7.7288,4.2643);
\draw[gray!65, line width=0.15pt] (7.6100,4.3390) -- (7.6151,4.4746);
\draw[gray!65, line width=0.15pt] (7.6151,4.4746) -- (7.5254,4.5522);
\draw[gray!65, line width=0.15pt] (7.7288,4.6012) -- (7.7069,4.6102);
\draw[gray!65, line width=0.15pt] (7.5254,4.6833) -- (7.6125,4.7458);
\draw[gray!65, line width=0.15pt] (7.7069,4.6102) -- (7.7288,4.6231);
\draw[gray!65, line width=0.15pt] (7.6125,4.7458) -- (7.5409,4.8814);
\draw[gray!65, line width=0.15pt] (7.5409,4.8814) -- (7.6055,5.0169);
\draw[gray!65, line width=0.15pt] (7.6055,5.0169) -- (7.5254,5.0703);
\draw[gray!65, line width=0.15pt] (7.7288,5.1125) -- (7.6687,5.1525);
\draw[gray!65, line width=0.15pt] (7.6687,5.1525) -- (7.7126,5.2881);
\draw[gray!65, line width=0.15pt] (7.7126,5.2881) -- (7.6808,5.4237);
\draw[gray!65, line width=0.15pt] (7.5254,5.5296) -- (7.5774,5.5593);
\draw[gray!65, line width=0.15pt] (7.6808,5.4237) -- (7.7288,5.4609);
\draw[gray!65, line width=0.15pt] (7.5254,5.6312) -- (7.5774,5.5593);
\draw[gray!65, line width=0.15pt] (7.7288,2.5391) -- (7.9322,2.5759);
\draw[gray!65, line width=0.15pt] (7.9322,2.8583) -- (7.7288,2.8875);
\draw[gray!65, line width=0.15pt] (7.7288,3.3769) -- (7.9322,3.3382);
\draw[gray!65, line width=0.15pt] (7.9322,3.4632) -- (7.7288,3.3988);
\draw[gray!65, line width=0.15pt] (7.7288,3.7357) -- (7.9322,3.6992);
\draw[gray!65, line width=0.15pt] (7.9058,3.9322) -- (7.7288,3.8642);
\draw[gray!65, line width=0.15pt] (7.9058,3.9322) -- (7.9058,4.0678);
\draw[gray!65, line width=0.15pt] (7.7288,4.1358) -- (7.9058,4.0678);
\draw[gray!65, line width=0.15pt] (7.9322,4.3008) -- (7.7288,4.2643);
\draw[gray!65, line width=0.15pt] (7.7288,4.6012) -- (7.9322,4.5368);
\draw[gray!65, line width=0.15pt] (7.9322,4.6618) -- (7.7288,4.6231);
\draw[gray!65, line width=0.15pt] (7.7288,5.1125) -- (7.9322,5.1417);
\draw[gray!65, line width=0.15pt] (7.9322,5.4241) -- (7.7288,5.4609);
\draw[gray!65, line width=0.15pt] (7.9322,2.5759) -- (8.1356,2.5505);
\draw[gray!65, line width=0.15pt] (8.1356,2.8795) -- (7.9322,2.8583);
\draw[gray!65, line width=0.15pt] (7.9322,3.3382) -- (8.1356,3.3852);
\draw[gray!65, line width=0.15pt] (8.1356,3.3933) -- (7.9322,3.4632);
\draw[gray!65, line width=0.15pt] (7.9322,3.6992) -- (8.1356,3.7418);
\draw[gray!65, line width=0.15pt] (8.1356,3.8585) -- (7.9681,3.9322);
\draw[gray!65, line width=0.15pt] (7.9681,3.9322) -- (7.9681,4.0678);
\draw[gray!65, line width=0.15pt] (7.9681,4.0678) -- (8.1356,4.1415);
\draw[gray!65, line width=0.15pt] (8.1356,4.2582) -- (7.9322,4.3008);
\draw[gray!65, line width=0.15pt] (7.9322,4.5368) -- (8.1356,4.6067);
\draw[gray!65, line width=0.15pt] (8.1356,4.6148) -- (7.9322,4.6618);
\draw[gray!65, line width=0.15pt] (7.9322,5.1417) -- (8.1356,5.1205);
\draw[gray!65, line width=0.15pt] (8.1356,5.4495) -- (7.9322,5.4241);
\draw[gray!65, line width=0.15pt] (8.3390,2.4149) -- (8.3241,2.4407);
\draw[gray!65, line width=0.15pt] (8.1356,2.5505) -- (8.1743,2.5763);
\draw[gray!65, line width=0.15pt] (8.3241,2.4407) -- (8.3390,2.4506);
\draw[gray!65, line width=0.15pt] (8.1743,2.5763) -- (8.1362,2.7119);
\draw[gray!65, line width=0.15pt] (8.1362,2.7119) -- (8.1914,2.8475);
\draw[gray!65, line width=0.15pt] (8.1914,2.8475) -- (8.1356,2.8795);
\draw[gray!65, line width=0.15pt] (8.3390,2.9484) -- (8.2944,2.9831);
\draw[gray!65, line width=0.15pt] (8.2944,2.9831) -- (8.3390,3.0912);
\draw[gray!65, line width=0.15pt] (8.3390,3.1436) -- (8.2885,3.2542);
\draw[gray!65, line width=0.15pt] (8.1356,3.3852) -- (8.1447,3.3898);
\draw[gray!65, line width=0.15pt] (8.2885,3.2542) -- (8.3390,3.2955);
\draw[gray!65, line width=0.15pt] (8.1447,3.3898) -- (8.1356,3.3933);
\draw[gray!65, line width=0.15pt] (8.3390,3.4748) -- (8.2878,3.5254);
\draw[gray!65, line width=0.15pt] (8.2878,3.5254) -- (8.2989,3.6610);
\draw[gray!65, line width=0.15pt] (8.1356,3.7418) -- (8.2161,3.7966);
\draw[gray!65, line width=0.15pt] (8.2989,3.6610) -- (8.3390,3.6870);
\draw[gray!65, line width=0.15pt] (8.2161,3.7966) -- (8.1356,3.8585);
\draw[gray!65, line width=0.15pt] (8.3390,3.9260) -- (8.3272,3.9322);
\draw[gray!65, line width=0.15pt] (8.3272,3.9322) -- (8.3272,4.0678);
\draw[gray!65, line width=0.15pt] (8.1356,4.1415) -- (8.2161,4.2034);
\draw[gray!65, line width=0.15pt] (8.3272,4.0678) -- (8.3390,4.0740);
\draw[gray!65, line width=0.15pt] (8.2161,4.2034) -- (8.1356,4.2582);
\draw[gray!65, line width=0.15pt] (8.3390,4.3130) -- (8.2989,4.3390);
\draw[gray!65, line width=0.15pt] (8.2989,4.3390) -- (8.2878,4.4746);
\draw[gray!65, line width=0.15pt] (8.1356,4.6067) -- (8.1447,4.6102);
\draw[gray!65, line width=0.15pt] (8.2878,4.4746) -- (8.3390,4.5252);
\draw[gray!65, line width=0.15pt] (8.1447,4.6102) -- (8.1356,4.6148);
\draw[gray!65, line width=0.15pt] (8.3390,4.7045) -- (8.2885,4.7458);
\draw[gray!65, line width=0.15pt] (8.2885,4.7458) -- (8.3390,4.8564);
\draw[gray!65, line width=0.15pt] (8.3390,4.9088) -- (8.2944,5.0169);
\draw[gray!65, line width=0.15pt] (8.1356,5.1205) -- (8.1914,5.1525);
\draw[gray!65, line width=0.15pt] (8.2944,5.0169) -- (8.3390,5.0516);
\draw[gray!65, line width=0.15pt] (8.1914,5.1525) -- (8.1362,5.2881);
\draw[gray!65, line width=0.15pt] (8.1362,5.2881) -- (8.1743,5.4237);
\draw[gray!65, line width=0.15pt] (8.1743,5.4237) -- (8.1356,5.4495);
\draw[gray!65, line width=0.15pt] (8.3390,5.5494) -- (8.3241,5.5593);
\draw[gray!65, line width=0.15pt] (8.3241,5.5593) -- (8.3390,5.5851);
\draw[gray!65, line width=0.15pt] (8.3390,2.4149) -- (8.3777,2.4407);
\draw[gray!65, line width=0.15pt] (8.3390,2.4506) -- (8.3777,2.4407);
\draw[gray!65, line width=0.15pt] (8.3390,2.9484) -- (8.4467,2.9831);
\draw[gray!65, line width=0.15pt] (8.4467,2.9831) -- (8.3390,3.0912);
\draw[gray!65, line width=0.15pt] (8.5424,3.0526) -- (8.3860,3.1186);
\draw[gray!65, line width=0.15pt] (8.3390,3.1436) -- (8.3860,3.1186);
\draw[gray!65, line width=0.15pt] (8.5424,3.2674) -- (8.3390,3.2955);
\draw[gray!65, line width=0.15pt] (8.3390,3.4748) -- (8.4339,3.5254);
\draw[gray!65, line width=0.15pt] (8.4339,3.5254) -- (8.3972,3.6610);
\draw[gray!65, line width=0.15pt] (8.3390,3.6870) -- (8.3972,3.6610);
\draw[gray!65, line width=0.15pt] (8.3390,3.9260) -- (8.3511,3.9322);
\draw[gray!65, line width=0.15pt] (8.3511,3.9322) -- (8.3511,4.0678);
\draw[gray!65, line width=0.15pt] (8.3390,4.0740) -- (8.3511,4.0678);
\draw[gray!65, line width=0.15pt] (8.3390,4.3130) -- (8.3972,4.3390);
\draw[gray!65, line width=0.15pt] (8.3972,4.3390) -- (8.4339,4.4746);
\draw[gray!65, line width=0.15pt] (8.3390,4.5252) -- (8.4339,4.4746);
\draw[gray!65, line width=0.15pt] (8.3390,4.7045) -- (8.5424,4.7326);
\draw[gray!65, line width=0.15pt] (8.3860,4.8814) -- (8.3390,4.8564);
\draw[gray!65, line width=0.15pt] (8.3390,4.9088) -- (8.4467,5.0169);
\draw[gray!65, line width=0.15pt] (8.3860,4.8814) -- (8.5424,4.9474);
\draw[gray!65, line width=0.15pt] (8.3390,5.0516) -- (8.4467,5.0169);
\draw[gray!65, line width=0.15pt] (8.3390,5.5494) -- (8.3777,5.5593);
\draw[gray!65, line width=0.15pt] (8.3390,5.5851) -- (8.3777,5.5593);
\draw[gray!65, line width=0.15pt] (8.5424,3.0526) -- (8.7458,3.0057);
\draw[gray!65, line width=0.15pt] (8.7458,3.2789) -- (8.5424,3.2674);
\draw[gray!65, line width=0.15pt] (8.5424,4.7326) -- (8.7458,4.7211);
\draw[gray!65, line width=0.15pt] (8.7458,4.9943) -- (8.5424,4.9474);
\draw[gray!65, line width=0.15pt] (8.7458,3.0057) -- (8.7722,3.1186);
\draw[gray!65, line width=0.15pt] (8.7722,3.1186) -- (8.7923,3.2542);
\draw[gray!65, line width=0.15pt] (8.7923,3.2542) -- (8.7458,3.2789);
\draw[gray!65, line width=0.15pt] (8.9492,3.3672) -- (8.9225,3.3898);
\draw[gray!65, line width=0.15pt] (8.9225,3.3898) -- (8.9492,3.4383);
\draw[gray!65, line width=0.15pt] (8.9492,3.8880) -- (8.9224,3.9322);
\draw[gray!65, line width=0.15pt] (8.9224,3.9322) -- (8.9224,4.0678);
\draw[gray!65, line width=0.15pt] (8.9224,4.0678) -- (8.9492,4.1120);
\draw[gray!65, line width=0.15pt] (8.9492,4.5617) -- (8.9225,4.6102);
\draw[gray!65, line width=0.15pt] (8.7458,4.7211) -- (8.7923,4.7458);
\draw[gray!65, line width=0.15pt] (8.9225,4.6102) -- (8.9492,4.6328);
\draw[gray!65, line width=0.15pt] (8.7923,4.7458) -- (8.7722,4.8814);
\draw[gray!65, line width=0.15pt] (8.7458,4.9943) -- (8.7722,4.8814);
\draw[gray!65, line width=0.15pt] (8.9492,3.3672) -- (8.9936,3.3898);
\draw[gray!65, line width=0.15pt] (8.9936,3.3898) -- (8.9492,3.4383);
\draw[gray!65, line width=0.15pt] (9.1525,3.5082) -- (9.1069,3.5254);
\draw[gray!65, line width=0.15pt] (9.1069,3.5254) -- (9.1122,3.6610);
\draw[gray!65, line width=0.15pt] (9.1122,3.6610) -- (9.1525,3.6640);
\draw[gray!65, line width=0.15pt] (8.9492,3.8880) -- (9.0524,3.9322);
\draw[gray!65, line width=0.15pt] (9.0524,3.9322) -- (9.0524,4.0678);
\draw[gray!65, line width=0.15pt] (8.9492,4.1120) -- (9.0524,4.0678);
\draw[gray!65, line width=0.15pt] (9.1525,4.3360) -- (9.1122,4.3390);
\draw[gray!65, line width=0.15pt] (9.1122,4.3390) -- (9.1069,4.4746);
\draw[gray!65, line width=0.15pt] (8.9492,4.5617) -- (8.9936,4.6102);
\draw[gray!65, line width=0.15pt] (9.1069,4.4746) -- (9.1525,4.4918);
\draw[gray!65, line width=0.15pt] (8.9492,4.6328) -- (8.9936,4.6102);
\draw[gray!65, line width=0.15pt] (9.1525,3.5082) -- (9.3559,3.5185);
\draw[gray!65, line width=0.15pt] (9.3559,3.6628) -- (9.1525,3.6640);
\draw[gray!65, line width=0.15pt] (9.1525,4.3360) -- (9.3559,4.3372);
\draw[gray!65, line width=0.15pt] (9.3559,4.4815) -- (9.1525,4.4918);
\draw[gray!65, line width=0.15pt] (9.3559,3.5185) -- (9.3633,3.5254);
\draw[gray!65, line width=0.15pt] (9.3633,3.5254) -- (9.3600,3.6610);
\draw[gray!65, line width=0.15pt] (9.3559,3.6628) -- (9.3600,3.6610);
\draw[gray!65, line width=0.15pt] (9.3559,4.3372) -- (9.3600,4.3390);
\draw[gray!65, line width=0.15pt] (9.3600,4.3390) -- (9.3633,4.4746);
\draw[gray!65, line width=0.15pt] (9.3559,4.4815) -- (9.3633,4.4746);
\draw[gray!75, line width=0.15pt] (4.2712,3.4972) -- (4.1973,3.5254);
\draw[gray!75, line width=0.15pt] (4.1973,3.5254) -- (4.2298,3.6610);
\draw[gray!75, line width=0.15pt] (4.2298,3.6610) -- (4.2712,3.6738);
\draw[gray!75, line width=0.15pt] (4.2712,4.3262) -- (4.2298,4.3390);
\draw[gray!75, line width=0.15pt] (4.2298,4.3390) -- (4.1973,4.4746);
\draw[gray!75, line width=0.15pt] (4.1973,4.4746) -- (4.2712,4.5028);
\draw[gray!75, line width=0.15pt] (4.2966,3.5254) -- (4.2712,3.4972);
\draw[gray!75, line width=0.15pt] (4.2966,3.5254) -- (4.2915,3.6610);
\draw[gray!75, line width=0.15pt] (4.2712,3.6738) -- (4.4192,3.7966);
\draw[gray!75, line width=0.15pt] (4.2915,3.6610) -- (4.4746,3.7579);
\draw[gray!75, line width=0.15pt] (4.4192,3.7966) -- (4.4746,3.8407);
\draw[gray!75, line width=0.15pt] (4.4746,4.1593) -- (4.4192,4.2034);
\draw[gray!75, line width=0.15pt] (4.4192,4.2034) -- (4.2712,4.3262);
\draw[gray!75, line width=0.15pt] (4.4746,4.2421) -- (4.2915,4.3390);
\draw[gray!75, line width=0.15pt] (4.2915,4.3390) -- (4.2966,4.4746);
\draw[gray!75, line width=0.15pt] (4.2712,4.5028) -- (4.2966,4.4746);
\draw[gray!75, line width=0.15pt] (4.6780,3.3538) -- (4.5301,3.3898);
\draw[gray!75, line width=0.15pt] (4.5301,3.3898) -- (4.6780,3.4609);
\draw[gray!75, line width=0.15pt] (4.6780,3.6599) -- (4.6774,3.6610);
\draw[gray!75, line width=0.15pt] (4.6774,3.6610) -- (4.4746,3.7579);
\draw[gray!75, line width=0.15pt] (4.6780,3.6662) -- (4.6497,3.7966);
\draw[gray!75, line width=0.15pt] (4.4746,3.8407) -- (4.6206,3.9322);
\draw[gray!75, line width=0.15pt] (4.6497,3.7966) -- (4.6780,3.8191);
\draw[gray!75, line width=0.15pt] (4.6206,3.9322) -- (4.6206,4.0678);
\draw[gray!75, line width=0.15pt] (4.6206,4.0678) -- (4.4746,4.1593);
\draw[gray!75, line width=0.15pt] (4.6780,4.1809) -- (4.6497,4.2034);
\draw[gray!75, line width=0.15pt] (4.4746,4.2421) -- (4.6774,4.3390);
\draw[gray!75, line width=0.15pt] (4.6497,4.2034) -- (4.6780,4.3338);
\draw[gray!75, line width=0.15pt] (4.6774,4.3390) -- (4.6780,4.3401);
\draw[gray!75, line width=0.15pt] (4.6780,4.5391) -- (4.5301,4.6102);
\draw[gray!75, line width=0.15pt] (4.5301,4.6102) -- (4.6780,4.6462);
\draw[gray!75, line width=0.15pt] (4.8814,3.3566) -- (4.6780,3.3538);
\draw[gray!75, line width=0.15pt] (4.6780,3.4609) -- (4.8814,3.4293);
\draw[gray!75, line width=0.15pt] (4.6787,3.6610) -- (4.6780,3.6599);
\draw[gray!75, line width=0.15pt] (4.6780,3.6662) -- (4.6969,3.7966);
\draw[gray!75, line width=0.15pt] (4.6787,3.6610) -- (4.8814,3.7306);
\draw[gray!75, line width=0.15pt] (4.6969,3.7966) -- (4.6780,3.8191);
\draw[gray!75, line width=0.15pt] (4.8814,3.8920) -- (4.7900,3.9322);
\draw[gray!75, line width=0.15pt] (4.7900,3.9322) -- (4.7900,4.0678);
\draw[gray!75, line width=0.15pt] (4.6780,4.1809) -- (4.6969,4.2034);
\draw[gray!75, line width=0.15pt] (4.7900,4.0678) -- (4.8814,4.1080);
\draw[gray!75, line width=0.15pt] (4.6969,4.2034) -- (4.6780,4.3338);
\draw[gray!75, line width=0.15pt] (4.8814,4.2694) -- (4.6787,4.3390);
\draw[gray!75, line width=0.15pt] (4.6780,4.3401) -- (4.6787,4.3390);
\draw[gray!75, line width=0.15pt] (4.8814,4.5707) -- (4.6780,4.5391);
\draw[gray!75, line width=0.15pt] (4.6780,4.6462) -- (4.8814,4.6434);
\draw[gray!75, line width=0.15pt] (5.0847,3.0201) -- (5.0307,3.1186);
\draw[gray!75, line width=0.15pt] (5.0307,3.1186) -- (5.0349,3.2542);
\draw[gray!75, line width=0.15pt] (5.0349,3.2542) -- (4.8814,3.3566);
\draw[gray!75, line width=0.15pt] (5.0847,3.2827) -- (4.9241,3.3898);
\draw[gray!75, line width=0.15pt] (4.8814,3.4293) -- (4.9241,3.3898);
\draw[gray!75, line width=0.15pt] (4.9520,3.7966) -- (4.8814,3.7306);
\draw[gray!75, line width=0.15pt] (4.8814,3.8920) -- (4.9520,3.7966);
\draw[gray!75, line width=0.15pt] (4.9520,4.2034) -- (4.8814,4.1080);
\draw[gray!75, line width=0.15pt] (4.8814,4.2694) -- (4.9520,4.2034);
\draw[gray!75, line width=0.15pt] (4.9241,4.6102) -- (4.8814,4.5707);
\draw[gray!75, line width=0.15pt] (4.8814,4.6434) -- (5.0349,4.7458);
\draw[gray!75, line width=0.15pt] (4.9241,4.6102) -- (5.0847,4.7173);
\draw[gray!75, line width=0.15pt] (5.0349,4.7458) -- (5.0307,4.8814);
\draw[gray!75, line width=0.15pt] (5.0307,4.8814) -- (5.0847,4.9799);
\draw[gray!75, line width=0.15pt] (5.2881,2.7982) -- (5.2457,2.8475);
\draw[gray!75, line width=0.15pt] (5.2457,2.8475) -- (5.2881,2.8644);
\draw[gray!75, line width=0.15pt] (5.1913,3.1186) -- (5.0847,3.0201);
\draw[gray!75, line width=0.15pt] (5.1913,3.1186) -- (5.1439,3.2542);
\draw[gray!75, line width=0.15pt] (5.0847,3.2827) -- (5.1439,3.2542);
\draw[gray!75, line width=0.15pt] (5.2881,3.7475) -- (5.1913,3.7966);
\draw[gray!75, line width=0.15pt] (5.1913,3.7966) -- (5.2881,3.8476);
\draw[gray!75, line width=0.15pt] (5.2881,4.1524) -- (5.1913,4.2034);
\draw[gray!75, line width=0.15pt] (5.1913,4.2034) -- (5.2881,4.2525);
\draw[gray!75, line width=0.15pt] (5.1439,4.7458) -- (5.0847,4.7173);
\draw[gray!75, line width=0.15pt] (5.1439,4.7458) -- (5.1913,4.8814);
\draw[gray!75, line width=0.15pt] (5.0847,4.9799) -- (5.1913,4.8814);
\draw[gray!75, line width=0.15pt] (5.2881,5.1356) -- (5.2457,5.1525);
\draw[gray!75, line width=0.15pt] (5.2457,5.1525) -- (5.2881,5.2018);
\draw[gray!75, line width=0.15pt] (5.4915,2.8199) -- (5.2881,2.7982);
\draw[gray!75, line width=0.15pt] (5.2881,2.8644) -- (5.4915,2.8610);
\draw[gray!75, line width=0.15pt] (5.4915,3.1183) -- (5.4899,3.1186);
\draw[gray!75, line width=0.15pt] (5.4899,3.1186) -- (5.4915,3.1192);
\draw[gray!75, line width=0.15pt] (5.4915,3.7576) -- (5.2881,3.7475);
\draw[gray!75, line width=0.15pt] (5.2881,3.8476) -- (5.4915,3.8342);
\draw[gray!75, line width=0.15pt] (5.4915,4.1658) -- (5.2881,4.1524);
\draw[gray!75, line width=0.15pt] (5.2881,4.2525) -- (5.4915,4.2424);
\draw[gray!75, line width=0.15pt] (5.4915,4.8808) -- (5.4899,4.8814);
\draw[gray!75, line width=0.15pt] (5.4899,4.8814) -- (5.4915,4.8817);
\draw[gray!75, line width=0.15pt] (5.4915,5.1390) -- (5.2881,5.1356);
\draw[gray!75, line width=0.15pt] (5.2881,5.2018) -- (5.4915,5.1801);
\draw[gray!75, line width=0.15pt] (5.5107,2.8475) -- (5.4915,2.8199);
\draw[gray!75, line width=0.15pt] (5.4915,2.8610) -- (5.6369,2.9831);
\draw[gray!75, line width=0.15pt] (5.5107,2.8475) -- (5.6949,2.9462);
\draw[gray!75, line width=0.15pt] (5.6369,2.9831) -- (5.4915,3.1183);
\draw[gray!75, line width=0.15pt] (5.6949,3.0998) -- (5.4993,3.1186);
\draw[gray!75, line width=0.15pt] (5.4915,3.1192) -- (5.5959,3.2542);
\draw[gray!75, line width=0.15pt] (5.4993,3.1186) -- (5.6949,3.1312);
\draw[gray!75, line width=0.15pt] (5.5959,3.2542) -- (5.6949,3.3013);
\draw[gray!75, line width=0.15pt] (5.6949,3.4608) -- (5.6212,3.5254);
\draw[gray!75, line width=0.15pt] (5.6212,3.5254) -- (5.6363,3.6610);
\draw[gray!75, line width=0.15pt] (5.6363,3.6610) -- (5.4915,3.7576);
\draw[gray!75, line width=0.15pt] (5.6949,3.6964) -- (5.5445,3.7966);
\draw[gray!75, line width=0.15pt] (5.4915,3.8342) -- (5.6617,3.9322);
\draw[gray!75, line width=0.15pt] (5.5445,3.7966) -- (5.6949,3.9116);
\draw[gray!75, line width=0.15pt] (5.6617,3.9322) -- (5.6617,4.0678);
\draw[gray!75, line width=0.15pt] (5.6617,4.0678) -- (5.4915,4.1658);
\draw[gray!75, line width=0.15pt] (5.6949,4.0884) -- (5.5445,4.2034);
\draw[gray!75, line width=0.15pt] (5.4915,4.2424) -- (5.6363,4.3390);
\draw[gray!75, line width=0.15pt] (5.5445,4.2034) -- (5.6949,4.3036);
\draw[gray!75, line width=0.15pt] (5.6363,4.3390) -- (5.6212,4.4746);
\draw[gray!75, line width=0.15pt] (5.6212,4.4746) -- (5.6949,4.5392);
\draw[gray!75, line width=0.15pt] (5.6949,4.6987) -- (5.5959,4.7458);
\draw[gray!75, line width=0.15pt] (5.5959,4.7458) -- (5.4915,4.8808);
\draw[gray!75, line width=0.15pt] (5.6949,4.8688) -- (5.4993,4.8814);
\draw[gray!75, line width=0.15pt] (5.4915,4.8817) -- (5.6369,5.0169);
\draw[gray!75, line width=0.15pt] (5.4993,4.8814) -- (5.6949,4.9002);
\draw[gray!75, line width=0.15pt] (5.6369,5.0169) -- (5.4915,5.1390);
\draw[gray!75, line width=0.15pt] (5.6949,5.0538) -- (5.5107,5.1525);
\draw[gray!75, line width=0.15pt] (5.4915,5.1801) -- (5.5107,5.1525);
\draw[gray!75, line width=0.15pt] (5.7611,2.9831) -- (5.6949,2.9462);
\draw[gray!75, line width=0.15pt] (5.6949,3.0998) -- (5.7286,3.1186);
\draw[gray!75, line width=0.15pt] (5.7611,2.9831) -- (5.8983,3.0804);
\draw[gray!75, line width=0.15pt] (5.7286,3.1186) -- (5.6949,3.1312);
\draw[gray!75, line width=0.15pt] (5.8983,3.1933) -- (5.8381,3.2542);
\draw[gray!75, line width=0.15pt] (5.6949,3.3013) -- (5.8381,3.2542);
\draw[gray!75, line width=0.15pt] (5.7713,3.5254) -- (5.6949,3.4608);
\draw[gray!75, line width=0.15pt] (5.7713,3.5254) -- (5.7513,3.6610);
\draw[gray!75, line width=0.15pt] (5.6949,3.6964) -- (5.8674,3.7966);
\draw[gray!75, line width=0.15pt] (5.7513,3.6610) -- (5.8983,3.7745);
\draw[gray!75, line width=0.15pt] (5.8674,3.7966) -- (5.6949,3.9116);
\draw[gray!75, line width=0.15pt] (5.8983,3.8160) -- (5.7239,3.9322);
\draw[gray!75, line width=0.15pt] (5.7239,3.9322) -- (5.7239,4.0678);
\draw[gray!75, line width=0.15pt] (5.6949,4.0884) -- (5.8674,4.2034);
\draw[gray!75, line width=0.15pt] (5.7239,4.0678) -- (5.8983,4.1840);
\draw[gray!75, line width=0.15pt] (5.8674,4.2034) -- (5.6949,4.3036);
\draw[gray!75, line width=0.15pt] (5.8983,4.2255) -- (5.7513,4.3390);
\draw[gray!75, line width=0.15pt] (5.7513,4.3390) -- (5.7713,4.4746);
\draw[gray!75, line width=0.15pt] (5.6949,4.5392) -- (5.7713,4.4746);
\draw[gray!75, line width=0.15pt] (5.8381,4.7458) -- (5.6949,4.6987);
\draw[gray!75, line width=0.15pt] (5.6949,4.8688) -- (5.7286,4.8814);
\draw[gray!75, line width=0.15pt] (5.8381,4.7458) -- (5.8983,4.8067);
\draw[gray!75, line width=0.15pt] (5.7286,4.8814) -- (5.6949,4.9002);
\draw[gray!75, line width=0.15pt] (5.8983,4.9196) -- (5.7611,5.0169);
\draw[gray!75, line width=0.15pt] (5.6949,5.0538) -- (5.7611,5.0169);
\draw[gray!75, line width=0.15pt] (6.1017,3.0804) -- (5.8983,3.0804);
\draw[gray!75, line width=0.15pt] (5.8983,3.1933) -- (6.1017,3.1933);
\draw[gray!75, line width=0.15pt] (6.1017,3.7745) -- (5.8983,3.7745);
\draw[gray!75, line width=0.15pt] (5.8983,3.8160) -- (6.1017,3.8160);
\draw[gray!75, line width=0.15pt] (6.1017,4.1840) -- (5.8983,4.1840);
\draw[gray!75, line width=0.15pt] (5.8983,4.2255) -- (6.1017,4.2255);
\draw[gray!75, line width=0.15pt] (6.1017,4.8067) -- (5.8983,4.8067);
\draw[gray!75, line width=0.15pt] (5.8983,4.9196) -- (6.1017,4.9196);
\draw[gray!75, line width=0.15pt] (6.3051,2.9462) -- (6.2389,2.9831);
\draw[gray!75, line width=0.15pt] (6.2389,2.9831) -- (6.1017,3.0804);
\draw[gray!75, line width=0.15pt] (6.3051,3.0998) -- (6.2714,3.1186);
\draw[gray!75, line width=0.15pt] (6.1017,3.1933) -- (6.1619,3.2542);
\draw[gray!75, line width=0.15pt] (6.2714,3.1186) -- (6.3051,3.1312);
\draw[gray!75, line width=0.15pt] (6.1619,3.2542) -- (6.3051,3.3013);
\draw[gray!75, line width=0.15pt] (6.3051,3.4608) -- (6.2287,3.5254);
\draw[gray!75, line width=0.15pt] (6.2287,3.5254) -- (6.2487,3.6610);
\draw[gray!75, line width=0.15pt] (6.2487,3.6610) -- (6.1017,3.7745);
\draw[gray!75, line width=0.15pt] (6.3051,3.6964) -- (6.1326,3.7966);
\draw[gray!75, line width=0.15pt] (6.1017,3.8160) -- (6.2761,3.9322);
\draw[gray!75, line width=0.15pt] (6.1326,3.7966) -- (6.3051,3.9116);
\draw[gray!75, line width=0.15pt] (6.2761,3.9322) -- (6.2761,4.0678);
\draw[gray!75, line width=0.15pt] (6.2761,4.0678) -- (6.1017,4.1840);
\draw[gray!75, line width=0.15pt] (6.3051,4.0884) -- (6.1326,4.2034);
\draw[gray!75, line width=0.15pt] (6.1017,4.2255) -- (6.2487,4.3390);
\draw[gray!75, line width=0.15pt] (6.1326,4.2034) -- (6.3051,4.3036);
\draw[gray!75, line width=0.15pt] (6.2487,4.3390) -- (6.2287,4.4746);
\draw[gray!75, line width=0.15pt] (6.2287,4.4746) -- (6.3051,4.5392);
\draw[gray!75, line width=0.15pt] (6.3051,4.6987) -- (6.1619,4.7458);
\draw[gray!75, line width=0.15pt] (6.1619,4.7458) -- (6.1017,4.8067);
\draw[gray!75, line width=0.15pt] (6.3051,4.8688) -- (6.2714,4.8814);
\draw[gray!75, line width=0.15pt] (6.1017,4.9196) -- (6.2389,5.0169);
\draw[gray!75, line width=0.15pt] (6.2714,4.8814) -- (6.3051,4.9002);
\draw[gray!75, line width=0.15pt] (6.2389,5.0169) -- (6.3051,5.0538);
\draw[gray!75, line width=0.15pt] (6.5085,2.8199) -- (6.4893,2.8475);
\draw[gray!75, line width=0.15pt] (6.4893,2.8475) -- (6.3051,2.9462);
\draw[gray!75, line width=0.15pt] (6.5085,2.8610) -- (6.3631,2.9831);
\draw[gray!75, line width=0.15pt] (6.3051,3.0998) -- (6.5007,3.1186);
\draw[gray!75, line width=0.15pt] (6.3631,2.9831) -- (6.5085,3.1183);
\draw[gray!75, line width=0.15pt] (6.5007,3.1186) -- (6.3051,3.1312);
\draw[gray!75, line width=0.15pt] (6.5085,3.1192) -- (6.4041,3.2542);
\draw[gray!75, line width=0.15pt] (6.3051,3.3013) -- (6.4041,3.2542);
\draw[gray!75, line width=0.15pt] (6.3788,3.5254) -- (6.3051,3.4608);
\draw[gray!75, line width=0.15pt] (6.3788,3.5254) -- (6.3637,3.6610);
\draw[gray!75, line width=0.15pt] (6.3051,3.6964) -- (6.4555,3.7966);
\draw[gray!75, line width=0.15pt] (6.3637,3.6610) -- (6.5085,3.7576);
\draw[gray!75, line width=0.15pt] (6.4555,3.7966) -- (6.3051,3.9116);
\draw[gray!75, line width=0.15pt] (6.5085,3.8342) -- (6.3383,3.9322);
\draw[gray!75, line width=0.15pt] (6.3383,3.9322) -- (6.3383,4.0678);
\draw[gray!75, line width=0.15pt] (6.3051,4.0884) -- (6.4555,4.2034);
\draw[gray!75, line width=0.15pt] (6.3383,4.0678) -- (6.5085,4.1658);
\draw[gray!75, line width=0.15pt] (6.4555,4.2034) -- (6.3051,4.3036);
\draw[gray!75, line width=0.15pt] (6.5085,4.2424) -- (6.3637,4.3390);
\draw[gray!75, line width=0.15pt] (6.3637,4.3390) -- (6.3788,4.4746);
\draw[gray!75, line width=0.15pt] (6.3051,4.5392) -- (6.3788,4.4746);
\draw[gray!75, line width=0.15pt] (6.4041,4.7458) -- (6.3051,4.6987);
\draw[gray!75, line width=0.15pt] (6.3051,4.8688) -- (6.5007,4.8814);
\draw[gray!75, line width=0.15pt] (6.4041,4.7458) -- (6.5085,4.8808);
\draw[gray!75, line width=0.15pt] (6.5007,4.8814) -- (6.3051,4.9002);
\draw[gray!75, line width=0.15pt] (6.5085,4.8817) -- (6.3631,5.0169);
\draw[gray!75, line width=0.15pt] (6.3051,5.0538) -- (6.4893,5.1525);
\draw[gray!75, line width=0.15pt] (6.3631,5.0169) -- (6.5085,5.1390);
\draw[gray!75, line width=0.15pt] (6.4893,5.1525) -- (6.5085,5.1801);
\draw[gray!75, line width=0.15pt] (6.7119,2.7982) -- (6.5085,2.8199);
\draw[gray!75, line width=0.15pt] (6.5085,2.8610) -- (6.7119,2.8644);
\draw[gray!75, line width=0.15pt] (6.5101,3.1186) -- (6.5085,3.1183);
\draw[gray!75, line width=0.15pt] (6.5085,3.1192) -- (6.5101,3.1186);
\draw[gray!75, line width=0.15pt] (6.7119,3.7475) -- (6.5085,3.7576);
\draw[gray!75, line width=0.15pt] (6.5085,3.8342) -- (6.7119,3.8476);
\draw[gray!75, line width=0.15pt] (6.7119,4.1524) -- (6.5085,4.1658);
\draw[gray!75, line width=0.15pt] (6.5085,4.2424) -- (6.7119,4.2525);
\draw[gray!75, line width=0.15pt] (6.5101,4.8814) -- (6.5085,4.8808);
\draw[gray!75, line width=0.15pt] (6.5085,4.8817) -- (6.5101,4.8814);
\draw[gray!75, line width=0.15pt] (6.7119,5.1356) -- (6.5085,5.1390);
\draw[gray!75, line width=0.15pt] (6.5085,5.1801) -- (6.7119,5.2018);
\draw[gray!75, line width=0.15pt] (6.7543,2.8475) -- (6.7119,2.7982);
\draw[gray!75, line width=0.15pt] (6.7119,2.8644) -- (6.7543,2.8475);
\draw[gray!75, line width=0.15pt] (6.9153,3.0201) -- (6.8087,3.1186);
\draw[gray!75, line width=0.15pt] (6.8087,3.1186) -- (6.8561,3.2542);
\draw[gray!75, line width=0.15pt] (6.8561,3.2542) -- (6.9153,3.2827);
\draw[gray!75, line width=0.15pt] (6.8087,3.7966) -- (6.7119,3.7475);
\draw[gray!75, line width=0.15pt] (6.7119,3.8476) -- (6.8087,3.7966);
\draw[gray!75, line width=0.15pt] (6.8087,4.2034) -- (6.7119,4.1524);
\draw[gray!75, line width=0.15pt] (6.7119,4.2525) -- (6.8087,4.2034);
\draw[gray!75, line width=0.15pt] (6.9153,4.7173) -- (6.8561,4.7458);
\draw[gray!75, line width=0.15pt] (6.8561,4.7458) -- (6.8087,4.8814);
\draw[gray!75, line width=0.15pt] (6.8087,4.8814) -- (6.9153,4.9799);
\draw[gray!75, line width=0.15pt] (6.7543,5.1525) -- (6.7119,5.1356);
\draw[gray!75, line width=0.15pt] (6.7119,5.2018) -- (6.7543,5.1525);
\draw[gray!75, line width=0.15pt] (6.9693,3.1186) -- (6.9153,3.0201);
\draw[gray!75, line width=0.15pt] (6.9693,3.1186) -- (6.9651,3.2542);
\draw[gray!75, line width=0.15pt] (6.9153,3.2827) -- (7.0759,3.3898);
\draw[gray!75, line width=0.15pt] (6.9651,3.2542) -- (7.1186,3.3566);
\draw[gray!75, line width=0.15pt] (7.0759,3.3898) -- (7.1186,3.4293);
\draw[gray!75, line width=0.15pt] (7.1186,3.7306) -- (7.0480,3.7966);
\draw[gray!75, line width=0.15pt] (7.0480,3.7966) -- (7.1186,3.8920);
\draw[gray!75, line width=0.15pt] (7.1186,4.1080) -- (7.0480,4.2034);
\draw[gray!75, line width=0.15pt] (7.0480,4.2034) -- (7.1186,4.2694);
\draw[gray!75, line width=0.15pt] (7.1186,4.5707) -- (7.0759,4.6102);
\draw[gray!75, line width=0.15pt] (7.0759,4.6102) -- (6.9153,4.7173);
\draw[gray!75, line width=0.15pt] (7.1186,4.6434) -- (6.9651,4.7458);
\draw[gray!75, line width=0.15pt] (6.9651,4.7458) -- (6.9693,4.8814);
\draw[gray!75, line width=0.15pt] (6.9153,4.9799) -- (6.9693,4.8814);
\draw[gray!75, line width=0.15pt] (7.3220,3.3538) -- (7.1186,3.3566);
\draw[gray!75, line width=0.15pt] (7.1186,3.4293) -- (7.3220,3.4609);
\draw[gray!75, line width=0.15pt] (7.3220,3.6599) -- (7.3213,3.6610);
\draw[gray!75, line width=0.15pt] (7.3213,3.6610) -- (7.1186,3.7306);
\draw[gray!75, line width=0.15pt] (7.3220,3.6662) -- (7.3031,3.7966);
\draw[gray!75, line width=0.15pt] (7.1186,3.8920) -- (7.2100,3.9322);
\draw[gray!75, line width=0.15pt] (7.3031,3.7966) -- (7.3220,3.8191);
\draw[gray!75, line width=0.15pt] (7.2100,3.9322) -- (7.2100,4.0678);
\draw[gray!75, line width=0.15pt] (7.2100,4.0678) -- (7.1186,4.1080);
\draw[gray!75, line width=0.15pt] (7.3220,4.1809) -- (7.3031,4.2034);
\draw[gray!75, line width=0.15pt] (7.1186,4.2694) -- (7.3213,4.3390);
\draw[gray!75, line width=0.15pt] (7.3031,4.2034) -- (7.3220,4.3338);
\draw[gray!75, line width=0.15pt] (7.3213,4.3390) -- (7.3220,4.3401);
\draw[gray!75, line width=0.15pt] (7.3220,4.5391) -- (7.1186,4.5707);
\draw[gray!75, line width=0.15pt] (7.1186,4.6434) -- (7.3220,4.6462);
\draw[gray!75, line width=0.15pt] (7.4699,3.3898) -- (7.3220,3.3538);
\draw[gray!75, line width=0.15pt] (7.3220,3.4609) -- (7.4699,3.3898);
\draw[gray!75, line width=0.15pt] (7.3226,3.6610) -- (7.3220,3.6599);
\draw[gray!75, line width=0.15pt] (7.3220,3.6662) -- (7.3503,3.7966);
\draw[gray!75, line width=0.15pt] (7.3226,3.6610) -- (7.5254,3.7579);
\draw[gray!75, line width=0.15pt] (7.3503,3.7966) -- (7.3220,3.8191);
\draw[gray!75, line width=0.15pt] (7.5254,3.8407) -- (7.3794,3.9322);
\draw[gray!75, line width=0.15pt] (7.3794,3.9322) -- (7.3794,4.0678);
\draw[gray!75, line width=0.15pt] (7.3220,4.1809) -- (7.3503,4.2034);
\draw[gray!75, line width=0.15pt] (7.3794,4.0678) -- (7.5254,4.1593);
\draw[gray!75, line width=0.15pt] (7.3503,4.2034) -- (7.3220,4.3338);
\draw[gray!75, line width=0.15pt] (7.5254,4.2421) -- (7.3226,4.3390);
\draw[gray!75, line width=0.15pt] (7.3220,4.3401) -- (7.3226,4.3390);
\draw[gray!75, line width=0.15pt] (7.4699,4.6102) -- (7.3220,4.5391);
\draw[gray!75, line width=0.15pt] (7.3220,4.6462) -- (7.4699,4.6102);
\draw[gray!75, line width=0.15pt] (7.7288,3.4972) -- (7.7034,3.5254);
\draw[gray!75, line width=0.15pt] (7.7034,3.5254) -- (7.7085,3.6610);
\draw[gray!75, line width=0.15pt] (7.7085,3.6610) -- (7.5254,3.7579);
\draw[gray!75, line width=0.15pt] (7.7288,3.6738) -- (7.5808,3.7966);
\draw[gray!75, line width=0.15pt] (7.5254,3.8407) -- (7.5808,3.7966);
\draw[gray!75, line width=0.15pt] (7.5808,4.2034) -- (7.5254,4.1593);
\draw[gray!75, line width=0.15pt] (7.5254,4.2421) -- (7.7085,4.3390);
\draw[gray!75, line width=0.15pt] (7.5808,4.2034) -- (7.7288,4.3262);
\draw[gray!75, line width=0.15pt] (7.7085,4.3390) -- (7.7034,4.4746);
\draw[gray!75, line width=0.15pt] (7.7034,4.4746) -- (7.7288,4.5028);
\draw[gray!75, line width=0.15pt] (7.8027,3.5254) -- (7.7288,3.4972);
\draw[gray!75, line width=0.15pt] (7.8027,3.5254) -- (7.7702,3.6610);
\draw[gray!75, line width=0.15pt] (7.7288,3.6738) -- (7.7702,3.6610);
\draw[gray!75, line width=0.15pt] (7.7702,4.3390) -- (7.7288,4.3262);
\draw[gray!75, line width=0.15pt] (7.7702,4.3390) -- (7.8027,4.4746);
\draw[gray!75, line width=0.15pt] (7.7288,4.5028) -- (7.8027,4.4746);
\draw[gray!85, line width=0.15pt] (4.8814,3.7870) -- (4.8545,3.7966);
\draw[gray!85, line width=0.15pt] (4.8545,3.7966) -- (4.8814,3.8105);
\draw[gray!85, line width=0.15pt] (4.8814,4.1895) -- (4.8545,4.2034);
\draw[gray!85, line width=0.15pt] (4.8545,4.2034) -- (4.8814,4.2130);
\draw[gray!85, line width=0.15pt] (4.8916,3.7966) -- (4.8814,3.7870);
\draw[gray!85, line width=0.15pt] (4.8814,3.8105) -- (4.8916,3.7966);
\draw[gray!85, line width=0.15pt] (4.8916,4.2034) -- (4.8814,4.1895);
\draw[gray!85, line width=0.15pt] (4.8814,4.2130) -- (4.8916,4.2034);
\draw[gray!85, line width=0.15pt] (5.2881,3.7720) -- (5.2397,3.7966);
\draw[gray!85, line width=0.15pt] (5.2397,3.7966) -- (5.2881,3.8221);
\draw[gray!85, line width=0.15pt] (5.2881,4.1779) -- (5.2397,4.2034);
\draw[gray!85, line width=0.15pt] (5.2397,4.2034) -- (5.2881,4.2280);
\draw[gray!85, line width=0.15pt] (5.4915,3.7814) -- (5.2881,3.7720);
\draw[gray!85, line width=0.15pt] (5.2881,3.8221) -- (5.4915,3.8112);
\draw[gray!85, line width=0.15pt] (5.4915,4.1888) -- (5.2881,4.1779);
\draw[gray!85, line width=0.15pt] (5.2881,4.2280) -- (5.4915,4.2186);
\draw[gray!85, line width=0.15pt] (5.6949,3.2363) -- (5.6805,3.2542);
\draw[gray!85, line width=0.15pt] (5.6805,3.2542) -- (5.6949,3.2611);
\draw[gray!85, line width=0.15pt] (5.6949,3.4931) -- (5.6581,3.5254);
\draw[gray!85, line width=0.15pt] (5.6581,3.5254) -- (5.6722,3.6610);
\draw[gray!85, line width=0.15pt] (5.6722,3.6610) -- (5.4915,3.7814);
\draw[gray!85, line width=0.15pt] (5.6949,3.6747) -- (5.5121,3.7966);
\draw[gray!85, line width=0.15pt] (5.4915,3.8112) -- (5.5121,3.7966);
\draw[gray!85, line width=0.15pt] (5.5121,4.2034) -- (5.4915,4.1888);
\draw[gray!85, line width=0.15pt] (5.4915,4.2186) -- (5.6722,4.3390);
\draw[gray!85, line width=0.15pt] (5.5121,4.2034) -- (5.6949,4.3253);
\draw[gray!85, line width=0.15pt] (5.6722,4.3390) -- (5.6581,4.4746);
\draw[gray!85, line width=0.15pt] (5.6581,4.4746) -- (5.6949,4.5069);
\draw[gray!85, line width=0.15pt] (5.6949,4.7389) -- (5.6805,4.7458);
\draw[gray!85, line width=0.15pt] (5.6805,4.7458) -- (5.6949,4.7637);
\draw[gray!85, line width=0.15pt] (5.7158,3.2542) -- (5.6949,3.2363);
\draw[gray!85, line width=0.15pt] (5.6949,3.2611) -- (5.7158,3.2542);
\draw[gray!85, line width=0.15pt] (5.7331,3.5254) -- (5.6949,3.4931);
\draw[gray!85, line width=0.15pt] (5.7331,3.5254) -- (5.7168,3.6610);
\draw[gray!85, line width=0.15pt] (5.6949,3.6747) -- (5.7168,3.6610);
\draw[gray!85, line width=0.15pt] (5.7168,4.3390) -- (5.6949,4.3253);
\draw[gray!85, line width=0.15pt] (5.7168,4.3390) -- (5.7331,4.4746);
\draw[gray!85, line width=0.15pt] (5.6949,4.5069) -- (5.7331,4.4746);
\draw[gray!85, line width=0.15pt] (5.7158,4.7458) -- (5.6949,4.7389);
\draw[gray!85, line width=0.15pt] (5.6949,4.7637) -- (5.7158,4.7458);
\draw[gray!85, line width=0.15pt] (6.3051,3.2363) -- (6.2842,3.2542);
\draw[gray!85, line width=0.15pt] (6.2842,3.2542) -- (6.3051,3.2611);
\draw[gray!85, line width=0.15pt] (6.3051,3.4931) -- (6.2669,3.5254);
\draw[gray!85, line width=0.15pt] (6.2669,3.5254) -- (6.2832,3.6610);
\draw[gray!85, line width=0.15pt] (6.2832,3.6610) -- (6.3051,3.6747);
\draw[gray!85, line width=0.15pt] (6.3051,4.3253) -- (6.2832,4.3390);
\draw[gray!85, line width=0.15pt] (6.2832,4.3390) -- (6.2669,4.4746);
\draw[gray!85, line width=0.15pt] (6.2669,4.4746) -- (6.3051,4.5069);
\draw[gray!85, line width=0.15pt] (6.3051,4.7389) -- (6.2842,4.7458);
\draw[gray!85, line width=0.15pt] (6.2842,4.7458) -- (6.3051,4.7637);
\draw[gray!85, line width=0.15pt] (6.3195,3.2542) -- (6.3051,3.2363);
\draw[gray!85, line width=0.15pt] (6.3051,3.2611) -- (6.3195,3.2542);
\draw[gray!85, line width=0.15pt] (6.3419,3.5254) -- (6.3051,3.4931);
\draw[gray!85, line width=0.15pt] (6.3419,3.5254) -- (6.3278,3.6610);
\draw[gray!85, line width=0.15pt] (6.3051,3.6747) -- (6.4879,3.7966);
\draw[gray!85, line width=0.15pt] (6.3278,3.6610) -- (6.5085,3.7814);
\draw[gray!85, line width=0.15pt] (6.4879,3.7966) -- (6.5085,3.8112);
\draw[gray!85, line width=0.15pt] (6.5085,4.1888) -- (6.4879,4.2034);
\draw[gray!85, line width=0.15pt] (6.4879,4.2034) -- (6.3051,4.3253);
\draw[gray!85, line width=0.15pt] (6.5085,4.2186) -- (6.3278,4.3390);
\draw[gray!85, line width=0.15pt] (6.3278,4.3390) -- (6.3419,4.4746);
\draw[gray!85, line width=0.15pt] (6.3051,4.5069) -- (6.3419,4.4746);
\draw[gray!85, line width=0.15pt] (6.3195,4.7458) -- (6.3051,4.7389);
\draw[gray!85, line width=0.15pt] (6.3051,4.7637) -- (6.3195,4.7458);
\draw[gray!85, line width=0.15pt] (6.7119,3.7720) -- (6.5085,3.7814);
\draw[gray!85, line width=0.15pt] (6.5085,3.8112) -- (6.7119,3.8221);
\draw[gray!85, line width=0.15pt] (6.7119,4.1779) -- (6.5085,4.1888);
\draw[gray!85, line width=0.15pt] (6.5085,4.2186) -- (6.7119,4.2280);
\draw[gray!85, line width=0.15pt] (6.7603,3.7966) -- (6.7119,3.7720);
\draw[gray!85, line width=0.15pt] (6.7119,3.8221) -- (6.7603,3.7966);
\draw[gray!85, line width=0.15pt] (6.7603,4.2034) -- (6.7119,4.1779);
\draw[gray!85, line width=0.15pt] (6.7119,4.2280) -- (6.7603,4.2034);
\draw[gray!85, line width=0.15pt] (7.1186,3.7870) -- (7.1084,3.7966);
\draw[gray!85, line width=0.15pt] (7.1084,3.7966) -- (7.1186,3.8105);
\draw[gray!85, line width=0.15pt] (7.1186,4.1895) -- (7.1084,4.2034);
\draw[gray!85, line width=0.15pt] (7.1084,4.2034) -- (7.1186,4.2130);
\draw[gray!85, line width=0.15pt] (7.1455,3.7966) -- (7.1186,3.7870);
\draw[gray!85, line width=0.15pt] (7.1186,3.8105) -- (7.1455,3.7966);
\draw[gray!85, line width=0.15pt] (7.1455,4.2034) -- (7.1186,4.1895);
\draw[gray!85, line width=0.15pt] (7.1186,4.2130) -- (7.1455,4.2034);
\draw[blue!80!black, line width=1.0pt] (5.9400,3.9600) -- (5.9401,3.9601) -- (5.9402,3.9602) -- (5.9404,3.9602) -- (5.9405,3.9603) -- (5.9406,3.9604) -- (5.9407,3.9605) -- (5.9408,3.9606) -- (5.9409,3.9606) -- (5.9411,3.9607) -- (5.9412,3.9608) -- (5.9413,3.9609) -- (5.9414,3.9609) -- (5.9415,3.9610) -- (5.9416,3.9611) -- (5.9418,3.9612) -- (5.9419,3.9613) -- (5.9420,3.9613) -- (5.9421,3.9614) -- (5.9422,3.9615) -- (5.9423,3.9616) -- (5.9425,3.9616) -- (5.9426,3.9617) -- (5.9427,3.9618) -- (5.9428,3.9619) -- (5.9429,3.9619) -- (5.9430,3.9620) -- (5.9431,3.9621) -- (5.9433,3.9622) -- (5.9434,3.9622) -- (5.9435,3.9623) -- (5.9436,3.9624) -- (5.9437,3.9625) -- (5.9438,3.9625) -- (5.9439,3.9626) -- (5.9440,3.9627) -- (5.9441,3.9628) -- (5.9443,3.9628) -- (5.9444,3.9629) -- (5.9445,3.9630) -- (5.9446,3.9631) -- (5.9447,3.9631) -- (5.9448,3.9632) -- (5.9449,3.9633) -- (5.9450,3.9634) -- (5.9451,3.9634) -- (5.9453,3.9635) -- (5.9454,3.9636) -- (5.9455,3.9636) -- (5.9456,3.9637) -- (5.9457,3.9638) -- (5.9458,3.9639) -- (5.9459,3.9639) -- (5.9460,3.9640) -- (5.9461,3.9641) -- (5.9462,3.9641) -- (5.9463,3.9642) -- (5.9464,3.9643) -- (5.9465,3.9644) -- (5.9467,3.9644) -- (5.9468,3.9645) -- (5.9469,3.9646) -- (5.9470,3.9646) -- (5.9471,3.9647) -- (5.9472,3.9648) -- (5.9473,3.9649) -- (5.9474,3.9649) -- (5.9475,3.9650) -- (5.9476,3.9651) -- (5.9477,3.9651) -- (5.9478,3.9652) -- (5.9479,3.9653) -- (5.9480,3.9653) -- (5.9481,3.9654) -- (5.9482,3.9655) -- (5.9483,3.9656) -- (5.9484,3.9656) -- (5.9485,3.9657) -- (5.9486,3.9658) -- (5.9487,3.9658) -- (5.9488,3.9659) -- (5.9489,3.9660) -- (5.9490,3.9660) -- (5.9491,3.9661) -- (5.9492,3.9662) -- (5.9493,3.9662) -- (5.9494,3.9663) -- (5.9495,3.9664) -- (5.9496,3.9664) -- (5.9497,3.9665) -- (5.9498,3.9666) -- (5.9499,3.9666) -- (5.9500,3.9667) -- (5.9501,3.9668) -- (5.9502,3.9668) -- (5.9503,3.9669) -- (5.9504,3.9670) -- (5.9505,3.9670) -- (5.9506,3.9671) -- (5.9507,3.9672) -- (5.9508,3.9672) -- (5.9509,3.9673) -- (5.9510,3.9674) -- (5.9511,3.9674) -- (5.9512,3.9675) -- (5.9513,3.9676) -- (5.9514,3.9676) -- (5.9515,3.9677) -- (5.9516,3.9677) -- (5.9517,3.9678) -- (5.9518,3.9679) -- (5.9519,3.9679) -- (5.9520,3.9680) -- (5.9521,3.9681) -- (5.9522,3.9681) -- (5.9523,3.9682) -- (5.9524,3.9683) -- (5.9525,3.9683) -- (5.9526,3.9684) -- (5.9527,3.9684) -- (5.9528,3.9685) -- (5.9529,3.9686) -- (5.9529,3.9686) -- (5.9530,3.9687) -- (5.9531,3.9688) -- (5.9532,3.9688) -- (5.9533,3.9689) -- (5.9534,3.9689) -- (5.9535,3.9690) -- (5.9536,3.9691) -- (5.9537,3.9691) -- (5.9538,3.9692) -- (5.9539,3.9693) -- (5.9540,3.9693) -- (5.9541,3.9694) -- (5.9542,3.9694) -- (5.9542,3.9695) -- (5.9543,3.9696) -- (5.9544,3.9696) -- (5.9545,3.9697) -- (5.9546,3.9697) -- (5.9547,3.9698) -- (5.9548,3.9699) -- (5.9549,3.9699) -- (5.9550,3.9700) -- (5.9551,3.9700) -- (5.9551,3.9701) -- (5.9552,3.9702) -- (5.9553,3.9702) -- (5.9554,3.9703) -- (5.9555,3.9703) -- (5.9556,3.9704) -- (5.9557,3.9705) -- (5.9558,3.9705) -- (5.9559,3.9706) -- (5.9559,3.9706) -- (5.9560,3.9707) -- (5.9561,3.9707) -- (5.9562,3.9708) -- (5.9563,3.9709) -- (5.9564,3.9709) -- (5.9565,3.9710) -- (5.9566,3.9710) -- (5.9566,3.9711) -- (5.9567,3.9712) -- (5.9568,3.9712) -- (5.9569,3.9713) -- (5.9570,3.9713) -- (5.9571,3.9714) -- (5.9572,3.9714) -- (5.9572,3.9715) -- (5.9573,3.9716) -- (5.9574,3.9716) -- (5.9575,3.9717) -- (5.9576,3.9717) -- (5.9577,3.9718) -- (5.9578,3.9718) -- (5.9578,3.9719) -- (5.9579,3.9719) -- (5.9580,3.9720) -- (5.9581,3.9721) -- (5.9582,3.9721) -- (5.9583,3.9722) -- (5.9583,3.9722) -- (5.9584,3.9723) -- (5.9585,3.9723) -- (5.9586,3.9724) -- (5.9587,3.9724) -- (5.9588,3.9725) -- (5.9588,3.9726) -- (5.9589,3.9726) -- (5.9590,3.9727) -- (5.9591,3.9727) -- (5.9592,3.9728) -- (5.9592,3.9728) -- (5.9593,3.9729) -- (5.9594,3.9729) -- (5.9595,3.9730) -- (5.9596,3.9730) -- (5.9596,3.9731) -- (5.9597,3.9732) -- (5.9598,3.9732) -- (5.9599,3.9733) -- (5.9600,3.9733) -- (5.9601,3.9734) -- (5.9601,3.9734) -- (5.9602,3.9735) -- (5.9603,3.9735) -- (5.9604,3.9736) -- (5.9604,3.9736) -- (5.9605,3.9737) -- (5.9606,3.9737) -- (5.9607,3.9738) -- (5.9608,3.9738) -- (5.9608,3.9739) -- (5.9609,3.9739) -- (5.9610,3.9740) -- (5.9611,3.9740) -- (5.9612,3.9741) -- (5.9612,3.9742) -- (5.9613,3.9742) -- (5.9614,3.9743) -- (5.9615,3.9743) -- (5.9615,3.9744) -- (5.9616,3.9744) -- (5.9617,3.9745) -- (5.9618,3.9745) -- (5.9618,3.9746) -- (5.9619,3.9746) -- (5.9620,3.9747) -- (5.9621,3.9747) -- (5.9621,3.9748) -- (5.9622,3.9748) -- (5.9623,3.9749) -- (5.9624,3.9749) -- (5.9624,3.9750) -- (5.9625,3.9750) -- (5.9626,3.9751) -- (5.9627,3.9751) -- (5.9627,3.9752) -- (5.9628,3.9752) -- (5.9629,3.9753) -- (5.9630,3.9753) -- (5.9630,3.9754) -- (5.9631,3.9754) -- (5.9632,3.9755) -- (5.9633,3.9755) -- (5.9633,3.9756) -- (5.9634,3.9756) -- (5.9635,3.9757) -- (5.9636,3.9757) -- (5.9636,3.9758) -- (5.9637,3.9758) -- (5.9638,3.9758) -- (5.9638,3.9759) -- (5.9639,3.9759) -- (5.9640,3.9760) -- (5.9641,3.9760) -- (5.9641,3.9761) -- (5.9642,3.9761) -- (5.9643,3.9762) -- (5.9643,3.9762) -- (5.9644,3.9763) -- (5.9645,3.9763) -- (5.9646,3.9764) -- (5.9646,3.9764) -- (5.9647,3.9765) -- (5.9648,3.9765) -- (5.9648,3.9766) -- (5.9649,3.9766) -- (5.9650,3.9767) -- (5.9651,3.9767) -- (5.9651,3.9767) -- (5.9652,3.9768) -- (5.9653,3.9768) -- (5.9653,3.9769) -- (5.9654,3.9769) -- (5.9655,3.9770) -- (5.9655,3.9770) -- (5.9656,3.9771) -- (5.9657,3.9771) -- (5.9657,3.9772) -- (5.9658,3.9772) -- (5.9659,3.9773) -- (5.9660,3.9773) -- (5.9660,3.9773) -- (5.9661,3.9774) -- (5.9662,3.9774) -- (5.9662,3.9775) -- (5.9663,3.9775) -- (5.9664,3.9776) -- (5.9664,3.9776) -- (5.9665,3.9777) -- (5.9666,3.9777) -- (5.9666,3.9777) -- (5.9667,3.9778) -- (5.9668,3.9778) -- (5.9668,3.9779) -- (5.9669,3.9779) -- (5.9670,3.9780) -- (5.9670,3.9780) -- (5.9671,3.9781) -- (5.9672,3.9781) -- (5.9672,3.9781) -- (5.9673,3.9782) -- (5.9673,3.9782) -- (5.9674,3.9783) -- (5.9675,3.9783) -- (5.9675,3.9784) -- (5.9676,3.9784) -- (5.9677,3.9784) -- (5.9677,3.9785) -- (5.9678,3.9785) -- (5.9679,3.9786) -- (5.9679,3.9786) -- (5.9680,3.9787) -- (5.9681,3.9787) -- (5.9681,3.9787) -- (5.9682,3.9788) -- (5.9683,3.9788) -- (5.9683,3.9789) -- (5.9684,3.9789) -- (5.9684,3.9790) -- (5.9685,3.9790) -- (5.9686,3.9790) -- (5.9686,3.9791) -- (5.9687,3.9791) -- (5.9688,3.9792) -- (5.9688,3.9792) -- (5.9689,3.9793) -- (5.9689,3.9793) -- (5.9690,3.9793) -- (5.9691,3.9794) -- (5.9691,3.9794) -- (5.9692,3.9795) -- (5.9692,3.9795) -- (5.9693,3.9795) -- (5.9694,3.9796) -- (5.9694,3.9796) -- (5.9695,3.9797) -- (5.9696,3.9797) -- (5.9696,3.9797) -- (5.9697,3.9798) -- (5.9697,3.9798) -- (5.9698,3.9799) -- (5.9699,3.9799) -- (5.9699,3.9799) -- (5.9700,3.9800) -- (5.9700,3.9800) -- (5.9701,3.9801) -- (5.9702,3.9801) -- (5.9702,3.9801) -- (5.9703,3.9802) -- (5.9703,3.9802) -- (5.9704,3.9803) -- (5.9705,3.9803) -- (5.9705,3.9803) -- (5.9706,3.9804) -- (5.9706,3.9804) -- (5.9707,3.9805) -- (5.9707,3.9805) -- (5.9708,3.9805) -- (5.9709,3.9806) -- (5.9709,3.9806) -- (5.9710,3.9807) -- (5.9710,3.9807) -- (5.9711,3.9807) -- (5.9712,3.9808) -- (5.9712,3.9808) -- (5.9713,3.9808) -- (5.9713,3.9809) -- (5.9714,3.9809) -- (5.9714,3.9810) -- (5.9715,3.9810) -- (5.9716,3.9810) -- (5.9716,3.9811) -- (5.9717,3.9811) -- (5.9717,3.9812) -- (5.9718,3.9812) -- (5.9718,3.9812) -- (5.9719,3.9813) -- (5.9720,3.9813) -- (5.9720,3.9813) -- (5.9721,3.9814) -- (5.9721,3.9814) -- (5.9722,3.9814) -- (5.9722,3.9815) -- (5.9723,3.9815) -- (5.9723,3.9816) -- (5.9724,3.9816) -- (5.9725,3.9816) -- (5.9725,3.9817) -- (5.9726,3.9817) -- (5.9726,3.9817) -- (5.9727,3.9818) -- (5.9727,3.9818) -- (5.9728,3.9819) -- (5.9728,3.9819) -- (5.9729,3.9819) -- (5.9729,3.9820) -- (5.9730,3.9820) -- (5.9731,3.9820) -- (5.9731,3.9821) -- (5.9732,3.9821) -- (5.9732,3.9821) -- (5.9733,3.9822) -- (5.9733,3.9822) -- (5.9734,3.9822) -- (5.9734,3.9823) -- (5.9735,3.9823) -- (5.9735,3.9824) -- (5.9736,3.9824) -- (5.9736,3.9824) -- (5.9737,3.9825) -- (5.9737,3.9825) -- (5.9738,3.9825) -- (5.9738,3.9826) -- (5.9739,3.9826) -- (5.9740,3.9826) -- (5.9740,3.9827) -- (5.9741,3.9827) -- (5.9741,3.9827) -- (5.9742,3.9828) -- (5.9742,3.9828) -- (5.9743,3.9828) -- (5.9743,3.9829) -- (5.9744,3.9829) -- (5.9744,3.9829) -- (5.9745,3.9830) -- (5.9745,3.9830) -- (5.9746,3.9830) -- (5.9746,3.9831) -- (5.9747,3.9831) -- (5.9747,3.9831) -- (5.9748,3.9832) -- (5.9748,3.9832) -- (5.9749,3.9832) -- (5.9749,3.9833) -- (5.9750,3.9833) -- (5.9750,3.9833) -- (5.9751,3.9834) -- (5.9751,3.9834) -- (5.9752,3.9834) -- (5.9752,3.9835) -- (5.9753,3.9835) -- (5.9753,3.9835) -- (5.9754,3.9836) -- (5.9754,3.9836) -- (5.9755,3.9836) -- (5.9755,3.9837) -- (5.9756,3.9837) -- (5.9756,3.9837) -- (5.9757,3.9838) -- (5.9757,3.9838) -- (5.9758,3.9838) -- (5.9758,3.9839) -- (5.9759,3.9839) -- (5.9759,3.9839) -- (5.9760,3.9840) -- (5.9760,3.9840) -- (5.9761,3.9840) -- (5.9761,3.9841) -- (5.9761,3.9841) -- (5.9762,3.9841) -- (5.9762,3.9842) -- (5.9763,3.9842) -- (5.9763,3.9842) -- (5.9764,3.9843) -- (5.9764,3.9843) -- (5.9765,3.9843) -- (5.9765,3.9843) -- (5.9766,3.9844) -- (5.9766,3.9844) -- (5.9767,3.9844) -- (5.9767,3.9845) -- (5.9768,3.9845) -- (5.9768,3.9845) -- (5.9769,3.9846) -- (5.9769,3.9846) -- (5.9769,3.9846) -- (5.9770,3.9847) -- (5.9770,3.9847) -- (5.9771,3.9847) -- (5.9771,3.9848) -- (5.9772,3.9848) -- (5.9772,3.9848) -- (5.9773,3.9848) -- (5.9773,3.9849) -- (5.9774,3.9849) -- (5.9774,3.9849) -- (5.9774,3.9850) -- (5.9775,3.9850) -- (5.9775,3.9850) -- (5.9776,3.9851) -- (5.9776,3.9851) -- (5.9777,3.9851) -- (5.9777,3.9851) -- (5.9778,3.9852) -- (5.9778,3.9852) -- (5.9778,3.9852) -- (5.9779,3.9853) -- (5.9779,3.9853) -- (5.9780,3.9853) -- (5.9780,3.9853) -- (5.9781,3.9854) -- (5.9781,3.9854) -- (5.9782,3.9854) -- (5.9782,3.9855) -- (5.9782,3.9855) -- (5.9783,3.9855) -- (5.9783,3.9856) -- (5.9784,3.9856) -- (5.9784,3.9856) -- (5.9785,3.9856) -- (5.9785,3.9857) -- (5.9785,3.9857) -- (5.9786,3.9857) -- (5.9786,3.9858) -- (5.9787,3.9858) -- (5.9787,3.9858) -- (5.9788,3.9858) -- (5.9788,3.9859) -- (5.9788,3.9859) -- (5.9789,3.9859) -- (5.9789,3.9860) -- (5.9790,3.9860) -- (5.9790,3.9860) -- (5.9791,3.9860) -- (5.9791,3.9861) -- (5.9791,3.9861) -- (5.9792,3.9861) -- (5.9792,3.9861) -- (5.9793,3.9862) -- (5.9793,3.9862) -- (5.9793,3.9862) -- (5.9794,3.9863) -- (5.9794,3.9863) -- (5.9795,3.9863) -- (5.9795,3.9863) -- (5.9796,3.9864) -- (5.9796,3.9864) -- (5.9796,3.9864) -- (5.9797,3.9864) -- (5.9797,3.9865) -- (5.9798,3.9865) -- (5.9798,3.9865) -- (5.9798,3.9866) -- (5.9799,3.9866) -- (5.9799,3.9866) -- (5.9800,3.9866) -- (5.9800,3.9867) -- (5.9800,3.9867) -- (5.9801,3.9867) -- (5.9801,3.9867) -- (5.9802,3.9868) -- (5.9802,3.9868) -- (5.9802,3.9868) -- (5.9803,3.9868) -- (5.9803,3.9869) -- (5.9804,3.9869) -- (5.9804,3.9869) -- (5.9804,3.9870) -- (5.9805,3.9870) -- (5.9805,3.9870) -- (5.9805,3.9870) -- (5.9806,3.9871) -- (5.9806,3.9871) -- (5.9807,3.9871) -- (5.9807,3.9871) -- (5.9807,3.9872) -- (5.9808,3.9872) -- (5.9808,3.9872) -- (5.9809,3.9872) -- (5.9809,3.9873) -- (5.9809,3.9873) -- (5.9810,3.9873) -- (5.9810,3.9873) -- (5.9810,3.9874) -- (5.9811,3.9874) -- (5.9811,3.9874) -- (5.9812,3.9874) -- (5.9812,3.9875) -- (5.9812,3.9875) -- (5.9813,3.9875) -- (5.9813,3.9875) -- (5.9813,3.9876) -- (5.9814,3.9876) -- (5.9814,3.9876) -- (5.9815,3.9876) -- (5.9815,3.9877) -- (5.9815,3.9877) -- (5.9816,3.9877) -- (5.9816,3.9877) -- (5.9816,3.9878) -- (5.9817,3.9878) -- (5.9817,3.9878) -- (5.9818,3.9878) -- (5.9818,3.9879) -- (5.9818,3.9879) -- (5.9819,3.9879) -- (5.9819,3.9879) -- (5.9819,3.9880) -- (5.9820,3.9880) -- (5.9820,3.9880) -- (5.9820,3.9880) -- (5.9821,3.9881) -- (5.9821,3.9881) -- (5.9822,3.9881) -- (5.9822,3.9881) -- (5.9822,3.9881) -- (5.9823,3.9882) -- (5.9823,3.9882) -- (5.9823,3.9882) -- (5.9824,3.9882) -- (5.9824,3.9883) -- (5.9824,3.9883) -- (5.9825,3.9883) -- (5.9825,3.9883) -- (5.9825,3.9884) -- (5.9826,3.9884) -- (5.9826,3.9884) -- (5.9826,3.9884) -- (5.9827,3.9885) -- (5.9827,3.9885) -- (5.9827,3.9885) -- (5.9828,3.9885) -- (5.9828,3.9885) -- (5.9829,3.9886) -- (5.9829,3.9886) -- (5.9829,3.9886) -- (5.9830,3.9886) -- (5.9830,3.9887) -- (5.9830,3.9887) -- (5.9831,3.9887) -- (5.9831,3.9887) -- (5.9831,3.9887) -- (5.9832,3.9888) -- (5.9832,3.9888) -- (5.9832,3.9888) -- (5.9833,3.9888) -- (5.9833,3.9889) -- (5.9833,3.9889) -- (5.9834,3.9889) -- (5.9834,3.9889) -- (5.9834,3.9890) -- (5.9835,3.9890) -- (5.9835,3.9890) -- (5.9835,3.9890) -- (5.9836,3.9890) -- (5.9836,3.9891) -- (5.9836,3.9891) -- (5.9837,3.9891) -- (5.9837,3.9891) -- (5.9837,3.9891) -- (5.9838,3.9892) -- (5.9838,3.9892) -- (5.9838,3.9892) -- (5.9839,3.9892) -- (5.9839,3.9893) -- (5.9839,3.9893) -- (5.9839,3.9893) -- (5.9840,3.9893) -- (5.9840,3.9893) -- (5.9840,3.9894) -- (5.9841,3.9894) -- (5.9841,3.9894) -- (5.9841,3.9894) -- (5.9842,3.9894) -- (5.9842,3.9895) -- (5.9842,3.9895) -- (5.9843,3.9895) -- (5.9843,3.9895) -- (5.9843,3.9896) -- (5.9844,3.9896) -- (5.9844,3.9896) -- (5.9844,3.9896) -- (5.9845,3.9896) -- (5.9845,3.9897) -- (5.9845,3.9897) -- (5.9845,3.9897) -- (5.9846,3.9897) -- (5.9846,3.9897) -- (5.9846,3.9898) -- (5.9847,3.9898) -- (5.9847,3.9898) -- (5.9847,3.9898) -- (5.9848,3.9898) -- (5.9848,3.9899) -- (5.9848,3.9899) -- (5.9849,3.9899) -- (5.9849,3.9899) -- (5.9849,3.9899) -- (5.9849,3.9900) -- (5.9850,3.9900) -- (5.9850,3.9900) -- (5.9850,3.9900) -- (5.9851,3.9900) -- (5.9851,3.9901) -- (5.9851,3.9901) -- (5.9852,3.9901) -- (5.9852,3.9901) -- (5.9852,3.9901) -- (5.9852,3.9902) -- (5.9853,3.9902) -- (5.9853,3.9902) -- (5.9853,3.9902) -- (5.9854,3.9902) -- (5.9854,3.9903) -- (5.9854,3.9903) -- (5.9854,3.9903) -- (5.9855,3.9903) -- (5.9855,3.9903) -- (5.9855,3.9904) -- (5.9856,3.9904) -- (5.9856,3.9904) -- (5.9856,3.9904) -- (5.9856,3.9904) -- (5.9857,3.9905) -- (5.9857,3.9905) -- (5.9857,3.9905) -- (5.9858,3.9905) -- (5.9858,3.9905) -- (5.9858,3.9905) -- (5.9858,3.9906) -- (5.9859,3.9906) -- (5.9859,3.9906) -- (5.9859,3.9906) -- (5.9860,3.9906) -- (5.9860,3.9907) -- (5.9860,3.9907) -- (5.9860,3.9907) -- (5.9861,3.9907) -- (5.9861,3.9907) -- (5.9861,3.9908) -- (5.9862,3.9908) -- (5.9862,3.9908) -- (5.9862,3.9908) -- (5.9862,3.9908) -- (5.9863,3.9908) -- (5.9863,3.9909) -- (5.9863,3.9909) -- (5.9864,3.9909) -- (5.9864,3.9909) -- (5.9864,3.9909) -- (5.9864,3.9910) -- (5.9865,3.9910) -- (5.9865,3.9910) -- (5.9865,3.9910) -- (5.9865,3.9910) -- (5.9866,3.9910) -- (5.9866,3.9911) -- (5.9866,3.9911) -- (5.9866,3.9911) -- (5.9867,3.9911) -- (5.9867,3.9911) -- (5.9867,3.9912) -- (5.9868,3.9912) -- (5.9868,3.9912) -- (5.9868,3.9912) -- (5.9868,3.9912) -- (5.9869,3.9912) -- (5.9869,3.9913) -- (5.9869,3.9913) -- (5.9869,3.9913) -- (5.9870,3.9913) -- (5.9870,3.9913) -- (5.9870,3.9913) -- (5.9870,3.9914) -- (5.9871,3.9914) -- (5.9871,3.9914) -- (5.9871,3.9914) -- (5.9871,3.9914) -- (5.9872,3.9914) -- (5.9872,3.9915) -- (5.9872,3.9915) -- (5.9872,3.9915) -- (5.9873,3.9915) -- (5.9873,3.9915) -- (5.9873,3.9915) -- (5.9873,3.9916) -- (5.9874,3.9916) -- (5.9874,3.9916) -- (5.9874,3.9916) -- (5.9875,3.9916) -- (5.9875,3.9917) -- (5.9875,3.9917) -- (5.9875,3.9917) -- (5.9876,3.9917) -- (5.9876,3.9917) -- (5.9876,3.9917) -- (5.9876,3.9918) -- (5.9877,3.9918) -- (5.9877,3.9918) -- (5.9877,3.9918) -- (5.9877,3.9918) -- (5.9877,3.9918) -- (5.9878,3.9918) -- (5.9878,3.9919) -- (5.9878,3.9919) -- (5.9878,3.9919) -- (5.9879,3.9919) -- (5.9879,3.9919) -- (5.9879,3.9919) -- (5.9879,3.9920) -- (5.9880,3.9920) -- (5.9880,3.9920) -- (5.9880,3.9920) -- (5.9880,3.9920) -- (5.9881,3.9920) -- (5.9881,3.9921) -- (5.9881,3.9921) -- (5.9881,3.9921) -- (5.9882,3.9921) -- (5.9882,3.9921) -- (5.9882,3.9921) -- (5.9882,3.9922) -- (5.9883,3.9922) -- (5.9883,3.9922) -- (5.9883,3.9922) -- (5.9883,3.9922) -- (5.9883,3.9922) -- (5.9884,3.9922) -- (5.9884,3.9923) -- (5.9884,3.9923) -- (5.9884,3.9923) -- (5.9885,3.9923) -- (5.9885,3.9923) -- (5.9885,3.9923) -- (5.9885,3.9924) -- (5.9886,3.9924) -- (5.9886,3.9924) -- (5.9886,3.9924) -- (5.9886,3.9924) -- (5.9886,3.9924) -- (5.9887,3.9924) -- (5.9887,3.9925) -- (5.9887,3.9925) -- (5.9887,3.9925) -- (5.9888,3.9925) -- (5.9888,3.9925) -- (5.9888,3.9925) -- (5.9888,3.9926) -- (5.9888,3.9926) -- (5.9889,3.9926) -- (5.9889,3.9926) -- (5.9889,3.9926) -- (5.9889,3.9926) -- (5.9890,3.9926) -- (5.9890,3.9927) -- (5.9890,3.9927) -- (5.9890,3.9927) -- (5.9890,3.9927) -- (5.9891,3.9927) -- (5.9891,3.9927) -- (5.9891,3.9927) -- (5.9891,3.9928) -- (5.9892,3.9928) -- (5.9892,3.9928) -- (5.9892,3.9928) -- (5.9892,3.9928) -- (5.9892,3.9928) -- (5.9893,3.9928) -- (5.9893,3.9929) -- (5.9893,3.9929) -- (5.9893,3.9929) -- (5.9894,3.9929) -- (5.9894,3.9929) -- (5.9894,3.9929) -- (5.9894,3.9929) -- (5.9894,3.9930) -- (5.9895,3.9930) -- (5.9895,3.9930) -- (5.9895,3.9930) -- (5.9895,3.9930) -- (5.9895,3.9930) -- (5.9896,3.9930) -- (5.9896,3.9931) -- (5.9896,3.9931) -- (5.9896,3.9931) -- (5.9896,3.9931) -- (5.9897,3.9931) -- (5.9897,3.9931) -- (5.9897,3.9931) -- (5.9897,3.9932) -- (5.9897,3.9932) -- (5.9898,3.9932) -- (5.9898,3.9932) -- (5.9898,3.9932) -- (5.9898,3.9932) -- (5.9898,3.9932) -- (5.9899,3.9932) -- (5.9899,3.9933) -- (5.9899,3.9933) -- (5.9899,3.9933) -- (5.9900,3.9933) -- (5.9900,3.9933) -- (5.9900,3.9933) -- (5.9900,3.9933) -- (5.9900,3.9934) -- (5.9901,3.9934) -- (5.9901,3.9934) -- (5.9901,3.9934) -- (5.9901,3.9934) -- (5.9901,3.9934) -- (5.9902,3.9934) -- (5.9902,3.9934) -- (5.9902,3.9935) -- (5.9902,3.9935) -- (5.9902,3.9935) -- (5.9902,3.9935) -- (5.9903,3.9935) -- (5.9903,3.9935) -- (5.9903,3.9935) -- (5.9903,3.9936) -- (5.9903,3.9936) -- (5.9904,3.9936) -- (5.9904,3.9936) -- (5.9904,3.9936) -- (5.9904,3.9936) -- (5.9904,3.9936) -- (5.9905,3.9936) -- (5.9905,3.9937) -- (5.9905,3.9937) -- (5.9905,3.9937) -- (5.9905,3.9937) -- (5.9906,3.9937) -- (5.9906,3.9937) -- (5.9906,3.9937) -- (5.9906,3.9937) -- (5.9906,3.9938) -- (5.9906,3.9938) -- (5.9907,3.9938) -- (5.9907,3.9938) -- (5.9907,3.9938) -- (5.9907,3.9938) -- (5.9907,3.9938) -- (5.9908,3.9938) -- (5.9908,3.9939) -- (5.9908,3.9939) -- (5.9908,3.9939) -- (5.9908,3.9939) -- (5.9909,3.9939) -- (5.9909,3.9939) -- (5.9909,3.9939) -- (5.9909,3.9939) -- (5.9909,3.9940) -- (5.9909,3.9940) -- (5.9910,3.9940) -- (5.9910,3.9940) -- (5.9910,3.9940) -- (5.9910,3.9940) -- (5.9910,3.9940) -- (5.9911,3.9940) -- (5.9911,3.9940) -- (5.9911,3.9941) -- (5.9911,3.9941) -- (5.9911,3.9941) -- (5.9911,3.9941) -- (5.9912,3.9941) -- (5.9912,3.9941) -- (5.9912,3.9941) -- (5.9912,3.9941) -- (5.9912,3.9942) -- (5.9912,3.9942) -- (5.9913,3.9942) -- (5.9913,3.9942) -- (5.9913,3.9942) -- (5.9913,3.9942) -- (5.9913,3.9942) -- (5.9914,3.9942) -- (5.9914,3.9942) -- (5.9914,3.9943) -- (5.9914,3.9943) -- (5.9914,3.9943) -- (5.9914,3.9943) -- (5.9915,3.9943) -- (5.9915,3.9943) -- (5.9915,3.9943) -- (5.9915,3.9943) -- (5.9915,3.9943) -- (5.9915,3.9944) -- (5.9916,3.9944) -- (5.9916,3.9944) -- (5.9916,3.9944) -- (5.9916,3.9944) -- (5.9916,3.9944) -- (5.9916,3.9944) -- (5.9917,3.9944) -- (5.9917,3.9944) -- (5.9917,3.9945) -- (5.9917,3.9945) -- (5.9917,3.9945) -- (5.9917,3.9945) -- (5.9918,3.9945) -- (5.9918,3.9945) -- (5.9918,3.9945) -- (5.9918,3.9945) -- (5.9918,3.9945) -- (5.9918,3.9946) -- (5.9919,3.9946) -- (5.9919,3.9946) -- (5.9919,3.9946) -- (5.9919,3.9946) -- (5.9919,3.9946) -- (5.9919,3.9946) -- (5.9920,3.9946) -- (5.9920,3.9946) -- (5.9920,3.9947) -- (5.9920,3.9947) -- (5.9920,3.9947) -- (5.9920,3.9947) -- (5.9920,3.9947) -- (5.9921,3.9947) -- (5.9921,3.9947) -- (5.9921,3.9947) -- (5.9921,3.9947) -- (5.9921,3.9948) -- (5.9921,3.9948) -- (5.9922,3.9948) -- (5.9922,3.9948) -- (5.9922,3.9948) -- (5.9922,3.9948) -- (5.9922,3.9948) -- (5.9922,3.9948) -- (5.9923,3.9948) -- (5.9923,3.9948) -- (5.9923,3.9949) -- (5.9923,3.9949) -- (5.9923,3.9949) -- (5.9923,3.9949) -- (5.9923,3.9949) -- (5.9924,3.9949) -- (5.9924,3.9949) -- (5.9924,3.9949) -- (5.9924,3.9949) -- (5.9924,3.9949) -- (5.9924,3.9950) -- (5.9925,3.9950) -- (5.9925,3.9950) -- (5.9925,3.9950) -- (5.9925,3.9950) -- (5.9925,3.9950) -- (5.9925,3.9950) -- (5.9925,3.9950) -- (5.9926,3.9950) -- (5.9926,3.9950) -- (5.9926,3.9951) -- (5.9926,3.9951) -- (5.9926,3.9951) -- (5.9926,3.9951) -- (5.9926,3.9951) -- (5.9927,3.9951) -- (5.9927,3.9951) -- (5.9927,3.9951) -- (5.9927,3.9951) -- (5.9927,3.9951) -- (5.9927,3.9952) -- (5.9927,3.9952) -- (5.9928,3.9952) -- (5.9928,3.9952) -- (5.9928,3.9952) -- (5.9928,3.9952) -- (5.9928,3.9952) -- (5.9928,3.9952) -- (5.9928,3.9952) -- (5.9929,3.9952) -- (5.9929,3.9953) -- (5.9929,3.9953) -- (5.9929,3.9953) -- (5.9929,3.9953) -- (5.9929,3.9953) -- (5.9929,3.9953) -- (5.9930,3.9953) -- (5.9930,3.9953) -- (5.9930,3.9953) -- (5.9930,3.9953) -- (5.9930,3.9953) -- (5.9930,3.9954) -- (5.9930,3.9954) -- (5.9931,3.9954) -- (5.9931,3.9954) -- (5.9931,3.9954) -- (5.9931,3.9954) -- (5.9931,3.9954) -- (5.9931,3.9954) -- (5.9931,3.9954) -- (5.9932,3.9954) -- (5.9932,3.9954) -- (5.9932,3.9955) -- (5.9932,3.9955) -- (5.9932,3.9955) -- (5.9932,3.9955) -- (5.9932,3.9955) -- (5.9933,3.9955) -- (5.9933,3.9955) -- (5.9933,3.9955) -- (5.9933,3.9955) -- (5.9933,3.9955);
\fill[blue!80!black] (5.9400,3.9600) circle (1.2pt);
\fill[blue!80!black] (5.9401,3.9601) circle (1.2pt);
\fill[blue!80!black] (5.9402,3.9602) circle (1.2pt);
\fill[blue!80!black] (5.9404,3.9602) circle (1.2pt);
\fill[blue!80!black] (5.9405,3.9603) circle (1.2pt);
\fill[blue!80!black] (5.9406,3.9604) circle (1.2pt);
\fill[blue!80!black] (5.9407,3.9605) circle (1.2pt);
\fill[blue!80!black] (5.9408,3.9606) circle (1.2pt);
\fill[blue!80!black] (5.9409,3.9606) circle (1.2pt);
\fill[blue!80!black] (5.9411,3.9607) circle (1.2pt);
\fill[blue!80!black] (5.9412,3.9608) circle (1.2pt);
\fill[blue!80!black] (5.9413,3.9609) circle (1.2pt);
\fill[blue!80!black] (5.9414,3.9609) circle (1.2pt);
\fill[blue!80!black] (5.9415,3.9610) circle (1.2pt);
\fill[blue!80!black] (5.9416,3.9611) circle (1.2pt);
\fill[blue!80!black] (5.9418,3.9612) circle (1.2pt);
\fill[blue!80!black] (5.9419,3.9613) circle (1.2pt);
\fill[blue!80!black] (5.9420,3.9613) circle (1.2pt);
\fill[blue!80!black] (5.9421,3.9614) circle (1.2pt);
\fill[blue!80!black] (5.9422,3.9615) circle (1.2pt);
\fill[blue!80!black] (5.9423,3.9616) circle (1.2pt);
\fill[blue!80!black] (5.9425,3.9616) circle (1.2pt);
\fill[blue!80!black] (5.9426,3.9617) circle (1.2pt);
\fill[blue!80!black] (5.9427,3.9618) circle (1.2pt);
\fill[blue!80!black] (5.9428,3.9619) circle (1.2pt);
\fill[blue!80!black] (5.9429,3.9619) circle (1.2pt);
\fill[blue!80!black] (5.9430,3.9620) circle (1.2pt);
\fill[blue!80!black] (5.9431,3.9621) circle (1.2pt);
\fill[blue!80!black] (5.9433,3.9622) circle (1.2pt);
\fill[blue!80!black] (5.9434,3.9622) circle (1.2pt);
\fill[blue!80!black] (5.9435,3.9623) circle (1.2pt);
\fill[blue!80!black] (5.9436,3.9624) circle (1.2pt);
\fill[blue!80!black] (5.9437,3.9625) circle (1.2pt);
\fill[blue!80!black] (5.9438,3.9625) circle (1.2pt);
\fill[blue!80!black] (5.9439,3.9626) circle (1.2pt);
\fill[blue!80!black] (5.9440,3.9627) circle (1.2pt);
\fill[blue!80!black] (5.9441,3.9628) circle (1.2pt);
\fill[blue!80!black] (5.9443,3.9628) circle (1.2pt);
\fill[blue!80!black] (5.9444,3.9629) circle (1.2pt);
\fill[blue!80!black] (5.9445,3.9630) circle (1.2pt);
\fill[blue!80!black] (5.9446,3.9631) circle (1.2pt);
\fill[blue!80!black] (5.9447,3.9631) circle (1.2pt);
\fill[blue!80!black] (5.9448,3.9632) circle (1.2pt);
\fill[blue!80!black] (5.9449,3.9633) circle (1.2pt);
\fill[blue!80!black] (5.9450,3.9634) circle (1.2pt);
\fill[blue!80!black] (5.9451,3.9634) circle (1.2pt);
\fill[blue!80!black] (5.9453,3.9635) circle (1.2pt);
\fill[blue!80!black] (5.9454,3.9636) circle (1.2pt);
\fill[blue!80!black] (5.9455,3.9636) circle (1.2pt);
\fill[blue!80!black] (5.9456,3.9637) circle (1.2pt);
\fill[blue!80!black] (5.9457,3.9638) circle (1.2pt);
\fill[blue!80!black] (5.9458,3.9639) circle (1.2pt);
\fill[blue!80!black] (5.9459,3.9639) circle (1.2pt);
\fill[blue!80!black] (5.9460,3.9640) circle (1.2pt);
\fill[blue!80!black] (5.9461,3.9641) circle (1.2pt);
\fill[blue!80!black] (5.9462,3.9641) circle (1.2pt);
\fill[blue!80!black] (5.9463,3.9642) circle (1.2pt);
\fill[blue!80!black] (5.9464,3.9643) circle (1.2pt);
\fill[blue!80!black] (5.9465,3.9644) circle (1.2pt);
\fill[blue!80!black] (5.9467,3.9644) circle (1.2pt);
\fill[blue!80!black] (5.9468,3.9645) circle (1.2pt);
\fill[blue!80!black] (5.9469,3.9646) circle (1.2pt);
\fill[blue!80!black] (5.9470,3.9646) circle (1.2pt);
\fill[blue!80!black] (5.9471,3.9647) circle (1.2pt);
\fill[blue!80!black] (5.9472,3.9648) circle (1.2pt);
\fill[blue!80!black] (5.9473,3.9649) circle (1.2pt);
\fill[blue!80!black] (5.9474,3.9649) circle (1.2pt);
\fill[blue!80!black] (5.9475,3.9650) circle (1.2pt);
\fill[blue!80!black] (5.9476,3.9651) circle (1.2pt);
\fill[blue!80!black] (5.9477,3.9651) circle (1.2pt);
\fill[blue!80!black] (5.9478,3.9652) circle (1.2pt);
\fill[blue!80!black] (5.9479,3.9653) circle (1.2pt);
\fill[blue!80!black] (5.9480,3.9653) circle (1.2pt);
\fill[blue!80!black] (5.9481,3.9654) circle (1.2pt);
\fill[blue!80!black] (5.9482,3.9655) circle (1.2pt);
\fill[blue!80!black] (5.9483,3.9656) circle (1.2pt);
\fill[blue!80!black] (5.9484,3.9656) circle (1.2pt);
\fill[blue!80!black] (5.9485,3.9657) circle (1.2pt);
\fill[blue!80!black] (5.9486,3.9658) circle (1.2pt);
\fill[blue!80!black] (5.9487,3.9658) circle (1.2pt);
\fill[blue!80!black] (5.9488,3.9659) circle (1.2pt);
\fill[blue!80!black] (5.9489,3.9660) circle (1.2pt);
\fill[blue!80!black] (5.9490,3.9660) circle (1.2pt);
\fill[blue!80!black] (5.9491,3.9661) circle (1.2pt);
\fill[blue!80!black] (5.9492,3.9662) circle (1.2pt);
\fill[blue!80!black] (5.9493,3.9662) circle (1.2pt);
\fill[blue!80!black] (5.9494,3.9663) circle (1.2pt);
\fill[blue!80!black] (5.9495,3.9664) circle (1.2pt);
\fill[blue!80!black] (5.9496,3.9664) circle (1.2pt);
\fill[blue!80!black] (5.9497,3.9665) circle (1.2pt);
\fill[blue!80!black] (5.9498,3.9666) circle (1.2pt);
\fill[blue!80!black] (5.9499,3.9666) circle (1.2pt);
\fill[blue!80!black] (5.9500,3.9667) circle (1.2pt);
\fill[blue!80!black] (5.9501,3.9668) circle (1.2pt);
\fill[blue!80!black] (5.9502,3.9668) circle (1.2pt);
\fill[blue!80!black] (5.9503,3.9669) circle (1.2pt);
\fill[blue!80!black] (5.9504,3.9670) circle (1.2pt);
\fill[blue!80!black] (5.9505,3.9670) circle (1.2pt);
\fill[blue!80!black] (5.9506,3.9671) circle (1.2pt);
\fill[blue!80!black] (5.9507,3.9672) circle (1.2pt);
\fill[blue!80!black] (5.9508,3.9672) circle (1.2pt);
\fill[blue!80!black] (5.9509,3.9673) circle (1.2pt);
\fill[blue!80!black] (5.9510,3.9674) circle (1.2pt);
\fill[blue!80!black] (5.9511,3.9674) circle (1.2pt);
\fill[blue!80!black] (5.9512,3.9675) circle (1.2pt);
\fill[blue!80!black] (5.9513,3.9676) circle (1.2pt);
\fill[blue!80!black] (5.9514,3.9676) circle (1.2pt);
\fill[blue!80!black] (5.9515,3.9677) circle (1.2pt);
\fill[blue!80!black] (5.9516,3.9677) circle (1.2pt);
\fill[blue!80!black] (5.9517,3.9678) circle (1.2pt);
\fill[blue!80!black] (5.9518,3.9679) circle (1.2pt);
\fill[blue!80!black] (5.9519,3.9679) circle (1.2pt);
\fill[blue!80!black] (5.9520,3.9680) circle (1.2pt);
\fill[blue!80!black] (5.9521,3.9681) circle (1.2pt);
\fill[blue!80!black] (5.9522,3.9681) circle (1.2pt);
\fill[blue!80!black] (5.9523,3.9682) circle (1.2pt);
\fill[blue!80!black] (5.9524,3.9683) circle (1.2pt);
\fill[blue!80!black] (5.9525,3.9683) circle (1.2pt);
\fill[blue!80!black] (5.9526,3.9684) circle (1.2pt);
\fill[blue!80!black] (5.9527,3.9684) circle (1.2pt);
\fill[blue!80!black] (5.9528,3.9685) circle (1.2pt);
\fill[blue!80!black] (5.9529,3.9686) circle (1.2pt);
\fill[blue!80!black] (5.9529,3.9686) circle (1.2pt);
\fill[blue!80!black] (5.9530,3.9687) circle (1.2pt);
\fill[blue!80!black] (5.9531,3.9688) circle (1.2pt);
\fill[blue!80!black] (5.9532,3.9688) circle (1.2pt);
\fill[blue!80!black] (5.9533,3.9689) circle (1.2pt);
\fill[blue!80!black] (5.9534,3.9689) circle (1.2pt);
\fill[blue!80!black] (5.9535,3.9690) circle (1.2pt);
\fill[blue!80!black] (5.9536,3.9691) circle (1.2pt);
\fill[blue!80!black] (5.9537,3.9691) circle (1.2pt);
\fill[blue!80!black] (5.9538,3.9692) circle (1.2pt);
\fill[blue!80!black] (5.9539,3.9693) circle (1.2pt);
\fill[blue!80!black] (5.9540,3.9693) circle (1.2pt);
\fill[blue!80!black] (5.9541,3.9694) circle (1.2pt);
\fill[blue!80!black] (5.9542,3.9694) circle (1.2pt);
\fill[blue!80!black] (5.9542,3.9695) circle (1.2pt);
\fill[blue!80!black] (5.9543,3.9696) circle (1.2pt);
\fill[blue!80!black] (5.9544,3.9696) circle (1.2pt);
\fill[blue!80!black] (5.9545,3.9697) circle (1.2pt);
\fill[blue!80!black] (5.9546,3.9697) circle (1.2pt);
\fill[blue!80!black] (5.9547,3.9698) circle (1.2pt);
\fill[blue!80!black] (5.9548,3.9699) circle (1.2pt);
\fill[blue!80!black] (5.9549,3.9699) circle (1.2pt);
\fill[blue!80!black] (5.9550,3.9700) circle (1.2pt);
\fill[blue!80!black] (5.9551,3.9700) circle (1.2pt);
\fill[blue!80!black] (5.9551,3.9701) circle (1.2pt);
\fill[blue!80!black] (5.9552,3.9702) circle (1.2pt);
\fill[blue!80!black] (5.9553,3.9702) circle (1.2pt);
\fill[blue!80!black] (5.9554,3.9703) circle (1.2pt);
\fill[blue!80!black] (5.9555,3.9703) circle (1.2pt);
\fill[blue!80!black] (5.9556,3.9704) circle (1.2pt);
\fill[blue!80!black] (5.9557,3.9705) circle (1.2pt);
\fill[blue!80!black] (5.9558,3.9705) circle (1.2pt);
\fill[blue!80!black] (5.9559,3.9706) circle (1.2pt);
\fill[blue!80!black] (5.9559,3.9706) circle (1.2pt);
\fill[blue!80!black] (5.9560,3.9707) circle (1.2pt);
\fill[blue!80!black] (5.9561,3.9707) circle (1.2pt);
\fill[blue!80!black] (5.9562,3.9708) circle (1.2pt);
\fill[blue!80!black] (5.9563,3.9709) circle (1.2pt);
\fill[blue!80!black] (5.9564,3.9709) circle (1.2pt);
\fill[blue!80!black] (5.9565,3.9710) circle (1.2pt);
\fill[blue!80!black] (5.9566,3.9710) circle (1.2pt);
\fill[blue!80!black] (5.9566,3.9711) circle (1.2pt);
\fill[blue!80!black] (5.9567,3.9712) circle (1.2pt);
\fill[blue!80!black] (5.9568,3.9712) circle (1.2pt);
\fill[blue!80!black] (5.9569,3.9713) circle (1.2pt);
\fill[blue!80!black] (5.9570,3.9713) circle (1.2pt);
\fill[blue!80!black] (5.9571,3.9714) circle (1.2pt);
\fill[blue!80!black] (5.9572,3.9714) circle (1.2pt);
\fill[blue!80!black] (5.9572,3.9715) circle (1.2pt);
\fill[blue!80!black] (5.9573,3.9716) circle (1.2pt);
\fill[blue!80!black] (5.9574,3.9716) circle (1.2pt);
\fill[blue!80!black] (5.9575,3.9717) circle (1.2pt);
\fill[blue!80!black] (5.9576,3.9717) circle (1.2pt);
\fill[blue!80!black] (5.9577,3.9718) circle (1.2pt);
\fill[blue!80!black] (5.9578,3.9718) circle (1.2pt);
\fill[blue!80!black] (5.9578,3.9719) circle (1.2pt);
\fill[blue!80!black] (5.9579,3.9719) circle (1.2pt);
\fill[blue!80!black] (5.9580,3.9720) circle (1.2pt);
\fill[blue!80!black] (5.9581,3.9721) circle (1.2pt);
\fill[blue!80!black] (5.9582,3.9721) circle (1.2pt);
\fill[blue!80!black] (5.9583,3.9722) circle (1.2pt);
\fill[blue!80!black] (5.9583,3.9722) circle (1.2pt);
\fill[blue!80!black] (5.9584,3.9723) circle (1.2pt);
\fill[blue!80!black] (5.9585,3.9723) circle (1.2pt);
\fill[blue!80!black] (5.9586,3.9724) circle (1.2pt);
\fill[blue!80!black] (5.9587,3.9724) circle (1.2pt);
\fill[blue!80!black] (5.9588,3.9725) circle (1.2pt);
\fill[blue!80!black] (5.9588,3.9726) circle (1.2pt);
\fill[blue!80!black] (5.9589,3.9726) circle (1.2pt);
\fill[blue!80!black] (5.9590,3.9727) circle (1.2pt);
\fill[blue!80!black] (5.9591,3.9727) circle (1.2pt);
\fill[blue!80!black] (5.9592,3.9728) circle (1.2pt);
\fill[blue!80!black] (5.9592,3.9728) circle (1.2pt);
\fill[blue!80!black] (5.9593,3.9729) circle (1.2pt);
\fill[blue!80!black] (5.9594,3.9729) circle (1.2pt);
\fill[blue!80!black] (5.9595,3.9730) circle (1.2pt);
\fill[blue!80!black] (5.9596,3.9730) circle (1.2pt);
\fill[blue!80!black] (5.9596,3.9731) circle (1.2pt);
\fill[blue!80!black] (5.9597,3.9732) circle (1.2pt);
\fill[blue!80!black] (5.9598,3.9732) circle (1.2pt);
\fill[blue!80!black] (5.9599,3.9733) circle (1.2pt);
\fill[blue!80!black] (5.9600,3.9733) circle (1.2pt);
\fill[blue!80!black] (5.9601,3.9734) circle (1.2pt);
\fill[blue!80!black] (5.9601,3.9734) circle (1.2pt);
\fill[blue!80!black] (5.9602,3.9735) circle (1.2pt);
\fill[blue!80!black] (5.9603,3.9735) circle (1.2pt);
\fill[blue!80!black] (5.9604,3.9736) circle (1.2pt);
\fill[blue!80!black] (5.9604,3.9736) circle (1.2pt);
\fill[blue!80!black] (5.9605,3.9737) circle (1.2pt);
\fill[blue!80!black] (5.9606,3.9737) circle (1.2pt);
\fill[blue!80!black] (5.9607,3.9738) circle (1.2pt);
\fill[blue!80!black] (5.9608,3.9738) circle (1.2pt);
\fill[blue!80!black] (5.9608,3.9739) circle (1.2pt);
\fill[blue!80!black] (5.9609,3.9739) circle (1.2pt);
\fill[blue!80!black] (5.9610,3.9740) circle (1.2pt);
\fill[blue!80!black] (5.9611,3.9740) circle (1.2pt);
\fill[blue!80!black] (5.9612,3.9741) circle (1.2pt);
\fill[blue!80!black] (5.9612,3.9742) circle (1.2pt);
\fill[blue!80!black] (5.9613,3.9742) circle (1.2pt);
\fill[blue!80!black] (5.9614,3.9743) circle (1.2pt);
\fill[blue!80!black] (5.9615,3.9743) circle (1.2pt);
\fill[blue!80!black] (5.9615,3.9744) circle (1.2pt);
\fill[blue!80!black] (5.9616,3.9744) circle (1.2pt);
\fill[blue!80!black] (5.9617,3.9745) circle (1.2pt);
\fill[blue!80!black] (5.9618,3.9745) circle (1.2pt);
\fill[blue!80!black] (5.9618,3.9746) circle (1.2pt);
\fill[blue!80!black] (5.9619,3.9746) circle (1.2pt);
\fill[blue!80!black] (5.9620,3.9747) circle (1.2pt);
\fill[blue!80!black] (5.9621,3.9747) circle (1.2pt);
\fill[blue!80!black] (5.9621,3.9748) circle (1.2pt);
\fill[blue!80!black] (5.9622,3.9748) circle (1.2pt);
\fill[blue!80!black] (5.9623,3.9749) circle (1.2pt);
\fill[blue!80!black] (5.9624,3.9749) circle (1.2pt);
\fill[blue!80!black] (5.9624,3.9750) circle (1.2pt);
\fill[blue!80!black] (5.9625,3.9750) circle (1.2pt);
\fill[blue!80!black] (5.9626,3.9751) circle (1.2pt);
\fill[blue!80!black] (5.9627,3.9751) circle (1.2pt);
\fill[blue!80!black] (5.9627,3.9752) circle (1.2pt);
\fill[blue!80!black] (5.9628,3.9752) circle (1.2pt);
\fill[blue!80!black] (5.9629,3.9753) circle (1.2pt);
\fill[blue!80!black] (5.9630,3.9753) circle (1.2pt);
\fill[blue!80!black] (5.9630,3.9754) circle (1.2pt);
\fill[blue!80!black] (5.9631,3.9754) circle (1.2pt);
\fill[blue!80!black] (5.9632,3.9755) circle (1.2pt);
\fill[blue!80!black] (5.9633,3.9755) circle (1.2pt);
\fill[blue!80!black] (5.9633,3.9756) circle (1.2pt);
\fill[blue!80!black] (5.9634,3.9756) circle (1.2pt);
\fill[blue!80!black] (5.9635,3.9757) circle (1.2pt);
\fill[blue!80!black] (5.9636,3.9757) circle (1.2pt);
\fill[blue!80!black] (5.9636,3.9758) circle (1.2pt);
\fill[blue!80!black] (5.9637,3.9758) circle (1.2pt);
\fill[blue!80!black] (5.9638,3.9758) circle (1.2pt);
\fill[blue!80!black] (5.9638,3.9759) circle (1.2pt);
\fill[blue!80!black] (5.9639,3.9759) circle (1.2pt);
\fill[blue!80!black] (5.9640,3.9760) circle (1.2pt);
\fill[blue!80!black] (5.9641,3.9760) circle (1.2pt);
\fill[blue!80!black] (5.9641,3.9761) circle (1.2pt);
\fill[blue!80!black] (5.9642,3.9761) circle (1.2pt);
\fill[blue!80!black] (5.9643,3.9762) circle (1.2pt);
\fill[blue!80!black] (5.9643,3.9762) circle (1.2pt);
\fill[blue!80!black] (5.9644,3.9763) circle (1.2pt);
\fill[blue!80!black] (5.9645,3.9763) circle (1.2pt);
\fill[blue!80!black] (5.9646,3.9764) circle (1.2pt);
\fill[blue!80!black] (5.9646,3.9764) circle (1.2pt);
\fill[blue!80!black] (5.9647,3.9765) circle (1.2pt);
\fill[blue!80!black] (5.9648,3.9765) circle (1.2pt);
\fill[blue!80!black] (5.9648,3.9766) circle (1.2pt);
\fill[blue!80!black] (5.9649,3.9766) circle (1.2pt);
\fill[blue!80!black] (5.9650,3.9767) circle (1.2pt);
\fill[blue!80!black] (5.9651,3.9767) circle (1.2pt);
\fill[blue!80!black] (5.9651,3.9767) circle (1.2pt);
\fill[blue!80!black] (5.9652,3.9768) circle (1.2pt);
\fill[blue!80!black] (5.9653,3.9768) circle (1.2pt);
\fill[blue!80!black] (5.9653,3.9769) circle (1.2pt);
\fill[blue!80!black] (5.9654,3.9769) circle (1.2pt);
\fill[blue!80!black] (5.9655,3.9770) circle (1.2pt);
\fill[blue!80!black] (5.9655,3.9770) circle (1.2pt);
\fill[blue!80!black] (5.9656,3.9771) circle (1.2pt);
\fill[blue!80!black] (5.9657,3.9771) circle (1.2pt);
\fill[blue!80!black] (5.9657,3.9772) circle (1.2pt);
\fill[blue!80!black] (5.9658,3.9772) circle (1.2pt);
\fill[blue!80!black] (5.9659,3.9773) circle (1.2pt);
\fill[blue!80!black] (5.9660,3.9773) circle (1.2pt);
\fill[blue!80!black] (5.9660,3.9773) circle (1.2pt);
\fill[blue!80!black] (5.9661,3.9774) circle (1.2pt);
\fill[blue!80!black] (5.9662,3.9774) circle (1.2pt);
\fill[blue!80!black] (5.9662,3.9775) circle (1.2pt);
\fill[blue!80!black] (5.9663,3.9775) circle (1.2pt);
\fill[blue!80!black] (5.9664,3.9776) circle (1.2pt);
\fill[blue!80!black] (5.9664,3.9776) circle (1.2pt);
\fill[blue!80!black] (5.9665,3.9777) circle (1.2pt);
\fill[blue!80!black] (5.9666,3.9777) circle (1.2pt);
\fill[blue!80!black] (5.9666,3.9777) circle (1.2pt);
\fill[blue!80!black] (5.9667,3.9778) circle (1.2pt);
\fill[blue!80!black] (5.9668,3.9778) circle (1.2pt);
\fill[blue!80!black] (5.9668,3.9779) circle (1.2pt);
\fill[blue!80!black] (5.9669,3.9779) circle (1.2pt);
\fill[blue!80!black] (5.9670,3.9780) circle (1.2pt);
\fill[blue!80!black] (5.9670,3.9780) circle (1.2pt);
\fill[blue!80!black] (5.9671,3.9781) circle (1.2pt);
\fill[blue!80!black] (5.9672,3.9781) circle (1.2pt);
\fill[blue!80!black] (5.9672,3.9781) circle (1.2pt);
\fill[blue!80!black] (5.9673,3.9782) circle (1.2pt);
\fill[blue!80!black] (5.9673,3.9782) circle (1.2pt);
\fill[blue!80!black] (5.9674,3.9783) circle (1.2pt);
\fill[blue!80!black] (5.9675,3.9783) circle (1.2pt);
\fill[blue!80!black] (5.9675,3.9784) circle (1.2pt);
\fill[blue!80!black] (5.9676,3.9784) circle (1.2pt);
\fill[blue!80!black] (5.9677,3.9784) circle (1.2pt);
\fill[blue!80!black] (5.9677,3.9785) circle (1.2pt);
\fill[blue!80!black] (5.9678,3.9785) circle (1.2pt);
\fill[blue!80!black] (5.9679,3.9786) circle (1.2pt);
\fill[blue!80!black] (5.9679,3.9786) circle (1.2pt);
\fill[blue!80!black] (5.9680,3.9787) circle (1.2pt);
\fill[blue!80!black] (5.9681,3.9787) circle (1.2pt);
\fill[blue!80!black] (5.9681,3.9787) circle (1.2pt);
\fill[blue!80!black] (5.9682,3.9788) circle (1.2pt);
\fill[blue!80!black] (5.9683,3.9788) circle (1.2pt);
\fill[blue!80!black] (5.9683,3.9789) circle (1.2pt);
\fill[blue!80!black] (5.9684,3.9789) circle (1.2pt);
\fill[blue!80!black] (5.9684,3.9790) circle (1.2pt);
\fill[blue!80!black] (5.9685,3.9790) circle (1.2pt);
\fill[blue!80!black] (5.9686,3.9790) circle (1.2pt);
\fill[blue!80!black] (5.9686,3.9791) circle (1.2pt);
\fill[blue!80!black] (5.9687,3.9791) circle (1.2pt);
\fill[blue!80!black] (5.9688,3.9792) circle (1.2pt);
\fill[blue!80!black] (5.9688,3.9792) circle (1.2pt);
\fill[blue!80!black] (5.9689,3.9793) circle (1.2pt);
\fill[blue!80!black] (5.9689,3.9793) circle (1.2pt);
\fill[blue!80!black] (5.9690,3.9793) circle (1.2pt);
\fill[blue!80!black] (5.9691,3.9794) circle (1.2pt);
\fill[blue!80!black] (5.9691,3.9794) circle (1.2pt);
\fill[blue!80!black] (5.9692,3.9795) circle (1.2pt);
\fill[blue!80!black] (5.9692,3.9795) circle (1.2pt);
\fill[blue!80!black] (5.9693,3.9795) circle (1.2pt);
\fill[blue!80!black] (5.9694,3.9796) circle (1.2pt);
\fill[blue!80!black] (5.9694,3.9796) circle (1.2pt);
\fill[blue!80!black] (5.9695,3.9797) circle (1.2pt);
\fill[blue!80!black] (5.9696,3.9797) circle (1.2pt);
\fill[blue!80!black] (5.9696,3.9797) circle (1.2pt);
\fill[blue!80!black] (5.9697,3.9798) circle (1.2pt);
\fill[blue!80!black] (5.9697,3.9798) circle (1.2pt);
\fill[blue!80!black] (5.9698,3.9799) circle (1.2pt);
\fill[blue!80!black] (5.9699,3.9799) circle (1.2pt);
\fill[blue!80!black] (5.9699,3.9799) circle (1.2pt);
\fill[blue!80!black] (5.9700,3.9800) circle (1.2pt);
\fill[blue!80!black] (5.9700,3.9800) circle (1.2pt);
\fill[blue!80!black] (5.9701,3.9801) circle (1.2pt);
\fill[blue!80!black] (5.9702,3.9801) circle (1.2pt);
\fill[blue!80!black] (5.9702,3.9801) circle (1.2pt);
\fill[blue!80!black] (5.9703,3.9802) circle (1.2pt);
\fill[blue!80!black] (5.9703,3.9802) circle (1.2pt);
\fill[blue!80!black] (5.9704,3.9803) circle (1.2pt);
\fill[blue!80!black] (5.9705,3.9803) circle (1.2pt);
\fill[blue!80!black] (5.9705,3.9803) circle (1.2pt);
\fill[blue!80!black] (5.9706,3.9804) circle (1.2pt);
\fill[blue!80!black] (5.9706,3.9804) circle (1.2pt);
\fill[blue!80!black] (5.9707,3.9805) circle (1.2pt);
\fill[blue!80!black] (5.9707,3.9805) circle (1.2pt);
\fill[blue!80!black] (5.9708,3.9805) circle (1.2pt);
\fill[blue!80!black] (5.9709,3.9806) circle (1.2pt);
\fill[blue!80!black] (5.9709,3.9806) circle (1.2pt);
\fill[blue!80!black] (5.9710,3.9807) circle (1.2pt);
\fill[blue!80!black] (5.9710,3.9807) circle (1.2pt);
\fill[blue!80!black] (5.9711,3.9807) circle (1.2pt);
\fill[blue!80!black] (5.9712,3.9808) circle (1.2pt);
\fill[blue!80!black] (5.9712,3.9808) circle (1.2pt);
\fill[blue!80!black] (5.9713,3.9808) circle (1.2pt);
\fill[blue!80!black] (5.9713,3.9809) circle (1.2pt);
\fill[blue!80!black] (5.9714,3.9809) circle (1.2pt);
\fill[blue!80!black] (5.9714,3.9810) circle (1.2pt);
\fill[blue!80!black] (5.9715,3.9810) circle (1.2pt);
\fill[blue!80!black] (5.9716,3.9810) circle (1.2pt);
\fill[blue!80!black] (5.9716,3.9811) circle (1.2pt);
\fill[blue!80!black] (5.9717,3.9811) circle (1.2pt);
\fill[blue!80!black] (5.9717,3.9812) circle (1.2pt);
\fill[blue!80!black] (5.9718,3.9812) circle (1.2pt);
\fill[blue!80!black] (5.9718,3.9812) circle (1.2pt);
\fill[blue!80!black] (5.9719,3.9813) circle (1.2pt);
\fill[blue!80!black] (5.9720,3.9813) circle (1.2pt);
\fill[blue!80!black] (5.9720,3.9813) circle (1.2pt);
\fill[blue!80!black] (5.9721,3.9814) circle (1.2pt);
\fill[blue!80!black] (5.9721,3.9814) circle (1.2pt);
\fill[blue!80!black] (5.9722,3.9814) circle (1.2pt);
\fill[blue!80!black] (5.9722,3.9815) circle (1.2pt);
\fill[blue!80!black] (5.9723,3.9815) circle (1.2pt);
\fill[blue!80!black] (5.9723,3.9816) circle (1.2pt);
\fill[blue!80!black] (5.9724,3.9816) circle (1.2pt);
\fill[blue!80!black] (5.9725,3.9816) circle (1.2pt);
\fill[blue!80!black] (5.9725,3.9817) circle (1.2pt);
\fill[blue!80!black] (5.9726,3.9817) circle (1.2pt);
\fill[blue!80!black] (5.9726,3.9817) circle (1.2pt);
\fill[blue!80!black] (5.9727,3.9818) circle (1.2pt);
\fill[blue!80!black] (5.9727,3.9818) circle (1.2pt);
\fill[blue!80!black] (5.9728,3.9819) circle (1.2pt);
\fill[blue!80!black] (5.9728,3.9819) circle (1.2pt);
\fill[blue!80!black] (5.9729,3.9819) circle (1.2pt);
\fill[blue!80!black] (5.9729,3.9820) circle (1.2pt);
\fill[blue!80!black] (5.9730,3.9820) circle (1.2pt);
\fill[blue!80!black] (5.9731,3.9820) circle (1.2pt);
\fill[blue!80!black] (5.9731,3.9821) circle (1.2pt);
\fill[blue!80!black] (5.9732,3.9821) circle (1.2pt);
\fill[blue!80!black] (5.9732,3.9821) circle (1.2pt);
\fill[blue!80!black] (5.9733,3.9822) circle (1.2pt);
\fill[blue!80!black] (5.9733,3.9822) circle (1.2pt);
\fill[blue!80!black] (5.9734,3.9822) circle (1.2pt);
\fill[blue!80!black] (5.9734,3.9823) circle (1.2pt);
\fill[blue!80!black] (5.9735,3.9823) circle (1.2pt);
\fill[blue!80!black] (5.9735,3.9824) circle (1.2pt);
\fill[blue!80!black] (5.9736,3.9824) circle (1.2pt);
\fill[blue!80!black] (5.9736,3.9824) circle (1.2pt);
\fill[blue!80!black] (5.9737,3.9825) circle (1.2pt);
\fill[blue!80!black] (5.9737,3.9825) circle (1.2pt);
\fill[blue!80!black] (5.9738,3.9825) circle (1.2pt);
\fill[blue!80!black] (5.9738,3.9826) circle (1.2pt);
\fill[blue!80!black] (5.9739,3.9826) circle (1.2pt);
\fill[blue!80!black] (5.9740,3.9826) circle (1.2pt);
\fill[blue!80!black] (5.9740,3.9827) circle (1.2pt);
\fill[blue!80!black] (5.9741,3.9827) circle (1.2pt);
\fill[blue!80!black] (5.9741,3.9827) circle (1.2pt);
\fill[blue!80!black] (5.9742,3.9828) circle (1.2pt);
\fill[blue!80!black] (5.9742,3.9828) circle (1.2pt);
\fill[blue!80!black] (5.9743,3.9828) circle (1.2pt);
\fill[blue!80!black] (5.9743,3.9829) circle (1.2pt);
\fill[blue!80!black] (5.9744,3.9829) circle (1.2pt);
\fill[blue!80!black] (5.9744,3.9829) circle (1.2pt);
\fill[blue!80!black] (5.9745,3.9830) circle (1.2pt);
\fill[blue!80!black] (5.9745,3.9830) circle (1.2pt);
\fill[blue!80!black] (5.9746,3.9830) circle (1.2pt);
\fill[blue!80!black] (5.9746,3.9831) circle (1.2pt);
\fill[blue!80!black] (5.9747,3.9831) circle (1.2pt);
\fill[blue!80!black] (5.9747,3.9831) circle (1.2pt);
\fill[blue!80!black] (5.9748,3.9832) circle (1.2pt);
\fill[blue!80!black] (5.9748,3.9832) circle (1.2pt);
\fill[blue!80!black] (5.9749,3.9832) circle (1.2pt);
\fill[blue!80!black] (5.9749,3.9833) circle (1.2pt);
\fill[blue!80!black] (5.9750,3.9833) circle (1.2pt);
\fill[blue!80!black] (5.9750,3.9833) circle (1.2pt);
\fill[blue!80!black] (5.9751,3.9834) circle (1.2pt);
\fill[blue!80!black] (5.9751,3.9834) circle (1.2pt);
\fill[blue!80!black] (5.9752,3.9834) circle (1.2pt);
\fill[blue!80!black] (5.9752,3.9835) circle (1.2pt);
\fill[blue!80!black] (5.9753,3.9835) circle (1.2pt);
\fill[blue!80!black] (5.9753,3.9835) circle (1.2pt);
\fill[blue!80!black] (5.9754,3.9836) circle (1.2pt);
\fill[blue!80!black] (5.9754,3.9836) circle (1.2pt);
\fill[blue!80!black] (5.9755,3.9836) circle (1.2pt);
\fill[blue!80!black] (5.9755,3.9837) circle (1.2pt);
\fill[blue!80!black] (5.9756,3.9837) circle (1.2pt);
\fill[blue!80!black] (5.9756,3.9837) circle (1.2pt);
\fill[blue!80!black] (5.9757,3.9838) circle (1.2pt);
\fill[blue!80!black] (5.9757,3.9838) circle (1.2pt);
\fill[blue!80!black] (5.9758,3.9838) circle (1.2pt);
\fill[blue!80!black] (5.9758,3.9839) circle (1.2pt);
\fill[blue!80!black] (5.9759,3.9839) circle (1.2pt);
\fill[blue!80!black] (5.9759,3.9839) circle (1.2pt);
\fill[blue!80!black] (5.9760,3.9840) circle (1.2pt);
\fill[blue!80!black] (5.9760,3.9840) circle (1.2pt);
\fill[blue!80!black] (5.9761,3.9840) circle (1.2pt);
\fill[blue!80!black] (5.9761,3.9841) circle (1.2pt);
\fill[blue!80!black] (5.9761,3.9841) circle (1.2pt);
\fill[blue!80!black] (5.9762,3.9841) circle (1.2pt);
\fill[blue!80!black] (5.9762,3.9842) circle (1.2pt);
\fill[blue!80!black] (5.9763,3.9842) circle (1.2pt);
\fill[blue!80!black] (5.9763,3.9842) circle (1.2pt);
\fill[blue!80!black] (5.9764,3.9843) circle (1.2pt);
\fill[blue!80!black] (5.9764,3.9843) circle (1.2pt);
\fill[blue!80!black] (5.9765,3.9843) circle (1.2pt);
\fill[blue!80!black] (5.9765,3.9843) circle (1.2pt);
\fill[blue!80!black] (5.9766,3.9844) circle (1.2pt);
\fill[blue!80!black] (5.9766,3.9844) circle (1.2pt);
\fill[blue!80!black] (5.9767,3.9844) circle (1.2pt);
\fill[blue!80!black] (5.9767,3.9845) circle (1.2pt);
\fill[blue!80!black] (5.9768,3.9845) circle (1.2pt);
\fill[blue!80!black] (5.9768,3.9845) circle (1.2pt);
\fill[blue!80!black] (5.9769,3.9846) circle (1.2pt);
\fill[blue!80!black] (5.9769,3.9846) circle (1.2pt);
\fill[blue!80!black] (5.9769,3.9846) circle (1.2pt);
\fill[blue!80!black] (5.9770,3.9847) circle (1.2pt);
\fill[blue!80!black] (5.9770,3.9847) circle (1.2pt);
\fill[blue!80!black] (5.9771,3.9847) circle (1.2pt);
\fill[blue!80!black] (5.9771,3.9848) circle (1.2pt);
\fill[blue!80!black] (5.9772,3.9848) circle (1.2pt);
\fill[blue!80!black] (5.9772,3.9848) circle (1.2pt);
\fill[blue!80!black] (5.9773,3.9848) circle (1.2pt);
\fill[blue!80!black] (5.9773,3.9849) circle (1.2pt);
\fill[blue!80!black] (5.9774,3.9849) circle (1.2pt);
\fill[blue!80!black] (5.9774,3.9849) circle (1.2pt);
\fill[blue!80!black] (5.9774,3.9850) circle (1.2pt);
\fill[blue!80!black] (5.9775,3.9850) circle (1.2pt);
\fill[blue!80!black] (5.9775,3.9850) circle (1.2pt);
\fill[blue!80!black] (5.9776,3.9851) circle (1.2pt);
\fill[blue!80!black] (5.9776,3.9851) circle (1.2pt);
\fill[blue!80!black] (5.9777,3.9851) circle (1.2pt);
\fill[blue!80!black] (5.9777,3.9851) circle (1.2pt);
\fill[blue!80!black] (5.9778,3.9852) circle (1.2pt);
\fill[blue!80!black] (5.9778,3.9852) circle (1.2pt);
\fill[blue!80!black] (5.9778,3.9852) circle (1.2pt);
\fill[blue!80!black] (5.9779,3.9853) circle (1.2pt);
\fill[blue!80!black] (5.9779,3.9853) circle (1.2pt);
\fill[blue!80!black] (5.9780,3.9853) circle (1.2pt);
\fill[blue!80!black] (5.9780,3.9853) circle (1.2pt);
\fill[blue!80!black] (5.9781,3.9854) circle (1.2pt);
\fill[blue!80!black] (5.9781,3.9854) circle (1.2pt);
\fill[blue!80!black] (5.9782,3.9854) circle (1.2pt);
\fill[blue!80!black] (5.9782,3.9855) circle (1.2pt);
\fill[blue!80!black] (5.9782,3.9855) circle (1.2pt);
\fill[blue!80!black] (5.9783,3.9855) circle (1.2pt);
\fill[blue!80!black] (5.9783,3.9856) circle (1.2pt);
\fill[blue!80!black] (5.9784,3.9856) circle (1.2pt);
\fill[blue!80!black] (5.9784,3.9856) circle (1.2pt);
\fill[blue!80!black] (5.9785,3.9856) circle (1.2pt);
\fill[blue!80!black] (5.9785,3.9857) circle (1.2pt);
\fill[blue!80!black] (5.9785,3.9857) circle (1.2pt);
\fill[blue!80!black] (5.9786,3.9857) circle (1.2pt);
\fill[blue!80!black] (5.9786,3.9858) circle (1.2pt);
\fill[blue!80!black] (5.9787,3.9858) circle (1.2pt);
\fill[blue!80!black] (5.9787,3.9858) circle (1.2pt);
\fill[blue!80!black] (5.9788,3.9858) circle (1.2pt);
\fill[blue!80!black] (5.9788,3.9859) circle (1.2pt);
\fill[blue!80!black] (5.9788,3.9859) circle (1.2pt);
\fill[blue!80!black] (5.9789,3.9859) circle (1.2pt);
\fill[blue!80!black] (5.9789,3.9860) circle (1.2pt);
\fill[blue!80!black] (5.9790,3.9860) circle (1.2pt);
\fill[blue!80!black] (5.9790,3.9860) circle (1.2pt);
\fill[blue!80!black] (5.9791,3.9860) circle (1.2pt);
\fill[blue!80!black] (5.9791,3.9861) circle (1.2pt);
\fill[blue!80!black] (5.9791,3.9861) circle (1.2pt);
\fill[blue!80!black] (5.9792,3.9861) circle (1.2pt);
\fill[blue!80!black] (5.9792,3.9861) circle (1.2pt);
\fill[blue!80!black] (5.9793,3.9862) circle (1.2pt);
\fill[blue!80!black] (5.9793,3.9862) circle (1.2pt);
\fill[blue!80!black] (5.9793,3.9862) circle (1.2pt);
\fill[blue!80!black] (5.9794,3.9863) circle (1.2pt);
\fill[blue!80!black] (5.9794,3.9863) circle (1.2pt);
\fill[blue!80!black] (5.9795,3.9863) circle (1.2pt);
\fill[blue!80!black] (5.9795,3.9863) circle (1.2pt);
\fill[blue!80!black] (5.9796,3.9864) circle (1.2pt);
\fill[blue!80!black] (5.9796,3.9864) circle (1.2pt);
\fill[blue!80!black] (5.9796,3.9864) circle (1.2pt);
\fill[blue!80!black] (5.9797,3.9864) circle (1.2pt);
\fill[blue!80!black] (5.9797,3.9865) circle (1.2pt);
\fill[blue!80!black] (5.9798,3.9865) circle (1.2pt);
\fill[blue!80!black] (5.9798,3.9865) circle (1.2pt);
\fill[blue!80!black] (5.9798,3.9866) circle (1.2pt);
\fill[blue!80!black] (5.9799,3.9866) circle (1.2pt);
\fill[blue!80!black] (5.9799,3.9866) circle (1.2pt);
\fill[blue!80!black] (5.9800,3.9866) circle (1.2pt);
\fill[blue!80!black] (5.9800,3.9867) circle (1.2pt);
\fill[blue!80!black] (5.9800,3.9867) circle (1.2pt);
\fill[blue!80!black] (5.9801,3.9867) circle (1.2pt);
\fill[blue!80!black] (5.9801,3.9867) circle (1.2pt);
\fill[blue!80!black] (5.9802,3.9868) circle (1.2pt);
\fill[blue!80!black] (5.9802,3.9868) circle (1.2pt);
\fill[blue!80!black] (5.9802,3.9868) circle (1.2pt);
\fill[blue!80!black] (5.9803,3.9868) circle (1.2pt);
\fill[blue!80!black] (5.9803,3.9869) circle (1.2pt);
\fill[blue!80!black] (5.9804,3.9869) circle (1.2pt);
\fill[blue!80!black] (5.9804,3.9869) circle (1.2pt);
\fill[blue!80!black] (5.9804,3.9870) circle (1.2pt);
\fill[blue!80!black] (5.9805,3.9870) circle (1.2pt);
\fill[blue!80!black] (5.9805,3.9870) circle (1.2pt);
\fill[blue!80!black] (5.9805,3.9870) circle (1.2pt);
\fill[blue!80!black] (5.9806,3.9871) circle (1.2pt);
\fill[blue!80!black] (5.9806,3.9871) circle (1.2pt);
\fill[blue!80!black] (5.9807,3.9871) circle (1.2pt);
\fill[blue!80!black] (5.9807,3.9871) circle (1.2pt);
\fill[blue!80!black] (5.9807,3.9872) circle (1.2pt);
\fill[blue!80!black] (5.9808,3.9872) circle (1.2pt);
\fill[blue!80!black] (5.9808,3.9872) circle (1.2pt);
\fill[blue!80!black] (5.9809,3.9872) circle (1.2pt);
\fill[blue!80!black] (5.9809,3.9873) circle (1.2pt);
\fill[blue!80!black] (5.9809,3.9873) circle (1.2pt);
\fill[blue!80!black] (5.9810,3.9873) circle (1.2pt);
\fill[blue!80!black] (5.9810,3.9873) circle (1.2pt);
\fill[blue!80!black] (5.9810,3.9874) circle (1.2pt);
\fill[blue!80!black] (5.9811,3.9874) circle (1.2pt);
\fill[blue!80!black] (5.9811,3.9874) circle (1.2pt);
\fill[blue!80!black] (5.9812,3.9874) circle (1.2pt);
\fill[blue!80!black] (5.9812,3.9875) circle (1.2pt);
\fill[blue!80!black] (5.9812,3.9875) circle (1.2pt);
\fill[blue!80!black] (5.9813,3.9875) circle (1.2pt);
\fill[blue!80!black] (5.9813,3.9875) circle (1.2pt);
\fill[blue!80!black] (5.9813,3.9876) circle (1.2pt);
\fill[blue!80!black] (5.9814,3.9876) circle (1.2pt);
\fill[blue!80!black] (5.9814,3.9876) circle (1.2pt);
\fill[blue!80!black] (5.9815,3.9876) circle (1.2pt);
\fill[blue!80!black] (5.9815,3.9877) circle (1.2pt);
\fill[blue!80!black] (5.9815,3.9877) circle (1.2pt);
\fill[blue!80!black] (5.9816,3.9877) circle (1.2pt);
\fill[blue!80!black] (5.9816,3.9877) circle (1.2pt);
\fill[blue!80!black] (5.9816,3.9878) circle (1.2pt);
\fill[blue!80!black] (5.9817,3.9878) circle (1.2pt);
\fill[blue!80!black] (5.9817,3.9878) circle (1.2pt);
\fill[blue!80!black] (5.9818,3.9878) circle (1.2pt);
\fill[blue!80!black] (5.9818,3.9879) circle (1.2pt);
\fill[blue!80!black] (5.9818,3.9879) circle (1.2pt);
\fill[blue!80!black] (5.9819,3.9879) circle (1.2pt);
\fill[blue!80!black] (5.9819,3.9879) circle (1.2pt);
\fill[blue!80!black] (5.9819,3.9880) circle (1.2pt);
\fill[blue!80!black] (5.9820,3.9880) circle (1.2pt);
\fill[blue!80!black] (5.9820,3.9880) circle (1.2pt);
\fill[blue!80!black] (5.9820,3.9880) circle (1.2pt);
\fill[blue!80!black] (5.9821,3.9881) circle (1.2pt);
\fill[blue!80!black] (5.9821,3.9881) circle (1.2pt);
\fill[blue!80!black] (5.9822,3.9881) circle (1.2pt);
\fill[blue!80!black] (5.9822,3.9881) circle (1.2pt);
\fill[blue!80!black] (5.9822,3.9881) circle (1.2pt);
\fill[blue!80!black] (5.9823,3.9882) circle (1.2pt);
\fill[blue!80!black] (5.9823,3.9882) circle (1.2pt);
\fill[blue!80!black] (5.9823,3.9882) circle (1.2pt);
\fill[blue!80!black] (5.9824,3.9882) circle (1.2pt);
\fill[blue!80!black] (5.9824,3.9883) circle (1.2pt);
\fill[blue!80!black] (5.9824,3.9883) circle (1.2pt);
\fill[blue!80!black] (5.9825,3.9883) circle (1.2pt);
\fill[blue!80!black] (5.9825,3.9883) circle (1.2pt);
\fill[blue!80!black] (5.9825,3.9884) circle (1.2pt);
\fill[blue!80!black] (5.9826,3.9884) circle (1.2pt);
\fill[blue!80!black] (5.9826,3.9884) circle (1.2pt);
\fill[blue!80!black] (5.9826,3.9884) circle (1.2pt);
\fill[blue!80!black] (5.9827,3.9885) circle (1.2pt);
\fill[blue!80!black] (5.9827,3.9885) circle (1.2pt);
\fill[blue!80!black] (5.9827,3.9885) circle (1.2pt);
\fill[blue!80!black] (5.9828,3.9885) circle (1.2pt);
\fill[blue!80!black] (5.9828,3.9885) circle (1.2pt);
\fill[blue!80!black] (5.9829,3.9886) circle (1.2pt);
\fill[blue!80!black] (5.9829,3.9886) circle (1.2pt);
\fill[blue!80!black] (5.9829,3.9886) circle (1.2pt);
\fill[blue!80!black] (5.9830,3.9886) circle (1.2pt);
\fill[blue!80!black] (5.9830,3.9887) circle (1.2pt);
\fill[blue!80!black] (5.9830,3.9887) circle (1.2pt);
\fill[blue!80!black] (5.9831,3.9887) circle (1.2pt);
\fill[blue!80!black] (5.9831,3.9887) circle (1.2pt);
\fill[blue!80!black] (5.9831,3.9887) circle (1.2pt);
\fill[blue!80!black] (5.9832,3.9888) circle (1.2pt);
\fill[blue!80!black] (5.9832,3.9888) circle (1.2pt);
\fill[blue!80!black] (5.9832,3.9888) circle (1.2pt);
\fill[blue!80!black] (5.9833,3.9888) circle (1.2pt);
\fill[blue!80!black] (5.9833,3.9889) circle (1.2pt);
\fill[blue!80!black] (5.9833,3.9889) circle (1.2pt);
\fill[blue!80!black] (5.9834,3.9889) circle (1.2pt);
\fill[blue!80!black] (5.9834,3.9889) circle (1.2pt);
\fill[blue!80!black] (5.9834,3.9890) circle (1.2pt);
\fill[blue!80!black] (5.9835,3.9890) circle (1.2pt);
\fill[blue!80!black] (5.9835,3.9890) circle (1.2pt);
\fill[blue!80!black] (5.9835,3.9890) circle (1.2pt);
\fill[blue!80!black] (5.9836,3.9890) circle (1.2pt);
\fill[blue!80!black] (5.9836,3.9891) circle (1.2pt);
\fill[blue!80!black] (5.9836,3.9891) circle (1.2pt);
\fill[blue!80!black] (5.9837,3.9891) circle (1.2pt);
\fill[blue!80!black] (5.9837,3.9891) circle (1.2pt);
\fill[blue!80!black] (5.9837,3.9891) circle (1.2pt);
\fill[blue!80!black] (5.9838,3.9892) circle (1.2pt);
\fill[blue!80!black] (5.9838,3.9892) circle (1.2pt);
\fill[blue!80!black] (5.9838,3.9892) circle (1.2pt);
\fill[blue!80!black] (5.9839,3.9892) circle (1.2pt);
\fill[blue!80!black] (5.9839,3.9893) circle (1.2pt);
\fill[blue!80!black] (5.9839,3.9893) circle (1.2pt);
\fill[blue!80!black] (5.9839,3.9893) circle (1.2pt);
\fill[blue!80!black] (5.9840,3.9893) circle (1.2pt);
\fill[blue!80!black] (5.9840,3.9893) circle (1.2pt);
\fill[blue!80!black] (5.9840,3.9894) circle (1.2pt);
\fill[blue!80!black] (5.9841,3.9894) circle (1.2pt);
\fill[blue!80!black] (5.9841,3.9894) circle (1.2pt);
\fill[blue!80!black] (5.9841,3.9894) circle (1.2pt);
\fill[blue!80!black] (5.9842,3.9894) circle (1.2pt);
\fill[blue!80!black] (5.9842,3.9895) circle (1.2pt);
\fill[blue!80!black] (5.9842,3.9895) circle (1.2pt);
\fill[blue!80!black] (5.9843,3.9895) circle (1.2pt);
\fill[blue!80!black] (5.9843,3.9895) circle (1.2pt);
\fill[blue!80!black] (5.9843,3.9896) circle (1.2pt);
\fill[blue!80!black] (5.9844,3.9896) circle (1.2pt);
\fill[blue!80!black] (5.9844,3.9896) circle (1.2pt);
\fill[blue!80!black] (5.9844,3.9896) circle (1.2pt);
\fill[blue!80!black] (5.9845,3.9896) circle (1.2pt);
\fill[blue!80!black] (5.9845,3.9897) circle (1.2pt);
\fill[blue!80!black] (5.9845,3.9897) circle (1.2pt);
\fill[blue!80!black] (5.9845,3.9897) circle (1.2pt);
\fill[blue!80!black] (5.9846,3.9897) circle (1.2pt);
\fill[blue!80!black] (5.9846,3.9897) circle (1.2pt);
\fill[blue!80!black] (5.9846,3.9898) circle (1.2pt);
\fill[blue!80!black] (5.9847,3.9898) circle (1.2pt);
\fill[blue!80!black] (5.9847,3.9898) circle (1.2pt);
\fill[blue!80!black] (5.9847,3.9898) circle (1.2pt);
\fill[blue!80!black] (5.9848,3.9898) circle (1.2pt);
\fill[blue!80!black] (5.9848,3.9899) circle (1.2pt);
\fill[blue!80!black] (5.9848,3.9899) circle (1.2pt);
\fill[blue!80!black] (5.9849,3.9899) circle (1.2pt);
\fill[blue!80!black] (5.9849,3.9899) circle (1.2pt);
\fill[blue!80!black] (5.9849,3.9899) circle (1.2pt);
\fill[blue!80!black] (5.9849,3.9900) circle (1.2pt);
\fill[blue!80!black] (5.9850,3.9900) circle (1.2pt);
\fill[blue!80!black] (5.9850,3.9900) circle (1.2pt);
\fill[blue!80!black] (5.9850,3.9900) circle (1.2pt);
\fill[blue!80!black] (5.9851,3.9900) circle (1.2pt);
\fill[blue!80!black] (5.9851,3.9901) circle (1.2pt);
\fill[blue!80!black] (5.9851,3.9901) circle (1.2pt);
\fill[blue!80!black] (5.9852,3.9901) circle (1.2pt);
\fill[blue!80!black] (5.9852,3.9901) circle (1.2pt);
\fill[blue!80!black] (5.9852,3.9901) circle (1.2pt);
\fill[blue!80!black] (5.9852,3.9902) circle (1.2pt);
\fill[blue!80!black] (5.9853,3.9902) circle (1.2pt);
\fill[blue!80!black] (5.9853,3.9902) circle (1.2pt);
\fill[blue!80!black] (5.9853,3.9902) circle (1.2pt);
\fill[blue!80!black] (5.9854,3.9902) circle (1.2pt);
\fill[blue!80!black] (5.9854,3.9903) circle (1.2pt);
\fill[blue!80!black] (5.9854,3.9903) circle (1.2pt);
\fill[blue!80!black] (5.9854,3.9903) circle (1.2pt);
\fill[blue!80!black] (5.9855,3.9903) circle (1.2pt);
\fill[blue!80!black] (5.9855,3.9903) circle (1.2pt);
\fill[blue!80!black] (5.9855,3.9904) circle (1.2pt);
\fill[blue!80!black] (5.9856,3.9904) circle (1.2pt);
\fill[blue!80!black] (5.9856,3.9904) circle (1.2pt);
\fill[blue!80!black] (5.9856,3.9904) circle (1.2pt);
\fill[blue!80!black] (5.9856,3.9904) circle (1.2pt);
\fill[blue!80!black] (5.9857,3.9905) circle (1.2pt);
\fill[blue!80!black] (5.9857,3.9905) circle (1.2pt);
\fill[blue!80!black] (5.9857,3.9905) circle (1.2pt);
\fill[blue!80!black] (5.9858,3.9905) circle (1.2pt);
\fill[blue!80!black] (5.9858,3.9905) circle (1.2pt);
\fill[blue!80!black] (5.9858,3.9905) circle (1.2pt);
\fill[blue!80!black] (5.9858,3.9906) circle (1.2pt);
\fill[blue!80!black] (5.9859,3.9906) circle (1.2pt);
\fill[blue!80!black] (5.9859,3.9906) circle (1.2pt);
\fill[blue!80!black] (5.9859,3.9906) circle (1.2pt);
\fill[blue!80!black] (5.9860,3.9906) circle (1.2pt);
\fill[blue!80!black] (5.9860,3.9907) circle (1.2pt);
\fill[blue!80!black] (5.9860,3.9907) circle (1.2pt);
\fill[blue!80!black] (5.9860,3.9907) circle (1.2pt);
\fill[blue!80!black] (5.9861,3.9907) circle (1.2pt);
\fill[blue!80!black] (5.9861,3.9907) circle (1.2pt);
\fill[blue!80!black] (5.9861,3.9908) circle (1.2pt);
\fill[blue!80!black] (5.9862,3.9908) circle (1.2pt);
\fill[blue!80!black] (5.9862,3.9908) circle (1.2pt);
\fill[blue!80!black] (5.9862,3.9908) circle (1.2pt);
\fill[blue!80!black] (5.9862,3.9908) circle (1.2pt);
\fill[blue!80!black] (5.9863,3.9908) circle (1.2pt);
\fill[blue!80!black] (5.9863,3.9909) circle (1.2pt);
\fill[blue!80!black] (5.9863,3.9909) circle (1.2pt);
\fill[blue!80!black] (5.9864,3.9909) circle (1.2pt);
\fill[blue!80!black] (5.9864,3.9909) circle (1.2pt);
\fill[blue!80!black] (5.9864,3.9909) circle (1.2pt);
\fill[blue!80!black] (5.9864,3.9910) circle (1.2pt);
\fill[blue!80!black] (5.9865,3.9910) circle (1.2pt);
\fill[blue!80!black] (5.9865,3.9910) circle (1.2pt);
\fill[blue!80!black] (5.9865,3.9910) circle (1.2pt);
\fill[blue!80!black] (5.9865,3.9910) circle (1.2pt);
\fill[blue!80!black] (5.9866,3.9910) circle (1.2pt);
\fill[blue!80!black] (5.9866,3.9911) circle (1.2pt);
\fill[blue!80!black] (5.9866,3.9911) circle (1.2pt);
\fill[blue!80!black] (5.9866,3.9911) circle (1.2pt);
\fill[blue!80!black] (5.9867,3.9911) circle (1.2pt);
\fill[blue!80!black] (5.9867,3.9911) circle (1.2pt);
\fill[blue!80!black] (5.9867,3.9912) circle (1.2pt);
\fill[blue!80!black] (5.9868,3.9912) circle (1.2pt);
\fill[blue!80!black] (5.9868,3.9912) circle (1.2pt);
\fill[blue!80!black] (5.9868,3.9912) circle (1.2pt);
\fill[blue!80!black] (5.9868,3.9912) circle (1.2pt);
\fill[blue!80!black] (5.9869,3.9912) circle (1.2pt);
\fill[blue!80!black] (5.9869,3.9913) circle (1.2pt);
\fill[blue!80!black] (5.9869,3.9913) circle (1.2pt);
\fill[blue!80!black] (5.9869,3.9913) circle (1.2pt);
\fill[blue!80!black] (5.9870,3.9913) circle (1.2pt);
\fill[blue!80!black] (5.9870,3.9913) circle (1.2pt);
\fill[blue!80!black] (5.9870,3.9913) circle (1.2pt);
\fill[blue!80!black] (5.9870,3.9914) circle (1.2pt);
\fill[blue!80!black] (5.9871,3.9914) circle (1.2pt);
\fill[blue!80!black] (5.9871,3.9914) circle (1.2pt);
\fill[blue!80!black] (5.9871,3.9914) circle (1.2pt);
\fill[blue!80!black] (5.9871,3.9914) circle (1.2pt);
\fill[blue!80!black] (5.9872,3.9914) circle (1.2pt);
\fill[blue!80!black] (5.9872,3.9915) circle (1.2pt);
\fill[blue!80!black] (5.9872,3.9915) circle (1.2pt);
\fill[blue!80!black] (5.9872,3.9915) circle (1.2pt);
\fill[blue!80!black] (5.9873,3.9915) circle (1.2pt);
\fill[blue!80!black] (5.9873,3.9915) circle (1.2pt);
\fill[blue!80!black] (5.9873,3.9915) circle (1.2pt);
\fill[blue!80!black] (5.9873,3.9916) circle (1.2pt);
\fill[blue!80!black] (5.9874,3.9916) circle (1.2pt);
\fill[blue!80!black] (5.9874,3.9916) circle (1.2pt);
\fill[blue!80!black] (5.9874,3.9916) circle (1.2pt);
\fill[blue!80!black] (5.9875,3.9916) circle (1.2pt);
\fill[blue!80!black] (5.9875,3.9917) circle (1.2pt);
\fill[blue!80!black] (5.9875,3.9917) circle (1.2pt);
\fill[blue!80!black] (5.9875,3.9917) circle (1.2pt);
\fill[blue!80!black] (5.9876,3.9917) circle (1.2pt);
\fill[blue!80!black] (5.9876,3.9917) circle (1.2pt);
\fill[blue!80!black] (5.9876,3.9917) circle (1.2pt);
\fill[blue!80!black] (5.9876,3.9918) circle (1.2pt);
\fill[blue!80!black] (5.9877,3.9918) circle (1.2pt);
\fill[blue!80!black] (5.9877,3.9918) circle (1.2pt);
\fill[blue!80!black] (5.9877,3.9918) circle (1.2pt);
\fill[blue!80!black] (5.9877,3.9918) circle (1.2pt);
\fill[blue!80!black] (5.9877,3.9918) circle (1.2pt);
\fill[blue!80!black] (5.9878,3.9918) circle (1.2pt);
\fill[blue!80!black] (5.9878,3.9919) circle (1.2pt);
\fill[blue!80!black] (5.9878,3.9919) circle (1.2pt);
\fill[blue!80!black] (5.9878,3.9919) circle (1.2pt);
\fill[blue!80!black] (5.9879,3.9919) circle (1.2pt);
\fill[blue!80!black] (5.9879,3.9919) circle (1.2pt);
\fill[blue!80!black] (5.9879,3.9919) circle (1.2pt);
\fill[blue!80!black] (5.9879,3.9920) circle (1.2pt);
\fill[blue!80!black] (5.9880,3.9920) circle (1.2pt);
\fill[blue!80!black] (5.9880,3.9920) circle (1.2pt);
\fill[blue!80!black] (5.9880,3.9920) circle (1.2pt);
\fill[blue!80!black] (5.9880,3.9920) circle (1.2pt);
\fill[blue!80!black] (5.9881,3.9920) circle (1.2pt);
\fill[blue!80!black] (5.9881,3.9921) circle (1.2pt);
\fill[blue!80!black] (5.9881,3.9921) circle (1.2pt);
\fill[blue!80!black] (5.9881,3.9921) circle (1.2pt);
\fill[blue!80!black] (5.9882,3.9921) circle (1.2pt);
\fill[blue!80!black] (5.9882,3.9921) circle (1.2pt);
\fill[blue!80!black] (5.9882,3.9921) circle (1.2pt);
\fill[blue!80!black] (5.9882,3.9922) circle (1.2pt);
\fill[blue!80!black] (5.9883,3.9922) circle (1.2pt);
\fill[blue!80!black] (5.9883,3.9922) circle (1.2pt);
\fill[blue!80!black] (5.9883,3.9922) circle (1.2pt);
\fill[blue!80!black] (5.9883,3.9922) circle (1.2pt);
\fill[blue!80!black] (5.9883,3.9922) circle (1.2pt);
\fill[blue!80!black] (5.9884,3.9922) circle (1.2pt);
\fill[blue!80!black] (5.9884,3.9923) circle (1.2pt);
\fill[blue!80!black] (5.9884,3.9923) circle (1.2pt);
\fill[blue!80!black] (5.9884,3.9923) circle (1.2pt);
\fill[blue!80!black] (5.9885,3.9923) circle (1.2pt);
\fill[blue!80!black] (5.9885,3.9923) circle (1.2pt);
\fill[blue!80!black] (5.9885,3.9923) circle (1.2pt);
\fill[blue!80!black] (5.9885,3.9924) circle (1.2pt);
\fill[blue!80!black] (5.9886,3.9924) circle (1.2pt);
\fill[blue!80!black] (5.9886,3.9924) circle (1.2pt);
\fill[blue!80!black] (5.9886,3.9924) circle (1.2pt);
\fill[blue!80!black] (5.9886,3.9924) circle (1.2pt);
\fill[blue!80!black] (5.9886,3.9924) circle (1.2pt);
\fill[blue!80!black] (5.9887,3.9924) circle (1.2pt);
\fill[blue!80!black] (5.9887,3.9925) circle (1.2pt);
\fill[blue!80!black] (5.9887,3.9925) circle (1.2pt);
\fill[blue!80!black] (5.9887,3.9925) circle (1.2pt);
\fill[blue!80!black] (5.9888,3.9925) circle (1.2pt);
\fill[blue!80!black] (5.9888,3.9925) circle (1.2pt);
\fill[blue!80!black] (5.9888,3.9925) circle (1.2pt);
\fill[blue!80!black] (5.9888,3.9926) circle (1.2pt);
\fill[blue!80!black] (5.9888,3.9926) circle (1.2pt);
\fill[blue!80!black] (5.9889,3.9926) circle (1.2pt);
\fill[blue!80!black] (5.9889,3.9926) circle (1.2pt);
\fill[blue!80!black] (5.9889,3.9926) circle (1.2pt);
\fill[blue!80!black] (5.9889,3.9926) circle (1.2pt);
\fill[blue!80!black] (5.9890,3.9926) circle (1.2pt);
\fill[blue!80!black] (5.9890,3.9927) circle (1.2pt);
\fill[blue!80!black] (5.9890,3.9927) circle (1.2pt);
\fill[blue!80!black] (5.9890,3.9927) circle (1.2pt);
\fill[blue!80!black] (5.9890,3.9927) circle (1.2pt);
\fill[blue!80!black] (5.9891,3.9927) circle (1.2pt);
\fill[blue!80!black] (5.9891,3.9927) circle (1.2pt);
\fill[blue!80!black] (5.9891,3.9927) circle (1.2pt);
\fill[blue!80!black] (5.9891,3.9928) circle (1.2pt);
\fill[blue!80!black] (5.9892,3.9928) circle (1.2pt);
\fill[blue!80!black] (5.9892,3.9928) circle (1.2pt);
\fill[blue!80!black] (5.9892,3.9928) circle (1.2pt);
\fill[blue!80!black] (5.9892,3.9928) circle (1.2pt);
\fill[blue!80!black] (5.9892,3.9928) circle (1.2pt);
\fill[blue!80!black] (5.9893,3.9928) circle (1.2pt);
\fill[blue!80!black] (5.9893,3.9929) circle (1.2pt);
\fill[blue!80!black] (5.9893,3.9929) circle (1.2pt);
\fill[blue!80!black] (5.9893,3.9929) circle (1.2pt);
\fill[blue!80!black] (5.9894,3.9929) circle (1.2pt);
\fill[blue!80!black] (5.9894,3.9929) circle (1.2pt);
\fill[blue!80!black] (5.9894,3.9929) circle (1.2pt);
\fill[blue!80!black] (5.9894,3.9929) circle (1.2pt);
\fill[blue!80!black] (5.9894,3.9930) circle (1.2pt);
\fill[blue!80!black] (5.9895,3.9930) circle (1.2pt);
\fill[blue!80!black] (5.9895,3.9930) circle (1.2pt);
\fill[blue!80!black] (5.9895,3.9930) circle (1.2pt);
\fill[blue!80!black] (5.9895,3.9930) circle (1.2pt);
\fill[blue!80!black] (5.9895,3.9930) circle (1.2pt);
\fill[blue!80!black] (5.9896,3.9930) circle (1.2pt);
\fill[blue!80!black] (5.9896,3.9931) circle (1.2pt);
\fill[blue!80!black] (5.9896,3.9931) circle (1.2pt);
\fill[blue!80!black] (5.9896,3.9931) circle (1.2pt);
\fill[blue!80!black] (5.9896,3.9931) circle (1.2pt);
\fill[blue!80!black] (5.9897,3.9931) circle (1.2pt);
\fill[blue!80!black] (5.9897,3.9931) circle (1.2pt);
\fill[blue!80!black] (5.9897,3.9931) circle (1.2pt);
\fill[blue!80!black] (5.9897,3.9932) circle (1.2pt);
\fill[blue!80!black] (5.9897,3.9932) circle (1.2pt);
\fill[blue!80!black] (5.9898,3.9932) circle (1.2pt);
\fill[blue!80!black] (5.9898,3.9932) circle (1.2pt);
\fill[blue!80!black] (5.9898,3.9932) circle (1.2pt);
\fill[blue!80!black] (5.9898,3.9932) circle (1.2pt);
\fill[blue!80!black] (5.9898,3.9932) circle (1.2pt);
\fill[blue!80!black] (5.9899,3.9932) circle (1.2pt);
\fill[blue!80!black] (5.9899,3.9933) circle (1.2pt);
\fill[blue!80!black] (5.9899,3.9933) circle (1.2pt);
\fill[blue!80!black] (5.9899,3.9933) circle (1.2pt);
\fill[blue!80!black] (5.9900,3.9933) circle (1.2pt);
\fill[blue!80!black] (5.9900,3.9933) circle (1.2pt);
\fill[blue!80!black] (5.9900,3.9933) circle (1.2pt);
\fill[blue!80!black] (5.9900,3.9933) circle (1.2pt);
\fill[blue!80!black] (5.9900,3.9934) circle (1.2pt);
\fill[blue!80!black] (5.9901,3.9934) circle (1.2pt);
\fill[blue!80!black] (5.9901,3.9934) circle (1.2pt);
\fill[blue!80!black] (5.9901,3.9934) circle (1.2pt);
\fill[blue!80!black] (5.9901,3.9934) circle (1.2pt);
\fill[blue!80!black] (5.9901,3.9934) circle (1.2pt);
\fill[blue!80!black] (5.9902,3.9934) circle (1.2pt);
\fill[blue!80!black] (5.9902,3.9934) circle (1.2pt);
\fill[blue!80!black] (5.9902,3.9935) circle (1.2pt);
\fill[blue!80!black] (5.9902,3.9935) circle (1.2pt);
\fill[blue!80!black] (5.9902,3.9935) circle (1.2pt);
\fill[blue!80!black] (5.9902,3.9935) circle (1.2pt);
\fill[blue!80!black] (5.9903,3.9935) circle (1.2pt);
\fill[blue!80!black] (5.9903,3.9935) circle (1.2pt);
\fill[blue!80!black] (5.9903,3.9935) circle (1.2pt);
\fill[blue!80!black] (5.9903,3.9936) circle (1.2pt);
\fill[blue!80!black] (5.9903,3.9936) circle (1.2pt);
\fill[blue!80!black] (5.9904,3.9936) circle (1.2pt);
\fill[blue!80!black] (5.9904,3.9936) circle (1.2pt);
\fill[blue!80!black] (5.9904,3.9936) circle (1.2pt);
\fill[blue!80!black] (5.9904,3.9936) circle (1.2pt);
\fill[blue!80!black] (5.9904,3.9936) circle (1.2pt);
\fill[blue!80!black] (5.9905,3.9936) circle (1.2pt);
\fill[blue!80!black] (5.9905,3.9937) circle (1.2pt);
\fill[blue!80!black] (5.9905,3.9937) circle (1.2pt);
\fill[blue!80!black] (5.9905,3.9937) circle (1.2pt);
\fill[blue!80!black] (5.9905,3.9937) circle (1.2pt);
\fill[blue!80!black] (5.9906,3.9937) circle (1.2pt);
\fill[blue!80!black] (5.9906,3.9937) circle (1.2pt);
\fill[blue!80!black] (5.9906,3.9937) circle (1.2pt);
\fill[blue!80!black] (5.9906,3.9937) circle (1.2pt);
\fill[blue!80!black] (5.9906,3.9938) circle (1.2pt);
\fill[blue!80!black] (5.9906,3.9938) circle (1.2pt);
\fill[blue!80!black] (5.9907,3.9938) circle (1.2pt);
\fill[blue!80!black] (5.9907,3.9938) circle (1.2pt);
\fill[blue!80!black] (5.9907,3.9938) circle (1.2pt);
\fill[blue!80!black] (5.9907,3.9938) circle (1.2pt);
\fill[blue!80!black] (5.9907,3.9938) circle (1.2pt);
\fill[blue!80!black] (5.9908,3.9938) circle (1.2pt);
\fill[blue!80!black] (5.9908,3.9939) circle (1.2pt);
\fill[blue!80!black] (5.9908,3.9939) circle (1.2pt);
\fill[blue!80!black] (5.9908,3.9939) circle (1.2pt);
\fill[blue!80!black] (5.9908,3.9939) circle (1.2pt);
\fill[blue!80!black] (5.9909,3.9939) circle (1.2pt);
\fill[blue!80!black] (5.9909,3.9939) circle (1.2pt);
\fill[blue!80!black] (5.9909,3.9939) circle (1.2pt);
\fill[blue!80!black] (5.9909,3.9939) circle (1.2pt);
\fill[blue!80!black] (5.9909,3.9940) circle (1.2pt);
\fill[blue!80!black] (5.9909,3.9940) circle (1.2pt);
\fill[blue!80!black] (5.9910,3.9940) circle (1.2pt);
\fill[blue!80!black] (5.9910,3.9940) circle (1.2pt);
\fill[blue!80!black] (5.9910,3.9940) circle (1.2pt);
\fill[blue!80!black] (5.9910,3.9940) circle (1.2pt);
\fill[blue!80!black] (5.9910,3.9940) circle (1.2pt);
\fill[blue!80!black] (5.9911,3.9940) circle (1.2pt);
\fill[blue!80!black] (5.9911,3.9940) circle (1.2pt);
\fill[blue!80!black] (5.9911,3.9941) circle (1.2pt);
\fill[blue!80!black] (5.9911,3.9941) circle (1.2pt);
\fill[blue!80!black] (5.9911,3.9941) circle (1.2pt);
\fill[blue!80!black] (5.9911,3.9941) circle (1.2pt);
\fill[blue!80!black] (5.9912,3.9941) circle (1.2pt);
\fill[blue!80!black] (5.9912,3.9941) circle (1.2pt);
\fill[blue!80!black] (5.9912,3.9941) circle (1.2pt);
\fill[blue!80!black] (5.9912,3.9941) circle (1.2pt);
\fill[blue!80!black] (5.9912,3.9942) circle (1.2pt);
\fill[blue!80!black] (5.9912,3.9942) circle (1.2pt);
\fill[blue!80!black] (5.9913,3.9942) circle (1.2pt);
\fill[blue!80!black] (5.9913,3.9942) circle (1.2pt);
\fill[blue!80!black] (5.9913,3.9942) circle (1.2pt);
\fill[blue!80!black] (5.9913,3.9942) circle (1.2pt);
\fill[blue!80!black] (5.9913,3.9942) circle (1.2pt);
\fill[blue!80!black] (5.9914,3.9942) circle (1.2pt);
\fill[blue!80!black] (5.9914,3.9942) circle (1.2pt);
\fill[blue!80!black] (5.9914,3.9943) circle (1.2pt);
\fill[blue!80!black] (5.9914,3.9943) circle (1.2pt);
\fill[blue!80!black] (5.9914,3.9943) circle (1.2pt);
\fill[blue!80!black] (5.9914,3.9943) circle (1.2pt);
\fill[blue!80!black] (5.9915,3.9943) circle (1.2pt);
\fill[blue!80!black] (5.9915,3.9943) circle (1.2pt);
\fill[blue!80!black] (5.9915,3.9943) circle (1.2pt);
\fill[blue!80!black] (5.9915,3.9943) circle (1.2pt);
\fill[blue!80!black] (5.9915,3.9943) circle (1.2pt);
\fill[blue!80!black] (5.9915,3.9944) circle (1.2pt);
\fill[blue!80!black] (5.9916,3.9944) circle (1.2pt);
\fill[blue!80!black] (5.9916,3.9944) circle (1.2pt);
\fill[blue!80!black] (5.9916,3.9944) circle (1.2pt);
\fill[blue!80!black] (5.9916,3.9944) circle (1.2pt);
\fill[blue!80!black] (5.9916,3.9944) circle (1.2pt);
\fill[blue!80!black] (5.9916,3.9944) circle (1.2pt);
\fill[blue!80!black] (5.9917,3.9944) circle (1.2pt);
\fill[blue!80!black] (5.9917,3.9944) circle (1.2pt);
\fill[blue!80!black] (5.9917,3.9945) circle (1.2pt);
\fill[blue!80!black] (5.9917,3.9945) circle (1.2pt);
\fill[blue!80!black] (5.9917,3.9945) circle (1.2pt);
\fill[blue!80!black] (5.9917,3.9945) circle (1.2pt);
\fill[blue!80!black] (5.9918,3.9945) circle (1.2pt);
\fill[blue!80!black] (5.9918,3.9945) circle (1.2pt);
\fill[blue!80!black] (5.9918,3.9945) circle (1.2pt);
\fill[blue!80!black] (5.9918,3.9945) circle (1.2pt);
\fill[blue!80!black] (5.9918,3.9945) circle (1.2pt);
\fill[blue!80!black] (5.9918,3.9946) circle (1.2pt);
\fill[blue!80!black] (5.9919,3.9946) circle (1.2pt);
\fill[blue!80!black] (5.9919,3.9946) circle (1.2pt);
\fill[blue!80!black] (5.9919,3.9946) circle (1.2pt);
\fill[blue!80!black] (5.9919,3.9946) circle (1.2pt);
\fill[blue!80!black] (5.9919,3.9946) circle (1.2pt);
\fill[blue!80!black] (5.9919,3.9946) circle (1.2pt);
\fill[blue!80!black] (5.9920,3.9946) circle (1.2pt);
\fill[blue!80!black] (5.9920,3.9946) circle (1.2pt);
\fill[blue!80!black] (5.9920,3.9947) circle (1.2pt);
\fill[blue!80!black] (5.9920,3.9947) circle (1.2pt);
\fill[blue!80!black] (5.9920,3.9947) circle (1.2pt);
\fill[blue!80!black] (5.9920,3.9947) circle (1.2pt);
\fill[blue!80!black] (5.9920,3.9947) circle (1.2pt);
\fill[blue!80!black] (5.9921,3.9947) circle (1.2pt);
\fill[blue!80!black] (5.9921,3.9947) circle (1.2pt);
\fill[blue!80!black] (5.9921,3.9947) circle (1.2pt);
\fill[blue!80!black] (5.9921,3.9947) circle (1.2pt);
\fill[blue!80!black] (5.9921,3.9948) circle (1.2pt);
\fill[blue!80!black] (5.9921,3.9948) circle (1.2pt);
\fill[blue!80!black] (5.9922,3.9948) circle (1.2pt);
\fill[blue!80!black] (5.9922,3.9948) circle (1.2pt);
\fill[blue!80!black] (5.9922,3.9948) circle (1.2pt);
\fill[blue!80!black] (5.9922,3.9948) circle (1.2pt);
\fill[blue!80!black] (5.9922,3.9948) circle (1.2pt);
\fill[blue!80!black] (5.9922,3.9948) circle (1.2pt);
\fill[blue!80!black] (5.9923,3.9948) circle (1.2pt);
\fill[blue!80!black] (5.9923,3.9948) circle (1.2pt);
\fill[blue!80!black] (5.9923,3.9949) circle (1.2pt);
\fill[blue!80!black] (5.9923,3.9949) circle (1.2pt);
\fill[blue!80!black] (5.9923,3.9949) circle (1.2pt);
\fill[blue!80!black] (5.9923,3.9949) circle (1.2pt);
\fill[blue!80!black] (5.9923,3.9949) circle (1.2pt);
\fill[blue!80!black] (5.9924,3.9949) circle (1.2pt);
\fill[blue!80!black] (5.9924,3.9949) circle (1.2pt);
\fill[blue!80!black] (5.9924,3.9949) circle (1.2pt);
\fill[blue!80!black] (5.9924,3.9949) circle (1.2pt);
\fill[blue!80!black] (5.9924,3.9949) circle (1.2pt);
\fill[blue!80!black] (5.9924,3.9950) circle (1.2pt);
\fill[blue!80!black] (5.9925,3.9950) circle (1.2pt);
\fill[blue!80!black] (5.9925,3.9950) circle (1.2pt);
\fill[blue!80!black] (5.9925,3.9950) circle (1.2pt);
\fill[blue!80!black] (5.9925,3.9950) circle (1.2pt);
\fill[blue!80!black] (5.9925,3.9950) circle (1.2pt);
\fill[blue!80!black] (5.9925,3.9950) circle (1.2pt);
\fill[blue!80!black] (5.9925,3.9950) circle (1.2pt);
\fill[blue!80!black] (5.9926,3.9950) circle (1.2pt);
\fill[blue!80!black] (5.9926,3.9950) circle (1.2pt);
\fill[blue!80!black] (5.9926,3.9951) circle (1.2pt);
\fill[blue!80!black] (5.9926,3.9951) circle (1.2pt);
\fill[blue!80!black] (5.9926,3.9951) circle (1.2pt);
\fill[blue!80!black] (5.9926,3.9951) circle (1.2pt);
\fill[blue!80!black] (5.9926,3.9951) circle (1.2pt);
\fill[blue!80!black] (5.9927,3.9951) circle (1.2pt);
\fill[blue!80!black] (5.9927,3.9951) circle (1.2pt);
\fill[blue!80!black] (5.9927,3.9951) circle (1.2pt);
\fill[blue!80!black] (5.9927,3.9951) circle (1.2pt);
\fill[blue!80!black] (5.9927,3.9951) circle (1.2pt);
\fill[blue!80!black] (5.9927,3.9952) circle (1.2pt);
\fill[blue!80!black] (5.9927,3.9952) circle (1.2pt);
\fill[blue!80!black] (5.9928,3.9952) circle (1.2pt);
\fill[blue!80!black] (5.9928,3.9952) circle (1.2pt);
\fill[blue!80!black] (5.9928,3.9952) circle (1.2pt);
\fill[blue!80!black] (5.9928,3.9952) circle (1.2pt);
\fill[blue!80!black] (5.9928,3.9952) circle (1.2pt);
\fill[blue!80!black] (5.9928,3.9952) circle (1.2pt);
\fill[blue!80!black] (5.9928,3.9952) circle (1.2pt);
\fill[blue!80!black] (5.9929,3.9952) circle (1.2pt);
\fill[blue!80!black] (5.9929,3.9953) circle (1.2pt);
\fill[blue!80!black] (5.9929,3.9953) circle (1.2pt);
\fill[blue!80!black] (5.9929,3.9953) circle (1.2pt);
\fill[blue!80!black] (5.9929,3.9953) circle (1.2pt);
\fill[blue!80!black] (5.9929,3.9953) circle (1.2pt);
\fill[blue!80!black] (5.9929,3.9953) circle (1.2pt);
\fill[blue!80!black] (5.9930,3.9953) circle (1.2pt);
\fill[blue!80!black] (5.9930,3.9953) circle (1.2pt);
\fill[blue!80!black] (5.9930,3.9953) circle (1.2pt);
\fill[blue!80!black] (5.9930,3.9953) circle (1.2pt);
\fill[blue!80!black] (5.9930,3.9953) circle (1.2pt);
\fill[blue!80!black] (5.9930,3.9954) circle (1.2pt);
\fill[blue!80!black] (5.9930,3.9954) circle (1.2pt);
\fill[blue!80!black] (5.9931,3.9954) circle (1.2pt);
\fill[blue!80!black] (5.9931,3.9954) circle (1.2pt);
\fill[blue!80!black] (5.9931,3.9954) circle (1.2pt);
\fill[blue!80!black] (5.9931,3.9954) circle (1.2pt);
\fill[blue!80!black] (5.9931,3.9954) circle (1.2pt);
\fill[blue!80!black] (5.9931,3.9954) circle (1.2pt);
\fill[blue!80!black] (5.9931,3.9954) circle (1.2pt);
\fill[blue!80!black] (5.9932,3.9954) circle (1.2pt);
\fill[blue!80!black] (5.9932,3.9954) circle (1.2pt);
\fill[blue!80!black] (5.9932,3.9955) circle (1.2pt);
\fill[blue!80!black] (5.9932,3.9955) circle (1.2pt);
\fill[blue!80!black] (5.9932,3.9955) circle (1.2pt);
\fill[blue!80!black] (5.9932,3.9955) circle (1.2pt);
\fill[blue!80!black] (5.9932,3.9955) circle (1.2pt);
\fill[blue!80!black] (5.9933,3.9955) circle (1.2pt);
\fill[blue!80!black] (5.9933,3.9955) circle (1.2pt);
\fill[blue!80!black] (5.9933,3.9955) circle (1.2pt);
\fill[blue!80!black] (5.9933,3.9955) circle (1.2pt);
\fill[blue!80!black] (5.9933,3.9955) circle (1.2pt);
\draw[red!80!black, line width=1.0pt] (5.9400,3.9600) -- (6.0000,4.0000);
\fill[red!80!black] (5.9400,3.9600) circle (1.2pt);
\fill[red!80!black] (6.0000,4.0000) circle (1.2pt);
\draw[->, thick] (0,0) -- (12.6000,0) node[right] {$x$};
\draw[->, thick] (0,0) -- (0,8.6000) node[above] {$y$};
\draw (0.0000,0) -- (0.0000,-0.12) node[below] {\small -100.00};
\draw (0,0.0000) -- (-0.12,0.0000) node[left] {\small -100.00};
\draw (2.4000,0) -- (2.4000,-0.12) node[below] {\small -60.00};
\draw (0,1.6000) -- (-0.12,1.6000) node[left] {\small -60.00};
\draw (4.8000,0) -- (4.8000,-0.12) node[below] {\small -20.00};
\draw (0,3.2000) -- (-0.12,3.2000) node[left] {\small -20.00};
\draw (7.2000,0) -- (7.2000,-0.12) node[below] {\small 20.00};
\draw (0,4.8000) -- (-0.12,4.8000) node[left] {\small 20.00};
\draw (9.6000,0) -- (9.6000,-0.12) node[below] {\small 60.00};
\draw (0,6.4000) -- (-0.12,6.4000) node[left] {\small 60.00};
\draw (12.0000,0) -- (12.0000,-0.12) node[below] {\small 100.00};
\draw (0,8.0000) -- (-0.12,8.0000) node[left] {\small 100.00};
\draw[fill=white, draw=black] (7.6,6.7) rectangle (11.8,7.9);
\draw[blue!80!black, line width=1.0pt] (7.9,7.45) -- (8.9,7.45);
\node[anchor=west] at (9.1,7.45) {\small gradDest};
\draw[red!80!black, line width=1.0pt] (7.9,7.05) -- (8.9,7.05);
\node[anchor=west] at (9.1,7.05) {\small fastGradDest};
\end{tikzpicture}

\caption{Функция Schaffer N.2.}
\end{figure}

\section{Обусловленность}
\begin{longtable}{|c|c|c|c|c|}
\caption{Число обусловленности в окрестности минимума.}\label{tab:conditioning}\\
\hline
Функция & Точка & Собственные числа & Число обусловленности & Комментарий \\
\hline
\endfirsthead
\hline
Функция & Точка & Собственные числа & Число обусловленности & Комментарий \\
\hline
\endhead
Adjiman & (2.000000, 0.105783) & полный гессиан неопределён & 1.000000 (редуц.) & Минимум лежит на границе $x=2$; для касательного направления по $y$ редуцированный гессиан одномерен, $H_{yy}=3.782735$. \\
Ackley 2 & (0.000000, 0.000000) & не определён & не определена & Функция негладкая в точке минимума, так как зависит от $\sqrt{x^2 + y^2}$. \\
Schaffer N.2 & (0.000000, 0.000000) & $\lambda_1=0.002000,\ \lambda_2=0.002000$ & 1.000000 & $H = \begin{pmatrix} 0.002000 & 0.000000 \\ 0.000000 & 0.002000 \end{pmatrix}$ \\
Rosenbrock & (1.000000, 1.000000) & $\lambda_1=0.399361,\ \lambda_2=1001.600641$ & 2508.007101 & $H = \begin{pmatrix} 802.000002 & -400.000000 \\ -400.000000 & 200.000000 \end{pmatrix}$ \\
\hline
\end{longtable}


Для Ackley 2 стандартный анализ через гессиан неприменим в точке минимума: функция зависит от $\sqrt{x^2+y^2}$ и не является гладкой в нуле. Поэтому число обусловленности по гессиану в этой точке не определено. Для Schaffer N.2 минимум локально близок к изотропному, поэтому число обусловленности оказывается близким к единице. Для Adjiman матрица вторых производных анизотропна, и это проявляется в более чувствительном выборе шага.
Для Adjiman нужен отдельный комментарий: глобальный минимум достигается на границе допустимой области, поэтому полный гессиан не описывает constrained-минимум так, как в обычной гладкой внутренней точке. В отчёте поэтому используется редуцированный анализ вдоль допустимого направления.

\section{Сравнение со \texttt{scipy.optimize.minimize}}
\begin{longtable}{|c|c|c|c|c|}
\caption{Сравнение с библиотечной реализацией scipy.optimize.minimize.}\label{tab:scipy}\\
\hline
Функция & Точка & $f(x)$ & Итерации & Вызовы $f$ \\
\hline
\endfirsthead
\hline
Функция & Точка & $f(x)$ & Итерации & Вызовы $f$ \\
\hline
\endhead
Adjiman & SciPy недоступен в текущем окружении & - & - & - \\
Ackley 2 & SciPy недоступен в текущем окружении & - & - & - \\
Schaffer N.2 & SciPy недоступен в текущем окружении & - & - & - \\
Rosenbrock & SciPy недоступен в текущем окружении & - & - & - \\
\hline
\end{longtable}


Если SciPy установлен, таблица выше заполняется реальными значениями библиотечного метода \texttt{L-BFGS-B} с теми же ограничениями на область поиска. В противном случае скрипт честно помечает отсутствие зависимости в окружении.

\section{Выводы}
\begin{itemize}
\item Исправленный \texttt{gradDest} теперь корректно реализует адаптацию шага и не зависает при неудачном пробном переходе.
\item \texttt{fastGradDest} делает меньше итераций за счёт одномерного поиска, но обычно тратит больше вычислений функции внутри каждой итерации.
\item На хорошо обусловленных функциях методы сходятся быстрее и стабильнее. На негладких задачах вроде Ackley 2 критерий по гессиану неприменим, а выбор шага становится более чувствительным.
\item Сравнение с эталонными минимумами удобно использовать как независимую проверку корректности реализации.
\end{itemize}

\section*{Воспроизводимость}
Для генерации таблиц и рисунков достаточно выполнить:
\begin{verbatim}
python3 lab1/experiments.py
\end{verbatim}
После установки SciPy дополнительно будет сформирована таблица сравнения с \texttt{scipy.optimize.minimize}.

\end{document}
